\documentclass[11pt,letterpaper,oneside,openany]{report}
\usepackage[T1]{fontenc}
\usepackage{lmodern,amsmath,amsthm,amssymb,mathtools,graphicx,array,xcolor}
\usepackage[left=0.94in,right=0.94in,top=0.86in,bottom=0.88in,headheight=14pt,headsep=13pt,footskip=25pt]{geometry}
\definecolor{ink}{HTML}{192522}
\definecolor{pine}{HTML}{20594C}
\definecolor{muted}{HTML}{4C5F59}
\usepackage[colorlinks=true,linkcolor=pine,urlcolor=pine,citecolor=pine,bookmarksopen=true,bookmarksopenlevel=0,bookmarksnumbered=true,pdfdisplaydoctitle=true]{hyperref}
\hypersetup{pdftitle={The Alp\"oge--Claude argument for Dirichlet GRH},pdfsubject={A mathematical exposition following grh-lean (72fead1) and the community revision (120962b4)}}
\urlstyle{same}
% Each source note is a group. Print its first source folder and each change of folder.
\newcommand{\lastsourcefolder}{}
\def\sourcefolder#1/#2\relax{%
  \def\sourcetail{#2}\ifx\sourcetail\empty\def\thissourcefolder{src}\else\def\thissourcefolder{#1}\fi
  \ifx\thissourcefolder\lastsourcefolder\else
    {\footnotesize\sffamily\color{muted}[\thissourcefolder]}\nobreak\hspace{.35em}%
    \let\lastsourcefolder\thissourcefolder
  \fi}
\newcommand{\src}[2]{\sourcefolder#1/\relax\href{https://github.com/lsernau/grh-lean/blob/72fead102f663d5c7a407964a13dc0cf9d977207/grh/src/#1}{#2}}
\newcommand{\comsrc}[2]{\sourcefolder#1/\relax\href{https://github.com/leanprover-community/riemann-hypothesis/blob/120962b4fc014406f9c76adc53b5cf64475143d3/src/#1}{#2}}
\newcommand{\rootlink}[2]{\href{https://github.com/lsernau/grh-lean/blob/72fead102f663d5c7a407964a13dc0cf9d977207/#1}{#2}}
\newcommand{\leanroute}[1]{\par\smallskip\noindent\textit{Comparison with Lean (optional).}\enspace #1}
\newcommand{\formalnote}[2]{\nobreak\hspace{.10em}\hypertarget{call.#1}{}\hyperlink{note.#1}{\textsuperscript{\normalfont\scriptsize #2}}}
\newcommand{\noteline}[3]{\begingroup\interlinepenalty=10000\noindent\hypertarget{note.#1}{}\hyperlink{call.#1}{\textbf{#2}}\enspace #3\par\endgroup}
\newcommand{\partheading}[1]{\gdef\waitingpart{#1}}
\newcommand{\waitingpart}{}
\newcommand{\currentpart}{}
\newtheorem{theorem}{Theorem}[chapter]
\newtheorem{lemma}[theorem]{Lemma}
\newtheorem{corollary}[theorem]{Corollary}
\newtheorem{proposition}[theorem]{Proposition}
\theoremstyle{definition}\newtheorem{definition}[theorem]{Definition}
\newtheoremstyle{workedexample}
  {8pt plus 2pt minus 2pt}{7pt plus 2pt minus 1pt}
  {\normalfont}{}{\bfseries}{}{\newline}
  {\thmname{#1}\thmnumber{ #2}.\thmnote{\enspace\normalfont\itshape #3}}
\theoremstyle{workedexample}\newtheorem{example}{Example}[chapter]
\AddToHook{env/example/begin}{\renewcommand{\qedsymbol}{\(\lozenge\)}\pushQED{\qed}}
\AddToHook{env/example/end}{\popQED}
\newcommand{\exampleentry}[2]{\noindent\hyperref[#1]{\makebox[3.5em][l]{\ref*{#1}}#2\nobreak\dotfill\makebox[2em][r]{\pageref*{#1}}}\par}
\theoremstyle{remark}\newtheorem{remark}[theorem]{Remark}
\theoremstyle{plain}
\newtheorem*{theorema}{Theorem A}\newtheorem*{theoremb}{Theorem B}\newtheorem*{theoremc}{Theorem C}
\renewcommand{\Re}{\operatorname{Re}}\renewcommand{\Im}{\operatorname{Im}}
\setlength{\emergencystretch}{2em}\setlength{\parindent}{1.25em}
\setlength{\parskip}{2pt plus .5pt minus .4pt}
\setcounter{tocdepth}{1}\setcounter{secnumdepth}{2}
\numberwithin{equation}{section}
\allowdisplaybreaks[2]\raggedbottom
\makeatletter
\newcommand{\lettertheoremlabel}[2]{\phantomsection\begingroup\def\@currentlabel{#1}\label{#2}\endgroup}
\newcommand{\editionchapter}[1]{\clearpage\thispagestyle{plain}%
 \ifx\waitingpart\@empty\else\refstepcounter{part}\addcontentsline{toc}{part}{\thepart\texorpdfstring{\quad}{. }\waitingpart}%
 \xdef\currentpart{Part \thepart\noexpand\enspace\noexpand\textperiodcentered\noexpand\enspace\waitingpart}\gdef\waitingpart{}\fi
 \refstepcounter{chapter}\def\@currentlabelname{#1}\addcontentsline{toc}{chapter}{\protect\numberline{\thechapter}#1}\chaptermark{#1}\@makechapterhead{#1}}
\def\@makechapterhead#1{\vspace*{2pt}{\parindent\z@\raggedright
  {\small\sffamily\color{muted}\currentpart\par}\vskip 9pt
  {\LARGE\bfseries\color{ink}\thechapter\quad #1\par}\nobreak\vskip 15pt}}
\def\@makeschapterhead#1{\vspace*{2pt}{\parindent\z@\raggedright\LARGE\bfseries\color{ink}#1\par}\nobreak\vskip 14pt}
\def\section{\@startsection{section}{1}{\z@}{-2.9ex\@plus -.7ex \@minus -.2ex}{1.1ex \@plus .2ex}{\normalfont\Large\bfseries\raggedright}}
\def\subsection{\@startsection{subsection}{2}{\z@}{-2.3ex\@plus -.5ex \@minus -.2ex}{.7ex \@plus .2ex}{\normalfont\large\bfseries\raggedright}}
\renewcommand*\l@section{\@dottedtocline{1}{1.65em}{3.35em}}
\renewcommand*\l@chapter[2]{\ifnum \c@tocdepth>\m@ne\addpenalty{-\@highpenalty}\vskip .36em plus .06em
  {\parindent\z@\rightskip\@pnumwidth\parfillskip-\@pnumwidth\leavevmode\bfseries
  \setlength\@tempdima{1.65em}\advance\leftskip\@tempdima\hskip-\leftskip #1\nobreak\hfill\nobreak\hb@xt@\@pnumwidth{\hss #2}\par}\fi}
\def\ps@paper{\def\@oddhead{\small\color{muted}\itshape The Alp\"oge--Claude argument for Dirichlet GRH\hfil\normalfont\small\leftmark}\let\@evenhead\@oddhead
  \def\@oddfoot{\hfil\small\thepage\hfil}\let\@evenfoot\@oddfoot}
\makeatother
\pagestyle{paper}\renewcommand{\chaptermark}[1]{\markboth{Chapter \thechapter}{}}
\begin{document}
\pagenumbering{roman}
\hypersetup{pageanchor=false}
\thispagestyle{empty}
\pdfbookmark[0]{Title and abstract}{title}
\vspace*{0.8in}
\begin{center}
{\LARGE\bfseries The Alp\"oge--Claude argument for Dirichlet GRH\par}
\vspace{.6em}
{\large From theta derivatives to a zero-free half-plane\par}
\vspace{1.2em}
{\normalsize An exposition of \texttt{grh-lean} with proofs from the community revision\par}
\end{center}
\vspace{1.4em}
\begin{quotation}
\noindent\textbf{Abstract.}
We present the argument of Alp\"oge and Claude for Dirichlet GRH in the form recorded in the readable Lean development. The arithmetic step gives a one-sided bound for a centered sum over prime powers in a progression. If a Dirichlet \(L\)-function had a zero to the right of the critical line, one could choose a quadratic character so that this zero gives a pole of the Mellin transform of the centered prime sum. The bound excludes that pole by Landau's theorem. The bound itself comes from estimating a Gram determinant at the archimedean and finite places. Divisibility of derivatives of theta functions with complex multiplication gives short vectors. A second-moment estimate shows that finitely many translates of the image of a Poincar\'e map span a subspace of nearly full dimension in the space of forms under consideration. Congruences between automorphic forms bound the denominators incurred by projecting onto the images of theta lifts. The main contribution per dimension is nonnegative and occurs in at least as many dimensions in the lower bound as in the upper bound. Comparing the bounds and dividing by the number \(m\asymp\sqrt x\) of theta lifts gives the square-root saving.
\end{quotation}
\vfill
\noindent\emph{Text and sources.} The \rootlink{README.md}{repository statement} gives the attribution and reports the machine checks. Superscript notes link to the \rootlink{README.md}{readable edition (\texttt{72fead1})} or, for the revised passages indexed in Appendix~\ref{sec:outstanding}, to the \href{https://github.com/leanprover-community/riemann-hypothesis/tree/120962b4fc014406f9c76adc53b5cf64475143d3}{community revision (\texttt{120962b4})}. Optional notes headed \emph{Comparison with Lean} identify changes of proof, including the use of published theorems. Appendix~\ref{sec:outstanding} indexes these notes.
\clearpage\hypersetup{pageanchor=true}
\pdfbookmark[0]{Contents}{contents}
\begingroup\small\setlength{\parskip}{0pt}\tableofcontents\endgroup
\chapter*{List of examples}\label{front:examples}\markboth{List of examples}{List of examples}
\addcontentsline{toc}{chapter}{List of examples}
\noindent Each example is optional. A diamond \(\lozenge\) marks its end.
\par\medskip
\begingroup\small\setlength{\parskip}{1.5pt}
\par\medskip\noindent\textbf{Part I.\enspace The main theorem and its reduction}\par\nobreak\smallskip
\exampleentry{ex:chebyshev-binomial}{Chebyshev's binomial coefficient}
\exampleentry{ex:lemniscate}{The lemniscate sine}
\exampleentry{ex:gaussian-derivative}{Gaussian growth and the size of a derivative}
\exampleentry{ex:integral-parallelogram}{The area of an integral parallelogram}
\exampleentry{ex:progression-filter}{Selecting a residue class by a change of weight}
\exampleentry{ex:quadratic-weight}{The quadratic weight at the first few prime powers}
\exampleentry{ex:mellin-positive}{Real and non-real poles of a positive transform}
\par\medskip\noindent\textbf{Part II.\enspace Short vectors and spanning by translates}\par\nobreak\smallskip
\exampleentry{ex:frobenius-monomial}{A divisibility condition on a single monomial}
\exampleentry{ex:inert-prime-ramp}{The contributions at one inert prime}
\exampleentry{ex:cyclotomic-frame}{A resolution from roots of unity}
\exampleentry{ex:cyclic-second-moment}{One line and its translates}
\exampleentry{ex:compact-coefficient-mass}{The mass of a matrix coefficient on \(U(2)\)}
\exampleentry{ex:radical-conjugation}{Unipotent conjugation in rank two}
\exampleentry{ex:elliptic-character-sign}{The sign of the character sum for \(U(1,1)\)}
\par\medskip\noindent\textbf{Part III.\enspace The determinant lower bound}\par\nobreak\smallskip
\exampleentry{ex:family-height}{A prime-dependent numerator}
\exampleentry{ex:projection-denominator}{The denominator of a projection onto a line}
\exampleentry{ex:line-gauss-three}{The Fourier transform over \(\mathbb Q_3\)}
\exampleentry{ex:plane-subspace-sum}{The subspace sum in a plane}
\exampleentry{ex:composite-gram}{The Gram determinant of a composite}
\exampleentry{ex:tame-tie}{An equation on tame local classes}
\exampleentry{ex:congruent-eigenvalues}{Two congruent eigenvalues}
\exampleentry{ex:weighted-group-norm}{A weighted group norm}
\par\medskip\noindent\textbf{Part IV.\enspace Galois representations and integral comparison}\par\nobreak\smallskip
\exampleentry{ex:local-lattice-extension}{Extending a map to a lattice of full rank}
\exampleentry{ex:small-height-lifting}{Exact lifting at small height}
\exampleentry{ex:transfer-swapped-trace}{The trace with a swap of factors}
\exampleentry{ex:quadratic-norm-congruence}{A norm congruent to minus one modulo sixteen}
\exampleentry{ex:torsion-sections-coalesce}{Torsion sections coalescing in characteristic \(p\)}
\exampleentry{ex:comparison-reduction-scalar}{Passing from residue fields to a quotient of a valuation ring}
\exampleentry{ex:period-linear-extension}{Linear extension of an action on invariant vectors}
\endgroup
\clearpage\pagenumbering{arabic}


\partheading{The main theorem and its reduction}

\editionchapter{Introduction}\label{ch:intro}



\section{Statement and scope}
For a Dirichlet character \(\chi\) modulo \(N\ge1\), let \(L(s,\chi)\) denote the meromorphic continuation of \(\prod_p(1-\chi(p)p^{-s})^{-1}\) from \(\Re s>1\). We give an exposition of the argument of Alp\"oge and Claude for the following theorem.
\begin{theorem}[Dirichlet GRH]\label{thm:main}
For every \(N\ge1\) and every Dirichlet character \(\chi\bmod N\), one has \(L(s,\chi)\ne0\) if \(\Re s>1/2\) and either \(\chi\) is nonprincipal or \(s\ne1\). Every zero in the open critical strip lies on \(\Re s=1/2\). In particular, every nontrivial zero of the Riemann zeta function has real part \(1/2\).
\end{theorem}
The principal character has a pole at \(s=1\). Passage from a primitive character to an induced one removes finitely many Euler factors, and the resulting additional zeros have real part zero. The functional equation gives the critical-strip assertion from nonvanishing to the right of \(1/2\). The character modulo one gives the zeta-function case.\formalnote{1.1}{1}

\paragraph{Reading the proof.}
From Theorem~\ref{thm:main}, go directly to \S\ref{sec:data}. Chapter~\ref{ch:volume} constructs the arithmetic data and combines Theorems A, B and C. Chapter~\ref{ch:analytic} deduces the assertion about zeros. Parts II--IV prove the three estimates used in that deduction.

The rest of this introduction is optional. Section~\ref{sec:motivation} illustrates the method with elementary prime sums and the lemniscate sine. Section~\ref{map:backward} starts from a proposed zero. Section~\ref{map:forward} identifies the estimates specific to the GRH argument and the previously known results on which they depend. A proof-only reading may omit every numbered example, the notes headed \emph{Comparison with Lean}, and the source guides in Appendices~\ref{sec:outstanding} and~\ref{app:formal}.

\section{From divisibility to a prime sum (optional motivation)}\label{sec:motivation}

\subsection{The factorial and a smoothed prime sum}
If every prime in an interval divides the same positive integer, the product of those primes cannot exceed that integer. Chebyshev used this observation to bound a sum over primes. It is also the starting point of the account prepared for M.~O.~Rubinstein \cite[\S\S I.2--I.4]{rubinstein}.

\begin{example}[Chebyshev's binomial coefficient]\label{ex:chebyshev-binomial}
Every prime \(n<p\le2n\) occurs in \((2n)!\) and is absent from \(n!\), so it divides \(\binom{2n}{n}\). At \(n=10\) one can see all four primes in the factorization
\[
                \binom{20}{10}=4\cdot11\cdot13\cdot17\cdot19 .
\]
In general \(\binom{2n}{n}\le 4^n\). Taking logarithms gives \(\sum_{n<p\le2n}\log p\le2n\log2\). Summing over dyadic intervals gives \(\sum_{p\le x}\log p\ll x\) \cite[\S2.2]{montgomery-vaughan}.
\end{example}

To count prime powers, write \(\Lambda(p^j)=\log p\) for \(j\ge1\) and set \(\Lambda(r)=0\) otherwise. For a positive integer \(n\), Legendre's formula gives
\[
  \log(n!)=\sum_{r\le n}\Lambda(r)\left\lfloor \frac nr\right\rfloor .
\]
The difference \(\lfloor n/r\rfloor-(n/r-1)=1-\{n/r\}\) lies in \((0,1]\). Summing this estimate loses \(O(n)\), by Chebyshev's bound \(\sum_{r\le n}\Lambda(r)\ll n\). To study fluctuations on the scale \(\sqrt n\), we work directly with the smoothed sum
\[
 S(x)=\sum_{r\le x}\Lambda(r)(x/r-1)\qquad (x\ge1).
\]
The weight vanishes at the cutoff and is linear in \(x/r\) up to that point. For \(\Re s>1\) its Mellin transform is
\[
 \int_1^\infty S(x)x^{-s-1}\,dx=\frac{-\zeta'(s)/\zeta(s)}{s(s-1)},
 \qquad
 \int_r^\infty (x/r-1)x^{-s-1}\,dx=\frac{r^{-s}}{s(s-1)}.
\]
Absolute convergence permits summation under the integral, and the Euler product gives the logarithmic derivative. At \(s=1\) the transform has principal part \((s-1)^{-2}-(1+\gamma)(s-1)^{-1}\), where \(\gamma\) is Euler's constant. Subtracting \(x\log x-(1+\gamma)x\) from \(S(x)\) removes that pole. At a zero \(\rho\) of \(\zeta\) of multiplicity \(m_\rho\) in the open critical strip the remaining pole has residue \(m_\rho/(\rho(1-\rho))\). In particular the residue is positive at a real zero. The pole at a zero explains the choice of smoothing, and the sign at a real zero is what allows the use of a one-sided estimate.

\subsection{Divisibility of Taylor coefficients}
\begin{example}[The lemniscate sine]\label{ex:lemniscate}
Let \(u(z)\) be the inverse near zero of \(t\mapsto\int_0^t(1-v^4)^{-1/2}\,dv\), with the square root equal to \(1\) at zero. Then \(u(0)=0\), \(u'(0)=1\), and \(u''=-2u^3\). Writing \(A_n=u^{(n)}(0)\), Leibniz's rule gives
\[
  A_{n+2}=-2\sum_{i+j+k=n}\frac{n!}{i!j!k!} A_i A_j A_k .
\]
Since \(A_0=0\) and \(A_1=1\), the recurrence shows that every \(A_n\) is an integer. The first terms of the Taylor series are \(z-z^5/10+z^9/120\), so \(A_5=-2^2\cdot3\) and \(A_9=2^4\cdot3^3\cdot7\). The primes \(3,7\) are inert in \(\mathbb Q(i)\), whereas \(5\) splits. Thus \(v_3(A_5)=1=\lfloor 5\cdot3/8\rfloor\) and \(v_3(A_9)=3=\lfloor9\cdot3/8\rfloor\). These are instances of the bound \(v_p(A_n)\ge\lfloor np/(p^2-1)\rfloor\) for inert \(p\) and \(A_n\ne0\). Since \(p/(p^2-1)=p^{-1}+p^{-3}+\cdots\), odd powers of \(p\) enter the lower bound. In \cite[\S III.12, Thm.~Le18]{rubinstein} the derivatives are taken for a composition \(f(u(z))\). The all-orders bound applies here with the odd polynomial \(f(t)=t\).
\end{example}

In the volume argument we need an upper bound for the complex absolute value of the same derivatives to which we apply divisibility. The theta functions used there are entire with Gaussian growth. Cauchy's formula turns this growth estimate into a bound on their derivatives.

\begin{example}[Gaussian growth and the size of a derivative]\label{ex:gaussian-derivative}
For \(F(z)=e^{z^2}\) the power series gives \(F^{(2k)}(0)=(2k)!/k!\). Its square is \((2k)!\binom{2k}{k}\le (2k)!\,4^k\), and all odd derivatives vanish. A single factorial therefore bounds the squared derivative up to an exponential factor.

Cauchy's formula explains why a single factorial suffices for any entire function of Gaussian growth. If \(|F(z)|\le A e^{C|z|^2}\) with \(A,C>0\), integration on the circle of radius \(\sqrt{n/(2C)}\) gives
\[
 |F^{(n)}(0)|^2\le A^2(n!)^2(2eC/n)^n\le A^2 n!(2eC)^n
 \qquad(n\ge1).
\]
The last inequality uses \(n!\le n^n\).
\end{example}

The factorial in Theorem A comes from applying Cauchy's formula to theta functions with Gaussian growth. For the divisibility estimate, the elliptic curve has complex multiplication by \(\mathbb Q(\sqrt{-d})\), where \(d\equiv7\bmod8\) is prime. Write \(\chi_d(n)=(n/d)\). The factor \(1-\chi_d(p^j)\) is \(2\) at odd powers of an inert prime and \(0\) at its even powers. It vanishes at every power of a split prime and equals \(1\) at powers of the ramified prime. For a derivative of order \(a\), the finite-place estimate uses terms of the form \((a/p^j-\kappa)_+\), where \(\kappa\) is the place weight used for that term and \(u_+=\max(u,0)\). Dividing by \(a\) gives \((1/p^j-1/y)_+\) with \(y=a/\kappa\). This is the weight multiplying \(\Lambda(p^j)\) in \(S(y)/y\). Section~\ref{sec:data} explains how the place weight depends on the prime power and the direction of differentiation. Section~\ref{sec:upper-A} proves the divisibility estimate.

\subsection{From one integer to a volume}
The determinant is the appropriate integer for a family of independent vectors. If \(v_1,\ldots,v_j\) are the columns of a nonsingular integral matrix \(A\), and \(p^{e_p}\mid\det A\), then
\[
              \sum_p e_p\log p\ \le\ \log|\det A|
                   \ \le\ \sum_{i=1}^j\log\|v_i\| .
\]
The first inequality follows from prime factorization, and the second is Hadamard's inequality for the Euclidean norm. For an earlier application to primes, Pritsker integrates a squared, weighted Vandermonde determinant with integer coefficients over a cube. Bounding its integral and clearing denominators gives a lower bound for the integral of Chebyshev's function \(\psi(t)=\sum_{n\le t}\Lambda(n)\) \cite[Thm.~1.1 and its proof in \S3]{pritsker}.

\begin{example}[The area of an integral parallelogram]\label{ex:integral-parallelogram}
Take the matrix \(A\) with columns \((1,1)^{\mathsf T}\) and \((1,7)^{\mathsf T}\). They coincide modulo \(2\) and modulo \(3\), so \(6\) divides their determinant. Their lengths have product \(\sqrt{2}\sqrt{50}=10\). The area \(|\det A|\) is a positive multiple of \(6\) at most \(10\), so it equals \(6\). The Gram determinant \(\det(A^{\mathsf T}A)\) is the squared area, namely \(36\). Its logarithm is \(2\log6\).
\end{example}

In the proof the vectors are automorphic forms, and we estimate their images under \(m\) successive theta lifts. Theorem A bounds the lengths of the lifted vectors. Theorem B gives a lower bound for the dimension spanned by the input vectors. Theorem C bounds the Gram determinant of the lifted vectors from below after accounting for denominators. Comparing the two estimates gives an upper bound for \(m\) times the prime remainder. Taking \(m\asymp\sqrt x\) gives the square-root saving in Chapter~\ref{ch:volume}.


\section{Roadmap from a proposed zero (optional)}\label{map:backward}
Following the roadmap in the guide to the first announcement \cite[p.~1]{oldguide}, we begin with a putative zero \(\rho\). Two choices determine which zeros the prime sum detects. For a proposed zero of a character modulo \(N\), put \(Q=18480N\). The progression \(1\bmod Q\) brings in all characters modulo \(Q\). Multiplying each character by the same quadratic character supplies a second family whose zeros enter with the opposite sign.

\paragraph{The pole that must survive.}
Suppose \(L(\rho,\chi)=0\) for a character \(\chi\bmod N\), with \(1/2<\Re\rho<1\). Fix \(1/2<\beta<\Re\rho\). For a quadratic character \(\chi_d(n)=(n/d)\), compare \(L(s,\psi)\) and \(L(s,\psi\chi_d)\) as \(\psi\) runs over the characters modulo \(Q\). The residue at \(\rho\) of the prime-sum transform below is proportional to the total multiplicity of \(\rho\) in the first family minus its total multiplicity in the second. For the induced character \(\chi_Q\) modulo \(Q\), one still has \(L(\rho,\chi_Q)=0\). An average over quadratic twists supplies a prime \(d\equiv Q/2-1\bmod Q\) for which every twisted \(L\)-function is zero-free in \(\Re s\ge\beta,\ |\Im s|\le|\Im\rho|\). Thus the transform has a pole at \(\rho\). At a real pole \(u\in(\beta,1)\) its residue is positive, since the factor multiplying the untwisted multiplicity is \(1/(\varphi(Q)u(1-u))\).

\paragraph{The prime statistic that excludes the pole.}
With the constant \(c_{Q,d}\) specified in \S\ref{sec:analytic}, the centered statistic is
\[
 R_{Q,d}(x)=\sum_{\substack{n\le x\\n\equiv1\pmod Q}}
    \Lambda(n)(1-\chi_d(n))\left(\frac1n-\frac1x\right)
    -\frac{\log x}{\varphi(Q)}-c_{Q,d}.
\]
Take \(\varepsilon=\beta-1/2\). The bound \(R_{Q,d}(x)\le x^{-1/2+\varepsilon}\) for large \(x\) and boundedness on compact intervals allow us to choose a constant \(C>0\) so that \(Cx^{\beta-1}-R_{Q,d}(x)\) is nonnegative for all \(x\ge1\). By Landau's theorem, the transform of this nonnegative function against \(x^{-s}\,dx\) must have a singularity at its real abscissa of convergence if that abscissa is finite. To the right of \(\beta\) the term \(Cx^{\beta-1}\) has holomorphic transform \(C/(s-\beta)\). Any real pole there would consequently be simple with negative residue, from the minus sign before \(R_{Q,d}\). Such a pole would make the integral negative just to the right of its abscissa. Hence the abscissa is at most \(\beta\), and the transform of \(R_{Q,d}\) is holomorphic at \(\rho\), a contradiction.

\paragraph{Admissible data at a prescribed cutoff.}
With \(Q,d\) fixed, construct the field, levels and weights for each large \(x\). The volume inequality involves remainders at \(x\) and a larger cutoff \(x_1\). At \(x_1\) the classical prime number theorem in progressions saves any fixed power of \(\log x_1\). At \(x\) the comparison of volumes gives \(R_{Q,d}(x)\le x^{-1/2+\varepsilon}\) for every \(\varepsilon>0\) and all sufficiently large \(x\) (Proposition~\ref{prop:an-one-sided}).

\section{Volume estimates and established inputs (optional roadmap)}\label{map:forward}
The readable edition's \rootlink{grh/INDEX.md}{index of results} distinguishes previously known mathematics formalized for this project from the constructions and estimates particular to the GRH argument. Theorem B supplies independent vectors. Theorems A and C bound the volume of their lifts from above and below.
\begin{description}
\item[{\hyperref[thm:short]{Theorem A} (short vectors).}] For a CM elliptic curve at an odd inert prime of good reduction, Katz proved congruences for normalized Bernoulli--Hurwitz numbers, which come from the Laurent expansion of the Weierstrass \(\wp\)-function \cite{katz-cartier}. Bannai and Kobayashi state and prove the congruences in this setting for a curve defined over its CM field \cite[\S5, Thms.~5.1 and 5.8]{bannai-integral}. Chapters~\ref{ch:upper}--\ref{ch:upperlift} calculate the divisibility of theta derivatives along products of matrix minors. At the marked primes they use scaled Hecke operators to recover divisibility on a sublattice and bound its index.
\item[{\hyperref[thm:spanning]{Theorem B} (spanning by translates).}] Schur orthogonality gives the squared norm of a discrete-series matrix coefficient in terms of the formal degree \cite[Thm.~B, (1)]{upper-moore-wolf}. Chapters~\ref{ch:span}--\ref{ch:trace} estimate the second moment of the translates of the Poincar\'e map and show that few translates span nearly all of the target. A comparison of traces across signatures transfers the dimension count from a definite unitary group.
\item[{\hyperref[thm:volume]{Theorem C} (the determinant).}] The Rallis inner product formula relates the squared norm of an individual theta lift to a doubling integral \cite[\S\S11.3--11.5]{gan-qiu-takeda-rallis}. Chapters~\ref{ch:lower}--\ref{ch:chain} compare determinants for a whole chain with and without a twist. To bound the heights of twisted Hecke values they use \S\ref{sec:frame-attainment}, where the local divisibility bound is attained in prescribed degrees up to one factor of controlled size. For each coefficient prime \(\ell\) and isomorphism \(\mathbb C\simeq\overline{\mathbb Q}_\ell\), Chapters~\ref{ch:selmer}--\ref{ch:deform} choose one twist whose first \(m\) powers control the projection denominators for the \(m\) lifts. Chapter~\ref{ch:local} bounds the loss at \(\ell\). Chapters~\ref{ch:galois}--\ref{ch:cohom} supply the Galois representations with the properties stated in \S\ref{sec:galois}.
\end{description}
\noindent Chapter~\ref{ch:volume} combines A, B and C. Chapter~\ref{ch:analytic} turns the resulting bound at \(x\) into a zero-free half-plane by choosing a quadratic twist and applying Landau's theorem.



\editionchapter{The arithmetic reduction}\label{ch:volume}

We begin the proof of Theorem~\ref{thm:main} by fixing a large integer \(x\). We seek a bound for a sum over prime powers \(n\le x\) in one residue class. Section~\ref{sec:data} encodes that class in a sum over places of a number field. Theorems A, B and C then give two estimates for the same Gram determinant. Comparing them bounds the remainder in the prime sum.

\section{Arithmetic data at a prescribed primary cutoff}
\label{sec:data}
Fix a positive multiple \(q\) of \(18480\) and a prime \(d\equiv q/2-1\bmod q\). In particular \(d\equiv7\bmod8\) and \((d,q)=1\). Write \(K=\mathbb Q(\sqrt{-d})\).

For each \(x\) we shall choose a second cutoff \(x_1>x\) and a positive integer \(a\), the derivative order used in \S\ref{sec:upper-A}. For a positive integer \(n\) and a positive weight \(w\), the factor
\[
                 (a/n-w)_+=\max\{a/n-w,0\}
\]
is positive exactly when \(n<a/w\). Thus \(w=a/x_1\) gives cutoff \(x_1\), and \(w=a/x\) gives cutoff \(x\). We shall index the complex places of a field by residue classes. At a place \(\sigma\), write \(\kappa_\sigma\) for the place weight, which takes one of these two values. The choice of places with \(\kappa_\sigma=a/x\) will select the class \(n\equiv1\bmod q\). We choose \(x_1\) large enough that the prime number theorem at \(x_1\) gives an error negligible on the scale of \(x\).

\begin{proposition}[Existence of admissible data]\label{prop:an-data}
Let \(0<\eta<1/100\). For every sufficiently large integer \(x\) there are admissible data as specified below, with a CM field \(E\supset K\), sets of infinite places \(\Phi_a,\Phi_b\), integers \(a,b,x_1\), and \(m\) theta lifts, with
\[
  \sqrt x\le m\le2\sqrt x,\qquad a=xb,\qquad
  x^\eta\le\log b\le x^\eta+1,\qquad
  x^{\eta/6}-1\le\log x_1\le x^{\eta/6}.
\]
The number \(f\) of infinite places equals \(\varphi(q)|\Phi_a|\), and \(0<|\Phi_a|/|\Phi_b|\le\log b\). There are also distinct primes \(p_0,p_1\) splitting completely in \(E\) whose product is the norm of an element of \(\mathcal O_K\), and a prime \(N\) splitting completely in \(E\) in the interval specified below.
\end{proposition}
The proof consists of the following three constructions. The threshold for \(x\) may depend on \(q,d,\eta\).

\subsection{The ranks, the field and the two weights}
Fix \(0<\eta<1/100\), and put \(W=x^{\eta/6}\). Take \(m\) to be the smallest power of two at least \(\sqrt x\), and \(r\) the smallest power of three at least \(x^5\). Consecutive powers differ by a factor two or three, hence
\[
 \sqrt x\le m\le2\sqrt x,\quad x^5\le r\le3x^5,\quad
 k=r+1-m,\quad t=5r .
\]
In \S\ref{sec:spaces} we construct a period map between spaces of forms on two unitary groups. Their underlying Hermitian spaces, denoted by \(X_0\) and \(Y_0\), have ranks \(k\) and \(r\), respectively. Choose \(s\) so \(2^s=m^{22}\), let \(l\) be the least power of three strictly greater than \(27s\), and set
\[
 q_1=3l\,2^s q,\quad q_2=(qd)^{s+6},\quad E=K(\zeta_{q_1}).
\]
Since \(m\le2\sqrt{x}\), we have \(s=22\log_2 m\le32\log x\) for sufficiently large \(x\). Also \(27s<l\le81s\). The discriminant of \(K\) is \(-d\) and \(d\nmid q_1\). The intersection of \(K\) with \(\mathbb Q(\zeta_{q_1})\) is \(\mathbb Q\), since \(d\) ramifies in \(K\) and is unramified in the cyclotomic field. Fix one of the two complex embeddings of \(K\). Each conjugate pair of embeddings of \(E\) has one member extending it. Write \(\Phi\) for these representatives and label them by the exponent \(\sigma\in(\mathbb Z/q_1)^\times\) in \(\zeta_{q_1}\mapsto\zeta_{q_1}^{\sigma}\). Write \(\Phi_a\) for the labels congruent to \(1\bmod q\). Increasing only the exponents of the primes \(2,3\) already dividing \(q\) gives
\[
 f=\varphi(q_1)=3l\,2^s\varphi(q),\qquad |\Phi_a|=3l\,2^s .
\]
Put \(x_1=\lfloor e^W\rfloor\) and take \(b\) to be the least multiple of \(q_1 x_1\) strictly above \(e^{W^6}\). Then \(e^{W^6}<b\le e^{W^6}+q_1x_1\). Since \(q_1\) grows polynomially in \(x\), \(q_1x_1\le e^{W^6}\) eventually. Thus \(W^6\le\log b\le W^6+1\), while \(W-1\le\log x_1\le W\). In particular \((\log b)^{100}\le x\) eventually since \(100\eta<1\). Set \(a=xb\). We take derivatives of order \(a\) in the directions indexed by \(\Phi_a\). Independently, choose a set \(\Phi_b\) of labels congruent to \(-1\bmod q\), constructed below. Set \(\kappa_\sigma=b=a/x\) for \(\sigma\in\Phi_b\) and \(\kappa_\sigma=a/x_1\) elsewhere. Both values are integers. For \(n\) prime to \(q_1\), the term indexed by \(\tau\in\Phi_a\) uses the place weight \(\kappa_{-n\tau}\), with the subscript reduced modulo \(q_1\).

Since \(W/\log x\to\infty\) and \(W^6/W\to\infty\), these choices also give \(x^{10^6}\le x_1\) and \(x_1^{10^6}\le b\) eventually. In particular \(1\le a/x_1\le b\). We have \((\log x_1)^6\le\log b\), and \(\log\log b\ge\log(x^\eta)=\eta\log x\) eventually exceeds \(10^{12}C(q,d)\), where \(C(q,d)\) is the fixed constant bound in the definition of admissibility. These are the scale conditions used by Theorems A--C.\formalnote{2.1}{1}

\begin{example}[Selecting a residue class by a change of weight]\label{ex:progression-filter}
Set \(a=60\). Use the units modulo \(20\) as indices. The two indices congruent to \(1\bmod5\) are \(1\) and \(11\). For each positive \(n\) coprime to \(20\), form the sum
\[
          \bigl(a/n-\kappa_{-n}\bigr)_++\bigl(a/n-\kappa_{-11n}\bigr)_+,
\]
where \(u_+=\max\{u,0\}\) and the subscripts are read modulo \(20\). First put \(\kappa_\sigma=4\) for every unit \(\sigma\). Both terms then have cutoff \(60/4=15\). Now change only \(\kappa_9\) to \(6\), giving cutoff \(60/6=10\) whenever a subscript is \(9\).

At \(n=11\) the subscripts are \(9\) and \(19\). The first term drops from \(16/11\) to zero, and the second is still \(16/11\). At \(n=3\) the subscripts are \(17\) and \(7\), so both terms are still \(16\). More generally, one of the subscripts is \(9\) precisely when \(-9n^{-1}\bmod20\) belongs to \(\{1,11\}\). Since \(-9\equiv1\bmod5\), this happens exactly when \(n\equiv1\bmod5\). In that class one summand uses cutoff \(10\) and the other keeps cutoff \(15\). Both use cutoff \(15\) in the other classes. The sum decreases precisely when \(n\equiv1\bmod5\) and \(n<15\). For instance, \(n=21\) uses the shorter cutoff, but both summands are already zero.
\end{example}

For the general count take \(n\) prime to \(q_1\). Multiplication by \(-n\) is a bijection on the units modulo \(q_1\). Every \(\sigma\in\Phi_b\) reduces to \(-1\bmod q\), so its unique preimage \(-n^{-1}\sigma\) lies in \(\Phi_a\) exactly when \(n\equiv1\bmod q\). This gives \(|\Phi_b|\) changes of cutoff for such \(n\) and none for the other units.

\subsection{A set of places with controlled prime-power overlaps}
We require \(\Phi_b\) in the fiber \(-1\bmod q\). For every prime \(p\le kt\) and \(0<i<j<l\) we prohibit triples \(\sigma,\tau,\rho\) in this set, with repeated places allowed, satisfying
\[
   \tau/\sigma\equiv \pm p^i,\qquad \rho/\sigma\equiv \pm p^j
\]
modulo the largest divisor of \(q_1\) coprime to \(p\). The signs are independent. For \(n\equiv1\bmod q\), the difference between the two positive parts for \(n\) occurs \(|\Phi_b|\) times. Recovering the prime sum in this class requires division by \(|\Phi_b|\). The lower bound \(|\Phi_b|\ge|\Phi_a|/\log b\) controls the resulting error.

Behrend used lattice points on one sphere to construct a set containing no three distinct terms in arithmetic progression \cite[pp.~331--332]{behrend}. Shapira records the weighted three-variable equation \cite[Thm.~1 of the preprint]{shapira-behrend}. Alon and Shapira construct one set of integers on which \(a u+b w=(a+b)v\) forces \(u=v=w\) simultaneously for all positive \(a,b\) up to a given bound \cite[Lemma~4.1 of the preprint]{alon-shapira}. We need the congruence modulo a specified integer, with an explicit count for a fixed number of digits. Suppose \((LH)^h\le o\) for positive integers \(L,H,h,o\) with \(L\ge2\). Encode \(v\in\{0,\ldots,H-1\}^h\) by its digits in base \(LH\). Pick one value of \(\|v\|^2\) attained by at least \(H^h/(hH^2)\) vectors, using that there are at most \(hH^2\) possible values. Let \(A\) be their encodings. If \(e_1,e_2,e_3\in A\) encode \(u,v,w\), and \(\alpha,\beta\) are positive integers with \(\alpha+\beta\le L\), then
\[
 (\alpha+\beta)e_2\equiv\alpha e_1+\beta e_3\pmod o
     \ \Longrightarrow\  (\alpha+\beta)v=\alpha u+\beta w .
\]
Indeed each digit of either side is less than \(LH\), so there are no carries, and both sides are below \(o\). Taking squared norms in the resulting convex combination gives
\[
 (\alpha+\beta)\{\alpha\|u\|^2+\beta\|w\|^2\}
   -\|\alpha u+\beta w\|^2=\alpha\beta\|u-w\|^2=0 .
\]
Thus \(u=v=w\). This argument handles every allowed pair of positive coefficients with the same set.

Take \(g=1+3lq\), \(v_2(q)=v\ge4\). Modulo \(2^{s+v}\), \(g\) has exact order \(o=2^s\). Indeed \(v_2(g-1)=v\) and successively squaring \(1+2^v u\), \(u\) odd, increases that valuation by one. The selected places are \(-g^e\) for \(e\in A\), using \(L=2l\). They are distinct and congruent to \(-1\bmod 3lq\). Suppose first that \(p\) is odd. Raise the two forbidden equalities to powers \(2j\) and \(2i\), respectively. The \(p\)-powers agree and the signs disappear, giving
\[
       2j e_2\equiv 2(j-i)e_1+2ie_3\pmod {2^s}.
\]
The additive construction forces equality of the three exponents. Then \(p^{2i}\equiv1\bmod 2^{s+1}\). The valuation identity
\(v_2(p^{2i}-1)=v_2(p-1)+v_2(p+1)+v_2(i)\)
implies \(2^s\le (p+1)2i\), since one of \(p-1,p+1\) has 2-adic valuation one and the other is at most \(p+1\). But \(2l(kt+1)<2^s\) eventually: the left side is \(O(x^{10}\log x)\) and \(2^s\ge x^{11}\). For \(p=2\) reduce the first relation modulo \(3l\) and square it. It gives \(4^i\equiv1\bmod 3l\), impossible since \(4=1+3\) has order \(l\) there by successive cubing.

For the size, choose a fixed positive \(h\), sufficiently large that \(2^s\le W^h\) eventually, and take a power of two \(H\) satisfying
\[
     H^h W\le 2^s < (2H)^hW .
\]
This is possible since \(2^s\ge W\). We have \(H\le W\). As \(l=O(\log x)\), eventually \((2l)^h\le W\) and \(3lh2^h\le W^3\). The cube fits, since \((2lH)^h\le W H^h\le 2^s\), and the count gives
\[
    |\Phi_a|=3l\,2^s
       \le \frac{H^h}{hH^2}\,(3lh2^h W H^2)
       \le |\Phi_b| W^6
       \le |\Phi_b|\log b .
\]
This proves both the required cardinality and sparseness.\formalnote{2.2}{2}

\subsection{Two norm primes and a level for the whole chain}
For each fixed \(A>0\), suppose that \(Q_{\rm lev}\le W^A\) and \(W\le U\le W^A\). For sufficiently large \(W\), there is a prime \(p\equiv1\bmod Q_{\rm lev}\) with \(U\le\log p\le U+W\). Indeed put \(z=e^U\). Apply Siegel--Walfisz \cite[Cor.~11.21]{montgomery-vaughan} to \(\vartheta(2z;Q_{\rm lev},1)-\vartheta(z;Q_{\rm lev},1)\), where \(\vartheta\) sums \(\log p\) over primes. Here \(Q_{\rm lev}\le(\log z)^A\), so the difference is \(z/\varphi(Q_{\rm lev})+O_A(z/(\log z)^{A+1})\). The main term is at least \(z/(\log z)^A\). The resulting prime has \(U<\log p\le U+\log 2\le U+W\).
\phantomsection\label{route:1}\leanroute{The formal result treats powers of a fixed modulus divisible by four. It descends the real characters to a fixed modulus and uses a zero-free region and Newman's contour argument for convergence of Dirichlet series \cite[pp.~693--695]{newman}. Siegel--Walfisz applies uniformly to the particular varying moduli \(Q_{\rm lev}=q_2^2\) used here because \(Q_{\rm lev}\) has polynomial size in \(W\). Its ineffective threshold suffices for existence at every sufficiently large \(x\).\formalnote{2.3}{3}}

Take \(Q_{\rm lev}=q_2^2\) and \(L_b=\log b\). Both \(Q_{\rm lev}\) and \(24mL_b\) are bounded by fixed powers of \(W\), because \(s=O(\log x)\), \(q,d,\eta\) are fixed, and \(W=x^{\eta/6}\). Write \(h_K\) for the class number of \(K\). The \(h_K+1\) disjoint logarithmic intervals of length \(W\) starting at \(10L_b+2iW\), \(0\le i\le h_K\), lie in \([10L_b,11L_b]\) for large \(x\). Each contains the logarithm of a different prime \(p_i\equiv1\bmod Q_{\rm lev}\). Since \(3l\le243s\le3^{s+5}\) and \(6\mid q\), we have \(3l\,2^s\mid(qd)^{s+5}\) and hence \(q_1d\mid q_2\). Thus primes congruent to \(1\bmod Q_{\rm lev}\) split completely in \(\mathbb Q(\zeta_{q_1d})\) and in its subfield \(E\). Factor \((p_i)=I_iJ_i\) into its two prime ideals in \(\mathcal O_K\). There are only \(h_K\) ideal classes, so \([I_i]=[I_j]\) for some \(i\ne j\). Since \([J_j]=[I_j]^{-1}\), \(I_iJ_j=(\alpha)\) is principal and its norm is \(p_ip_j=\alpha\bar\alpha\). Call this product \(M\), and relabel the selected primes \(p_0,p_1\). In particular \(b^{20}\le M\le b^{22}\).

Write \(\mathrm{rd}_E=|\operatorname{disc}(E)|^{1/[E:\mathbb Q]}\) for the root discriminant of \(E\) and \(\bar\kappa\) for the geometric mean of the place weights, so \(\log\bar\kappa=f^{-1}\sum_{\tau\in\Phi}\log\kappa_\tau\). Set \(G=\log(M\,\mathrm{rd}_E/(2\pi))\) and \(U=m\log\bar\kappa+(m-1)G-\frac{15}{2}L_b\). The field \(E\) lies in \(\mathbb Q(\zeta_{q_1d})\), since \(K\) has conductor \(d\). The cyclotomic discriminant formula and the tower formula for discriminants give \(\mathrm{rd}_E\le q_1d\le b\). With \(1\le\kappa_\tau\le b\) this yields \(mL_b\le U\le24mL_b\) for large \(x\). For the lower bound use \(G\ge20L_b-\log(2\pi)\) and \(m\to\infty\). Choose a prime \(N\equiv1\bmod Q_{\rm lev}\) with
\[
             U\le\log N\le U+W\le U+\sqrt{L_b}.
\]
Since \(N\equiv1\bmod4\), a sum-of-two-squares representation gives \(\pi_N\in\mathbb Z[i]\subset\mathcal O_E\) with \(\pi_N\bar\pi_N=N\). Finally
\[
   m(\log\bar\kappa+G)-\log N
       =G+\tfrac{15}{2}L_b-(\log N-U)=O(L_b).
\]
The two quantities \(m(\log\bar\kappa+G)\) and \(\log N\) are of order \(mL_b\). Choosing \(N\) in this interval keeps their difference within \(O(L_b)\).\formalnote{2.4}{4}


\section{The spaces and the invariant being estimated}
\label{sec:spaces}
For a Hermitian space \(Z\) over \(E\), let \(U(Z)\) be the group of \(E\)-linear isometries of its form, viewed over the totally real subfield of \(E\). Thus \(U(Z)\) has finite adelic points, which record all finite completions simultaneously, and real points, which are products of unitary groups of prescribed signatures. A compact open subgroup at the finite places fixes the level of an automorphic form. On a definite group, fixing the weight at infinity identifies a space of such forms with vector-valued functions on finitely many adelic classes, with the stabilizer at each class acting on the values. This makes the integral structure concrete. The modules of interest distinguish a source vector's integrality at a finite prime from the complex size of its lifted image.

There are two chains. For \(0\le j\le m\) the ranks of their spaces are
\[
             n_j=\operatorname{rank}_E X_j=k+j(t+1),
                 \qquad \operatorname{rank} Y_j=r+jt .
\]
The \(X_j\) are definite. The group of \(Y_0\) has a Bessel subgroup associated with \(X_0\); its Hermitian form is the orthogonal sum of \(X_0\), \((m-2)/2\) hyperbolic planes, and a line of norm \(-N\). The ranks agree because \(k+(m-2)+1=r\). At the far end \(Y_m\) is a definite subspace of \(X_m\) of codimension one, since \(k+m-r=1\). These two endpoint relations give a period map from forms on \(X_0\) to forms on \(Y_0\), and a restriction from forms on \(X_m\) to forms on \(Y_m\).

Recall that \(M=p_0p_1=\alpha\bar\alpha\). Write \(\bar\pi_M=\alpha\) and \(\pi_M=\bar\alpha\). The chosen places at \(p_0,p_1\) are those dividing \(\pi_M\), one from each conjugate pair. Their prime ideals have product \((\pi_M)\), since \(p_0p_1\) is squarefree and splits completely in \(E\). On either group the principal level at these places consists of the finite integral elements reducing to the identity at every \(w\mid\pi_M\).

The source space \(S_X\) consists of algebraic modular forms on \(U(X_0)\) with highest weight \((a,0,\ldots,0)\) at each \(\tau\in\Phi_a\) and zero weight elsewhere. Its finite level consists of the elements of principal level at the chosen places whose last row is congruent to \((0,\ldots,0,1)\) modulo \(\pi_N\). For each place \(w\mid\pi_N\), choose an integral norm-one finite idele \(z_w\), equal to \(1\) off \(w,\bar w\), whose component at \(w\) reduces to a generator of \(\mathbb F_N^\times\). Its scalar matrix acts by right translation. Its \(i\)-th power belongs to the source level exactly when \(N-1\mid i\). The ramified source space \(S_{\rm ram}\) consists of the forms whose average over \(1,z_w,\ldots,z_w^{N-2}\) vanishes for every such \(w\).

At infinity we also fix a type on \(U(Y_0)\). At the place \(\tau\in\Phi\) its Harish--Chandra parameter consists of the two ordered lists
\[
       \bigl(k/2-i\bigr)_{1\le i\le k-1},\qquad
       \bigl(-\kappa_\tau-2jr^8\bigr)_{1\le j\le m}.
\]
The basis directions belonging to the first list have positive sign in the Hermitian form. In the second list the signs alternate, starting with a positive sign. For this discrete series let \(c_\tau\) be the matrix coefficient of a unit highest weight vector in its lowest compact type, and let \(d_\tau\) be its formal degree. Write \(e_\infty=\prod_{\tau\in\Phi}d_\tau\overline{c_\tau}\) for the projection kernel on the product of the real groups (\S\ref{sec:real-pairing}). Put \(G_\infty=U(Y_0)_\infty\), the product of the archimedean groups. Let \(\mathcal A(Y_0)\) consist of the bounded automorphic forms of the principal level above satisfying
\[
       F(y)=\int_{G_\infty}e_\infty(g)F(yg)\,dg .
\]

At the finite places the group \(\Gamma=\prod_{w\mid\pi_M}\mathrm{GL}_r(k_w)\), with \(k_w\) the residue field at \(w\), acts on \(\mathcal A(Y_0)\) by translation. The good space \(Y=e_0\mathcal A(Y_0)\) is the image of the central idempotent \(e_0\in\mathbb C[\Gamma]\) specified in \S\ref{sec:B}. Write \(D=\dim_{\mathbb C}Y\).

For a pair of groups a theta kernel is obtained by summing a Weil-representation test function over rational vectors. Pairing the theta kernel with an automorphic form over the source quotient produces a form on the other group. Schematically, with the conjugation determined by the pairing convention,
\[
      \phi\longmapsto
      \left(h\longmapsto\int_{[U(Z)]}
              \mathcal K(g,h)\,\phi(g)\,dg\right).
\]
At the definite places the resulting entries can be expressed using Taylor coefficients. Composing the \(m\) maps along \(Y_0,\ldots,Y_m\) gives the composite theta lift \(\Theta_\xi\). The kernel at step \(j\) uses the power \(\xi^j\) of an automorphic character \(\xi\) of the norm-one torus. For an \(m\)-fold twist, \(\xi\) has weight zero and its tame part is supported on a finite set of places of \(E\) disjoint from its conjugate and avoiding the places above \(MN\). At each place of this set the residue character has order greater than \(m\), so \(\xi,\ldots,\xi^m\) have the same conductor. The trivial twist has empty support and gives \(\Theta_1\).

For a source form \(\phi\in S_X\), its Bessel period against a test form on \(U(Y_0)\) defines a conjugate-linear functional of that test form. The Poincar\'e map sends \(\phi\) to the Riesz representative of the functional. The seesaw identities express the result of applying \(\Theta_\xi\) to the projected Poincar\'e vector in terms of lifting along the definite \(X\) chain and then restricting to \(Y_m\). The choice of local data and all normalizing scalars in that identity matter for integrality. Section~\ref{sec:upper-A} records them in the size and index calculations.\formalnote{2.5}{5}

To measure a family at infinity we apply each automorphism \(\sigma\) of \(\mathbb C\) to the lift coordinates in a fixed rational Gram model. The corresponding squared volume is the Gram determinant of these transformed vectors. At a coefficient prime \(\ell\), integrality is tested for the twists whose conductors are prime to \(\ell\). The index is a bound for the determinant valuation of a matrix of integral linear combinations that makes all these lifts integral at once. The upper estimate controls both the length and the index for explicitly constructed vectors. The lower estimate compares the determinants of \(\Theta_1^*\Theta_1\) and \(\Theta_\xi^*\Theta_\xi\) on the entire target \(Y\). Completing the family to one basis lets us compare these determinants with a different twist at each prime while keeping the same family at infinity.

\begin{samepage}
\subsection*{Notation for the volume comparison}
\begin{center}\begin{tabular}{@{}p{.19\linewidth}p{.72\linewidth}@{}}
\(q,d,\chi_d\) & Progression modulus \(q=Q\), auxiliary prime \(d\), and quadratic character \(\chi_d(n)=(n/d)\).\\[3pt]
\(x,x_1,a,b\) & The two cutoffs \(x<x_1\), the derivative order \(a=xb\), and the place weight \(b\) on \(\Phi_b\).\\[3pt]
\(E,\Phi,f\)& CM field, its type, and the number \(f=|\Phi|\) of complex places.\\[3pt]
\(m,k,r,t\)& Number of lifts, initial definite rank, initial target rank, and rank increment on the target chain.\\[3pt]
\(M,N\)& Principal level \(M=p_0p_1\) and level prime \(N\). In Chapter~\ref{ch:analytic}, \(N\) is the modulus of the original Dirichlet character.\\[3pt]
\(S_{\rm ram},Y,D\)& Ramified source space \(S_{\rm ram}\), good space \(Y\), and \(D=\dim_{\mathbb C}Y\).\\[3pt]
\(\Theta_\xi,H,L\)& Composite theta lift \(\Theta_\xi\), main term \(H\), and error bound \(L\) defined below.
\end{tabular}\end{center}
\end{samepage}


\section{The three volume estimates}\label{sec:three-estimates}
Fix admissible data as in \S\ref{sec:data}. The normalized Poincar\'e map \(\mathfrak p:S_{\rm ram}\to Y\) includes the scalar normalization of \S\ref{sec:upper-A} and the projection onto the good space \(Y\). For \(g\in\Gamma\), write \(g\mathfrak p\) for the composition with translation by \(g\) on \(Y\).

We measure the same family at infinity and at the finite primes. For an automorphism \(\sigma\) of \(\mathbb C\), the squared norm \(\|\cdot\|_\sigma^2\) is obtained by applying \(\sigma\) to the coordinates of the untwisted lift and evaluating its rational Gram form. At the finite primes, a logarithmic index bound \(\delta\) means that there are \(e_\ell\ge0\) with \(\sum_\ell e_\ell\log\ell\le\delta\) and the following property. For every field isomorphism \(\iota:\mathbb C\to\overline{\mathbb Q}_\ell\), one can take linear combinations of the family whose lifts are integral for every twist allowed at \(\iota\). The coefficient matrix lies in the valuation ring for \(v_\ell\circ\iota\), and its determinant has \(v_\ell\circ\iota\)-valuation at most \(e_\ell\), where \(v_\ell(\ell)=1\). This matrix may depend on \(\iota\) and must serve all the allowed twists there. Section~\ref{sec:lower-C} gives the local conditions and the pairing.

For \(\tau\in\Phi\) and an integer \(0\le y\le a\), the complex and finite-place estimates for a derivative of order \(y\) involve the quantities
\[
\begin{aligned}
 s_\infty(y)&=\log(y!)+y\log(M/\Omega),\\
 s_{\rm fin}(y,\tau)&=\sum_{1\le n\le a}\Lambda(n)(1-\chi_d(n))
                         (y/n-\kappa_{-n\tau})_+ .
\end{aligned}
\]
Here \(\Omega\) is the positive CM period in \S\ref{sec:upper-A}, \(\chi_d(n)=(n/d)\), and the von Mangoldt function \(\Lambda\) equals \(\log p\) at powers \(p^j\) with \(j\ge1\) and vanishes elsewhere. If \((n,q_1)>1\), the convention is \(\kappa_{-n\tau}=\kappa_\tau\). A prime power \(n=p^j\) in this case has \(p\mid q\), so \(\chi_d(n)=1\) and the term vanishes. For \(p=2\) this follows from \(d\equiv7\bmod8\). For odd \(p\mid q\) it follows from \(d\equiv-1\bmod p\) and quadratic reciprocity, which gives \((p/d)=(-d/p)=1\). The first quantity enters the complex norm estimate. The second records the prime divisibility to be deducted from it. Put \(U_0=fa(\log b)^3\), \(L=U_0\log b\), and set
\[
 H=mk\sum_{\tau\in\Phi}\!
       \bigl(s_\infty(\kappa_\tau)-s_{\rm fin}(\kappa_\tau,\tau)\bigr).
\]
Using the derivative order \(a\) at every \(\tau\in\Phi_a\), define the remainder
\[
 \mathcal R=\sum_{\tau\in\Phi_a}\left[s_{\rm fin}(a,\tau)-a\{\log(a/\bar\kappa)+c_1\}\right].
\]
Here \(c_1=(L'/L)(1,\chi_d)-\gamma-1\), with \(\gamma\) Euler's constant, centers the unrestricted prime sum. The place count in Chapter~\ref{ch:analytic} identifies \(\mathcal R\) with a combination of centered prime sums.

\begin{theorema}[Short vectors and their index]\leavevmode\lettertheoremlabel{A}{thm:short}
The ramified source space \(S_{\rm ram}\) is spanned by integral \(\phi\) for which, simultaneously for all \(g\in\Gamma\) and all \(\sigma\),
\[
       \log\|\Theta_1(g\mathfrak p\phi)\|_\sigma^2
                \le H-m\mathcal R+L
\]
whenever the image is nonzero. For \(T\subset\Gamma\), every finite family of images under \(g\mathfrak p\), \(g\in T\), of vectors in \(S_{\rm ram}\) integral at every finite coefficient embedding has index bound \(2|T|(\dim S_{\rm ram})U_0\).
\end{theorema}
The two conclusions concern the same normalization of the Poincar\'e map. The proof begins in \S\ref{sec:upper-A}.

\begin{theoremb}[Spanning by few translates]\leavevmode\lettertheoremlabel{B}{thm:spanning}
The dimension \(D=\dim Y\) is positive. Some \(T\subset\Gamma\) satisfies
\[
 \dim\sum_{g\in T}g\mathfrak p(S_{\rm ram})\ \ge\ (1-4/b^2)D,\qquad
 |T|\dim S_{\rm ram}\ \le\ 5D\log(1+2b^2).
\]
\end{theoremb}
The first inequality is needed to apply Theorem C, and the second controls the index charged by Theorem A. The proof begins in \S\ref{sec:B}.

\begin{theoremc}[Lower bound for a lifted volume]\leavevmode\lettertheoremlabel{C}{thm:volume}
Let \(y_1,\ldots,y_n\in Y\) be independent, with \(n\ge(1-4/b^2)D\). If \(c\) bounds their logarithmic squared covolume after applying any \(\sigma\in\operatorname{Aut}(\mathbb C)\) to the coordinates of their untwisted lifts, and \(\delta\) bounds their logarithmic index at the finite places, then
\[
                       D(H-3L)\le c+2\delta .
\]
\end{theoremc}
The covolume here is that of the untwisted lifted family, and the index includes the twisted integrality conditions in \S\ref{sec:lower-C}. The factor \(2\) multiplying \(\delta\) comes from the two arguments of a Hermitian pairing. The proof occupies Parts III and IV. We now apply the three estimates to one family.


\clearpage
\section{Comparing the volume bounds}
\label{sec:comparison}
\begin{proposition}[The arithmetic inequality]\label{prop:arithmetic-inequality}
For the admissible data of Proposition~\ref{prop:an-data}, with \(x\) sufficiently large,
\begin{equation}
                       m\mathcal R\le fa(\log b)^5. \label{eq:star}
\end{equation}
\end{proposition}
\begin{proof}
We first show that \(H\ge0\). For integral \(0\le y\le a\) only \(n\le x_1\) contribute to \(s_{\rm fin}(y,\tau)\), because every weight is at least \(a/x_1\). Each summand is bounded by \(2y\Lambda(n)/n\). For an integer \(z\ge2\), use \(\lfloor z/n\rfloor\ge z/(2n)\) for \(n\le z\) and the divisor identity \(\sum_{n\mid j}\Lambda(n)=\log j\). It follows that
\[
 \sum_{n\le z}\frac{\Lambda(n)}{n}
 \le\frac2z\sum_{n\le z}\Lambda(n)\lfloor z/n\rfloor
 =\frac2z\log(z!)\le 2\log z .
\]
Consequently \(0\le s_{\rm fin}(y,\tau)\le4y\log x_1\le y\log b\) for large \(x\). Since \(\Omega=e^{\gamma-\operatorname{Re}(L'/L)(1,\chi_d)}2\pi/\mathrm{rd}_E\), \(M\ge b^{20}\), and \(\mathrm{rd}_E\ge1\), we have \(s_\infty(y)\ge19y\log b\) for large \(x\) with \(d\) fixed. Apply this at \(y=\kappa_\tau\) and sum over \(\tau\).\formalnote{2.6}{6}

By Theorem B, \(D>0\) and there are translates \(T\) whose images of \(S_{\rm ram}\) span at least \((1-4/b^2)D\) dimensions and satisfy \(|T|\dim S_{\rm ram}\le5D\log(1+2b^2)\). Let \(\mathcal S\) be the spanning set of short integral source vectors from Theorem A. The vectors \(g\mathfrak p u\), with \(u\in\mathcal S\) and \(g\in T\), span the same space. Choose a basis \(y_1,\ldots,y_n\) from these images. Then
\[
       (1-4/b^2)D\le n\le D .
\]
At each conjugate, Theorem A bounds the squared norm of the untwisted lift of each basis vector by \(e^{H-m\mathcal R+L}\). Hadamard's inequality gives the logarithmic squared covolume bound \(V=n(H-m\mathcal R+L)\). The index estimate in Theorem A gives \(\delta=DI\) as an index bound for the same basis, where
\[
 I=10\log(1+2b^2)U_0 .
\]
Theorem C applied to this family gives \(V+2\delta\ge D(H-3L)\), so
\[
 D(H-3L)\le n(H+L)-nm\mathcal R+2DI .
\]
Since \(n\le D\) and \(H+L\ge0\), we may bound \(n(H+L)\le D(H+L)\), and \(nm\mathcal R\le D(4L+2I)\). For \(b\) large, \(n\ge D/2>0\). Put \(B_0=\log b\); then \(\log(1+2b^2)\le3B_0\) for \(b\ge3\), so
\[
  2(4L+2I)\le128faB_0^4\le faB_0^5\quad(B_0\ge128).
\]
Dividing by \(n\) gives \eqref{eq:star}.\formalnote{2.7}{7}
\end{proof}


\editionchapter{Prime sums and zeros of Dirichlet L-functions}\label{ch:analytic}

The arithmetic inequality of Chapter~\ref{ch:volume} bounds a sum over prime powers in one residue class. We take its Mellin transform to recover logarithmic derivatives of Dirichlet \(L\)-functions. A zero of one of these functions can produce a pole in the transform. The quadratic character in the prime sum may be chosen after the proposed zero is fixed, so that no twisted \(L\)-function cancels this pole.

\section{The analytic deduction of the zero-free half-plane}\label{sec:analytic}

Fix a positive multiple \(Q\) of \(18480\). For the moment \(d\) is any
prime in the class \(Q/2-1\bmod Q\). This class is prime to \(Q\) since
its representative is odd and is relatively prime to \(Q/2\).
In particular \(d\equiv7\bmod8\). Put
\(\chi_d(n)=(n/d)\), the quadratic character of \(\mathbb Q(\sqrt{-d})\).
We shall work with the moduli \(q=Q\) and \(q=1\), always with this same
\(d\). A subscript \(q\) on a prime sum suppresses its dependence on \(d\).
For each such \(d\), the geometric estimates give a bound at every sufficiently large argument. Given a proposed zero \(\rho\), we will choose \(d\) so that every twisted \(L\)-function is nonzero at \(\rho\).

\subsection{The smoothed prime sum and its Mellin transform}

For \(x\ge1\) define
\[
 \lambda_q(n)=\mathbf1_{n\equiv1\bmod q}\Lambda(n)(1-\chi_d(n)),
 \qquad S_q(x)=\sum_{n\le x}\lambda_q(n)(1/n-1/x).
\]
Here \(\Lambda(p^j)=\log p\) for \(j\ge1\), and \(\Lambda\) vanishes
elsewhere. A split prime contributes zero. At an inert prime only the odd
powers contribute, with factor two. At the ramified prime the factor is
one. In each case only powers congruent to \(1\bmod q\) enter \(S_q\).

\begin{example}[The quadratic weight at the first few prime powers]\label{ex:quadratic-weight}
The definition of \(S_1(x)\) makes sense for every odd prime \(d\). Take \(d=7\) and \(x=12\). The nonzero squares modulo \(7\) are \(1,2,4\). Thus the powers of \(2\) and the prime \(11\equiv4\bmod7\) contribute zero. The primes \(3,5\) are nonsquares, \(3^2\) is a square, and \(7\) is ramified. The three surviving terms give
\[
 S_1(12)=2(\tfrac13-\tfrac1{12})\log3+
         2(\tfrac15-\tfrac1{12})\log5+
           (\tfrac17-\tfrac1{12})\log7
 =\tfrac12\log3+\tfrac7{30}\log5+\tfrac5{84}\log7 .
\qedhere\]
\end{example}

We seek the Mellin transform of \(S_q\) after subtracting its main term. In the Dirichlet series for its coefficients \(\lambda_q(n)\), that main term comes from the pole of the principal \(L\)-function at \(1\). Writing \(h=\varphi(q)\), let \(F_\psi=L(\,\cdot\,,\psi)\) for a
nonprincipal character and let \(F_\psi\) be the analytic continuation of
\((s-1)L(s,\psi)\) for a principal character. Define
\begin{equation}\label{eq:an-log-derivative}
A_q(s)=\frac1h\left[
 \sum_{\psi\bmod q}
 \left(\frac{F'_{\psi\chi_d}}{F_{\psi\chi_d}}(s)
               -\frac{F'_\psi}{F_\psi}(s)\right)-1\right].
\end{equation}
None of the twists \(\psi\chi_d\) is principal. Indeed the Chinese
remainder theorem supplies a unit equal to \(1\bmod q\) whose image
modulo \(d\) is a quadratic nonresidue. Dirichlet's nonvanishing theorem
at \(1\) and the nonzero residue of the principal pole make \(A_q\)
holomorphic there \cite[Thm.~4.9 and (4.29)]{montgomery-vaughan}. Pairing characters with their conjugates
shows that \(c_q=A_q(1)\) is real. The centered error is
\begin{equation}\label{eq:an-centered-error}
R_q(x)=S_q(x)-h^{-1}\log x-c_q.
\end{equation}

To derive the normalization, take \(\Re s>1\). Orthogonality gives
\(\sum_{\psi\bmod q}\psi(n)=h\,\mathbf1_{n\equiv1\bmod q}\), including
the case \((n,q)>1\), when both sides vanish. The Euler product for each
character gives \(-L'/L=\sum_n\psi(n)\Lambda(n)n^{-s}\).
Replacing \(L_\psi\) by \(F_\psi\) adds \(1/(s-1)\) to the logarithmic
derivative for exactly one untwisted character. It follows that
\begin{equation}\label{eq:an-dirichlet-series}
D_q(s):=\sum_n\lambda_q(n)n^{-s}
     = A_q(s)+\frac1h\left(1+\frac1{s-1}\right)
     = A_q(s)+\frac{s}{h(s-1)}.
\end{equation}
The quadratic twist enters the prime weight with
coefficient \(-1\), but enters \(A_q\) with coefficient \(+1\).

For \(n\ge1\) and \(\Re s>1\), direct integration gives
\[
 \int_n^\infty(1/n-1/x)x^{-s}\,dx
       =n^{-s}\left(\frac1{s-1}-\frac1s\right).
\]
The absolute integrals sum to at most a constant times
\(\sum_n\Lambda(n)n^{-\Re s}\), which converges for \(\Re s>1\).
Substitution of \eqref{eq:an-dirichlet-series} gives
\begin{equation}\label{eq:an-mellin-transform}
\begin{aligned}
\Phi_q(s)
   &:=\int_1^\infty R_q(x)x^{-s}\,dx
     =\int_0^\infty R_q(e^v)e^{(1-s)v}\,dv\\
   &=\frac{D_q(s)}{s(s-1)}-\frac1{h(s-1)^2}
                                   -\frac{c_q}{s-1}
     =\frac{A_q(s)/s-c_q}{s-1}.
\end{aligned}
\end{equation}
The right side continues meromorphically to \(\Re s>0\), with a removable
singularity at \(1\). At any other \(\rho\) in this half-plane its residue is
\begin{equation}\label{eq:an-residue}
r_q(\rho)=\frac{n_q(\rho)-n_{q,d}(\rho)}
                 {h\rho(1-\rho)},\qquad
 n_q(\rho)=\sum_{\psi\bmod q}\operatorname{ord}_{\rho} F_\psi,\quad
 n_{q,d}(\rho)=\sum_{\psi\bmod q}\operatorname{ord}_{\rho} F_{\psi\chi_d}.
\end{equation}
Indeed \(F(s)=(s-\rho)^e U(s)\) with \(U(\rho)\ne0\) gives
\(F'/F=e/(s-\rho)+U'/U\). Thus every possible pole is simple,
independently of the multiplicity of the zero. At regular points we put
\(r_q=0\).\formalnote{3.1}{1}

At modulus one the Laurent expansion
\((s-1)\zeta(s)=1+\gamma(s-1)+O((s-1)^2)\) gives
\(c_1=(L'/L)(1,\chi_d)-\gamma-1\) in \eqref{eq:an-log-derivative},
where \(\gamma\) is Euler's constant. In the one-sided estimate below,
\(R_1\) occurs at the auxiliary cutoff \(x_1\), and \(R_Q\) is the
remainder to be bounded at the primary cutoff \(x\).

\paragraph{Comparison with the earlier zeta statistic (optional).}
At \(q=1\), one may write \(xR_1(x)\) as a difference of two centered sums.
Let \(T_\chi(x)=\sum_{n\le x}\Lambda(n)\chi(n)(x/n-1)\). With \(1\) the
character of modulus one and \(\gamma\) Euler's constant, set
\[
G(x)=T_1(x)-x(\log x-1-\gamma),\qquad
G_{\chi_d}(x)=T_{\chi_d}(x)+x(L'/L)(1,\chi_d).
\]
The first is the centered zeta statistic of \S\ref{sec:motivation}.
These are the normalizations in the account prepared for
Rubinstein \cite[eqs.~(I.2) and (II.7)]{rubinstein}. Substitution of the value of \(c_1\) gives
\[
                     xR_1(x)=G(x)-G_{\chi_d}(x).
\]

\subsection{The one-sided estimate at the primary cutoff}
\begin{proposition}[One-sided estimate]\label{prop:an-one-sided}
Fix \(Q\) and \(d\) as above. For each \(\varepsilon>0\), the inequality \(R_Q(y)\le y^{-1/2+\varepsilon}\) holds for all sufficiently large real \(y\).
\end{proposition}
\begin{proof}

By Propositions~\ref{prop:an-data} and \ref{prop:arithmetic-inequality}, at each
sufficiently large \emph{integer} \(x\), for fixed \(Q,d\) and
\(0<e<1/100\), there are admissible data with scales \(x_1,b\), derivative
order \(a=xb\), and \(m\) lifts, such that
\begin{equation}\label{eq:an-scales}
m\ge\sqrt x,\qquad \log b\le2x^e,\qquad
       \log x_1\ge x^{e/12}.
\end{equation}
The places in \(\Phi_a\) are all those with residue label \(1\bmod Q\), where derivatives of order \(a\) are taken. The places in \(\Phi_b\) are the selected ones with residue label \(-1\bmod Q\), where the weight is \(b\). Write \(U=|\Phi_a|\), \(V=|\Phi_b|>0\), and \(f=\varphi(Q)U\). The
cardinality bound is \(U/V\le\log b\). The main geometric inequality
and the identity for its remainder read
\begin{equation}\label{eq:an-geometric-remainder}
\begin{aligned}
m\mathcal R&\le fa(\log b)^5,\\
      \mathcal R&=a\{V[R_Q(x)-R_Q(x_1)]+U R_1(x_1)\}.
\end{aligned}
\end{equation}
To see the origin of the difference, first put weight \(a/x_1\) at every
place. For each of the \(U\) places in \(\Phi_a\), the positive part in the prime sum is then
\((a/n-a/x_1)_+=a(1/n-1/x_1)_+\), independent of the residue of \(n\).
Changing the place weight to \(a/x\) on \(\Phi_b\) replaces precisely \(V\) contributions when \(n\equiv1\bmod Q\), and none otherwise. Indeed, the labels of \(\Phi_a\) reduce to \(1\bmod Q\), and those of \(\Phi_b\) reduce to \(-1\bmod Q\). For \(n\) coprime to \(q_1\), multiplication by \(-n\) permutes the place labels modulo \(q_1\). Its preimage of each place in \(\Phi_b\) lies in \(\Phi_a\) exactly when \(n\equiv1\bmod Q\). Every prime dividing \(q_1\) already divides \(Q\).
For an odd prime \(p\mid Q\), we have \(-d\equiv1\bmod p\), so
quadratic reciprocity gives \(\chi_d(p)=(-d/p)=1\).
Also \(\chi_d(2)=1\) since \(d\equiv7\bmod8\). Thus every prime
dividing \(q_1\) has \(1-\chi_d(p^j)=0\) at each of its powers \(p^j\). The sum before centering
is \(a\{US_1(x_1)+V[S_Q(x)-S_Q(x_1)]\}\). The two weights and
\(f=\varphi(Q)U\) give
\[
 U\log(a/\bar\kappa)
     = U\log x_1 + \frac{V}{\varphi(Q)}\log(x/x_1).
\]
This gives the logarithms in the centering term. The two occurrences
of \(c_Q\) cancel, and the remaining constant is \(Uc_1\).
Solving for \(R_Q(x)\) and using the cardinality bound gives
\begin{equation}\label{eq:an-three-errors}
R_Q(x)\le\frac{\varphi(Q)(\log b)^6}{m}
           +|R_Q(x_1)|+\log b\,|R_1(x_1)|.
\end{equation}
We bound the two auxiliary terms using the prime number theorem in fixed progressions.\formalnote{3.2}{2}

\paragraph{The auxiliary cutoff.}

For either modulus \(q\) put
\(C_q(t)=\sum_{n\le t}\lambda_q(n)-t/\varphi(q)\).
The equality \(1/n-1/y=\int_n^y t^{-2}\,dt\) and a finite interchange give
\begin{equation}\label{eq:an-tail-integral}
R_q(y)=\int_1^y C_q(t)t^{-2}\,dt-c_q.
\end{equation}
Let \(\Psi(t)=\sum_{n\le t}\Lambda(n)\) be Chebyshev's function. For an integer \(u\ge1\), Legendre's formula for \(\log\binom{2u}{u}\) counts each prime power \(u<p^j\le2u\) with coefficient \(\log p\). All other contributions are nonnegative. Thus \(\Psi(2u)-\Psi(u)\le2u\log 2\). Summing at successive powers of two up to the first one at least \(t\) gives \(\Psi(t)\le4t\log 2\le6t\) for \(t\ge1\). Since \(0\le\lambda_q(n)\le2\Lambda(n)\), it follows
that \(|C_q(t)|\le13t\). In particular
\begin{equation}\label{eq:an-continuity}
|R_q(z)-R_q(y)|\le13\log(z/y)\le13(z-y)/y\quad(z\ge y\ge1).
\end{equation}
In particular, \(v\mapsto R_q(e^v)\) is Lipschitz on \(v\ge0\).
The fixed-modulus prime number theorem needed here is
\begin{equation}\label{eq:an-fixed-modulus-pnt}
\text{for every integer }B\ge0,\qquad
       |R_q(y)|\le(\log y)^{-B}\quad\text{for all sufficiently large }y.
\end{equation}
Apply Siegel--Walfisz to the fixed modulus \(qd\) \cite[Cor.~11.19]{montgomery-vaughan}. Among the reduced classes with \(a\equiv1\bmod q\), the sum of \(1-\chi_d(a)\) is \(d-1\) by character orthogonality modulo \(d\). Their main term is \(t/\varphi(q)\). Of the remaining prime powers, only powers of \(d\) can contribute, with total weight at most \(\log t\). Thus \(C_q(t)\ll_{q,d,B}t(\log t)^{-B-2}\) for \(t\ge2\). The centering constant in \eqref{eq:an-tail-integral} is important. By partial summation and \eqref{eq:an-dirichlet-series},
\[
  \int_1^\infty C_q(t)t^{-s-1}\,dt
       =D_q(s)/s-\frac{1}{\varphi(q)(s-1)}=A_q(s)/s \qquad(s>1).
\]
Dominated convergence at \(s=1\) gives \(c_q=\int_1^\infty C_q(t)t^{-2}\,dt\). Subtracting this constant leaves the tail
\[
  R_q(y)=-\int_y^\infty C_q(t)t^{-2}\,dt
          =O_{q,d,B}\bigl((\log y)^{-B-1}\bigr).
\]
The extra power of the logarithm absorbs the implied constant and gives \eqref{eq:an-fixed-modulus-pnt}.
\phantomsection\label{route:2}\leanroute{The formal proof obtains a zero-free region of logarithmic width and applies a quantitative Tauberian argument with a gamma kernel. Siegel--Walfisz supplies exactly the required error after the partial summation above. Both \(q\) and \(d\) are fixed when \(y=x_1\) tends to infinity.\formalnote{3.3}{3}}

In \eqref{eq:an-scales} choose an integer \(B\ge24/e\). For \(q=1,Q\), \eqref{eq:an-fixed-modulus-pnt} yields
\(|R_q(x_1)|\le(\log x_1)^{-B}\le x^{-2}\).
Using this in \eqref{eq:an-three-errors} gives, for sufficiently large integers \(x\),
\begin{equation}\label{eq:an-integer-bound}
R_Q(x)\le 64\varphi(Q)x^{-1/2+6e}
            +(1+2x^e)x^{-2}
       \le (64\varphi(Q)+3)x^{-1/2+6e}.
\end{equation}
For a real \(y\ge2\) put \(x=\lfloor y\rfloor\). By \eqref{eq:an-continuity} the floor
costs at most \(26/y\). Given \(\varepsilon>0\), choose
\(0<e<1/100\) with \(7e\le\varepsilon\). Since \(x\ge y/2\) and
\(-1\le-1/2+6e\le0\), \(x^{-1/2+6e}\le2y^{-1/2+6e}\).
The constant in \eqref{eq:an-integer-bound} and the floor error are absorbed in the spare
factor \(y^e\). We have proved
\begin{equation}\label{eq:an-one-sided-bound}
R_Q(y)\le y^{-1/2+\varepsilon}
               \quad\text{for all sufficiently large real }y.
\end{equation}
The threshold may depend on \(d,Q,\varepsilon\).
\end{proof}

\subsection{A theta-series criterion for a zero}

For \(\beta>1/2\) and \(H_0\ge0\) let
\(\mathcal B(\beta,H_0)=\{s:\Re s\ge\beta,\ |\Im s|\le H_0\}\).
Fix \(1/2<\beta<1\) and \(H_0\ge0\). We seek one prime \(d\equiv Q/2-1\bmod Q\) for which all the twists \(L(s,\psi\chi_d)\), \(\psi\bmod Q\), are zero-free in \(\mathcal B(\beta,H_0)\). To detect a zero of a character \(L\)-function, we use two ordinary theta series, one formed from its Dirichlet coefficients and one from the coefficients of its reciprocal. The lemma below shows that a zero forces \(J_j(\chi,T)\), the sum of the absolute-value tail integrals of the two series, to be large. The next subsection bounds the mean of \(J_j\) as the twisting prime \(d\) varies.

For a sequence \(|c_n|\le1\) and \(j=0,1\) define
\[
 w_j(z)=z^j e^{-\pi z^2},\qquad
 \theta_{c,j}(u)=u^{1/2}\sum_{n\ge1}c_n w_j(nu),\qquad
 \Gamma_j(s)=\tfrac12\pi^{-(s+j)/2}
                  \Gamma\bigl((s+j)/2\bigr).
\]
Substitution \(z=\pi u^2\) in
\(\int_0^\infty u^{s+j}e^{-\pi u^2}\,du/u\) gives \(\Gamma_j(s)\).
Thus on \(\Re s>1\),
termwise integration followed by \(t=nu\) gives
\begin{equation}\label{eq:an-theta-mellin}
\int_0^\infty\theta_{c,j}(u)u^{s-1/2}\,\frac{du}{u}
                   =\Gamma_j(s)\sum_{n\ge1}c_n n^{-s}.
\end{equation}
Let \(\chi\) be primitive and nonprincipal of conductor \(f\), and choose the parity \(j\) by \(\chi(-1)=(-1)^j\). The classical theta transformation, in the present scaling, is
\[
 \theta_{\chi,j}(u)=(-i)^j\frac{\tau(\chi)}{\sqrt f}
                  \theta_{\bar\chi,j}(1/(fu)),\qquad |\tau(\chi)|=\sqrt f .
\]
To compare normalizations in \cite[Thms.~9.7 and 10.6]{montgomery-vaughan}, substitute \(z=fu^2\) in its bilateral series and multiply by \(u^{1/2+j}/2\). Since \(w_j\) is real, we obtain
\begin{equation}\label{eq:an-theta-symmetry}
|\theta_{\chi,j}(u)|=|\theta_{\chi,j}(1/(fu))|.
\end{equation}
Gaussian decay at the two ends now continues the Mellin integral \eqref{eq:an-theta-mellin} for \(c=\chi\) to every bounded vertical strip, by dominated convergence on compact subsets.
\phantomsection\label{route:3}\leanroute{The code derives the theta transformation in absolute value from Poisson summation. The cited classical transformation also gives its phase. Taking absolute values yields the same identity used here.\formalnote{3.4}{4}}

For \(T\ge1\), let \(J_j(c,T)\) be the sum of absolute-value tail integrals of \(\theta_{c,j}\) and \(\theta_{\mu c,j}\), where \((\mu c)_n=\mu(n)c_n\) and \(\mu\) is the M\"obius function. Let \(\mathcal M(s)\) be the truncated Mellin integral of \(\theta_{\mu\chi,j}\). Explicitly,
\begin{align*}
 J_j(c,T)&=\int_{1/T}^\infty
          (|\theta_{c,j}(u)|+|\theta_{\mu c,j}(u)|)\frac{du}{u},\\
 \mathcal M(s)&=\int_{1/T}^\infty\theta_{\mu\chi,j}(u)
                          u^{s-1/2}\frac{du}{u}.
\end{align*}
\begin{lemma}[A zero forces a large theta tail]\label{lem:an-detector}
Fix \(0<\delta<1/18\) and \(H_0\ge0\). For every sufficiently large \(T\), every primitive nonprincipal \(\chi\) of conductor \(f\le T^2\) having a zero in \(\mathcal B(1/2+9\delta,H_0)\) satisfies \(J_j(\chi,T)>T^{2\delta}\), where \(j\) is the parity of \(\chi\).
\end{lemma}
\begin{proof}
The full integral for \(c=\mu\chi\) equals \(\Gamma_j(s)/L(s,\chi)\)
for \(\Re s>1\), by absolute convergence of the reciprocal Euler product.
We estimate \(\mathcal M(s)\) and \(\Gamma_j(s)L(s,\chi)\) on the two boundaries of
\(1/2\le\Re s\le2\). On the line \(\Re s=2\) we have
\[
 \left|\Gamma_j(s)/L(s,\chi)-\mathcal M(s)\right|\le3/T,
       \qquad |\Gamma_j(s)L(s,\chi)|\le3 .
\]
Indeed \( |w_j(t)|\le3e^{-2t}\) for \(t>0\), and
\(\sum_{n\ge1}e^{-2nu}=e^{-u}/(2\sinh u)\le e^{-u}/(2u)\).
The absolute integrand in \eqref{eq:an-theta-mellin}, when measured against \(du\), is
therefore at most \(3e^{-u}\) on this line. Integrate over \(0<u<1/T\)
for the first claim and \(0<u<\infty\) for the second.

On \(\Re s=1/2\), \(|\mathcal M(s)|\le J_j(\chi,T)\) directly. If
\(f\le T^2\), fold the integral for \(\Gamma_j L\) at \(u=f^{-1/2}\)
using \eqref{eq:an-theta-symmetry}. Since \(f^{-1/2}\ge1/T\), this gives
\[
                  |\Gamma_j(s)L(s,\chi)|\le2J_j(\chi,T).
\]
Suppose \(J_j(\chi,T)\le T^{2\delta}\) and set
\[
          P(s)=\Gamma_j(s)L(s,\chi)\mathcal M(s)-\Gamma_j(s)^2.
\]
The preceding estimates give \(|P(s)|\ll_j T^{4\delta}\) on \(\Re s=1/2\) and \(|P(s)|\le9/T\) on \(\Re s=2\). The functions \(\Gamma_j L\), \(\mathcal M\) and \(\Gamma_j\) defining \(P\) are bounded on this closed strip for fixed \(T,\chi\). Indeed,
one uses \eqref{eq:an-theta-mellin}--\eqref{eq:an-theta-symmetry} for \(\Gamma_j L\), the cutoff and Gaussian
decay for \(\mathcal M\), and the Mellin integral for \(\Gamma_j\).
They are holomorphic in its interior. The constants in the two
\emph{boundary estimates} are uniform in \(T,\chi\).
Hadamard's three-lines theorem, with
\(\alpha=(\Re s-1/2)/(3/2)\), gives
\[
              |P(s)|\ll_j T^{4\delta(1-\alpha)-\alpha}.
\]
At a zero with \(\Re s\ge1/2+9\delta\) one has \(\alpha\ge6\delta\),
so the exponent is at most \(-2\delta\), whereas
\(P(s)=-\Gamma_j(s)^2\).
The Euler product and the classical zero-free line exclude zeros with
\(\Re s\ge1\) \cite[Thms.~4.9, 6.6 and 11.3]{montgomery-vaughan}.
On the compact rectangle \(1/2\le\Re s\le1,\ |\Im s|\le H_0\)
the continuous function \(|\Gamma_j|\) has a positive minimum.
Consequently, for all sufficiently large \(T\) depending only on
\(\delta,H_0\) and the parity, every primitive nonprincipal
\(\chi\) of conductor at most \(T^2\) with a zero in
\(\mathcal B(1/2+9\delta,H_0)\) satisfies
\begin{equation}\label{eq:an-zero-detector}
J_j(\chi,T)>T^{2\delta}.
\end{equation}
The lower bound is uniform in the character with the height box fixed.\formalnote{3.5}{5}
\end{proof}

\subsection{A prime giving simultaneous nonvanishing}
Jutila obtains a zero-density estimate for real characters by bounding mean squares of Dirichlet polynomials \cite[Thms.~1--2]{jutila}. Here the twists of each fixed character modulo \(Q\) must be handled, including when that character is complex. The large sieve used below allows its values to be put into the coefficients.
\begin{proposition}[Simultaneous nonvanishing]\label{prop:an-good-prime}
For a positive multiple \(Q\) of \(18480\), a real number \(1/2<\beta<1\) and a height \(H_0\ge0\), there exists a prime \(d\equiv Q/2-1\bmod Q\) such that \(L(s,\psi\chi_d)\ne0\) throughout \(\mathcal B(\beta,H_0)\) for every character \(\psi\bmod Q\), where \(\chi_d(n)=(n/d)\).
\end{proposition}
\begin{proof}
Put \(\delta=(\beta-1/2)/9\).
Let \(\xi\) be a fixed primitive character of conductor \(M\mid Q\), including the trivial character of conductor one. The quadratic large sieve of Heath-Brown gives, for every \(\eta>0\), all complex coefficients \(a_1,\ldots,a_N\) and positive integers \(D,N\),
\[
 \sum_{\lambda\in\mathcal X_{\rm real}(D)}
     \left|\sum_{n\le N}a_n\lambda(n)\right|^2
   \ll_\eta (DN)^\eta (D+N)N\max_{n\le N}|a_n|^2 ,
\]
where the sum is over primitive real characters of conductor at most \(D\) \cite[Cor.~3, p.~238]{heath-brown-sieve}. Each \(\chi_d=(\cdot/d)\) for an odd prime \(d\le D\) is one such character. The fixed character \(\xi\) is absorbed into \(a_n\), which is why this bound applies also to complex \(\xi\).

Here are the cutoff and the loss in this application. Set \(D=\lfloor T^2\rfloor\), \(T\ge 3\). For \(1/T\le u\le1\) truncate the theta sum at \(N=\lceil u^{-1}\log T\rceil\). For any fixed \(A>0\) its tail is \(O_{A,j}(T^{-A})\), uniformly in the characters, since the absolute tail is \(O_j(u^{-1/2}e^{-c(\log T)^2})\). Use \(a_n=\xi(n)w_j(nu)\) and then \(a_n=\mu(n)\xi(n)w_j(nu)\). Both satisfy the coefficient bound. For either resulting theta sum, the factor \(\sqrt u\) contributes \(u\) to its square. We have \(uN\ll\log T\), \(N\ll T\log T\), \(D+N\ll T^2\) and \(DN\ll T^4\). Thus the sum of the squares over the prime conductors \(d\le T^2\) is \(O_{\eta,j}(T^{2+4\eta}\log T)\). Cauchy--Schwarz bounds the sum of the absolute values by \(O_{\eta,j}(T^{2+2\eta}\sqrt{\log T})\). Integration against \(du/u\) costs \(\log T\). For \(u\ge1\) the Gaussian gives an integrable majorant independent of \(d\), so that part costs \(O_j(T^2)\). Taking \(\eta=\delta/4\) and absorbing the logarithms gives
\begin{equation}\label{eq:an-large-sieve-mean}
\sum_{\substack{d\le T^2\\ d\ {\rm odd\ prime}}}
          J_j(\xi\chi_d,T) \ll_{\delta,j} T^{2+\delta}.
\end{equation}
\phantomsection\label{route:4}\leanroute{The development extends the twisting parameter from primes to every positive integer, uses quadratic reciprocity, and estimates the off-diagonal Jacobi sums by P\'olya--Vinogradov. Heath-Brown's inequality allows arbitrary complex coefficients and all positive indices \(n\). It gives the prime-conductor first moment \eqref{eq:an-large-sieve-mean} directly, with the same power needed by the zero detector.\formalnote{3.6}{6}}

For \(d\equiv Q/2-1\bmod Q\), the product \(\xi\chi_d\) is primitive of conductor \(Md\) by coprimality of the conductors. Its parity is opposite to that of \(\xi\), since \(d\equiv7\bmod8\). Thus \(j\) is fixed over the primes in this progression. Each such prime \(d\) with \(Md\le T^2\) for which \(L(s,\xi\chi_d)\) vanishes in \(\mathcal B(\beta,H_0)\) contributes more than \(T^{2\delta}\) to \eqref{eq:an-large-sieve-mean}, by \eqref{eq:an-zero-detector}. There are \(O(T^{2-\delta})\) such primes. A character modulo \(Q\) is induced by a primitive \(\xi\) of conductor \(M\mid Q\). The Euler factors removed in the induced twist are nonzero for \(\Re s>0\). Taking \(T=\sqrt{QX}\) and summing over the finite list of primitive characters gives
\begin{equation}\label{eq:an-bad-primes}
\#\left\{ d\le X:
    \begin{array}{l}
     d\text{ prime},\ d\equiv Q/2-1\bmod Q,\\
     L(s,\psi\chi_d)=0 \text{ for some }\psi\bmod Q
                     \text{ and }s\in\mathcal B(1/2+9\delta,H_0)
    \end{array}\right\}
 \ll_{Q,H_0,\delta} X^{1-\delta/2}.
\end{equation}
The fixed-modulus prime number theorem gives \(\sim X/(\varphi(Q)\log X)\) primes in this reduced class \cite[Cor.~11.21]{montgomery-vaughan}. This exceeds \eqref{eq:an-bad-primes} for all sufficiently large \(X\). A single surviving prime \(d\) gives simultaneous nonvanishing for every \(\psi\bmod Q\) throughout the prescribed box.
\phantomsection\label{route:5}\leanroute{Lean completes this selection by comparing sums of \((\log p)/p\). The sum over primes in the chosen progression diverges, while the sum over bad primes converges. Here the comparison of the two counting functions suffices.}

\end{proof}
\subsection{Landau's theorem and the sign of a real residue}

For a nonnegative integrand, the real boundary point of the region of
convergence plays a special role, even if the pole one wishes to exclude
is non-real.

\begin{example}[Real and non-real poles of a positive transform]\label{ex:mellin-positive}
For \(\alpha\in\mathbb R\) and \(t\in\mathbb R\setminus\{0\}\), consider
\(x^{\alpha-1}(1+\cos(t\log x))\) on \(x\ge1\). It is nonnegative. On
\(\Re s>\alpha\) its transform against \(x^{-s}\,dx\) is
\[
 \frac1{s-\alpha}+
 \frac1{2(s-\alpha-it)}+\frac1{2(s-\alpha+it)}.
\]
At \(s=\alpha\), put \(v=\log x\). The integral becomes
\(\int_0^\infty(1+\cos(tv))\,dv\), whose integral over each period
is \(2\pi/|t|\). The abscissa of absolute convergence is therefore
\(\alpha\). The oscillation produces the two non-real poles, and the
constant term gives the real pole of residue \(1\).
\end{example}

Fix \(1/2<\beta<1\) and \(H_0\ge0\). Take \(d\) as in Proposition~\ref{prop:an-good-prime}, so that every \(L(s,\psi\chi_d)\) is zero-free in \(\mathcal B(\beta,H_0)\). Assume \(R_Q(y)\le y^{\beta-1}\) for all sufficiently large \(y\). This is the hypothesis needed from the one-sided estimate at \(\varepsilon=\beta-1/2\).

Put \(A_\beta(x)=x^{\beta-1}-R_Q(x)\) for \(x\ge1\). This function is continuous and eventually nonnegative. Choose \(X_0\ge1\) so that \(A_\beta(x)\ge0\) for \(x\ge X_0\). For \(\Re s>1\) its Mellin integral is
\begin{equation}\label{eq:an-landau-transform}
\mathcal L_\beta(s)=\int_1^\infty A_\beta(x)x^{-s}\,dx
                  =\frac{1}{s-\beta}-\Phi_Q(s).
\end{equation}
Let \(\sigma_c\) be the abscissa of convergence of the integral over \([X_0,\infty)\). The formula for \(R_Q\) gives \(R_Q(x)=O(1+\log x)\), so \(\sigma_c\le1\). The integral over \([1,X_0]\) is entire. Landau's theorem for Mellin integrals makes \(\mathcal L_\beta\) holomorphic on \(\Re s>\sigma_c\) and singular at any finite \(\sigma_c\) \cite[Lemma~15.1]{montgomery-vaughan}. The integral is entire if \(\sigma_c=-\infty\). On the common half-plane it agrees with the meromorphic continuation in \eqref{eq:an-landau-transform}.

\phantomsection\label{route:6}\leanroute{Lean derives this singularity theorem by expanding the transform about a real point just right of the abscissa and interchanging the positive Taylor terms with the integral. The Mellin form of Landau's theorem supplies the same boundary singularity. In the Laplace--Stieltjes notation of \cite[Ch.~II, \S1]{widder-laplace}, the same tail is \(\int_0^\infty e^{-sv}\,dF(v)\), where \(F(v)=0\) for \(0\le v\le\log X_0\) and \(F(v)=\int_{X_0}^{e^v}A_\beta(x)\,dx\) thereafter. Nonnegativity of \(A_\beta\) on \([X_0,\infty)\) is exactly what makes this \(F\) nondecreasing.\formalnote{3.7}{7}}

If \(\sigma_c>\beta\), \eqref{eq:an-landau-transform} makes the singularity a simple pole. The singularity of \(\Phi_Q\) at \(1\) is removable, and the Euler product excludes zeros of the \(L\)-functions to the right of \(1\). Thus \(\beta<\sigma_c<1\). This real point lies in \(\mathcal B(\beta,H_0)\), where every twisted \(L\)-function is zero-free. By \eqref{eq:an-residue} the residue of \(\Phi_Q\) at \(\sigma_c\) is \(n_Q(\sigma_c)/(\varphi(Q)\sigma_c(1-\sigma_c))\ge0\). Hence the residue of \(\mathcal L_\beta\) there is nonpositive. For real \(\sigma>\sigma_c\), the integral over \([X_0,\infty)\) is nonnegative. On multiplying by \(\sigma-\sigma_c\) and letting \(\sigma\downarrow\sigma_c\), the integral over \([1,X_0]\) contributes zero. The residue is therefore also nonnegative and must be zero. This contradicts the simple pole. Therefore \(\sigma_c\le\beta\) and \eqref{eq:an-landau-transform} makes \(\Phi_Q\) holomorphic on \(\Re s>\beta\).\formalnote{3.8}{8}

\begin{samepage}
\subsection{The zero-free criterion and the functional equation}
\begin{proposition}[One-sided zero-free criterion]\label{prop:an-criterion}
Fix \(N\ge1\) and put \(Q=18480N\). Suppose that for every prime \(d\equiv Q/2-1\bmod Q\) and every \(\varepsilon>0\) the centered sum \(R_{Q,d}(y)\) in \eqref{eq:an-centered-error} is at most \(y^{-1/2+\varepsilon}\) for all sufficiently large \(y\). For every Dirichlet character \(\chi\bmod N\), \(L(s,\chi)\) has no zero in \(\Re s>1/2\), with the principal pole excepted.
\end{proposition}
\end{samepage}
Suzuki proves a related criterion for logarithmically weighted sums over \(n\equiv1\bmod q\). For \(1/2\le\beta<1\), he assumes that the \(L\)-functions of all characters modulo \(q\) have no real zero above \(\beta\). Eventual constancy of sign then gives nonvanishing throughout \(\Re s>\beta\) \cite[Thm.~6(1)]{suzuki}. Here the choice of the quadratic twist after the proposed zero is fixed gives the required residue sign at a real pole.
\begin{proof}

Let \(\chi\) be a character modulo \(N\).
Suppose \(L(\rho,\chi)=0\) with \(\Re\rho>1/2\) away from a principal
pole. The case \(\Re\rho\ge1\) is already excluded by the Euler product
and the classical nonvanishing theorem on the line \(1\), with the pole
excepted. Thus assume \(\Re\rho<1\). Set \(Q=18480N\),
choose \(1/2<\beta<\Re\rho\), and take \(H_0=|\Im\rho|\).
Raise \(\chi\) to level \(Q\). Its \(L\)-function is multiplied by a
finite product of Euler factors, so the zero persists and
\(n_Q(\rho)\ge1\). For \(d\) supplied by Proposition~\ref{prop:an-good-prime} at these \(Q,\beta,H_0\),
simultaneous nonvanishing gives \(n_{Q,d}(\rho)=0\). Formula \eqref{eq:an-residue} therefore gives
\[
                 r_Q(\rho)=\frac{n_Q(\rho)}
                             {\varphi(Q)\rho(1-\rho)}\ne0.
\]
The assumed bound at \(\varepsilon=\beta-1/2\) applies to this \(d\). The preceding Landau argument gives holomorphy of \(\Phi_Q\) on \(\Re s>\beta\), hence \(r_Q(\rho)=0\). This contradiction proves the
right-half-plane result for every character. Taking \(H_0=0\)
is permitted, so the same proof rules out a real zero in this range.\formalnote{3.9}{9}
\end{proof}

\begin{proof}[Conclusion of the proof of Dirichlet GRH]
Proposition~\ref{prop:an-one-sided} supplies the hypothesis of Proposition~\ref{prop:an-criterion} for every \(N\).

If \(\chi^*\) is the primitive character inducing \(\chi\), the
imprimitive Euler factors have modulus bounded away from zero at each
point of \(\Re s>0\), since \(|\chi^*(p)p^{-s}|<1\). The gamma and
conductor factors of the completed primitive \(L\)-function are finite
and nonzero inside \(0<\Re s<1\). The functional equation \cite[Corollaries~10.3 and~10.8]{montgomery-vaughan} takes a zero
at \(\rho\) to a zero for the conjugate character at \(1-\rho\).
Both must have real part at most \(1/2\). It follows that
\(\Re\rho=1/2\), for primitive and imprimitive characters alike.
For \(\zeta\), the functional equation reflects the nonvanishing on
\(\Re s\ge1\) to \(\Re s\le0\), except at the zeros
\(-2,-4,\ldots\) from the gamma factor. At the exceptional point \(0\)
one has \(\zeta(0)=-1/2\)
\cite[Corollaries~10.3--10.4 and (10.11)]{montgomery-vaughan}.
This proves the zeta assertion as well as the half-plane and strip
statements for Dirichlet \(L\)-functions.\formalnote{3.10}{10}
\end{proof}


\partheading{Short vectors and spanning by translates}

\editionchapter{Divisibility and complex multiplication}\label{ch:upper}

Theorem A requires vectors that are both integral and small in the complex norm. We construct them from derivatives of theta functions attached to elliptic curves with complex multiplication. If a derivative is divisible by a power of \(p\), dividing by that power reduces its complex size while preserving integrality at \(p\). The calculation below gives the power available at each prime and explains the dependence on the two places \(\tau\) and \(-p^e\tau\) in the prime sum.

\section{Proof of Theorem A}\label{sec:upper-A}

Fix admissible data as in \S\ref{sec:data}. The imaginary quadratic field is
\(K=\mathbb Q(\sqrt{-d})\), and \(E=K(\zeta_{q_1})\) has \(f\) infinite places,
indexed by \(\Phi\simeq(\mathbb Z/q_1)^\times\). For this and the next chapter write \(q_0=q\) for the progression modulus of \S\ref{sec:data}. As in that construction, \(a=xb\) is the derivative order at the places in \(\Phi_a\). The set
\(\Phi_a\) consists of the places over \(1\bmod q_0\), while \(\Phi_b\) is
a chosen subset over \(-1\bmod q_0\). The place weight \(\kappa_\tau\) is \(a/x=b\)
on \(\Phi_b\) and \(a/x_1\) elsewhere. There are \(m\) lifts and the source
rank is \(k\). Write
\[
 U=fa(\log b)^3,\qquad s_\infty(y)=\log(y!)+y\log(M/\Omega)
 \quad(y\in\mathbb Z_{\ge0}),
\]
where \(M=p_0p_1\) is the product of the norm primes and \(\Omega\) is the positive CM period. The level prime \(N\) controls the norm of the final restriction to \(Y_m\).

The two assertions of Theorem A concern the same normalized images.
Dividing a form by \(p^g\) divides its squared complex norm by \(p^{2g}\).
To preserve integrality we must first establish \(p^g\)-divisibility
in the lattice of forms. At \(p\) the required divisibility comes from
theta derivatives.
The finite-place calculation must bound the index needed when this
normalization is applied to all the images in a chosen family. In particular,
the same correction at a given coefficient embedding must work for all the
twists allowed there. Throughout, \(v_p(p)=1\) and \(t_+=\max(t,0)\).

\subsection{The formal-group calculation behind a local term}

Fix an embedding into \(\mathbb C_p\) and a nonsplit \(p\). Put \(q_p=p^2\)
at an inert prime and \(q_p=p\) at \(p=d\). Take
\(\pi_p=p\) in the first case and \(\pi_p=\sqrt{-d}\) in the second, and set
\[
 v=v_p(\pi_p),\qquad \rho=\frac{v}{q_p-1},\qquad
 \gamma_p=\frac1{p-1}-\rho .
\]
The expansion of the theta series will contain products of a coordinate at one place, raised to \(p^e\), and a coordinate at a second place. The indices \(e\) of the Frobenius powers \(p^e\) form a set \(I_p^\circ\). It comprises all positive odd
integers at an inert prime and all nonnegative integers at \(p=d\). Write
\(I_p=I_p^\circ\cap\mathbb Z_{>0}\).
For a fixed \(e\), call \(\tau'=-p^e\tau\bmod q_1\) the partner of \(\tau\). Every nonsplit \(p\) is prime to \(q_1\), so this multiplication is defined on the unit labels. In the prime sum over all \(n\), we interpret \(-n\tau\) as \(\tau\) when \(n\) is a nonunit modulo \(q_1\).
A geometric series gives
\[
                     \gamma_p=v\sum_{e\in I_p}p^{-e}.
\]
For a derivative of total order \(B\), the term \(B\gamma_p\) will therefore be a sum of contributions \(vB/p^e\). From each we must deduct the cost of a change of variables that makes the corresponding pairing divisible by \(\pi_p\).

Let \(F\) be the Taylor series at zero of a theta function on a standard positive definite Hermitian space, with standard level-one data. Such data are
integral linear combinations at \(p\) of translates of the integral vacuum
by points integral above \(p\). The variables in \(F\) have been multiplied by \(c\,a_{\rm jet}\), where \(c\) is integral at \(p\) and \(a_{\rm jet}=P_0 A_0/\Omega_0\) is the Taylor scale of Chapter~\ref{ch:upperlift}. Here \(P_0\) is its period factor, \(A_0=\sqrt{2\pi/\sqrt d}\), and \(\Omega_0\mathcal O_K\) is the period lattice of the normalized elliptic curve. There exist
\(\varpi\) with \(v_p(\varpi)=\rho\), integral coefficients \(\lambda_{ej}\)
and an integral formal series \(W\) such that
\begin{align}
 F(\varpi z)
   &=W(z)\exp\left(\frac{\varpi^{p+1}}p
            \sum_{e\in I_p^\circ}\sum_{j\ge0}
                           \lambda_{ej}\Gamma_e(z)^{p^j}\right),\label{eq:Athetaform}\\
 \Gamma_e(z)&=\sum_{\tau,i}z_{\tau,i}^{\,p^e}z_{-p^e\tau,i}. \nonumber
\end{align}
The sum makes sense degree by degree. The proof of \eqref{eq:Athetaform} uses the one-curve CM theta expansion, then a \(p\)-integral resolution
\(\sum_j m_j u_j u_j^\dagger=1\). Substituting \(u_j^\dagger z\) into the
one-curve factors and applying the Frobenius congruence reduces the sum of
their exponents, modulo \(\pi_p\), to the couplings \(\Gamma_e\). The
difference in the exponent contributes an integral factor to \(W\).
Section~\ref{sec:frame-proof} constructs the resolution. The size estimate also requires the archimedean bound for the derivatives.\formalnote{4.1}{1}

\subsection{A derivative along a product of minors}

For a homogeneous polynomial \(P\) write \(J_P(F)=(P(\partial)F)(0)\).
If \(F=\sum_\alpha F_\alpha z^\alpha\), then
\(J_P(F)=\sum_\alpha P_\alpha\alpha!F_\alpha\).
For integral \(P_\alpha,F_\alpha\) of fixed degree \(B\), Legendre's formula
gives a lower bound \(B/(p-1)-D(1+\lfloor\log_p B\rfloor)\) on its
valuation for \(B\ge1\), where \(D\) counts variables. Indeed
\(v_p(n!)=(n-s_p^{\rm dig}(n))/(p-1)\), and the digit sum of each nonzero
index is at most \((p-1)(1+\lfloor\log_p B\rfloor)\).
It remains to obtain integral degree-\(B\) coefficients in enough
directions while keeping track of the cost of the change of variables.

An entry of the Schur derivative is controlled by derivatives along products
of minors. A filling specifies rows for each column of the Young diagram;
a column of height \(h\) takes a minor on these rows and the first \(h\)
columns of \(z_\tau y_\tau\), for integral directions \(y_\tau\). For a rank
bound \(R\), only column heights \(h<R\) enter. At \(\tau\) let \(t_\tau\)
count the height-one columns, \(c_\tau\) the other columns, and \(b_\tau\)
the boxes of the other columns. Thus \(B=\sum_\tau(t_\tau+b_\tau)\).
For a partner pair in layer \(e\) define the deduction
\[
 D_e(\tau)=
 \min\{t_\tau/p^e,t_{\tau'}+c_{\tau'}\}
       +c_\tau/p^e+\min\{b_\tau/p^e,b_{\tau'}\}.
\]
\begin{example}[A divisibility condition on a single monomial]\label{ex:frobenius-monomial}
Consider two integral variables \(X,Y\) over \(\mathbb C_p\) and a monomial \(X^{p^e}Y\), \(e\ge1\), from \(\Gamma_e\). We want to make this monomial divisible by \(\pi_p\) while losing as little as possible in a derivative of fixed degree in each direction. Choose \(\alpha\) with
\(\alpha^{p^e}=\pi_p\). Either change of variable
\(X\mapsto\alpha X\) or \(Y\mapsto\pi_pY\) makes the term divisible
by \(\pi_p\). A product of minors using the first direction \(t\) times and the second \(t'\) times is multiplied in the first case by \(\alpha^t\), in the second by \(\pi_p^{t'}\). Their valuations are \(vt/p^e\) and \(vt'\), where \(v=v_p(\pi_p)\). Choosing the cheaper
of the two explains the appearance of \(\min(t/p^e,t')\).
\end{example}

For the full pairing \(\sum_i y_i^{p^e}y'_i\), one must impose the
congruences simultaneously in all column heights. To impose a linear
congruence over the valuation ring, choose a coefficient of least
valuation, use it to eliminate the other coefficients, and scale the
remaining direction. Only this last operation contributes to the
determinant valuation. The simultaneous choice makes all
required pairings vanish modulo \(\pi_p\), and the valuation of the
product of the change-of-direction determinants, counted with column
multiplicity, is at most \(v\sum_{e,\tau}D_e(\tau)\). Here it suffices to
impose \(e\in I_p^\circ\) with \(e\le\lfloor\log_p B\rfloor\).
Section~\ref{sec:ab-closure} constructs the simultaneous bases, including the two minimum-cost choices and the marked height-one exceptions.\formalnote{4.2}{2}

After the substitution, the part of each relevant \(\Gamma_e\) is divisible
by \(\pi_p\). In \eqref{eq:Athetaform},
\(v_p(\varpi^{p+1}\pi_p/p)\ge 1/(p-1)\); this makes the exponential integral
through the required degree. There are at most \(fR^3\) direction variables.
In the chain rule the derivative gains a factor \(\varpi^B\) and the
column determinants. Removing both factors from the integral derivative
yields, in the original directions,
\[
 v_p(J_P(F))\ \ge\
 B\gamma_p-v\sum_{\substack{e\in I_p^\circ\\e\le\lfloor\log_p B\rfloor}}
                         \sum_\tau D_e(\tau)
              -fR^3(1+\lfloor\log_p B\rfloor).
\]
This estimate explains both the positive part and the ramified exception.
Extend the nonnegative deductions to all layers and write
\(u-\min(u,u')=(u-u')_+\). The contribution of \(B\gamma_p\) in
each positive layer is \(v\sum_\tau(t_\tau+b_\tau)/p^e\).
Subtracting \(vD_e(\tau)\) leaves the explicit prediction
\begin{align}
 P_p={}&v\sum_{e\in I_p}\sum_\tau
 \left[\left(t_\tau/p^e-t_{\tau'}-c_{\tau'}\right)_+
       +\left(b_\tau/p^e-b_{\tau'}\right)_+ -c_\tau/p^e\right]\nonumber\\
 &-\mathbf1_{p=d}\,\frac12\sum_\tau D_0(\tau).\label{eq:Apred}
\end{align}
The final term occurs because the ramified normal form contains \(e=0\)
but the expansion of \(\gamma_p\) starts at \(e=1\). The proved bound is
\[
 v_p(J_P(F))\ge P_p-fR^3\bigl(1+\log_p\max(B,1)\bigr).
\]
In the charged version, mark places \(S\) and attach powers \(p^h\) to rows
of the fillings. Each pair involving a marked height-one column must have a
positive charge at one end (with the paired height and row as in the
theorem). Add the total filling charge to the left, compute \(P_p\) with
the marked height-one columns erased, and add \(\gamma_p\) on the right for
each erased box. This is the precise variant used after Hecke transfer.\formalnote{4.3}{3}

\subsection{From the diagrams to an arithmetic sum}

At lift \(j<m\) take the rank bound \(R=n_{j+1}=\dim_E X_{j+1}=k+(j+1)(t+1)\) in the preceding derivative estimate. The long row at an
unmarked \(\tau\in\Phi_a\) consists of \(a\) height-one columns. The other
places have none. The column estimates at each place give
\[
 c_\tau\le \kappa_\tau+x^{60},\qquad
            k\kappa_\tau\le b_\tau\le k\kappa_\tau+x^{60}.
\]
Thus the second positive part in \eqref{eq:Apred} contributes at least
\(k(\kappa_\tau/p^e-\kappa_{\tau'})_+-x^{60}\). For a surviving long row, the first
positive part contributes at least
\((a/p^e-\kappa_{\tau'})_+-x^{60}\) unless the partner also has a surviving long
row. In that case its extra \(a\) may absorb the entire layer \(a/p^e\).
This is the collision the proof has to control.

Define the local contribution \(s_p(y,\tau)\) and the predicted valuation gain \(G_p\) for the composite lift by
\begin{align}
 s_p(y,\tau)&=v\sum_{e\in I_p}(y/p^e-\kappa_{-p^e\tau})_+,\nonumber\\
 G_p&=m\left[\sum_{\tau\in\Phi_a}s_p(a,\tau)
                       +k\sum_{\tau\in\Phi}s_p(\kappa_\tau,\tau)\right].
 \label{eq:Aasked}
\end{align}
At a split prime both are defined to be zero. Let
\(C_p(S)=\sum_{\tau,e}a/p^e\), summing over \(e\in I_p\) for which
\(\tau,\tau'\in\Phi_a\setminus S\). Denote by \(V_{p,j}(S)\) the prediction
in \eqref{eq:Apred} after erasing the marked height-one columns, plus
\(\gamma_p\) per erased box, less the logarithmic error in the derivative
count. For a nonsplit prime \(p<x_1\), these estimates give
\[
 V_{p,j}(S)\ \ge\ G_p/m-C_p(S)+v|\Phi_a\cap S|\,a/x_1
                         -\frac{fa(\log b)^2}{mp}.                 \tag{A1}
\]
The errors \(2x^{60}\) in comparing the two positive parts
are summed over at most \(\lceil\log_p\lceil x_1\rceil\rceil\) layers whose
prime powers are below the cutoff. Admissibility makes this product at most
\(a/(pm)\). The bound for \(R^3(1+\log_p\max(B,1))\) is also \(a/(pm)\);
the deduction at \(e=0\) and the term \(c_\tau/p^e\) are accounted for in
the same comparison. The surplus \(va/x_1\) for an erased long row uses the
common lower bound \(\kappa_{\tau'}\ge a/x_1\).\formalnote{4.4}{4}

In \eqref{eq:Aasked} the factor \(v\,\mathbf1_{e\in I_p}\) equals
\((1-\chi_d(p^e))/2\) for \(e\ge1\), including at \(p=d\).
Since \(s_p(y,\tau)=0\) for \(p\ge x_1\) when \(0\le y\le a\), regrouping
the finite sum over powers gives exactly
\begin{equation}
 s_{\rm fin}(y,\tau)=2\sum_p s_p(y,\tau)\log p
 =\sum_{1\le n\le a}\Lambda(n)(1-\chi_d(n))(y/n-\kappa_{-n\tau})_+ .\label{eq:Aprime}
\end{equation}
The positive part in \eqref{eq:Aprime} restricts the prime powers in a term
with partner weight \(w=a/z\) and \(y=a\) to \(n<z\). Dividing that
positive part by \(a\) gives the weight \((1/n-1/z)_+\) in the
centered prime sum.

\begin{example}[The contributions at one inert prime]\label{ex:inert-prime-ramp}
Fix an inert prime \(p\) and write \(w=a/x_1\) for the smaller weight.
Fix one place \(\tau\). In the sum \(s_{\rm fin}(y,\tau)\) just obtained, take \(y=p^4w\), so a power \(p^e\) is tested against the weight at its own partner, \(-p^e\tau\). For sufficiently large \(x\) we have \(y\le a\) and
\(p^3w\le b\), since \(x_1/x\to\infty\).
Among the powers of this prime, only \(p\) and \(p^3\) can contribute. Indeed \(1-\chi_d(p^e)\) is
zero for even \(e\) and is two for odd \(e\), while \(y/p^e<w\) for \(e\ge5\).
If both partners \(-p\tau,-p^3\tau\) have weight \(w\), their total
contribution to \eqref{eq:Aprime} is
\[
  2\log p\,\bigl((p^3w-w)+(pw-w)\bigr)
      =2w(p^3+p-2)\log p.
\]
A partner in \(\Phi_b\) has the larger weight \(b\), so its term
vanishes. Thus each of the two powers is counted according to its
own partner.
\end{example}

The centered sum at modulus \(q\) is
\[
 R_q(z)=\sum_{\substack{n\le z\\n\equiv1\bmod q}}
 \Lambda(n)(1-\chi_d(n))(1/n-1/z)
                  -\frac{\log z}{\varphi(q)}-c_q .
\]
For \(n\) prime to \(q_1\), as \(\tau\) runs over
\(\Phi_a\), precisely \(|\Phi_b|\) partners have weight \(a/x\) when
\(n\equiv1\bmod q_0\), and none otherwise. The other partners have weight
\(a/x_1\). A prime power with prime divisor dividing \(q_1\) has
weight zero, since that prime splits in \(K\), as in \S\ref{sec:analytic}.
Starting with the common cutoff \(x_1\), making those replacements, and subtracting \(s_{\rm main}=a\{\log(a/\bar\kappa)+c_1\}\) for each \(\tau\in\Phi_a\) gives the identity for the arithmetic remainder \(\mathcal R\) of \S\ref{sec:three-estimates}. Here \(\bar\kappa\) is the geometric mean of the place weights and \(c_1\) is the constant in the centered sum \(R_1\). Thus
\[
 \mathcal R/a=|\Phi_a|R_1(x_1)+|\Phi_b|\{R_{q_0}(x)-R_{q_0}(x_1)\}.
\]
Both the local regrouping \eqref{eq:Aprime} and the displayed expression for \(\mathcal R\) are identities.\formalnote{4.5}{5}

\subsection{The additional long row}

Write \(\mathcal A\) for all source forms and \(D_X=\dim_\mathbb C\mathcal A\).
At an embedding over \(p\), \(L_p\subset\mathcal A\) denotes the module
whose values at each finite adelic point are integral in the basis obtained
by translating the standard lattice to that point. The gauge of a value is
the largest absolute value of its coordinates in this basis. An estimate
\(\mathrm{gauge}\le p^{-v_0}\) at every point expresses a valuation of at
least \(v_0\).

For \(p<x_1\) nonsplit define
\[
 \ell_p=\frac{fa(\log b)^2}{p}
 \begin{cases}m,&p\equiv\pm1\pmod m,\\1,&\text{otherwise}.\end{cases}
\]
With empty marking, (A1) and the ultrametric inequality give a valuation
of at least \(G_p-\ell_p-mC_p(\varnothing)\) for the composite lift on integral source forms. At each step, (A1) gives the bound for every allowed translate of the vacuum. The valuation bounds add under composition and still hold after taking integral linear combinations. The twist conductor is
assumed prime to \(p\) here and below.

Put \(r_p=\lceil\log_p m\rceil\) and let \(o\) be the order of the class
named by \(-p\) modulo \(q_0\). A small nonsplit prime is exceptional if
it is inert, \(p\not\equiv\pm1\bmod m\), and \(o\) is odd with \(o\le r_p\).
Outside the exceptional cases the collision term satisfies
\(m C_p(\varnothing)\le\ell_p\). To see the saving when \(o>r_p\),
observe that an odd layer joining two \(a\)-places requires \(o\mid e\).
Hence the geometric tail
begins past \(p^{r_p}\ge m\), absorbing the factor \(m\) from the
number of stages. At \(p\equiv\pm1\bmod m\)
the definition of \(\ell_p\) already pays the extra factor \(m\).

At an exceptional prime choose a full mark \(\mathcal M\), one prime of
\(E\) from each conjugate pair induced by the \(a\)-places. The set \(S\)
of infinite places inducing a member of \(\mathcal M\) meets every
collision, so \(C_p(S)=0\). A mark must be realized by an operator on
the source. For each \(\lambda\in\mathcal M\) set
\(A_\lambda=a\,\#\{\tau\in\Phi_a:\tau\text{ induces }\lambda\}\).
Let \(T_{\mathcal M}\) be the product of the index-\(r_p\) Hecke operators at these places, first scaling each by \(p^{A_\lambda}\). Passing a
Hecke operator through a lift expands it in integral combinations of
lower-index operators after scaling. Evaluation at the resulting
neighboring lattices supplies the charges in the derivative estimate.
The extra term \(|\Phi_a\cap S|a/x_1\) in (A1) and \(\ell_p/m\)
cover both the transfer and the charges on the low columns. It follows
that the valuation on \(T_{\mathcal M}^m L_p\) is at least \(G_p-2\ell_p\).
It is at least as large on \(p^{j_p}L_p\), where
\(j_p=\lceil2mfa/p\rceil\), by the estimate before marking.\formalnote{4.6}{6}

We now bound the colength of \(T_{\mathcal M}^mL_p+p^{j_p}L_p\) in \(L_p\). The integral
coefficient ring is a valuation ring that need not be discrete. Say a
submodule has \emph{colength dividing \(p^j\)} if, for every finite
family of \(n\) vectors of the ambient module, there is an integral
\(n\times n\) matrix \(A\) with \(\det A\mid p^j\) whose row combinations
belong to that submodule. The Frobenius cycle of the places over
\(\lambda,\bar\lambda\) gives an elementary-divisor estimate for the
scaled Hecke operator. Its cost per dimension modulo \(p^{j_p}\) is at most
\[
               \frac{12A_\lambda}{mp}+\frac{j_p}{m^2f}.
\]
Here the underlying edge-operator theorem constructs an approximate inverse
to the map that sums values on edges ending at a lattice. Congruences on
the edge functions and a bound for the number of additional generators
control the image of functions depending only on the origin. The
comparisons of integral lattices at the two ends control the remaining
elementary divisors. This is where the estimate uses more than a
pointwise bound for a theta kernel. Summing the costs over \(\mathcal M\)
and the \(m\) factors gives at most \(\ell_p D_X\).\formalnote{4.7}{7}

There is one further local calculation. At \(p=N\ge x_1\) the predicted
gain is zero. Ramification of the source allows each glued pair of lifts
to be replaced by direct lifts after the operator \(\mathrm{Up}\), with a
scalar of controlled valuation. Use the inverse image under
\(\mathrm{Up}^{m/2}\) of the module integral for the modified lattice
\(L_0\). Comparison with the standard and dual lattices costs at most
\(mfbD_X\) in colength. The valuation of the composite lift is at least
\(-2mfb\) there. At every other prime outside the small nonsplit set,
the standard data preserve integrality and the predicted valuation gain \(G_p\) is zero.\formalnote{4.8}{8}

We can now take \(\epsilon_p=2\ell_p\) at small nonsplit primes,
\(\epsilon_N=2mfb\), and zero elsewhere. For any finite set of primes,
\(\sum_p\epsilon_p\log p\le U\). To see the summation, a prime of the
form \(mc\pm1\) has \(m/p\le 2/c\). For the other small nonsplit primes use \(c=p\). In each case, \(\epsilon_p\log p\le 4fa(\log b)^2\log x_1/c\). For \(p=N\) use \(c=1\). The bounds \(m\le2\sqrt{x}\) and \(\log N\le25m\log b\) give \(2mfb\log N\le200fa\log b\), within the same bound for large \(x\). At most four primes are assigned to each \(c\), and \(\sum_{1\le c\le x_1}c^{-1}\le1+\log x_1\). Hence the sum is at most \(16fa(\log b)^2\log x_1(1+\log x_1)\le U\) for large \(x\).
In particular, the same \(\epsilon_p\) works for every embedding above \(p\).
For each embedding there is a submodule of the integral ramified source,
of colength dividing \(p^{j}\), \(j\le\epsilon_pD_X\), on which
every admissible composite lift of conductor prime to \(p\) has valuation
at least \(G_p-\epsilon_p\).\formalnote{4.9}{9}

\subsection{The Poincar\'e map, the seesaw and the scale}

Let \(S\subset\mathcal A\) be the ramified source. For each \(w\mid\pi_N\), use the scalar \(z_w\) of \S\ref{sec:spaces}, whose component at \(w\) generates the residue-field units modulo \(w\). Then \(S\) is defined by
\(\sum_{i=0}^{N-2}z_w^i\phi=0\) for all \(w\).
Recall from \S\ref{sec:spaces} that \(\mathcal A(Y_0)\) has prescribed finite level and archimedean type, and that \(Y=e_0\mathcal A(Y_0)\). Write \(\overline G\) for the finite group of level translates denoted by \(\Gamma\) in \S\ref{sec:spaces}. The Poincar\'e map
\(\mathcal P:\mathcal A\to\mathcal A(Y_0)\) represents the Bessel
period. For its chosen point \(g\) and linear functional \(t\), this means
\[
 \langle h,\mathcal P\phi\rangle_{\rm Pet}
        =\int_{[R_0]}\mathcal W_{\phi,t}(r)\,\overline{h(rg)}\,dr .
\]
Here \(R_0\) is the Bessel subgroup and \(\mathcal W_{\phi,t}\) uses its
unipotent character and \(t\) applied to the transformed values of \(\phi\).
Finite-dimensional Riesz representation gives this map. Its estimates
require information about how the period behaves under lifting.

Let \(\iota_\vartheta\) be the composite theta lift on \(\mathcal A(Y_0)\). Its restriction to \(Y\) is \(\Theta_\vartheta\). On the other chain let \(B_{\vartheta,\Psi}\) be the composite theta lift from \(X_0\) to \(X_m\), followed by restriction to \(Y_m\), with Schwartz data \(\Psi\).
The seesaw identity of \S\ref{sec:seesaw-levels} gives a single \(c_{\rm ss}\in
\mathbb C^\times\) such that for every twist \(\vartheta\) there is an
admissible, finitely supported \(\Psi\) with
\[
           \iota_\vartheta(\mathcal P\phi)=
                      c_{\rm ss}B_{\vartheta,\Psi}\phi
                      \qquad(\phi\in S).                         \tag{A2}
\]
Pair the left side with a test form on the definite \(Y_m\). Moving its
Bessel period through the stages leaves a real integral and a weighted
finite coset sum at each stage. At the top, restriction to the definite
subspace realizes the resulting integral as the inner product against the
restricted form. The coset calculation and the expression of the resulting
Schwartz tensor in translates of the vacuum give control of the
coefficients of the chain datum. They are integral away from the twist modulus,
and in the untwisted case their conjugates have a controlled sum of absolute
values after expansion. Those properties allow (A2) to be used in the
local and complex estimates.\formalnote{4.10}{10}

Let \(\mathcal P_0\) be the set of small nonsplit primes and put
\[
 h_p=\lfloor G_p-\epsilon_p\rfloor,\qquad
 c_B=\prod_{p\in\mathcal P_0\cup\{N\}}p^{-h_p}\in\mathbb Q_{>0},
 \qquad s\mathfrak p\phi=c_B|\overline G|\,e_0s\,c_{\rm ss}^{-1}
                                         \mathcal P\phi .          \tag{A3}
\]
Centrality shows this is the \(s\)-translate of the map with \(s=1\).
Every prime outside the product has \(h_p=0\). Hence \(v_p(c_B)=-h_p\)
at every prime. At the real place the floor inequality and the error bound
give
\[
       \log c_B\le-\sum_{p\in\mathcal P_0}G_p\log p+U+b.          \tag{A4}
\]
Indeed the sum of the extra \(\log p\)'s from rounding is at most
\(\lceil x_1\rceil x_1+\log N\le b\). The arbitrary complex scalar
\(c_{\rm ss}\) disappears through (A2); the remaining scale in (A4) is
positive rational and has the same value under all complex automorphisms.\formalnote{4.11}{11}

\subsection{The complex estimate and the ramified spanning set}

Recall the geometric mean \(\bar\kappa\) of the place weights and the constant \(c_1\) of the centered sum \(R_1\). With \(s_{\rm main}=a\{\log(a/\bar\kappa)+c_1\}\), the definitions of \(H\) and \(\mathcal R\) are
\[
 H=mk\sum_\tau(s_\infty(\kappa_\tau)-s_{\rm fin}(\kappa_\tau,\tau)),\quad
 \mathcal R=\sum_{\tau\in\Phi_a}(s_{\rm fin}(a,\tau)-s_{\rm main}).
\]
We have \(\Omega=\exp(\gamma-\Re(L'/L)(1,\chi_d))\,2\pi/\mathrm{rd}_E\),
where \(\mathrm{rd}_E\) is the root discriminant and \(\gamma\) is Euler's
constant. In particular
\(\log(M/\Omega)=1+c_1+\log(M\,\mathrm{rd}_E/(2\pi))\).
Stirling's upper bound and the lower end of the prescribed window for
\(\log N\) give
\begin{align*}
 s_\infty(a)-s_{\rm main}
   &\le a\log(\bar\kappa M\,\mathrm{rd}_E/(2\pi))+1+\tfrac12\log a,\\
 m(s_\infty(a)-s_{\rm main})-a\log N&\le32a\log b.
\end{align*}
For the second line, put \(g_0=\log(M\,\mathrm{rd}_E/(2\pi))\). Admissibility
gives \(\log N\ge m(\log\bar\kappa+g_0)-(g_0+\tfrac{15}2\log b)\),
\(g_0\le23\log b\), and \(m(1+\tfrac12\log a)\le a\log b\).
These three inequalities give the displayed constant.

Let \(K_{\rm size}=m(|\Phi_a|s_\infty(a)+k\sum_\tau s_\infty(\kappa_\tau))\).
Equations \eqref{eq:Aasked} and \eqref{eq:Aprime} imply
\[
 H-m\mathcal R+2\sum_p G_p\log p
       =m|\Phi_a|s_{\rm main}+mk\sum_\tau s_\infty(\kappa_\tau).
\]
It follows that \(K_{\rm size}-|\Phi_a|a\log N\) is at most
\(H-m\mathcal R+2\sum_pG_p\log p+32fa\log b\). The term with \(\log N\) on the
left is a saving present in the actual restriction map. The following
details explain how it enters the norm calculation.

At stage \(j\), let \(u_j\) be the sum of \(s_\infty\) of the first
\(n_j\) row lengths over all places. The complex-conjugation and vacuum
estimates for the Schur kernel bound its operator norm by
\(\exp(x^{50}+u_j/2)\), uniformly in the translate and the finite
adelic arguments. Section~\ref{sec:ab-closure} proves this by conjugating the finite vacuum and integrating the Schur polynomial against its real Gaussian. Main row lengths are
\(a\) or zero on the first row, \(\kappa_\tau\) on the next \(k-1\) rows,
and zero beyond them. A surplus box costs at most \(25\log b\), since
the row length is at most \(2a\). Summing the surplus gives
\(\sum_{j<m}u_j\le K_{\rm size}+U/2\).
The weak chain expansion has at most \(\exp(x^{30})\) classes per stage;
its untwisted coefficients have conjugate absolute sum at most
\(\exp(x^{60})\). At the final restriction the squared norm is multiplied
by at most
\[
\exp\{x^{70}-(|\Phi_a|a-fm\Delta/2)\log N\}.
\]
Here \(\Delta=2r^8\) is the gap between consecutive lift weights, with \(r=\dim_E Y_0\). The term \(|\Phi_a|a-fm\Delta/2\) is the total
branching exponent obtained from the diagrams. Multiplying the kernel
norms, summing the class and coefficient bounds, squaring, and applying
the restriction gives
\[
 \|\sigma(B_{1,\Psi}\phi(y))\|^2\le
 \exp\!\left(B_0+H-m\mathcal R+2\sum_pG_p\log p+U\right),                 \tag{A5}
\]
whenever \(\|\sigma(\phi(g))\|^2\le e^{B_0}\) at all input points.
In this estimate \(U/2\) paid the row surplus. The other half pays the
conjugation, class, coefficient, branching-gap and \(32fa\log b\)
terms.\formalnote{4.12}{12}

It remains to produce inputs satisfying (A5) with \(B_0=U\).
At the free source level, class-supported delta forms identify each value
space separately. A short basis for the lattice at that class, combined
with Weyl evaluations at roots of unity, supplies spanning values integral
at all finite places and bounded in every conjugate. Quantitatively, if
\(D_0=q_1d\) is in the different, the bound for their squared norms is at
most
\[
        (6D_0(fk)^2)^{6(kf+fa)}\le e^U/4^f.
\]
Section~\ref{sec:ab-closure} derives this estimate at each class and for every determinant weight. Here \(fa\) bounds the number of boxes at source rank. For each \(w\mid
\pi_N\), the image of \(1-z_w\) equals the kernel of
\(\sum_{i=0}^{N-2}z_w^i\). The commuting finite-order operators therefore
give \(S=\operatorname{im}\prod_{w\mid\pi_N}(1-z_w)\). Each difference
preserves integrality by the ultrametric inequality and multiplies the
squared complex norm bound by at most four. Applying all \(f\) differences
to the delta forms therefore spans \(S\) by integral forms whose squared norms at every conjugate are at most \(e^U\).\formalnote{4.13}{13}

To pass from an individual value to the good target norm, expand
\(|\overline G|e_0s\) in group elements. After lifting, each term is an
integral right translation with coefficient \(|\overline G|e_0(t)\), an
algebraic integer; its absolute value under every conjugate is at most
\(|\overline G|\). Thus the norm before squaring gains at most
\(|\overline G|^2\). The class sum and this fourth power together are
at most \(\exp(x^{99})\). Now apply (A2), (A5) with \(B_0=U\), and the
factor \(\exp(2\log c_B)\) from (A3). By (A4) the positive
\(2\sum_pG_p\log p\) cancels. The remaining exponent is at most
\[
 H-m\mathcal R+x^{99}+4U+2b\ \le\ H-m\mathcal R+fa(\log b)^4.
\]
This proves Theorem A(a). The integral \(\phi\in S\) satisfying
\(\|\Theta_1(s\mathfrak p\phi)\|_\sigma^2
\le\exp(H-m\mathcal R+fa(\log b)^4)\) for \emph{every} \(s\in\overline G\) and
\(\sigma\in\operatorname{Aut}(\mathbb C)\) span \(S\).
Here the norm is the class sum of squared Weyl Gram norms.\formalnote{4.14}{14}

\subsection{The same map has controlled index}

At an embedding \(\iota\) over \(p\), a vector \(y\) is integral in the
twisted target lattice if \(\Theta_\vartheta y\) is integral for
\emph{every} allowed \(m\)-fold twist of level prime to \(p\). A finite
family \(y_i\) has logarithmic index at most \(\delta\) if there are
numbers \(e_p\ge0\) with \(\sum_{p\in F}e_p\log p\le\delta\) for every
finite set \(F\), and at each \(\iota\) one integral matrix \(A_\iota\)
with \(v_p(\det A_\iota)\le e_p\) whose row combinations of \(y_i\)
are integral. The same matrix must work for all such twists.

Fix \(s\) and take \(y_i=s\mathfrak p\phi_i\) with the \(\phi_i\)
globally integral in \(S\). At \(\iota\), apply the colength result to the
family \(\phi_i\). This gives a matrix of determinant valuation at most
\(\epsilon_pD_X\) carrying it into the local good submodule. For each
resulting combination \(\psi\), (A2) replaces the lifted Poincar\'e image
by \(B_{\vartheta,\Psi}\psi\); its valuation is at least \(G_p-\epsilon_p\).
Multiplication by \(c_B\) preserves a nonnegative valuation because
\(v_p(c_B)=-\lfloor G_p-\epsilon_p\rfloor\). Projection and translation
preserve integrality by the group-element expansion above. Linearity
therefore gives the needed matrix on the target, simultaneously for all
allowed twists. The total logarithmic index is at most \(D_X U\), where \(D_X=\dim_{\mathbb C}\mathcal A\) is the dimension of the full source space and \(U=fa(\log b)^3\).

Finally the invariant part for each \(z_w\) has dimension at most
\(D_X/(N-1)\), by comparison with the level enlarged by that cyclic
subgroup. The codimension of the image of the product of the \(1-z_w\)
is at most the sum of their nullities. Consequently
\[
 \dim S\ge(1-f/(N-1))D_X,\qquad f/(N-1)\le b^{-8}\le\tfrac12.
\]
For one translate the index cost is at most \(2(\dim S)U\).
For a family drawn from \(T\subset\overline G\), assign each vector a
translate from which it came. Grouping the correction matrices into a
block diagonal matrix adds the determinant valuations. We obtain
Theorem A(b) in the form used in the determinant comparison,
\[
              \operatorname{IndexLe}\bigl((y_i),\,
                         2|T|\,(\dim S)\,fa(\log b)^3\bigr).
\]
The index estimate applies to every finite family of these images,
including dependent families. Its cost depends on the number of
translates that the spanning theorem needs.\formalnote{4.15}{15}


\section{The integral elliptic theta coordinate and its Frobenius term}
\label{sec:cm-integral}
We construct a local coordinate on the CM elliptic curve in which the theta function can be expanded integrally after a change of scale. The calculation of \S\ref{sec:cm-detail} will use it to determine which powers occur in the exponent of the series. We also prove an identity for six copies of the function that will allow us to isolate a single Taylor coefficient in \S\ref{sec:frame-attainment}. Throughout, \(K=\mathbb Q(\sqrt{-d})\), \(d\) is prime and \(d\equiv7\bmod8\). Write \(\beta=(1+\sqrt{-d})/2\). Fix the comparison \(\iota:\mathbb C\hookrightarrow\mathbb C_p\), the valuation ring \(O\) and its maximal ideal \(\mathfrak m\). Complex identities below are transported by applying \(\iota\) to their Taylor coefficients.

\subsection{The normalized cubic and its good reduction}
Let \(\Lambda\) be the CM lattice, \(A\) its area, and \(t\) the chosen point of order two for which \(\beta t,\bar\beta t\) are both nonzero modulo \(\Lambda\). Normalize by \(\wp''(t)=2\). Write \(a^2=12\wp(t)\), and put
\[
      x(z)=\wp(z+t)-\wp(t),\qquad
      y(z)=\tfrac12\bigl(\wp'(z+t)-a x(z)\bigr).
\]
The Weierstrass differential equation, expanded about the root \(\wp(t)\), becomes
\[
        y^2+a x y=x^3+x,\qquad
        x'=2y+ax,\quad y'=1+3x^2-a y.                     \tag{CM1}
\]
In particular \(y=z+O(z^2)\) and \(x=z^2+O(z^3)\). We first prove \(|a|_p\le1\) and \(|a^4-64|_p=1\), including \(p=2\).

For each of the three nonzero points \(t_i\) of order two put \(e_i=\wp(t_i)\). For a cyclic ordering let \(\delta_i=e_{i+1}-e_{i+2}\), and define
\[
 s(t_i)=\frac{288e_i^2}{\wp''(t_i)}-64
       =\frac{-16\delta_i^2}{\delta_{i+1}\delta_{i+2}} .
\]
Indeed the cubic has roots \(e_i\) of sum zero and \(\wp''(t_i)=2(e_i-e_{i+1})(e_i-e_{i+2})\). The second expression gives \(\sum_i s(t_i)=-48\) and \(\prod_i s(t_i)=-2^{12}\), since \(\sum_i\delta_i=0\) and \(\sum_i\delta_i^3=3\prod_i\delta_i\). Consider an isogeny \(z\mapsto c z\) with kernel \(\{0,t_1\}\), and use a subscript \(D\) for its target. The two-coset formula gives
\[
     c^2\wp_D(cz)=\wp(z)+\wp(z+t_1)-e_1.
\]
Both sides have the same poles and principal parts and the same constant term at zero. At \(t_2\) it gives \(c^2\wp_D(ct_2)=-2e_1\); differentiating twice gives \(c^4\wp_D''(ct_2)=2(e_2-e_3)^2\). Substitution in the formula for \(s\) gives \(s(t_1)s_D(ct_2)=2^{12}\).

Choose normalized CM lattices in the ideal classes of \(K\). Taking the quotient by \(b t\), for \(b=\beta\) or \(\bar\beta\), permutes the classes by multiplication by a fixed class \(g\). Indeed the enlarged lattice is \(\Lambda+\mathbb Z bt\), the product of \(\Lambda\) with the inverse of a prime above two: on \(\Lambda/2\Lambda\) the idempotents \(b,1-b\bmod2\) select the two factors of \(\mathcal O_K/2\). This also identifies the surviving point \((1-b)t\) upstairs with \((1-b)t\) of the normalized target modulo its lattice. If \(s_C=s_C(bt)\), \(s'_C=s_C((1-b)t)\), \(h_C=s_C(t)\), the three-point and two-coset identities give
\[
 s_C+s'_C+h_C=-48,\quad s_Cs'_{gC}=2^{12},\quad
                  s_{gC}h_{gC}=-s_C.
\]
If \(|s_C|_p>1\), then \(|s'_{gC}|_p\le1\), while the last product has norm greater than one. The first identity on \(gC\) forces \(|s_{gC}|_p=|h_{gC}|_p=|s_C|_p^{1/2}\). Iterating around the finite orbit contradicts \(|s_C|_p>1\). Applying this to both choices of \(b\) proves the two corresponding entries integral, and then \(h_C\) is integral. Multiplying \(s_{gC}h_{gC}=-s_C\) over all classes gives \(\prod_C h_C=(-1)^{h_K}\), where \(h_K\) is the class number, since every \(s_C\ne0\). Thus each \(h_C\) is a unit. In the chosen normalization \(h_C=a^4-64\), which proves both assertions.\formalnote{4.16}{16}

\subsection{The local height and an integral basis of sections}
Let \(\sigma\) be the odd sigma section with \(\sigma'(0)=1\), and \(\theta\) its even twisted translate by \(t\), with \(\theta(0)=1\). Write \(|b|_p=|\iota(b)|_p\) for complex \(b\). The local height \(h(z)=\max(|\sigma(z)|_p,|\theta(z)|_p)\) is positive since the two sections have disjoint zero sets. The sigma addition formula gives \(x\theta^2=\sigma^2\), \(\theta(2z)=\theta^4-\sigma^4\) and \(\sigma(2z)=x'\theta^4\) off the poles of \(x\). By (CM1) the pair \((\theta^2,\sigma^2)\) transforms under doubling by the quartics
\[
 D(U,V)=\big((U^2-V^2)^2,\ UV(4U^2+a^2 UV+4V^2)\big).
\]
They have no common projective zero after reduction. Indeed, a zero of the first coordinate has \(V=\pm U\ne0\), and vanishing of the second would give \(a^4=64\), contrary to the unit discriminant. Homogeneity now gives \(h(2z)=h(z)^4\), also at \(p=2\).

Put \(G=y\theta^3\), which extends across the poles of \(y\). The equation \(G^2+a\theta\sigma^2G=\sigma^6+\theta^4\sigma^2\) gives \(|G(z)|_p\le h(z)^3\): a larger left-hand root would make \(G^2\) strictly larger in norm than every other term. Thus a monomial \(\theta^j\sigma^{2k}G^e\) of degree \(m=j+2k+3e\) has norm at most \(h^m\). The complex spanning result of \S\ref{sec:sigma-identities} gives \(|f|_p\le C_f h^m\) for every section of degree \(m\) with multiplier that of \(\theta^m\).

For an isogeny \(z\mapsto\alpha z\) of degree \(n\) to a lattice with height \(h'\), the ratio \(r(z)=h'(\alpha z)/h(z)^n\) is bounded: the pullbacks of \(\theta'\) and \((\sigma')^2\) have degrees \(n,2n\). Doubling gives \(r(z)^4=r(2z)\), hence \(r\le1\). For translation by a torsion point \(u\), put \(b_u(z)=\exp[-\pi(\bar u z+|u|^2/2)/A]\), where \(A\) is the area of the lattice. Thus \(T_u^{(m)}f(z)=b_u(z)^m f(z+u)\). The ratio \(s(z)=|b_u(z)|_p h(z+u)/h(z)\) is bounded since the \(2N\)-th powers of the two translated functions are sections of degree \(2N\) when \(Nu\in\Lambda\). The identity \(b_u(z)^4=b_u(2z+u)b_u(2z)\) and doubling give \(s(z)^4=s(2z+u)s(2z)\). Taking suprema gives \(\sup s\le1\). Applying this also to the inverse translate gives equality pointwise. In particular the normalized torsion value \(f[u]=(T_u^{(m)}f)(0)\) is integral for each of the above monomials, since \(h(0)=1\), and \(|\sigma[t]|_p=1\) since \(\theta(t)=0\).\formalnote{4.17}{17}

Let \(f\) be a section of degree \(m\) with the multiplier of \(\theta^m\). Suppose that \(|f(z)|_p\le h(z)^m\) for every \(z\). Off the zeros of \(\theta\) the complex basis gives
\[
 f=\theta^m\{P(x)+yQ(x)\},\qquad
 2k\le m\text{ in }P,\quad 2k+3\le m\text{ in }Q,                  \tag{CM2}
\]
where \(k\) denotes an exponent with nonzero coefficient. The coefficients of \(P,Q\) are integral. To see this, put \(D_1(X)=4X^3+a^2X^2+4X\), so \(x'^2=D_1(x)\). Choose an integral \(c\) at which \(P,Q,D_1\) each attain their largest coefficient norm. Such a \(c\) exists: normalize each nonzero polynomial by a largest coefficient and avoid the finitely many roots of the reductions in the infinite residue field of the algebraically closed valued field \((\mathbb C,|\cdot|_p)\). At least one of \(4,a^2\) is a unit by the discriminant, so \(|D_1(c)|_p=1\). Choose \(z\) with \(x(z)=c\) by elliptic uniformization. At both \(z\) and \(-z\) we have \(h=|\theta|_p\) since \(\sigma^2=x\theta^2\). Thus \(P(c)+y(\pm z)Q(c)\) are integral. Their difference is \(x'(z)Q(c)\), with \(x'(z)\) a unit. Also \(y(z)\) is integral by its monic equation (CM1). Hence \(Q(c),P(c)\) are integral, as are all their coefficients by the choice of \(c\). For even \(f\), subtraction at \(z,-z\) gives \(x'Q(x)=0\) identically, hence \(Q=0\).\formalnote{4.18}{18}

\subsection{Pullback, the CM phase and the first unit}
The cubic reads \(x=y^2+axy-x^3\). Since \(x(0)=y(0)=0\), solving recursively expresses \(x\) as an integral series in \(y\), with leading term \(y^2\). Put \(\mathcal I=O[[y]]\). For an isogeny \(z\mapsto\alpha z\) from \(\Lambda\) to \(\Lambda'\) of degree \(n\), one has \(|\alpha|_\infty^2 A=n A'\). For a period \(l\) the theta multiplier is \((-1)^{|l|^2\sqrt d/(2A)}\). This exponent is the integral alternating pairing of \(\bar\beta l\) with \(l\). The sign formula follows from the theta semicharacter identity applied to \(l=\beta l+\bar\beta l\), whose two summands have multiplier one. The area identity matches both the multiplier and the exponential in the twist under pullback. The height comparison gives the required bounds on the pullbacks. Applying (CM2) to the target sections \(\theta'\) and \(G'\) composed with the isogeny gives
\[
 \theta'(\alpha z)=\theta(z)^n R_1,\qquad
 G'(\alpha z)=\theta(z)^{3n}R_3,\qquad R_1,R_3\in\mathcal I .
\]
Here the prime on \(\theta',G'\) denotes the target. Since \(R_1(0)=1\), \(y_{\Lambda'}(\alpha z)=R_3/R_1^3\) belongs to \(\mathcal I\), and its linear coefficient \(\iota(\alpha)\) is integral. At a nonsplit \(p\) take \(\pi=p,\ n=p^2\) in the inert case and \(\pi=\sqrt{-d},\ n=d\) in the ramified case. The pullback of \(\theta\) is even, and \(n=q\) is odd, so (CM2) gives the polynomial of degree \(<q/2\) used in (F1).\formalnote{4.19}{19}

Write \(g=-\log\theta\). Differentiating twice gives \(g''=g''(0)+x\), since the logarithmic second derivative of the normalized translated sigma is \(\wp(z+t)\) plus a constant with this sign for \(g\). For \(\alpha\in\mathcal O_K\), the first pullback formula gives
\[
        \exp\!\bigl(\alpha\bar\alpha g(z)-g(\alpha z)\bigr)\in\mathcal I.
\]
Its linear term is zero, so its quadratic term in \(z\) equals its quadratic term in \(y\). This proves integrality of \(\alpha(\bar\alpha-\alpha)[z^2]g\). Subtract the assertion for \(\beta\) from the one for \(\bar\beta\). Since \(\beta+\bar\beta=1\), the difference proves \(\sqrt{-d}\,[z^2]g\in O\).

The two pullback statements extend to limits \((\delta,s)\) of the pairs \((\alpha,\alpha\bar\alpha)\) in \(\mathbb C_p^2\). Indeed coefficients at any fixed degree vary continuously, and re-expansion in \(y\) also depends on finitely many coefficients since \(y=z+\cdots\). The condition that each resulting coefficient lie in the closed set \(O\) passes to the limit. If \(\zeta^{q-1}=1\), choose \(\alpha\) whose residue is \(\bar\zeta\) (here meaning reduction of \(\zeta\)) and whose complex conjugate reduces to \(\bar\zeta^p\). In the inert case conjugation is Frobenius on \(\mathbb F_{p^2}\); for \(q=p\), a rational integer lifting the residue suffices. The powers \(\alpha^{q^j}\) and \(\bar\alpha^{q^j}\) tend to \(\zeta\) and \(\zeta^p\): their ratios to these roots reduce to \(1\) and successive \(p\)-powering tends to \(1\). Thus \(y(\zeta z)\in\mathcal I\). Apply the phase formula at \((\zeta^{-1},\zeta^{-p-1})\) and substitute \(\zeta z\) to obtain
\[
           \exp\bigl(\zeta^{-p-1}g(\zeta z)-g(z)\bigr)\in\mathcal I .
\]
This gives the integrality needed to project onto the exponent in (F2).\formalnote{4.20}{20}

At a place of good reduction, Deuring's criterion says that a CM elliptic curve has ordinary reduction when \(p\) splits in the CM field and supersingular reduction when \(p\) is inert or ramified \cite[Thm.~1.1]{blake-deuring}. We prove directly that the coefficient of \(T^q\) in the multiplication series \(T(\pi z)\), expressed in the coordinate \(T=T(z)\) constructed below, is a unit. For \(\alpha\in\mathcal O_K\setminus\{0\}\) let \(o(\alpha)\) be the order in \(y\) of the reduction of \(y(\alpha z)\) modulo \(\mathfrak m\). Apply the even expansion (CM2) to \(\sigma(\alpha z)^2\), of degree \(2n\), \(n=\alpha\bar\alpha\). This gives \(\theta^{2n}R(x)\) with \(R\) integral, \(\deg R\le n\). At \(z=t/\alpha\) the left side has normalized torsion value \(\sigma[t]^2\), a unit. The right side is a sum of the integral monomial sections \(\theta^{2n-2i}\sigma^{2i}\) multiplied by \(R_i\). Thus some \(R_i\) with \(i\le n\) is a unit. Using the first pullback and \(x\theta^2=\sigma^2\) gives \(x(\alpha z) R_1^2=R(x)\). After reduction, the left side has order \(2o(\alpha)\), since \(x=y^2+\cdots\) in \(O[[y]]\) and \(R_1\) is a unit. The right side has order \(2\operatorname{ord}\bar R\le2n\). Consequently \(o(\pi)\le q\) for nonsplit \(p\). If \(p\) splits, take a power of the prime over \(p\) which is principal, generated by \(\alpha\) with \(\alpha\bar\alpha=p^h\) and \(|\bar\alpha|_p=1\) at the chosen embedding. Composition multiplies orders, and \(o(\bar\alpha)=1\) by its linear coefficient. Hence \(o(p)^h=o(\alpha)\le p^h\), giving \(o(\pi)\le q=p\) here as well.\formalnote{4.21}{21}

Take \(T=(q-1)^{-1}\sum_{\zeta^{q-1}=1}\zeta^{-1}y(\zeta z)\). This is integral in \(y\), tangent to \(y\), and therefore another generator. Its equivariance and the integral pullback give
\[
       T(\pi z)=\pi T(z)+\sum_{j\ge1} b_j T(z)^{1+j(q-1)},\qquad b_j\in O .
\]
Orders after reduction are independent of the integral generator tangent to \(y\). The linear coefficient \(\pi\) reduces to zero, and the only remaining possible degree at most \(q\) is \(q\). The bound \(o(\pi)\le q\) thus proves \(b_1\) is a unit. Choose \(\varpi^{q-1}=-\pi/b_1\), put \(Y=T(\varpi z)/\varpi\), and substitute:
\[
 \pi^{-1}Y(\pi z)=Y-Y^q+\pi E(Y),\qquad E\in O[[Y]],\quad E\in(Y^2).
\]
In degree \(n\ge2\), moving the linear term to the left leaves the unit \(\pi^{n-1}-1\) times \([z^n]Y\); the right side uses only earlier coefficients. With \([z]Y=1\) this proves \(Y\in O[[z]]\). The left side reduces to \(z\bmod\pi\), giving \(Y-Y^q\equiv z\bmod\pi\). Its unique zero-constant solution is \(\sum_{j\ge0}z^{q^j}\) by successive substitution in characteristic \(p\). In particular every series integral in \(y\) is integral after scaling by \(\varpi\); its terms of degrees \(<q\) in the unscaled \(z\) are integral since \(T=z+O(z^q)\).\formalnote{4.22}{22}

\subsection{Integrality after six copies}
Equation (CM1) also gives the integral two-variable addition law used for the one-curve derivative bound. To see the absence of small denominators, set \(x_i=x(z_i), y_i=y(z_i)\), \(h=x_0x_1\), \(u=(1-h)^{-1}\). Euler's expression for \(y(z_0+z_1)\) is
\[
\begin{split}
 u^2\big((1+h)(y_0+y_1)-a y_0y_1\big)
  +u^3\big((3+h)(x_0+x_1)+a^2 h\big)
                    (x_0y_1+x_1y_0+a h).
\end{split}
\]
For verification apply \(\partial_{z_0}-\partial_{z_1}\), using (CM1) and \(D u=u^2 D h\). Substituting \(y_i^2=x_i^3+x_i-a x_iy_i\) cancels the terms. At \(z_0=0\) the expression is \(y_1\). In characteristic zero the equation \(D F=0\) and this boundary value determine \(F\), by comparison in the coordinates \(z_0,z_0+z_1\). The sole inverse \(u\) is integral in \(O[[y_0,y_1]]\) because \(h\) has constant coefficient zero. Differentiation also preserves \(\mathcal I\) by (CM1). These facts justify the submultiplicativity and derivative integrality in \S\ref{sec:cm-detail}.\formalnote{4.23}{23}

To obtain four \(p\)-adic integers whose squares sum to \(-1\), write each \(p^n-1\) as a sum of four integer squares by Lagrange's theorem. The resulting quadruples have a convergent subsequence in \(\mathbb Z_p^4\), and its limit has the required sum of squares. Use these four, \(1\) in the first slot, and \(0\) in the last to form \(\delta_0,\ldots,\delta_5\). For any \(\delta\in\mathbb Z_p\) the pair \((\delta,\delta^2)\) is a limit of \((n,n^2)\), \(n\in\mathbb N\). The integral phase just proved therefore gives
\[
   \prod_c\theta(\delta_c z)=
   \prod_c\exp\{\delta_c^2 g(z)-g(\delta_c z)\}\ \in\ O[[y(z)]]^\times .
\]
The equality uses \(\sum_c\delta_c^2=0\); the constant coefficient is one, which proves invertibility. For the frame product \(\Theta=\prod_i\theta(\mu_i)^{m_i}\), substitute each linear form \(\mu_i\) into the identity and its inverse, and multiply with exponents \(m_i\). This proves the unit assertion in the exact integral ring \(O[[y(\mu_i):i]]\). Section~\ref{sec:frame-attainment} uses this unit and its inverse to obtain the six-copy coefficient bounds \textup{(V3)}--\textup{(V4)}.\formalnote{4.24}{24}

\subsection{A period integral at every finite place}
Finally we justify the finite-unit normalization in \S\ref{sec:frame-proof}. Put \(h=h_K\). The ideal class group has odd order by the genus theorem for the prime fundamental discriminant \(-d\). In particular squaring permutes its elements. Choose representatives \(\mathfrak b_C\) and replace them by \(\mathfrak c_C=\mathfrak b_{C^2}/\mathfrak b_C\). These represent the same classes and have product the unit ideal. Choose normalized CM lattices \(\Lambda_C=\mu_C\mathfrak c_C\Lambda\), and let \(\Delta_C=a_C^4-64\) be their unit discriminants. Define \(u\) by
\[
                          u^{12h}=\prod_C\mu_C^{12}\Delta_C .
\]
Every number of the lattice \(\mu_C\mathfrak c_C\) sends \(\Lambda\) into \(\Lambda_C\) and is integral at the chosen embedding by the linear term of the isogeny pullback of \(y\). Their products and sums are integral. Since \(\prod_C\mathfrak c_C=(1)\), it follows that \(\prod_C\mu_C\) is integral. Applying the argument to the inverse ideal multipliers, which send \(\Lambda_C\) into \(\Lambda\), proves integrality of the reciprocal. Hence \(|\iota u|_p=1\) for every \(p,\iota\).

To identify the archimedean normalization put \(a_d=2/\sqrt d\), and
\[
                 B_\sigma=A(\Lambda^\sigma)|\sigma u|^2
\]
for \(\sigma\in\operatorname{Aut}(\mathbb C/K)\); \(\Lambda^\sigma\) denotes the lattice of the conjugate cubic. The conjugation theorem for normalized CM lattices gives \(\Lambda_C^\sigma=(\sigma\mu_C)\mathfrak c_C\Lambda^\sigma\). Using covolumes and the defining equation of \(u\),
\[
 \prod_C A(\Lambda_C^\sigma)\,|\sigma\Delta_C|^{1/6}=B_\sigma^h .
\]
Indeed each area contributes \(|\sigma\mu_C|^2 N(\mathfrak c_C) A(\Lambda^\sigma)\), the product of ideal norms is one and \(|\sigma u|^{2h}=\prod_C|\sigma\mu_C|^2|\sigma\Delta_C|^{1/6}\). The classes of these conjugate lattices still run through the class group.

For a representative \(L\) set \(F_L(s)=(s-1)\sum_{0\ne l\in L}(|l|^2/(a_d A(L)))^{-s}\). The first Kronecker limit formula gives \(F_L(1)=\pi a_d\), and
\[
 \frac{F_L'(1)}{F_L(1)}=
   2\gamma-\log\!\left(\frac{A(L)}{a_d\pi^2}\right)
                    -\tfrac16\log|\Delta_L|.
\]
On summing the underlying series over a full set of ideal classes one obtains the number of units of \(K\) times \(\zeta(s)L(s,\chi_d)\). Indeed, after the covolume normalization the summands correspond to norms of integral ideals in the inverse class, with one generator per unit. After multiplying by \(s-1\) and dividing the derivative by the value at one, the right-hand side gives \(\gamma+\operatorname{Re}(L'/L)(1,\chi_d)\). The left-hand side is the mean of the displayed logarithmic derivatives since the values \(F_L(1)\) are equal. Substitution of the product formula for \(B_\sigma^h\) yields
\[
 B_\sigma=a_d\pi^2\exp\{\gamma-\operatorname{Re}(L'/L)(1,\chi_d)\}.
\]
This identifies the common normalization at all embeddings and proves the finite-unit claim for the same \(u\). The finite norm of the residual Hecke-character root in \(P_0\) is then the calculation in \S\ref{sec:frame-proof}. Sections~\ref{sec:sigma-identities}--\ref{sec:kronecker-limit} derive the doubling identities, construct the conjugate period and evaluate the lattice limit in these normalizations. The argument uses complex uniformization, the CM isogeny and multiplier laws, and the indicated local-field facts.\formalnote{4.25}{25}


\section{From the period lattice to the two doubling polynomials}
\label{sec:sigma-identities}
The finite-place proof in \S\ref{sec:cm-integral} iterates a pair of homogeneous quartics. We derive these quartics over \(\mathbb C\) first. The constants in them depend on the normalization of a half-period. In the following calculation \(A\) is the Euclidean area of the lattice and \(\sigma'(0)=1\).

\subsection{The normalized section and the addition law}
Write \(\Lambda=a(\mathbb Z+\tau\mathbb Z)\), \(y=\Im\tau>0\), \(A=|a|^2y\), and let
\(\vartheta(w)=\theta_1(\pi w\mid\tau)\). The odd Jacobi theta function has simple zeros at \(\mathbb Z+\tau\mathbb Z\). Its multipliers on \(1,\tau\) are respectively \(-1\) and \(-e^{-\pi i\tau-2\pi iw}\). In the gauge needed here the section is
\[
 \sigma(z)=a\,e^{\pi (z/a)^2/(2y)}\frac{\vartheta(z/a)}{\vartheta'(0)} .
\]
Thus \(\sigma'(0)=1\). Substitution in the two generator shifts gives, for \(l=a(m+n\tau)\),
\[
 \sigma(z+l)=(-1)^{m+n+mn}
   \exp\!\left(\frac{\pi\bar l z}{A}+\frac{\pi |l|^2}{2A}\right)\sigma(z).
\]
Uniqueness follows by dividing any such odd section by \(\sigma\). It vanishes at all the lattice points, so the quotient is entire and elliptic. It is constant, and the derivative normalization gives its value one at zero. The classical Weierstrass sigma function also equals a constant times an exponential in \(z^2\) times \(\vartheta(z/a)\). Both sigma functions have derivative one at zero, so their ratio is \(e^{cz^2}\). This factor cancels in the sigma addition formula because \((z+v)^2+(z-v)^2=2z^2+2v^2\). Thus for \(z,v\notin\Lambda\) the formula in our gauge is
\[
 \sigma(z+v)\sigma(z-v)=(\wp(v)-\wp(z))\sigma(z)^2\sigma(v)^2 .
 \tag{SG1}
\]
The Jacobi multipliers and the simple zeros follow from (20.2.6) and the product (20.5.1) of \cite{theta-dlmf}. The sigma--theta identity is (23.6.9) there, and the addition formula in the Weierstrass gauge is (23.10.3).
\phantomsection\label{route:7}\leanroute{Lean constructs the normalized section in a space of theta functions and derives addition by evaluation in degree two. The classical identity applies after the quadratic change of gauge displayed above. The half-period and the integral quartics below retain the normalization used by Lean.}\formalnote{4.26}{26}

\subsection{The two quartics}

For \(z\notin\Lambda\), differentiate with respect to \(v\) at \(v=z\). The factor \(\sigma(z-v)\) has derivative \(-1\) there, which fixes the sign in \(\sigma(2z)=-\wp'(z)\sigma(z)^4\).\formalnote{4.27}{27}

Write \(H(t,z)=\pi\bar t z/A\). Let \(t\notin\Lambda\) with \(2t\in\Lambda\), and normalize \(\wp''(t)=2\). Put \(e=\wp(t)\), \(b^2=12e\), and
\[
 \theta(z)=e^{-H(t,z)}\frac{\sigma(z+t)}{\sigma(t)},\qquad
 x(z)=\wp(z+t)-e .
\]
Because \(2t\notin2\Lambda\), its sigma multiplier has sign \(-1\). The shift by \(2t\) and oddness give \(\theta(-z)=\theta(z)\). Apply (SG1) to \(t,z\). Multiplication of the two exponential factors at \(z,-z\) gives one, so
\[
                 \theta(z)^2=(\wp(z)-e)\sigma(z)^2 .
\]
The function \((\wp(z)-e)(\wp(z+t)-e)\) has possible poles only at \(0\) and \(t\) modulo periods. At zero its factors have leading terms \(z^{-2}\) and \(z^2\), since \(\wp'(t)=0,\ \wp''(t)=2\). Translation by \(t\) interchanges the factors. Their product is therefore constant and has value one. We have proved
\[
                   x=(\wp-e)^{-1},\qquad x\theta^2=\sigma^2 .
 \tag{SG2}
\]
Apply (SG1) again, now at \(z+t,v\), and divide by the normalization constants for \(\theta\). This gives
\[
   \theta(z+v)\theta(z-v)=\theta(z)^2\{\theta(v)^2-x(z)\sigma(v)^2\}.
\]
Taking \(v=z\) and using \(\theta(0)=1\) gives \(\theta(2z)=\theta(z)^4-\sigma(z)^4\). Differentiating the reciprocal in (SG2) gives \(x'\theta^4=-\wp'\sigma^4=\sigma(2z)\). Identities of entire products proved away from the exceptional points extend there by continuity.\formalnote{4.28}{28}

The Weierstrass equation gives \(4e^3-g_2e-g_3=0\) and \(12e^2-g_2=2\wp''(t)=4\). Expanding it at \(e+x\) gives
\[
                       (x')^2=x(4x^2+b^2x+4).
\]
Write \(U=\theta(z)^2,\ V=\sigma(z)^2\), so \(x=V/U\). Squaring the two doubling formulas produces
\[
 \big(\theta(2z)^2,\sigma(2z)^2\big)
       =\big((U^2-V^2)^2,\ UV(4U^2+b^2UV+4V^2)\big).
 \tag{SG3}
\]
The Gaussian at \(2z\) for a squared section equals the fourth power of the Gaussian at \(z\). Hence (SG3) also holds for the normalized torsion evaluations in \S\ref{sec:cm-integral}. That section uses the unit \(b^4-64\) to show that reduction preserves the maximum norm of this pair to the fourth power.

\subsection{The three generators}
Let \(Y=(x'-bx)/2\) and \(G=Y\theta^3\). The last differential equation gives \(Y^2+bxY=x^3+x\). Near \(z=-t\), \(x\) has pole order two, \(Y\) has pole order three, and \(\theta\) has a simple zero. Thus \(G\) is entire. Dividing any degree-\(m\) section with multiplier that of \(\theta^m\) by \(\theta^m\) gives an elliptic function with just one possible pole, at \(-t\), of order at most \(m\). For every integer \(2\le j\le m\) write \(j=2h+3c\) with \(h\ge0\) and \(c\in\{0,1\}\). A multiple of \(x^hY^c\) cancels a leading pole of order \(j\). Descending in \(j\) leaves at worst a single simple pole, whose residue must vanish by the residue theorem on a period parallelogram. The remainder is constant. Multiplying back by \(\theta^m\) gives a spanning family
\[
            \theta^{m-2h-3c}(\sigma^2)^hG^c,\qquad
            c=0,1,\quad 2h+3c\le m .
\]
Its distinct pole orders before multiplication also give linear independence. For \(m=0\) or \(1\) only the constant survives before multiplication. This establishes the basis over \(\mathbb C\). The local height criterion of \S\ref{sec:cm-integral} gives integrality of the coefficients in \textup{(CM2)}.\formalnote{4.29}{29}


\section{Transporting a CM period and its theta sections}
\label{sec:cm-conjugation}
In \S\ref{sec:cm-integral} the periods of a conjugate cubic are compared with ideal multiples of a reference lattice. For an arbitrary \(\alpha\in\operatorname{Aut}(\mathbb C/K)\), this comparison must be made through equations: such an automorphism need not preserve analytic limits. The classical CM theorem describes the period lattice and torsion of the conjugate curve by the action of a finite idele \cite[Theorem~9.17]{upper-milne-cm}. Here we also need the scalar by which the normalized theta section changes. Write \(K=\mathbb Q(\sqrt{-d})\), \(d\equiv7\bmod8\) prime, and \(\omega=(1+\sqrt{-d})/2\).

\subsection{The ideal classes and the common scale}
The class group of \(K\) has no element of order two. To prove it, let \(I\) be a nonzero integral ideal with \(I^2=(u)\) and put \(n=NI\). Then \(u\bar u=n^2\). Choose the sign of \(u\) so that \(\bar u+n\ne0\), and set \(J=(\bar u+n)I\). Since
\[
     (\bar u+n)^2u=(\bar u+n)(u+n)n=NJ ,
\]
we have \(J^2=(NJ)\). For every unramified rational prime \(p\), the exponents of the primes above \(p\) in \((NJ)\) are \(v_p(NJ)\). They must be even. The discriminant of \(K\) is \(-d\), so \(NJ=s^2\) or \(s^2d\) for an integer \(s>0\). The unique prime above \(d\) is \((\sqrt{-d})\), since \((\sqrt{-d})^2=(d)\). Thus \(J^2\) is the square of the principal ideal \((s)\) or \((s\sqrt{-d})\). Unique factorization of ideals identifies \(J\) with that principal ideal, so \(I\) is principal as well. Squaring is therefore a bijection of the finite class group.\formalnote{4.30}{30}

Every lattice stable under \(\mathcal O_K\) is \(w\mathfrak a\) for a fractional ideal: after dividing by a nonzero period its rational span is \(K\), since it has dimension two over \(\mathbb Q\) and is stable under \(\omega\). Clearing one denominator gives an integral ideal. Bijectivity of squaring now permits the form \(w\mathfrak c^2\) with \(\mathfrak c\) integral, absorbing a principal factor into \(w\). Set \(e=N\mathfrak c\) and \(\Omega=we\). Then
\[
     w\mathfrak c^2=\Omega\,\mathfrak c^2/e,\qquad
     A(w\mathfrak c^2)=|\Omega|^2\sqrt d/2,\qquad
     \Omega e\in w\mathfrak c^2\subseteq (\Omega/e)\mathcal O_K .
 \tag{CG1}
\]
The last inclusion uses \(\mathfrak c^2\subseteq\mathcal O_K\), and the membership uses \(e\in\mathfrak c\) since \(|\mathcal O_K/\mathfrak c|=e\). This supplies both the area and a common denominator.\formalnote{4.31}{31}

\subsection{Action of the conjugated isogeny on periods}
Use \(b,x,Y,G,\theta,\sigma\) as in \S\ref{sec:sigma-identities}, and subscript \(C\) for a curve. If \(b_{C'}=\alpha b_C\), the Taylor series \(x_{C'},Y_{C'}\) at zero are the coefficientwise conjugates of the old ones. In fact \(x(0)=Y(0)=0\), and
\[
    x'=2Y+bx,\qquad Y'=1+3x^2-bY .
\]
The second equation follows by differentiating the cubic and canceling \(x'\) on its nonzero set, then extending by analyticity. Equating the coefficient of \(z^j\) recursively determines those of \(z^{j+1}\), since the characteristic is zero. The theta unit
\[
             Q_C=\frac{\operatorname{Tay}_0\theta_{C'}}
                           {\alpha(\operatorname{Tay}_0\theta_C)}
                  \in 1+z\mathbb C[[z]]
\]
appears for each factor of degree one. Indeed a monomial of weight \(m\) in \(\theta,\sigma^2,G\) has series \(\theta^m x^hY^c\). Its transported series equals \(Q_C^m\) times its coefficientwise conjugate. Expanding in the basis from \S\ref{sec:sigma-identities} extends this rule to every section.\formalnote{4.32}{32}

We also need to transport vanishing at the marked point \(t\). At ordinary points of the chart \(\theta\ne0\), transport \((x,Y)\) by \(\alpha\) on the cubic, use complex elliptic uniformization for the resulting point and compare the values of \(\theta\) by a single factor \(\mu\ne0\). The generators \(\sigma^2\) and \(G\), of degrees two and three, acquire the factors \(\mu^2\) and \(\mu^3\), respectively. At \(\theta=0\) use the homogeneous equation from \S\ref{sec:cm-integral}, which there reads \(G^2=(\sigma^2)^3\) with both entries nonzero. If \(u,v\) are the ratios of the transported weight-two and weight-three coordinates, then \(v^2=u^3\); taking \(\mu=v/u\) gives both \(\mu^2=u\) and \(\mu^3=v\). Consequently evaluation of a transported degree-\(m\) section is \(\mu^m\) times the conjugate of the old evaluation, also at \(t\).\formalnote{4.33}{33}

Suppose \(a\Lambda_C\subseteq\Lambda_D\) is an isogeny of degree \(N>0\) between the marked data. The pullback rule for sections at this level gives degrees \(2N,N\) for \(\sigma_D(az)^2,\theta_D(az)\). Transport them to sections \(U,V\) on \(C'\); \(U(0)=0\), \(V(0)=1\). If \(D'\) has parameter \(\alpha b_D\), the recursion for \(x=\sigma^2/\theta^2\) gives, as an identity of entire functions,
\[
       U(z)\theta_{D'}((\alpha a)z)^2
                 =\sigma_{D'}((\alpha a)z)^2 V(z)^2 . \tag{CG2}
\]
To see the equality without interchanging conjugation and evaluation, take the two Taylor series at zero. The factors \(Q_C^{2N}\) and the target theta units on the two sides cancel using \(\sigma_D^2=x_D\theta_D^2\). Equal series give equal entire functions. At \(l\in\Lambda_{C'}\) the multipliers give \(U(l)=0\), \(V(l)\ne0\). Hence \((\alpha a)l\in\Lambda_{D'}\). Similarly if \((\alpha a)t_{C'}\in\Lambda_{D'}\) then \(\theta_{D'}\) is nonzero there and \(U(t_{C'})=0\), so \(at_C\in\Lambda_D\) by transport of evaluation.\formalnote{4.34}{34}

For existence of the conjugate datum take a uniformizing lattice for the nonsingular cubic with invariants \(\alpha g_2,\alpha g_3\). The conjugate of \((\wp(t),0)\) supplies a marked half-period, still satisfying \(\wp''(t')=2\). The preceding period test can first be applied to these normalized marked cubics before imposing CM on the new lattice. Applied to each \(a\in\mathcal O_K\), which is fixed by \(\alpha\), it gives \(\mathcal O_K\)-stability. The marked-point test for \(a=\omega,\bar\omega\) preserves the two nonperiod conditions. In particular if \(\Lambda_D=u\mathfrak b\Lambda_C\), applying the period test to every \(u b\) with \(b\in\mathfrak b\) yields \((\alpha u)\mathfrak b\Lambda_{C'}\subseteq\Lambda_{D'}\). The reverse inclusion follows by applying the same test to \(u^{-1}\mathfrak b^{-1}\), since the fractional ideal is invertible. This proves the lattice equality used when the CM areas are multiplied over ideal classes.\formalnote{4.35}{35}

\subsection{Eliminating the unknown theta unit}
We only need the following rescaling for an isogeny whose conjugate also has degree \(N\). At a zero of \(\theta_{D'}((\alpha a)z)\), the other target section \(\sigma_{D'}\) does not vanish, so (CG2) forces \(V=0\). These zeros are simple. Division gives an entire elliptic function since the two sections on \(C'\) have the same multipliers. Their values at zero are one. Thus \(V(z)=\theta_{D'}((\alpha a)z)\). Comparing its two Taylor series gives
\[
                        Q_D((\alpha a)z)=Q_C(z)^N .
\]
Take \(C=D\), \(a=\omega\). Its degree on both lattices is \(N=\omega\bar\omega\) by the area ratio, and \(\alpha\omega=\omega\). If \(Q_C=1+c z^j+\cdots\) first differs from one in degree \(j\), this gives \(\omega^j c=Nc\). Now \(|\omega|^2=(1+d)/4>1\), so \(j=2\). But \(\omega^2=\omega\bar\omega\) would make \(\omega\) real. Therefore \(Q_C=1\), proving exact conjugation of every unshifted Taylor series of the generators, with the same formal coordinate on both sides.\formalnote{4.36}{36}

\subsection{One finite adele for the torsion translates}
Suppose \(\Lambda_C=\Omega_0\mathcal O_K\). For a translate on \(C\) put
\[
 H_C(u,z)=\frac{\pi\bar u z}{A(\Lambda_C)},\qquad
 (T_u\theta_C)(z)=e^{-H_C(u,z)-H_C(u,u)/2}\theta_C(z+u).
\]
Here \(A(\Lambda_C)\) is the area of a fundamental parallelogram. Translates on the transported curve \(C'\) use \(H_{C'}\) with its own area \(A(\Lambda_{C'})\). For a finite set \(S\subset K\) choose a positive integer \(n\) with \(nS\subset\mathcal O_K\), and write \(\Omega_0=2n^2w\). Apply the point comparison in the reverse direction at \(w\), obtaining \(w'\) on \(C'\). Evaluations of sections at \(w,w'\) match up to the factor by degree. Pulling a section back by \(a\in\mathcal O_K\) gives the same rule at \(aw,aw'\) with the degree multiplied by \(N(a)\). In particular the zeros of the three generators there match. Apply this to \(\sigma^2\) and \(a=2n^2\): \(v=2n^2w'\) is a period of \(C'\).

For \(s\in S\), put \(p=n^2 s w,\ p'=n^2s w'\). The pulled-back functions \((T_{2p}\theta_C)(nz)\), \((T_{2p'}\theta_{C'})(nz)\) belong to the degree \(n^2\) spaces by the torsion multiplier law. If the first vanishes, \(nz+2p+t_C\in\Lambda_C\); twice this equality shows \(z=(b-2ns)w\) for some \(b\in\mathcal O_K\). At this point the theta argument before adding \(t_C\) is \(nbw\). Transport of zeros at \(nbw,nbw'\) now shows that the transported second function vanishes at each of the simple zeros of the first. Their ratio extends to an entire elliptic function and is constant, say \(c\). To determine it, evaluate at \(z=-nsw\), so \(nz=-p\). The Gaussian in \(T_{2p}\theta(-p)\) is one and the value is \(\theta(p)\); the unshifted \(\theta(nz)\) at the same point also gives \(\theta(p)\) by evenness. On the conjugate curve the same comparison gives \(\theta(p')\), and the two degree-\(n^2\) sections have the same evaluation factor. When \(\theta(p)\ne0\) cancellation gives \(c=1\). When \(\theta(p)=0\), the matching value \(\theta(p')\) also vanishes, so \(p,p'\) represent the marked half-periods. The normalized value of each translate by \(2p,2p'\) at zero is \(-1\) by its multiplier, again giving \(c=1\). Cancel the invertible formal rescaling \(z\mapsto nz\). Thus the \emph{one} period \(v\) works for all \(s\in S\).\formalnote{4.37}{37}

Take \(\Omega,e\) for \(C'\) from (CG1). Each such period is \(\Omega b_S\) with \(b_S\in e^{-1}\mathcal O_K\). The net indexed by finite subsets ordered by inclusion has a cluster point \(\beta\) in the compact set \(e^{-1}\widehat{\mathcal O}_K\). Fix \(s\in K\) and take \(m>0\) with \(ms\in\mathcal O_K\). A coset of \(2me\,\widehat{\mathcal O}_K\) around \(\beta s\) contains \(b_Ss\) for some \(S\ni s\), by the cluster property. Every rational \(s'\) in that coset satisfies \(s'-b_Ss\in 2me\,\mathcal O_K\). The difference of the shifts \(\Omega s'\) and \(\Omega b_Ss\) is twice \(m\) times a period, and \(m\Omega b_Ss\) is a period. The semicharacter and the commutator phase in the twisted translation law are then both one. Thus \(T_{\Omega s'}\theta_{C'}\) has precisely the conjugated Taylor series of \(T_{\Omega_0s}\theta_C\), for every such \(s'\). This is the finite-adelic transport in the construction. Its background includes uniformization of a nonsingular complex cubic and the isogeny pullback rule for theta multipliers at the marked level.\formalnote{4.38}{38}


\section{The CM multiplier and conjugation of every finite translate}
\label{sec:cm-multiplier}

The elliptic calculation in \S\ref{sec:cm-conjugation} transports a one-variable Taylor series. The derivative estimate uses a stronger statement. A single finite adele must transport every translation in every standard hermitian space over the same field, and the transported series must have the period prescribed by the estimate. We derive both statements here.

\subsection{Pullback in the normalized theta trivialization}
Let \(A\) be the Euclidean covolume of the period lattice \(\Lambda\), put \(H(u,v)=\pi\bar u v/A\), and put \(E_\Lambda(u,v)=\operatorname{Im}(\bar u v)/A\). Thus \(E_\Lambda(\ell,m)\) is integral for \(\ell,m\in\Lambda\), by the determinant formula for area. For a weight \(r\) section use the translation
\[
    T_u^{(r)}f(z)=e^{-r(H(u,z)+H(u,u)/2)}f(z+u).
\]
Appell--Humbert classifies the line bundle by a Hermitian form and a semicharacter. In the cited convention the term \(\pi H_{\mathrm{BK}}(z,u)\) in the factor of automorphy is our \(H(u,z)\). A line bundle of degree one on an elliptic curve has a one-dimensional space of holomorphic sections by Riemann--Roch. The reduced section of the divisor \([0]\) is odd and satisfies \(\sigma_\Lambda'(0)=1\) after normalization \cite[Thm.~1.7 and Ex.~1.9, eq.~(8)]{legacy-bannai-kobayashi}. We normalize its translate below by dividing by its value at the marked point.

Write \(K=\mathbb Q(\sqrt{-d})\), \(d\equiv7\bmod 8\) prime, and \(\omega=(1+i\sqrt d)/2\). The marked CM data have \(\mathcal O_K\Lambda\subseteq\Lambda\), \(2t\in\Lambda\), and \(t,\omega t,\bar\omega t\notin\Lambda\). For the normalized odd reduced sigma function \(\sigma_\Lambda\), its lattice sign \(\epsilon(\ell)\) is \(1\) on \(2\Lambda\) and \(-1\) elsewhere in \(\Lambda\). Translating it to the marked point gives
\[
   \theta(z)=e^{-H(t,z)}\frac{\sigma_\Lambda(z+t)}{\sigma_\Lambda(t)},
   \qquad
   T_\ell^{(1)}\theta=\chi(\ell)\theta,\qquad
   \chi(\ell)=\epsilon(\ell)\exp(2\pi i E_\Lambda(\ell,t)).
\]
The sign rule for \(\epsilon(\ell+2t)\), together with \(\epsilon(2t)=-1\), gives \(\chi(\ell)=-\epsilon(\ell+2t)\). In particular \(\chi(\ell)=-1\) exactly when \(t+\ell/2\in\Lambda\).\formalnote{4.39}{39}

For \(a=\omega,\bar\omega\) we have \(1-a=\bar a\) and \(a\bar a/2=(d+1)/8\in\mathbb Z\). If \(t+a\ell/2\) were a period, then
\[
   (1-a)(t+a\ell/2)-(a\bar a/2)\ell=\bar a t
\]
would also be a period. Thus \(\chi(\omega\ell)=\chi(\bar\omega\ell)=1\). Apply the semicharacter rule to \(\ell=\omega\ell+\bar\omega\ell\). Alternation gives \(E_\Lambda(\bar\omega\ell,\omega\ell)=E_\Lambda(\bar\omega\ell,\ell)\), whence
\[
   \chi(\ell)=(-1)^{n_\Lambda(\ell)},\qquad
   n_\Lambda(\ell)=E_\Lambda(\bar\omega\ell,\ell)
                  =\frac{|\ell|^2\sqrt d}{2A}\in\mathbb Z .
   \tag{MT1}
\]
This calculation applies separately to the original marked curve and to the marked conjugate constructed in \S\ref{sec:cm-conjugation}.\formalnote{4.40}{40}

Let \(a\ne0\) satisfy \(a\Lambda\subseteq\Lambda'\) with \(N=[\Lambda':a\Lambda]\). Comparing covolumes gives \(|a|^2A=NA'\), hence \(H'(au,av)=NH(u,v)\) and \(n_{\Lambda'}(a\ell)=Nn_\Lambda(\ell)\). Substitution proves the two exact pullback laws
\[
 \chi'(a\ell)=\chi(\ell)^N,\qquad
 T_u^{(rN)}[f(a\,\cdot)](z)=(T_{au}^{\prime(r)}f)(az).
 \tag{MT2}
\]
In particular the degree becomes \(rN\). For the torsion shift used in interpolation take \(a=\bar\delta\), an endomorphism, and \(\delta u\in\Lambda\). The commutator of \(T_{\bar\delta\ell}^{(1)}\) with \(T_u^{(1)}\) is \(\exp(2\pi i E_\Lambda(\bar\delta\ell,u))=1\), since \(E_\Lambda(\bar\delta\ell,u)=E_\Lambda(\ell,\delta u)\). Consequently \(z\mapsto(T_u^{(1)}\theta)(\bar\delta z)\) also has degree \(N=\delta\bar\delta\) and character \(\chi^N\). Evaluating (MT2) at zero commutes with transport of the normalized values. The isogeny and torsion-translation bounds of \S\ref{sec:cm-integral} therefore apply to the shifted pullback as well. Formula \textup{(CM2)} gives its integral power-series factor.\formalnote{4.41}{41}

\subsection{From the elliptic translates to a single adelic multiplier}
Take \(E=K(\zeta)\), \(q_1=\operatorname{ord}(\zeta)\), and a standard positive hermitian space of rank \(n\). For \(0\leq k<q_1\), a coordinate \(j\), and the chosen complex embeddings \(\Phi_w\) extending the given embedding of \(K\), put
\[
 v_k(y)=-\frac{\operatorname{Tr}_{E/K}(y\bar\zeta^k)}{q_1},
 \qquad
 \lambda_{kj}(\Omega,c)=\frac{\Omega c}{q_1 A_0}
         \sum_w\overline{\Phi_w(\zeta^k)}Z_{wj},\qquad
 A_0=\sqrt{2\pi/\sqrt d}.
\]
Write \(J_c f\) for the Taylor series of the entire theta function of \(f\) at scale \(c\). For a finite vector \(x\) let \(\rho_f(x)\) be its Heisenberg translation. In this paragraph a rational \(\gamma\in E^n\) denotes its image in the finite translation coordinates of the chosen adelic frame. The product construction of \S\ref{sec:frame-proof} yields a vector \(F\) whose translates satisfy
\[
    J_c(\rho_f(\gamma)F)
       =\prod_{k,j}
         \left[(\operatorname{Tay}_0 T_{\Omega v_k(\gamma_j)}^{(1)}\theta)
                        \bigl(\lambda_{kj}(\Omega,c)\bigr)\right]^{q_1}.
 \tag{MT3}
\]
Indeed the frame uses \(\zeta^k/q_1\) repeated \(q_1\) times, the adjoint coordinate of \(-\gamma_j\) is \(v_k(\gamma_j)\), and differentiation along a frame vector substitutes the indicated linear form. At \(\gamma=0\) each factor has value \(1\). The finite Fourier decomposition of this product has a nonzero coefficient, since its theta value is \(1\). Projection to that character expresses a nonzero multiple of the vacuum as a sum of translates of \(F\). A further translation and the Heisenberg composition law express every translate of the vacuum in the same span. Thus \(F\) generates the entire finite part.\formalnote{4.42}{42}

Fix \(\eta\in\operatorname{Aut}(\mathbb C/K)\). Use \(\Omega_0,\Omega'\) for the periods of the original and conjugate curves, \(P_0\) for the period factor of \S\ref{sec:frame-proof}, and \(F'\) for the conjugate curve's product. The Taylor scale is \(a_{\rm jet}=P_0 A_0/\Omega_0\), and its counterpart for the conjugate curve is \(c_\eta=\eta(P_0)A_0/\Omega'\).
The scalar in \(\lambda_{kj}\) is then \(P_0/q_1\) on one side and \(\eta(P_0)/q_1\) on the other. Moreover \(\overline{\Phi_w(\zeta^k)}=\Phi_w(\zeta^{-k})\). Hence applying \(\eta\) to these coefficients just permutes the embeddings. Write \(D_\eta\) for the variable renaming that expresses the new series in the old order. By the area and period formula in \S\ref{sec:kronecker-limit},\formalnote{4.43}{43}
\[
    |c_\eta|^2=\mathcal B^{-1},\qquad
    \mathcal B=\frac{2\pi}{\operatorname{rd}_E}
      \exp\bigl(\gamma_{\rm Euler}-\operatorname{Re}(L'/L)(1,\chi_d)\bigr).
\]

From \S\ref{sec:cm-conjugation} there is a finite adele \(\beta\) of \(K\) with the following property. For each \(v\in K\), the \(\eta\)-image of the elliptic Taylor series at shift \(\Omega_0 v\) equals the new elliptic Taylor series at shift \(\Omega' v'\) for every \(v'\in K\) in a suitable adelic neighborhood of \(\beta v\). Let \(b\) be the image of \(\beta\) in \(\mathbb A_{E,f}\). For a rational \(\gamma\), impose these neighborhood conditions on \(a v_k(\gamma_j)\), simultaneously for all \(k,j\), with \(a\) varying near \(\beta\). Also require \(\rho_f(a\gamma)F'=\rho_f(b\gamma)F'\). This last condition is open by smoothness. Their intersection contains a rational \(a\in K\) by density in the finite adeles. Since each \(v_k\) is \(K\)-linear, the elliptic comparison applies to every factor on the right of (MT3) with new argument \(a\gamma\). Each linear form has zero constant term, so coefficientwise transport commutes with formal substitution. We conclude
\[
       \eta\bigl(J_{a_{\rm jet}}(\rho_f(x)F)\bigr)
          =D_\eta J_{c_\eta}(\rho_f(bx)F')\qquad (x\in\mathbb A_{E,f}^{\,n}).
 \tag{MT4}
\]
We have proved the identity on rational vectors. Their images are dense and both translated vectors are locally constant in \(x\), proving (MT4) everywhere. The choice of \(b\) came from the elliptic curve and works in every rank.\formalnote{4.44}{44}

\subsection{The finite Heisenberg phase and the derivative}
Define \(L\) on the finite part by extending \(\rho_f(x)F\mapsto\rho_f(bx)F'\) \(\eta\)-semilinearly. Any linear relation between these generators maps, by (MT4), to a relation between their new jets. Jets at nonzero scale determine a finite vector. Indeed the zero series restricts to zero on each complex line. The one-variable identity theorem makes the entire theta function zero there, and its values on the archimedean slice determine the vector by rational translation and density. Since \(c_\eta\ne0\), the proposed map is well-defined and
\(\eta(J_{a_{\rm jet}}f)=D_\eta J_{c_\eta}(Lf)\) on the finite part.\formalnote{4.45}{45}

For translation of an arbitrary generator one further scalar must agree. Write \(\psi_f\) for the normalized finite additive character in the Heisenberg product. Use (MT4) in rank one with the same \(b\). Smoothness gives an integer \(m>0\) with \(\rho_f(m)F=F\) and \(\rho_f(bm)F'=F'\). Given \(z\in\mathbb A_{E,f}\), write \(z=mx\). In the original model compare \(\rho_f(2x)\rho_f(m)F\) with \(\rho_f(2x+m)F\). The product law supplies the factor \(\psi_f(mx)\). In the new model the corresponding factor is \(\psi_f(b\bar b\,mx)\). Equation (MT4) for \(2x\) and \(2x+m\) relates the jets in both comparisons. The translate of \(F\ne0\) has a nonzero jet, so we can cancel it coefficientwise. We obtain
\[
                  \eta(\psi_f(z))=\psi_f(b\bar b\,z).
 \tag{MT5}
\]
Now the product law gives \(L\rho_f(x)=\rho_f(bx)L\) on each generator in every rank, hence on the whole finite part. Taking the constant coefficient of the translated jet gives \(\eta(f(\operatorname{fin}x))=(Lf)(\operatorname{fin}(bx))\), where \(\operatorname{fin}x\) has vanishing infinite part. Taking the coefficient of any monomial gives its Taylor coefficient at the prescribed conjugate scale, with places permuted. Multiplication by the multi-index factorial converts these coefficients to ordinary partial derivatives, and \(\eta\) fixes that integer.\formalnote{4.46}{46}

Finally let \(g\) act at the finite places, write \(\omega_f(g)\) for its Weil operator and \(\Delta_f(g)\) for its vacuum multiplier. We claim that on the finite part
\[
  L\omega_f(g)=r_\eta(g)\,\omega_f(g)L,\qquad
  r_\eta(g)=\frac{\eta(\Delta_f(g))}{\Delta_f(g)},\qquad |r_\eta(g)|=1 .
  \tag{MT6}
\]
The self-dual lattice for \(\psi_f\) is \(\widehat{\mathcal O}_E\). By (MT5) and injectivity of \(\eta\), \(b\bar b\,z\) is integral exactly when \(z\) is integral, by testing their characters against every integral adele. Thus \(b\bar b\) is a local unit everywhere. Write \(b=ua\) with \(u\bar u=1\) and \(a\) an integral adelic unit. At a place fixed by conjugation \(b\) itself is a unit. On an exchanged pair take the two components of \(u\) to be \(b_w\) and \(1/\overline{b_w}\), where the bar identifies the two local fields. The resulting \(b/u\) is a unit at both places. Replace \(L\) temporarily by \(L_0=\omega_f(u)^{-1}L\), which scales translations by \(a\). Formula (MT5) also transports the quadratic phase in the normalized finite translation \(H(x)\), giving \(L_0 H(x)=H(ax)L_0\). The lattice translations fix the vacuum line and \(a\) permutes the lattice, so \(L_0v_0\) lies on that line.\formalnote{4.47}{47}

In the finite sum (VP2) of \S\ref{sec:vacuum-projection}, \(a\) permutes the cosets of the intersection lattice and commutes with \(g\). The factor \(\Delta_f(g)/h\), with \(h=|\Delta_f(g)|^2\) the integral lattice index, acquires precisely the ratio \(r_\eta(g)\). This proves (MT6) on the vacuum for \(L_0\), and covariance with \(H(x)\) extends it to the span. Since \(u\) is central in the finite unitary group, the identity also holds for \(L\). For the modulus use \(\overline{\Delta_f(g)}=\upsilon(\det g)\Delta_f(g)\) and \(\eta(h)=h\). Their quotient gives \(r_\eta(g)^2=\upsilon(\det g)/\eta(\upsilon(\det g))\). Every conjugate of the finite value of this algebraic unitary character has modulus one by its weight law at points with identity infinite component. This proves the last assertion and accounts for the scalar in the derivative norm estimate.\formalnote{4.48}{48}


\section{The constant term of the lattice zeta function}
\label{sec:kronecker-limit}
The common period in \S\ref{sec:cm-integral} involves a logarithmic derivative. To see why the area and the discriminant enter it with their particular exponents, consider a complex lattice \(\Lambda\) of area \(A\), its Weierstrass discriminant \(\Delta_\Lambda=g_2^3-27g_3^2\), and an arbitrary normalization \(k>0\). For \(1<s<2\) put
\[
 F(s)=(s-1)\sum_{0\ne l\in\Lambda}\left(\frac{|l|^2}{kA}\right)^{-s}.
\]
\subsection{The Laurent coefficient in the period normalization}
Write \(\Lambda=a(\mathbb Z+\tau\mathbb Z)\), \(y=\Im\tau>0\), so \(A=|a|^2y\). Set
\[
 Z_\tau(s)=\sum_{(m,n)\ne(0,0)}
           \left(\frac{|m+n\tau|^2}{y}\right)^{-s}.
\]
This is the Epstein zeta function of the positive binary quadratic form of determinant one. Kronecker's first limit formula gives its residue \(\pi\) and its constant term
\[
 2\pi\left(\gamma-\log 2-\log(\sqrt y\,|\eta(\tau)|^2)\right)
\]
at \(s=1\) \cite[Ch.~I, \S1, Theorem~1]{theta-siegel}. Both signs of every nonzero lattice vector occur in this convention. Since \(F(s)=(s-1)k^s Z_\tau(s)\), multiplication of Laurent series gives \(F(1)=k\pi\) and
\[
  F'(1)/F(1)=\log k+2\gamma-2\log2-\log y-4\log|\eta(\tau)|.
\]
For the Weierstrass invariants of the lattice the discriminant is
\(\Delta_\Lambda=(2\pi/a)^{12}\eta(\tau)^{24}\), by the normalized modular discriminant product. Substituting this identity and using \(A=|a|^2y\) yields
\[
 \boxed{\displaystyle
     F(1)=k\pi,\qquad
     \frac{F'(1+)}{F(1)}
       =2\gamma-\log\frac{A}{k\pi^2}-\frac16\log|\Delta_\Lambda| .}
 \tag{KL4}
\]
The combination of area and discriminant is unchanged on scaling \(\Lambda\) by a complex scalar. This checks that the normalized Epstein sum depends only on the shape of the lattice.
\phantomsection\label{route:8}\leanroute{Lean derives the constant term by subtracting an integral from each row of the lattice sum, expanding the remaining cotangent series and differentiating the beta factor at one. The limit formula gives exactly that Laurent coefficient for the same sum. Multiplication by \(k^s\) and the discriminant conversion above are still needed for the arithmetic application.}\formalnote{4.49}{49}

\subsection{The mean over ideal classes}

In the application \(d\) is prime, \(d\equiv7\pmod8\), \(K=\mathbb Q(\sqrt{-d})\), and \(k=2/\sqrt d\). Thus \(-d\) is the field discriminant. A fractional ideal \(\mathfrak a\) has covolume \(\sqrt d\,N(\mathfrak a)/2\) in \(\mathbb C\), by computing the area of the integral basis \(1,(1+\sqrt{-d})/2\) and scaling by the fractional ideal norm. For \(\Lambda=c\mathfrak a\), a nonzero vector \(l=c\alpha\) thus contributes the inverse \(s\)-th power of \(N((\alpha)\mathfrak a^{-1})\). The latter ideals are precisely the integral ideals in the inverse class, each represented once per unit. The Euler product of their sum over classes is \(\zeta(s)L(s,\chi_d)\) times the number of units: the local factors of \(\zeta_K\) are \((1-p^{-s})^{-2}\), \((1-p^{-2s})^{-1}\), and \((1-p^{-s})^{-1}\) at split, inert and ramified primes. All products here converge absolutely. After multiplication by \(s-1\), the logarithmic derivative on the right is \(\gamma+\Re(L'/L)(1,\chi_d)\). The values on the left are all \(k\pi\), so their logarithmic derivatives are averaged with equal weight. This is the logarithmic form of the Chowla--Selberg class product. Gross \cite[\S3]{upper-gross-periods} gives the discriminant product by taking the logarithmic derivative of \(\zeta_K\) at \(s=0\). The calculation here compares the same ideal classes at \(s=1\), giving the common archimedean scale in terms of \((L'/L)(1,\chi_d)\).\formalnote{4.50}{50}


\section{Local exponents for translated theta functions}\label{sec:cm-detail}
\subsection{The Frobenius congruence for the translated theta function}

Let \(C\) be the normalized CM lattice for \(\mathbb Q(\sqrt{-d})\),
\(d\equiv7\bmod8\) prime, and fix a nonsplit \(p\). Work over the valuation
ring \(\mathcal O\) of \(\mathbb C_p\). At an inert prime take \((\pi,q)=(p,p^2)\), and at the ramified prime \(p=d\) take \((\pi,q)=(\sqrt{-d},p)\). Thus \(\pi=\pi_p\), \(q=q_p\) and \(\rho=v_p(\pi)/(q-1)\), as in the derivative estimate. In both the inert and the ramified case
\(|q|=|\pi|^2\), and \(p\) is odd. Write \(\theta(0)=1\),
\(g=-\log\theta\), and \(X(z)=\wp(z+t)-\wp(t)=z^2+\cdots\), where \(t\)
is the chosen point of order two. Write \(y\) for the translated formal
coordinate and \(\mathcal I=\mathcal O[[y]]\), as a subring of series in
\(z\). Membership in \(\mathcal I\) implies that coefficients of degrees
less than \(q\) are integral. In all degrees it implies integrality
after \(z\mapsto\varpi z\) with \(v_p(\varpi)\ge\rho\).

We use two inputs from the CM curve. First, multiplication by \(\pi\)
pulls back \(\theta\) to an even section bounded by the \(q\)-th power of the local height of \S\ref{sec:cm-integral}. Its degree is \(q\), so (CM2) gives a polynomial
\(D\) such that
\[
 \theta(\pi z)=\theta(z)^qD(X(z)),\qquad 2\deg D<q .
\]
Second, \(X\in\mathcal I\), \(g''(z)=g''(0)+X(z)\), and the normalization
gives \(\sqrt{-d}\,[z^2]g\in\mathcal O\). These facts come from elliptic
functions, their CM periods and the integral model at \(p\).\formalnote{4.51}{51}

Here is a proof of the stronger assertion \(D-1\in\pi\mathcal O[X]\).
Put \(g_n=[z^n]g\). For \(3\le n<q\),
\(n(n-1)g_n=[z^{n-2}]X\) is integral. In the inert case \(n<p^2\):
at most one of \(n,n-1\) is divisible by \(p\), to order at most one.
At \(p=d\), both are units since \(n<p\). Hence \(|\pi g_n|\le1\).
The quadratic coefficient has the same bound by the normalization above, and
the two lower coefficients vanish. Since \(|q-\pi^n|\le|\pi|^2\) for
\(n\ge2\), all coefficients of \(qg(z)-g(\pi z)\) below degree \(q\)
are divisible by \(\pi\). Exponentiating preserves this congruence:
for \(k\ge1\), \(v_p(\pi^{k-1}/k!)\ge0\) by Legendre's formula and
\(v_p(\pi)\ge1/(p-1)\). Thus \(D(X(z))-1\) is divisible by \(\pi\)
in the same degrees. If \((D-1)_i\) is its \(i\)-th polynomial
coefficient, the coefficient of \(z^{2i}\) is \((D-1)_i\) plus an
integral combination of its predecessors. Induction covers the whole
polynomial because \(2\deg D<q\). Taking the formal logarithm of
\(1+\pi H\) with \(H\in\mathcal I\) and zero constant term now yields
\[
                     qg(z)-g(\pi z)\in\pi\mathcal I.             \tag{F1}
\]
Indeed \(v_p(\pi^{k-1}/k)\ge0\). The polynomial degree bound was what
allowed the low-degree argument to prove a statement about the full series.

\subsection{Projecting the exponent, including the ramified layer}

Let \(g^\flat\) contain exactly the terms of \(g\) of degree
\(p+1\bmod(q-1)\). It equals the average of
\(\zeta^{-p-1}g(\zeta z)\) over \(\zeta^{q-1}=1\). The CM transformation
law gives \(\exp(\zeta^{-p-1}g(\zeta z)-g(z))\in\mathcal I\).
Their product is \(\exp(g^\flat-g)^{q-1}\). The unique root with constant
coefficient one belongs to \(\mathcal I\) since \(q-1\) is prime to \(p\).
This step uses the transformation law for the CM theta function under these
roots of unity. The proof obtains it by approximating the roots and their
norm characters by CM integers.

Recall the parameter \(T=z+O(z^q)\) of \S\ref{sec:cm-integral}, supported in degrees \(1\bmod(q-1)\). The coefficient \(b_1\) of \(T^q\) in \(T(\pi z)\) is a unit. With \(\varpi^{q-1}=-\pi/b_1\), the scaled parameter \(Y=T(\varpi z)/\varpi\) belongs to \(\mathcal O[[z]]\) and reduces to \(\sum_{j\ge0}z^{q^j}\) modulo \(\pi\mathcal O\).\formalnote{4.52}{52}

The least nonnegative representative \(i\) of \(p+1\bmod(q-1)\) is
\(p+1\) if \(p\) is inert and \(2\) if \(p=d\). Project (F1)
to this class in \(\mathcal O[[T]]\). Each \(T^h\) has degrees
\(h\bmod(q-1)\), so the result has the form
\[
 qg^\flat(z)-g^\flat(\pi z)=
                \pi\{uT^i+T^{i+q-1}J\},\qquad u\in\mathcal O,\ J\in\mathcal I .
\]
Set \(c_\varpi=\varpi^{p+1}/p\), \(\Xi=-g^\flat(\varpi z)/c_\varpi\).
The scalar \(-\pi\varpi^i/(qc_\varpi)\) is a unit. For inert \(p\) it
is \(-1\), and at ramified \(p\) its norm is one by the choice of \(\varpi\).
Substituting and dividing shows that for some \(\lambda\in\mathcal O\),
\[
 (1-\pi^n/q)[z^n]\Xi\ \equiv\
                   \lambda[z^n]Y^i\pmod{\pi\mathcal O}.          \tag{F2}
\]
The factor \(\varpi^{q-1}\) makes the higher \(T\)-terms vanish modulo
\(\pi\). For \(n>2\), \(\pi^n/q\in\pi\mathcal O\), so the multiplier is a
unit congruent to one. Terms of \(\Xi\) below degree \(i\) are already
zero. At \(p=d,n=2\) the multiplier is \(2\), since \(\pi^2=-p\), and
\([z^2]Y^2\equiv1\). Thus (F2) proves integrality of \(\Xi\) and
\[
 \Xi\equiv
 \begin{cases}
      \lambda Y^{p+1}&p\ \text{inert},\\
      \lambda(Y^2-z^2/2)&p=d
 \end{cases}\quad\bmod\pi .                                      \tag{F3}
\]
At an inert prime \(Y\) has one term for each even Frobenius exponent modulo \(\pi\), and
\(Y^p\) has one for each odd exponent. Every pair of such exponents has
a unique smaller exponent \(j\) and a positive odd difference \(e\).
Their product is \((z^{p^e}z)^{p^j}\). In the ramified case both lists
contain all exponents. Unequal pairs occur twice, equal pairs once,
and the first diagonal term has coefficient \(1/2\) after the correction
in (F3). This proves the support and coefficients in the one-curve normal
form. Indeed \(\theta(\varpi z)=W\exp(c_\varpi\Xi)\) with
\(W=\exp(g^\flat-g)(\varpi z)\) integral. Replacing \(\Xi\) by its explicit
sum modulo \(\pi\) changes \(W\) by an integral series because
\(v_p(c_\varpi\pi)=1/(p-1)\).\formalnote{4.53}{53}

\subsection{The unit coefficient and the limit of one-curve divisibility}

In the inert case the value of \(\lambda\bmod\pi\) matters for sharpness.
Write \(h_j=[z^{p-1}]X(z)^j\). From the Weierstrass equations
\(X'^2=4X^3+a_1^2X^2+4X\) and \(X''=6X^2+a_1^2X+2\), differentiation
gives, modulo the maximal ideal of \(\mathcal O\),
\[
 (4j+6)h_{j+2}+(j+1)a_1^2h_{j+1}+(4j+2)h_j=0 .
\]
The expression before reduction is the coefficient of \(z^{p-1}\) in
\((X^jX')'\). It equals \(p[z^p](X^jX')\), and the coefficient \([z^p](X^jX')\) is integral because \(X^jX'\in\mathcal I\) and \(p<q=p^2\).
Let \(p=2s+1\). We have \(h_0=0\) and \(h_s=1\), since \(X=z^2+\cdots\).
If \(h_1\) reduced to zero, the recurrence would force the same for all
\(h_j,\ j\le s\), as \(4j+6=2(2j+3)\) is a unit for \(j+2\le s\).
Thus \(h_1\) is a unit. The identity for \(g''\) gives
\([z^{p+1}]\Xi=-p\,g_{p+1}=-h_1/(p+1)\), also a unit.
In (F3) the coefficient of \(z^{p+1}\) in \(Y^{p+1}\) is \(1\) modulo
\(\pi\). This proves that \(\lambda\), hence every leading odd coupling
coefficient in the one-curve expansion, is a unit.

The \rootlink{grh/src/New/Derivatives/Sharpness/BoundAttained.lean\#L34-L42}{source note on one-curve sharpness} records that for one curve the inert-prime slope is attained in infinitely many degrees. The following Lubin--Tate calculation exhibits the slope explicitly. Let \(\mathcal F\) be a Lubin--Tate formal group for the unramified quadratic extension of \(\mathbb Q_p\), with formal logarithm \(\log_{\mathcal F}(T)\). There is a period \(\omega_p\) of valuation \(1/(p-1)-1/(p^2-1)=p/(p^2-1)\) for which the character \(\exp(\omega_p\log_{\mathcal F}(T))\) has coefficients in the valuation ring of \(\mathbb C_p\) \cite[Appendix, theorem~(a), (c)]{upper-fourier}. In the logarithm coordinate its \(n\)-th derivative at zero is \(\omega_p^n\). We now need a bound in each prescribed degree for the ring \(\mathcal I\) of the chosen elliptic coordinate.
Let \(C_n=\max_{F\in\mathcal I}|n![z^n]F|_p\) and put
\(\gamma_p=1/(p-1)-\rho\). The maximum exists: only the finitely many
monomials \(1,y,\dots,y^n\) can affect the coefficient and no integral
linear combination has larger norm. The integral group law writes
\(y(z_1+z_2)\) as an integral series in \(y(z_1),y(z_2)\). Taking the
\((m,n)\)-coefficient, multiplying by \(m!n!\), and applying the
ultrametric inequality therefore gives \(C_{m+n}\le C_mC_n\).
The integral differential structure gives \(C_n\le1\), \(C_0=1\).
For \(Q=q^j\), the coefficient \([z^Q]Y\equiv1\bmod\pi\) gives
\[
 |[z^Q]T|_p=p^{\rho(Q-1)},\qquad
 C_Q\ge |Q![z^Q]T|_p=p^{-\gamma_p(Q-1)} .
\]
The equality uses that \(Q\) is a power of \(p\). For any fixed
\(n>0\), write \(Q=hn+r\), \(0\le r<n\). Then
\(p^{-\gamma_p Q}\le C_Q\le C_n^h C_r\le C_n^h\). Letting \(j\) grow
gives \(C_n\ge p^{-\gamma_p n}\). The series attaining the bound can depend on \(n\).\formalnote{4.54}{54}

Section~\ref{sec:frame-attainment} uses this one-variable bound on the degrees left after the Frobenius pairs have been chosen. For twelve copies of the frame this gives the valuation bound \textup{(V5)}. Fourier inversion then gives the indicator estimate \textup{(V8)}, with one scalar valid at every prime.\formalnote{4.55}{55} Chapter~\ref{ch:lower} uses this attainment estimate to \hyperlink{call.9.14}{bound the heights of twisted Hecke values} in the proof of Theorem C.



\editionchapter{Integral frames and the definite theta lift}\label{ch:upperlift}

To apply the elliptic-curve calculation in many variables, we write the Hermitian form as a sum of rank-one forms. An integral choice of vectors keeps track of the denominators in this change of coordinates. Example~\ref{ex:cyclotomic-frame} shows how roots of unity give such a sum. The estimates at each prime concern forms on the definite spaces \(X_j\), whereas Theorem A concerns images in the spaces of forms on \(Y_j\). A period identity for the two chains, called the seesaw identity, lets us use the definite estimates on those images.

\section{Constructing the theta frame and paying the factorial losses}
\label{sec:frame-proof}
The normal form of the theta series is obtained by expressing the
Hermitian form as a sum of rank-one forms and applying the calculation
for a single CM curve to each summand. An arbitrary orthonormal basis
need not belong to the prescribed lattice. We therefore seek lattice
vectors whose rank-one operators sum to \(M\) times the identity,
where \(M=p_0p_1\) is fixed by the level. At each prescribed finite embedding, at least one of the resulting linear forms must have unit coefficients in every coordinate. This unit condition is needed to attain the predicted valuation on a coefficient of the product of one-variable theta functions. The other possible loss comes from distributing the orders of differentiation among the factors of their product. We bound the resulting factorials by keeping track of the pairs \(\tau,-p^e\tau\) at which both weights are large.

\begin{example}[A resolution from roots of unity]\label{ex:cyclotomic-frame}
Let \(\zeta\) be a primitive cube root of unity, and let
\(v_j=(\zeta^j,\zeta^{2j})^T\) for \(j=0,1,2\). The two coordinates
are the two cyclotomic embeddings of \(\zeta^j\). Each diagonal entry
of \(v_jv_j^*\) is \(1\). The upper off-diagonal entry is
\(\zeta^{-j}\), so that summing over \(j\) gives
\[
             \sum_{j=0}^2 v_jv_j^* =3I_2.
\]
Thus \(3I_2\) is a sum of rank-one matrices with integral entries in \(\mathbb Z[\zeta]\). Using \(t\) copies gives \(3tI_2\). More generally,
the powers of a primitive \(Q\)-th root give \(QI_f\), where
\(Q=q_1\) and \(f=\varphi(Q)\). A second family with sum
\(hI_f\) and \((h,Q)=1\) allows us to reach the prescribed large \(M\)
by combining copies of the two families.
\end{example}

\subsection{A trace-dual frame with length coprime to the cyclotomic level}
Write \(Q=q_1,\ \zeta=\zeta_Q,\ f=\varphi(Q)\), and let \(P(X)\) be the
\(Q\)-th cyclotomic polynomial. The progression modulus is \(q_0\), as in Chapter~\ref{ch:upper}. Put \(D=P'(\zeta)\) and
\(\beta=(1+\sqrt{-d})/2\). The discriminants of \(K\) and
\(\mathbb Q(\zeta)\) are coprime. The coprime-discriminant compositum
theorem therefore gives \(\mathcal O_E=\mathcal O_K[\zeta]\), and the
power-basis formula gives the relative different \((D)\).
Moreover \(D\mid Q\): differentiating \(X^Q-1=P(X)H(X)\) at \(\zeta\)
gives \(Q\zeta^{Q-1}=DH(\zeta)\). Put \(j=\lceil\log_2 Q\rceil\).
Every rational prime dividing \(Q\) divides
\(\beta\bar\beta=(d+1)/4\). Indeed \(d+1\equiv q_0/2\bmod q_0\),
\(16\mid q_0\), and \(Q/q_0\) has only the primes \(2,3\), already
dividing \(q_0\). The exponent of any prime in \(Q\) is at most \(j\).
Hence
\[
                   D\mid Q\mid(\beta\bar\beta)^j .
\]

At a place dividing \(\beta\) the lattice \(\mathfrak a_0\) is
\(D^{-1}\mathcal O_{E,w}\); elsewhere it is \(\mathcal O_{E,w}\).
It contains \(1\) and is stable under \({\rm Gal}(E/K)\), since that
group fixes \(\beta\) and preserves the different. Its product with its
conjugate is the codifferent. To check at a place where \(D\) is a
nonunit, its rational prime divides \(Q\), so it divides
\(\beta\bar\beta\); it cannot divide both \(\beta\) and \(\bar\beta\)
since \(\beta+\bar\beta=1\). Conjugation preserves the ideal of \(D\).
At every other place both sides are the unit ideal. In particular, if
\(y\in\mathfrak a_0\), its trace norm
\(\nu(y)=\operatorname{Tr}_{E/K}(y\bar y)=\sum_{\tau\in\Phi}|\tau(y)|^2\)
is integral in \(K\). It is real, so it is a nonnegative rational integer.\formalnote{5.1}{1}

For each \(w\in E\) let \(v(w)\) be its column of CM embeddings.
A family of length \(h\) here means one satisfying
\(\sum v(w)v(w)^*=hI_f\), counting multiplicity. The family of
all \(Q\) powers of \(\zeta^{-1}\) has length \(Q\). A diagonal
entry is \(Q\), and every off-diagonal entry is the sum of a
nontrivial geometric progression of \(Q\)-th roots. More generally,
for any \(y\in\mathfrak a_0\), the family
\[
             \zeta^a\gamma(y),\qquad 0\le a<Q,\quad
                       \gamma\in{\rm Gal}(E/K),
\]
has length \(Q\nu(y)\). Summing over \(a\) kills the off-diagonal
entries, and summing over \(\gamma\) gives \(\nu(y)\) on each diagonal.
Every member still belongs to the lattice.

Write \(b_i=[X^i]\bigl(P(X)/(X-\zeta)\bigr)\),
\(0\le i<f\), and put
\(e_i=\bar\beta^j b_i/D\). These vectors belong to \(\mathfrak a_0\).
At a place over \(\beta\), \(De_i\) is integral. At every other place
\(\beta\) is a unit and \(D\mid(\beta\bar\beta)^j\), so \(e_i\)
itself is integral there. Lagrange interpolation of \(X^k\), for
\(0\le k<f\), at the conjugates \(\zeta_\tau=\tau(\zeta)\) gives
\[
     X^k=\sum_{\tau\in\Phi}\zeta_\tau^k
           \frac{P(X)}{(X-\zeta_\tau)P'(\zeta_\tau)}.
\]
Comparing coefficients gives
\(\operatorname{Tr}_{E/K}(\zeta^k b_i/D)=\delta_{ik}\).
The two square matrices of these embeddings are inverses in either
order. It follows that
\[
   \sum_{i=0}^{f-1}v(e_i)v(\zeta^{-i})^*=\bar\beta^j I_f .
\]
Start with all \(Q\) powers of \(\zeta^{-1}\) and all the orbits just
constructed for \(y=e_i\). For each \(i<f\), replace the member
\(\zeta^{-i}\) and the member \(e_i\) of its indicated orbit by
\(\zeta^{-i}+e_i\). The change in the sum of outer products is the
last displayed matrix plus its adjoint. The new frame therefore has
length
\[
       h=Q\left(1+\sum_{i<f}\nu(e_i)\right)+
                    \beta^j+\bar\beta^j.                         \tag{T1}
\]
The last two terms sum to an integer. At any prime of \(K\) over a
prime divisor of \(Q\), the reductions of \(\beta,\bar\beta\) are
\(0,1\) in some order. Hence \(h\equiv1\) modulo each prime divisor
of \(Q\). This gives \((h,Q)=1\), and positivity of the outer products
gives \(h>0\). The member \(\zeta e_{f-1}\) survives the replacement:
it occurs in the orbit at exponent \(a=1\), and the replacements used
only \(a=0\).\formalnote{5.2}{2}

Here the sizes can be checked directly along the family of
\S\ref{sec:data}, with \(q_0,d\) fixed. Write \(Q=Tq_0\).
As every prime of \(T\) already divides \(q_0\),
\(P(X)=\Phi_{q_0}(X^T)\): the roots of the right side are precisely
the primitive \(Q\)-th roots, and its degree is \(\varphi(Q)\).
Thus \(|\tau(D)|\ge T\min_{\xi}|\Phi'_{q_0}(\xi)|\), where the minimum
over primitive \(q_0\)-th roots is positive. In particular it exceeds
one for sufficiently large \(x\). All coefficients of \(P\) have
absolute value at most \(2^{q_0}\) by expanding \(\Phi_{q_0}\) in
its roots. Synthetic division at a point on the unit circle gives
\(|\tau(b_i)|\le Q2^{q_0}\). From (T1) we obtain
\[
    h\le Q\{1+f^2(d^j Q2^{q_0})^2\}+2d^j .
\]
Here \(Q\) is polynomially bounded in \(x\) and \(j=O(\log x)\).
Consequently this expression grows at most polynomially in \(x\),
whereas \(b>\exp(x^\eta)\) for a fixed \(\eta>0\). Hence \(h\le b\).

Take \(R=6fb\). Since \(M\ge b^{20}\), eventually
\(M-R(Q+h)\ge Qh\). Choose \(0\le v<Q\) with
\(vh\equiv M-R(Q+h)\pmod Q\); after subtracting \(vh\), what remains
is a nonnegative multiple \(uQ\). Use \(R+u\) copies of each member of
the \(Q\)-frame and \(R+v\) copies of each member of the \(h\)-frame.
This gives (V1) with length exactly \(M\). Both \(1\) and
\(\zeta e_{f-1}\) occur.

For the period factor \(P_0=s_{\rm vac}/u\), the finite norm formula at a fixed
embedding \(\iota:\mathbb C\to\mathbb C_p\) is
\[
 |\iota P_0|_p=
 \begin{cases}
   |\phi_\tau(D)|_p &\text{if the place of }\phi_\tau\text{ divides }\beta,\\
   1&\text{otherwise},
 \end{cases}
 \qquad \phi_\tau(y)=\iota(\overline{\tau(y)}).
\]
Here the Chowla--Selberg unit \(u\) has finite norm one
and \(s_{\rm vac}^f=\chi_{\rm fin}(\alpha)^{-1}\). For this algebraic
Hecke character, the finite norm of \(\chi_{\rm fin}(\alpha)\) at
\(\iota\) is the product of the norms of the components of \(\alpha\)
under the \(\phi_\tau\). These norms are independent of \(\tau\):
all these embeddings agree on \(\beta\), and the \(K\)-conjugates
of \(D\) generate the same ideal. Taking the positive \(f\)-th root
gives the displayed formula from the specified components
\(D^{-1}\) and \(1\).\formalnote{5.3}{3}
Which case occurs
depends only on \(|\iota\bar\beta|_p\), since every CM embedding
restricts to the chosen embedding of \(K\). At places away from \(\beta\), the frame vector \(1\) gives a unit coefficient at every \(\tau\). At places dividing \(\beta\), the element \(\bar\beta\) is a unit and the
monic quotient of \(P\) has top coefficient \(b_{f-1}=1\).
Thus \(|\phi_\tau(\zeta e_{f-1})|_p=|\phi_\tau(D)|_p^{-1}\), which
shows that the frame vector \(\zeta e_{f-1}\) gives a unit coefficient for every \(\tau\) in this case.\formalnote{5.4}{4}

\subsection{Frobenius on the linear forms}
All prime divisors of \(Q\) split in \(K\): at such a prime the two
roots \(0,1\) of the reduction of the minimal polynomial of \(\beta\)
are distinct. For nonsplit \(p\) we may therefore use the unit class
\(-p^e\bmod Q\) to define \(t_e(\tau)=-p^e\tau\). Set
\(\pi=p\) for inert \(p\), and \(\pi=\sqrt{-d}\) when \(p=d\).
In \(\mathcal O_{\mathbb C_p}/(\pi)\) the two maps
\[
        y\longmapsto\iota(\tau(y))^{p^e},\qquad
        y\longmapsto\iota(\overline{t_e(\tau)(y)})
\]
agree on \(\mathcal O_E\) for every permissible exponent: odd \(e\)
in the inert case and every \(e\ge0\) at \(d\).
Indeed this quotient has characteristic \(p\), so both are ring maps.
On \(\zeta\), the minus sign naming the partner cancels complex
conjugation. On \(\sqrt{-d}\), at an inert prime the \(p\)-power
acts as the nontrivial automorphism of the quadratic residue field,
and an odd iterate still negates the square root. At the ramified
prime the square root reduces to zero. These checks determine the
maps on \(\beta\), since nonsplit primes are odd, and hence on the
ring generated by \(\beta,\zeta\). For a \(y\) integral at all
places above \(p\), clear denominators by a rational integer prime
to \(p\) and cancel its unit reduction. Applying the congruence to
\(\bar y\) gives
\[
      \iota(\overline{\tau(y)})^{p^e}
                 \equiv\iota(t_e(\tau)(y))\pmod{\pi}.
\]
At nonsplit \(p\), the vectors \(u_i=w_i/\bar\pi_M\) are integral
above \(p\): the possible denominators of \(w_i\) divide \(D\), and
the rational primes in \(M\) split in \(K\). Their resolution has
length one. After Frobenius, multiplying the linear form at the
first endpoint by the ordinary form at the second gives a coefficient
\(\sum_i c_i\,t_e(\tau)(u_i)\overline{\sigma(u_i)}
=\mathbf1_{t_e(\tau)=\sigma}\). This is precisely the coupling
\(\Gamma_e=\sum_\tau z_\tau^{p^e}z_{t_e(\tau)}\) in the residue
ring. At \(p=d\) the reduction was by the element of valuation
\(1/2\); retaining this modulus gives the factor \(1/2\) in the
local prime sum below.\formalnote{5.5}{5}

\subsection{The matching at small nonsplit primes}
Use \(\ell_0\) for the least power of three above \(27s\), and
\(l=(l_\tau)\) for a multi-index satisfying
\(\kappa_\tau\le l_\tau\le\kappa_\tau+x^{42}\).
The place weight \(\kappa_\tau\) is \(b\) on \(\Phi_b\) and \(a/x_1\)
elsewhere. For a nonsplit \(p\le k\) call \((\tau,e)\) a near pair
if \(e\in I_p\), \(2e<\ell_0\), and both \(\tau,t_e(\tau)\) belong
to \(\Phi_b\). Here \(I_p\) consists of the positive exponents allowed
by Frobenius. All \(b\)-places lie over \(-1\bmod q_0\), so the
existence of such a pair at \(d\) would imply
\(-d^e\equiv1\bmod q_0\). Since \(d\equiv q_0/2-1\) and \(4\mid q_0\),
\(d^2\equiv1\bmod q_0\). Neither \(1\) nor \(q_0/2-1\) equals
\(-1\bmod q_0\), proving that the near-pair set at \(d\) is empty.

At an inert prime the exponents are odd. Put \(g=-p\bmod Q\);
then \(t_e(\tau)=g^e\tau\) for these exponents. The sparseness proved
in \S\ref{sec:data} forbids
\(\tau,g^u\tau,g^v\tau\) from all being \(b\)-places whenever
\(0<u<v<\ell_0\), including coincidences among the three places.
Indeed the two ratios are \((-1)^u p^u,(-1)^v p^v\) in the prime-to-\(p\)
part of \(Q\), and \(p\le k\le kt\). Thus two near pairs with the same
start have the same exponent. If pairs of lengths \(e<e'\) had the
same end, then, starting from the start of the longer pair, the
points at distances \(0,e'-e,e'\) would contradict sparseness. Equal
lengths and equal ends instead force equal starts by cancellability.
Finally if an end were also a start, there would be three such places
at distances \(0,e,e+e'\), with \(e+e'<\ell_0\). These arguments also
exclude a loop. Each selected start and each selected end therefore
occurs in exactly one selected pair and the two sets are disjoint.\formalnote{5.6}{6}

There is a slightly stronger uniqueness needed for (V3).
Suppose an odd exponent \(e'\le e_\tau\) takes another selected start
to the end of the selected edge \((\tau,e_\tau)\). That edge of
length \(e'\) has two \(b\)-endpoints and \(2e'<\ell_0\), so it is
itself near and must be the very same edge by the preceding argument.
Now set \(j_\tau=\lfloor l_\tau/p^{e_\tau}\rfloor\) at selected starts
and \(j_\tau=0\) elsewhere. Allocate the degrees \(p^{e_\tau}j_\tau\)
to the start and \(j_\tau\) to the end. At a start the allocated degree
is at most \(l_\tau\). At an end it is at most
\(2b/p^{e_\tau}\le b\le l_{t_{e_\tau}(\tau)}\), since every
\(l_\tau\le2b\) for sufficiently large \(x\). Thus the allocation
fits coordinatewise. Call its two parts \(\nu',\nu''\), and call the
remainder \(\eta=l-\nu'-\nu''\).

Let \(s_0=\lceil6b/k\rceil,\ t_\tau=(l_\tau-s_0)_+\) for the next
calculation only, and define
\[
                B_\tau=\prod_{u=1}^{s_0}(t_\tau+u).
\]
We use \(v=v_p\) and the convention \(v(l!)=\sum_\tau v(l_\tau!)\).
The quotient \(B_\tau/s_0!=\binom{t_\tau+s_0}{s_0}\) is an integer.
Also \(l_\tau!\mid t_\tau!B_\tau\), with equality if \(l_\tau\ge s_0\).
Consequently \(v(B_\tau)\ge v(s_0!)\ge s_0/p-1\).

For \(e\in I_p\) let
\[
   R_\tau=\sum_{e\in I_p}\min\{\kappa_\tau/p^e,\kappa_{t_e(\tau)}\},
   \qquad S_p=v(\pi)\sum_{\tau\in\Phi}\sum_{e\in I_p}
                  (\kappa_\tau/p^e-\kappa_{t_e(\tau)})_+ .
\]
Both series over \(e\) converge by geometric comparison.
Since \(\gamma_p=v(\pi)\sum_{e\in I_p}p^{-e}\), subtracting the
positive parts gives the exact equality
\[
             \gamma_p\sum_\tau\kappa_\tau-S_p
                         =v(\pi)\sum_\tau R_\tau.                \tag{T2}
\]
The sum \(R_\tau\) has at most one near term, at most \(j_\tau+1\)
because \(\kappa_\tau\le l_\tau\). For any other \(e\) with
\(2e<\ell_0\), either the start has small weight, or its end does.
The corresponding minimum is at most \(a/x_1\). For the tail put
\(n=\lceil\ell_0/2\rceil\). Its sum is at most
\(2b/p^n\le2b/(kp)\), since \(p^{n-1}\ge2^s\ge k\).
Hence
\[
 R_\tau\le j_\tau+1+\ell_0\,a/x_1+2b/(kp).                    \tag{T3}
\]
This bound holds at \(p=d\) as well, with all \(j_\tau=0\).

For \(p\ge3,\ u\le2b\), the base-\(p\) digits of \(u\) give
\[
       \frac{u}{p-1}-(\log b+2)\le v(u!)\le\frac{u}{p-1}.
\]
Indeed \(v(u!)=(u-s_p(u))/(p-1)\), and \(s_p(u)/(p-1)\) is bounded
by the number of digits, at most \(1+\log_p(2b)<\log b+2\) for \(u>0\).
Apply the upper bound to \(l_\tau!\) and the lower bound to the
factorials of its three parts at \(\tau\). The resulting multinomial
valuation is at most \(3(\log b+2)\). Formula (V4) for \(\nu',\nu''\)
and the one-curve estimate for \(\eta\) therefore show
\[
      v(l!)+v_{12}(l)
              \le\gamma_p|l|-\sum_\tau j_\tau+3f(\log b+2).
\]
At \(p=d\) the two selected degrees are zero, so the same calculation
uses only the one-curve estimate. All nonsplit primes here exceed two.

To compare this inequality with \(S_p\), note the following direct
consequences of the scales from \S\ref{sec:data}:
\[
   3(\log b+2)+x^{42}+1\le a/x_1,\qquad
       1\le(\ell_0+1)a/x_1\le b/k^2 .
\]
For the first inequality, \(a/x_1=xb/x_1\) grows faster than every
power of \(x\). For the second, \(x_1\) eventually exceeds
\((\ell_0+1)xk^2\). For \(p\le k\) these bounds, followed by
\(s_0\ge6b/k\), give
\[
  3(\log b+2)+x^{42}+1+\ell_0a/x_1+2b/(kp)
             \le3b/(kp)\le s_0/p-1\le v(B_\tau).
\]
Here \(b/(kp)\ge1\) gives the middle inequality. In (T2) both
\(\gamma_p\) and \(v(\pi)\) are at most one. The change from
\(\kappa_\tau\) to \(l_\tau\) consequently costs at most \(x^{42}\)
per place. Inserting (T3) now proves
\[
                 v(l!)+v_{12}(l)\le S_p+\sum_\tau v(B_\tau)
\]
at every small nonsplit prime.\formalnote{5.7}{7}

\subsection{The large-prime allocation}
For \(p>k\) nonsplit, \(d\le k\) eventually, so \(p\) is inert.
Use every place with exponent one. The uniqueness for (V3) then
follows from the bijectivity of \(\tau\mapsto t_1(\tau)\) and the
fact that one is the smallest positive odd integer. Put
\[
     i_\tau=\lfloor t_\tau/p\rfloor,\qquad
     j_\tau=\min\{i_\tau,l_{t_1(\tau)}\}.
\]
We have \(i_\tau\le2b/p\le s_0\).
At place \(\sigma\) the outgoing allocation is at most
\((l_\sigma-s_0)_+\) and its unique incoming allocation is at most
\(\min(s_0,l_\sigma)\), whose sum with the first bound is \(l_\sigma\).
Use \(v_6(\nu'+\nu'')\le-\sum j_\tau\) from (V4), and the bound
\(v_6(\eta)\le0\) at the remaining inert degrees. It remains to
bound \(\sum_\tau[v(l_\tau!)-j_\tau]\).

Removing the first factor \(p\) from each multiple of \(p\) in
\(t_\tau!\) gives \(v(t_\tau!)=i_\tau+v(i_\tau!)\).
Since \(i_\tau\le s_0\), it follows that
\[
      v(l_\tau!)\le v(t_\tau!)+v(B_\tau)\le i_\tau+2v(B_\tau).
\]
Also \(s_0\ge x^{42}\) eventually, so \(t_\tau\le\kappa_\tau\).
Consequently
\[
 i_\tau-j_\tau=(i_\tau-l_{t_1(\tau)})_+
          \le(\kappa_\tau/p-\kappa_{t_1(\tau)})_+.
\]
The last expression is the first term of the local sum defining
\(S_p\), where \(v(\pi)=1\) at an inert prime. Summing proves
\(v(l!)+v_{12}(l)\le S_p+2\sum_\tau v(B_\tau)\).
At a split prime \(\gamma_p=S_p=0\), and the unit member of the frame
allows the one-curve estimate to supply the whole degree. This
establishes the arithmetic case split used in (V5).\formalnote{5.8}{8}

The finite absolute value of the period factor \(P_0\), the one-curve coefficient estimate and
the six-copy extraction in (V4) enter this derivation.
Section~\ref{sec:cm-integral} derives the period normalization, the inert and split local calculations and the extraction.
The frame construction and the treatment of near pairs above use
neither a uniform choice of the unit member over \(p\) nor a common
choice of the multi-index witnessing the local estimate.


\section{Attainment of the local exponent}\label{sec:frame-attainment}
\subsection{From the elliptic unit to a coefficient in twelve variables per place}

Write \(v=v_p\) with \(v(p)=1\), let \(\Phi\) be the CM type, and put
\(f=\#\Phi\). We prove how the unit \(\lambda\) calculated above is
detected by an actual theta datum. Three operations can affect its
valuation: substitution of the linear forms of a lattice frame,
extraction of an individual coefficient from an exponential, and
passage to the finite level at which the theta datum is specified.
We calculate them in that order.

There is a fixed family \(w_i\in\mathfrak a_0=E\cap\alpha\widehat{\mathcal
O}_E\), with integral multiplicities \(c_i\ge6fb\), such that
\[
 \sum_i c_i\Phi_\sigma(w_i)\overline{\Phi_\tau(w_i)}
                          =M\mathbf1_{\sigma=\tau}.
 \tag{V1}
\]
At each \(p\)-adic embedding some member \(w_{i_0}\) makes all
coefficients of \(\mu_{i_0}\) units, where
\(\mu_i(z)=P_0\sum_\tau\overline{\Phi_\tau(w_i)}z_\tau\) and \(P_0\)
is the normalized period factor. The underlying frame is independent
of the embedding. Its construction uses the roots-of-unity frame of
length \(q_1\) and a trace-dual frame of coprime length \(h\le b\).
Indeed, if \(R=\lceil6fb\rceil\) and
\(M\ge q_1h+R(q_1+h)\), subtract \(R(q_1+h)\) from \(M\) and choose
\(0\le v<q_1\) solving \(vh\equiv M-R(q_1+h)\pmod {q_1}\).
The remainder after subtracting \(vh\) is a nonnegative multiple of
\(q_1\). This supplies multiplicity at least \(R\) for both frames.
Section~\ref{sec:frame-proof} constructs the trace-dual frame and identifies a unit member. The possible unit member is either \(1\) or
a root of unity times the last adjusted trace-dual basis vector.\formalnote{5.9}{9}

Let \(\delta=\bar\pi_M\), so \(\delta\bar\delta=M\), and put
\(u_i=w_i/\delta\). Dividing (V1) by \(M\) gives a resolution of the
identity. At an inert prime \(p\), the \(u_i\) are integral and
\(\mu_i=s u_i^\dagger\), with \(s=P_0\bar\delta\) a \(p\)-adic
unit. For odd \(e\), reduction modulo \(p\) changes the
\(p^e\)-power of \(\overline{\Phi_\tau(u_i)}\) into
\(\Phi_{-p^e\tau}(u_i)\). This is the CM Frobenius congruence on the
cyclotomic extension. Setting \(p_e(\tau)=-p^e\tau\), we get
\[
 \begin{split}
 \sum_i c_i(u_i^\dagger z)^{p^e}(u_i^\dagger z)
 &\equiv \sum_{\tau,\sigma}
     \left(\sum_i c_i\Phi_{p_e(\tau)}(u_i)
                   \overline{\Phi_\sigma(u_i)}\right)
                    z_\tau^{p^e}z_\sigma\\
 &=\sum_\tau z_\tau^{p^e}z_{p_e(\tau)}=:\Gamma_e(z)\pmod p .
 \end{split}
\]
The congruence also applies after higher Frobenius powers. Multiplying
the one-curve normal forms, their exponents add, and these congruences
replace that sum by the indicated \(\Gamma_e\)'s. The discrepancy
belongs to \(p\mathcal O[[z]]\) and has constant term zero. Its
exponential after multiplication by \(\varpi^{p+1}/p\) is integral:
\(v(\varpi^{p+1})=1/(p-1)\) and
\(v(n!)\le n/(p-1)\). It can therefore be absorbed into the integral
unit. The extra homothety \(s\) multiplies the coefficient of
\(\Gamma_e^{p^a}\) by the unit \(s^{(p^e+1)p^a}\). Hence
\[
 F(z):=\prod_i\theta(\varpi\mu_i(z))^{c_i}
       =W(z)\exp(\lambda_p\Xi(z)),\qquad
 \Xi=\sum_{\substack{e\ {\rm odd}\\a\ge0}}
           \lambda_{e,a}\Gamma_e^{p^a},\quad
 \lambda_p=\frac{\varpi^{p+1}}p .
 \tag{V2}
\]
Here \(v(\varpi)=\rho=1/(p^2-1)\), \(W\) is an integral unit,
all \(\lambda_{e,a}\) are integral, and every
\(\lambda_{e,0}\) is a unit. Thus the original unit controls a specific
coefficient even after summing over a whole frame.\formalnote{5.10}{10}

\paragraph{Isolating the unit.}
Write \(\Theta(z)=F(z/\varpi)\) and
\(\mathcal I_\mu=\mathcal O[[y(\mu_i(z)):i]]\) for the image
of the formal series ring under the indicated substitution.
Choose six \(p\)-adic integers \(\delta_0,\ldots,\delta_5\) with
\(\delta_0=1\) and \(\sum\delta_c^2=0\): take four whose squares
sum to \(-1\), and use \(0\) for the last entry. The elliptic
zero-square-sum product identity places \(\prod_c\theta(\delta_c z)\)
in \(\mathcal O[[y(z)]]\). Its constant term is one, so its inverse
also belongs to that ring. Substitution into every frame form and
multiplication show that \(\prod_c\Theta(\delta_c z)\) is an
invertible element of \(\mathcal I_\mu\). After scaling by \(\varpi\),
\[
 U(z)=\prod_{c=0}^{5}F(\delta_c z)\in\mathcal O[[z]],\qquad U(0)=1.
\]
These facts use integrality of the elliptic formal multiplication
law and the theta product for a zero sum of squares.\formalnote{5.11}{11}

Choose starting places \(A\subseteq\Phi\). To each \(\tau\in A\)
assign an odd \(e_\tau\) such that the only pair
\((\sigma,a)\) with \(\sigma\in A\), \(a\) odd,
\(a\le e_\tau\), and \(p_a(\sigma)=p_{e_\tau}(\tau)\) is
\((\tau,e_\tau)\). Write \(t_\tau=p_{e_\tau}(\tau)\).
Prescribe nonnegative counts \(j_\tau\). The exponents that we seek
in two sets of variables are
\[
 \nu'=\sum_{\tau\in A}p^{e_\tau}j_\tau[\tau],
 \qquad \nu''=\sum_{\tau\in A}j_\tau[t_\tau],
 \qquad J=\sum_{\tau\in A}j_\tau .
\]
Here \([\tau]\) is the unit multi-index. Consider
\[
 G(z',z'')=F(z'+z'')\prod_{c=1}^5F(\delta_c z').
\]
Only \(F(z'+z'')\) depends on \(z''_t\).
By (V2), \(\lambda_p^{-1}\partial_t F/F
=\partial_t\Xi+\lambda_p^{-1}(\partial_t W)/W\) is integral, and
reduces to \(\partial_t\Xi\) since \(v(\lambda_p)=1/(p-1)-1<0\).
Iterating, \(\lambda_p^{-J}\partial_{z''}^{\nu''}G=gG\), where
\(g\) is integral and \(g(z',0)\) reduces to
\(\prod_{\tau\in A}(\partial_{t_\tau}\Xi(z'))^{j_\tau}\).
Indeed, when the next derivative hits the previous integral multiplier,
its contribution contains \(\lambda_p^{-1}\), so vanishes in the residue
field.

Set \(z''=0\), so \(G=U\), and in computing the coefficient
at \(\nu'\) set \(z'_\sigma=0\) for \(\sigma\notin A\).
In the residue field differentiating \(\Xi\) kills the \(p^a\)-powers
with \(a>0\) and the differentiated \(p^e\)-power in every
\(\Gamma_e\). The least-degree term remaining in
\(\partial_{t_\tau}\Xi\) is therefore
\(\bar\lambda_{e_\tau,0}(z'_\tau)^{p^{e_\tau}}\), uniquely by the
condition on \(A\). The product has a unique initial monomial at
\(\nu'\); multiplication by \(U\) leaves its coefficient unchanged
modulo the maximal ideal. This proves the exact valuation
\[
 v\!\left((\nu''!)\,[{z'}^{\nu'}{z''}^{\nu''}]G\right)
                          =J\bigl(1/(p-1)-1\bigr).
 \tag{V3}
\]
In particular the selected terms cannot cancel among the different
places.

Let \(v_n(\nu)\le r\) mean that there are \(n\) factors
\(H_c\Theta\), \(H_c\in\mathcal I_\mu\), and multi-indices
\(\alpha_c\) with \(\sum_c\alpha_c=\nu\) for which
\(\prod_c[z^{\alpha_c}](H_c\Theta)\) has valuation at most \(r\).
Undoing the scaling in (V3) subtracts
\(\rho|\nu'+\nu''|\) from the valuation. Expanding the two-group
substitution by the multinomial theorem uses coefficients in
\(\mathcal O\); a sum with the norm in (V3) requires at least
one summand of that norm or larger. Thus for \(\nu=\nu'+\nu''\),
\[
 v_6(\nu)\le J/(p-1)-J-\rho|\nu|-v(\nu''!).
 \tag{V4}
\]
As \(|\nu|\ge(p+1)J\), this gives \(v_6(\nu)\le-J\). Also
\(v(\nu'!)\le|\nu'|/(p-1)\) and \(|\nu''|=J\) give
\(v(\nu'!\nu''!)+v_6(\nu)\le\gamma_p|\nu|-J\).\formalnote{5.12}{12}

For comparison, the unit-coefficient form \(\mu_{i_0}\) realizes
the one-variable derivative in every multi-index \(\eta\). Choose
\(h\in\mathcal O[[y]]\) with
\(v(h^{(|\eta|)}(0))\le\gamma_p|\eta|\) as proved above. Then
\[
 [z^\eta]h(\mu_{i_0}(z))
    =\frac{h^{(|\eta|)}(0)}{\eta!}
      \prod_\tau\mu_{i_0,\tau}^{\eta_\tau}.
\]
In copy zero take \(H_0=h(\mu_{i_0})/\prod_c\Theta(\delta_c z)\),
and take \(H_c=1\) in the other five. On the diagonal
\(z_c=\delta_c z\) the product of these six \(H_c\Theta\)'s is
\(h(\mu_{i_0})\), giving \(v(\eta!)+v_6(\eta)\le\gamma_p|\eta|\).
At an inert place the frame resolution also expresses
\(s z_\tau\) as an integral linear combination of the \(\mu_i(z)\).
Replace each \(\mu_i\) in this combination by \(y(\mu_i)\). Since
\(y(z)=z+O(z^2)\), the combinations start with \(s z_\tau\).
Their product in the prescribed powers, multiplied by \(\Theta\),
has a unit coefficient of degree \(\eta\). This proves
\(v_6(\eta)\le0\) there.

The two sets of six copies are independent, so their bounds add.
For inert \(p>k\), take \(A=\Phi,e_\tau=1\); uniqueness follows
because \(\tau\mapsto-p\tau\) is a permutation. With
\(s_0=\lceil6b/k\rceil\) take
\(j_\tau=\min(\lfloor(l_\tau-s_0)_+/p\rfloor,l_{-p\tau})\).
The explicit degree inequality gives \(\nu\le l\) coordinatewise.
Use (V4) on \(\nu\) and the zero upper bound on \(l-\nu\).
For nonsplit \(p\le k\), first consider the case in which the near-pair set is empty.
This includes the ramified prime \(p=d\). Take the constant coefficient \(1\) in the
first six copies and apply the one-variable derivative bound to the full degree in the other six.
If a near pair occurs, \(p\) is inert. The uniqueness condition for the near pairs
between the \(b\)-places then permits \(j_\tau=\lfloor l_\tau/p^{e_\tau}\rfloor\).
Use the second bound from (V4) on their degrees and the one-variable bound on the remainder.
At split \(p\) the one-variable bound alone applies with \(\gamma_p=0\).
Section~\ref{sec:frame-proof} proves the simultaneous degree inequalities, sparsity of near pairs, and the comparison of their factorial valuations with \(S_p\). With \(B_\tau\) the
loss product of the preceding subsection, they give in all cases\formalnote{5.13}{13}
\[
          v_{12}(l)\le S_p+2\sum_\tau v(B_\tau)-v(l!).
 \tag{V5}
\]

\subsection{Realization by an adelic theta vector and Fourier inversion at the level}

Put \(r=S_p+2\sum_\tau v(B_\tau)-v(l!)\). Choose the factors and
degrees witnessing (V5). The \(y(\mu_i)\) have zero constant term,
so only monomials of degree at most \(|\alpha_c|\) in the series
defining \(H_c\) affect \([z^{\alpha_c}](H_c\Theta)\). The coefficients
of those series are integral. In each copy, one such monomial gives
a coefficient of norm at least the norm for \(H_c\Theta\).
Thus there are exponents \(e_{ci}\) with \(\sum_i e_{ci}\le|\alpha_c|\)
such that
\[
 v\!\left(\prod_{c=0}^{11}[z^{\alpha_c}]
      \left(\prod_i y(\mu_i(z))^{e_{ci}}\Theta(z)\right)\right)\le r.
 \tag{V6}
\]
Index \(c=0,\dots,11\) labels the copies here and \(c_i\) always
denotes the multiplicity of a frame vector. Since \(l_\tau\le2b\),
\(3 e_{ci}\le3|l|\le6fb\le c_i\).
This is where a lower bound on \emph{every} frame multiplicity is
needed. The section on the \(i\)-th elliptic factor may be written
\[
    s_{ci}=(y\,\theta^3)^{e_{ci}}\theta^{\,c_i-3e_{ci}}.
\]
The elliptic section \(y\theta^3\) is entire, including at the
poles of \(y\). Hence \(s_{ci}\) is an entire section of degree
\(c_i\); near zero its Taylor series is \(y^{e_{ci}}\theta^{c_i}\).
Although (V6) was obtained by using \(p\)-adic formal series, these
exponents and the actual elliptic sections now define complex
holomorphic data.

Here are the scale and lattice of the realization. Let
\(\Omega_0\mathcal O_K\) be the period lattice of the normalized
elliptic curve, let \(A_0=\sqrt{2\pi/\sqrt d}\), and set \(a_{\rm jet}=P_0 A_0/\Omega_0\), with the period factor \(P_0\) that occurs in the linear forms \(\mu_i(z)=P_0\sum_\tau\overline{\Phi_\tau(w_i)}z_\tau\) of (V1). On the standard space
\(V=\operatorname{Hom}(E,E^{12})\), pull the \(i,c\) section back by
\[
 L_{ic}(z)=\frac{\Omega_0}{\bar\delta A_0}
                     \sum_\tau\overline{\Phi_\tau(w_i)}z_{\tau c}.
\]
The CM Hermitian form on the elliptic line is
\((\pi/\operatorname{area})\,\bar u v\) with
\(\operatorname{area}=|\Omega_0|^2\sqrt d/2\).
Summing the pullbacks weighted by \(c_i\) gives the standard Fock
Hermitian form: the scalar multiplying each coordinate pairing is
\[
   \frac{\pi}{|\Omega_0|^2\sqrt d/2}\,
   \frac{|\Omega_0|^2}{M A_0^2}\,M=1
\]
by (V1). Moreover \(L_{ic}\) sends the real image of the lattice
\(\bar\delta\,\mathfrak a_0^{12}\) into the elliptic period lattice.
This uses the trace pairing of \(\mathfrak a_0\) with itself and
\(w_i\in\mathfrak a_0\). Consequently the translation laws of the
sections multiply to the signed theta translation law at this
lattice: the exponents in the Gaussian factors sum by the computed
Hermitian identity, and the lattice multipliers are the canonical
ones for the chosen trace-one element. Thus the entire product
\(\prod_{c,i}s_{ci}\circ L_{ic}\) is the theta function of a
finite-part vector \(v\) in the standard adelic Heisenberg model.
The existence of such a vector also follows from the finite inversion
below, using the one-dimensional space of theta functions at the
self-dual lattice.

For every finite argument \(\operatorname{fin}(x)\), the value
\(v(\operatorname{fin}(x))\) is integral. Indeed the finite part is
smooth, so a rational \(\gamma\) approximating \(x\) closely enough
gives exactly the same translate of \(v\). Evaluation is then a
product of twisted translates of \(s_{ci}\) at zero. Their shifts
are torsion: clear the denominators of \(\gamma\) by \(L\ne0\);
\(ML\) times each shift belongs to the period lattice by the same
trace pairing. Apply the CM integrality theorem for the normalized values of these monomial sections at torsion, including the translation factor, as derived in \S\ref{sec:cm-integral}. This proves the claim for
every finite argument, as opposed to a bound solely on the Taylor
coefficients. Finally,
\(L_{ic}(\bar\delta a_{\rm jet}z)=\mu_i(z_{\bullet c})\).
Let \(\mathcal T(v)\) denote the Taylor series of the theta function
at exactly this scale, and put
\(\beta_{\tau c}=(\alpha_c)_\tau\). Its variables on the twelve
copies are disjoint, so the coefficient of the product is the product
in (V6). We have proved\formalnote{5.14}{14}
\[
  v_p([z^\beta]\mathcal T(v))\le r,\qquad
  \sum_c\beta_{\tau c}=l_\tau .
\]

To extract a datum of the precise residue level let
\(\Lambda=\mathfrak a_0^{12}\) and use the finite groups
\[
    Y=\delta^{-1}\Lambda/\Lambda,\qquad
    \Gamma=\Lambda/\bar\delta\Lambda .
\]
The Heisenberg commutator \(\langle \gamma,u\rangle\) gives a perfect
pairing of these groups by self-duality of \(\Lambda\). Their orders
are equal; their product is
\([\delta^{-1}\Lambda:\bar\delta\Lambda]
=[\Lambda:M\Lambda]=M^{24f}\). Thus each has order \(M^{12f}\).
For a representative \(u\in Y\) form the unnormalized sum of the
signed \(\Gamma\)-translates of the translate of \(v\) by \(-u\).
It is invariant under the signed full lattice, hence lies in the
vacuum line; write it \(a_u v_0\), with \(v_0\) the ordinary vacuum.
Translate it back by \(u\) and sum over \(Y\). Commuting the two
translations multiplies each \(\gamma\)-term by
\(\langle\gamma,u\rangle\) or its inverse according to the sign
convention. The sum over \(u\) vanishes for every nonzero class
\(\gamma\), since translating that character sum by an element
where its character differs from one leaves the sum unchanged.
The trivial class contributes \(\#Y\,v\), proving
\[
     M^{12f}v=\sum_{u\in Y}a_u\,\rho(u)v_0.
 \tag{V7}
\]
The sum over \(\Gamma\) defining \(a_u v_0\) also gives an integral bound. At any finite argument \(x\),
\(a_u v_0(\operatorname{fin}(x))\) is a sum of the integral values
of \(v\) times roots of unity. Therefore
\(\lvert a_u v_0(\operatorname{fin}(x))\rvert_p\le1\).\formalnote{5.15}{15}

For \(u\in\delta^{-1}\Lambda\) the remaining translated vacuum has
the expansion
\[
       \rho(u)v_0=\sum_{y_0}\zeta_{u,y_0} v_{y_0},
       \qquad \lvert\zeta_{u,y_0}\rvert_p=1.
\]
Here \(y_0\) runs through \(12\times1\) residue matrices modulo
\(\pi_M=\bar\delta\). The datum \(v_{y_0}\) has the Gaussian at
infinity, the class indicator at the places of the modulus, and the
vacuum elsewhere. On each residue class the additive character of
the allowed translate is constant, since multiplying its parameter
by \(\delta\) makes it integral. This gives the indicated roots of
unity. Consequently for any linear functional \(L\), (V7) implies
\[
 \left|M^{12f} L(v) v_0(\operatorname{fin}(x))\right|_p
                  \le\max_{y_0}|L(v_{y_0})|_p .
\]
Thus a bound on \(L(v)\) yields an indicator \(v_{y_0}\) with \(|L(v_{y_0})|_p\) at least the left-hand side. The argument uses only the integrality of \(a_u v_0(\operatorname{fin}(x))\) and the unit absolute values of the phases \(\zeta_{u,y_0}\).\formalnote{5.16}{16}

Take \(L=\mathrm{vacuumScale}^{12}[z^\beta]\mathcal T\), and choose
one \(x\) so that \(c_0=\mathrm{vacuumScale}^{12}
v_0(\operatorname{fin}(x))\ne0\). This choice can be made before
choosing \(p\). For some \(y_0\) we now have
\[
 v_p\!\left(\mathrm{vacuumScale}^{12}[z^\beta]\mathcal T(v_{y_0})\right)
      \le r+v_p(M^{12f}c_0).
\]
Put \(t=M^{12f}c_0\prod_\tau B_\tau^2\). Multiplication by \(l!\)
and substitution of \(r\) from (V5) prove
\[
 v_p\!\left(\frac{\mathrm{vacuumScale}^{12}}{t}\,
                    l!\,[z^\beta]\mathcal T(v_{y_0})\right)\le S_p .
 \tag{V8}
\]
To see the global size cost, let \(s_0=\lceil6b/k\rceil\).
Each \(B_\tau\) is a product of \(s_0\) positive integers bounded by
\(2b+s_0\). The normalized vacuum value has
\(\log|\sigma c_0|\le12 C_P q_1^2\), by the vacuum-height estimate
with its field constant \(C_P\). It follows that
\[
 \log|\sigma t|\le12f\log M+2f s_0\log(2b+s_0)+12 C_P q_1^2
                   \le fmb/k
\]
by the admissible scale inequalities. The same nonzero \(t\) works
at all primes and embeddings. The level class \(y_0\) and the
distribution \(\beta\) of \(l\) across the copies are chosen locally.
Thus (V8) gives sharpness for the family of indicator data, with an
explicit common allowance in height.\formalnote{5.17}{17}

Section~\ref{sec:frame-proof} constructs the frame and estimates the
loss from places paired across its two ends.
Section~\ref{sec:hecke-average} proves the interpolation statement in
(H1), the flag count and volume estimate in (H2), and the invariant
averaging statement in (H3).


\section{The invariant average at a marked prime}\label{sec:hecke-average}

At an exceptional inert prime, \S\ref{sec:upper-A} uses a scaled Hecke operator on the source forms. To bound the colength of its image, we construct an averaging operator on the local polynomial coordinates. This operator removes monomials whose defect exceeds a specified bound. The defect is the weighted sum of exponents \(\Delta_i\) defined below. On the pairings used to compare integral lattices, the averaging operator acts as a scalar modulo a power of the uniformizer.
We construct the averaging functional first. We then insert it into the
Frobenius substitutions one matrix at a time. Two dimension estimates control how many directions can be lost. The first uses orders of vanishing of polynomials at residue-field points. The second counts monomials by their exponents and bounds the count by the volume of an interlacing cone. Section~\ref{sec:marked-divisors} combines the resulting estimates (H1) and (H2) with the scalar congruence (H3).

\subsection{From a marked place to the Frobenius cycle}

Fix an exceptional \(p\) and a marked place \(\lambda\). It is unramified
in \(E=K(\zeta_{q_1})\), since \(p\) is inert in \(K\) and prime to the
two discriminants here. Let \(z=-p\bmod q_1\), \(M=\operatorname{ord}(z)\)
and \(o=\operatorname{ord}(-p\bmod q_0)\). View the class \(z\) as an
automorphism fixing \(K\). If \(\kappa\) is complex conjugation,
\(z\kappa\) is arithmetic Frobenius at \(p\): on roots of unity it
is the \(p\)-power map, and on \(K\) it is the nonidentity automorphism.
The decomposition group is generated by this element. In the coordinates
\(\operatorname{Gal}(E/K)\times\langle\kappa\rangle\) it follows that
the elements of \(\operatorname{Gal}(E/K)\) carrying a place to itself
or to its conjugate are exactly the powers \(z^j\); even \(j\) give
the original place, odd \(j\) its conjugate. The primes of a mark are
distinct from their conjugates. In particular the two possibilities
are disjoint and \(M\) is even. This also follows from
\(m\mid M\) for exceptional \(p\) with \(m\) even. The residue degree
is the order of \(z\kappa\), namely \(\mathrm{lcm}(M,2)=M\).

Choose an \(a\)-place \(\tau\) inducing \(\lambda\), as supplied by the
mark. In the orbit \(z^j\tau\) the \(a\)-places are exactly those for
which \(o\mid j\), by reduction modulo \(q_0\). Thus
\(t_i=z^{oi}\tau\) gives a bijective numbering of the required
\(a\)-places by \(\mathbb Z/h\) with \(h=M/o\). Since \(o\) is odd,
the inducing primes alternate between \(\lambda\) and \(\bar\lambda\).
Put \(P=p^o\), \(q=P^h\). Write \(W=(\mathcal O_E)_\lambda\). For
even \(i\) take \(\iota_i=\Phi(t_i)\) on \(W\), and for odd \(i\) take
\(\iota_i=\Phi(t_i)\circ\kappa\). Both take values in the same
valuation ring \(O\) under the chosen \(p\)-adic embedding, because
the conjugation has made their inducing place \(\lambda\).
The residue field of \(W\) has cardinality \(q\), and, on integers
of \(E\), \(\iota_{i+1}=\iota_i\circ(z^o\kappa)\).
Since \(o\) is odd, \(z^o\kappa=(z\kappa)^o\).
The arithmetic Frobenius congruence therefore gives
\(\iota_{i+1}(u)\equiv\iota_i(u)^P\bmod pO\). To see that this holds
modulo \(p\) itself, it holds on the cyclotomic generator and on
\(\mathcal O_K\), and these generate \(\mathcal O_E\) by the
coprime-discriminant identity. The two maps to \(O/(p)\) remain equal
on localization since denominators outside \(\lambda\) are units.
Taking \(w\in O\) with \(w^2=p\) supplies the cycle used here.
It gives exactly \(h/2\) even letters and the Hecke scale
\(p^{a h/2}\), since \(a\) is the symmetric degree at those places.\formalnote{5.18}{18}

Locally let \(d=\operatorname{diag}(pI_r,I_{k-r})\). In the double
coset of \(d\), two elements \(\beta d,\gamma d\) have the same right
integral coset precisely when \(d^{-1}\gamma^{-1}\beta d\) is integral
with integral inverse. Conjugation divides the upper-right block by
\(p\) and multiplies the lower-left block by \(p\). The determinant
remains a unit, so integrality of the entries suffices for the inverse.
Thus \(\gamma^{-1}\beta\) must be in \(J\), the lower block parahoric.
Reduction of integral matrices onto the residue group is surjective
both in \(W\) and its completion, by lifting entries and the unit
determinant. This gives precisely the coset set
\(\mathrm{GL}_k(W)/J\) appearing in the finite counts. Section~\ref{sec:marked-intertwining} carries these local cosets
to the two endpoint maps on the integral source by changing the adelic
representatives and computing the transition on the multiplicity spaces.\formalnote{5.19}{19}

\subsection{An integral trace on one matrix variable}

Let \(W\) be a local ring with residue field of order \(q\) and let
\(\iota:W\longrightarrow O\), where \(O\) is a valuation domain of
characteristic zero in which the prime \(p\) is a nonunit. Fix a decomposition
of the \(k\) indices into blocks and let \(G(W)\) be the invertible
block-diagonal matrices. Give index \(i\) an integer weight \(\rho_i\).
In a monomial \(X^\alpha\) in the entries of one \(k\)-by-\(k\) matrix,
the row weight is \(u(\alpha)=\sum_{ij}\rho_i\alpha_{ij}\).
Call a monomial high if \(u(\alpha)>t_0\). Suppose
\[
                      D<p^b,\qquad D+k<q,\qquad C=p^{k^2 b}.
\]
There is an \(O\)-linear functional \(I\) on matrix polynomials, invariant
under \(X\mapsto\iota(g)^{-1}X\iota(g)\) for \(g\in G(W)\), which kills
every high monomial. Its further property concerns any ideal
\(\mathfrak j\subset O\) and any polynomial \(\psi\) of degree at most
\(D\). If every high coefficient in every such conjugate of \(\psi\)
lies in \(\mathfrak j\), then
\[
                         I(\psi)\equiv C\psi(1)\pmod{\mathfrak j}.
                         \tag{HA1}
\]
Here \(1\) is the identity matrix. The three properties will be needed
simultaneously.

Fix a degree \(d\le D\) and let \(U=(O^k)^{\otimes d}\), with basis words
\(e_j=e_{j_1}\otimes\cdots\otimes e_{j_d}\).
If \(j,i\) are two words, set \(\alpha(j,i)_{ab}=
\#\{\nu:j_\nu=a,\ i_\nu=b\}\). For \(|\alpha|=d\) the quantity
\(\eta_\alpha=C/\binom d\alpha\) belongs to \(O\). Indeed the digit-sum
formula gives
\[
 v_p\binom d\alpha
  =\frac{\sum_{ab}s_p(\alpha_{ab})-s_p(d)}{p-1}\le k^2 b .
\]
Every rational prime different from \(p\) is a unit in \(O\). For a
homogeneous \(\psi\) form the operator \(A_\psi\) with entry in row \(i\),
column \(j\) equal to \(\eta_{\alpha(j,i)}[X^{\alpha(j,i)}]\psi\) if
\(i_\nu,j_\nu\) belong to the same block for every \(\nu\), and zero
otherwise. For a generic block-diagonal matrix \(Z\),
\[
                  \operatorname{tr}(A_\psi Z^{\otimes d})=C\psi(Z).
                  \tag{HA2}
\]
In fact every surviving monomial occurs in the trace exactly
\(\binom d\alpha\) times. In particular
\(\operatorname{tr}_U A_\psi=C\psi(1)\).

For a single block over the residue field, the operators \(g^{\otimes d}\), with \(g\in{\rm GL}_k(\mathbb F_q)\), span the endomorphisms commuting with permutations of the \(d\) tensor factors whenever \(q>d\) \cite[Theorem~4.3]{upper-benson-doty}. We need the integral span of the operators \(\iota(g)^{\otimes d}\) for the prescribed blocks.
It consists of all matrices supported on matching block words whose
entries on that support depend only on \(\alpha(j,i)\). To see this,
list all degree-\(d\) monomials in the allowed block entries. Their
vectors of simultaneous values at \(g\in G(W)\) span the free \(W\)-module
on the list. Otherwise Nakayama's lemma gives a nonzero linear functional
on its reduction killing these evaluations. Its coefficients would give
a nonzero degree-\(d\) polynomial \(F\) over the residue field vanishing
on every invertible block matrix, since each such matrix lifts. Multiplying
by the product \(H\) of the block determinants gives a polynomial vanishing
on \emph{all} block matrices. It has total degree \(d+k<q\).
Successive one-variable interpolation forces \(FH=0\) as a polynomial.
Since \(H\ne0\), this contradicts \(F\ne0\). Taking transposes if needed,
and applying \(\iota\) to a spanning expression, proves the assertion
over \(O\). This span is an algebra since \(G(W)\) is a group.

In particular \(A_\psi\) belongs to this algebra. The left side of (HA2)
determines its operator within the algebra, since its coefficient at
\(Z^\alpha\) recovers the associated entry after multiplication by the
nonzero integer \(\binom d\alpha\). Cyclic permutation under the trace now
gives \(A_{\psi(g^{-1}Xg)}=g^{\otimes d} A_\psi (g^{-1})^{\otimes d}\),
where the matrices act through \(\iota\). Let \(U_+\) be the span
of all \(g^{\otimes d}e_j\) with \(\sum_\nu\rho_{j_\nu}>t_0\).
It is stable under the algebra and finitely generated, by the finite list
of word vectors and the finite set of coefficient matrices spanning the
algebra. A finitely generated torsion-free module over a valuation domain
is free. Consequently we can set
\[
       I_d(\psi)=\operatorname{tr}_U A_\psi-
                       \operatorname{tr}_{U_+}(A_\psi|_{U_+}).
\]
The covariance just proved makes both traces invariant.

If \(\psi\) is a high monomial, \(A_\psi e_j\) can be nonzero only
on high input words. All those vectors \(e_j\) belong to \(U_+\), so
\(A_\psi\) factors as \(U\to U_+\hookrightarrow U\). Interchanging the
two factors preserves the trace, even when the two free modules have
different ranks. Thus \(I_d(\psi)=0\). Under the hypotheses of (HA1),
apply \(A_{\psi(g^{-1}Xg)}\) to a high word. Only high monomials can
contribute, their coefficients are in \(\mathfrak j\), and each acts
in \(U_+\). Covariance, allowing \(g\) to vary over the group, gives
\(A_\psi U_+\subset\mathfrak j U_+\). Its restricted trace belongs to
\(\mathfrak j\). Summing \(I_d\) over the homogeneous components up to
degree \(D\), and setting it equal to zero in higher degrees, proves (HA1).\formalnote{5.20}{20}

\subsection{Inserting the average around the Frobenius cycle}

Number the positions of the Frobenius cycle by \(\mathbb Z/h\). We use the monomial pairing \(\langle\cdot,\cdot\rangle_b\) of \S\ref{sec:marked-divisors}, with \(h=2\ell\),
\(P=p^o,\ q=P^h,\ w^2=p\). The two blocks meet at rows \(a=r-1\) and \(r\). Call the corresponding variables at position \(s\) of the cycle \(X_s,Y_s\), taking \(Y_s\) from row \(r\) at even positions and from row \(a\) at odd positions. For a monomial with exponent vector \(\beta\), write \(\beta(X_s),\beta(Y_s)\) for its two exponents. Call \(\Delta_i(\beta)\) the defect of this monomial at cycle position \(i\) for the pair of rows \(a,r\). It is
\[
 \Delta_i(\beta)=\beta(Y_i)+\frac{1}{q-1}\sum_{s\bmod h}
 P^{(s-i)\bmod h}\bigl(\beta(Y_s)-\beta(X_s)\bigr).
\]
Here the exponent of \(P\) lies between \(0\) and \(h-1\). Let \(G\) consist of the two blocks of sizes \(r,k-r\).
The size condition used here is \(hk(\ell+2)n+n+k<P^\ell\). Hence
\[
          D=hnk(\ell+2)P^\ell,\quad D+k<q,\quad
          c_0=k^2oh,\quad L=hc_0,\quad c=5L,\quad t=3kc .
\]
At letter \(j\) use its map \(\iota_j:W\to O\) in the trace construction,
take \(b=oh\), and use row weights \(+1\) at \(r\), \(-1\) at \(a\)
and zero elsewhere. Set \(t_0=2tP^\ell\) and denote its functional by
\(I_j\). Write \(E(\psi)=p^{c_0}\psi(1)\).

Here is the polynomial being averaged. Introduce independent matrix
variables \(Z_j\) indexed by \(j\bmod h\). In letter \(s\) take the
product of the matrices \(Z_{s-d}^{[P^d]}\) for \(d\) in descending
order from \(\ell-1\) to zero, where brackets mean powers entry by entry.
For a deep row start at \(d=\ell\). In each case use the product in an
even letter and its transposed adjugate in an odd letter, taking the
specified row to form the letter matrix. Take the \(n\)-th symmetric
power in each letter, their tensor product over the letters, and apply
the resulting operator to a coefficient vector \(v\). Denote the result
by \(F(v)\). At \(Z_j=1\) for all \(j\), \(F(v)=v\); in each individual
copy its degree is at most \(D\), by multiplication of the factors and
the degree bound \(k-1\) for cofactors. Write \(V_Jv\) for the contraction
of \(F(v)\) with \(I_j\) at \(j\in J\) and \(E\) at \(j\notin J\).
In particular \(V_\varnothing v=p^L v\).

We spell out the two arithmetic estimates needed here. In a Levi frame
\(g\), express the same substitution on \(gv\) by transforming its
matrix factors, and contract in every copy except \(j\). Denote the
polynomial in coordinate \(\beta\) that remains by \(H^g_\beta(Z_j)\).
Conjugating the last variable in \(H^1\) gives the vector \(g^{-1}H^g\).
Indeed before the contractions, this is change of variables in each
matrix product; the even matrices transform by \(g\), the transposed
adjugates by \(g^{-t}\). Block diagonality preserves the deep and shallow
rows. Invariance of the functionals in the other copies permits the
change of their matrix variables. The estimates are
\begin{align*}
 [Z_j^\alpha]H^g_\beta\notin(w^c)
   &\ \Longrightarrow\
       \big|u(\alpha)/P^\ell-\Delta_{j+\ell}(\beta)\big|\le t,\\
 \Delta_{j+\ell}(\beta)>t,\quad v'\text{ in the dual all-edge lattice}
   &\ \Longrightarrow\
         \langle g^{-1}e_\beta,v'\rangle_b\in(w^c).
\end{align*}
If \([Z_j^\alpha]H^g_\beta\) survives modulo \(w^c\), then a monomial of the uncontracted substitution must survive there,
with factors surviving their functionals in each other copy.
Invariance forces the row and column weights in each such copy to agree.
Indeed, diagonal conjugation by
\(\operatorname{diag}((1+p)^{\rho_i})\) multiplies a monomial by
\((1+p)^{\mathrm{column\ weight}-\mathrm{row\ weight}}\).
A nonzero exponent gives a factor different from one in characteristic zero.
Here \(1+p\) is a unit in \(W\), whose residue characteristic is \(p\).
Thus the full monomial's copy weight, which includes weighted row-minus-column
terms from the other copies, reduces to \(u(\alpha)/P^\ell\).
The Frobenius-factor coefficient bound gives the first inequality in this
form. The second estimate is the adjoint root estimate for the halfway
lattice, with the two monomial pairings satisfying \(b_\beta b'_\beta=p^L\)
and the actions of \(g\) and \(g^t\) adjoint. These are the two specialized
inputs used below.\formalnote{5.21}{21}

Take \(V=V_{\mathbb Z/h}\). Whenever \(\Delta_{j+\ell}(\beta)>3t\),
the first estimate makes every coefficient of \(H^g_\beta\) outside
\((w^c)\) high. The remaining functional \(I_j\) kills all such monomials.
Thus every required coefficient of \(gVv\) lies in \((w^c)\), for every
Levi frame. For the pairing, fix \(j\notin J\), contract the other
copies as prescribed by \(J\), and form the scalar polynomial
\(\psi=\sum_\beta b_\beta v'_\beta H^1_\beta\). By change of frame, the
coefficient of \(Z_j^\alpha\) in its conjugate is
\[
             \sum_\beta [Z_j^\alpha]H^g_\beta\,
                               \langle g^{-1}e_\beta,v'\rangle_b .
\]
For high \(\alpha\), either \([Z_j^\alpha]H^g_\beta\in(w^c)\) or \(\Delta_{j+\ell}(\beta)>t\). In the latter case \(\langle g^{-1}e_\beta,v'\rangle_b\in(w^c)\). Formula (HA1) proves \((I_j-E)\psi\in(w^c)\), precisely the
change in the pairing on passing from \(V_J\) to \(V_{J\cup\{j\}}\).
Adding these congruences from the empty set to the full cycle gives
\[
              \langle Vv,v'\rangle_b\equiv
                       p^L\langle v,v'\rangle_b\pmod{w^c}.
\]
This yields the operator used in the factorization (H3).\formalnote{5.22}{22}

\subsection{The two coefficient estimates}

To compare the substitution polynomial with the defect \(\Delta_{i+\ell}(\beta)\) of the original monomial, we first assign weights to its variables. Fix a cycle position \(i\), use \([s-i]\in\{0,\ldots,h-1\}\) for cyclic differences,
and write \(\rho_b={\bf1}_{b=r}-{\bf1}_{b=a}\). Give entry \(Z_{j,bc}\)
weight
\[
 \lambda(j,b,c)=
 \begin{cases}
   \rho_b/P^\ell,&j=i,\\
   (\rho_b-\rho_c)P^{[j-i]}/(qP^\ell),&j\ne i .
 \end{cases}
\]
The weight of a monomial with exponent vector \(\nu\) is \(\lambda\cdot\nu\).
Take \(\epsilon=1/2\) or \(-1/2\) and use the following filtration on
polynomials: level \(v\) means that each coefficient at \(\nu\) is
divisible by \(w^{\max(0,\lceil v-\epsilon\lambda\cdot\nu\rceil)}\).
In a product the levels add, by termwise multiplication.

In a changed frame the \(d\)-th factor in letter \(s\) is
\[
 B_{s,d}=\iota_s(g)
   \big(\iota_{s-d}(g)^{-1}Z_{s-d}\iota_{s-d}(g)\big)^{[P^d]}
   \iota_s(g)^{-1}.
\]
Modulo \(p\) the entrywise power is a ring endomorphism and carries
\(\iota_{s-d}(g)\) to \(\iota_s(g)\) by the Frobenius cycle.
Thus \(B_{s,d}\equiv Z_{s-d}^{[P^d]}\pmod p\). Each entry is homogeneous
of degree \(P^d\). Since \(|\epsilon\lambda|\le P^{-\ell}\) and \(d\le\ell\),
all these monomials have weights in \([-1,1]\). Suppressing the main
monomial costs at most two levels, paid exactly by \(p=w^2\). More
precisely the entry \(bc\) has level \(P^d\epsilon\lambda(s-d,b,c)\).

Put \(u=(2P^\ell)^{-1}\) and introduce
\[
 A_{s,d}(b)=
      \frac{\epsilon\rho_b P^{d+[s-d-i]}}{qP^\ell}.
\]
These quantities sum to zero over \(b\), and \(|A_{s,0}(b)|\le u\).
One step from \(d\) to \(d+1\) has
\[
 A_{s,d+1}(b)-A_{s,d}(c)-u
           \le P^d\epsilon\lambda(s-d,b,c).
\]
If \(s-d\ne i\), the cyclic exponent drops by one and the formula
before subtraction of \(u\) is an equality. At \(s-d=i\), it wraps
from \(0\) to \(h-1\), yielding the main term
\(P^d\epsilon\rho_b/P^\ell\) and an error bounded by \(u\), since
\(P^d\le q\). The intermediate \(A\)'s cancel in a matrix product.
For \(m=\ell\) or \(\ell+1\), each entry of row \(b\) in
\(B_{s,m-1}\cdots B_{s,0}\) thus has level at least
\(A_{s,m}(b)-(\ell+2)u\). In the transposed adjugate, a row \(b\)
is formed from minors omitting row \(b\). The sum of the potentials
in the remaining rows is \(-A_{s,m}(b)\), and the total loss is at
most \(k(\ell+2)u\).

For each variable \((b,s)\) take \(m=\ell+{\bf1}_{(b,s)\ {\rm deep}}\)
and put \(\varepsilon_s=1\) at even and \(-1\) at odd letters.
Then \(\varepsilon_s A_{s,m}(b)\) is \(\epsilon(1-q^{-1})\) times
the defect weight of that variable at \(i+\ell\). For rows outside
\(a,r\) both sides vanish. For the two remaining rows the factor
\(P^\ell\) cancels, leaving \(P^{[s-i-\ell]}/q\) in the shallow
row and \(P^{1+[s-1-i-\ell]}/q\) in the deep row, with their
respective signs. At \(s=i+\ell\) the deep exponent wraps to \(h\),
giving the extra term of the defect. Passing to symmetric degree \(n\)
in every letter adds at most \(hnk(\ell+2)u\le1/2\) to the loss. Thus
the coefficient of a full monomial \(\nu\) in output coordinate \(\beta\)
is divisible by the nonnegative power of \(w\) with exponent
\[
 \max\!\left(0,\left\lceil
       \epsilon\big((1-q^{-1})\Delta_{i+\ell}(\beta)-\lambda\cdot\nu\big)
                       -hnk(\ell+2)u\right\rceil\right).
\]
If it survives modulo \(w^c\), applying both signs gives
\(|(1-q^{-1})\Delta_{i+\ell}(\beta)-\lambda\cdot\nu|\le 2c+1\).
The defect definition gives \(|\Delta_{i+\ell}(\beta)|\le2n\le q\).
Combining these inequalities proves the desired \(2c+2\le3kc=t\)
bound for \(c\ge1,\ k\ge2\); for \(c=0\) no coefficient survives.\formalnote{5.23}{23}

To bound the pairing \(\langle g^{-1}e_\beta,v'\rangle_b\), let \(d(\beta)\) be the ordinary deep degree
of a monomial, and write \(D_w\) for the diagonal operator
\(e_\beta\mapsto w^{d(\beta)}e_\beta\). A vector \(v'\) in the
all-edge lattice has a representative \(f\) and a nonnegative integer
\(m\) with \(w^m v'=D_w f\) and \(D_w\gamma f\in w^m O^I\) for every
\(\gamma\in\mathrm{GL}_k(W)\). Thus in every translate each term has
order in integral powers of \(w\), plus deep degree, at least \(m\).
Fix \(i\). Replace the height on each variable by its deep degree
minus one third of its defect weight at \(i\). From \(X_s\) to \(Y_s\)
the deep degree rises by one and the defect weight rises by
\({\bf1}_{s=i}+2P^{[s-i]}/(q-1)\), a number in \([0,3/2]\) in the
admissible range. Hence the new step is between \(1/2\) and \(1\).
If any rank-two translate violated level \(m\) for the new height,
take a term of smallest order plus new height among such translates.
There are only finitely many possible exponents, and only finitely
many integral orders below the putative cutoff for each exponent.
The extreme-term root lemma forces this term's defect to be
nonpositive. Its order plus ordinary deep degree is then at most its
order plus new height, which is below \(m\), a contradiction.
The argument applies also after any given Levi translation of \(f\).

The Levi action commutes with \(D_w\), since it preserves deep degree.
If \(\Delta_i(\beta)\ge3c\), the improved level for \(gf\) says that
its coefficient times \(w^{d(\beta)}\) is divisible by \(w^{m+c}\).
Cancel \(w^m\) to get \((gv')_\beta\in(w^c)\). Finally adjointness gives
\(\langle g^{-1}e_\beta,v'\rangle_b=b_\beta((g^{-1})^t v')_\beta\).
The transpose of \(g^{-1}\) is still in the Levi. This proves the
second estimate whenever \(\Delta_i(\beta)>3kc\). Here the remaining
rank-two input is precisely the extreme-term assertion before (H2).\formalnote{5.24}{24}

\subsection{Multiplicity interpolation in the depth estimate}

We give the algebraic implication used in (H1). Let \(K\) be a field,
\(R=K[Y_1,\ldots,Y_m]\), and assign \(Y_i\) weight \(1/e_i\) with
\(e_i\) positive integers. Suppose a finite set \(Z\subset K^m\) has
at least \(\prod_i e_i\) points. Assume that polynomials
\(H_i=Y_i^{e_i}-G_i\) vanish at every point of \(Z\), with every monomial
of \(G_i\) of weight strictly below one. If a nonzero \(F\) vanishes
to order at least \(\delta\) at each point of \(Z\), then a monomial of
\(F\) has weight at least \(\delta\).

Indeed, as \(a\) ranges over nonnegative integer vectors and
\(0\le b_i< e_i\), the family \(Y^b H^a\) is a \(K\)-basis of \(R\).
Each term has a single top-weight monomial \(Y^{b+(e_i a_i)_i}\);
those monomials are distinct and exhaust all monomials. Induction on
weight after clearing denominators proves spanning, and inspecting the
largest weight proves independence. The ideal \(J=(H_1,\ldots,H_m)\)
has \(J^\delta\) spanned by the basis terms with \(|a|\ge\delta\);
expand the coefficient of each product of \(\delta\) generators in the
basis to see the inclusion in this span. In particular
\(\dim R/J=\prod e_i\). Simultaneous evaluation maps \(R/J\) onto \(K^Z\)
by the Chinese remainder theorem. The cardinality assumption implies
equality of dimensions, so \(J=\bigcap_{z\in Z}\mathfrak m_z\).
The point ideals are pairwise comaximal, as are their powers. Thus
\(\bigcap_{z\in Z}\mathfrak m_z^\delta=J^\delta\). Expanding \(F\) in
the displayed basis now proves the claim: every nonzero basis term has
top weight \(\sum_i b_i/e_i+|a|\ge\delta\), and distinct top terms of
the maximal weight cannot cancel.

\subsection{The points and their equations}

For (H1), take \(K\) to be the rational function field in the shallow
variables over the residue field of \(O\), and write \(x_{b,s}\) for
those independent indeterminates. The \(Y_{b,s}\) are the deep variables.
Their \(e_{b,s}\) are \(P\) for \(b<2r\) and \(P^{2r}\) for \(b\ge2r\).
Counting the deep positions in the \(h=2\ell\) letters gives
\(\prod e_{b,s}=P^{hr+h r(k-2r)}=q^{r(k-2r+1)}\).
Put \(m_0=k-2r\). Introduce one residue parameter \(y_{al}\in\mathbb F_q\)
for each pair \(0\le a<r,\ 0\le l\le m_0\).
Place its lift in a matrix \(A\) in row \(a\) and column
\(d(a,0)=a+r\) or \(d(a,l)=k-l\) for \(l>0\). All other entries are
zero. Since \(A^2=0\), \(1+A\) and \(1-A\) are inverses.
Write \(y^{(s)}_{al}\) for the image of the lift under the residue
embedding in letter \(s\); the cycle gives
\((y^{(s)}_{al})^{P^j}=y^{(s+j)}_{al}\).

Define a point \(z(y)\) with one coordinate for each deep position by
\[
 z_{b,s}(y)=
 \begin{cases}
  \displaystyle\sum_{a,l:\ d(a,l)=b} x_{a,s}\,y^{(s)}_{al},
                 &s\ \text{even},\ b\ge r,\\[3pt]
  \displaystyle-\sum_l x_{d(b,l),s}\,y^{(s)}_{bl},
                 &s\ \text{odd},\ b<r .
 \end{cases}
\]
Indeed the action of \((1+A)^{-1}\) on a deep variable sends it to
\(Y_{b,s}-z_{b,s}(y)\) after reading the shallow variables in \(K\).
In an even letter the columns of \(1-A\) give the first formula, and
in an odd letter the columns of \((1+A)^t\) give the second.
Shallow variables are fixed by this substitution. Now \(F\in F_\delta\)
means in particular that its translate by \(1+A\) belongs to the
\(\delta\)-th power of the ideal generated by deep variables.
Translating back shows that \(F\), over \(K\), belongs to
\(\mathfrak m_{z(y)}^\delta\).

Fix a letter \(i\). We claim that every linear form in the
\(y^{(i)}_{al}\), with coefficients in \(K\), is a \(K\)-linear
combination of functions
\[
               z_{b,s}(y)^{P^j},\qquad
               s+j=i\bmod h,\qquad P^j<e_{b,s}.
               \tag{HA3}
\]
For \(l>0\), take \(b=k-l\) and all \(j\) in \(0,\ldots,2r-1\)
for which \(s=i-j\) is even. There are exactly \(r\) choices. These
\(r\) functions involve only the \(r\) parameters \(y^{(i)}_{al}\)
for the chosen \(l\); their coefficient matrix is
\((x_{a,i-j}^{P^j})_{j,a}\).
Its determinant is nonzero. In any term of the determinant there is
exactly one factor whose variable has first index \(a\); the exponent
\(P^j\) recovers the row \(j\). Distinct permutations give distinct
monomials, even if letters repeat. This recovers all the parameters with
\(l>0\). For \(l=0\) use \(j=0,\ s=i\). If \(i\) is even, \(z_{a+r,i}\)
equals \(x_{a,i}y^{(i)}_{a0}\). If \(i\) is odd, \(z_{a,i}\) contains
\(-x_{a+r,i}y^{(i)}_{a0}\) and terms with \(l>0\) already recovered.
The coefficients are nonzero in \(K\), proving (HA3).

There are two consequences. First, knowing \(z(y)\) gives all the
\(y^{(0)}_{al}\), hence \(y\), by injectivity of the residue embedding.
There are exactly \(q^{r(m_0+1)}=\prod e_{b,s}\) distinct points.
Second, writing \(e_{b,s}=P^d\), the value \(z_{b,s}(y)^{e_{b,s}}\)
is a linear form in the parameters with superscript \(s+d\). Apply
(HA3) at \(i=s+d\). Its expression uses only powers of point coordinates
whose weight \(P^j/e_{b',s'}\) is strictly less than one. This gives
the polynomial \(G_{b,s}\) required in the interpolation lemma. Reading
the shallow variables in the rational function field cannot send a
nonzero \(F\) to zero, by their algebraic independence. A surviving
deep monomial of weight at least \(\delta\) lifts to a monomial of \(F\)
of the same cost. This proves the interpolation bound in (H1).\formalnote{5.25}{25}

\subsection{Packing the bounded-defect space}

Let \(K\) now denote the coefficient residue field, \(\mathcal L\)
the two-block Levi over \(W\), and \(B\) the subgroup of matrices lower
triangular modulo the maximal ideal. Take a set \(\mathcal R\) such
that every Levi element can be written \(bu\), with
\(b\in B\cap\mathcal L,\ u\in\mathcal R\), uniquely. Its size is the
number of pairs of full flags in the two residue blocks. Let \(\Omega\)
be the cone of real arrays of nonpositive defect at every letter for
every adjacent root. Write \(\Delta_{i,b}\) for the defect at letter
\(i\) and root \((b,b+1)\). For real upper degrees \(D_s\), let
\(\mathcal D(D)\) be the nonnegative arrays with
\(D_s-1\le\sum_b\beta_{b,s}\le D_s\) for every \(s\).

First observe that a nonzero polynomial \(f\) has a translate \(uf\),
\(u\in\mathcal R\), with a nonzero monomial of nonpositive defect at
every root \emph{inside} either block, at all letters. Give variable
\((b,s)\) height \(b\) at even letters and \(-b\) at odd letters. Over
all Levi translates choose a nonzero term of least height. Only finitely
many exponents can occur since the total degree is bounded. The height
increases by one from \(X_s\) to \(Y_s\) at every simple root. The
residue-field root lemma proved through the degree comparison before
(H2) gives nonpositive defect for the chosen term at each within-block
root, since that rank-two group lies in the Levi. Write the translating
element \(bu\). On each variable \(b\) acts by a nonzero diagonal
multiple of the same variable plus variables of strictly greater height:
at odd letters the action is upper triangular and height has reversed
sign. Taking a coefficient of least height therefore just multiplies
the corresponding coefficient in \(uf\) by a nonzero scalar. The term
already occurs at the representative \(u\).

For an integer \(M=q^j\) assign to every residue array
\(u\in\{0,\ldots,M-1\}^{kh}\) a polynomial \(F_u\in S_0\). In letter
\(s\) set \(d_s=\sum_b u_{b,s}\) and add \(M(k-\lfloor d_s/M\rfloor)\)
at row \(a=r-1\), leaving the other entries of \(u\) alone. Call the
result \(\gamma(u)\). Its letter degrees lie in \([Mk,M(k+1)]\).
The polynomial
\[
                       f=\sum_u x^{\gamma(u)}F_u(x^M)
\]
encodes the full family, since its summands have disjoint exponent
residues modulo \(M\). In every letter the acting matrices have entries
fixed by the \(q\)-power map. Consequently substitution \(x\mapsto x^M\)
commutes with their action. A monomial in any Levi translate of \(f\)
therefore has exponent \(\alpha=M\beta+\epsilon\), with \(\beta\)
occurring in a translate of some \(F_u\), and with the degrees of
\(\epsilon\) in \([Mk,M(k+1)]\). For a nonnegative array of letter
degrees at most \(d\), the defect formula bounds its absolute defect
by \(d+d\sum_{v=0}^{h-1}P^v/(q-1)\le2d\).
Hence the boundary defect of \(\alpha/M\) is at most
\(3t+2(k+1)\le4t\); its letter degrees are in
\([n+k,n+k+1]\). At a representative and exponent of least height as
above all the other adjacent-root defects are already nonpositive.

To remove the boundary defect use the following explicit integral
translation. Put \(\mu_b=P^{2(r-1-b)}\) for \(b<r\) and zero otherwise.
Set \(\sigma_{b,s}=(P^2-1)\mu_b\) at even letters and zero at odd ones.
Its even letter degree is \(P^{2r}-1\). Its defects at \((b,b+1)\) are
\(P^2\mu_{b+1}-\mu_b\) at even letters and
\(P(\mu_{b+1}-\mu_b)\) at odd letters. This follows by substituting
in the defect formula and summing the geometric series over even
letters. These numbers are nonpositive, and at \(b=r-1\) are at most
\(-1\). Thus \(4t\sigma+\alpha/M\) belongs to
\[
 U=\Omega\cap\mathcal D(D),\qquad
 D_s=n+k+1+4t\sum_b\sigma_{b,s}\le n+4tP^{2r}.
\]
The lattice \(M^{-1}\mathbb Z^{kh}\) contains this translated point.
Record for each representative the coefficients of its translate of
\(f\) at all nonnegative exponents \(\alpha\) so obtained, allowing
every point of \(U\cap M^{-1}\mathbb Z^{kh}\). These coefficients give
an injective linear map: every nonzero family has a nonzero packed
polynomial, and the least-height argument detects it at one of those
entries. It follows that
\[
 M^{kh}\dim S_0\le|\mathcal R|\,\#(U\cap M^{-1}\mathbb Z^{kh}).
\]
As \(j\to\infty\), the normalized count on the right tends to
\(\operatorname{vol}(U)\). The bounded polytope has boundary of measure
zero, including the case in which it has lower dimension. This proves
\(\dim S_0\le|\mathcal R|\operatorname{vol}(U)\).\formalnote{5.26}{26}

\subsection{The exact volume of the Pieri cone}

Write \(\varepsilon_s=1\) in even and \(-1\) in odd letters. For a tuple
\(\pi=(\pi_s)\) of permutations of the \(k\) rows define the linear map
\[
 (T_\pi v)_{b,s}=\varepsilon_s(Pv_{b,s+1}-v_{\pi_s(b),s}).
\]
We use the identity tuple \(1\) first. Write \(\beta=T_1v\).
The sum over letters in \(\Delta_{i,b}(\beta)\) telescopes. Indeed if
\(u_s=v_{b,s}-v_{b+1,s}\), its summand before division by \(q-1\) is
\(P^{(s-i)\bmod h}(u_s-Pu_{s+1})\), giving \(-(q-1)u_i\) in all.
Consequently
\[
                \Delta_{i,b}(\beta)=\beta(Y_i)-
                                      (v_{b,i}-v_{b+1,i}).
\]
Away from equality hyperplanes, the conditions \(\beta>0,\ \beta\in\Omega\)
therefore interlace two descending lists at each letter. For an even
letter the lists are \(x_b=Pv_{b,s+1}\) and \(y_b=v_{b,s}\) and the
condition is \(x_0>y_0>x_1>y_1>\cdots>x_{k-1}>y_{k-1}\).
At an odd letter exchange the lists. These interlacings already force
each \(v_{\bullet,s}\) into descending order.

For descending lists \(x,y\), the determinant of
\(({\bf1}_{x_b>y_c})_{b,c}\) equals one if they interlace and zero
otherwise, up to equality hyperplanes. Indeed each row is a final
segment of ones, nested as one moves down the rows. A nonzero determinant
requires all \(k\) final segments to be distinct and nonempty. They then
start successively at \(0,\ldots,k-1\), giving an upper triangular matrix
with diagonal entries one.

Take the product of these determinants over the \(h\) letters, transposing
the matrix at the odd letters if necessary. Its expansion is
\(\sum_{\pi}\operatorname{sgn}(\prod_s\pi_s){\bf1}_{T_\pi v>0}\).
Sorting each \(v_{\bullet,s}\) preserves this product, since each
permutation of one letter contributes its sign twice. For almost every
\(v\) there is a unique ordering of each letter into descending order.
The product therefore also equals the sum over all within-letter
reorderings \(v'\) of the indicator that \(T_1v'>0\) and
\(T_1v'\in\Omega\). The letter-sum constraints in \(\mathcal D(D)\)
depend only on the sums of the \(v_{\bullet,s}\), so impose them in
every term of this identity and integrate over \(v\).

Put \(Q=q^{-1}\) and let \(D(\tau)=\det(1-Q\tau)>0\) on row space.
We have \(|\det T_\pi|=q^k D(\tau)\) with \(\tau\) the cyclic product
of the row permutations (with the ordering determined by the shift).
In fact one factors out the signed shift \(v_s\mapsto
\varepsilon_s P v_{s+1}\), with determinant of absolute value \(P^{hk}\).
Successive block elimination in the remaining matrix \(1-P^{-1}S\),
with \(S\) the permutation-weighted cyclic shift, gives identity diagonal
blocks except for the last, \(1-P^{-h}\tau\). A length-\(j\) cycle in
\(\tau\) contributes \(1-Q^j>0\), proving both positivity and
invertibility. The integration on the interlacing side has \((k!)^h\)
equal terms with Jacobian \((q^kD(1))^{-1}\). On the determinant side
there are \((k!)^{h-1}\) tuples for each product \(\tau\), with Jacobian
\((q^kD(\tau))^{-1}\). It follows that
\[
 \operatorname{vol}(U)=\operatorname{vol}(\mathcal D(D))\,
 \frac{(1-Q)^k}{k!}\sum_{\tau\in S_k}
                              \frac{\operatorname{sgn}\tau}{D(\tau)} .
\]
This identity also holds when the slab has zero volume.

For completeness the finite permutation sum can be evaluated directly.
The exponential generating function of these sums, using the cycle
decomposition and the sign \((-1)^{j-1}\) of a \(j\)-cycle, is
\[
 \exp\!\left(\sum_{j\ge1}
       \frac{(-1)^{j-1}z^j}{j(1-Q^j)}\right)
                      =\prod_{m\ge0}(1+zQ^m).
\]
The equality follows by expanding the logarithm for \(|z|<1\).
If \(a_k\) is the coefficient on the right, removing its first factor
gives \((1-Q^k)a_k=Q^{k-1}a_{k-1}\), \(a_0=1\). Multiplication by
\((1-Q)^k\) now gives the required volume fraction
\[
        (1-Q)^k\frac{Q^{k(k-1)/2}}{\prod_{j=1}^k(1-Q^j)}
                         =\prod_{j=1}^k\frac{q-1}{q^j-1}.
\]
The source evaluates the same signed sum by comparison with the
one-letter cone.\formalnote{5.27}{27}

Finally a full flag over \(\mathbb F_q\) can be chosen by successively
choosing a line in spaces of dimensions \(k,k-1,\ldots,1\). Its count
is the reciprocal of the displayed fraction. Reduction modulo the maximal
ideal identifies the parahoric cosets with the partial flags of its
block type, and the refinements of a partial flag with the two full
flags counted by \(|\mathcal R|\). Thus the two indices cancel the
volume fraction. For \(E=4tP^{2r}\), a letter slab of upper degree
\(D_s\le n+E\) has volume
\((D_s^k-(D_s-1)^k)/k!\le (n+E)^{k-1}/(k-1)!\), since \(D_s\ge1\).
Multiplying over letters costs at most \(\exp(h(k-1)E/n)\) times
\((n^{k-1}/(k-1)!)^h\). The latter is at most the number \(|I|\)
of exponent arrays. From \(2hkE\le\eta n,\ 0<\eta\le1\) we obtain
\(\exp(h(k-1)E/n)\le e^{\eta/2}\le1+\eta\). This proves (H2).\formalnote{5.28}{28}

\subsection{The rank-two auxiliary polynomial}

We finish the extreme-term argument invoked in the packing and adjoint
estimates. Fix a root, a letter \(i\), and weights \(c\) with
\(\delta_s=c(Y_s)-c(X_s)\in[1/2,1]\). Suppose a monomial with exponent
\(\alpha\) in \(f\) has coefficient nonzero modulo \(w^{l+1}\), and
every term in every translate by the embedded \(\mathrm{SL}_2(W)\) has
integral \(w\)-order plus \(c\)-weight at least
\(B=l+c\cdot\alpha\). The assertion is \(\Delta_i(\alpha)\le0\).
This formulation also applies over the coefficient residue field, with
\(w=0,\ l=0\).

Put \(P_s=P^{[s-i]}\), \(\rho=\delta_i/q\), and give variables two
integer weight vectors \(d,e\):
\[
 d(X_s)=qP_s,\qquad e(Y_i)=q,\qquad e(Y_s)=P_s\quad(s\ne i);
\]
all other entries vanish. The parity sign \(\xi_s\) is \(-1\) when
\(s\) and \(i\) have the same parity, and \(+1\) otherwise. With an
auxiliary variable \(T\) define a substitution \(S\) by
\begin{align*}
 S(X_i)&=T^q X_i, &S(Y_i)&=X_i,\\
 S(X_s)&=T^{qP_s}X_s-\xi_s T^{P_s}Y_s,
 &S(Y_s)&=Y_s-\xi_s X_s\qquad(s\ne i).
\end{align*}
Leave the other variables fixed. Denote by \(S_{\rm main}\) the first summand
in each formula, including the whole formula at \(i\). Let \(m(\alpha)\)
be the exponent after merging \(Y_i\) into \(X_i\). Over
\(O/(w^{l+1})\), take \(\Phi(T)\) to be the coefficient of
\(x^{m(\alpha)}\) in \(S(f)\).

At infinity give \(T\) weight \(-\rho\) and \(z\) weight
\(c(z)+d(z)\rho\). Every main term \(S_{\rm main}(z)\) has exactly weight
\(c(z)\), since the weight of \(X_i\) now equals \(c(Y_i)\).
Each cross term has weight strictly greater than \(c(z)+1/20\).
Indeed \(P_s\rho\le1/P\le1/5\), and \(qP_s\rho\ge P/2\ge5/2\) for
\(s\ne i\). For the cross term of \(X_s\) the gain is
\(\delta_s-P_s\rho\); for that of \(Y_s\) it is
\(qP_s\rho-\delta_s\). Thus \(S(f)\) has level at least \(B\), and
\(S(f)-S_{\rm main}(f)\) has level at least \(B+1/20\), by expansion of each
monomial. The coefficient at \(x^{m(\alpha)}T^j\) in the former must
vanish modulo \(w^{l+1}\) for \(j>d\cdot\alpha\). In the latter it
vanishes already at \(j=d\cdot\alpha\), since any strictly positive
gain beyond order \(l\) forces integer order at least \(l+1\).
In \(S_{\rm main}(f)\) that coefficient is exactly \([x^\alpha]f\): the merged
exponent fixes all variables except the split between \(X_i,Y_i\);
\(d\cdot\alpha\) fixes this split since their pole weights differ by
\(q\). We have proved
\(\deg\Phi=d\cdot\alpha\), including nonvanishing of its top coefficient.

At a finite point take \(v\in W,\ v_s=\iota_s(v)\). Give \(T\)
weight \(+\rho\) and \(z\) weight \(c(z)-e(z)\rho\). In the root
rows take the matrix \(g_v=\left(\begin{smallmatrix}0&1\\-1&x
\end{smallmatrix}\right)\), with \(x=-v\) for even \(i\) and \(x=v\)
for odd \(i\). Reading its columns and those of its inverse transpose
shows that, up to a sign on each input variable, it sends \(X_s\) to
\(Y_s\) and \(Y_s\) to \(X_s+\xi_s v_s Y_s\). Write \(\tau\) for the
swap of \(X_s\) and \(Y_s\) at all letters. First \(S(\tau z)\) has
level at least \(c(z)\). At \(s=i\) this follows from \(q\rho=\delta_i\).
At \(s\ne i\) use \(P_s\rho\le1/5\le\delta_s\) and
\(qP_s\rho\ge5/2\ge\delta_s\) in the two formulas for \(S\).

Apply \(g_v\), then \(S\), then replace \(T\) by \(T+v_i\). On \(X_s\)
this gives \(S(Y_s)\), up to its input sign. On \(Y_s\), direct
substitution gives \(S(X_s)\) up to the input sign and errors
\[
 \big((T+v_i)^{qP_s}-T^{qP_s}-v_s\big)X_s
   -\xi_s\big((T+v_i)^{P_s}-T^{P_s}-v_s\big)Y_s .
\]
At \(s=i\) the second error vanishes and the first is
\(((T+v_i)^q-T^q-v_i)X_i\). Both bracketed polynomials are multiples of
\(p=w^2\); in characteristic \(p\) the powers distribute over sums
and \(v_i^{P_s}\equiv v_s\equiv v_i^{qP_s}\bmod p\) by the cycle.
At the finite point \(T\) has positive weight, so the error levels
are at least \(2+c(X_s)\) and \(2+c(Y_s)-e(Y_s)\rho\).
Both exceed \(c(Y_s)\), since \(\delta_s\le1\) and \(e(Y_s)\rho\le1\).
The composite substitution therefore takes every \(z\) to level at
least \(c(z)\), including variables fixed by the root. Apply it to
\(g_v^{-1}f\), which has level at least \(B\) by hypothesis.

It follows that \(S(f)(T+v_i)\) has level at least \(B\). In the
finite-point weights, \(X_i,Y_i\) have the same weight. The contribution
of \(x^{m(\alpha)}T^j\) to that weight is
\(c\cdot\alpha+(j-e\cdot\alpha)\rho\). Hence its coefficient is zero
modulo \(w^{l+1}\) whenever \(j<e\cdot\alpha\). This says exactly
that \((T-v_i)^{e\cdot\alpha}\) divides \(\Phi(T)\).
Take lifts of all \(q\) residue classes of \(W\). Their differences
are units, whose images in the coefficient quotient are also units,
so the \(q\) powers of the monic linear factors are pairwise comaximal.
Their product divides the nonzero \(\Phi\). A monic divisor gives a
degree bound even over this coefficient ring with zero divisors. Thus
\[
 q\left(q\alpha(Y_i)+\sum_{s\ne i}P_s\alpha(Y_s)\right)
       =q\,e\cdot\alpha\le d\cdot\alpha
       =q\sum_s P_s\alpha(X_s).
\]
Division by \(q(q-1)\) and rearrangement give \(\Delta_i(\alpha)\le0\).
This completes the rank-two input in both (H2) and the construction
leading to (H3).\formalnote{5.29}{29}


\section{The marked Hecke operator in coordinates}\label{sec:marked-intertwining}

Section~\ref{sec:hecke-average} constructs the local estimate for a
Frobenius cycle. To apply it to the source forms we must specify where
each local lattice change goes in the global class set and how it acts
on every factor of the weight. Fix the marked place \(\lambda\mid p\).
Let \(k=\dim X_0\), \(r=\lceil\log_p m\rceil\),
\(W=\mathcal O_{E,\lambda}\subset E\), and
\(O=\{z\in\mathbb C:|\iota_p(z)|_p\le1\}\).
Write \(K\) for the source level, \(\Gamma=U(I_k)(F)\), and
\(\mathcal X=\Gamma\backslash U(I_k)(\mathbb A_{F,f})/K\).
Choose a numbering of the basis of \(X_0\). The prime \(p\) is
a uniformizer in \(W\), the pair \(\lambda,\bar\lambda\) is split
over \(F\), and \(K\) contains the full maximal compact supported at
this pair. Put \(d=\operatorname{diag}(pI_r,I_{k-r})\) and let
\(J\subset\mathrm{GL}_k(W)\) be the inverse image of the lower block
parabolic modulo \(p\).

\subsection{An edge and its reverse coordinates}
For a unimodular lattice at a split place of odd residue characteristic, the right cosets in the Hecke double coset of \(d=\operatorname{diag}(pI_r,I_{k-r})\) parametrize lattice neighbours of type \(r\) \cite[Thm.~5.11]{upper-greenberg-voight}. The coordinates below also record the map on the coefficient representation.

The stabilizer \(\Gamma\cap xKx^{-1}\) of each class representative
is trivial. Indeed a stabilizing rational \(\gamma\) is integral in
some frame at every finite place, so \(\operatorname{tr}\gamma-k\)
is an algebraic integer. At each place above the principal level prime \(p_0\) where the level condition requires \(\gamma\equiv 1\) in integral coordinates, \(\operatorname{tr}\gamma-k\) vanishes in the residue field. Every complex conjugate
of this difference has absolute value at most \(2k\), as every
archimedean component is unitary for a definite form. The product
of the residue norms at the chosen places is \(p_0^f>(2k)^{[E:\mathbb Q]}\).
Hence \(\operatorname{tr}\gamma-k=0\). A unitary matrix of trace \(k\) has all
eigenvalues one and is diagonalizable, proving \(\gamma=1\).
Choose a representative \(\rho_x\) of every class \(x\) integral at
all places over \(p\). This uses weak approximation of the rational
unitary group at this finite set of places, applied to the open
condition \(\gamma x\in K_p^{\max}\).\formalnote{5.30}{30}

Let \(b\in\mathrm{GL}_k(W)/J\) and choose a representative \(\beta_b\).
For the edge \(e=(x,b)\), call \(x\) its terminus. Change the
\(\lambda\)-component of \(\rho_x\) to \(\beta_b d\), changing the
conjugate component by the unitary condition. Equivalently, multiply
\(\rho_x\) on the right by the split adelic embedding of
\[
                  h_e=(\rho_x)_\lambda^{-1}\beta_b d .
\]
As \(b\) varies, these are representatives of all right integral
cosets in the double coset at \(\lambda\). Equality of two cosets
is equivalent to \(d^{-1}\beta_b^{-1}\beta_{b'}d\) being integral,
which is exactly \(\beta_b^{-1}\beta_{b'}\in J\). Surjectivity follows
because the left integral factor of a double-coset representative,
after absorbing \((\rho_x)_\lambda\), needs only to be matched modulo
the lower block parabolic. Matrices over \(W\) lift every invertible
residue matrix by lifting entries and using the unit determinant.

Write \(\widetilde h_e\) for the split adelic embedding of \(h_e\).
The modified point has a unique class \(o(e)\), the origin, and
choose factors satisfying
\[
             \rho_x\widetilde h_e=\gamma_e\,\rho_{o(e)}u_e ,
             \qquad \gamma_e\in\Gamma,\quad u_e\in K.           \tag{MI1}
\]
Evaluating at \(\lambda\) shows that \(g_e=\gamma_e^{-1}\beta_b d\)
belongs to \(\mathrm{GL}_k(W)\), since its completion is
\((\rho_{o(e)})_\lambda(u_e)_\lambda\), an integral invertible
matrix. Here all factors defining \(g_e\) before completion are
rational over \(E\), so both \(g_e\) and its inverse have entries
in the localization \(W\). In particular,
\[
                    \gamma_e^{-1}=g_e d^{-1}\beta_b^{-1}.       \tag{MI2}
\]
Away from \(\lambda,\bar\lambda\), equation (MI1) gives
\((\gamma_e)_v=(\rho_x)_v(u_e)_v^{-1}(\rho_{o(e)})_v^{-1}\);
it is integral and integrally invertible at all other \(v\mid p\).

The endpoint data needed in the elementary divisor estimate have
no multiplicities. Explicitly the map
\[
       (x,b)\longmapsto
       \left(o(e),[(g_e^t)^{-1}]\right)
       \quad\hbox{on }\mathcal X\times\mathrm{GL}_k(W)/J
                                                               \tag{MI3}
\]
is a permutation. To prove it, consider edges \(e,e'\) with equal
images, and let \(b,b'\) be their terminus cosets. The origin classes
are equal and \(g_e^t(g_{e'}^t)^{-1}\in J\). By the parahoric
criterion and inverse transpose, \(d g_e^{-1}g_{e'}d^{-1}\) is
integral at \(\lambda\) with integral inverse. Consequently
\[
      \eta=\gamma_e\gamma_{e'}^{-1}
             =\beta_b d g_e^{-1}g_{e'}d^{-1}\beta_{b'}^{-1}
\]
is integral there. In (MI1) cancel the common origin on comparing
the two edges. With \(u=u_{e'}^{-1}u_e\in K\), the point
\(\rho_x^{-1}\eta\rho_{x'}u\) is trivial off
\(\{\lambda,\bar\lambda\}\) and integral there. Maximality of \(K\)
on this pair puts \(\rho_x^{-1}\eta\rho_{x'}\) in \(K\). Hence
\(x=x'\), and the stabilizer calculation forces \(\eta=1\).
Equation (MI2) now makes \(d^{-1}\beta_b^{-1}\beta_{b'}d\) integral,
so \(b=b'\). Injectivity gives bijectivity on the finite set.

If \(R_{\rm wt}\) is the source weight representation, the unscaled
Hecke operator consequently satisfies
\[
  (T_{\lambda,r}\phi)(\rho_x)
    =\sum_{b\in\mathrm{GL}_k(W)/J}
          R_{\rm wt}(\gamma_{(x,b)})\,\phi(\rho_{o(x,b)}).       \tag{MI4}
\]
Indeed the indicator of the double coset acts by the sum over its
right integral cosets, each once. Right \(K\)-invariance eliminates
\(u_e\) from (MI1), and left rational equivariance contributes exactly
\(R_{\rm wt}(\gamma_e)\).\formalnote{5.31}{31}

\subsection{The transposed symmetric powers and the remaining weight}

Retain the \(h\) marked letters \(s\) and their embeddings
\(\iota_s:W\to O\) from Section~\ref{sec:hecke-average}. Write
\(\varphi_s\) for the corresponding complex CM embedding, so that
\(\iota_s\) on \(W\) equals \(\varphi_s\) on an even letter and
\(\varphi_s\circ\mathrm{conj}\) on an odd letter. For
\(u\in\mathrm{GL}_k(W)\), let \(L_s(u)=\iota_s(u)\) on even
letters and \(L_s(u)=\iota_s(u)^{-t}\) on odd letters. Define
\[
 A_s=\begin{cases}
          \operatorname{diag}(I_r,p I_{k-r}),&s\text{ even},\\
          \operatorname{diag}(p I_r,I_{k-r}),&s\text{ odd}.
       \end{cases}
\]
For each edge write \(g=g_e\), \(\beta=\beta_b\) and
\(\gamma=\gamma_e\). With \(c_s=p\) on even letters and \(c_s=1\)
on odd ones, the identity in degree one is
\[
        c_s\,\varphi_s(\bar\gamma)^t
                     = L_s(g) A_s L_s(\beta^{-1}).             \tag{MI5}
\]
On even letters \(\bar\gamma^t=\gamma^{-1}\) and \(p d^{-1}=A_s\),
so (MI2) proves the formula. On an odd letter one transposes
\(\gamma=\beta d g^{-1}\) and applies \(\iota_s\). This gives
\(\iota_s(g)^{-t}d\,\iota_s(\beta)^t\), as asserted.

The source determinant exponent is zero. Its diagram has one row of
length \(a\) at the \(a\)-places and an empty diagram elsewhere.
For a row of length \(a\), a semistandard filling is uniquely
specified by the counts of the \(k\) entries. Its polynomial is the
corresponding monomial, since every column minor has size one.
Consequently the Weyl basis index splits as
\[
  I=\mathcal B\times H,\qquad
  \mathcal B=\left\{ (\alpha_{is}):
       \alpha_{is}\in\mathbb N,\ \sum_i\alpha_{is}=a
                              \text{ for every }s\right\}.
\]
Here \(H\) indexes the Weyl factors at all places outside the marked
letters. Let \(m_e\) be the Schur matrix of \(\gamma_e\) on these
factors, using the embeddings \(\varphi_\tau\circ\mathrm{conj}\).
At any remaining \(a\)-place the induced place and its conjugate
are different from \(\lambda,\bar\lambda\). Thus (MI1) makes
\(\gamma_e\) integral with integral inverse there. The corresponding
Schur matrix and its inverse are integral because they come by
base change from the integral Weyl representation. At an empty
diagram the action is \(1\). Tensoring gives \(m_e\in\mathrm{GL}_H(O)\).

Put \(\rho(u)=\bigotimes_s\operatorname{Sym}^a L_s(u)\), and let
\(\Delta\) be the diagonal matrix on \(\mathcal B\) whose entry
is \(p^{\sum_{\rm deep}\alpha_{is}}\), with the deep indices those
where \(A_s\) has diagonal \(p\). It equals
\(\bigotimes_s\operatorname{Sym}^a A_s\). In the source Schur
matrix a row gives the coefficients of the transformed basis
polynomial. Hence for a one-row place its \(i,j\) entry, after
scaling by \(c_s^a\), is the \(j,i\) entry of the symmetric power
of \(c_s\varphi_s(\bar\gamma)^t\). Formula (MI5) and multiplicativity
now give the full matrix on \(I\):
\[
 p^{ah/2}R_{\rm wt}(\gamma_e)
   =\bigl(\rho(g_e)\Delta\rho(\beta_b^{-1})\bigr)^t\otimes m_e .
                                                               \tag{MI6}
\]
Exactly \(h/2\) even letters each contribute \(p^a\). On the
monomials, the transpose here corresponds to working on coefficients
of the scaled monomial basis \(\prod_s(a!/\prod_i\alpha_{is}!)z^\alpha\).
The one-row Schur formula verifies the identity of matrix entries
integrally, including when some multinomial coefficient is divisible
by \(p\).\formalnote{5.32}{32}

\subsection{Returning to the integral source}

Let \(f\) send a family \(F=(F_x)_{x\in\mathcal X}\), \(F_x\in O^I\),
to the form with values \(F_x\) on \(\rho_x\). Explicitly,
\(f(F)(\gamma\rho_x u)=R_{\rm wt}(\gamma)F_x\).
Freeness proved above makes this well-defined. At \(\rho_x\), the
Weyl gauge of \(F_x\) is at most one exactly when all its entries
are in \(O\), since \(\rho_x\) is integral with integral inverse
above \(p\). At \(\gamma\rho_x u\) the inverse frame in the gauge
cancels the left weight action, while \(u\) and \(u^{-1}\) preserve
the integral lattice. Thus \(f\) identifies these families with the
whole integral source \(\mathcal M\). Over \(\mathbb C\) the identical
class evaluation gives \(\dim\mathrm{SourceForm}=|\mathcal X|\,|I|\).

Let \(\mathcal E\) act on families by summing the matrices on the
right of (MI6) from \(o(e)\) to \(x\) over all \(e=(x,b)\).
Combining (MI4) and (MI6) gives
\[
                    p^{ah/2}T_{\lambda,r}\, f=f\,\mathcal E.
\]
The origin and terminus conditions for the finite edge estimate are
precisely (MI3) and the fact that \(\beta_b\) represents \(b\).
The other factor is the invertible matrix \(m_e\) at each edge.
It follows that \(Y=p^{ah/2}T_{\lambda,r}\) preserves \(\mathcal M\);
for every \(j_0\ge0\) the colength divisor of
\(Y\mathcal M+p^{j_0}\mathcal M\) in \(\mathcal M\) divides \(p^u\)
with
\[
             u\le(3c_{\rm div}+\eta j_0)\dim\mathrm{SourceForm}.
\]
This is the colength notion over the coefficient valuation ring
used in the marked estimate. The argument needs weak approximation
for unitary groups at the places over \(p\) and the local elementary
divisor estimate proved in the preceding marked-prime arguments.\formalnote{5.33}{33}


\section{The marked Hecke operator and its elementary divisors}\label{sec:marked-divisors}
\subsection{A single scaled Hecke operator}

Fix an exceptional inert prime \(p\) and a marked prime \(\lambda\) above it.
Let \(r=\lceil\log_p m\rceil\), let \(o\) be the order of \(-p\bmod q_0\),
and write \(P=p^o\). The \(a\)-places inducing \(\lambda\) or
\(\bar\lambda\) form one cycle under multiplication by \(-p^o\).
Write \(h\) for its length, to distinguish it from the level \(N\).
The letters \(s\in\mathbb Z/h\) alternate between the two primes, so
\(\lambda\) receives \(h/2\) letters. Their embeddings from the local
ring \(W\) above \(\lambda\) to the coefficient valuation ring \(O\) obey
\(\iota_{s+1}(z)\equiv\iota_s(z)^P\bmod p\). The residue field of \(W\)
has \(q=P^h\) elements. This is the arithmetic Frobenius-cycle input.
We work with pairs of neighbouring lattices, called edges, so that summing over neighbours represents the local Hecke action. The coordinate map of \S\ref{sec:marked-intertwining} intertwines the operator on edges with the operator on integral source forms. A bound on the elementary divisors upstairs will thus give a bound in this quotient.
\formalnote{5.34}{34}

At each letter take degree \(n=a\) in \(k\) variables. The set \(I\) of
exponent arrays has cardinality \(\binom{n+k-1}{k-1}^{h}\).
The action of \(\mathrm{GL}_k(W)\) uses the embedding in the even
letters and its inverse-transpose action in the odd letters. An edge
is a change of lattice of type \((r,k-r)\), indexed at each origin by
a coset of the parahoric subgroup \(J\). Set
\[
 c_{\rm div}=hn\left(P^{-2r}+\frac{r^2}{kP}\right),\qquad
 \eta=\frac1{m^2f}.
\]
The cycle satisfies \(2r\le k,\ 0<\eta\le1,\ n^4\le q\) and
\((hk\log q/\eta)^5P^{3r}\le n\). Their verification uses
\(13\le p<x_1,\ 0<o\le r,\ m\le oh\), and the admissible scale
inequalities. In particular all the quantitative hypotheses of the
edge theorem hold. Write \(D_*\) for the number of coordinates, including
the origins and the additional multiplicity indices. We derive a cost
of at most \((3c_{\rm div}+\eta j_0)D_*\) modulo \(p^{j_0}\).

\subsection{The two finite-field counts}

The first count bounds how deep the lattice of vectors depending only on
the origin can lie inside the lattice of all edge values. Over the residue
field of \(O\), let \(F_\delta\) consist of forms all of whose translates
have degree at least \(\delta\) in the deep variables. A variable indexed by
\((b,s)\), \(0\le b<k\), is deep if \(s\) is even and \(b\ge r\), or \(s\)
is odd and \(b<r\). Give a shallow variable cost zero, a deep variable
with \(b\ge2r\) cost \(P^{-2r}\), and every other deep variable cost \(P^{-1}\).

Each nonzero \(F\in F_\delta\) has a monomial of cost at least \(\delta\).
To see this, regard the shallow variables as coefficients in a rational
function field. Unipotent substitutions \(1+A\), with entries only in
rows \(<r\) and columns \(\ge r\), supply points of multiplicity at least
\(\delta\). The equations of these points express each deep variable raised
to the reciprocal of its cost in terms of monomials of cost \(<1\).
The injectivity and number of those points enter the polynomial
vanishing-ideal estimate. Section~\ref{sec:hecke-average} proves the interpolation statement by constructing the points, their equations, and the injection into the multiplicity quotient.

Projection to monomials of cost \(\ge\delta\) is therefore injective on
\(F_\delta\). Sum the dimension bound over \(\delta\ge1\):
each monomial \(\beta\) is counted at most \(\lfloor\mathrm{cost}(\beta)\rfloor\)
times. Symmetry gives the average exponent of each variable at a letter
as \(n/k\). For every row index exactly half the letters are deep. The
sum of the variable costs is consequently at most
\(h(kP^{-2r}+r^2/P)\). We obtain
\[
                  \sum_{\delta\ge1}\dim F_\delta
                                 \le |I|c_{\rm div}.               \tag{H1}
\]
The integral lattice comparison factors each edge restriction as \(Qd\),
where \(Q\) has a left inverse and spans the functions depending only
on the origin. Its proof bounds the exponent of \(\det d\) in units of
\(w=\sqrt p\) by the sum of the layer dimensions in (H1) times the
number of origins. The dual restriction factors as \(Q'd'\). Thus for
the two ends together \(v_p(\det d\det d')\le c_{\rm div}D_*\).\formalnote{5.35}{35}

The second count controls the number of edge directions lost between
the two ends. For the root between consecutive rows \(a,a+1\), call the
two variables \(X_s,Y_s\), with \(Y_s\) the second at even \(s\) and the
first at odd \(s\). For a letter \(i\) put \(P_s=P^{(s-i)\bmod h}\).
The defect of a monomial at that root and letter is
\[
 \beta(Y_i)+\frac1{q-1}\sum_sP_s\bigl(\beta(Y_s)-\beta(X_s)\bigr).
\]
The root lemma bounds this by zero for a term of least order plus weight
among all rank-two translates, when the weight rises from \(X_s\) to
\(Y_s\) by between \(1/2\) and \(1\). Here is the degree comparison
in the lemma. For the leading monomial the auxiliary polynomial over
\(O/(w^{l+1})\) has a nonzero coefficient in degree
\(q\sum_sP_s\beta(X_s)\), and no higher coefficients, where \(l\) is
the order of the leading coefficient in units of \(w\). At each lifted
residue class of \(W\) it vanishes to order at least
\[
                   v=q\beta(Y_i)+\sum_{s\ne i}P_s\beta(Y_s).
\]
For distinct residue classes the linear factors
have unit differences, so their \(v\)-th powers are pairwise comaximal.
Their product of degree \(qv\) divides the polynomial. Comparing degrees
gives \(v\le\sum_sP_s\beta(X_s)\). Since \(P_i=1\), rearranging gives
exactly the stated defect inequality.\formalnote{5.36}{36}

Let \(c_0=k^2oh,\ L=hc_0,\ c=5L,\ t=3kc\). At the boundary root
\((r-1,r)\), take the residual space \(S_0\) of forms all of whose
block-Levi translates have zero coefficient at defect \(>3t\) at every
letter. The dimension estimate needed is
\[
             [\mathrm{GL}_k(W):J]\dim S_0\le(1+\eta)|I|.           \tag{H2}
\]
Its main geometric input is the flag-and-volume estimate. Among all
Levi translates choose a monomial of least height. The root lemma forces
the Pieri inequalities at the other simple roots. Packing forms using
substitution by high powers of \(q\) and counting exponent arrays in
the translated cone bounds the dimension by the number of Levi flags
times the cone volume. The enlarged degree bounds are
\(n+E_0,\ E_0=4tP^{2r}\). The relative volume of the zero-defect cone in
a degree slab is
\(\prod_{j=1}^{k}(q-1)/(q^j-1)\), the reciprocal of the number of
complete flags. Thus the parahoric index times the Levi-flag factor
is absorbed by this relative volume. The enlargement costs at most
\(\exp(h(k-1)E_0/n)\le1+\eta\), since \(8hktP^{2r}\le\eta n\).
The slab volume at degree \(n\) is bounded by the monomial count. This
proves (H2) upon applying the packing estimate and the interlacing-volume
formula derived in \S\ref{sec:hecke-average}. The latter follows by integrating a determinant of interlacing indicators and summing over permutations.\formalnote{5.37}{37}

\subsection{The pairing, its approximate inverse, and the resulting index}

Pair the edge and dual-edge coordinates with diagonal monomial weights
\(b_\beta\), and take complementary weights \(b'_\beta\) with
\(b_\beta b'_\beta=p^L\). In each edge there is also an invertible
transition matrix on the multiplicity indices. We need the following
specialized averaging result. There is an integral matrix \(V\) whose
values satisfy the above defect conditions modulo \(w^c\). The lattice
\(C_0\) of forms satisfying these conditions contains the forms belonging
to the halfway lattices at all edges. For each such dual form \(h'\) and
each \(h\),
\[
       \langle Vh,h'\rangle_b\equiv
                          p^L\langle h,h'\rangle_b\bmod w^c.
\]
Moreover \(w^{c-1}h\in C_0\) implies that the reduction of \(h\) lies in
\(S_0\): cancel \(w^{c-1}\) in each forbidden coefficient, and then
reduce modulo the maximal ideal. The averaging theorem acts on one
generic-matrix variable at a time. In each letter its invariant linear
functional kills high row-weight monomials and agrees with \(p^{c_0}\)
times evaluation at the identity in pairings of the indicated precision.
The unchanged pairings telescope through the \(h\) letters. The result
uses the controlled average and the adjoint monomial pairing constructed in \S\ref{sec:hecke-average}.\formalnote{5.38}{38}

Let \(A\) send all edge values to the pairings with the columns of \(Q'\),
and let \(R'Q'=1\). Build \(B\) by applying \(R'{}^t\), inverse transition
on each edge, multiplication by \(b'\), and finally \(V\) in the polynomial
variables. Omitting \(V\) would give \(AB=p^LI\), since transitions cancel,
\(b_\beta b'_\beta=p^L\), and \((R'Q')^t=I\). Restoring \(V\) gives
\[
                   AB\equiv p^{2L}I\bmod w^c .                  \tag{H3}
\]
Both the columns of \(Q\) and those of \(B\) have their edge values in
\(C_0\). Put \(u=w^{c-1}\), \(a_0=p^{2L}\). There is room to lift (H3)
because \(a_0w\mid u\) by \(w^2=p,\ c=5L,\ L\ge2\).

For clarity, take Smith form of the columns of \(B\) in the complement
of the split summand \(\operatorname{im}Q\). This gives basis vectors
\(v_j\) and coefficients \(t_j\) such that \(\operatorname{im}Q+
\operatorname{im}B\) is spanned by \(\operatorname{im}Q\) and \(t_jv_j\).
For \(0\ne t_j\mid u\), the values of \(uv_j\) lie in \(C_0\), so these
\(v_j\) and the columns of \(Q\) reduce into \(S_0\) edge by edge.
They remain independent on reduction: \(Q\) is split and the \(v_j\)
are basis vectors of its complementary summand. Their number is at most
the number of edge slots times \(\dim S_0\), which is at most
\((1+\eta)D_*\) by (H2). Since \(Q\) already has \(D_*\) columns, there
are at most \(\eta D_*\) such additional directions.

Let \(G\) be the span of \(AQ\) and the \(Av_j\) for these additional
directions. At each remaining index \(u\mid t_j\), by total divisibility
in the valuation ring. Hence (H3) gives
\(a_0O^{D_*}\subset G+uO^{D_*}\). To remove the last term, pass to the
quotient by \(G\). The image of \(a_0O^{D_*}\) is finitely generated and
is contained in \(w\) times itself, since \(a_0w\mid u\). Nakayama's
lemma makes it zero. Thus \(a_0O^{D_*}\subset G\).

The actual edge operator in these coordinates has factoring matrices
\(d'{}^t(AQ)d\) on the paired input. Its image on the full input contains
the corresponding image. Inserting \(d,d'\) costs their determinants.
To remove the extra generators at precision \(p^{j_0}\), multiply their
coefficients by \(p^{j_0}\), costing at most one such factor per generator.
The determinant bound for the final image plus \(p^{j_0}O^{D_*}\)
therefore divides
\[
     \det(d')\det(d)\, (p^{2L})^{D_*}
                                  (p^{j_0})^{\#\{\text{extra directions}\}} .
\]
Each end of the edge restriction has already been counted in (H1).
The cycle size hypotheses give \(L\le c_{\rm div}\). Combining the
three costs gives \(v_p\le(3c_{\rm div}+\eta j_0)D_*\).
Mapping the associated integral combinations to the actual source
preserves this bound. At the free source level its coordinate count
is \(D_X\).

Finally \(P\ge p,\ p^r\ge m,\ r\ge1\), and the scale assumptions give
\(mr^2\le k\). Thus \(P^{2r}\ge mp\) and
\(r^2/(kP)\le1/(mp)\). With \(A_\lambda=a h/2\) we obtain
\[
           3c_{\rm div}\le\frac{6ha}{mp}
                         =\frac{12A_\lambda}{mp}.
\]
At the indicated precision the scaled operator at \(\lambda\) consequently
has colength exponent at most
\((12A_\lambda/(mp)+j_0/(m^2 f))D_X\). The two inputs carrying the most
specialized geometry in this derivation are the integral
Frobenius-compatible average and the leading-term interpolation and
flag-volume estimates behind (H1)--(H2). Their exact hypotheses control
the permissible degree and precision.\formalnote{5.39}{39}



\section{The period identity on the prescribed levels}
\label{sec:seesaw-levels}

The covolume argument takes place on the target chain \(Y_j\).
The integrality and divisibility estimates apply to theta kernels on
the definite chain \(X_j\), where the lifts are finite class sums.
Kudla's seesaw identity compares two pairings by interchanging the two integrations of a common theta kernel \cite[\S1, pp.~247--249]{upper-kudla}. The period identity needed here relates the two chains at their prescribed levels. It rewrites the
lift of a Poincar\'e image on the \(Y_j\) as a finite linear
combination of definite lifts, followed by restriction to \(Y_m\).
The coefficients in this combination determine whether the
finite-place estimates can be used for all twists.

Fix the admissible data, a coefficient functional \(t\) at the source, and
\(g\in U(Y_0)_\infty\). Write \(\mathcal P=\mathcal P_{g,t}\),
\(m=2\varrho+2\), \(n_j=\dim_E X_j\), \(G_j=U(Y_j)\), and
\(H_j=U(X_j)\). Let \(F\) be the totally real subfield of \(E\).
The assertion to be proved is that there exists \(c_{\rm ss}\in
\mathbb C^\times\), depending on the fixed data, \(g,t\), such that for
each permitted twist \(\vartheta\) there is a finite set of coefficients
\(\Psi_\vartheta\) with
\[
 (T_{m-1}\cdots T_0)\mathcal P\phi
    =c_{\rm ss}\operatorname{Res}_{Y_m}
       \Theta_{\rm def}\!\left(\sum_x\Psi_\vartheta(x)V_x\right)\phi
 \qquad(\phi\ \text{in the ramified source space}).
\]
Here \(T_j\) is the normalized lower theta lift, \(V_x\) is the tensor
of the translated integral vacua glued in adjacent pairs, and
\(\Theta_{\rm def}\) contracts this tensor with the upper lifts.
Both sides take values in the same definite algebraic forms at \(Y_m\).
The arithmetic assertions on \(\Psi_\vartheta\) follow below.

For each \(Y_j\), let \(K_j\) be the stabilizer of the standard finite
lattice, \(J_{M,j}\) the subgroup reducing to the identity at the chosen
places above \(p_0,p_1\), and \(J_j\subset J_{M,j}\) the subgroup also
reducing to the identity at every \(w\in Q\), the tame conductor of
\(\vartheta\). Put \(M=p_0p_1\), \(q=\prod_{w\in Q}\mathrm Nw\), and
\(I_j=[K_j:J_j]\). On \(X_j\) use the chosen level further restricted by
the condition that the last row is the last basis covector modulo
\(\pi_N\), where \(\pi_N\bar\pi_N=N\). In particular the conditions at
the marked prime \(N\) and the tame conductor have different origins.
For \(j<m\), the grading of \(Y_j\) has definite part \(X_j\) and a
string of length \(m-1-j\). Write \(R_j\) for its Bessel group and
\([R_j]=R_j(F)\backslash R_j(\mathbb A_F)\). At the top use the trivial
grading on the definite space \(Y_m\), of dimension \(n_m-1\).\formalnote{5.40}{40}

\subsection{Unfolding a single period and the initial index}

Write \(f_j\) for the Schwartz datum of \(T_j\) and
\(\theta_j(g',h)=\theta(\omega_j(g',h)f_j)\). Here \(\omega_j\) uses the
rank twists of the prescribed chain character at step \(j+1\). In the
bilinear pairing against a conjugate-type form \(z\), the transpose is
\[
 T_j^\vee z(g')=I_j\int_{[G_{j+1}]}\theta_j(g',h)z(h)\,dh .
\]
The lift in the other direction carries the same \(I_j\).
The interchange of the two kernel integrals gives
\(\int_{[G_{j+1}]}z\,T_jh=\int_{[G_j]}(T_j^\vee z)h\). The level of
\(T_j^\vee z\) contains \(J_j\) by invariance of \(f_j\). This pairing
puts all level factors on the same side before unfolding.

Let \(W_0\) be the function on \(R_0\) with the prescribed unipotent
character and with coefficient \(t(\phi)\) on the definite subgroup
(evaluated at the point of that subgroup). For outer variables
\(g_j=(\xi_j,y_j)\), put \(f_j^{g_j}=\omega_j(g_j,1)f_j\).
Denote by \(C_j(r,h;f_j^{g_j})\) its constant term at the map sending
one graded string to the next. Set
\[
 W_{j+1}(h)=\int_{[R_j]}W_{j,\mathrm{on}}(r)
                          C_j(r,h;f_j^{g_j})\,dr,\qquad
 \Xi_{j+1}(z)=\int_{R_{j+1}(F)\backslash G_{j+1}(\mathbb A_F)}
                                  W_{j+1}(h)z(h)\,dh .
\]
The notation \(W_{j,\mathrm{on}}\) means the preceding unfolded
function on the Bessel group, with the translations by the preceding
outer points as in the recursion. More explicitly, after \(W_j\) is
formed its input to the next step on \(R_j\) is \(r\mapsto W_j(rg_j)\)
when \(j>0\). The first step takes \(W_{0,\mathrm{on}}=W_0\).
Then
\[
 \int_{[R_j]}W_{j,\mathrm{on}}(r)(T_j^\vee z)(rg_j)\,dr
       = I_j\Xi_{j+1}(z).
\]
One can see the orbit calculation behind this assertion by expanding
the theta kernel in constant terms indexed by rational maps modulo
maps of nonnegative degree. Terms in classes incompatible with the
unipotent character integrate to zero. Only \(G_{j+1}(F)\)-translates
of the classes agreeing with the prescribed string map can contribute.
Their translates are disjoint modulo \(R_{j+1}(F)\).
Summing over these cosets unfolds the domain to
\(R_{j+1}(F)\backslash G_{j+1}(\mathbb A_F)\).
Section~\ref{sec:orbit-unfolding} proves the orbit identification by
two expansions, checks which unipotent characters vanish on the
contributing terms, and gives the absolute-convergence bound. The required
rapidity for conjugate test forms on this chain is established before
applying it.\formalnote{5.41}{41}

The bottom outer sum has a distinct normalization. The defining
Poincaré pairing is at level \(J_{M,0}\). Tracing the test form, which
has level \(J_0\), gives
\[
 [J_{M,0}:J_0]\int_{[G_1]}z\,T_0\mathcal P\phi
 =\sum_{\delta\in J_{M,0}/J_0}
  \int_{[R_0]}W_0(r)(T_0^\vee z)(r(g,\delta))\,dr .
\]
Substituting the one-step identity leaves
\([K_0:J_{M,0}]\sum_\delta\Xi_1(z;g,\delta)\). Thus
\(\mathcal C_0=J_{M,0}/J_0\), with multiplicity one. The constant
\([K_0:J_{M,0}]\) is independent of the twist. Treating the bottom
sum as an unnormalized trace would leave a spurious factor
\([J_{M,0}:J_0]\).

\subsection{The other outer variables and their finite measures}

Fix \(0<j<m\). Put \(H=G_j(\mathbb A_{F,f})\),
\(S=R_j(\mathbb A_{F,f})\), \(A=K_j\), \(J=J_j\), and \(S_A=S\cap A\).
The support calculation supplies a compact subset of \(H\) meeting
every relevant \(S\)-orbit modulo \(J\). Its construction uses a
frequency constraint: whenever a string constant term is nonzero, the
string map carries the scaled projected lattice of the preceding
grading inside the lower lattice of the new grading. This propagates
lower and upper lattice bounds. At even intermediate stages a lattice
split along the definite summand at a place dividing \(\pi_N\) forces
vanishing. A scalar fixing such a lattice acts through the level,
and ramification kills the resulting scalar-invariant period. The
lattice classification then supplies finitely many representatives
modulo \(K_j\), supported at \(Q\cup\bar Q\). Passing to the open
subgroup \(J\) preserves finiteness. These support and classification
statements are essential inputs to the argument that follows.\formalnote{5.42}{42}

Choose one point \(xJ\) in the compact set modulo \(J\) for each
represented \(S\)-orbit, and let \(\mathcal C_j\) be the union of their
\(S_A\)-orbits. For \(t\in H/J\) put
\[
 \mu_j(t)=\#\{sS_A:s^{-1}t\in\mathcal C_j\}.
\]
On the cut this number is finite and positive. At a chosen
representative \(t=xJ\) let \(Z_x=S\cap xJx^{-1}\). If
\(s^{-1}xJ\in\mathcal C_j\), it lies in the \(S_A\)-orbit of \(xJ\),
since exactly one representative of the entire \(S\)-orbit was
chosen. Hence the counted classes are the \(Z_x\)-orbit of \(S_A\)
in \(S/S_A\), and
\[
 \mu_j(xJ)=[Z_x:Z_x\cap A].
\]
Left translation shows that \(\mu_j\) is constant on the entire
\(S\)-orbit. Away from \(Q\cup\bar Q\) the chosen \(x\) is integral.
At \(w\in Q\), the elements of \(J\) reduce to the identity. A
sufficiently high congruence subgroup remains integral after
conjugation by \(x_w\), and taking \(q^\nu\)-th powers puts \(J_w\)
inside it for some common \(\nu\). Unitarity gives the analogous
assertion at \(\bar w\). Consequently \(z^{q^\nu}\in Z_x\cap A\)
for every \(z\in Z_x\). The same powers lie in the normal core since
the assertion holds for all conjugates of \(z\). In the resulting finite
quotient every element order divides \(q^\nu\), so every prime dividing
its order divides \(q\). Its order is a multiple of \(\mu_j\).
In particular \(\mu_j\mid q^e\) for some \(e\). This proof allows
\(q\) to have several prime factors.

Normalize finite Haar measure on \(H\) by \(\operatorname{vol}(A)=1\)
and on \(S\) by \(\eta(S_A)=1\). The function
\(\psi(t)=\mathbf 1_{\mathcal C_j}(t)/\mu_j(t)\), set to zero off
\(\mathcal C_j\), satisfies
\[
 \int_S \psi(s^{-1}t)\,d\eta(s)=1
\]
on each orbit meeting the cut. Indeed the integrand is
\(1/\mu_j(t)\) on exactly \(\mu_j(t)\) cosets of mass one. To see the
archimedean normalization, choose a nonnegative measurable cutoff
\(\beta\) on \(G_j(\mathbb A_F)\) with integral one on each orbit
under the adelic \(R_j\), and set
\(\alpha(\xi)=\int_S\beta(\xi,u)\,d\eta(u)\). Fubini and translation
invariance show that \(\alpha(\xi)\psi(y)\) has orbit integral equal
to \(\int_S\psi(s^{-1}y)\,d\eta(s)\). Inserting this density in
quotient integration recovers the original integral of any
integrable, \(R_j(F)\)-invariant function supported on these orbits.
In particular, if on the cut its \([R_j]\)-period equals
\(I_j a(\xi)b(yJ)\), its integral becomes
\[
 I_j\left(\int a(\xi)\alpha(\xi)\,d\xi\right)
      \sum_{t\in\mathcal C_j}\frac{b(t)}{\mu_j(t)}
                               \operatorname{vol}(J)
 =\left(\int a(\xi)\alpha(\xi)\,d\xi\right)
      \sum_{t\in\mathcal C_j}\frac{b(t)}{\mu_j(t)} .
\]
Indeed \(\operatorname{vol}(J)=I_j^{-1}\). The archimedean operation
depends on \(a\), the grading, and \(A\), hence a twist-independent
archimedean factor remains twist-independent at the preceding step.
Iterating this calculation from a factorized terminal period and
using the bottom formula gives the finite functional
\[
 \mathcal C(F)=
 \sum_{c_0\in\mathcal C_0,\ldots,c_{m-1}\in\mathcal C_{m-1}}
   \frac{F(c_0,\ldots,c_{m-1},1)}
        {\prod_{j=0}^{m-1}\mu_j(c_j)}.
\]
All functions of \(c_j\) here descend to the quotient by \(J_j\).\formalnote{5.43}{43}

\subsection{A pair of constant terms and the top branching scalar}

It remains to establish the factorization used under the outer
integrals. Consider a pair \(V\to V'\to V''\) starting at \(j=2i\).
Use the adjoint for the Hermitian forms to identify the joined arrow
space with \(\operatorname{Hom}(V',V^-\oplus V'')\). Write \(u\) and
\(u'\) for the two string maps and \(U=(u^*,u')\) for the matrix formed
by stacking them. The product \(f\boxtimes f'\) here means the
Schwartz arrow product in these coordinates. For suitably matched
theta normalizations there is a fixed \(z_\theta\ne0\) such that
\[
 \theta(f\boxtimes f')=z_\theta\theta(f)\theta(f').
\]
For translations by \(w,w'\), the product of the two constant terms
is \(z_\theta^{-1}\) times the value at zero of the grade-zero
projection of \(\rho(w)f\boxtimes\rho(w')f'\), where \(\rho\) is
Heisenberg translation. To verify this, expand both constant terms as
probability means on additive adelic quotients. The map
\((a,a')\mapsto(a^*,a')_{>0}\) onto the positive-degree entries of
the joined space is additive and sends rational entries to rational
entries. It is surjective, since arbitrary specified entries can be
extended to a matrix, unstacked, and the first block adjointed back.
The pushforward of normalized Haar on the product of the two compact
quotients is therefore normalized Haar on the target. Apply the theta
product formula inside this mean. This proves the asserted constant
term identity, including the inverse power of \(z_\theta\).

There are two unipotent integrations in the pair. On the odd starting
string the character of \(W_{2i,\mathrm{on}}\) cancels the character of
the first kernel. Its unipotent mean leaves the integral over
\([H_{2i}]\) times a measure scalar \(\kappa\). In the middle group,
translation by the unipotent changes the frequency of both kernels.
After stacking, these changes are the coordinates of the isotropic
part of the remaining planes. Averaging over them gives the zero term
for the full splitting times a second measure scalar \(\kappa'\).
In the source proof these steps use the mean formula for each Bessel
group and the explicit coordinate identity for the stacked shifts.
The surviving definite theta kernel on
\(\operatorname{Hom}(X_{2i+1},X_{2i}^-\oplus X_{2i+2})\) has datum
\[
 \operatorname{ZT}\!\left(\rho(U)(f\boxtimes f')\right)
\]
and scalar \(\kappa\kappa'/z_\theta\). Its two remaining group
integrals are over \([H_{2i}]\) and \([H_{2i+1}]\), with the triple
Weil action of the endpoint and middle groups. Linearity and the
theta composition formula identify them with the matrix coefficient
of the two consecutive definite lifts.\formalnote{5.44}{44}

For \(i<\varrho\) the two outside strings have middle vectors \(v,v'\)
of norms \(N,-N\) in \(V^-\oplus V''\). Their sum and
\((v-v')/(2N)\) form a hyperbolic pair: each is isotropic and their
mutual pairing is one. Removing this plane gives the stated definite
space. For \(i=\varrho\) the string lengths are \(1,0,0\).
The sole string vector \(v\) has norm \(N\) in \(V^-\), so \(v/\pi_N\)
has norm one. It becomes the last basis vector of \(X_m\). Thus the
terminal \(H\)-argument for \(k\in G_m\) is \(\operatorname{diag}(k,1)\).
The real and finite parts of the translated zero term separate.
The characters entering the real Weil action have weights independent
of \(\vartheta\). For each tuple \(\xi\) of real outer points the
coefficient functional at the top is consequently a functional
\(\ell_\xi\) independent of \(\vartheta,y,\phi\). Concretely, if
\(e_i(y)\) is the finite zero-term evaluation of the joined translates
on pair \(i\), then at \(y_m=1\) the top stage is the coefficient at
\(\ell_\xi\) of the definite lift with tensor
\(\bigotimes_{i=0}^{\varrho} e_i(y)\), evaluated at
\(\operatorname{diag}(k,1)\). This follows by induction because an
equivariant function on the Bessel group of an odd string is determined
by its restriction to the definite subgroup. The initial functional
is \(t\).\formalnote{5.45}{45}

To normalize the top pairing let \(x\) be a test form on \(Y_m\) and
\(h\) the definite lift just obtained. Let \(B\) denote the chosen
integral branching matrix, and set
\[
 d_Y=\dim W_{Y_m},\qquad
 b_\pi=\prod_{\tau}\Phi_\tau(\bar\pi_N)^{e_{\rm br}(\tau)} .
\]
Here \(W_{Y_m}\) is the top coefficient representation and
\(e_{\rm br}(\tau)\) is the branching determinant exponent. Average
the pointwise pairing of the algebraic form \(x\) with
\(Bh(\operatorname{diag}(k,1))\) over the finite class integral and
call it \(A(x,h)\). The archimedean average of the two matrix
coefficients on the compact top group is
\(C(\ell_\xi)A(x,h)\), by the multiplicity-one property of the
chosen branching map. With the algebraic inner product and normalized
restriction used here one has
\[
 \langle x,\operatorname{Res}_{Y_m}h\rangle
             =d_Y^{-1}b_\pi^{-1}A(x,h).
\]
The dimension is Schur's normalization, and the restriction uses
\(B/b_\pi\). Thus the archimedean coefficient for the last period is
\(a(\xi)=d_Y b_\pi C(\ell_\xi)\). With this \(a\), the preceding
cut argument gives, for all \(x\), the same scalar \(c_0\) times the
pairing against the restricted chain with finite tensor
\[
 D_\vartheta=\mathcal C\!\left(c\longmapsto
                       \bigotimes_{i=0}^{\varrho} e_i(c)\right).
\]
Nondegeneracy of the top pairing gives the equality of forms.
Moreover, a tuple \(b_j\) in the upper levels with \(b_m=1\)
permutes each cut via the inclusion of the definite subgroup into
the Bessel group. It preserves multiplicity. Equivariance of finite
evaluation in the three groups of each pair turns this change of cut
variables into the joint Weil action of \(b\), also for the last
embedding with \(b_m=1\). Hence \(D_\vartheta\) is jointly invariant.\formalnote{5.46}{46}

\subsection{Clearing the denominator of the glued tensor}

Set \(B_0=M^{f n_m^2}\), where \(M=p_0p_1\) is the product of the two principal level primes, \(f=[F:\mathbb Q]\) and \(n_m=\dim_E X_m\) is the rank of the last definite space.
At a cut point all \(y_j\) are integral away from \(q\).
The finite datum on one lower arrow is the indicator of the inclusion
matrix modulo the chosen and tame places, applied to the vacuum.
Expanding a congruence indicator as a finite additive Fourier average
produces roots of unity divided by the order of its matrix residue
group. After ignoring denominators over \(q\), the denominator on
edge \(j\) divides \(\mathrm N(\pi_M)^{\dim Y_j\,\dim Y_{j+1}}\).
Moving by the finite outer variables preserves integrality away from
\(q\). When \(Q\) is empty the sum of absolute values of the expansion
coefficients is at most one after every embedding. The joined
translation has allowed denominators \(\pi_M\) on its first block and
\(\bar\pi_M\) on its second, the conjugation coming from adjunction.

The finite zero term requires an additional lattice calculation. If
\(I\) is the isotropic space of the splitting and \(\Lambda\) the
vacuum lattice, the image of \(\Lambda\cap I^\perp\) in \(I^\perp/I\)
is the neighbour lattice for the glued pair. The calculation constructs
an integral lift of its vectors whose adjoint kills the shift and
keeps the allowed denominators. A translated vacuum either has zero
zero-term or gives, up to a root of unity, a translate of the zero
term of the vacuum. The latter is a fixed nonzero scalar times a sum
over neighbour-lattice representatives. Section~\ref{sec:neighbour-lattices} constructs these representatives and lifts on the graded lattices. In the resulting tensors the two
upper translations become integral away from \(q\) after multiplication
by \(\bar\pi_M\). A uniform bound for the averaging denominator in this
evaluation is given by the number of matrices modulo \(\pi_M\) of shape
\((X_j\oplus X_{j+2})\times X_{j+1}\). The needed middle and last
neighbour-lattice identities are proved with the respective splittings
described above.\formalnote{5.47}{47}

For \(j=2i\), the three matrix counts together give a denominator
dividing \(\mathrm N(\pi_M)^{a_i}\), where
\[
 a_i=\dim Y_j\dim Y_{j+1}+\dim Y_{j+1}\dim Y_{j+2}
                  +(n_j+n_{j+2})n_{j+1}\le 4n_m^2 .
\]
Since \(\mathrm N(\pi_M)=M^f\), this divides \(B_0^4\).
One fixed constant \(C_i\) suffices for each pair for all twists and
cut points. To express the result with the integral vacuum used in
\(V_x\), divide \(C_i\) by \(\kappa_{2i}\kappa_{2i+1}\), where
\(\kappa_j\) is the integral vacuum's fixed scale on edge \(j\).
Tensoring multiplies the new constants into \(C\) and the
denominators into
\[
 d_0=(B_0^4)^{\varrho+1}=B_0^{2m}.
\]
For \(Q=\varnothing\) the absolute coefficient sum of each pure
tensor is at most one. In that case \(q=1\) forces \(\mu_j=1\) on
the cut. Also \(\#\mathcal C_j\le B_0\): representatives are integral
everywhere, and the principal level modulo \(\pi_M\) is contained
in \(J_j\), so the cut injects into the cosets of that principal
level inside \(K_j\). Their number is bounded by the number of all
residue matrices. For general \(q\), the factors \(1/\mu_j\) have
denominators supported at \(q\). Thus we can write
\[
 D_\vartheta=C\sum_x r_xV_x,\qquad
 d_0r_x\text{ integral away from }q,\qquad
 \sum_x|\sigma(r_x)|\le B_0^m\quad(Q=\varnothing)
\]
for every automorphism \(\sigma\) of \(\mathbb C\).

If \(c_0C\ne0\), set \(c_{\rm ss}=c_0 C/d_0\) and take the new
coefficients to be \(d_0r_x\). If \(c_0C=0\), take \(c_{\rm ss}=1\)
and the zero tensor. Equal translated vacuum data may be represented
by translations vanishing at every place prime to \(Mq\). Replacing
the representatives and collecting coefficients preserves the bounds
and joint invariance of the tensor. This gives the prescribed
\(\Psi_\vartheta\), whose values are integral away from \(q\), with
\[
 \sum_x|\sigma(\Psi_\vartheta(x))|
       \le B_0^{3m}=M^{3fm n_m^2}\qquad(Q=\varnothing).
\]
The same nonzero scalar \(c_{\rm ss}\) works for every permitted twist \(\vartheta\).
The proof permits a zero map for particular choices of the original
coefficient functional \(t\).\formalnote{5.48}{48}

\subsection{The ramified calculation at the marked split prime}

At the split level prime \(N\), changing the last row and column coordinates in two consecutive lifts contributes a scalar \(G_i\) for the pair beginning at \(j=2i\). We determine \(G_i\) before taking a norm. Choose the set
\(\Sigma_N\) of non-induced places, one in each conjugate pair above
\(N\), and put \(A_N=\prod_{w\in\Sigma_N}\mathbb F_N\). The original
gluing uses the places \(P_\pi\) dividing \(\pi_N\). Passing to
coordinates on \(\Sigma_N\) uses the marked lattice on \(X_j\), whose
last coordinate is divisible by \(\pi_N\), and changes the last
coordinate on the dual lattice accordingly. Write \(D_j(S)\) for
the diagonal matrix on \(X_j\) with last entry \(N\) at the places
of \(S\cap\Sigma_N\) and unit entries elsewhere. Put
\(\delta_j=D_j(\Sigma_N)^{-1}\).

Let \(e\) denote the last basis vector of each relevant residue module
and \(\psi_N\) the fixed additive character of \(A_N\), given by the
self-dual character at the negative simple-pole lift of its argument.
For integral arrow matrices \(T:X_j\to X_{j+1}\) and
\(R:X_{j+1}\to X_{j+2}\), define
\[
 P_y(T)=\mathbf1_{\bar T e=y},\qquad
 Q_y(R)=\psi_N(e^{\mathsf t}\bar R y).
\]
Extend both functions by zero off the integral matrices. The gluing
kernel in the marked coordinates is then
\[
 \mathbf1_{\rm int}(T)\mathbf1_{\rm int}(R)
       \psi_N(e^{\mathsf t}\overline{RT}e)
       =\sum_{y\in A_N^{n_{j+1}}}P_y(T)Q_y(R)
\]
on integral inputs, with both sides zero otherwise. This identity
just sorts the product by the last column of \(T\). The direct datum
is \(P_e\), and \(Q_e\) is the mirror datum on the next arrow.
The local Schrödinger calculation identifies the original sum over
simple poles with this kernel after the marked changes of basis and
vacuum projection at \(N\). Write \(L_j(P)\) for the one-step lift
with the indicated matrix datum there, holding its remaining finite
data fixed. The changes of basis cancel at the middle group and give
the operator for a glued pair with \(j=2i\) in the form
\[
  G_i\,\rho_{j+2}(\delta_{j+2})
       \sum_y L_{j+1}(Q_y)L_j(P_y).
\]
To compute \(G_i\), let \(\chi_{s,r}\) denote the determinant
character called \(\mathrm{liftCharAtN}(\vartheta,s,r)\) in the
construction. It is evaluated on the indicated group, with \(r\)
the rank parameter of the opposite space. Then
\[
 G_i=\chi_{j+2,n_{j+1}}\!\left(D_{j+2}(\Sigma_N\setminus P_\pi)\right)
       \chi_{j+1,n_{j+1}}\!\left(D_j(P_\pi)\right).
\]
The factor involving \(D_{j+2}\) comes from scaling the target rows of \(R:X_{j+1}\to X_{j+2}\). The factor involving \(D_j\) comes from scaling the source column of \(T:X_j\to X_{j+1}\). This
uses the actual Schrödinger change-of-variables formula for these
characters.\formalnote{5.49}{49}

For a ramified input, at each \(w\in\Sigma_N\) write
\(\phi=(1-\rho(s_w))v_w\), where \(s_w\) is an integral scalar unit
reducing to one at the other marked places. If \(y_w=0\), the
indicator \(P_y\) is invariant under that scalar since \(s_w y=y\).
Moving \(s_w\) into the datum by integral equivariance therefore
shows that \(L_j(P_y)\phi=0\). If \(y_w\ne0\) at every \(w\), choose
an invertible residue matrix carrying \(e\) to \(y\), by completing
the column separately at each place, and lift it to an integral
invertible matrix. In the two lifts its actions on \(P_e\) and
\(Q_e\) cancel at the middle group. Every surviving summand equals
\(L_{j+1}(Q_e)L_j(P_e)\phi\). The number of summands is
\[
     h_i=\prod_{w\in\Sigma_N}(N^{n_{j+1}}-1).
\]

For a single lift \(L_j\) put
\[
 \gamma_j=
 \frac{\chi_{j+1,n_{j+1}}(D_j(\Sigma_N))}
      {\chi_{j+1,n_j}(D_{j+1}(\Sigma_N))}.
\]
Its second local equivariance law moves \(\delta_{j+1}\) past the
last-row indicator, turning it into the direct last-column indicator
and introducing \(\gamma_j\). The associated marked operator on a
representation \(\rho\) is
\[
 \operatorname{Up}_\rho
    =\rho(\delta)\!\!\sum_{\substack{y\\y_{\rm last}\in A_N^\times}}
        \psi_N(y_{\rm last})\,\rho(k_y).
\]
Here \(k_y\) is integral and its last column and last row reduce to
\(y\) and \(y^{\mathsf t}\). For a unit corner \(a=y_{\rm last}\)
and other coordinates \(v\), a possible residue matrix has upper
left block \(I+vv^{\mathsf t}/a\), last column \((v,a)\) and last
row \((v^{\mathsf t},a)\). Its Schur complement has determinant one
after taking out the corner \(a\), so it is invertible. The same
scalar argument as before kills any prescribed corner value vanishing
at one place. Sorting the mirror datum by its last column gives
\(\rho_{j+1}(\delta_{j+1})L_j(Q_e)=\operatorname{Up}_{j+1}L_j(P_e)\)
on ramified vectors. Sorting by the last row and moving \(\delta\)
through the last-row indicator gives instead
\(\gamma_j L_j(P_e)\operatorname{Up}_j\). In particular a direct
lift commutes with \(\operatorname{Up}\) at the price of \(\gamma_j\),
and preserves ramification.\formalnote{5.50}{50}

Apply the mirror identity to step \(j+1\), then commute its resulting
\(\operatorname{Up}_{j+1}\) left through \(L_j(P_e)\). The operator
for pair \(i\) on a ramified vector is
\[
  G_i h_i\,\gamma_{j+1}\gamma_j\,
             L_{j+1}(P_e)L_j(P_e)\operatorname{Up}_j .
\]
Each pair creates one marked operator. Commuting the one created at
pair \(i\) through the earlier \(2i\) direct lifts introduces
\(\gamma_0\cdots\gamma_{2i-1}\). Thus the full glued lift equals
the \(m\) direct lifts on \(\operatorname{Up}_0^{m/2}\phi\) times
\[
       c_N=\prod_{i=0}^{\varrho}
               \left(G_i h_i\prod_{l=0}^{2i+1}\gamma_l\right).
\]
This is an identity on the vacuum projections of the finite inputs
above \(N\). At the \(N\)-adic embedding used for the upper estimate,
the rank monotonicity of the determinant-character norm gives
\(\lvert\gamma_j\rvert_N\le1\). The integer \(h_i\) also has norm
at most one. Each place contributes to one of the two factors in
\(G_i\), and the local weight bound gives
\(\lvert G_i\rvert_N\le N^{2fb}\). Consequently
\[
     \lvert c_N\rvert_N\le N^{(2fb)(\varrho+1)}=N^{mfb}.
\]
This is the scalar loss before estimating the direct lifts
themselves. The scalar \(c_N\) uses the lift characters at the chosen
twist and is controlled at the marked prime.\formalnote{5.51}{51}


\section{The splitting character and the integral gauge in a pair of lifts}
\label{sec:weil-stages}

In \S\ref{sec:seesaw-levels} the product of two consecutive kernels becomes one kernel on a smaller space. There are three points at which a scalar can enter this calculation. Joining the two theta functionals gives a fixed complex scalar, restricting the Weil operator to a smaller space gives a determinant character, and translating the finite vacuum gives coefficients which must be integral away from the conductor. We calculate each at the place where it occurs.

\subsection{The reproducing integral in the adelic model}
Write \(W\) for the additive adelic phase space, \(\beta\) for its unit-modulus Heisenberg cocycle, and \(e\) for the vacuum coefficient. In the chosen coordinates it is a Gaussian at infinity and a semicharacter times the indicator of the self-dual lattice at the finite places. Haar measure is the product of the vacuum measures, with that finite lattice of mass one. The normalized vacuum has \(|e(x)|\le1\), \(e(0)=1\) and \(\int |e|^2=1\). In the coherent-state realization the vectors are the \(L^1\) functions reproduced by the kernel
\[
       K(y,x)=\beta(y,x)e(x-y),\qquad
       f(x)=\int_W f(y)K(y,x)\,dy .
\]
The equality picks a continuous representative and gives \(\|f\|_\infty\le\|f\|_1\). The Gaussian and the finite indicator produce the vacuum vector \(u=e\) in this model. Equivalently, in a Schrödinger polarization the orthogonality calculation is Plancherel: for a unit vacuum \(a\), Fourier transform in the modulation variable gives
\[
 \int |\langle M_\eta T_q a,b\rangle|^2\,d\eta
       =\int |a(x-q)|^2|b(x)|^2\,dx .
\]
Integration in \(q\) gives \(\|b\|^2\). Polarization gives the reproducing kernel on the matrix coefficients. Unit-modulus phase choices multiply the kernel by the same Heisenberg factors as the translated vacuum. This calculation works for the real and finite variables together, using additive dual Haar measures.\formalnote{5.52}{52}

In the coherent-state model the translate of the vacuum \(u=e\) is the function \(x\mapsto(\varrho(y)u)(x)=K(y,x)\). The scalar product is \(\int\bar f_1 f_2\), conjugate-linear in its first variable, and the reproduction identity and Hermitian symmetry of \(K\) give \(\langle\varrho(y)u,f\rangle=f(y)\). The model has its \(L^1\)-norm. Translations act isometrically and continuously in this norm, so \(\int_W f(y)\varrho(y)u\,dy\) exists as a Banach-space integral: its norm integrand is \(|f(y)|\|u\|_1\). Evaluation at each \(x\) is bounded, and evaluating the integral gives \(f(x)\) by the displayed identity. Applying any bounded linear map \(T\) to this equality gives
\[
                Tf=\int_W\langle\varrho(y)u,f\rangle\,
                                 T(\varrho(y)u)\,dy .
\]
This proves both the superposition used in \S\ref{sec:vacuum-projection} and the determination of a bounded operator by its values on all vacuum translates. The separate rational invariance used to join two theta functionals enters below.\formalnote{5.53}{53}

\subsection{Two determinant calculations}
Let \(V,V',V''\) be three consecutive graded spaces, of ranks \(d,e,f\). Their string lengths are \(s,s'=s-1,s''=\max(s'-1,0)\), with \(s\) odd. The outer ranks have even sum. The part of weight zero in \(V^-\oplus V''\) is \(X^-\oplus X''\) together with \(h\) hyperbolic planes. Write \(D,E_0,F_0\) for the ranks of \(X,X',X''\), and \(P\) for the number of positive-weight basis vectors in \(V^-\oplus V''\). Every positive weight pairs with its negative, and the middle string has \(s'/2\) negative weights. Hence
\[
      d+f=D+F_0+2h+2P,\qquad e=E_0+s'.
\]
The three definite groups act through the grade-zero inclusions. At the last pair the third group acts on \(X''\) by \(k\mapsto\operatorname{diag}(k,1)\), so its determinant still agrees with the one on the large space.\formalnote{5.54}{54}

Write \(x,y,z\) for the three determinants in \(U(1)(\mathbb A_F)\), and \(u\) for the restriction of the fixed splitting character \(\upsilon\) to this group. On a space of rank \(j\) the character for one arrow is \(a_j=a\,u^{\lfloor j/2\rfloor}\). The extra scalar on an arrow of ranks \(d,e\) is \(a_d(y)/a_e(x)\). Join the first arrow to the second by applying the scaled adjoint to the first map and taking the orthogonal sum. The Weil operators then act on \(\operatorname{Hom}(V',V^-\oplus V'')\) by its ordinary Weil operator times
\[
 \begin{split}
 M(d,e,f)&=\frac{a_d(y)}{a_e(x)}\frac{b_e(z)}{b_f(y)}
                         \frac{u(x)^e}{u(y)^d}\\
 &=\frac{a(y)b(z)}{a(x)b(y)}\,
   u(x)^{\lceil e/2\rceil}u(z)^{\lfloor e/2\rfloor}
   u(y)^{-(d+f)/2}.
 \end{split}                                                    \tag{WS1}
\]
The last factor in the first line is the character of the adjoint flip. Indeed negating all signs replaces the number \(n_-\) of negative signs by \(n-n_-\), so the standard character \(u^{n_-}\) becomes \(u^n/u^{n_-}\). The two determinant powers produced on transposing an arrow are \(u(x)^e/u(y)^d\). The coefficient of \(u(y)\) in the second line uses the fact that \(d,f\) have the same parity. The orthogonal sum itself introduces no character, since the two Weil operators act on its two blocks separately.\formalnote{5.55}{55}

For the comparison with the smaller ranks we need the effect of a zero term. If an element acts by a matrix \(A\) on the isotropic frame and induces \(g_0\) on its orthogonal quotient, that effect in the chosen normalization is
\[
   Z\omega(g)=\upsilon(\det A)^{-1}|\det A|_{\mathbb A_E}^{1/2}
                           \omega(g_0)Z.                         \tag{WS2}
\]
In mixed coordinates the Levi factor takes the remaining slice parameter \(v\) to \(A^*v\). Its Jacobian gives the half power: preservation of the integrated squared norm of the slices gives \(|c|^2=|\det A|_{\mathbb A_E}\) for the scalar on the zero slice. Its argument is obtained by pairing an adapted vacuum with its Levi translate. In adapted coordinates the overlap factors into the hyperbolic and the small-space overlap. For the hyperbolic multiplier \(m\), the chosen Weil normalization gives \(m\upsilon(\det A)>0\); the mixed pairing gives \(c\,J\,\bar m=I\) with \(I>0\) and \(J\ge0\) integrals of products of vacuum Gaussians and finite vacuum indicators. Thus \(c\upsilon(\det A)>0\), which determines \(c\). A radical element acts on \(Z\) by a scalar since it fixes the isotropic frame and the quotient and hence preserves the Heisenberg transformation law of \(Z\). This law determines the slices on translates of an adapted vacuum, which span, by \S\ref{sec:vacuum-projection}. The scalar of the radical element is one: it is a commutator of a Levi scalar dilation and a radical element. In the two-step radical a dilation may be chosen to multiply its first block by \(2\) and its central block by \(4\); solving for a prescribed commutator first in the first block and then in the centre uses the invertibility of \(2-1\) and \(4-1\). Scalars have trivial commutators. This proves (WS2) for the parabolic elements used here.\formalnote{5.56}{56}

The first zero term discards the positive weights in the space of maps. On their coordinate subspace the determinant of the matrix action is \((xz)^{s'/2}y^{-P}\). Indeed on each negative-weight column of the middle string the variable part of the left action has determinant \(xz\), supported in weight zero of the outer space. On each positive-weight row of the outer space the variable part of the right action on the weight-zero columns has determinant \(y^{-1}\). All remaining diagonal blocks have determinant one. The second zero term discards the isotropic halves of the \(h\) planes. Here the middle group acts on the columns by inverse transpose, giving \(y^{-h}\). All determinants have module one. The two factors from (WS2) therefore multiply to
\[
                   u(xz)^{-s'/2}u(y)^{P+h}.
\]
By the two rank equalities this is exactly \(M(D,E_0,F_0)/M(d,e,f)\). In particular the full zero term intertwines the large and small \emph{twisted} triple actions exactly. The two zero terms compose because the large isotropic frame consists of the first frame and the image of the second under the complementary-space inclusion. Successive integrals over the two compact additive quotients are the product probability integral.\formalnote{5.57}{57}

\subsection{Translations and the definite matrix coefficient}
For the standard global input we use Weil's rational invariance of the theta sum on a finite-dimensional vector space over a number field, for its adelic Schwartz functions and the rational splitting. Here the symplectic space is obtained from the hermitian pairing by taking its alternating trace with the fixed imaginary scalar. Its rational vectors are the rational matrices in the given frames. The character on the additive adeles is trivial on the rational field. The kernel with the unitary splitting retains rational invariance since its additional determinant characters are automorphic \cite[no.~41, eq.~(39) and Theorem~6]{legacy-weil}. The tensor theta functional in the joined model differs from the product functional by the fixed nonzero scalar \(z_\theta\) in \S\ref{sec:seesaw-levels}. In the orthogonal sum, the functional factors on pure data. To compare the two functionals under the flip, use the normalization for which the mean of \(|\theta(\rho(x)v_0)|^2\) is one, where \(\rho(x)\) is Heisenberg translation and \(v_0\) is the unit vacuum. The mean is taken on the additive adelic quotient by the rational vectors with its Haar probability measure. The flip transports the rational vectors, the translations and the vacuum, so both means are one. The two functionals invariant under the rational Heisenberg operators are proportional by uniqueness, and their equal means give a proportionality scalar of modulus one.\formalnote{5.58}{58}
\phantomsection\label{route:34}\leanroute{Weil's rational splitting gives invariance of the adelic theta sum in the Schrödinger model. The formal proof works with a functional invariant under the rational Heisenberg operators and normalizes its value on translates of the vacuum by a mean square. The comparison under the flip uses this latter normalization, as above.}

Let \(T,T'\) be the two string maps and \(U=\binom{T^*}{T'}\). The common generic constant of the two gradings makes \(TT^*=T'^*T'\) as endomorphisms of \(V'\). With the minus sign on \(V\) this says \(U^*U=0\). Moving a middle unipotent \(n\) across the two kernels replaces \(T,T'\) by \(n^{-1}T,T'n\), and the joined frequency is \(Un\). Its displacement is \(U(n-1)\). The string maps have degree at least \(-1\), while \(n-1\) raises degree, so the displacement has no negative-degree coordinate. The positive-degree part is absorbed by the first zero-term average. In degree zero the frame formula leaves precisely the matrix of entries \((n-1)_{\lambda i}\) from grade-zero \(i\in X'\) to the distinguished grade-one vectors \(\lambda\), multiplied by the common generic constant. To see this, on every degree-zero block of \(U(n-1)\) only the degree \(-1\) block of \(U\) and degree \(+1\) block of \(n-1\) contribute. The string formula identifies its output with the isotropic half of the \(h\) planes. This also covers \(h=0\) in the last pair.\formalnote{5.59}{59}

In the Heisenberg multiplication law the positive-degree frame is orthogonal to the grade-zero inclusion, so factoring these two translations introduces no phase. Composing with the original translation introduces the bicharacter of \(U(n-1)\) and \(U\). Its trace pairing factors through \(U^*U\), hence it too is one. The middle unipotent average now takes exactly the second zero term. The first unipotent average cancels against the character of the incoming Bessel coefficient as in \S\ref{sec:orbit-unfolding}. With the two quotient-measure constants \(\kappa,\kappa'\), the result is \(\kappa\kappa'/z_\theta\) times the double integral on the two definite quotients with datum \(Z\rho(U)(f\boxtimes f')\). All three constants are independent of the two twisting characters. The definite quotients are compact, so the continuous integrands are integrable in either order.\formalnote{5.60}{60}

One more uniformity matters in iterating this formula. The linear functional defining the outgoing coefficient must work for every finite datum, tame character and incoming form. For real polynomial vectors split a monomial into its two disjoint blocks of variables. The double integral is then an iterated ordinary theta lift, with \(z_\theta\) absorbed into the second polynomial; extend by linearity in the polynomial and finite tensor. For a general real vector \(v\), send \(v\) to the family of all these integral values indexed by the finite data, twist, input form and evaluation point. Each coordinate of this map is continuous by integration over the compact definite quotients. The families obtained by evaluating the composite lift at a single outgoing functional form the image of a finite-dimensional vector space in this Hausdorff product, hence are closed. Polynomial vectors are dense, so the entire family for \(v\) comes from one functional. In the fixed-real-point family used for the chain the archimedean character weights are independent of the tame twist. Thus the real vector used here is common to all finite data. The induction on the number of pairs can use this one functional at its next stage.\formalnote{5.61}{61}

\subsection{The finite frame and the gauge of an algebraic form}
Let \(L=\widehat O_E^n\) in orthonormal finite coordinates, \(v_0\) the finite vacuum, \(H(x)\) the finite Heisenberg translation including its lattice semicharacter, and \(G\) a finite unitary matrix. Choose representatives \(b\) for \(L/(L\cap G^{-1}L)\), of cardinal \(I\). The lattice projection formula of \S\ref{sec:vacuum-projection}, followed by the Heisenberg multiplication law, gives
\[
 \omega(G)H(x)v_0 =
       \sum_{b}\frac{\Delta_f(G)}{I}\,\zeta_b\,
                          H\bigl(G(x+b)\bigr)v_0,\qquad
 I=|\Delta_f(G)|^2.                                            \tag{WS3}
\]
The \(\zeta_b\) are finite additive character values and hence roots of unity. For integral \(d_0\), the condition that \(d_0 x\) be integral away from the conductor norm \(q\) passes to \(G(x+b)\) whenever \(G\) is integral invertible away from \(q\). Moreover \(\Delta_f(G)/I=1/\Delta_f(G^{-1})\). To determine its arithmetic size at \(\ell\nmid q\), use the matrix \(B\) spanning \(L\cap GL\) in the multiplier of \(G^{-1}\). At every place above \(\ell\) this intersection equals \(L\). Hence \(\det B\) is a unit there. The multiplier is the finite value of the algebraic character \(\upsilon|\cdot|^{-1/2}\) at \(\det B\), whose \(\ell\)-adic norm is one for a unit at all those places. Thus every coefficient in (WS3) is a unit at \(\ell\) for every coefficient embedding. The determinant characters of a twisted arrow have the same unit property: their algebraic weights express their \(\ell\)-adic norms as products of powers of local norms of the two endpoint determinants above \(\ell\). For empty conductor, \(G\) preserves \(L\) everywhere, \(I=1\) and \(\Delta_f(G)=1\); the determinant characters here have modulus one under all complex automorphisms by the weight formula.\formalnote{5.62}{62}

There is one finite-coordinate difference from the adelic calculation. Writing \(s_L\) for the semicharacter of the standard lattice model and \(\eta\) for its finite integral unit with \(\bar\eta\eta=-1\), translation of a joined datum by two row blocks \((A,B)\) gives the product of translations by \(\eta A^*\) and \(B\) times \(s_L(A)^2\); here \(A^*\) denotes conjugate transpose in the orthonormal coordinates. Indeed \(H\) includes the semicharacter of both blocks. In pulling the first block through the flip its norm form changes sign, so the ratio of the two semicharacters on that block is its square. This is another root of unity. Conjugation exchanges \(\pi_M\) and \(\bar\pi_M\), explaining the distinct clearing factors on the two blocks before the neighbour-lattice calculation. Section~\ref{sec:neighbour-lattices} computes the resulting lattice after the zero term, including the exceptional top row above \(N\). The counts of the two input residue matrices and one output residue matrix give the per-pair denominator dividing \(B_0^4\), \(B_0=M^{[F:\mathbb Q] n_m^2}\), as in \S\ref{sec:seesaw-levels}. Taking the tensor over \(m/2\) pairs gives \(B_0^{2m}\) and keeps all scalar changes fixed across twists or integral away from \(q\) as specified in that section.\formalnote{5.63}{63}

Finally consider a definite algebraic form of weight \((\lambda,e_0)\). Fix \(\iota:\mathbb C\to\mathbb C_\ell\). If \(\xi\) is its vector of values at the finite point \(g\), its gauge there is
\[
       \left\|R_{\lambda,e_0}\big((g_\tau^{-1})_\tau\big)
                       \,\iota(\xi)\right\|_\infty ,
\]
where \(g_\tau\) is evaluated through the completion induced by \(a\mapsto\iota(\Phi_\tau(\bar a))\). A rational change \((g,\xi)\mapsto(\gamma g,R(\gamma)\xi)\) leaves the gauge equal by cancellation inside \(R\). A right change \(g\mapsto gk\) integral invertible above \(\ell\) also leaves it equal. At each relevant embedding the Weyl lattice is preserved by the integral matrix \(k_\tau\) and its inverse; determinant twists act by units. Applying both matrices proves equality of the coordinate norms. Thus the bound may be tested at points integral above \(\ell\): to move a prescribed \(g\) there by rational \(\gamma\), approximate its inverse at the finitely many places of \(F\) over \(\ell\). For the unitary group this weak approximation follows directly from the Cayley transform. Its chart consists of group elements \(h\) for which \(1+h\) is invertible, with coordinates \((1-h)(1+h)^{-1}\) in the \(F\)-vector space of skew-adjoint endomorphisms. In each prescribed open neighborhood a scalar of norm one sufficiently close to one moves a local element into the chart: only finitely many scalar values are excluded by its eigenvalues, and the scalars \((1+j t)/(1-j t)\), \(\bar j=-j\ne0\), give infinitely many values near one. Weak approximation in an \(F\)-vector space, followed by the inverse Cayley transform, now supplies the rational element. At integral points the gauge is the ordinary sup norm. Its ultrametric inequality shows that a sum with integral coefficients preserves a common coordinate bound, including a divisibility bound \(\ell^{-v}\). This proves that the coordinate divisibility estimate implies the gauge bound in Theorem A.\formalnote{5.64}{64}


\section{Orbit selection in a single theta lift}
\label{sec:orbit-unfolding}
In the period identity, each theta kernel is summed over rational
linear maps. Integrating against a Bessel character selects terms
whose unipotent phase agrees with that character. The principle is
the same as in a Fourier coefficient on \(\mathbb R/\mathbb Z\):
\(\int_0^1e^{2\pi i(a-b)t}\,dt\) vanishes for distinct integers
\(a,b\). Here the phases depend on the Gram matrices of the
rational maps. The main step is to show that agreement of the
tested entries singles out the required orbit. For an orthogonal--symplectic pair, Wang unfolds a period of a theta lift along a nilpotent orbit into periods of the original form \cite[Proposition~3.2 and Theorem~4.5]{upper-wang}. For unitary groups, Furusawa and Morimoto obtain a Whittaker period from Bessel periods in the equal-dimensional case \cite[Proposition~3.1]{upper-furusawa-morimoto}. Here we follow the orbit at each step between the prescribed unitary spaces.

Consider two consecutive spaces in \S\ref{sec:seesaw-levels}. Write \(G'=U(V')\) for the target. The source grading presents \(V\) as a definite space \(X\) of weight zero and a string \(e_0,\ldots,e_{s-1}\) of weights \(g(a)=s-1-2a\). The target has its own definite part and a string of length \(s-1\), with weights \(g'\). Its form pairs opposite weights. In these bases the rational map \(T:V\to V'\) takes each string vector except the last to a specified nonzero multiple of the same-index target vector, and kills the last vector and \(X\). Its degree is \(-1\). The operator \(T^\dagger T\) is the source shift, sending a string vector to a nonzero multiple of its successor wherever a successor exists. The dagger is the adjoint for the two forms. Write \(B,B'\) for the Bessel groups associated with the two gradings and \([B]=B(F)\backslash B(\mathbb A_F)\).

\subsection{The phase detects a Gram entry}
For a target grading \(D\), a matrix \(A\) has degree at least \(i\) when \(A_{ba}=0\) for \(D(b)<g(a)+i\). In the expansion of the Weil theta kernel collect the rational maps by their classes modulo maps of degree at least zero. Let \(f_w(r,h)\) be \(W(r)\) times the resulting constant term at \(w\), where \(r\in B(\mathbb A_F)\), \(h\in G'(\mathbb A_F)\), and \(W\) has the prescribed Bessel equivariance. The phase attached to a map \(y\) and a source unipotent \(u\) is
\[
 P_y(u)=\psi_E\!\left(\frac{j}{2}\operatorname{tr}(y^\dagger y u)\right).
\]
Here \(j\) is the fixed scalar of the Weil datum and \(\psi_E\) is its standard adelic additive character. Take \(u\) with \(u-1\) of degree at least two and \(w(u-1)\) of degree at least one. The change of variable on the constant term and the equivariance of \(W\) give \(f_w(ur,h)=P_T(u)P_w(u)^{-1}f_w(r,h)\). Say \(w\) survives the test when the two phases agree for every such \(u\).

Here is how to read matrix entries from this equality. Write \(\bar k\) for the basis index paired with \(k\) by the source Gram matrix \(J\), and \(J'\) for the target matrix. For a vector \(v\) and an index \(k\) impose the following support conditions:
\[
\begin{array}{ll}
 v_a\ne0 &\Longrightarrow\quad g(a)\ge\max(1,2+g(k)),\\
 (wv)_b\ne0 &\Longrightarrow\quad D(b)>\max(g(k),-1),\\
 w_{b,\bar k}v_a\ne0 &\Longrightarrow\quad D(b)+g(a)>0.
\end{array}
\]
Take \(x\) supported at \(\bar k\). Use the unitary Eichler transformation
\[
 u(z)=z+\langle x,z\rangle v-\langle v,z\rangle
           \bigl(x+\tfrac12\langle x,x\rangle v\bigr).
\]
The supports force the stated two degree bounds on \(u-1\) and \(w(u-1)\). Substitution into the trace defining \(P_y\), for \(y=T,w\), gives
\[
                 P_y(u)=\psi_E\!\left(j x^*y^*J'yv\right).
\]
Indeed the two mixed rank-one terms pair under conjugation in the additive character. The term involving \(\langle x,x\rangle/2\) and the constant trace give zero in its symmetrized argument. Vary the one coordinate of \(x\). Since \(J^{-1}_{k,\bar k}\ne0\) and \(y^*J'y=J y^\dagger y\), nondegeneracy of the additive pairing gives
\[
                  (T^\dagger T v)_k=(w^\dagger w v)_k . \tag{U1}
\]
This tests exactly the entries permitted by the three support conditions.\formalnote{5.65}{65}

\subsection{Two gradings select a single block}
Take first \(D=0\). For source vectors \(e_i,e_k\) of positive weight, apply (U1) to \(e_i\) and the coordinate \(\bar k\). The inequalities hold since \(g(i),g(k)\ge1\). Multiplying by \(J_{k,\bar k}\) gives equality of the two Gram matrices of the positive columns. Those columns are independent for \(w\). Indeed, extend a nonzero linear relation by zero to a vector \(v\) on all of \(V\), and let \(a\) be the largest string index with \(v_a\ne0\). Then \(wv=0\). At \(k=a+1\), one has \(g(k)=g(a)-2\), and all nonzero coordinates of \(v\) have weight at least \(g(a)\ge1\). Thus all three tests in (U1) hold. Its right side is zero while the left side is a nonzero multiple of \(v_a\). The same independence holds for \(T\) by its explicit columns. The isometry between the two spans, \(we_i\mapsto Te_i\), now extends by Witt's theorem to a rational unitary \(\delta\). This use of Witt extension also covers a degenerate restriction of the ambient form to the column span.

Call a class \emph{agreeing} if \(w_{ba}=T_{ba}\) whenever \(g(a)>0\) and \(g'(b)<g(a)\), and denote the set by \(\mathcal C_D\). This condition is independent of representative for \(D=0,g'\). We have just shown that a survivor for \(D=0\) has a rational translate in \(\mathcal C_0\). We need the stabilizer of this set, since distinct translates must be counted once. If \(\gamma\in B\) and \(\delta\in B'\), then \(\delta T\gamma-T\) has nonnegative degree for the two string gradings. The matrices of \(\gamma,\delta\) preserve the filtrations and their prescribed graded string vectors, so \(\delta w\gamma\) again agrees. Conversely assume \(w\gamma\) and \(\delta w\gamma\) agree. In each row of \(\delta-1\) determine the entries in decreasing order of target weight. The tested columns \(k\) are the string columns of nonnegative weight, and the tested rows satisfy \(g'(b)\le g'(k)\). For such a column the string supplies a source column of weight \(g'(k)+1\) whose \(k\)-entry is a unit. In this column the tested entry of \((\delta-1)w\gamma\) is the unit times \((\delta-1)_{bk}\): lower-weight contributions vanish by agreement, and higher ones by induction. Agreement makes the product zero. These vanishings are the entrywise description of \(B'\), with the definite weight-zero block allowed its unitary action. Thus a target element taking an agreeing class to an agreeing class lies in \(B'\). In particular the distinct translates of \(\mathcal C_0\) are disjoint.

Now use \(D=g'\). For an agreeing survivor we claim \(w-T\) has degree at least zero. Proceed by increasing row weight \(g'(b)\). The only unproved entries satisfy \(g'(b)<g(a)\le0\). The dual target index \(\bar b\) has positive weight. Take a source column \(i\) of weight \(1+g'(\bar b)\) with unit entry \(T_{\bar b,i}\). Apply (U1) to \(v=e_i,\ k=\bar a\). To check the supports, \(g(i)=1-g'(b)\ge2-g(a)\), and agreement gives \(w e_i=T e_i\) in every row below \(g(i)\). Hence the only possible nonzero entries of \(we_i\) have weight at least \(g(i)-1=-g'(b)>-g(a)\). Finally, if \(g'(p)+g(i)\le0\), then \(g'(p)<g'(b)\). The induction hypothesis gives \(w_{pa}=T_{pa}=0\), the last equality following from the degree of \(T\) and \(g'(p)<g'(b)<g(a)\). This verifies all three conditions. In the product \(w^\dagger w e_i\) the entries paired with rows of weight above \(g'(\bar b)\) are already known, since their dual rows have weight below \(g'(b)\). The entries below \(g'(\bar b)\) in column \(i\) vanish by agreement. The remaining term has the unit factor \(T_{\bar b,i}\). Cancelling it in (U1) gives \((w^\dagger)_{\bar a,\bar b}=(T^\dagger)_{\bar a,\bar b}\), hence \(w_{ba}=T_{ba}\). This finishes the induction. Modulo nonnegative maps the only possible agreeing survivor for the second expansion is therefore the class of \(T\).\formalnote{5.66}{66}

\subsection{From character cancellation to the unfolded integral}
Consider a family of killed classes stable under \(B(F)\). Group the terms into \(B(F)\)-orbits. The sum for the orbit of \(w\), integrated over \([B]\), unfolds to the integral of \(f_w\) over the quotient by its rational stabilizer \(\Gamma_w\). A failed phase test gives a left translation \(u\) with \(f_w(ur,h)=c f_w(r,h)\), \(|c|=1,\ c\ne1\). One has to justify cancellation on this quotient even when left translation by \(u\) need not preserve its cosets. For fixed \(h\) put \(f(r)=f_w(r,h)\). Form the closed subgroup \(S\) of translations preserving, as a function of \(t,y\), the defect
\[
                     \big|f(tuy)-c f(ty)\big|.
\]
It contains \(\Gamma_w\) by invariance and \(u\) since \(|c|=1\). The defect at \(t=1\) vanishes, hence \(f(suy)=c f(sy)\) for \(s\in S\). In the quotient integration formula through \(S\), right translation by \(u\) in the \(S\) variable scales its left Haar measure by a positive real factor \(d\). The two evaluations of the inner mean give \(c\) and \(d\) times that mean. Since a positive real of modulus one equals \(1\), \(c\ne d\). The quotient integral vanishes. This proves cancellation for the orbit after unfolding to its stabilizer.\formalnote{5.67}{67}

Apply this twice. At \(D=0\), classes outside the union of translates of \(\mathcal C_0\) form a stable killed family. Inside one translate it suffices to calculate the sum on \(\mathcal C_0\). To pass from \(\mathcal C_0\) to \(\mathcal C_{g'}\), both sums are the same Fourier coefficient. Indeed let \(v:V'\to V\) range over all adelic maps modulo rational maps. Keep the entries \(v_{ab}\) for \(g(a)>\max(0,g'(b))\), putting the others equal to zero, and call the resulting map \(y\). The coefficient for either grading is the mean of the Heisenberg translate by \(y^\dagger\) evaluated at zero, divided by \(\psi_E(j\operatorname{tr}(yT))\). In the Fourier expansion, averaging each extra coordinate kills every nonmatching rational frequency by character orthogonality. This gives precisely the agreement restriction in each grading. In the second grading the complement of the class of \(T\) inside \(\mathcal C_{g'}\) is a stable killed family by the row induction. The whole block therefore integrates to the \(T\)-term \(W'(h)=\int_{[B]}W(r)\operatorname{CT}_T(\omega(r,h)\Phi)\,dr\).\formalnote{5.68}{68}

Accounting for the translates identified above yields
\[
 \int_{[B]} W(r)\theta(\omega(r,h)\Phi)\,dr
       =\sum_{\delta\in B'(F)\backslash G'(F)}W'(\delta h).
\]
For convergence, the sum of absolute constant terms of a Schwartz datum is bounded by a fixed multiple of its defining norm. The Bessel quotient has a compact set of representatives, so the same bound after the source translation is uniform in \(r\) up to the target Weil operator norm. In target grading zero every rational target translate is allowed. Reduction gives one such translate at which that norm is at most \(2^b H(h)^b\), with \(b\) fixed. For a rapidly decreasing automorphic test function \(z\) on \(G'(F)\backslash G'(\mathbb A_F)\), \(|z(h)|H(h)^b\) is bounded. Consequently the sum of absolute values after multiplication by \(z(h)\) is bounded uniformly in both variables. The original target quotient has finite volume, and \([B]\) has finite mass. Tonelli and Fubini now give absolute convergence and permit the final unfolding over the disjoint translates of a target fundamental domain:
\[
 \int_{[B]}W(r)\!\int_{[G']}\theta(\omega(r,h)\Phi)z(h)\,dh\,dr
     =\int_{B'(F)\backslash G'(\mathbb A_F)} W'(h)z(h)\,dh.
\]
Section~\ref{sec:bessel-bounds} derives the absolute-sum estimate and compactness of the Bessel quotient, and identifies the height reduction used in the final integral. The local integral lattice bounds used in later stages of the seesaw are separate.\formalnote{5.69}{69}


\section{Extending an isometry of primitive integral frames}
\label{sec:integral-witt}
In \S\ref{sec:neighbour-lattices} the form on the common string is nondegenerate over the local field. Its Gram determinant in an integral basis can nevertheless be a nonunit. At such a place the Gram matrix has singular reduction, so the field formula for orthogonal projection along the string involves an inverse with nonintegral entries. We use an extension argument that requires only primitive integral frames. Let \(O=\widehat{\mathcal O}_E\) for the CM field \(E\) of the construction, write \(M^*=\bar M^{\,t}\), and let \(X,Y\) be \(n\times k\) integral matrices with integral left inverses. We prove
\[
                X^*X=Y^*Y\quad\Longrightarrow\quad
         \exists g\in\mathrm{GL}_n(O),\quad g^*g=1,\quad gX=Y . \tag{WT1}
\]
The hypothesis on the field used below is the existence of \(c\in O\) with \(c-\bar c\) a unit at every place. In the construction \(E\) contains the \(5040\)-th roots of unity. It contains \(i\) and \(\omega=(1+\sqrt{-7})/2\), an integer with \(\omega+\bar\omega=1\); thus \(c=i\omega\) gives \(c-\bar c=i\). In particular
\(\lambda=c/(c-\bar c)\) is integral and \(\lambda+\bar\lambda=1\), also at the primes over two.\formalnote{5.70}{70}

\subsection{An invertible matrix with two prescribed actions}
Choose \(L,M\) with \(LX=MY=1\). The matrix
\[
                a_0=Y L+M^*X^*(1-XL)
\]
satisfies \(a_0X=Y\) and \(a_0^*Y=X\). The first follows since \((1-XL)X=0\). For the second,
\(a_0^*Y=L^*Y^*Y+(1-L^*X^*)XMY=L^*X^*X+X-L^*X^*X\).
For any integral \(C\) put
\[
          a=a_0+(1-a_0L^*Y^*)\,C\,(1-XMa_0).
\]
Since \((1-XMa_0)X=0\) and \((1-a_0L^*Y^*)^*Y=0\), every such \(a\) satisfies \(aX=Y\) and \(a^*Y=X\). We claim \(C\) can be chosen to make \(a\) invertible.\formalnote{5.71}{71}

In general, given square \(D,U,V\), one can make \(D+(1-DU)C(1-VD)\) invertible over \(O\). First reduce at one finite place. Changing the source and target bases independently puts \(D\) in the form \(\operatorname{diag}(D_1,0)\), with \(D_1\) invertible. In detail, if the original \(D=P D_0 Q\), set \(C=P C_0 Q\), \(U_0=QUP\), \(V_0=QVP\), and factor \(P,Q\) out on the two sides. Take \(C_0=\operatorname{diag}(0,1)\) in the new bases. The correction has the form
\[
       \begin{pmatrix} B\\1\end{pmatrix}
       \begin{pmatrix} C_1&1\end{pmatrix},
       \qquad B=-D_1(U_0)_{12},\quad C_1=-(V_0)_{21}D_1 .
\]
Subtracting \(B\) times the lower block row from the upper one yields a triangular matrix with diagonal \(D_1,1\). This also covers either empty block. At each finite place lift the resulting entries of \(C\) arbitrarily to its valuation ring. They form a matrix over the full product \(O\), and the resulting determinant is a unit in every component. Apply this to the displayed correction of \(a_0\). We have found an invertible \(a\) with \(aX=Y,\ a^*Y=X\). In particular \(h=a^*a\) is an invertible hermitian matrix satisfying \((h-1)X=0\).\formalnote{5.72}{72}

\subsection{An orthonormalization that fixes the prescribed columns}
For every invertible hermitian \(h\) we prove a relative form of orthonormalization:
\[
             g_0=1+W(h-1),\qquad g_0^* h g_0=1 . \tag{WT2}
\]
An arbitrary change to an orthonormal basis would lose control of \(X\). To obtain the required congruence, put \(H=h-1\), \(Q(T)=\det(TI-H)\). Its coefficients are fixed by conjugation because \(H=H^*\). In the commutative algebra \(S=O[T]/(Q)\), conjugate coefficients and fix \(T\). The class \(1+T\) is a unit since \(Q(-1)=(-1)^n\det(h)\) is a unit. We will solve
\[
      u\bar u=(1+T)^{-1},\qquad u=1+Tv,\qquad v\in S. \tag{WT3}
\]
Evaluation \(T\mapsto H\) is legitimate by Cayley--Hamilton. With \(g_0=u(H)\), the adjoint becomes coefficientwise conjugation in this commuting algebra. Then \(W=v(H)\) gives (WT2); its left inverse \(g_0^*h\) makes \(g_0\) invertible by the determinant identity.\formalnote{5.73}{73}

Here is a proof of (WT3) in the untwisted case used here. Start in a finite conjugation-stable quotient of \(O\) and the corresponding finite quotient \(S_0\) of \(S\). Every finite commutative ring is a product of finite local rings: its Jacobson radical is nilpotent and sufficiently high powers of its distinct maximal ideals are comaximal, so the Chinese remainder theorem gives the factors. Conjugation permutes these factors. On a pair of exchanged factors, any fixed unit \(b\) has the form \((b_1,\bar b_1)\). Choose \(u=(b_1,1)\); for \(b=(1+T)^{-1}\) this satisfies the required congruence because \(b_1-1=-T_1(1+T_1)^{-1}\).

On a fixed local factor \(A\) with maximal ideal \(\mathfrak m\), the residue field has nontrivial involution because \(c-\bar c\) is a unit. The norm on residue field units maps onto the units of the fixed field: on \(\mathbb F_{q^2}^\times\) it is the power \(q+1\), with image of order \(q-1\). If \(T\) is a unit in \(A\), the congruence \(u\equiv1\bmod T\) is automatic. Take any lift of a residue norm solution. If \(T\) is a nonunit, start with \(u=1\) modulo the nilpotent ideal \(J=(T)\); \(b=(1+T)^{-1}\) already equals one there. In the first case set \(J=\mathfrak m\). In either case we have a unit \(u\) with \(u\bar u\equiv b\bmod J\). Suppose the congruence holds modulo \(J^j\) for \(j\ge1\), and put \(e=b/(u\bar u)-1\in J^j\). This is fixed. Replacing \(u\) by \(u(1+\lambda e)\) changes its norm by the factor
\[
             (1+\lambda e)(1+\bar\lambda e)\equiv 1+e\pmod{J^{j+1}} .
\]
Iteration terminates since \(J\) is nilpotent, and preserves \(u\equiv1\bmod T\) whenever this condition has content. This solves (WT3) in every finite quotient.

To pass to \(O\), use the coefficients of \(v\) in the basis \(1,\ldots,T^{n-1}\) as unknowns. For any finite quotient, the tuples solving (WT3) after reduction form a closed nonempty subset of compact \(O^n\). Finitely many quotient conditions hold simultaneously by solving in a single common refinement; conjugation-stable finite quotients are cofinal by intersecting an ideal with its conjugate. Compactness gives a common tuple, which satisfies (WT3) in \(S\) since the quotients separate coefficients. Rank zero is immediate. This argument proves (WT2) without dividing by two at any place.

Finally \((h-1)X=0\) gives \(g_0X=X\). Set \(g=ag_0\). Then \(gX=aX=Y\) and \(g^*g=g_0^*h g_0=1\), proving (WT1). The argument also supplies an orthonormal basis for each unimodular integral hermitian matrix by (WT2) alone.\formalnote{5.74}{74}

\subsection{Orbit determination by the common lattice}
Let \(F,F'\) be two orthonormal adelic frames in rational coordinates and \(\mathcal L=FO^n,\ \mathcal L'=F'O^n\) their lattices. Suppose their intersections with the fixed string space \(S_{\rm str}\) agree. In local coordinates each intersection is saturated in its containing lattice: if \(0\ne r\) and \(rv\) is in the vector subspace \(S_{\rm str}\), then \(v\) is already in it. Over a DVR the Smith form makes a saturated submodule a direct summand. Choose local bases for the common intersection; in the adelic frames these give integral primitive matrices \(X=F^{-1}B\) and \(Y={F'}^{-1}B\) with integral left inverses assembled componentwise. Equivalently one can first choose \(B\) for \(F\); equality of the intersections says the coefficients for which \(Bc\) belongs to the second lattice are again exactly the integral coefficients, proving primitivity there. The two Gram matrices both equal \(B^*J B\) for the ambient rational form \(J\). Equation (WT1) thus gives \(gX=Y\) with \(g\) an integral unitary. The isometry \(F'gF^{-1}\) fixes \(B\) pointwise and takes \(\mathcal L\) to \(\mathcal L'\). The whole string space is nondegenerate in the grading and is spanned by \(B\). Its pointwise stabilizer acts on its orthogonal complement, the definite part, and acts there without changing the string coordinates. This is the orbit identification used in the neighbour calculation.\formalnote{5.75}{75}


\section{The finite lattices in the seesaw}\label{sec:neighbour-lattices}

Two calculations in Section~\ref{sec:seesaw-levels} use the same
pair of basis vectors. One bounds the support of a period as the string
gets shorter. The other computes the lattice in a constant term after
two successive lifts. In both calculations the last definite vector
is glued to the middle string vector when the string has odd length.
Write \(O=\widehat{\mathcal O}_E\), let \(P\) contain the chosen half
of the places over \(N\), and put \(\pi=\pi_N\), \(N=\pi\bar\pi\).
The element \(N\) is a uniformizer at every place over \(N\) in this
construction and the residue fields there are \(\mathbb F_N\).

\subsection{Intersection, projection, and the induction on support}

Split a graded space as \(D\oplus S\), with \(D\) its definite part
and \(S\) the span of the string \(z_0,\ldots,z_{s-1}\). For the
finite point \(x\) set \(L_x^- = L_x\cap S\) and
\(L_x^+=\operatorname{pr}_S L_x\), and write \(L^\pm=L_1^\pm\).
Use \(e\) for the last vector of \(D\). Let \(a\) be the idele
generating the chain ideal, \(u=(a/\bar a)^{\lfloor(s+2)/2\rfloor}\),
and \(l_b=u\) for odd \(s\). For even \(s\) set
\(l_b=u/a\) in the first half and \(l_b=u\bar a\) in the second.
The standard lattice has generators \(l_b z_b\) and the definite
basis, with a single replacement for odd \(s\). Writing
\(z=z_{(s-1)/2}\), the replacement for \(e\) is
\(\bar\pi^{-1}e+(u/N)z\). In particular the exact membership test
on the coordinates \(y_e,y_z\) in this plane is
\[
       \bar\pi y_e\in O,\qquad  u^{-1}y_z-\pi^{-1}y_e\in O .
       \tag{NL1}
\]
Indeed the two expressions are the coefficients of the replacement
and of \(uz\). Putting \(y_e=0\) gives the intersection. Allowing
\(y_e\) to vary gives the projection. Thus \(L^-=\bigoplus_b O l_b z_b\),
and the projection enlarges just the middle ideal \(Ou\) to \(N^{-1}Ou\).
In particular \(L^+/L^-\) is zero for even \(s\), and for odd \(s\)
its localization at each place above \(N\) has dimension one over
the residue field.

There is a convenient way to compare strings of consecutive lengths.
Write their vectors in the normalized string basis
\(\mathfrak e^i\) of the string map, so the map \(h\) carries
\(\mathfrak e^i\) to the vector of the same index on the shorter
string and kills the last vector. Put
\(\mathfrak l_{s,i}=a^{s-i}\bar a^{-1-i}O\).
For \(\sum_i t_i\mathfrak e^i\) the condition for membership in
\(L^-\) is \(t_i\in\mathfrak l_{s,i}\) when \(2i<s\) and
\(t_i\in N^{-1}\mathfrak l_{s,i}\) otherwise. The condition for
\(L^+\) has the cutoff \(2i+1<s\). To check the change of basis,
put \(r=i-\lfloor s/2\rfloor\) if \(2i<s\), and put
\(r=i-\lfloor(s-1)/2\rfloor\) otherwise. Direct substitution gives
\(l_i O=(a\bar a)^r\mathfrak l_{s,i}\). The coefficient of \(z_i\)
in the string power \(\mathfrak e^i\) generates
\((a\bar a)^rO\) in the first case and \(N(a\bar a)^rO\) in the
second, since the generic string constant generates \(a\bar aO\).
Cancellation gives the two tests. For \(i<s-1\), the cutoff for
\(L^+\) agrees with that for the shorter intersection, and
\(\mathfrak l_{s,i}=a\mathfrak l_{s-1,i}\). Extending a tuple by
a zero last coordinate gives surjectivity. Hence\formalnote{5.76}{76}
\[
                          hL^+=a(L')^- .                   \tag{NL2}
\]

Let \(c_M\) be the chosen idele with ideal \((\pi_M)\), let
\(b\) generate the full modulus of the lower datum, and put
\(d=b/c_M\). Thus \(d\) generates the tame conductor, and both
\(d,\bar d\) are units over \(N\). The finite Weil character is
normalized by an idele \(\alpha\), and \(a=\alpha/\bar c_M\).
We claim that wherever the \(j\)-th stage is nonzero, every graded
Bessel translate of its finite lattice satisfies
\[
           \bar d^{\,j}L^-\ \subseteq\ L_x^-\ \subseteq\ d^{-j}L^- .
           \tag{NL3}
\]
At stage zero there is an integral point in the orbit. For the
induction a nonzero integral has a point where both its factors,
the preceding stage and the single string constant term, are nonzero.
Choose a graded representative of the preceding bounded orbit.

Here is the frequency constraint supplied by the single constant
term. The modulus datum is invariant under the Heisenberg translations
by \(b\) times an integral matrix. Indeed its additive Fourier expansion
has shifts \(t\) with \(\bar b t\) integral, so all these translations
commute with them and fix the vacuum. After changing frames by a
unitary \(g\) the period lattice becomes \(M O^t\), \(M=bg\), with
dual matrix \(M'=g/\bar b\). Let \(k\) span the isotropic subspace
of positive-degree maps. Nonvanishing of the finite constant term at
the shift \(y\) implies
\[
                           k^*(y-M'\ell)=0\quad(\ell\in O^t).
\]
To see this, change the columns of \(M'^* k\) by an invertible
matrix \(T\) so \(M'^*kT\) is integral and has an integral left
inverse \(H\). For every integral \(v\), the vector \(kTv\) is both
a period and an isotropic translation, under which the constant-term
operator is invariant. Its commutator with \(y\) must therefore act
trivially on the nonzero value. Self-duality of the additive character
gives \(t_0=T^*k^*y\) integral. Taking \(\ell=H^*t_0\) proves the
claim. In the original frames \(\ell\) defines an integral lattice
map \(W:L_x\to L_{x'}\). The orthogonality just proved says its
negative-degree entries are those of \((\bar b/\alpha)h\).
A Bessel change of frames preserves its degree \(-1\) part, since
it preserves the flag and the prescribed string shift on the associated
graded spaces. In graded lattices each homogeneous piece of a
lattice map is integral, by projecting each input and each output
to its grade. Since \(h\) kills \(D\) and has image in the next
string, this proves\formalnote{5.77}{77}
\[
                     (\bar d/a) h L_x^+ \subseteq L_{x'}^- .
                     \tag{NL4}
\]

Under the hermitian pairing, the dual of the intersection with
the string is the projection to it, since the ambient lattice is
self-dual and projection is adjoint to inclusion. Dualizing
\(d^jL_x^-\subseteq L^-\) therefore gives
\(\bar d^{\,j}L^+\subseteq L_x^+\). Now (NL2) and (NL4) give
\(\bar d^{\,j+1}(L')^-\subseteq L_{x'}^-\). Its dual gives
\(d^{j+1}L_{x'}^+\subseteq (L')^+\).
If the upper bound at this new stage fails, some \(v\in L_{x'}^-\)
has \(d^{j+1}v\) nonzero in \((L')^+/(L')^-\). By (NL1) the new
string is odd and this element generates the quotient at a place
over \(N\). At this place, \(d,\bar d\) are units, so
\[
 L_{x'}^+ \ \subseteq\ (L')^+=(L')^-+O_v v\
 \subseteq\ L_{x'}^- .
\]
Consequently the lattice there splits between the definite and
string summands. Duality gives the split at the conjugate place
as well. At this even chain stage the chosen ramification scalar
at the place in \(P\) acts on the definite summand and fixes the
string. Its action preserves the split lattice, is trivial away
from the pair, and belongs after conjugation to the imposed level
since the pair is disjoint from the level places. Moving the scalar
through the chain fixes the last lift datum. The ramification
projection of the source then makes the chain zero, contradicting
nonvanishing. This completes the induction.\formalnote{5.78}{78}

For completeness, (NL3) leaves only finitely many possibilities for
the string intersection. The quotient of its two bounding lattices
is a finite module supported on the tame conductor and its conjugate.
Away from those places the intersection equals the standard one, so
one may truncate each representative to those places. Two self-dual
point lattices with the same intersection are equivalent under the
unitary group of the definite part. Indeed a basis of the common
intersection is primitive in both frames and has equal Gram
matrices. The integral Witt extension theorem identifies the two
embeddings by an integral isometry. In the original space the
resulting isometry fixes the string pointwise and acts on its
orthogonal complement.

Finally, each orbit has a graded point. Choose by local Hermite
reduction a basis \(b\) adapted to the grade flag and let \(b_0\)
denote its block diagonal part. Put \(G_b=b^*Jb\), and let \(G_0\)
consist of the entries with weight sum zero. The entries of \(G_b\)
of positive weight sum, and those of \(G_b^{-1}\) of negative
weight sum, vanish by triangularity and the pairing of opposite
grades in \(J\). Hence taking weight-sum zero in \(G_b^{-1}G_b=1\)
shows that \(G_0\) is integrally unimodular. Over the chosen CM field
its unimodular hermitian form admits an integral orthonormal basis.
The forms \(G_b\) and \(G_0\) agree on the primitive columns of
nonnegative weight. Integral Witt extension therefore gives an
integral \(n\) fixing those columns with \(n^*G_b n=G_0\).
We claim \(n-1\) raises weights. Its nonnegative-weight columns
vanish. For a remaining column \(j\) and \(w_i\le w_j<0\), expand
\((n-1)_{ij}\) using \(G_b^{-1}(G_bn-G_b)\). A possible summation
index \(h\) satisfies \(w_h\ge-w_i\ge-w_j\), so its weight is
nonnegative. The corresponding row of \(G_bn\) equals that of
\(G_0\) because \(n\) fixes column \(h\). At \((h,j)\) this
equals \(G_b\), since \(w_h+w_j\ge0\). This proves the claim.
Since \(G_0=b_0^*Jb_0\), the matrix \(bn b_0^{-1}\) is a
unipotent isometry. Its inverse carries the lattice to the graded
lattice of \(b_0\). The integral hermitian theorems used here apply
to this CM field containing a primitive \(5040\)-th root.\formalnote{5.79}{79}

\subsection{A matrix for the neighbour}

Consider two consecutive arrows \(V\to V'\to V''\). Write
\(D_-,H,D_+\) for the three definite index sets on the upper chain.
At the end \(D_+\) includes the line coming from the last negative
string vector. The reduced finite coordinates are matrices on
\((D_-\amalg D_+)\times H\); use \(e,e''\) for the last basis vector
of each row block. All adjoints below are taken in the finite
standard coordinates furnished by the Weil frames.
Write \(j_P\) for the characteristic idempotent of \(P\),
\(\eta\in O^\times\) for the chosen finite unit with
\(\bar\eta\eta=-1\), and put
\[
 \rho=e+\bar\eta\bar j_P e'',\qquad
 A=j_P(\pi-1)+\bar j_P(\bar\pi^{-1}-1),\qquad
 G=1+A\rho\rho^* .
\]
Because \(j_P\bar j_P=0\), \(\rho^*\rho=1\). The scalar \(1+A\)
is \(\pi\) on \(P\), \(\bar\pi^{-1}\) on \(\bar P\), and \(1\)
elsewhere. Its norm is one, so \(G\) is unitary. Set
\(\beta=\pi_M\), \(\delta=c_M/\bar c_M\), and let
\(\kappa_b=1\) on all rows except the last target row of the final
pair, where \(\kappa_{e''}=\delta\). Then \(\delta=1\) above \(N\)
and \((\bar\beta\delta)=(\beta)\) as fractional ideals.
The neighbour lattice is the image of the standard matrix lattice
\(L_0\) under \(g=(\operatorname{diag}(\kappa)G)\otimes 1_H\).
The two factors on the rows commute since their nontrivial supports
are disjoint.

For an integral input column, the two exceptional output coordinates
of \(G\), in the order \(e,e''\), are
\[
\begin{array}{c|cc}
 &\text{first coordinate}&\text{second coordinate}\\ \hline
 P      &\pi x_e+(\pi-1)\eta x_{e''} & x_{e''}\\
 \bar P &\bar\pi^{-1}x_e & x_{e''}+
                       \bar\eta(\bar\pi^{-1}-1)x_e .
\end{array}                                                     \tag{NL5}
\]
On \(P\) the output satisfies \(v_e+\eta v_{e''}\equiv0\bmod\pi\).
On \(\bar P\) it satisfies \(\bar\pi v_e\in O\) and
\(v_{e''}-\bar\eta v_e\in O\). In each case solving the two equations
in (NL5) gives the converse. In particular the nonintegral part at
\(\bar P\) is one residue per column and per place, and it has a
representative with rows
\[
                n(y)_e=\eta\bar y/N,\qquad n(y)_{e''}=-\bar y/N ,
                \qquad y\in (\mathbb F_N^{P})^{H}.             \tag{NL6}
\]
Here \(y\) denotes integral representatives on \(P\), extended by
zero elsewhere. Multiplication by \(\eta\), conjugation, and reduction
give a bijection between the residues and the classes of the first
row. The second row has the required class because
\(\bar\eta\eta=-1\). Thus there are \(N^{|P||H|}\) local terms.\formalnote{5.80}{80}

\subsection{Integral lifting in the middle and at the end}

Write \(i\) for the inclusion of the reduced space in \(I^\perp\)
and \(k\) for the frame of its isotropic subspace \(I\).
We construct an integral lifting matrix \(B=ig+kC\). The verification
is columnwise. For each column \(x\) of \(G\) there is a coefficient
\(c\), fixed by \(\delta\), such that
\[
 \bar\pi x_e,\bar\pi x_{e''}\in O,\qquad
 c-x_e/\pi\in O,\qquad \bar\eta c-x_{e''}/\pi\in O .           \tag{NL7}
\]
Other rows are integral. The coefficients for the two exceptional
columns are respectively \(A\bar j_P/\pi\) and \(A\eta j_P/\pi\);
every other coefficient is zero. To check the signs, on \(P\) the
\(e''\)-column gives \(c=(\pi-1)\eta/\pi\), so the last difference
in (NL7) equals \(-1\). On \(\bar P\) the \(e\)-column gives
\(c=(\bar\pi^{-1}-1)/\pi\), so the first difference is
\(-1/\pi\), integral there, and the last difference vanishes. The
other entries follow directly from (NL5). All these \(c\)'s are
supported over \(N\).

Write \(\lambda=\bar\eta\bar\alpha/\alpha\) for the change of finite
frame on negating \(V\), and \(u,u''\) for the two odd-string
scales. They satisfy \(\lambda u=\bar\eta\delta u''\), since
\(u=(a/\bar a)u''\) and \(a=\alpha/\bar c_M\).
In a middle pair lift \(x\) to the column with definite coordinates
\(\lambda x_-\) and \(x_+\), and put the same coefficient
\(\lambda u c\) on both middle string vectors. The sum of those
two middle vectors is isotropic in \(V^-\oplus V''\). After undoing
the frame change on the first block, (NL1) asks exactly for the
first congruence of differences in (NL7). On the last block the
coefficient is \(u''\bar\eta c\) because \(\delta c=c\), giving
the second congruence. This proves integrality in both blocks.

In the final pair the string lengths are \(1,0,0\). The remaining
string vector \(z\) has norm \(-N\) in \(V\), and \(z/\pi\) has norm
one in \(V^-\). Thus \(e''\) lifts to \(z/\pi\), with the additional
row factor \(\delta\), and there is no middle plane. Here \(u''=1\)
and \(\lambda^{-1}\delta=-\eta u\). In the original source frame
the coefficient at \(z\), divided by \(u\), is \(-\eta x_{e''}/\pi\).
Subtract \(x_e/\pi\) as (NL1) prescribes. The result is integral, since
\[
 -x_e/\pi-\eta x_{e''}/\pi
        =(c-x_e/\pi)+\eta(\bar\eta c-x_{e''}/\pi).
\]
The target is standard in this case. The denominator on its lifted
last row matches the old source denominator because
\((\bar\beta\delta)=(\beta)\).

Call the resulting matrix of lifted columns \(A_0\), expressed in
the point lattice of \(V^-\oplus V''\). If \(R\) consists of the
definite rows of the middle point, it is integral by the even-string
membership test. On matrix spaces set \(B(Z)=A_0 ZR\). The column
construction reads \(B=ig+kC\) with \(C\) zero in positive grades
and given by the plane coefficients in grade zero. Also \(B^*\tau=0\),
where \(\tau\) is the stacked pair of string shifts. Indeed \(R^*\)
in the middle rational frame selects the definite summand, and
both the forward shift and the adjoint of the preceding shift vanish
on it. The quotient lattice of the finite zero term is therefore
exactly \(gL_0\). For an integral vector \(v\), \(Bv\) is an integral
lift of \(gv\). Conversely, if \(iw+kb\) is integral, then
\(B^*(iw+kb)=g^{-1}w\) is integral, since \(i^*i=1\),
\(i^*k=k^*k=0\). This proves both inclusions.\formalnote{5.81}{81}

\subsection{The vacuum, the gluing sum, and its denominator}

Let \(H(v)\) denote finite Heisenberg translation and \(Z_I\) the
finite zero-term operator. For any set of representatives \(r\) of
\(gL_0/(gL_0\cap L_0)\) one has
\[
                   Z_I({\bf1})=c_0\sum_r H(r){\bf1}_0 ,
                   \qquad c_0\ne0 .                           \tag{NL8}
\]
Indeed \(iw+kb\) integral makes \(H(w)\) fix \(Z_I({\bf1})\).
The \(gL_0\)-fixed line is spanned by the Weil translate of the
vacuum. Its expansion follows by averaging. Give \(L_0\) volume one
and denote the finite Weil multiplier by \(\Delta_f(g)\). Conjugating
the vacuum projection by \(g\) gives
\(\omega(g){\bf1}_0=\Delta_f(g)\int_{L_0}H(gv){\bf1}_0\,dv\).
The integrand is constant on the cosets of \(L_0\cap g^{-1}L_0\),
whose volume is \(1/[L_0:L_0\cap g^{-1}L_0]=|\Delta_f(g)|^{-2}\).
Consequently each term in the expansion of \(\omega(g){\bf1}_0\)
has the same coefficient \(\Delta_f(g)/|\Delta_f(g)|^2\).
In (NL8) this is multiplied by the scalar comparing the zero-term
vacuum with the Weil translate. To see that \(Z_I({\bf1})\ne0\),
apply the commutator argument above
with \(M=M'=1\). A translated vacuum can survive only for a shift
of the form \(iw+kb+\ell\), with \(\ell\) integral. Its image is
then \(H(w)Z_I({\bf1})\) times a root of unity by the Heisenberg
law and invariance along \(I\). If the vacuum died, so would all
its translates, which span the finite representation, whereas the
constant term is nonzero on a vacuum adapted to the splitting.
More generally \(gB^*(iw+kb+\ell)=w+gB^*\ell\), and
\(gB^*\ell\in gL_0\). Thus a surviving shift \(y\) may be represented
by \(gB^*y\) up to a root of unity. The extra shift \(\tau\)
in the evaluation disappears because \(B^*\tau=0\).\formalnote{5.82}{82}

Let \(n\) be prime to \(N\), and suppose a translation is cleared away
from \(n\) by \(\beta\) on the incoming rows and \(\bar\beta\) on the
outgoing rows. It is integral over \(N\). Subtracting its component
there changes the translated vacuum only by a root of unity. On the
remaining places the correction \(C\) vanishes, so \(gB^*y=i^*y\).
The latter has the same allowed row denominators. In detail, \(B\)
is integral and \(g\) there is row scaling by \(\kappa\); whenever an
input index \(l\) contributes to output row \(b\), its clearing ideal
is \((d_b\kappa_b)\) with \(d_b=\beta\) or \(\bar\beta\) according
to that row. Each term in \(d_b(gB^*y)_b\) is therefore integral
away from \(n\).

Off \(N\) the representatives in (NL8) are classes in a finite group
\(\mathcal A\), the scaled row lattice supported off \(N\) modulo its
integral elements. Each full representative is \(n(y)+a\), \(a\)
a representative of \(\mathcal A\). Put \(K=\#\mathcal A\).
The split-support commutators between \(n(y)\) and away vectors
are trivial. There is also no phase in factorizing the \(n(y)\)
operator into the tensor product: its two rows are supported on
\(\bar P\), so their self-pairings vanish. On the first arrow the
adjoint takes \(v_-\) to \(\eta v_-^*\), and on the second the
coordinate is \(v_+\). Thus (NL6) gives exactly the last column
\(-y/N\) on the first arrow and last row \(-\bar y/N\) on the
second. Their sum over \(y\) is the prescribed gluing operator,
with coefficient one for every \(y\). For a general away translate
\(v\), the two translations are \((\eta v_-^*,v_+)\), so both
are cleared away from \(n\) by \(\bar\beta\).

Finally write the sum over \(\mathcal A\) as \(K\) times its average
and absorb \(c_0K\) into the fixed scalar for this pair. Each summand
has coefficient a root of unity divided by \(K\), or zero; collecting
equal terms therefore gives total absolute value at most one under
each complex automorphism. To bound the denominator, send an away
vector \(v\) to
\[
  \bigl(\overline{\bar\kappa_b v_{bh}}\bmod\beta\bigr)_{b,h}
       \ \in (O/\beta O)^{(D_-\amalg D_+)\times H}.
\]
The entries being reduced are integral. If the reduction is zero,
\(v_{bh}\in\kappa_b\bar\beta O\subset O\). Thus the kernel is
contained in the subgroup of integral vectors, and \(K\) divides
the size of this residue matrix group. In fact \(K=1\) for every
middle pair because all \(\kappa_b=1\). In the last pair only row
\(e''\) contributes: at each place where \(\delta\) has a pole it
contributes \(\#(\delta O_v/O_v)^{|H|}\). Since \((\beta)\) and
\((\bar\beta)\) have disjoint supports, their product over these
places is \(\mathrm N(\beta)^{|H|}\). The manuscript's uniform
denominator in Section~\ref{sec:seesaw-levels} uses the full matrix
count in both cases. Section~\ref{sec:vacuum-projection} derives the lattice multiplier and the nonzero constant term of an adapted vacuum. The lattice classification above uses
the integral orthonormalization and extension of primitive frames proved in \S\ref{sec:integral-witt}.\formalnote{5.83}{83}


\section{The lattice vacuum and its constant term}
\label{sec:vacuum-projection}

Let \(v_0\) be the normalized vacuum, \(L\) its finite lattice and \(\omega(g)\) the Weil operator for a finite unitary change of frame \(g\). We compute the coefficient \(\langle v_0,\omega(g)v_0\rangle\), whose squared absolute value is \(1/[L:L\cap g^{-1}L]\). A second scalar \(c\) occurs in the constant term (VP4) along a null space. Its value can be absorbed into the choice of theta data once we prove \(c\ne0\).

\subsection{Averaging over the lattice}
Let \(L=\widehat{\mathcal O}_E^d\subset\mathbb A_{E,f}^d\) in the standard finite coordinates of the Weil model. Give \(L\) mass one. Write \(v_0\) for the normalized vacuum, \(H(y)\) for the finite Heisenberg translation with its semicharacter correction, and \(\psi_0\) for the finite additive character normalized to make \(L\) self-dual. The correction matters even when \(y\) is integral. In the raw model the finite vacuum at \(y\) is the semicharacter times \({\bf1}_L(y)\). Pairing a raw translate with the vacuum evaluates it at that point. The correction on the translate is conjugated in the first slot and cancels the semicharacter there. Thus
\[
      \langle H(y)v_0,v_0\rangle={\bf1}_L(y),\qquad
      H(x)H(y)=\psi_0(y^*x)H(y)H(x).
\]
The second equality follows by dividing the two product laws in opposite orders: the two trace-one-adele terms add to \(1\) in the character. On \(L\) the full multiplication cocycle evaluates an integral argument and is one. Hence \(x\mapsto H(x)\) is a representation on \(L\). For \(y\in L\) the unit vectors \(H(y)v_0,v_0\) have inner product one, and are equal.

Let \(T\) consist of a set of places and their conjugates; let \(1_T\) denote its adelic idempotent. Put \(P_T=\int_L H(1_Tx)\,dx\). Translation invariance fixes its image under every \(H(1_Tx)\), so \(P_T^2=P_T\). Inversion invariance and unitarity imply \(P_T=P_T^*\). On \(P_Tv=v\) the commutation law gives
\[
       P_TH(y)v=H(y)v\int_L\psi_0((1_Ty)^*x)\,dx
               ={\bf1}_L(1_Ty)\,H(y)v .                          \tag{VP1}
\]
The last equality is additive character orthogonality and self-duality. Choose representatives \(s(q)\) for \(\mathbb A_{E,f}^d/L\). Translating both entries in an inner product by \(-s(q')\) reduces orthogonality for distinct \(q,q'\) to the zero vacuum coefficient outside \(L\). Another choice of representative differs by a lattice translate and a scalar of modulus one. The \(H(s(q))v_0\) thus give an orthonormal basis of the algebraic finite part (the span of finite translates). On that part \(P_{\rm all}\) projects onto the single vacuum line. At some places only, (VP1) selects integral components and preserves all unrestricted components.\formalnote{5.84}{84}

\subsection{The phase and the index of a neighbour}
Take a finite unitary \(g\) and let \(L_g=L\cap g^{-1}L\). Locally use elementary divisors to choose a matrix \(D\) with \(DL=L_g\): for an elementary divisor \(\varpi^a\) of \(g\), the required exponent for \(D\) on that input coordinate is \(\max(0,-a)\). Only finitely many places contribute. Define \(\chi_f(b)=\upsilon(b)|b|_f^{-1/2}\), with \(\upsilon\) the unitary splitting character of the Weil datum, trivial on integral finite units in this construction. The finite multiplier is \(\Delta_f(g)=\chi_f(\det D)\). The formula is independent of the integral bases, and
\[
  \operatorname{vol}(L_g)=|\det D|_f,\qquad
  h=[L:L_g]=|\Delta_f(g)|^2 .
\]
This also proves the volume formula when \(g\) has nontrivial factors on both members of a conjugate pair.

The vacuum coefficient of the normalized operator \(\omega(g)\) is
\(\langle v_0,\omega(g)v_0\rangle=\overline{\Delta_f(g)}^{-1}\). Indeed the Weil operator on the vacuum is its multiplier times the integrated translates, and evaluation at zero gives \(\overline{\Delta_f(g)}\cdot(\omega(g)v_0)(0)=1\). By unitarity \(P_{\rm all}\omega(g^{-1})v_0=\Delta_f(g)^{-1}v_0\); \(g\) acts only on the finite factor. Conjugate this average by \(g\) and use covariance of \(H\) under the unitary action. It yields
\[
 \omega(g)v_0
   =\Delta_f(g)\int_L H(gy)v_0\,dy
   =\frac{\Delta_f(g)}{|\Delta_f(g)|^2}
       \sum_{y\in L/L_g} H(gy)v_0.                             \tag{VP2}
\]
There is no choice of phase left at the finite sum: for \(z\in L_g\), apply \(\omega(g)\) to the law \(H(y+z)=H(y)H(z)\) on \(L\) and use that \(H(gz)\) fixes \(v_0\). This proves constancy on each coset, including the phase. Terms from different cosets are mutually orthogonal by (VP1). Their common coefficient has modulus \(h^{-1/2}\), as also follows by taking the norm of (VP2). For the neighbour in \S\ref{sec:neighbour-lattices} the residues indexing this sum, and any additional factor at the last edge, were counted in that section. Here the same representative-independent scalar has been derived before counting residues.\formalnote{5.85}{85}

\subsection{The constant term of an adapted vacuum}
In the global model write a Witt decomposition as
\(V=lE^p\oplus iV_0\oplus kE^p\), with \(l,k\) mutually dual null frames. Use the corresponding adelic coordinates in evaluations. Define \(Zf(y)\) as the mean of \(f(iy+kb)\) on \(\mathbb A_E^p/E^p\). This function of \(b\) is periodic since the rational translation phase on \(kE^p\) is trivial against \(iV_0+kE^p\). One can see boundedness in the model by first splitting off the hyperbolic coordinates. On a pure tensor the mean at zero, considered as a functional of the \(V_0\)-factor, is invariant under its rational translations. Uniqueness of the rational theta functional makes it \(c(w)\theta_{V_0}\), with \(c\) bounded on hyperbolic data \(w\). The map on products is then \(c\otimes1\): on translated vacua its norm is uniformly bounded, and the coherent-state reproducing integral extends it to the full model. Evaluating after an \(iV_0\)-translation gives the stated mean at \(y\). Put \(S_x f=Z(\rho(lx)f)\). Commuting two translations and absorbing \(kb\) in the mean gives
\[
 S_x\rho(lc+iw+kb)f=\beta(2x+c,b)\,\rho_0(w)S_{x+c}f             \tag{VP3}
\]
with \(\beta\) the normalized phase of the null pair. Successive zero terms on compatible splittings are a product of quotient averages, by Fubini.\formalnote{5.86}{86}

Suppose \(g\) is adapted to \(t\in{\rm GL}_p(\mathbb A_E)\): the columns of \(g^{-1}(l(t^{-1})^*,kt,i)\) are integral in the finite standard frame and orthonormal at infinity. Let \(u=\omega(g)v_0\). We first prove with a possibly zero \(c\) that
\[
                     S_x u=c\,|\operatorname{vac}(t^*x)|^2 v_{0,V_0}.
                                                                    \tag{VP4}
\]
In scaled coordinates write \(t^*x=(a,z)\) with real Fock coordinate \(a\) and finite lattice coordinate \(z\). Explicitly for \(y=t^*x\) one has \(z=\alpha^{-1}y_f\) for the finite scale idele \(\alpha\), and \(\|a\|_2^2=2\pi\sum_{i,\tau}\Im(\tau(\jmath))|y_{i,\tau}|^2\) for the fixed imaginary Weil scalar \(\jmath\). Thus \(z\) integral means \(y_f\in\alpha\widehat{\mathcal O}_E^p\). At finite places adaptivity makes integrality equivalent to integrality in the three frame blocks. In (VP3) with \(c=w=0\), translating by an integral \(b\) leaves the input vacuum fixed. A nonzero slice must therefore have trivial phase for every \(b\), forcing \(z\in L_p=\widehat{\mathcal O}_E^p\). For integral \(z\) translate by \(-z\) in the first block and average the middle block. This identifies the finite slice with \(P_{{\rm all},V_0}\) of the slice with that finite coordinate zero.

At infinity insert the real vacuum projection \(\int e(y)\rho(y)\,dy\), with \(e(y)=e^{-\|y\|_2^2/2}\) in normalized coordinates, where the squared norm here is the sum of the squared coordinate moduli. The adapted frame preserves the measure and Gaussian. Write its variables as \((c,b,w)\). Integration in \(w\) gives the smaller real vacuum projection. In \(b\), the Fourier transform of the Gaussian at \(2a+c\) is \(2^j e(2a+c)\), where \(j\) is the number of complex coordinates in the null block and volume is normalized as in the Fock model. Multiplication by \(e(c)\) gives
\[
       2^j e(c)e(2a+c)=2^j e^{-\|a\|_2^2} e(a+c)^2 .
\]
Shifting the remaining \(c\) integral removes \(a\) from it. Comparison at \(a=0\) leaves the factor \(e^{-\|a\|_2^2}\). The two projections together extract \(\langle v_{0,V_0},S_0u\rangle v_{0,V_0}\). The squared vacuum function in (VP4) is exactly \(e^{-\|a\|_2^2}{\bf1}_{L_p}(z)\) in these coordinates, proving (VP4).\formalnote{5.87}{87}

\subsection{Nonvanishing of the scalar}
First prove that \(\sum_{\xi\in E^p}\|S_\xi f\|<\infty\), without assuming \(c\ne0\). In the adapted frame with \(t=1\), (VP4) and (VP3) bound the norm of the zero term on any translate of \(u\) along \(lx\) by a translate of an integrable Gaussian times a lattice indicator, say \(b(x)\). More precisely the norm as a function of \(x\) is precisely such a translate, times a fixed scalar in modulus; the other frame directions introduce an isometry and a phase. The reproducing integral for \(u\) gives a local average bound. To spell it out, for \(z\in V(\mathbb A)\) in frame coordinates put
\[
 G(s)=\int |f(z)|\,\|Z\rho(z+ls)u\|\,dz .
\]
Tonelli gives \(\int G(s)\,ds=\|b\|_1\|f\|_1\) when \(b\) includes the fixed scalar. Reproducing \(f\) against \(\rho(ls)u\), then applying \(Z\rho(lx)\), gives
\[
              |u(-ls)|\,\|S_x f\|\le G(x+s).
\]
The nonzero vacuum coefficient of \(\omega(g)\) gives \(u(0)\ne0\), irrespective of \(c\). Choose a compact neighborhood \(B\) of zero with \(d=\int_B |u(-ls)|\,ds>0\). Summing the inequality integrated on \(B\) shows that
\[
     d\sum_{\xi\in E^p}\|S_\xi f\|
       \le \#\bigl(E^p\cap(B-B)\bigr)\,\|b\|_1\|f\|_1 .
\]
Indeed a point belonging to several translates \(\xi+B\) determines, by subtracting any one of their indices, distinct rational vectors in \(B-B\). The rational lattice is discrete and closed, so this number is finite.

Now \(b'\mapsto f(kb')\) is continuous and periodic. On the rational index \(\xi\) its Fourier coefficient is \(\theta_{V_0}(S_\xi f)\): commute \(\rho(l\xi)\) past \(\rho(kb')\) in the average. The commutator is the Fourier character, and rational translation leaves the theta functional invariant. The rational \(\xi\)'s give the full dual of the compact additive quotient by adelic self-duality. Absolute summability allows evaluation of the Fourier series at zero, giving
\[
                       \theta_V(f)=\sum_{\xi\in E^p}\theta_{V_0}(S_\xi f).
\]
If \(Z=0\), every slice on the right vanishes for every \(f\). Translating \(v_0\) and using evaluation by \(\theta_V\) would give \(v_0(y)=0\) for every \(y\), contradicting its unit norm. Thus \(Z\ne0\). Finally if \(c=0\) in (VP4), (VP3) shows that \(Z\) vanishes on every Heisenberg translate of \(\omega(g)v_0\). Those translates determine a bounded operator by the vacuum reproducing formula and invertibility of \(\omega(g)\). This again forces \(Z=0\). This proves \(c\ne0\) for every adapted \(g,t\), and justifies division by the adapted-vacuum scalar in the lattice descent.\formalnote{5.88}{88}


\section{Simultaneous divisibility and integral spanning vectors}\label{sec:ab-closure}
Section~\ref{sec:upper-A} reduced the derivative estimate to congruences on pairs of directions. Here we choose a single set of integral changes of basis satisfying all the congruences at once. Their determinants measure the divisibility lost in replacing the original directions.

\subsection{Paying for all pairings with one change of basis per height}
Fix \(p,\pi=\pi_p\), a finite set of exponents and the height bound of \S\ref{sec:upper-A}. Put \(O=\mathcal O_{\mathbb C_p}\). For a place \(s\) and a height \(h\), the original frame consists of \(h\) vectors \(y_{s,k}\), \(0\le k<h\). We replace them by the rows \(\beta_{s,h,k}=\sum_l g_{s,h;kl}y_{s,l}\) of an \(O\)-linear combination. A height-\(h\) minor acquires the factor \(\det g_{s,h}\). Its valuation must be charged once for each column of this height, irrespective of how many couplings visit \(s\).

Consider a single arrow \(s\to s'=-p^e s\), put \(P=p^e\), and fix a \(P\)-th root \(\delta\) of \(\pi\). The pairing to annihilate modulo \(\pi\) is \(\sum_a\beta(a)^P\beta'(a)\). For a single source row facing a partner of height one, impose the linear congruence
\[
       \sum_a\beta(a)y_{s',0}(a)^{1/P}=0\bmod\delta .
\]
Any integral choices of the roots work. Raising to the \(P\)-th power modulo \(\pi\) gives the required relation, since \(O/(\pi)\) has characteristic \(p\) and the new height-one partner is a multiple of \(y_{s',0}\). In the opposite direction, impose \(\sum_a y_{s,0}(a)^P\beta'(a)=0\bmod\pi\) when the source has height one. A third option on either side is to make all \(h\) coefficients of the new row divisible by that side's modulus. This kills its pairing with every opposing row. The costs in units of \(v_p(\pi)\) for one pairing congruence are \(1/P\) at the source and \(1\) at the partner.

Let \(t_s\) count height-one columns after erasing the marked ones, \(c_s\) count columns of height other than one, and \(b_s\) count their boxes. In the Young diagrams all columns have positive height. If zero-height columns are formally allowed, their determinant is one and the same upper bound holds. Choose the congruences for this arrow as follows.
\begin{enumerate}
\item On every source row of height greater than one, impose the one source pairing congruence with the height-one partner. This costs \(c_s/P\).
\item For the unmarked height-one sources, choose the cheaper of making their entries divisible by \(\delta\) and imposing one partner pairing congruence for every partner column except the marked height-one columns. The cost is \(\min(t_s/P,c_{s'}+t_{s'})\).
\item For pairs of heights both greater than one, choose between imposing all the entry congruences on the source and doing so on the partner. The cost is \(\min(b_s/P,b_{s'})\).
\end{enumerate}
These cases cover every pair except those for which the source has height one and either the source is marked or the partner is both marked and of height one. The sum is exactly the deduction \(D_e(s)\) with the marked columns erased.

For each fixed \((s,h)\) now collect the forms from \emph{all} arrows, on both sides. They involve only the original frames and unknown coefficients of a row of \(g_{s,h}\). Starting with the identity matrix, suppose the values of the next form on the current rows are \(a_1,\ldots,a_h\) and its modulus is \(\alpha\). If one is nonzero, choose one of least valuation, \(a_i\). Subtract \((a_j/a_i)\) times row \(i\) from row \(j\ne i\), then multiply row \(i\) by \(\alpha\). The new determinant is the old one times \(\alpha\), the form takes values in \(\alpha O\), and previous congruences survive integral row operations. A zero form costs nothing. This bounds the valuation of the product of the determinants, with column multiplicity, by \(v_p(\pi)\sum_{e,s}D_e(s)\). The proof also applies to a loop \(s=s'\). A marked height-one exemption is harmless after multiplying a charged row variable by \(p^{w}\): the required positive charge at one end makes its contribution to the pairing divisible by \(p\), hence by \(\pi\).\formalnote{5.89}{89}

For completeness, introduce variables \(s_{h,k,j}\) for a row direction \(k<h\) and a coordinate \(j\) of the target at each place. Substitute into the theta variables the linear combinations \(\sum_{h,k}p^{w_{h,j}}\beta_{h,k}(a)s_{h,k,j}\), separately for each target coordinate \(j\). Expansion of the Frobenius power modulo \(\pi\) shows that every coefficient of a relevant coupling is one of the pairings just killed, multiplied by row-charge factors. Exponents \(p^e\) exceeding the total degree \(B\) cannot contribute to the jet. The exponential in \eqref{eq:Athetaform} consequently has integral coefficients through degree \(B\), by \(v_p(n!)\le n/(p-1)\). Differentiating the substituted series along the determinant in the \(h\) row directions of each column extracts precisely the original product of minor derivatives times \(\varpi^B\), the product of the determinants of \(g_{s,h}\), and \(p\) to the sum of charges over the boxes. This follows by expanding each alternating determinant of first derivatives and commuting constant-coefficient derivatives. The number of variables is at most \(fr^3\). The factorial estimate in \S\ref{sec:upper-A} now proves the charged jet inequality there, also when a relevant coupling has \(e=0\).\formalnote{5.90}{90}

\subsection{The complex norm of a Schur jet}
At a lift from rank \(r\) to \(r'\), let \(a_{\tau i}\) be the lengths of the first \(r\) rows in its polynomial weight, and put \(d_0=rr'\), \(D=fd_0\). The translated integral vacuum, after conjugating the finite data, satisfies a uniform bound \(\exp(C(q_1d_0+1)^2)\). We first account for this bound.

In rank one the finite lattice projection of a translated product of the elliptic theta sections is a nonzero multiple of the vacuum. Each summand is a product of \(q_1^2\) torsion values of the fixed CM elliptic theta function, times a root of unity, and the finite projection has \((q_1^2)^f\) summands. To see the cardinality, the product has translation level \(q_1\), and the symplectic quotient of a self-dual lattice by this level has \(q_1^{2f}\) elements. Finite-character orthogonality projects onto the invariant vacuum line, and translating first permits a nonzero projection since translates of the line span. The torsion values are integral by \S\ref{sec:cm-integral}. Their conjugate moduli are uniformly bounded because \S\ref{sec:cm-conjugation} transports each to a normalized translate on one of the finitely many ideal-class lattices, where the modulus is continuous and periodic. Hence the finite values of the chosen integral multiple are bounded by \(\exp(C_0 q_1^2)\) in rank one, using \(f\le q_1\). Taking tensor powers gives \(\exp(C_0 q_1^2 d_0)\). The map conjugating the jets sends this multiple to a scalar times a unit vacuum moved by a norm-one idele, by \S\ref{sec:cm-multiplier}. Its squared norm is the mean of the squared absolute values on the compact additive quotient. Finite points and their rational translates are dense there, and rational translates preserve absolute values. The same bound therefore controls the scalar.\formalnote{5.91}{91}

The unit vacuum after any finite motion also has a uniform bound. Expand the square of its theta value as the sum of its translated matrix coefficients over rational vectors. This identity follows by Fourier expansion on the compact rational additive quotient, or by unfolding the square and grouping by the difference. In absolute value each coefficient on its finite support is at most \(\exp(-\pi\|z\|^2/\sqrt d)\) in the archimedean coordinates of the moved lattice, with \(d\) the fixed quadratic discriminant parameter. The squared length of a difference of distinct support vectors is a positive integer. It is the trace to \(K\) of a Hermitian norm lying in the inverse relative different by the finite support condition, hence a real algebraic integer in the imaginary quadratic field. The support is thus one-separated in real dimension \(2D\). In the shell \(j\le\|z\|<j+1\), disjoint balls of radius \(1/2\) bound its size by \((2j+3)^{2D}\). Completing the square gives \(-\pi v^2/\sqrt d\le-(2D+1)v+\sqrt d(2D+1)^2\), and \(2j+3\le3e^j\). The sum is at most \(2\cdot3^{2D}\exp(\sqrt d(2D+1)^2)\). Together with the scalar bound and the unit-modulus multipliers, this proves the claimed bound for the transported vacuum \(F_\sigma\).\formalnote{5.92}{92}

For a polynomial \(Q\) representing a matrix coefficient against a unit dual vector, the conjugate jet is the integral of \(Q^\sigma(t\bar y)\exp(-\|y\|^2/2)F_\sigma(y)\) against \(dy/\pi^D\), with the places permuted as in \S\ref{sec:cm-multiplier}. Here \(Q^\sigma\) also includes the change of scale: the squared modulus of the argument scale is \(T=M/\Omega\). The formula follows termwise from the complex Gaussian moments \(\int z^a\bar z^b e^{-|z|^2}\,dz/\pi=\mathbf1_{a=b}a!\), applied to the Fock evaluation functional. The extra Gaussian factor is in its vacuum measure.

Put \(L=a+1\), where \(a\) is the large weight of \S\ref{sec:upper-A}, \(u=\log T\), and \(R=\|y\|^2/(2L+2)\). Let \(\lambda_{\tau i}\) be the squared singular values, in decreasing order, of the matrices at the scaled argument. Their sum is at most \(e^u(2L+2)R\). On the polynomial representation of highest weight \((a_{\tau i})_i\), the squared operator norm is at most \(\prod_i\lambda_{\tau i}^{a_{\tau i}}\). Indeed the two unitary factors of the singular-value decomposition preserve the norm, and on a diagonal matrix the largest weight monomial is the highest one by dominance. Apply \(z^l\le l!\exp(lv+ze^{-v})\) for \(z\ge0\), \(v=u+1/L\), obtained from the exponential series. Since \(e^{1/L}\ge1+1/L\), the sum of the last exponents is at most \(2LR\). Thus, for a unit input and a unit dual vector, the polynomial factor has absolute value at most \(\exp(w+LR)\), where
\[
                     2w=\sum_{\tau,i}
                      \left(\frac{a_{\tau i}}{L}
                           +\log(a_{\tau i}!)+a_{\tau i} u\right).
\]
After multiplication by the real vacuum the remaining Gaussian is \(e^{-R}\), whose normalized integral is \((2L+2)^D\). The logarithm of the resulting norm bound is consequently
\[
          C(q_1rr'+1)^2+D\log(2L+2)+w .
\]
The last two terms inside \(2w\) sum to the row-size sum \(u_j\) of \S\ref{sec:upper-A}. The other terms are bounded by \(x^{50}\) by the rank and level bounds at the prescribed cutoff, using \(a_{\tau i}\le2a\). Taking the supremum over unit dual vectors proves the kernel factor \(\exp(x^{50}+u_j/2)\) in (A5). All finite translations have the same bound. In a finite linear combination the triangle inequality supplies precisely the conjugate coefficient sum.\formalnote{5.93}{93}

\subsection{An integral spanning family with a uniform bound}
Write \(k\) for the definite rank, \(f=[F:\mathbb Q]\), and \(N_0\) for a positive integer in the different of \(E\). For a representative \(g\) of a finite class let \(Y(g)\) have as columns the independent vectors in \S\ref{sec:height-reduction}, for the lattice \(L_g=E^k\cap g\widehat{\mathcal O}_E^k\). In the definite case the shortest nonzero length is at least one. Indeed, the Hermitian norm of a nonzero element of \(L_g\) is a totally positive algebraic integer. The product of its real conjugates is at least one, so some summand in the squared Euclidean length is at least one. Therefore \(g^{-1}Y(g)\) is integral and every archimedean entry of \(Y(g)\) has modulus at most \(C=6N_0(fk)^2\).

At a free level a form is determined by one arbitrary Weyl-module vector per class. For the class \(g\) take the usual frame evaluation vector \(v(y)\) whose coordinates are the products of flag minors of shape \(\kappa\), with the entries of \(y\) ranging independently through roots of unity. Transform it by the fixed invertible operator
\[
  A_g=\left(\prod_\tau\det\bigl(\tau Y(g^{e_\tau})\bigr)\right)
       \operatorname{Schur}_{\kappa}\bigl((\tau Y(g))_\tau\bigr).
\]
Using a fresh short frame at \(g^{e_\tau}\) is essential for a determinant weight of arbitrary size or sign. In the finite gauge at \(g\), its determinant becomes that of the integral matrix \(g^{-e_\tau}Y(g^{e_\tau})\). The Schur factor of \(A_g\) uses \(g^{-1}Y(g)\), and the flag minors and Schur action have integral coefficients in the tableau lattice. Thus all evaluations are integral at every finite embedding.

Under any complex embedding the entries of \(y\) still have modulus one. The vector \(v(y)\) is obtained from the unit highest-weight frame vector by applying the Schur action of the square frame of \(y\). Each resulting matrix \(\tau Y(g)y_\tau\) has squared Frobenius norm at most \(k^4 C^2\le C^6\). The operator norm in the tensor representation of degree \(d\) is at most the \(d\)-th power of the underlying operator norm, and this bound descends to the unitary Schur component. Summing \(d\) over the embeddings gives the number \(d_\kappa\) of boxes in the polynomial part. Each determinant has squared norm at most \(C^{6k}\), by Hadamard. Thus the squared Weyl norm is at most \(C^{6(d_\kappa+kf)}\). Finally the evaluation vectors span. A linear functional annihilating them gives a polynomial in the entries of \(y\) vanishing on a product of infinite sets. Induction on the variables makes it the zero polynomial. Its coordinates are standard tableau polynomials, which are linearly independent by Howe's basis theorem \cite[Theorem~5.1.5 and Corollary~5.2.9]{theta-howe}. Invertibility of \(A_g\) preserves spanning. Extending each vector by the delta form on its class preserves integrality and its archimedean norm, by equivariance and freeness. Taking the union over the classes proves the asserted short spanning result.\formalnote{5.94}{94}
\phantomsection\label{route:9}\leanroute{Howe's basis theorem supplies the linear independence of the products of flag minors in the last step. Lean develops the tableau polynomials and their straightening rules internally. The chosen vectors, their integrality and their norm bounds above still use the explicit tableau lattice.}



\editionchapter{The second moment and the orbital sieve}\label{ch:span}

To use the determinant comparison one needs independent vectors. Theorem B obtains them from translates of a single linear map. We begin with the second moment of those translates and the dimension gained by adding one of them.

\section{Proof of Theorem B}
\label{sec:B}
To apply the volume lower bound, the span of the translated source
must occupy almost all of the target. A lower bound for the norm of
one translated map alone gives no control of the overlap between its
image and the images of the other translates. The overlap is measured
by a second moment, which bounds how much of the translated source can
fall into a subspace already spanned.

Let \(S=S_X\) be the full source space and \(Y\) the good space on \(Y_0\). Put \(s_0=\dim S\) and \(D=\dim Y\). The source has coordinates indexed by pairs consisting of a definite lattice class and a Bessel coordinate. Both indexing sets are nonempty, so \(s_0>0\). Write \(\Gamma\) for the product of residue-field general linear groups at the chosen places above \(p_0,p_1\). Let \(\rho\) be the unitary action of \(\Gamma\) on \(\mathcal A(Y_0)\) from \S\ref{sec:spaces}. Let \(P\) be the scaled Poincar\'e map from \(S\) to \(\mathcal A(Y_0)\). The central idempotent \(e_0\) from \S\ref{sec:spaces} acts as the orthogonal projection \(e=\rho(e_0)\) onto \(Y\), so \(e\) commutes with the action of \(\Gamma\). We use the Hilbert--Schmidt norm \(\|A\|_{\mathrm{HS}}^2=\operatorname{tr}(A^*A)\), written \(\|A\|_2^2\) throughout this section.\formalnote{6.1}{1}

\begin{example}[One line and its translates]\label{ex:cyclic-second-moment}
Let \(S=\mathbb C\), \(Y=\mathbb C^3\), \(e=1\), and let the cyclic
group of order three permute the coordinate vectors. If \(P(1)=e_1\),
the three translates of its image are the three coordinate lines.
The squared overlaps \(\lvert\langle P(1),\rho(g)P(1)\rangle\rvert^2\)
are \(1,0,0\), with mean \(1/3=1/\dim Y\). Now take
\(P(1)=(e_1+e_2+e_3)/\sqrt3\). All three translates give the same
unit vector, so the squared overlaps have mean \(1\).
The norm of \(P(1)\) is the same in both cases. The mean squared
overlap distinguishes a span of dimension three from a span of
dimension one.
\end{example}

\subsection{From squared overlaps to a new direction}
The argument uses two estimates. The first measures how much squared norm remains after projection onto the good space \(Y\). The second measures the overlap of the image of \(P\) with its translates. For an orthonormal basis \((u_i)\) of \(S\), the entries of \(P^*\rho(g)P\) are the inner products of \(Pu_i\) with \(\rho(g)Pu_j\). Suppose \(0<\epsilon<1/16\) and assume the inequalities below, the first for every subspace \(U\) of \(S\), with \(p_U\) its orthogonal projection.
\begin{align}
 \|ePp_U\|_2^2&\ge\dim U-\epsilon s_0, \label{eq:near}\\
 D\sum_{g\in\Gamma}\|P^*\rho(g)P\|_2^2
      &\le |\Gamma|(1+\epsilon)s_0^2. \label{eq:overlap}
\end{align}
Taking \(U=S\) shows that \(D>0\). Let \(F\subset Y\) occupy the share \(\alpha=\dim F/D\). Form the positive operator
\[
    C=\sum_{g\in\Gamma} C_g,\qquad C_g=\rho(g)PP^*\rho(g)^*.
\]
Expanding the squared norm and using cyclicity of trace gives
\[
 \langle C_g,C_h\rangle_{\mathrm{HS}}
    =\|P^*\rho(g^{-1}h)P\|_2^2,\qquad
 \|C\|_2^2=|\Gamma|\sum_{g\in\Gamma}\|P^*\rho(g)P\|_2^2.
\]
Indeed every difference \(g^{-1}h\) occurs exactly \(|\Gamma|\) times. Since \(\|p_F\|_2^2=\dim F\), Cauchy--Schwarz for \(\operatorname{tr}(p_F C)\) and \eqref{eq:overlap} give
\[
 \frac1{|\Gamma|}\sum_{g\in\Gamma}\|p_F\rho(g)P\|_2^2
 \le s_0\sqrt{\alpha(1+\epsilon)}
 \le \tfrac12(\alpha+1+\epsilon)s_0 .
\]
Choose \(g\) whose summand does not exceed this average. For \(U=(e\rho(g)P)^{-1}F\), the maps \(e\rho(g)P\) and \(p_F\rho(g)P\) agree on \(U\). The component perpendicular to \(Y\) is perpendicular to \(F\), and the component in \(Y\) already belongs to \(F\). Applying \eqref{eq:near} to this \(U\), commuting \(e\) and \(\rho(g)\), and using unitarity yields
\[
 \dim U-\epsilon s_0\ \le\ \|e\rho(g)Pp_U\|_2^2
       \ \le\ \|p_F\rho(g)P\|_2^2
       \ \le\ \tfrac12(\alpha+1+\epsilon)s_0 .
\]
The rank of \(e\rho(g)P\) modulo \(F\) is \(s_0-\dim U\). It is at least \(\tfrac12(1-\alpha-3\epsilon)s_0\), regardless of how the earlier translates were chosen.

\subsection{The logarithmic number of translates}
If \(1-\alpha>16\epsilon\), this gain is at least \(\tfrac25(1-\alpha)s_0\), since \(1/2-3/32>2/5\). Put \(c=2s_0/(5D)>0\). While the missing dimension exceeds \(16\epsilon D\), the gain reduces it by a factor at most \(1-c\). The first step from \(F=0\) also shows \(c\le1\), since a codimension is nonnegative. Induction and \(1-c\le e^{-c}\) bound the missing share after \(j\) steps by \(\max(16\epsilon,e^{-cj})\). Take \(j=\lceil c^{-1}\log(1/(16\epsilon))\rceil\). The set \(\mathcal T\) of translates chosen has
\[
 |\mathcal T|s_0\le\tfrac52D\bigl(1+\log(1/(16\epsilon))\bigr),\qquad
 \dim\sum_{g\in\mathcal T} e\rho(g)P(S)\ge(1-16\epsilon)D .
\]
The chosen translates are distinct while the dimension gain is positive. All previously used maps have image contained in the current \(F\). If the construction stops early, the same bound on the number of translates applies.\formalnote{6.2}{2}

\subsection{The two automorphic estimates}
Let \(d_\infty\) be the product of the local Bessel formal degrees and let \(\mu_r=\sum_{[L]}|\operatorname{Aut}L|^{-1}\) be the mass of the standard rank \(r\) lattice genus. Comparing the traces on the good subspaces for unitary groups differing in signature gives equality of their dimensions. On a definite group, evaluation at class representatives identifies the space of algebraic modular forms with the direct sum of the weight spaces fixed by the respective stabilizers \cite[Proposition~2.3 and Lemma~2.4]{upper-greenberg-voight}. At our free principal level these are copies of the whole weight space. The class count is \(|\Gamma|\mu_r\), since passing from the standard level multiplies the mass by its index \(|\Gamma|\), and at the smaller level all stabilizers are trivial. The dimension of the weight is \(d_\infty\) by the Weyl degree formula. Since the good forms are a subspace,
\[
                        D\le |\Gamma|\mu_r d_\infty .
\]
The signature comparison uses a cuspidal test at a split place in the trace formula. Its applicability to the good idempotent is one of the specialized facts here.\formalnote{6.3}{3}

Write \(P_g=P^*\rho(g)P\). The two scales for the second moment are \(\mathcal D_{\rm main}=s_0^2/(\mu_r d_\infty)\) and \(\mathcal D_{\rm id}=|\bar R_0|s_0\), where \(\bar R_0\) is the image of the Bessel subgroup in \(\Gamma\). Unfolding the Poincar\'e series expresses each matrix block of \(P_g\) as a sum over rational isometries \(\gamma\) with fixed integral reduction \(g\). A block joins two class representatives \(h_c,h_{c'}\). The element \(h_c^{-1}\gamma h_{c'}\) must be integral and have the prescribed reduction. Each amplitude is a Bessel integral with the prescribed finite and infinite place data. Divide the sum of Bessel integrals by the common unipotent covolume to obtain the block of \(P_g\). Summing the squared block norms gives \(\beta(g)=\|P_g\|_2^2\). For a subset \(A'\subset\Gamma\), put \(\beta(A')=\sum_{g\in A'}\beta(g)\). Its positive type follows from \(\beta(g^{-1}h)=\langle C_g,C_h\rangle_{\mathrm{HS}}\).

Fix a subset \(J\) of the chosen places above \(p_0\), with a nonzero test vector or covector at each \(w\in J\). Write \(H_J\) for the simultaneous stabilizer of these vectors and covectors. At a chosen place \(w\), write \(g_w\) for the component of \(g\in\Gamma\) and \((f_i)\) for the reduced frame there. We call \(g\) \emph{open} if, at some chosen place \(w\) above \(p_0\), all leading principal minors of order less than the frame length are nonzero in the matrix \((\langle f_i,g_w f_j\rangle)_{ij}\). If \(J\) omits a chosen place above \(p_0\), the open part of \(H_J\) has mass at most \((1+4b^{-4})\,b^{-6k|J|}\mathcal D_{\rm main}\), where the mass of a subset is its \(\beta\)-sum. In the expansion of \(\|P_g\|_2^2\), a pair of rational isometries occurs in each term. For one such isometry \(\gamma\), let \(A\) be the isotropic span of the first part of the frame. Take the middle coordinates of \(t\) and \(\gamma^{-1}t\) for every \(t\) perpendicular to both \(A\) and \(\gamma A\). The resulting linear relation is the \emph{link} of \(\gamma\) (\S\ref{sec:frame-lifts}). For pairs with the same link, the open terms are estimated by successively summing their common frame coordinates. Fixing \(J\) introduces the factor \(b^{-6k|J|}\) in the main class-sum estimate even before adding the tails. Bounds for pairs with different links and for pairs with a shared link depend on the lattice on their common subspace. Both the local integrals and the uniform count of those lattices are needed for this input.\formalnote{6.4}{4}

\subsection{A weighted determinant test over a finite field}
Let a finite group \(H\) act through \(\pi:H\longrightarrow\mathrm{GL}_n(\mathbb F_Q)\). Let \(A_0,A_1\) be \(n\)-by-\(h\) matrices. Suppose \(\theta:H\to\mathbb R_{\ge0}\) has positive type, written \(\theta(a^{-1}b)=\langle u_a,u_b\rangle\). For each nonzero \(v\in\mathbb F_Q^h\), assume the \(\theta\)-masses of the stabilizers of \(A_0v\) under \(\pi\) and \(A_1v\) under \(\pi^{-T}\) are at most \(Y_0\). Then, writing \(Z_g=A_1^T\pi(g)A_0\),
\[
 Q\sum_{\det Z_g=0}\theta(g)\ \le\
 2\sum_g\theta(g)+Q^{h+n/2}Y_0 .
\]
Choose a nontrivial additive character \(\psi\) of the field. In \(\ell^2(\mathbb F_Q^n)\) let \(e_z\) be the delta vector at \(z\), and \(c_x(z)=Q^{-n/2}\psi(-x\cdot z)\). Both families are orthonormal. For \(x=A_1y,\ z=A_0v\) with \(y,v\ne0\), form
\[
       X=\sum_{a\in H}c_{\pi(a)^{-T}x}\otimes u_a,\qquad
       Z=\sum_{b\in H}e_{\pi(b)z}\otimes u_b .
\]
Expanding \(\|X\|^2\) leaves only pairs for which \(a^{-1}b\) fixes \(x\) under the contragredient action. Hence \(\|X\|^2\le |H|Y_0\), and likewise \(\|Z\|^2\le |H|Y_0\). Expanding their inner product and grouping by \(g=a^{-1}b\) bounds \(\big|\sum_g\psi(y^TZ_gv)\theta(g)\big|\) by \(Q^{n/2}Y_0\). Now sum over \(y\) and average over \(v\in\mathbb F_Q^h\). Orthogonality makes the result \(\sum_g|\ker Z_g^T|\theta(g)\). The axes \(y=0\) and \(v=0\) contribute at most \(2\sum_g\theta(g)\). The remaining terms contribute at most \(Q^{h+n/2}Y_0\). A singular matrix has at least \(Q\) left-kernel vectors. This proves the inequality.\formalnote{6.5}{5}

Maximize \(b^{6k|J|}\beta(H_J)\) over \(J\) and its vector choices. For any unused place, adjoining a fixed-vector condition bounds its stabilizer mass inside \(H_J\) by \(b^{-6k}\beta(H_J)\). Apply the determinant test to the failure of each leading-minor condition. The residue-field and frame sizes give total nonopen mass at most \(b^{-4}\beta(H_J)\). For a full collection the separate terminal estimate is \(\beta(H_J)\le b^{-6k|J|}\mathcal D_{\rm main}\). Otherwise put \(z=b^{-4}\) and \(M_J=b^{-6k|J|}\mathcal D_{\rm main}\). Adding the open and nonopen estimates gives
\[
           \beta(H_J)\le (1+4z)M_J+z\beta(H_J).
\]
Since \(z\le1/10\), \((1-z)(1+9z)\ge1+4z\), giving \(\beta(H_J)\le(1+9z)M_J\). Comparison with the empty \(J\) proves \(\sum_g\beta(g)\le(1+9b^{-4})\mathcal D_{\rm main}\). Multiplication by \(D\) and the dimension bound cancel \(\mu_r d_\infty\) exactly. This proves \eqref{eq:overlap} for \(\epsilon\ge9b^{-4}\).\formalnote{6.6}{6}

\subsection{Counting the finite types that survive the projection}
At one principal-level prime take \(f\) fields of order \(Q\) and groups \(\mathrm{GL}_r(\mathbb F_Q)\), where \(1<r<Q\). At either prime \(p_i\) take \(Q=p_i\) and \(J=q_2=(qd)^{s+6}\), with \(s\) as in \S\ref{sec:data}. A local polynomial is allowed when it is monic of degree \(r\), squarefree, and has no root at zero or at a \(J\)-th root of unity in the field. A tuple is good when all its polynomials are allowed and at least one has a primitive root of \(\mathbb F_{Q^r}^{\times}\) as a zero and is irreducible of degree \(r\).

Begin with the \(Q^r\) monic polynomials. At most \((1+J)Q^{r-1}\) have a forbidden linear factor, by dividing by each such factor. A polynomial with a square factor has the form \(P^2B\) with \(P,B\) monic and \(e=\deg P\ge1\). For fixed \(e\) there are at most \(Q^e Q^{r-2e}\le Q^{r-1}\) choices. Summing gives at most \((1+J+r)Q^{r-1}\) exceptions in total. Give each polynomial weight \(v=(1-r/Q)Q^{-r}\). The allowed ones have total weight at least \(1-\varepsilon\), where \(\varepsilon=(1+J+2r)/Q\), by multiplying the surviving fraction by \(1-r/Q\).

Write \(n=Q^r-1\). The multiplicative group of the degree-\(r\) field has \(\varphi(n)\) generators. Each has degree \(r\), since it cannot lie in a field of order \(Q^e\) with \(e<r\). Its minimal polynomial accounts for at most \(r\) generators. If \(p_1<\cdots<p_t\) are the primes dividing \(n\), then \(p_i\ge i+1\), and
\[
 \frac{\varphi(n)}{n}=\prod_i(1-1/p_i)\ge\prod_{i=1}^t i/(i+1)
       =1/(t+1),\qquad t\le \log n/\log 2\le2\log n.
\]
As \(n\ge Q^r/2\), the primitive fraction among all monic polynomials is at least \(u=1/[2r(1+2r\log Q)]\). The nonprimitive polynomials thus have total weight at most \(1-u\le e^{-u}\).

For a squarefree label \(c\) with nonzero constant coefficient, the central idempotent of its irreducible character \(\Xi_c\) has coefficient \(\Xi_c(1)^2/|\mathrm{GL}_r(\mathbb F_Q)|\) at the identity. If the irreducible factors of \(c\) have degrees \(d_j\), Green's degree formula \cite[Theorem~14, pp.~443--444]{theta-green} gives \(\Xi_c(1)=\prod_{a=1}^r(Q^a-1)/\prod_j(Q^{d_j}-1)\). Dividing its square by the group order gives \(Q^{-r}\prod_{a=1}^r(1-Q^{-a})/\prod_j(1-Q^{-d_j})^2\ge Q^{-r}(1-r/Q)=v\). For the counting weight \(v^f\) per tuple, the total over all allowed tuples is at least \((1-\varepsilon)^f\), and the total over all tuples with no primitive place is at most \(e^{-fu}\). Subtraction gives a lower bound for the counting weight of the good tuples. Each such tuple contributes at least \(v^f\) to the actual retained coefficient, since tensor products multiply the local coefficients. For \(0\le\varepsilon\le1\), the coefficient discarded in the good-type projection is therefore at most
\[
 1-(1-\varepsilon)^f+e^{-fu}\ \le\ f\varepsilon+e^{-fu}.
\]
At either principal-level prime, \(b^{10}\le Q\le b^{11}\) by \S\ref{sec:data}. The numerator \(f(1+J+2r)\) grows at most polynomially in \(x\) for fixed \(q,d\), so \(f\varepsilon=o(b^{-8})\). Also \(f\ge x^{11}\) and \(r\le3x^5\), giving \(fu\gg x/\log b\). Since \(\log b\le x^\eta+1\) with \(\eta<1/100\), \(fu/\log b\to\infty\). Thus \(f\varepsilon+e^{-fu}\le b^{-8}/2\) at each prime for sufficiently large \(x\). This bounds the trace lost in the regular representation. Passing to an automorphic space requires the estimate below.\formalnote{6.7}{7}

\subsection{Projecting to good forms and restricting to the ramified source}
At the identity translate \(g=1\), the estimate for \(P^*\rho(g)P=P^*P\) in \S\ref{sec:bessel-projection} gives \(\|P^*P-1\|_2^2\le b^{-12}s_0\). For any \(U\),
\[
  \big|\|Pp_U\|_2^2-\dim U\big|=
  |\operatorname{tr}(p_U(P^*P-1))|\le b^{-6}\sqrt{s_0\dim U}\le b^{-6}s_0 .
\]
The two factors \(e_{p_0},e_{p_1}\) of the central idempotent \(e_0\) each discard a share at most \(b^{-8}/2\) in the regular representation, measured by their complement's coefficient at the identity. Passing from that share to \(P\) requires the Bessel trace estimate. It yields \(\|(1-e)P\|_2^2\le4b^{-4}s_0\). Orthogonality gives
\[
    \|ePp_U\|_2^2=\|Pp_U\|_2^2-\|(1-e)Pp_U\|_2^2
                 \ge \dim U-5b^{-4}s_0 ,
\]
which proves \eqref{eq:near} for \(\epsilon\ge5b^{-4}\). Section~\ref{sec:bessel-projection} derives the identity-block bound and the Bessel trace bound from the restricted orbital second moment.\formalnote{6.8}{8}

Take \(\epsilon=1/(8b^2)\), which exceeds \(5b^{-4}\) and \(9b^{-4}\) for our data. Applying \eqref{eq:near} to \(U=S\) gives \(D>0\). The general lemma gives \(16\epsilon=2/b^2\) and \(|\mathcal T|s_0\le5D\log(1+2b^2)\), using \(1\le\log(1+2b^2)\) and \(\log(b^2/2)\le\log(1+2b^2)\). On the ramified source space \(S_{\rm ram}\), each map \(e\rho(g)P\) has the same image as the translate \(g\mathfrak p\) in Theorem A. Here \(\mathfrak p\) is the Poincar\'e map to \(Y\). The normalizing scalars of the two maps are nonzero. The mean-zero conditions for the scalings at the primes above \(N\) have total codimension at most \(fs_0/(N-1)\le b^{-8}s_0\). The image of a complement to \(S_{\rm ram}\) accounts for at most this many directions per translate. The total loss is at most \(5b^{-8}D\log(1+2b^2)\le 2D/b^2\). Thus\formalnote{6.9}{9}
\[
  \dim\sum_{g\in\mathcal T}g\mathfrak p(S_{\rm ram})\ge(1-4/b^2)D,\qquad
  |\mathcal T|\dim S_{\rm ram}\le5D\log(1+2b^2).
\]


\section{The orbital sieve in the second moment}
\label{sec:orbital-sieve}

We expand the open-cell estimate in Section~\ref{sec:B}. Write
\(n=\rho+1\), where \(\rho\) is the number of hyperbolic planes,
\(k\) for the rank of the definite middle block, and \(f=[E:\mathbb Q]/2\).
Put \(h=|J|\), the number of chosen places above \(p_0\) at which a
nonzero vector or covector is fixed. The scale for the sum before division by \(V_U^2\) is
\[
      B_{\rm open}=b^{-6kh}\mathcal D_{\rm main}V_U^2,
\]
where \(V_U\) is the unipotent covolume. Expanding the square gives two lists of frame coordinates. Holding their initial entries fixed, we sum over the possible next common entry. Poisson summation separates this sum into the integral of the kernel (frequency zero) and an error from nonzero frequencies. Decay with exponent proportional to \(k\) makes the accumulated error at most the allowed fraction of the main term.

\subsection{Summing one common entry}

Let \(T\) be a finite conjugation-stable set of finite places containing
the denominators of the data under consideration. Write \(E_T=\prod_{w\in
T}E_w\) and \(X=\mathbb C^f\times E_T\). A nonzero rational entry
\(t\) has coordinates
\(x_a(t)=((r_{va}\Phi_v(t))_v,(t)_T)\) in position \(a\), where
\[
 u_v=\frac{M\Phi_v(\varpi)}{2\pi i},\qquad
 r_{va}=\frac{|u_v|^{2(\rho-a)}}{N}\quad(0\le a\le\rho),\qquad
 c_a=\frac{\operatorname{vol}(\widehat{\mathcal O}_E)}
                {\prod_v r_{va}^2}.
\]
Here \(\varpi\) generates the different, and Haar measures have the
adelic normalization of the source. The square in the denominator
is the real Jacobian of dilation of a complex coordinate.

Let \(A\subset E^\times\) be a finite set of possible entries, with
\(A,A^{-1},M^{-1}\) integral off \(T\). For \(0\le i\le n\) let
\(J_i(p,q)\) be the product of the normalized local orbital kernels
with \(i\) entries fixed in each of the two lists \(p,q\). These initial lists are called the \emph{heads}, and the stage with \(i\) entries fixed is called floor \(i\). Write \((p,x)\) and \((q,x)\) for the heads obtained by appending the same local entry \(x\). The local
integration formula gives
\[
  \int_X J_{i+1}((p,x),(q,x))\,dx=J_i(p,q)\qquad(i<n).
  \tag{B1}
\]
The transform in \(x\) at a nonzero rational dual frequency \(\xi\)
has absolute value at most
\(\beta(\xi)\Lambda_i(p)\Lambda_i(q)\), where the head weights
\(\Lambda_i\) are nonnegative. For different last rational entries
\(t,s\), the same bound applies directly to \(J_i\), with weights
of length \(i\) and with frequency coming from
\((t^{-1}-s^{-1})/M\). These are two separate uses of decay.

For each \(x\in X\) set
\(F_i(x)=\sum_{p,q\in A^i}J_{i+1}((p,x),(q,x))\), evaluating all
rational entries at their prescribed scales. This is real and
nonnegative. On the top floor the matrix with common last entry is
positive semidefinite, so the sum is its quadratic form at the vector
of all ones. On lower floors the claim follows by (B1), integrating
the same quadratic form with further common coordinates. Products of
local positive semidefinite matrices preserve positivity, since their
entries are inner products of tensor products. Continuity of the top
form at nonzero entries and Fatou's lemma on the successive integrals
give lower semicontinuity at each \(x_i(t)\).

The classical adelic Poisson formula identifies the sum over \(E\) with the dual sum for Schwartz--Bruhat functions \cite[Lemma~4.2.4]{tate1950}. The following one-sided form uses lower semicontinuity only at the finitely many entries being counted. If \(F\ge0\) is integrable and
lower semicontinuous at these finitely many points, then
\[
 \sum_{t\in A} F(x_i(t))
 \ \le\ c_i\left(\int_X F+\sum_{0\ne y\in\mathcal O_{E,T}}
                              |\widehat F(\xi_i(y))|\right).
 \tag{B2}
\]
Our transform uses the product of the characters
\(\exp(z-\bar z)\) at infinity and the self-dual character at \(T\)
normalized by \(-\varpi^{-1}\). With these conventions
\(\xi_i(y)\) has components
\(2\pi i\Phi_v(y)/(\Phi_v(\varpi)r_{vi})\) and \(y\) at \(T\).
To see (B2), extend the rescaled \(F\) to the adeles by multiplying
by the indicator of integral components off \(T\). Its Fourier
coefficient at \(y/\varpi\) vanishes unless \(y\) is integral off
\(T\), by invariance under those integral translations. At the
remaining frequencies it equals \(c_i\widehat F(\xi_i(y))\).
Periodize the extension on \(\mathbb A_E/E\). If the displayed
absolute sum is finite, the periodization agrees almost everywhere
with its continuous, absolutely convergent Fourier series. Convolve
with a nonnegative approximate identity of integral one at the
origin. Each of the finitely many translates belonging to \(A\)
contributes in the limit inferior at least its value there by lower
semicontinuity. The other translates contribute nonnegative terms.
The Fourier series is bounded by the sum of the absolute values
of its coefficients. This proves (B2), including the inequality for
a subset of rational points.\formalnote{6.10}{10}

Define \(D_0=\Re J_0\) and
\(D_{i+1}=\sum_{t\in A}F_i(x_i(t))\). Put
\(W_i=\sum_{p\in A^i}\Lambda_i(p)\). Suppose the decay weights
sum to at most \(\epsilon\) over nonzero rational frequencies at
the indicated scales, including the scales for inverse differences.
Apply (B2) to \(F_i\). The nonzero frequencies cost at most
\(\epsilon W_i^2\), by the triangle inequality and expansion of
\(W_i^2\). The integral in (B2) equals
\(\Re\sum_{p,q\in A^i}J_i(p,q)\) by (B1). The portion with equal
last entries equals \(D_i\). Each of the other terms has a nonzero
frequency, because inversion on \(E^\times\) is injective. Even the
singleton frequency bound is at most \(\epsilon\), so this portion
costs at most another \(\epsilon W_i^2\). When \(i=0\) there is
only the empty pair. Consequently
\[
 D_{i+1}\le c_i(D_i+2\epsilon W_i^2),\qquad
 D_n\le\prod_{a<n}c_a\,\bigl(B_0+2n\Gamma\epsilon\bigr)
 \tag{B3}
\]
provided \(D_0\le B_0\) and
\(W_i^2\le\Gamma\prod_{a<i}c_a\). The second inequality follows by
dividing the first by \(\prod_{a\le i}c_a\) and summing. Positivity
allows the final common entry to range over any subset of \(A\).

\subsection{Summing the frequencies and counting the heads}

The local estimates feeding (B3) have specific exponents. Put
\(K_*=k-n^7\), \(D_K(z)=z^K\min(1,z)^6\) for \(z>0\), and
\(V_*=\pi^{(2\rho+1)k}\operatorname{vol}U(k)/\operatorname{vol}U(r)\).
At infinity the bound for a local transform or for two unequal
last entries contains \(V_*b^{2000n^9}D_{K_*}(|\xi_v|^{-1})\).
At \(T\) the factor is
\[
 b^{n^9 f}\frac{[\mathcal O_T:\xi_T\mathcal O_T]_{\rm gen}^{K_*/2}}
 {[\mathcal O_T:\mathcal O_T\cap\xi_T^{-1}\mathcal O_T]^3}.
 \tag{B4}
\]
Here the generalized index is the ratio of the two indices over
the intersection. For the archimedean bottom mass, split the integral
of the squared cut Bessel pairing into a position ball and its
complement. The matrix-coefficient estimate and the cutoff error on
the ball give \((1+b^{-12})^2/\deg_v\). The shell estimate outside gives
\(b^{-12}/\deg_v\). Their sum is at most \((1+b^{-10})/\deg_v\)
for these data. The underlying oscillatory shell bounds constitute
local inputs to the following summation. Sections~\ref{sec:local-orbital} and \ref{sec:shell-measure} establish them at infinity. Section~\ref{sec:bessel-projection} establishes the finite shell, denominator and frequency bounds.\formalnote{6.11}{11}

Here is why summing the bound over frequencies preserves its large
power. For \(y\ne0\) integral off \(T\), let \(\mathfrak d_y\)
be its denominator ideal. The finite part of the product formula bounds
the square root of the generalized index by \(\prod_v|\Phi_v(y)|\),
and the denominator index is at least \(\mathrm N\mathfrak d_y\).
At a scaled frequency \(c'_v\Phi_v(y)\), suppose
\(|c'_v|^{-1}\le x_v\le B\) with \(B\ge1\). Then
\[
 D_{K_*}(|c'_v\Phi_v(y)|^{-1})|\Phi_v(y)|^{K_*}
 \le x_v^{K_*}B^6\max(1,|\Phi_v(y)|)^{-6}.
\]
Indeed the powers of \(|\Phi_v(y)|\) cancel exactly. The displayed
bound on the remaining minimum follows by separating
\(|\Phi_v(y)|\le1\) and \(|\Phi_v(y)|\ge1\).

The sum of the residual height weight
\((\mathrm N\mathfrak d_y)^{-3}\prod_v\max(1,|\Phi_v(y)|)^{-6}\)
is at most \((64\cdot38e^4)^f\). For a fixed denominator
\(\mathfrak d\), the eligible elements are in the
\(\mathrm N\mathfrak d\) cosets of
\(\mathfrak d^{-1}/\mathcal O_E\). In one coset there is at most
one element in any product of half-open squares of side \(1/2\), since
two such elements would differ by a nonzero algebraic integer with
all conjugates of modulus less than one. The weight at a point is
at most \(2^6\) times the weight at the corner of its square.
The corner sum is at most \(38\). Indeed the half-unit lattice points
within radius \(R\ge1\) number at most \((1+4R)^2\le25R^2\)
by disjoint disks of radius \(1/4\). Integrating this count against
\(6R^{-7}dR\) on \([1,\infty)\) gives at most \(37.5\).
Summing now over ideals costs
\(\zeta_E(2)\le\zeta(2)^{2f}\le e^{4f}\).
Both Poisson and inverse-difference scales permit
\(x_v=|u_v|^{2\rho+1}/(MN)\). Using \(1\le|u_v|\le b^{23}\)
in the factors \(B^6\), one obtains\formalnote{6.12}{12}
\[
 \epsilon=b^{2002n^9 f}\Delta,\qquad
 \Delta=\prod_v V_*\left(\frac{|u_v|^{2\rho+1}}{MN}\right)^{K_*}.
 \tag{B5}
\]

The finite weight vanishes unless every successive product of
the head entries is integral at \(T\). It is already integral
elsewhere. The archimedean cutoffs and the inverse dilations bound
every conjugate of these algebraic integers by \(R=b^{120n^2}\).
The successive products determine the head by division. A packing
argument in real dimension \(2f\) bounds the possible lists by
\((1+2R)^{2fi}\). Images of distinct integers are at distance
at least one in the maximum of the \(f\) complex moduli, by their
nonzero integral norm. The balls of radius \(1/2\) in that norm are
disjoint. The product formula applied to each successive product
also bounds its weight by \(b^{120n^8f}\), and hence gives
\(W_i\le b^{130n^9f}\) for \(i\le\rho\). The frame scalings give
\((\prod_{a<i}c_a)^{-1}\le b^{100n^2f}\), so we may take
\(\Gamma=b^{400n^9 f}\). Substitution into (B3) costs at most
\(b^{2403n^9f}\Delta\) beyond the bottom mass.\formalnote{6.13}{13}

\subsection{Saving from each fixed vector}

At the pair of conjugate places belonging to a specified vector,
let \(K\) be the integral unitary group, \(K_M\) its congruence
subgroup and \(Z\) the inverse image of the vector stabilizer.
Let \(\mu\) be Haar measure and set
\[
 I(Z)=\int\left|\int_Z \chi_M(yxy^{-1})\,d\mu(y)\right|^2d\mu(x),
\]
where \(\chi_M\) is the extended Bessel phase. The residue general
linear group acts transitively on nonzero vectors and on nonzero
covectors. Thus \([K:Z]=p_0^r-1\). Write \(B=b^9\) and choose
\(t\) with \(6k\le9t,\ t+1<r,\ t+1+n+2\cdot10^9n^4\le k\),
as the rank hypotheses allow. Since \(B\le p_0\),
\(\mu(Z)\le B^{-(t+1)}\mu(K)\).

To obtain the required relative mass bound, separate the outer
integral over \(K_M\) and its complement. On \(K_M\) the squared
integral contributes exactly \(\mu(Z)^2\mu(K_M)\). Outside \(K_M\)
the change of frame has a vector of depth at least \(B\), because
an integral element with nonintegral inverse has modulus at most
\(B^{-1}\). In the orbital majorant this gives a power saving.
Precisely, the shell-volume estimate bounds the integral of the
majorant outside \(K_M\) by
\(B^{-m}\mathfrak m_Z2^{L}\) for \(m+10^8n^4\le k\), with
\(L=10^9n^4\) and \(\mathfrak m_Z=\mu(Z)\mu(K_M)/V_{U,T}^2\),
where \(V_{U,T}\) is the local unipotent mass. Taking \(m=t+1+n+L\)
and using \(2\le B\) absorbs \(2^L\). The squared phase integral is
bounded on this complement by \(V_{U,T}^2\) times the majorant,
so the local unipotent mass cancels. Therefore
\begin{align*}
 I(Z)&\le \mu(Z)\mu(K_M)\bigl(\mu(Z)+B^{-(t+1)}\mu(K)\bigr),\\
 I(K)&\ge \mu(K)^2\mu(K_M).
\end{align*}
For the first inequality the shell error before simplification is
\(B^{-(t+1+n)}\). The lower volume estimate
\(\mu(K)\ge2^{-n}\ge B^{-n}\) absorbs it. For the second the phase
on \(K_M\) has no oscillatory loss in the inner integral. Since
\(2/B\le1\), multiplying the first bound by \(B^t/\mu(Z)\)
and using the index gives \(B^t I(Z)/\mu(Z)\le I(K)/\mu(K)\).
Decompose \(T\) into the \(h\) disjoint conjugate pairs and its
complement. The group, phase and mass integrate as products on these
blocks. Consequently the gain is \(b^{-6kh}\) in normalized mass.
At the full integral group the bottom mass is at most \(1+b^{-6}\)
in the normalization by \(\mu(Z)\mu(K_M)\) and the unipotent
volume. Combining all places gives\formalnote{6.14}{14}
\[
 B_0=\frac{1+b^{-6}}{b^{6kh}}\,
                      \prod_v\frac{1+b^{-10}}{\deg_v}.
 \tag{B6}
\]

\subsection{From the class sum to the density on the last floor}

For an open block the nonzero leading minors allow successive
elimination in the two unipotent groups. For a fixed link this
writes the class sum as \(\sum_{p\in A^\rho} a_{p,i}\), where
\(i\) indexes the block, the entries in \(p\) are Weyl parameters,
and the last entry \(t_i\) is common. The cut terms under consideration
have \(|\Phi_v(t_i)|>N\) for every \(v\). The definite coordinate
\(\omega\) of the last frame vector satisfies
\[
 \omega^*\omega=\alpha,\qquad \alpha=(t_i\bar t_i-N^2)/N.
\]
In particular \(\alpha\) is totally positive. This is a one-vector
representation problem. The integral at each archimedean place of
the radial density is the derivative of the volume of a complex
\(k\)-ball, \(\pi^k\alpha_v^{k-1}/(k-1)!\).

For clarity, apply the one-vector Siegel formula to a nonnegative
locally constant compactly supported function \(\Psi\) on the finite
vectors. It includes the class representative, the shift by the last
frame vector and the indicator of integrality off \(T\). With
\(\Psi_T\) denoting the resulting test at \(T\), the finite density
factors into its norm-fiber integral
\(\int_{\omega^*\omega=\alpha}\Psi_T(\omega)d\nu_\alpha(\omega)\)
and a factor off \(T\). This follows by writing the norm density as
the Fourier integral for the equation \(\omega^*\omega=\alpha\)
and splitting the variables according to \(T\). Outside \(T\) the
test is the lattice indicator. The integrated absolute Gauss-sum
bound there is
\(\operatorname{vol}(\widehat{\mathcal O}_E)^k\zeta_E(3\rho+4)\)
in the integral-adele ideal-sum normalization, and is at most
\(\operatorname{vol}(\widehat{\mathcal O}_E)^k(1+b^{-6})\).
The admissible degree bound \([E:\mathbb Q]\le2^{2\rho}\) and
\(2(3\rho+5)\le k\) are used here. The genus average, rather than
a formula for an individual lattice, gives the required upper
bound for the selected class points by positivity.

At fixed \(t\) take as coefficients the Gram matrix of the
archimedean orbital-integral vectors indexed by \(p\). At a
finite norm-fiber point the finite pairing of the corresponding
vectors produces a nonnegative quadratic form after summing over
\(p,q\). A compactly supported step function agrees there with
the positive part of its real value. This makes it a valid Siegel
test. The head-density factor in each amplitude times its
archimedean normalization is
\(r_0(\prod_v(|\Phi_v(t)|^2-N^2))^{(k-1)/2}\), where \(r_0>0\)
is the common scalar of the opening chart. On squaring, the norm
power from the Siegel formula cancels. Group by \(t\), applying
the formula once in each fiber, to obtain
\[
       r_0^2\sum_i\left\|\sum_p a_{p,i}\right\|^2
                         \le (1+b^{-6})C_S D_n .       \tag{B7}
\]
Here \(C_S\) is the scaled Siegel constant including the class
count, the archimedean sphere constants, and the chart Jacobians.
The outer indexing by class pair, reduction and representative
is injective, so no extra multiplicity enters (B7).\formalnote{6.15}{15}

Write \(w=(\prod_{a<n}c_a)/r_0^2\).
The volume formulas give
\[
 C_Sw\le(1+b^{-9})\mathcal D_{\rm main}d_\infty V_U^2,\qquad
 C_Sw\Delta\le b^{n^9f}\,\#I\,\mathcal D_{\rm id}V_U^2,
\]
where \(I\) is the set of Bessel coordinates and \(\mathcal D_{\rm id}=|\bar R_0|s_0\) is the identity second moment. Admissibility gives
\(\#I\le b^{2kf}\), \(b^{11kf}\mathcal D_{\rm id}\le\mathcal D_{\rm main}\)
and \(512000n^9\le k\). Since \(h\le f\),
\[
 2404n^9f+2kf+10+6kh\le11kf.
\]
Thus the error after (B5) contributes at most \(b^{-10}B_{\rm open}\).
For the mass, \(d_\infty=\prod_v\deg_v\) and
\((1+b^{-10})^f\le1+b^{-9}\). The coefficient of \(B_{\rm open}\) in (B7),
after division by \(r_0^2\), is therefore at most
\[
 (1+b^{-6})\big((1+b^{-9})^2(1+b^{-6})+b^{-10}\big)
 \le 1+b^{-4}\quad(b\ge1000).
\]
This proves the main-link bound, prior to adding the other pairs.\formalnote{6.16}{16}

\subsection{The discriminant sum for the remaining pairs}

Write \(\tau=b^{kf/4}\). Let \(d_L\) and \(\ell\) be the discriminant
and corank of the common lattice of a near pair. Put
\(P_*=(\#\text{labels})(\#\text{classes})/N^{kf}\).
For each \(0\le\ell\le n^3\) the count takes the weighted form
\[
    \sum_{\text{eligible pairs of corank }\ell}d_L^{-k}
                          \le \tau^{\ell+2}P_* .          \tag{B8}
\]
Eligible pairs here have \(d_L\le b^{n^3f}\) and lie over an agreeing
block. The proof first assigns a common lattice and a reduction
label. At each such choice the number of pairs is at most \(\tau\).
Both terms are bounded by the near-term count over that lattice,
and the square of that count fits in \(\tau\). The genus-weighted
sum of \(d_L^{-k}\) over lattices of corank \(\ell\) is at most
\((\#\text{classes})N^{-kf}\tau^{\ell+1}\). Multiplying these
bounds gives (B8).

To locate the factor \(N^{-kf}\), write \(s=k-\ell\) and
\(H=b^{n^3f}\). Above a unimodular lattice the finer classes vary
the distinguished row modulo \(\bar\pi_N\). There are at least
\((N^k/b)^f\) cosets. A lattice \(L\) in the class occurs only
when the row annihilates \(L\bmod\bar\pi_N\). Let \(j\) be its index
in the saturation of its rational span, and \(d\ge1\) the
discriminant of that saturation. Then \(d_L=dj\). The saturation is
a direct summand locally. Its annihilator modulo \(\bar\pi_N\) thus
has at most \((N^f)^\ell\) row vectors. Passing to \(L\) enlarges
this by at most \(j\): in a local diagonal basis each elementary
divisor contributes at most its index to the number of additional
solutions, and these indices multiply to \(j\).

Put \(\mu=2\rho+4\). The condition \(dj\le H\) and \(s\le k\)
imply
\[
 j(dj)^{-k}\le H^{\mu+2}d^{-(k+1)}j^{-(s+\mu)} .
\]
Indeed replacing \(H\) by \(dj\) makes the right side divided
by the left \(d^{\mu+1}j^{k-s+1}\ge1\). For a fixed saturation,
the sum of \(j^{-(s+\mu)}\) is bounded by the \(s\)-th power of the
ideal zeta sum at \(\mu\), hence by \(2^k\). The subspace-discriminant
estimate in the corank calculation bounds the sum of \(d^{-(k+1)}\)
by \((200fk\,\mathrm{rd}_E)^{f(k+1)\ell}\), where \(\mathrm{rd}_E=|\operatorname{disc}(E)|^{1/[E:\mathbb Q]}\) is the root discriminant. Thus the row sum above a
given unimodular lattice is at most
\[
 (N^f)^\ell H^{2\rho+6}(200fk\,\mathrm{rd}_E)^{f(k+1)\ell}2^k .
\]
Admissibility bounds this expression by \(\tau^{\ell+1}/b^f\).
Dividing by the number of row cosets now gives
\(\tau^{\ell+1}/N^{kf}\) per class on average, as asserted.

Let \(G,G'\) be the two positive level gauges and let \(D\) bound
both cell second moments. Each individual orbit size is bounded by
\(\tau^2\sqrt{D/(P_*G^k)}\) with its corresponding gauge.
If \((d_L\eta)^2\le GG'\), multiplication bounds the product by
\(\tau^4 D\, d_L^{-k}/(P_*\eta^k)\). By (B8), any bound
\(\tau^{\ell+6}D/\eta^k\le B\) supplies at most \(B\) for this
corank. There are \(n^3+1\le\tau\) possible coranks, so the sum
is at most \(\tau B\). For different links or vanishing term
discriminant the congruence and lower-cell estimates give
\(B=\mathcal D_{\rm id}/\tau^{10}\), and
\(\tau B\le b^{-kf}\mathcal D_{\rm id}\). For equal nondegenerate
links in blocks meeting at places of total norm at least \(z>0\),
the corank is at most one and \(\eta=z\) is admissible. This gives
\(\tau^8\mathcal D_{\rm main}/z^k
=(b^{2f}/z)^k\mathcal D_{\rm main}\).

After the cutoff and remaining-pair estimates, the square root of
the open second moment on the prescribed fixer is bounded by
\(\sqrt{(1+b^{-4})b^{-6kh}\mathcal D_{\rm main}}+
b^{-15}\sqrt{\mathcal D_{\rm id}}\) in the nonterminal case.
Using \(b^{11kf}\mathcal D_{\rm id}\le\mathcal D_{\rm main}\)
to absorb the second term, and \(\sqrt{1+x}\le1+x/2\), gives
the coefficient \((1+b^{-4})^2\le1+4b^{-4}\) in Section~\ref{sec:B}.\formalnote{6.17}{17}

Section~\ref{sec:bessel-period} derives the Bessel majorant, and
\S\ref{sec:lattice-heights} proves the weighted genus estimate in
(B8) and accounts for the successive frame-count losses. The local
Jacobian and one-vector Siegel formulas, and the frame congruence
counts specified there, enter the argument at the points indicated.
The Poisson inequality, the two errors at
each step of descent, the degree of frequency summability, the
per-vector mass gain and the final discriminant-weight summation
follow from those inputs by the displayed arguments.



\section{Pairs of frames and their lattice determinants}
\label{sec:cell-reduction}
The second moment uses the same integer \(j\) to choose the rows of a product matrix, the size of a minor, and the integral equations for a frame lift. We derive that choice from the character in the Bessel period. We then calculate the determinant of the pair of lattices in (FL2) and construct the box in which the lifts are counted.

\subsection{Detecting an intersection by two minors}
Decompose the Hermitian space as \(A\oplus(X_0\oplus Ez)\oplus A^*\), with paired isotropic spaces \(A,A^*\) of dimension \(\rho\), a definite middle space \(X_0\) of rank \(k\), and \(z\) perpendicular to \(X_0\) with \([z,z]=-N\). Write \([x,y]=x^*Jy\) for the pairing, \(n=\rho+1\), \(F=(f_0,\ldots,f_\rho)=(e_0,\ldots,e_{\rho-1},z)\), and \(F'=\gamma F\). Put \(Z=F^*JF'\), \(\delta_i=\det Z_{<i,<i}\) for \(0\le i\le n\), and \(\delta_0=1\). The quantity used to detect an intersection is the element of the real subfield
\[
  g=\delta_n\bar\delta_n-N^2\delta_\rho\bar\delta_\rho.
\]
Suppose \(Zy=D e_\rho\), where here \(e_\rho\) denotes the last \emph{coordinate} vector in \(E^n\). The first \(\rho\) products say that \(F'y\) has zero \(A^*\)-coordinate. Its product with \(z\) says that its \(z\)-coordinate is \(-D/N\). Its length for the rational form is \([Fy,Fy]=-N\bar y_\rho y_\rho\) by isometry. Consequently, at each embedding \(v\),
\[
 N\| (F'y)_{X_0}\|_v^2=|D|_v^2-N^2|y_\rho|_v^2.                \tag{CR1}
\]
For the last column \(c\) of \(\operatorname{adj}(Z)\), we have \(Zc=\delta_n e_\rho\) and \(c_\rho=\delta_\rho\). Hence \(g\) is totally nonnegative. If \(\delta_n=0\), then \(g=0\) by (CR1). Choose a nonzero vector \(y\) of \(\ker Z\). The same identity, now with \(D=0\), gives \(y_\rho=0\) and \((F'y)_{X_0}=0\). Thus \(F'y\ne0\) lies in the first frame. If \(\delta_n\ne0\), then \(c\ne0\), and \(g=0\) forces \((F'c)_{X_0}=0\), again giving a shared vector. Conversely a common nonzero vector \(Fx=F'y\) satisfies \(Zy=-Nx_\rho e_\rho\). In the invertible case \(y\) is a nonzero multiple of \(c\), so the vanishing of the \(X_0\)-coordinate forces \(g=0\). This proves that the two frame spans meet precisely when \(g=0\), including the singular case where the cofactor column could vanish.\formalnote{6.18}{18}

\subsection{Selection of the leading minor by character agreement}
Let \(U\) be the rational Bessel unipotent. Say that \(\gamma\) agrees if \(\chi(u)=\chi(u')\) whenever \(u,u'\in U\) and \(u\gamma=\gamma u'\). In the unfolded kernel a translation by such a pair multiplies the stabilizer quotient integral by \(\chi(u)^{-1}\chi(u')\). A factor different from one forces the integral to vanish. Unfolding by the class unipotents on each side applies this observation to a stable collection of terms. We can therefore classify the agreeing terms.\formalnote{6.19}{19}

A pair \((\mu,\nu)\) occurs at \(a\) if \(Fx=F'y\) for coordinates supported in \(0,\ldots,a\) with \(x_a=\mu,\ y_a=\nu\). We prove that it then occurs at \(a+1\) when \(a<\rho\). Let \(m\) be any linear form killing \(f_i,f'_i\) for \(i\le a\). The assertion amounts to
\[
                     \mu m(f_{a+1})=\nu m(f'_{a+1}).              \tag{CR2}
\]
Indeed, if this holds for every such \(m\), the difference \(\mu f_{a+1}-\nu f'_{a+1}\) belongs to the span of the earlier columns on both sides together. Writing it in these columns supplies the new \(x,y\).

Represent \(m\) as \([w,\cdot]\), put \(v=Fx=F'y\), and choose any \(r\in E\). The vector \(h=rv\) is isotropic and \([w,h]=0\). For \(s+\bar s=[w,w]\), the map
\[
 T(t)=t+h[w,t]-(w+sh)[h,t]
\]
preserves the form. In the expansion of \([T(t),T(t')]\) the terms in the product of the two \(h\)-pairings cancel by \(s+\bar s=[w,w]\), and the cross terms cancel in conjugate pairs. The form \([h,\cdot]\) kills both entire frames. On the first frame \(T\) therefore changes \(f_i\) by \(rv\,m(f_i)\), zero for \(i\le a\) and in the span through \(a\) for \(i>a\). The same formula in the primed frame gives the same triangular action there. In the ambient Witt coordinates, the two terms with \([h,\cdot]\) have only dual-isotropic input, and the term \(h[w,\cdot]\) has only isotropic output. Together with the support through \(a\) and the vanishing of \(m\) through \(a\), these are the block conditions for \(T\in U\) and \(\gamma^{-1}T\gamma\in U\). The sum of the immediately superdiagonal entries on the first frame is \(r\sum_{i<\rho}x_i m(f_{i+1})=r\mu m(f_{a+1})\), and the primed calculation gives \(r\nu m(f'_{a+1})\). Taking \(r=-M\delta_E t\), where \(\delta_E\) normalizes the character by the different, agreement gives the same finite-adelic additive character on \(t\mu m(f_{a+1})\) and \(t\nu m(f'_{a+1})\) for every \(t\in E\). Density of \(E\) in its finite adeles and nondegeneracy of this character pairing prove (CR2).\formalnote{6.20}{20}

At the last index, a common vector with coefficients \(\mu,\nu\) has norm both \(-N\bar\mu\mu\) and \(-N\bar\nu\nu\). By (CR2) the equality of these two scalars holds for a pair occurring at \emph{any} index. Thus at the highest occupied index of a nonzero relation both coefficients are nonzero, and we can rescale them to \((\zeta,1)\) with \(\zeta\bar\zeta=1\).

Let \(j\) be the least integer from which a single pair \((\zeta,1)\) occurs at every index through \(\rho\), allowing \(j=n\) for the empty interval. If \(g\ne0\), the frame spans are disjoint and \(j=n\), with \(\delta_n\ne0\). If \(g=0\), a shared vector and (CR2) give \(j\le\rho\). Every nonzero common vector then has its highest occupied index at least \(j\), with a nonzero coefficient there in both frames. The witnesses for \(b=j,\ldots,\rho\) have successively larger highest coordinate and nonzero diagonal coefficients on both sides. Back substitution constructs a common vector with any prescribed tail of coordinates from \(j\) onward in either one of the two frames.\formalnote{6.21}{21}

We claim \(\delta_j\ne0\). For \(j\le\rho\), singularity would give \(0\ne y\) supported below \(j\) with the first \(j\) entries of \(Zy\) zero. For every common vector \(Fx=F'y'\) one has
\[
          x^* Zy=[F'y',F'y]=-N\bar y'_\rho y_\rho=0.
\]
The tail of \(x\) is arbitrary, so the remaining entries of \(Zy\) also vanish. By (CR1), \(F'y\) is a nonzero common vector. Its second list of coefficients has no occupied index \(\ge j\), a contradiction. The bound \(\delta_j\ne0\) holds in the original ordering of the flags.

Write \(I=A\cap\gamma A\) and \(A_{<j}=\langle f_0,\ldots,f_{j-1}\rangle\) for \(j\le\rho\). The witnesses at indices \(j,\ldots,\rho-1\) lie in \(I\). Their tails span all isotropic tail coordinates, and a vector in \(I\) supported below \(j\) in either frame must vanish by the highest-index argument. Hence
\[
        A=A_{<j}\oplus I,\qquad
        \gamma A=\gamma A_{<j}\oplus I,\qquad \dim_E I=\rho-j.     \tag{CR3}
\]
This also gives the intrinsic rank formula \(j=\rho-\dim_E I\) on the zero-discriminant cells. For the open cell one has \(I=0\) and \(g\ne0\). Thus the discriminant distinguishes the two possibilities when \(I=0\). For completeness, another index \(i\le\rho\) satisfying the tail and pivot conditions must equal \(j\). Necessarily \(i\ge j\) by minimality. If \(i>j\), a witness at \(i-1\) is isotropic with both coordinate lists supported below \(i\). Its primed coordinates give a nonzero kernel vector of \(Z_{<i,<i}\), since its unprimed vector lies in \(A\), contradicting \(\delta_i\ne0\).\formalnote{6.22}{22}

\subsection{The last-vector minor and the complementary lattice}
Let \(a=\min(j,\rho)\) and \(S_j=\{0,\ldots,a-1\}\cup\{\rho\}\). For \(j\le\rho\), form \(Y\) by taking the first \(j\) rows of \(Z\) and the remaining rows of the identity. Put \(w=Y^{-1}e_\rho\). The tail prescription above gives a common vector \(Fu=F'y\) with \(y_b=(e_\rho)_b\) for \(b\ge j\). Since \(Zy=-Nu_\rho e_\rho\), this vector satisfies \(Yy=e_\rho\), hence \(y=w\). The last coefficient of \(w\) is one, so \(u_\rho\bar u_\rho=1\). Replacing the last row of \(Y\) by the last row of \(Z\) multiplies its determinant by \(Z_{\rho,*}Y^{-1}e_\rho=-Nu_\rho\), by expansion in cofactors. Expanding along the remaining identity rows gives
\[
 \det Z_{S_j,S_j}=-Nu_\rho\,\delta_j \quad(j\le\rho).
\]
For \(j=n\), \(S_j\) is the full index set. With \(f=[E:\mathbb Q]/2\) and \(\mu_\gamma=|\operatorname{Nm}_{E/\mathbb Q}\delta_j|/N^{2f{\bf1}_{j=n}}\), both cases give
\[
            |\operatorname{Nm}_{E/\mathbb Q}\det Z_{S_j,S_j}|
                              =N^{2f}\mu_\gamma.                 \tag{CR4}
\]
Passing from \(\gamma\) to \(\gamma^{-1}\) conjugate-transposes \(Z\), preserving the absolute norms of these minors.\formalnote{6.23}{23}

Here is the determinant comparison that turns (CR4) into a lattice discriminant. Work in finite-adelic integral bases of \((\Lambda\cap I^\perp)/(\Lambda\cap I)\). Its form \(J_0\) is unimodular. Indeed, a primitive basis of \(\Lambda\cap I\) extends locally to a lattice basis and has integral pairing-dual vectors by unimodularity. The paired block has matrix \(\left(\begin{smallmatrix}0&1\\1&*\end{smallmatrix}\right)\) and splits off integrally. Write \(D\) for the classes of the two lists of \(a\) isotropic vectors, \(p\) for the two last vectors, and \(V\) for a primitive integral basis of the orthogonal complement of \(D\). The columns of \(D\) are independent by (CR3) in the closed cells and by disjointness in the open cell. Saturate their span, giving \(D_0\) with \(D=D_0 R\), and put \(q=[\det R]\), where brackets denote the inverse Haar modulus. Form the bordered Gram matrix
\[
 G(D)=(D\mid p)^*J_0(D\mid p)+\operatorname{diag}(0,NI_2).
\]
The class of each of the two vectors in \(p\) pairs to zero with its own isotropic list and has norm \(-N\). Reordering \(G(D)\) into the two frame lists therefore gives zero diagonal blocks and off-diagonal blocks the subframe matrix and its adjoint, in the finite class coordinates. Thus \([\det G(D)]=(N^{2f}\mu_\gamma)^2\).

To evaluate the same determinant on the complement, introduce the single matrix
\[
 \mathcal K=\begin{pmatrix}
       J_0&J_0p&J_0D_0\\ p^*J_0&-NI_2&0\\ D_0^*J_0&0&0
      \end{pmatrix}.
\]
Eliminate its upper left block. The Schur complement is \(-G(D_0)\) with the block order reversed, so \([\det\mathcal K]=[\det G(D_0)]\). Eliminate the \(-NI_2\) block first instead. This gives \([N]^2\) times the index of the determinant of
\[
   \begin{pmatrix}J_+&J_0D_0\\ D_0^*J_0&0\end{pmatrix},
                 \qquad J_+=J_0+N^{-1}J_0pp^*J_0 .
\]
Because \(D_0\) is primitive and \(J_0\) unimodular, the map \(D_0^*J_0\) has an integral right inverse \(Q\). Its kernel has basis \(V\), so \((Q\mid V)\) is an integral change of basis. After that change the last block row is \((1\ \ 0\ \ 0)\) and the last block column is its adjoint. Expanding along them leaves precisely \(\det(V^*J_+V)\) up to a sign. Finally \(G(D)\) is obtained from \(G(D_0)\) by the basis change \(\operatorname{diag}(R,I_2)\) in both arguments. Equating the two evaluations gives
\[
       [\det(V^*J_+V)]\,q^2[N]^2=[\det G(D)]
                                  =(N^{2f}\mu_\gamma)^2.          \tag{CR5}
\]
The metric on the graph sends the two last-coordinate pairings of a lift to the correction \(N^{-1}J_0pp^*J_0\). Thus \([\det(V^*J_+V)]\) is its squared lattice discriminant. Since \([N]=N^{2f}\), (CR5) gives \(q\,d(B_\gamma^1)=\mu_\gamma\) in the notation of (FL2). The calculation only inverts \(N\) in the finite-adelic algebra, so it applies at primes dividing \(N\) as well.\formalnote{6.24}{24}

\subsection{The intersection test at a level prime}
For the finite analogue, set \(M=p_0p_1\), take one of the specified split places \(w\) above either prime, and write \(s\) for the action of a lattice stabilizer element on the residue space. Agreement here compares the characters of integral unipotents \(u,u'\) whenever \(us=su'\) after level reduction. Suppose \(Fx=sFy\) modulo \(w\), with largest occupied coordinate \(a<\rho\), and \(m\) annihilates the two lists of columns through \(a\). Lift the residue coordinates of \(x,y,m,ms\) independently to \(\widehat{\mathcal O}_E\) with zero residues at every other specified place, using the Chinese remainder theorem and choosing zero as the lift of zero. Since \(w\ne\bar w\), the conjugates of these lifts reduce to zero at \(w\).

Use the same Eichler formula over the integral adeles with an integral \(\tau\) of trace one and \(s_0=\tau[w_0,w_0]\) for a representing vector \(w_0\) of the lifted linear form. Such an integral \(\tau\) exists since \(E\) contains a primitive sixth root of unity. The integral lattice form in these coordinates is unimodular. For any \(r\in\widehat{\mathcal O}_E\), the two maps reduce at the chosen place to \(1+r(Fx)m\) and \(1+r(Fy)(ms)\), since all terms involving a conjugated lifted frame coefficient vanish there. At the chosen orientation of every other pair the maps reduce to the identity. Their reductions therefore intertwine with \(s\) at the full level, the conjugate actions being determined by unitarity. Agreement equates their characters for all \(r\). Self-duality of \(\widehat{\mathcal O}_E\) for the normalized character implies that the difference of the two lifted simple-root pairings is divisible by \(M\). Reduction at \(w\) gives (CR2) over its residue field.\formalnote{6.25}{25}

Consequently any nonzero intersection modulo \(w\) yields a relation at the last index with coefficient pair \((\mu,\nu)\ne(0,0)\), by iteration. If two integral forms kill both isotropic lists, their values on the two last vectors form a \(2\)-by-\(2\) matrix whose reduction kills \((\mu,-\nu)^t\). Its determinant vanishes modulo \(w\). When the determinant is a finite-adelic unit, the quotient of the integral adeles by its ideal surjects onto the product of the residue fields at all meeting places. Taking orders gives the factor \(\mathrm{meetNorm}(s)\) used in (FL1). This argument requires only that one last coefficient survive in the residue field.\formalnote{6.26}{26}

\subsection{A bounded representative for the triangular coordinates}
At each complex place apply ordinary Euclidean Gram--Schmidt to the moved normalized frame \(y_v^{-1}F_{\rm norm}\), retaining each orthogonal vector \(h_i\) with its length. Write the matrix of the original columns in these vectors as \(C_v\), upper unitriangular. A unipotent change of frame adds earlier columns to later ones and leaves every orthogonal vector unchanged. Thus the triangle transforms by right multiplication by the inverse frame matrix. In rational coordinates put \(T_{ab,v}=C_{ab,v}d_{a,v}/d_{b,v}\), where \(d_{a,v}=\kappa_v^{\rho-a}/\sqrt N\) and \(|\kappa_v|\ge1\). The change \(T\mapsto T\theta(u)^{-1}\) now uses the rational frame entries of \(u\), embedded at every complex place.\formalnote{6.27}{27}

Let \(D_{\rm res}\) represent \(\mathbb C^f/M\mathcal O_E\). We can choose it measurably inside \(\max_v|w_v|\le b^{26}\). Indeed, for any \(w\), the covering estimate for the different gives \(l\in M\mathcal O_E\) with \(|w_v+l_v|\le M(2f)^2|\delta_E|_v\). Admissibility bounds these three factors by \(b^{23},b^2,b\). There are finitely many such translates in the closed box, by discreteness of the embedded lattice. For each coset take the lexicographically least point of its intersection with the box in fixed real coordinates. This is measurable, since the set of lattice translations is countable and the comparison of two points is Borel.\formalnote{6.28}{28}

The image of the class unipotents on the rational frame has entries above the diagonal in \(M\mathcal O_E\) and contains every \(1+lE_{ab}\) for \(a<b,\ l\in M\mathcal O_E\). For the latter assertion take the coordinate dual of \(f_b\) on the class lattice, prescribing its entries off \(X_0\). The possible entries on \(X_0\) giving an integral dual form constitute a nonempty open set at the finite adeles, and density supplies rational entries. The integral Eichler transformation in row \(a\) formed from this dual form and the scalar \(l\) acts on the frame by the required elementary matrix.\formalnote{6.29}{29}

Reduce \(T\) in order of increasing \(b'-a\). Right multiplication by \(1+lE_{a b'}\) adds \(l\) to entry \((a,b')\) and can otherwise change only entries \((i,b')\) with \(i<a\). Their distance from the diagonal is greater. The unique translate of \(T_{a b'}\) into \(D_{\rm res}\) therefore preserves all preceding choices. If two triangular group elements give reduced matrices, their quotient has a first nonzero superdiagonal. There the corresponding two entries would be distinct representatives differing by a lattice translation, impossible. Thus the acting \emph{frame matrix} is unique, and the elements \(u\) yielding it form one coset of the subgroup fixing the frame. This gives multiplicity one on the quotient used in the box partition.\formalnote{6.30}{30}

Finally, for \(a\le b'\), \(|d_{b',v}/d_{a,v}|\le1\) gives \(|C_{ab',v}|\le b^{26}\) on a reduced point. Its Gram--Schmidt vectors have length at most the Bessel length \(L(y_v)\). Indeed they are invariant under a left Bessel move, orthogonal projection decreases norm, and each moved frame vector has length at most the operator norm of the inverse of that representative. In the normalized unitary group the inverse has the same norm as the representative because \(J^{-1}=J=J^*\) is Euclidean unitary. Taking the infimum in \(L\) proves the bound. Expanding the original vector in the \(h_i\) costs at most \(n b^{26}L(y_v)\le b^{27}L(y_v)\). The cut \(L(y_v)\le b^{6n+2}\) gives \(b^{27+6n+2}\le b^{37n}\), the radius used for both frames at a near point.\formalnote{6.31}{31}

The local inputs to these calculations are additive approximation, self-duality for the character normalized by the different, and the Euclidean lattice covering estimate used to choose \(D_{\rm res}\). The subsequent count of the three frame-lift fibers, the index cost of the level-\(M\) condition, and the use of the meeting factor for a pair with common graph are given in \S\ref{sec:frame-lifts}. The averaging over classes and the height sum are in \S\ref{sec:lattice-heights}.


\section{The one-frame measure and the depth of a remote position}
\label{sec:shell-measure}
We prove the real-place shell estimate (L5) of \S\ref{sec:local-orbital}. Write \(R\) for the radius of the Bessel cut, \(L\ge R\) for the scale of the two heads, \(A=L^2\), \(m=k-5=s+1+3n\), and \(P=2^h\) for a shell height. Throughout this section size \(u\) at scale \(L\) means that each coordinate has absolute value at most \(L^u\). Write \(|c|_{\rm mod}=|c|^2\) for the modulus of multiplication by a complex scalar on planar area. Lengths and radii without this subscript use the ordinary complex absolute value.

\subsection{The reciprocal volume that measures depth}
Let \(A_a=\langle e_0,\ldots,e_{a-1}\rangle\) and \(Y_a=A_a^\perp/A_a\). At a group element \(g\) give this quotient the norm
\[
       |[v]|_{g,a}=\inf_{w\in A_a}\|g^{-1}(v+w)\|_2 .
\]
For a pair \(g,gx\), define
\[
 d_a(g,x)=\inf_{\substack{v\in A_a^\perp\\
                    \langle e_a,v\rangle=1}}
          (|[v]|_{g,a}^2+|[v]|_{gx,a}^2),\qquad D_a(g,x)=d_a(g,x)/\pi.
\]
The product map \([v]\mapsto\langle e_a,v\rangle\) on the open unit ball for the sum of squared norms has image the disc \(\{c:|c|^2d_a<1\}\). Indeed a lift of a point \(c\ne0\) can be divided by \(c\), and any lift over \(1\) can be multiplied by \(c\). The two norms can use different representatives of the same class. This disc has area \(1/D_a\). At successive floors the weight is
\[
                  w_a(g,x)=\prod_{b<a}D_b(g,x)^{-s},\qquad
                         w_a/w_{a+1}=D_a^s.                           \tag{SH1}
\]

We need a uniform moment bound for these depths. Let \(X'\cong\mathbb C^m\) be the free orthonormal coordinates and let \(J\) be the Hermitian involution of the real group. Its negative eigenspace has dimension \(n\). In \(X'\) impose the at most \(n\) linear conditions of a frame chart and require the projections to the negative eigenspace of both \(g^{-1}w\) and \((gx)^{-1}w\) to vanish. There are at most \(3n\) conditions in all. Choose \(s+1\) orthonormal columns \(O\) in their kernel. On the image of \(O\), both Euclidean lengths after transport equal the original length: there \(J\) acts as \(1\) after transport and the Hermitian form is preserved. For \(u\in X'\) put \(q=\|O^*u\|^2\). Products with vectors short at both points and satisfying the chart conditions fill a disc of radius \(\sqrt q/2\). In fact for \(|z|\le\sqrt q/2\) take \(w=(z/q)OO^*u\). It has product \(z\) with \(u\) and sum of squared lengths at most \(1/2\). Consequently
\[
                           \Phi(u)=36/\|O^*u\|^2
\]
bounds the reciprocal area of the products. It also bounds the area of the scalars \(c\) for which all the products by \(cu\) lie in the unit disc: testing the boundary point forces \(|c|\le 2/\sqrt q\). Set \(\Phi=+\infty\) on the null space of \(O^*\).

Here is the estimate that removes an exponential dependence on \(s\). Extend \(O\) to a unitary basis. On the coordinate box of size \(u_0\) each transformed coordinate has modulus at most \(L^{u_0+1}\), since \(m\le L\). There are \(3n\) remaining complex coordinates. The radial variable \(q\) satisfies \(q\le Q=L^{2u_0+3}\) and has Lebesgue density \(\pi^{s+1}q^s/s!\) in the first \(s+1\) coordinates. Integrating gives
\[
 \int_{\text{box of size }u_0}\Phi(u)^s\,du
  \le \frac{\pi(36\pi)^s}{s!}\,Q\,
                   (\pi L^{2u_0+2})^{3n}
  \le A^{99n(u_0+1)}.                                                 \tag{SH2}
\]
Indeed \(\pi(36\pi)^s/s!\le 4e^{144}\le4\cdot3^{144}\), and \(145+(2u_0+3)(3n+1)\le198n(u_0+1)\). The extra complex dimension \(s+1\) makes the radial integrand constant.\formalnote{6.32}{32}

\subsection{Integrating a single pivot}
Let \(\theta_a\) be the measure on floor \(a\). A point is near at \(a\) if it can be moved into the box by a transformation preserving the first \(a\) frame vectors modulo their predecessors in the prescribed partial unitriangular group. The orbit map \(\eta\mapsto\eta^{-1}e_a\) carries \(\theta_a\) to the quadric measure for an isotropic floor; use the last frame vector on the final floor. If \(\eta g\) is near at \(a+1\), then \(g\) is near at \(a\). A change of variable therefore permits us to take \(g\) in the box. The relevant orbit points have size \(7\), by multiplying their boxed representatives. The coordinates in the floor basis then have size \(7+4=11\): its inverse reads products with its columns and the Gram matrix merely permutes each hyperbolic pair.

For the quadric or the ordinary vector measure, the chart disintegration in the floor basis gives, for nonnegative functions \(h\) of the free coordinates alone,
\[
 \int_{\mathrm{size}\ b}h(v_{X'})\,d\nu(v)
       \le A^{99n(b+6)}\int_{\mathrm{size}\ b+4}h(u)\,du .            \tag{SH3}
\]
Fix the coordinates off one reserved block of two positive complex directions. Their squared length in the equation of the quadric is then fixed. The block gives the radial sphere measure, of density proportional to its squared radius. Integrate this bounded fiber and then the bounded hyperbolic and reserved coordinates. The floor basis costs four in the size exponent, and the uniform fiber and coordinate-volume bound is \(A^{99n(b+6)}\). The free coordinates are disjoint from the solved block, so Fubini leaves their full integral. This uses the normalized quadric and section Jacobians of the source.\formalnote{6.33}{33}

For \(h=1\), take \(b=7\) and take the gauge at the cut radius \(R\). The free volume is at most \((R^{24})^m\): one coordinate has area \(\pi R^{22}\le R^{24}\). Thus at every pivot and every \(g\) the mass of \(\{\eta:\eta g\ \text{near at }a+1\}\) is at most \(R^{2574n}(R^{24})^m\le V\), where \(V=(\pi R^{24})^kR^{9999n}\). For \(h=\Phi^s\), the reciprocal-area comparison bounds \(D_a(\eta g,x)^s\) by \(h\) on the same orbit. Indeed the common short vectors are orthogonal to the earlier frame vectors, and the product of the orbit point with such a vector is its free-coordinate pairing. Combining (SH2) and (SH3) yields
\[
 \int_{\eta g\ \mathrm{near\ at}\ a+1}
                 D_a(\eta g,x)^s\,d\theta_a(\eta)
           \le A^{99n(13+12)}\le A^{10^4 n}.                          \tag{SH4}
\]

The reciprocal power occurs when integrating in the opposite order, over a lower Heisenberg group. Write its coordinates as \((\omega,\beta)\), with \(\beta\) skew. For a product \(c=\langle e_a,\zeta\rangle\) from a stable pair \(\zeta,\zeta'\) and a common short vector \(u\), subtraction of the pairings before and after the Heisenberg translation gives \(\bar c\langle\omega,u\rangle\). The other terms vanish because \(\zeta-\zeta'\in A_a\) and \(u\perp A_{a+1}\). The two short pairs bound this difference by \(1\). Thus the volume of permitted \(c\), namely \(D_a^{-1}\), is at most \(\Phi(\omega_{X'})\) by the second area estimate. On a position ball of loss \(A'\), one can center \(\omega\) in a box of size \(17\) modulo the preceding frame vectors. The skew differences have size \(18\). Applying (SH3) to ordinary volume and (SH2) to \(\omega_{X'}\), and integrating the skew coordinate, bounds this integral by
\[
  \int_{\mathrm{close}}D_a(g,g^{-1}\eta y)^{-s}\,d\eta
     \le (A')^{99n(17+6)+99n(17+4+1)+19}
     \le (A')^{4600n}.                                                 \tag{SH5}
\]
The archimedean level unit is \(1\), so its Heisenberg modulus is \(1\). For the corresponding head section \(\langle e_a,v\rangle=\bar p_a\), scaling its vector and skew coordinates contributes \(|p_a|_{\rm mod}^{-(r_a+1/2)}\), \(r_a=k+1+2(\rho-a-1)\). Using the section box of size \(20n\) in (SH2)--(SH3) gives a bound \(A^{10^6 n^2}|p_a|_{\rm mod}^{-(r_a+1/2)}\) for the section integral of \(D_a^s\). Iterating along at most \(n\) floors gives \(C=A^{10^6n^3}\) times the product of these head moduli. Equation (SH1) makes successive weight factors telescope.\formalnote{6.34}{34}
\subsection{The shell bound from two second moments}
For either frame kernel \(K(g,y)=\langle F(g),F(wy)\rangle\), write \(K_{a+1}\) for its average over the floors above \(a\), and write \(N_j(K;x)\) for the sum over \(a=j,\ldots,\rho\) of \(w_a(g,x)|K_{a+1}(g,gx)|^2\) integrated in \(g\) through floors \(0,\ldots,a\). The one-frame estimate on a ball of loss \(A'\) is
\[
 \int_{\rm ball}N_j(K;x)\,dx
       \le n\left(\pi^{2m}(A')^{4900n}\right)^\rho
                      \int |K(g,y)|^2\,d\bar g\,d\bar y .             \tag{SH6}
\]
Here \(d\bar g\) denotes the floor-normalized quotient measure. For one pivot \(a\), disintegrate Haar measure in \(x\) by writing \(x=g^{-1}\nu y\), with \(\nu\) in the unipotents below \(a\). The absolute value of \(K_{a+1}\) is unchanged on this \(\nu\)-fiber by its character law. The integral of \(w_a\) along the close part of the fiber is bounded by (SH5) iterated \(a\) times. It remains to integrate \(|K_{a+1}|^2\), cut to those \(g,y\) for which the fiber of close points has positive measure. One floor upward replaces the integral of this cut square over a Heisenberg variable by the integral of the next cut square at \((\eta g,\eta y)\), with factor at most \(\pi^{2m}(A')^{4900n}\). Its character has unit frequency, (L3) gives the squared-norm bound \(\pi^{2m}\) on the free coordinates, and on centering the support the coordinate bounds in the Heisenberg Parseval inequality give exponent \(100n(1+17+31)=4900n\). If the next fiber of close points has measure zero, disintegration gives measure zero almost everywhere for the smaller fibers. This justifies keeping the cutoff through the iteration. At the compact endpoint Cauchy--Schwarz has constant one. The remaining double quotient integral is exactly the norm on the right. Each lower-floor cost in (SH5) is at most the same Parseval cost. Thus the total product has \(\rho\) factors regardless of \(a\). Summing at most \(n\) possible pivots proves (SH6).\formalnote{6.35}{35}

For the first kernel take \(w=1\) and start the sum at floor zero. For the second use the relative Weyl element of \(p,q\) and start at \(j\). Call their moments \(N_1,N_2\). The weighted square product used in (L4) is
\[
    \mathcal Q=D\,C^3N_1N_2,\qquad
               D=\prod_{a<j}|p_a|_{\rm mod}^{-(2r_a+1)} .
\]
The powers of \(C\) allow the inverse-weight integrations through the head and the pivot tail sum. In particular, with the bound \(V\) for unweighted pivot mass, two applications of Cauchy--Schwarz on the pivot give the no-gain bound \(\phi^2\le V^2\mathcal Q\). To track a gain, suppose at each active point some \(a\ge j\) has \(D_a\ge\sigma^{-2}\). Cover the first pivot integral by its pieces indexed by such an \(a\). On each piece \(1\le \sigma^{2s}D_a^s\). By (SH1) the squared absolute correlation there is bounded by \(\sigma^{2s}w_{a+1}^{-1}\) times the corresponding weighted square. Cauchy--Schwarz bounds the integral squared by the product of the two integrals. The first is at most \(A^{10^4 n\rho}w_j^{-1}\), by (SH4) one floor at a time with the near indicators. The second is bounded by the sum of the weighted squares above \(j\). The at most \(\rho\) choices cost \(\rho^2\) after squaring, with \(\rho^2 A^{10^4n\rho}\le C\). The other pivot uses the mass bound \(V\), and \(V\ge1\). This proves
\(\phi^2\le (V\sigma^s)^2\mathcal Q\) for the absolute majorant and hence for the smaller oscillatory integrand.\formalnote{6.36}{36}

Let \(E_0\) be any pointwise gain satisfying \(\phi^2\le E_0^2\mathcal Q\). Denote \(\mathcal B=\int|\langle F(g),F(y)\rangle|^2\,d\bar g\,d\bar y\). For the relative Weyl element the quotient change of variables multiplies this norm by \(J_h=\prod_{a<j}|p_a/q_a|_{\rm mod}^{r_a+1}\). The square of the head density \(d_h\) is \(\prod_{a<j}|p_aq_a|_{\rm mod}^{r_a}\), so
\[
           d_h^2 D J_h=\prod_{a<j}|q_a|_{\rm mod}^{-1}\le A^2 .
\]
All exponents of \(p_a\) cancel. Cauchy--Schwarz in \(x\) and (SH6) applied to both kernels now give
\[
 d_h\int_{\rm ball}\phi(x)\,dx
 \le E_0\, A^{1500000n^3+1}(A')^{25000n^2}
                                  \pi^{2m\rho}\mathcal B .
\]
Here \(1500000\) is \(3/2\) times the exponent coefficient in \(C\), and \(n(A')^{4900n\rho}\le (A')^{25000n^2}\). For a dyadic shell of height \(P\), a single archimedean ball of scale \(LP\) covers the shell, so take \(A'=AP^2\). Since \(A\ge2\), \(P\ge1\), the last bound is at most
\[
          E_0 A^{2\cdot10^6n^3}P^{10^5n^2}\pi^{2m\rho}\mathcal B.
\]
This is (L5), with \(E_0=Vg^s\) and \(\pi^{2m\rho}\mathcal B\) the shell normalization.\formalnote{6.37}{37}

\subsection{A vector of large depth at a remote position}
For a dyadic shell the height \(P\) is the least power of two among \(1,2,4,\ldots\) which bounds all entries of \(x-1\). Fix \(T\ge L^{11n}\). Suppose \(F(g),F(gx)\) have nonzero inner product and the Weyl translates for the two heads are near at \(j\). If every \(d_a(g,x)<T^2\) for \(a\ge j\), the stable product disc at each such floor contains \(T^{-1}\). Dividing the corresponding two stable lifts by \(T^{-1}\) gives dual vectors of product \(1\), congruent modulo the earlier \(e_b\), short at their respective points. At \(a<j\) the heads already supply these vectors. Cramer's rule on the first \(a+1\) frame products uses the inverse partial products bounded by \(L\), and the resulting standard duals have size \(8\rho+7\) at scale \(L\). Scaling to a common short pair costs two further powers, with \(8\rho+9\le10n\). Thus the inverses of the stabilizing scalars have size one at scale \(T\). At that scale one can take size \(5\) for all the normalized short duals.

To pass from the duals to a bound on \(x-1\), first compare the gauges of \(g,gx\) on \(Y_\rho\). Moving both into the cut box by Bessel unipotents acts trivially on \(Y_\rho\), and multiplication of the difference of the boxed points by the bounded inverse of the second gives a loss of size \(9\). This means a vector of size \(u\) at the first point has a lift whose coordinate difference at the second has size \(u+9\). Suppose the corresponding loss on \(Y_{a+1}\) is \(t_0\). For \(v\) on \(Y_a\) write \(v=\alpha\zeta+v_0\), where \(\zeta\) is the short dual and \(\alpha=\langle e_a,v\rangle\), so \(v_0\perp e_a\). The comparison on \(Y_{a+1}\) gives a second lift \(v_0'\) with bounded difference. The coefficient \(\beta\) of \(e_a\) in \(v_0'-v_0\), modulo earlier directions, is read by pairing with the short dual. If \(\zeta'\) is its second lift and \(f'\) a short lift of \(e_a\) at the second point, then \(\alpha\zeta'+v_0'-\beta f'\) is a lift of \(v\) on \(Y_a\). Bounding the scalar products, differences and the two scalar multiples with a common size \(c\) for the short vectors increases the loss by at most \(4c+12\). Here \(c=5\), so all \(\rho\) steps cost at most \(9+32\rho\). On the whole space the lift is unique. Testing on the coordinate basis and taking the Hermitian adjoint to convert an inverse to the original isometry costs four more. We obtain
\[
                     \max_{u,v}|(x-1)_{uv}|\le T^{99n}.
\]
Consequently, outside this ball an active point has some \(a\ge j\) with \(d_a\ge T^2\). This also covers \(j=\rho\), since a remote position with no later isotropic vector cannot be active.\formalnote{6.38}{38}

For \(0<\sigma<1\) assume \(A^{1200n^2}\le \sigma^{400n^2}P\). Choose \(T=L^{11n+1}/\sigma\). Then \(2T^{99n}\le P\), since \(2\le L\), \(99n\le400n^2\), and \(1+(11n+1)99n\le2400n^2\). By minimality of \(P>1\), at least one entry of \(x-1\) exceeds \(P/2\). Thus some \(D_a=d_a/\pi\ge T^2/\pi\ge\sigma^{-2}\). Choosing \(\sigma=P^{-1/(800n^2)}\) satisfies the assumption whenever \(P>A^{2400n^2}\), exactly the remote range of (L5). The shell summation in \S\ref{sec:local-orbital} now applies. The local disintegration into the normalized orbit, section, and Heisenberg measures and its Parseval formula enter (SH3), (SH5), and (SH6); their Haar normalizations determine the stated power of \(\pi\).\formalnote{6.39}{39}


\section{Counting the frame lifts and gaining from a common block}
\label{sec:frame-lifts}
The lattice sum (G3) weights a common lattice \(L\) by \(d_L^{-k}\). Two facts make it applicable to the second moment. First, the number of frames above one such \(L\) grows with its \emph{corank}, even when its rank is large. Second, the two orbital denominators pay for \(d_L\), with an extra gain when the two links differ. We give the finite-place and archimedean arguments that produce these two bounds.

The minor tests determining the cell and the reduction giving the radius \(b^{37n}\) for these lifts are proved in \S\ref{sec:cell-reduction}.

\subsection{The three measurements of a frame}
Use the coordinates \(A\oplus(X_0\oplus Ez)\oplus A^*\), where \(A\) and \(A^*\) are paired isotropic spaces of dimension \(\rho\), \(X_0\) is definite of dimension \(k\), and \(\langle z,z\rangle=-N\). Write \(F_0=(e_1,\ldots,e_\rho,z)\), \(n=\rho+1\), \(f=[E:\mathbb Q]/2\), and \(F=\gamma F_0\). Let \(Z=F_0^*JF\) be the matrix of products in the ambient form \(J\), \(P\) its first \(\rho\) rows, and \(H\) the \(X_0\)-coordinates of \(F\). Thus \(P\) reads its \(A^*\)-coordinates and the last row of \(Z\) reads \(-N\) times its \(z\)-coordinates. A move is a rational unipotent fixing \(F_0\) and congruent to the identity at level \(M\) in the lattice of the source class. It preserves \(Z\). A near term admits an archimedean point \(y\) at which \(y,\gamma^{-1}y\) are both in the box of the cut. We use indices \(0,\ldots,n-1\) for columns below.

Fix the source class and the cell rank \(0\le j\le n\). Let \(Y\) have the first \(j\) rows of \(Z\) followed by the corresponding rows of the identity. Its determinant \(\delta\) is the leading \(j\)-minor of \(Z\), nonzero by the cell condition (the empty determinant is \(1\)). For every \(a\ge j\), \(F Y^{-1}e_a\) lies in the span of \(F_0\), by the shared-vector condition of the cell. The matrices of shared vectors used below record the \(A\)-coordinates of \(F\operatorname{adj}(Y)e_a\) for these \(a\), with zero at the other columns. Let \(O=\widehat{\mathcal O}_E\) and \(\Lambda_0=\Lambda\cap X_0(\mathbb A_{E,f})\) in the lattice of the source class. As in \S\ref{sec:lattice-heights}, \(L=B O^r\subset\Lambda_0\), with \(B=ya\) for a rational \(k\)-by-\(r\) matrix \(y\) of rank \(r=k-\ell\) and \(a\in\mathrm{GL}_r(\mathbb A_{E,f})\). In what follows \(y\) denotes this rational matrix. The volume convention is \(d_L^2=[\det(B^*B)]\).

\paragraph{The lift of \(L\).}
The condition \(L\subset L_\gamma^M\) means that each basis vector of \(L\) has a lift obtained by adding a vector of \(M(\Lambda\cap A)\) to make it orthogonal to \(F\). In coordinates these conditions give a \(\rho\)-by-\(r\) matrix \(w\) such that
\[
                  H^*y+MP^*w=0,\qquad wa\in M_{\rho,r}(O).
\]
One can choose \(w\) over \(E\). The individual adelic lifts give an adelic solution to the rational linear system. Row reduction gives a rational particular solution: every consistency row is a rational zero since it is zero over the adeles. The remaining parameters lie in a rational vector space. Its rational points are dense in its finite adeles by additive approximation, and the condition \(wa\) integral is open. Thus an integral solution in this sense exists among the rational ones.

Partition the frames by the rows of \(wa\) modulo the row lattice of \(B^*B\). This costs at most \(d_L^{2\rho}\), since the integral Gram matrix has full rank and its row lattice has index \([\det(B^*B)]\). For lifts \(w,w'\) in the same class write \((w'-w)a=t^* B^*B\) with \(t\) integral. Then
\[
       \eta=y(y^*y)^{-1}(w'-w)^*,\qquad \eta^*y=w'-w.
\]
Over the finite adeles \(\eta=Bt\): the two matrices have their columns in the span of \(y\) and the same products with \(y\), and \(y^*y\) is invertible by definiteness. In particular the columns of \(\eta\) lie in \(\Lambda_0\).

The requisite move sends \(H\) to \(H-M\eta P\). Here is an integral realization, which also works at places above \(2\). Denote by \(X\) the columns \(e_a\) and by \(V\) the columns of \(\eta\) embedded in \(X_0\). Choose \(\tau\in\mathcal O_E\) with \(\tau+\bar\tau=M\): the field contains a primitive sixth root \(\zeta_6\) with trace \(1\) over the real subfield, so \(\tau=M\zeta_6\) works. Put
\[
          U=1+M\bigl(XV^*J-VX^*J-X\tau(V^*JV)X^*J\bigr).
\]
Since \(X^*JX=0\) and \(X^*JV=0\), the quadratic term in \(U^*JU-J\) is \(M^2 JX(V^*JV)X^*J\) and the linear terms sum to its negative. Thus \(U\) is unitary. The added operator is nilpotent by the same two orthogonalities, and all its terms have integral lattice coefficients and carry the factor \(M\). It fixes \(F_0\); its \(X_0\)-row action on \(F\) is as asserted. Applying this to \(w'-w\) makes the new matrix of products \(H^*y\) equal to \(-MP^*w'\). In each class choose one reference \(w'\); these products then agree for all its frames.\formalnote{6.40}{40}

\paragraph{Removing differences in the definite coordinates.}
Write \(\Pi=\Lambda_0\cap\ker(y^*)\), a saturated lattice of rank \(\ell\). After the first moves, the columns of \(V_0=H'-H\) are orthogonal to \(y\). They also lie in \(\Lambda_0\). Indeed, pass from each term to \(h_c^{-1}\gamma h_{c'}\) in the source and target lattice bases. The middle blocks of these frames are integral. On the difference, the \(z\)-coordinate is zero because the last rows of \(Z\) agree, leaving exactly the lattice coordinate of \(H'-H\). Hence measuring all \(n\) columns relative to one reference modulo \(d\Pi\), where \(d=M\delta\), costs at most \([d]^{\ell n}\). This follows in local bases of \(\Pi\) by taking the index of multiplication by \(d\).

Inside one residue class all columns of \(\eta'=V_0\operatorname{adj}(Y)/(M\delta)\) are integral in \(\Lambda_0\). We show that \(M\eta'=V_0Y^{-1}\) vanishes on the last column and at all columns \(a\ge j\). For \(a\ge j\) the shared vectors of both frames lie in \(\operatorname{span} F_0\), which has zero \(X_0\)-coordinates. In the remaining case \(j=n\) take \(t=Z^{-1}e_{n-1}\). The vector \(Ft\) has zero \(A^*\)-coordinates and \(z\)-coordinate \(-1/N\). Its norm equals the norm of \(F_0t\). At every complex place this gives
\[
                 N\|(Ft)_{X_0}\|^2=1-N^2 |t_{n-1}|^2\le1.
\]
It holds for both frames. Their difference in that column has squared length at most \(4/N\le4/5<1\). An integral column of \(\Lambda_0\) with this property at every complex place is zero: its Hermitian norm is an algebraic integer, and all conjugates have modulus below one, so its absolute field norm is an integer below one in modulus. Definiteness gives the assertion. Now multiply \(\eta'\) by \(Y\). Its vanishing columns leave only the rows of \(Y\) that coincide with the corresponding rows of \(P\). Thus \(V_0=M\eta P\), with \(\eta\) the first \(\rho\) columns of \(\eta'\). The move just constructed makes \(H'\) agree exactly with \(H\). No factor growing with the ambient dimension \(k\) has occurred.\formalnote{6.41}{41}

\paragraph{Removing differences in the isotropic coordinates.}
Suppose \(Z,H\) and the shared-vector matrix coincide. These shared vectors are unchanged by the preceding moves because the columns concerned belong to \(\operatorname{span}F_0\), which is fixed pointwise. Only the top \(\rho\) rows of \(F\) can differ. They are integral in the class coordinates, which preserve these rows. Measure them modulo \(d'=M\delta\bar\delta\). There are at most \([d']^{n\rho}\) possibilities. To prove injectivity modulo moves, write the difference in those rows as \(d' D\) with \(D\) integral. Equality of the two frame Gram matrices gives \(D^*P=-P^*D\), since a difference valued in \(A\) has zero self-pairing.

Let \(\alpha\) consist of the first \(\rho\) columns of \(\operatorname{adj}(Y)\). We claim \(D\alpha P=\delta D\). The columns \(a\ge j\) of \(D\operatorname{adj}(Y)\) are zero by equality of shared vectors. At \(j=n\) we still need its last column to vanish. Here \(Y=Z\) so \(P\operatorname{adj}(Y)e_{n-1}=0\). Skew-Hermitian symmetry gives \(P^*D\operatorname{adj}(Y)e_{n-1}=0\); the map \(P^*\) is injective because \(P\alpha=\delta I_\rho\). Consequently all the possibly omitted columns vanish. Multiplication by \(Y\) now proves the claim. Define the integral skew-Hermitian matrix
\[
         \sigma=\bar\delta D\alpha-\delta(D\alpha)^*
                         -\alpha^*(P^*D)\alpha .
\]
Multiplying by \(P\) on the right, the first term yields \(\delta\bar\delta D\). The remaining terms cancel using \(D^*P=-P^*D\) and \(D\alpha P=\delta D\). The central unipotent \(1+M X\sigma X^*J\) fixes the standard frame and adds \(M\sigma P=d'D\) to its transformed frame's isotropic rows, so it is the desired move. This argument explains the two distinct moduli \(M\delta\) and \(M\delta\bar\delta\).\formalnote{6.42}{42}

\subsection{Bounding the two matrices at infinity}
Two counts remain before multiplying the residue-class bounds. A nonzero algebraic integer has modulus at least one at some complex place, by the product formula. Distinct integral matrices in a box of radius \(R\) are therefore at distance at least one in the sup norm over entries and places. Disjoint balls of radius \(1/2\) in the ball of radius \(R+1/2\) show that the number is at most \((1+2R)^{2f a b'}\) for \(a\)-by-\(b'\) matrices.

At a boxed point each frame vector has length at most \(b^{37n}\). This follows from its Gram--Schmidt expansion: the reduced triangular coefficients are bounded by \(b^{26}\), their sum costs at most \(b^{27}\), and the cut bounds the orthogonal-vector length by \(b^{6n+2}\). The near-term comparison for the pair of boxed points gives \(|Z_{ab'}|\le b^{134n}\), in the rational normalization at every complex place. The entries of \(Z\) are integral because they pair the two integral frames in lattice coordinates. Hence at most \(b^{270 n^3 f}\) matrices \(Z\) occur. Each leading minor satisfies \(|\delta|^2\le b^{270n^2}\) by the determinant bound from the entry norm and \(n\le b\). With \(M\le b^{23}\), the two generalized indices are therefore bounded by \(b^{2n^3 f}\), by taking the product over the complex places at the admissible scales.

For the shared vectors we also need a lower bound on the size of a linear combination in terms of its coefficients. At a boxed point \(h\), choose a unipotent \(u\) with \(\|uh\|\le b^{7n}\), using the cut. It acts on the frame by a unitriangular matrix \(T\). With the normalized frame vectors \(v_i=h^{-1}f_i\) at this place,
\[
       |(Tx)_i|\le b^{7n}\Big\|\sum_a x_av_a\Big\|,\qquad
       |T_{ij}|\le b^{44n}.
\]
Since \(\det T=1\), \(x=\operatorname{adj}(T)Tx\). A cofactor of a matrix with entries at most \(b^e\) has size at most \(b^{(2+e)\rho}\) by the determinant bound and \(\rho\le b\). Taking the sum over \(n\) columns gives \(|x_i|\le b^{46n^2}\|\sum_a x_av_a\|\).

For \(a\ge j\) write \(F\operatorname{adj}(Y)e_a=F_0w\). The entries of \(\operatorname{adj}(Y)\) are bounded by \(b^{(2+134n)\rho}\). Express this vector at the two boxed points. After dividing by the scale of the \(i\)-th frame vector, its coefficient at the first point is \(\sigma(w_i)\) divided by that scale at the embedding \(\sigma\). The length of the combination at the second point is at most the maximum cofactor times \(b^{1+23\rho+37n}\) divided by the same scale: the \(n\) terms cost \(b\), the change in absolute scale across the frame costs \(b^{23\rho}\), and each normalized vector has length at most \(b^{37n}\). The two expansions give the same vector because the points differ by \(\gamma\). The preceding inverse bound now gives
\[
               |\sigma(w_i)|\le
         b^{46n^2+(2+134n)\rho+1+23\rho+37n}\le b^{182n^2}.
\]
The shared entries are integral, being \(A\)-coordinates of an integral frame times an integral adjugate. Packing gives \(b^{366n^4f}\) possibilities. With \(d_L\le b^{n^3 f}\), the three successive measurements bound the number at fixed \(j,Z\) by
\[
                  d_L^{2\rho}\,[M\delta]^{\ell n}\,
                   b^{366n^4 f}\,[M\delta\bar\delta]^{n\rho}.
\]
There are \(\rho+2\le b^{2f}\) possible \(j\). Substitution gives the exponent \(2000n^6(\ell+1)f\) of \S\ref{sec:lattice-heights}. Within one reduction label, equality of frames after moves puts the reductions in the same coset of the reduced frame stabilizer; their common position within this coset gives equality of reductions. The class and double-coset comparison then identifies the two chosen terms. Thus this count applies once for each term of a near pair, as required by (B8).\formalnote{6.43}{43}
To check the exponent \(134n\), the change from the real normalized Gram entry to \(Z_{ij}\) gives
\[
 |M_v|^{(\rho-i)+(\rho-j)}\,|\sigma_v Z_{ij}|
       \le N\,\|v_i(y)\|\,\|v_j(\gamma^{-1}y)\|
       \le b^{60n+37n+37n}.
\]
Here \(M_v\) is the scale in the real frame change and \(|M_v|\ge1\). The box supplies both length bounds.\formalnote{6.44}{44}

\subsection{The discriminant bound from one common block}
Let \(g=h_c^{-1}\gamma h_{c'}\) and \(g'=h_c^{-1}\gamma' h_{c'}\) in lattice coordinates. Both are integral isometries, with the same reduction \(s\bmod M\). In rational coordinates the link \(\Gamma_\gamma\) consists of the two middle coordinates of \(t,\gamma^{-1}t\) as \(t\) ranges over \((A+\gamma A)^\perp\). For a pair \((u+\lambda' z,u'+\lambda z)\) in the link, regard \(u+\lambda z^+\) as input and \(u'+\lambda' z^+\) as output, where \(z^+\) has norm \(N\). Switching the two scalar coordinates in this fashion makes the link the graph of an isometry \(U_\gamma\) of the definite space \(W=X_0\perp Ez^+\). Indeed the difference of the two original middle forms vanishes on the link by invariance of the original form. After the switch it becomes the difference of two positive forms. Projection to input is injective by definiteness. Its dimension is \(k+1\), since the kernel of the map \((A+\gamma A)^\perp\to\Gamma_\gamma\) is \(A\cap\gamma A\). Thus projection is onto \(W\). Give each graph the metric of its input.

Write \(C=\Gamma_\gamma\cap\Gamma_{\gamma'}\), \(i=k+1-\dim C\), and let \(C_0\subset X_0\) be the inputs of elements of \(C\) with both scalar coordinates zero. Put \(P_\gamma=L_\gamma^M\cap (C_0\otimes\mathbb A_{E,f})\), and similarly for \(P_{\gamma'}\). The common lattice is \(L_{12}=P_\gamma\cap P_{\gamma'}\). Let \(t_\gamma\) be the index of \((\Lambda\cap A)+g(\Lambda\cap A)\) in the saturation of that sum, now in lattice coordinates. Define
\[
 \begin{aligned}
 \mu_\gamma&=\frac{|\mathrm{Nm}_{E/\mathbb Q}\delta|}
                  {N^{2f{\bf1}_{j=n}}},&
 \Delta_\gamma&=\prod_{v\mid\infty}|\mathrm{termDisc}(\gamma)|_v,\\
 H_\gamma&=\max(N^f\mu_\gamma,\Delta_\gamma),&
 G_\gamma&=\frac{M^{2a f}H_\gamma}{I_g I_{g^{-1}}},\qquad a=\min(j,\rho).
 \end{aligned}
\]
Here \(I_g\) is the index of the products of the first \(a\) basis vectors of \(A\) with \(g F_0\) in the module of products of all of \(\Lambda\cap A\) with \(g F_0\). Take \(1\le x\le\mathrm{meetNorm}(s)\), the product of residue-field orders at level places where the reduced frame spans share a nonzero vector. Set \(X=x\) if \(i=0,j=n\), and \(X=1\) otherwise. We prove
\[
             (X d_{L_{12}})^2
                  \le (2/M)^{2if}G_\gamma G_{\gamma'}.                 \tag{FL1}
\]

\paragraph{The discriminant of one full link.}
We first construct a lattice \(B_\gamma^+\) of full rank in the graph admitting integral lifts \(t\) with \(A\)-part in \(M(\Lambda\cap A)\), such that
\[
           t_\gamma\,d(B_\gamma^+)
                   \le M^{2af}\mu_\gamma/(I_g I_{g^{-1}}).            \tag{FL2}
\]
Here is the determinant calculation. Pass to the quotient of \(\Lambda\cap I^\perp\) by \(\Lambda\cap I\), with \(I=A\cap gA\). It is unimodular locally: complete a primitive basis of \(\Lambda\cap I\) by dual vectors using unimodularity, then take the orthogonal complement of these paired vectors. Their Gram matrix has unit determinant. Use \(J_0\) for the resulting quotient form and \(D\) for the columns of the first \(a\) vectors of each isotropic frame in the quotient. A basis of \(I\) can be chosen with identity entries in the isotropic coordinates of index \(\ge j\), by the triangular shared-vector condition. Subtracting these basis elements from orthogonal lifts kills the corresponding coordinates without changing either middle projection.

Let \(q\) be the index of the columns of \(D\) in their saturation. Choose a primitive basis \(V\) of the integral orthogonal complement of \(D\) in the quotient. Two additional columns \(p\) are the classes of \(z,gz\). On \(V\) the positive metric has Gram matrix \(V^*(J_0+N^{-1}J_0pp^*J_0)V\). In fact each negative scalar coordinate removed from the original form contributes the square of its product with its frame's \(z\) divided by \(N\). Adjoin two independent vectors of norm \(N\), and add one to each column of \(p\). Each list (the first \(a\) isotropic columns, and its corrected last column) is now totally isotropic. Its pairing with the other list is the subframe minor on these indices, so the absolute norm of the determinant of the two lists together is \((N^{2f}\mu_\gamma)^2\). At \(j=n\) this is the full frame determinant squared in absolute norm. At smaller ranks the shared last vector gives the additional factor of modulus \(N\) at each embedding in the subframe minor.

For clarity, the complementary-minor identity used in this step is
\[
 [\det V^*(J_0+N^{-1}J_0pp^*J_0)V]\ q^2[N]^2
       =\left[\det\left((D\mid p)^*J_0(D\mid p)
                                    +\operatorname{diag}(0,NI_2)\right)\right].
\]
One obtains it by using a basis of the saturated columns of \(D\), completing their pairing forms to a dual basis of the quotient, and expanding the bordered determinant along those forms. Replacing saturated columns by \(D\) contributes \(q\) from rows and \(q\) from columns. Eliminating the two norm-\(N\) entries on the complementary side contributes \([N]^2\) and leaves the displayed rank-two correction. Since \([N]=N^{2f}\), the lattice \(B_\gamma^1\) of images of these primitive lifts has \(q\,d(B_\gamma^1)=\mu_\gamma\).

Enlarge the module spanned by \(D\) in three steps: add the remaining first-frame columns, then the remaining second-frame columns, then saturate. The index at the first step is divisible by \(I_g\), by taking products with \(g F_0\), which kills the second isotropic list. Products with \(F_0\) give the divisor \(I_{g^{-1}}\) at the second step. The final index is divisible by \(t_\gamma\): projection from \(I^\perp\) to the quotient carries the saturation of the full double list into its saturation downstairs, and its kernel on the integral vectors is contained in the double list. Thus \(I_g I_{g^{-1}}t_\gamma\mid q\). Finally require the \(A\)-part of a lift to be a multiple of \(M\). Its products with \(gF_0\) lie in the row module of the full isotropic list, of rank \(a\). Requiring these products to belong to \(M\) times this module cuts an index at most \([M]^a=M^{2af}\). If they belong to that multiple, subtract the corresponding difference of \(A\)-vectors, which has zero products with \(gF_0\) and lies in \(I\) by the cell condition. This gives a lift with the required \(A\)-part. Taking the kernel of this residue map in \(B_\gamma^1\) proves (FL2).\formalnote{6.45}{45}

\paragraph{The two scalar coordinates.}
Split the rational graph orthogonally into three summands \(W_0,W_1,W_2\). The first is the part of \(C\) with both scalar coordinates zero, the first two together give \(C\), and \(\dim W_2=i\). Let \(D_0,D_1,D_2\) be the discriminants on the three successive quotients of the induced filtration of \(B_\gamma^+\), with the orthogonal metrics. An integral basis adapted to the filtration is block triangular, place by place. Its determinant is the product of the three diagonal determinants, and the rational orthogonal Gram determinant also factors. Hence \(D_0D_1D_2=d(B_\gamma^+)\). The first lattice projects isometrically into \(P_\gamma\), giving \(d(P_\gamma)\le D_0\).

Write \(e=\dim W_1\). The two scalar forms have zero common kernel on \(W_1\), so their restrictions span its dual. Choose \(e\) of them giving a nonsingular square pairing with \(W_1\). Their representing vectors in input coordinates are, up to signs and order, \(z^+\) and \(U_\gamma^{-1}z^+\), when one uses the products with \(z,gz\) for scalar forms. Thus \(e\le2\). At nontop rank the shared last vector makes these forms proportional. At top rank they are independent on the whole graph. In particular \(e=2\) if \(X>1\), and the corank of \(L_{12}\) in \(X_0\) is \(\ell=i+e-1\le i+{\bf1}_{j=n}\).

The pairing determinant in an adelic basis of the second quotient is integral: lift each column to the full lattice, changing its representative only by \(W_0\), which both forms kill. Its generalized index is at least \(X\). To see the nontrivial case, at a meeting place the agreement of characters propagates a relation of the reduced flags to a relation with at least one nonzero last coefficient (the flag-climbing argument). Apply the two integral pairing rows to this relation. All isotropic columns are killed, so the \(2\)-by-\(2\) matrix on the two last columns has a nonzero kernel vector modulo that place. Its determinant is divisible by the place. Multiplication over the meeting places proves the claim.\formalnote{6.46}{46}

At infinity apply Cauchy--Schwarz to the two exterior products defining this determinant, and multiply over embeddings. The Gram principal minor of the representers is \(1\) in rank zero, \(N\) in rank one, and in rank two has modulus \(N^2|\mathrm{termDisc}(\gamma)|_v/|\delta|_v^2\). The last equality follows by taking the determinant of their two-by-two Gram matrix, whose diagonal entries are \(N\), with the mixed entry furnished by the ratio of the consecutive frame minors in the cell. Its absolute field norm is \((\Delta_\gamma/\mu_\gamma)^2\) in that rank. In every case the square root of the absolute norm is at most \(H_\gamma/\mu_\gamma\). The finite basis changes multiply the exterior-product bound by their determinant indices, giving \(X\mu_\gamma\le D_1 H_\gamma\). Combining this with (FL2) yields
\[
                     X\, t_\gamma\, d(P_\gamma)\,D_{2,\gamma}
                             \le G_\gamma.                             \tag{FL3}
\]
If both graphs agree, the two scalar forms have the same rank on both, so the exceptional top-rank gain is available on both sides of the pair.\formalnote{6.47}{47}

\paragraph{The mixed determinant.}
The signed difference of the two middle forms vanishes on each graph. Between the two graphs its radical on either side is exactly \(C\), since each graph is maximal isotropic for that difference form. Thus the induced \(i\)-dimensional mixed pairing between the two complements is nondegenerate. For integral lifts \(t,t'\), its value in lattice coordinates is
\[
             t^*J(1-g(g')^{-1})t'\ \in M O.
\]
Removing their components in \(C\) has no effect. Consequently the mixed determinant on the projected lattices has generalized index at least \([M]^i=M^{2if}\). At an archimedean place the same form in positive input coordinates is \(\langle u,(1-U_\gamma^{-1}U_{\gamma'})v\rangle\). The operator has norm at most \(2\), so its \(i\)-th exterior power has norm at most \(2^i\). Applying this to rational bases, multiplying over the embeddings and inserting the finite determinant changes as above gives
\[
                     D_{2,\gamma}D_{2,\gamma'}\ge (M/2)^{2if}.
\]
This exhibits separately the origin of the level in the numerator and of the power of two in the denominator.

\paragraph{Passing to the same lattice on both sides.}
For \(v\in P_\gamma\) take \(t=v-Mb\) perpendicular to \(gF_0\) with \(b\in\Lambda\cap A\). Set \(m=g'g^{-1}\), so \(m\equiv1\bmod M\). Then \(mt\) is perpendicular to \(g'F_0\), and \(w=(t-mt)/M\) is integral. It belongs to \(A+g'A\): since the pair of middle coordinates of \(t,g^{-1}t\) belongs to both graphs, there is a lift \(t'\) of the same pair on the second graph. Then \(t-t'\in A\) and \(t'-mt\in g'A\), giving the assertion. Now \(v=Mb+Mw+mt\). Test the products of \(v\) against \(g'F_0\) modulo \(M\) times the products of \(\Lambda\cap A\) against that frame. Only \(Mw\) contributes. Changing \(w\) by an integral member of \(A\) or \(g'A\) leaves its residue unchanged. The test on \(P_\gamma\) therefore takes values in an image of the torsion module of order \(t_{\gamma'}\), and its kernel is \(P_\gamma\cap P_{\gamma'}\) by the definition of a level lift. Hence
\[
              d(L_{12})\le t_{\gamma'}d(P_\gamma),\qquad
              d(L_{12})\le t_{\gamma}d(P_{\gamma'}).
\]
Multiply these inequalities and apply (FL3) for the two orders of the pair and the mixed determinant bound. This gives (FL1), including \(i=0\), for which the empty mixed determinant equals one. The per-frame count above now applies to \(L_{12}\) with its actual corank \(\ell\).\formalnote{6.48}{48}


\section{Summing the small-discriminant lattices in the second moment}
\label{sec:lattice-heights}
In (B8) the number of near pairs is first bounded for each common lattice \(L\), and then \(d_L^{-k}\) is summed over \(L\). The ambient rank \(k\) is large. The estimate must save \(N^{kf}\) from the level even for corank zero. We prove the two zeta bounds needed in that sum, indicate exactly where the row condition costs the index of a nonsaturated lattice, and track the factors for the near-pair count.

Let the CM field \(E\) have \(f=[E:\mathbb Q]/2\) complex places, root discriminant \(\mathrm{rd}_E\), and the primitive \(5040\)-th root of unity used for the local integral bases. Write \(O=\widehat{\mathcal O}_E\) and take the self-dual lattice \(\Lambda=uO^k\), where \(u\) is finite adelic unitary. All rank-\(s\) lattices in the sum have the form \(Y A O^s\), with \(Y\) a \(k\)-by-\(s\) matrix of rank \(s\) over \(E\) and \(A\in\mathrm{GL}_s(\mathbb A_{E,f})\). Their adelic spans therefore come from rational subspaces. For a rational subspace \(V\), let \(S_V\) be the intersection of its finite adelic span with \(\Lambda\). Write \(d(V)\) for its discriminant, the square root of the generalized index of the determinant of an adelic basis Gram matrix. The generalized index \([a]\) for a finite idele is the reciprocal of its finite absolute value. A change from an integral basis of \(S_V\) to a sublattice \(L\) of index \(i\) multiplies the Gram determinant by \(\overline{\det A}\det A\) with \([\det A]=i\). Thus \(d_L=i\,d(V)\). Moreover \(d(V)\ge1\), since the Gram matrix of an integral basis relative to the unitary lattice is integral at every finite place.

\subsection{The zeta bound inside one saturation}
For a free rank-\(s\) \(O\)-module \(P\), let \(Z_s(a)=\sum_{M\subset P}[P:M]^{-a}\) over finite-index submodules. Take the projection \(P=O^{s-1}\oplus O\to O\). Its image on \(M\) is an ideal \(J\) of finite index, and \(M_0=M\cap O^{s-1}\) has finite index \(i_0\) in the first summand. The index of \(M\) is \(i_0[O:J]\), by the exact sequence on quotients. With \(J,M_0\) fixed, \(M\) is the inverse image of the graph of a homomorphism \(J\to O^{s-1}/M_0\). Each finite-index ideal in \(O\) is generated by an idele integral at every finite place: use local uniformizer powers and units at all but finitely many places. Thus there are at most \(i_0\) choices for this homomorphism. Summing positive terms gives
\[
 Z_s(a)\le Z_{s-1}(a-1)\zeta_E(a),\qquad
 Z_s(s+\mu)\le\prod_{j=1}^s\zeta_E(\mu+j)\le\zeta_E(\mu)^s .
\]
Here the finite-index ideals are the finite tuples of prime powers, so their index series has the usual ideal Euler product. This argument applies in local integral bases of a saturated lattice as well. Choose \(\mu=2\rho+4,\ n=\rho+1\). From the local Euler factors \(\zeta_E(\mu)\le\zeta(\mu)^{2f}\); above a rational prime there are at most \(2f\) factors, each with residue degree at least one. Since \(\sum_{m\ge2}m^{-\mu}\le2^{-\mu}(1+2/(\mu-1))\le 2^{1-\mu}\) and \(2fn\le2^{2\rho}\), we have \(\zeta_E(\mu)^n\le e^{1/8}<2\). In particular the sum of \(i^{-s-\mu}\) in a fixed saturation is at most \(2^k\) for \(s\le k\).\formalnote{6.49}{49}

\subsection{Heights of rational kernels}
Take \(V\) of corank \(\ell\). We first construct vectors \(x_1,\ldots,x_\ell\) of \(\Lambda\cap E^k\) that cut out \(V\) by Hermitian orthogonality and satisfy
\[
     \left(\prod_{j=1}^{\ell}\|x_j\|_\infty\right)^{2f}
      \le (8f\ell\,\mathrm{rd}_E)^{f\ell}\,d(V).
      \tag{G1}
\]
The length is the Euclidean norm over all complex places together. By the number-field form of Siegel's lemma due to Bombieri and Vaaler, an \(E\)-subspace of \(E^k\) has an \(E\)-basis of vectors with integral coordinates whose product of heights is bounded in terms of the subspace height \cite{upper-bombieri-vaaler}. In (G1) the vectors must lie in the specified lattice \(\Lambda\), and we compute the covolume to obtain the bound in terms of \(d(V)\). If \(\ell=0\), then \(V=E^k\) and the empty choice gives the assertion with constant one.

For \(\ell>0\), the determinant identity used here is \(d(V^\perp)=d(V)\). Locally use \(u^{-1}\) to put \(\Lambda\) in standard position. Complete a primitive integral basis of \(S_V\) to an integral invertible matrix \(B\). The last \(\ell\) columns of \((B^*)^{-1}\) form an integral basis of the orthogonal annihilator. Their Gram matrix is the last principal block of \((B^*B)^{-1}\). The complementary-minor formula gives its determinant as that of the first block of \(B^*B\) divided by \(\det(B^*B)\). The latter determinant has index one. At conjugate split places this same calculation uses mutually adjoint bases at the two places. Multiplication of the finite factors proves the identity.

To see the covolume constant, take a rational basis \(Y\) of \(V^\perp\), written as columns, and an adelic change of basis \(A\) for which \(YA\) spans \(S_{V^\perp}\) over \(O\). In complex Euclidean coordinates the volume factor for \(Y\) is \(\prod_v\det(Y_v^*Y_v)=|\mathrm{Norm}_{E/\mathbb Q}\det(Y^*Y)|^{1/2}\). The finite volume factor for the lattice of coordinates is \([\det A]\). This follows by comparing with \(\mathcal O_E^\ell\) through their intersection of finite index and taking local indices on both sides. The product of these factors is \(d(V^\perp)=d(V)\). Since \(\mathcal O_E\) has real covolume \(2^{-f}|\Delta_E|^{1/2}\), the lattice of rational points of \(S_{V^\perp}\) in its Euclidean span has covolume \(2^{-f\ell}|\Delta_E|^{\ell/2}d(V)\).

Write \(D=2f\ell\) and \(\lambda_1\le\cdots\le\lambda_D\) for the real successive minima of the unit ball. One can choose the \(j\)-th \(E\)-independent vector with length \(\le\lambda_{2f(j-1)+1}\): the \(E\)-span of the first \(j-1\) vectors has rational dimension \(2f(j-1)\). Among \(2f(j-1)+1\) real-linearly independent lattice vectors at least one lies outside that span. Hence the left side of (G1) is at most \(\prod_{i=1}^D\lambda_i\). By Minkowski's second theorem \cite[Theorem~1.3]{legacy-henk} this product is at most \(2^D\,\mathrm{covol}/\mathrm{vol}(B_D)\), where \(B_D\) is the Euclidean unit ball. The inscribed cube gives \(\mathrm{vol}(B_D)\ge(2/\sqrt D)^D\). Substitution proves (G1), even with a smaller constant. The selected vectors span \(V^\perp\), so they determine \(V\).\formalnote{6.50}{50}

Nonzero points of \(\Lambda\cap E^k\) have length at least one. Indeed \(x^*x\) is integral at every finite place, totally positive and a nonzero algebraic integer. The product of its values across the CM type is at least one. At least one archimedean length is therefore at least one. Taking differences gives separation by one. In real dimension \(D'=2fk\), packing disjoint half-unit balls bounds the number of such points of norm at most \(R\ge1\) by \((1+2R)^{D'}\). With \(a_0=2f(k+1)\), layer integration gives
\[
 \sum_{x\ne0}\|x\|_\infty^{-a_0}
  \le a_0\int_1^\infty (3R)^{D'}R^{-a_0-1}\,dR
  = 3^{D'}(k+1)\le 5^{a_0}.
\]
Choose one ordered tuple in (G1) for each \(V\). Two different subspaces cannot have the same tuple, because it determines its common orthogonal. Raising (G1) to \(k+1\), putting the product of vector norms on the other side and summing over all tuples therefore gives
\[
  \sum_{\operatorname{codim} V=\ell}d(V)^{-(k+1)}
       \le(200 f k\,\mathrm{rd}_E)^{f(k+1)\ell}.
       \tag{G2}
\]
The factor \(200\) is \(8\cdot5^2\). The spare exponent \(+1\) on \(d(V)\) makes the packing integral converge.\formalnote{6.51}{51}

\subsection{The genus sum and its normalization}
Let \(H_0=b^{n^3 f}\), \(\tau=b^{kf/4}\), \(s=k-\ell\). As in \S\ref{sec:orbital-sieve} vary the distinguished row modulo the idele \(\bar\pi_N\), of index \(N^f\). For a lattice of index \(i\) in a saturation of rank \(s\), at most \(i(N^f)^\ell\) rows annihilate the lattice. To check this prime by prime, extend a basis of the saturation to one of \(\Lambda\), and diagonalize the sublattice inside the first \(s\) coordinates. If its local diagonal entry is a uniformizer to the power \(e\) and the modulus is the uniformizer to the power \(a\), a row coefficient has at most \(q^{\min(e,a)}\) choices. Their product is bounded by \(i\). Each of the last \(\ell\) coordinates is unrestricted. Rows coming from different stabilizer cosets are distinct modulo the level, so this counts an upper bound even when the admissible rows have additional conditions.

For each \(L\) with \(di\le H_0\), \(d=d(V)\), multiply the row bound by \((di)^{-k}\). Using \(d,i\ge1\) gives
\[
     i(di)^{-k}\le H_0^{\mu+2}d^{-(k+1)}i^{-(s+\mu)} .
\]
We may now sum over the lattice and then over \(V\) by positive summation. Our two bounds give the following upper bound for the sum over row cosets above a given unimodular representative:
\[
    (N^f)^\ell\,H_0^{2\rho+6}\,2^k
               (200fk\,\mathrm{rd}_E)^{f(k+1)\ell}.
\]
To divide by the number of row cosets without imposing a uniform class size at a coarser level, integrate the sum of right translates on the adelic quotient. Haar invariance makes the integral the number of translates times the integral of the original function. At the finer source level the rational stabilizers are trivial, so its normalized integral is the sum over classes divided by their number. This yields precisely the desired average over fine classes. Their row-coset index is at least \((N^k/b)^f\).

The scale comparisons \(b^f H_0^{2\rho+6}2^k\le\tau\) and \(N^f(200fk\,\mathrm{rd}_E)^{f(k+1)}\le\tau\) absorb the two factors. For the former, \(2^k\le b^f\) and \(2+n^3(2\rho+6)\le k/4\). For the latter take thousandth powers and use \(N\le b^{60n}\), \((200fk\,\mathrm{rd}_E)^{1000}\le b\), and \(60000n+k+1\le250k\). Thus, with \(\mathcal C\) the set of classes and \(\Lambda_c\) their row-constrained lattices, we obtain
\[
 \sum_{c\in\mathcal C}\ 
 \sum_{\substack{L\subset\Lambda_c,\ \mathrm{rank}\,L=k-\ell\\d_L\le H_0}}
          d_L^{-k}
       \le \frac{|\mathcal C|}{N^{kf}}\tau^{\ell+1}.
       \tag{G3}
\]
The local basis completion, trivial stabilizers and the level-index bound here refer to the actual levels in the construction.\formalnote{6.52}{52}

\subsection{Counting frames on a fixed lattice}
For clarity we distinguish the other estimate used with (G3). Fix a near term represented by a frame of \(n=\rho+1\) vectors with cell rank \(j\); let \(Z\) be its matrix of pairings with the standard frame. The leading \(j\)-minor is nonzero. The frame-counting calculation chooses the cell rank in at most \(b^{2f}\) ways and \(Z\) in at most \(b^{270n^3 f}\) ways. For fixed \(j,Z,L\) write \(a=M\det Z_{[j]}\) and let \(a'\) be the nonzero integral modulus for the isotropic part of the frame. Both have generalized index at most \(b^{2n^3f}\). Three successive counts are needed. The choices of integral lifts of \(L\) modulo admissible unipotent moves contribute at most \(d_L^{2\rho}\). With their pairings fixed, the \(X_0\) coordinates vary in a space of corank \(\ell\), contributing \([a]^{\ell n}\) residue classes together with the archimedean count. After those coordinates and the shared vectors are fixed, the remaining \(\rho\)-by-\(n\) entries determine the frame modulo moves when read modulo \(a'\), contributing \([a']^{n\rho}\). The archimedean and shared-vector estimates contribute \(b^{366n^4 f}\) in total. Section~\ref{sec:frame-lifts} derives these three bounds, including the congruence moduli and the gain for a common block.

To combine the counts, use that if there are \(A\) labels and each label has at most \(B\) elements, then there are at most \(AB\) elements. Apply this successively to the three stages and to \(j,Z\). Substituting \(d_L^2,[a],[a']\le b^{2n^3 f}\) gives an exponent of \(b\) at most
\[
 f\{2+270n^3+366n^4+2n^3(\rho+\ell n+n\rho)\}
       \le2000n^6(\ell+1)f .
\]
At fixed reduction label this is also the bound for one term over \(L\), using the identification with its frame modulo moves. The square bounds two terms. For \(\ell\le n^3\) its exponent is at most \(10^5 n^9 f\), hence the square is at most \(\tau\) at the prescribed scales. Multiplication by (G3) and by the number of reduction labels gives exactly (B8). Section~\ref{sec:cell-reduction} gives the common cell and the box for its frames. Section~\ref{sec:frame-lifts} proves the three lifting bounds, the level index and the gain relating \(d_L\) to the two orbital denominators.\formalnote{6.53}{53}


\section{Projection losses and finite-place shell bounds}\label{sec:bessel-projection}
\subsection{The near-identity orbital estimate and the two discarded projections}
Write \(\Gamma=G_0\times G_1\), where \(G_i\) is the product of residue-field groups at the chosen places above \(p_i\), and \(H=\bar R_0\) for the image of the Bessel subgroup. The source of \(P\) has dimension \(s_0=D_X\), and the identity second moment is \(\mathcal D_{\rm id}=|H|s_0\). Write \(P_s=P^*\rho(s)P\), and let \(r(s)\) be the operator obtained by retaining just the rational Bessel-subgroup terms in each unfolded block. Thus \(r(s)=0\) off \(H\). On \(H\) it acts by changing the definite source class and its Bessel coordinates, with the unipotent phase. The class permutation and coordinate action are unitary. More explicitly, multiplying the rational summation variable by a Bessel representative gives the corresponding permutation and unitary factor on the source index. At \(1\) the unipotent integral supplies precisely the normalizing covolume and freeness removes the definite stabilizer. Hence \(r(1)=1\), \(r(hs)=r(h)r(s)\), and the same change of variables in the full amplitude gives \(P_{hs}=r(h)P_s\), for \(h\in H\). Put \(D_s=P_s-r(s)\).\formalnote{6.54}{54}

Fix \(i\), and let \(\Omega_i=HG_i\). For \(s\in\Omega_i\), \(s\) meets the Bessel condition at all \(f\) places above the other prime \(p_j\) and their conjugates. Its meeting norm, the product of the norms there and at any additional meetings, is at least \(p_j^{2f}\ge(M/b)^f\). Also \(\Omega_i\) avoids \(\zeta H\) for each nontrivial admissible scalar \(\zeta\). The first frame coordinate has diagonal entry \(1\) at the other level, so such a scalar would reduce to \(1\) there, which detects the identity on this finite scalar group. Apply the square expansion of \S\ref{sec:orbital-sieve} to the amplitude outside the Bessel subgroup. Since \(kf\ge10\), the equal-link contribution on this set is at most
\[
 \left(\frac{b^{2f}}{(M/b)^f}\right)^k\mathcal D_{\rm main}
 \ \le\ \left(\frac{b^3}{M}\right)^{kf}
           \left(\frac{M}{b^6}\right)^{kf}\mathcal D_{\rm id}
 \ \le\ b^{-30}\mathcal D_{\rm id}.
\]
The unequal-link contribution is at most \(b^{-40}\mathcal D_{\rm id}\), and removing the cutoffs costs at most \(b^{-20}\sqrt{\mathcal D_{\rm id}}\) in the square-root norm. The triangle inequality and \(b\ge1000\) give
\[
             \sum_{s\in\Omega_i}\|D_s\|_{\rm HS}^2
                   \le b^{-12}|H|s_0.                     \tag{AB1}
\]
Indeed the left side sums exactly the squared norms of the off-Bessel amplitudes, since the Bessel blocks cancel before squaring. For \(h\in H\), \(D_h=r(h)(P^*P-1)\). Restricting (AB1) to \(H\) and cancelling \(|H|\) gives \(\|P^*P-1\|_{\rm HS}^2\le b^{-12}s_0\).\formalnote{6.55}{55}

For a central idempotent \(x=\sum_{t\in G_i}x(t)t\), right convolution \(F(s)\mapsto\sum_t x(t)F(st)\) on square-summable Hilbert-space-valued functions on \(\Gamma\) is an orthogonal projection. This follows by decomposing the finite unitary right translation representation into irreducibles, on each of which \(x\) acts as zero or one. Apply the contraction to \(F=1_{\Omega_i}D\). At every \(h\in H\) the result equals \(r(h)\sum_t x(t)D_t\), since \(ht\in\Omega_i\). Summing just over \(H\) on the output side and using (AB1) on the input gives
\[
              \left\|\sum_t x(t)D_t\right\|_{\rm HS}
                     \le b^{-6}\sqrt{s_0} .
\]
Its trace has modulus at most \(b^{-6}s_0\), by Cauchy--Schwarz against the identity matrix.\formalnote{6.56}{56}

For \(t\in G_i\) embedded with identity at the other prime, \(\operatorname{tr}r(t)=s_0\psi_U(t)\), where \(\psi_U\) is the phase on the reduced Bessel unipotent group, extended by zero. To see this, take a diagonal class block contributing to the trace. The definite component of the rational Bessel element belongs to the stabilizer of that class and reduces to the identity at the other prime. This stabilizer with principal level there is trivial. Indeed definiteness makes the stabilizer finite, and the first congruence subgroup at a completely split odd prime is torsion-free, as follows from expanding powers of a matrix \(1+p_j^aX\) with \(X\) nonzero modulo \(p_j\). Thus only a unipotent lift can contribute. Translation by such a lift multiplies the identity block by its phase. Summing the traces gives the formula.

Let \(\psi_U^\natural\) be the average of \(\psi_U\) over conjugation in \(\Gamma\). In the finite local orbital integral \(I(K)\) of \S\ref{sec:orbital-sieve}, use the product \(K\) of full integral groups at all the level places. Reduction maps \(K\) onto \(\Gamma\) with kernel \(K_M\). On the coset above \(s\) the inner integral is \(\mu(K)\psi_U^\natural(s)\), since reduction carries Haar probability to uniform probability. These cosets are disjoint, giving
\[
 \mu(K_M)\mu(K)^2\sum_{s\in\Gamma}|\psi_U^\natural(s)|^2
 \le I(K)\le\mu(K_M)\mu(K)(\mu(K)+2^{-n}).
\]
The upper inequality is the finite mass estimate with \(Q=2\) and \(t=n\), allowed by \(n+2\cdot10^9n^4\le k\). Since \(\mu(K)\ge2^{-n}\), the square sum is at most \(2\).\formalnote{6.57}{57}

For \(x=1-e_{p_i}\in\mathbb C[G_i]\), the discarded idempotent at \(p_i\), we have \(\Re x(1)\le b^{-8}/2\) by the finite type count in \S\ref{sec:B}. The identity coefficient in \(x^*x=x\) gives \(\sum_t|x(t)|^2=\Re x(1)\). As \(x(t)\) is a class function on \(G_i\), its pairing with \(\psi_U\) is unchanged by replacing the latter by \(\psi_U^\natural\). Cauchy--Schwarz therefore bounds \(|\sum_t x(t)\operatorname{tr}r(t)|\) by \(b^{-4}s_0\). Adding the off-Bessel trace bound gives \(\|\rho(x)P\|_{\rm HS}^2=\operatorname{tr}(P^*\rho(x)P)\le 2b^{-4}s_0\). In the group algebra, \(1-e_0=(1-e_{p_0})+e_{p_0}(1-e_{p_1})\). The two summands act as orthogonal projections with mutually orthogonal images. Their estimates give \(\|(1-e)P\|_{\rm HS}^2\le4b^{-4}s_0\), as required in \S\ref{sec:B}.\formalnote{6.58}{58}

\subsection{The finite-place shell estimate used in the trace and frequency sums}
Work over \(A=\prod_{v\in T}E_v\), with \(R=\prod_v\mathcal O_v\) of additive volume one. Write \(|c|\) for the modulus of multiplication on \(A\), \(g=M\) for the local level, and \(n=\rho+1\). A self-conjugate scale \(\alpha\in R\cap A^\times\) defines size \(j\) by \(\alpha^jz\in R\) and has loss \(L=2/|\alpha|\). All common denominators below are ideals in this product ring. Put \(m=3\rho+5\). The sum \(\zeta_R(m)=\sum_J[R:J]^{-m}\) over finite-index ideals satisfies \(\zeta_R(m)^n\le1+2^{-n}\). Indeed its Euler factors occur among those of \(\zeta_E(m)\). For \(m\ge4\) the Euler product gives \(\log\zeta_E(m)\le 2[E:\mathbb Q]2^{2-m}\), by bounding the factors over a rational prime by the corresponding factors for \(\zeta(m)^{[E:\mathbb Q]}\) and comparing the sum over integers \(\ge2\) with its first term and a decreasing integral. Here \(m=n+2\rho+4\) and \(n[E:\mathbb Q]\le2^{2\rho}\), so \(n\log\zeta_E(m)\le2^{-n-1}\le\log(1+2^{-n})\).\formalnote{6.59}{59}

We describe the finite version of the one-frame estimate in \S\ref{sec:shell-measure}. At a pivot, stable pairs \(w,w'\) agree modulo the earlier frame directions. Their transported coordinates at the two points are integral and agree modulo \(g\). The possible values of \(\langle e_a,w\rangle\) form a fractional ideal \(d_aR\): these conditions define a full lattice in the linear subspace of admissible pairs, and the indicated nonzero coordinate projects it to a full lattice in \(A\). This may be checked one discrete valuation ring at a time, by clearing denominators and projecting in an adapted basis. Its reciprocal volume \(D_a=|d_a|^{-1}\) is the depth. At a near point the coordinate is integral by moving to the integral box with a triangular isometry fixing this coordinate. The difference of the two Hermitian pairings, divided by \(g\), is integral for close pairs, by expanding the pairings using their congruences.

Here is the replacement for integration of the reciprocal area of a disc. At an active relative position \(z=y^{-1}uyk\), with \(u\) Bessel unipotent and \(k\in K_M\), forget \(3\rho\) free directions. The remaining ones can be taken as integral common vectors, primitive together. Indeed impose the \(\rho\) chart equations, the vanishing of the last \(\rho\) rows of the transported columns, and that of the first \(\rho\) rows after applying \(u^{-1}-1\). Integral column operations at each valuation ring put these at most \(3\rho\) rows on the omitted columns. The shape of \(u\) then fixes every remaining transported column, and \(k\equiv1\bmod g\) supplies closeness. Let \(Y\) collect the remaining columns of the resulting matrix in \(\mathrm{GL}(R)\). For a test vector \(v\), let \(tR\) be the ideal of the entries of \(Y^*v\). Products with common vectors contain \(\bar t R\). Conversely, if \(\langle cv,w\rangle\) is integral for every integral common vector \(w\), then testing the columns of \(Y\) gives \(c(Y^*v)_i\in R\) for each \(i\). Hence \(c\in t^{-1}R\). Thus \(\Phi(v)=|t|^{-1}\) controls both volume factors in the two pivot integrals.

There are \(s+m+5\) such columns, where \(s\) is the volume exponent of the finite frame. For integral, uniformly distributed coordinates, the event that their content lies in a particular ideal \(J\) has probability \([R:J]^{-(s+m+5)}\). Sum over the possible contents and bound the event of equality by that of inclusion. The \(s\)-th moment of \(\Phi\) is at most \(\zeta_R(m)\). A coordinate change in \(\mathrm{GL}(R)\) preserves the box. Rescaling a box of size \(j\) to size zero costs \(|\alpha^j|^{-(3\rho+m+5)}\) after cancellation against the \(s\)-th power, giving a bound \(L^{99n(j+1)}\).\formalnote{6.60}{60}

The remaining block has \(\ell=2m+4\) coordinates. To put it in hyperbolic form,
choose \(\epsilon\in R^\times\) with \(\epsilon\bar\epsilon=-1\). Such a unit exists
because each finite place of \(E/E^+\) is split or unramified
(\S\ref{sec:split-vacuum}) and the norm on units is surjective in either case.
A primitive sixth root \(\tau\in R\) has \(\tau+\bar\tau=1\).
For each pair of orthonormal basis vectors \(e,f\), the vectors
\(e+\epsilon f,\ \bar\tau e-\tau\epsilon f\) are isotropic, have mutual pairing \(1\),
and have determinant \(-\epsilon\) in this basis. Reserve one of these planes for the quadric chart. Fix one half \(v\in R^{m+1}\) of the remaining
hyperbolic pairs. Almost surely its entries generate a finite-index ideal \(cR\).
Write \(c=\sum_i r_i v_i\) with \(r_i\in R\). Translating the coordinates opposite \(v\)
by \(c\bar r_i t\), \(t\in R\), preserves the box and changes the Hermitian norm by
\(c\bar c(t+\bar t)\). The integral trace is onto by the choice of \(\tau\).
After multiplication by the integral skew unit, these shifts comprise \(c\bar c\) times the
integral skew elements, a subgroup of measure \(|c|\). Thus the conditional mass on any one quadric,
multiplied by \(|c|\), is at most the mass integrated over all quadrics. The latter is the
ambient volume by disintegration. The resulting weight is \(|c|^{-1}=[R:cR]\).
For each finite-index ideal \(I\), the probability that \(cR\subseteq I\) is
\([R:I]^{-(m+1)}\). Summing the index over possible contents gives
\[
 \mathbb E[R:cR]\le\sum_I[R:I]\Pr(cR\subseteq I)=\zeta_R(m).
\]
This bounds the mass of the unit box uniformly in the value of the form.
Scaling to \(u^{-1}R^\ell\) multiplies ambient volume by \(|u|^{-\ell}\). The skew parameter
scales by \((u\bar u)^{-1}\), of modulus \(|u|^{-1}\) on the skew line. The fibre bound is
therefore \(\zeta_R(m)|u|^{-(\ell-1)}\).
Together with the preceding moment bound this supplies the two finite-place estimates
in \S\ref{sec:shell-measure}. The unweighted pivot mass at unit scale is at most
\(V=2^{99\cdot13n}\), from the box of size \(7\) in the quadric chart, whose free
integral volume at this scale is one.\formalnote{6.61}{61}

For each position \(x\), let \(J=cR\) be the common denominator ideal of \((x-1)/g\), with \([R:J]=|c|^{-1}\). Its generator is obtained by taking the largest denominator valuation in each component. Let \(P=2^h\) be the first power of two at least this index. There are at most \(2P^{9n}\) such ideals, by the zeta bound and \(m<9n\). For each ideal its entire shell lies in the position box at scale \(\bar c c\), with loss \(2[R:J]^2\le LP^2\).

Suppose every remaining stable product ideal at an active point has volume greater than \(\sigma^2\), with \(0<\sigma<1\). It contains an integral generator or an integral element \(d_a\) with \(|d_a|>\sigma^2\): on the active supports the product ideals are contained in \(R\), by the integral box and the Bessel triangular action. Take \(C=\prod_{a\ge j}d_a\bar d_a\), so \(|C|\ge\sigma^{4n}\). At scale \(\alpha^{10n}C\), each \(d_a^{-1}\) has size one. The comparison on \(A_a^\perp/A_a\) in \S\ref{sec:shell-measure} applies to these common pairs with differences divided by \(g\). Each preceding-head inverse has size one in the original scale and the cofactors of its frame have size at most \(10n\). The bottom comparison costs size \(9\), and the dual-vector step on passing to a preceding floor adds at most \(32\). Thus \((\alpha^{10n}C)^{99n}\) clears \((x-1)/g\). Since \(c\) divides every such denominator,
\[
       \sigma^{396n^2}\le |C|^{99n}
                     \le |c|\,L^{990n^2}.
\]
Consequently \(L^{990n^2}\le \sigma^{400n^2}[R:J]\) forces a remaining depth at least \(\sigma^{-2}\). At floor zero, at unit scale, there is a second implication. If each stable product contains a unit of \(R\), the same comparison gives \((x-1)/g\) integral, hence \(x\in K_M\). If every proper denominator ideal has index at least \(Q\), then at every active point outside \(K_M\) some depth is at least \(Q\).\formalnote{6.62}{62}

\clearpage
\paragraph{Summation over the shells.}
Write \(e^s\), \(0<e\le1\), for the pointwise gain relative to the square root of the product of weighted moments, taking \(e=1\) for the absolute majorant. The two pivot integrations give the shell bound
\[
 V h_0^s L^{2\cdot10^6n^3}P^{10^5n^2}\,\mathfrak m_Z
\]
after multiplication by the head density. Here \(\mathfrak m_Z=\mu(Z)\mu(K_M)/V_{U,T}^2\). Indeed the per-ball loss in \S\ref{sec:shell-measure} at \(L'\le LP^2\) contributes at most \(P^{50000n^2}\), and \(2P^{9n}\) balls fit within the displayed bound. The shell normalization is at most \(\mathfrak m_Z\) by orthogonality of distinct translates of the level vector. Take \(h_0=e\) up to \(P=L^{2400n^2}\). Above this threshold take \(h_0=\min(e,P^{-1/(800n^2)})\). The latter power meets the deep-vector criterion since \([R:J]\ge P/2\); the depth gain comes from the pivot Cauchy--Schwarz argument in \S\ref{sec:shell-measure}. When \(s\ge r+9\cdot10^7n^4\), factor out \(e^r\). On a remote shell the spare powers contribute at most \(P^{-112500n^2}\), absorbing \(P^{1+10^5n^2}\). On a close shell the latter is at most \(L^{2400n^2(1+10^5n^2)}\). Thus the shell costs at most \(P^{-1}V e^r L^{3\cdot10^8n^4}\mathfrak m_Z\). Sum \(\sum_h2^{-h}=2\) and absorb \(2V\) into \(L^{10^9n^4-3\cdot10^8n^4}\). The rank hypothesis \(r+10^8n^4\le k\) supplies the required bound on \(s\) by its definition. For the floor-zero majorant outside \(K_M\) both pivots coincide, so a depth \(Q\) supplies \(e=Q^{-1}\) by applying the gain at both pivots. Here \(L=2\). Taking \(r=t+10^9n^4\) absorbs \(2^{10^9n^4}\) and yields the error \(Q^{-t}\mathfrak m_Z\) used above. To obtain \(\mu(K)\ge2^{-n}\) at the selected level places, the integral frame chart at each of the \(n\) floors has mass at least \(\operatorname{vol}(R^\times)\). The union of the residue nonunit conditions has volume at most \(\sum_{v\in T}q_v^{-1}\le1/2\) by the size of the level primes. Multiply the \(n\) lower bounds.\formalnote{6.63}{63}

\clearpage
\subsubsection*{The finite Fourier transform}

Finally consider the finite transform in (B4). With the two heads \(p,q\) fixed, take common multiples \(\beta_p,\beta_q\) of their integral successive products by multiplying all those successive products in each head. They also clear the inverses of the single entries, since an entry is a quotient of successive products. Set \(\alpha=g\,\overline{\beta_p\beta_q}\,\beta_p\beta_q\) and \(b_0=\alpha^{200n^2}\). For a free frequency \(\xi\) choose an element \(\eta\) of the same valuation as \(\gcd(\xi,b_0)\) in every component and congruent to \(\xi\) modulo \(b_0R\). One may use \(\xi\) where its valuation is strictly smaller than that of \(b_0\), and a generator of \(b_0R\) elsewhere. The phase is unchanged on arguments of size \(200n^2\). Put \(w=b_0/\eta\). For an inverse-head difference \(\xi=(p_a^{-1}-q_a^{-1})/g\), take \(\eta=\xi,\ w=b_0\), since \(b_0\xi\in R\). In both cases \(w,\xi w,\xi/\eta\) are integral. With \(e=|\eta|^{-1/2}\), the free-coordinate Fourier estimate gives the gain \(\min(1,e)^s\) (the Plancherel constant is one in this normalization). Write \(d=[R:R\cap \xi^{-1}R]\). From \(wR\subset R\cap\xi^{-1}R\) one gets \(|w|\le d^{-1}\), while \(\min(1,e)^2\le |w|/|b_0|\). Therefore
\[
     \min(1,e)^{t+12}
          \le |\xi|^{-t/2}d^{-6}|b_0|^{-6}.
\]
In the shell bound the extra \(|b_0|^{-6}\) costs at most \(L^{1200n^2}\). Since \(L=2/(|g||\beta_p\beta_q|^2)\), all resulting powers \(L^E\) with \(4E\le K\) are paid by \((2/|g|)^K\) and the two head weights \(|\beta_p\beta_q|^{-K/2}\). The assumed \(K\ge10^{10}n^4\) permits this. This proves the finite transform bound, with even six denominator powers before the weakening to three in (B4).\formalnote{6.64}{64}


\editionchapter{Archimedean Bessel periods}\label{ch:bessel}

The second moment involves integrals of matrix coefficients at each real place. We construct a nonzero Bessel function \(W\) at the prescribed weight and compute the integral of the squared absolute value of its correlation with its translates. Estimates away from the flags control the error when \(W\) is cut off. The chapter begins with a compact calculation of the normalization of these integrals.

\section{The real-place orbital estimate}
\label{sec:local-orbital}

Fix one complex place of \(E\) over a real place of its totally real
subfield. The group \(G\) preserves a Hermitian form with \(\rho\)
hyperbolic planes, \(k\) positive middle directions and one additional
negative vector. Write \(n=\rho+1\), \(r=k+2\rho+1\), and let \(U\) be
the Bessel unipotent subgroup. The second moment in
\S\ref{sec:orbital-sieve} requires two integrations. At fixed \(g\) one first pairs a Bessel function \(W(y)\) with its translate \(W(yg)\) and integrates in \(y\) over \(U\backslash G\). The resulting correlation is a function of \(g\). Its squared absolute value is then integrated over \(G\). For a Bessel function of unit \(L^2\)-norm, the scale \(1/d\) comes from discrete-series orthogonality. For a unit vector \(v\) in the discrete series \(\pi\), one has \(\int_G|\langle v,\pi(g)v\rangle|^2\,dg=1/d\), where \(d\) is the formal degree for the measure \(dg\). This is the discrete-series orthogonality relation \cite[Theorem~B, (1)]{upper-moore-wolf}, with probability Haar measure on the compact centre and the corresponding quotient measure. We also need decay
after taking Fourier transforms in the frame coordinates. The
quantitative bounds are \((1+b^{-10})/d\) for the normalized squared correlation integral and an
exponent \(k-n^7\) for the Fourier estimate.

\subsection{The coefficient of the full Bessel function and its cutoff}

\begin{example}[The mass of a matrix coefficient on \(U(2)\)]\label{ex:compact-coefficient-mass}
Give \(U(2)\) Haar measure one. For the vector-valued function
\(W(y)=y e_1\) the correlation of \(W\) with its right translate is
\[
 \int_{U(2)} \langle y e_1, y g e_1\rangle\,dy
       =\langle e_1,g e_1\rangle=g_{11}.
\]
By Haar invariance the entries \(g_{11}\) and \(g_{21}\) have the
same mean square. Since \(|g_{11}|^2+|g_{21}|^2=1\), one has
\(\int_{U(2)}|g_{11}|^2\,dg=1/2\), the reciprocal of the dimension
of the standard representation. For \(h>0\), replacing \(W\) by \(\sqrt h\,W\) multiplies the correlation at \(g\) by \(h\), so the integral of its squared absolute value becomes \(h^2/2\). In the discrete-series calculation below the formal
degree \(d\) takes the place of \(2\).
\end{example}

Let \(c_\lambda\) be the normalized discrete-series coefficient,
\(c_\lambda(1)=1\), and put \(k_\lambda=d\,c_\lambda\).
The Bessel construction supplies a nonzero average of \(c_\lambda\)
against the Bessel character and an irreducible unitary type of
\(U(k)\). Normalize the resulting function \(W:G\to\mathbb C^I\) to
\(\|W\|_{L^2(U\backslash G)}^2=h=|I|\). The unipotent acts through
a character of absolute value one, so its norm and the inner
product of simultaneous right translates descend to the quotient.
With the convention
\[
        (W*k_\lambda)(y)=\int_G k_\lambda(g^{-1})W(yg)\,dg ,
\]
the operator \(W\mapsto W*k_\lambda\) acts by a rank-one projection on each copy of the discrete series of parameter \(\lambda\). It reproduces \(W\), and compressing a right
translation by this kernel multiplies the projection by \(c_\lambda(g)\).
In a discrete-series model the latter assertion is the identity
\(P_v\pi(g)P_v=\langle v,\pi(g)v\rangle P_v\) for a unit vector \(v\).
It applies on each copy of the represented summand. Thus
\[
 C_{W,W}(g):=\int_{U\backslash G}\langle W(y),W(yg)\rangle\,dy
       =h c_\lambda(g),\qquad
 \int_G|C_{W,W}(g)|^2\,dg=\frac{h^2}{d}.
 \tag{L1}
\]
The last equality uses \(d^{-1}=\int_G|c_\lambda|^2\). Independently,
the matrix \(\int W(y)W(y)^*\,dy\) commutes with \(U(k)\), by the
left transformation law. Its trace is \(h\), so irreducibility makes
it the identity. In particular the \(h\) coordinates are orthonormal.\formalnote{7.1}{1}

Write \(\ell(y)=\inf_{a\in U\,U(k)}(1+\|ay\|_{\rm op})\), put
\(R=b^{6n+2}\), and define
\(F(y)=\max(0,\min(1,2-2\ell(y)/R))W(y)\).
Where this cutoff changes \(W\), the scaled length
\(\max(1,\ell(y)/b^{6n})\) is at least \(b\) for \(b\ge2\).
Use \(f=|\Phi|\), the number of archimedean places of \(E\), as the decay parameter. The Bessel average supplies a function
\(0\le A(y)\le a_0\) with
\[
 \max(1,\ell(y)/b^{6n})^{3r^3f}|W(y)|\le A(y),\qquad
 a_0\int_{U\backslash G}A(y)\,dy\le b^{400r^2}h.
\]
Squaring the cutoff difference, using \(A^2\le a_0A\), and integrating
therefore gives \(\|F-W\|_2^2\le b^{400r^2-6r^3f}h
\le b^{-2r^3f}h\) since \(r\ge100\), \(f\ge1\).
For later use \(\|F\|_2\le\sqrt h\), and its covariance matrix
\(\int FF^*\) is bounded by the identity. In particular its squared
Hilbert--Schmidt norm is at most \(h\).\formalnote{7.2}{2}

\subsection{The local mass, with its three Haar factors}

Let \(\lambda_H\) be the ratio of the orbital-coordinate Haar measure
to \(dg\). Put \(V_j=\prod_{i=1}^j \frac{2\pi^i}{(i-1)!}\), the product
of the areas of the unit spheres in \(\mathbb C,\ldots,\mathbb C^j\).
In rank \(j\), division of the Cayley-coordinate density by \(V_j\)
gives the Haar normalization with compact dual \(U(j)\) of mass one.
Unfolding the unipotent coordinates identifies their
quotient measure with \(\lambda_H\,dy\) on invariant integrands. The
coordinate calculation gives \(0<\lambda_H\le V_r/V_k\).
The normalizing constant of the floor forms is
\[
                 c_F=\frac{V_k/V_r}{(\lambda_H h)^2}.
\]
At floor zero the two correlations are complex conjugates. Each uses
the quotient measure \(\lambda_H\,dy\), while the position integral
uses \(\lambda_H\,dg\). The resulting density is
\(c_F\lambda_H^3|C_{F,F}(g)|^2\,dg\). The bound
\(c_F\lambda_H^3h^2=(V_k/V_r)\lambda_H\le1\) will turn (L1)
into precisely the required formal-degree factor.\formalnote{7.3}{3}

Define the near positions by
\[
  \max_{i,j}|(g-1)_{ij}|
     \le (b^{6r}R^{11n})^{99n}.
\]
On this set \(\|g\|_{\rm op}\le b^{700nr}\). Indeed an \(r\)-by-\(r\)
matrix of entries bounded by \(Z\) has norm at most \(r^2Z\).
Apply \(r^2+1\le R^2,\ R\le b^{7n}\) and \(10^9n^3\le r\).
For \(T\ge1\) the parameter and Haar estimates give
\(d\,\operatorname{vol}\{\|g\|\le T\}\le T^{9r^2}b^{99r^2}\).
Indeed on this ball \(1\le (T/2)^{9r^2}(2/\|g\|)^{9r^2}\), so the
volume estimate follows by integrating the right-hand side.
Right translation is unitary on the quotient. Cauchy--Schwarz
applied to the two terms of
\(C_{F,F}-C_{W,W}=C_{F-W,F}+C_{W,F-W}\) gives the pointwise
bound \(2b^{-r^3f}h\). Here \(n=m/2\), whereas \(f=3l\,m^{22}\varphi(q)\) by \S\ref{sec:data}. Thus \(f\ge4000n\) for sufficiently large \(x\), and the near-position volume
absorbs its squared size, since
\[
 4 b^{-2r^3f+6300nr^3+99r^2}\le b^{-24}.
\]
Minkowski's inequality and (L1) bound the near contribution by
\((1+b^{-12})^2h^2/d\) before multiplication by \(c_F\lambda_H^3\).

For the complement put \(\Theta=\pi^{(2\rho+1)k}V_k/V_r\) and
\(s=k-3n-6\). The remote-position shell bound established below
gives, including the normalization, at most
\[
   \Theta R^{10^9n^4+24k}
                  (b^{6r})^{10^8n^4-s}.
\]
The Weyl degree estimate \(d\Theta\le b^{4r^2}\) and \(R\le b^{7n}\)
make this at most \(b^{-12}/d\). Explicitly the exponent inequality is
\[
 4r^2+7n(10^9n^4+24k)+6r\,10^8n^4+12\le 6rs.
\]
It follows on substituting \(k=s+3n+6,\ r=s+5n+5\) and
\(s\ge10^{10}n^4\). Adding the two contributions proves\formalnote{7.4}{4}
\[
                  |J_0|\le\frac{(1+b^{-12})^2+b^{-12}}d
                          \le\frac{1+b^{-10}}d.
 \tag{L2}
\]

\subsection{Fourier gain in each free complex coordinate}

The additive character is \(\operatorname{wave}(z)=e^{z-\bar z}\).
For \(m_0\) free complex coordinates the essential Fourier integral
is
\[
 \int_{\mathbb C^{m_0}}\!\!\int_{\mathbb C^{m_0}}
       e^{\,2i\operatorname{Im}(\xi\, x^*y)}
                    g(x)\overline{h_1(y)}\,dx\,dy .
\]
In one complex coordinate write \(x=a+ib\), \(y=c+id'\).
The real bilinear form \(\operatorname{Im}(\bar xy)=ad'-bc\)
is ordinary dot product after a rotation. For complex \(\xi\ne0\)
we may rotate once more and dilate by \(|\xi|\). With the usual
unitary transform \(e^{-2\pi i x\cdot\eta}\), the new real frequency
is a rotation of \(y\) dilated by \(|\xi|/\pi\). By Plancherel and
Cauchy--Schwarz, the square of the displayed integral is at most
\[
          (\pi^2/|\xi|^2)^{m_0}\|g\|_2^2\|h_1\|_2^2.    \tag{L3}
\]
Both powers of \(|\xi|\) in a single coordinate come from its two
real dimensions. The two uses of (L3) involve the diagonal Gram
entry, with phase \(\operatorname{wave}(\xi\langle v,w\rangle)\),
and a pair of different last entries \(p_i,q_i\). In the latter case
the Heisenberg translations of the two lifted frame vectors produce
the phase with frequency \(p_i^{-1}-q_i^{-1}\), up to sign.
Its absolute bound is identical.

On the relevant frame varieties the solved-block coordinates
leave \(m_0=k-5\) free complex coordinates in the resulting Fourier
bound. The remaining fiber measures and the other coordinates
are bounded by box volumes. Fubini followed by (L3) on the
free block and Cauchy--Schwarz on the other variables gives, for the
inner absolute correlation \(\phi(x)\) at a position \(x\),
\[
  \phi(x)^2\le(\pi^2/|\xi|^2)^{k-5}\mathcal Q(x).
 \tag{L4}
\]
Here \(\mathcal Q\) is the product of the two weighted second moments
of the frame correlations, including the solved-block volume
constants. The local quadric disintegration and the separation of
the frame variables justify precisely this reduction on each floor.\formalnote{7.5}{5}

\subsection{Loss of powers in the position integral}

For a pair of heads choose \(L\ge R\) large enough to bound the
inverses of all their partial products. Write \(A=L^2\) for the
associated volume parameter. Let \(H\) be the product of both head
weights, each being the product of
\(\max(1,|p_0\cdots p_j|^{-1})^{n^6}\) over the specified entries.
The weight is set to zero if a partial product exceeds \(b^{60n^2}\).
In that case the support estimate makes the floor form zero, since
the cut bounds the leading minors by \(R^{9n}\le b^{60n^2}\).
Otherwise, taking \(L\) to be the maximum of \(R\) and the inverse
partial-product moduli gives \(L^{n^6}\le R^{n^6}H\).

Set \(t=\max(1,R^{24}|\xi|)\), \(e=t^{-1}\) and
\(V=(\pi R^{24})^kR^{9999n}\). For \(t>1\), (L4) implies
\(\phi^2\le(Ve^s)^2\mathcal Q\), since
\[
   (\pi/|\xi|)^{k-5}=(\pi R^{24})^{k-5}t^{-(k-5)}
          \le (\pi R^{24})^k t^{-s}.
\]
When \(t=1\), the estimate for the absolute majorant without
oscillation gives the same bound with \(e=1\).
Partition the position variable according to the least dyadic
height \(P=2^j\ge1\) bounding all entries of \(x-1\). We spell
out the shell summation. Suppress the common shell normalization
\(\mathcal N\), which at this place is at most
\(\pi^{2(k-5)\rho}\lambda_H^2 h\). The power of \(\pi\) comes from
the successive Fourier kernels. The remaining pairing norm is
bounded by \(\lambda_H^2 h\) using \(\int FF^*\le1\).

On a shell, the one-frame volume bound and Cauchy--Schwarz give,
including the head density, the upper bound
\[
       Vg^s A^{2\cdot10^6 n^3}P^{10^5n^2}\,\mathcal N.
 \tag{L5}
\]
One may take \(g=e\) from (L4). When \(P>A^{2400n^2}\) the frame
geometry also allows \(g=\sigma=P^{-1/(800n^2)}\), since
\(A^{1200n^2}\le\sigma^{400n^2}P\). In that range at least one
remaining vector has depth at least \(\sigma^{-2}\) on the active
frame chart, giving the gain \(\sigma^s\) in its majorant. Thus use
the smaller of \(e,\sigma\) on the remote shells. The factor
\(A^{2\cdot10^6 n^3}P^{10^5n^2}\) in (L5) is a geometric
one-frame estimate. The summation that removes its \(P\)-dependence
is elementary.

Reserve \(9\cdot10^7n^4\) powers of \(g\). For a remote shell
these powers contribute at most
\(\sigma^{9\cdot10^7n^4}=P^{-112500n^2}\), enough for
\(P^{-(1+10^5n^2)}\). For a near shell,
\(P^{1+10^5n^2}\le A^{2400n^2(1+10^5n^2)}
\le A^{29\cdot10^7n^4}\).
In both ranges (L5) is at most
\(P^{-1} V e^{s-10^8n^4}A^{3\cdot10^8n^4}\mathcal N\).
The sum of \(P^{-1}\) is two. Since \(A\ge2\), this factor fits
within \(A^{4\cdot10^8n^4}\). Restoring normalization gives
\[
 h\,|c_F\,\mathrm{floorForm}|
    \le\Theta L^{10^9n^4}R^{24k}t^{10^8n^4-s}.
 \tag{L6}
\]
Indeed \(c_Fh\mathcal N\le(V_k/V_r)\pi^{2(k-5)\rho}\).
The \(\pi^k\) in \(V\) raises this at most to \(\Theta\), and
\(R^{9999n}A^{4\cdot10^8n^4}\le L^{10^9n^4}\).
The same summation bounds the absolute majorant outside the position
ball of scale \(t L^{11n}\): every active frame there already has
a vector of depth at least \(t^2\). Setting \(L=R,t=b^{6r}\)
recovers the remote estimate used in (L2).\formalnote{7.6}{6}

Finally put \(K=k-n^7\). The parameter regime gives
\(10^9n^4\le n^6\) and \(10^8n^4+K+6\le s\). Rewrite the right
side of (L6) using these inequalities as an upper bound
\[
 \Theta H R^{n^6+24n^7}
           (R^{24}/t)^K t^{-6}
 \le \Theta H b^{2000n^9}
            |\xi|^{-K}\min(1,|\xi|^{-1})^6.          \tag{L7}
\]
Here \(R^{24}/t\le|\xi|^{-1}\) and
\(t^{-1}\le\min(1,|\xi|^{-1})\). This establishes the local
transform and unequal-entry bounds in Section~\ref{sec:orbital-sieve}
with the stated frequency exponent. The stronger left side in
(L6) contains \(h\ge1\), which can now be dropped.

For the top form, the orbital integral of the cut vectors at a
Weyl element is multiplied by
\[
 \frac{\sqrt{c_G}}h
       \big(\max(0,|p_\rho|^2-1)\big)^{(k-1)/2}
       \prod_{a<\rho}|p_a|^{k+1+2(\rho-a-1)} .
\]
Here \(c_G=V_k^2/(V_rV_{k-1})\) is the open-cell Haar constant. Pairing two such
tensors gives a positive semidefinite kernel. The factor
\((|p_\rho|^2-1)^{k-1}\) on a common last entry is the sphere
density on the last quadric. Integrating in that entry, with
the chart Jacobian, recovers \(c_F\) times the preceding floor form.
The other floors descend by integration of the common Gram entry.
Consequently (L2) and (L7) apply to the same system whose top
kernel occurs in the Siegel step.\formalnote{7.7}{7}

Section~\ref{sec:bessel-period} constructs the nonzero Bessel average
and gives its decay majorant, including a lower bound that controls
the cost of normalization. That calculation uses the lowest-type
coefficient estimates stated there. The Haar and quadric Jacobians
enter the explicit inequalities above. Section~\ref{sec:shell-measure}
derives the one-frame volume and depth estimates in (L5). In particular (L3) accounts for the Fourier
constant, and the shell summation accounts for the loss between
\(k-5\) free coordinates and the exponent \(k-n^7\) in (L7).\formalnote{7.8}{8}


\section{Producing a nonzero Bessel average at a real place}
\label{sec:bessel-period}
The function normalized in \S\ref{sec:local-orbital} comes from an integral with a varying complex phase. A nonzero value of the integrand near the identity does not by itself give a nonzero integral. We construct a position in which the two representation-theoretic phases cancel the Bessel character and bound the remaining cancellation. We also keep the lower bound on the period: its size controls the normalization of the cutoff.

\subsection{An indefinite Hermitian space made from two sets of real nodes}
We seek an isotropic flag in the larger space and an isometric copy of the smaller space perpendicular to its isotropic directions. In these coordinates the first derivative of the phase can be read off from matrix entries. Let \(X\) index the common positive coordinates of weight zero in the two spaces. On the complementary coordinates take integers
\[
 u_0>\cdots>u_{L-1},\qquad t_0<\cdots<t_{L'-1},
 \qquad u_i<t_j,\qquad L=L'+2d+1.
\]
The signs in each complement alternate, starting with \(+1\). In the application the \(u_i\) lie in \([-3B,-B]\), the \(t_j\) in \([0,B']\), and \(B\ge1\). Define \(p(z)=\prod_i(z-u_i)\), \(q(z)=\prod_j(z-t_j)\). To a real polynomial \(F\) associate the vector
\[
     [F]_i=\frac{F(u_i)}{\sqrt{|p'(u_i)q(u_i)|}} .
\]
Extend the scalar product with diagonal entries \((-1)^i\) to a Hermitian form. Since the signs of \(p'(u_i)\) and \(q(u_i)\) are \((-1)^i\) and \((-1)^{L'}\), for real \(F_1,F_2\) the form is
\[
 ([F_1],[F_2])=(-1)^{L'}\sum_i
       \operatorname*{Res}_{z=u_i}\frac{F_1(z)F_2(z)}{p(z)q(z)}.
\]
Multiplication of a vector by the diagonal weight \(D=\operatorname{diag}(u_i)\) takes \([F]\) to \([zF]\).

Take \(v_a=[z^a q]\), \(0\le a\le d\). These vectors are independent: dividing a vanishing linear combination pointwise by \([q]_i\) would give a polynomial of degree at most \(d<L\) vanishing at all the \(u_i\). For every polynomial \(H\) of degree at most \(L-1\),
\[
        \sum_i H(u_i)/p'(u_i)=[z^{L-1}]H .
\]
This is the leading coefficient in the Lagrange interpolation formula. Apply it to \(H=z^{a+b}q\). We get \((v_a,v_b)=0\) except when \(a=b=d\), where the value is \((-1)^{L'}\).

Put \(q_j=q/(z-t_j)\) and \(w_j=\sqrt{|p(t_j)/q'(t_j)|}\,[q_j]\). The inner product of \([q_j]\) with \(v_a\) is zero by the same interpolation, since \(\deg(z^a q_j)<L-1\). For \(j\ne j'\), the finite poles of \(q_jq_{j'}/(pq)\) are all at the \(u_i\). Its decay at infinity gives sum of residues zero there. For \(j=j'\) there is one additional pole at \(t_j\), with residue \(q'(t_j)/p(t_j)\). Because \(p(t_j)>0\) and \(q'(t_j)\) has sign \((-1)^{L'-1-j}\), the normalization gives \((w_j,w_{j'})=(-1)^j\delta_{jj'}\). Consequently the map \(I\) with columns \(w_j\), extended as the identity on \(\mathbb C^X\), is an isometry orthogonal to the \(v_a\). Extend \(D\) by zero on \(X\). The identity \(zq_j=t_jq_j+q\) proves
\[
 Dv_a=v_{a+1}\quad(a<d),\qquad
 Dw_j=t_j w_j+\sqrt{|p(t_j)/q'(t_j)|}\,v_0 .
\]
In particular \(I^\dagger D I\) is exactly the second weight \(D'\); the dagger denotes the adjoint for the signed Hermitian forms.

All node differences are integers of absolute value at least one. Set \(D_0=3B+B'\). The coordinates of \(D^a v_0\), for \(a\le2d\), have absolute value at most \(D_0^{a+L'/2}\le D_0^L\), by cancelling one square root of \(|q(u_i)|\). The squared modulus of a coordinate of \(w_j\) is at most \(|p(t_j)q_j(u_i)|\le D_0^{2L}\), using \(|q'(t_j)|,|p'(u_i)|,|u_i-t_j|\ge1\). This bounds the cost of the change of position. For the Bessel group in question \(L=2n,\ L'=1,\ d=n-1\). The last vector has signed square norm \(-1\), and the smaller compact group's remaining weight is the one nonnegative integer \(t_0\).\formalnote{7.9}{9}

\subsection{The phase at a stationary flag}
Let \(S\) preserve the position, with \(g-1\) sending \(v_a\) into the span of its predecessors and \(g'\) its induced action on the smaller space. The successive first-superdiagonal entries of \(g\) on the flag add under multiplication, as is seen from multiplying two upper unitriangular matrices. Write \(\sigma(g)\) for their sum. In this position the trace identities are
\[
 \operatorname{tr}(g-1)=\operatorname{tr}(g'-1),\qquad
 H(g-1,g'-1)=\sigma(g)-\overline{\sigma(g)}=:\phi(g),
 \quad H(A,A')=\operatorname{tr}(DA)-\operatorname{tr}(D'A').
\]
The flag trace calculation uses the dual isotropic flag as well: its first-superdiagonal contribution to the weighted trace is the negative conjugate of the one on the original flag. The induced middle block contributes \(\operatorname{tr}(D'(g'-1))\). We use the signed flag trace identity here. Thus \(\phi\) is imaginary and additive. The image of the joint map \(s\mapsto(g,g')\) contains \(U(X)\) acting diagonally on both common blocks.\formalnote{7.10}{10}

There is an exact identity for logarithms. Write \(A=(g-1,g'-1)\) and compute powers componentwise. For every integer \(j\ge0\),
\[
             H\bigl((1+A)^j-1\bigr)=j\phi(g).
\]
Taking \(j=2,3,\ldots\) in turn and expanding by the binomial theorem gives \(H(A^j)=0\) for all \(j\ge2\). For \(\|A\|<1\) the logarithm series now gives
\[
            H(\log g,\log g')=\phi(g).
\]
Thus the scalar phase removed from the product of lowest-type coefficients at \(g,g'\) is exactly the Bessel character, even though the individual coefficients have nonlinear phases.\formalnote{7.11}{11}

\subsection{A cone estimate for the two coefficients}
Let \(\lambda,\nu\) be the Harish--Chandra parameters in the two groups of ranks \(r,k\). Their doubles are \(2\lambda,2\nu\). Write \(c_\lambda,c_\nu\) for the normalized coefficients. Put
\[
 Z(s)=c_\lambda(g)\overline{c_\nu(g')}\,e^{-\phi(s)},\quad
 N_0=(r+k)^2,\quad T=N_0\log\beta .
\]
Use \(\beta\ge10^{12}\) and \(\beta\ge r+k+2\max_{i,j}(|\lambda_i|,|\nu_j|)\). In either rank \(s_0\), the gap of the undoubled parameters between unlike signs is at least \(500s_0^3+h\), with \(h\ge10^{30}T^3\). These conditions hold for the chosen parameters with \(\beta=b^3\). The rank, gap and weight checks are in the local construction.\formalnote{7.12}{12}

First consider \(g,g'\) within \(\delta\) of \(1\) in Frobenius norm, where \(\delta T\le10^{-12}\). For either parameter the coefficient divided by its trace phase has distance at most \(a/10+e^{-100T}\) from a real number \(a\in[0,1]\). Here is how this error arises. In the Flensted--Jensen integral at \(g=\exp Y\), put \(q_0=Q_{-(\lambda+\rho)}(g-1)/4\), where \(\rho\) is the half-sum of positive roots in parameter order and
\[
 Q_w(U)=\sum_{i,j}|w_i-w_j|\,|U_{ij}|^2 .
\]
Use \(a=e^{-q_0}\). On the maximal compact group let \(m(k_0)\) be the squared modulus of the weight minor of the lowest type \(\mu\), \(Q=Q_\mu(k_0)\), and \(K_\mu=\int m(k_0)\,dk_0\) for probability Haar measure. The minor expansion expresses the error after removing the phase as the integral of \(m(k_0)e^{-q_0}(e^{R(k_0)}-1)/K_\mu\). Its quantitative inputs are
\[
 |R(k_0)|\le10^6\delta(q_0+Q+s_0^2),\qquad m(k_0)\le e^{-Q/2},\qquad
                  K_\mu\ge e^{-5000T}.
\]
The middle bound has an elementary interpretation. On a coordinate subspace with projection \(P\), the squared singular values \(s_i^2\) of \(Pk_0P\) lie in \([0,1]\). The loss \(1-s_i^2\) sums to \(\|(1-P)k_0P\|_{\rm HS}^2\). Since \(u\le e^{-(1-u)}\) for \(u\in[0,1]\), the squared determinant is at most the exponential of minus that sum. The two opposite crossing blocks of a unitary matrix have equal squared Hilbert--Schmidt norms. Thus multiplying over the threshold projections \(P=\mathbf1_{\{\mu_i\ge t\}}\) counts each ordered crossing \(|\mu_i-\mu_j|\) times with coefficient \(1/2\). Determinants of the whole compact element have modulus one. This proves the asserted bound on \(m\).

Set \(S_0=q_0+Q+s_0^2\). When \(S_0\le50000T\) the remainder is at most \(0.05\) and \(|e^R-1|\le1/10\). When \(S_0>50000T\), use \(|R|\le S_0/4\), \(s_0^2\le T\), and \(|e^R-1|\le e^{|R|}\) (from its series). The integrand in absolute value before dividing by \(K_\mu\) is bounded by
\[
          e^{-q_0-Q/2+S_0/4}\le e^{-S_0/4+T/2}\le e^{-5100T}.
\]
Integrating the two bounds and dividing by \(K_\mu\) proves the claimed \(a/10+e^{-100T}\). This calculation uses the uniform second-order remainder for the principal minors and the lower estimate on \(K_\mu\).\formalnote{7.13}{13}

Use the parameter \(-\nu\) for the complex conjugate coefficient. The logarithmic identity above cancels the two phases. If \(u,v\) are the resulting ratios, close to \(a,c\in[0,1]\) respectively, then expanding the product yields
\[
 |uv-ac|\le0.21ac+4e^{-100T},\qquad
             |uv|\le2\operatorname{Re}(uv)+12 e^{-100T}.
\]
For the second inequality use \(|uv|\le ac+|uv-ac|\) and \(\operatorname{Re}(uv)\ge ac-|uv-ac|\).

Now use the compact smaller group of the application, with \(L'=1\). For any group element \(s\), put \(\rho_0=10^{-13}/T\). The noncompact part of a matrix consists of its entries joining unlike signs. If the sum of the squared norms of the two noncompact parts is at least \(\rho_0^2\), the coefficient bound
\(
 |c_\lambda(g)|^4\le(1+\|g_{\rm nc}\|^2)^{-h}
\)
and its companion give \(|Z|^4\le (1+\rho_0^2)^{-h}\le e^{-400T}\). Here \(\log(1+\rho_0^2)\ge\rho_0^2/2\). Otherwise the compact-block localization gives a common diagonal \(U(X)\) translate with the same value of \(Z\) and with both matrices at Frobenius distance \(\le4\rho_0\) from \(1\). The preceding near-identity calculation applies. We have used the discrete-series gap estimate and the compact-block localization at this step.\formalnote{7.14}{14}

\subsection{Integrating a near-positive function}
Let \(J(s)=g\oplus g'\), using the operator norm from now on, \(M_0=9N_0\), \(G(s)=(2/\|J(s)\|)^{M_0}\). Discrete-series decay gives \(|Z|\le G\) and \(|Z|\le G^2\), using \(2M_0\le h\); apply the decay estimate to the larger component and the unit bound to the other. The positive part of \(|Z|-2\operatorname{Re}Z\) is at most both \(12e^{-100T}\) and \(3G^2\). Taking their geometric mean gives an integrable error:
\[
           |Z|\le 2\operatorname{Re} Z+6\beta^{-50N_0}G .
\]

Here are the requisite volume bounds. In the real normed vector space of complex \((r+k)\)-by-\((r+k)\) matrices the dimension is \(2N_0\). A maximal \(\alpha\)-separated subset of a ball of radius \(R\) has at most \((1+2R/\alpha)^{2N_0}\) points, by taking the disjoint \(\alpha/2\)-balls about them. Take centers from the image being covered. If \(\|J(s)\|,\|J(s_i)\|\le R\) and \(\|J(s)-J(s_i)\|\le\alpha\), then \(\|J(ss_i^{-1})-1\|\le R\alpha\), since the inverse of a signed unitary matrix has the same norm. Thus a right Haar measure of the ball of radius \(R\) in norm is bounded by \((1+2R^2/\epsilon)^{2N_0}\) times that of the \(\epsilon\)-ball about the identity. For \(R=\beta+1\) and \(\epsilon=(200\beta^5)^{-1}\) this factor is at most \(\beta^{16N_0}\). Write \(V,W\) for the respective two volumes.

The same covering with \(\epsilon=1\) gives integrability at infinity. Since \(\|J\|\ge1\), on \(2^j\le\|J\|\le2^{j+1}\), \(j\ge0\), the integrand \(G\) is at most \(2^{M_0}2^{-jM_0}\). The covering factor there is at most \(9^{2N_0}2^{4jN_0}\). The geometric series is bounded by \(16^{M_0}\) after including the factor \(2^{M_0}9^{2N_0}\). Therefore \(\int G\le16^{M_0}V\), since the unit ball about the identity is contained in the large norm ball. Properness of the first component and continuity of the second ensure the covering sets have finite measure. On the small ball each coefficient has absolute value at least \(1/2\) by its quantitative identity-neighborhood estimate. Writing \(I=|\int Z|\) gives
\[
       W/4\le 2I+6\beta^{-50N_0}16^{9N_0}V,\qquad
       V\le\beta^{16N_0}W .
\]
Because \(48\beta^{16N_0}16^{9N_0}\le\beta^{50N_0}\), substitution yields
\[
                 V\le16\beta^{16N_0}I\le\beta^{20N_0}I .
\]
In particular \(I>0\), because \(V\) contains an identity neighborhood of positive Haar measure. The same covering proof bounds translated matrix balls uniformly: apply the separation argument to \(J(s)y\), whose pairwise quotients are \(J(ss_i^{-1})\).\formalnote{7.15}{15}

\subsection{From a scalar period to a normalized vector with a decay bound}
Let \(R=U H\) be the Bessel group with compact factor \(H=U(k)\). In its unitary value representation \(\pi\) choose a unit vector \(v\) whose compact matrix coefficient is \(c_\nu\); on the unipotent factor use the prescribed character. Matching the stationary flag with the standard Bessel frame gives \(x\) with Bessel length at most \(b^{5n}\) when \(3B+B'\le b^2\) and \(8nr\le b\). The frame matching with this metric bound is a further input. Let
\[
 W_0(y)=\int_R c_\lambda(x^{-1}sy)\pi(s)^{-1}v\,ds,\qquad
 V_x=\operatorname{vol}_R\{s:\|x^{-1}sx-1\|\le1\}.
\]
The compact component has norm one, so this ball is contained in the large norm ball of the conjugated pair. The scalar period just bounded is \(\langle v,W_0(x)\rangle\). Let \(d_\lambda\) be the formal degree, and use compact Haar mass one. The convolution kernel \(d_\lambda c_\lambda\) reproduces \(W_0\) by discrete-series orthogonality. To see its normalization, insert the integral defining one factor of \(\|W_0\|^2\), change variable \(g=sy\) in the unfolded integral, and use \(W_0(sy)=\pi(s)W_0(y)\). This gives
\[
 d_\lambda\|W_0\|_{L^2(U\backslash G)}^2
 =\int_G d_\lambda\overline{c_\lambda(x^{-1}g)}
                   \langle v,W_0(g)\rangle\,dg
 =\langle v,W_0(x)\rangle .
\]
The Haar decomposition and absolute convergence use the quotient measure of \S\ref{sec:local-orbital} and the large-gap coefficient decay. The stationary bound with \(\beta=b^3\) proves \(V_x\le b^s d_\lambda\|W_0\|_2^2\) for \(s=60(r+k)^2\). In particular the constant \(c>0\) determined by \(c^2\|W_0\|_2^2=\dim\pi\) exists.\formalnote{7.16}{16}

Set \(m_0=9r^2\), \(\psi(g)=(2/\|g\|)^{m_0}\), \(\Phi(y)=\int_R\psi(x^{-1}sy)\,ds\). If \(\ell\) is the quotient length, its invariance under \(s\) and submultiplicativity give \(\ell(y)/b^{6n}\le\|x^{-1}sy\|/2\). Use the coefficient decay exponent \(m_0+i\) on the part where \(\ell(y)>b^{6n}\) and exponent \(m_0\) elsewhere. For every \(i\) up to the available excess gap this yields
\[
 \max(1,\ell(y)/b^{6n})^i\|c W_0(y)\|\le c\Phi(y),\qquad
 \Phi(y)\le16^{m_0}V_x .
\]
The latter inequality is the translated covering bound. Unfolding gives \(\int_{U\backslash G}\Phi=\int_G\psi\). The group normalization estimate for the actual formal degree is \(16^{m_0}d_\lambda\int_G\psi\le b^{99r^2}\). The bound for \(V_x\) consequently gives
\[
   (c16^{m_0}V_x)\int c\Phi
        \le b^{s+99r^2}\,c^2\|W_0\|_2^2
        \le b^{400r^2}\dim\pi .
\]
Here \(k\le r\) implies \(s+99r^2\le339r^2\). Taking the prescribed exponent \(i=3r^3f\) gives the majorant and the nonzero normalized function used in the cutoff argument. Section~\ref{sec:bessel-minors} derives the minor remainder, compact localization and small-ball coefficient bound from the lowest-type formula. The stationary signed-flag and frame comparisons and the Haar/formal-degree normalization are used in this deduction.\formalnote{7.17}{17}


\section{The lowest-type coefficient and localization of its phase}
\label{sec:bessel-minors}
In \S\ref{sec:bessel-period} the coefficient integral is approximated by a positive Gaussian after division by a trace phase. The quadratic form of that Gaussian has weights of the size of the Harish--Chandra parameter. A bound of order \(\|g-1\|_F^3\) per minor would accumulate an uncontrolled multiple of the largest weight. We prove the stronger bound that charges a minor only for matrix entries crossing its two blocks, then derive the reduction to the identity neighborhood.

\subsection{A dimension-independent expansion of a minor}
Write \(\|\cdot\|\) for the Frobenius norm. Suppose \(P,Q\) are orthogonal projections,
\(\|X\|\le\delta\le1/2\), and \(Y\) satisfies
\[
 \|[P,X]-Y\|\le\delta c,\qquad
 \max(\|Y\|,\|P-Q\|)\le c,\qquad \operatorname{tr}([P,Q]Y)=0 .
\]
Put \(Z=\log(1+X)\), using its convergent series. Then
\[
 \det(1+PX)=\exp\{\operatorname{tr}(QZ)-\tfrac14\operatorname{tr}(Y^2)+r\},
 \qquad |r|\le100\delta c^2 .                                      \tag{M1}
\]
Indeed, set \(A=[P,X]\). Expand the trace logarithm on \(\operatorname{im}P\) and subtract \(\operatorname{tr}(PZ)\). The linear terms cancel, and the quadratic difference is \(-\operatorname{tr}(A^2)/4\). For degree \(j\ge3\), telescope \(PX^jP-(PXP)^j\) at the first occurrence of \(1-P\). Its trace is the sum of the traces of
\((PXP)^kPX(1-P)X^{j-1-k}P\), \(0\le k\le j-2\), on \(\operatorname{im}P\). The two factors crossing between the blocks have Frobenius norms at most \(\|A\|\) and \((j-1-k)\delta^{j-2-k}\|A\|\), respectively, by the commutator formula for a power. Cauchy--Schwarz bounds the sum by \(\binom j2\delta^{j-2}\|A\|^2\). Dividing by \(j\) and summing gives at most \(3\delta\|A\|^2\).

To replace \(P\) by \(Q\), split \(P-Q\) into blocks for \(P\). Its two diagonal blocks have opposite signs and their trace norms sum to \(\|P-Q\|^2\). Their contribution to \(\operatorname{tr}((P-Q)Z)\) is at most \(2\delta c^2\). The off-diagonal contribution is \(\operatorname{tr}([P,Q][P,Z])\). We have \(\|[P,Q]\|\le\|P-Q\|\), may subtract \(Y\) from \([P,Z]\) by hypothesis, and the logarithm series gives \(\|[P,Z]-Y\|\le\delta c+\delta\|A\|/(1-\delta)\le4\delta c\). Finally \(\|A\|\le(1+\delta)c\), so replacing \(A^2\) by \(Y^2\) in the trace costs at most \(\delta c(\|A\|+\|Y\|)\). These bounds give (M1). The formal calculation integrates the two logarithmic derivatives on the segment from \(0\) to \(X\); the series above estimates the same trace logarithms term by term.\formalnote{7.18}{18}

\subsection{From one minor to the discrete-series coefficient}
For one group of rank \(d\), let \(J=\operatorname{diag}(\epsilon_i)\), \(g\in U(J)\), \(X=g-1\), and \(\delta=\|X\|<10^{-3}\). Write \(Z=\log g\). For the undoubled Harish--Chandra parameter \(\lambda\), put \(v=\lambda+\rho_{\rm roots}\), where \(\rho_{\rm roots}\) is the half-sum of positive roots in parameter order. The highest weight of the lowest compact type is \(\mu=v-2\rho_c\), with \(\rho_c\) the half-sum of positive compact roots. For an integer weight \(w\) define \(m_w(M)=\det(M)^b\prod_{b<t\le B}\det M_{P_t,P_t}\), where \(P_t=\mathbf1_{\{w_i\ge t\}}\) and \(b\le w_i\le B\). This is independent of the bounds.

Let \(u,u'\) retain the entries \(X_{ij}\) with \(\epsilon_i=\epsilon_j\) and \(\lambda_i<\lambda_j,\ \lambda_i>\lambda_j\), respectively, and vanish elsewhere. Both \(1-u\) and \(1-u'\) are triangular in parameter order with diagonal one. Applying the Borel transformation law on both sides of \(W=(1-u)g(1-u')\) gives
\[
 c_\lambda(g)=\frac1{K_\mu}\int_K\frac{|m_\mu(k)|^2}{m_{-v}(k^*Wk)}\,dk,
 \qquad K_\mu=\int_K |m_\mu(k)|^2\,dk .                            \tag{M2}
\]
Here \(K=U(d_+)\times U(d_-)\) has Haar mass one.\formalnote{7.19}{19}

Fix a threshold \(P=P_t\) of \(-v\), so that the entries of \(\lambda\) on \(P\) are lower than those on \(1-P\). Write \(A=[P,X]\), \(E=A-[P,u]-[P,u']\), \(Q=kPk^*\), \(a=\|A\|\), and \(s=\|[P,k]\|=\|Q-P\|\). The matrix \(E\) consists of the entries of \(A\) between opposite signs, whereas \(P,Q\) preserve signs, so \(\operatorname{tr}([Q,P]E)=0\). Expanding the three factors gives \(\|W-1\|\le4\delta\) and \(\|[P,W-1]-E\|\le8\delta a\): each remaining product in the commutator has a factor of norm at most \(\delta\) and a commutator of norm at most \(a\). Replacing \(P\) by \(Q\) in this last commutator costs at most \(8\delta s\).

Apply (M1) to the complementary minors of \(g,W\), using \(1-P\) for both projections and \(-A,-E\) for \(Y\), respectively. Apply it once more to the minor of \(k^*Wk\) on \(P\), using projections \(Q,P\) and \(Y=E\). The respective pairs of bounds \((\delta,c)\) in (M1) are \((\delta,a),(4\delta,4a),(4\delta,4(a+s))\). Triangularity gives equal complementary minors for \(W,g\) and also \(\det W=\det g\). Combining the three exponentials therefore leaves the linear term \(\operatorname{tr}(PZ)\) and the quadratic term \(\{\operatorname{tr}(A^2)-2\operatorname{tr}(E^2)\}/4\). Since \(g^*Jg=J\), we have \(A^*=Jg^{-1}Ag^{-1}J\). Thus \(\|A^*-JAJ\|\le5\delta a\), using \(\|g^{-1}-1\|\le2\delta\). Since \(JAJ=A-2E\), the quadratic term differs from \(a^2/4\) by at most \(5\delta a^2/4\). In total the error per minor is at most \(2\cdot10^5\delta(a^2+s^2)\).

Multiply over all thresholds and restore the determinant power in \(m_{-v}\). Each ordered pair \(i,j\) crosses exactly \(|v_i-v_j|\) thresholds, so the quadratic terms sum to \(q=Q_{-v}(X)/4\), with \(Q_w\) as in \S\ref{sec:bessel-period}. The linear terms sum to \(-\operatorname{tr}(\operatorname{diag}(v)Z)\). Using \(|2\rho_{c,i}|\le d\), \(\|Z\|\le2\delta\) and \(\|k\|^2=d\) gives \(Q_{-v}(k)\le Q_\mu(k)+2d^2\). Hence the denominator in (M2) is
\[
 \exp\{q-\operatorname{tr}(\operatorname{diag}(\mu)Z)-r(k)\},
 \qquad |r(k)|\le10^6\delta\{q+Q_\mu(k)+d^2\}.
\]
This is the coefficient estimate used in \S\ref{sec:bessel-period}.\formalnote{7.20}{20}

\subsection{The compact normalization and the identity neighborhood}
Let \(L\) bound the absolute values of the doubled parameters on both groups, with \(\beta\ge r+k+L\) and \(T=(r+k)^2\log\beta\). The assumptions include \(\beta\ge10^{12}\). Within each sign block, order the coordinates by decreasing parameter. The weight \(\mu\) is dominant in that order. For two coordinates within a block, the difference of the \(\lambda\)-coordinates is at least their distance in the full order, and the compact \(\rho_c\)-difference is at most that distance. The principal minors defining \(m_\mu\) on \(K=U(d_+)\times U(d_-)\) are coefficients of the leading wedge vectors in the exterior powers. Their product, with the indicated determinant power, is the coefficient of the tensor of these unit highest vectors, which generates the irreducible highest-weight representation \(V_\mu\). Since \(K\) has Haar mass one, Schur orthogonality \cite[Corollary~4.10]{theta-knapp} gives \(K_\mu=(\dim V_\mu)^{-1}\).

Each \(|\mu_i|\le2\beta\), since \(\mu=\lambda+\rho_{\rm roots}-2\rho_c\). In Weyl's dimension formula \cite[Theorem~5.84]{theta-knapp} the factor for a pair \(i<j\) in a block, indexed within that block, is \((\mu_i-\mu_j+j-i)/(j-i)\le5\beta\). There are at most \(d^2/2\) pairs, including the empty product for a torus. Schur orthogonality and Weyl's dimension formula therefore give
\[
       K_\mu\ge (5\beta)^{-d^2/2}\ge e^{-T}\ge e^{-5000T}.
\]
\phantomsection\label{route:10}\leanroute{Lean proves the compact mass bound by applying its quantitative derivative and covering estimate to \(K\). The classical dimension formula bounds the same squared norm by Schur orthogonality. The full-group estimate below still uses the derivative bound.}\formalnote{7.21}{21}

For the full group, \(\beta\ge d+L+1\) because the other group has positive rank. The real derivative estimate of \S\ref{sec:real-decay} gives
\[
 |c_\lambda(g)-1|\le100(d+L+1)^5\|g-1\|_{\rm op}.
\]
Hence \(|c_\lambda(g)|\ge1/2\) when \(\|g-1\|_{\rm op}\le (200\beta^5)^{-1}\). This bound will supply positive mass on the identity neighborhood in \S\ref{sec:bessel-period}.\formalnote{7.22}{22}

\subsection{Decay from the spectral gap}
For \(y>0\) in \(U(J)\), write \(a=(y+JyJ)/2\) for its sign-block diagonal part, which is positive definite. We claim
\[
                     |\det a|^h\le |m_{-v}(y)|^2                       \tag{M3}
\]
for any nonnegative integer \(h\) at most the gap between unlike signs in the parameter. For each threshold \(t\) let \(S=\{i:v_i\ge t\}\), \(S'=\{i:v_i+h\ge t\}\). The difference \(C=S'\setminus S\) has only one sign by the gap, so the principal blocks of \(a,y\) on \(C\) coincide. For a positive matrix the determinant on a union of disjoint coordinate sets is at most the product of their determinants, by Schur complementation. Thus \(\det a_{S'}\le\det a_S\det y_C\). Since \(y^{-1}=JyJ\) and \(\det y=1\), Jacobi's complementary-minor identity gives \(\det y_C=\det y_{C^c}\). Splitting \(C^c=S\sqcup(S')^c\) a second time gives
\[
                   \det a_{S'}\le\det a_S\,\det y_S\,\det y_{(S')^c}.
\]
Multiply over the thresholds. The ratio of the weight products for \(v+h\) and \(v\) on \(a\) is \(\det(a)^h\). The other two products are \(m_v(y)\) and \(m_{v+h}(y^{-1})\), with determinant factors of modulus one. Complementation and \(y^{-1}=JyJ\) give the same absolute value \(|m_{-v}(y)|\) for both. This proves (M3).

Polar decomposition writes general \(g\in U(J)\) as \(k y\), \(k\in K,\ y>0\). In the polar Flensted--Jensen formula the denominator is \(m_{-v}(u^* y u)\), \(u\in K\), while the numerator is a product of compact coefficients on \(ku\) and \(u\). The compact-block determinant has the same modulus for \(g,y,u^*yu\). Bound the reciprocal denominator by (M3) and integrate the numerator using Cauchy--Schwarz. Its two squared integrals both equal \(K_\mu\), by Haar invariance. Therefore \(|c_\lambda(g)|^2|\det a_g|^h\le1\). Writing \(b_g=g-a_g\) for the opposite-sign part, the sign-block diagonal equations in \(g^*Jg=J\) imply \(a_g^*a_g=1+b_g^*b_g\). The eigenvalues of \(b_g^*b_g\) are nonnegative and \(\prod(1+t_i)\ge1+\sum t_i\). In particular
\[
                    |c_\lambda(g)|^4\le(1+\|b_g\|_F^2)^{-h},          \tag{M4}
\]
the estimate used outside the compact tube.\formalnote{7.23}{23}

\subsection{Localization modulo the shared zero-weight subspace}
For the pair \(g,g'\) at a stationary position let \(b,b'\) be their opposite-sign parts and let \(t^2=\|b\|_F^2+\|b'\|_F^2\). Put \(A=g-1,\ A'=g'-1\), and write \(P\) for projection off the common positive subspace \(X\). The lowest weights on its complement satisfy \(B\le-\mu_i\le3B\). The smaller weights \(\mu'_i\) are nonnegative. Their upper bound on coordinates meeting an opposite sign is \(3B\); in the compact application every entry of \(b'\) is zero. The row equations in \(gJg^*=J\) give
\[
                    \sum_j|A_{ij}|^2
                    =2\sum_j|b_{ij}|^2-2\operatorname{Re}A_{ii}.
\]
Multiply by \(-\mu_i\) and sum. The analogous sum with \(\mu'_i\) in the second group has a nonnegative left side. On adding, the two trace terms cancel because their difference is the stationary phase, which is purely imaginary. Division by \(B\) gives \(\|PA\|_F^2\le6t^2\). Apply the same estimate to the inverse pair, using \(g^{-1}=Jg^*J\), and get \(\|AP\|_F^2\le6t^2\).

Let \(a\) be the compression of \(g\) to \(X\). Its column Gram equations give \(\|a^*a-1\|_F\le\|PA\|_F^2\): the missing off-\(X\) block contributes a signed Gram matrix of Frobenius norm bounded by the squared Frobenius norm of that block. For \(0<t\le1/6\), \(6t^2\le t<1\); the polar decomposition \(a=u z\) has \(u\) unitary and \(z>0\). Since \(|s-1|\le|s^2-1|\) for \(s\ge0\), \(\|z-1\|_F\le6t^2\le t\). Translate the pair on the left by \(u^{-1}\) in the shared diagonal \(U(X)\) of the position group. Both lowest weights vanish there, so left translation fixes both coefficients (the compact highest-weight vector is fixed on its equal zero-weight block). The character is also fixed because its trace phase on this element is zero, and the norms of \(b,b'\) are unchanged. The new \(X\)-block is \(z\). Thus for the translated pair
\[
     \|g-1\|_F^2\le6t^2+6t^2+t^2=13t^2.
\]
On summing the unweighted row equations and using the equality of the traces of \(g-1\) and \(g'-1\), we obtain
\[
   \|g'-1\|_F^2=\|g-1\|_F^2-2\|b\|_F^2+2\|b'\|_F^2\le15t^2 .
\]
The polar argument also works for \(t=0\), with \(z=1\). Both distances are thus at most \(4t\), with exactly the same value of the period integrand. For \(t\ge10^{-13}/T\), (M4) gives the decay used in \S\ref{sec:bessel-period}; for \(t<10^{-13}/T\) the translated pair is in the Gaussian neighborhood. The analytic inputs to this calculation are the polar and near-identity Flensted--Jensen formulas for the normalized lowest-type coefficient and its triangular covariance. Section~\ref{sec:flag-measure} constructs the stationary flag and proves the trace identities assumed of the pair here.\formalnote{7.24}{24}


\section{The signed flag, its Bessel orbit and the Haar mass}\label{sec:flag-measure}
Three geometric facts enter the real second moment in \S\ref{sec:bessel-period}. The linear phase is a character on an entire subgroup. A point identifying that subgroup with the Bessel group can be chosen at bounded quotient length. Finally, the Haar measure used for the period gives the required quantitative bound when multiplied by the formal degree. We prove these facts with the norms and measures used in the period estimate.

\subsection{A trace identity on an isotropic flag}
Let \([x,y]=x^*Jy\) be a nondegenerate hermitian form and \(T^\dagger\) the adjoint for it, using the corresponding form on the domain if \(T\) is rectangular. Use the independent columns \(v_0,\ldots,v_d\) of the stationary position. Their only nonzero pairing is \([v_d,v_d]=\pm1\). Write \(V=(v_0,\ldots,v_d)\), \(A=(v_0,\ldots,v_{d-1})\), and \(I\) for the isometric inclusion of the smaller space. The image of \(I\) is orthogonal to every \(v_j\). The ambient dimension is the smaller dimension plus \(2d+1\). Consequently
\[
     (\operatorname{im}A)^\perp
               =\operatorname{im}I\ \oplus\ \operatorname{im}V .
\]
Indeed the right side is contained there, its summands meet trivially because \(I^\dagger I=1\) and \(I^\dagger V=0\), and it has the requisite dimension.

Let \(g\) preserve the form, with \(Y=g-1\) carrying \(v_j\) into the span of its predecessors. Then \(Y^\dagger=g^{-1}-1\) lowers the same flag, since the inverse of an upper unitriangular matrix is upper unitriangular. Write \(YV=AN\), \(Y^\dagger V=AN_-\). The matrices \(N,N_-\) have zero entries in positions \((i,j)\) with \(j\le i\). The columns of \(YI\) are orthogonal to all of \(V\), by adjunction. The decomposition above and \([v_d,v_d]\ne0\) therefore give
\[
      YI=IY'+AC,\qquad Y'=I^\dagger YI .
\]
In particular \(g'=1+Y'\) is the action on the nondegenerate smaller summand modulo \(\operatorname{im}A\). It respects its form. This compression is multiplicative, since a product preserves \(\operatorname{im}A\) and that space is killed by \(I^\dagger\). The subgroup is closed: fix any left inverse of \(A\); the requirement that the columns of \(YV\) equal \(A\) times their coordinates there, with the indicated zero entries, is a system of continuous matrix equalities.\formalnote{7.25}{25}

Let \(D\) be self-adjoint, with \(Dv_i=v_{t(i)}\) for \(0\le i<d\), and put \(D'=I^\dagger D I\). We claim
\[
    \operatorname{tr}(DY)-\operatorname{tr}(D'Y')
       =\sum_{i<d}N_{i,t(i)}
           +\overline{\sum_{i<d}(N_-)_{i,t(i)}} .                  \tag{FM1}
\]
Choose \(Z\) with \(ZA=1,\ ZI=0\) by extending the two independent sets of columns to a basis. Set \(T=DAZ\). Then \(T\) agrees with \(D\) on \(\operatorname{im}A\), kills \(\operatorname{im}I\), and has image in \(\operatorname{im}V\). Hence \(T^\dagger\) kills both \(\operatorname{im}A\) and \(\operatorname{im}I\). The self-adjoint remainder \(M=D-T-T^\dagger\) kills \(A\), has image in \(A^\perp\), and has \(I^\dagger M I=D'\). Cyclicity gives
\[
     \operatorname{tr}(TY)=\operatorname{tr}(ZYDA)
          =\sum_{i<d}N_{i,t(i)}.
\]
Taking adjoints gives the conjugated sum in (FM1) for \(T^\dagger Y\). To compute the remaining trace, put \(C_0=(I\mid V)\), write \(M=C_0B\) by the description of \(A^\perp\), and use \(BA=0\) because \(C_0\) is injective. Then \(BYC_0=(BIY'\mid0)\). The trace of this rectangular cyclic product, taken on the domain of \(C_0\), is the trace of the top block of \(BIY'\). The top block of \(B\) equals \(I^\dagger M\), proving (FM1).\formalnote{7.26}{26}

For \(D=1\) take \(t(i)=i\); both displayed sums vanish and \(\operatorname{tr}Y=\operatorname{tr}Y'\). For the stationary lowest weight take \(t(i)=i+1\). Its compression \(D'\) is the smaller weight by the residue calculation of \S\ref{sec:bessel-period}. On the whole flag \(g=1+N_{\rm full}\), where \(N_{\rm full}\) has the rows of \(N\) followed by a zero row. All powers of \(N_{\rm full}\) beyond the first vanish on its first superdiagonal. Thus the adjacent entries in \(g^{-1}-1\) are the negatives of those in \(g-1\). If \(\sigma(g)=\sum_{i<d}N_{i,i+1}\), (FM1) now gives exactly \(\sigma(g)-\overline{\sigma(g)}\). Multiplication adds adjacent entries, so the exponential of this trace difference is a character. The formulas include \(d=0\), where both sums are empty.\formalnote{7.27}{27}

\subsection{An isometry with bounded quotient length}
Use the Bessel space \(G=U(J)\) with \(d=\rho\) hyperbolic planes, \(k\) positive middle coordinates \(X_0\) and the negative unit vector \(z\). Put \(n=d+1,\ r=k+2d+1\), and let \(F=(e_0,\ldots,e_{d-1},z)\) be the standard frame. The standard ordered basis is Euclidean orthonormal, and \(J=J^*=J^{-1}\) interchanges \(e_i,e_i^*\), is \(+1\) on \(X_0\) and \(-1\) on \(z\). The sign-coordinate change \(p\) satisfies \(p^*Jp=J_{\rm sign}\) and is Euclidean unitary. In this application \([v_d,v_d]=-1\). The columns of \(p(I\mid V)\) have the same Gram matrix as the corresponding standard columns \((X_0\mid F)\). The induced partial isometry extends by Witt's theorem to \(x\in G\) identifying these lists, including when their span is degenerate. Concretely, for each isotropic radical vector choose a dual orthogonal to the nondegenerate columns and subtract half the matrix of their mutual pairings in the radical directions. The duals are then isotropic and give hyperbolic pairs that complete the partial isometry.

Witt extension alone gives no numerical bound on \(x\). Put \(w_a=D^av_0\) in sign coordinates and \(V_a=pw_a\) for \(0\le a\le2d\). The first \(d+1\) of these are the columns of \(x^{-1}F\). Moving \(D\) from one side of the pairing to the other shows that its value depends only on \(a+b\). For \(a+b\le2d\), rewrite \(a+b=a'+b'\) with \(a',b'\le d\). The known Gram matrix of the frame therefore gives
\[
       [V_a,V_b]=0\ (a+b<2d),\qquad [V_a,V_b]=-1\ (a+b=2d).
\]
By the node bound in \S\ref{sec:bessel-period}, every coordinate of \(w_a\) has modulus \(\le(3B+B')^{2n}\). In Euclidean norm the \(V_a\) have length at most \(M=\sqrt r\,(3B+B')^{2n}\). Orthogonalize the first \(n\) vectors in Euclidean norm without normalizing, obtaining \(f_a=V_a+\sum_{i<a}c_iV_i\). Orthogonal projection gives \(\|f_a\|\le M\). Pairing against \(V_{2d-a}\) gives \(-1\), so Cauchy--Schwarz and \(\|J\|=1\) give \(\|f_a\|\ge M^{-1}\).\formalnote{7.28}{28}

Let \(U\) be the Bessel unipotent group and \(R=U\,U(k)\). Every upper unitriangular change of frame is realized by \(U\). In fact if it sends the columns \(E=(e_i)_{i<d}\) to \(ET\) and \(z\) to \(z+Ec\), take the dual columns with bottom block \((T^*)^{-1}\), and choose their \(z\)-coordinates to make them orthogonal to \(z+Ec\). Their mutual Gram matrix can then be made zero by adding linear combinations of \(ET\), since its change is the sum of the correcting matrix and its adjoint. Fixing the \(X_0\) columns completes a block unitriangular isometry. Use this to choose a left translate \(y\in Ux\) with \(y^{-1}F=(f_0,\ldots,f_d)\).

For \(i<d\) the vectors \(b_i=Jf_i/\|f_i\|^2\) satisfy \([f_a,b_i]=\delta_{ai}\) and are mutually isotropic. Their norms are at most \(M\). A further translation, now by the part of \(U\) fixing \(F\) pointwise, arranges that the inverse columns at \(e_i^*\) equal \(b_i\). To verify this explicitly, write \(yb_i\) together in the basis \(E,X_0,z,E^*\). Duality forces their bottom block to be \(1\) and their \(z\)-row to vanish. Write their other blocks as \(C,B_0\), respectively. Isotropy reads \(C+C^*=-B_0^*B_0\). Thus an isometry fixing \(F\) sends \(E^*\) to these columns by sending each \(X_0\)-column to itself minus the corresponding column of \(E B_0^*\). This is an element of \(U\), as all its nonidentity entries strictly raise the flag. Left multiplication of \(y\) by its inverse has the asserted effect.\formalnote{7.29}{29}

Call the resulting matrix \(y_0\), and set \(A_0=y_0^{-1}\). All its columns at \(E,z,E^*\) have length at most \(M\). For the other columns use \(A_0 J A_0^*=J\), which gives
\[
      \|A_0\|^2\le1+\|A_0(1-J)\|\ \|A_0\|.
\]
The nonzero columns of \(1-J\) are \(2z\) and the two copies with opposite signs of \(e_i-e_i^*\). Summing the norms of the resulting rank-one column maps bounds \(\|A_0(1-J)\|\) by \((4d+2)M\). Since \(\|A_0\|=\|y_0\|\ge1\), and \(M\ge1\), this proves
\[
     \ell(x)=\inf_{s\in R}(1+\|sx\|)
       \le 8n\sqrt r\,(3B+B')^{2n}\le b^{5n}.                    \tag{FM2}
\]
The last inequality uses \(3B+B'\le b^2,\ 8nr\le b,\ b\ge1\), giving \(b^{4n+1}\le b^{5n}\). The stationary isometry is \(x\) itself; the left translations bound its quotient length.\formalnote{7.30}{30}

Conjugation by \(xp\) identifies the flag-lowering group with \(R\): after the frame and the middle subspace have been identified, an element acts by \(h\in U(k)\) on the positive quotient, and removal of \(h\) leaves the defining block upper unitriangular isometry. In the other direction every \(hu\) lowers the standard frame after subtraction of \(1\). On the graph \(\{(hu,h)\}\subset G\times U(k)\) its projection to \(G\) is proper, because it is closed and the second factor is compact. Thus the pushforward of the graph measure is finite on compact sets and right invariant on the full position group, by properness and surjectivity. The first superdiagonal sum and the action of \(h\) are preserved by the identification. This justifies applying the signed-flag period inequality to the Bessel integral with its actual character.\formalnote{7.31}{31}

\subsection{Unfolding and the formal-degree estimate}
Give \(H=U(k)\) probability Haar measure and \(U\) its Haar measure in real unipotent coordinates. In the subgroup \(U^0\) fixing the frame, the coordinates \((\omega,\mu)\) satisfy \(\mu+\mu^*=\omega^*\omega\), so only the skew-Hermitian part of \(\mu\) is free. Multiplication sends these parameters and \((\omega',\mu')\) to \((\omega+\omega',\mu+\mu'+\omega^*\omega')\). Left and right translations therefore have real Jacobian one in the free coordinates. For \(U\) the remaining coordinates are the strictly upper triangular frame entries. Ordering them by the difference in flag levels makes the free-entry changes under multiplication triangular with diagonal one. Conjugation by \(H\) preserves \(du\) because the scale factor is a continuous positive character of the compact group \(H\), hence equals one. The graph measure \(ds\) is the image of \(du\,dh\) under \((u,h)\mapsto(hu,h)\). Its right invariance also follows immediately from \((hu)(h'u')=hh'(h'^{-1}uh')u'\).\formalnote{7.32}{32}

Use the quotient measure \(d\dot y\) on \(U\backslash G\) corresponding to \(du,dg\). For every nonnegative measurable \(F_0\),
\[
     \int_{U\backslash G}\int_R F_0(sy)\,ds\,d\dot y
       =\int_H\int_G F_0(hg)\,dg\,dh=\int_G F_0(g)\,dg .
\]
The inner integral is independent of the representative of \(\dot y\). This identity explains both the quotient by \(U\) used for the squared norm and the absence of a second compact volume factor. For the normalized coefficient \(c_\lambda\), put \(d_\lambda=(\int_G |c_\lambda|^2\,dg)^{-1}\). Schur orthogonality gives \(c_\lambda*(d_\lambda c_\lambda)=c_\lambda\). Translating and averaging this identity reproduces the Bessel function \(W_0\). Inserting its defining integral into \(\|W_0\|^2\) and using \(W_0(sy)=\pi(s)W_0(y)\) now unfolds to \(d_\lambda\|W_0\|^2=\langle v,W_0(x)\rangle\), as in \S\ref{sec:bessel-period}. The gap gives \(c_\lambda\in L^1(G)\) and bounds the integral defining \(W_0\) uniformly. Their product after unfolding is absolutely integrable.\formalnote{7.33}{33}

For the cutoff, let \(L\) bound the absolute values of the entries of the doubled Harish-Chandra parameter. Put \(D=r+L+1\) and \(K=2000D^6\). At the gap used in the real construction, the normalized moment bound (RD2) of \S\ref{sec:real-decay} gives, for \(a>8r^2\),
\[
      d_\lambda\int_G(2/\|g\|)^a\,dg\le8 K^{2r^2}2^a .
\]
With \(b\ge2000,\ D\le b^4,\ a=9r^2\), one has \(K\le b^{25}\), \(8\le b\) and \(2,16\le b\). Multiplying by \(16^a\) gives the required bound \(b^{50r^2+1+2a}\le b^{99r^2}\). The moment uses \(d_\lambda=(\int|c_\lambda|^2)^{-1}\). The identification of \(d_\lambda\) with a Weyl dimension in \S\ref{sec:elliptic-character} uses the compact-dual normalization and the stronger gap \(500r^3+10r^5\).\formalnote{7.34}{34}


\section{Decay and variation of the real discrete-series coefficient}\label{sec:real-decay}
The trace calculation uses a pointwise bound for a normalized discrete-series coefficient, and the volume calculation in \S\ref{sec:flag-measure} uses a ball on which that coefficient is bounded below. Put \(J=\operatorname{diag}(\epsilon_1,\ldots,\epsilon_n)\), \(G=U(J)\), and \(K=G\cap U(n)\), with \(n\ge1\). Use the Euclidean operator norm on matrices. Let \(l=2\lambda\) be regular, with entries congruent to \(n-1\bmod 2\). Assume \(|\lambda_i-\lambda_j|\ge h\) for opposite signs, where \(h\) is a nonnegative integer. In parameter order let \(\rho,\rho_K\) be the half-sums of all positive roots and the compact positive roots, respectively, and put \(\Lambda=\lambda+\rho-2\rho_K\). The construction uses a compact integral of Flensted--Jensen type. For its origin in the theory of symmetric spaces see \cite{theta-flensted}. In the polar integral below, setting \(g=1\) makes the integrated numerator equal to the normalization and the denominator inside the integral equal to one. Hence \(c_\lambda(1)=1\).\formalnote{7.35}{35}

\subsection{Pointwise decay from principal minors}
Put \(a(g)=(g+JgJ)/2\), the two diagonal blocks of \(g\). In the polar decomposition \(g=kz\) the positive factor satisfies \(JzJ=z^{-1}\) and \(k\in K\). Indeed, \(g^*Jg=J\) gives \(J(g^*g)J=(g^*g)^{-1}\). Taking the positive square root proves the identity for \(z=(g^*g)^{1/2}\). The eigenvalues of \(a(z)\) are \((s+s^{-1})/2\) for the eigenvalues \(s\) of \(z\). Each is at least one, and the extreme eigenvalues of \(z\) are \(r=\|g\|\) and \(r^{-1}\). When \(r>1\) these contribute two factors of \((r+r^{-1})/2\). When \(r=1\), \(z=1\). Thus\formalnote{7.36}{36}
\[
                  |\det a(g)|=|\det a(z)|\ \ge\ (r/2)^2.
\]

For any integer \(0\le c\le h\), apply \textup{(M3)} and the ensuing polar-integral estimate of \S\ref{sec:bessel-minors} with gap exponent \(c\). The weight denoted there by \(v\) is \(\lambda+\rho\), and the block determinant is \(\det a(g)\). Together with the preceding eigenvalue bound this gives\formalnote{7.37}{37}
\[
              |c_\lambda(g)|^2 |\det a(g)|^c\le1,\qquad
              |c_\lambda(g)|\le(2/\|g\|)^c .                         \tag{RD1}
\]
This is the exponent used to dominate the real orbital integrals in \S\ref{sec:elliptic-character}.\formalnote{7.38}{38}

For the representation used below, suppose \(h\ge500n^3\) and write \(c=c_\lambda\). First \(c\in L^1(G)\cap L^2(G)\). Indeed, put \(B_R=\{g\in G:\|g\|\le R\}\). The set \(B_R\) can be covered by \(O_n(R^{4n^2})\) right translates of the fixed neighborhood \(\{\|g-1\|\le1\}\). To see this, take a maximal \(1/R\)-separated subset of \(B_R\) in the real matrix space of dimension \(2n^2\), and use \(\|x^{-1}\|=\|x\|\le R\) at each center \(x\). This neighborhood is compact and has finite Haar measure. Summing (RD1) on the sets \(B_{2^{j+1}}\setminus B_{2^j}\) gives integrability on taking any integer exponent between \(8n^2\) and \(h\), and (RD1) with exponent zero gives \(|c|\le1\).

Put \(M=\int_K|m_\Lambda(k)|^2\,dk\). The compact averaging identity for the normalized integral is
\[
  \int_K\overline{m_\Lambda(k)}c(xky)\,dk=M c(x)c(y).
\]
For the compact-type calculation, replace the conjugated minor on the left by \(\overline{m_\nu}\), where \(\nu\) is compact dominant. Suppose that \(\nu=\Lambda\) or that \(\nu-\Lambda\) lies outside the positive-root cone. Take an integral vector \(o\) strictly reversing the parameter order such that \(N=o\cdot(\nu-\Lambda)\) is zero in the first case and positive in the second. Such an \(o\) exists by separation of the cone, with a small perturbation and clearing denominators. Initially take \(y\) near the identity. Write \(a=a(x)\), \(w=-\lambda-\rho\), \(W_j=j^*a^{-1}xj\), \(Y_k=k^*yk\). Substituting the defining integral for \(c(xuy)\) and changing variables \(j=uk\) gives
\[
  M\int_K\overline{m_\nu(u)}c(xuy)\,du
  =\int_{K^2}
   \frac{m_\nu(kj^{-1})m_\Lambda(aj)\overline{m_\Lambda(k)}}
        {m_w(W_jY_k)}\,dj\,dk .
\]
Use Haar invariance to average the right side under \(k\mapsto kt_s\), where \(t_s=\operatorname{diag}(s^{o_i})\), \(|s|=1\). Cauchy--Binet on the ordered flag minors shows that \(s^{-o\cdot\nu}m_\nu(kt_sj^{-1})\) has nonnegative powers of \(s\) and constant term \(m_\nu(k)m_\nu(j^{-1})\). Conjugating \(m_\Lambda(kt_s)\) contributes \(s^{-o\cdot\Lambda}\overline{m_\Lambda(k)}\). For the denominator, factor \(Y_k\) by triangular elimination in the order \(o\). In \(t_s^{-1}Y_kt_s\), the unipotent factor contracts to the identity as \(s\to0\). The complementary triangular factor acts on the right of the flag minors by its diagonal character \(m_w(Y_k)\). The remaining denominator has no zeros on the closed disc when \(y\) is close to \(1\), by continuity and compactness; at \(y=1\) the conjugates \(W_j\) have positive definite Hermitian part. The circle mean is therefore the constant term of \(s^N\) times a holomorphic function. It is zero for \(N>0\). For \(\nu=\Lambda\) its value is
\[
   \frac{m_\Lambda(aj)\overline{m_\Lambda(j)}}{m_w(W_j)}
   \frac{|m_\Lambda(k)|^2}{m_w(Y_k)} .
\]
The two integrals are \(M c(x)\) and \(M c(y)\), respectively. Division by \(M\) proves the asserted identity near \(y=1\). Both sides are real analytic in \(y\) on the connected group \(G\), so both the identity and the vanishing hold throughout \(G\).

Write \(I(x,y)=\int_G c(xgy)\overline{c(g)}\,dg\). Since \(c(g^{-1})=\overline{c(g)}\) by the polar integral, averaging \(c(xkgy)c(g^{-1})\) against \(\overline{m_\Lambda(k)}\) gives \(M c(x)I(1,y)\). The change \(g\mapsto k^{-1}g\), followed by the same compact identity applied to \(c(g^{-1}k)\), gives \(M I(x,y)\). Thus \(I(x,y)=c(x)I(1,y)\). Inverting \(g\) and translating gives \(I(1,y)=I(y,1)=c(y)I(1,1)\). With \(d_\lambda=1/I(1,1)\), the vector \(v_0=\sqrt{d_\lambda}\,c\) has norm one. For right translation \(R(y)f(g)=f(gy)\) its coefficient is
\[
   \langle v_0,R(y)v_0\rangle=d_\lambda I(1,y)=c(y).
\]
On \(K\) the compact integral gives \(c(k)=m_\Lambda(k)\). Thus the translates of \(v_0\) by \(K\) have the same Gram matrix as the translates of a unit highest vector of weight \(\Lambda\). Their span is a copy of \(V_\Lambda\), with \(v_0\) its highest vector. The orbit map \(x\mapsto R(x)v_0\) is smooth in the Hilbert norm. Indeed its Gram kernel is \(c(x^{-1}y)\), smooth in both variables. Inner products of divided differences of the orbit vectors are the corresponding mixed divided differences of this kernel. In local coordinates their limits show that the vector divided differences are Cauchy at every order and that the resulting derivatives are continuous. This justifies differentiating \(v_0\) in the following estimate.

\subsection{A ball of known radius}
Now assume \(h\ge500n^3+B\), \(B\ge0\), and \(|l_i|\le L\). Let \(v_0\) be the unit vector giving the coefficient and \(E_{ij}\) the derived elementary matrix. For a compact-sign pair the adjoint of \(E_{ij}\) is \(E_{ji}\), and for opposite signs it is \(-E_{ji}\). If \(\Omega=\sum E_{ij}E_{ji}\) and \(\Omega_K\) is the sum on compact pairs, it follows that
\[
        \sum_{ij}\|E_{ij}v_0\|^2
        =2\langle v_0,\Omega_K v_0\rangle-\langle v_0,\Omega v_0\rangle .
\]
The compact scalar on the highest vector is \(\sum_i\Lambda_i(\Lambda_i+(2\rho_K)_i)\). Its absolute value is at most \(nN(N+n)\), where \(N=L+n\), since \(|\Lambda_i|\le L+n\) and \(|(2\rho_K)_i|\le n\). The full scalar is \(\sum_i(l_i^2-(2\rho)_i^2)/4\), hence at least \(-n^3/4\). This eigenvalue follows by differentiating the explicit Flensted--Jensen integral, before identifying the cyclic representation. At each evaluation point, freeze the block-diagonal factor in the numerator and differentiate the principal-minor product in the denominator. For fixed \(k\in K\) the resulting germ is a lowest-weight vector in the conjugated basis \(kE_{ij}k^{-1}\). The Casimir acts by the same scalar on each germ, and differentiation commutes with the compact average.\formalnote{7.39}{39}
Since \(|X_{ij}|\le\|X\|\), Cauchy--Schwarz over \(n^2\) entries yields
\[
 \|d\pi(X)v_0\|
 \le n\sqrt{2nN(N+n)+n^3/4}\,\|X\|
 \le(n+1)^4(n+L+1)\|X\|.
\]
For the last bound use \(N+n\le2N\), \(n\le N\), and \(3n^{3/2}\le(n+1)^4\). Matrix coefficients against a unit vector are bounded by the norm of their second vector. This proves the uniform right derivative estimate at every \(g\).\formalnote{7.40}{40}

Set \(D_0=n+L+1\). For \(\|g-1\|<1\), write \(A=g-1\), \(u=1+(1-t)A/2\) and \(\gamma(t)=1+tAu^{-1}\). The difference \((\gamma u)^*J(\gamma u)-u^*Ju=t(g^*Jg-J)\) vanishes, so this curve joins \(1\) to \(g\) inside \(G\). On \(0\le t\le1\), the Neumann series gives \(\|u^{-1}\|\le2\), hence \(\|\gamma^{-1}\|=\|\gamma\|\le3\) and \(\|\gamma'\|\le4\|A\|\). Integration of the right derivative gives \(|c_\lambda(xg)-c_\lambda(x)|\le 100D_0^5\|A\|\). The same bound for \(\|A\|\ge1\) follows from \(|c_\lambda|\le1\). In particular \(|c_\lambda|\ge1/2\) on the ball of radius \(\delta=(200D_0^5)^{-1}\) about \(1\).

Write \(d_\lambda^{-1}=\int_G|c_\lambda|^2\,dg\). This ball has mass at most \(4/d_\lambda\). In the \(2n^2\)-dimensional real normed matrix space, any maximal \(\delta/R\)-separated set of points of \(G\) of norm at most \(R\), \(R\ge1\), has size at most \((1+2R^2/\delta)^{2n^2}\) by packing disjoint balls of half that separation. The same centers cover the set by balls of radius \(\delta/R\). If \(x\) is a center and \(\|g-x\|\le\delta/R\), then \(\|gx^{-1}-1\|\le\delta\) because \(\|x^{-1}\|=\|x\|\). Right invariance gives
\[
 d_\lambda\,\operatorname{vol}\{\|g\|\le R\}
       \le 4(2000D_0^6)^{2n^2}R^{4n^2}.
\]
Finally \(\|g\|\ge1\), and \(r^{-c}=\int_r^\infty cR^{-c-1}\,dR\). Integrating the ball estimate costs \(c/(c-4n^2)\le2\) when \(c>8n^2\). We obtain the normalized moment used in the tail estimate,\formalnote{7.41}{41}
\[
     \int_G d_\lambda(2/\|g\|)^c\,dg
            \le 8(2000D_0^6)^{2n^2}2^c.                              \tag{RD2}
\]


\subsection{Irreducibility and the lowest compact type}
In the range \(h\ge 500n^3+10n^5\) used for \S\ref{sec:elliptic-character}, write \(v=v_0\) for the unit generator constructed above, \(d=d_\lambda\), and \(\pi\) for the cyclic representation generated by \(v\). The coefficient identity
\[
 d\int_G c_\lambda(x^{-1}hy)\overline{c_\lambda(h)}\,dh
        =c_\lambda(x^{-1})c_\lambda(y)
\]
identifies \(\pi(d\,\overline{c_\lambda})\) with orthogonal projection onto \(\mathbb C v\), by testing on the dense span of translates of \(v\). Any bounded operator commuting with \(\pi\) preserves \(\mathbb C v\) and therefore acts by a scalar on this dense span. This proves irreducibility in \(L^2(G)\).\formalnote{7.42}{42}

To identify its parameter, for each dominant compact weight \(\nu\) let \(m_\nu\) be the normalized highest-vector coefficient on \(K\) and \(V_\nu\) the associated compact representation. By compact Schur orthogonality \cite[Corollary~4.10]{theta-knapp}, \(\dim V_\nu\int_K\overline{m_\nu(k)}\pi(k)\,dk\) projects onto the highest-weight lines of type \(\nu\). Its matrix entries on translates of \(v\) are given by the integral of \(\overline{m_\nu(k)}c_\lambda(x^{-1}ky)\). The compact averaging identities for the Flensted--Jensen integral make this projector zero unless \(\nu-\Lambda\) lies in the cone spanned by all positive roots, and give the projection onto \(\mathbb C v\) at \(\nu=\Lambda\). Since \(\Lambda+2\rho_K=\lambda+\rho\) pairs positively with all the positive roots, \(\Lambda\) is the unique weight of an occurring compact type minimizing \(\|\nu+2\rho_K\|\).\formalnote{7.43}{43}

By the discrete-series classification for \(U(p,q)\), the irreducible representation \(\pi\) has a regular parameter \(\lambda'\), chosen dominant for the same compact positive roots \cite[\S2.2]{theta-baldoni-vergne}. Let \(\rho'\) be the half-sum of all roots positive for \(\lambda'\), and put \(\rho_n'=\rho'-\rho_K\). For the compact highest weight \(\nu\), the compact Harish--Chandra parameter denoted by \(\mu\) in the cited formula is \(\nu+\rho_K\). Each nonzero summand in the Blattner formula counts expressions for \(w(\nu+\rho_K)-(\lambda'+\rho_n')\) as sums of positive noncompact roots \cite[\S2.2, eq.~(1)]{theta-baldoni-vergne}. Because \(\nu+\rho_K\) is compact dominant, \((\nu+\rho_K)-w(\nu+\rho_K)\) is a nonnegative sum of compact positive roots. Adding gives \(\nu-(\lambda'+\rho'-2\rho_K)\) in the full positive-root cone of \(\lambda'\).

The weight \(\lambda'+\rho'-2\rho_K\) occurs with multiplicity one. At that weight the identity summand counts the empty sum, and each nonidentity \(w\) gives a nonzero negative sum of compact roots because \(\lambda'+\rho_n'\) is strictly dominant on the compact roots. Indeed on each sign block the coordinates of \(\rho_n'\) are nonincreasing in parameter order. Thus \(\lambda'+\rho'-2\rho_K\) is the unique minimizer of the same shifted norm, and equals \(\Lambda\). The strictly ordered coordinates of \(\Lambda+2\rho_K=\lambda'+\rho'\) determine the full chamber and hence \(\rho'\). Subtracting it gives \(\lambda'=\lambda\), with the central coordinate included.
\phantomsection\label{route:11}\leanroute{Lean obtains the multiplicity and root-cone support by analyzing the cyclic representation and uses them to identify a polynomial for the orbital integral. Here the compact averaging identity identifies the lowest type of the actual coefficient, and the Blattner formula recovers its Harish--Chandra parameter for use in the published character formula. The pointwise and moment bounds (RD1)--(RD2) still use the actual Flensted--Jensen coefficient.}\formalnote{7.44}{44}


\section{Short frames and reduction at the height of a point}\label{sec:height-reduction}
The absolute sum in \S\ref{sec:bessel-bounds} and the kernel bound in \S\ref{sec:cusp-trace} use two different conclusions from reduction theory. Finite covolume suffices to integrate a bounded function on the quotient. A bound that varies with the point also requires a short frame for its finite lattice and control of the real factors that move that point into a compact set. We give both conclusions with the height used by the proof.

\subsection{The dual of the finite lattice}
Write \(E/F\) for the CM extension, \(f=[F:\mathbb Q]\), \(n>0\) for the rank, \(h\) for the nonsingular Hermitian form, and \(x\) for an adelic frame. In its coordinates \(h\) has Gram matrix \(1\) at finite places and a diagonal sign matrix \(S\) at the infinite places. Set
\[
 L=E^n\cap x_f\widehat{\mathcal O}_E^n,\qquad
 \iota(b)=(x_w^{-1}b_w)_{w\mid\infty},\qquad
 \lambda=\min\bigl(1,\min_{0\ne b\in L}\|\iota(b)\|_2\bigr).
\]
The real Euclidean dimension is \(D=2fn\), with one embedding for each conjugate pair. The lattice \(L\) lies between two rational integral multiples of \(\mathcal O_E^n\) allowing a common denominator, so \(\iota(L)\) is a full real lattice. Its completions are the prescribed \(x_v\mathcal O_{E,v}^n\) by approximation. With the involution pairing the places above a split place, these finite lattices together are self-dual for the \emph{Hermitian} pairing. For the \emph{trace} pairing there is an inverse different. In fact if \(\operatorname{Tr}_{E/\mathbb Q}h(b,a)\in\mathbb Z\) for every \(b\in L\), the same condition holds with \(ub\) for all \(u\in\mathcal O_E\). By the definition of the inverse different this puts \(h(L,a)\) in \(\mathfrak D_E^{-1}\). Thus any positive integer \(N\in\mathfrak D_E\) satisfies \(Na\in L\), by the finite self-duality.\formalnote{7.45}{45}

Consider the real dual of \(\Lambda=(2/N)\iota(L)\), with the dual defined by integral Euclidean pairings. Given \(X\in\Lambda^*\), its pairing with \((2/N)\iota(b)\) extends to a rational linear form on \(E^n=L\otimes_{\mathbb Z}\mathbb Q\). The trace pairing of a nonsingular Hermitian form is nondegenerate, hence the form is \(b\mapsto\operatorname{Tr}h(b,a)\) for a unique \(a\in E^n\). Comparing the two expressions at infinity gives
\[
  2\sum_{w}\Re\langle\iota(b)_w,S_w\iota(a)_w\rangle
     = (2/N)\sum_{w}\Re\langle\iota(b)_w,X_w\rangle .
\]
The real span of \(\iota(L)\) is full, so \(X=N S\iota(a)\). Since \(Na\in L\) and \(S\) is Euclidean unitary, every nonzero such \(X\) has length at least \(\lambda\). This accounts for the \(2\) arising from complex conjugate pairs and the \(N\) required by trace duality.\formalnote{7.46}{46}

Banaszczyk's Euclidean transference theorem, applied to \(\Lambda^*\), gives \(\mu(\Lambda)\lambda_1(\Lambda^*)\le D/2\), and hence the required upper bound \(D^2\) \cite[Theorem~2.2]{theta-banaszczyk}. Put \(r=D^2/\lambda\). The vectors of \(\Lambda\) of length at most \(3r\) span the real space. Otherwise choose a unit vector perpendicular to their span and approximate twice \(r\) times this vector by a lattice point within \(r\). The point has length at most \(3r\) and has positive projection in the chosen direction. Rescaling shows that the elements of \(L\) of length at most
\[
                       \frac{3ND^2}{2\lambda}
                          =\frac{6N(fn)^2}{\lambda}
\]
span \(E^n\). Indeed an \(E\)-linear form annihilating them would annihilate their real span after extending to the archimedean coordinates. Select \(n\) independent columns. These give an \(E\)-basis contained in \(L\), exactly what the small-box argument of \S\ref{sec:cusp-trace} uses.
\phantomsection\label{route:12}\leanroute{Lean proves the required weak transference bound by successive projections and a Minkowski estimate. Banaszczyk's Euclidean transference theorem applies to the full lattice in the real trace-dual coordinates computed above. Weakening its constant gives exactly the bound used to select the frame.}\formalnote{7.47}{47}

\subsection{Halving the height in a mixed signature}
Write \(H(g)=1/\lambda\) for the frame moved by \(g\). If \(H(g)\le2\), the preceding construction gives uniformly bounded independent columns of the finite lattice. In the moved coordinates these columns are integral at the finite places. Their rational Gram matrices belong to a compact subset of the adele matrices, hence to a finite set. Choosing one rational basis for each occurring Gram matrix supplies a compact set of representatives for these \(g\) modulo the rational isometry group, as in \S\ref{sec:bessel-bounds}. When the fixed form is given by a rational Gram matrix \(J\), the rational frame isometry for two bases \(Y,Y_0\) with the same Gram matrix is \(Y_0Y^{-1}\). The transformed inverses are bounded adelically in a closed unitary group (or its translate by the initial frame). This proves the height-two compactness needed below.\formalnote{7.48}{48}

At an arbitrary point the finite pairing of two lattice columns is integral. If both columns have length \(<1\), every infinite conjugate of their pairing has absolute value \(<1\) by Cauchy--Schwarz. An algebraic integer with this property is zero. At each infinite place the columns of this short set therefore span a totally isotropic space \(A_w\). Inside \(\mathbb C^n\) with its usual positive Euclidean product, \(A_w\) and \(S_w A_w\) are orthogonal. On their direct sum take \(y_w\) to be \(2\) on \(A_w\) and \(1/2\) on \(S_w A_w\), and take \(1\) on the orthogonal complement. The sign involution exchanges the two summands and preserves the complement, so \(y_w^*S_w y_w=S_w\). Both \(y_w\) and its inverse have norm at most two.

At the new point \(g(y,1)^{-1}\) the finite lattice is unchanged and the infinite coordinates are multiplied by \(y\). Each formerly short vector doubles in length. The others still have length at least \(1/2\). Thus
\[
      H\bigl(g(y,1)^{-1}\bigr)\le\max\{H(g)/2,\,2\}.
\]
For \(2^k\le H(g)<2^{k+1}\) perform \(k\) steps and use height-two compactness. We have written a rational translate of \(g\) as \(z(y_k\cdots y_1,1)\), with \(z\) in a fixed compact set and every factor of norm at most two. If an action has operator norm bounded by \(2^a\) on this compact set and on the real norm-two ball, the product has norm at most \(2^{a(k+1)}\le 2^aH(g)^a\). These are precisely the factors used before integrating the orbit sum.\formalnote{7.49}{49}

Borel's adelic finiteness theorem gives the finite volume. The connected reductive group here is \(U(h)\) over the number field \(F\). Over an algebraic closure its character group is generated by the determinant. The nontrivial automorphism of \(E/F\) sends this generator to its negative in the character group, so there is no nontrivial \(F\)-rational character. Borel's criterion \cite[Theorem~5.6(i)]{theta-borel} gives finite invariant volume to \(U(h)(F)\backslash U(h)(\mathbb A_F)\). If the form is definite at infinity, a rational unipotent must be the identity, since its image in a compact unitary group is diagonalizable. The compactness criterion \cite[Theorem~5.6(ii)]{theta-borel} then gives the definite compact quotient used in \S\ref{sec:bessel-bounds}. The shrinking of a radical at fixed finite level in the mixed case uses the further column and approximation estimates in \S\ref{sec:cusp-trace}.
\phantomsection\label{route:13}\leanroute{The formal reduction theory proves finite volume by integrating a count of rational matrices and specializing its column number to zero. Borel's theorem gives the finite-covolume and compactness assertions for the indicated rational unitary group. The pointwise short-frame and height-halving estimates above supply additional quantitative information.}\formalnote{7.50}{50}


\section{Absolute convergence of the Bessel unfolding}
\label{sec:bessel-bounds}

In \S\ref{sec:orbit-unfolding} a rational orbit is discarded when its character average vanishes. This procedure is legitimate only if one can group and integrate the orbits in the order used. We give the required bound uniformly in the subset of rational maps being summed. It uses two properties of the construction: the positive-degree entries are a null frame, and the definite quotient in the source Bessel group has bounded height.

\subsection{From map classes to rational slices}
Let \(g_i\) be the source grades and \(d_j\) the target weights, assumed odd under reversal of the target strings. Give the \((j,i)\)-entry of a map \(V\to V'\) degree \(d_j-g_i\). Write \(J\) for the Gram matrix on \(\operatorname{Hom}(V,V')\) in these rational coordinates. In the source and target a basis vector pairs only with its reversed string index or the matching definite coordinate. The product Gram matrix therefore has a zero entry between two coordinates whenever their degrees do not sum to zero. Its positive-degree coordinate vectors \(k\) are isotropic. Extend \(k\) to a Witt frame \((l,k,i_0)\) with \(k^*Jl=1\) and middle space \(V_0\).

Two rational maps give the same class modulo maps of nonnegative degree just when their negative-degree entries agree. The coordinate projector onto those entries is \(P=J^{-1}kk^*J\): \(kk^*\) is the positive-coordinate projector and \(J\) interchanges opposite degrees. Every class has a representative \(l\xi\) with unique \(\xi\) in the rational coordinate space of the frame. Indeed for any \(w\) take \(\xi=k^*Jw\); then \(P l\xi=J^{-1}k\xi=Pw\). If \(P l\xi=0\), injectivity of \(k\) gives \(\xi=0\).

The Fourier term of the map \(w\) is the theta of the zero term after Heisenberg translation by \(w\). On rational maps the functional is unchanged when a map of nonnegative degree is added: exchange the translation past the averaging variable of positive degree. The extra phase is trivial since there are no opposite degrees to pair, and the rational translation on the theta functional is trivial. At the representative \(l\xi\) the term is therefore precisely \(\theta_{V_0}(S_\xi f)\) in the notation of \S\ref{sec:vacuum-projection}. By that section its sum is \(f(0)\), and its absolute sum is at most
\[
                 \|\theta_{V_0}\|\sum_{\xi}\|S_\xi f\|
                        \le C\|f\|.
\]
The same estimate applies to any subset of the classes, before an orbit has been selected.\formalnote{7.51}{51}

\subsection{Compact representatives for the Bessel group}
Write \(B\) for the source Bessel group and \(X\) for the definite part of the grading. Projection onto \(U(X)\) has a section \(a\mapsto i_X(a)\) over the ground field, and the kernel consists of grade-raising unipotents. First recall the bounded-frame argument. Work in rational coordinates, so the matrix of a unitary point \(g\) preserves the rational form \(J_V\). Suppose an invertible rational frame \(r\) in a prescribed subgroup \(H\subset{\rm GL}(V)(E)\) has \(g^{-1}r\) in a compact set of adelic matrices \(C\). Then \(r^*J_Vr=(g^{-1}r)^*J_V(g^{-1}r)\) belongs to a compact adelic set. Its entries are rational and rational points are closed and discrete, so there are only finitely many such Gram matrices. For each occurring matrix \(T\) choose one rational \(r_T\in H\) of that Gram matrix. Now \(\gamma=r_T r^{-1}\) is a rational isometry in \(H\), and \((\gamma g)^{-1}=(g^{-1}r)r_T^{-1}\) lies in a finite union of compact sets. The unitary group is closed in the adelic matrix space (the inverse is given by the adjoint), so \(\gamma g\) lies in a compact set of unitary points.

For \(U(X)\), a nonzero vector of the integral finite lattice of archimedean length strictly less than one would have algebraic integral self-pairing of absolute value less than one at every embedding. Its self-pairing would be zero by taking the field norm. Definiteness rules this out. Thus the first minimum cut off at one is always one, and the short-basis bound supplies a rational frame of uniformly bounded archimedean length with integral finite part. Applying the bounded-frame argument gives a compact set \(D\) of representatives for the definite quotient.\formalnote{7.52}{52}

For the unipotent part, choose a compact \(A\subset\mathbb A_E\) and rational shifts \(q(c)\) such that \(c+q(c)\in A\). Let \(m\) be a grade-raising unipotent, let \(n=m^{-1}\) in rational coordinates, and order the grades using \(w=-g\). Now \((n-1)_{ij}\) can be nonzero only for \(w_i<w_j\). Construct a rational unitriangular matrix \(r\) by putting \(r_{ij}=\delta_{ij}\) for \(w_i\ge w_j\), and otherwise
\[
                            r_{ij}=q\bigl(((n-1)r)_{ij}\bigr).
\]
In a fixed column, the right side uses only rows of strictly larger \(w\); descend through the grades. Consequently \(nr=r+(n-1)r\) has entries in \(A\cup\{0,1\}\), and \(\det r=1\). The constructed \(r\) need only lie in the rational unitriangular group \(H\). The bounded-frame argument with this \(H\) gives an actual unitary \(\gamma\in H\) with \(\gamma m\) in a fixed compact set \(K\). Such a \(\gamma\) lies in \(B\) by the grade condition. For an arbitrary \(b\), first reduce its definite component into \(D\) by the section applied to a rational element of \(U(X)\), then reduce its unipotent component into \(K\). The resulting rational translate belongs to \(K i_X(D)\cap B\), compact since \(B\) is closed. This proves compactness of \([B]\) and finiteness of its quotient measure.\formalnote{7.53}{53}

\subsection{The uniform sum on the two quotients}
For equivariant \(W\), the class term at \((h,b)\) uses the datum \(W(b)\omega(b,h)f\), with the prescribed unitary twists understood. Rational translations of \(b\) and rational target parabolic translations of \(h\) just permute the class terms; both the phase of \(W\) and the Weil phase are accounted for by its equivariance. Reduce \(b\) to the compact set above. Continuity bounds \(\|W(b)\omega(b,1)f\|\) there. Apply the preceding absolute slice bound to \(\omega(1,\delta h)W(b)\omega(b,1)f\). For every subset \(D_1\) of classes and every rational target parabolic element \(\delta\) it gives
\[
            \sum_{w\in D_1}|\operatorname{term}_w(h,b)|
                     \le C_f\|\omega(1,\delta h)\|.
\]
When the target weight \(d\) is identically zero the allowed parabolic is the whole rational unitary group. Write \(H(h)\) for the reciprocal of its truncated first minimum. The height reduction gives a rational translate \(\delta h=z x_1\cdots x_k\), with \(z\) in a fixed compact set, each \(x_i\) a real unitary point of operator norm at most two at every place, and \(2^k\le H(h)\). On the two compact sets of factors the Weil operators have norm at most \(2^a\) for a fixed integer \(a\). Multiplication bounds \(\|\omega(1,\delta h)\|\) by \(2^{a(k+1)}\le 2^aH(h)^a\). For the rapidly decreasing automorphic target test \(y\), \(\sup_h |y(h)|H(h)^a<\infty\). The full absolute sum is invariant under the rational groups by reindexing. On representatives, all its partial sums weighted by \(|y|\) are therefore uniformly bounded on the product of finite-volume quotients. This proves absolute integrability, permits termwise integration over \(B\), and permits unfolding the final sum over target translates as an ordinary quotient integral. The quantitative reduction by real factors and the short-basis transference theorem are the lattice results used here.\formalnote{7.54}{54}


\editionchapter{Signature transfer and the target dimension}\label{ch:trace}

Theorem B needs an upper bound for the dimension of its target space. On a definite unitary group one can bound the corresponding dimension by counting lattice classes. To transfer the count, we compare traces in the two signatures with the same finite-level projector.

The background result compares traces of test functions that are cuspidal at a split place. Labesse's transfer theorem preserves the finite component on changing the unitary group under its hypotheses on cuspidal base change and the local components \cite[Theorem~5.5]{upper-labesse}. Here we compare the specified projected spaces directly. Each summand of the projector onto the good forms gives a test cuspidal at some split place, which may depend on the summand. Applying the trace comparison with spherical Hecke operators then gives the same multiplicity for each eigensystem on the two spaces of good forms. See the \rootlink{grh/src/Lib/Adelic/TraceFormula/CuspidalComparison.lean}{trace comparison} and its \rootlink{grh/src/New/Spanning/Multiplicities/EqualMultiplicities.lean}{application to good forms}.

\section{The signature comparison in the dimension formula}\label{sec:sieve-transfer}
\subsection{Equality of target dimensions across signatures}

The two Hermitian spaces \(V,V'\) may have different signatures at infinity. They are identified at every finite place, so we may use the same finite adelic group, principal level \(K\), and central idempotent \(e_0\) on both. In the application to Theorem B one of the spaces is definite. At infinity their kernels use the same doubled Harish--Chandra parameter
with the two choices of signs. For these data \(E\) contains a primitive
\(5040\)-th root of unity, its number of infinite places is even,
\(r\ge2\), and the large-gap hypothesis of the comparison holds for
both sign choices, with threshold \(500r^3+10r^5\) for the gaps between opposite-sign entries of the undoubled parameter. Section~\ref{sec:elliptic-character} gives the normalization.

Write \(S_V\) for the principal-level forms of this archimedean type, \(\pi_V\) for the
action of the residue groups, and \(Y_V=\pi_V(e_0)S_V\).
The central idempotent \(e_0\) is a sum of orthogonal idempotents
\(e_c=\bigotimes_w e_{c_w}\) indexed by the good families of
polynomial labels. A spherical Hecke operator \(T\) away from
the level commutes with \(e_0\), so
\[
       \operatorname{tr}(T\mid Y_V)
       =\operatorname{tr}(T\pi_V(e_0)\mid S_V)
       =\sum_c\operatorname{tr}(T\pi_V(e_c)\mid S_V).
\]
The first equality follows by taking the trace on the image and
kernel of the projection. Fix one \(c\). At some split place
\(w\) above \(p_0\) its label is primitive and in particular
irreducible. The corresponding finite general linear group
representation \(\sigma\) is cuspidal by Green's character formula.
For a proper parabolic radical \(U_P\) its averaging operator
\(\sum_{u\in U_P}\sigma(u)/|U_P|\) is zero, since \(\sigma^{U_P}=0\).
Multiplying on the left and right by \(\sigma(a),\sigma(b)\)
and taking traces gives \(\sum_u\chi_\sigma(aub)=0\).
The coefficients of \(e_{c_w}\) are proportional to
\(\overline{\chi_\sigma}\), so the same cancellation holds for all
its translates. Inflating the tuple to \(K\), multiplying by the
common finite Hecke element representing \(T\), and averaging over
\(K\) produces a common test \(\phi\) cuspidal at \(w\).\formalnote{8.1}{1}

We compare the terms \(I_R(e,\phi)\) of the trace expansion in \S\ref{sec:signature-trace}, indexed by rational characteristic polynomials \(R\).
Choose a second split place \(v\) away from \(w,\bar w\), where
the level is maximal and the finite group-algebra element is
integral. These requirements exclude only finitely many split
places. Cut the level at \(v\) to the stabilizer of a residue line.
At this level positive powers of the line Hecke operator \(T^+\) force zero to be a simple root of the characteristic polynomial modulo \(v\) on the support. In fact, at this split place \(T^+\) is supported on the double coset of \(\operatorname{diag}(\varpi_v,1,\ldots,1)\) for the line stabilizer. The reduction of every matrix in this double coset kills the chosen line and induces an invertible map on the quotient. Products preserve both properties. Cuspidality at \(w\) persists. The
geometric side, equal to the level volume times the trace, is
grouped by rational characteristic polynomials \(R\). The two
finite test conditions eliminate different terms. If \(R\) is a proper power, its monic root is integral at \(v\), and each root of the reduction of \(R\) has multiplicity at least two. If the characteristic polynomial has two coprime factors, integration along a proper unipotent radical at \(w\) gives zero by cuspidality. If \(R\) is
irreducible with some archimedean root off the unit circle,
the discrete-series orbital integral vanishes. For an irreducible
\(R\) with all roots on the circle, the class matching and the
archimedean character identity give
\[
        C(R,\phi')\prod_{u\mid\infty}
                      \overline{\chi^{\rm comp}_{\lambda_u}(z_u)} .
\]
Here \(\phi'\) is the test at the cut level, \(z_u\) is the diagonal
torus element of the numbered roots, and the compact character is
that of the lowest compact type for the all-positive signature.
The matching constant \(C(R,\phi')\) is independent of the real
signs. Summing over \(R\) yields equality of the cut-level mixed
traces \(\operatorname{tr}(A^j(T^+)^m)\) for all \(j,m\ge1\),
where \(A\) represents \(T e_c\) off the auxiliary place.\formalnote{8.2}{2}

To recover the original trace, let \(S\) be the cut-level space
and let \(e\) average over the full maximal compact at \(v\). On \(S\),
\(T^+(T^+)^*=d(1+c'e)\) for local constants \(c',d>0\), and
\(A\) commutes with both operators. On the generalized
\(\alpha\)-eigenspace of \(A\), the mixed trace is
\(\alpha^j\operatorname{tr}((T^+)^m)\), because the other
terms in the binomial expansion contain a commuting nilpotent.
Vandermonde inversion in \(j\) separates the nonzero \(\alpha\)'s.
The power sums in \(m\) then determine the eigenvalues of
\(T^+\) with multiplicity on that space. It is invertible by the
displayed local relation, so they determine its dimension \(a\)
and its determinant. Taking determinants in that relation gives
\[
            |\det T^+|^2=d^a(1+c')^{\operatorname{rank} e}.
\]
Thus the rank of the full-level projection agrees on each nonzero
generalized eigenspace for \(V,V'\). Multiplying each rank by
\(\alpha\) and summing recovers \(\operatorname{tr}(A\mid eS)\)
at the original level, with the zero eigenspace contributing zero.
The trace of \(Te_c\) therefore agrees. Summing over \(c\) and
taking \(T=1\) proves \(\dim Y_V=\dim Y_{V'}\), in particular for
the definite target and each \(W_j\). Section~\ref{sec:signature-trace} derives the line Hecke relation,
the contribution of the elliptic classes and the use of the two auxiliary
places. It gives the local and analytic inputs to this comparison.\formalnote{8.3}{3}


\section{Transfer of the trace between signatures}
\label{sec:signature-trace}
On a definite unitary group, a form of fixed weight and level is
determined by its values on a finite set of lattice classes. We use
this description to bound the dimension of the good space in \S\ref{sec:sieve-transfer}. On each of the two unitary groups
the dimension is the trace of the projector onto the chosen forms.
We compare the orbital terms in these two traces.

Let \(V,V'\) have rank \(n\ge2\), the same finite adelic group, the same finite level and the same doubled integral parameters, with gaps large enough for the corresponding Harish--Chandra parameters in both signatures. The CM field \(E\) contains the \(5040\)-th roots of unity and has an even number \(f\) of complex places. Two local conditions are used in the comparison. An auxiliary split place \(v\) excludes characteristic polynomials that are proper powers. The split place \(w\), where the finite test function is cuspidal, excludes the remaining reducible polynomials by averaging along a unipotent radical.

\subsection{The local operator and its two eigenvalues after taking a square}
At \(v\), put \(K=\mathrm{GL}_n(O_v)\), let \(q\) be the size of its residue field \(k\), and put \(J=\{g\in K:\bar g\,k e_1=k e_1\}\). Let \(D=\operatorname{diag}(\varpi,1,\ldots,1)\). In a unitary representation write \(T^+,T^-\) for the sums of right-coset operators for \(J D J\) and \(J D^{-1}J\) on the \(J\)-fixed vectors. They are adjoint. In fact the coset sums are integrations over the two inverse double cosets for Haar measure normalized by \(\operatorname{vol}J=1\), and inversion interchanges these integrals upon taking adjoints. Let \(e_K\) be averaging over \(K\) with mass one. We claim
\[
          T^+T^-=q^{n-2}\bigl(1+(q^n-1)e_K\bigr).       \tag{TR1}
\]
In particular \(T^+\) is surjective on each finite-dimensional invariant space where (TR1) holds and both operators act. The two eigenvalues of the positive product on the nonspherical and spherical parts are \(q^{n-2}\) and \(q^{2n-2}\) respectively, when the corresponding parts occur.

For \(y\) fixed by \(J\) and \(z\in k^n\setminus\{0\}\), denote by \(y[z]\) the translate by any member of \(K\) carrying the line of \(e_1\) to that of \(z\). Two such elements differ on the right by \(J\). Cosets upstairs are indexed by forms \(\xi\) with \(\xi(e_1)=1\) and cosets downstairs by vectors \(z\) with first entry \(1\). Explicit representatives are \(B_\xi^{-1}D\) and \(D^{-1}C_z\), where \(B_\xi\) and \(C_z\) are integral lifts of the identity matrix with the first row replaced by \(\xi\) and the first column replaced by \(z\), respectively. To check the indexing, intersect \(J\) with its conjugate by \(D\) (respectively \(D^{-1}\)). The remaining datum for the raising coset is the kernel of a form nonzero on \(e_1\); for the lowering coset it is a line not in the complementary hyperplane. Normalizing the form on \(e_1\), or the vector at its first entry, gives the indicated unique parameter. The product of the chosen representatives is \(B_\xi^{-1}C_z\). Its image of the first basis vector runs over the vectors \(z'\) with \(\xi(z')=1\). Thus
\[
            T^+T^-y=\sum_{\xi(e_1)=1}\ \sum_{\xi(z)=1}y[z].
\]
For \(z=e_1\) the number of forms is \(q^{n-1}\); for the other vectors on its line it is zero. For any \(z\) independent of \(e_1\) the two linear conditions leave \(q^{n-2}\) forms. The \(q-1\) nonzero multiples of \(e_1\) all give \(y\). Factoring out \(q^{n-2}\) consequently leaves \(y+\sum_{z\ne0}y[z]\). The group \(\mathrm{GL}_n(k)\) is transitive on nonzero vectors, so this last sum equals \((q^n-1)e_Ky\), proving (TR1). Modulo \(v\), all members of \(JDJ\) kill the first line and induce an invertible map on the quotient. Both properties persist in a nonempty product. Its reduced characteristic polynomial has the factor \(X\) exactly once. This explains why every positive power of \(T^+\) imposes the same support condition on the test.\formalnote{8.4}{4}

\subsection{The trace and the two sorts of reducible polynomial}

\begin{example}[Unipotent conjugation in rank two]\label{ex:radical-conjugation}
Take distinct \(a,b\in F^\times\) in a local field \(F\). Put
\(y=\operatorname{diag}(a,b)\) and write \(n(z)\) for the upper
triangular unipotent with upper right entry \(z\). Then
\[
 \begin{pmatrix}1&z\\0&1\end{pmatrix}
 \begin{pmatrix}a&0\\0&b\end{pmatrix}
 \begin{pmatrix}1&-z\\0&1\end{pmatrix}
       =\begin{pmatrix}a&(b-a)z\\0&b\end{pmatrix}.
\]
The matrix on the right is \(y n(t)\) with \(t=(b-a)z/a\). For
additive Haar measure and a compactly supported test function \(\phi\),
\[
 \int_F\phi\bigl(n(z)y n(-z)\bigr)\,dz
   =\left|\frac{a}{b-a}\right|_F\int_F\phi\bigl(y n(t)\bigr)\,dt .
\]
Thus a vanishing integral along \(yN\), with \(N=\{n(t):t\in F\}\),
gives a vanishing integral over these conjugates of \(y\).
With two diagonal blocks \(y_1,y_2\), the upper right entry is
\(Zy_2-y_1Z\). Disjoint spectra make this linear map invertible.
\end{example}

Use \(K'\) for the global level cut to \(J\) at \(v\). If \(a\) is the corresponding group-algebra operator, use at finite places \(\phi=a*1_{K'}\), where the indicator is unnormalized. Let \(e\) be the product of the discrete-series idempotent kernels at infinity. It has decay order \(B=10n^5\) with \(B\ge13n^3\) in the prescribed gap range. In the automorphic kernel
\[
    {\cal K}(x,x')=\sum_{\gamma\in U(V)(F)}
       (e\otimes\phi)(x^{-1}\gamma x')
\]
we take the right-action convention in which translating the output \(x\) represents \(a\). For \(\phi=1_{K'}\) integration against a vector of level \(K'\) and type \(e\) reproduces the vector times \(\operatorname{vol}K'\): unfold the rational sum, apply \(e*e=e\) on its image and integrate the finite-level indicator. The conjugate kernel with one point fixed belongs to this same space by the kernel bound and the compact reproducing function. Expanding it in any orthonormal basis \(b_i\) of that space gives \({\cal K}(x,x')=\operatorname{vol}K'\sum_i b_i(x)\overline{b_i(x')}\). By linearity in finite translates the right side for \(\phi=a*1_{K'}\) has \((a b_i)(x)\) in place of \(b_i(x)\). On setting \(x'=x\) and integrating, orthonormality gives
\[
   \operatorname{vol}K'\operatorname{tr}(a)
     =\sum_R I_R(e,\phi),                                \tag{TR2}
\]
where \(I_R\) integrates those rational terms with characteristic polynomial \(R\).

The convergence issue in (TR2) is control uniformly over partial polynomial sums. In bounded height the sum of absolute kernel terms is at most \(C H(x)^{4fn^2}\). In a cusp choose an isotropic frame with \(\delta^nH(x)=\kappa\), a constant independent of \(x\). Split the sum into those rational points preserving its span and those moving the span. The reduction and cusp-kernel estimates bound the first partial kernel in modulus by \(C_2(\delta^nH(x))^{4fn^2}\), and the sum for the movers by \(C_1\kappa^{4fn^2}\), uniformly in the set of polynomials. If \(\delta\ge1\) the bounded-height estimate already suffices. Thus all partial sums of polynomial terms are bounded by a constant on a quotient of finite measure. At a fixed point the rational series converges absolutely. Dominated convergence gives (TR2), including the grouping by \(R\). The reduction and stabilizer bounds, and the fact that the reproducing kernel is a form, enter at these indicated points.\formalnote{8.5}{5}

Suppose \(R\) contributes. The support at \(v\) forces a simple zero on reduction. Consequently \(R\) cannot be a proper power. Indeed a monic divisor of the characteristic polynomial of an integral matrix has integral coefficients over a discrete valuation ring, by integrality of all its roots. If \(R=P^j\) with \(j>1\) we can normalize \(P\) to be monic; a root of its reduction would contribute multiplicity at least \(j\). A reducible monic \(R\) left over has two coprime factors \(A,C\) of positive degree: split off an entire power of one irreducible factor. They stay coprime over the local field at \(w\). At a matrix \(y\) with characteristic polynomial \(R\) the Bezout identity gives a polynomial projector onto the summand \(\ker A(y)\) along \(\ker C(y)\). Thus \(y\) has two diagonal blocks \(y_1,y_2\) with disjoint spectra. Put \(N=\{1+Z:Z\in\operatorname{Hom}(\ker C(y),\ker A(y))\}\), the corresponding abelian unipotent radical. Conjugating \(y\) by \(1+Z\) fills \(yN\): the map \(Z\mapsto Zy_2-y_1Z\) is invertible since its kernel would intertwine modules with coprime annihilating polynomials.

The functional \(I_R\) is invariant under rational-sum change of variables and finite adelic conjugation, and depends only on the local fibre with polynomial \(R\). Average the test over conjugation by the integral compact at \(w\). On the slice where the polynomial projector is the fixed coordinate projector, integration along each radical as above is zero: the cuspidal test has sum zero over the finite radical translates, each of which can be effected by conjugation in \(N\); right Haar translation by these unipotents preserves measure. The local descent formula obtains the original functional from this slice by conjugation and a nonzero compact-centralizer volume. It therefore vanishes. Section~\ref{sec:projector-descent} constructs the descent to the projector slice and verifies compact support of its lifted tests. In particular we have specified how cuspidality at \(w\) deals with the case the simple-zero test at \(v\) leaves open.\formalnote{8.6}{6}

\subsection{Vanishing for eigenvalues off the unit circle}
Assume \(R\) irreducible and let \(\gamma\) have polynomial \(R\). Its eigenvalues are distinct in characteristic zero. At a complex place where an eigenvalue \(r\) has \(|r|\ne1\), take a nonzero left eigenvector \(v_r\) of \(\gamma\). Define
\[
          H(x)=\log\|v_r x_\infty\|,\qquad \ell=\log|r|.
\]
Only this place is used in \(H\). Then \(H(\gamma x)=H(x)+\ell\). Every element of the rational centralizer acts by a scalar on the same left eigenline, so also shifts \(H\) by a constant. Right multiplication by \(b\) changes \(H\) by at most \(\max(\log\|b\|,\log\|b^{-1}\|)\) at that place; the two operator norms coincide in a signed unitary group. This bound gives an integrable weight against the triple kernel below by the moment estimate for the cusp kernel.

Write \(p(x,y)=e(x^{-1}y)\) at infinity and use \(e=e*e*e\) in the orbital integral of \(\gamma\). With \(x\) modulo its rational centralizer and two real intermediate points \(x_1,x_2\), the real factor becomes
\[
   p(x,x_1)p(x_1,x_2)p(x_2,\gamma x).
\]
Let \(a_1=H(x_1)-H(x)\), \(a_2=H(x_2)-H(x_1)\) and \(a_3=H(\gamma x)-H(x_2)\). Their sum is \(\ell\). The three integrals against these weights are equal: rotate the triple to \((x_1,x_2,\gamma x)\) in its real coordinates and keep the finite coordinate of \(x\). The real difference weights descend under the centralizer, the finite orbital factor is unchanged, and Haar changes of variables give the trace property on this quotient. Moreover the integral for \(a_1+a_2\) equals that for \(a_1\). To see this, in the former the weight is \(H(x_2)-H(x)\), so convolution integrates out \(x_1\). In the latter integrate out \(x_2\). Both leave the same two factors and the same difference after renaming the intermediate point. Hence the integral for \(a_2\) vanishes, and the two others vanish with it. Thus \(\ell\) times the original orbital integral is zero. As \(\ell\ne0\), every such class gives zero.\formalnote{8.7}{7}

\subsection{Changing the signs of elliptic eigenlines}
For irreducible \(R\) with all roots on the circle, let \(L=E[X]/(R)\) and let \(\theta\) be the image of \(X\). On this field complex conjugation extends that on \(E\) by \(\theta\mapsto\theta^{-1}\); at every complex embedding the roots have absolute value one. Thus \(L\) is a CM field. The space of an elliptic \(y\) with polynomial \(R\) is a one-dimensional \(L\)-space through \(\theta\mapsto y\). Its \(L\)-eigenlines at each complex place are orthogonal for the given form: invariance under \(y\) gives \((1-\bar r_i r_j)\langle v_i,v_j\rangle=0\). Fix a numbering of the roots. To change the assigned signs, let \(\eta_{ui}\in\{\pm1\}\) describe the required changes, and suppose their total product is \(1\).

Quadratic idelic reciprocity for \(L/L^+\) says that the idele of \(L^+\) with these real signs and finite part \(1\) is principal times a norm precisely when its quadratic character is \(1\). Its character is the product of the signs. Consequently there are \(a\in (L^+)^\times\) with these signs and a finite idele \(h\) of \(L\) with \(a=h\bar h\) at the finite places (invert the finite norm factor in the principal-times-norm expression). Represent \(a\) by a polynomial in \(\theta\), and \(h\) by a polynomial with finite adelic coefficients. Set \(Z=h(y)\). For a Gram matrix \(\beta\) preserved by \(y\), the adjoint of \(y\) is \(y^{-1}\). Evaluating the norm identity therefore gives \(Z^*\beta Z=\beta a(y)\). Multiplication of \(\beta\) by \(a(y)\) changes its signs exactly as prescribed. The two pairs consisting of the old form with frame \(Zx\) and the changed form with frame \(x\) have identical finite Gram matrix. Also \(Z\) commutes with \(y\) and with its rational centralizer, since it is a polynomial in \(y\). The assignment is compatible with rational conjugacy and is invertible by using \(Z^{-1}\). This gives a bijection of the corresponding finite frame classes that preserves the represented finite point and the embedded centralizer. The even number \(f\) ensures the sign condition when passing from definite signs to any fixed signature repeated at every complex place: their product is \((\prod_i s_i)^f=1\).\formalnote{8.8}{8}

For an elliptic class its finite centralizer image is discrete in the setup of the comparison. One can integrate over it using a finite-adelic cutoff whose translates over the centralizer sum to one. Landherr's local-to-global isometry theorem identifies the representative of a specified signature with a point on the given global space. With this cutoff, Fubini expresses its orbital integral as the product of the finite orbital integral and one real orbital integral at each complex place. The matched class uses the same finite point and centralizer, hence the same cutoff and finite factor. Summing over classes with one fixed assignment of signs to the numbered roots gives a factor \(C(R,\phi)\) independent of that assignment.

For completeness, write \(I_u(\sigma)\) for the local orbital integral with the roots permuted by \(\sigma\), and \(W_s\) for the permutations preserving the signs. If the signature is \((p_+,p_-)\), then \(|W_s|=p_+!p_-!\). Each occurring class appears exactly \(|W_s|^f\) times as the orderings vary independently: at every place only the orders within the two signs are free. The local discrete-series character identity with the stated gap gives
\[
 (-1)^{p_+p_-}\sum_{\sigma\in S_n}I_u(\sigma)
     =|W_s|\,\overline{\chi^{\rm comp}_{\lambda_u}(z_u)}
\]
on the regular diagonal torus element \(z_u\) of numbered roots. Multiply over \(u\), use the evenness of \(f\) to cancel the sign, and divide by the overcount \(|W_s|^f\). This gives \(I_R(e,\phi)=C(R,\phi)\prod_u\overline{\chi^{\rm comp}_{\lambda_u}(z_u)}\) for either space. Together with the two vanishing arguments this proves equality of all the cut-level mixed traces in \S\ref{sec:sieve-transfer}. Its determinant argument using (TR1) now recovers the original spherical trace and the equality of good-target dimensions. Sections~\ref{sec:elliptic-character} and~\ref{sec:cusp-trace} derive the elliptic sum and the cusp estimate in (TR2). The finite type and the projector descent are treated in \S\S\ref{sec:green-cuspidal} and~\ref{sec:projector-descent}. Section~\ref{sec:classical-matching} applies Landherr's theorem to the prescribed rational hermitian classes and global reciprocity to the signs.\formalnote{8.9}{9}


\section{The finite character and the Hermitian forms in the trace comparison}\label{sec:classical-matching}
Three classical results enter the matching in \S\ref{sec:signature-trace}. Green's theorem identifies a particular function on a finite general linear group as an irreducible character. Global reciprocity allows an elliptic eigenline to change sign at infinity with a prescribed finite norm. Landherr's theorem puts the resulting form back on the space whose trace is being computed. We check the input to each result.

\subsection{Identifying the finite type from its values}
Let \(q\) be the residue cardinality and \(P\) a monic irreducible label of degree \(n\ge1\) with \(P(0)\ne0\). Choose a root \(a\in\mathbb F_{q^n}^{\times}\). The fixed faithful multiplicative character \(T\) of the algebraic closure identifies this cyclic group with the complex \((q^n-1)\)-st roots of unity. Write \(T(a)=\exp(2\pi i j/(q^n-1))\) and put \(\theta=T^j|_{\mathbb F_{q^n}^\times}\). The map \(a\mapsto\theta\) is a Frobenius-equivariant isomorphism with the character group, using the operation \(\theta\mapsto\theta^q\) there. Hence \(\theta\) has orbit of length \(n\) and changing the root of \(P\) just permutes the orbit. This is the regularity condition in Green's classification.\formalnote{8.10}{10}

To check the character used in the cuspidality argument, evaluate the finite class function specified by the label. Regard a matrix \(g\in{\rm GL}_n(\mathbb F_q)\) as a finite \(\mathbb F_q[t]\)-module, where \(t\) acts by \(g\). Write \(\mu(U)\) for the Möbius number of its submodule lattice from \(0\) to \(U\) and \(\beta(U)\) for the sum of \(T^j\) over the eigenvalues on \(U\), counted with multiplicities. The function at \(g\) is
\[
             (-1)^n\sum_{U\subseteq V\ {\rm stable}}\mu(U)\,\beta(U).
\]
In a direct sum of distinct primary components all submodules split, \(\mu\) is multiplicative and \(\beta\) is additive. Each summand grouped by a chosen component of \(\beta\) then has a factor \(\sum_{U'\subset V'}\mu(U')=0\) in a different nonzero component, where the sum includes \(V'\). Thus only primary matrices contribute.

Suppose the sole primary polynomial is \(f\), with root \(b\), degree \(m\mid n\), and let \(r=\dim_{\mathbb F_{q^m}}\ker f(g)\). Only submodules in \(\ker f(g)\) have nonzero Möbius number. Indeed, induct on the length of a primary module. For a module strictly larger than its socle the induction restricts the sum over proper submodules to the full nonzero socle, where the Möbius sum vanishes. On a semisimple module of dimension \(i\) over \(\mathbb F_t\), with \(t=q^m\), the Möbius number is \((-1)^it^{i(i-1)/2}\). This follows by induction from the finite Gaussian binomial identity
\[
   \sum_{i=0}^r \genfrac{[}{]}{0pt}{}{r}{i}_t(-1)^it^{i(i-1)/2}u^i
                    =\prod_{k=0}^{r-1}(1-t^ku).
\]
Here \(\genfrac{[}{]}{0pt}{}{r}{i}_t\) counts \(i\)-dimensional subspaces; the identity follows by separating subspaces contained in a fixed hyperplane from those meeting it in codimension one. On each such submodule \(\beta=i\sum_{k=0}^{m-1}\theta(b^{q^k})\). Differentiate at \(u=1\). The derivative of the factor \(1-u\) is the only surviving term, and the resulting value is
\[
      (-1)^{n-1}\left(\prod_{k=1}^{r-1}(1-q^{mk})\right)
                   \sum_{k=0}^{m-1}\theta(b^{q^k}).
\]
In particular its degree is \(\prod_{k=1}^{n-1}(q^k-1)>0\). This is Green's single-orbit character formula (a single part of size \(n\) in the basic character, with the partition \((1)\) at the regular orbit). His irreducibility theorem applies to this orbit \cite[\S\S4--7, Theorem~12]{theta-green}. The sign and degree above fix the convention even for \(n=1\). For a squarefree label the irreducible factors give distinct orbits and Green's product theorem gives their parabolic-induction product \cite[Theorems~13 and~14]{theta-green}. In \S\ref{sec:green-cuspidal} one marked factor has a single orbit; the ensuing zero average on radical translates uses that specific factor.\formalnote{8.11}{11}

\subsection{From prescribed signs to a rational isometry}
For a surviving polynomial, \(L=E[g]\) is a field. Its CM involution \(g\mapsto g^{-1}\) acts as ordinary conjugation under each complex embedding, since the eigenvalues have absolute value one. Each chosen sign on a numbered eigenline specifies a sign \(\eta_w\) at a real place \(w\) of \(L^+\), the fixed field. In the comparison their product is \(1\). Take the idele \(u\) of \(L^+\) with component \(\eta_w\) at \(w\mid\infty\) and \(1\) at every finite place. The local Artin maps for \(\mathbb C/\mathbb R\) send a negative number to the nonidentity and a positive number to the identity. The finite symbols are trivial. Thus the global symbol for the quadratic extension \(L/L^+\) is \(1\). Global reciprocity gives \(u=aN_{L/L^+}(h)\) for some \(a\in(L^+)^\times\), \(h\in\mathbb A_L^\times\). In particular \(a\) has the prescribed signs and \(a=N(h_f^{-1})\) at the finite places. Conversely, principal ideles and norms have trivial global symbol, so these conditions force \(\prod_w\eta_w=1\). This proves the exact sign obstruction \cite[\href{https://www.jmilne.org/math/CourseNotes/CFTc.pdf}{V, Theorem 5.3, and VII, §8}]{arith-cft}.\formalnote{8.12}{12}

The regular \(L\)-module \(E^n\) identifies \(h_f^{-1}\) with an invertible matrix \(Z\) commuting with \(g\). The adjoint on that algebra for the given form \(J\) is the CM involution, by its value on \(g\). Hence \(Z^*JZ=J a(g)\). This supplies local isometries between the new and the old form at every finite place. At infinity the new form has the requested signs. Compare the new form with the fixed target form of this signature. The finite frames give local isometries there and diagonalization gives local isometries at every infinite place. Apply Landherr to the matrix algebra over \(E/F\) with its involution of the second kind. Both forms are nonsingular of the same dimension. For their Gram matrices \(J_1,J_2\), put \(I(X)=J_1^{-1}X^*J_1\) and \(Y=J_1^{-1}J_2\). Then \(I(Y)=Y\) and a local isometry gives \(Y=I(a_v)a_v\) at each place. Taking \(x_v=I(a_v)\) puts this in the form \(Y=x_vI(x_v)\) of Landherr's theorem \cite[\S5.5, Landherr's theorem, second formulation]{legacy-kneser}. It gives a rational solution and hence a rational isometry. More explicitly, if \(a_0^*J_1a_0=J_2\) and \(x_i\) are the two frames, \(g_0=x_1^{-1}a_0x_2\) belongs to the common adelic unitary group. For rational \(\gamma\in U(J_2)\),
\[
       x_1^{-1}(a_0\gamma a_0^{-1})x_1
                     =g_0(x_2^{-1}\gamma x_2)g_0^{-1}.
\]
This equality transports the rational subgroup together with the point. The separate bijection matching the finite centralizers is the one given by \(Z\) in \S\ref{sec:signature-trace}.
\phantomsection\label{route:14}\leanroute{Landherr's theorem gives the rational isometry from the local isometries just exhibited. Lean proves the case of adelically framed spaces over a CM field containing a primitive \(5040\)-th root of unity. It uses the Hasse norm theorem and strong approximation to represent one, and then constructs the rational frame one column at a time. The displayed conjugation is needed with either proof.}\formalnote{8.13}{13}

\subsection{The use of Landherr in cyclic descent}
\noindent\textit{This subsection may be deferred when reading the proof of Theorem B. The descent argument in Chapter~\ref{ch:galois}, \S\ref{sec:weak-transfer}, uses the rational coboundary constructed here.}

Let \(E'/E\) be the cyclic CM extension of prime degree with generator \(\sigma\). Suppose \(s\in U(I_n)=\{s\in {\rm GL}_n(E'):s^*s=I_n\}\) has twisted norm one and there is a finite adelic unitary \(u\) with \(s=u\sigma(u)^{-1}\). Matrix Hilbert 90 gives \(b\in{\rm GL}_n(E')\) with \(s=b\sigma(b)^{-1}\). Equivalently \(v\mapsto s\sigma(v)\) is a semilinear cyclic action and a basis of its invariant space gives the columns of \(b\).
Since \(s^*s=1\) and \(\sigma\) commutes with CM conjugation, \(J=b^*b\) is fixed by \(\sigma\) and descends to \(E\). The matrix \(z=u^{-1}b\) is fixed at the finite places and likewise descends to the finite adeles of \(E\). There \(z^*z=J\). At each infinite place choose an embedding extending it to \(E'\); the matrix \(b\) there exhibits \(J\) as positive definite. Landherr gives \(l\in{\rm GL}_n(E)\) with \(l^*Jl=1\). Now \(g=bl\) is unitary and \(\sigma(l)=l\) gives \(s=g\sigma(g)^{-1}\), exactly the rational coboundary used to descend the pair. The level-preserving finite solution \(u\) is furnished separately by the division and norm calculations of \S\S\ref{sec:tame-root} and~\ref{sec:global-root}.\formalnote{8.14}{14}


\section{Vanishing on a proper unipotent radical}
\label{sec:green-cuspidal}

Let \(k=\mathbb F_q\). At the place used for cuspidality the label \(P\) is irreducible of degree \(n\) with \(P(0)\ne0\). Its root in \(k_n^\times=\mathbb F_{q^n}^\times\) specifies a regular character \(\theta\) of that cyclic group. The identification with its character group uses a fixed faithful character and carries the \(q\)-power operation to the \(q\)-power operation. Hence its orbit has exactly \(n\) elements. The signed finite class function used in the trace is Green's character of that orbit. The comparison of its values, including the degree and sign, is given in \S\ref{sec:classical-matching}.

Cuspidality can be read directly from Green's parametrization. In his character ring, multiplication is parabolic induction. For a fixed Frobenius orbit, the characters indexed by partitions form the images, up to their prescribed signs, of Schur functions under a ring homomorphism. Products with the same orbit label therefore expand using only that label. Products from distinct labels give Green's factorization of irreducibles \cite[\S7, Definition~7.3 and Theorem~12, and \S8, Theorem~13]{theta-green}. A representation of a proper block Levi has factors on groups of ranks strictly less than \(n\), so all their orbit labels have degree strictly less than \(n\). Its induced character consequently has zero inner product with the character labeled by the single degree-\(n\) orbit. Frobenius reciprocity gives cuspidality. In the elliptic-torus description the corresponding criterion is that \(T^F=k_n^\times\) preserves no nonzero proper \(k\)-subspace and \(\theta\) has trivial rational Weyl stabilizer. The signed torus character is then irreducible and cuspidal \cite[Theorem~6.8, Proposition~7.4 and Theorem~8.3]{theta-deligne-lusztig}. If \(U\) is the radical of a proper \(k\)-parabolic, the corresponding representation \(\pi\) has \(\pi^U=0\). The uniform mean of \(\pi(u)\) is the projection onto this invariant space. For arbitrary \(y,z\in\mathrm{GL}_n(k)\) it follows that
\[
 \sum_{u\in U}\chi_P(yuz)
  =\operatorname{tr}\!\left(\pi(y)\sum_{u\in U}\pi(u)\,\pi(z)\right)=0 .
\]
This proves the translated cancellation needed for the central idempotent, whose coefficient at \(g\) is \(\chi_P(1)\overline{\chi_P(g)}/|\mathrm{GL}_n(k)|\). It also covers parabolics of longer flags.

\phantomsection\label{route:15}\leanroute{The finite development proves the block mean by counting stable subspaces in the Hall algebra, and constructs the irreducible character using Brauer lifting and its squared norm. Green's regular-orbit theorem supplies the representation and its cuspidality at this one marked place. Taking the trace of its zero averaging operator gives the stronger, two-sided translate identity needed by the adelic calculation.}\formalnote{8.15}{15}

Lift the product of the residue idempotents at the marked places by a section of reduction on the adelic maximal compact \(K_0\), and convolve with the indicator of the kernel of reduction. The resulting function on \(K_0\) is its coefficient at the reduced element. For a basis point \(s\) of the residue group the indicator is one exactly when \(g\) and the lift of \(s\) are congruent. Linearity gives the claim. Its support is contained in \(K_0\). In the sum over the finitely many lifts of a radical translate at the chosen place all other residue components are fixed. For \(g\in K_0\) the sum is the finite sum just calculated, times the other components. For \(g\notin K_0\) multiplication by these integral lifts stays outside \(K_0\), so each value is zero. Arbitrary finite left convolution is a linear combination of translates and preserves the same cancellation. This proves the marked cuspidality used in \S\S\ref{sec:signature-trace} and~\ref{sec:cusp-trace}.\formalnote{8.16}{16}


\section{The elliptic orbital character in the signature comparison}\label{sec:elliptic-character}
Fix one archimedean place in \S\ref{sec:signature-trace}. Write \(J=\operatorname{diag}(\epsilon_i)\) with \(p\) plus and \(q\) minus signs, \(n=p+q\ge1\), \(G=U(J)\), \(K=U(p)\times U(q)\), and \(W=S_p\times S_q\), acting on the given sign blocks. The integral is over \(G\) itself. Its Haar measure is the compact-dual measure of \S\ref{sec:flag-measure}, which gives \(U(n)\) mass one. Integrals on \(K\) use Haar probability. The parameter in the source is \(l=2\lambda\). The \(l_i\) are distinct integers of parity \(n-1\), and
\[
   |\lambda_i-\lambda_j|\ \ge h,\qquad
   h=500n^3+10n^5,\qquad \epsilon_i\ne\epsilon_j .
\]
Orient a root \(e_i-e_j\) positively if \(l_i>l_j\). Let \(\rho\) be the half sum of positive roots, and \(\rho_c\) the half sum over pairs of equal sign. The lowest \(K\)-type has highest weight \(\Lambda=\lambda+\rho-2\rho_c\). The compact representation that appears after transfer has highest weight \(\nu_0=\lambda-\rho\). These weights coincide in the definite case. For a diagonal matrix \(t\) we also use \(t\) for the vector of its entries. Write \(t^\alpha=\prod_i t_i^{\alpha_i}\), and put
\[
\begin{aligned}
  V(t)&=\prod_{l_i>l_j}(t_i-t_j),&
  V_c(t)&=\prod_{\substack{l_i>l_j\\ \epsilon_i=\epsilon_j}}(t_i-t_j),\\
  D(t)&=\prod_{\substack{l_i<l_j\\ \epsilon_i\ne\epsilon_j}}(t_i-t_j),&
  A_H(\alpha)(t)&=\sum_{w\in H}\operatorname{sgn}(w)t^{\alpha\circ w}.
\end{aligned}
\]
There are \(pq\) reversed factors in \(D\), so \(V_cD=sV\), where \(s=(-1)^{pq}\). The integer numerator exponent is \(\kappa_i=(l_i+n-1)/2\).

Let \(c_l(1)=1\) be the normalized lowest-type coefficient, let \(d^{-1}=\int_G|c_l(g)|^2\,dg\), and put \(e_l=d\,\overline{c_l}\). With \(l=2\lambda\) we have \(c_l=c_\lambda\), so \(e_l(g)=k_\lambda(g^{-1})\) in the convolution formula of \S\ref{sec:local-orbital}. We shall prove that for \(|t_i|=1\) with the \(t_i\) pairwise distinct,
\[
 \Theta(t):=d\int_G c_l(g^{-1}tg)\,dg
       =\frac{A_W(\kappa)(t)}{V_c(t)D(t)} .                  \tag{EC1}
\]
The normalized coefficient in \S\ref{sec:real-decay} belongs to the discrete series of parameter \(\lambda\) and lowest compact weight \(\lambda+\rho-2\rho_c\). These data specify the member of the packet, including its central character. To apply the semisimple character theorems, use the finite central covering \(U(1)\times SU(p,q)\to U(p,q)\), \((z,g)\mapsto zg\). The central circle acts by a character \(\chi\), so the pullback is \(\chi\otimes\pi_0\), with \(\pi_0\) irreducible on \(SU(p,q)\). Give the circle mass one and choose Haar measure on \(SU(p,q)\) so that product measure pushes forward to the specified measure on \(G\). The squared norm of a coefficient is then unchanged on restriction to \(SU(p,q)\), since its absolute value is constant along the circle. In particular \(\pi_0\) is square integrable and has the same formal degree. For \(t=zt_0\), both the coefficient orbital integral and the character acquire the factor \(\chi(z)\). The root system ignores the scalar coordinate, which is restored by this factor. For \(n=1\) the representation is a character of the compact circle and (EC1) follows by direct integration.

On the special unitary group the lowest-type vector is \(K\)-finite. Harish-Chandra's orbital-integral formula for its normalized matrix coefficient gives \(d_\lambda\int c_\lambda(g^{-1}tg)\,dg=\Theta_\pi(t)\) on the regular compact Cartan \cite[\S40, Lemma~80 and Theorem~14]{theta-hc-orbital}. In Arthur's formulation the integral is over \(T_{\mathbb R}\backslash G\), where \(T_{\mathbb R}\) is the vector part of the Cartan centralizing \(t\). That part is trivial here, and the integral converges absolutely \cite[p.~205, eq.~(1)]{theta-arthur}. Equivalently one can integrate over \(T\backslash G\) with probability Haar measure on the diagonal torus \(T\). The sign of the coefficient convention agrees on the diagonal torus, since \(c_l(t)\) has the weight of the lowest-type vector. The conjugate is introduced only when passing to \(e_l\).

Harish-Chandra's compact-Cartan character formula is the compact-Weyl alternant of \(\lambda\), divided by the Weyl denominator for all positive roots, with sign \((-1)^{\dim(G/K)/2}\) \cite[Theorem~8.1]{theta-atiyah-schmid}. The special unitary group used here is linear in its defining representation, so the original index argument applies, as explained in the erratum. Here \(\dim_{\mathbb R}G/K=2pq\). To convert this formula into integer powers of \(t_i\), factor each root term as
\[
   (t_i/t_j)^{1/2}-(t_j/t_i)^{1/2}
          =\frac{t_i-t_j}{(t_it_j)^{1/2}} .
\]
Every \(t_i\) occurs \(n-1\) times under the square roots. Multiplication by the common monomial thus changes \(\lambda_i\) in the numerator to \(\kappa_i=\lambda_i+(n-1)/2\) and leaves \(V(t)\) in the denominator. The result is \(sA_W(\kappa)/V=A_W(\kappa)/(V_cD)\), proving (EC1). The semisimple formulas apply to \(U(p,q)\) by the finite central covering \(SU(p,q)\times U(1)\to U(p,q)\). On the covering group the circle contributes the single central character and has mass one. Pushing the product Haar measure to the quotient on both real forms preserves that normalization. For \(p q=0\) this reduces to compact Schur orthogonality and the Weyl character formula, and also covers \(n=1\).

In the compact-dual Haar measure of \S\ref{sec:real-orbit-density}, the formal-degree formula of the same theorem gives
\[
   d=\prod_{\alpha>0}\frac{(\lambda,\alpha)}{(\rho,\alpha)}
    =\frac{\prod_{l_i>l_j}(\lambda_i-\lambda_j)}
            {\prod_{j=0}^{n-1}j!}
    =\dim \nu_0 .
\]
The denominator is the product of all positive differences of the integer positions in the chosen order. The last equality is the Weyl dimension formula of \(U(n)\). The product formula for the formal degree uses the normalization of the individual measure. The identity (EC1) itself is unchanged upon rescaling \(dg\), because \(d\) scales inversely.

\phantomsection\label{route:16}\leanroute{The character formula and the formal degree in the formal development are deduced from compact multiplicities, a polynomial on the torus and an integral on an open flag orbit. Here the orbital-integral theorem and the compact-Cartan formula supply the individual character directly. The Haar comparison in \S\ref{sec:real-orbit-density} identifies the scalar needed for the formal degree. The permutation sum below still has to be done for the signature comparison.}\formalnote{8.17}{17}

\subsection{The permutation sum in the trace}
Put \(I(\sigma)=\int_G e_l(g^{-1}\operatorname{diag}(t\circ\sigma)g)\,dg\). From (EC1) and \(V_cD=sV\), the unconjugated expression \(s\sum_{\sigma}\Theta(t\circ\sigma)\) equals
\[
        \frac{1}{V(t)}\sum_{\sigma\in S_n}
             \operatorname{sgn}(\sigma)A_W(\kappa)(t\circ\sigma).
\]
For each fixed \(w\in W\) the inner sum over \(\sigma\), including the sign of \(w\), is exactly \(A_{S_n}(\kappa)(t)\) by reindexing the permutations. Weyl's compact formula and the conjugation in \(e_l\) now give
\[
          (-1)^{pq}\sum_{\sigma\in S_n} I(\sigma)
                 =p!\,q!\,\overline{\chi_{\nu_0}(t)}.           \tag{EC6}
\]
\begin{example}[The sign of the character sum for \(U(1,1)\)]\label{ex:elliptic-character-sign}
For \(U(1,1)\), take doubled parameters \(l_1>l_2\) in the range above and put \(\kappa_i=(l_i+1)/2\). The two permutations assign the two unit-circle eigenvalues to the positive and negative directions in opposite orders. The compact Weyl group
\(S_1\times S_1\) is trivial. For distinct \(t_1,t_2\) on the unit
circle, put \(A=t_1^{\kappa_1}t_2^{\kappa_2}\) and
\(B=t_2^{\kappa_1}t_1^{\kappa_2}\). Formula (EC1) gives
\[
 -\bigl(\Theta(t_1,t_2)+\Theta(t_2,t_1)\bigr)
   =-\left(\frac{A}{t_2-t_1}+\frac{B}{t_1-t_2}\right)
   =\frac{A-B}{t_1-t_2}.
\]
This is the compact \(U(2)\) character at \((t_1,t_2)\). The minus
sign is \((-1)^{pq}\) in (EC6), and conjugating the identity gives
the formula with the kernel \(e_l\).
\end{example}

For the definite group this just counts the \(n!\) orderings of a single class. In the mixed signature it counts each sign assignment \(p!q!\) times. The global comparison uses the same signature at each of \(f\) archimedean places. After dividing by \((p!q!)^f\), the product of local signs is \(1\) because \(f\) is even. This proves the real identity used there for both signatures on its required regular domain.\formalnote{8.18}{18}


\section{Haar normalization on the two real forms}\label{sec:real-orbit-density}
The elliptic character in \S\ref{sec:elliptic-character} is independent of a rescaling of Haar measure. Its formal degree depends on that scale. We identify the measure in the unitary-group integrals with the one in the discrete-series formal-degree formula.

Let \(G=U(p,q)\), \(K=U(p)\times U(q)\) and \(G_c=U(n)\), where \(n=p+q\). Write \(J=\operatorname{diag}(1_p,-1_q)\). The compact-dual convention takes probability Haar measure on \(G_c\) and transports its invariant density at the identity to \(G\) using
\[
  \mathfrak g=\mathfrak k\oplus\mathfrak p,\qquad
  \mathfrak g_c=\mathfrak k\oplus i\mathfrak p .
\]
The compact factor has the same density on both sides and \(X\mapsto iX\) on the noncompact directions identifies the absolute Riemannian densities. This is the choice in the formal-degree theorem \cite[\S3, eqs.~(3.8)--(3.11), and Theorem~8.1]{theta-atiyah-schmid}. The central circle occurs on both sides. The covering by the circle times the special unitary group has the same finite kernel on both sides.

To compare with the code, put \(Y=JX\) for \(X\in\mathfrak g\). Then \(Y^*=-Y\). In left logarithmic coordinates the Haar density for any signature is \(C_n^{-1}dY\) at the identity, where
\[
 C_n=\prod_{j=1}^n\frac{2\pi^j}{(j-1)!}.
\]
Indeed the Cayley chart \(Y\mapsto(1+JY/2)(1-JY/2)^{-1}\) has differential \(Y\mapsto JY\) at zero, and the chart density there is \(C_n^{-1}\). In compact signature this is precisely the probability Haar measure by the compact unitary volume formula. On \(\mathfrak k\) the substitution \(X\mapsto JX\) changes only row signs. On \(\mathfrak p\) the independent entries form a complex rectangular block. Changing between \(JX\) and the compact-dual coordinates of \(iX\) only multiplies those entries by a unit of modulus one. Both changes have real absolute determinant one. The two invariant densities therefore agree at the identity, and hence at every point by translation. This proves the measure assertion used in (EC1) and in the separate formal-degree calculation.\formalnote{8.19}{19}

\phantomsection\label{route:17}\leanroute{For the direct computation of the character, Lean transports densities over an entire open flag orbit by signed Gram--Schmidt and computes the triangular Jacobian. The published character formula removes the need to integrate over that orbit. The remaining comparison is between two invariant Haar densities and can be made in one tangent space.}\formalnote{8.20}{20}


\section{Convergence in the cusp for the trace by characteristic polynomial}
\label{sec:cusp-trace}

Use the unitary group \(G\), rational subgroup \(\Gamma=U(V)(F)\),
real kernel \(e\), and finite test \(\phi\) of Section~\ref{sec:signature-trace}.
Write \(n=\dim_E V>0,\ f=[E:\mathbb Q]/2\), and
\(L(t)=\prod_{w\mid\infty}\|t_w\|_{\rm op}\), taking one complex embedding in
each conjugate pair. Each local norm is at least one and equals the norm
of the inverse, since \(t_w^{-1}=J_w t_w^*J_w\) for a diagonal sign matrix.
For a set \(\mathcal P\) of polynomials in \(E[X]\) put
\[
 K_{\mathcal P}(x)=
  \sum_{\substack{\gamma\in\Gamma\\\det(X-\gamma)\in\mathcal P}}
    e((x^{-1}\gamma x)_\infty)\phi((x^{-1}\gamma x)_f).
\]
The assertion needed for the diagonal trace is
\(\sup_{x,\mathcal P}|K_{\mathcal P}(x)|<\infty\).
The sum also converges absolutely at each fixed \(x\). The proof of the
uniform bound uses cancellation among the elements preserving an isotropic
subspace. For those moving it one can take absolute values term by term.\formalnote{8.21}{21}

\subsection{A bound for a sum with no cancellation}

Let \(\Lambda_x\) be the lattice in \(E^n\) defined by the finite
coordinates of the moved space, with lengths measured in its archimedean
coordinates. Put \(H(x)=1/\min(1,\lambda_1(\Lambda_x))\). Choose one
compact open subgroup \(J\) of the integral level so that \(\phi\) is
invariant on both sides. Set
\[
 A_0=8n^2+1,\qquad d=4fn^2,\qquad
 M_A(t)=|\phi(t_f)|L(t_\infty)^{-A}.
\]
We first prove a bound uniform even in a second point \(g\):
\[
             \sum_{\gamma\in\Gamma}M_{A_0}(g^{-1}\gamma x)
                          \le C H(x)^d .                         \tag{CT1}
\]
The finite factor has compact support. At each real factor the Haar
volume bound makes \(\|t_w\|^{-A_0}\) integrable for \(A_0>8n^2\);
the integral of the product is therefore finite. Around \(1\) take a
box with finite part \(J\) and archimedean radii
\(\varepsilon=(cH(x)^2)^{-1}\). The translates \(g^{-1}\gamma x\)
of this box are pairwise disjoint for one fixed large enough \(c\).
In fact \(\Lambda_x\) contains an \(E\)-basis of vectors of length at
most \(bH(x)\), with \(b\) independent of \(x\). If two boxes meet,
\(x^{-1}\gamma'^{-1}\gamma x\) is integral at the finite places and
within \(4\varepsilon\) of \(1\) at infinity. Subtracting a basis vector
from its image gives a lattice vector of length at most
\(4b\varepsilon H(x)<H(x)^{-1}\). It must be zero, so
\(\gamma=\gamma'\).

The real norm-two ball can be covered by at most
\((12/\varepsilon)^{2n^2}\) group translates of the single-place
\(\varepsilon\)-ball, by packing in the \(2n^2\) real matrix coordinates.
The identity box in the product therefore has volume at least
\(c_J\varepsilon^{2fn^2}\). On multiplying on the right by a point
of the box, \(M_{A_0}\) decreases by at most the factor \(2^{fA_0}\),
by level invariance and submultiplicativity. Integrating this inequality
over the disjoint boxes and dividing by the box volume proves (CT1).
Its height exponent is \(2\cdot 2fn^2\).\formalnote{8.22}{22}

The cusp kernel has \(|e(t)|L(t)^B\le C_e\), with \(B\ge13n^3\).
For any subset of \(\Gamma\), (CT1) now implies absolute convergence.
If on that subset \(1\le hL((g^{-1}\gamma x)_\infty)\) whenever the
finite factor is nonzero, and \(A_0+j\le B\), the same argument bounds
its sum of moduli by \(C h^j H(x)^d\). Indeed \(L^{-B}\le h^jL^{-A_0}\).
We use precisely \(j=4n^3\) below.\formalnote{8.23}{23}

\subsection{The isotropic flag and the elements that move it}

Fix an integer \(M>0\) with principal level \(J(M)\subseteq J\).
Let \(D\) be the absolute norm of the different of \(E\), and
\[
       c_0=8fn\,{\rm rd}(E)+48 n^6 f^2 D+1,\qquad
       \kappa=2(Mc_0^2)^n,\qquad
       \delta^n H(x)=\kappa .
\]
When \(\delta\le1\), the lattice supplies a nonzero isotropic subspace
\(A\), and columns of integral matrices in the moving lattice which
span \(A,A^\perp\) with length bounds \(l_X\le l_Y\), respectively.
In particular \(l_Xl_Y\le\delta\). Here is the argument for these
columns. If \(\lambda_1^{\rm cap}<\theta^n\) and \(\theta c_0\le1\),
Minkowski reduction gives an ordered rational basis with lengths
\(\alpha_1\le\cdots\le\alpha_n\), product at most \(c_0^n\).
Take a nonzero lattice vector of length less than \(\theta^n\).
Some member of this basis pairs nontrivially with it. Their hermitian
pairing is an algebraic integer, so the product formula yields
\(1\le\|\xi\|\alpha_n\). Thus \(\prod_{i<n}\alpha_i<
(\theta c_0)^n\). Pair these \(n-1\) lengths in reverse order.
Some \(1\le i\le n-i\) satisfies \(\alpha_i\alpha_{n-i}<\theta c_0\);
otherwise squaring the product would give at least
\((\theta c_0)^{n-1}\). The first \(i\) vectors pair to zero with the
first \(n-i\), by the same product formula and the strict bound \(<1\).
Their spans are therefore \(A,A^\perp\) by the dimension count. In
rank one the short-vector premise itself is impossible.\formalnote{8.24}{24}

The radical \(N_A\) must also have sufficiently small representatives.
Apply the preceding construction with \(\theta=\delta/(Mc_0^2)\).
The required strict inequality holds because
\(\theta^n=2/H(x)\). The stronger bound \(l_Xl_Y\le\delta/(Mc_0)\)
lets one cover
\[
           N_A(\mathbb A_F)x
             \subseteq N_A(F)x\,Z_\delta J(M),                 \tag{CT2}
\]
where \(Z_\delta\) is the real ball of radius \(\delta\) at \(1\).
In the moved frame \(X,Y\) denote the matrices of the two sets of
columns. In the source's radical coordinates, a matrix \(T\) gives
\(1+2(T-T^*-T^*J_0 T)J_0\), with \(J_0\) denoting the standard form at
each place. A rational radical translation brings the parameter of
each point to \(MT\), with
\[
          T=\bigl(Y(C-C_0)-X(D'-D'_0)\bigr)X^* ,
\]
where the subscript \(0\) denotes entrywise adelic rounding to \(E\).
This parametrization uses that \({\rm im}\,Y=\ker(X^*J)\).
Resolving the identity as \(Yh+2MbX^*J=1\) by a rational linear
splitting permits both rounding corrections within the radical. The
rounding errors have integral finite components and real norm at most
\(n^2(2f)^2D\). Bounding \(X,Y\) by \(n^2l_X,n^2l_Y\) gives
\[
       \|T_w\|\le 4n^6 f^2D(l_Y+l_X)l_X
                    \le \frac{c_0 l_X l_Y}{6}
                    \le \frac{\delta}{6M}.
\]
Consequently the group output is congruent to the identity modulo
\(M\) at the finite places and within \(\delta\) at the real ones
by its quadratic formula. This proves (CT2). For \(\delta\ge1\)
we already have \(H(x)\le\kappa\).\formalnote{8.25}{25}

If \(\gamma A\ne A\), some columns \(y\in A,\ y'\in A^\perp\)
have \(u=\langle y',\gamma y\rangle\ne0\). On the finite support
of \(\phi((x^{-1}\gamma x)_f)\) a single integer \(c>0\) clears all
matrix denominators, so \(cu\) is an algebraic integer. At each
archimedean place Cauchy--Schwarz in the moved coordinates bounds
\(|(cu)_w|\le c\delta\|(x^{-1}\gamma x)_w\|\). The norm of the
nonzero algebraic integer has modulus at least one. Taking its square
root, since \(E\) is CM, gives
\[
                  1\le(c\delta)^fL((x^{-1}\gamma x)_\infty).
\]
Use \(h=(c\delta)^f,\ j=4n^3\) after (CT1). Since \(fj=nd\), the sum
of absolute values over \emph{any} selection of movers is at most
\[
          C(c\delta)^{4fn^3} H(x)^d
                       = C c^{nd}(\delta^n H(x))^d .
\]
If the characteristic polynomial is irreducible, every element moves
\(A\), since a matrix preserving the proper nonzero \(E\)-space \(A\)
has reducible characteristic polynomial. Thus this absolute uniform
bound alone applies to the irreducible orbital terms.\formalnote{8.26}{26}

\subsection{The sum over the stabilizer}

Fix \(x,A,\mathcal P\). Let \(Z\) consist of the rational elements
preserving \(A\) with characteristic polynomial in \(\mathcal P\).
An element \(\nu\in N_A(F)\) acts as the identity on the graded pieces
of \(0\subset A\subset A^\perp\subset E^n\). The last quotient is
dual to \(A\) through the form. Since \(\gamma\) preserves this flag
whenever \(\gamma A=A\), its matrix and that of \(\nu\gamma\) have
the same three diagonal blocks. Thus \(\nu Z=Z\) for \emph{every}
\(\mathcal P\). Form the function of a new variable
\[
       F_x(g)=\sum_{\gamma\in Z}
                 e((g^{-1}\gamma x)_\infty)\phi((g^{-1}\gamma x)_f).
\]
It is left \(N_A(F)\)-invariant by reindexing, right \(J\)-invariant
by left invariance of \(\phi\), and bounded uniformly in \(g\) by
\(T=C H(x)^d\) using (CT1). Convolution on the right by \(e\) fixes
it. One can integrate term by term since the sum of moduli inside
the convolution is dominated by \(T|e|\), and then use \(e*e=e\).
The convolution is continuous by continuity of translation in \(L^1\)
and level invariance.\formalnote{8.27}{27}

The mean of \(F_x\) along \(N_A(F)\backslash N_A(\mathbb A_F)\) is
zero, at every real right translate. Indeed (CT2) gives finite
covolume. Averaging \(F_x\) over that quotient produces a bounded
continuous function \(\Psi\) on the real group, left invariant by the
real radical. Fubini and reproduction of \(F_x\) show \(T_e\Psi=\Psi\),
where \(T_e b(y)=\int e(t)b(yt)\,dt\). We prove \(T_e\Psi=0\).

For a real isotropic column matrix \(X\ne0\) put
\(h(y)=\log\|y_w^{-1} X_w\|\) at one real group factor. Its right
increment has modulus at most \(L(t)^{B+1}\). A real dilation \(a\)
on the flag with scalar \(s>1\) at every place and acting trivially
on \(A^\perp/A\) changes \(h\) by \(-\log s\). We can choose \(s\)
close enough to one that for every \(y\) such a dilation is available
in the form \(y b_y y^{-1}\), with
\(\|e(b_y^{-1}\cdot)-e\|_1<1/2\). In fact project orthogonally onto
\(y_w^{-1}A_w\) in the positive definite Euclidean metric, with
projector \(P\), and take \(Q=J_wPJ_w\). The two ranges are orthogonal
by isotropy. The matrix \(1+(s-1)P+(s^{-1}-1)Q\) is unitary for
\(J_w\), acts on those ranges by \(s,s^{-1}\), and is within
\(s-1\) of \(1\), uniformly in \(y\).

Let \(\mathcal V\) be the bounded continuous real-group functions both
fixed by \(T_e\) and left invariant under the real radical. A fixed
choice \(a\) of the dilation acts on \(\mathcal V\) by
\(U b(y)=b(ay)\). The two dilations \(a,yb_y y^{-1}\) differ by a
radical element, so \(b(ay)=b(yb_y)\). Reproduction and the \(L^1\)
bound show \(\|1-U\|\le1/2\). For an integer \(m\) put
\(\theta=-\pi/(m\log s)\) and \(S b=T_e(e^{i\theta h}b)\). This maps
\(\mathcal V\) into itself by idempotence and radical invariance.
The shift of \(h\) gives \(US=e^{i\pi/m}SU\), so \(US^m=-S^mU\).
Writing \(2S^m=(1-U)S^m+S^m(1-U)\) forces \(S^m=0\).
On the other hand if \(Sb=0\), subtract
\(e^{i\theta h(y)}T_e b(y)\) inside the integral and use the bound on
the right increment of \(h\). Then
\[
          \|b\|_\infty\le |\theta|\!\int |e(t)|L(t)^{B+1}\,dt
                               \ \|b\|_\infty .
\]
For large \(m\) the coefficient is less than one, making \(S\) injective.
Thus \(\mathcal V=0\), proving the asserted vanishing of the mean.\formalnote{8.28}{28}

By (CT2) and its zero mean, \(|F_x(xz)|\) is bounded by the maximum
of \(|F_x(xvz)-F_x(xz)|\) over \(v\in Z_\delta\), where we omit the
inclusion of real factors into the adeles. Indeed average the
difference from \(F_x(xz)\) using quotient probability. A compactly
supported continuous \(\alpha\) reproduces \(e\), and therefore also
\(F_x\). For each fixed integer \(k\ge0\) there is a \(\delta_0>0\)
such that, whenever \(|F_x(gw)|\le C' L(w)^{2k}\),
\[
         |F_x(gu)-F_x(g)|\le C'/2
                    \quad\text{if }d(u,1)\le\delta_0 .
\]
To see this, on all supports of the translates of \(\alpha\) by the
unit ball the factor \(L(w)^{2k}\) has a fixed maximum. Choose
\(\delta_0\) so the \(L^1\)-distance between \(\alpha\) and its
translate is less than the reciprocal of twice that maximum.

Set \(M_*=\sup_z |F_x(xz)|/L(z)^{2k}\le T\).
If \(\delta L(z)^2\le\delta_0\), then \(z^{-1}v z\) is within
\(\delta_0\) of \(1\) for \(v\in Z_\delta\). Apply the last estimate
at \(g=xz\), using \(C'=M_* L(z)^{2k}\), and use the zero-mean
bound by oscillation. The quotient at this \(z\) is at most
\(M_*/2\). For the other \(z\), it is at most \(T(\delta/\delta_0)^k\).
Taking the supremum and then \(z=1\) yields
\[
         |F_x(x)|\le T(\delta/\delta_0)^k .
\]
Take \(k=nd\). This pays for the factor \(H(x)^d\) in \(T\), giving
\(|F_x(x)|\le C\delta_0^{-nd}\kappa^d\). The same \(\delta_0\)
works for all the polynomial sets because only \(e\) and \(k\)
entered the translation estimate.\formalnote{8.29}{29}

\subsection{Grouping the trace and the triple integral}

The bounds in the two height ranges now prove the stated supremum.
At fixed \(x\), absolute convergence gives the pointwise limit as finite
sets of polynomials exhaust \(E[X]\). On the finite-volume quotient the
constant bound gives dominated convergence for the \emph{polynomial}
partial sums. Finite additivity and passage to that limit identify
the integral of the full kernel with the sum of the integrated terms
indexed by polynomials. For irreducible polynomials the absolute
estimate over all movers also permits unfolding each conjugacy class
to the quotient by its centralizer.\formalnote{8.30}{30}

For such a \(\gamma\) put \(t=x^{-1}\gamma x\), let \(y,u\) be the
two real-group increments, and put \(z=u^{-1}y^{-1}t_\infty\).
For a measurable weight \(a\) with \(|a|\le L(y)L(u)\), the absolute
value of the weighted triple kernel is at most
\[
   C\,|\phi(t_f)|L(t_\infty)^{-B}\,
               \bigl(|e(y)|L(y)^{B+1}\bigr)
               \bigl(|e(u)|L(u)^{B+1}\bigr).
\]
Indeed \(t_\infty=yuz\) and \(|e(z)|L(z)^B\le C\).
The first factor is integrable in \(x\) modulo the centralizer by the
irreducible absolute bound applied to the real kernel \(L^{-B}\).
The other two integrals are the moments in the cusp-kernel hypothesis.
This accounts for the extra \(1\) in \(B+1\), even for a logarithmic
increment. It justifies Fubini for the first two increments and
their sum, with an inessential constant if necessary. The third
increment is the total constant \(\log |r|\) minus the first two,
where \(r\) is the eigenvalue used in Section~\ref{sec:signature-trace}.
It is consequently integrable against the triple kernel as well.
Cyclic rotation now compares the three integrals.\formalnote{8.31}{31}

\subsection{The exact cancellation at the marked finite place}

Uniform convergence above holds for every locally constant compactly
supported \(\phi\). To eliminate polynomials with coprime factors,
Section~\ref{sec:signature-trace} uses an additional identity for the
specific test at its marked split place. In \(G_0={\rm GL}_n(k)\) fix
a nonempty proper coordinate set \(T\) and its projection \(P\).
The upper radical \(U\) consists of \(1+PZ(1-P)\) as \(Z\) runs over all
matrices. Each member is repeated the same number of times. For the
irreducible polynomial label \(c\) of degree \(n\), let \(\rho\) be the
representation with Green character \(\chi_c\). One needs the finite sum
\[
   \sum_Z [e_c]\bigl(y(1+PZ(1-P))z\bigr)=0
                        \qquad(y,z\in G_0),
\]
where \([e_c](g)=\chi_c(1)\overline{\chi_c(g)}/|G_0|\).
The average \(P_U=|U|^{-1}\sum_{u\in U}\rho(u)\) is a projection.
Its trace is the mean of \(\chi_c\) on \(U\). That trace is zero, as
follows from the block-average formula in Section~\ref{sec:green-cuspidal}
with both diagonal blocks equal to the identity. Thus
\(P_U=0\), and the trace of \(\rho(y)P_U\rho(z)\) vanishes for every
pair \(y,z\), proving the displayed assertion including all translates.\formalnote{8.32}{32}

For several marked places the other tensor factors are constant in
the chosen radical variable. Lift the tuple idempotent to the
finite adelic group and convolve with the indicator of the reduction
kernel. On the maximal compact this gives precisely the coefficient
function of the reduction. Outside the maximal compact all the terms
remain outside after right multiplication by any integral lift, so are
zero. Left convolution by any finite group-algebra element preserves
this identity by linearity and translation. This is the finite-place
cuspidality used in the local projector-slice descent. The separate
simple-root condition at the second place removes proper powers.\formalnote{8.33}{33}

The quantitative inputs to these estimates are the short-basis and
Minkowski product bounds, adelic rounding by the different, the
archimedean volume-growth bound, and the decay, moment and compact
reproduction of the chosen real kernel. The marked test also uses
Green's character realization and the finite block-average identity of
Section~\ref{sec:green-cuspidal}. After grouping the kernel, apply the real character calculation in \S\ref{sec:elliptic-character} and the projector descent in \S\ref{sec:projector-descent}.


\section{Descent to the fibre of a polynomial projector}
\label{sec:projector-descent}
Let \(D\) be the finite adelic group of \S\ref{sec:signature-trace}, \(w\) its split cuspidal place, \(F=E_w\), \(G=\mathrm{GL}_n(F)\), and \(K=\mathrm{GL}_n(O_F)\). Embed \(G\) at the split place and write \(g\cdot y=j(g)y j(g)^{-1}\) for its action on \(D\). Write \(\mathcal S(Y)\) for the locally constant, compactly supported complex functions on \(Y\). Fix coprime monic polynomials \(A,B\in F[t]\) of positive degree. Suppose \(I\) is a linear form on \(\mathcal S(D)\), invariant under \(G\)-conjugation, that kills every function vanishing on
\[
       C=\{x\in D:\det(t-x_w)=A(t)B(t)\}.
\]
We prove \(I(\phi)=0\) for a test cuspidal at \(w\) in the sense used in \S\ref{sec:green-cuspidal}. The assertion is immediate when \(C\) is empty. The proof for nonempty \(C\) also applies to nonsemisimple \(x_w\).\formalnote{8.34}{34}

\subsection{Choosing coordinates for the primary summands}
Choose \(uA+vB=1\) and set \(p=vB\). The polynomial \(p^2-p=-uvAB\) gives \(e_x=p(x_w)\) with \(e_x^2=e_x\), by Cayley--Hamilton. Its image is \(\ker A(x_w)\), since \(A(x_w)e_x=0\) and \(e_x\) acts as the identity on that kernel. The complementary statement for \(1-e_x\) gives \(\ker B(x_w)\). Over an algebraic closure triangularize \(x_w\). The diagonal of \(p(x_w)\) consists of one \(1\) for every root of \(A\) and one \(0\) for every root of \(B\), with multiplicities. An idempotent has trace equal to its rank in the characteristic-zero field \(F\), so \(d=\deg A\) is the common rank. Fix the coordinate projector \(\epsilon\) of rank \(d\) and put
\[
 X=\{y\in C:p(y_w)=\epsilon\},\qquad
 H=Z_G(\epsilon),\qquad
 N=\{1+Z:Z=\epsilon Z(1-\epsilon)\}.
\]
Here \(H=\mathrm{GL}_d(F)\times\mathrm{GL}_{n-d}(F)\) in the chosen coordinates. Its elements preserve both coordinate summands and may change the matrix \(y_w\) inside them. In particular the Haar measure used below belongs to the centralizer of the \emph{projector}.\formalnote{8.35}{35}

Every rank-\(d\) projector \(e\) has the form \(g\epsilon g^{-1}\) with \(g\in KN\). Indeed integral row and column operations put \(e\) in diagonal form over \(F\). At each step take a nonzero entry of smallest valuation, put it at the pivot, and use the integral ratios of entries to the pivot to clear its row and column. Therefore some \(k\in K\) carries \(\operatorname{im}\epsilon\) onto \(\operatorname{im}e\). In these coordinates \(e'=k^{-1}ek=\left(\begin{smallmatrix}1&C_0\\0&0\end{smallmatrix}\right)\). The matrix \(r=1+\epsilon-e'\in N\) satisfies \(r\epsilon r^{-1}=e'\), so take \(g=kr\). Applying this to \(e=a\epsilon a^{-1}\) for each \(a\in G\) also gives \(G=KNH\).\formalnote{8.36}{36}

Conjugation \(\pi(g,y)=g\cdot y\) now maps \(G\times X\) onto \(C\). Its fibres are exactly the sets \(\{(gh,h^{-1}\cdot y):h\in H\}\), by equivariance of \(e_x\). Both \(C\) and \(X\) are closed by continuity of the matrix coefficients involved. There are local continuous choices of coordinates. For \(y\in X\) and \(x\in C\) close to \(y\), take \(s(x)=e_x\epsilon+(1-e_x)(1-\epsilon)\). This tends to \(1\), hence is invertible nearby, and \(s(x)\epsilon=e_xs(x)\). Thus \(s(x)^{-1}\cdot x\in X\). Conjugating this construction gives a local section near each point of \(C\).\formalnote{8.37}{37}

\subsection{Descent to the stabilizer}
The equations above present \(C\) as the associated space \(G\times_H X\). All spaces are locally compact and totally disconnected. Both \(G=\mathrm{GL}_n(F)\) and its block-diagonal subgroup \(H\) are unimodular. Give \(G/H\) the invariant measure determined by Haar measures \(\mu\) and \(\nu\) on \(G\) and \(H\), respectively. A compactly supported locally constant function on the closed subspace \(C\) extends to \(D\) by lifting its finitely many nonzero level sets to compact open sets. The support assumption on \(I\) thus makes it an invariant distribution on \(C\).

Frobenius descent for invariant distributions \cite[\S1.5]{theta-bernstein} gives an \(H\)-invariant distribution \(J\) on the fibre \(X\). For \(\Phi\in\mathcal S(G\times D)\) put
\[
 F_\Phi(g\cdot y)=\int_H\Phi(gh,h^{-1}\cdot y)\,d\nu(h),
 \qquad y\in X .
\]
This is the usual pushforward of the restriction to \(G\times X\). Compact support and the local sections above make \(F_\Phi\) a test on \(C\). The explicit induction formula in Frobenius descent reads
\[
 \Lambda(\Phi):=I(F_\Phi)
       =J\!\left(y\longmapsto\int_G\Phi(g,y)\,d\mu(g)\right).
\]
Indeed the integral over \(G/H\) in the formula is applied to the \(H\)-average defining \(F_\Phi\). Invariance of \(J\) absorbs \(h^{-1}\) in the fibre and quotient integration combines the two measures into \(\mu\). The identities can first be checked on compact open rectangles, since every test is a finite sum of such functions. In particular \(\Lambda(\Phi)=0\) whenever the inner integral vanishes on \(X\).

\phantomsection\label{route:18}\leanroute{Lean proves the needed implication by extending fibre averages and decomposing tests into finitely many cosets of a compact open subgroup. Bernstein's descent theorem gives the same functional, since both groups are unimodular and the projector map has a homogeneous base. The compact lift and its centralizer volume below still use the specific projector and level.}\formalnote{8.38}{38}

\subsection{A compact lift and its centralizer volume}
Average \(\phi\) over conjugation under \(K\), with probability Haar measure, and denote the result by \(\psi\). This is a finite average with multiplicities after passing to open cosets fixing the test, so \(I(\psi)=I(\phi)\). For \(y\in X\) we want a lift whose value at \((g,y)\) is
\[
       1_{KN}(g)\,\psi(g\cdot y).
\]
The set \(KN\) itself is generally unbounded. Put \(r(Z)=1+\epsilon Z(1-\epsilon)\) and
\[
 \rho(x)=r\bigl(\epsilon-p(x_w)\bigr),\qquad
 V=K\,\rho(\operatorname{supp}\psi)\subset KN .
\]
The map \(\rho\) is continuous, so \(V\) is compact and, as a union of right cosets of \(K\), open. If \(g=ka\in KN\) and \(\psi(g\cdot y)\ne0\), then \(x=a\cdot y\in\operatorname{supp}\psi\) by \(K\)-invariance. Writing \(a=1+Z\) with \(Z\) in the upper right block gives \(a\epsilon a^{-1}=\epsilon-Z\). Hence \(\rho(x)=a\) and \(g\in V\). Thus on \(G\times X\) the desired formula agrees with the restriction of the ambient test
\[
        \Psi(g,z)=1_V(g)\psi(g\cdot z).
\]
Its support is contained in the compact image of \(V\times\operatorname{supp}\psi\) under \((g,x)\mapsto(g,g^{-1}\cdot x)\).\formalnote{8.39}{39}

For \(x=g\cdot y\), the fibre integral of \(\Psi\) is \(\psi(x)\) times \(\nu\{h\in H:gh\in KN\}\). Choose \(h_0\in H\) such that \(gh_0\in KN\). The set measured is \(h_0(H\cap K)\). Indeed, for \(g_0=ka\in KN\), \(h\in H\), membership \(g_0h\in KN\) is equivalent to \(h\in K\). If \(h\in K\), use \(g_0h=kh(h^{-1}ah)\) and \(h^{-1}ah\in N\). Conversely if \(g_0h=k'a'\), then \(h(h^{-1}ah)(a')^{-1}=k^{-1}k'\in K\). The diagonal blocks of this upper triangular matrix and its inverse are those of \(h\) and \(h^{-1}\), so both are integral and \(h\in K\). Consequently
\[
          \Lambda(\Psi)=c_H I(\psi),\qquad c_H=\nu(H\cap K)>0 .
\]
The inequality and finiteness follow because \(H\cap K\) is compact open in \(H\) and contains the identity. No matching of the normalizations of \(\nu\) and \(\mu\) is required.\formalnote{8.40}{40}

\subsection{Radical translations force the other integral to vanish}
Let \(k_w\) be the residue field. For each \(i\in\operatorname{Mat}_n(k_w)\) choose an integral lift \(\widetilde i\), and put \(n_i=r(\widetilde i)\in K\cap N\). Finite cuspidality in the lifted-coset form gives
\(\sum_i\phi(xj(n_i k))=0\) for every \(x\in D,\ k\in K\). Every right translate of the finite radical occurring here has the same multiplicities since the upper right block is a surjective linear projection. Substitute \(j(k)^{-1}x\) for \(x\), then average over \(k\) to obtain \(\sum_i\psi(xj(n_i))=0\). The original vanishing over radicals of the finite group is supplied by \S\ref{sec:green-cuspidal}.\formalnote{8.41}{41}

Fix \(y\in X\). For each \(i\) there is \(m_i\in N\) with \(m_i\cdot y=yj(n_i)\). Here is the polynomial formula. Write \(b=y_w,\ x=bn_i,\ W=x-b\), so \(b\) is block diagonal, \(W\) is upper right and \(D_i=p(x)-p(b)\) is upper right, as is seen by taking powers of the block triangular matrix \(x\). Since \(p(b)=\epsilon\), the equation \(x(\epsilon+D_i)=(\epsilon+D_i)x\) gives \(bD_i-D_i b=W\). Indeed \(W\epsilon=WD_i=D_iW=0\) and \(\epsilon W=W\). Therefore \((1+D_i)x=b(1+D_i)\) and \(m_i=(1+D_i)^{-1}\) does the job at \(w\). This gives the asserted adelic equality since \(y j(b)^{-1}\) is trivial at the split factor and commutes with the image of \(G\).\formalnote{8.42}{42}

If \(g=ka\in KN\), then \(gm_i\in KN\) and
\[
       \Psi(gm_i,y)=\psi\bigl((a\cdot y)j(n_i)\bigr).
\]
Indeed \(K\)-invariance removes \(k\), and \(a,n_i\) commute in \(N\). If \(g\notin KN\), all \(gm_i\) are outside \(KN\). Thus \(\sum_i\Psi(gm_i,y)=0\) for every \(g\). Unimodularity of \(G\) makes \(\mu\) right invariant. Integrating the finite sum now gives \(|\operatorname{Mat}_n(k_w)|\int_G\Psi(g,y)\,d\mu(g)=0\). This proves the required zero integral for every \(y\in X\). Descent gives \(c_H I(\phi)=\Lambda(\Psi)=0\).\formalnote{8.43}{43}

In the trace comparison, \(I=I_R\) has the required invariance by right translation in the adelic quotient, and has the required support because each summand with rational polynomial \(R\) has local polynomial \(R_w\). A nonzero contribution forces \(R\) to be monic and to occur at the second place with simple zero on reduction. A proper power would give a multiple zero there, since its monic factor has integral coefficients. If \(R\) is reducible and is no proper power, split off the largest power of one of its irreducible factors. The two nonconstant factors are coprime and their B\'ezout identity persists at \(w\). The calculation just proved eliminates this last reducible case. The analytic input to the application is the existence of \(I_R\) as a linear functional for the allowed cusp kernels and their decay bound, discussed in \S\ref{sec:cusp-trace}.\formalnote{8.44}{44}


\partheading{The determinant lower bound}

\editionchapter{Volumes of single lifts and of the full chain}\label{ch:lower}

Theorem C is a lower bound for the volume of the same family to which Theorem A gives an upper bound. An earlier comparison for a single eigenform is Prasanna's integrality theorem for the quotient of the Petersson norms of a newform and its transfer to a Shimura curve, away from an explicit set of primes \cite[Theorem~2.4(1)]{prasanna-norms}. Here the quantity to be bounded is a determinant on the entire space of lifted forms. The product formula supplies the lower bound if one controls the determinant at the finite places. At each of these places one can choose a twist of the lifts, provided the resulting change in determinant is also bounded. This chapter makes that comparison.

\section{Proof of Theorem C}\label{sec:lower-C}

Fix the arithmetic data of \S\ref{sec:three-estimates}. In this chapter \(V=Y\) denotes the good space of forms on \(Y_0\), and \(d_T=\dim_{\mathbb C}V=D\). Write \(\Theta_\xi\) for the composite theta lift with finite-order twisting character \(\xi\). Its \(m\) steps use \(\xi,\ldots,\xi^m\) in order. The chain calculation in \S\ref{sec:chain-squares} proves injectivity. The terminal space carries an inner product and lattices of integral forms at finite places. Put
\[
 \begin{split}
 J&=\sum_\tau\big(\log(\kappa_\tau!)+
                  \kappa_\tau\log(M/\Omega)\big),\qquad
 F=\sum_\tau s_{\rm fin}(\kappa_\tau,\tau),\\
 U&=J-F,\qquad H=mkU,\qquad
 B=\frac{fa(\log b)^3}{mk},\qquad L=fa(\log b)^4.
 \end{split}
\]
Here \(\tau\) runs over the CM type and \(f\) is its cardinality. The integer \(\kappa_\tau\) is the place weight from Chapter~\ref{ch:upper}, equal to \(b\) on \(\Phi_b\) and \(a/x_1\) elsewhere. The term in \(J\) is the factorial and period contribution at that weight, with \(M=p_0p_1\) the principal level and \(\Omega\) the CM period. The term \(s_{\rm fin}(\kappa_\tau,\tau)\) in \(F\) is the corresponding finite-prime divisibility sum. The main term is \(H\), and \(B\) is the logarithmic budget for denominators. Thus \(mk(\log b)B=L\).

For independent \(y_1,\ldots,y_n\in V\) with \(n\ge(1-4/b^2)d_T\), let \(\delta\) and \(c\) be the upper bounds of Theorem C for their logarithmic index and their logarithmic squared covolume under \(\Theta_1\), respectively. The claim is
\[
                         c+2\delta\ \ge\ d_T(H-3L).
\]
The proof compares the Gram determinant at the trivial twist with its determinant at a different twist for each finite embedding. To do this, first bound the ratio of the Gram determinants of \(\Theta_1\) and \(\Theta_{\xi_\iota}\). Then complete the \(y_i\) to one basis that can be used at all the chosen twists and apply the product formula.\formalnote{9.1}{1}

\subsection{A height that permits a different twist at each prime}
For a fixed nonzero linear map between adelically normed vector spaces, Gaudron defines its height as a weighted sum of the logarithms of its local operator norms. The product formula makes this height invariant under multiplication by a nonzero scalar from the base field \cite[Def.~6.1 of the preprint]{gaudron}. For an injection into a filtered target, Bost bounds the normalized Arakelov degree of the source using the heights of the induced maps and the maximal slopes of the successive quotients of the target \cite[Prop.~4.6]{bost}. Here the maps at two different primes may be different theta lifts.

The product formula compares the absolute values of a single number.
When the twist varies with the prime, the determinant ratios vary too.
We charge one common multiplier at infinity and bound the remaining
denominators separately at the finite places.

An index \(\iota\) consists of a rational prime \(\ell\) and a field isomorphism \(\mathbb C\simeq\overline{\mathbb Q}_\ell\). Normalize \(|\ell|_\ell=\ell^{-1}\). For a family \(u=(u_\iota)_\iota\) of complex numbers, write \(\mathcal H(u)\le C\) if there exist \(h\in\mathbb C^\times\), \(\alpha\in\mathbb R\), and \(e_\ell\ge0\) satisfying
\[
 |\sigma h|\le e^\alpha,\qquad
 |\iota(h u_\iota)|_\ell\le\ell^{e_\ell},\qquad
 \alpha+\sum_{\ell\in S}e_\ell\log\ell\le C.
\]
These inequalities hold for every \(\sigma\in{\rm Aut}(\mathbb C)\), every \(\iota\), and every finite set \(S\) of primes, respectively. One multiplier \(h\) serves all indices. Multiplying two witnesses adds their logarithmic bounds and gives the product and power rules. If \(|\sigma z|\le e^v\) for every \(\sigma\) and \(z\ne0\), the constant family \(z^{-1}\) has bound \(v\) with multiplier \(h=z\). A factor \(N_\iota^{a/b'}\), for \(N_\iota\) a positive integer prime to \(\ell\), \(a\in\mathbb Z\), \(b'>0\) integral, has finite norm one because its \(b'\)-th power does. It can be inserted at no cost, even if the conductors \(N_\iota\) grow without a uniform Archimedean bound.

\begin{example}[A prime-dependent numerator]\label{ex:family-height}
Let \(u_\iota=N_\ell/6\), where \(N_\ell\) is any positive integer
prime to \(\ell\). The \(N_\ell\) can grow arbitrarily quickly.
Take \(h=6\). At every \(\iota\) above \(\ell\) the product
\(\iota(h u_\iota)=N_\ell\) has \(\ell\)-adic absolute value \(1\),
so \(e_\ell=0\) suffices at every prime. The archimedean cost is
\(\log 6\), giving \(\mathcal H(u)\le\log 6\).
For the reciprocal family \(6/N_\ell\), take \(h=1/6\). The finite
costs are still zero and the bound is now \(-\log 6\). These
two bounds are sharp. The two families multiply to \(1\), so any
improvement in either bound would give the constant family \(1\) a
negative bound, contrary to the product formula proved below.
\end{example}

The needed positivity assertion is \(\mathcal H(1)\le C\Rightarrow C\ge0\). For a witness \(h\), the bound at all isomorphisms above \(2\) implies \(h\) is algebraic. A transcendental \(h\) could be sent, by choosing images of a transcendence basis and extending to algebraic closures, to elements of arbitrarily large \(2\)-adic norm. Set \(n_h=[\mathbb Q(h):\mathbb Q]\) and \(q=N_{\mathbb Q(h)/\mathbb Q}(h)\in\mathbb Q^\times\). Each of the \(n_h\) factors of \(q\) has modulus at most \(e^\alpha\) at infinity and norm at most \(\ell^{e_\ell}\) above \(\ell\). Take \(S\) to contain the primes dividing the numerator or denominator of \(q\). Rational prime factorization gives
\[
 0=\log|q|+\sum_{\ell\in S}\log|q|_\ell
   \le n_h\left(\alpha+\sum_{\ell\in S}e_\ell\log\ell\right)\le n_h C.
\]
Thus the product formula applies to a single rational number. Intermediate families may have negative upper bounds. This is how the term \(-H\) will survive.\formalnote{9.2}{2}

\section{Regular twists and the cost of an orthogonal projection}\label{sec:regular-twists}

For a fixed \(j\), \(1\le j\le m\), let \(T_\eta\) be the theta lift from the selected forms on \(W_j\) to forms on \(V_j\), with twist \(\eta\). Put \(S_\eta=\operatorname{im}T_\eta\), and let \(P_\eta\) be orthogonal projection onto \(S_\eta\). Orthogonal projection can introduce a denominator even when the subspace is spanned by integral vectors.

\begin{example}[The denominator of a projection onto a line]\label{ex:projection-denominator}
Take \(v=(1,2)^{\mathsf T}\) with the usual inner product on \(\mathbb Q^2\).
Projection onto \(\mathbb Qv\) has matrix
\[
 P=\frac{vv^{\mathsf T}}{v^{\mathsf T}v}
   =\frac15\begin{pmatrix}1&2\\2&4\end{pmatrix}.
\]
Thus \(5P\) preserves \(\mathbb Z^2\). Since \(P(1,0)^{\mathsf T}=v/5\), every integer that clears the denominator is divisible by \(5\).
\end{example}

For the lattice of integral automorphic forms we seek an operator \(\Lambda\) on the lifted space \(S_\eta\) such that \(\Lambda P_\eta\) sends the integral test vectors specified below to integral vectors. The size of the denominator cleared in this way is measured by the valuation of \(\det\Lambda\). We construct \(\Lambda\) from Hecke operators, so this valuation depends on congruences between the eigenvalues on \(S_\eta\) and those on its complement.

For a twist \(\eta\) with auxiliary modulus \(Q_\eta\), write \(\operatorname{Nm}Q_\eta=\prod_{w\in Q_\eta}q_w\), with the empty product equal to one. A regular family chooses \(\xi_\iota\) with \(\ell\nmid\operatorname{Nm}Q_{\xi_\iota}\). This coprimality condition concerns the new twisting places. For each \(j\) the family supplies nonnegative logarithmic losses \(\beta_\ell\), independent of the embedding above \(\ell\), with \(\sum_{\ell\in S}\beta_\ell\le B\) for all finite sets \(S\) of primes. With \(\eta=\xi_\iota^j\), there is an operator \(\Lambda\) on \(S_\eta\) satisfying
\[
 |\iota(\det\Lambda)|_\ell\ge \exp(-d_T r\beta_\ell).
\]
When \(u\) is integral and orthogonal to the false lifts, \(\Lambda P_\eta u\) is integral. The false lifts lie in the orthogonal complement of the genuine lifts. They have the selected type and are killed by the test value of the Hecke polynomial at some split place away from the level. There is also a bound for a single lifted line. A critical pair \((\nu,z)\) consists of an integer \(1\le z\le t=5r\) and an automorphic character \(\nu\) of the norm-one torus. Its weight at each infinite place \(\tau\) is \(-c(\tau)\), with \(z\le c(\tau)\le4b\), and its tame residue characters are trivial off the chosen places over \(p_0,p_1\). There are finitely many such pairs. The integer, weights and residue characters range over finite sets, and two characters with the same weight and tame part differ by a character of the finite class group of the torus. The partial Hecke \(L\)-value is evaluated at \(z\). A \emph{main} pair at step \(j\) gives a lift from \(U(1)\) to \(U(2z)\), with \(m<z\le r\). Its tame character is the \(j\)-th power of the fixed residue character at the chosen level. At each infinite place \(\tau\), put \(w_j(\tau)=\kappa_\tau+(2j+1)r^8\), the negative of the weight of the \(j\)-th lifting character. The main-pair condition is \(w_j(\tau)+r-t\le c(\tau)\le w_j(\tau)\). The separate rank-one bound for these pairs is
\[
 \big|\iota\langle P\Xi,P\Xi'\rangle\big|_\ell\le e^{\beta_\ell},
\]
where \(P\) is the orthogonal projection onto the line spanned by the rank-one theta lift. The test forms \(\Xi,\Xi'\) must be integral at every embedding \(\mathbb C\hookrightarrow\mathbb C_\ell\) and orthogonal to the false lifts in the kernel of the rank-one descent, using the selected type with label \(X-1\). Finally, for each \(j\), each ramified critical pair, and either sign of the twist \(\xi_\iota^j\), the family of inverse ratios of twisted and untwisted absolute partial \(L\)-values has height at most \(x_1\). A ramified critical pair has nontrivial tame character exactly at the original chosen places. The auxiliary primes in \(Q_\xi\) are additional.\formalnote{9.3}{3}

\subsection{One character for a finite set of Selmer problems}

Fix \(\ell,\iota\). Only finitely many Hecke eigenvalue systems occur in the finite-dimensional good space on \(W_m\). For each system and each \(j\), take the rank-\(r\) Galois representation attached to the source eigensystem. In the normalized trace formula of \S\ref{sec:congruence-detail}, write \(R_x\) for this main summand after the normalization there. The trace of the lift is the sum of \(\psi\operatorname{tr}R_x\) and the \(t'\) terms \(\epsilon^{-e}\) for \(0\le e<t'\), where \(t'\) is the rank increase, \(\psi\) is the inverse Galois character of the auxiliary twist, and \(\epsilon\) is the cyclotomic character. For each main critical pair take an additional summand of rank one. Chapter~\ref{ch:galois} constructs the attached representations and establishes the required local properties. Over a sufficiently large coefficient ring \(O\), choose a lattice model of the summand twisted by the inverse Galois character corresponding to the twist. The module whose length is bounded consists of continuous cocycles with coefficients in this model tensored with \({\rm Frac}(O)/O\). They are unramified outside \(\ell p_0p_1\) and \(Q\cup\bar Q\), satisfy the specified crystalline \(H^1_f\) relations above \(\ell\), and vanish on a distinguished element acting by a scalar on the twisted lattice. The length is normalized by \(\operatorname{len}(O/\ell O)=1\).\formalnote{9.4}{4}

The general twisting theorem treats finitely many such lattices \(T_i\) together with finitely many Stickelberger systems. A system supplies elements \(\Theta_Q\) of group rings of ray-type class groups, compatible under norm maps with Euler factors. In the application, evaluating a suitable system at a character controls its algebraic \(L\)-value ratio. The simultaneous bounds concern a twisted Selmer length, the valuation of the evaluation of \(\Theta_Q\), and the sum of \(v_\ell(\psi(\gamma_w)-1)\) at tame generators of the primes of \(Q\) and their conjugates. The last enters the ramification term of the projection estimate.\formalnote{9.5}{5}

Work first with \(T_i/\varpi^kT_i\) over a discrete valuation ring and bounded \(k\). At a new prime \(w\) and its conjugate impose a relation on cocycles of the form
\[
                       z(\gamma_w)=s_i z(\varphi_w),
\]
where \(\gamma_w\) is a tame generator and \(\varphi_w\) a Frobenius element. Slope \(s_i=0\) imposes unramifiedness there. For one selected datum a prime and a slope can reduce the total length whenever a shortening step is available. Slope zero for the other data preserves their unramified condition at this prime. The sum of ordinary composition lengths over the finitely many data and powers drops by at least one. The descent terminates. At the terminal stage reciprocity bounds the dual group at the exponent of the tied group, and comparison of lengths bounds the tied group. This duality estimate and the auxiliary-prime theorem prescribing Frobenius are inputs to the argument. Termination alone would give no numerical bound.\formalnote{9.6}{6}

To pass from slopes to characters one constructs homomorphisms \(\lambda_w\) with values in \(\mathbb Z/\ell^N\mathbb Z\). The constraints are represented by finitely many polynomials in variables indexed by the auxiliary primes. Simultaneous specialization at roots of unity produces \(\zeta_w\) for which the single character
\[
                      \psi(g)=\prod_{w\in Q}\zeta_w^{\,\lambda_w(g)}
\]
has the required bounds for its first \(m\) powers. In the application the character comes from \(U(1)\), has \(\ell\)-power order, is trivial on decomposition groups above \(\ell\), and ramifies at the auxiliary primes. These primes avoid the bad set and its conjugate. The resulting bounds for the relevant divisible-coefficient Selmer lengths and inverse ratios use
\[
               e'_\ell=(4f+3)r\,v_\ell(N_{\rm Aut})+2^{-\ell},
\]
where \(N_{\rm Aut}=\prod_{w\mid\pi_M}|\mathrm{GL}_{rt}(k_w)|\) is a common multiple of the orders of the class stabilizers. Taking logarithms of finite norms converts valuation \(e'_\ell\) into \(e'_\ell\log\ell\). In particular,
\[
 \sum_{\ell\in S} e'_\ell\log\ell
   \le (4f+3)r\log N_{\rm Aut}+2\le x^{999}\le x_1.
\]
The first inequality uses prime factorization and the convergent sum of \(2^{-\ell}\log\ell\). The multiplier one gives the promised height bound \(x_1\) for ramified pairs.\formalnote{9.7}{7}

\subsection{From cocycles to a denominator}

To estimate the denominator cleared by \(\Lambda P_\eta\), we use a quantitative version of Ribet's method. A two-dimensional calculation explains why cohomology appears. If a triangular representation has diagonal characters \(\chi_1,\chi_2\) and off-diagonal entry \(b(g)\), multiplication of matrices gives
\[
 c(gh)=c(g)+(\chi_1/\chi_2)(g)c(h),
 \qquad c(g)=b(g)/\chi_2(g).
\]
Changing the splitting changes \(c\) by a coboundary. At finite coefficient depth, possible off-diagonal entries therefore yield classes with prescribed local conditions. The actual argument treats a main summand and a string of cyclotomic characters in a conjugate self-dual representation, allowing congruent constituents. The two-dimensional identity explains the coefficient module. The higher-rank separating-element theorem supplies the numerical estimate.

At auxiliary primes the construction uses relations in a polynomial algebra of the form
\[
 \gamma_w-1-X_w(\varphi_w-1)H_w(\varphi_w).
\]
Here \(H_w\) is a polynomial from the tied local data. The resulting separating Hecke element, evaluated on a lift, has zero coefficients at the proper divisors of the squarefree monomial \(\prod_w X_w\), while the coefficient of that monomial has bounded valuation. To obtain this coefficient, even Tate classes are tested on a triangular sequence of tame generators and Frobenius elements. A new row controls classes satisfying the vanishing conditions of the previous rows, with a bounded power-of-\(\ell\) loss. This avoids charging the full sizes of the even groups. The odd classes enter through one positive integer \(X\) whose valuations bound the lengths of the specified odd cocycle modules. Write \(N_E=q_1d\) for the modulus with \(E\subseteq\mathbb Q(\mu_{N_E})\), and let \(h_E\) be the class number of \(E\). For twist range \(t\), the arithmetic input gives
\[
               X\le 2^{\,4(N_E+t+3)^{27}}h_E .
\]
Its scope is the odd module and local conditions defined in the cyclotomic argument. Using a common integer bounds the global sum of its valuations with weights \(\log\ell\) by \(\log X\). Section~\ref{sec:congruence-detail} gives the bound for the odd cocycles and explains how to choose the successive Frobenius tests for the even cocycles.\formalnote{9.8}{8}

For source rank \(r'\ge1\) with rank difference \(t'\le t\), the bound on a separating eigenvalue's valuation, before clearing stabilizers, is
\[
     {\rm sel}+2t'{\rm ram}+r'{\rm loc}_\ell+e_\ell^{\,0}.
\]
Here \({\rm sel}\) bounds the twisted Selmer length, \({\rm ram}\) bounds tame ramification, \({\rm loc}_\ell\) measures the defect of the local lattice comparison at primes above \(\ell\), and \(e_\ell^{\,0}\) collects valuations from the odd classes and finite group orders, as specified in \S\ref{sec:congruence-detail}. The twist gives \({\rm sel},{\rm ram}\le e'_\ell\). The separating operators clear the projected integral test forms. To choose one integral combination for all eigenlines, normalize the list of eigenvalues on each line by its largest \(\ell\)-adic absolute value, dividing by one eigenvalue of that size. The reductions of the combination coefficients must avoid the corresponding cancellation hyperplane on each line. There are finitely many such hyperplanes and the algebraic residue field is infinite. The combination keeps the largest absolute value on every line. Multiplying by \(N_{\rm Aut}\) clears stabilizer denominators in the inner products. Set
\[
       C_\ell={\rm loc}_\ell+(1+2t)e'_\ell
                     +e_\ell^{\,0}+v_\ell(N_{\rm Aut}).
\]
The resulting operator \(\Lambda\) has \(v_\ell(\iota\mu)\le r'C_\ell\) for every eigenvalue \(\mu\). Taking their product with multiplicities gives determinant norm at least \(\ell^{-d'r'C_\ell}\) on a space of dimension \(d'\). This proves the required estimate when \(\beta_\ell\ge C_\ell\log\ell\). For a main-pair line the pairing estimate also involves the number \(h_U\) of classes of the finite \(U(1)\) level quotient. The common allowance is consequently\formalnote{9.9}{9}
\[
                    \beta_\ell=(C_\ell+v_\ell(h_U))\log\ell .
\]

The contribution of the crystalline local comparisons is bounded by \(4\cdot10^6 e^{180} f(m/k)b\log b\). Summing the remaining bounds gives, for every finite \(S\),
\[
 \begin{split}
 \sum_{\ell\in S}\beta_\ell
 &\le 4\cdot10^6 e^{180} f(m/k)b\log b+(1+2t)x^{999}\\
 &\hspace{20mm}{}+x^{4000}+\log N_{\rm Aut}+x^{30}\ \le B .
 \end{split}
\]
For the last inequality use \(a=xb\). Dividing the first term by \(B\) gives \(4\cdot10^6 e^{180}m^2/[x(\log b)^2]\), which tends to zero since \(m\le2\sqrt{x}\). All remaining terms grow at most polynomially in \(x\), including \(\log N_{\rm Aut}\) by the class stabilizer estimate. The scale \(b>\exp(x^\eta)\) absorbs any such polynomial for fixed \(\eta>0\). This proves the required sum bound simultaneously for all finite \(S\). The crystalline lattice and global duality estimates behind the first inequality are proved in \S\S\ref{sec:kisin-height}, \ref{sec:hodge-cup} and~\ref{sec:selmer-formula}; the finite-depth operator is constructed in \S\S\ref{sec:trace-ideal} and~\ref{sec:deformed-norm}.\formalnote{9.10}{10}

\section{The determinant of a single theta lift}\label{sec:single-lift-norm}

\subsection{Comparison by a positive operator}

For one \(j\), keep the notation \(T_\eta,S_\eta,P_\eta\). The multiplicity theorem identifies the dimension of its source with \(d_T\). Write
\[
 Z_j(\eta)=\det(T_\eta^*T_\eta)>0,\qquad
 \rho_\eta=\frac{\mathcal L_j(\eta)}{\mathcal L_j(1)},
\]
where \(\mathcal L_j(\eta)\) is the product of the three factors evaluated in \S\ref{sec:three-exchanges}. There is a gluing map \(g_\eta\) from tuples of forms on a smaller definite space to forms on \(V_j\). The theta identities for exchanging the order of lifts give a common period map \(T^\circ\), independent of the twist, and scalars \(c'_\eta\ne0\) such that
\[
       g_\eta^*T_\eta=(c'_\eta)^{-1}T^\circ,\qquad
       \rho_\eta=|c'_1|^2/|c'_\eta|^2 .
\]
In particular the ratio is nonzero. If \(h_j\) is the hook number of the auxiliary weight diagram, put \(A_\eta=N_{\rm Aut}h_j^2g_\eta g_\eta^*\). Section~\ref{sec:common-period} constructs an integral frame for this operator, bounds its norm under every conjugation at the trivial twist, and proves that its value on a genuine lift is orthogonal to false lifts. The same section proves injectivity of \(T^\circ\) by testing a nonzero form at a peak.\formalnote{9.11}{11}

Compress the comparison operator \(A_\eta\) to the lifted space \(S_\eta\) by putting \(D_\eta=P_\eta A_\eta|_{S_\eta}\). Write \(\Delta_\eta=\det D_\eta\). Since \(T_\eta^*=T_\eta^*P_\eta\), factor \(T_\eta\) through its image as an isomorphism \(\widetilde T_\eta\) followed by inclusion. Taking determinants of the resulting square maps gives
\[
 \det(T_\eta^*A_\eta T_\eta)
 =\det(\widetilde T_\eta^*)\det(D_\eta)\det(\widetilde T_\eta)
 =Z_j(\eta)\Delta_\eta .
\]
On the other hand \(T_\eta^*A_\eta T_\eta=N_{\rm Aut}h_j^2|c'_\eta|^{-2}(T^\circ)^*T^\circ\). The untwisted operator on the source is positive definite. Taking its determinant and dividing the identity for \(\eta=1\) by the identity for \(\eta\) yields
\[
 \boxed{\quad \frac{Z_j(1)}{Z_j(\eta)}
       =\frac{\Delta_\eta}{\Delta_1}\rho_\eta^{-d_T}. \quad}
\]
This is where the dimension enters the exponent on the \(L\)-value side.

Let \(A=fmb(\log b)^2\). For an eigenvector of \(D_1\) with eigenvalue \(\lambda\), the norm bound for \(A_1\), followed by contraction of the projection on its false-orthogonal image in the conjugated norm, gives \(|\sigma\lambda|\le e^{3A}\). Taking the product gives \(|\sigma\Delta_1|\le e^{3d_T A}\). At the finite place, \(\Lambda D_\eta\) preserves integral vectors in \(S_\eta\) because \(A_\eta\) preserves integrality and maps them into the subspace where \(\Lambda P_\eta\) clears the denominator. To see why the determinant of the composite is integral, scale a nonzero eigenvector so that its coordinates are integral and one has maximal norm one (in the integral-form coordinates). Any scalar multiplying this vector into an integral vector has norm at most one. Apply this to the eigenvalue and take products. Thus
\[
 |\iota(\Delta_\eta)|_\ell\le|\iota(\det\Lambda)|_\ell^{-1}
                   \le e^{d_T r\beta_\ell}.
\]
Taking \(h=\Delta_1\) in the definition of family height, and using \(\sum\beta_\ell\le B\), proves\formalnote{9.12}{12}
\[
             \mathcal H\big((\Delta_{\xi_\iota^j}/\Delta_1)_\iota\big)
                          \le d_T(3A+rB).
\]

\subsection{The gain from one critical value}

For a main critical pair \((\nu,z)\), let \(R(\eta)\) be the ratio of the absolute partial Hecke values at \(z\) for \((\nu\eta)^{-1}\) and \(\nu^{-1}\). Both values omit the Euler factors at the chosen places and their conjugates. The numerator also omits the factors at the places in \(Q_\eta\) and their conjugates. Write \(\phi_\eta\) for the lift from \(U(1)\) to \(U(2z)\). Its algebraic part is its squared norm divided by \((\operatorname{Nm}Q_\eta)^{2z-1}\). The ratio formula gives
\[
                R(\eta)^{-1}=(\operatorname{Nm}Q_\eta)^{2z-1}
                             \frac{\|\phi_1\|^2}{\|\phi_\eta\|^2}.
\]
Let \(\phi_0\) be the untwisted lift for an empty row diagram of the
same rank \(2z\). The common denominator for its values is \(M^{4fz}\).
The class weights in its squared norm have denominators dividing
\(N_{\rm Aut}\). Thus
\[
                 a_0=N_{\rm Aut}M^{8fz}\|\phi_0\|^2
\]
is integral at every prime. The period constant in the archimedean
norm formula depends on the rank, so it cancels on comparing the
two lifts. The jet estimate then gives
\(\lvert\sigma(a_0/\|\phi_1\|^2)\rvert\le\exp(B-J)\) for every
\(\sigma\). With \(h=a_0/\|\phi_1\|^2\), the product
\(h\|\phi_1\|^2=a_0\) is integral, whence
\(\mathcal H(\|\phi_1\|^2)\le B-J\).
Section~\ref{sec:common-period} constructs the empty row and computes
the archimedean bound.\formalnote{9.13}{13}

For the twisted norm a bound in the other direction is needed. Let \(\phi=\phi_\eta\) and let \(P\) project onto \(\mathbb C\phi\). Apply the construction of test forms twice, once to make the pairing large at \(\iota\) and once at \(\iota\) composed with complex conjugation. The resulting forms \(\Xi,\Xi'\) both satisfy the integrality and orthogonality conditions of the rank-one bound. Set \(D=\langle\Xi,\phi\rangle\langle\phi,\Xi'\rangle\). With \(r_0(\eta)\) the ratio for one auxiliary ramified empty-row pair and \(c_0\ne0\) independent of \(\iota\), the two estimates give
\[
 |\iota(D)|_\ell\ge
 |\iota(c_0\overline{c_0}\, r_0(\eta)^2)|_\ell
                         \ell^{-2P_\ell},\qquad |\sigma c_0|\le e^B,
\]
where \(P_\ell\) is the local prime sum and \(2\sum_{\ell\in S}P_\ell\log\ell\le F\). In particular \(D\ne0\). The tests are built from class sums. By Fourier orthogonality on the finite abelian \(U(1)\) quotient, the normalized sums twisted by all its characters add to a selected value of the original function. The ultrametric triangle inequality gives at least one character whose sum has norm as large as that value. Applying this selection to the theta-jet construction uses the sharpness of the finite-place derivative estimate. The resulting lower bound on each pairing gives the displayed inequality for \(D\).\formalnote{9.14}{14}

The formula for projection onto a line gives \(\langle P\Xi,P\Xi'\rangle=D/\|\phi\|^2\). Apart from the conductor unit, \(R(\eta)^{-1}\) therefore factors as
\[
 \|\phi_1\|^2\,
 \langle P\Xi,P\Xi'\rangle\,
 r_0(\eta)^{-2}\,
 \frac{r_0(\eta)^2}{D}.
\]
The four height bounds in this order are \(B-J,\ B,\ 2x_1,\ 2B+F\). For the last take multiplier \(c_0\overline{c_0}\) and use the displayed test inequality. For the second use the line-pairing clause of regularity. Since \(x_1\le B\) the sum is at most \(6B-(J-F)=6B-U\). The possible saving \(U\) has thus been identified for a single critical value before counting how many times it occurs.

\subsection{Counting the critical values in one step}

The three strings in \(\mathcal L_j\) consist of squared absolute partial \(L\)-values at integers \(1\le z\le z_i\), times powers of the norms of omitted places. Their endpoints are \(7m/2+k/2\), \(m\), and \(m/2\). The upper \(k/2\) positions of the first string are main pairs. The remaining positions number \(7m/2+m+m/2=5m\), and each admits the ramified-pair bound. Here \(m,k\) are even. The level factors give the exponent \(e=\sum_i z_i^2\); absorb the signs of the two nonprimary twists into the appropriate ratios \(R'_v\). Then
\[
  \rho_\eta=(\operatorname{Nm}Q_\eta)^e
  \left(\prod_{u=1}^{k/2} R_u(\eta)
        \prod_{v=1}^{5m} R'_v(\eta)\right)^2 .
\]
Taking inverse heights and remembering the square gives
\[
       \mathcal H\big((\rho_{\xi_\iota^j}^{-1})_\iota\big)
           \le 2\big((k/2)(6B-U)+5mx_1\big)
           =k(6B-U)+10mx_1 .
\]
The \(k\) multiplying \(U\) is thus an exact count. Combining this with the compression formula gives a bound per dimension of \(3A+rB+k(6B-U)+10mx_1\). The scale inequalities give \(3A+rB+6kB+10mx_1\le k(\log b)B\). We have proved for every \(j\) that\formalnote{9.15}{15}
\[
 \boxed{\quad
  \mathcal H\!\left(\left(\frac{Z_j(1)}{Z_j(\xi_\iota^j)}\right)_\iota\right)
     \le d_T\big(L/m-kU\big).\quad}
\]

\section{Chain determinants and completion of a basis}

\subsection{Multiplication of single-step determinants}\label{sec:chain-step-norms}

Write \(Z(\xi)=\det(\Theta_\xi^*\Theta_\xi)>0\) for the Gram determinant of the composite theta lift from \(Y_0\) to \(Y_m\). The single lifts considered above have distinct sources \(W_j\) and targets \(V_j\), so multiplying their matrices would not calculate this determinant. The connection is obtained by transferring the same spherical Hecke operators to a common source. In the Rallis inner-product formula for step \(j\), write \(\mathcal R_j(\eta)\) for the resulting Hecke series acting on \(V\). Equality of the selected Hecke multiplicities on \(W_j\) and \(Y_0\) gives
\[
                Z_j(\eta)=(a_j c_Q)^{d_T}
                                \det\nolimits_V\mathcal R_j(\eta).
\]
Here \(a_j\ne0\) is independent of \(\eta\), and \(c_Q\) is the local normalizing factor for the modulus. On each common eigenspace the determinant raises the Hecke eigenvalue to the \emph{multiplicity} of that eigenspace.

Index the intermediate lifts by consecutive strings of blocks \(i+1,\ldots,i+h\). The empty string gives the base space, and adjoining its first block gives a lift to which the convergent inner-product formula applies. Exchanging adjacent lifts reduces the Gram operator of the full chain to the operators of these single lifts. On \(V\) the result is
\[
 \Theta_\xi^*\Theta_\xi
   =c(\xi)c_Q^{\,m}\,\mathcal R_1(\xi)
                        \mathcal R_2(\xi^2)\cdots\mathcal R_m(\xi^m).
\]
For all powers of one allowed \(m\)-fold twist the local normalizing factor in this identity agrees with the step factors. The exchange theorem evaluates the absolute value of \(c(\xi)/c(1)\) up to an integral half-power of the conductor. Positivity of the Gram determinants then gives the following exact relation for some integer \(e\) independent of \(\xi\).
\[
 \begin{split}
 \frac{Z(1)}{Z(\xi)}
   &=A(\xi)^{-d_T}\prod_{j=1}^{m}\frac{Z_j(1)}{Z_j(\xi^j)},\\
 A(\xi)&=(\operatorname{Nm}Q_\xi)^{e/2}
       \prod_{1\le i<j\le m} R_{ij}(\xi).
 \end{split}
\]
Indeed, taking determinants replaces the scalar \(c_Q^m\) by \(c_Q^{md_T}\), which also occurs in the product of the \(m\) step formulas. The fixed \(a_j\)'s disappear after division by the trivial-twist formula. Each \(R_{ij}\) is a product of \(t\) ratios of critical values for the quotient of the \(j\)-th and \(i\)-th untwisted lift characters, at \(\xi^{j-i}\). The exchange theorem and multiplicity comparison justify this factorization. Section~\ref{sec:chain-squares} gives the square laws, the induction through the triangle, and the cancellation of the level indices.\formalnote{9.16}{16}

Each elementary inverse ratio in \(A\) satisfies the ramified bound \(x_1\), and the conductor term has height zero at the chosen finite embeddings. Hence
\[
 \mathcal H\big((A(\xi_\iota)^{-1})_\iota\big)
       \le \binom{m}{2}t x_1\le m^2t x_1\le b.
\]
Adding the \(m\) single-lift estimates now gives the determinant estimate for the entire chain,\formalnote{9.17}{17}
\[
  \mathcal H\!\left(\left(\frac{Z(1)}{Z(\xi_\iota)}\right)_\iota\right)
 \le d_Tb+md_T(L/m-kU)=d_T(b+L-H).
\]

\subsection{A simultaneous completion with controlled cost}

The determinant \(Z\) concerns a whole basis, whereas Theorem C starts with only \(n\) independent vectors. A single completion must serve every auxiliary conductor so that the same complex determinant occurs in the product formula. We use the following spanning result about the theta kernels. There is a function \(\gamma_\xi\ne0\), \(\gamma_1=1\), and a spanning set \(\mathcal A\) of \(V\) such that, for each \(z\in\mathcal A\),
\[
 \begin{gathered}
       \gamma_\xi\Theta_\xi z\text{ is integral at every embedding above }
                      \ell\nmid \operatorname{Nm}Q_\xi,\\
       \|\sigma\Theta_1 z\|^2\le e^{h_s}\quad\hbox{for every }\sigma,
                 \qquad h_s=2x_1^2b^2 .
 \end{gathered}
\]
The notation \(\sigma\Theta_1z\) means applying \(\sigma\) to its coordinates in the fixed rational Gram model. The scalar \(\gamma_\xi\) works simultaneously for all members of \(\mathcal A\).

To obtain one spanning family for every twist, start with the periods on the two paths of the grid in \S\ref{sec:grid-completion}. On the last path the period is a finite sum over tuples of residue matrices, paired with the lifts along its two halves. Powers of the fixed level and field parameters clear the denominators of each kernel away from the conductor. The number and sizes of the summands bound the untwisted conjugate norms. The first path recovers the coordinates of the lifted form by pairing with delta forms. Comparing the two paths by \(m^2\) exchanges transfers these bounds to the coordinates, with a scalar \(s_\xi\) independent of the input form. A nonzero multiple of the projection onto \(V\) sends forms on \(Y_0\) to \(V\) and preserves the required integrality. Its archimedean loss gives \(x^{99}+x^{202}b^2\le h_s\).\formalnote{9.18}{18}

Normalize \(\gamma_\xi=s_\xi/s_1\). The exchange formula gives
\[
 |\gamma_\xi|=(\operatorname{Nm}Q_\xi)^{e'/2}\prod_{0\le i,j<m} R^{\rm turn}_{ij}(\xi)
\]
for another fixed integer \(e'\). Each turn contributes \(t\) elementary ratios with the ramified height bound. Squaring and inverting therefore gives
\[
       \mathcal H\big((|\gamma_{\xi_\iota}|^{-2})_\iota\big)
                 \le 2m^2 t x_1\le b .
\]
The scalar \(\gamma_\xi\) from the square grid is applied to the vectors added to complete the basis. Choose \(z_1,\ldots,z_{n'}\) from \(\mathcal A\) successively outside the span of the preceding vectors and the \(y_i\), until they complete a basis. Then \(n'=d_T-n\le4d_T/b^2\).

\subsection{Mixed Gram determinants at complex and finite embeddings}

Let \(w=(y_1,\ldots,y_n,z_1,\ldots,z_{n'})\) be the fixed completed basis, and put \(G_\xi=\det(\langle\Theta_\xi w_i,\Theta_\xi w_j\rangle)_{ij}\). It is positive because \(\Theta_\xi\) is injective. In an orthonormal basis of \(V\), let \(W\) be the coordinate matrix of \(w\). The formula
\[
 G_\xi=\det(W^*(\Theta_\xi^*\Theta_\xi)W)
        =|\det W|^2 Z(\xi)
\]
proves \(G_1/G_\xi=Z(1)/Z(\xi)\). It also explains why an arbitrary completion that depends on \(\xi\) would be unusable.

The covolume hypothesis gives bound \(e^c\) for the positive Gram determinant of the conjugated coordinates of the first \(n\) vectors at twist one, for every \(\sigma\). Appending each \(z_i\) multiplies that determinant by at most \(e^{h_s}\), since the factor contributed by a new vector is the squared length of its orthogonal component to the previous span. This is the Gram form of Hadamard's inequality and extends by continuity to dependent lists. Applying \(\sigma\) directly to \(G_1\) deserves care. If \(j_{\mathbb C}\) denotes complex conjugation, put \(\sigma'=j_{\mathbb C}\sigma j_{\mathbb C}\). Rationality of the Gram matrix makes \(\sigma G_1\) the mixed Gram determinant between the \(\sigma'\)-conjugated and \(\sigma\)-conjugated coordinate lists. On the top exterior power, Cauchy--Schwarz bounds its modulus by the geometric mean of their two positive Gram determinants. Hence
\[
                       |\sigma G_1|\le e^{c+n'h_s}.
\]
The multiplier \(G_1\) bounds the constant family \(G_1^{-1}\) by \(c+n'h_s\).

For the index hypothesis choose \(a_\ell\ge0\) such that \(\sum_{\ell\in S}a_\ell\log\ell\le\delta\) for each finite \(S\). At each \(\iota\) there is a matrix \(U_\iota\) over its valuation ring with \(|\iota(\det U_\iota)|_\ell\ge\ell^{-a_\ell}\). Its combinations of the \(y_i\) lift to integral forms for \emph{all} twists allowed at \(\ell\). Fix \(\xi=\xi_\iota\) and set \(\bar\iota=\iota\circ j_{\mathbb C}\). Use \(U_\iota\) on the linear side of the inner product and \(U_{\bar\iota}\) on the conjugate-linear side. Both work for the same \(\xi\), and they may be different matrices. Extend them on the last \(n'\) coordinates by the block \(\gamma_\xi I\). The twisted lattice pairs the resulting two integral lists integrally under \(\iota\). Taking determinants of their mixed Gram matrix gives
\[
 |\bar\iota(\det U_{\bar\iota})|_\ell\,
 |\iota(\det U_\iota)|_\ell\,
 |\iota((\gamma_\xi\overline{\gamma_\xi})^{n'}G_\xi)|_\ell\le1 .
\]
Consequently the last norm is at most \(\ell^{2a_\ell}\), and the multiplier one proves
\[
             \mathcal H\big((|\gamma_{\xi_\iota}|^{2n'}G_{\xi_\iota})_\iota\big)
                        \le 2\delta .
\]
The index is doubled because the pairing uses two integral lists.\formalnote{9.19}{19}

\subsection{The product formula closes the inequality}

The basis and \(G_1\) are fixed, while \(\xi=\xi_\iota\) can vary. At each embedding the four nonzero factors satisfy
\[
       1=G_1^{-1}\,
            \frac{Z(1)}{Z(\xi)}\,
            (|\gamma_\xi|^{-2})^{n'}\,
            |\gamma_\xi|^{2n'}G_\xi .
\]
The height bounds for these factors and nonnegativity for the constant family \(1\) give
\[
 0\le c+n'h_s+d_T(b+L-H)+n'b+2\delta.
\]
Only \(n'\le4d_Tb^{-2}\) vectors were added, so the completion costs
\[
 n'(h_s+b)
 \le 4d_Tb^{-2}(2x_1^2b^2+b)
 =4d_T(2x_1^2+b^{-1})\le d_Tb .
\]
By the scale bounds, the charges \(d_Tb\) for the chain exchanges and for completion are each at most \(d_TL\).\formalnote{9.20}{20}
This gives Theorem C,
\[
                c+2\delta\ge d_T(H-L-2b)\ge d_T(H-3L).
\]


\editionchapter{Theta exchange, local factors and the norm formula}\label{ch:exchange}

For the theta lift of the constant function \(1\) from \(U(1)\) to a definite \(U(2k)\), \(k\ge2\), the squared norm at the prescribed level is an explicit multiple of a partial Hecke \(L\)-value. The Hecke character enters through the splitting used in the Weil representation. The Rallis formula relates the inner product of two theta lifts to a doubling integral \cite[\S\S11.3--11.5]{gan-qiu-takeda-rallis}. For the pair \(U(1)\times U(1)\), Yang evaluates a central Hecke value as a squared norm, with the local data and root number specified in his formula \cite[Theorem~0.3]{yang-hecke}. For the definite target \(U(2k)\) used here the source computes the \rootlink{grh/src/Lib/SiegelWeil/RallisInner/SharpData.lean}{local sum and Euler product} in \texttt{Lib} and states the \rootlink{grh/src/New/OneLift/HeckeValues/CharacterLiftNorm.lean}{norm formula with the chosen levels and measures} in \texttt{New}. Chapter~\ref{ch:lower} needs the change in the determinant when the twisting character is changed. For this we choose matched test functions for three successive exchanges of theta integrals and keep track of all three scalar factors.

\section{The three theta exchanges}
\label{sec:three-exchanges}
To compare twisted and untwisted norms in \S\ref{sec:single-lift-norm},
we express both lifts using the common period map \(T^\circ\), which is independent of the twist.
In a theta exchange, one sums a kernel over the lattice classes on
a definite group, weighting each class by the reciprocal of its
automorphism-group order. For matched test functions this sum is
a scalar times an integral of the partner kernel. The local matching
and Rallis norm calculations give its squared modulus. In the three
exchanges below, the first partner contributes the constant function,
the second gives a fixed theta lift of that function, and the third
gives the period map. The scalar factors record all the variation
with the twist. Fix a lift number \(1\le j\le m\) and put \(e=m/2\).

For a set \(S\) of marked places, a character \(\alpha\), and an integer \(z\ge1\), use
\[
  \mathcal L_S(\alpha,z)=
  \left(\prod_{w\in S} q_w\right)^{z^2}
             \prod_{u=1}^z |L^S(u,\alpha)|^2 .
\]
Here \(q_w\) is the residue cardinality at \(w\). The partial Hecke value \(L^S(u,\alpha)\) omits the Euler factors at \(S\) and its conjugate \(\bar S\), with one factor \(q_w^{z^2}\) outside the \(L\)-product for each \(w\in S\). Let \(\upsilon\) be the Weil splitting character on \(E^\times\backslash\mathbb A_E^\times\). In products of characters on the norm-one torus, \(\upsilon\) denotes its restriction to that torus. For a tame weight-zero character \(c\), let \(\chi_c\) denote its finite residue character.\formalnote{10.1}{1}

\paragraph{Turning the first class sum into the constant function.}
Take the standard positive definite space \(I_m\) twice. The row set of the partner in the first exchange is empty. Its finite datum is \(1\) and its archimedean polynomial is \(1\), since it is an empty product of minors. The polynomial on the definite side specializes to the prescribed product of powers of the determinant. Twisting both splitting characters cancels the determinant character outside the kernel on the rank-zero partner. The integral on that space is therefore the constant \(1\) with the fixed measure normalization. For the character \(c^{-1}\) the first exchange supplies a scalar \(\gamma_1\), making the normalized class sum \(1\), with
\[
               |\gamma_1|^2=C_1/\mathcal L_S(c\upsilon^e,e),
\]
where \(C_1\) is independent of \(c\). Negating the form on \(I_m\) conjugates the first identity, introducing a constant of norm one and the value of \(\chi_c\) at a unit. Both disappear from the squared norm.\formalnote{10.2}{2}

\paragraph{Turning the second class sum into one fixed form.}
Let \(\phi^\circ\) be the theta lift of \(1\) on the negative definite space of rank \(m\) to the space of signature \((2m,m)\). Use the second exchange with character \(c\upsilon^m\) and the direct-sum decomposition in its kernel. For block variables the Schwartz function is a product, and summing the theta series over the direct-sum lattice factors it into the two component kernels. Substitute the conjugate first identity into the component on the negative definite space. Integration of the other component gives the fixed theta lift \(\phi^\circ\). Thus the weighted class sum of the second kernel is \(\phi^\circ/\gamma_{12}\). If \(\gamma_2\) is the second exchange scalar, then \(\gamma_{12}\) is \(\gamma_2\overline{\gamma_1}\) times a number of modulus one. The second character factor is \((c\upsilon^m)^{-1}\upsilon^{2m}=c^{-1}\upsilon^m\), which gives
\[
    |\gamma_{12}|^2=
       \frac{C_1C_2}{\mathcal L_S(c\upsilon^e,e)\,
                          \mathcal L_S(c^{-1}\upsilon^m,m)} .
\]
Thus the form \(\phi^\circ\) is the same for all the tame characters \(c\) at this weight. Its rapid decrease, holomorphy, and invariance at full finite level allow its use in the next exchange.\formalnote{10.3}{3}

\paragraph{The last averaging factor.}
Write \(W_j=E^{3m}\oplus W'\). Its partner at this exchange has \(3m+2z_*\) indices, where \(z_*=7m/2+k/2\). Recall from \S\ref{sec:spaces} that \(\pi_M\) generates the product of the chosen prime ideals. At the original level \((\pi_M^2)\), the finite test function in the integral against \(\phi^\circ\) is the sum of the translates of a residue-matrix additive character \(\psi_\Lambda\) by \(A=\mathrm{GL}_{3m}(\mathcal O_E/(\pi_M^2))\). Each translate gives the same integral: change the group variable by its integral lift, use invariance of the integration measure and the full level of \(\phi^\circ\). Thus the integral with this test function is \(|A|F_i(g)\), where \(F_i(g)\) is the integral of \(\phi^\circ\) against the theta kernel for datum \(i\) with the original finite test function \(\psi_\Lambda\). Replace \(\phi^\circ\) by the class sum just computed and apply the third matched exchange to each term. The resulting double class sum \(T_{i,\xi}(g)\) therefore satisfies
\[
             c_j(\xi)T_{i,\xi}(g)=F_i(g),\qquad
             c_j(\xi)=\gamma_3\gamma_{12}/|A|.
\]
The cardinality is the same for every twist \(\xi\), because the averaging group was taken at the original level. This equality applies to all choices of the row datum \(i\) and all \(g\) with the \emph{same} \(c_j(\xi)\).\formalnote{10.4}{4}

To identify the resulting product, let \(\chi\) be the fixed weight-zero character at the chosen places, \(\mu_j^0\) the untwisted lift character, and \(S\) the union of the chosen places and the places of \(\xi\). Put \(s_j=j+(t-11m-1)/2\), the exponent of \(\upsilon\) in the third critical pair. The three pairs \((\nu,z)\) are
\[
   (\upsilon^{-e}\chi^{-j},e),\qquad
   (\upsilon^{-m}\chi^j,m),\qquad
   (\mu_j^0\upsilon^{s_j},z_*).
\]
In these pairs the tame twisting powers are respectively \(-1,+1,+1\). In the first two scalar formulas set \(c=\chi^j\xi\). The factors there become \(\mathcal L_S((\nu\xi^{-1})^{-1},e)\) and \(\mathcal L_S((\nu\xi)^{-1},m)\) for the respective \(\nu\). The third matching contributes \(\mathcal L_S((\mu_j^0\upsilon^{s_j}\xi)^{-1},z_*)\). Calling the product of the three factors \(\mathcal L_j(\xi)\), we obtain
\[
              |c_j(\xi)|^2=\frac{C_j}{\mathcal L_j(\xi)},\qquad
      \frac{|c_j(\xi)|^2}{|c_j(1)|^2}=\frac{\mathcal L_j(1)}{\mathcal L_j(\xi)}
\]
whenever both twists are evaluated with their respective sets \(S\) and the nonzero constants of the construction. The fixed factor \(C_j=C_1C_2C_3/|A|^2\) cancels. This explains why a single untwisted family of functions \(F_i\) can test all the twisted square kernels, and why taking a quotient of determinants introduces exactly these Hecke-value ratios. Section~\ref{sec:one-exchange} derives the one-exchange identity and its normalization from the local matching and period calculations specified there.\formalnote{10.5}{5}


\section{The matched theta exchange and its normalization}
\label{sec:one-exchange}

In an orthogonal case, Snitz matches the two lifts of prescribed characters by an explicit map between the Schwartz spaces of quadratic spaces of dimensions three and one \cite[Theorem~2]{snitz-matching}. The unitary exchange used in Section~\ref{sec:three-exchanges} compares
two integrals on the same adelic group \(U(W)\). One integrates over a
positive definite space \(I_a\) of rank \(a\), the other over a space
\(V'\) of rank \(b\). Write the signature of \(W\) as \((m,l)\) and that
of \(V'\) as \((p,l)\). The dimension conditions are
\[
 n=\dim W=m+l=a+b,\qquad m=p+a,\qquad b=p+l,\qquad
 a=b+2d,\qquad n\ge 2l+2.
\]
In the following weight formulas, \(\upsilon\) is restricted to the norm-one torus.
Choose \(t_\tau\) frames \(\eta_{\tau i}\) of size \(p\times m\) at each
infinite place; their entries can be arbitrary. Let \(\alpha\) be tame
of weight \(l-t_\tau\), ramified exactly at the nonempty set \(S\) of
places of the modulus \(\mathfrak m\), with residue character \(\chi\).
The modulus is prime to its conjugate.

\subsection{Finite data and factorization on a direct sum}

Write \(R=\mathcal O_E/\mathfrak m\). Split the \(n\) indices into an
ordered block of size \(b\) and one of size \(a\). For \(k\in\mathrm{GL}_n(R)\)
let \(X_k\) be the last \(a\) rows of \(k^{-1}\), and \(Y_k\) the first
\(b\) columns of \(k\). Functions \(\Phi(X)\) and \(\Psi(Y)\) are
matched for \(\chi^{-1}\) if they vanish outside their respective sets
of such matrices and satisfy
\[
                 \Phi(X_k)=\chi(\det k)^{-1}\Psi(Y_k).
\]
The exchange uses the following pair attached to a function \(F\) on \(a\times b\) matrices.
Writing \(X=(X_1\ X_2)\), \(Y=\binom{Y_1}{Y_2}\), put
\[
 \Phi(X)=\chi(\det X_2)F(X_2^{-1}X_1),\qquad
 \Psi(Y)=\chi(\det Y_1)F(-Y_2Y_1^{-1}).
\]
Each function is defined to be zero if its indicated square block is
singular over \(R\), meaning that its determinant is a nonunit. On their
supports the matrices complete to invertible block triangular matrices.
For \(X=X_k,\,Y=Y_k\), the complementary-minor identity gives
\(\det Y_1=(\det X_2)\det k\), and \(XY=0\) gives
\(X_2^{-1}X_1=-Y_2Y_1^{-1}\) wherever the blocks are invertible. This
proves matching, including simultaneous vanishing on the complement.
If \(F\) is right \(\mathrm{GL}_b(R)\)-invariant, the last argument can
be replaced by \(Y_2\). Thus the column datum is the product
\(\chi(\det Y_1)F(Y_2)\).

For a corresponding orthogonal decomposition \(W=I_b\oplus W'\), take
the archimedean polynomial to depend only on \(W'\). The two factors of
the Schwartz function, the block diagonal group action, and the sum over
rational vectors all separate. The theta kernel on \(U(V')\) becomes
the product of the determinant kernel for \((V',I_b)\) and the kernel
for \((V',W')\) with datum \(F\). On the support, \(Y_1\) is invertible;
hence \(Y\) has a left inverse. The primitive-support convergence
estimate gives rapid decay and integrability. Consequently one may
multiply the exchange by the coefficients of another finite class
sum and integrate the sum term by term. For a subgroup \(H\) stabilizing
the chosen rows, cutting the row datum down to \(X_2\in H\) gives the
level-\(H\) version. The twisted translates of this restriction over
cosets partition the original datum according to \(X_2\), and their
kernels are translates by integral lifts. Summing first over the
fibers of the level map gives the original full-level class sum, with
the inverse automorphism weights. These observations account for the
product kernel and level changes in successive exchanges.\formalnote{10.6}{6}

\subsection{Conductor powers and critical values}

Set \(s_0=n/2-l=b+d-l\ge1\), and introduce the positive-weight
character \(\lambda=\alpha^{-1}\upsilon^{b+d}\), of weight
\(t_\tau+s_0\) at \(\tau\). The splitting character prescribed for \(W\)
in the general complementary-space theorem is then
\(\lambda^{-1}\upsilon^{n/2}=\alpha\). Put \(Q_S=\prod_{w\in S}q_w\) and
\[
 L_*(h)=Q_S^{h^2}\prod_{z=1}^h|L^S(z,\lambda)|^2,\qquad
 L_*(0)=1,\qquad L_*(-h)=L_*(h)^{-1}\quad(h>0).
\]
The partial \(L\)-function uses base change to \(E\) and omits \(S\)
and its conjugate.

The reduction in hyperbolic planes explains this product. Start with
standard determinant row and column functions at the ramified places.
For a large triple, unfolding along an isotropic frame expresses the
theta integral as Fourier terms in the center of the unipotent radical.
The terms use rational frames of prescribed Gram matrix. Terms whose
Gram matrix is excluded by the signatures disappear. On an orthogonal
sum the surviving term factors into a smaller theta functional and a
period of the equal-rank pair. This yields an identity for the two
integrals on the standard functions, measured by that period.
One then removes hyperbolic planes by taking the constant term at the
new isotropic line. The nonzero line orbit unfolds into a Tate integral.
On the large side its value at the step indexed by \(z\) has modulus,
up to a fixed normalization, \(Q_S^{z-1/2}|L^S(z,\lambda)|\).
The finite integral in the region \(z>1\) is found by expanding valuations
away from \(S\cup\bar S\) into the Euler factors; at the marked places
the determinant character cancels the unit character, and the change
of scale and the line Gauss sum give the norm power. Poisson summation
gives the opposite-side calculation, and continuation of the line
integral supplies the endpoint where the absolutely convergent
calculation does not apply directly. In squared norm the ratio of
successive factors is therefore
\[
       Q_S^{\,2z-1}|L^S(z,\lambda)|^2=L_*(z)/L_*(z-1).
\]
Summing the exponents \(1,3,\ldots,2d-1\) gives \(d^2\). Nonvanishing
of the positive-weight Hecke \(L\)-function on the real interval
\([1,\infty)\) makes each division legitimate. Sections~\ref{sec:plane-period} and~\ref{sec:hecke-endpoint} derive the two formulas for the single line orbit and prove the nonvanishing at one.\formalnote{10.7}{7}

It remains to determine the dependence of the equal-rank period on \(\lambda\).
Put \(k=s_0\), let \(H\) be a hyperbolic plane and take \(J=I_{2k}\) positive definite.
Then \(X=H^{\oplus l}\oplus J\) has the signature of \(W\).
For an archimedean Schwartz vector \(v\) on the pair \((Y,Y)\), where \(Y=H\oplus X\), set
\[
 r_\lambda(v)=|I_{Y,Y}(v)(1)|^2/L_*(k+l+1).
\]
At the ramified places the finite function on an invertible integral matrix \(y\) is the inverse of the local unit character of \(\lambda\) evaluated at \(\det y\). It is zero off the integral general linear group.
Remove the \(l+1\) planes above \(J\) in the exchange identity, first with complementary
partners \(J,0\) and then with the partners interchanged. Use the same \(v\) in both
identities and take the intermediate vector on \((J,J)\) from the first as datum for the second. This gives
\[
 B|c_1|^2=L_*(k)r_\lambda(v),\qquad D|c_2|^2=L_*(-k)r_\lambda(v).
\]
Here \(c_1,c_2\) are the exchange scalars. Integrating the resulting common kernel on
\([U(J)]\times[U(J)]\) in the two orders gives \(c_1c_2=\xi_v\) after cancellation of the
rank-zero volume. The exchange identities fix \(B,D,\xi_v\) independently of \(\lambda\).
Since \(L_*(k)L_*(-k)=1\), the nonnegative number \(r_\lambda(v)\) has constant square
\(BD|\xi_v|^2\).

Fix one admissible character. If its ratio vanished for every \(v\), the equal-rank
periods in the Fourier coefficients for \((H^{\oplus 2k}\oplus Y,J^-\oplus Y)\) would all vanish,
where \(J^-\) carries the negative of the form on \(J\). The theta integral for this pair would
therefore vanish on all standard data. By \S\ref{sec:det-cone}, vanishing on determinant
vectors descends through \(2k+l+1\) planes to \((J,0)\), where the vacuum integral is a
positive volume. Thus some \(v\) has positive ratio for the chosen character and hence for
every admissible character. Replacing \(Y\) by \(H\oplus W\) translates the argument of the
constant equal-rank theta integral and multiplies its value by a scalar of modulus one.
This gives the same ratio for the required space.

As a result, for any prescribed real Schwartz datum \(v'\) on
\((V',W)\) there is a real datum \(v\) on \((W,I_a)\), independent of
the tame character, so that on matched finite functions
\[
           I_{W,I_a}(v,\Phi)=c_0 I_{V',W}(v',\Psi),
           \qquad |c_0|^2=C_0L_*(d),\quad C_0>0 .
\]
The extension from standard functions uses equality of their induced
sections on complementary frames. A matched pair with locally constant
column function determines the same sections after multiplying by the
determinant character. The equivariant linear functional on row data
then factors through the section on column data. Reduction from the
product of the marked local integer rings to \(R\) preserves the
residue character and detects units, so functions matched modulo
\(\mathfrak m\) give these same local data.\formalnote{10.8}{8}

\subsection{Projection to the determinant vectors}

At \(\tau\), let \(X_\tau\) be the \(a\times m\) positive-column block
of the row variables. In the column variables write \(Y^+_\tau\)
for the \(m\times p\) block pairing positive signs and \(Y^-_\tau\)
for the \(l\times l\) block pairing negative signs. The required Fock
polynomials are
\[
 P_\eta=\prod_{\tau,i}\det\binom{\eta_{\tau i}}{X_\tau},
 \qquad
 Q_\eta=\prod_\tau\det(Y^-_\tau)^{t_\tau+n-2l}
             \prod_{\tau,i}\det(\eta_{\tau i}Y^+_\tau).
\]
The columns of the first determinant are ordered by \(m=p+a\).
Use \(v'\) equal to the Gaussian Fock vector of \(Q_{\eta^0}\) for
the standard frame \(\eta^0=(1_p\ 0)\) repeated in every factor.
For a fixed adelic point \(g\) the integrals above define linear
functionals on the two sets of real data. The first is invariant under
the compact \(U(a)\) and equivariant for translations by \(U(m,l)\)
on \(g\). Since the equality holds at every \(g\), it holds for
translates of the standard-frame data.

To isolate the desired component of \(v\), act in the positive \(m\)
block by the circle \(\operatorname{diag}(z^{-2}1_p,z^{-1}1_a)\),
and in \(U(a)\) by \(z^{-1}1_a\). Multiply by
\(z^{-a\sum_\tau t_\tau}\) and average over \(|z|=1\). On a Fock
polynomial the remaining circle weights on variables outside the last
\(a\) positive columns are strictly positive; variables inside that
square block have weight zero. Integration selects the constant term
in \(z\), which is exactly restriction to that block. Project under
the action of the compact \(U(a)\) with determinant character at each
\(\tau\). On the surviving \(a\times a\) matrix the resulting polynomial
transforms by \(\det^{\,t_\tau}\). This transformation law extends from
the unitary group to \(\mathrm{GL}_a(\mathbb C)\) by Zariski density.
Its value at the identity therefore determines the polynomial on
invertible matrices as a scalar multiple of \(\det^{\,t_\tau}\).
Polynomiality gives the same identity on all matrices. Density of the polynomial Fock vectors
and continuity of the mean extend this calculation to the vector
\(v\). The one-dimensional space of resulting functionals is closed.
Thus a scalar \(c'\), depending only on the fixed real data, converts
evaluation on \(P_{\eta^0}\) into this circle mean. The explicit
transformation of \(Q_{\eta^0}\) under the positive block and the
determinant twist give the same mean on the other side. This proves
the exchange for \(\eta^0\) and its unitary translates.

For general frames the two polynomials obey the same linear relations.
Indeed products \(\prod_{\tau,i}\det(\eta_{\tau i}Y_\tau)\) belong to the
representation of rectangular highest weight generated by the product
at the standard frame. Irreducibility gives one finite combination
of its unitary translates, valid for all \(Y_\tau\), that represents
the prescribed product. This immediately works for \(Q_\eta\).
For \(P_\eta\) specialize the entries of \(X_\tau\). There are then
a scalar \(d_\tau\) and an \(m\times p\) matrix \(Z_\tau\) with
\(\det\binom{\eta}{X_\tau}=d_\tau\det(\eta Z_\tau)\) for every
\(\eta\). If \(X_\tau\) has full rank, choose the columns of \(Z_\tau\)
as a basis of its null space and complete to a basis of \(\mathbb C^m\);
column operations give the formula with one determinant factor.
If its rank is smaller, every stacked determinant vanishes. The
same unitary linear combination therefore holds pointwise for
\(P_\eta\), hence polynomially. This proves the exchange for all the
frames with one scalar \(c'\).

On \(I_a\) the Fock vector just obtained is fixed by the real group
in the twisted Weil representation. The residue character of
\(\alpha\) cancels the change of \(\Phi\) under the integral finite
group. The resulting kernel is a function of the definite double
coset, and its integral is
\(\sum_{[h]}K(h)/|\operatorname{Aut}(h)|\) in the prescribed measure.
Introducing the free splitting character \(\beta\) multiplies the two
sides by the same factor \(\beta(\det g)^{-1}\), with the partner
characters becoming \(\alpha/\beta,\alpha\). If \(\Theta_{S,T}^{u,v}\)
denotes the kernel of the ordered pair with the displayed splitting
characters, the outcome is
\[
 \begin{aligned}
 \gamma_\alpha\sum_{[h]}
   \frac{\Theta_{W,I_a}^{\alpha,\beta}(\Phi,P_\eta;g,h)}
        {|\operatorname{Aut}(h)|}
    &=\int_{[U(V')]}
      \Theta_{V',W}^{\alpha/\beta,\alpha}(\Psi,Q_\eta;x,g)\,dx,\\
 |\gamma_\alpha|^2
    &=\frac{C}{Q_S^{d^2}\prod_{z=1}^d
                  |L^S(z,\alpha^{-1}\upsilon^{b+d})|^2},
 \qquad C=\frac{|c'|^2}{C_0}.
 \end{aligned}
\]
Indeed the projected equality has \(c'\) on the class-sum side and
\(c_0\) on the integral side, so \(\gamma_\alpha=c'/c_0\). Only the
nonzero \(c_0\) is divided out. The theorem permits \(C=0\); showing
that a particular resulting function is nonzero forces \(c'\ne0\)
in that application. In the first exchange of Section~\ref{sec:three-exchanges}
the partner integral is \(1\). In the second the required nonvanishing
is the nonvanishing of the fixed theta lift \(\phi^\circ\).\formalnote{10.9}{9}

\subsection{The sharp test vector in the lift norm}

The Petersson norm calculation uses a further test-vector identity.
For a lift from rank \(r\) to rank \(N>2r\), a sharp finite datum
\(\varphi_f\) for a level \(J\) is the vacuum away from the marked
places. At finite group elements integral away from the marked set,
its matrix coefficient is \(\|\varphi_f\|^2\) on \(J\cap K_0\) and zero
off the subgroup, with \(K_0\) the full integral level. Its mass is
\(\|\varphi_f\|^2/[K_0:J\cap K_0]\). For the real vector \(v\) the
Rallis doubling integral then factors into this mass, the quotient
volume, and the integral over the real group of the product of the
Weil coefficient \(\langle v,\omega(x)v\rangle\) with the coefficient
of the original form acted on by the sum of unramified vacuum Hecke
operators. For a joint-harmonic vector joining the specified compact
types, the real zeta integral
\(\int\langle v,\omega(x)v\rangle\overline{e(x)}/e(1)\,dx\)
is nonzero, where \(e\) is the source type function. This yields the
norm pairing. For injectivity, one may impose integrality at successively
more unramified places in the doubling integral. The limiting integral
is the sharp mass times the nonzero real factor times the squared norm
of the source form. This argument also applies in mixed signature: the source form
of the specified cuspidal type is rapidly decreasing, so doubling is absolutely
convergent for \(N>2r\). In the archimedean pairing its real translates contribute
the coefficient \(\overline{e(x)}/e(1)\), giving the nonzero integral above
(\S\ref{sec:real-pairing}). The finite-place test and limit are described in
\S\ref{sec:odd-lifts}.\formalnote{10.10}{10}

For the character lift \(U(1)\to U(2k)\), \(k\ge2\), each chosen place contributes \(1\) to the product of the local mean square, target level index, and missing Euler factor. Each twisting place of norm \(q\) contributes \(q^{2k-1}\). Thus for twisting modulus \(Q\) the conductor contributes \(N(Q)^{2k-1}\). Formula \textup{(SV3)} in \S\ref{sec:split-vacuum} gives the full squared norm of the class-sum lift.\formalnote{10.11}{11}


\section{Doubling the theta kernel in the convergent range}
\label{sec:doubling-integral}
To compute the squared Petersson norm of a theta lift from \(S\), one pairs its theta kernel with the complex conjugate kernel from a second copy of \(S\). Their product is a single theta kernel for \(S\oplus S^-\), where \(S^-\) has the negative Hermitian form. This is the doubling construction \cite[\S11.3]{gan-qiu-takeda-rallis}. Unfolding its integral gives the operator on the smaller group used for nonvanishing in \S\ref{sec:odd-lifts}. The Siegel--Weil formula identifies the convergent theta integral with an Eisenstein series \cite[no.~52, Theorem~5]{weil-siegel-formula}. We give the derivation in the normalization of this paper, including the strict inequality at a split prime.

\subsection{The mean and the sum over Lagrangians}
Let \(V\) be a split hermitian space consisting of \(r\) hyperbolic planes over the CM field \(E\), and let \(T\) be a hermitian space of rank \(m>2r\), with the unimodular adelic data of the construction. Put \(F=E^+\), \(H=U(T)\), and \([H]=H(F)\backslash H(\mathbb A_F)\). For a rational Lagrangian \(L\subset V\) let \(Z_L f\) be the zero term of the datum \(f\), namely its value at zero in the Schr\"odinger model for a complete splitting along \(L\). In the common model it is a normalized mean of Heisenberg translates over adelic vectors modulo rational ones. The desired formula is
\[
       \int_{[H]}\theta(\omega(1,h)f)\,dh
              =\operatorname{vol}[H]\sum_{L}Z_L f .                \tag{RI1}
\]
Here the splitting and the Weil representation in (RI1) are those of \(V,T\). We describe the estimates required to identify these two continuous functionals.

Fix a split frame. For \(x\in\mathbb N^r\), take the Levi torus element with diagonal entries \(2^{x_j}\) at the infinite places and identity at every finite place. Normalize its action by \(2^{-d m\sum_j x_j}\), where \(d=[F:\mathbb Q]\). The estimates on the rational Siegel set give a summable majorant for the norms of the zero terms after every such action. More explicitly, each \(L\) is assigned a vector \(t\in\mathbb N^r\). After fixed shifts have been absorbed into constants, a term has norm at most \(A(2^d)^{-m\sum t_j}\). The number with a fixed \(t\) is at most \(C(2^d)^{\sum_{i,j}(t_i+t_j)}\). In the reduction giving this count, one records the rational hermitian part of a unipotent parameter, whose \(ij\) entry has scale \(2^{t_i+t_j}\), and a rational group element in a fixed compact difference set; the latter has finitely many choices. Thus the sum is bounded by
\[
        AC\prod_{j=1}^r\sum_{u\ge0} 2^{-d(m-2r)u}<\infty .
\]
On the integral side the theta norm at each normalized torus translate is bounded by \(B(h)\|f\|\), where \(B\) is integrable on \([H]\); this is the separate reduction estimate for theta. These two bounds justify the sum and the integral defining the difference
\(\ell(f)=\operatorname{vol}[H]\sum_L Z_L f-\int_{[H]}\theta(\omega(1,h)f)\,dh\).\formalnote{10.12}{12}

Induct on \(r\). For \(r=0\) both sides give the unique zero term times the volume. Fix a column \(j\) of the split frame for \(r>0\), and write \(D\) for its line. The Lagrangians containing \(D\) correspond to those of \(D^\perp/D\). Their zero terms factor through the mean along \(D\). The induction hypothesis therefore identifies their sum times \(\operatorname{vol}[H]\) with the theta integral on \(D^\perp/D\) of that mean.

We compare this integral with the theta integral on \(V\), both at the normalized torus translate. In the mean along \(D\), conjugation by the torus divides the archimedean Heisenberg shifts by \(2^{x_j}\). To control the finite shifts, let the original data range over a compact set \(C\), and choose a compact set \(\Omega\) representing every class of the additive adelic quotient for the line mean. The largest change in a datum under translation by \(a\Omega\) tends to zero for each \(h\) as the adelic scalar \(a\) tends to zero. This change is bounded by \(2\sup_{f\in C}\|f\|\). Multiplying by the integrable \(B(h)\) and applying dominated convergence gives the same smallness after integration in \(h\). For any prescribed error choose an integer \(Q>0\) with \(Q\widehat{\mathcal O}_E\) inside the resulting finite neighbourhood of zero. Every quotient class has a representative \(Qb\), \(b\in\Omega\), because multiplication by the rational scalar \(Q\) is an automorphism of the quotient. Periodicity permits this change of representative in the mean. The scalar acting on \(b\) now has finite part \(Q\) and infinite part \(Q2^{-x_j}\), which lies in the required adelic neighbourhood for large \(x_j\). This proves convergence of the difference uniformly in \(f\in C\) and in all other dyadic coordinates.\formalnote{10.13}{13}

For each \(L\) outside the family containing \(D\), the normalized zero term tends to zero as \(x_j\to\infty\). The summable majorant allows summation of these limits. Thus \(\ell\) on a normalized translate tends to zero uniformly when any one dyadic coordinate tends to infinity.\formalnote{10.14}{14}

\subsection{The strict inequality that kills the difference}
Write \(\omega_g=\omega(g,1)\) and let \(v_0\) be the vacuum. Define \(M(g)=\sum_L|Z_L(\omega_g v_0)|\). For each \(f\) the function \(D_f(g)=|\ell(\omega_g f)|-M(g)\) is bounded above. In fact reduction puts \(g\) in a rational translate of a dyadic torus point times an element of a fixed compact set. Divide the compact family of translates of \(f\) by the absolute zero term of the corresponding vacuum along the standard Lagrangian. This denominator is continuous and strictly positive, so the rescaled family is still compact. Outside a finite box the normalized difference has absolute value at most one by the preceding limit. Undoing the normalization then bounds \(|\ell(\omega_g f)|\) by just the translated standard Lagrangian term in \(M(g)\). Inside the box the boundedness follows from compactness and continuity. Rational translations permute the zero terms.\formalnote{10.15}{15}

It suffices by density to test \(f\) that is a polynomial Fock vector times a finite translate of the vacuum. Choose a split prime where that translate is integral at both conjugate places. Let \(q\) be its residue cardinality, \(s=\sqrt q\), and \(k=m-2r>0\). At this prime the split Hecke identity at parameter \(-1\) for any \(H(\mathbb A)\)-invariant functional gives
\[
       B\,|\ell(\omega_g f)|
          =\left|\sum_U c_U\ell(\omega_{g g_U} f)\right|,
       \qquad B=\prod_{j=0}^{m-1}(1+q^j).
\]
Here \(U\) runs over the subspaces of the \(2r\)-dimensional residue space and \(g_U\) is its lattice modification. For a fixed \(L\) and \(g\), put \(u=\dim U\) and \(\eta=\upsilon_w(\varpi_w)^m\), where \(\upsilon_w\) is the splitting character. The local vacuum calculation gives \(\eta^u Z_L(\omega_{gg_U}v_0)=r_U v\), with \(r_U\ge0\) and one complex \(v\) independent of \(U\). The weight has the form \(c_U=c\,a_U\eta^u\), where \(a_U\ge0\) and \(c\) is fixed. Hence the weighted summands are all nonnegative real multiples of \(cv\). Their absolute values sum to the absolute value of their sum, which the Hecke identity for \(Z_L\) evaluates as \(B|Z_L(\omega_g v_0)|\).
Summing over \(L\) gives \(\sum_U|c_U|M(gg_U)=BM(g)\). Apply the triangle inequality to the Hecke identity for \(\ell\) and subtract this equality to obtain
\(B D_f(g)\le\sum_U|c_U|D_f(g g_U)\).\formalnote{10.16}{16}

The absolute weight of a subspace of dimension \(u\) is
\(\prod_{j<k}(1+q^j s^{2r})\,s^{ku}q^{\binom u2}\).
The finite-field binomial identity in \S\ref{sec:period-insertion} gives
\[
            \sum_U|c_U|
           =\prod_{j<k}(1+q^{r+j})\prod_{j<2r}(1+q^js^k)< B .
 \tag{RI2}
\]
For the strict sign, remove the \(k\) middle factors on both sides. Pair the two remaining factors with indices \(j,r+j\) on the left and \(j,r+k+j\) on the right for \(0\le j<r\). The difference between the right and left products for that pair is \(s^{2j}(s^k-1)(s^{2r+k}-1)>0\). There is at least one pair since \(r>0\). Now taking the finite upper bound \(\sup_g D_f(g)\) gives \(\sup D_f\le0\) by (RI2). The bound applies also to every scalar multiple of \(f\), with the same \(M(1)\). Evaluation at \(g=1\) forces \(\ell(f)=0\). Density and continuity prove (RI1). The analytic inputs remaining in this proof are Siegel-set reduction, the mixed-model theta and zero-term estimates used for domination and the boundary limit, nonvanishing of the translated vacuum zero term, alignment of the weighted local phases, and the local Hecke transfer.\formalnote{10.17}{17}

\subsection{Identification of the surviving Lagrangians}
Now take \(V=S\oplus S^-\), with \(S\) of rank \(r\), and tensor a datum \(f\) on \(\operatorname{Hom}(S,T)\) with the conjugate of \(f'\) on the negative copy. The negative-model factor has modulus one and is chosen so that its theta is the conjugate theta for \(f'\). The product rule for the tensor theta identifies the integrand on the left of (RI1) with the product of theta kernels at \(g,g'\in G=U(S)(\mathbb A_F)\). The two target character factors cancel under conjugation. Write the frame of a Lagrangian in the sum on the right as \(\binom{K}{K'}\), with \(r\) independent columns. Isotropy means \(K^*J K={K'}^*J K'\).

Suppose \(K'c=0\). When \(S\) is definite, \(Kc\) has zero norm and therefore vanishes. The other hypothesis allowed in the doubling theorem is that \(f=\operatorname{Sch}_A(\Psi)f_0\), where \(A\) is a set of finite places disjoint from its conjugate, \(\Psi\) is a Schwartz--Bruhat multiplier in the coordinates above \(A\), and at some split place \(w\) every \(y\) with \(\Psi(y)\ne0\) is an injective map \(y_w:S_w\to T_w\). If \(Kc\ne0\), the vector \((Kc,0)\) in the Lagrangian makes the zero-term functional invariant under its associated Heisenberg shifts. Moving these shifts through the \(g\)-action expresses \(Kc\) in the source coordinates at \(g\). Those coordinates differ from \(Kc\) by an invertible linear map at \(w\), so \(y_w\) sends the resulting vector to a nonzero vector. The change to the target dual frame is invertible as well. Nondegeneracy of the local additive pairing now gives a shift acting on the multiplier at \(y\) by a scalar different from one. By local constancy the scalar stays fixed on a neighbourhood of \(y\). Decompose the multiplier into indicators of finitely many compact open pieces in its support coordinates subordinate to these neighbourhoods. On each piece invariance equates the value of the functional to that scalar times its value, forcing it to vanish. Consequently \(K'c=0\) implies \(Kc=0\) for every nonzero term. Independence of the stacked columns gives \(c=0\).\formalnote{10.18}{18}

Consequently \(K'\) is invertible and the Lagrangian is exactly the graph \(\{(\delta y,-y)\}\) with \(\delta=-K{K'}^{-1}\) and \(\delta^*J\delta=J\). Conversely every rational isometry gives a graph, and its lower coordinate determines the upper map uniquely. In this graph splitting, the zero term is the quotient mean of the product of \(f'\) conjugated and \(f\) translated by \(\delta^{-1}\). Indeed the two blocks of the graph frame restrict the two data to the same variable via the isometry; the change to this variable is an invertible adelic linear map preserving the rational quotient and hence its normalized Haar measure. With the initial translates \(g,g'\) and reindexing \(\delta^{-1}\) by \(\delta\), the zero term becomes \(C(g'^{-1}\delta g)\), where \(C(u)=\langle f',\omega_{1,b}(u,1)f\rangle\) and \(b\) is the source-side splitting character in this convention. Thus the doubled integral is \(\operatorname{vol}[H]\sum_{\delta\in U(S)(F)} C(g'^{-1}\delta g)\).\formalnote{10.19}{19}

\subsection{Unfolding the Petersson product}
Integrate the identity against \(\overline{y'(g')}y(g)\) on the two copies of \([U(S)]\), where \(y,y'\) are rapidly decreasing automorphic forms. For fixed \(g'\) the rational translates \(\delta g\) of a fundamental domain tile \(G\). Put \(u=g'^{-1}\delta g\); Haar invariance and automorphy give \(y(g)=y(g'u)\). The integral over \(g'\) is now the source Petersson product \(\langle y',R(u)y\rangle\). This proves the Rallis identity in the normalization used for (OL2):
\[
 \langle\Theta_{f'} y',\Theta_f y\rangle_{[H]}
   =\operatorname{vol}[H]\int_G
              \langle f',\omega_{1,b}(u,1)f\rangle\,
              \langle y',R(u)y\rangle_{[U(S)]}\,du .             \tag{RI3}
\]
For the exchanges of integrals, the absolute theta integral is bounded by the integrable majorant for the doubled datum; the sum of absolute graph coefficients is bounded by the sum of zero-term operator norms. The datum norm is bounded by a constant times the product of the two translated norms. Reduction by rational elements bounds each translated norm by a fixed power of the source height. The absolute values of \(y,y'\) absorb these powers, and the product quotient has finite measure. Tonelli's theorem followed by Fubini's theorem therefore justifies both the unfolding and the absolute convergence of (RI3). In the odd lifts the spaces are definite and \(m=2r+1\), which meets the strict range by exactly one. Section~\ref{sec:odd-lifts} applies this identity in odd rank. For the determinant character lift see \S\ref{sec:real-pairing} for the real projection and \S\ref{sec:split-vacuum} for the finite Euler product.\formalnote{10.20}{20}


\section{The real pairing in doubling and in the lift from a circle}\label{sec:real-pairing}
The archimedean factor of (OL2) in \S\ref{sec:odd-lifts} is a projection of the real Weil vector onto a specified type. For the character lift \(U(1)\to U(N)\) in \S\ref{sec:one-exchange}, the same factor is an explicit positive integer. We compute the projection, its norm and the additional dimension factor introduced by passage to algebraic forms. All inner products below are conjugate-linear in the first variable. The real groups carry compact-dual Haar measures, hence each compact unitary group and each circle has mass one.

\subsection{Compression of a real translate to the source type}
At one infinite place let \(G=U(J)\), \(K=G\cap U(r)\), and let \(l=2\lambda\) be a discrete-series parameter of rank \(r\ge1\). Write \(c=c_l\) for the unit lowest-type coefficient and \(d=(\int_G|c|^2)^{-1}\). Then \(c(1)=1\), \(c(g^{-1})=\overline{c(g)}\), and the type kernel is \(e(g)=d\,\overline{c(g)}\). In the noncompact argument take \(l_i\) distinct of parity \(r-1\) and \(|\lambda_i-\lambda_j|\ge500r^3+A\) for opposite signs. The nonnegative integer \(A\) bounds the degree of operator-norm growth in whichever representation is to be integrated. For the oscillator model below \(A\) can be the number of matrix coordinates. With all signs equal the separation across opposite signs is vacuous.

Write \(m(k)\) for the unit highest-vector coefficient of the lowest \(K\)-type \(\Lambda=\lambda+\rho-2\rho_c\), and \(a=\int_K|m|^2>0\). The lowest-type projection gives
\[
      \int_K\overline{m(k)}\,c(xkb)\,dk=a\,c(x)c(b).
\]
In representation terms, the average against \(\bar m\) is \(a\) times projection onto the highest line of that type. The lowest type occurs once. The coefficient on the left inserts this projection between \(x\) and \(b\); both ends of the projected coefficient use the same unit vector.

This compact projection gives the corresponding Schur integral on \(G\). Put \(w(g)=c(g^{-1})\). The displayed identity with \(x=g^{-1}, b=1\) gives
\(\int_K\overline{m(k)}w(k^{-1}g)\,dk=a w(g)\).
Multiply by \(c(xgb)\), integrate, and in the inner integral put \(g=kh\). Applying the display again, now to \(c(xkhb)\), gives
\[
 \int_G w(g)c(xgb)\,dg
      =c(x)\int_G w(g)c(gb)\,dg .
\]
The exchanges use \(w\in L^1\) and \(|c|\le1\), which follow from coefficient decay. In the last integral put \(g=h^{-1}b^{-1}\). Invariance under translation and inversion gives
\[
 \int_G c(g^{-1})c(gb)\,dg
   =\int_G c(h^{-1})c(bh)\,dh .
\]
Applying the preceding factorization with \(x=b\) and last argument \(1\) gives \(c(b)/d\). We have proved\formalnote{10.21}{21}
\[
       d\int_G c(xgb)\overline{c(g)}\,dg=c(x)c(b).             \tag{RP1}
\]

Now let \(R(x)y(g)=y(gx)\), with \(x\) in the real factor, and put \(P_e=\int_G e(u)R(u)\,du\). On a right unitary action this is self-adjoint since \(e(u^{-1})=\overline{e(u)}\). The kernel of \(P_eR(x)P_e\), at a point \(h\) in the convolution variable, is
\[
 \int_G e(u)e(x^{-1}u^{-1}h)\,du
   =d^2\int_G\overline{c(u)}c(h^{-1}ux)\,du
   =c(x)e(h)
\]
by (RP1). In particular \(P_e^2=P_e\). The same computation for a product of real groups factors place by place. If \(y,z\) are automorphic forms of type \(e\), then \(P_ey=y,\ P_ez=z\), so
\[
       \langle y,R(x)z\rangle
         =\langle y,P_eR(x)P_ez\rangle
         =\frac{\overline{e(x)}}{e(1)}\langle y,z\rangle .      \tag{RP2}
\]
The Petersson pairing uses a finite-volume quotient. The bounded forms and the integrable type kernels justify these convolutions by Fubini. In the sharp-data doubling formula \(z=\mathcal R y\), with \(\mathcal R\) the finite vacuum operator, belongs to the same real type. Substitution of (RP2) leaves exactly
\[
    Z_\infty=\int_{G_\infty}
        \langle v,\Omega(x,1)v\rangle\,
        \frac{\overline{e(x)}}{e(1)}\,dx
\]
times \(\langle y',\mathcal R y\rangle\). The target quotient volume and the sharp finite mass multiply this expression as in (OL2).\formalnote{10.22}{22}

\subsection{Positivity of the joint harmonic pairing}
Let the source parameter at place \(\tau\) be \(l_\tau\). In the source oscillator representation \(\pi_\tau\), the harmonic \(v_\tau\) has the lowest-type highest weight for \(L_\tau=-l_\tau\). This sign follows from the source acting on the columns of the space of maps. The kernel identity \(\overline{e_{l_\tau}}=e_{L_\tau}\) and factorization of the pure tensor \(v=\bigotimes_\tau v_\tau\) therefore give
\[
 \begin{aligned}
 Z_\infty&=
   \frac{\prod_\tau\langle v_\tau,
                   \pi_\tau(e_{L_\tau})v_\tau\rangle}
        {\prod_\tau d_{l_\tau}},\\
 \pi_\tau(e_L)v_\tau&=\int_{G_\tau}e_L(g)\pi_\tau(g)v_\tau\,dg .
 \end{aligned}                                                   \tag{RP3}
\]
The matching hypotheses give nonzero polynomials \(v_\tau\) that are highest and termwise harmonic for both groups. They also give the large gap needed to integrate \(e_L(g)\|g\|^A\). The tensor product of these integrable bounds permits the factorization. For the definite groups in the two odd lifts this already finishes the positive-pairing calculation. Indeed \(G_\tau=K_\tau\) and \(c_L|_K=m\). Averaging \(\pi_\tau(k)v_\tau\) against \(\overline{m(k)}\) gives \(a v_\tau\), also when there are several copies of the type, since \(v_\tau\) lies in their highest lines. Here \(d_L=1/a\), so the numerator factor in (RP3) is \(\|v_\tau\|^2>0\).\formalnote{10.23}{23}

The source also proves positivity for mixed signature, where \(\pi(e_L)v\) can be a proper projection of \(v\). We detail this argument to specify what enters the real pairing. Work at one place and write \(L=2\lambda_0,\ \mu=\lambda_0+\rho,\ \Lambda=\lambda_0+\rho-2\rho_c\) for this parameter. For an integer weight \(\alpha\) write \(m_\alpha(Y)\) for its product of principal weight minors, as in \S\ref{sec:real-decay}, so \(m=m_\Lambda\). The polar expression for the coefficient at \(k_1 z k_2^{-1}\), where \(z\) is positive definite in \(G\), is
\[
 c_L(k_1 z k_2^{-1})
  =\frac1a\int_K
     \frac{m(k_1 u)\,\overline{m(k_2u)}}{m_{-\mu}(u^*zu)}\,du .
\]
Average the integrand defining \(\langle v,\pi(e_L)v\rangle\) over \(g\mapsto k_1^{-1}g k_2\). This preserves its integral. For \(g=\kappa z\) in the polar decomposition, insert the polar expression for \(\bar c_L\). At fixed \(u\), translate \((k_1,k_2)\) by \((\kappa u,u)\) on the left. The two numerator factors become \(m(k_1)\overline{m(k_2)}\), and the operator inside the pairing becomes \(\pi(k_1^{-1}u^{-1}zu k_2)\). Integration on each side projects the highest vector with factor \(a\), cancelling one power against the polar denominator. Thus the compact average of the original integrand is
\[
       a\,d_L\int_K \frac{\langle v,\pi(u^{-1}zu)v\rangle}
                                  {m_{-\mu}(u^*zu)}\,du .       \tag{RP4}
\]
The denominator is positive, since all principal minors and the determinant of a positive-definite matrix are positive.\formalnote{10.24}{24}

It remains to see why the numerator of (RP4) is positive. Write \(P\) for the nonzero harmonic polynomial of \(v\). The Fock inner product on polynomial vectors is
\[
       (P,Q)_{\mathcal F}=(\overline P(\partial)Q)(0),
       \qquad (X^\alpha,X^\beta)_{\mathcal F}
                     ={\bf1}_{\alpha=\beta}\prod_j\alpha_j! .
\]
Here the bar conjugates coefficients. Multiplication by \(X_j\) is adjoint to \(\partial_j\) by the oscillator raising and lowering identities, and the normalized vacuum paired with \(Q\) gives \(Q(0)\). Iterating proves both assertions.

For \(z>0\) in \(G\), its linear action \(U\) on the matrix coordinates is positive definite and preserves the spectator labels, namely the columns for a row action or the rows for a column action. In the two conventions this action is \(z\otimes I\) or \(I\otimes(z^{-1})^{\mathsf T}\). Its determinant twist is one since \(\det z>0\) and \(|\det z|=1\). Denote the coordinate-sign projections by \(P_\pm\), put
\[
        C=P_+UP_++P_-\overline U P_-,\qquad
        b b^*=C^{-1},
\]
and write \(S_M\) for the substitution \(X_j\mapsto\sum_i M_{ij}X_i\). Both diagonal blocks of \(C\) are positive definite, so \(b\) is invertible. For harmonic polynomials \(Q,P\), the oscillator coefficient equals
\[
   (Q,\omega(U)P)_{\mathcal F}
      =(\det U_{--})^{-1}(S_{C^{-1}}Q,P)_{\mathcal F}.          \tag{RP5}
\]
Indeed termwise harmonicity means that the mixed derivatives in opposite signs of the same label vanish. Commuting a creation operator through \(\omega(U)\) and applying the annihilation operators gives, for the coefficient with \(X_jQ\) on the left, the sum of the coefficients with \(Q,\partial_i P\) multiplied by \(\overline{(C^{-1})_{ij}}\). In this step a monomial \(X_jQ\) being harmonic forces each prohibited opposite-sign derivative of \(Q\) to vanish. The formula on the right of (RP5) has exactly the same recursion by the Fischer adjunction. The initial value with \(Q=1\) is \((\det U_{--})^{-1}P(0)\), obtained from the vacuum overlap and the same commuting relations. Monomial induction proves (RP5). For \(Q=P\), moving the substitution \(S_b\) to the other side gives \((S_{b^*}P,S_{b^*}P)_{\mathcal F}\) in its numerator. Invertibility of \(b\) and \(P\ne0\) make this strictly positive, and \(\det U_{--}>0\). Now (RP4) is positive at every \(g\). Its integral is finite by the weighted kernel bound, and is strictly positive. This proves positivity of every numerator in (RP3).\formalnote{10.25}{25}

\subsection{Evaluation for the character lift and the dimension multiplier}
For \(U(1)\to U(N)\) both groups are definite and the source automorphic form is \(1\). At each infinite place \(\tau\) the weight diagram has one row of length \(r_\tau\), possibly zero. Its highest polynomial in the coordinates of the one-column matrix is \(Z_{\tau,i_\tau}^{r_\tau}\) with \(i_\tau\) the first row in parameter order. Indeed each of the \(r_\tau\) columns of the diagram supplies the same one-by-one highest minor. The polynomial used in the lift is the product of these polynomials, with the variables renamed to the matrix entries of \(\operatorname{Hom}(W,V)\). To verify its normalization in the Weyl module coordinates, the highest vector \(\eta\) is the evaluation vector on the permutation array determined by the order of the parameter. Its Gram norm is one. Its polynomial is the evaluation of the Weyl kernel at that array, giving the same product of first-row minors, with coefficient one.\formalnote{10.26}{26}

The source splitting exponent in the notation \(\Omega_{\gamma,\gamma'}\) is \(\gamma'_\tau=-r_\tau\), with \(\gamma_\tau\) the target determinant exponent. On a polynomial homogeneous of degree \(r_\tau\) in that column, the Fock substitution at \(x_\tau\in U(1)\) contributes \(x_\tau^{-r_\tau}\), and the determinant twist \((\det x_\tau)^{-\gamma'_\tau}=x_\tau^{r_\tau}\) cancels it. Thus \(\Omega_{\gamma,\gamma'}(x,1)v=v\). The Fischer pairing gives its squared norm \(H=\prod_\tau r_\tau!\), equal to the hook number. Integrating over circles of mass one, while retaining the Petersson product of the constant functions on the adelic source, gives
\[
 \int_{U(1)_\infty}\langle v,\Omega(x,1)v\rangle
                       \langle1,1\rangle_{[U(1)]}\,dx
          =\operatorname{vol}[U(1)]\prod_\tau r_\tau! .           \tag{RP6}
\]
In the Rallis norm formula the finite series acting on \(1\) is the constant vacuum sum, so it can be taken outside this integral. The convergent-range hypothesis here is \(N>2\), and in the special-value application \(N=2k,\ k>1\).\formalnote{10.27}{27}

For clarity, the additional dimension in the squared norm in \S\ref{sec:one-exchange} comes from a second compact average on the \emph{target}. Let \(\rho\) be the target compact Weyl representation with its unitary Gram form, \(D_\rho=\dim\rho\), and \(v^+\) the Gram-unit highest vector. The function attached to an algebraic form \(\phi\) at finite point \(g\) is \(F_\phi(x,g)=\langle v^+,\rho(x^{-1})\phi(g)\rangle\). Averaging the conjugates of \(v^+\langle v^+,\,\cdot\,\rangle\) over the compact group gives an intertwiner. Irreducibility and trace one give the scalar \(1/D_\rho\). Hence
\[
 \int_{U(N)_\infty}\overline{F_\phi(x,g)}F_{\phi'}(x,g)\,dx
      =D_\rho^{-1}\langle\phi(g),\phi'(g)\rangle_{\rm Gram}.
\]
Let \(J_f\subset K_0\) be the target level, with \(\operatorname{vol}(K_0)=1\), and let the two forms have the same finite-order level character. Their Gram pairing is then constant on \(J_f\)-cosets. Each rational double coset represented by \(g\) contributes \(\operatorname{vol}(J_f)/|\Gamma\cap gJ_fg^{-1}|\) times its compact average to the Petersson integral, where \(\Gamma\) is the diagonal image of the rational unitary group in its finite adelic points. The stabilizer \(\Gamma\cap gJ_fg^{-1}\) is finite. The algebraic class sum uses weight \(1/|\Gamma\cap gJ_fg^{-1}|\). Thus it equals \([K_0:J_f]D_\rho\langle F_\phi,F_{\phi'}\rangle_{\rm Pet}\). For these one-row weights \(D_\rho=\prod_\tau\binom{N+r_\tau-1}{r_\tau}\), since at each place the degree-\(r_\tau\) monomials in \(N\) variables form a basis of the symmetric power. This accounts for both archimedean factors \(H D_\rho\) in the algebraic norm.\formalnote{10.28}{28}

The local inputs used here are the polar and lowest-\(K\)-type identities for the normalized discrete-series coefficient, its weighted integrability, the existence of the matched joint-harmonic polynomials, and the oscillator commutation and vacuum formulas. The global factorization uses the sharp-data Rallis identity of \S\ref{sec:doubling-integral} and \S\ref{sec:odd-lifts}. Formula (RP6) requires only the compact Fock calculation once the specified vector and splitting exponents have been identified.


\section{The finite vacuum sum for the rank-one Rallis norm}
\label{sec:split-vacuum}

Let \(F=E^+\), let the source be the definite line and the target have rank \(N=2k\), \(k\ge2\). The source form is the constant \(1\). Its arithmetic character enters the finite Weil representation through
\[
      b=\nu\,(\upsilon|_{U(1)})^k,\qquad
      \widetilde\nu(x)=\nu(x/\bar x).
\]
Here \(\nu\) is the unitary character of the adelic norm-one torus in the lift data. For a critical pair and twist it is the product of the pair's character and the tame twist. Let \(S=C\sqcup Q\) consist of the chosen places \(C\) and the designated twisting places \(Q\). They are distinct split places of \(E\), and \(S\) is disjoint from its conjugate. Put \(D=S\cup\bar S\) and let \(S_F\) be the places of \(F\) below \(S\). The tame part of \(\nu\) is trivial off \(S\). We write \(L_E^D\) for the finite Hecke \(L\)-function with all factors in \(D\) removed.\formalnote{10.29}{29}

Section~\ref{sec:doubling-integral} expresses the norm of the lift by an integral of matrix coefficients. Section~\ref{sec:finite-period} computed the support and mass of the marked test vector. We evaluate the complementary finite sum and obtain
\[
     \Sigma=\frac{L_E^D(k,\widetilde\nu^{-1})}
                    {\zeta_F^{S_F}(2k)}
       =\frac{L_E^D(k,\widetilde\nu^{-1})}{\zeta_F(2k)}
                    \prod_{w\in S}(1-q_w^{-2k})^{-1}.       \tag{SV1}
\]
The deleted set includes the auxiliary places even when a residue character there happens to be trivial. There is one correction per conjugate pair.

\subsection{The finite integral as a sum}

Write \(K_A\) for the finite points of \(U(1)\) that are integral at every place of \(A=D^c\), and \(K_0\) for the integral points at all places. Each coset modulo \(K_A\) has a representative \(u\) equal to \(1\) on \(D\). The vacuum coefficient in rank \(N\) is
\[
                c_N(u)=\bigl(b(\det u)\,m_f(u)^N\bigr)^{-1},
\]
where \(m_f\) is the finite Weil multiplier. On the translate \(u K_A\), the coefficient of the sharp vector at \(ug\) is \(c_N(u)\|\Psi\|^2\) for \(g\in J\cap K_0\) and is zero elsewhere. This follows from disjoint support of the vacuum change at \(A\) and the marked datum at \(D\). The character is trivial on integral points supported in \(A\), so the expression is independent of the supported representative. Normalizing \(K_0\) to have measure one, the integral over this coset is \(c_N(u)\|\Psi\|^2/[K_0:J\cap K_0]\). The source function \(1\) has the same value after every right translation. Thus the finite operator contributes the marked mass times \(\Sigma=\sum_u c_N(u)\).\formalnote{10.30}{30}

\subsection{One split place and one nonsplit place}

First take a split place \(v\notin S_F\), with places \(w,\bar w\) above it in \(E\) and common residue cardinality \(q\). Evaluation at \(w\) identifies the local norm-one torus with \(E_w^\times\), its integral subgroup with \(\mathcal O_w^\times\), and the quotient with \(\mathbb Z\). Represent \(z\in\mathbb Z\) by the component \(\varpi_w^z\) at \(w\) and its conjugate inverse at \(\bar w\). Put \(z_+=\max(z,0)\), \(z_-=\max(-z,0)\), \(a=\widetilde\nu(\varpi_w)\), and \(\alpha=\upsilon(\varpi_w)\), using a local idele in each character value. Both have absolute value one. As \(\widetilde\nu\) is unramified at both places, its value at any uniformizer of \(\bar w\) is \(a^{-1}\).\formalnote{10.31}{31}

The intersection \(\mathcal O_w\cap\varpi_w^{-z}\mathcal O_w\) equals \(\varpi_w^{z_-}\mathcal O_w\). Its index supplies the square root of the norm in the multiplier, and its generator supplies the splitting character. The one-sided multiplier is therefore \((\alpha\sqrt q)^{z_-}\). At the other place the Weil conjugation rule gives the complex conjugate of \((\alpha\sqrt q)^{z_+}\). Multiplication on disjoint supports gives
\[
          m_f(u_z)=
          (\alpha\sqrt q)^{z_-}
          (\alpha^{-1}\sqrt q)^{z_+}
                    =\alpha^{-z}q^{|z|/2}.
\]
On a norm-one element \(x/\bar x\) the splitting character takes the value \(\upsilon(x)^2\). Hence \(b(\det u_z)=(a\alpha^{2k})^z\), and the two powers of \(\alpha\) cancel exactly in the coefficient. With \(x=q^{-k}\) we obtain
\[
 c_{2k}(u_z)=a^{-z}x^{|z|},\qquad
 \sum_{z\in\mathbb Z}c_{2k}(u_z)
      =\frac{1}{1-a^{-1}x}+\frac{ax}{1-ax}
      =\frac{1-x^2}{(1-a^{-1}x)(1-ax)}.                   \tag{SV2}
\]
Exactly one of the two geometric series includes \(z=0\). The denominator consists of the two unramified Hecke factors in inverse form. The numerator is the inverse local factor of \(\zeta_F(2k)\).\formalnote{10.32}{32}

For \(v\notin S_F\) with a single \(w=\bar w\), every element \(g\in U(1)(F_v)\) satisfies \({\rm ord}_w(g)=-{\rm ord}_w(g)\). The integral subgroup is the whole local group and there is one coset of coefficient \(1\). Moreover \(t/\bar t\) is a local unit for every \(t\in E_w^\times\). Since the tame part of \(\nu\) is trivial there, \(\widetilde\nu\) is trivial on all of \(E_w^\times\). In the inert case, if \(q_v\) is the residue cardinality of \(F_v\), then \(q_w=q_v^2\). Consequently the numerator \(L\)-factor \((1-q_w^{-k})^{-1}\) cancels the local \(\zeta_F(2k)\)-factor. This proves (SV1) also at inert places.\formalnote{10.33}{33}

The residue-degree assertion uses an actual property of \(E\). The Weil data require a primitive \(5040\)-th root of unity in \(E\). In particular \(\mathcal O_E\) contains \(i\) and \(\omega\) with \(i^2=-1,\ \bar i=-i,\ \omega+\bar\omega=1\). Then \(y=i\omega\) satisfies \(y-\bar y=i\), a unit at every finite place. At a ramified finite place of the quadratic extension \(E/F\), conjugation would act trivially on the residue field and \(y-\bar y\) would vanish there. Thus every nonsplit finite place here is inert. In a general ramified quadratic extension the torus and local base-change arguments would still give \(1\) under the same character hypothesis, but \(q_w=q_v\) would change the identification of the denominator with \(\zeta_F(2k)\).\formalnote{10.34}{34}

\subsection{Absolute convergence and the omitted factors}

An adelic coset amounts to finitely many nonzero integer valuations at split \(v\notin S_F\), with the representative above at each such \(v\). Every such finite choice occurs by multiplying the supported local ideles. The absolute local sum is \((1+x)/(1-x)\) at a split \(v\) and is \(1\) at an inert \(v\). The product of the absolute local sums is finite: for \(k\ge2\), the logarithm of \((1+x)/(1-x)\) is \(O(x)\) uniformly for \(0\le x\le 1/4\), and
\[
   \sum_{v\ {\rm finite}}q_v^{-k}
         \le [F:\mathbb Q]\sum_{\ell\ {\rm prime}}\ell^{-k}<\infty .
\]
Thus finite products exhaust the cosets in an absolutely convergent sum. This argument also permits reindexing the Hecke factors by \(w\mapsto\bar w\), giving the same partial \(L\)-value for \(\widetilde\nu\) and \(\widetilde\nu^{-1}\). At a split pair the resulting factor in (SV2) is \((1-x^2)/|1-ax|^2>0\), and its logarithm is \(O(x)\), so the full finite scalar is positive.\formalnote{10.35}{35}

One can also write the product over places of \(E\). Give the two places of a split pair opposite signs \(\epsilon_w\) and put \(\epsilon_w=1\) at a fixed place. For \(w\notin D\) its factor is
\[
    \frac{1-\epsilon_w q_w^{-k}}
         {1-\widetilde\nu(\varpi_w)^{-1}q_w^{-k}} .
\]
The choice of which place gets the minus sign disappears in the product. Omitting \(D\) removes both signs at each marked pair, replacing the full denominator \(\zeta_F(2k)\) by \(\zeta_F(2k)\prod_{w\in S}(1-q_w^{-2k})\). The Hecke factors at \(D\) stay omitted even if the character is unramified there. This proves (SV1) with the imprimitive \(L\)-value that occurs in Section~\ref{sec:critical-values}.\formalnote{10.36}{36}

\subsection{Putting the finite and real calculations together}

Let \(I_C=\prod_{w\in C}(q_w-1)\), \(N(Q)=\prod_{w\in Q}q_w\). The mean square of the marked vector was computed in Section~\ref{sec:finite-period}. The target integral-level index counts nonzero residue columns at a chosen place and residue hyperplanes at a twisting place. Multiplying mean square, target index, and correction in (SV1) gives, respectively,
\[
\begin{aligned}
   q^{-N}(q^N-1)(1-q^{-N})^{-1}&=1,\\
   \frac{q-1}{q}\frac{q^N-1}{q-1}(1-q^{-N})^{-1}&=q^{N-1}.
\end{aligned}
\]
The marked mass divides the mean square by the source index \(I_C\). The class-sum lift \(\Phi_X\) is \(I_C\) times the theta integral of the constant \(1\), since each source level class in the quotient integral carries the inverse level index. Squaring leaves one factor \(I_C\) after the marked mass has been included. For this specific definite line the real Weil action fixes the chosen Fock vector, whose squared norm is its hook number \(h_\lambda\). Its integral with the source constant contributes \(h_\lambda\,{\rm vol}[U(1)]\). Passing from the target Petersson norm to the algebraic norm contributes the target level index above and the dimension \(d_\lambda\) of its weight representation. With the volume normalization of Section~\ref{sec:doubling-integral}, we have proved
\[
  \|\Phi_X\|^2=
      \frac{{\rm vol}[U(1)]\,{\rm vol}[U(N)]}{\zeta_F(N)}
       I_C\,h_\lambda d_\lambda\,
       L_E^D(k,\widetilde\nu^{-1})\,N(Q)^{N-1}.              \tag{SV3}
\]
Section~\ref{sec:real-pairing} describes the real pairing in the oscillator representation. The critical pair's integral normalization further multiplies \(\Phi_X\) by the specified vacuum and derivative scales; the squared norm therefore acquires their squared absolute value. Formula (SV3) is for the class-sum lift before this multiplication. It uses the definite block and sharp mass, the convergent doubling identity and the real pairing. The split-place calculation above uses the finite Weil action on a lattice vacuum. The hypothesis \(k>1\) explains where absolute convergence enters this calculation; the endpoint at \(k=1\) in Section~\ref{sec:finite-period} has a different proof.\formalnote{10.37}{37}


\section{Removing a plane by Poisson summation on one line}
\label{sec:plane-period}

Let \(W'=H\perp W,\ V'=H\perp V\) be the two adapted splittings used in a single exchange, over a CM extension \(E/F\). Put \(n=\dim_E W\), \(r=\dim_E V\), and \(e=r-n/2\), with \(n\) even. The constant term on \(U(W')\) is taken along the added line. In the unitary tower property for a cuspidal input, taking a constant term on the target and restricting to the smaller group gives its theta lift with the added Schwartz variable set to zero \cite[Proposition~3.1, (3.2)]{wu-unitary-tower}. Here the source of the lift also contains an added plane. The large theta integral integrates over \([U(V')]\), where brackets mean quotient by rational points. We follow both integrations. They reduce to the theta integral on \((W,V)\) times one incomplete Hecke value. The exponent at which that value must be evaluated is \(e+1\) for \(e\ge0\) and \(-e\) for \(e<0\).

Fix the row labels for the standard determinant datum: a common block of size \(b\) together with the \(r\) variable rows labels the \(n\) small columns; with the \(r+2\) large rows it labels the \(n+2\) large columns. Thus \(n=b+r\). The split tame character has a nonempty set \(S\) of marked places disjoint from \(\bar S\), and its residue character is nontrivial at each of them. These conditions will eliminate the zero terms.\formalnote{10.38}{38}

\subsection{The rational orbit and its measure}
Write a theta sum in source-line coordinates, with rational column \(t\in V'(E)\). On this term the central coordinate \(b\) of the source radical acts by the standard additive character of \(\operatorname{tr}(b\,t^*J_{V'}t)\). The average of a nontrivial character of a compact group is zero: translation by a point where the character differs from one both preserves the average and multiplies it by that value. Thus only columns with \(t^*J_{V'}t=0\) survive. The remaining radical coordinates still have to be averaged.

The column \(t=0\) also disappears, for a different reason. At a marked place the two row blocks in the standard finite datum must stack to a matrix having an integral left inverse. In the zero slice the variable block kills the column \(K\) of the adapted frame. The fixed block \(yI^\vee\) kills \(K\) since \(I^\vee K=0\). A left inverse would then imply \(K=0\), contradicting \(K^\vee K=1\). This proves vanishing before integrating the smaller theta kernel.\formalnote{10.39}{39}

Group the remaining columns by their rational spans. Every isotropic line is rationally isometric to the fixed line \(E k_V\): scale generators to identify the two null lines and extend the isometry by Witt's theorem. The orbit sum thus unfolds \([U(V')]\) to the quotient by the rational stabilizer \(P(F)\) of \(E k_V\). Over this one line sum over \(t=\xi k_V,\ \xi\in E^\times\). The remaining radical average in the source is the zero term in the target complement. Indeed its translation directions there are the rows of \(\xi k_V y^*J_W\); the change from \(y\) to this full set of coordinates is invertible rationally on the parameter space, so preserves its quotient probability measure. This leaves the theta of a line slice on \(\operatorname{Hom}(W,V)\).\formalnote{10.40}{40}

For the unfolded measure write a point in radical, Levi and compact coordinates. The target radical has coordinates \((C,b)\) with \(C\in\mathbb A_E^r\) and \(b\) a self-adjoint scalar. Its group law translates \(C\) and translates \(b\) by a quantity depending on \(C\) and the left factor. Integrating in \(b\) first proves left invariance of the product additive measure. Conjugation by the Levi point \(m(a,h)\) sends these coordinates to
\[
        (hC\bar a^{-1},\ (a\bar a)^{-1}b).
\]
Consequently the integral of \(f(m(a,h)u m(a,h)^{-1})\) over the radical is \(|a|_E^{r+1}\int f(u)\,du\): the \(r\) column entries contribute \(|a|_E^r\) since the determinant of \(h\) has module one, and the self-adjoint scalar contributes \(|N_{E/F}a|_F=|a|_E\). The rational radical has compact quotient, and the summed line term is invariant under its adelic radical after averaging. The Iwasawa integral therefore has a fixed positive Haar factor and remaining variables \(a\in E^\times\backslash\mathbb A_E^\times,\ h\in[U(V)],\ k\in K_{V'}\), with the factor \(|a|_E^{r+1}\). This is where the extra full power of \(|a|\) enters.\formalnote{10.41}{41}

\subsection{The two functions on the line}
Fix \(g\in U(W)(\mathbb A_F)\) and the source Levi scalar \(u=1\), or any \(u\) when \(n=r\). For fixed \(h,k\), denote by \(\Phi_+(z)\) the smaller theta kernel at \((g,h)\) of the following datum. Act by \(k\), take the source slice in direction \(zk_V\), then the target zero term. For \(\Phi_-(z)\), translate in the opposite pair of null coordinates, take the target zero term first, then the source zero term. These orders give two polarizations of the remaining adelic line.

Here is the summation identity. For either polarization, first take the common zero term in the directions perpendicular to the residual plane. On that plane the rational slice sum equals evaluation of the original theta functional. To see this, expand on the compact additive quotient by rational points in the other null direction. Its Fourier coefficient at the character indexed by \(\xi\in E\) is the zero-term average of the translate by \(\xi\); hence its value at the identity is the sum of those coefficients. The rapid slice bound ensures absolute summability. Applying this in the two orders gives the same theta value. It applies also after translating by a Levi point. The two new arguments are respectively \(\xi a\) and \(\xi\bar a^{-1}\). The ratio of their scalar multipliers is \(|a|\). Explicitly, with the above parabolic Jacobian included, the two factors are \(\widetilde\nu(a)|a|^{e+1}\) and \(\widetilde\nu(a)|a|^{e}\). The \(n\) half powers from the source Weil action give \(|a|^{-n/2}\), and the opposite line model has one more inverse power. Here \(\widetilde\nu(a)=\nu(a/\bar a)\). At general allowed \(u\) the resulting arguments and factors use \(au\). At zero both line values vanish: the finite one-dimensional factor is supported on units on one set of marked components, and the transformed nontrivial residue character has sum zero at frequency zero. Equivalently apply the finite factorization of \S\ref{sec:finite-period} on each side. With \(u=1\) the identity becomes\formalnote{10.42}{42}
\[
    |a|\sum_{\xi\in E^\times}\Phi_+(\xi a)
        =\sum_{\xi\in E^\times}\Phi_-(\xi\bar a^{-1}).                 \tag{PP1}
\]

Let \(\mathcal Z(x)\) be the integral in \(a,h,k\) of \(\widetilde\nu(au)|au|^x\sum_{\xi\ne0}\Phi_+(\xi au)\), using the quotient measures just described. The constant term is \(c_0\mathcal Z(e+1)\) for a fixed nonzero \(c_0\), provided the unfolded integral converges absolutely. In (PP1) replace \(a\) by \(au\). Conjugate inversion preserves the idele-class Haar measure, fixes \(\widetilde\nu\), and inverts the module. Translation by \(u\) also preserves the measure. It follows that \(\mathcal Z(1-x)\) has the minus-side expression with exponent \(x\). Thus on \emph{both} sides we can now unfold the sum over \(E^\times\) against the idele-class integral at a positive exponent.

\subsection{Convergence and the scalar}
Write \(M_N(v)=\int(1+\|y\|)^N|v(y)|\,dy\) for the weighted absolute integral in the Gaussian model, with its fixed measure and the maximum norm in the polynomial weight. In Gaussian exponents use \(\|y\|_2^2=\sum_j|y_j|^2\). A Schwartz vector has all these moments finite. If \(S(a)\) is a slice in an adapted mixed model and \(M,N\ge0\) are integers, then
\[
       (1+\|a\|)^M M_N(S(a)v)\ \le\ C_{N,M}M_{N+M}(v).
\]
Indeed resolve \(v\) into translates of the vacuum with coefficient \(v(y)\) in the reproducing integral. Decompose \(y=L A y+I B y+K D y\) using the three bounded frame-coordinate maps. The slice of that translate is, up to phase, \(\exp(-\|a+Ay\|_2^2)\) times the small vacuum translated by \(By\). The \(N\)-th moment of the translate is at most \((1+\|By\|)^NM_N(v_0)\). The triangle inequality bounds \(1+\|a\|\) by \((1+\|a+Ay\|)(1+\|Ay\|)\). The Gaussian absorbs the first polynomial and the other two are bounded by \(C(1+\|y\|)^{M+N}\). Integration proves the estimate. Restricting the full slice parameter to the line preserves this decay since the adapted line-coordinate embedding bounds its argument from below.\formalnote{10.43}{43}

The real cone integrand is the slice of \(kv\) times \(\prod_w(a_w/|a_w|)^{2\mu_w}|a_w|^{2x}\), with \(a_w\in\mathbb C^\times\). For \(x>0\), \(t^{2x}(1+\log^2 t)\) is bounded near zero and grows at most polynomially at infinity. The slice estimate thus bounds any required moment by
\[
                C(k)\prod_w (1+\log^2|a_w|)^{-1}.
\]
Here \(C(k)\) is bounded on the compact group, by the Fock moment bound for the real Weil operator and continuity of its matrix multiplier. Multiplicative Haar measure in polar coordinates is a constant times \(d\theta\,d\log t\), so the displayed majorant is integrable. For the remaining \(h\) integral use the primitive-support theta bound: at fixed \(g\) the kernel is bounded by \(C_g M_N\), uniformly on \([U(V)]\), which has finite volume. On the support the variable rows are rows of an invertible stacked matrix, so the variable map is onto. Section~\ref{sec:primitive-decay} derives the resulting estimate with arbitrary target height decay. These two bounds justify changing order and passing the vector-valued cone integral inside the smaller theta kernel.\formalnote{10.44}{44}

For \(x>1\), unfolding and separating finite from infinite variables now gives, for \(t_+=x,\ t_-=1-x\),
\[
 \mathcal Z(t_\pm)=
 \epsilon_\pm\,\alpha_0^x Q_S^{\,x-1/2}L^S(x,\nu)\,
     \Theta_{W,V}\bigl(\operatorname{Cone}_\pm(x)v,\,f_y\bigr)(g).
                                                                    \tag{PP2}
\]
Here \(\alpha_0>0\) is the module of the scale idele, \(Q_S=\prod_{w\in S}q_w\), \(y\) is the fixed integral row datum, and \(|\epsilon_\pm|>0\) depends only on the fixed data and the side. The factor \(Q_S^{-1/2}\) is the modulus of the finite Fourier scalar, \(Q_S^x\) is the finite translation in \S\ref{sec:finite-period}, and the remaining finite shell sum is the incomplete Euler product. Absolute convergence also holds at \(R\) and \(1-R\) for every \(R>1\). Given \(t\in\mathbb R\), take \(R=|t|+2\). Powers of \(\log|au|\) on an interior subinterval are absorbed by the two endpoint powers. Mellin integration gives absolute convergence at \(t\) and a real-analytic function there. In (PP2) let \(x\downarrow1\): the cone integral is continuous by domination between its integrable endpoint weights at 1 and 2. The Hecke factor extends continuously and is nonzero by \S\ref{sec:hecke-endpoint}.

For \(e\ge0\) take \(x=e+1\) on the plus side. For \(e<0\) take \(x=-e\) on the minus side so that \(1-x=e+1\). In either case \(x\ge1\) and the scalar in the plane constant term has modulus \(C Q_S^{x-1/2}|L^S(x,\nu)|\), with \(C>0\) independent of the twisting character at fixed weight and ranks. In particular, the first negative step \(e=-1\) and the middle step \(e=0\) both call for the endpoint \(x=1\).\formalnote{10.45}{45}

\subsection{Descent on determinant vectors}\label{sec:det-cone}
To descend vanishing of a theta integral through one hyperbolic plane, we need only one
family of test vectors in \textup{(PP2)}. Here \(k=s_0\) and \(\mu_w\) is the weight of
\(\lambda\) at \(w\). Put \(m_w=\mu_w-k\ge0\).
For a matrix \(F_w\) of frame parameters at each infinite place, let \(u_{+,w}\) be the
positive-coordinate part of the variable \(u\) in the Fock model. Multiply the Gaussian by the
complex conjugate of the product, over \(w\), of the two principal minors of \(u_{+,w}F_w\)
on the positive and negative sign blocks of the target, raised to \(m_w\) and \(m_w+2k\), respectively.
Denote this vector by \(v_F\). On extending \(F\) by the adapted plane frame, each sign block
acquires a scalar diagonal block in the line slice. The two minors thus contribute
\(\bar a_w^{\,2\mu_w}\) in either model. The vector transforms under the compact group by the
character canceled in the cone integral, and \((a_w/|a_w|)^{2\mu_w}\) cancels the phase of this power,
so the cone integrand for the extended frame is \(v_F\) times
\[
                 \prod_{w\mid\infty} e^{-|a_w|^2}|a_w|^{2(x+\mu_w)}.
\]
For \(x>0\) its scalar integral is positive and finite, as is seen by putting \(t=|a_w|^2\)
in each radial integral. Rescaling the extended vector gives any prescribed multiple of \(v_F\)
downstairs. For the value \(x\ge1\) used in \textup{(PP2)}, the Hecke factor is nonzero, and
the adapted integral frame carries each standard row datum to an admissible datum upstairs.
Consequently vanishing on these vectors upstairs implies vanishing on them downstairs.
When the second space has rank zero, the two minors are empty determinants and \(v_F\) is
the vacuum.\formalnote{10.46}{46}

\section{The ramified line in the theta exchange}
\label{sec:finite-period}

The scalar for removing a hyperbolic plane in \S\ref{sec:one-exchange}
contains a familiar local calculation. The Fourier transform of a
nontrivial residue character is a Gauss sum and is supported on a
single valuation shell. Its absolute value supplies \(Q_S^{-1/2}\).
The shell forces a translation at the conjugate places, contributing
\(Q_S^z\) to the multiplicative integral. Together these give
\(C Q_S^{\,z-1/2}|L^S(z,\nu)|\), with \(Q_S\) the product of residue
cardinalities at the marked places.

\subsection{The Fourier transform on the ramified line}

Let \(\nu\) be the positive-weight character of the norm-one torus in the plane-removal formula (\(\lambda\) in \S\ref{sec:one-exchange}), with tame ramification at the nonempty set \(S\), disjoint from \(\bar S\). Its base change is \(\widetilde\nu(a)=\nu(a/\bar a)\). At a place \(w\in S\), write \(F_w\) for the local field, \(O_w\) for its integers, \(\varpi_w\) for a uniformizer, \(k_w\) for its residue field and \(q_w=|k_w|\). The scalar character \(\chi_w\) on this factor of \(\prod_{w\in S}F_w^\times\) restricts to a nontrivial residue character on \(O_w^\times\). Define
\[
  b_w(x)=
  \begin{cases}\chi_w(x)^{-1}&x\in O_w^\times,\\0&\text{otherwise}.\end{cases}
\]
The additive character \(\psi_w\) in the integral frame has conductor \(O_w\). Thus \(\bar\psi_w(r)=\psi_w(\varpi_w^{-1}\tilde r)\) is well-defined and nontrivial on \(k_w\), where \(\tilde r\) is any lift. Normalize the additive measure by \(\operatorname{vol}(O_w)=1\).\formalnote{10.47}{47}

\begin{example}[The Fourier transform over \(\mathbb Q_3\)]\label{ex:line-gauss-three}
Choose \(\psi\) of conductor \(\mathbb Z_3\) with
\(\psi(1/3)=\omega=e^{2\pi i/3}\). Let \(b\) be the quadratic
character on \(\mathbb Z_3^\times\), extended by zero off the units.
The classes \(1+3\mathbb Z_3\) and \(2+3\mathbb Z_3\) have measure
\(1/3\) and carry the values \(1\) and \(-1\). For their Fourier
integral \(I(t)=\int_{\mathbb Z_3}b(x)\psi(tx)\,dx\) this gives
\[
 I(1/3)=(\omega-\omega^2)/3=i/\sqrt3,\qquad I(2/3)=-i/\sqrt3 .
\]
Adding any \(3\)-adic integer to \(t\) leaves the integral
unchanged. At integral \(t\) the two unit classes cancel. If
\(t=3^{-k}u\), \(u\in\mathbb Z_3^\times\), \(k\ge2\), translation
of \(x\) by \(3^{k-1}\) preserves \(b\) and multiplies
\(\psi(tx)\) by \(\psi(u/3)\ne1\), so the integral vanishes there
too. Its support is exactly \(3^{-1}\mathbb Z_3^\times\), and its
absolute value throughout that shell is \(3^{-1/2}\).
\end{example}

For \(t\in F_w\), put \(I_w(t)=\int_{O_w}b_w(x)\psi_w(tx)\,dx\). If \(v_w(t)<-1\), translating by a suitable \(\varpi_w h\), \(h\in O_w\), preserves \(b_w\) and multiplies the integrand by a nontrivial scalar. Thus \(I_w(t)=0\). If \(t\in O_w\), including \(t=0\), the additive character is trivial throughout \(O_w\). The sum of the nontrivial multiplicative character over \(k_w^\times\) vanishes. In the remaining case \(t=\varpi_w^{-1}u\) with \(u\) a unit. Integrating residue class by residue class and substituting \(s=\bar u r\) in \(k_w\) gives
\[
 I_w(\varpi_w^{-1}u)=\chi_w(u)\gamma_w,\qquad
 \gamma_w=q_w^{-1}\sum_{s\in k_w^\times}
                      \chi_{w,\mathrm{res}}(s)^{-1}\bar\psi_w(s).
 \tag{FP1}
\]
To compute its size, take the squared absolute value of the unnormalized sum and put \(s=h r\) in the double sum over \(s,r\ne0\). The result is
\[
 \sum_{h\in k_w^\times}\chi_{w,\mathrm{res}}(h)^{-1}
                 \sum_{r\in k_w^\times}\bar\psi_w((h-1)r)
       =(q_w-1)-\sum_{h\ne0,1}\chi_{w,\mathrm{res}}(h)^{-1}=q_w .
\]
Consequently \(|\gamma_w|^2=q_w^{-1}\). Over \(F_S=\prod_{w\in S}F_w\) these integrals factor. With \(\chi=\prod\chi_w\), \(\varpi=(\varpi_w)_w\), \(O_S=\prod O_w\) and \(Q_S=\prod q_w\), the transform is supported where \(\varpi t\in O_S^\times\) and equals \(\gamma\chi(\varpi t)\) there, with \(|\gamma|^2=Q_S^{-1}\). Primitivity at each marked place is what removes the frequencies integral at that place.\formalnote{10.48}{48}

\subsection{The transform at the prescribed determinant datum}

In the smaller pair, stack an integral fixed row block \(P\) above the variable row block \(Y\), in the prescribed ordering, to obtain a square matrix \(M\). Its standard row function is \(\chi(\det M)^{-1}\) on \(\mathrm{GL}(O_S)\) and zero elsewhere. Split a hyperbolic plane from each space and use integral bases together with their integral duals. In these coordinates the larger row function is the standard determinant function evaluated on \(G T R\). Here \(R\) is the inverse of the source frame, and \(G\) inserts the two target isotropic columns and the remaining frame columns, together with the identity on the fixed rows. The dual rows give an integral left inverse for \(G\). The combined square change \(RG\) has unit determinant. Recovering \(T\) with the integral left inverse and the source frame shows that \(GTR\) is integral exactly when \(T\) is integral. Its determinant is \(\det(RG)\det T\). Thus the basis change only multiplies the row function by the unit-modulus scalar \(\epsilon=\chi(\det(RG))^{-1}\).\formalnote{10.49}{49}

In the enlarged matrix the diagonal blocks of \(T\) are \(a_S,M,t\) in this order on one side of Poisson summation, and \(t,M,a_S\) on the other. All blocks below the diagonal are zero. The free off-diagonal entries are coordinates of the slice, with some positions fixed to zero. Here \(a_S\) is the marked part of the line variable after division by the fixed Weil scale \(\alpha_0\). In both orders, \(\det T=a_S t\det M\). On integral matrices each factor must be a unit for the determinant function to survive. Hence the three factors in its value are \(b(a_S)\), \(b(t)\) and the smaller row function at \(Y\), where \(b=\prod b_w\). Every other free slice coordinate carries the indicator of \(O_S\). The mixed-model Fourier character has frequency zero in all these other coordinates and frequency \((\bar a)_S\) in the \(t\) coordinate. The integrals of the indicators are one. Formula (FP1) evaluates the one remaining integral.\formalnote{10.50}{50}

The resulting finite function on an adele \(a\), apart from a scalar of modulus \(Q_S^{-1/2}\) and the fixed vacuum scalar of the chosen side, is
\[
 H_{\varpi}(a)=
 \begin{cases}
   \chi(a_S)^{-1}\chi(\varpi(\bar a)_S)&
      \begin{array}{l} a\text{ integral off }S\cup\bar S,\\
         a_S,\ \varpi(\bar a)_S\in O_S^\times,\end{array}\\
   0&\text{otherwise}.
 \end{cases}                                                     \tag{FP2}
\]
The vacuum outside \(S\cup\bar S\) gives the first support condition, and \(\epsilon\) is absorbed into \(\gamma\). In particular the same computation kills the zero line on both sides. A marked place \(w\) and its conjugate play different roles. In (FP2) the valuation of \(a_w\) is zero and that of \(a_{\bar w}\) is \(-1\).\formalnote{10.51}{51}

\subsection{The translation and the finite idele integral}

Let \(H_1\) be (FP2) with \(\varpi\) replaced by \(1\). Take the finite idele \(c\) equal to one away from \(\bar S\), with \((\bar c)_w=\varpi_w\) for \(w\in S\). Then
\[
        H_{\varpi}(a)=H_1(ac),\qquad |c|_f=Q_S^{-1}.
\]
Multiplication by \(c\) preserves the components at \(S\) and makes the shifted components at \(\bar S\) units. It also preserves the integrality conditions elsewhere. Write \(U_f=\prod_vO_v^\times\). On \(u\in U_f\), the base-change character is \(\chi(u_S)\chi((\bar u)_S)^{-1}\), so \(H_1(u)\widetilde\nu_f(u)=1\).\formalnote{10.52}{52}

For a positive real exponent \(x>1\) take any multiplicative Haar measure \(d^\times a\) on the finite ideles and set \(V_f=\operatorname{vol}(U_f)\). The support of \(H_1\) on ideles is indexed by integral ideals \(\mathfrak a\) prime to \(S\cup\bar S\). For each choose \(a_{\mathfrak a}\) with the prescribed valuations and components \(1\) at the marked and conjugate places. The fiber is \(a_{\mathfrak a}U_f\) and has measure \(V_f\). By the unit cancellation its contribution is
\[
             V_f\,\widetilde\nu_f(a_{\mathfrak a})(N\mathfrak a)^{-x}.
\]
The character is unitary, hence the sum of absolute values is bounded by \(V_f\zeta_E(x)\). Expanding the geometric series independently at each unmarked place now gives
\[
 \int H_1(a)\widetilde\nu_f(a)|a|_f^x\,d^\times a
    = V_f \prod_{v\notin S\cup\bar S}
         \left(1-\widetilde\nu_v(\varpi_v)q_v^{-x}\right)^{-1}
    = V_f L^S(x,\nu).
\]
Changing variables \(b=ac\) multiplies this answer for \(H_\varpi\) by \(\widetilde\nu_f(c)^{-1}Q_S^x\). There is no change in multiplicative Haar measure. The absolute value of the extra character is one. Together with (FP1) this gives the asserted \(Q_S^{x-1/2}\) in a single plane removal.\formalnote{10.53}{53}

More precisely, the direct Poisson variable is \(a\) and the dual variable is \(\bar a^{-1}\). The base-change character takes the same value on the two, and their modules are reciprocal. After quotient unfolding, if \(I_\beta(x)\) denotes the integral of the smaller theta kernel against the real cone on the chosen side, the two identities have the form
\[
 \mathcal Z(s_\beta(x))
  =e_\beta\,|\alpha_0|^x Q_S^{\,x-1/2} L^S(x,\nu)\, I_\beta(x),
 \qquad s_+(x)=x,\quad s_-(x)=1-x .
\]
Here \(\mathcal Z\) includes the fixed unfolding constant. The modulus of \(e_\beta\) depends on the side and the fixed spaces and measures, uniformly in the tame character. Its phase includes the normalized Gauss sum and the character at \(\alpha_0c^{-1}\). The rank of the smaller source is even, the splittings are adapted to the integral vacuum, the real vector is Schwartz, and the Levi coordinate is either \(1\) or the two smaller spaces have equal rank. These are the hypotheses of the constant-term identity used in the exchange.\formalnote{10.54}{54}

The endpoint argument following \textup{(PP2)} in \S\ref{sec:plane-period} applies at \(s=1\) on the direct side and \(s=0\) on the Poisson side. Thus at the first step the squared modulus contains \(Q_S|L^S(1,\nu)|^2\), up to the fixed factor used in \S\ref{sec:critical-values}.\formalnote{10.55}{55}

\subsection{Support of the definite character-lift vector}

For the other critical values the definite pair \(U(1)\times U(N)\), \(N=2z\ge4\), has a different finite vector. At a chosen place of norm \(q\) its column function \(\Psi\) is supported on the distinguished basis column \(e_i\) modulo the prime. At a twisting place it is the tame character at the distinguished entry \(x_i\), extended by zero when \(x_i=0\). Elsewhere it is the vacuum. Write \(J\) for the subgroup of the integral source group that reduces to \(1\) at each chosen place. We verify the support of its matrix coefficient, the property used in the definite norm calculation in \S\ref{sec:one-exchange}.

Every supported column \(x\) has \(x_i\) a unit at each marked place and \(x_i=1\) in the residue field at each chosen place. Thus it has an integral left inverse at the marked places. To lift its distinguished entries in the full residue ring to an element of \(J\), lift the units at one of each pair of split places and use their conjugate inverses at the other places to obtain a norm-one idele. The finite Weil transformation law fixes the vector on \(J\). Indeed at the chosen places its distinguished column is preserved, and at the twisting places the determinant character of the action cancels the character acquired by the block function. These use the prescribed splitting character of the pair.\formalnote{10.56}{56}

Take a finite group point \(g\) integral off the marked places and their conjugates. A nonzero matrix coefficient requires that some column \(x\) and \(xg\) both meet the support in the local Schrödinger coordinates. Their primitivity forces \(g\) to have valuation zero at every marked place, since \(x_i\) and \(x_i g\) are both units. The norm-one condition gives the same at the conjugates. If this now integral \(g\) fails to belong to \(J\), at some chosen place the distinguished entry of \(xg\) differs from \(1\), contradicting its support. Thus the coefficient vanishes off \(J\) in this region. On \(J\) it is the squared norm of the vector. With the integral vacuum normalized, this is the mean of \(|\Psi|^2\) on the residue columns. There is one supported column out of \(q^N\) at a chosen place, and \((q-1)q^{N-1}\) out of \(q^N\) at a twisting place, each of squared absolute value one. Dividing the product of these means by \([U(1)_{\rm int}:J]=\prod_{\rm chosen}(q-1)\) gives the mass fed into the Rallis pairing. This argument includes auxiliary moduli at which a twisting residue character is trivial, since the block function still requires its distinguished entry to be a unit.\formalnote{10.57}{57}

The finite calculations above use the mixed-model slice formula and its self-dual measure, the integral adapted frames, and the finite Weil action of the splitting character on the block datum. Sections~\ref{sec:plane-period} and~\ref{sec:hecke-endpoint} prove the global plane identity and the endpoint used here. Sections~\ref{sec:doubling-integral}, \ref{sec:real-pairing} and~\ref{sec:split-vacuum} derive the definite norm after inserting this marked coefficient.


\section{Inserting a prime into the determinant period}
\label{sec:period-insertion}
In the endpoint case of \S\ref{sec:critical-values}, the character stays fixed while its declared modulus grows. Removing a pair of Euler factors from an absolutely convergent series would suffice above \(1\). At \(1\) we shall instead calculate the change in a definite theta period from its test function. The local calculation works in every rank and explains both factors in rank two.

Section~\ref{sec:finite-period} computes the finite Gauss and translation factors, and the support of the character-lift vector used here.

\subsection{A sum over subspaces removes exactly the singular matrices}
Let \(n\) be the common rank in the definite pair \((I_n,I_n)\). Let \(\mathfrak m,\mathfrak m'\) be moduli prime to their conjugates, with the support of \(\mathfrak m'\) obtained from that of \(\mathfrak m\) by adjoining \(v\). Fix the Weil splitting character \(\upsilon\). On both sides of the pair use the same automorphic character \(\alpha\) of the norm-one torus, of weight \((-c_\sigma)_\sigma\) with \(c_\sigma\ge0\). At both levels use the polynomial \(\prod_\sigma\det(Z_\sigma)^{c_\sigma}\), where \(Z_\sigma\) is the matrix of variables at \(\sigma\), and the tame residue characters of \(\alpha\). Away from \(\mathfrak m\) each residue character is trivial on units and has value zero at zero. In particular \(v\) is split and neither \(v\) nor \(\bar v\) was in the old support. Denote its residue field by \(k\), its order by \(q\), its uniformizer by \(\varpi\), and its local integer ring by \(O_v\). Write
\[
 c_v=\alpha_{\rm sp}(\varpi)q^{n/2},\qquad X=c_v^{-1},
\]
where \(\alpha_{\rm sp}(a)=\alpha(a/\bar a)\upsilon(a)^{-n}\) is the local splitting-character normalization. At \(v\) the old datum is the vacuum. The new datum multiplies the characteristic function of integral matrices by the indicator of invertibility on reduction, since the trivial residue character of the determinant has value zero on every singular matrix.

Given \(U\subset k^n\) of dimension \(i\), choose \(\gamma_U\) with elementary divisors \((\varpi\,1_i,1_{n-i})\) and reduction kernel \(U\). To construct one, choose an integral basis change carrying \(U\) on reduction to the first \(i\) coordinates and multiply it on the left by this diagonal matrix. For every matrix \(y\) over the local field one has
\[
    y\gamma_U^{-1}\text{ integral}\quad\Longleftrightarrow\quad
         y\text{ integral and }U\subset\ker(\bar y).
\]
The forward direction follows by multiplying on the right by the integral \(\gamma_U\) and reducing. For the converse use the chosen basis: the first \(i\) columns are divisible by \(\varpi\) and can be divided by it while keeping all entries integral. In the split Schrödinger model the target action by \(\gamma_U^{\,t}\) on the vacuum contributes the local scalar \(c_v^i\) and the indicator of integrality of \(y\gamma_U^{-1}\), with \(y\) the transpose of the argument matrix. Indeed the action changes the argument by the inverse transpose and multiplies by the unramified determinant character with the half-density factor. The determinant has valuation \(i\), so that multiplier is \(\alpha_{\rm sp}(\varpi)^i q^{ni/2}\). Neither the kernel dimension nor invertibility changes on transposing.\formalnote{10.58}{58}

Weight the translated terms by \(h_U=(-X)^i q^{\binom i2}\). For integral \(y\), let \(d=\dim\ker(\bar y)\). The sum evaluated at \(y\) is
\[
      \sum_{i=0}^d \genfrac{[}{]}{0pt}{}{d}{i}_q\,q^{\binom i2}(-Xc_v)^i
              =\prod_{a=0}^{d-1}(1-q^a Xc_v).
\]
The bracket counts the \(i\)-planes in \(k^d\). To see the equality, fix a hyperplane. The \(i\)-planes it contains give the corresponding count in dimension \(d-1\). Every other plane meets it in an \((i-1)\)-plane, and choosing a lift of a generator of the transverse line then gives \(q^{d-i}\) possibilities. Substitution of this recurrence multiplies the degree-\((d-1)\) product by \(1-q^{d-1}Xc_v\). Starting at \(d=0\) proves the identity. Since \(Xc_v=1\), the product is \(1\) when \(d=0\), and otherwise its \(a=0\) factor is zero. Every translated term is zero for nonintegral \(y\). The old factors at the other places multiply this indicator. Consequently the linear combination of Weil translates is the desired new datum, including its normalization.\formalnote{10.59}{59}

\subsection{The product after integration}
Let \(K\) stabilize the standard lattice and let \(F_{\mathfrak m}\) denote the old kernel. The theta functional is linear. Applying it to the datum identity gives the equality of functions
\[
        F_{\mathfrak m'}(g)=\sum_{U\subset k^n}h_U
                 F_{\mathfrak m}(g\gamma_U^{\,t}),
\]
where the transpose indicates the finite adelic translate at the split pair. The automorphic transformation at rational \(\eta\) contributes the scalar \(\prod_\sigma\sigma(\det\eta)^{c_\sigma}\); substituting \(\eta^{-1}\) in \(\prod_\sigma\det^{c_\sigma}\) contributes its inverse. Thus the kernel is left invariant under rational points. For \(k_0\in K\) the finite datum changes by the factor \(\alpha(\det k_0)\) and the substitution \(y\mapsto k_0^{-1}y\). The latter multiplies \(\chi_{\mathfrak m}(\det y)\) by \(\chi_{\mathfrak m}(\det k_0)^{-1}\), which cancels the former by the given tame part of \(\alpha\). Both full kernels therefore have right level \(K\).

Normalize Haar so that a \(K\)-invariant kernel integrates on the definite quotient to the sum of its class values with weights \(1/|\operatorname{Aut}|\). For each \(U\) the translated old kernel has the same integral by right invariance of quotient measure. Individual translates can have different levels. Integrating their sum before replacing the two full integrals by \(K\)-class sums gives
\[
  \Lambda_{\mathfrak m'}(\alpha)=
       \left(\sum_{U\subset k^n}h_U\right)\Lambda_{\mathfrak m}(\alpha)
       =\prod_{a=0}^{n-1}(1-q^aX)\,\Lambda_{\mathfrak m}(\alpha).
                                                        \tag{PI1}
\]
The sum in \textup{(PI1)} runs over all subspaces of \(k^n\). The factor \(c_v^i\) belongs to the translated kernel, whose integral is unchanged.

\begin{example}[The subspace sum in a plane]\label{ex:plane-subspace-sum}
Let \(\bar y\) be a two-by-two matrix over the field \(k\) of \(q\) elements. To test whether it is invertible, sum over subspaces of its kernel, assigning \(1\) to the zero space, \(-1\) to each line and \(q\) to the plane. These weights are \(h_U c_v^{\dim U}\) in the preceding calculation, since \(Xc_v=1\). The plane \(k^2\) has \((q^2-1)/(q-1)=q+1\) lines. If \(\ker(\bar y)\) is a line, its two subspaces give \(1-1=0\). If \(\bar y=0\), all subspaces occur, with total \(1-(q+1)+q=0\). Only an invertible \(\bar y\) gives the value \(1\).

After integration each translated kernel has the same integral. The coefficients that multiply it are \(1\) for the zero subspace, \(-X\) for each line, and \(qX^2\) for the plane. Their sum is
\[
                   1-(q+1)X+qX^2=(1-X)(1-qX).
\qedhere\]
\end{example}

Apply \textup{(PI1)} at rank two with \(\alpha=\nu\tau\upsilon\). In the notation \(a_v=(\nu\tau)_{\rm BC}(\varpi)\) of \textup{(CV2)}, the splitting normalization gives \(X=(qa_v)^{-1}\). Since \(\overline{a_v}=a_v^{-1}\), squaring the absolute value of the two linear factors gives the factors of \textup{(CV2)} with exponents \(1\) and \(0\).\formalnote{10.60}{60}

\subsection{A common scalar for the integral values}
The scale \(N_0\) used at \(z=1\) can also be specified. Write \(f=[E:\mathbb Q]/2\), let \(M\) be the principal level and \(d_K\) the prime of the quadratic CM field. At the chosen tower step let \(a_0\) be the rank-one vacuum scale and \(b_0\) the Taylor-variable scale. For definite ranks \(m,n\), put \(t=mn\), \(D=M^2\), and let \(h\) be the homogeneous degree of the derivative polynomial. Take
\[
             C_{m,n,h}=D^{2ft}a_0^t(D d_K b_0)^h.       \tag{PI2}
\]
Every factor is nonzero. Assume all values of the residue datum \(\Psi\) are integral away from \(NQ\), and \((NQ)D\in\mathfrak m\). Fix \(p\nmid NQ\) and an embedding at which the derivative polynomial has integral coefficients. Finite Fourier inversion writes the \(a_0^t\)-scaled datum as a sum of translates of the integral vacuum. The denominators in that inversion divide the order of the finite character group used for the expansion. The \(p\)-part is bounded by that of \(|\mathcal O_E/D\mathcal O_E|^t=D^{2ft}\), since the residue modulus divides \((NQ)D\) and \(NQ\) is a \(p\)-unit. The translation parameters become integral over \(p\) when multiplied by \(D\), by the same divisibility applied to the dual residue group in the vacuum coordinates. Multiplying by \(D^{2ft}\) therefore gives an integral combination of the standard level-\(D\) translates. This argument uses the identification of the finite additive characters with the Heisenberg translates in the chosen self-dual coordinates.\formalnote{10.61}{61}

The integral Weil-translation and Taylor estimates say that at integral invertible adelic arguments over \(p\), the Taylor series of such a combination has integral coefficients with variable scale \(D d_K b_0\). Apply the polynomial in derivatives at zero. For a monomial \(\partial^\beta\) the coefficient extracted is multiplied by \(\prod_i\beta_i!\), an integer, and \(|\beta|=h\) for every monomial. Scaling the variables multiplies the jet by \((Dd_Kb_0)^h\); linearity supplies the datum scale \(D^{2ft}a_0^t\). This proves that (PI2) times the tested kernel is integral. For the endpoint period take \(m=n=2\), \(h=2\sum_\sigma(c_\sigma-1)\). Each class point in the finite measure of \S\ref{sec:critical-values} may be moved to an integral point above the selected \(p\) by rational translation. The determinant kernel is rationally invariant as checked above. Thus the same \(C_{2,2,h}\) works at every allowed \(p\), for every auxiliary modulus.\formalnote{10.62}{62}

For \(z\ge2\) the point kernels in \S\ref{sec:critical-values} already include the vacuum and Taylor scales. Their integral multiplier can be taken to be the integer
\[
                     D_0=M^{4fz}(M d_K)^{h'},\qquad
                  h'=\text{number of boxes of the lift weight}.
\]
Indeed the finite datum has \(2z\) entries and uses level \(M\), giving \(M^{2f(2z)}\) by the same inversion. A jet with \(h'\) boxes costs at most \((Md_K)^{h'}\) relative to its prescribed scale: the additional factor \(\pi_M\) (the conjugate of the chosen generator on the elliptic side) in the variable of that jet is integral. Moving the output back to the Weyl lattice uses the integral gauge transformation for the kernel, applied to the vector \(1\) on \(U(1)\), whose gauge is one. This yields exactly the pointwise integral vectors needed before taking their class Gram pairing. The remaining CM input for both scales is the integral Taylor estimate at the specified period, together with preservation of standard data by integral Weil translations and the transformation of the Weyl gauge.\formalnote{10.63}{63}


\section{The partial Hecke value at one}\label{sec:hecke-endpoint}
Let \(K\) be the CM field in the determinant period of \S\ref{sec:critical-values}. Let \(X\) be a unitary character of \(K^\times\backslash\mathbb A_K^\times\), with infinite component \(\prod_{w\mid\infty}(z_w/|z_w|)^{n_w}\). Write \(S\) for a finite set of finite places containing its ramification.
\begin{proposition}\label{prop:hecke-one}
If some \(n_w\ne0\), the Euler product \(L^S(\sigma,X)\), initially defined for real \(\sigma>1\), has a finite, nonzero right limit at \(1\).
\end{proposition}
\begin{proof}
We use the adelic form of Hecke's continuation theorem \cite{hecke1920}. For an idele class quasicharacter \(c\), Tate's global zeta integral can have poles only at \(c=1\) and \(c=|\cdot|_{\mathbb A_K}\) \cite[Thm.~4.4.1 and \S4.5]{tate1950}. Put \(c=X^j|\cdot|_{\mathbb A_K}^{s}\), \(j=1,2\). On the unit circle at the specified place its exponent is \(jn_w\ne0\), so neither exceptional character occurs.

To pass to the partial Euler product, take the characteristic function of \(\mathcal O_v\) at every finite \(v\notin S\), with multiplicative measure giving \(\mathcal O_v^\times\) volume one. Its local zeta integral is the Euler factor. At the remaining places take a nonnegative smooth test function of compact support in \(K_v^\times\), positive near \(1\), times \(X_v^{-j}\). The resulting local integrals are entire in \(s\) and positive at \(1\). Dividing Tate's integral by their product gives the holomorphy of \(L^S(s,X^j)\) at \(1\). For the trivial character the usual finite-place indicators and archimedean Gaussians give a nonzero residue in Tate's formula. Their archimedean factors are finite and nonzero at \(1\), and omitting the factors at \(S\) multiplies the residue by \(\prod_{v\in S}(1-q_v^{-1})>0\). Thus \(\zeta_K^S\) has a simple pole at \(1\).

For \(\sigma>1\) expand the logarithm of
\[
   \zeta_K^S(\sigma)^3\,|L^S(\sigma,X)|^4\,|L^S(\sigma,X^2)|.
\]
For each \(v\notin S\) and each \(j\ge1\), the coefficient of \(q_v^{-j\sigma}/j\) is \(3+4\Re z+\Re(z^2)=2(1+\Re z)^2\), where \(z=X(\varpi_v)^j\). The logarithm is nonnegative, so the product is at least one. A zero of \(L^S(s,X)\) at \(1\) would give a factor of order at least four against a pole of order three, while the last factor stays bounded. The product would tend to zero. Holomorphy and nonvanishing give the asserted right limit.
\end{proof}
In the application the Hecke character comes from \(b\mapsto\nu(b/\bar b)\). Its infinite type is \(2\mu\ne0\). Inversion negates the type, and finite-order twists preserve it. Thus the proposition applies at the endpoint of each string. It evaluates the limit from \(\sigma>1\).\formalnote{10.64}{64}

\phantomsection\label{route:19}\leanroute{The formal proof establishes an analytic germ at \(1\) by unfolding a compactly supported adelic test function and estimating the difference of two of its translates at equal module. It obtains the principal residue by an ideal count. Hecke's continuation theorem supplies both analytic inputs here. The Euler-factor inequality then proves exactly the nonvanishing required by the period.}\formalnote{10.65}{65}


\section{Decay of theta kernels with primitive support}
\label{sec:primitive-decay}

In the plane integral of \S\ref{sec:plane-period}, the archimedean vector varies with the integration variables. We need a bound uniform over the target unitary quotient, with the dependence on this vector made explicit. Write \(S,T\) for the source and target Hermitian spaces over \(E/F\), of ranks \(n,r\), and \(X=\operatorname{Hom}_E(S,T)\). Let \(\omega(g,h)\) be the pair Weil action with its two fixed automorphic twists, and let \(\theta\) be the rational theta functional on the data space of \(X\). Write \(P\) for the set of marked finite places of \(E\), with \(P\cap\bar P=\varnothing\). In the Schrödinger coordinates at \(P\), take a multiplier \(\Psi\) locally constant and compactly supported in those coordinates. Apply it to the finite vacuum to obtain a fixed datum \(F_\Psi\). We write \(J_\Psi v=v\otimes F_\Psi\) for the integrated datum with archimedean Fock vector \(v\).

The hypothesis on \(\Psi\) is the following. There is a unital homomorphism \(\phi:\mathbb A_{E,f}\to R\) to a nonzero commutative ring such that whenever \(\Psi(Y)\ne0\), the \(r\times n\) matrix obtained from \(Y\) by applying \(\phi\) has a right inverse. The application uses projection to the product of the marked completions. Put \(d=[F:\mathbb Q]nr\) and
\[
 M_j(v)=\int_{\mathbb C^d}(1+\|z\|_\infty)^j|v(z)|\,\frac{dz}{\pi^d}.
\]
The Gaussian belongs to the vector in this convention. In particular \(M_0(v)=\|v\|\). For some integer \(A\ge0\) depending on the pair and its action, each integer \(N\ge0\) has a constant \(C_N\) such that
\[
  |\theta(\omega(g,h)J_\Psi v)|\,H_T(h)^N
       \le C_N\,M_{N+A}(v)\,H_S(g)^{N+2A}.                       \tag{PD1}
\]
It suffices that the displayed moment be finite. The constant is uniform in \(g,h,v\) for the fixed data. There is a companion estimate with source and target interchanged when the matrices on the support have a left inverse.\formalnote{10.66}{66}

\subsection{A short isotropic family in the space of maps}

The adelic frame \(x_S\) takes the rational Hermitian matrix to a diagonal matrix of signs at infinity and to the identity at the finite places. At \(g\), the lattice consists of the rational columns \(c\) for which \((x_S g)^{-1}c\) is integral at the finite places. Its length is the Euclidean norm over all archimedean coordinates. The height \(H_S(g)\) is the reciprocal of the first minimum cut off above at \(1\). Define \(H_T(h)\) in the same manner. Set \(\rho=H_S(g)\) and \(H=H_T(h)\).

First, the full functional \(L=\theta\circ\omega(g,h)\) satisfies
\[
                         \|L\|\le C(\rho H)^A .                    \tag{PD2}
\]
Here is the height reduction behind this estimate. Two lattice vectors of lengths less than \(1\) have zero Hermitian pairing. The pairing is an algebraic integer and each archimedean absolute value is less than \(1\), so its field norm forces it to vanish. The short columns therefore span an isotropic space at each infinite place. A unitary dilation doubles this space, halves its dual isotropic space and fixes the remaining summand. Its operator norm and inverse norm are at most \(2\). In the moved lattice the short lengths double, and all other lengths remain at least \(1/2\). Thus each step replaces a height \(H_0\) by a height at most \(\max(H_0/2,2)\). After \(k=\lfloor\log_2 H_0\rfloor\) steps Mahler compactness gives a representative in a fixed compact set modulo rational points. The Weil operators on that set and on the compact set of real factors of norm at most \(2\) have bounded norm, say at most \(2^A\), increasing \(A\) to accommodate both groups. The \(k\) factors cost at most \(H_0^A\). Rational invariance of \(\theta\) and reduction in each group give (PD2).\formalnote{10.67}{67}

Suppose \(H>2\). The definition of the minimum gives a nonzero rational target column \(u\) in the moved lattice with length less than \(2/H\). Since \(2/H<1\), its Hermitian norm is zero. In the source there is a rational basis of columns \(c_1,\ldots,c_n\) belonging to the moved lattice and of lengths at most \(Z_S\rho\). To see the second assertion, put \(f=[F:\mathbb Q]\) and take a positive integer \(D_E\) in the different of \(E/\mathbb Q\). The trace-dual lattice calculation gives first minimum at least \(1/\rho\) for the Euclidean dual of the source lattice rescaled by \(2/D_E\). Weak lattice transference in real dimension \(2fn\) places every point within \(R_0=(2fn)^2\rho\) of this rescaled lattice. Its vectors of length at most \(3R_0\) span. Otherwise a point of length \(2R_0\) orthogonal to their span would have a lattice point within \(R_0\), necessarily of length at most \(3R_0\), a contradiction. Real spanning of the archimedean images implies rational \(E\)-spanning, since a rational linear form vanishing there would also vanish on every vector in any one complex realization. Rescaling and selecting a basis give \(Z_S=6D_E(fn)^2\) for \(n>0\). In rank zero the empty family is the basis.\formalnote{10.68}{68}

For \(1\le j\le n\) form the maps \(B_j(s)=u\langle c_j,s\rangle_S\). Their span \(Z\subset X\) is totally isotropic. Indeed the Hermitian pairing of two rank-one maps has a target factor \(\langle u,u\rangle_T=0\). Their moved coordinates are integral at the finite places. At an infinite place the Hilbert--Schmidt norm of a rank-one matrix is the product of the two column norms, since the diagonal signature matrix preserves Euclidean length. Summing the squares over places bounds the total length of each map by\formalnote{10.69}{69}
\[
                             \delta=2Z_S\rho/H .                   \tag{PD3}
\]

\subsection{The mean killed by the support condition}

Write \(\varrho(z)\) for Heisenberg translation. On \(Z\) the cocycle is trivial by isotropy. Consequently the rational translation invariance of \(\theta\) defines the probability mean
\[
   \theta_0(F)=\int_{Z(\mathbb A_E)/Z(E)}\theta(\varrho(z)F)\,dz ,
\]
where the coordinates of \(Z\) are taken in the fixed adelic frame of \(X\). Translation invariance of Haar measure makes \(\theta_0\) invariant under the whole adelic space \(Z(\mathbb A_E)\).

Let \(U=(x_T h)^{-1}u\) and \(C=(x_Sg)^{-1}[c_1\,\cdots\,c_n]\). The finite pairings of \(Y\) with the moved columns \(B_j\) are the entries of \(U_f^*Y C_f\), since the standard finite Gram matrices are identities. This row cannot vanish where \(\Psi(Y)\ne0\). Indeed \(C_f\) is invertible. After applying \(\phi\), cancel \(C_f\) on the right and then use a right inverse to \(\phi(Y)\). Vanishing would force \(\phi(U_f^*)=0\). But \(U_f^*=u^*((x_T h)^{-1}_f)^*\) is a nonzero rational row times an invertible matrix. The unital map from the field \(E\) to \(R\) is injective.

Pull \(\theta_0\) back through \(\omega(g,h)\). Weil covariance makes this functional invariant under translations by the moved \(Z\). On the marked coordinates, a translation supported on the conjugates \(\bar P\) multiplies Schrödinger data by an additive character. At a point of the support take its representative with complementary coordinates zero. Some pairing with the moved \(Z\) is nonzero there. Self-duality of the finite additive character supplies a scalar multiple in that adelic subspace whose character has value different from \(1\) there. Project this translation to \(\bar P\) by the corresponding adelic idempotent. Its multiplier still has that value near the point, by local constancy. Subdivide the compact marked support into small open cosets with these constant multipliers. On each piece, invariance equates its value under the functional with the same value multiplied by a scalar different from \(1\). Each value vanishes. Summing over the pieces gives\formalnote{10.70}{70}
\[
                         \theta_0(\omega(g,h)J_\Psi v)=0 .         \tag{PD4}
\]

\subsection{Recovering an arbitrary power of the short length}

Set \(D=(\theta-\theta_0)\circ\omega(g,h)\circ J_\Psi\). This equals the original value on every \(v\) by (PD4). The difference \(\theta-\theta_0\) has mean zero, so its value on any datum is the mean of its differences under translation by \(z\). Moreover \(\|D\|\le2\|J_\Psi\|\|L\|\), since the Heisenberg translates pulled back through \(\omega(g,h)\) are isometries.

The finite datum belongs to the span of finite Heisenberg translates of the vacuum. Its finite translates are locally constant, so some integer \(m>0\) makes every shift by \(m\widehat{\mathcal O}_E^{nr}\) fix it in the finite coordinates. For \(z=\sum b_j B_j\), reduce each adelic \(b_j\) modulo \(E\) to \(mw_j\), with \(w_j\) finite integral and every archimedean component bounded in absolute value by a field constant. This follows by choosing a bounded archimedean fundamental domain for the integral representatives of \(\mathbb A_E/E\), applied to \(b_j/m\). The rational changes do not affect the translated theta value. After moving the remaining shift back through \(\omega(g,h)\), its finite part acts trivially by the choice of \(m\), and its archimedean part acts by a Fock translation \(T(y)\) with \(\|y\|_\infty\le t=C_0\delta\). Here \(C_0\) accounts for \(m\), the number of generators, the fundamental domain and the fixed Weil scaling. We have proved, for every Fock vector \(w\),\formalnote{10.71}{71}
\[
         |D(w)|\le\sup_{\|y\|_\infty\le t}|D(w)-D(T(y)w)| .       \tag{PD5}
\]

For \(t\le1\), a polynomial times the signed-model Gaussian, denoted \(v_Q\), satisfies
\[
     |D(T(z)v_Q)|\le C_{Q,j}\|D\|\,t^j(1+\|z\|_\infty)^j         \tag{PD6}
\]
for every \(j\ge0\). At \(j=0\) this is the isometry of \(T(z)\). Suppose the bound known at \(j\) for every polynomial. To estimate a difference in (PD5), differentiate \(D(T(sy)T(z)v_Q)\) on \(0\le s\le1\). The derivative replaces \(Q\) inside the two translations by the polynomial
\[
  (z^*J_0 y-y^*J_0z)Q+
     \sum_{i=1}^d\bigl(y_i A_iQ-\overline{y_i\epsilon_i}\,B_i'Q\bigr),
\]
where \(J_0=\operatorname{diag}(\epsilon_i)\), and \(A_i,B_i'\) are the fixed raising and lowering operators in the signed Fock model. The first coefficient comes from commuting the two Heisenberg translations. Its modulus is at most \(2d\,t\|z\|_\infty\), and all other coefficients have modulus at most \(t\). Also \(T(sy)T(z)\) is a unit scalar times \(T(z+sy)\), with \(1+\|z+sy\|_\infty\le2(1+\|z\|_\infty)\). The induction hypothesis on the finitely many resulting polynomials bounds the derivative by \(C\|D\|t^{j+1}(1+\|z\|_\infty)^{j+1}\). Integration from \(0\) to \(1\) gives the same bound on every difference. Equation (PD5) transfers it to the value, proving (PD6).

Use \(Q=1\) in (PD6) and integrate against \(v(z)\,dz/\pi^d\). The Fock reproducing identity expresses \(v\) as this superposition of translated vacua. Truncating at radius \(R\) leaves an error at most
\[
       C_j\|D\|\,t^j
       \int_{\|z\|_\infty>R}(1+\|z\|_\infty)^j |v(z)|\,\frac{dz}{\pi^d},
\]
which tends to zero for \(M_j(v)<\infty\). For \(t\ge1\) the direct bound \(\|D\|M_0(v)\le\|D\|t^j M_j(v)\) suffices. Thus for all short families (PD3) and finite \(j\)-th moments we have proved\formalnote{10.72}{72}
\[
                |\theta(\omega(g,h)J_\Psi v)|
                    \le C_j'\|L\|\delta^j M_j(v).                 \tag{PD7}
\]

\subsection{The powers of the heights and the application to the plane}

Choose \(j=N+A\) in (PD7). Combining (PD2) and (PD3) gives the factor
\[
             (\rho H)^A(2Z_S\rho/H)^{N+A}
                =(2Z_S)^{N+A}\rho^{N+2A}H^{-N}.
\]
For \(H\le2\), use \(|L(J_\Psi v)|\le\|L\|\|J_\Psi\|M_0(v)\), \(H^{N+A}\le2^{N+A}\), and \(\rho\ge1\). This proves (PD1) everywhere. The geometry-of-numbers inputs here are Mahler compactness, weak lattice transference and the trace-dual calculation. The Heisenberg covariance, finite character duality and Fock reproducing identity are the representation-theoretic inputs.\formalnote{10.73}{73}

For the plane period take \(S=W,T=V\). Its fixed \(b\times n\) row block \(y\) and the \(r\times n\) variable \(Y\) have \(b+r=n\). Over the marked completions the multiplier equals the inverse tame character of \(\det M\) when \(M\), formed by stacking and reordering these two blocks, is integral with unit determinant. Elsewhere it is zero. This function is locally constant and compactly supported in the marked variables. On its support let \(P_r\) select the variable rows from \(M\). Then
\[
            Y=P_r M,\qquad Y(M^{-1}P_r^{\mathsf T})=I_r .
\]
This gives even an integral right inverse. The marked set is nonempty, so its product of completions is a nonzero ring as required. At fixed \(g\), (PD1) with \(N=0\) now bounds the smaller theta kernel by \(C_g M_A(v)\), uniformly in \(h\). The real cone estimate from \S\ref{sec:plane-period} supplies an integrable bound for this moment of the sliced vector. Multiplication by the finite volume of \([U(V)]\) proves the absolute integrability used to interchange the cone and group integrals. For \(N>0\) the same argument bounds the integral over \(H_T(h)>R\) by \(\operatorname{vol}[U(V)]\,C_N\rho^{N+2A}R^{-N} M_{N+A}(v)\) for a fixed \(v\).\formalnote{10.74}{74}


\editionchapter{The comparison period and composition of lifts}\label{ch:chain}

Two determinant identities are needed in Chapter~\ref{ch:lower}. For a single lift we compare its norm at two twists by testing both against the same period. In a chain the next map is applied only to the image of the preceding maps. Its Gram form has to be restricted to that subspace. Example~\ref{ex:composite-gram} calculates this restriction in dimension two. We use identities between pairs of lifts and their adjoints to compute the restricted forms.

\section{The common period and the integral comparison operator}\label{sec:common-period}

Fix a step \(j\). As in \S\ref{sec:single-lift-norm}, \(T_\eta\) is the theta lift with twist \(\eta\) from the good space on \(W_j\) to the forms on \(V_j\). The gluing map \(g_\eta\) takes the tuples described below to the forms on \(V_j\). Write \(T^\circ\) for the common period map. The tuples have one Weyl-module value for each choice \(i\) of finite and archimedean data and each class of the definite second factor \(V^\circ\). The seesaw gives \(c'_\eta g_\eta^* T_\eta=T^\circ\). We prove that \(T^\circ\) detects every nonzero source form. We also bound the comparison operator \(A_\eta=N_{\rm Aut}h_j^2 g_\eta g_\eta^*\) at the finite and infinite places, where \(h_j\) is the hook number of the auxiliary weight diagram. For \(y\) in the good space on \(W_j\), we show that \(A_\eta T_\eta y\) is orthogonal to the false lifts.

\subsection{Detecting a form by its common periods}
With the convention that the inner product is linear in the second variable, the common period against a tuple is a sum over \(i\) of nonzero constants times
\[
   \int_{[U(W_j)]} y(g)F_i(g)\Theta^\circ_{\zeta_i}(g)\,dg .
\]
Here \(F_i\) is the integral of the fixed theta lift \(\phi^\circ\) against the theta kernel for datum \(i\) with its original finite test function, as in \S\ref{sec:three-exchanges}. The form \(\phi^\circ\) is the theta lift of \(1\) from the negative definite group of rank \(m\) to the group of signature \((2m,m)\). Both \(F_i\) and \(\phi^\circ\) are independent of the twist. The factor \(\Theta^\circ_{\zeta_i}\) results from pairing the theta kernel on the definite group with the class values \(\zeta_i\) from the tuple. The other slots can be set to zero. Suppose all these integrals vanish. Write \(u\) for the holomorphic weight of \(\bar y\), \(u'\) for the kernel weight and \(m_v\) for the normalized holomorphic matrix coefficient of weight \(v\) (\(m_v(1)=1\)). For any \(g\) and constants \(a_i\) the vanishing implies
\begin{equation}
 \int_{U(W_j)_\infty} y(gh)\left(\sum_i a_iF_i(gh)\right)m_{u'}(h)\,dh=0.
 \label{eq:period-real-test}
\end{equation}
Indeed \(yF_i\) is rapidly decreasing of level \((\pi_M^2)\). Apply the theta reproducing formula to this product with the complementary definite space, whose rank is greater than \(2\dim W_j\). In that formula the real translates of the highest weight polynomial are linear combinations of the basis polynomials in the negative letters. Its finite sum ranges over a definite class and one such polynomial. Each term is the original common period for a tuple supported at that class, with the vector dual to the basis polynomial there. This proves \eqref{eq:period-real-test}. This use of reproduction accounts for the second factor in the products.\formalnote{11.1}{1}

Suppose \(y\ne0\), and call \(g\) a peak if \(|y(g)|>\tfrac12\sup|y|\). Such points exist by boundedness. Put \(s=2ft^3\), \(B_*=b^s>10\), and let \(|h|\) denote the larger of one and the archimedean operator norms used in the matrix coefficient estimates. At every peak the cut-off first minimum of the moved lattice is at least \((5b^{24})^{-1/2}\). To see the mechanism, a shorter vector \(v\) is isotropic by the field norm of its integral self-pairing. The transvection \(1-pvv^*J\), for the purely imaginary level element \(p\), fixes the form. The restriction to \(1-zpvv^*J\) is consequently periodic of period one. Unitary times triangular factorization continues it to \(\Im z>-1\): completion of the square gives a positive block precisely on the half-plane needed under the shortness assumption. The triangular weight factor is bounded there and tends to zero at \(i\infty\) by the positive weight gap. Passing to the disk by \(e^{2\pi i(z+i)}\) and applying Schwarz's lemma gives \(|y(g)|\le \exp(2-2\pi)\sup|y|<\tfrac12\sup|y|\). This contradicts the peak inequality. The bound is \((AA')^{-1/2}\) for weight entries bounded by \(A\) and imaginary level coordinates bounded by \(A'\); here \(A=5b,\ A'=b^{23}\). The same argument for the fixed theta lift \(\phi^\circ\) gives a peak \(g_1\) with minimum at least \((5b^{3/2})^{-1/2}\). The padded space at \((g,1)\) has minimum at least the cut-off minimum at \(g\), since the added standard definite lattice has minimum one. The product of the two minima is at least \(b^{-13}\) in the scale of the construction.\formalnote{11.2}{2}

We next construct a test at these peaks. Let \(P\) be the product of the prescribed minors with the padding columns adjoined, and let \(c_P(h',h)\) be its archimedean oscillator matrix coefficient, with \(h\) embedded in the padded block. The indices of each signed minor are distinct, so its determinant is a nonzero polynomial. The product is nonzero and \(c_P(1,1)\) is its positive Fock norm squared (in the chosen normalization). Unitarity bounds \(c_P\) on the whole real group. Write \(u_1\) for the weight of the kernel used to define the fixed theta lift \(\phi^\circ\) and \(u_3\) for the antiholomorphic weight after the minor is inserted. Set
\[
 F_0(h)=\int \phi^\circ(g_1h')c_P(h',h)m_{u_1}(h')\,dh',
 \qquad J=\int |m_{u_3}(h')|^2\,dh' .
\]
The weight gaps make \(J\) positive and finite. At \(h=1\) compact torus averaging leaves the highest weight coefficient in the fixed theta lift \(\phi^\circ\) and in the padded minor. Their combined character with the kernel weight \(u_1\) gives \(u_3\). Holomorphic reproduction then gives \(F_0(1)=J\phi^\circ(g_1)c_P(1,1)\ne0\). The integrable majorant for \(m_{u_1}\) also proves boundedness of \(F_0\), uniformly in \(h\).

For every peak \(g\) there are constants \(a_i(g)\) such that \(K_g(h)=\sum_i a_i(g)F_i(gh)\) satisfies
\begin{equation}
       |K_g(h)-F_0(h)|\le \delta |h|^{6s},
       \qquad \delta=2B_*^{-19}|F_0(1)| .
 \label{eq:period-approx}
\end{equation}
The construction uses an average over additive characters of the matrix residue ring at \((\pi_M^2)\). By character orthogonality the correlation of a translated finite vacuum with the average is supported on zero residue. The zero rational matrix contributes exactly \(c_P\). For a remaining matrix \(v\), the multiple \((q_1/\pi_M^2)v\) is in the moved Hom lattice. A short basis of the source, as in \S\ref{sec:primitive-decay}, contains a vector on which \(v\ne0\); its image has length at least the target minimum. This lower-bounds the Hom minimum by the product of the minima divided by \(6\operatorname{cond}(E)(ft)^2t^2\). Rescaling gives separation at least \(b^6\) in the original sum.

Here is how the separation enters the error. If \(d\le b^{3/2}\) is the degree of \(P\), \(D\) the complex dimension of the matrix variables, and \(a\) a positive integer, the absolute translated coefficient at length \(l>0\), divided by \(c_P(1,1)\), is at most
\[
  \|\sigma\|^{D+2a}\,[4(d+D+a)]^a l^{-2a},
\]
where \(\sigma\) is the real action on matrices. This follows by dilating the polynomial times the Gaussian in the Fock integral. At least one of the two translated arguments has length comparable to \(l/(1+\|\sigma\|)\); removing its Gaussian tail and applying Cauchy--Schwarz to the remaining factors gives the displayed inverse power. Disjoint balls of half the separation bound the number of points in each successive annulus by the ratio of the volumes of balls in real dimension \(2D\). Taking \(a=2s>D\) bounds the relative error by
\[
 (|h'||h|)^{3a}
       \left(\frac{100(d+2a)}{b^{12}}\right)^a
 \le (|h'||h|)^{6s}B_*^{-20}.
\]
Indeed \(d\le b^{3/2}\), \(s\) has polynomial growth in \(x\), and \(100(d+4s)\le b^2\) at the chosen scale. Replacing \(c_P\) by this theta kernel in the integral for \(F_0\) gives a combination of the \(F_i\): expand the finite residue function in point masses, split the polynomial into its padded and original letters, and expand on the definite classes. Finally \(|\phi^\circ|\le2|\phi^\circ(g_1)|\), and the weighted integral of \(|m_{u_1}(h')||h'|^{6s}\) is at most \(B_*J\). These inequalities give \eqref{eq:period-approx} including the factor 2.\formalnote{11.3}{3}

Each \(K_g\) is holomorphic of weight \(u-u'\), as every \(F_i(g\,\cdot)\) has that weight. Rotate a peak \(g\) by the diagonal compact circle defined by the weight order. Its translates are peaks since the form changes there by a unitary character. At each rotated point \eqref{eq:period-real-test} holds. Replace its \(K_{gk}\) by \(K_g\); by \eqref{eq:period-approx} the cost inside the integral is at most \(2\delta |h|^{6s}\). Remove the rotation character and average over the circle. The holomorphic circle mean of \(\bar y(gkh)\) is \(\bar y(g)m_u(h)\) after the character is removed. Reproduction evaluates the new integral at \(K_g(1)\) times \(y(g)\int|m_u|^2\). The change in the integrand is bounded by \(2|y(g)|\) times the kernel error and \(|m_{u'}(h)|\), and the weighted moment of \(|m_{u'}|\) is at most \(B_*\int|m_u|^2\). Division gives \(|K_g(1)|\le4B_*\delta\). Consequently
\[
 |F_0(1)|\le (4B_*+1)\,2B_*^{-19}|F_0(1)| .
\]
The multiplier is strictly below one, contradicting \(F_0(1)\ne0\). This proves injectivity and hence positive definiteness of \((T^\circ)^*T^\circ\).\formalnote{11.4}{4}

\subsection{An integral frame for the adjoint}
In the standard tableau coordinates let \(F\) be the integral filling matrix, with rows indexed by all increasing-column fillings. The Gram matrix of the Weyl module is \(G=F^tF\). If \(H\) is the matrix of the integral Fock differential pairing on the products of minors, then \(GH=hI\), with \(h=h_j\) the hook number: differentiating the reproducing kernel against a basis minor gives \(h\) times that minor. In particular for each coordinate vector \(e_a\) of the column-filling space the vector \(v_a=HF^t e_a\) is integral, and \(Fv_a=hP_F e_a\), where \(P_F=F G^{-1}F^t\) is orthogonal projection. Thus \(\|v_a\|\le h\), and, for the Weyl inner product,
\[
           \sum_a v_a\langle v_a,z\rangle=h^2 z .
\]
Indeed apply \(F\) to both sides and use \(P_F^2=P_F\). For each definite class \(c\) choose a rational unitary \(\gamma_c\) taking its local representative into the integral gauge at the prime \(\ell\). Transport the integral frame by \(v\mapsto\overline{\rho(\gamma_c^{-1})\bar v}\) and place it in the single slot \((i,c)\). Rational unitary transport preserves the Gram identity even after a complex field automorphism (which permutes the CM places). We obtain tuples \(f_a\) of squared conjugate norms at most \(h^2\), satisfying the same frame identity, with \(\rho(\gamma_c)\overline{f_a(i,c)}\) integral in tableau coordinates.\formalnote{11.5}{5}

Recall that \(N_{\rm Aut}\) is the common multiple of stabilizer orders in \S\ref{sec:regular-twists}. The comparison operator is \(A_\eta=N_{\rm Aut}h^2g_\eta g_\eta^*\), with \(h=h_j\). Fix a prime \(\ell\) prime to the norm of the twist modulus. Applying \(g_\eta\) to the frame identity for \(g_\eta^*z\) gives
\[
        A_\eta z= N_{\rm Aut}\sum_a \langle g_\eta f_a,z\rangle\,g_\eta f_a .
\]
The gluing sum clears the divided splitting by its prescribed power of \(M\); the factors from the twist modulus are units at the prime under consideration. Thus its terms are integral in the class gauges, both at \(\iota\) and at \(\iota\circ{\rm conj}\). For an integral \(z\) the scalar \(N_{\rm Aut}\langle g_\eta f_a,z\rangle\) is integral at \(\iota\): each summand in the class pairing has denominator the stabilizer order, which divides \(N_{\rm Aut}\). The conjugated first input to the pairing is the reason the second embedding is needed. Every summand in \(A_\eta z\) is now integral.

For the trivial twist put \(A=fmb(\log b)^2\), \(s_0=\#\{a\}\), \(n_0=\#\{(i,c)\}\). In the conjugate squared norms gluing costs at most \(e^{2A}\). Explicitly the class sums cost at most \(e^b\) classes on either definite factor and \(N_{\rm Aut}\) to clear the stabilizers. If \(e\) is the divided-splitting exponent, their contribution is bounded by \(M^{2e} e^b(n_0 e^b N_{\rm Aut})^2\le e^A\), with the other \(e^A\) allowing the norms of the first-factor tests. The number of increasing-column fillings and the hook number satisfy \(n_0(s_0N_{\rm Aut}h^2)^2\le e^{2A}\) by bounding every filling column by its list of entries and every hook by the number of boxes. To track the two applications of gluing put \(S=e^{2A}\|z\|_\sigma^2\). Cauchy--Schwarz for the conjugated pairing gives \(|\sigma(N_{\rm Aut}\langle g_1 f_a,z\rangle)|\le N_{\rm Aut}h\sqrt S\) (apply the first-input norm estimate at \({\rm conj}\,\sigma\,{\rm conj}\)). The tuple with these coefficients on the \(f_a\)'s has squared norm at most \(n_0(s_0N_{\rm Aut}h^2)^2S\). Applying gluing gives \(\|A_1z\|_\sigma^2\le e^{6A}\|z\|_\sigma^2\), the required unsquared bound \(e^{3A}\).\formalnote{11.6}{6}

\subsection{Orthogonality to false lifts}
Let \(\Xi\) be a false lift of the specified type. Its projection by the self-adjoint type idempotent is tested against all genuine lifts. Adjointness and orthogonality give zero normalized descent on \(W_j\). The good type is annihilated by the prescribed hyperplane sums at the chosen places, and the tame character is nontrivial at every twist place. These are two separate inputs. At the highest coefficient the normalized descent is a nonzero multiple of the conjugate descent for the harmonic datum. Descent for any Fock polynomial with the same finite block datum is a convolution of this zero function on the real group, hence vanishes. For the finite test it suffices to take indicators of arbitrarily small residue classes about a matrix \(x\). At a marked place write \(x_0\) for its distinguished square block of rows and \(x_1\) for the remaining rows. If \(x_0\) is invertible and \(x_1x_0^{-1}\) is integral at all marked places and is zero modulo each chosen place, averaging the indicator over the level gives a multiple of a source translate of the block datum. Otherwise an integral annihilating row produces shears fixing \(x\) and a sufficiently small class about it. Their sum contributes zero by the hyperplane relation at the chosen places or by the nontrivial character at a twist place. This proves vanishing for all finite tests vacuum outside the marked places and their conjugates. Other weight coefficients follow by translating the highest coefficient.

One must also test a slice on a block of the definite space. Adjoin \(a\) hyperbolic planes to \(W_j\), where \(2r+2a+3<n\) and \(n\) is the definite rank. Induct on the number of planes. The constant term along the new isotropic line is a descent for the smaller member of the tower, hence vanishes; applying this to translates gives vanishing of its slices. The remainder of the theta sum is bounded after reduction, by the rapid decay of the nonzero terms. At the split place specifying the false lift, spherical Hecke transfer gives an annihilating polynomial for the raised descent. Its evaluation point is in the small-root range \(q^{d+1}|X|\le1\) on the group of rank \(d\) (inserting the lift root and the rank inequality gives precisely this range). There the absolute sum of all the coefficients is \(\prod_{j=0}^{d-1}(1+q^j|X|)\) by the finite subspace count. It is at most \(\exp(|X|\sum_{j<d}q^j)\le e^{1/2}<2\). Isolating the coefficient of modulus one and taking the supremum over translates then forces any bounded spherical function killed by the polynomial to vanish. This gives zero throughout the required part of the tower.\formalnote{11.7}{7}

Split the definite target as \(I_a\oplus I_l\). The nonsingular Fourier coefficient for the \(a\) hyperbolic planes unfolds on frames with prescribed Gram matrix. Witt extension makes these frames one rational orbit; its stabilizer is the unitary group of the complementary block. Arbitrary Schwartz tests on the frame orbit therefore give the block period. More explicitly, multiply the coefficient of \(\Xi\) by a theta kernel on the original \(W_j\); tensor product of the two data gives a descent in the raised tower and hence zero. Choose a test on the first block whose theta kernel has a nonzero value at the prescribed point (a translated vacuum and finite test isolate such a value). Factorization of kernels on \(I_a\oplus I_l\) then leaves zero for the complementary block integral. On the definite quotient the latter is a finite sum over classes with their stabilizer weights. Equivariance lets one take any polynomial coordinate and transport the level by the divided splitting. These are exactly the slice pairings \(\Theta^\circ_{(g_\eta^*\Xi)_i}\). They vanish for each \(i\), so the common-period identity gives \(\langle \Xi,g_\eta T^\circ y\rangle=0\). Substituting \(T^\circ=c'_\eta g_\eta^*T_\eta\) proves the claimed orthogonality of \(A_\eta T_\eta y\).\formalnote{11.8}{8}

\subsection{Cancellation at the empty row}
For a main critical pair \(p=(\nu,z)\), let \(\phi_p\) be the form obtained by lifting the corresponding character at the trivial auxiliary twist. We seek a multiplier taking \(\|\phi_p\|^2\) to an algebraic integer, with logarithmic absolute value at most \(B-J\) at every complex conjugate. Here \(B\) is the denominator budget and \(J\) is the sum of the factorial and CM-period terms in \S\ref{sec:lower-C}. The term \(-J\) comes from comparing \(\phi_p\) with a lift of the same rank and archimedean polynomial of degree zero. On a definite space of the same dimension \(2z\) take \(p_{\varnothing}=(\upsilon^{-z},z)\), with trivial tame part and negative weight \(z\) at every place. Its rank twist is trivial and every row degree is zero. For any embedding \(\sigma\) fixing the imaginary quadratic field, the real pairing and the transformation of the elliptic period give, with \(d_\tau(p)=\operatorname{negWeight}_p(\tau)-z\),
\[
 \log |\sigma(\|\phi_p\|^2)|
  = c_{\sigma,2z}+r_p+\sum_\tau\operatorname{jetSize}(d_\tau(p)).
\]
The finite vacuum constant, the transformed elliptic multiplier at row zero and the common Haar factor depend only on the rank and embedding. For a row of length \(d\) the hook contributes \(d!\) and the integral elliptic normalization supplies \((M/\Omega)^d\), where \(M=p_0p_1\) and \(\Omega\) is the CM period. Thus the jet size is \(\log(d!)+d\log(M/\Omega)\). The remaining archimedean representation dimension is at least one and equals one on the empty row. At \(z\ge2\) comparison of the Euler factors with \(1\pm (Nv)^{-z}\) bounds the log absolute partial value between \(-4f\) and \(4f\). Thus \(r_p\ge-4f\) and \(r_{p_{\varnothing}}\le4f\). The same \(c_{\sigma,2z}\) occurs in both formulas.

For the empty row the lift denominator is \(M^{4fz}\). At the
trivial twist the conductor has norm one, so the integral algebraic
part is \(a_0=N_{\rm Aut}M^{8fz}\|\phi_{p_{\varnothing}}\|^2\).
The class pairing proves its integrality at every prime, and
\(\log(N_{\rm Aut}M^{8fz})\le x^{33}\). For a main pair at step \(j\), the bounds \(c(\tau)\ge w_j(\tau)+r-t\) and \(z\le r\) from \S\ref{sec:regular-twists} give \(d_\tau(p)=c(\tau)-z\ge w_j(\tau)-t\). Since \(w_j(\tau)-\kappa_\tau=(2j+1)r^8\ge t\), the row length \(d_\tau(p)\) is at least the place weight \(\kappa_\tau\). Also \(M/\Omega>1\) by the scale bounds in \S\ref{sec:three-estimates}, so \(\operatorname{jetSize}(d)\) increases with \(d\). The sum over the rows of \(p\) is therefore at least
\(J=\sum_\tau\operatorname{jetSize}(\kappa_\tau)\) from \S\ref{sec:lower-C}.
Subtracting the two logarithms gives
\[
   \log|\sigma(a_0/\|\phi_p\|^2)|\le x^{33}+8f-J\le B-J.
\]
For an automorphism acting nontrivially on the imaginary quadratic field, composition on the left with complex conjugation preserves the absolute value and fixes that field. Multiplication by the displayed quotient takes the constant family \(\|\phi_p\|^2\) to the integral \(a_0\). This proves the height inequality used in the one-step determinant.\formalnote{11.9}{9}



\section{The Gram determinant of the theta chain}
\label{sec:chain-squares}
For a single cuspidal lift, the Rallis formula expresses an inner product in terms of a standard \(L\)-function and local zeta integrals. Gan, Qiu and Takeda give the formula in the second-term range under their first-occurrence hypothesis \cite[Theorem~1.3]{gan-qiu-takeda-rallis}. In the calculation here, exchanging adjacent lifts will reduce the chain to pairings in the convergent range. For a composite \(BA\) the Gram
operator is \(A^*B^*BA\). Thus the quadratic form of \(B^*B\) must
be known on the image of \(A\).

\begin{example}[The Gram determinant of a composite]\label{ex:composite-gram}
Give \(\mathbb R\) and \(\mathbb R^2\) their Euclidean inner products.
Let \(A(x)=(x,x)\) and \(B(u,v)=(2u,3v)\). On the full intermediate
space the Gram determinant of \(B\) is \(4\cdot9=36\), while the Gram
determinant of \(A\) is \(2\). The composite sends \(1\) to \((2,3)\),
so its Gram determinant is \(13\). Indeed the restriction of the
quadratic form \(4u^2+9v^2\) to the unit direction
\((1,1)/\sqrt2\) gives \(13/2\), and \(A\) contributes the other
factor of \(2\). The map \(B'(u,v)=(u,6v)\) also has Gram
determinant \(36\). Its composite sends \(1\) to \((1,6)\), of
squared length \(37\). The same individual determinants have thus
given two different composite determinants.
\end{example}

In the theta chain the intermediate lifts may also leave the range
of convergence of the single-lift Rallis integral. We arrange the lifts by the number \(h\) of blocks already adjoined. A lift that adds the first block starts from the fixed base rank, where the integral converges. Exchanging adjacent lifts and passing to adjoints expresses the required inner products in terms of these lifts from the base rank. This gives one scalar for each pair
of distinct lift indices.

Fix the twist \(\xi\). Write \(R\) for the product of the residue
fields at the chosen and twisting places. At position \((i,h)\),
\(i+h\le m\), the hermitian space consists of a fixed block of rank
\(b_0\), called the stem, and the \(h\) blocks \(i+1,\ldots,i+h\).
Its rank is \(b_h=b_0+ht\). Write \(F_{i,h}\) for the space of
automorphic forms at this position. At \(h=0\), changing \(i\)
preserves the stem up to negation. Identifying the unitary groups
of a form and its negative gives isometric copies of the base
space, with embeddings \(j_i\).
Let \(a_{i,h}:F_{i,h}\to F_{i,h+1}\) add a block on the right,
using the delta mass at the standard inclusion as finite test function.
Set \(w_h=|\mathrm{GL}_{b_h}(R)|\). Multiplication by \(w_h\)
changes the source measure from volume one for the integral group to
volume one for its principal congruence subgroup. Write \(Q_\xi\) for the auxiliary twisting modulus and \(D=\operatorname{Nm}Q_\xi\) for its norm.
For a scalar function \(c=c(\xi)\) and a nonnegative function \(f\),
write \(c\sim f\) when \(|c|^2=C D^k f^2\) for some
twist-independent \(C\ge0\) and \(k\in\mathbb Z\). These relations
multiply. The injectivity argument below will establish the
nonvanishing needed for division.

\subsection{The two paths round one square}
The four corners are the source \(U_1\) and target \(U_2\) of ranks \(b_{h+1}\), the large intermediate space \(A\) of rank \(b_{h+2}\), and the small intermediate space \(A'\) of rank \(b_h\). The first path appends the last block and then removes the first. The second path removes the first and then appends the last. Retain the finite matrix on each edge as a variable by taking its delta mass as test function. Test the two resulting maps on the same input vector and against the same output vector. This gives functions \(L(X)\) and \(L'(Y)\), where \(X\) has \(b_{h+2}\) rows, \(Y\) has \(b_h\) columns, and their common ambient space has rank \(2b_{h+1}\).

We apply the complementary-partner exchange of Section~\ref{sec:one-exchange} on \(U_1\oplus U_2\). For completeness, its dimension and weight tests here are the following. The ranks of \(A,A'\) sum to \(2b_{h+1}\) and differ by \(2t\). Their negative-sign counts sum to the number \(l\) of negatives on \(U_1\oplus U_2\), by counting the alternating signs of the same stem entries on the two diagonals. Added blocks are positive, and each node has at most \(b_0\) negative entries. As \(b_0<t\), this gives \(l<b_{h+1}\). The absolute negative weight of the quotient of the lifted characters is \(\mu=(h+1)\operatorname{gap}\), which is at least \(b_{h+1}\) by the choice of the weight gap. Hence \(2l+2\le2b_{h+1}\le 2\mu+2l\). The quotient character for indices \(u=i+1,\ v=i+h+2\) is \(\nu=\chi_v/\chi_u\), since the untwisted characters form a geometric progression. The product of splitting characters carries precisely the rank twist of \(\nu\xi^{v-u}\). Its inverse has nontrivial tame part at each chosen or twisting place by the nontrivial-power condition through \(m\). Denote its residue character by \(\omega\).

Splitting a finite delta mass on \(U_1\oplus U_2\) into its two blocks splits the theta sum into the product of its two edge kernels. Inserting the two test forms and integrating first in \(U_1,U_2\) yields the two paths just described. The common determinant twist from this splitting is the quotient of the splitting characters at the intermediate group and is the same on both paths. The cusp forms have rapid decay, and the finite functions in a matched pair are supported on split injections and surjections, so the moment bounds of the exchange justify the change of integration order. On the small partner use the product of the two prescribed real polynomial vectors. On the large partner first take a polynomial vector. Expand it into products of polynomials on the two blocks. On each product the lift of each factor is an archimedean convolution of the prescribed edge lift. Schur orthogonality for the specified type on their common group makes the integrated product of these two convolutions a scalar multiple of the original pairing. Add over the polynomial terms. The scalars involve only the real data. Approximation of the exchanged Schwartz vector by polynomial vectors in the moment norms extends this assertion to it, since the one-dimensional span of the family of original pairings is closed in the topology of pointwise convergence. Denote the scalar by \(k_0\). Thus for all pairs \(\Phi,\Psi\) matched by \(\omega\),
\[
 \sum_Y\Psi(Y)L'(Y)=s\sum_X\Phi(X)L(X),\qquad
 s=k_0/c,\qquad |c|^2=C\,\mathcal L_S((\nu\xi^{v-u})^{-1},t)>0.
\]
Here \(\mathcal L_S\) is the product of \(t\) squared incomplete Hecke values and the residue norm to power \(t^2\) from Section~\ref{sec:three-exchanges}. Matching means that for an invertible \(g\) of the ambient space the functions at the last rows of \(g^{-1}\) and the complementary first columns of \(g\) differ by \(\omega(\det g)\), with both functions supported on these frames. Changing the frame in a path gives \(L(kX)=\omega(\det k)L(X)\) and \(L'(Yk)=\omega(\det k)L'(Y)\).\formalnote{11.10}{10}

To evaluate the path identity on the prescribed frames, temporarily write \(b,a\) for their column and row sizes, in the ambient order \(b,a\). For an \(a\times b\) matrix \(Z\), use the matched pair
\[
 \Phi(X_1,X_2)=\omega^{-1}(\det X_2)\mathbf1_{X_2^{-1}X_1=Z},\quad
 \Psi\binom{Y_1}{Y_2}=\omega^{-1}(\det Y_1)\mathbf1_{-Y_2Y_1^{-1}=Z}.
\]
A singular pivot gives value zero. The columns of \(\binom{1}{-Z}\) form a basis of the kernel of \((Z,1)\). More explicitly, for the matched blocks of \(g\) we have \(X_1Y_1+X_2Y_2=0\) and the complementary-minor identity \(\det Y_1=\det g\det X_2\). Thus the pair is matched, including at singular pivots. Summation over the two possible changes of frame gives
\[
              |\mathrm{GL}_b(R)|\,L'\binom{1}{-Z}
              =s\,|\mathrm{GL}_a(R)|\,L(Z,1).
\]
Define \(p^+\) on block matrices by
\[
 p^+\begin{pmatrix}\eta&B\\0&1_t\end{pmatrix}
       =p(-\eta)/\#\{B\},\qquad p^+=0\ \text{elsewhere}.
\]
Set \(Z=\binom{\eta}{0}\), permute the blocks back to the order of the edges, and left-multiply \(X\) by the shear with upper right block \(B\). Its determinant is one. Average over \(B\) and sum over \(\eta\) against \(p(-\eta)\). This gives the first square identity for a descent \(d_{i,h}:F_{i,h+1}\to F_{i+1,h}\) using datum \(p\), with \(d_{i,h+1}\) using \(p^+\):
\[
                 w_h a_{i+1,h}d_{i,h}
                 =s_{i,h} w_{h+2}d_{i,h+1}a_{i,h+1}.                 \tag{CS1}
\]
The datum for \(d_{i,0}\) is chosen by taking the adjoint of the first edge on the base. Iterating \(p\mapsto p^+\) defines the descents on the entire row. If \(s_{i,h}\ne0\), set \(\beta_{i,h}=1/\bar s_{i,h}\). Its relation is \(\beta_{i,h}\sim R_{u v}\). Indeed the squared modulus of \(c\) above, divided by its trivial-twist value, gives precisely the product of the \(t\) squared critical-value ratios \(R_{uv}^2\) times a fixed power of \(D\); all factors from the chosen places are fixed. Here and below \(R_{uv}\) has the same orientation as in the determinant formula.\formalnote{11.11}{11}

For the adjoint identity reverse the left edges. Reversal conjugates the archimedean polynomial and takes its signed transpose. At a marked place it replaces \(p\) by the finite Fourier transform of \(\bar p\), by Poisson summation interchanging the two polarizations. Its scalar consists of fixed real and measure factors, a power of \(\#R\), and a central unitary translation. The latter commutes through the right edge up to the value of its splitting character, of modulus one. In this square the two intermediate ranks are equal, so the preceding complementary-partner calculation has \(t=0\) in its value product and contributes no Hecke values.

We detail the only extra finite factor. With the unnormalized matrix Fourier transform \(\widehat p(X)=\sum_Yp(Y)\psi(\operatorname{tr}XY)\) for a primitive additive character \(\psi\), character orthogonality in \(B\) gives
\[
  \widehat{p^+}\begin{pmatrix}A&B'\\C&D'\end{pmatrix}
        =\mathbf1_{C=0}\widehat p(-A)\psi(\operatorname{tr}D').
\]
When summed against a left \(\omega\circ\det\)-equivariant function, integration over the \(t\times t\) block \(D'\) replaces it by \(1_t\) at the cost of a Gauss sum \(G\), with \(|G|^2=|R|^{t^2}\). Indeed on each field factor a singular block has a stabilizer with arbitrary determinant on a complementary line, whereas \(\omega\) is nontrivial on that factor. Thus an equivariant function vanishes at singular matrices and on the units equals its value at \(1_t\) times \(\omega(\det)\). The same reasoning for its Fourier transform gives \(\omega^{-1}(\det)\) there. Parseval for the unnormalized transform gives the claimed squared modulus. The free upper block contributes \(|R|^{b_h t}\), so reversing the order of Fourier transformation and \(p\mapsto p^+\) costs an inverse scalar of squared modulus \(|R|^{-t(t+2b_h)}\). As the chosen primes are fixed, \(|R|\) is a constant times \(D\). The resulting adjoint identity, with the pairing linear in its second variable, is
\[
 \langle d_{i,h}a_{i,h+1}^*z,y\rangle
       =k_{i,h}\langle d_{i,h+1}z,a_{i+1,h}y\rangle,\qquad k_{i,h}\sim1 .
                                                                    \tag{CS2}
\]
Here \(a^*z\) denotes the adjoint vector obtained by transposing the kernel against \(z\); central translations and the reversal scalars have been included in \(k_{i,h}\).\formalnote{11.12}{12}

\subsection{Peeling one block off the Gram form}
Write \(X_{i,0}=1\), \(X_{i,h+1}=w_h a_{i,h}X_{i,h}\), and let
\(j_{i+1}B_i x=d_{i,0}a_{i,0}j_i x\).
To see this last definition is possible, pair the right side with an arbitrary base form transported by negation. The first-edge adjoint gives the Gram pairing of \(a_{i,0}\). The base level preserves that pairing by its action on the finite point datum, so the represented vector has the base level. For the first-edge lift \(a_{i,0}\), the ranks satisfy \(2b_0<b_1\), the point datum is sharp for a level fixing \(j_i x\), and the real vector is the joint harmonic for the two prescribed types. Hence the injectivity argument of \S\ref{sec:one-exchange} applies in the signature of this edge. It follows that \(B_i\) has zero kernel, since \(\langle B_i x,x\rangle=\|a_{i,0}j_i x\|^2\).\formalnote{11.13}{13}

Applying (CS1) successively gives, on the base,
\[
 X_{i+1,h}j_{i+1}B_i
  =\frac{w_{h+1}}{w_1w_0}
       \left(\prod_{\ell=0}^{h-1}s_{i,\ell}\right)
                  d_{i,h}X_{i,h+1}j_i .                                  \tag{CS3}
\]
The empty product gives the defining identity for \(B_i\). To pass to \(h+1\), multiply by \(w_h a_{i+1,h}\), use (CS1), and replace \(w_{h+1}a_{i,h+1}X_{i,h+1}\) by \(X_{i,h+2}\). If \(X_{i,h+1}j_ix=0\), (CS3) gives \(X_{i+1,h}j_{i+1}B_ix=0\). Induction and injectivity of \(B_i\) imply \(x=0\). Thus all chains on the base are injective. For a nonzero base the left side of (CS3) is injective, so its scalar product never vanishes. All occurring \(s_{i,\ell}\) are consequently nonzero.

The first-edge adjoint is \(\langle z,a_{i,0}j_i y\rangle=\langle d_{i,0}z,j_{i+1}y\rangle\). Iterating (CS2) gives
\[
 \langle z,a_{i,h}X_{i,h}j_i y\rangle
   =K_{i,h}\langle d_{i,h}z,X_{i+1,h}j_{i+1}y\rangle,\qquad
 K_{i,h}=\prod_{\ell<h} k_{i,\ell}\sim1.
\]
For the induction move the last \(a\) across the pairing, apply the preceding identity to \(a^*z\), and then use (CS2). Write \(G_{i,h}(x,y)=\langle X_{i,h}j_ix,X_{i,h}j_i y\rangle\). Apply the last display with \(z=X_{i,h+1}j_i x\), insert (CS3) into its first argument, and conjugate that scalar. One obtains
\[
 G_{i,h+1}(x,y)
  =K_{i,h}\Big(\prod_{\ell<h}\beta_{i,\ell}\Big)
       \frac{w_0w_1w_h}{w_{h+1}}\,G_{i+1,h}(B_i x,y).                    \tag{CS4}
\]
This exhibits both the conjugation in \(\beta\) and the two end weights.

Write \(V\) for the selected target space in the base. The single-edge Rallis formula on each row gives, for \(y\in V\) and any base \(x\),
\[
 w_0\langle B_i x,y\rangle=C_i\langle x,S_i y\rangle,\qquad
 S_i=\frac{\mathcal R_{i+1}(\xi^{\,i+1})}{|R|^{b_0 b_1}},
\]
where \(C_i\) is independent of \(\xi\). Induct on \(h=m-i\) in (CS4). If the next row contributes \(w_0w_1^{h-1}/w_{h-1}\), multiplication by \(w_0w_1w_{h-1}/w_h\) and the use of the first-edge identity, with its factor \(1/w_0\), leave \(w_0w_1^h/w_h\). At row \(i\) the new \(\beta\)'s give \(R_{i+1,v}\) for \(i+1<v\le m\), each once. The other scalars satisfy \(\sim1\). The length-zero form is the original inner product. Thus\formalnote{11.14}{14}
\[
 G_{0,m}(x,y)=K\,\frac{w_0w_1^m}{w_m}
               \langle x,S_0 S_1\cdots S_{m-1}y\rangle,\qquad
 K\sim\prod_{1\le u<v\le m}R_{uv}.                                      \tag{CS5}
\]

\subsection{The last level and the determinant}
The norm on forms at the last level is \(G_{0,m}\) times the level index and the fixed dimension of the archimedean weight space. At each twist prime of residue order \(q\), the flag stabilizer for \(j\) blocks has order
\[
 |\mathrm{GL}_{b_0}(q)|\,|\mathrm{GL}_{t}(q)|^j\,
             q^{b_0tj+t^2\binom j2}.
\]
Indeed its diagonal blocks are invertible and all entries in the strictly upper blocks are free. Taking the quotient of \(|\mathrm{GL}_{b_j}(q)|\) by this order and comparing the last flag with \(m\) copies of the first cancels \(w_0w_1^m/w_m\) from (CS5), up to a fixed constant, an integral power of \(D\), and \(m\) copies of the first-edge normalization. More explicitly, let \(\delta\) be the first-edge Gram factor (the first flag index times the lower-rank group order and the measure power, divided by the fixed chosen-place group order). The last index is \(C D^k\delta^m w_m/(w_0w_1^m)\); this follows place by place by the flag count. Dividing \(\delta\) by \(|R|^{b_0b_1}\) gives the fixed chosen-place factor times
\[
       c_Q=D^{rt}\prod_{w\mid Q}\prod_{t<j\le t+r}(1-q_w^{-j}).
\]
Indeed \(r=b_0\). The first flag index times \(|\mathrm{GL}_r(q)|\) counts full-rank \((r+t)\times r\) matrices. Dividing by \(q^{(r+t)r}\) gives their proportion \(\prod_{j=t+1}^{t+r}(1-q^{-j})\). The measure factor is \(q^{rt}\) at each twist prime, as asserted. The operators in (CS5) consequently give
\[
  \Theta_\xi^*\Theta_\xi=c(\xi)c_Q^m\prod_{j=1}^m\mathcal R_j(\xi^j),
       \qquad c\sim\prod_{u<v}R_{uv}.
\]
On the selected space of positive dimension, chain injectivity gives \(c(1)\ne0\). The critical-value ratios are nonnegative and equal one at the trivial twist. Hence
\(
 |c(\xi)/c(1)|=D^{e/2}\prod_{u<v}R_{uv}(\xi)
\)
for a fixed \(e\in\mathbb Z\).

Write \(Z_j(\xi^j)\) for the Gram determinant of the \(j\)-th single lift and \(Z(\xi)\) for the Gram determinant of the composite theta lift \(\Theta_\xi\). With \(\mathcal R_j(\xi^j)\) the Hecke series transferred to the base space \(V\), the multiplicity calculation in \S\ref{sec:chain-step-norms} gives \(Z_j(\xi^j)=(a_jc_Q)^{d_T}\det_V\mathcal R_j(\xi^j)\), where \(a_j\) is independent of the twist. The same \(c_Q\) occurs because all the powers through \(m\) have the same ramified places. Take determinants, cancel \(c_Q^{md_T}\) and divide by the corresponding equality at the trivial twist. The Gram determinants are positive, so taking absolute values gives
\[
 \frac{Z(\xi)}{Z(1)}
    =\left(D^{e/2}\prod_{u<v}R_{uv}(\xi)\right)^{d_T}
                   \prod_{j=1}^m\frac{Z_j(\xi^j)}{Z_j(1)}.
\]
If the selected space has dimension zero, every determinant here is the empty determinant and the same identity holds.\formalnote{11.15}{15}


\section{Simultaneous basis completion for the twisted chains}\label{sec:grid-completion}
The product formula in \S\ref{sec:lower-C} requires one completed basis of the base space for all primes and permitted twists. We therefore construct a single spanning set on that space. Any given independent family can then be extended by selecting vectors from this set. The norm and integrality bounds for the added vectors must hold at all the required twists. Put \(\mathfrak b=(\pi_M^b)\). The two arrays of spaces used for this purpose are \(Y^i_j\) and \(\widetilde Y^j_i\); along their common diagonal the latter has one more orthogonal line. A path pairs forms at its two ends by integrating the kernels on its edges. Initially the two halves run along the first columns. Finally they run along the last rows, where their integrals are sums on definite groups. Moving from the first path to the last interchanges every horizontal step with every vertical step once.

\paragraph{The restriction at the first path.}
The transpose of the untwisted auxiliary chain maps onto its deep-level space. Indeed the second triangle starts with the stem plus one line, adds blocks of length \(t-1\), and has the same first-edge injectivity and forward square identity (CS1) as in Section~\ref{sec:chain-squares}, at its deep level. Induction using (CS3) proves injectivity of its full weighted chain. If a form is orthogonal to the transposed image, pair against the conjugate of its own weighted lift. Transposing the kernel identifies that pairing with the squared norm of the lift. It vanishes, forcing the form to vanish. Finite dimensionality gives the asserted surjectivity. For the next operation choose a single archimedean translate \(\omega\), and take the restriction to \(Y_0\) followed by the type projection. To see that this is onto, suppose \(y\) is orthogonal to all these restrictions and let \(S\) bound \(|y|\). If the first minimum is small, a short isotropic line gives a unipotent integration with vanishing constant term. Subtracting the translated kernel from the original kernel and integrating its derivative bounds \(|y|\) by \(S/2\) in this range. At larger first minimum use instead a Poincaré series on \(Y_0\oplus\langle1\rangle\) of principal level \(\mathfrak b\), with translated discrete-series coefficient at infinity. Its period against \(y\) vanishes by assumption. Unfold the period by the possible rational last columns. The identity column gives the convolution of \(y\) with the translated coefficient. The remaining columns agree with the distinguished one modulo \(\mathfrak b\). They are separated in archimedean norm by the level and the lower bound for the minimum. Summing the decaying coefficient over shells bounds the convolution by the tail. Convolve once more with the projector of \(y\)'s own type. The rank-one coefficient identity gives back \(\eta y\) where \(\eta\) is the product of the local pairings with the translated coefficient. At each real place the two lists of lowest compact weights have opposite signs outside their common zero block. With the alternating signs of the Hermitian form, polynomial interpolation realizes the smaller weight matrix as the compression of the larger one to the orthogonal complement of a vector of Hermitian square \(1\). This gives the local components of \(\omega\). The local pairing estimates give \(|\eta|>0\), and the level bound makes the tail after convolution at most \(S|\eta|/2\). Hence \(|y|\le S/2\). Since every uniform bound is halved, \(y=0\). The composite \(r\) of the transposed chain and the translated restriction is therefore onto the principal forms on \(Y_0\). The choice of \(\omega\) and the map \(r\) are independent of the twist.\formalnote{11.16}{16}

\paragraph{Following the finite data around the grid.}
Use residue matrices at \(\mathfrak b Q_\xi\) on the edges. A change of basis \(u\) at the nodes sends a matrix on \(p\to q\) to \(u_q^{-1}S_{pq}u_p\). For equivariance the changes of basis lie in the levels at the two initial nodes, and on a diagonal node they act by \(u\oplus1\) on the side with the added line. Under such changes a path datum transforms by the product of the edge characters and is supported where the vertical matrices are split injections. Starting from the first column imposes these conditions by orbit summation. At an off-diagonal turn, fix the two exterior matrices, sum over the intermediate one, and use the one-square exchange of \S\ref{sec:one-exchange}. Reindexing by the new intermediate basis gives equivariance for the new path and the same prescribed injection condition. On the diagonal the extra line must first be reversed. Move its kernel to the common space by the seesaw formula, Fourier transform the column belonging to that line, apply the square exchange, and move the reversed line to the other pair of edges by the same seesaw. Orthogonality of additive characters gives the inverse transform and its cardinality factor. For each turn its twist-dependent denominator is nonzero. More precisely, if the new period is proportional to a family \(\Xi\), the old period is proportional to \(\Xi\) divided by that denominator. The constant of proportionality may depend on the real test point, but is independent of the twist and of both end forms. At most the number of vectors in the line transform multiplies the allowed number of roots of unity in each coefficient. Induction on the remaining inversions of the path yields one scalar \(c_\xi\), with one factor for each of the \(m^2\) squares. The exchange formula gives \(|c_\xi/c_1|=(NQ_\xi)^{e/2}\prod_{0\le i,j<m}R^{\rm turn}_{ij}(\xi)\) for a fixed integer \(e\). The final datum has coefficients that are sums of at most \(V^{m^2}U^{2m}\) roots of unity, where \(V\) and \(U\) bound, respectively, the number of vectors in a turn and the number of matrices on an edge.\formalnote{11.17}{17}

\paragraph{An integer independent of the twist.}
Let \(\Xi_\xi\) be the last-path sum. For each tuple of residue matrices one lifts each of two input vectors \(m\) times, takes the branching pairing at their common end, and sums over all the intermediate classes. Write \(n_m\) for the last chain rank. One lift is integral away from \(NQ_\xi\) after multiplication by
\[
 D_0=(M^b)^{2f(2n_m)^2}(M^b d)^{br^4}.
\]
Indeed \(NQ_\xi M^b\) belongs to the modulus; it clears the translated finite datum on the \(2f(2n_m)^2\) scalar entries. The second factor clears the coordinates in each of at most \(br^4\) boxes in the polynomial. The extra determinant from reversing the line may have either sign of exponent. If \(t_0\) generates the modulus and \(\alpha\) is the fixed finite scaling idele, the flip scalar is \(\overline{t_0\alpha}/(t_0\alpha)\). Multiplication by \(NQ_\xi M^b q_1\) makes that scalar and its inverse integral, since \(q_1\alpha\) and \(\alpha^{-1}\) are integral. If \(L_\tau\) is the branching determinant exponent at \(\tau\), one can take \(D_1=\prod_\tau(M^bq_1)^{m|L_\tau|}\). Indeed the line coordinate is the inverse \(m\)-th power of the flip scalar and \(NQ_\xi\) is a unit at the primes under consideration. The bounds \(|L_\tau|\le br^4+\lceil20b\rceil\) give \(\log D_1\le x^{98}b^2\). Each half uses \(D_0^m\). The finite path sum is an integral combination of such pairings (the level indices supply the class weights). Consequently \(N_*=D_1D_0^{2m}\) clears \(\Xi_\xi\) for all integral starting values and \(\log N_*\le x^{99}b^2\). For the archimedean estimate at trivial twist take \(G=\exp(x^{99}b^2)\). Both the norm cost of one edge kernel and the number of definite classes at a node are at most \(G\). The two lifted vectors have norms at most \(G^m\) times their original norms. There are \(2m+1\) class sums. Each coefficient contributes at most \(V^{m^2}U^{2m}\) roots of unity, and the \(2m\) residue matrices can themselves be chosen in at most \(U^{2m}\) ways. Thus before the class sums there are at most \(V^{m^2}U^{4m}\) terms. Including the norm of the branching pairing gives
\[
 |\sigma\Xi_1|\le\exp(x^{200}b^2)\sqrt{B_-B_+},
\]
where \(B_\pm\) bound the two starting squared norms.\formalnote{11.18}{18}

\paragraph{Recovering the vector.}
In the period take a delta form at the last class \(x_0\), with test vector \(v\) in the dual Weyl module, and the deep-level input \(\phi\) on the other side. The sum at that class collapses. In the integral dual pairing \(B(v,z)\), for a constant \(\beta\) independent of \(\xi,x_0,v,\phi\), this reads
\[
 B\bigl(v,(c_\xi/c_1)\Theta_\xi(r\phi)(x_0)\bigr)
       =\beta\,\Xi_\xi(\phi,x_0,v).
\]
The pairing matrix has an integer left multiple equal to the dual hook number \(h^\vee\) times the identity, by the Fock pairing identity. Choose one nonzero \(s\) with \(s\beta=h^\vee N_*\) if \(\beta\ne0\); if \(\beta=0\) take \(s=1\). Testing with the rows of that integer matrix recovers the integral coordinates of \((c_\xi/c_1)\Theta_\xi(sr\phi)\). At an infinite embedding use the norm dual vector instead. The preceding bound, \(\log N_*\le x^{99}b^2\), and the hook bound yield squared norm at most \(\exp(x^{202}b^2)\) for short \(\phi\) with initial squared norm at most \(G\). The principal level here lies in the chosen level and has trivial rational class stabilizers. At each class choose short integral bases for the representative lattices, and apply their invertible Schur and determinant matrices to the Weyl evaluation vectors with roots of unity as entries. The resulting delta forms are integral in the class gauge. The short-basis bound \(C=6\operatorname{cond}(E)(f n)^2\) in rank \(n\) bounds the squared conjugate norms by \(C^{6H}\), where \(H\) is the number of boxes in columns of height below \(n\) plus \(fn\). The size conditions give \(C^{6H}\le G\). These evaluations span because a polynomial vanishing on all root-of-unity tuples vanishes identically, by the one-variable root bound successively in each variable. Thus these short forms span the initial level space. Their images \(sr\phi\) span since \(s\ne0\) and \(r\) is onto. Applying the surjective multiple of the good projection to this family preserves integrality at the allowed primes by its level normalization, and costs at most \(x^{99}\) in log squared norm. As \(x^{99}+x^{202}b^2\le2x_1^2b^2\), this gives the simultaneous completion with precisely the cost charged in Theorem C.\formalnote{11.19}{19}


\editionchapter{Auxiliary characters and global Selmer bounds}\label{ch:selmer}

Fix a coefficient prime \(\ell\) and a field isomorphism \(\iota_\ell:\mathbb C\simeq\overline{\mathbb Q}_\ell\). We seek one finite-order character \(\xi\) such that each of \(\xi,\ldots,\xi^m\) satisfies two bounds. The first controls the lengths of the Selmer groups that enter the projection estimate. The second controls how much twisting changes the critical Hecke \(L\)-values in the norm formula. The product formula in Chapter~\ref{ch:lower} allows a separate choice of \(\xi\) at each such isomorphism \(\iota_\ell\). The \rootlink{grh/src/New/Congruences/FitTwistsExist.lean}{simultaneous twisting theorem} bounds the total cost over the coefficient primes.

\section{The group-ring elements for the critical Hecke values}
\label{sec:critical-values}
The interpolation of twisted Hecke \(L\)-values by integral group-ring elements, together with their norm relations, appears among the \rootlink{grh/INDEX.md\#known-results}{known results in the source index}. For an ordinary CM type, Katz constructs a measure on a ray class group and gives a formula for its moments in terms of Hecke values and periods \cite[\S5.5, Theorem~5.5.8, and (5.8.10)]{katz-cm-lfunctions}. Chida and Hsieh construct group-ring elements for the anticyclotomic twists of a modular form over \(\mathbb Q\) after base change to an imaginary quadratic field \cite[Definition~4.1 and Theorem~4.6]{chida-hsieh-theta}. The present construction varies the tame modulus away from \(\ell\). We specify how adding a prime changes the Euler factors. At the critical integer \(1\), the norm relation has two factors for each of the exponents \(0\) and \(1\), indexed by \(\epsilon=\pm1\) in \textup{(CV2)}.

We need a lower bound for certain twisted critical values in
\(\ell\)-adic absolute value. An analytic Euler product gives no such
bound directly. Instead, the norm identity expresses the critical-value
ratio in terms of the squared norm of a theta lift. At fixed level each
character enters that norm through a finite sum. This gives a group-ring
element whose specialization at a character can be controlled by the
polynomial argument of \S\ref{sec:specialization}.

Fix one of the ramified critical pairs \(p=(\nu,z)\) of \S\ref{sec:three-exchanges}. Here \(\nu\) is the character and \(z\) the critical integer at which the \(L\)-function is evaluated. Fix also a place \(w_*\) chosen in the construction away from \(\ell\) and a sign \(\delta=\pm1\). Write \(Q\) for an auxiliary modulus disjoint from its conjugate, and \(NQ=\prod_{v\in Q}q_v\) for its norm, where \(q_v\) is the residue cardinality at \(v\). Let \(S\) be the set of places chosen over \(p_0,p_1\). The superscript \(S\cup Q\) on an \(L\)-function omits the Euler factors at \(S\cup Q\) and their conjugates. For a twist \(\xi\) ramified at every place of \(Q\), put
\[
    r_\delta(\xi)=
    \frac{|L^{S\cup Q}(z,(\nu\xi^\delta)^{-1})|}
         {|L^S(z,\nu^{-1})|}.
\]
A \emph{tame datum of modulus \(Q\)} also allows the residue character at an element of \(Q\) to be trivial. This distinction lets us enlarge a modulus while keeping the character fixed: that operation supplies the Euler relation. Let \(I_z\) index the exponents \(k\) in the Euler factors of \textup{(CV2)}, and let \(e_z\) be the exponent of \(NQ\) in \textup{(CV1)}. Set \((I_z,e_z)=(\{z\},0)\) for \(z\ge2\) and \((I_z,e_z)=(\{0,1\},1)\) for \(z=1\).

\subsection{An integral finite sum on each modulus}
We seek a single function \(F\) of tame data with the following properties. At a fixed modulus it is a sum \(\sum_i c_i\tau(u_i)\) with \(u_i\) adelic norm-one points and coefficients integral away from \(NQ\). This means \(|j(c_i)|_l\le1\) for every prime \(l\nmid NQ\) and every field embedding \(j:\mathbb C\to\mathbb C_l\). Its value \(F(1)\) at trivial character and empty modulus is nonzero. At a ramified twist one has
\[
                 F(\xi)=F(1)(NQ)^{e_z} r_+(\xi)^{e_z+1}. \tag{CV1}
\]
Inserting a good place \(v\) such that \(v,\bar v\) avoid the old modulus and chosen places, and \(Q\cup\{v\}\) is auxiliary, into the modulus of a fixed character \(\tau\) should multiply \(F\) by
\[
     \prod_{k\in I_z}\ \prod_{\epsilon=\pm1}
               (1-a_v^\epsilon q_v^{-k}),\qquad
      a_v=(\nu\tau)_{\rm BC}(\varpi_v),                       \tag{CV2}
\]
where \({\rm BC}\) is base change from the norm-one torus.

First suppose \(z\ge2\). Write \(D_0\) for the common integral denominator of the point kernels, \(A\) for the group order divisible by the orders of their class stabilizers, and \(T(\tau)\) for the definite theta lift to \(U(2z)\). Use
\[
               F(\tau)=\frac{A D_0^2}{(NQ)^{2z-1}}\|T(\tau)\|^2.
\]
At fixed modulus the expansion of \(T(\tau)\) has the same finite set of indices for all \(\tau\). Denote its points by \(g_i\) and its untwisted point kernels by \(K_i\), in the normalization incorporating the weights of the source sum. The scalar multiplying \(K_i\) in the lift is \(\tau(\det g_i)^{-1}\). The Gram pairing is conjugate-linear in its first variable, so expanding the squared norm writes \(F(\tau)\) as
\[
          \sum_{i,j}
          \frac{A D_0^2\langle K_j,K_i\rangle_Q}{(NQ)^{2z-1}}\,
                         \tau\bigl(\det(g_j/g_i)\bigr).
\]
For this modulus choose one representative \(h\) of each target class \([h]\) and use the same representatives for every \(i,j,\tau\). The coefficient pairing is
\[
        \langle K_j,K_i\rangle_Q=
        \sum_{[h]}\frac{\langle K_j(h),K_i(h)\rangle_{\rm Gram}}
                            {|\operatorname{Aut}([h])|}.
\]
Individual point kernels can depend on the representative. The full lift has the prescribed level, so its squared norm is independent of the choice. At a prime away from \(NQ\) the kernel lattice estimate puts both \(D_0K_i\) and \(D_0K_j\) in the integral gauge, at every point, so the pointwise Gram values after scaling are integral. The factor \(A\) clears each stabilizer in the class sum, and \(NQ\) is a unit. These are the claimed integral coefficients. The norm formula for the character lift gives \(F(\tau)=C L^{S\cup Q}(z,(\nu\tau)^{-1})\) with \(C\ne0\) independent of \(\tau\). The positive value at the empty twist and the same identity at \(\xi\) give (CV1). At \(v,\bar v\) the base change has reciprocal uniformizer values since \(\mu_{\rm BC}(\bar y)=\mu_{\rm BC}(y)^{-1}\) for a character \(\mu\) of the norm-one torus and the character is unramified at both places. Erasing those two Euler factors from the absolutely convergent series gives (CV2).\formalnote{12.1}{1}

The endpoint \(z=1\) has a different integral formula. Let \(\Lambda_Q(\tau)\) be the rank-two determinant period of \(\nu\tau\upsilon\) with archimedean determinant polynomial \(\prod_\sigma\det^{c_\sigma-1}\), where \(-c_\sigma\) is the weight of \(\nu\). The modulus at the chosen places has exponent two, and has exponent one at \(Q\). The rank-zero exchange in \S\ref{sec:three-exchanges} gives, at each actual twist, \(\gamma_\xi\Lambda_Q(\xi)=1\) and \(|\gamma_\xi|^2=C/\mathcal L_{S\cup Q}((\nu\xi)^{-1},1)\), with \(C\) fixed and nonzero. In particular
\[
   \frac{|\Lambda_Q(\xi)|^2}{|\Lambda_\emptyset(1)|^2}
                   = (NQ)\,r_+(\xi)^2 .
\]
Choose the common nonzero kernel scale \(N_0\) for this homogeneous polynomial, of degree \(2\sum_\sigma(c_\sigma-1)\), and set \(F(\tau)=|A N_0\Lambda_Q(\tau)|^2\). This proves (CV1) including nonvanishing at \(1\).

For integrality, expand the period over definite lattice classes and residue-matrix translates. At \(g\), the kernel for \(\nu\tau\upsilon\) is the kernel for \(\nu\upsilon\) times \(\tau(\det g)\). For a fixed \(Q\) the indices \(g\) are independent of \(\tau\). Their coefficients in \(A N_0\Lambda_Q\) are the scaled kernel at \(g\) times \(A/|\operatorname{Aut}([g])|\). At each prime away from \(NQ\) a rational translate makes the point integral, and the homogeneous kernel bound then makes the scaled kernel integral. The second factor is an integer. Taking the squared modulus gives the products \(c_i\overline{c_j}\) at the points \(u_i/u_j\). Integrality at every embedding over a given prime also holds after complex conjugation. Finally, the local rank-two period relation upon inserting \(v\) is multiplication by the polynomial \((1-X)(1-q_vX)\) in the inverse normalized Frobenius value. In its squared modulus the two conjugate factors for the first term are \(1-a_v^{\pm1}/q_v\), and those for the second are \(1-a_v^{\pm1}\), in either ordering. This gives (CV2) with \(I_1=\{0,1\}\) without needing absolute convergence of the Euler product at one. Section~\ref{sec:period-insertion} derives both the insertion formula and the common scale, tracing them to the unramified vacuum action and the integral Taylor bound.\formalnote{12.2}{2}

\subsection{Transport to the group ring and the cost at a tame generator}
Fix the field isomorphism \(\iota_\ell:\mathbb C\simeq\overline{\mathbb Q}_\ell\). At a fixed auxiliary modulus \(Q\), take any of the finite sums just constructed. For each \(u_i\) take a finite idele \(y_i\) for which the finite part of \(y_i/\bar y_i\) agrees with the finite part of \(u_i\). This is local Hilbert 90, with unit representatives away from the finitely many exceptional places. Weight-zero tame characters have trivial archimedean component, so \(\tau(u_i)=\tau_{\rm BC}(y_i)\). Write \(\mathrm{Cl}_Q\) for the group of norm-one finite ideles modulo rational norm-one points and its level-\(Q\) units. Set
\[
       \Theta_Q=\sum_i \iota_\ell(c_i)
          \left[\bigl(y_i/\bar y_i\bigr)^{-\delta}\right]
                   \ \in\ \overline{\mathbb Z}_\ell[\mathrm{Cl}_Q].
\]
The coefficients are integral because \(\ell\nmid NQ\). To evaluate a character \(\psi\) on \([y/\bar y]\), take \(\psi(\operatorname{Art}_{\rm geom}y)\). Admissibility means that the latter factors through the class of \(y/\bar y\), so it gives a well-defined character of \(\mathrm{Cl}_Q\). When \(\psi\) is the inverse of the \(\ell\)-adic Galois character attached to \(\tau\) by reciprocity (the inverse \(\ell\)-adic avatar in \S\ref{sec:hecke-avatars}), evaluating the \(i\)-th class therefore amounts to applying \(\psi\) to the Artin symbol of \(y_i^{-\delta}\), and the value is \(\iota_\ell(\tau^\delta(u_i))\). This proves \(\psi(\Theta_Q)=\iota_\ell(F(\tau^\delta))\), with the modulus \(Q\) on the right. At the empty modulus this is the nonzero integral scalar \(\mathbf1(\Theta_\emptyset)=\iota_\ell F(1)\). One can choose zero at nonauxiliary sets, since the relations use only auxiliary sets.\formalnote{12.3}{3}

Let \(\vartheta\) be the \(\ell\)-adic avatar of the base change of \(\nu^{-\delta}\). For an arithmetic Frobenius at a good newly inserted place \(v\) with \(Q\cup\{v\}\) auxiliary, \(\vartheta(\mathrm{Fr}_v)\psi(\mathrm{Fr}_v)\) is the comparison image of the local parameter for \(\nu\tau^\delta\) to the power \(\delta\), with the arithmetic Frobenius convention used here. The product over both signs in (CV2) is invariant under inversion. Every admissible finite Galois character in the norm relation comes from such a tame \(\tau\) by class field theory, and on enlarging the modulus its evaluation at the old classes is the same. Hence
\[
 \psi(\Theta_{Q\cup\{v\}})=\psi(\Theta_Q)
    \prod_{k\in I_z,\ \epsilon=\pm1}
       \left(1-(\vartheta(\mathrm{Fr}_v)\psi(\mathrm{Fr}_v))^\epsilon
                  q_v^{-k}\right).
\]
At an actual twist \(\xi\), equation (CV1) and \(|NQ|_\ell=1\) give
\[
   |\psi(\Theta_Q)|_\ell=
              |\mathbf1(\Theta_\emptyset)|_\ell\,
                        |\iota_\ell r_\delta(\xi)|_\ell^{e_z+1}. \tag{CV3}
\]
Thus the two distinct critical cases give exactly the shape assumed by \S\ref{sec:specialization}.\formalnote{12.4}{4}

To price this system, let \(g_0\) be a tame generator at the chosen \(w_*\). Its image \(u=\vartheta(g_0)\) is a nontrivial root of unity of order dividing \(q_{w_*}-1\), because this pair is ramified there and the local residue group has that order. Put \(n=q_{w_*}-1\). In \(\overline{\mathbb Z}_\ell\),
\(\prod_{\zeta^n=1,\ \zeta\ne1}(1-\zeta)=n\), by evaluating \((X^n-1)/(X-1)\) at \(1\). Every factor is integral, so \(v_\ell(u-1)\le v_\ell(n)\). In the system bound there are \(2|I_z|\le4\) Euler factors. In particular
\[
    \beta\le v_\ell(\mathbf1(\Theta_\emptyset))+4v_\ell(q_{w_*}-1).
\]
The first term cancels on taking valuations in (CV3). For a chosen twist with \(v_\ell(\psi(\Theta_Q))\le\beta+\varepsilon\) we obtain
\[
   \max\{0,v_\ell(\iota_\ell r_\delta(\xi))\}
           \le \frac{4v_\ell(q_{w_*}-1)+\varepsilon}{e_z+1}.
\]
This bound can be applied simultaneously to every critical pair and sign for which the specialization was made. For \(z=1\), the exponent of the ratio \(r_\delta(\xi)\) in \textup{(CV3)} is \(e_1+1=2\).\formalnote{12.5}{5}


\section{The character of the critical-value system and its Frobenius signs}\label{sec:hecke-avatars}
Two reciprocal conventions meet in the lower determinant bound. The Hecke character is read at a uniformizer by geometric reciprocity, whereas the norm relation in \S\ref{sec:critical-values} is evaluated at arithmetic Frobenius. At the coefficient prime there is also an algebraic factor that must be removed to read Frobenius on crystalline periods. We compute the three evaluations.

\subsection{From a finite idele to a Galois character}
Let \(E\) be totally complex. An algebraic Hecke character in the convention used here consists of a finite-idele character \(\theta_f\) with open kernel and integers \(n_\sigma\) such that
\(\theta_f(a)=\prod_{\sigma:E\hookrightarrow\mathbb C}\sigma(a)^{-n_\sigma}\) for \(a\in E^\times\). Fix the comparison of scalar fields \(\iota:\mathbb C\simeq\overline{\mathbb Q}_\ell\), and extend each \(\iota\sigma\) to the completion \(E_{v(\sigma)}\) it induces. For a finite idele \(x\) put
\[
 A(x)=\prod_{\sigma}(\iota\sigma)(x_{v(\sigma)})^{n_\sigma},
 \qquad r_\theta(\operatorname{Art}_{\rm geom}x)=\iota(\theta_f(x))A(x).
\]
The right side kills principal ideles by the type identity. The factor \(\iota(\theta_f(x))\) is locally constant because \(\theta_f\) has open kernel, and \(A(x)\) is continuous by continuity of the local embeddings. The product \(\iota(\theta_f(x))A(x)\) takes values in the units of a finite extension of \(\mathbb Q_\ell\). On finite unit ideles the first factor has finite image by compactness and the second takes unit values in the compositum of the finitely many conjugates of the \(E_v\) above \(\ell\). Adjoin the values of the product at representatives of the finitely many ideal classes. A power of each representative is a principal idele times a unit idele, so its value under the product also has valuation zero. Passing to finite quotients of these local coefficient units and applying global reciprocity, then taking the inverse limit, gives a unique continuous character \(r_\theta:G_E\to\overline{\mathbb Q}_\ell^\times\). Complex archimedean components do not contribute to finite Artin quotients \cite[\href{https://www.jmilne.org/math/CourseNotes/CFTc.pdf}{V, §§4--5, especially Theorems 5.3 and 5.5}]{arith-cft}.\formalnote{12.6}{6}

For a finite place \(v\nmid\ell\) the correction \(A\) on the local factor is identically one. If \(\theta_f\) is unramified there and \(\phi_v\) is arithmetic Frobenius, the two values are
\[
       r_\theta(\phi_v^{-1})=\iota\theta_f(\pi_v),\qquad
       r_\theta(\phi_v)=\iota\theta_f(\pi_v)^{-1}.
\]
The choice of \(\pi_v\) is immaterial since a unit is in the kernel. At a ramified place the Artin images of the units give the inertia values. These assertions concern the abelianization and therefore apply to every lift of the indicated unramified Frobenius.\formalnote{12.7}{7}

\subsection{The finite Fourier sum that defines the norm relation}
Use the ramified pair \((\nu,z)\) and the sign \(\delta=\pm1\) from \S\ref{sec:critical-values}. Base change a norm-one character by evaluating it at \(x/\bar x\); indicate this by a tilde. The system character is \(\vartheta=r_{\widetilde\nu^{-\delta}}\) and the auxiliary character is \(\psi=r_{\widetilde\tau^{-1}}\), where \(\tau\) has weight zero and tame modulus \(Q\). At a new good split prime \(v\), both local base changes are unramified, \(v\nmid\ell\), and \(v,\bar v\) lie outside \(Q\). Thus
\[
 \vartheta(\phi_v)\psi(\phi_v)
    =\iota\!\left(\widetilde\nu(\pi_v)^\delta\widetilde\tau(\pi_v)\right)
    =\iota\!\left(\widetilde{\nu\tau^\delta}(\pi_v)\right)^\delta .
\]
Taking a product over \(\epsilon=\pm1\) removes the final \(\delta\) in the Euler factors \(1-(\vartheta\psi)(\phi_v)^\epsilon q_v^{-e}\). At a chosen split \(w\nmid\ell\) where the residue character is nontrivial and the conjugate residue character is trivial, the value of \(\vartheta\) at the stipulated tame generator is \(\iota(\nu_w(\bar u))^{-\delta}\) for a generator \(\bar u\in k_w^\times\). This value has order dividing \(q_w-1\) and differs from one. These two computations account for the Euler factor and, separately, for the distinguished nonzero denominator.\formalnote{12.8}{8}

To check that the allowed class characters suffice, write \(I\) for an auxiliary modulus. In the norm relation admissibility means that a finite Galois character has a realization
\(\psi(\operatorname{Art}x)=\eta(x/\bar x)\), with the stated unit and level conditions on \(\eta\), including triviality at the primes above \(\ell\). Restrict \(\eta\) to the finite parts of norm-one ideles. It kills the rational norm-one subgroup by principality and the integral norm-one subgroup congruent to \(1\) at \(I\) by the local conditions. Hence \(\iota^{-1}\eta\) there defines a weight-zero character \(\mu\) of the norm-one class group. The conductor condition at a marked component makes the character on local units factor through reduction to the residue units. Thus \(\tau=\mu^{-1}\) is a tame character with the allowed level \(I\), including when some residue factors are trivial. Its inverse avatar at \(\operatorname{Art}x\) is \(\iota\mu(x/\bar x)=\psi(\operatorname{Art}x)\); reciprocity gives equality on \(G_E\).\formalnote{12.9}{9}

For the measure at level \(I\), let \(u_j\) denote the finitely many norm-one adelic points and \(c_j\) the coefficients of \S\ref{sec:critical-values}. Choose a finite idele \(y_j\) with \(y_j/\bar y_j\) equal to the finite part of \(u_j\). Hilbert 90 gives local solutions. Off a finite set the nonsplit extensions are unramified, so the solutions may be made units by multiplying by an element of the base field of the same valuation. At split places one takes one component to be one. The characters \(\tau\) have weight zero, so their value on \(u_j\) is their base change on \(y_j\). The group-ring element has coefficients \(\iota c_j\) at the classes of \((y_j/\bar y_j)^{-\delta}\). In the evaluation defined in \S\ref{sec:critical-values}, the value of \(\psi\) on this class is its value on the geometric Artin symbol of \(y_j^{-\delta}\). The sum is thus \(\sum_j\iota(c_j\tau(u_j)^\delta)\). Since \(\ell\nmid\prod_{v\in I}q_v\), each coefficient is \(\ell\)-integral by the integral measure bound.\formalnote{12.10}{10}

\subsection{The local character at the coefficient prime}
Fix \(v\mid\ell\), \(L=E_v\), and let \(A_v\) be the product of the factors of \(A\) inducing this place. In arithmetic local reciprocity the character on \(L^\times\) is \((\iota\theta_{f,v} A_v)^{-1}\). On an open subgroup of units this agrees with the algebraic torus character \(A_v^{-1}\). If \(\theta_{f,v}\) is unramified, it agrees on all units. Conrad's rank-one crystalline criterion therefore applies over \(L\) with a finite coefficient field containing the local embeddings and character values. For \(a=[k_v:\mathbb F_\ell]\) his Frobenius formula on the covariant \(D_{\rm cris}\) gives
\[
      \varphi^a=
       r_\theta(\operatorname{Art}_{\rm arith}\pi_v)^{-1} A_v(\pi_v)^{-1}
                        =\iota\theta_{f,v}(\pi_v).
\]
The algebraic factor cancels only after passing to periods. If \(\theta_{f,v}\) is ramified, extend its finite character on units to \(L^\times\) by assigning the value \(1\) to a chosen uniformizer. Removing this finite-order factor leaves the preceding crystalline character. Adjoining a field that kills the finite character proves potential crystallinity. On putting it back, the finite factor acts on the linear Weil operator by its Galois value, and the scalar on the Weil element represented geometrically by \(x\), with valuation \(0\) or \(1\), is \(\iota\theta_{f,v}(x)\). Changing the extension off units changes the two factors by inverse unramified characters. In the filtration convention of the proof the single jump at \(\sigma\) above \(v\) is \(-n_\sigma\). Indeed Conrad decomposes the algebraic character into labelled Lubin--Tate factors and computes their crystalline Frobenius \cite[Appendix B, Proposition B.4 and its proof]{arith-conrad}. For a Lubin--Tate formal group the action on its one-dimensional tangent space is the structural embedding of \(L\). The Hodge--Tate decomposition of its Tate module thus has cyclotomic filtration jump \(-1\) at that embedding and jump zero at each other embedding \cite[\S4, Corollary 2 to Theorem 3]{arith-tate-pdiv}. Taking tensor powers and an unramified or finite-order twist gives the asserted jumps.\formalnote{12.11}{11}

\phantomsection\label{route:20}\leanroute{The formal proof first obtains the character of the prescribed local infinity type at a split place by an odd-rank attachment and a determinant. It divides by that character, treats the resulting weight-zero twist, and descends from a split CM step at the other places. The rank-one classification applies directly to the local field \(E_v\); its Frobenius and filtration conventions give the same scalar and jump displayed above.\formalnote{12.12}{12}}


\section{The split-prime relation bound by global duality}\label{sec:selmer-formula}\label{sec:split-duality}
Allow a global cocycle to ramify at a new set of split primes \(P\). Its values on the tame inertia generators at those primes must satisfy linear relations coming from global duality. We need relations that work for every later choice of \(P\) outside a single finite exceptional set. Inequality \textup{(SD2)} below bounds the size of the module of such relations.

\subsection{Statement of the relation bound}
Let \(K\) be totally complex with \(r_2\) infinite places, \(G=G_K\), and \(M\) a finite discrete \(G\)-module killed by \(m=\ell^r\), where \(r>0\). An open normal subgroup \(U\) acts trivially on \(M\) and on \(\mu_m\). Write \(F=\overline K^U\). Suppose \(\sigma\in G\) is central modulo \(U\), fixes all \(\ell\)-power roots of unity, and \(\rho(\sigma)-1\) is injective on \(M\). At places \(v\) of a finite set \(S\) choose a decomposition group and an additive group \(A_v\) of allowed continuous cocycles containing its coboundaries. Inertia acts trivially on \(M\) outside \(S\). Let \(Z(P)\) consist of the global cocycles satisfying these local conditions and vanishing on inertia away from \(S\cup P\). Write \(M^\vee=\operatorname{Hom}(M,\mathbb Q/\mathbb Z)\) and give this group the action \(g\lambda=\chi_m(g)\lambda\circ\rho(g)^{-1}\), where \(\chi_m\) is cyclotomic.

For finite \(T\supset S\), a relation is a locally constant function \(y:G\to M^\vee\), zero off \(U\), additive on \(U\), equivariant there under conjugation, and zero on inertia away from \(T\). It also satisfies the following test. Take any family \(P\) of distinct primes outside \(T\) split in \(F\), with Frobenius lifts \(\phi_v\in U\) and tame generators \(\tau_v\) normalized by one primitive \(m\)-th root of unity. For every \(z\in Z(P)\) one must have
\[
                    \sum_{v\in P}y(\phi_v)(z(\tau_v))=0.       \tag{SD1}
\]
Call the relation module \(\mathcal D_T\). It is finite by finiteness of bounded-ramification cocycles with finite coefficients over \(F\), applied to its restriction to \(U\). Fix a positive integer \(p\) prime to \(m\), with every place above \(p\) outside \(S\) and unramified in \(F\). The split-prime relation bound asserts that there is a \(T\) avoiding all the places dividing \(p\), for which
\[
       |Z(\varnothing)|\,|M|^{|S|+r_2}
           \ \le\ |\mathcal D_T|\,|M|\prod_{v\in S}|A_v| .       \tag{SD2}
\]
The freedom to avoid \(p\) allows the descent to preserve unramifiedness at the first tied prime and its conjugate.\formalnote{12.13}{13}

\subsection{Duality and the relation module}
\phantomsection\label{route:21}\leanroute{The formalization builds global relations by detecting roots of units at an auxiliary finite set of primes. It counts trivializations on successive coefficient layers and extends their relations to every later choice of split primes. We take the dual Selmer group as the source of universal relations. The Greenberg--Wiles formula supplies its order, global reciprocity supplies the relation, and the central-element hypothesis makes restriction injective. These three steps give (SD2) for the full coefficient module in one calculation.\formalnote{12.14}{14}}

Identify \(a/m\in \frac1m\mathbb Z/\mathbb Z\) with \(\zeta_m^a\) for the primitive root used to normalize tame generators. This identifies \(M^\vee\) with the Cartier dual \(M^*=\operatorname{Hom}(M,\mu_m)\). Set \(\mathcal L_v=A_v/B^1(K_v,M)\) for \(v\in S\), take the unramified subgroup at the other finite places, and the zero group at complex places. The Selmer group with these conditions is \(Z(\varnothing)/B^1(K,M)\). In fact a class unramified at \(v\notin S\) has every representative zero on inertia because inertia acts trivially on \(M\). At \(v\in S\) adding a coboundary preserves membership in \(A_v\).\formalnote{12.15}{15}

Write \(D=\operatorname{Sel}_{\mathcal L^\perp}(K,M^*)\), using the exact annihilator at each place under the local invariant pairing. The finite-module Greenberg--Wiles formula \cite[Theorem 2.19]{arith-ddt}, from Poitou--Tate duality and the local and global Euler characteristics \cite[I, Corollary 2.3 and Theorems 2.8, 4.10 and 5.1]{arith-milne}, gives
\[
 \frac{|\operatorname{Sel}_{\mathcal L}(K,M)|}{|D|}
  =\frac{|H^0(K,M)|}{|H^0(K,M^*)|}
     \prod_v\frac{|\mathcal L_v|}{|H^0(K_v,M)|}.
\]
All but finitely many local terms are one. More precisely, at a finite \(v\notin S\) an unramified class is given by an arbitrary value in \(M/(\phi_v-1)M\). Its order equals \(|\ker(\phi_v-1)|\) since \(M\) is finite. At \(v\in S\) the factor is \(|A_v|/|M|\) by counting coboundaries. Each of the \(r_2\) complex places contributes \(1/|M|\). Multiplying by \(|B^1(K,M)|=|M|/|H^0(K,M)|\) therefore gives
\[
 |Z(\varnothing)|\,|M|^{|S|+r_2}
    =\frac{|D|}{|H^0(K,M^*)|}\,|M|\prod_{v\in S}|A_v|.              \tag{GW1}
\]
This uses orders of finite abelian groups and applies without freeness over the coefficient ring or a restriction to exponent \(\ell\).

We now embed \(D\) in the relation module. Take for \(T\) the union of \(S\), the places dividing \(m\), and the places ramified in \(F/K\). These sets avoid the auxiliary integer \(p\) by the hypotheses of (SD2). Take a cocycle \(b\) representing a member of \(D\) and restrict it to \(U\). Since \(U\) acts trivially on \(M^*\), the restriction is a homomorphism and is independent of the representative. Extend it by zero off the open and closed subgroup \(U\) and call the function \(y\). The cocycle identity gives \(b(gug^{-1})=g b(u)\), \(u\in U\), which is exactly the conjugation rule in (SD1). The function is locally constant. For \(v\notin T\) the finite modules are unramified and of order prime to the residue characteristic. Local Tate duality identifies the annihilator of the unramified subgroup with the dual unramified subgroup. Thus \(b\), and consequently \(y\), vanishes on inertia there.

Let \(P\) be any of the families in (SD1) and take \(z\in Z(P)\). The sum of the local invariants of \(b\cup z\) is zero by reciprocity for the Brauer group of \(K\). Away from \(P\) the two classes are locally orthogonal by construction, including at the places dividing \(m\) and at infinity. At a \(v\in P\) the decomposition group acts trivially on \(M\) and \(\mu_m\), since \(v\) splits in \(F\), and \(b\) is unramified. The local term, with the common sign prescribed by the order of cup factors, is
\[
                        y(\phi_v)\bigl(z(\tau_v)\bigr).
\]
To check the normalization, split \(M\) into cyclic groups of orders \(d\mid m\) at this place. In one such coordinate use \(\zeta_d=\zeta_m^{m/d}\). The normalized generator \(\tau_v\) sends \(\pi_v^{1/d}\) to \(\zeta_d\,\pi_v^{1/d}\), so the Kummer class of a uniformizer has tame coordinate one. The unit part of a Kummer class is unramified here since \(d\) is invertible and \(\mu_d\subset K_v\). Two unramified classes have zero cup product. The invariant of the cup of an unramified character with the class of \(\pi_v\) is its value on arithmetic Frobenius by local reciprocity. Bilinearity gives the displayed evaluation in each coordinate and hence on \(M\). Reversing the cup order reverses every sign simultaneously. Global reciprocity therefore gives exactly (SD1), for every such \(P\) at once.
For the local invariant and global reciprocity see \cite[III, \S\S1, 3--4, and VIII, Theorem 4.2]{arith-cft}.\formalnote{12.16}{16}

For injectivity, put \(\Delta=G_K/U\) and take the element \(\sigma\) of (SD2). Its image is central in \(\Delta\), it fixes \(\mu_m\), and \(\rho(\sigma)-1\) is injective on the finite group \(M\), hence invertible. On the dual, \(\sigma-1\) is also invertible, by dualizing \(\rho(\sigma)^{-1}-1\). In particular \(|H^0(K,M^*)|=1\). A class with zero restriction to \(U\) is represented by a cocycle \(w\) on \(\Delta\). Comparing \(w(\sigma g)\) and \(w(g\sigma)\) gives
\[
                 (\sigma-1)w(g)=(g-1)w(\sigma).
\]
The inverse of \(\sigma-1\) commutes with \(g\). Thus \(w\) is the coboundary of \((\sigma-1)^{-1}w(\sigma)\), proving injectivity. With \(|D|\le|\mathcal D_T|\), equation (GW1) is (SD2). This argument supplies the terminal inequality used in \S\ref{sec:tied-primes}. The excluded set avoids the first tied prime and its conjugate by the stated hypotheses on \(p\).\formalnote{12.17}{17}


\section{Choosing the tied primes by Selmer descent}\label{sec:tied-primes}

At a new split prime we evaluate a cocycle on a tame inertia generator and on a Frobenius lift. We require the inertia value to equal a chosen scalar times the Frobenius value, and impose the corresponding equation at the conjugate prime. We call these the tied conditions. Mazur and Rubin use auxiliary ramification to decrease the lengths of a pair of dual residual Selmer groups for an abelian variety \cite[Proposition~9.17]{mazur-rubin-stability}. Here the lengths to be bounded come from several lattices and several powers of one character. A suitable prime and scalar can shorten the finite module of global cocycles satisfying all the local conditions. We iterate this step and use reciprocity to bound the module when the descent stops. For the specialization in \S\ref{sec:specialization}, the number of primes used must stay bounded as the orders of the twisting characters increase.

\begin{example}[An equation on tame local classes]\label{ex:tame-tie}
Take trivial coefficients \(\mathbb F_5\) at the local field
\(\mathbb Q_{11}\). A continuous cocycle is an additive homomorphism.
It kills wild inertia, which is pro-\(11\), and is determined by its values \(a,b\)
on an arithmetic Frobenius lift and a tame generator, respectively. The relation
\(\phi\gamma\phi^{-1}=\gamma^{11}\) in the tame quotient imposes \(10b=0\), automatic in
\(\mathbb F_5\). Thus the two values are initially arbitrary.
Imposing \(b=0\) retains the unramified line. Imposing \(b=2a\)
retains a different line. To see the effect on a global condition,
suppose that the image of a space of global classes at this place is
the unramified line. All its classes satisfy the first equation.
Only those restricting to zero satisfy the second, since
\((a,0)=(a,2a)\) forces \(a=0\). The new condition cuts the global
space by one dimension even though the two local lines have the same
dimension.
\end{example}

\subsection{The finite problem and the choice of Frobenius}

In the Selmer descent the restrictions are tested at a prime and its
conjugate together. The local conditions must also be compatible as
the coefficient depth varies.

Let \(G=G_E\) for the CM field \(E\), and let \(c\) be complex conjugation. For \(g\in G\) put \(g^c=cgc\) and \(g^{-c}=(g^c)^{-1}\). Fix the ring of integers \(O\) of a finite extension of \(\mathbb Q_\ell\), with uniformizer \(\varpi\) and residue field \(k\). Write \(\lambda\) for ordinary \(O\)-module length. The data are finitely many lattices \(T=O^d\) with continuous actions \(\rho\), unramified outside \(S\subset S_0\), where \(S\) contains all places above \(\ell\). At \(w\mid\ell\) each datum specifies a finite-rank module \(Z_w\) of integral local cocycles containing the coboundaries. Its finite-level local conditions are the matrix-linear relations satisfied by \(Z_w\). The hypotheses used below are
\[
  {\rm span}_O\rho(G)=M_d(O),\qquad
  \sum_{w\mid\ell}{\rm rank}_O Z_w
          \le (\#\{w\mid\ell\}+r_2(E))d .
\]
For each datum there is an inertia element \(\gamma_*\) at a place of \(S\) away from \(\ell\) such that both \(\rho(\gamma_*)\) and \(\rho(\gamma_*^{-c})\) are the scalar \(\zeta\), with \(\zeta^2-1\in O^\times\). In particular \(\zeta-1\) is a unit. Each cohomology class has a unique representative \(z\) with \(z(\gamma_*)=0\), since subtracting a coboundary changes this value by \((\zeta-1)m\).

Fix a depth \(M\), put \(e=v_\varpi(\ell^M)\), and assume \(\#k>2de+1\) for each datum. Fix also the element \(g_0\) in inertia at a place of \(S_0\) away from \(\ell\) and a neighborhood of it. For \(\mathfrak b=(\varpi^j)\) put \(V=T/\mathfrak b T\). Given primes \(J\), the tied module \(\mathcal S_J(V)\) consists of normalized cocycles obeying the old local conditions, unramified off \(S\cup J\cup\bar J\), with the two additional equations at each \(P\in J\)
\[
 z(\gamma_P)=s_{T,P}z(\phi_P),\qquad
 z(\gamma_P^{-c})=s_{T,P}z(\phi_P^{-c}).
\]
Here \(\gamma_P\) is a tame generator and \(\phi_P\) a Frobenius lift at a chosen place over \(P\). Our target for every \(N>M\) is a single \(J\) and slopes with
\[
 \lambda\mathcal S_J(T/\ell^M T)\ \le\ g_T+l_T .
\]
Write \(V^*={\rm Hom}_{\mathbb Z}(V,\mathbb Q/\mathbb Z)\) with the Cartier dual action, which includes the cyclotomic factor, and twist this action by \(c\) for \(V^{*c}\). The global allowance \(g_T\) is the supremum of \(\lambda{\rm Hom}_{O[G]}(T/aT,(T/aT)^{*c})\) over nonzero \(a\in O\). The term \(l_T\) sums, over the old places in \(S\) away from \(\ell\), the supremum of the excess of the unrestricted local cocycle length over \(\lambda(T/aT)\), with negative excess counted as zero. Both allowances are finite in the Selmer data. The first-prime determinant terms are charged later when the ties are removed.\formalnote{12.18}{18}

The primes are chosen with the following conditions. There is a distinguished \(P_0\) with Frobenius in the given neighborhood of \(g_0\), with zero slopes for all data. Each \(P\) has prime norm \(1\bmod\ell^N\), \(P\ne\bar P\), and a character \(\chi_P:G\to\mathbb Z/\ell^N\) coming from \(U(1)\), with \(\chi_P(\gamma_P)=1\), unramified away from \(P,\bar P\), trivial on decomposition groups above \(\ell\). Write these characters additively. At the Frobenius of a later prime, \(\chi_{P_0}\) has value \(1\) and all other chosen characters have value \(0\). Where the slope for \(T\) is nonzero both \(\phi_P,\phi_P^{-c}\) act trivially on \(T/\ell^M T\).

We recall why such primes can be prescribed for the reciprocity argument. Let \(H\) be the subgroup generated by the commutators \([g,\gamma_*]\), \([g,\gamma_*^c]\) for \(g\in G\). It is normal, since \(b[g,\gamma]b^{-1}=[bg,\gamma][b,\gamma]^{-1}\), and stable under \(c\). Its elements act trivially on \(T\), and all abelian characters kill it. For \(h\in H\), Frobenius at a new prime can be prescribed in any neighborhood of \(\gamma_{P_0}h\), with \(\gamma_{P_0}\) the tame inertia generator at the first tied prime, outside any fixed finite exclusion set and its conjugate. Indeed the target belongs to the normal closure of old inertia at \(P_0\), the old ramified places and their conjugates, all away from \(\ell\) for the inertia used here. It fixes \(T/\ell^M T\) and \(\mu_{\ell^N}\), as does its conjugate, and its old character values are those at \(\gamma_{P_0}\).

For the number-theoretic construction choose an auxiliary set \(D\) of places containing the required decomposition-triviality places and their conjugates. Enlarge it using primes representing the ideal classes so that principal ideles and ideles unit off \(D\) generate the finite ideles. Choose the added places away from the inertia used to form the target. Fix \(\ell^N\)-th roots of a finite set of generators of the \(D\)-units and the \(\ell^N\)-th roots of unity. At the indicated inertia places the units and \(\ell\) are invertible, so these root extensions are unramified there, including after taking their normal closure. The target therefore fixes the roots. The Chebotarev density theorem \cite[Chapter~V, Theorem~3.23]{arith-cft} gives positive density to the Frobenius conjugacy class of the target in the finite Galois quotient imposing these conditions. Primes of absolute residue degree at least two have Dirichlet density zero, since their reciprocal norms sum to at most \([E:\mathbb Q]\sum_p p^{-2}\). We may exclude the finitely many ramified and previously chosen primes. An absolute degree-one prime left in this class splits over \(E^+\), since an unramified inert prime has even absolute residue degree. Choosing the place above it fixes the Frobenius representative in the prescribed neighborhood. In the residue field an additive discrete logarithm modulo \(\ell^N\) kills the reductions of the \(D\)-units. It defines a character on the ideles unit off \(D\) through the component at the new prime and extends trivially across principal ideles since the formula agrees on their intersection. Evaluating this character on \(b/\bar b\) and applying global Artin reciprocity gives the stated anticyclotomic character. Normalizing its value at the new tame generator to \(1\) allows multiplication of the Frobenius lift by a tame power to make its own character value \(0\). The Chebotarev prescription may be made in a finite quotient unramified at the new prime, so this correction preserves the prescribed neighborhood. At the first prime the same construction uses the target \(g_0\) directly. At later primes the new character kills decomposition at the old primes other than \(P_0\), as required in the displayed character pattern.\formalnote{12.19}{19}

\subsection{A prime that shortens a group}

Fix a datum, a level \(1\le j\le e\), and a candidate new prime \(P\) with the stated prescriptions for this datum. Temporarily allow ramification at \(P,\bar P\) without imposing their ties. Call the resulting finite module \(B\). Let
\[
    A z=(z(\gamma_P),z(\gamma_P^{-c})),\qquad
    F z=(z(\phi_P),z(\phi_P^{-c})),\qquad A,F:B\longrightarrow V^2 .
\]
The old group is \(\ker A\). In fact inertia acts trivially on \(V\) at both places, and for these finite \(\ell\)-primary coefficients, vanishing at tame generators implies vanishing on inertia. Adding slope \(u^{-1}\) for a unit \(u\) gives \(\ker(uA-F)\).

\paragraph{The slope lemma.}
Outside at most \(2d\) residues of \(u\) one has
\[
          \lambda\ker(uA-F)\le
          \lambda\{z\in\ker A:Fz\in A(B)\}.
\]
Let \(B'\) be the sum of the kernels \(\ker(vA-F)\) over all \(v\in O\). Then \(F(B')\subset A(B')\). Take the inverse image \(\Lambda\subset O^{2d}\) of \(A(B')\subset (O/\varpi^j)^{2d}\). It is free of rank at most \(2d\), with surjection \(p:\Lambda\to A(B')\). By freeness lift \(p\) to \(s:\Lambda\to B'\) with \(As=p\), and lift \(Fs\) to \(U:\Lambda\to\Lambda\) with \(pU=Fs\). If \(u\) avoids the roots of the characteristic polynomial of \(U\bmod\varpi\), then \(u-U\) is invertible. Hence \((uA-F)s=p(u-U)\) surjects onto \(A(B')\). Its image on \(B'\) has at least the length of \(A(B')\). Comparing kernels on \(B'\) gives the inequality because the whole kernel on the left lies in \(B'\), and \(\ker(A|_{B'})\subset\ker A\cap F^{-1}A(B)\).

If some \(z\in\mathcal S_J(V)\) has \(Fz\notin A(B)\), the bound is strictly smaller than the old length. Avoid the bad residues for all \(1\le j\le e\) at once. At most \(2de\) residues are excluded, and \(\#k>2de+1\) leaves the residue of a unit. Give this datum that inverse unit slope and all other data slope zero. No tied module grows, and at least one shrinks. The sum of the lengths over data and levels is a nonnegative integer. Starting at \(J=\{P_0\}\) it is bounded by the sum for the old groups without auxiliary ramification, since all first slopes are zero. Those groups are finite by finiteness of Galois cohomology with bounded ramification and finite coefficients. Thus the final number of primes is bounded by one plus this initial sum, independently of \(N\). At termination we have for \emph{every} datum, level, admissible new prime and old class \(z\)\formalnote{12.20}{20}
\[
                 \text{there exists }z'\in B\text{ with }Az'=Fz . \tag{D1}
\]

\subsection{The size of a module when descent stops}

Fix a datum and use the least \(j\le e\) for which \(\varpi^j\) kills its final module at level \(e\). Multiplication by \(\varpi^{e-j}\) identifies \(O/\varpi^j\) with the \(\varpi^j\)-torsion in \(O/\varpi^e\). The cocycle and tie equations commute with this injection, so the level-\(j\) module has the same length as the level-\(e\) module. For \(j=0\) there is nothing to prove. Otherwise minimality supplies a cocycle \(z\) whose value at some \(g\) has a unit coordinate in \(V=(O/\varpi^j)^d\). From the cocycle law and \(z(\gamma_*)=0\) one obtains
\[
                         z([g,\gamma_*])=(1-\zeta)z(g).
\]
Consequently \(z(H)\) has a vector with a unit coordinate. Its span is stable under every \(\rho(g)\), because \(z(ghg^{-1})=\rho(g)z(h)\) for \(h\in H\). The matrix-span assumption lets us apply every matrix unit to the chosen vector. Thus \(z(H)\) spans \(V\).

The split-prime relation bound (SD2) of \S\ref{sec:selmer-formula} applies to the local conditions of the descent. We use its cardinality factors to cancel the local contributions. Put \(S'=S\cup (J\setminus\{P_0\})\cup\overline{J\setminus\{P_0\}}\), \(D_0=|S'|+r_2(E)\), and \(v_0=\lambda V\). Denote by \(L_w\) the allowed local cocycles at \(w\), and by \(Z_{\rm gl}\) the unnormalized cocycles unramified outside \(S'\) with these local conditions. For a suitable finite excluded set the relation theorem gives a finite \(O\)-module \(\mathcal D\) with
\[
             \lambda Z_{\rm gl}+D_0v_0
                \le\lambda\mathcal D+v_0+\sum_{w\in S'}\lambda L_w .
\]
A member \(y\) of \(\mathcal D\) is an equivariant, locally constant homomorphism on \(U=\ker(\rho\bmod\ell^M)\cap\ker(\epsilon\bmod\ell^N)\) into \(V^*\), extended by zero off \(U\), with the stated reciprocity conditions at split primes. It is unramified at \(P_0,\bar P_0\), since the excluded set avoids the rational norm of \(P_0\).

At \(w\mid\ell\) the linear relations give \(\lambda L_w\le\operatorname{rank}_O(Z_w)\lambda(O/\varpi^j)\). The rank hypothesis bounds the sum of these lengths by \((\#\{w\mid\ell\}+r_2(E))v_0\). At each added place, the Frobenius value determines the tame value by the tie, so it determines the cocycle and \(\lambda L_w\le v_0\). At the old places away from \(\ell\) the sum of excesses over \(v_0\) is at most \(l_T\) by definition. Hence
\[
               \sum_{w\in S'}\lambda L_w\le D_0v_0+l_T .
\]
Adding coboundaries to the normalized cocycles gives \(\lambda Z_{\rm gl}\ge v_0+\lambda\mathcal S_J(V)\). A coboundary obeys each tie, since either the slope is zero or Frobenius acts trivially. Evaluation at \(\gamma_*\) distinguishes the translates by distinct coboundaries. Cancelling the two contributions of \(v_0\) in the count yields\formalnote{12.21}{21}
\[
                         \lambda\mathcal S_J(V)\le\lambda\mathcal D+l_T . \tag{D2}
\]

It remains to bound \(\mathcal D\). For \(y\in\mathcal D\), \(t\in O\), \(h,h'\in H\), form the additive pairing
\[
     B_t(h,h')=y(h')\bigl(tz(h)\bigr)
                            -y(h^c)\bigl(tz((h')^c)\bigr)\in\mathbb Q/\mathbb Z .
\]
Its additivity in each group argument follows because \(H,H^c\subset U\) act trivially on \(V\) and its character module. We first prove \(B_t(h,h)=0\). Both \(z\) and \(y\) vanish at \(\gamma_{P_0}\) and its conjugate, so the diagonal value is unchanged on replacing \(h\) by \(\gamma_{P_0}h\). By local constancy and the prime prescription choose a new \(P\) outside the exclusions with the same diagonal value at \(\phi_P\), where \(\phi_P,\phi_P^c\in U\). Use (D1) to find \(z'\) ramified additionally only at the new pair and satisfying
\[
 z'(\gamma_P)=z(\phi_P),\qquad
 z'(\gamma_P^{-c})=z(\phi_P^{-c})=-z(\phi_P^c).
\]
Reciprocity for \(t z'\) and the relation \(y\) gives
\[
            y(\phi_P)\bigl(t z'(\gamma_P)\bigr)
              +y(\phi_P^c)\bigl(t z'(\gamma_P^{-c})\bigr)=0 .
\]
The normalized tame generator at the conjugate place is the inverse conjugate: complex conjugation inverts the chosen roots of unity. Passing to the prescribed generator multiplies both tame values by the same integer. This proves the required diagonal identity, including its sign.

Apply the diagonal identity to \(hh'\) to get \(B_t(h,h')+B_t(h',h)=0\). Conjugating the first argument by \(\gamma_*\) multiplies the parameter \(t\) by \(\zeta\). Conjugating the second divides it by \(\zeta\). Indeed \(z\) transforms by \(\rho\), \(y\) by its Cartier dual, \(\rho(\gamma_*^c)=\zeta^{-1}\), and \(\epsilon(\gamma_*)=1\) at this inertia place away from \(\ell\). Apply antisymmetry with parameter \(\zeta t\) to \(\gamma_*h\gamma_*^{-1},h'\), and subtract its value with parameter \(t\) at \(h,h'\). This gives \(B_{(\zeta^2-1)t}(h,h')=0\). The unit hypothesis now implies \(B_t=0\) for every \(t\), including at \(\ell=2\).

Define \(\eta_y(z(h))=y(h^c)\). To check relations take \(\sum a_i z(h_i)=0\). By \(B_{t a_i}=0\), the character \(\sum a_i y(h_i^c)\) vanishes at each \(t z((h')^c)\), since its value equals \(y(h')(t\sum a_i z(h_i))\). Those elements generate \(V\) additively, since \(c\) preserves \(H\) and \(z(H)\) spans \(V\) over \(O\). Thus \(\eta_y:V\to V^{*c}\) is well-defined and \(O\)-linear. The conjugation formulas on \(H\) give its \(G\)-equivariance. It determines \(y\): if \(y|_H=0\), then for \(g\in U\) the commutator \([g,\gamma_*]\) lies in \(H\). Hence \(y(g)=y(\gamma_*g\gamma_*^{-1})\), which by equivariance equals the character \(m\mapsto y(g)(\zeta^{-1}m)\). Since \(1-\zeta^{-1}\) is a unit, \(y(g)=0\); outside \(U\) the function is zero by definition. We have proved an injection
\[
          \mathcal D\ \lhook\joinrel\longrightarrow\
                {\rm Hom}_{O[G]}(V,V^{*c}),\qquad y\longmapsto\eta_y .
\]
Combining this with (D2) proves the asserted bound \(g_T+l_T\). Division by \(v_\varpi(\ell)\) gives the normalized lengths of Section~\ref{sec:specialization}. That section uses the bound at \(X=0\) for its line of twists and includes the two first-prime determinant allowances when it removes the ties.\formalnote{12.22}{22}


\section{One character for all the Selmer and Hecke-value bounds}
\label{sec:specialization}

Fix a prime \(\ell\). The Selmer descent of \S\ref{sec:tied-primes} bounds cocycle lengths. Below we encode the bound for each Galois lattice as a polynomial evaluated at a point that may depend on the lattice. The critical-value group-ring elements \(\Theta_Q\) of \S\ref{sec:critical-values} give a second family of polynomials, with their own evaluation points. These points have to be replaced by one finite-order character, and each of its powers \(\xi^1,\ldots,\xi^m\) must satisfy both bounds. There is an order-of-choice constraint. The auxiliary primes depend on the conductor exponent \(N\), while the character used at those primes must have order dividing \(\ell^N\). We first prove a specialization theorem in which the order bound is independent of the polynomials, their degrees and their test points. We then apply it to the two kinds of arithmetic data.

\subsection{The specialization theorem}
Write \(K=\overline{\mathbb Q}_\ell\), let \(\mathcal O\) be its valuation ring and use the ordinary additive valuation \(v(\ell)=1\), \(v(0)=+\infty\). All evaluations we bound lie in \(\mathcal O\). Fix a bound \(s_0\) on the number of variables, a positive rational radius \(\rho\), a finite \(\beta\ge0\), and a real error \(\eta>0\). There is an integer \(H\) with the following property. Let \(P_i\in\mathcal O[X_1,\ldots,X_s]\), \(s\le s_0\), be any finite family, and let the point \(x_i\in\mathcal O^s\) depend on \(i\). Suppose
\[
        v(x_{ij})\ge\rho,\qquad
        \sum_i v(P_i(x_i))\le\beta .
\]
There are roots of unity \(\zeta_1,\ldots,\zeta_s\) of orders dividing \(\ell^H\) such that \(v(\zeta_j-1)\le\rho\) for every \(j\), and
\[
       v(P_i(\zeta_1-1,\ldots,\zeta_s-1))
                     \le v(P_i(x_i))+\eta\qquad\text{for every }i.
       \tag{S1}
\]
The finite-degree polynomials here suffice for the Galois application. The proof gives a uniform \(H\) even when those degrees increase after a larger conductor exponent has been fixed.\formalnote{12.23}{23}

\paragraph{Gauss valuations.}
For \(g(Y)=\sum_j a_j Y^j\ne0\) and \(r\ge0\) rational, put
\[
                    w_r(g)=\min_j\{v(a_j)+jr\}.
\]
Choosing \(t\in K\) with \(v(t)=r\), this is the minimum of the coefficient valuations of \(g(tY)\). After division by a coefficient realizing the minimum, its reduction is a nonzero polynomial over the residue field. Products of two such reductions are nonzero, so \(w_r(gg')=w_r(g)+w_r(g')\). Also \(w_r(g)\le v(g(y))\) whenever \(v(y)\ge r\). For finitely many \(g\) one can attain all these minima at the same \(y\) of valuation \(r\). Choose a nonzero residue avoiding the roots of their reduced polynomials and lift it to a unit \(u\); then take \(y=tu\). The residue field of \(K\) is the infinite field \(\overline{\mathbb F}_\ell\).\formalnote{12.24}{24}

Factor \(g=a\prod_\alpha(Y-\alpha)\) in \(K[Y]\), with roots repeated according to multiplicity. Then
\[
  w_r(g)=v(a)+\sum_\alpha\min\{r,v(\alpha)\}
        =w_0(g)+\sum_{v(\alpha)>0}\min\{r,v(\alpha)\}.
        \tag{S2}
\]
This accounts also for nonintegral roots, since their summands at \(0\) and at \(r\) coincide. For integral coefficients \(w_0(g)\ge0\). If \(v(y)=r\), ultrametricity gives equality \(v(y-\alpha)=\min\{r,v(\alpha)\}\) whenever \(v(\alpha)\ne r\). Therefore the excess of \(v(g(y))\) over \(w_r(g)\) comes entirely from roots with \(v(\alpha)=r\).

\paragraph{Finding a cyclotomic radius with few roots.}
Put \(r_h=1/\varphi(\ell^h)\) for \(h\ge1\). Every primitive \(\ell^h\)-th root \(\zeta\) has \(v(\zeta-1)=r_h\). Indeed its residue is \(1\), and for \(a\) prime to \(\ell\), \((\zeta^a-1)/(\zeta-1)\) is a unit, as its residue is \(a\) for positive \(a\). All primitive powers thus have the same valuation, and their differences from \(1\) have product up to sign \(\Phi_{\ell^h}(1)=\ell\). The numbers \(r_h\) decrease strictly to zero.

For fixed \(0<\delta<1\) and finite \(B\ge0\), choose \(H\) with
\[
    B<\delta\sum_{\substack{1\le h\le H\\r_h\le\rho}}\frac{1}{h+1}.
    \tag{S3}
\]
This is possible by divergence of the harmonic series. Let \(g\in\mathcal O[Y]\) satisfy \(w_\rho(g)\le B\), and write \(R_h\) for its multiset of roots of valuation \(r_h\). By (S2) and disjointness of the radii,
\[
         \sum_{\substack{h\le H\\1\le h,\ r_h\le\rho}}
               |R_h|r_h\le w_\rho(g)\le B .
\]
Some allowed \(h\) consequently satisfies \(|R_h|r_h\le\delta/(h+1)\).

\paragraph{Finding a point on that radius.}
Fix this \(h\). A root \(\alpha\in R_h\) could be arbitrarily close to a particular \(\zeta-1\), including equal to it. To average, truncate at \(1\). For \(F(Y)=(Y+1+\alpha)^{\ell^h}-1\), all numbers \(\zeta-1-\alpha\) with \(\zeta^{\ell^h}=1\) are integral roots. The coefficient of \(Y\) has valuation \(h\) since \(1+\alpha\) is a unit. The root formula for this monic polynomial gives
\[
 \sum_{\zeta\ {\rm primitive\ of\ order}\ \ell^h}
       \min\{v(\zeta-1-\alpha),1\}
       \ \le\ w_1(F)\ \le\ h+1 .
\]
Sum over \(\alpha\in R_h\). The number of primitive roots is \(1/r_h\), so for at least one \(\zeta\) their truncated distances, summed over \(\alpha\), are at most \(|R_h|r_h(h+1)\le\delta\). Every summand is nonnegative and \(\delta<1\); each valuation itself is therefore below \(1\) and truncation can be removed. With \(y=\zeta-1\), subtract the contributions in (S2) for \(r=r_h\) to obtain
\[
          v(g(y))-w_{r_h}(g)=
           \sum_{\alpha\in R_h}(v(y-\alpha)-r_h)\le\delta .
           \tag{S4}
\]
Both the choice of \(h\) among finitely many radii and the averaging at fixed \(h\) are needed: the first controls how many roots lie on the radius, and the second controls how closely one test point approaches them.\formalnote{12.25}{25}

\paragraph{Changing one coordinate at a time.}
Take \(\delta=\min(1,\eta)/(s_0+2)\), \(B=\beta+1\), and choose \(H\) by (S3). At an intermediate stage let \(b_i\) be the old value of \(v(P_i)\). To change a still untouched coordinate \(j\), specialize all other variables of \(P_i\) at their current, possibly \(i\)-dependent, values and call the resulting polynomial \(g_i(Y)\). The old \(j\)-th coordinate still has valuation at least \(\rho\), so \(w_\rho(g_i)\le b_i\). Provided \(\sum_i b_i\le B\), the product \(g=\prod_i g_i\) satisfies \(w_\rho(g)\le B\).

Apply (S4) to this product, giving a single \(y=\zeta_j-1\) with \(r_h=v(y)\le\rho\). Each difference \(v(g_i(y))-w_{r_h}(g_i)\) is nonnegative, and the sum of these differences is at most \(\delta\). Thus each new valuation is at most \(w_{r_h}(g_i)+\delta\le b_i+\delta\), and their sum is at most \(\sum_i b_i+\delta\). The sum costs only one \(\delta\) even if there are many constraints. After \(t\le s\) coordinates the sum is at most \(\beta+t\delta\le B\). After all coordinates, each increase is at most \(s\delta\le\eta\). This proves (S1).

\subsection{All powers of the same character}
Suppose that the individual initial bounds are \(v(P_i(x_i))\le b_i<\infty\) and that \(v(x_{iu})\ge r>0\). To control all \(1\le j\le m\) for \(m\ge1\), replace the family by
\[
      P_{ij}(X)=P_i\big((1+X_1)^j-1,\ldots,(1+X_s)^j-1\big).
\]
For each \(x=x_{iu}\) there exists \(x'\) with \((1+x')^j=1+x\) and \(v(x')\ge r/m\). For \(x=0\) take \(x'=0\). Otherwise the monic polynomial \((1+Y)^j-(1+x)\) has \(j\) integral roots whose product has valuation \(v(x)\). One root has valuation at least \(v(x)/j\). Choose such preimages separately for all \(i,j,u\). At the resulting points \(P_{ij}\) takes the original value \(P_i(x_i)\). Apply (S1) with \(\rho=r/m\) and \(\beta=m\sum_i b_i\). For a common tuple of roots of order dividing \(\ell^H\) this gives
\[
 v(\zeta_u-1)\le r/m,\qquad
 v\!\left(P_i((\zeta_u^{\,j}-1)_u)\right)
                \le v(P_i(x_i))+\eta\quad (1\le j\le m).
                \tag{S5}
\]
The exponent \(H\) is fixed by the numerical inputs. In particular none of the original \(P_i\) need yet have been chosen.

When \(r<1\), these powers retain bounded distance from \(1\). If \(\zeta_u\) has order \(\ell^h\), then \(m/\varphi(\ell^h)\le r<1\). Thus \(\ell^{v_\ell(j)}\le j\le m<\varphi(\ell^h)\le\ell^h\). Its \(j\)-th power has order \(\ell^{h-v_\ell(j)}>1\) and
\[
    v(\zeta_u^j-1)=\ell^{v_\ell(j)}v(\zeta_u-1)\le r .
    \tag{S6}
\]
This last bound controls the loss at the tame generators.\formalnote{12.26}{26}

\subsection{Obtaining the test point for a Galois lattice}
Let \(O\) be the integers of a finite Galois extension of \(\mathbb Q_\ell\). Normalize \(O\)-length by multiplying ordinary composition length by \(v(\varpi_O)\); then \(O/\ell O\) has length \(1\). Use the same normalization after extending \(O\). Fix the finite set \(S_0\) of bad places containing those above \(\ell\), the finite list of Selmer data and Stickelberger systems, and the element \(g_0\) of inertia above a bad place away from \(\ell\) used to define their respective allowances \(\beta_i\). For a Selmer lattice \(T\) of rank \(d\) with representation \(\sigma\), take a finite testing depth \(M>\beta_T+1\). The tied descent supplies one distinguished prime \(w_0\) in a set \(Q\) of auxiliary primes, characters \(\lambda_w:G_E\longrightarrow\mathbb Z/\ell^N\), and slopes \(s_w\) at the other primes that may depend on \(T\). The construction imposes \(q_w\equiv1\bmod\ell^N\) on residue cardinalities and the indicated Frobenius and inertia values below. Denote by \(s_0\) a uniform bound with \(2d|Q|\le s_0\) for every lattice, and with the same bound using the number of exponents for each system in place of \(d\).

Lift the values of the \(\lambda_w\) to exponents in \([0,\ell^N)\) and form
\[
             U_X(g)=\prod_{w\in Q}(1+X_w)^{[\lambda_w(g)]}.
\]
On specializing \(X_w=\zeta_w-1\) with \(\zeta_w^{\ell^N}=1\) it is a character, since reducing exponents modulo \(\ell^N\) then preserves products. At arbitrary \(X=x\) use the matrices \(U_x(g)\sigma(g)\) in the defining linear equations: the cocycle identity for each pair \(g,h\), zero on inertia outside the allowed set, and the specified class relations.

All these equations can be placed in a single matrix \(A(X)\) with finitely many columns. To see the uniformity, take the common open kernel \(D\) of all \(\lambda_w\) and of \(\sigma\bmod\ell^M\). On \(D\), a solution with coefficients killed by \(\ell^M\) is additive. There are only finitely many continuous characters \(D\to\mathbb Z/\ell^M\) unramified outside the fixed finite set upstairs, by finiteness of abelian extensions of bounded exponent unramified outside a given finite set (equivalently here, the finite-coefficient \(H^1\) finiteness theorem). Take an open subgroup of \(D\) killing them all, and its normal core in \(G_E\). Every permitted solution vanishes there, since homomorphisms of an abelian group killed by \(\ell^M\) into \(\mathbb Z/\ell^M\) separate its points. For the latter assertion extend a nonzero character on a cyclic subgroup to the whole group using the injectivity of \(\mathbb Q/\mathbb Z\); its image lies in the \(\ell^M\)-torsion. The cocycle rule now makes the solution a function on the same finite quotient of \(G_E\) for every \(x\). Listing its \(d\) coordinates on each coset gives the columns of \(A\). The rows are the linear equations above after combining all terms referring to the same coset. There can be infinitely many rows.\formalnote{12.27}{27}

For a fixed specialization into a discrete valuation ring \(O'\), the length of the solution module in \((O'/\ell^M)^{k}\), where \(k\) is the number of columns, is the least valuation of a \(k\)-row minor of \(\binom{A(x)}{\ell^M I_k}\). Indeed the row span cuts out a finite-length quotient \(C\) of \(O'^k\) killed by \(\ell^M\). The solutions are \(\operatorname{Hom}_{O'}(C,O'/\ell^M)\), which has the same length as \(C\) on each cyclic summand. Smith normal form identifies that length with the valuation of the ideal of maximal minors. Even with an infinite row family the span is finitely generated, and multilinearity identifies the ideal with the ideal generated by the original row minors. Thus any minor polynomial gives an upper bound at every specialization, and some minor gives equality at a chosen one. The same reasoning applies after adding the tied linear equations.\formalnote{12.28}{28}

Let \(a_{w_0}=1\), \(a_w=s_w\) elsewhere, and specialize on the line \(X_w=a_wY\). At \(Y=0\) impose all ties except the one at \(w_0\) and its conjugate. The bound \(\beta_T\) includes the cost of omitting these two conditions. Take a minor for this tied system attaining the length at zero and denote its restriction to the line by \(p_0(Y)\). Then \(w_r(p_0)\le v(p_0(0))\le\beta_T\).

At \(w\ne w_0\) let \(\varphi,\gamma\) be the chosen Frobenius and tame inertia generator, and use their conjugate inverses at \(\bar w\). On the line the twisting scalars on these elements are \(1+y\) and \(1+s_wy\). For \(v(y)\ge r\), \(0<r\le1\), \(M\le Nr\), their reduction modulo \(\ell^M\) defines a character. Indeed for \(v(\alpha-1)\ge r\), the geometric sum \(\sum_{i<\ell^N}\alpha^i\) is divisible by \(\ell^M\). Factor it into the \(N\) sums \(\sum_{i<\ell}(\alpha^{\ell^j})^i\). Each such sum has valuation at least \(r\) because its terms are all congruent to \(1\) modulo the ideal of elements of valuation at least \(r\), which contains \(\ell\). Multiplying by \(\alpha-1\) also gives \(\alpha^{\ell^N}=1\bmod\ell^M\).

Write \(\Phi=(1+y)\sigma(\varphi)\). Local unramifiedness of \(\sigma\) makes the twisted action of inertia scalar. For every finite-depth cocycle \(z\) the affine inertia action \(u\mapsto\alpha u+z(\gamma)\) has exponent dividing \(\ell^N\), by the same geometric sum. The image of wild inertia in this affine action is trivial: the action is continuous on a finite module and wild inertia is pro-\(p\) for the residue characteristic \(p\ne\ell\). On the tame quotient conjugation by Frobenius raises inertia to \(q_w\) (or its inverse in the opposite Frobenius convention), hence acts identically on this image since \(q_w\equiv1\bmod\ell^N\). Commuting the two affine actions and evaluating at \(0\) gives
\[
                  (\Phi-I)z(\gamma)=s_wy\,z(\varphi).
\]
At a nonzero slope the chosen Frobenius acts trivially on \(T/\ell^M\), so \(s_w(\Phi-I)z(\varphi)=s_w y z(\varphi)\) in both cases \(s_w=0,\ s_w\ne0\). The defect \(z(\gamma)-s_wz(\varphi)\) thus lies in \(\ker(\Phi-I)\). The conjugate-inverse computation gives the same conclusion. The quotient by all ties injects into the direct sum of these kernels. The normalized length of each kernel is bounded by \(v(\det(\Phi-I))\) by Smith normal form. At the first prime the corresponding calculation at \(y=0\) gives the two determinant allowances included in \(\beta_T\), using the prescribed congruence to \(g_0\) and to its conjugate inverse.\formalnote{12.29}{29}

Multiply \(\det((1+Y)\sigma(\varphi)-I)\) over both choices above each \(w\ne w_0\) and call the product \(p_1(Y)\). The coefficient of \(Y^d\) of each factor is \(\det\sigma(\varphi)\), a unit. Thus \(w_r(p_1)\le2d|Q|r\le s_0r\). Choose a single \(y\) of valuation \(r\) attaining the two Gauss valuations \(w_r(p_0)\) and \(w_r(p_1)\) by residue avoidance. Enlarge the coefficient ring to contain \(y\), and take a minor \(P(X)\) attaining the length of the \emph{untied} system at \(x_w=a_wy\). By length additivity and the preceding kernel bound,
\[
      v(P(x))\le v(p_0(y))+v(p_1(y))\le\beta_T+s_0r.
      \tag{S7}
\]
The original integral polynomial \(P\) still bounds the finite-depth length at every root-of-unity specialization. Whenever that bound is strictly below \(M\) it also bounds the divisible-coefficient Selmer group. Indeed each locally constant cocycle takes finitely many values and so is torsion. Its subgroup killed by \(\ell^M\) is the finite-depth group by the coefficient identification \((\operatorname{Frac}(O')/O')[\ell^M]\simeq O'/\ell^M O'\). An element not killed by \(\ell^M\) would generate a cyclic subgroup whose \(\ell^M\)-torsion already has normalized length \(M\). This contradiction proves the passage to divisible coefficients.\formalnote{12.30}{30}

\subsection{The test point for a Stickelberger system}
Let \(\Theta_Q\) be the group-ring element of one of the systems, with its character \(\vartheta\), and a list of \(d'\) integer exponents. Extend \(U_X\) linearly with the representatives used for the class-group elements and write \(P(X)=U_X(\Theta_Q)\). At every \(\ell^N\)-torsion tuple this evaluates to the actual character value on \(\Theta_Q\). On the line on which only \(X_{w_0}=Y\) can be nonzero, call its restriction \(P_*(Y)\), and let \(R_0(Y)\) use \(\Theta_{\{w_0\}}\) on that same line.

By the system norm relation at the trivial character,
\[
  v(R_0(0))\le
     v(\mathbf 1(\Theta_\emptyset))+2d'\,v(\vartheta(g_0)-1)=:\beta_Y .
\]
In detail, take a depth \(M'>\beta_Y+1\) with \(M'\le N\) and take the first Frobenius congruent to \(\vartheta(g_0)\) in its \(\vartheta\)-value modulo \(\ell^{M'}\). Each Euler factor at the first prime is a product of two terms \(1-\vartheta(\varphi)^{\pm1}q_{w_0}^{-e}\), one of each sign. Both terms are congruent to \(1-\vartheta(g_0)^{\pm1}\) modulo \(\ell^{M'}\) since \(q_{w_0}\equiv1\bmod\ell^N\). For a nonempty exponent list the valuation of each latter term is below \(M'\) and equals \(v(\vartheta(g_0)-1)\); an empty list contributes \(1\).

At each auxiliary prime \(w\) other than the distinguished prime \(w_0\), write \(\varphi_w\) for its Frobenius and \(q_w\) for its residue cardinality. Put \(a=\vartheta(\varphi_w)\) and \(c=q_w^{-e}\). Both are \(\ell\)-adic units. At a root \(\zeta\) with \(\zeta^{\ell^N}=1\), the first auxiliary character has exponent \(1\) at \(\varphi_w\). Clearing the factor \(\zeta^{-1}\) in the paired Euler factors gives the integral polynomial
\[
        (1-ac(1+Y))(1+Y-a^{-1}c).
\]
Its value at \(\zeta-1\) is \(\zeta(1-ac\zeta)(1-a^{-1}c\zeta^{-1})\), and its top coefficient is a unit. Taking the product \(R_1\) over all pairs \((w,e)\) yields \(w_r(R_1)\le 2d'(|Q|-1)r\le s_0r\). The norm relation, now applied at each of the \(\ell^N\) torsion points on the line, gives the polynomial congruence
\[
    (1+Y)^{d'(|Q|-1)}P_*(Y)\equiv R_0(Y)R_1(Y)
                       \pmod{(1+Y)^{\ell^N}-1}.
\]
Indeed the difference vanishes at all \(\ell^N\) distinct roots, and the divisor is monic, so division preserves integral coefficients. Assume \(M'\le Nr\) and \(s_0r<1\). Choose \(y\) of valuation \(r\) attaining \(w_r(R_0R_1)\le\beta_Y+s_0r<M'\). The modulus at \(y\) has valuation at least \(Nr\) by the geometric-sum argument. The two sides of the congruence have equal valuations at \(y\) and \(1+y\) is a unit. Thus \(v(P_*(y))\le\beta_Y+s_0r\). Take all other coordinates zero. This gives precisely the second kind of initial test point required by (S5).\formalnote{12.31}{31}

\subsection{Choosing the exponent before the auxiliary primes}
The common tied-prime descent at finite coefficient depth supplies \(s_0\ge1\) before choosing the radius \(r\). It gives the prime data used in the preceding two subsections for every \(0<r\) with \(s_0r<1\) and every sufficiently large \(N\), including \(N\) large enough for \(M,M'\le Nr\). It is here that the auxiliary-prime existence, the duality bound for the terminal tied module, and the Frobenius prescriptions enter. For the Selmer data the initial level \(M\) also requires that the residue field of \(O\) have more than \(2d\,\operatorname{length}_O(O/\ell^M O)+1\) elements for each rank \(d\); this length in the cardinality condition is unnormalized.\formalnote{12.32}{32}

Given an error \(0<\varepsilon<1\), choose rational \(r>0\) with \(s_0r\le\varepsilon/2\). For \(s\) running over all data and systems, their numbers \(\beta_s\) are finite and already fixed. Use (S5) with original bounds \(b_s=\beta_s+1\), number of powers \(m\), and loss \(\eta=\varepsilon/2\) to fix \(H\). Only now take \(N\ge H\) sufficiently large and choose \(Q,\lambda_w,P_s,x_s\). Since \(v(P_s(x_s))\le\beta_s+s_0r\), (S5) gives
\[
       v(P_s((\zeta_w^j-1)_{w\in Q}))
                \le \beta_s+s_0r+\varepsilon/2\le\beta_s+\varepsilon
               \qquad(1\le j\le m).
\]
Set \(\psi(g)=\prod_{w\in Q}\zeta_w^{\lambda_w(g)}\). The formula is well defined because all the orders divide \(\ell^N\). The local class-field conditions on \(\lambda_w\) give an anticyclotomic character trivial above \(\ell\). On inertia at \(w\) only \(\lambda_w\) can contribute, and it is surjective. The inertia image under \(\psi^j\) is consequently generated by \(\zeta_w^j\); the conjugate prime has the same image. Their residue characteristics are different from \(\ell\), so any tame generator maps to a primitive generator of that image. Its difference from \(1\) has valuation \(v(\zeta_w^j-1)\le r\) by (S6). This also proves ramification at each of the prescribed primes for every \(j\). The sum over \(Q\cup\bar Q\) is at most \(2|Q|r\le s_0r\le\varepsilon\). The valuations of the system values have the required bound, and the strict inequalities \(\beta_T+\varepsilon<M\) allow the length conclusion from (S7) at all the same character powers. If the requested \(\varepsilon\) is at least \(1\), apply the argument with \(1/2\) and weaken the result.

For fixed auxiliary data this proves the simultaneous specialization, the polynomial comparisons and the tame-generator bound. Section~\ref{sec:tied-primes} constructs the bounded set of primes and their characters, and derives the length at the origin from the stated duality formula. Section~\ref{sec:critical-values} constructs the systems and proves their norm and value relations for the critical Hecke values.\formalnote{12.33}{33}


\section{The prime-field Jacobi sums and their tame valuations}
\label{sec:jacobi-valuations}

We need cyclotomic numbers with prescribed orders at the places over a prime \(p\). We construct them from Jacobi sums. In \S\ref{sec:derivative-cocycles} weighted products of their Galois conjugates will give cocycles with controlled ramification.
Fix a positive divisor \(M\) of \(p-1\), a place \(B\) of \(\overline{\mathbb Q}\) above \(p\), and an integer \(c\ge0\). The index \(x\) ranges over the classes in \(\mathbb Q/\mathbb Z\) killed by \(M\). Write \(v_B(p)=1\) and
let \(\{x\}\in[0,1)\) be the representative of \(x\in\mathbb Q/\mathbb Z\).
We construct elements \(J(x)\in\mathbb Q(\mu_M)^\times\) such that
\[
     v_B(J(x))=s_c(x),\qquad
     s_c(x)=\lfloor c\{x\}\rfloor-\lfloor c\{-x\}\rfloor .
\]
The same elements work simultaneously for all \(M\mid p-1\). We will
compute both their principal divisors and the inertia action on the roots
used in the derivative. The divisor records integral orders, while the exponent on an
\(m\)-th root is read modulo \(m\).\formalnote{12.34}{34}

\subsection{The valuation coordinate}
Let \(n=p-1\). At \(B\) choose the Teichm\"uller lift \(\omega(u)\in\mu_n\) of \(u\in\mathbb F_p^\times\), and choose a primitive \(p\)-th root \(\zeta_p\). For \(t\in\mathbb Z/n\), set
\[
             G(t)=\sum_{u\in\mathbb F_p^\times}\omega(u)^{-t}\zeta_p^u,
             \qquad S_B(t)=t/n \pmod{\mathbb Z}.
\]
For \(1\le t<n\), the prime-field Gross--Koblitz formula \cite[Theorem 1.7]{arith-gross-koblitz} gives
\[
            -G(t)=\pi^t\Gamma_p(t/n),\qquad
             \pi^{p-1}=-p,\quad \pi\equiv\zeta_p-1\pmod{(\zeta_p-1)^2}.
\]
The residue field in the theorem is \(\mathbb F_p\) because \(n\mid p-1\). The factor \(\Gamma_p(t/n)\) is a \(p\)-adic unit, so \(v_B(G(t))=\{S_B(t)\}\). Also \(G(0)=-1\); this covers the remaining character and the entire case \(p=2\). All sums are units away from \(p\): for \(t\ne0\), multiplicative character orthogonality gives \(G(t)\overline{G(t)}=p\), with both factors algebraic integers.

\phantomsection\label{route:22}\leanroute{Lean starts with an arbitrary generator of the multiplicative characters of \(\mathbb F_p\). It proves that \(t\mapsto v_B(G(t))\bmod\mathbb Z\) is an isomorphism by showing additivity and trivial kernel. Fixing the inverse Teichm\"uller generator identifies that isomorphism explicitly as \(t\mapsto t/(p-1)\). Only its fractional parts enter the subsequent valuations.\formalnote{12.35}{35}}

If \(\sigma\) acts on \(\mu_n\) by \(a\) and on \(\zeta_p\) by \(b\in\mathbb F_p^\times\), changing \(u\) to \(b^{-1}u\) gives \(\sigma G(t)=\omega(b)^{at}G(at)\). The sign in this formula and the sign in \(S_B\) both correspond to the exponent \(-t\) in \(G(t)\).

\subsection{The principal divisor, distribution and reduction at an auxiliary prime}

For \(x=S_B(t)\) define
\[
          j(t)=\frac{G(t)^c}{G(ct)},\qquad
          J(x)=\frac{j(t)}{j(-t)} .
\]
In particular \(J(0)=1\). The root-of-unity factor in the
Galois action on \(G\) cancels in \(j\), so \(\sigma j(t)=j(at)\).
If \(Mx=0\), then \(Mt=0\) in \(\mathbb Z/n\). Thus \(J(x)\) lies
in \(F_M=\mathbb Q(\mu_M)\), and \(\sigma_a J(x)=J(ax)\) with
\(\sigma_a\zeta_M=\zeta_M^a\). We also have
\[
       v_B(j(t))=c\{x\}-\{cx\}=\lfloor c\{x\}\rfloor,\qquad
       v_B(J(x))=s_c(x).
\]
This includes the case \(ct=0\). If \(\mathfrak p=B|_{F_M}\), then
\(p\) splits completely in \(F_M\), and these equalities give the
fractional ideal identity
\[
         (J(x))=\prod_{a\in(\mathbb Z/M)^\times}
                    (\sigma_a\mathfrak p)^{\,s_c(a^{-1}x)} .
\]
The inverse occurs because the order at \(\sigma_a\mathfrak p\) is
the order at \(\mathfrak p\) of \(\sigma_a^{-1}J(x)\). Negative
exponents can occur. Complex conjugation sends \(j(t)\) to \(j(-t)\)
after every embedding, so every conjugate of \(J(x)\) has absolute
value one.\formalnote{12.36}{36}

For \(r\mid M\) let \(b\) run through \(\mathbb Z/r\). The \(r\)
representatives \(\{x+b/r\}\) are the numbers
\((j+\{rx\})/r,\ 0\le j<r\), in some order. Hence their sum is
\(\{rx\}+(r-1)/2\). If \((c,r)=1\), multiplication by \(c\)
permutes the translates for \(cx\). Writing
\(\lfloor c\{y\}\rfloor=c\{y\}-\{cy\}\) and subtracting the same
formula for the negatives shows that
\[
                  \sum_b s_c(x+b/r)=s_c(rx).
\]
Apply the ideal formula at each \(\sigma_a\mathfrak p\), where
\(a^{-1}\) also permutes \(b\). The quotient
\(D=\prod_b J(x+b/r)/J(rx)\) has trivial fractional ideal. Its
conjugates have modulus one. For completeness, every \(D^k\) is an
algebraic integer of degree at most \([F_M:\mathbb Q]\) with all
conjugates of modulus one. The coefficients of its minimal polynomial
are bounded by the binomial coefficients of that degree bound, so only
finitely many such polynomials and roots occur. Two positive powers
of \(D\) coincide, giving the root-of-unity distribution relation.\formalnote{12.37}{37}

Fix a prime \(q\mid p-1\), a place \(A\mid q\), and \(b\) with
\(qb=0\) in \(\mathbb Q/\mathbb Z\). If \(b=S_B(t')\), then
\(\omega(u)^{-t'}\in\mu_q\) has residue one at \(A\) for every
\(u\ne0\). Thus \(G(t+t')\equiv G(t)\) modulo its maximal ideal.
These are \(A\)-units, since they divide \(p\). The same holds after
multiplying the indices by \(c\), or negating them. Taking the four
ratios therefore gives \(J(x+b)/J(x)\equiv1\) at \(A\).
This even applies if \(q\nmid M\), since both indices still occur
at level \(p-1\). The congruence compares unit values at the new
prime, while the displayed fractional ideal formula counts their
zeros and poles over \(p\).\formalnote{12.38}{38}

\subsection{Two actions of inertia on roots}

Let \(K\) be a local field with uniformizer \(\varpi\), residue
characteristic \(l\), and let \(I\) be its inertia group. For any
\(m\) prime to \(l\) and \(z\) with \(z^m\) an inertia-fixed unit,
every \(\tau\in I\) fixes \(z\). Indeed \(z\) is a unit, hence
\(\tau z/z\) is both an \(m\)-th root of unity and congruent to one,
which forces it to equal one. In contrast, for \(r^m=\varpi\) and
any \(\zeta\in\mu_m\), some \(\tau\in I\) satisfies \(\tau r=\zeta r\).
To prove this, the inertia orbit of \(r\) is a set of \(d\le m\)
roots, each of normalized order \(1/m\). Their product is fixed by
inertia and has order \(d/m\), an integer. Hence \(d=m\).
This orbit version of Eisenstein's argument works before choosing
a generator of the tame quotient. With \(m\) prime to \(l\), inertia
fixes \(\mu_m\), so \(\tau\mapsto\tau r/r\) is a surjective
homomorphism \(I\to\mu_m\), independent of the chosen root \(r\).\formalnote{12.39}{39}

Apply these two calculations at \(B\) with \(K=\mathbb Q_p\),
\(m\mid M\mid p-1\). Fix a primitive \(\zeta_m\) and choose
\(\tau_0\in I_B\) with \(\tau_0p^{1/m}=\zeta_m p^{1/m}\).
For any \(y(x)^m=J(x)\) put \(s=s_c(x)\). The \(m\)-th power of
\(y(x)/(p^{1/m})^s\) is a unit of \(\mathbb Q_p\) since
\(F_M\) embeds in \(\mathbb Q_p\) at \(B\). Hence
\(\tau_0 y(x)=\zeta_m^{s_c(x)}y(x)\), for every \(x\) at that level
and for every choice of the root. If \(\gamma\) acts on \(\mu_M\)
by \(a\), then \(\gamma y(x)\) is an \(m\)-th root of \(J(ax)\).
Consequently
\[
       \frac{\gamma y(x)}{\tau_0\gamma\tau_0^{-1}y(x)}
          =\zeta_m^{\,a s_c(x)-s_c(ax)} .
\]
Applying \(\gamma\) to \(\zeta_m\) supplies the coefficient \(a\).
The minus sign comes from the denominator. This identity gives the
bottom term (KC5) with the character and commutator convention of
Section~\ref{sec:derivative-cocycles}.\formalnote{12.40}{40}

At an auxiliary \(A\mid q\), \((q,M)=1\), the neighboring derivative
needs an element acting on \(\mu_q\) by a prescribed \(b\) and
trivially on \(\mu_M\). Here \(\Phi_q(1+X)\) is Eisenstein of
degree \(q-1\) over \(\mathbb Q_q\). Its roots \(\zeta_q^u-1\) for
\(u\ne0\) have order \(1/(q-1)\). The same orbit-product argument
shows they form a single inertia orbit, whereas inertia fixes
\(\mu_M\) by the separate-reductions argument above. Taking
\(b=v^{-1}\) gives exactly the action \(x+vu/q\mapsto x+u/q\)
used to compare the neighboring derivative with arithmetic Frobenius.
At higher \(q\)-power levels one uses \(\Phi_{q^j}(1+X)\), whose
reduction is \(X^{\varphi(q^j)}\) and whose constant term is \(q\).
Its orbit has \(\varphi(q^j)\) elements, so each unit class occurs
at each finite level. The inverse images in the compact inertia
group of the actions by a fixed \(q\)-adic unit at successive levels
are nested nonempty closed sets, and have a common element. In
particular this supplies the local lift used to choose \(\gamma\)
in the descent.\formalnote{12.41}{41}

The local inputs here are the valuation of a uniformizer, the
unramified fixed-field description of inertia, and the identification
of a global decomposition group with the absolute group of its local
field. The finite-field character sums, floor valuations, principal
units and prescribed inertia exponents were computed above for any
prime \(p\) with \(M\mid p-1\). Choosing a \(p\) whose Frobenius also
detects a specified global cocycle, and using global reciprocity
in the ensuing descent, are the separate steps of
Section~\ref{sec:derivative-cocycles}.


\section{The derivative cocycles and the number of odd Tate classes}
\label{sec:derivative-cocycles}

One input to the determinant in Section~\ref{sec:trace-ideal} is an upper bound for the length of the odd Tate cocycles. In the classical ideal class group setting, Rubin bounds the order of an odd character component by a generalized Bernoulli number under the local hypotheses of \cite[Chapter~III, Theorem~4.6]{rubin-euler-systems}. We need a bound on the number of generators for the Tate-twisted cocycles of Section~\ref{sec:congruence-detail}. Let \(O\) be the integers of a finite extension of \(\mathbb Q_\ell\), let \(N_E\) be a positive integer (taken to be \(q_1d\) in the application), let \(\psi\) be an \(O\)-valued Dirichlet character of conductor dividing \(N_E\), and put \(\alpha=\epsilon^e\psi\), \(\Upsilon=\epsilon\alpha\), where \(\epsilon\) is the cyclotomic character. Assume \(e>0\) and \(\alpha(c_\infty)=-1\) at complex conjugation. The component \(Z_\nu\) consists of continuous cocycles for \(\alpha^{-1}\) in \(O/\ell^\nu\) that vanish at \(c_\infty\), are unramified away from \(\ell\) over \(\mathbb Q(\mu_{N_E})\), and vanish at the indicated local cyclotomic kernels above \(\ell\). Let \(X\) be the positive integer with the finite Stickelberger congruence in Section~\ref{sec:congruence-detail}. We will prove that, for \(\ell>N_E+e+3\), there exist \(1+v_\ell(X)\) elements of \(G_{\mathbb Q}\) whose evaluations detect every member of \(Z_\nu\).

The word \emph{derivative} here refers to weighted translates of Jacobi sums. In the trace argument the variables \(X_j\) occur in relations on an abstract group algebra and in operators on an integral lattice. The finite quotient measuring odd Tate cocycles enters that argument through an upper bound for the colength of a row space. The construction below applies to the cyclotomic and Dirichlet characters themselves.

\subsection{A Jacobi family with a specified valuation}

Section~\ref{sec:jacobi-valuations} computes the leading term with a choice of generator and checks the inertia sign directly.

Fix an integer \(c\ge1\), \(m=\ell^n\), a finite set \(P\) of primes \(q\equiv1\bmod cNm\), a prime \(p\equiv1\bmod N_Em\prod_{q\in P}q\), and a place \(B\) above \(p\). At the levels \(M\mid p-1\) we shall use a family \(J(x)\) for \(x\in(\mathbb Q/\mathbb Z)[M]\). Write \(\{x\}\in[0,1)\) for the fractional representative and \(s_c(x)=\lfloor c\{x\}\rfloor-\lfloor c\{-x\}\rfloor\). Here and in the Frobenius calculation below, Frobenius acts as the residue cardinality power on residue classes.

Choose a primitive additive character of \(\mathbb F_p\) and a generator of its multiplicative character group. Let \(G(t)\) be the Gauss sum for the \(t\)-th power of that generator, and normalize \(v_B(p)=1\). The map \(t\mapsto v_B(G(t))\bmod\mathbb Z\) identifies \(\mathbb Z/(p-1)\) with the \((p-1)\)-torsion in \(\mathbb Q/\mathbb Z\). Indeed inertia acts on \(G(t)\) by a root of unity multiplicative in \(t\), so \(G(t+u)/(G(t)G(u))\) is fixed by inertia and has integral valuation. A nontrivial character sum has positive valuation because its additive character is \(1\) in the residue field and the sum of the nontrivial multiplicative character is zero. Its inverse-character Gauss sum, with inverse additive character, also has positive valuation, and their product is \(p\). Hence \(0<v_B(G(t))<1\) for \(t\ne0\). For \(t=0\) the Gauss sum is \(-1\). This gives additivity, injectivity, and the claimed choice of representatives.

Write \(x=v_B(G(t))\bmod\mathbb Z\), set \(j(t)=G(t)^c/G(ct)\) and \(J(x)=j(t)/j(-t)\). The order of \(j(t)\) is \(c\{x\}-\{cx\}\), hence \(v_B(J(x))=s_c(x)\). Changing the additive character multiplies \(G(t)\) by a factor multiplicative in \(t\), which cancels in \(j\). Thus an element acting as \(a\) on \(\mu_M\) sends \(J(x)\) to \(J(ax)\). The Gauss sums divide \(p\), so \(J\) is a unit away from \(p\). Complex conjugation interchanges \(j(t)\) and \(j(-t)\), so every complex absolute value of \(J\) is one.

For \(r\mid M,\ (r,c)=1\), the quotient
\[
                    \frac{\prod_{b\bmod r}J(x+b/r)}{J(rx)}
                    \quad\hbox{is a root of unity}.                    \tag{KC1}
\]
At \(B\) this follows at the level of valuations from
\(\sum_b s_c(x+b/r)=s_c(rx)\), the floor distribution identity. The same holds at every place over \(p\) by Galois equivariance and reindexing \(b\). At all other places it is a unit, and all its conjugates have absolute value one, so Kronecker's theorem \cite[Corollary~5.6]{legacy-milne-ant} applies. Also \(J(x+b/q)/J(x)\equiv1\) at every place over \(q\mid p-1\) when both arguments occur at the chosen levels. In the Gauss sums the added \(q\)-torsion character reduces to the trivial character at that place, and the denominators are units there.\formalnote{12.42}{42}

\subsection{Commutator logarithms and the cocycle identity}

Choose \(m\)-th roots \(y(x)\) of \(J(x)\), a primitive \(m\)-th root \(\zeta\), and \(\gamma\in G_{\mathbb Q}\). For each \(g\) define \(F_g(x)\bmod m\) by
\[
                 \zeta^{F_g(x)}
                  =\frac{\gamma y(x)}{(g\gamma g^{-1})y(x)} .
\]
The two automorphisms agree on \(J(x)\) since their actions on \(\mu_M\) agree. They also agree on \(\mu_m\), so changing the chosen root \(y(x)\) leaves the ratio unchanged. Equation (KC1) applied to the roots gives
\[
                    \sum_{b\bmod q}F_g(x+b/q)=F_g(qx).
\]
The discrepancy in the product of the roots has finite order, on which \(\gamma\) and \(g\gamma g^{-1}\) act alike.

For \(q\in P\) choose a surjection \(L_q:\mathbb F_q^\times\to\mathbb Z/m\). Allow a shift \(d_q\) and form the commuting translation operators
\[
 \begin{aligned}
 D_{q,d_q}F(x)&=\sum_{u\in\mathbb F_q^\times}(L_q(u)-d_q)F(x+u/q),\\
 \Lambda F&=\sum_{b\in(\mathbb Z/N_Em)^\times}b^e\psi(b)F(b/N_Em).
 \end{aligned}
\]
Put \(\kappa_S(g)=\Lambda\prod_{q\in S}D_{q,0}F_g\) for \(S\subset P\). In this formula all shifts may be changed arbitrarily. To prove it, induct on \(|S|\) and change one \(d_q\) at a time. Its change contributes a scalar times the remaining derivative applied to \(F_g(qx)-F_g(x)\), by the distribution identity and omission of the zero residue in \(D\). In the \(qx\) term reindex the remaining translates by multiplication by \(q\). On the base character sum \(q\equiv1\bmod N_Em\) and on the other translates this changes only the shifts, since \(L_r(qu)=L_r(q)+L_r(u)\). The smaller-set induction kills the difference.

For \(g\) acting as \(a\) on \(\mu_{N_Em\prod_{q\in S}q}\), \(g y(x)/y(ax)\) is an \(m\)-th root of unity. Applying \(g\) to the defining ratio and cancelling equal actions on this root gives
\[
                   a F_h(x)=F_{gh}(ax)-F_g(ax).
\]
To see the direction of the action, set \(\Lambda_{S,d}=\Lambda\prod_{q\in S}D_{q,d_q}\). For \((F\circ a)(x)=F(ax)\), replacing \(b,u\) by \(ab,au\) gives
\[
       \alpha(g)\Lambda_{S,d}(F\circ a)=
           \Lambda_{S,d+L(a)}F,\qquad L(a)_q=L_q(a\bmod q).
\]
Apply this identity to the right-hand side of the preceding difference and use shift independence. It follows that
\[
           \kappa_S(gh)=\kappa_S(g)+\Upsilon(g)\kappa_S(h),
           \qquad \kappa_S(\gamma)=0.                                  \tag{KC2}
\]
Each finite sum is locally constant because it uses finitely many algebraic numbers. At a prime away from \(\ell,p,S\), inertia over \(\mathbb Q(\mu_{N_E})\) fixes the relevant roots of unity and the \(J(x)\) are units. Their \(m\)-th roots lie in the maximal unramified extension there after the \(J(x)\) are fixed by inertia. This follows from Hensel's lemma for units, since \(m\) is prime to the residue characteristic. Applying the same fact at conjugate places makes each commutator logarithm vanish. Thus the only allowed additional ramification over \(\mathbb Q(\mu_{N_E})\) is at \(p\) and \(S\).\formalnote{12.43}{43}

\subsection{A Frobenius value as a tame inertia value}

Let \(q\in P\setminus S\), \(M=N_Em\prod_{r\in S}r\), \(A\) a place over \(q\), and \(\phi\) an arithmetic Frobenius at \(A\). Choose \(v\in\mathbb F_q^\times\) with \(L_q(v)=1\). There is \(\tau\in I_A\) fixing \(\mu_M\) and acting on \(\mu_q\) by \(v^{-1}\), since \((q,M)=1\). For \(Mx=0\) we claim
\[
          \sum_{u\in\mathbb F_q^\times}L_q(u)F_\tau(x+u/q)=F_\phi(qx).
                                                                        \tag{KC3}
\]
Put \(Y_b=y(x+b/q)\), \(y_*=y(qx)\), and choose integer lifts \(l(u)\) of \(L_q(u)\). All \(Y_b^m\) are units congruent to \(Y_0^m\) at \(A\). For \(\delta\in G_{\mathbb Q}\) form
\[
 W_\delta=
 \frac{\prod_{b\bmod q}Y_b}{(\delta\phi\delta^{-1})Y_0}
 \prod_{u\in\mathbb F_q^\times}
 \left(\frac{(\delta\tau\delta^{-1})Y_{vu}}{Y_u}\right)^{l(u)} .
\]
At \(\delta A\) the inertia acts trivially on residues, and the Frobenius on a unit raises the residue to its \(q\)-th power. The resulting product telescopes to \(1\): reindex \(u\mapsto vu\) and use \(l(vu)-l(u)=1+mk_u\). The remaining powers of \(Y_b^m/Y_0^m\) are principal units, and summing over the \(q-1\) nonzero residues accounts for the \(q\)-th power of \(Y_0\). Hence \(W_\delta\equiv1\) at \(\delta A\). In addition \(W_1\) has finite order. Indeed \(\prod_bY_b/y_*\) has finite order by (KC1), and both \(\phi Y_0/y_*\) and \(\tau Y_{vu}/Y_u\) are \(m\)-th roots of unity. Consequently \(\gamma W_1\) is a power of \(W_1\) and is also \(1\) modulo \(A\).

The quotient of the two roots of unity asserted equal in (KC3) is \(\gamma W_1/\gamma W_{\gamma^{-1}}\). This follows by expanding the displayed product: in the factors \(\tau Y_{vu}/Y_u\) and \(\phi Y_0/y_*\), conjugating \(\gamma\) leaves its action unchanged on \(m\)-th roots of unity. The quotient is therefore a principal unit at \(A\). Roots of unity of order prime to \(q\) have distinct reductions there, proving (KC3). Apply \(\Lambda\) and the derivatives at \(S\); reindexing multiplication by \(q\) removes shifts as above. We obtain
\[
                          \kappa_S(\phi)=\kappa_{S\cup\{q\}}(\tau).      \tag{KC4}
\]

At \(B\) over the distinguished prime \(p\), choose \(\tau_0\) sending \(p^{1/m}\) to \(\zeta p^{1/m}\). If \(\gamma\) acts as \(a\) on \(\mu_{N_Em}\), divide each \(y(x)\) by \((p^{1/m})^{s_c(x)}\). Its \(m\)-th power is a unit at \(B\). Applying the same inertia computation gives \(F_{\tau_0}(x)=a s_c(x)-s_c(ax)\). Thus, if \(\kappa_\varnothing\) vanishes on \(I_B\), reindexing the base sum gives
\[
           (\Upsilon(\gamma)-1)\sum_{u\in(\mathbb Z/N_Em)^\times}
                     u^e\psi(u)s_c(u/N_Em)=0\quad\hbox{in }O/mO .       \tag{KC5}
\]
The auxiliary relation (KC4) applies at every choice of \(A\mid q\), with an inertia element for that choice. The terminal relation (KC5) uses the single place \(B\mid p\) used to label the Gauss sums. Its conjugate valuations are determined by the Galois action on \(x\).\formalnote{12.44}{44}

\subsection{Detection by global reciprocity}

Write \(F=\mathbb Q(\mu_{N_E\ell})\), take \(q\equiv1\bmod N_E\ell\), and choose a Frobenius \(\phi\) at \(A\mid q\). Suppose \(x\in Z_\nu\) and a continuous \(\Upsilon\)-cocycle \(k\) are both killed by \(\ell\). Allow \(k\) to ramify over \(\mathbb Q(\mu_{N_E})\) only at \(\ell\) and a finite set \(T\ni q\). Assume \(x\) vanishes on all decomposition groups over \(T\setminus\{q\}\). Choose \(a\equiv1\bmod N_E\) prime to \(\ell\) and suppose \(2,\varphi(N_E\ell),a^e-1\) are units in \(O\). We prove
\[
                          x(\phi)\ne0\quad\Longrightarrow\quad
                          k|_{I_A}=0 .                                  \tag{KC6}
\]

Fix \(\tau\in I_A\) and let \(v=k(\tau)\). Choose an \(\mathbb F_\ell\)-basis \(b_i\) of the cyclic \(O\)-module \(Ov\). Read each coordinate on this module as an \(\ell\)-th root of unity, and extend it to an additive map from the coefficient module to \(\overline{\mathbb Q}^{\times}\), using divisibility of that group. On \(G_F\) the \(\ell\)-torsion coefficients and \(\mu_\ell\) have trivial action, so each coordinate of \(k\) is a character. By Hilbert 90 it is \(g\mapsto g(r_i)/r_i\), with \(r_i^\ell\in F^\times\). For lifts \(\sigma_u\) of \(u\in(\mathbb Z/N_E\ell)^\times\), set \(m_{iu}=v_A(\sigma_u r_i^\ell)\), normalized at \(q\). These are integers since \(q\) splits in \(F\).

Apply global Artin reciprocity for \(F\) to \(r_i^\ell\) and the character obtained by restricting \(2(a^e-1)x\) to \(G_F\), testing its values by additive functionals if necessary. At \(\ell\) the cyclotomic kernel is killed by definition of the component. On the residual abelian cyclotomic quotient, commuting with an inertia element of cyclotomic value \(a\) makes the factor \(a^e-1\) kill the restrictions over \(F\). The factor \(2\) kills contributions at real places. At a finite place away from \(\ell,q\), either \(x\) is zero on the decomposition group or the Kummer character is unramified and \(r_i^\ell\) has order divisible by \(\ell\). The coefficient at that place is then zero. At the places over \(q\), the conjugation law for \(x\) contributes precisely the weight \(\alpha(u)=u^e\psi(u)\). The product formula in additive notation gives
\[
         \left(\sum_u m_{iu}\alpha(u)\right)2(a^e-1)x(\phi)=0.          \tag{KC7}
\]

To recover \(v\), choose \(\pi^\ell=q\) and define \(j\bmod\ell\) by \(\tau\pi/\pi=\zeta_\ell^j\). The quotient \(\sigma_u r_i/\pi^{m_{iu}}\) has valuation zero, so its inertia difference is a principal unit. Its difference is also an \(\ell\)-th root of unity, hence trivial. The cocycle conjugation law and the action \(u\) on \(\mu_\ell\) give the coordinate identity
\[
                   v=\alpha(u)\sum_i (j m_{iu}\bmod\ell)b_i
                   \qquad \text{for each }u .
\]
If \(x(\phi)\ne0\), every scalar \(\sum_u m_{iu}\alpha(u)\) belongs to the maximal ideal of \(O\), by (KC7). Sum the coordinate identity over \(u\). Its left-hand side is \(\varphi(N_E\ell)v\) and its right-hand side is in the maximal ideal times \(Ov\). The totient is a unit, so Nakayama applied to \(Ov\) forces \(v=0\). This proves (KC6) for every inertia element, including when \(O/\ell O\) has nilpotents.\formalnote{12.45}{45}

\subsection{Descent on the number of auxiliary primes}

Take \(k=1+v_\ell(X)\). Since \(\ell>N_E+e+3\), choose \(a\equiv1\bmod N_E\) whose residue generates \(\mathbb F_\ell^\times\). Then \(a^e-1\) and \(a^{e+1}-1\) are units, as are \(2\) and \(\varphi(N_E\ell)\). Choose \(\gamma\) over \(\ell\) acting on cyclotomic roots with \(\Upsilon(\gamma)=a^{e+1}\). In (KC1)--(KC5) take \(c=a,\ m=\ell^k\). Suppose for contradiction that no list of \(k\) group elements detects \(Z_\nu\).

For any set \(T\) of at most \(k\) primes \(q\equiv1\bmod N_E\ell\), choose one Frobenius at each. By assumption some nonzero component cocycle vanishes at all those elements. Multiplying it by the largest power of \(\ell\) that keeps it nonzero produces an \(\ell\)-torsion witness \(x\). Such an \(x\) also vanishes on inertia at those primes. Their Frobenii act trivially on its coefficients. Its vanishing therefore extends to conjugate Frobenii by the cocycle identity, and then to every decomposition group over \(T\).

Call \(T\) tight when its primes are \(1\bmod N_E\ell\) and every \(\ell\)-torsion \(\Upsilon\)-cocycle in \(O/\ell^k\) unramified off \(\ell,T\) over \(\mathbb Q(\mu_{N_E})\) is already unramified off \(\ell\). The empty set is tight. If \(T\) is tight, a witness \(x\) vanishes over \(T\), and \(x(\phi_r)\ne0\) at a new split prime \(r\), (KC6) shows \(T\cup\{r\}\) is tight. Indeed each eligible dual cocycle must vanish on inertia over \(r\) by (KC6), so the original tightness applies.

There is a finite reservoir \(P\) of primes \(1\bmod aN_E\ell^k\) which detect any pair consisting of a nonzero \(x\in Z_\nu\) and a \(\Upsilon\)-cocycle \(y\) into \(O/\ell^k\), unramified off \(\ell\), visible on \(G_{\mathbb Q(\mu_{aN_E\ell^k})}\). To see that \(x\) is visible there as well, suppose it factors through the abelian cyclotomic quotient. Comparing \(x(c_\infty g)\) and \(x(gc_\infty)\), using \(x(c_\infty)=0\), gives \(2x(g)=0\), hence \(x=0\). Thus the zero subgroups of both restrictions are proper. Their union does not cover the group. Apply Chebotarev on a common finite quotient to find a split prime with a Frobenius outside both kernels, choosing the place above it to obtain the required representative. At fixed coefficient depths the cocycle spaces with the stated ramification are finite, by finiteness of bounded-degree extensions with fixed ramification, so choosing one prime per pair gives a finite \(P\). Next choose a nonzero \(\ell\)-torsion witness and use the same argument with modulus \(N_E\ell^k\prod_{r\in P}r\) to find \(p\) whose Frobenius detects it. In particular \(\{p\}\) is tight. All primes now satisfy the hypotheses for (KC1)--(KC5).

We prove by descending induction on \(i=k,k-1,\ldots,0\) that
\[
 \begin{gathered}
   S\subset P,\quad |S|\le i,\quad S\cup\{p\}\ \text{tight}\\
   {}\Longrightarrow \ell^i\kappa_S
            \ \text{is unramified off }\ell\text{ over }\mathbb Q(\mu_{N_E}).
 \end{gathered}                                                         \tag{KC8}
\]
The case \(i=k\) follows because the coefficients are killed by \(\ell^k\). Suppose the assertion holds at \(i+1\) and take \(S\) at \(i\). By tightness it suffices to prove \(\ell^{i+1}\kappa_S=0\). Otherwise \(y=\ell^{i+1}\kappa_S\) is nonzero and unramified off \(\ell\) by the induction hypothesis applied to \(S\). It is visible on the required cyclotomic kernel, since \(y(\gamma)=0\) and \(\Upsilon(\gamma)-1\) is a unit. This follows by the same commuting-elements argument as for \(x\). Since \(|S\cup\{p\}|\le i+1\le k\), choose a witness \(x\) trivial above this set. The reservoir gives \(r,\phi_r\) with both \(x(\phi_r)\ne0\) and \(y(\phi_r)\ne0\), so \(r\) is new and \(S\cup\{p,r\}\) is tight. The induction hypothesis at \(S\cup\{r\}\) makes \(\ell^{i+1}\kappa_{S\cup\{r\}}\) unramified at \(r\). But (KC4) identifies \(y(\phi_r)\) with one of its inertia values. The contradiction proves (KC8).

At \(i=0\) the empty derivative is unramified at \(p\). Equation (KC5) gives zero for the floor sum in \(O/\ell^k\), since \(\Upsilon(\gamma)-1\) is a unit. The finite Stickelberger congruence then gives \(\ell^k\mid(a^{e+1}-1)X\), hence \(\ell^k\mid X\) in \(O\). But \(v_O(\ell^k)=k v_O(\ell)>v_O(X)\). This proves the promised detection bound. Notice that (KC8) descends both in the annihilating power and through the possible sets of auxiliary primes. The final assertion at the empty set alone would not propagate through (KC4).\formalnote{12.46}{46}

\subsection{Factors of the determinant}

Section~\ref{sec:congruence-detail} gives \(4X^2Z_\nu=0\). A detected component embeds in the product of \(k\) copies of \((O/\ell^\nu)[4X^2]\). Thus with ordinary composition length and \(v_O\) normalized at the uniformizer,
\[
       {\rm len}_O Z_\nu
            \le(1+v_\ell(X))\min\{\nu v_O(\ell),\,v_O(4X^2)\}.
\]
For \(G_E\), if \(E\subseteq\mathbb Q(\mu_{N_E})\) is CM, the integral character transform gives \(Y=16\varphi(N_E)X^2\) as annihilator. When \(\ell>(N_E+t+3)^2\), \(e<t\), the transform uses at most \(\varphi(N_E)^2(1+v_\ell(X))\) evaluations on \(G_E\), by indexing a residue class and a character for each component test. The same proof bounds its length by this number times \(\min\{\nu v_O(\ell),v_O(Y)\}\). The source takes the looser uniform bound \(v_O(Y^{(N_E+t+3)^{20}})\), which also works in the small-prime range by the evaluation bound of \S\ref{sec:ray-detection}. At \(e=0\) class-field theory supplies the class number. Multiplying the positive bounds over \(e\) gives one integer suitable for all the character rows.\formalnote{12.47}{47}

On the main representation the rows instead use the crystalline Selmer length and the local loss. In the notation of Section~\ref{sec:trace-ideal}, the exact cost inequality before normalization is
\[
 \begin{aligned}
 v_O(d_\pi) \le {}& b_0+\sum_{e<t} b_e+
                  t\sum_{v\in W}v_O(o_v(1))\\
                 &+(2t|s|+|J|+1)v_O(C).
 \end{aligned}
\]
The \(b_e\) use the odd length just bounded and the local crystalline loss for \(O(-e)\). The automorphic Galois representations enter in realizing these character rows and the main block for the same family of traces, in obtaining the specified action of inertia and Frobenius, and in bounding \(b_0\). The powers of distinct auxiliary variables are combined by squarefree-coefficient extraction. The local root-chain argument supplies the polynomial \(H_j\) without a further cocycle-length term. These are the links between the cyclotomic derivative bound and the integral projection.\formalnote{12.48}{48}


\section{Ray-class detection at small coefficient primes}\label{sec:ray-detection}
Fix \(E\subset\mathbb Q(\mu_{N_E})\), a rational prime \(\ell\), and \(G_E=\operatorname{Gal}(\overline{\mathbb Q}/E)\). Let \(O\) be the integers in a finite extension of \(\mathbb Q_\ell\). Write \(A=O/\ell^\nu O\), with \(G_E\) acting by \(\chi_\ell^{-e}\) for \(e\ge0\). We will choose a single set of elements of \(G_E\) at which evaluation detects every continuous \(A\)-valued cocycle unramified away from \(\ell\). The set depends only on \(E,N_E,\ell\). If \(N_E\le y\) and \(\ell\le y^2\), its cardinality is at most \(y^{20}\). The objects being counted here are Galois elements. The bound is uniform in \(\nu,e,O\) and involves no norm bound for a prime ideal.

\subsection{The abelian quotient that suffices for one layer}
Put \(M=N_E\ell^2\), \(K=\mathbb Q(\mu_M)\), \(d=[K:\mathbb Q]=\varphi(M)\), and \(\Sigma=\{v:v\mid\ell\}\) in \(K\). Since \(M>2\), all archimedean places of \(K\) are complex. In the finite ideles divide by the principal ideles, by the units outside \(\Sigma\), and by \(K_v^{\times M}\) at \(v\in\Sigma\). Call the resulting abelian group \(B\). The image of the full integral idele units has cardinality dividing \(\prod_{v\in\Sigma}[K_v^\times:K_v^{\times M}]\). After dividing by this image, the quotient is a quotient of the ordinary ideal class group. Consequently
\[
         |B|\ \mid\ h_K\prod_{v\in\Sigma}[K_v^\times:K_v^{\times M}].
\]
Artin reciprocity and the existence theorem identify this finite idele quotient with a Galois quotient. A continuous additive homomorphism from \(G_K\) killed by \(M\) and unramified off \(\ell\) factors through it. Indeed its Artin pullback kills the outside units by the inertia condition and kills the inside \(M\)-th powers by additivity. Archimedean reciprocity contributes nothing at the complex places. Choose one lift \(b_i\in G_K\) for each generator in a decomposition \(B=\prod_{i=1}^s\mathbb Z/m_i\), \(m_i\ge2\). The values at the \(b_i\) detect every such homomorphism \cite[\href{https://www.jmilne.org/math/CourseNotes/CFTc.pdf}{V, Theorems 5.3 and 5.5}]{arith-cft}.\formalnote{12.49}{49}

Here is a bound for \(s\) including \(\ell=2\). For any finite completion \(K_v\), the normalized absolute value and multiplication by \(M\) give
\[
 [K_v^\times:K_v^{\times M}]
     =M\,|\mu_M(K_v)|\,|M|_v^{-1}\le M^2|M|_v^{-1}.
\]
The valuation group contributes the first \(M\). On a sufficiently deep group of principal units, the logarithm changes the power map into multiplication by \(M\) on an open additive ideal, whose cokernel has order \(|M|_v^{-1}\) and whose kernel is trivial. On the finite group of remaining unit cosets the cokernel and kernel of an endomorphism have the same order. Multiplicativity of the cokernel-to-kernel ratio in this short exact sequence gives the formula for all units, with their actual kernel \(\mu_M(K_v)\). One may take the logarithm arbitrarily deep at residue characteristic two. Since \(\prod_{v\in\Sigma}|M|_v^{-1}\le M^d\), the product of the displayed indices is at most \(M^{2|\Sigma|+d}\).\formalnote{12.50}{50}

For the class number the community proof invokes the effective bound \(h_K\le4^d|\Delta_K|\).
Indeed Minkowski's theorem \cite[\href{https://www.jmilne.org/math/CourseNotes/ANT301.pdf}{Theorem 4.3}]{legacy-milne-ant}
gives each class an integral ideal of norm at most \(B_K=(4/\pi)^{r_2}(d!/d^d)\sqrt{|\Delta_K|}\),
where \(r_2\) counts complex places. Summing \((B_K/N\mathfrak a)^2\) over nonzero integral ideals gives
\(h_K\le B_K^2\zeta_K(2)\le4^d|\Delta_K|\), since \(2r_2\le d\) and the Euler products give
\(\zeta_K(2)\le\zeta(2)^d<2^d\). Thus \(h_K\le(2^d|\Delta_K|)^{d+1}\). To bound the discriminant, write \(X^M-1=\Phi_M(X)H(X)\) and let \(\zeta\) be a primitive \(M\)-th root. Differentiation gives \(\Phi_M'(\zeta)H(\zeta)=M\zeta^{-1}\). The absolute discriminant of the power basis is \(\lvert N_{K/\mathbb Q}\Phi'_M(\zeta)\rvert\), which divides \(M^d\), and the field discriminant divides that of the power basis. Using \(d,|\Sigma|\le M\) and \(M\le 2^M\) gives\formalnote{12.51}{51}
\[
 2^s\le |B|\le (2^M(2^M)^M)^{M+1}\,(2^M)^{3M}
               =2^{M^3+5M^2+M}.
\]


\subsection{Restoring the twist without paying for its depth}
Add to \(\{b_i\}\) one representative per coset of \(G_K\) in \(G_E\). There are at most \(\varphi(M)\) of these. Suppose \(z\) vanishes at the resulting set, and induct on \(j\) such that \(\ell^j z=0\). For \(j=0\) this already says \(z=0\). For \(j\ge1\), the cocycle \(Mz\) vanishes at the same points and is killed by \(\ell^{j-1}\), since \(\ell\mid M\). Induction gives \(Mz=0\).

Now \(\chi_\ell(g)\equiv1\bmod\ell^{v_\ell(M)}\) for \(g\in G_K\), since \(K\) contains the corresponding roots of unity. This makes \(\chi_\ell(g)^{-e}\) trivial on \(A[M]\), which contains all the values of \(z\). The cocycle equation gives \(z(gh)=z(g)+z(h)\) for \(g\in G_K,\ h\in G_E\). In particular the restriction to \(G_K\) is an \(M\)-torsion homomorphism. Inertia away from \(\ell\) restricts into inertia away from \(\ell\) in \(G_E\), so the preceding abelian test kills the restriction. Writing any \(h\in G_E\) as \(g c\) for a chosen coset representative \(c\) gives \(z(h)=0\). The same test works for every \(j\), including \(j=\nu\).\formalnote{12.52}{52}

In the stated small-\(\ell\) range \(M\le y^5\), and \(M\ge4\) gives \(M^3+5M^2+2M\le M^4\). The number of test points is therefore at most \(y^{20}\). If the cyclotomic image has a topological generator represented by \(g_0\), each test point can also be written \(h c\) with \(h\in\ker\chi_\ell\) and \(c\in\overline{\langle g_0\rangle}\): the compact group on the right maps onto the cyclotomic image. For cocycles already zero at \(g_0\), the zero set is a closed subgroup and contains every \(c\). Vanishing at the corresponding \(h\)'s suffices. This last variant assumes the stated generator; the preceding unconditional test works in particular for \(\ell=2\).\formalnote{12.53}{53}

\clearpage
\paragraph{The bound for the determinant.}
Assume that \(E\) is CM and that \(O\) contains the roots of unity of order \(\exp((\mathbb Z/N_E\mathbb Z)^\times)\) at each coefficient prime. For \(0<e<t\), take the odd cocycle module of \S\ref{sec:derivative-cocycles} and the integer \(Y_e=16\varphi(N_E)X_e^2\) that annihilates it there. Evaluation embeds this module into at most \(y^{20}\) copies of \(A[Y_e]\), each of length \(\min\{\nu v_O(\ell),v_O(Y_e)\}\). For \(e=0\) the everywhere unramified homomorphisms factor through \(\mathrm{Cl}(E)\) in the totally complex case. Decomposing the class group into cyclic groups of orders \(d_i\) gives the bound \(\sum_i v_O(d_i)=v_O(h_E)\). Putting \(y=N_E+t+3\), the same integer
\[
              h_E\left(\prod_{0<e<t}Y_e\right)^{y^{20}}
\]
bounds all the lengths by its \(O\)-valuation, using the component tests of \S\ref{sec:derivative-cocycles} for \(\ell>y^2\). To check the size, \(X_e\le2^{y^6}\), \(\varphi(N_E)\le N_E\le y\), and \(4+\log_2 y+2y^6\le4y^6\) for \(y\ge4\). There are at most \(y\) factors, so this integer is at most \(h_E2^{4y^{27}}\). This is the absolute cost used for the odd rows of the determinant.\formalnote{12.54}{54}


\editionchapter{Congruences and integral Hecke operators}\label{ch:deform}

A projection onto selected Hecke eigenspaces can require division by differences of eigenvalues. The estimates of Chapter~\ref{ch:selmer} control the losses from those congruences. This chapter constructs an operator at the original level that clears the resulting denominators.

\section{The Galois congruence estimate at finite coefficient depth}
\label{sec:congruence-detail}

To bound the denominator in \S\ref{sec:regular-twists}, we need the
lengths of two modules of Tate cocycles. The odd cocycles are killed
by an integer constructed from Bernoulli polynomials. The bound
also counts the evaluations needed to determine a cocycle, since an
annihilator alone need not bound the length. For the even cocycles
the evaluations at suitably chosen Frobenius elements satisfy a
triangular system of relations. Both bounds include the prime \(2\).

\subsection{The Galois representation seen by a lift}

For the CM field \(E\) of \S\ref{sec:data}, set \(N_E=q_1d\), so that \(E\subseteq\mathbb Q(\mu_{N_E})\). Write \(G_E={\rm Gal}(\overline{\mathbb Q}/E)\) and let \(\epsilon\) be the \(\ell\)-adic cyclotomic character and \(c\) a fixed complex conjugation. Fix an auxiliary inverse character \(\psi\) from the simultaneous twisting theorem, a lift with source rank \(r'\), and rank increase \(1\le t'\le t\). Both \(r'\) and \(t'\) are odd in this setting. Index the Hecke eigensystems at the chosen level by \(x\). Chapter~\ref{ch:galois} attaches a Galois representation to each such system. After the normalizing twist, we use the resulting finite family of conjugate self-dual representations \(u_x\) with the prescribed Hodge--Tate jumps, tame types at the chosen places over \(p_0,p_1\), and local conditions at \(\ell\). At inertia over the conductor of \(\psi\) and its conjugate,
\[
                     (u_x(g)-\psi(g)I)(u_x(g)-I)=0 .
\]
When \(x\) is an eigensystem arising from the lift, the normalized trace has the form
\[
 {\rm tr}\,u_x(g)=\psi(g){\rm tr}\,R_x(g)+
                           \sum_{e=0}^{t'-1}\epsilon(g)^{-e}.
\]
Here \(R_x\) is the main summand. At the chosen lift its twist has a lattice with a specified scalar at a distinguished inertia element and the divisible-coefficient Selmer bound from the twisting theorem. The proof uses these assertions about the representations attached to the actual family of Hecke systems. A decomposition of their reductions at a single prime would give weaker information.\formalnote{13.1}{1}

Over a sufficiently large coefficient ring \(O\), the trace relation admits a lattice model with the main block and the \(t'\) character blocks on the lifts. The lattice lemma uses nondegeneracy of the trace form on the span of the image and the full matrix-span property for the chosen main block. Its integral conclusion is needed because the sizes of congruences are measured over \(O\). A family of first-order relations among \(u_x(g)-1\) gives cocycles with coefficients in each block. The separate contributions to their length are the main-block Selmer length, the loss in imposing the crystalline relations at \(\ell\), and cocycles of \(O(-e)\) for \(e<t'\). We treat these last modules. In what follows \(O\) is the integer ring of a finite extension of \(\mathbb Q_\ell\) containing enough roots of unity for the Dirichlet characters modulo \(N_E\). Use \(v_O(\varpi)=1\) for a uniformizer and \(v_\ell(\ell)=1\), so \(v_O(a)=v_O(\ell)v_\ell(a)\) on rational \(a\).\formalnote{13.2}{2}

\subsection{The integer for the odd cocycles}

Put \(A_n=O/\ell^nO\). For \(e>0\) let \(M_e(n)\) be the \(O\)-module of locally constant functions \(z:G_E\to A_n\) satisfying
\[
 z(gh)=z(g)+\epsilon(g)^{-e}z(h),\qquad z(cgc)=-z(g).
\]
They vanish on inertia away from \(\ell\) and on \(D_w\cap\ker\epsilon\) for every decomposition group above \(\ell\). These are actual cocycles with specified values. For \(e=0\) the definition changes to all everywhere unramified continuous homomorphisms to \(A_n\). By class field theory those homomorphisms factor through the ideal class group of \(E\). If its cyclic factors have orders \(d_i\), evaluation on generators embeds the homomorphisms into \(\prod_i A_n[d_i]\). Since \(\operatorname{len}_O A_n[d_i]\le v_O(d_i)\), their total length is at most \(v_O(h_E)\). This accounts for the class-number factor below.\formalnote{13.3}{3}

Fix \(1\le e<t\). There is a concrete integer \(X_e\) involved in the construction. For \(0\le s<N_E\) put
\[
 b(s)=(e+2)!\,N_E^{e+1}B_{e+1}(s/N_E),\quad
 g(s)=2b(s)+(-1)^e{\bf1}_{s=1}+{\bf1}_{s=-1},
\]
with the two indicators read modulo \(N_E\). A power-sum identity for Bernoulli polynomials makes \(b(s)\) integral. Form the integral matrix \(T=(g(v/u))_{u,v\in(\mathbb Z/N_E)^\times}\), taking least nonnegative representatives in \(b\). Write \(D(a)=(a^e-1)(a^{e+1}-1)\). One choice is
\[
                X_e=D(1+N_E)D(1+2N_E)\,|\det T|.
\]
To see why the determinant is nonzero, apply \(T\) to the column \((\chi(v))_v\). Its eigenvalue is
\[
                2\sum_{s\bmod N_E}\chi(s)b(s)+(-1)^e+\chi(-1).
\]
For \(\chi(-1)=(-1)^e\), the sum is a multiple of \(L(-e,\chi)=0\) and the last two terms give \(2(-1)^e\). In the other parity the last two terms cancel and the multiple of \(L(-e,\chi)\) is nonzero. This uses the functional equation at the negative integer and nonvanishing at \(1+e>1\), with the finite factors for imprimitive characters included. Complex characters give a basis, so \(\det T\ne0\). For each \(\ell\) at least one of \(1+N_E,1+2N_E\) is prime to \(\ell\), as their greatest common divisor is one. This gives the factor \(D(a)\mid X_e\) for a suitable \(a\equiv1\bmod N_E\) prime to \(\ell\).

To see that each character eigenvalue divides \(\det T\) in the character's coefficient ring, apply \({\rm adj}(T)\) to the eigenvector equation and read the component at \(1\). The Bernoulli power-sum congruence then gives, for \(\chi(-1)=-(-1)^e\) and \(a\equiv1\bmod N_E\), a coefficient \(w\) independent of \(m\) such that for \((m,a)=1\)
\[
  (a^{e+1}-1)X_e\equiv
    2w\sum_{0\le s<N_Em}\left\lfloor\frac{as}{N_Em}\right\rfloor
                              s^e\chi(s)\pmod{m}.
\]
The congruence, together with \(D(a)\mid X_e\) for a suitable \(a\) at each prime, is the finite Stickelberger property used here. For the size, the Bernoulli estimate bounds every entry of \(T\) by \(64(e+2)!N_E^{e+1}(e+1)!\). Bounding the determinant and the four elementary factors gives\formalnote{13.4}{4}
\[
                        1\le X_e\le 2^{(N_E+e+3)^6}.
\]

\paragraph{From the floor sum to annihilation.}
For a Dirichlet character \(\chi\bmod N_E\) put \(\alpha=\epsilon^e\chi\). Its component over \(\mathbb Q\) consists of \(\alpha^{-1}\)-cocycles with values in \(A_n\) that vanish at \(c\), are unramified off \(\ell\) over \(\mathbb Q(\mu_{N_E})\), and are locally trivial at \(\ell\) on the indicated cyclotomic kernel. For \(\alpha(c)=-1\), the component reciprocity input is as follows. At a split prime \(p\equiv1\bmod N_E\ell^n\), a Jacobi sum is a unit away from \(p\) and its conjugates' valuations at \(p\) are explicit differences of floors. Global Artin reciprocity applied to that sum and a cocycle \(z\) gives a relation on \(z({\rm Frob}_p)\). The floor-sum congruence and \(D(a)\mid X_e\) imply \(2X_e^2z({\rm Frob}_p)=0\) at the split Frobenius elements used here. If \(4X_e^2z\ne0\), the facts \(\alpha(c)=-1\) and \(z(c)=0\) force \(2X_e^2z\) to be visible on the cyclotomic kernel at level \(N_E\ell^n\). Chebotarev then gives a split Frobenius with \(2X_e^2z({\rm Frob}_p)\ne0\), a contradiction. This gives the component annihilator \(4X_e^2\). The Jacobi-sum valuation identity and its reciprocity relation are the number-theoretic inputs to this deduction.\formalnote{13.5}{5}

The factor \(2X_e^2\) can be seen directly in the residue ring \(A_n\). Let \(S\) be the Stickelberger sum there, with its floor differences evaluated at the conjugates of the Jacobi sum. The reciprocity identity is \(S\,2(a^e-1)z({\rm Frob}_p)=0\). The congruence gives \(S k=(a^{e+1}-1)X_e\) for some \(k\in A_n\). Write \(X_e=(a^e-1)(a^{e+1}-1)b\). Multiplication of the reciprocity identity by \(kb\) now gives \(2X_e^2 z({\rm Frob}_p)=0\), with no division in \(A_n\). Oddness supplies the additional factor \(2\) in passing from the split Frobenius assertion to all group elements.

To pass back to \(G_E\) choose lifts \(\sigma_u\) of \(u\in(\mathbb Z/N_E)^\times\) to \(G_{\mathbb Q}\). For \(s,g\in G_{\mathbb Q}\), write \(b(s,g)=s(cgc^{-1}g^{-1})s^{-1}\in G_E\), using the abelian cyclotomic field containing \(E\). The character transform of a cocycle \(z\) on \(G_E\) is
\[
          z_\chi(g)=\sum_{u\in(\mathbb Z/N_E)^\times}
                           \alpha(\sigma_u)\,z(b(\sigma_u,g)).
\]
Reindexing the cosets and using the cocycle identity shows that \(z_\chi\) is a component cocycle. Character orthogonality gives the integral implication that if all odd \(z_\chi\) vanish, then \(4\varphi(N_E)z=0\). The integer is retained because the group order and \(2\) may fail to be units in \(O\). Applying the implication to \(4X_e^2z\) proves that\formalnote{13.6}{6}
\[
                       Y_e=16\,\varphi(N_E)X_e^2
                    \quad\hbox{annihilates }M_e(n).
\]

\subsection{From an annihilator to the length needed for a determinant}

An annihilator alone controls the exponent of a finite module. An arbitrary number of copies of \(O/\varpi O\) have the same exponent. For \(M_e(n)\) the missing input is a set of evaluations that detects a cocycle. If at most \(k\) elements \(g_i\in G_E\) suffice, evaluation gives
\[
  M_e(n)\hookrightarrow
       \prod_{i=1}^{k} A_n[Y_e],\qquad
   \operatorname{len}_O M_e(n)\le k v_O(Y_e).
\]
Indeed the kernel is zero by detection and each factor has length \(\min(nv_O(\ell),v_O(Y_e))\). This is the precise way an annihilator becomes a sum in the logarithm of the determinant.

Set \(y=N_E+t+3\). If \(\ell\le y^2\), the ray-class detection theorem supplies at most \(y^{20}\) testing points for cocycles unramified off \(\ell\). Its proof passes to the cyclotomic field \(K=\mathbb Q(\mu_M)\), with \(M=N_E\ell^2\). There \(h_K\) times the idele power index at \(\ell\) is at most \(2^{M^3+5M^2+M}\). This bounds the order of the relevant ray quotient. In a finite group, adjoining an element outside the subgroup generated so far at least doubles its order. A group of order at most \(2^D\) therefore needs at most \(D\) generators. Section~\ref{sec:ray-detection} uses this count to bound the number of evaluation points and treats the twist and the cyclotomic generator.

If \(\ell>y^2\), the Kolyvagin derivative-cocycle theorem supplies at most \(1+v_\ell(X_e)\) testing elements for each odd component over \(\mathbb Q\). Its arithmetic mechanism can be stated precisely. The derivative cocycles \(\kappa_S\) for \(\epsilon^{e+1}\chi\) are indexed by subsets of auxiliary split primes. For \(q\notin S\), evaluation of \(\kappa_S\) on a Frobenius at \(q\) equals an inertia value of \(\kappa_{S\cup\{q\}}\). If \(1+v_\ell(X_e)\) arbitrary tests never sufficed, one could find nonzero \(\ell\)-torsion classes trivial at the required sets. Reciprocal pairings against these classes, using the Frobenius--inertia identity inductively, make the required multiple of \(\kappa_\varnothing\) unramified at the distinguished split prime. Its prescribed Jacobi-sum valuation would imply \(\ell^{1+v_\ell(X_e)}\mid X_e\). The contradiction establishes detection. Section~\ref{sec:derivative-cocycles} constructs the derivative classes and their Frobenius--inertia relation; the resulting triangular tests and their local reciprocity induction are detailed in the sections on odd classes.\formalnote{13.7}{7}

To count tests on \(G_E\), use the points \(b(\sigma_u,g_{\chi,i})\) for every \(u,\chi,i\). Their number is at most \(\varphi(N_E)^2(1+v_\ell(X_e))\le y^{20}\), since \(v_\ell(X_e)\le y^6\). If \(z\) vanishes there, every odd \(z_\chi\) vanishes at its detecting points and is zero. Hence \(4\varphi(N_E)z=0\). Since \(4\varphi(N_E)<\ell\) here, it follows that \(z=0\). Both prime ranges now give \(\operatorname{len}_O M_e(n)\le y^{20}v_O(Y_e)\). One and the same integer serving all \(e<t,\ell,n\) is
\[
             X=h_E\left(\prod_{e=1}^{t-1}
                          16\varphi(N_E)X_e^2\right)^{y^{20}} .
\]
Its size follows directly from \(X_e\le2^{y^6}\), \(\varphi(N_E)\le y\), and at most \(t-1\le y\) factors. The base-two logarithm of the product to the power \(y^{20}\) is at most \(4y^{27}\). We have proved, with the two stated detection theorems and reciprocity inputs,
\[
   \operatorname{len}_O M_e(n)\le v_O(X),
              \qquad 0<X\le 2^{4(N_E+t+3)^{27}}h_E .
\]
Dividing by \(v_O(\ell)\) gives the normalized bound \(v_\ell(X)\) independently of the coefficient extension. It applies to the explicitly defined module \(M_e(n)\). This argument makes no identification with a full Tate--Shafarevich group and uses only the displayed finite-level consequences of the arithmetic inputs.\formalnote{13.8}{8}

\subsection{Even classes and a triangular family of Frobenius elements}

For \(e>0\) an even cocycle satisfies \(z(cgc)=z(g)\). It vanishes on \(D_w\cap\ker\epsilon\) above \(\ell\) and may ramify at auxiliary places and their conjugates. We explain how the proof handles these allowed ramification values without summing the lengths of the full even modules. The hypotheses here include \(\mu_{5040}\subset E\), at least four infinite places in the CM type, and \(E\subset\mathbb Q(\mu_{N_E})\). Fix a coefficient depth \(\nu>0\). Take \(y_0>N_E+t'\) and an element \(\gamma\) generating the cyclotomic image topologically, with \(u=\epsilon(\gamma)\) and \((u^k-1)\mid\ell^h\) in \(\mathbb Z_\ell\) for \(0<k<t'\). Assume \(\ell\le y_0^2\) or \(h=0\). In the application \(N_E=O_{q,d}(x^{11}\log x)\) and \(t'\le t=5r\), so \(y_0=r^8>N_E+t'\) for all sufficiently large \(x\). The integer \(N_{\rm Aut}\) is the common multiple of the class stabilizer orders. It gives the bound \(h\le v_\ell(N_{\rm Aut})\), and one can take \(h=0\) for \(\ell>y_0^2\). Indeed for such primes \(E\) lies in a cyclotomic field unramified at \(\ell\), and an inertia element can act on \(\mu_\ell\) by a generator of \(\mathbb F_\ell^\times\). Its powers with \(0<k<t'\) differ from \(1\) by units.

There can also be prescribed restrictions \(W\) on Frobenius data, used to retain the automorphic level and Hecke conditions. The prime-selection hypothesis is that, near a translate \(\delta^{\ell^\nu}g\) of each prescribed \(g\), with \(\epsilon(\delta)=1\), the permissible data occupy a nonempty open part of every sufficiently small Frobenius neighborhood, apart from finitely many primes. The compatible restrictions for the lift are verified in the projection argument. The conclusion chooses distinct degree-one primes \(w_j\), with pairwise distinct primes below them in \(E^+\), avoiding the excluded places and with \(Nw_j\equiv1\bmod\ell^\nu\). Denote their geometric Frobenius and tame generator by \(\phi_j,\tau_j\).\formalnote{13.9}{9}

\paragraph{Constructing row \(j\).}
The rows run over pairs of a positive exponent \(e<t'\) and a detecting point \(g\in\ker\epsilon\), ordered first by exponent. A Kummer map \(K\) is a function to \(\mathbb Z/\ell^\nu\) whose restriction to \(G_\nu=\ker(\epsilon\bmod\ell^\nu)\) is an odd homomorphism. At exponent \(e\), conjugation by \(a\in G_E\) acts on this restriction by \(\epsilon(a)^{e+1}\). Its global reciprocity relation comes from the maximal real subfield. The local/global relation-space dimension theorem there supplies more independent Kummer reductions modulo \(\ell\) than the number of earlier rows of exponent \(e\). Take a new map \(K_j\) outside their span. The map evaluating all the Kummer maps of this exponent on \(G_\nu\) is onto \((\mathbb Z/\ell^\nu)^s\). Indeed linear independence of the reductions gives surjectivity modulo \(\ell\), and the full image then follows by Nakayama's lemma on the finite quotient. Choose an element on which the earlier maps of this exponent are zero and \(K_j\) is one.

To treat the other exponents apply the separating polynomial in conjugation by \(\gamma\), rescaled by \(u^{-1}\) to match the exponents \(k\) instead of \(k+1\). Conjugation and products in \(G_\nu\) realize the linear combinations on its abelian images. This gives \(v\in G_\nu\) such that all earlier maps kill \(v\) and
\[
         K_j(v)=\prod_{\substack{0\le k<t'\\ k\ne e}}(u^e-u^k)
                   \pmod{\ell^\nu}=:\Delta_e .
\]
Each difference is a unit times \(\pm(u^{|e-k|}-1)\), so the product divides \(\ell^{(t'-1)h}\) in \(\mathbb Z_\ell\). The factor \(4\) needed below divides \(\ell^{3h}\) in the case at hand. At \(2\) the inclusion of fourth roots of unity ensures \(4\mid u^e-1\), and the hypothesis \(u^e-1\mid\ell^h\) bounds its valuation. Thus \(4\Delta_e\) divides \(\ell^{y_0h}\) modulo \(\ell^\nu\).

The actual Frobenius prescription explains both appearances of \(2\). Write \(a^c=cac\), and select Frobenius near
\[
               \delta^{\ell^\nu}\,g g^c\,v(v^c)^{-1}.
\]
The factor \(\delta^{\ell^\nu}\) contributes zero to the additive evaluations modulo \(\ell^\nu\). On \(G_\nu\) the odd Kummer maps are additive and the even Tate cocycles of exponent \(e\) are additive there too. An odd map takes the value \(2K(v)\) on the product, and an even cocycle takes the value \(2z(g)\). Local constancy and the Frobenius selection theorem give a new \(\phi_j\) for which the earlier Kummer maps vanish, \(K_j(\phi_j)=2\Delta_e\), and \(z(\phi_j)=2z(g)\) for all even cocycles of that exponent unramified off \(\ell\). The statement for all such cocycles at once uses the uniform-neighborhood lemma in the prime construction. Distinctness and avoidance are obtained by excluding finitely many primes at each step.\formalnote{13.10}{10}

\paragraph{How the rows eliminate the even values.}
Take \(z\) of exponent \(e\) allowed to ramify at the new places. Suppose it vanishes at \(\tau_i\) for all \(i<j\) and at the \(\phi_i\) of the same exponent for \(i<j\). Triviality on earlier tame inertia, its conjugate, and wild inertia removes those places from its ramification set. The reciprocity law for \(K_j\), with generators normalized over \(E^+\), is
\[
                     \sum_{i\ge j}K_j(\phi_i)\,2z(\tau_i)=0 .
\]
Each later Frobenius was chosen to kill the earlier maps. The only term left is \(2K_j(\phi_j)z(\tau_j)\). Since \(2K_j(\phi_j)=4\Delta_e\) divides \(\ell^{y_0h}\), the conclusion is
\[
                          \ell^{y_0h}z(\tau_j)=0 .
\]
Changing to the chosen local tame generator multiplies the additive inertia value by a scalar and preserves this conclusion. This is triangular elimination with the actual diagonal divisibility. For \(h=0\) it gives vanishing. At the small primes it gives the specific torsion bound above.

Let \(s\) be the set of test points, which together with \(\gamma^{-1}\) detects cocycles unramified away from \(\ell\). The construction bounds \((\#\{j\}+\#s+1)y_0h\le y_0^{22}h\). In the case \(h>0\) this follows from the polynomial count at \(\ell\le y_0^2\) and the at most \(t'-1\) exponents per test point. For \(h=0\) this particular cost vanishes even though testing points are still required. This total bound is independent of \(\nu\), a necessary feature for passing to divisible coefficients.\formalnote{13.11}{11}

\subsection{Returning the two parity calculations to the projection}

For any cocycle \(z\) write \(z^c(g)=z(cgc)\). Imposing the prescribed crystalline relations forces vanishing on the local set defining \(M_e(n)\). The cost of replacing a Ribet cocycle module \(Z\) by one with this local vanishing is bounded by the length of \(Z/(Z\cap Z_f)\), where \(Z_f\) obeys the crystalline relations. For positive exponent \(z+z^c\) is an even cocycle. At the tame and Frobenius test elements, which act trivially on \(A_\nu(-e)\), the cocycle equation gives \(z((g^c)^{-1})=-z(g^c)\). Thus equality of the paired evaluations \(z(g)\) and \(z((g^c)^{-1})\) is exactly vanishing of \(z+z^c\) at \(g\). The triangular argument just proved applies to these equalities. For the residual submodule whose paired tame values are zero and whose same-exponent Frobenius pairs agree, the triangle cost theorem bounds the length modulo the specified \(\ell^{y_0h}\)-torsion by the odd contribution, besides the cost of the local vanishing. Its additional hypothesis is that on members unramified off \(\ell\), \((\epsilon(\gamma^{-1})^{-e}-1)z(\gamma^{-1})=0\). The separation products and the tests give control of the even residual term even at \(2\). The resulting admissible length cost for each character is
\[
          \operatorname{len}_O\bigl(Z/(Z\cap Z_f)\bigr)
                +\operatorname{len}_O M_e(\nu).
\]
At \(e=0\) the residual term goes directly into the unramified module. This use of \(z+z^c\) avoids dividing by \(2\) in the coefficient ring.\formalnote{13.12}{12}

For clarity, the precise algebraic output at the chosen lift \(\pi\) is a single element \(f\) of a Hecke trace algebra with polynomial variables \(X_j\). The trace-relation module contains \(\ell^\nu\) times the traces and the traces of the ties \(\tau_j-1-X_j(\phi_j-1)H_j(\phi_j)\), with \(H_j(1)=1\), including the conjugate relations. The element \(f\) sends every determinant \(\det(1-u_x(g))\) into that module on multiplication. These determinants vanish on true lifts because the summand \(e=0\) is trivial. On the lift locus every coefficient of \(f\) at a proper divisor of \(\prod_jX_j\) vanishes. Its coefficient \(d_\pi\) at \(\prod_jX_j\), evaluated at \(\pi\), satisfies
\[
 v_\ell(d_\pi)+2v_\ell(N_{\rm Aut})
       \le {\rm sel}+2t'{\rm ram}
                       +r'{\rm loc}_\ell+e_\ell^0.
\]
The underlying Ribet matrix theorem glues the tied rows of the main block and the character blocks by a Fitting-ideal determinant. Section~\ref{sec:trace-ideal} shows how this turns the lengths above into the coefficient of a single polynomial serving all determinants. The crystalline comparison bounds the sum of the local colengths and contributes \(r'{\rm loc}_\ell\). The term \(2t'{\rm ram}\) bounds the losses at the conjugate auxiliary-conductor places. With \(t\) the ambient upper bound and \(y_0=r^8\), the remaining term in the readable edition is explicitly
\[
 \begin{split}
 e_\ell^0={}&t\,v_\ell(X)\\
 &+\bigl(8ft+2+(4t+1)y_0^{22}+y_0\bigr)v_\ell(N_{\rm Aut}).
 \end{split}
\]
The term \(t\,v_\ell(X)\) sums the length bound \(v_\ell(X)\) for at most \(t\) odd cocycle modules. Here \(X\) is the common integer from the odd-cocycle construction above. The terms with \(y_0^{22}\) reflect the entire triangular cost, using \(h\le v_\ell(N_{\rm Aut})\). No sum of orders of the even cocycle modules appears in this expression.\formalnote{13.13}{13}

For a character of exponent \(e\) the weighted determinant sum agrees with a scalar \(\Gamma_e\) times the monomial assigned to \(e\), modulo a colon ideal. After multiplication by the product of the earlier variables its error is zero modulo \((X_j^2)_j\). Multiplication in increasing order of \(e\) now gives vanishing at every proper divisor of the distinguished monomial and computes its top coefficient. Section~\ref{sec:trace-ideal} proves this congruence and the valuation bound.\formalnote{13.14}{14}

To recover an integral operator, deform the auxiliary-place norm operators to first order. The coefficient of \(\prod_jX_j\) in their product, multiplied by a tied relation, vanishes modulo the chosen depth by the norm and derivative congruences. On a true lift the lower coefficients of \(f\) vanish, so the squarefree coefficient controls the surviving scalar. Take the depth large enough that the comparison modulus is divisible by the scalar to be cancelled. The finite-level congruences of \S\ref{sec:deformed-norm} then give the operator on the projected test forms of the desired Hecke type. Multiplication by \(N_{\rm Aut}^2\) accounts for the type and the change of level. The two finite group-order factors in the coefficient inequality already allow for this loss. A further \(N_{\rm Aut}\) clears the stabilizer denominators in pairings. Taking one integral linear combination that works on all the finitely many eigenlines and multiplying its eigenvalues gives the determinant bound of \S\ref{sec:regular-twists}.\formalnote{13.15}{15}


\section{From cocycle lengths to an integral projection}\label{sec:trace-ideal}

The projection in \S\ref{sec:regular-twists} must separate the lifted
systems of Hecke eigenvalues from the other systems. The potential
denominator is visible already for two eigenvalues.

\begin{example}[Two congruent eigenvalues]\label{ex:congruent-eigenvalues}
If \(T\) is diagonalizable and has exactly two eigenvalues \(u\) and \(v\), the operator \((T-vI)/(u-v)\) is the identity on the \(u\)-eigenspace and zero on the \(v\)-eigenspace. The full denominator can occur on a lattice. Over \(\mathbb Z_\ell\), take \(s\ge1\) and
\[
 T=\begin{pmatrix}\ell^s&1\\0&0\end{pmatrix}.
\]
The eigenvalues are \(\ell^s\) and \(0\). The projection onto the \(\ell^s\)-eigenspace is \(T/\ell^s\), which takes the second basis vector to \(\ell^{-s}\) times the first. Thus clearing the projection costs exactly \(s\) powers of \(\ell\).
\end{example}

In Ribet's method a congruence between characteristic polynomials gives a Galois cohomology class. In dimension two, Dasgupta, Kakde, Silliman and Wang construct a class whose coefficient module has zeroth Fitting ideal contained in the congruence ideal, even when the two characters agree in the residue field \cite[Theorem~1.1]{dksw-ribet}. For the higher-rank family of eigenvalues here, the Galois cocycles measure the congruences. Bounds for their lengths control the valuation of
a polynomial coefficient. After a change of level this coefficient
clears the denominator of the projection.

Fix a prime \(\ell\), a finite family \(Y\) of systems of Hecke eigenvalues, and a distinguished lift \(\pi\) in the subset \(\mathcal X\subset Y\) of lifts. We seek a polynomial \(f\) in variables indexed by the auxiliary primes, with coefficients evaluated on the family of Hecke systems. The coefficient of \(X_J=\prod_{j\in J}X_j\), the product of the variables at the auxiliary primes indexed by \(J\), will give a multiplier on the lifts. To bound its valuation we first derive a relation among traces valid at every Galois element.

\subsection{Relations in the trace algebra}

Work over a discrete valuation ring \(O\) with fraction field \(K\). Write \(v\) for its valuation, normalized at a uniformizer. The systems give representations \(u_x:G\to {\rm GL}_n(O)\). At \(x\in\mathcal X\) choose bases in which
\[
                    u_x=\rho_x\oplus\bigoplus_{e=0}^{t-1}\xi_e,
                    \qquad t>0,\quad \xi_0=1.
\]
For a finite ordered set \(J\) put \(R=O[X_j:j\in J]\), and take the ring of families \(\mathcal R=\prod_{x\in Y}L[X_j:j\in J]\) for a coefficient field \(L\supset K\). A subalgebra \(S\subset\mathcal R\) contains \(R\) diagonally and all coefficient families of \(\det(T-u_x(g))\). Let \(U:S[G]\to M_n(\mathcal R)\) denote simultaneous evaluation. At depth \(p=\ell^\nu\) consider an \(S\)-submodule \(\mathcal J\subset\mathcal R\) containing \(p\,{\rm tr}\,U(g)\) and \({\rm tr}\,U(t_{j,\pm}g)\) for every \(g\), where
\[
 t_{j,\pm}=\tau_{j,\pm}-1-X_j H_j(\phi_{j,\pm})(\phi_{j,\pm}-1).
\]
The two signs denote the chosen place and its conjugate. Define
\[
 \mathfrak a=\{w\in S[G]:{\rm tr}\,U(wv')\in\mathcal J
                           \text{ for every }v'\in S[G]\}.
\]
Linearity, associativity and \({\rm tr}(ABC)={\rm tr}(BCA)\) make this a two-sided ideal. In particular it contains \(\ker U,\ pS[G]\) and every \(t_{j,\pm}\).

Choose generators \(g_i\) so that the families \(U(g_i-1)\) span all \(U(g-1)\) over \(O\). These families lie in a finite free \(O\)-module, hence finitely many suffice. Include the old inertia generators \(\delta_v\) and the detector set \(s\), allowing repeated coordinates. Put \(y_i=g_i-1\) and \(I=\sum_i S[G]y_i\). For each lift suppose \(\rho_x(\gamma_x)=u'_x\,1\), with \(u'_x-1\) a unit, and \(\xi_e(\gamma_x)=1\) for all \(e\). With \(a\in G\) as below put
\[
 n=(u'_\pi-1)^{-1}(\gamma_\pi-1),\qquad n_0=n+(1-n)(a-1).
\]
Thus \(\rho_\pi(n_0)=1\) and \(\xi_0(n_0)=0\). A row \(b=(b_i)\) over \(R[G]\) is tied if \(\sum_i b_i y_i\) belongs to the left ideal generated by the common kernel, \(p\) times the augmentation ideal, the ties \(t_{j,\pm}\), and the products \(w n_0\) for \(w\) in the augmentation ideal. Every tied square matrix \(B\) therefore satisfies, for some \(B_1\),
\[
                    By=B_1 y\, n_0\pmod{\mathfrak a}.
\]
Indeed all but the last generators belong to \(\mathfrak a\). Each augmentation element is a linear combination of \(g-1\)'s, and each \(g-1\) belongs to \(I+\ker U\) by the choice of coordinates.\formalnote{13.16}{16}

\paragraph{The determinant calculation.}
Put \(D(T)=\det U_{\rm flat}(B-TB_1)\), where each matrix entry is replaced by its \(n\)-by-\(n\) image, giving one ordinary square matrix. Its characteristic coefficients, in particular \(D(T)\), lie in \(S[T]\). Here is an integral-coefficient argument for this assertion. For finitely many matrices \(A_i\) and a formal variable \(z\), the identity
\[
 \det(1-z\textstyle\sum_i A_i)
       =\prod_{[w]} \det(1-z^{|w|}A_w)
\]
runs over cyclic equivalence classes of nonempty primitive words in the indices, where \(A_w\) is the ordered matrix product. Only finitely many factors affect a given degree. For generic matrices over \(\mathbb Q\), take formal logarithms. In degree \(m\) the left side gives \(-{\rm tr}((\sum_i A_i)^m)/m\). An \(m\)-letter word repeating a primitive word of length \(h\) has \(h\) cyclic rotations, giving its trace the coefficient \(-h/m\), exactly the logarithm of its indicated factor. Each coefficient before taking logarithms is an integer polynomial in the entries, so the identity specializes to any commutative coefficient ring. Expand each group-algebra block in group elements times scalars in \(S[T]\), with each term placed in its own matrix-unit block. A word of these terms either has determinant factor \(1\) (zero product or an off-diagonal matrix unit), or has one diagonal block, a scalar times \(U(g)\) for some \(g\in G\). Its factor has coefficients in \(S[T]\) by hypothesis. This proves the assertion and covers the single-block case used below.

The adjugate is a polynomial in the flattened matrix with these characteristic coefficients, by Cayley--Hamilton. Consequently it lifts to a block matrix over \(S[G][T]\), and its product with \(B-TB_1\) is \(D(T)\,1\) modulo \(\ker U\), coefficient by coefficient. In \((S[G]/\mathfrak a)^{\#\{i\}}\), let the coefficients act on the left and \(T\) act by \emph{right} multiplication by \(n_0\). Those actions commute, so the tied-row equality and the adjugate identity give
\[
                               I D(n_0)\subset\mathfrak a.
\]
For \(w\in I\), the same adjugate argument in one block gives some \(v'\) with \(U(v'w)=\det U(w)\,1\). For any \(n_1\in S[G]\), the polynomial family
\[
                      f_B={\rm tr}\,U\bigl(n_1 D(n_0)\bigr)
\]
belongs to \(S\) by the characteristic-coefficient assertion. Its product with \(\det U(w)\) is the trace of \(n_1v'wD(n_0)\), which belongs to \(\mathcal J\). This works for every \(w\in I\) using the same \(f_B\).\formalnote{13.17}{17}

\subsection{Producing tied rows with small determinants}

We spell out the two length inputs. First, the matrices \(\rho_\pi(g)\) span \(M_{\dim\rho_\pi}(O)\) over \(O\). Let \(b_0\) bound the length of the \(\rho_\pi\)-cocycles with values in \((K/O)^{\dim\rho_\pi}\) that vanish at all \(\tau_{j,\pm}\) and at \(\gamma_\pi\), and respect every \(O\)-linear relation among the families \(U(g-1)\). Evaluate tied rows on \(\rho_\pi\) at \(X=0\) and take the \(O\)-span of their ordinary rows. The annihilator of this row lattice in the corresponding power of \(K/O\) injects into the indicated cocycle module. In fact a tuple of values on the \(g_i\) extends by the chosen linear coordinates to a function of \(g\). The identity \(gh-1=(g-1)+g(h-1)\) gives the cocycle identity. Ties at zero give vanishing at the tame generators. The normalizer gives vanishing at \(\gamma_\pi\): the row representing \((\gamma_\pi-1)n_0\) can be taken as \((\gamma_\pi-1)\) times the coordinate row for \(n_0\), and on the main block this gives the coordinate row for \(\gamma_\pi-1\). Indeed, writing \(n_0=n+(a-1)-n(a-1)\), the last two terms contribute zero under this evaluation since \(\rho_\pi(n)=1\); the first contributes \((u'_\pi-1)^{-1}\) times the coordinate row of \(\gamma_\pi-1\). Evaluation on the \(g_i\) recovers the tuple. Smith normal form now bounds the colength by \(b_0\) and produces an ordinary square matrix of rows with determinant of valuation at most \(b_0\). These ordinary rows can be packed into block rows of tied elements: the evaluation algebra contains every matrix unit by the spanning hypothesis, so left linear combinations select a single row and place it in any row position. We obtain a tied matrix \(B_0\). Write \(D_x=\det(\rho_{x,{\rm flat}}(B_0))\). Then \(v(D_\pi(0))\le b_0\).\formalnote{13.18}{18}

For the characters, assign each \(j\) a nondecreasing type \(e(j)\in\{0,\ldots,t-1\}\). Set \(J_e=\{j:e(j)=e\}\) and \(X_A=\prod_{j\in A}X_j\). Assume \(H_j(\xi_e(\phi_{j,\pm}))\equiv1\bmod p\). Use the character cocycles modulo \(p\) which respect the linear relations on \(U(g-1)\), vanish on the \(\delta_v\), and satisfy \(c(\gamma_\pi)=0\Rightarrow(\xi_e(g)-1)c(a)=0\) for every \(g\). Assume \(s\) detects these cocycles. For each submodule \(Z\) of them take the triangle data of the earlier parity calculation: a subgroup \(Z_0\) defined by vanishing on some test elements, such that for \(j\in J_e\) the earlier equalities of paired tame values and the earlier same-type equalities of paired Frobenius values force \(C(c(\tau_{j,+})-c(\tau_{j,-}))=0\) on \(Z_0\). Suppose
\[
  {\rm len}(Z/Z_0)+{\rm len}(Z_1/Z_1[C])\le b_e
\]
whenever \(Z_1\subset Z_0\) vanishes at every tame generator and satisfies all the same-type paired Frobenius equalities. Here \(C\ne0\) is also divisible by \(\prod_{e'\ne e}(\xi_e(a)-\xi_{e'}(a))\) for every \(e\).

Group the coordinates into old inertia \(W\), extra generators \(s_0\), and detectors \(s\). On a fixed \(\xi_e\), tied rows supply unit rows on the \(s_0\)-coordinates. To see this, project to those coordinates and test their image on constant coordinate vectors by linear forms modulo \(p\). Such forms give character cocycles by the coordinate relations and kill the detector coordinates by construction. Detection forces every form to vanish. Duality for finite modules over \(O/(p)\) then places each constant unit row in the image. One can use the same rows on every character. If \(u\) already gives unit row \(i\) on earlier characters, and \(v_j\) gives unit row \(j\) on the next character, replace \(u\) by
\[
                         u+\sum_{j\in s_0}(\delta_{ij}-u_j)v_j.
\]
The correction vanishes on previous characters and gives the requested row on the next. Use these common rows to eliminate the \(s_0\)-coordinates. Let \(K_e\subset R^s\) be the resulting character image of tied rows after forgetting \(W\).

Let \(\mathfrak n=(X_j^2:j\in J)\) and \(\mathfrak q_e=(\mathfrak n:X_{\{j:e(j)<e\}})\). Compute the Fitting ideal first with row space \(K_e+\mathfrak q_e R^s\). It contains \(\Gamma_e X_{J_e}\) with
\[
              v(\Gamma_e(0))\le b_e+(|J_e|+|s|)v(C).
\]
Here are the determinant operations behind the estimate. For a row space \(T\), write \({\rm Fit}(T)\) for the ideal generated by its full-size minors. Adjoining a row \(w\) with \(a'w\in T\) costs at most \(a'\): by multilinearity \(a'{\rm Fit}(T+Rw)\subset {\rm Fit}(T)\). Constant coordinate rows are tested by their \(O/(p)\)-valued duals, since their quotient is killed by \(p\). To impose a subgroup on those duals, add the constant-coordinate images annihilating that subgroup. The added constant module has length equal to the codimension of the subgroup in the dual by elementary divisors. Lifting generators of a composition series and applying the one-row rule costs a scalar of valuation at most that length. Apply this first to \(Z_0\), costing \({\rm len}(Z/Z_0)\).

Next adjoin in order the coordinate differences \(w_j\) at the two Frobenii for \(j\in J_e\). Modulo the current rows the tame difference equals \(X_jw_j\). Earlier Frobenius differences of the same type are already rows. For earlier \(i\) of smaller type, \(X_i\) belongs to \(\mathfrak q_e\). Thus the dual forms satisfy the earlier paired tame equalities as well. The triangle implication gives \(C X_jw_j\) in the current rows by duality applied to the constant tame difference. Each step therefore costs \(CX_j\). Finally specialize the enlarged row space at zero. The surviving duals \(Z_1\) kill both tame generators and agree at the required Frobenii. Restricting them to \(Z_1[C]\) costs at most \({\rm len}(Z_1/Z_1[C])\). At that stage the relations contain \(C\) times each of the \(|s|\) coordinate vectors, so their determinant is \(C^{|s|}\). Lifting its determinant expression from \(X=0\) leaves its constant term unchanged. Multiplying the three costs proves the displayed estimate.\formalnote{13.19}{19}

Expanding the determinants by rows shows that
\({\rm Fit}(K_e+\mathfrak q_e R^s)\subset{\rm Fit}(K_e)+\mathfrak q_e\). Consequently there are polynomials \(d_e\in{\rm Fit}(K_e)\) and \(\zeta_e\in\mathfrak q_e\) with
\[
                          d_e+\zeta_e=\Gamma_e X_{J_e}.
\]
Write \(d_e=\sum_i h_{ei}\det B_{ei}\) using matrices of rows in \(K_e\), and lift their rows to tied rows with zero \(s_0\)-coordinates on \(\xi_e\). The group-algebra element \(q_e\) obtained from \(\prod_{e'\ne e}(a-\xi_{e'}(a))\) by the divisibility assumption evaluates to \(C\) on \(\xi_e\) and zero on all other characters. For each tuple \(i=(i_e)_e\), sum the lifted rows over \(e\) after multiplying by \(q_e\), and restore the common unit rows. For \(v\in W\) assume \(\xi_e(\delta_v)=1\) and \((u_x(\delta_v)-1)o_v(u_x(\delta_v))=0\) for all \(x\). The row with sole entry \(o_v(\delta_v)\) at \(v\) is tied by this identity. Restoring these rows gives a tied matrix \(M_i\) whose diagonal blocks on character \(e\) have determinants \(\prod_{v\in W}o_v(1)\), \(1\), and \(C^{|s|}\det B_{e i_e}\). In particular put \(c_0=C^{|s|}\prod_{v\in W}o_v(1)\).\formalnote{13.20}{20}

\subsection{The squarefree coefficient and the projection}

Take \(k=\prod_{x\in\mathcal X}(u'_x-\gamma_x)/(u'_x-1)\) in any fixed order. Each factor evaluates to \(1\) on all the characters, and at \(\rho_x\) the factor indexed by \(x\) vanishes. Use the tied matrix \(B_i=(1-k)B_0+kM_i\) and set \(n_1=kq_0\) in the trace construction. On lifts \(U_x(n_1)\) is \(C\) on the single trivial character and zero elsewhere, while \(U_x(n_1n_0)=0\). Thus \(f_{B_i}(x)=C D(0)(x)\). Splitting this last determinant by the constituent representations gives \(C D_x\prod_e\det \xi_e(M_i)\). Sum over \(i\) with weight \(\prod_e h_{e i_e}\). By distributivity, one \(f\in S\) still multiplies every \(\det U(w)\), \(w\in I\), into \(\mathcal J\), and on lifts
\[
                           f(x)=C c_0^t D_x\prod_e d_e.
\]
In \(R/\mathfrak n\) successively replace \(d_e\) by \(\Gamma_e X_{J_e}\) in increasing order. At the \(e\)-th step all earlier variables already divide the product, and their product kills \(\zeta_e\) in this quotient. Therefore
\[
             f(x)\equiv C c_0^t D_x\Bigl(\prod_e\Gamma_e\Bigr)X_J
                             \pmod{\mathfrak n}.
\]
Every proper squarefree coefficient on every lift vanishes. The top coefficient at \(\pi\) is
\[
                         d_\pi=C c_0^t D_\pi(0)\prod_e\Gamma_e(0).
\]
Counting the powers gives
\[
 v(d_\pi)\le b_0+\sum_e b_e+
     t\left(\sum_{v\in W}v(o_v(1))+2|s|v(C)\right)+(|J|+1)v(C).
\]
This is the bound over \(O\). Divide by \(v(\ell)\) to obtain the normalization \(v_\ell(\ell)=1\). Taking \(w\) representing \(g-1\) supplies \(\beta_g(x)=\det(1-u_x(g))\) up to the fixed sign \((-1)^n\). All \(\beta_g\) vanish on the lifts because \(\xi_0\) is trivial.\formalnote{13.21}{21}

For the last passage take \(S\) to be the image of the polynomial Hecke algebra and \(\mathcal J\) the image of the ideal generated by the stated trace relations. Coefficientwise evaluation on all eigensystems is injective on the Hecke algebra of functions, so \(f\beta_g\) lies in the tied ideal there. For each non-weak system of base level choose \(g_k\) such that \(\beta_{g_k}(x)=a_k\ne0\) at that system. There are finitely many. Then
\[
        \delta_0=\prod_k(a_k-\beta_{g_k}),\qquad n_*=\prod_k a_k
\]
give \(\delta_0=0\) on these systems and \(\delta_0=n_*\ne0\) on the lifts. Expansion proves \((\delta_0-n_*)f\) belongs to the tied ideal. Let \(A\) be the Hecke indicator of the congruence class of \(\pi\), so \(A(\pi)=1\). Put \(c=N_{\rm Aut}\) and choose depth with \(2n!(n_*c^2)^{|J|+1}\mid\ell^\nu\) in the valuation ring.

Section~\ref{sec:deformed-norm} proves the root-chain formula under the specified local hypotheses, the norm and derivative congruences, and the integral change-of-level norm lift. For precision, let \(l\) have base level, with \(d\,l\) integral, and let \(y\) be its prescribed lift component, so that \(\delta_0 y=n_*y\) and \(A\delta_0 l=n_* A y\). The norm lift supplies an integral \(u\) with \({\rm Nm}\,u=dc\,l\). In the product \(\mathcal N=\prod_j({\rm Nm}_j+X_j\partial_j)\), take \(G_*\) to be the operator \([X_J](\mathcal N f)\). The first-order congruences kill the tied ideal at this coefficient and factor each derivative term through the base norm after multiplying by the prescribed powers of \(c^2 A\delta_0\). Thus, with \(h=|J|+1\), modulo \(\ell^\nu/(2n!)\) the operator \((n_*c^2)^h A^h G_*\) applied to \(u\) equals
\[
                         (n_*c^2)^h\,dc\, A^h [X_J]f\, y .
\]
Indeed a term indexed by \(T\subset J\) contains \([X_T]f\) acting on \(y\); all terms with \(T\ne J\) vanish by the coefficient calculation. In the surviving term the norm sends \(u\) to \(dc\,l\), and \(A\delta_0 l=n_* A y\) supplies the indicated scalar. The modulus is divisible by \((n_*c^2)^h\), so cancellation inside the lattice proves \(dc\,A^h[X_J]f\,y\) integral. This is a finite-depth calculation. The set of auxiliary places, and hence the level used for this calculation, may depend on the exponent \(\nu\) of \(\ell^\nu\). At \(\pi\) the diagonal multiplier \(A^h[X_J]f\) takes the value \(d_\pi\).\formalnote{13.22}{22}

For a fixed lift generator of the chosen type, type projection gives \(d=c\). Orthogonality of the original form to false lifts gives the agreement of the two projected forms required above: at a suitable split place the Hecke polynomial evaluated at the lift root kills the competing weak eigenspaces. This produces an operator integral on the projected test forms, diagonal on the lift generators and with the bound for \(v_\ell(c^2 d_\pi)\) on the selected generator. To use one operator on \(Z\), choose a basis of lift generators and an operator for each. On each basis line scale its list of eigenvalues by one of largest \(\ell\)-adic absolute value. All entries are then integral and at least one has nonzero residue. Over the infinite residue field of \(\overline{\mathbb Q}_\ell\), finitely many proper linear hyperplanes cannot cover coefficient space. Choose integral coefficients outside their union after reduction. On every basis line the sum has the maximal absolute value in its list, so each valuation satisfies the required bound. The sum clears the projected forms. Multiplication by the last \(N_{\rm Aut}\) clears their pairings by the stabilizer formula. The determinant is the product over the basis lines, giving the factor \(\dim Z\) in its logarithmic bound.\formalnote{13.23}{23}


\section{Deforming the norm on the auxiliary level lattice}
\label{sec:deformed-norm}
At an auxiliary prime, applying the group norm, the sum of the
translations by all residue-group elements, returns a form to the
original level. The trace relations in \S\ref{sec:trace-ideal}
also involve tame generators. The following calculation shows why a
weighted sum of translations is useful for imposing these relations.

\begin{example}[A weighted group norm]\label{ex:weighted-group-norm}
Let \(g\) generate a cyclic group of order three. In its integral group
ring put \(D=1+g+g^2\). The identity \((g-1)D=0\) expresses the
invariance of the norm. The exponents of \(1,g,g^2\) modulo three
have least absolute representatives \(0,1,-1\). Weight the same
summands by these integers to obtain \(B=g-g^2\). Then
\[
                  (g-1)B=2g^2-g-1=3g^2-D\equiv-D\pmod3.
\]
Indeed the coefficient at \(g^j\) in \(gB\) is the weight at
\(g^{j-1}\). It differs from the weight at \(g^j\) by \(-1\)
modulo three for every \(j\).
\end{example}

For the matrix group at an auxiliary prime, weight each translation by a discrete logarithm of its determinant. Translating this weighted sum by an element whose determinant has logarithm \(1\) subtracts one from every coefficient modulo the chosen depth. The difference is therefore the negative of the ordinary norm. We prove the corresponding parabolic identity in \textup{(DN2)}.

Write \(n\) for the rank, \(m=\ell^\nu\) for the depth, and \(w_j\), \(j\in J\), for the auxiliary split places. These places and their conjugates are pairwise disjoint, avoid \(\ell\) and the base level, and satisfy \(m\mid q_j-1\), where \(q_j\) is the residue cardinality. Write \(\Gamma_j=\mathrm{GL}_n(\mathbb F_{q_j})\), \(\Gamma=\prod_j\Gamma_j\), and \(\rho\) for their action on forms of the specified type with principal level at the \(w_j\). The residue-group actions preserve the integral lattice \(L\) defined by the coefficient gauge. The base-level space is the subspace fixed by \(\Gamma\), since the full integral group modulo principal level is \(\Gamma_j\) at \(w_j\). Let \(\mathcal E\) be the ring of endomorphisms of the complex form space that preserve \(L\), with scalars in \(O=\overline{\mathbb Z}_\ell\) acting through the chosen comparison with complex numbers. The diagonal Hecke algebra \(\mathbb T\) commutes with the local translations.

Write \(D_j=\sum_{s\in\Gamma_j}\rho(s)\) for the group norm at \(w_j\) and \(D=\prod_jD_j\) for the group norm over all the auxiliary primes. Let \(c=N_{\rm Aut}\) be the common multiple of the stabilizer orders of the original level. Throughout the calculation \(A\equiv B\) means equality after a ring map killing every \(Z\in\mathcal E\) such that \(2n!\,Z\in m\mathcal E\). At the depth used for the denominator, \(\mu=m/(2n!)\) belongs to \(O\), and we can use the actual quotient \(\mathcal E/\mu\mathcal E\). Indeed multiplication by \(2n!\) is injective on these endomorphisms in characteristic zero. This choice makes every final congruence a divisibility statement on the lattice.\formalnote{13.24}{24}

\subsection{The lift through the norm}
For each subgroup \(\Delta\subseteq\Gamma\) we need an integral operator \(\Psi_\Delta\) satisfying
\[
                \sum_{g\in\Delta}\rho(g)^{-1}\Psi_\Delta\rho(g)=c\,{\rm Id}.       \tag{DN1}
\]
Take the finite adelic classes modulo the kernel of reduction of the standard level. This preserves the invariance under the original type. The preimage of \(\Delta\) in the standard level moves these classes through its quotient \(\Delta\). Choose one point per orbit. If its stabilizer in \(\Delta\) has order \(s\), give that point weight \(c/s\), and give the other points of the orbit weight zero. The stabilizer is the image of a subgroup of the rational automorphism group at coarse level, since a lift fixing a class modulo the reduction kernel can be multiplied by a kernel element to fix a representative modulo rational points. Conversely a coarse automorphism in the preimage acts trivially on that class. Thus \(s\mid c\), including for the orbits with nontrivial stabilizers. Each point encounters the chosen point exactly \(s\) times under translation by \(\Delta\).

Multiply a form pointwise by this integer weight, using the finite inverse argument if classes are written as right quotients. The weight is invariant under rational points and the kernel of auxiliary reduction. This preserves the principal level and commutes with the type action supported at the old places, whose elements are trivial under auxiliary reduction. It also preserves integrality at each point. Summing the conjugated multipliers now gives (DN1). For a base-level vector \(l\) with \(d\,l\in L\), set \(u=\Psi_\Gamma(d\,l)\). Since every \(\rho(g)\) fixes \(l\), (DN1) gives
\[
                             u\in L,\qquad D u=dc\,l .
\]
This construction pays for the actual stabilizers. In particular it does not require division by \(\#\Gamma\) in \(O\).\formalnote{13.25}{25}

\subsection{A logarithm on a line and its transfer}
Temporarily fix \(w=w_j\) with residue field \(k\), of order \(q\), and write \(G=\mathrm{GL}_n(k)\), \(P=\operatorname{Stab}_G(k e_1)\). The fixed tame generator \(\gamma\) has the local Artin class of a unit whose residue \(z\) admits a homomorphism \(\lambda:k^\times\to\mathbb Z/m\) with \(\lambda(z)=1\). Indeed \(k^\times\) is cyclic, the tame-generator condition makes the residue primitive on its \(m\)-quotient, and \(m\mid q-1\). A uniformizer can be chosen whose Artin class is that of the specified \emph{geometric} Frobenius \(\phi\) (its inverse is arithmetic Frobenius). Set
\[
 a_P(p)=\lambda(p_{11}),\qquad a_G(g)=\lambda(\det g),\qquad d_z=\operatorname{diag}(z,1,\ldots,1).
\]
For integer functions \(f\) on \(\mathbb Z/m\) write \(W_P(f)=\sum_{p\in P}f(a_P(p))\rho(p)\), and define \(W_G(f)\) using \(a_G\). For each residue class \(s\), choose an integral representative \(b(s)\). Choose representatives \(s_L\) for \(G/P\), indexed by the lines \(L\), and write \(C=\sum_L\rho(s_L)\), \(C'=\sum_L\rho(s_L^{-1})\).

The transfer from \(P\) of \(p\mapsto p_{11}\) is \(\det\). To see this, decompose the action of \(g\) on lines into cycles. A cycle through \(L\) of length \(e\) contributes the scalar by which \(g^e\) acts on \(L\), by cancellation of the choices of bases along the cycle. For \(g=\operatorname{diag}(v,1,\ldots,1)\) only the first coordinate line can contribute anything other than \(1\): on any other line a nonzero component outside the first coordinate is fixed under every power of \(g\). Thus the transfer takes value \(v\). Conjugation gives the result at every one-coordinate diagonal. On each elementary transvection subgroup the transfer is trivial, since the subgroup has order \(q\) while its target has order \(q-1\). Row reduction generates \(G\) by these transvections and the diagonal subgroup. Composing with \(\lambda\) proves the assertion for the logarithms.

At a given \(g\), the coefficient in \(CW_P(b)C'\), before applying \(\rho\), is
\(\sum_L b(a_P(s_{gL}^{-1}g s_L))\). Modulo \(m\) this is the transferred logarithm \(a_G(g)\). Moreover the coefficient of \(g\) in \((\rho(d_z)-1)W_G(b)\) comes from \(b(a_G(g)-1)-b(a_G(g))\equiv-1\pmod m\). We have proved\formalnote{13.26}{26}
\[
 CW_P(b)C'\equiv W_G(b),\qquad
 (\rho(d_z)-1)W_G(b)\equiv-D_j .                           \tag{DN2}
\]

\subsection{The tame and Frobenius traces}
On a Hecke system \(x\) write \(N_x\) for its attached Galois representation twisted by the character \(a\), with \(a(\gamma)=1\). For \(g\in G_E\) and \(p\in O[X]\) define diagonal operators with values
\[
 S_p(g)(x)=\operatorname{tr}\!\left(p(N_x(\phi))N_x(g)\right),\qquad
 T_p(g)(x)=\operatorname{tr}\!\left(p(N_x(\phi))(N_x(\gamma)-1)N_x(g)\right).
\]
They are integral Hecke operators once the characteristic-polynomial coefficients at Galois elements belong to \(\mathbb T\). Indeed the trace is minus the coefficient of degree \(n-1\); expanding \(p\) expresses the displayed functions as integral sums of traces at \(\phi^i g\), or differences of those traces at \(\phi^i\gamma g\) and \(\phi^i g\).

We first take the determinant character \(\psi\circ a_G\) of the full residue group. Put \(\zeta=\psi(1)\), its value on the chosen residue generator. For this parahoric there is one block, so every diagonal slot of the generator has character value \(\zeta\). At the split place \(w\nmid\ell\), parahoric compatibility says that any polynomial vanishing on the set of these distinct unit-character values annihilates inertia. To obtain this assertion from \S\ref{sec:tame-labels}, use the annihilating polynomial with multiplicities there and finite inertia: in characteristic zero the inertia element of finite order is semisimple, hence its minimal polynomial is squarefree. In particular \(X-\zeta\) annihilates \(N_x(\gamma)\), because the twist is trivial on \(\gamma\). The action of \(\rho(d_z)\) is the same scalar. The image of \(W_G(f)\) decomposes into such subspaces: averaging over the kernel of \(a_G\) makes the kernel act trivially, and the remaining finite cyclic action diagonalizes over \(\mathbb C\). Thus
\[
             W_G(f)T_p(g)=(\rho(d_z)-1)W_G(f)S_p(g),\qquad D_jT_1(g)=0.       \tag{DN3}
\]
The equality for the norm is the constant-weight case, since \(\rho(d_z)D_j=D_j\).

For the line stabilizer put \(U\) equal to the unnormalized Hecke operator for \(\operatorname{diag}(\varpi,1,\ldots,1)\) at pro-unipotent Iwahori level and put \(\widetilde U=a(\phi)U\). Then
\[
                         B(f)=p(\widetilde U)W_P(f)\quad\text{belongs to }\mathcal E .
\]
In fact the weighted sum takes values fixed by the pro-unipotent Iwahori: modulo the maximal ideal that group is upper unitriangular and its first entry is \(1\). On this space \(U\) is a finite sum of local coset translations. At a place away from \(\ell\) the local group algebra with integral coefficients preserves the integral gauge. It keeps the fixed type at the old places because the supports are disjoint. The principal level at \(w\) is contained in the pro-unipotent Iwahori. Thus each composition defining \(B(f)\) gives a local integral endomorphism on the asserted domain. The image of \(B(1)\) is fixed by \(P\), since the unweighted input is parahoric-fixed and \(U\) preserves the parahoric-fixed space.\formalnote{13.27}{27}

For a line character \(\psi\circ a_P\) with \(\psi(n!)\ne1\), the local trace identity on an eigenform of this type is
\[
                      p(\widetilde U)\, T_1(g)y=T_p(g)y .
\]
We use the parahoric exponents at the unramified and ramified blocks. Suppress the embeddings, let \(F,T\) be untwisted Frobenius and inertia, and use the line character \(\zeta\) with \(\zeta^{n!}\ne1\). On a nonzero vector in the parahoric character space one can choose a nonvanishing exponent functional \(l\): the Jacquet exponents detect Iwahori vectors, and sorting characters within equal unit-character blocks preserves the parahoric character. Write \(\alpha_i\) for the normalized uniformizer values of its exponents and number the marked line \(0\). The split-place compatibility gives commuting \(F,T\), with \((T-1)(T-\zeta)=0\); for each integer \(b\ge0\), the roots of an annihilating polynomial of \(FT^b\) are \(\alpha_0\zeta^b,\alpha_1,\ldots,\alpha_{n-1}\), with the indicated multiplicities. Put \(Q(X)=\prod_{i>0}(X-\alpha_i)\). Restricting the annihilator to the image of \(T-1\) and cancelling \(\zeta^b\) gives
\[
  Q(\zeta^b F)(F-\alpha_0)(T-1)=0 .
\]
As a polynomial in \(\zeta^b\) this has degree at most \(n-1\). The \(n\) values with \(0\le b<n\) are distinct since every positive difference less than \(n\) divides \(n!\). Taking the constant coefficient and using \(Q(0)\ne0\) gives \(F(T-1)=\alpha_0(T-1)\). The same functional satisfies \(l(Uy)=\alpha_0 l(y)\), because the coset degree \(q^{n-1}\) cancels the modulus in the normalized torus character. Apply this argument to a functional detecting the difference of the two sides of the desired degree-one trace identity, if that difference is nonzero. The displayed matrix equality makes its value under \(l\) zero, a contradiction. Induction in the polynomial degree gives the assertion for \(p\); twisting \(F\) and \(U\) by the same scalar gives the version above. This calculation explains the factorial in the line-character hypothesis.\formalnote{13.28}{28}

For the other line characters one can still prove the needed congruence over the integral ring. Let
\(\mathcal E_g(f)=B(f)T_1(g)-W_P(f)T_p(g)\). It is additive in the weight. Finite Fourier decomposition shows that \(\mathcal E_g(f(\,\cdot-n!)-f)=0\) exactly. In this difference the components with \(\psi(n!)=1\) cancel, and the trace identity applies on every other component. With \(k_0=n!\), take \(f(x)=b(x)(b(x)+k_0)\). Reduction modulo \(m\) gives
\[
 f(x-k_0)-f(x)+2k_0 b(x)\equiv0\pmod m
\]
at every \(x\). Applying the additive operator to the integer identity yields \(2n!\mathcal E_g(b)\in m\mathcal E\), hence
\[
                            B(b)T_1(g)\equiv W_P(b)T_p(g).                   \tag{DN4}
\]
This explains the factor two even when \(\ell=2\). No projector onto an individual finite character has been applied to the integral lattice.\formalnote{13.29}{29}

\subsection{The polynomial in the local Hecke operator}
We explain why a single polynomial works on the relevant base systems at \(w\). Write \(r\) for the rank of the main summand and \(t=n-r>0\). A weak lift means here that every group element has eigenvalue \(1\) on its normalized Galois representation. For a base system of this kind let \(P_x\) be its spherical Hecke polynomial. The local annihilating polynomial for geometric Frobenius, and the eigenvalue \(1\) after twisting, force \(b_0=a(\phi)^{-1}\), read in \(\mathbb C\), to be a root of \(P_x\). Its ordinary complex absolute value satisfies \(|b_0|^2=q^r\) in the normalization of the lift.

Here are the two spherical calculations used with this root. Let \(y_0\) be a nonzero spherical eigenvector with eigenvalues \(t_i\) in a smooth unitary representation of \(\mathrm{GL}_n(E_w)\), and let \(P(X)=\sum_{i=0}^n(-1)^i q^{\binom i2}t_iX^{n-i}\). If \(b\) is a root and \(qb\) is not, then \(|b|^2\ge q^{n-2}\). If \(b\ne0\) and \(qb\) is also a root, then the cofactor \(P(X)/(X-b)\) evaluated at \(U\) kills \(y_0\).

Put \(d_0=n-1\) and let \(M\) be the mean over the pro-unipotent Iwahori. Write \(U'\) for the adjoint of the compressed operator
\[
                 U=q^{d_0}M\rho(\operatorname{diag}(\varpi,1,\ldots,1))M.
\]
Here are the finite counts for \(U\) and \(U'\). Write \(O_w\) for the valuation ring, \(k\) for its residue field, \(\Lambda=O_w^n\), \(I\) for the pro-unipotent Iwahori, and \(a=\operatorname{diag}(\varpi,1,\ldots,1)\). To a lattice \(L=g\Lambda\) attach \([L]=\rho(g)y_0\). This is well defined since \(\mathrm{GL}_n(O_w)\) fixes \(y_0\). Subspaces \(W\subset\Lambda/\varpi\Lambda\) of codimension \(i\) correspond to the lattices \(L\) between \(\varpi\Lambda\) and \(\Lambda\) of colength \(i\). The sum of \([L]\) over all of them is \(t_i y_0\). Let \(v_i\) be the sum restricted to \(e_1\in L\). It is \(I\)-fixed since \(I\) preserves \(k e_1\) modulo \(\varpi\). In particular \(v_0=y_0\) and \(v_n=0\).

For a row \(u\in k^{d_0}\) choose a lift \(\widetilde u\in O_w^{d_0}\), and let \(n_u\) be the identity with \(\widetilde u\) inserted to the right of its first diagonal entry. Representatives of \(I a I/I\) are \(n_u^{-1}a\). Indeed \(I\cap a I a^{-1}\) is cut out in \(I\) by vanishing of those first-row entries modulo \(\varpi\). For \(h\in I\) exactly one \(u\) kills them in \(n_u h\), since the lower-right block of \(h\bmod\varpi\) is invertible. Thus the index is \(q^{d_0}\); averaging \(\rho(h a)\) over \(h\in I\) gives the uniform mean of these translates on \(I\)-fixed vectors. This proves the stated normalization of \(U\). The mean \(M\) is an orthogonal projection by unitarity and invariance of Haar probability under inversion. Hence \(U'=q^{d_0}M\rho(a^{-1})M\).

Each lattice contributing to \(v_i\) is \(O_w e_1\oplus\widehat S\), where \(\widehat S\) is the inverse image in \(O_w^{d_0}\) of a codimension-\(i\) subspace \(S\subset k^{d_0}\). After applying \(n_u^{-1}a\) the lattice still contains \(\varpi\Lambda\) and reduces to the graph of \(-u|_S:S\to k\) in \(k\oplus k^{d_0}\). Every codimension-\((i+1)\) subspace missing \(e_1\) is such a graph, uniquely determining \(S\) and the functional on it. That functional has \(q^i\) extensions to \(k^{d_0}\). Subtracting the terms containing \(e_1\) from the full Hecke sum therefore gives \(Uv_i=q^i(t_{i+1}v_0-v_{i+1})\), including \(i=d_0\).

For the adjoint use \(\Lambda_+=a^{-1}\Lambda=O_w\varpi^{-1}e_1\oplus O_w^{d_0}\). The group \(I\) preserves both this lattice and the hyperplane \(H=\Lambda/\varpi\Lambda_+\) in \(V_+=\Lambda_+/\varpi\Lambda_+\). In the indicated basis of \(\Lambda_+\), the orbit of the first residue line under \(I\) consists of all \(q^{d_0}\) lines \(k(1,b)\), \(b\in k^{d_0}\). Indeed the first column below its top entry in the original basis belongs to \(\varpi O_w^{d_0}\) and can be chosen arbitrarily there. Thus \(q^{d_0}M\rho(a^{-1})v_{i+1}\) counts each codimension-\((i+1)\) subspace \(W\subset V_+\) once per such line in \(W\). If \(W\subset H\), it contains none. Otherwise the first coordinate on \(W\) is surjective and its fibre at \(1\) has \(q^{d_0-i-1}\) elements, one per line. The full sum over \(W\) is \(t_{i+1}[\Lambda_+]\) by translating the spherical sum. Those \(W\) in \(H\) correspond to lattices between \(\varpi\Lambda_+\) and \(\Lambda\) of colength \(i\) in \(\Lambda\). Since \(\varpi\Lambda_+=O_we_1+\varpi\Lambda\), their sum is \(v_i\).

Finally \([\Lambda_+]\) is \(I\)-fixed. The top spherical lattice is \(\varpi\Lambda\), so \(\rho(\varpi\,{\rm Id})v_0=t_n v_0\) with \(t_n\ne0\). The unique lattice for \(v_{d_0}\) is \(\varpi\Lambda_+\), whence \(v_{d_0}=t_n[\Lambda_+]\) by centrality. In particular \(U'v_0=q^{d_0}[\Lambda_+]\). Collecting the two counts gives
\[
\begin{split}
 &v_0=y_0,\quad v_n=0,\qquad Uv_i=q^i(t_{i+1}v_0-v_{i+1}) \quad(0\le i<n),\\
 &U'v_0=q^{d_0}t_n^{-1}v_{d_0},\\
 &U'v_{i+1}=q^{d_0-i-1}\bigl(t_{i+1}q^{-d_0}U'v_0-v_i\bigr)
                          \quad(0\le i<d_0).
\end{split}
\]
The last case of the forward recurrence gives \(Uv_{d_0}=q^{d_0}t_n v_0\), so \(UU'v_0=q^{2d_0}v_0\).\formalnote{13.30}{30}
Let \(P_h(X)=\sum_{i=0}^h(-1)^i q^{\binom i2}t_i X^{h-i}\) and \(D_h=(P_h(X)-P_h(b))/(X-b)\). The first recurrence gives \(P_h(U)v_0=(-1)^h q^{\binom h2}v_h\). In particular \(P(U)v_0=0\). Substitution in the second recurrence, using \(D_{h+1}=bD_h+P_h\), gives for \(0\le h\le n\)
\[
 U'\bigl((bD_h+P_h)(U)v_0\bigr)
       =q^{-h}P_h(qb)\,U'v_0+q^{n-2}D_h(U)v_0 .
\]
For the step from \(h\) to \(h+1\), multiply the identity by \(b\) and add
\[
                U'P_{h+1}(U)v_0=(-1)^{h+1}q^{\binom{h+1}2}U'v_{h+1}.
\]
At the last step use \(v_n=0\) and the formula for \(U'v_0\). The identity at \(h=0\) has \(D_0=0,t_0=1\).

Take \(h=n\), \(P(b)=0\), \(z_b=D_n(U)v_0\), \(A_b=P(qb)\). Then \(Uz_b=bz_b\) and \(bU'z_b=q^{-n}A_bU'v_0+q^{n-2}z_b\). Pair this last equality first with \(U'v_0\), using \(UU'v_0=q^{2(n-1)}v_0\). It gives \(b\langle v_0,z_b\rangle(q^n-1)=A_b\|v_0\|^2\), with the pairing linear in its second variable. Next pair with \(z_b\) and use \(Uz_b=bz_b\); elimination yields
\[
 (|b|^2-q^{n-2})\|z_b\|^2
                =\frac{|A_b|^2\|v_0\|^2}{q^n(q^n-1)} .
\]
This proves the lower bound when \(A_b\ne0\). When \(A_b=0\) and \(b\ne0\), \(z_b\) is orthogonal to \(v_0\) and \(U'z_b=q^{n-2}b^{-1}z_b\). It is therefore orthogonal to all \(U^i v_0\), including their combination \(z_b\). Positivity proves \(z_b=0\).\formalnote{13.31}{31}

The adjoint duality between spherical Hecke operators of degrees \(i\) and \(n-i\) implies that the roots of \(P\) are symmetric under \(b\mapsto q^{n-1}/\bar b\). In fact the dual of the degree-\(i\) double coset is, up to the inverse central uniformizer, the degree-\((n-i)\) coset. Pairing spherical eigenvectors gives \(\bar t_i=t_{n-i}/t_n\) and \(|t_n|=1\). Substituting these in the coefficients yields \(P(X)=(-1)^n t_n q^{-\binom n2}X^n\overline{P(q^{n-1}/\bar X)}\). For \(P_x\), the involution interchanges \(b_0q^i\) and \(b_0q^{t-1-i}\). Starting at the root \(b_0\), one obtains successive roots \(b_0q^i\) up to the middle of the list. If \(b_0q^{i-1}\) were the last root for \(1\le i\le(t-1)/2\), its squared modulus \(q^{n-t+2i-2}\) would be strictly smaller than \(q^{n-2}\). Reflection then gives the other half, with the same reasoning for the central two roots when \(t\) is even. The \(t\) roots are distinct.

The same cofactor calculation singles out the top root when pairing with \(y_0\). The polynomials killing \(v_0\) under \(U\) form an ideal. This ideal contains \(P\) and \(P/(X-b)\) for every distinct nonterminal root \(b\) in the chain. Their greatest common divisor is \(P/\prod_{i<t-1}(X-b_0 q^i)\), even if the roots have higher multiplicities. Denote this cofactor by \(R\). It kills the whole span of the \(v_i\), since each \(v_i\) is a polynomial in \(U\) applied to \(v_0\). That span is stable under \(U'\), so \(\overline R(U')v_0\) lies in the span and pairs to zero with it. Thus \(\overline R(U')v_0=0\), and \(v_0\) pairs to zero with \(R(U)y\) for every \(y\) fixed by the pro-unipotent Iwahori. This is the orthogonality needed for the norm.

For all selected systems choose \(H\in O[X]\), \(H(1)=1\), so that after rescaling \(X\) by \(a(\phi)^{-1}\) the factor complementary to the chain divides \(H\). To see the normalization, write \(B_0=\prod_{i<t}(X-q^i)\) over the \(\ell\)-adic field, and let \(A_x\) be the complementary polynomial after correcting the leading scalar to make \(A_xB_0\) monic. The scaled Hecke coefficients are congruent modulo the maximal ideal to the coefficients of the normalized Galois characteristic polynomial at the selected Frobenius. This uses the attached representation on the prescribed set of Frobenii. On a system congruent to the selected lift, these coefficients are in turn congruent to those of \(S B_0\), where \(S\) is the characteristic polynomial of the twisted main summand: the exact decomposition at the lift follows from equality of the traces of all powers and Newton's identities. One arranges \(|S(1)|_\ell\ge1\) at the auxiliary place. Then \(A_x\) is integral and \(A_x-S\) has all coefficient norms \(<1\). Indeed use the nonarchimedean Gauss norm and \(\|B_0\|_{\rm Gauss}=1\), since \(B_0\) is integral and monic. The inequalities for the products transfer to the factors. Thus \(|A_x(1)|_\ell=1\), and \(A_x/A_x(1)\) is an integral multiple of the complementary factor with value \(1\) at \(1\). Multiply these polynomials over the finitely many relevant systems.

For the requirement \(|S(1)|_\ell\ge1\), the scalar-inertia hypothesis on the main summand supplies an inertia element whose scalar \(y\) remains of residual order \(>4r+2\) after taking its \(m\)-th power, and on which the twisting and cyclotomic characters are trivial. For a prescribed \(g\), the \(r+1\) scalings \(1,y^m,\ldots,(y^m)^r\) are distinct in the residue field. No eigenvalue of the twisted main summand at \(g\) can have distance \(<1\) from more than one of their inverses. There are at most \(r\) eigenvalues. Some choice of scaling therefore gives \(|\det(1-(y^m)^i N_{\mathrm{main}}(g))|_\ell\ge1\). This persists on a neighborhood by continuity. Translating the Frobenius target by this \(m\)-th power leaves the required cyclotomic condition unchanged. The prime-triangle construction invoked in \S\ref{sec:congruence-detail} allows precisely these translated targets together with the required Frobenius compatibility conditions. This choice is necessary for dividing by \(A_x(1)\).\formalnote{13.32}{32}

\subsection{The single-place derivative and the many-place factorization}
Take \(p(X)=(X-q^{t-1})H(X)\) and define
\[
                             \partial_j=-C B(b) C' .
\]
Using (DN4), then transfer and (DN3), gives for each \(g\in G_E\)
\[
 \partial_jT_1(g)\equiv -CW_P(b)C'T_p(g)
      \equiv-W_G(b)T_p(g)\equiv D_jS_p(g) .
\]
Local operators commute with the diagonal traces. As \(q\equiv1\pmod m\), \(S_p(g)\equiv S_{(X-1)H}(g)\) by linearity. This proves the asserted positive sign in the Frobenius tie, including the case \(t=1\).

The construction uses one \(\Gamma_j\) and one variable \(X_j\) at each split pair. To obtain the trace relation at its conjugate, write \(s^*=(c_\infty s c_\infty^{-1})^{-1}\) and \(\chi=\epsilon^{-1}\). Conjugate self-duality of the normalized \(N_x\) and cyclic invariance of trace give, for \(s=\gamma_j\) or \(\phi_j\),
\[
 {\rm tr}\,N_x((s^*)^i g)=
      \chi((s^*)^i g)^{t-1}{\rm tr}\,N_x(s^i g^*) .
\]
We have \(\chi(\gamma_j^*)=1\) and \(\chi(\phi_j^*)\equiv1\bmod m\). Expand \((X-1)H_j(X)\) in powers of \(X\). The conjugate trace tie is therefore \(\chi(g)^{t-1}\) times the original tie at \(g^*\), up to a linear combination of \(m\) times traces with coefficients in the polynomial algebra. Those terms already belong to the trace ideal. This supplies the two Galois relations per \(X_j\) used in Section~\ref{sec:trace-ideal}.\formalnote{13.33}{33}

For a Hecke operator \(\delta'\) supported at base level only on the selected systems (its values at higher level can be arbitrary), we have \(\delta'D B(1)=0\). Indeed \(DB(1)\) has image at base level. Each base test vector in a contributing system is \(\Gamma\)-invariant, so pairing it with \(DB(1)y\) multiplies the pairing with \(B(1)y\) by \(\#\Gamma\). The latter pairing vanishes by the top-root orthogonality just proved, since \(p(a(\phi)X)\) is a multiple of the cofactor \(R\). Decomposing into systems and using a positive invariant pairing proves the vanishing.

For \(\sigma\in G\) put \(h_L=s_{\sigma L}^{-1}\sigma s_L\in P\). Right multiplication of \(C'\) by \(\rho(\sigma)\) permutes the summands up to these \(h_L\). Moreover \(B(f)\rho(h)=B(f(\,\cdot-a_P(h)))\). Thus
\[
 B(b) C'(\rho(\sigma)-1)\equiv B(1)Q_\sigma,\qquad
 Q_\sigma=-\sum_L b(a_P(h_L))\rho(s_L^{-1}) .
\]
As \(B(1)\) is fixed on the left by \(P\), (DN1) for \(P\) gives
\(c B(1)=(\sum_{h\in P}\rho(h))\Psi_P B(1)\). Multiplying by \(-C\) gives
\[
 c\partial_j(\rho(\sigma)-1)\equiv D_j Z_\sigma,\qquad
 Z_\sigma=-\Psi_P B(1)Q_\sigma,\qquad \delta'D Z_\sigma=0 .             \tag{DN5}
\]
For the last equality, apply \(\delta'D\) to the preceding exact identity for \(cB(1)\). The left side vanishes. The right side is \(\#P\,\delta'D\Psi_P B(1)\) since \(D\rho(h)=D\). We may cancel \(\#P\) in the characteristic-zero endomorphism ring before reducing.

Now use local operators at all \(j\). Those at different places commute by disjoint support. In the polynomial ring with central variables modulo \(X_j^2\) write
\[
       \prod_j(D_j+X_j\partial_j)=\sum_{S\subseteq J}N_S X_S,\quad
        X_S=\prod_{j\in S}X_j,\qquad N_\varnothing=D .
\]
It annihilates the tied ideal modulo \(\mu\). Indeed multiplying the factor at \(j\) by \(T_1(g)-X_j S_{(X-1)H_j}(g)\) gives a zero constant term by (DN3) and a zero linear term by the derivative identity. Terms having \(X_j^2\) give no squarefree coefficient, and the depth generators are multiples of \(m\). Taking further Hecke-polynomial multiples and summing preserves these assertions because the \(N_S\) commute with \(\mathbb T\).

Put \(\Delta=c^2\delta'\). We prove that \(\Delta^i N_S\equiv Y_S D\) for some integral \(Y_S\) whenever \(i\ge|S|\). If \(S=\varnothing\), take \(Y_S=\Delta^i\). For \(S\ne\varnothing\) put \(V_0=c\delta'\Delta^{i-1}N_S\). We show that \(V_0\) is right invariant modulo \(\mu\) under the whole \(\Gamma\). For \(j\notin S\) the factor \(D_j\) makes this exact. For \(j\in S\), (DN5) replaces \(cN_S(\rho(\sigma)-1)\) by \(N_{S\setminus\{j\}}Z_\sigma\). Move \(\delta'\) across \(\Delta^{i-1}N_{S\setminus\{j\}}\) and apply induction to the latter. The result contains \(Y_{S\setminus\{j\}}\delta'DZ_\sigma=0\). Thus right invariance holds for each factor, hence for all \(\Gamma\). Formula (DN1) for \(\Gamma\) on the right of \(V_0\) now gives
\[
                    \Delta^i N_S=c V_0\equiv V_0\Psi_\Gamma D .
\]
Taking \(i=|J|\) gives the congruence for every coefficient. This congruence and the integral norm lift \(Du=dc\,l\) give the cancellation in \S\ref{sec:trace-ideal}.\formalnote{13.34}{34}


\partheading{Galois representations and integral comparison}

\editionchapter{Local cohomology and the integral period lattice}\label{ch:local}

The preceding chapter reduces a denominator bound to lengths of Galois cohomology groups with finite coefficients. Here we estimate the part at the coefficient prime. Example~\ref{ex:local-lattice-extension} illustrates the obstruction to extending a map from a sublattice.

The background calculation constructs an integral frame for a crystalline lattice, giving the part of \rootlink{grh/src/Lib/PadicHodge/Kisin/FrameTheorem.lean}{Breuil--Kisin theory} used below. The estimate specific to the determinant argument bounds the length of the crystalline cocycles on the Tate dual lattice modulo those whose finite cup products vanish against every allowed crystalline cocycle. This \rootlink{grh/src/New/Congruences/Lifting/CupOrthogonalityBounds.lean}{local length bound} enters the sum over coefficient primes.

\section{The loss at places above the coefficient prime}
\label{sec:crystalline-loss}

Let \(v\mid\ell\) be a finite place of \(E\), and let \(L\) be a Galois-stable integral lattice. A trace congruence gives cocycles with values in \(L/\ell^nL\). The local Selmer condition at \(v\) consists of reductions of integral cocycles \(z\) in \(L\) that become coboundaries after extending scalars to the crystalline period ring \(B_{\max}\). Thus there must be a period vector \(y\) with \(z(g)=gy-y\) for every \(g\in G_{E_v}\). We bound the length contributed by the congruence cocycles modulo those that already satisfy this condition. First we realize each congruence cocycle by reduction from a second integral lattice. We then use the finite cup product and local duality to compare the two reductions. Summing the resulting local bounds over coefficient primes costs at most a constant times \(f(m/k)b\log b\).

\subsection{An integral lift of a finite-coefficient cocycle}

Let \(O\) be the integer ring of a sufficiently large finite Galois extension of \(\mathbb Q_\ell\), \(q=\ell^n\), \(n\ge1\). Let \(u_x\) be the finite family of integral representations used for the trace congruence, and let \(L\) be an integral summand of one family member, with action \(\sigma\). We write \(G\) for the global Galois group and \(\Gamma=G_{E_v}\) for the local group. For \(A=\prod_x\operatorname{Mat}_d(O)\), put \(U_g=(u_x(g)-1)_x\). Give \(A\) the action by left multiplication with \((u_x(g))_x\). The defining condition on a Ribet cocycle \(w:G\to L/qL\) is
\[
 w(gh)=w(g)+\sigma(g)w(h),\qquad
 \sum_g c_g U_g=0\quad\Longrightarrow\quad\sum_g c_g w(g)=0
\]
for every finitely supported family \((c_g)\) in \(O\). Thus \(U_g\mapsto w(g)\) gives a linear map \(\phi:I\to L/qL\), where \(I=\sum_g O U_g\). It is equivariant since \(gU_h=U_{gh}-U_g\).

\begin{example}[Extending a map to a lattice of full rank]\label{ex:local-lattice-extension}
Take \(A=\mathbb Z_\ell e_1\oplus\mathbb Z_\ell e_2\),
\(I=\ell^2\mathbb Z_\ell e_1\), and \(q=\ell\). Set
\(\phi(\ell^2 a e_1)=a\bmod\ell\). A linear map
\(A\to\mathbb F_\ell\) sends \(\ell^2e_1\) to zero, so it cannot
extend \(\phi\). We can, however, enlarge \(I\) to a lattice of full
rank. Prescribing zero on \(\ell^2 A\) would contradict
\(\phi(\ell^2 e_1)=1\). Prescribing zero on \(\ell^3 A\) works
because \(I\cap\ell^3 A=\ell I\). The extension is
\[
 I+\ell^3 A\longrightarrow\mathbb F_\ell,\qquad
      \ell^2 a e_1+\ell^3 b e_2\longmapsto a\bmod\ell .
\]
Its graph is free on \((\ell^2 e_1,1)\) and \((\ell^3 e_2,0)\).
The first projection realizes it inside \(A\), and the second
recovers the prescribed values on \(I\).
\end{example}

For the general construction apply the Artin--Rees lemma to the finite \(O\)-modules \(I\subseteq A\) and the ideal \((q)\) \cite[Tag~00IN]{arith-stacks}. It gives an \(s\ge1\) with
\[
 I\cap q^s A=q(I\cap q^{s-1}A)\subseteq qI .
\]
Thus \(\phi\) kills the intersection, and extending it by zero on \(q^sA\) gives an equivariant map \(\widetilde\phi:I+q^s A\to L/qL\). Let \(Y\) be its graph. The first projection identifies \(Y\) with a finite free \(O\)-lattice in \(A\). Reading the matrix columns gives an equivariant embedding into a sum of copies of the \(u_x\).

In \(Y\) consider \(z(g)=(U_g,w(g))\) and \(y_0=(q^s{\bf1},0)\), where \({\bf1}\) denotes the family of identity matrices. Both coordinates of \(z\) satisfy the cocycle identity. Moreover
\[
                          g y_0-y_0=q^s z(g).
\]
This identity proves that \(z|_\Gamma\) becomes a coboundary in \(B_{\max}\otimes Y\), since \(q\) is invertible in \(B_{\max}\). Thus it is an integral crystalline cocycle. The local class of lattices \(\mathcal C_v\) specifies a period condition, allowed Hodge--Tate jumps at each embedding, and at the level primes annihilating polynomials on inertial periods. It is closed under equivariant embeddings into sums of its members. Consequently \(Y|_\Gamma\in\mathcal C_v\), and the second projection sends its crystalline cocycle to \(w|_\Gamma\). This proves a lifting assertion for every congruence cocycle, without assuming that its value in \(L/qL\) lifts to a crystalline cocycle of \(L\) itself.\formalnote{14.1}{1}

\subsection{The finite pairing and the rank that it costs}

Write \(Z_f(T)\) for the \(O\)-module of cocycles with values in \(T\) which become coboundaries over \(B_{\max}\). For \(T=L^\vee(1)\) let \(H\) consist of the integral cochains whose finite Tate cup products with crystalline cocycles of every \(Y\in\mathcal C_v\) are coboundaries, for all integral pairings allowed in the class. Define
\[
 d_{\mathcal C_v}(T)=
 \operatorname{len}_*\bigl(Z_f(T)/(Z_f(T)\cap H)\bigr),
 \qquad\operatorname{len}_*(O/\ell O)=1 .
\]
The normalization divides usual \(O\)-length by the absolute ramification index of \(O\). For a module \(N\) of restrictions constructed above let \(A_n\) denote the reductions of \(Z_f(L)\). Then
\[
                    \operatorname{len}_*(N/(N\cap A_n))\le d_{\mathcal C_v}(T).
                    \tag{1}
\]

Here is the use of duality in (1). Given a class lattice \(Y\) and an equivariant \(t:Y\to L/qL\), lift \(t\) as an \(O\)-linear map \(Q:Y\to L\). The differences \(gQ-Q\) lie in \(q\operatorname{Hom}_O(Y,L)\). Dividing the differences by \(q\) gives an integral cocycle which is a coboundary after \(q\) is inverted. For the evaluation form \(\Psi(y,\varphi)=\varphi(Qy)\), the same calculation shows that \((g\Psi-\Psi)/q\) is integral and crystalline. Composing evaluation with the integral functional on \(O\) giving a perfect \(\mathbb Z_\ell\)-pairing yields exactly an allowed pairing on \(Y\times T\) modulo \(q\). Consequently every \(s=t\circ z\) in \(N\) cups trivially with \(H\). The same statement holds for \(A_n\) by taking \(Y=L\) and \(t\) reduction.

We give the index calculation. Set \(V=L/qL\), \(V'=T/qT\). The evaluation just chosen identifies \(V'\) with the finite Cartier dual of \(V\). Put \(d_v=[E_v:\mathbb Q_\ell]\). For any subgroup \(J\subset Z^1(\Gamma,V')\) and its cup annihilator \(J^\perp\subset Z^1(\Gamma,V)\) we have
\[
                        |J^\perp|\,|J|\le |V|^{2+d_v}.
\]
Indeed let \(\overline J\) be the image of \(J\) in \(H^1(\Gamma,V')\). By finite local Tate duality the annihilator of \(\overline J\) in \(H^1(\Gamma,V)\) has order \(|H^1(\Gamma,V)|/|\overline J|\). The fibers on passing from cocycles to classes have size \(|B^1(\Gamma,V)|\) on the first side and at most \(|B^1(\Gamma,V')|\) in \(J\). Now
\[
 |H^1(\Gamma,V)|=|H^0(\Gamma,V)|\,|H^0(\Gamma,V')|\,|V|^{d_v},
 \qquad |B^1(\Gamma,V)|=|V|/|H^0(\Gamma,V)|.
\]
The first equality is the local Euler characteristic formula and duality in degree two. Substitution gives the claimed bound.

Let \(A'_n\) be the reduction of \(Z_f(T)\). Reduction of integral crystalline cocycles has kernel exactly \(qZ_f\): if every value of a cocycle is divisible by \(q\), dividing by \(q\) preserves the cocycle rule and the coboundary condition over \(B_{\max}\). The Bloch--Kato rank formula says that the \(O\)-rank of \(Z_f\) is the rank of its lattice plus the number of negative jumps summed over embeddings at \(v\). Each integer \(u\) in \(L\) and its mate \(-u-1\) in \(T\) contribute exactly one negative jump. Thus
\[
                        |A_n|\,|A'_n|=|V|^{2+d_v}.
\]
Take \(J\) to be the reduction of \(H\cap Z_f(T)\). Both \(N\) and \(A_n\) annihilate \(J\). Since \(N/(N\cap A_n)\simeq (A_n+N)/A_n\), the preceding estimates give
\[
 [N:N\cap A_n]\le |J^\perp|/|A_n|
 \le |A'_n|/|J|\le[Z_f(T):Z_f(T)\cap H].
\]
The last comparison is the quotient map induced by reduction. All finite indices are powers of \(\ell\), giving the divisibility and the normalized length inequality (1). Imposing all local linear relations on \(Z_f(L)\) can only decrease this colength, as every reduction already obeys those relations.\formalnote{14.2}{2}

Put \(\rho=\operatorname{rank}_O T\). The same rank formula gives \(\operatorname{rank}_O Z_f(T)\le (1+d_v)\rho\). The cocycles form a finite torsion-free \(O\)-module. Hence a cup exponent \(s\), meaning \(\ell^s Z_f(T)\subset H\), gives
\[
 d_{\mathcal C_v}(T)\le\operatorname{len}_*\bigl(Z_f(T)/\ell^sZ_f(T)\bigr)
                  \le 2d_v\rho s.                                      \tag{2}
\]
The degree factor here is local, which matters when many places lie above \(\ell\).\formalnote{14.3}{3}

\subsection{The integer which records the exceptional primes}

Write \(h=r^8\), \(k=r+1-m\), \(t=5r\), \(Q_0=kt\), \(M=p_0p_1\), and let \(b\) denote the integral large weight. Define
\begin{align*}
 I_b&=\{j\in\mathbb Z_{>0}:\ b-30h^3<j\le b+30h^3\},\\
 B&=(2b)!\left(\prod_{j\in I_b}j\right)^{h^3},\qquad
 \lambda=v_\ell(B),\qquad K=\left\lfloor\frac{2\lambda}{Q_0}\right\rfloor,\\
 D_\ell&=10^6(m/k)\bigl(b\,v_\ell(M)+\lambda\bigr).
\end{align*}
The number \(D_\ell\) is an upper bound, and need not be integral. For a main summand set \(x=1\); for a one-dimensional Tate summand set \(x=t\). We prove
\[
                         d_{\mathcal C_v}(T)\le d_v\,\rho\,D_\ell/x.      \tag{3}
\]
Two observations account for both factors of \(B\). The factorial contains \(\lfloor2b/\ell\rfloor\) multiples of \(\ell\), so \(2b<\ell(\lambda+1)\). If \(\ell\mid j\) for any \(j\in I_b\), then \(\lambda\ge h^3\). This last threshold will allow one to exclude a short interval of multiples simultaneously for all allowed weights.

Here are the weight data used below. Write \(\eta=h^2+1\), \(W=4t\), \(G_0=a/x_1\), and \(\delta=2ih+h-2t\), where \(0\le i\le m\) is the comparison index. If the embedding \(\sigma\) induces the infinite place \(\tau\), put \(a_\sigma=\kappa_\tau+\delta\). Call \(\sigma\) thick when \(\tau\in\Phi_b\), so \(\kappa_\tau=b\), and thin otherwise, when \(\kappa_\tau=G_0\). The allowed jumps \(U\) of the class and the actual jumps \(U'\) of \(T\) both lie in \([-\eta,\eta]\) together with a far interval. Their far intervals are respectively \(-\epsilon_\sigma[a_\sigma,a_\sigma+W]\) and \(+\epsilon_\sigma[a_\sigma,a_\sigma+W]\). This reversal comes from the rule \(u\mapsto-u-1\) for the Tate dual, with a margin built into the interval. For a main summand, every negative element of \(U\) or \(U'\) is at most \(-h/2\), and at most \(\mu=\lfloor\rho m/(kx)\rfloor\) jumps of \(T\), with multiplicity, lie in the low interval. For a Tate summand \(U'\) is the singleton \(\{-c-1\}\subset[-\eta,0]\). We still use the same symbol \(\mu\) in the bound. The admissible scales give
\begin{align*}
 99\bigl(Q_0^2+(W+1)Q_0+Q_0+W+m\bigr)&\le h, &
 99h^3&\le G_0,\\
 99(G_0h^3+G_0)&\le b, &
 2\log(4b)&<m,\\
 99\bigl((\eta+W)Q_0+\eta+h\bigr)&\le h^3, &
 \delta&<\eta .
\end{align*}
For crystalline places the local comparison uses a Kummer chart of degree \(d_v\) and ramification index at most \(h\). A letter is a residue-field unit indexing a set of embeddings by the Teichmüller lift. If \(\ell>Q_0\) the chart is unramified and \(\ell\ge7\). For \(\ell\le Q_0\) the Behrend condition supplies sparsity: along \(z\mapsto z^\ell\), no three thick letters occur with both successive gaps between \(1\) and \(\lfloor\log_\ell(h^3)\rfloor\). The chart is unramified with \(\ell\ge7\), or the signs on the letters agree. The source constructs these charts after enlarging \(O\).\formalnote{14.4}{4}

\paragraph{Crystalline places with small prime or large valuation.}
For \(\ell\le Q_0\) one may widen the low interval to \([-2G_0,2G_0]\), absorbing each thin far interval. The sparse-letter comparison theorem for these windows gives a cup exponent
\[
                        10^4K(\lfloor\log_\ell h^3\rfloor+1)
\]
provided \(2(b+\delta)<h^2K\). For this condition, \(2b<\ell(\lambda+1)\) gives \(K>4b/Q_0^2-2/Q_0-1\). Since \(b\ge 2Q_0^2\), this implies \(K>b/Q_0^2\); now \(h\ge99 Q_0^2\) and \(\delta<b\) suffice. Also \(\lfloor\log_\ell h^3\rfloor<m\), since \(\log 2>1/2\), \(h^3<4b\), and \(2\log(4b)<m\).

For \(\ell>Q_0\) and \(\lambda\ge h^3\) the chart is unramified. The unramified comparison for any pattern of thick letters has hypotheses
\[
        b+\delta<K\ell^2,\qquad 2(G_0+\delta)<K\ell,\qquad W+4\eta\le K .
\]
Indeed \(K\ge \lambda/Q_0\), since \(\lambda\ge h^3\ge Q_0\), and \(K\ge h^3/Q_0\ge99(\eta+W)\). Also \(K\ell\ge\lambda\ell/Q_0>b/Q_0\), using \(\lambda\ge1\) and \(2b<\ell(\lambda+1)\). As \(b\ge99G_0h^3\), this dominates \(2(G_0+\delta)\). Finally \(\delta<b/2\) and \(\lambda\ge3\) give \(b+\delta<3\ell(\lambda+1)/4\le\lambda\ell<K\ell^2\). The resulting cup exponent is
\(10^4(K+\lfloor\log_\ell(2(b+\delta+W+4\eta))\rfloor)\).
The logarithm has floor strictly less than \(m\), by the scale inequalities and \(2\log(4b)<m\), and \(K\ge1\). Both cases therefore give an exponent at most \(10^4 K m\). By (2) their defect is at most \(2\cdot10^4 d_v\rho Km\).\formalnote{14.5}{5}

\paragraph{Crystalline places with small valuation and large prime.}
Suppose \(\ell>Q_0\) and \(\lambda<h^3\). Put \(D=2(\lambda+1)\), \(n_0=6(D+\eta+W)\). Since \(\lambda+1\le h^3\), the two inequalities \(2b<\ell(\lambda+1)\) and \(99G_0h^3\le b\) imply \(\ell>198G_0\). Thus \(G_0+\delta+n_0<\ell\). No nonzero multiple of \(\ell\) lies within \(n_0\) of \(0\) or \(\pm(G_0+\delta)\). Since \(n_0+\delta<30h^3\), every positive integer within \(n_0\) of \(b+\delta\) belongs to \(I_b\). A multiple there would force \(\lambda\ge h^3\), a contradiction. Reflection gives the assertion at \(-b-\delta\) also.

The no-nearby-multiple comparison theorem applies with height \(D\). Its window hypotheses include \(a_\sigma+W+\eta\le(\ell-3)D\) and \(n_0<a_\sigma\). Indeed \(a_\sigma\le b+\delta\), \(\ell D>4b\), and \(3D+\delta+W+\eta<b\); also \(n_0<G_0\). The theorem gives zero defect if \(U'\subset[-\eta,0]\), or if both sets of jumps avoid \([-D,0)\). In the other case, when the low multiplicity is at most \(\mu\), it gives
\[
                    11d_v\{\mu(D+1)+\rho(W+1)\}.
\]
For the Tate summand the zero criterion always applies in this range. For the main summand it applies when \(\lambda<\lfloor h/5\rfloor\), because \(D<h/2\). Otherwise \(\lambda\ge\lfloor h/5\rfloor\), so \(W+1\le K\) by \(99(W+1)Q_0\le h\), and \(D+1\le6\lambda\). This proves in all three crystalline cases that
\[
 d_{\mathcal C_v}(T)\le 10^5 d_v(\rho K m+\mu\lambda)
       \le 3\cdot10^5 d_v\frac{\rho m\lambda}{kx}
       \le d_v\frac{\rho D_\ell}{x}.
\]
The middle inequality uses \(x\le t\) and drops the floors in \(K,\mu\).\formalnote{14.6}{6}

\paragraph{The two possible tame level primes.}
When the local class is of tame type, \(\ell\) is one of the two split level primes. For the action of each inertial Weil element on periods the type supplies an annihilator \(X^{\operatorname{lcm}_{1\le j\le r}(\ell^j-1)}-1\) and a monic annihilator over \(\mathbb Z_\ell\) of degree at most \(N_0=(r+t)r\). Splitting makes the local group the full group at \(\ell\), so the signs and window center are fixed at that place. The comparison requires \(\ell>a_\sigma+6(W+\eta)+r^2(N_0+1)^2+2\), which follows from \(p_i\ge b^{10}\) and the scale bounds. Its tame chart has ramification \(\operatorname{lcm}_{1\le j\le r}(\ell^j-1)\le\ell^{r^2}\); over this chart the class becomes crystalline. The tame comparison gives a cup exponent
\[
                 s=3r^2\bigl((r+t)r+1\bigr)^2+150(\eta+W+1).
\]
The scale bounds give \(2s\le 10^6 mb/(kt)\). As \(v_\ell(M)\ge1\) and \(x\le t\), (2) again gives (3).\formalnote{14.7}{7}

\subsection{Summing the local constraints}

The intersection of the kernels of the local relation maps is the kernel of their direct sum. The colength of the global module on imposing all local relations is therefore at most the sum of the local colengths. There is one main summand of rank \(r'\ge1\) and \(t'\le t\) Tate summands. The local degrees sum to \([E:\mathbb Q]=2f\), so (1) and (3) give
\[
          \sum_{v\mid\ell}\sum_{\rm summands}d_{\mathcal C_v}(T)
         \le 2f(r'+t'/t)D_\ell\le r'\,\operatorname{loc}_\ell,
         \qquad\operatorname{loc}_\ell=4fD_\ell.
\]
In particular adding all the Tate characters costs at most one additional unit of main rank. Finally \(\sum_\ell v_\ell(A)\log\ell=\log A\) for every positive integer \(A\). There are at most \(60h^3\) factors in \(\prod_{j\in I_b}j\), and \(M\le b^{22}\). Consequently
\begin{align*}
 \sum_\ell \operatorname{loc}_\ell\log\ell
 &\le 4\cdot10^6 f(m/k)
 \{b\log M+\log((2b)!)+60h^6\log(b+30h^3)\}\\
 &\le 4\cdot10^6 f(m/k)\,28b\log b .
\end{align*}
For the last inequality use \(60h^6\le b\) and \(30h^3\le b\) from the scale bounds, \(\log((2b)!)\le2b\log(2b)\), and \(b\ge2\). This implies the manuscript's allowance \(4\cdot10^6 e^{180} f(m/k)b\log b\).\formalnote{14.8}{8}

Section~\ref{sec:local-lengths} applies finite local Tate duality, the local Euler characteristic formula and the crystalline cocycle rank formula with their normalizations here. The argument further uses closure of the prescribed period class under embeddings and the integral chart estimates that enter the four cup comparisons derived in \S\ref{sec:cup-comparisons}. In particular the vanishing in the small-valuation case requires the integral cup comparison. The weight counts and sparse local charts are the arithmetic inputs furnished by the admissible data. The graph construction, finite-pairing index argument, deduction of (3) from the comparisons, and summation of the resulting losses yield the term used by the global congruence estimate.


\section{From local duality to the integral cocycle count}
\label{sec:local-lengths}

The estimate of \S\ref{sec:crystalline-loss} uses an annihilator under local duality. To evaluate its size we need the ranks of the integral crystalline cocycles and their duals. Fix \(v\mid p\), write \(K=E_v,\ n=[K:\mathbb Q_p]\), and identify the decomposition group with \(G_K\) through the open embedding into \(G_{\mathbb Q_p}\). Let \(O\) be the ring of integers in a finite extension of \(\mathbb Q_p\), \(L\) the free \(O\)-lattice, and \(T=\operatorname{Hom}_O(L,O)(1)\). Write \(d=\operatorname{rank}_{\mathbb Z_p}L\). The local class supplies potential crystallinity of \(L[1/p]\) and its Hodge jumps. The dual also has potential crystallinity, with jumps obtained by duality and the Tate twist.\formalnote{14.9}{9}

\subsection{Identifying the finite condition and recovering the lattice}

For \(V=L[1/p]\), let \(Z_f(V)\) denote the continuous cocycles \(c:G_K\to V\) admitting a \(y\in B_{\max}\otimes_{\mathbb Q_p}V\) with \(c(g)=gy-y\). The project uses the \(p\)-adic completion of \(A_{\inf}[\xi/p]\), inverts \(p\), and then inverts \(\mu=[\epsilon]-1\). Its cyclotomic period is a unit multiple of \(p\mu\) in \(A_{\max}\). For Fontaine's classical period \(t=\log[\epsilon]\), Colmez proves directly that \(t/\mu\) is an \(A_{\max}\)-unit \cite[Lemma~III.3.9]{arith-colmez}. Thus localizing at \(\mu\) gives Colmez's \(B_{\max}\). Colmez's inclusions \cite[III, \S2]{arith-colmez}
\[
                \varphi(B_{\max})\subset B_{\rm cris}\subset B_{\max}
\]
identify the two splitting conditions directly. If \(y\) is a primitive over \(B_{\max}\), apply \(\varphi\otimes1\). Frobenius commutes with Galois and fixes \(\mathbb Q_p\), so \(g(\varphi\otimes1)y-(\varphi\otimes1)y=c(g)\) in \(B_{\rm cris}\otimes V\). The reverse implication is inclusion. It follows that \(Z_f(V)/B^1(K,V)\) is precisely the Bloch--Kato group \(H_f^1(K,V)\). This uses the local absolute Galois group at the coefficient prime.\formalnote{14.10}{10}

\phantomsection\label{route:23}\leanroute{The formalization computes the rank of the space of crystalline cocycles from a period basis and its filtration. Once the splitting condition has been identified with \(H^1_f\), the Bloch--Kato dimension formula gives the rational rank immediately. Their local duality gives the corresponding orthogonality over the coefficient field. The integral reduction and the loss for a comparison between two different lattices are calculated below.\formalnote{14.11}{11}}

The potential crystalline periods span over a finite extension, since a basis uses only finitely many periods fixed by open subgroups. They give a de Rham basis there. After replacing that extension by a finite Galois extension, semilinear descent on its de Rham invariants gives a de Rham basis over \(K\). Bloch--Kato's dimension formula therefore applies at \(K\) itself. Put \(D=D_{\rm dR}(V)\), \(t_V=D/\operatorname{Fil}^0 D\). In the second exact sequence of \cite[Corollary~3.8.4]{arith-bloch-kato}, the middle map sends \((x,y)\) to \(((1-\varphi)x,x-y)\), with \(y\in\operatorname{Fil}^0D\). Quotienting the second target factor \(D\) by \(\operatorname{Fil}^0D\) and eliminating \(y\) gives
\[
 0\longrightarrow V^{G_K}\longrightarrow D_{\rm cris}(V)
 \xrightarrow{x\mapsto((1-\varphi)x,\ x\bmod\operatorname{Fil}^0)}
 D_{\rm cris}(V)\oplus t_V\longrightarrow H_f^1(K,V)\longrightarrow0 .
\]
It follows that \(\dim_{\mathbb Q_p} H_f^1(K,V)=\dim_{\mathbb Q_p}V^{G_K}+\dim_{\mathbb Q_p}t_V\). The map \(x\mapsto(g\mapsto gx-x)\) has kernel \(V^{G_K}\), so restoring the coboundaries replaces the first term by \(\dim_{\mathbb Q_p}V=d\). In the sign convention of this manuscript, \(\mathbb Q_p(1)\) has jump \(-1\). In a de Rham basis adapted to the filtration, the quotient by \(\operatorname{Fil}^0\) has one \(K\)-dimension for each negative jump in the list of the underlying \(\mathbb Q_p\)-representation. If \(d^-\) counts these jumps with multiplicity, then
\[
                \dim_{\mathbb Q_p}Z_f(V)=d+n d^- .                \tag{LC1}
\]\formalnote{14.12}{12}

For a class lattice the integral period definition gives continuous cocycles. One way to check continuity is to form the extension of the trivial rank-one lattice by \(L\) determined by the cocycle. Its period splitting preserves potential crystallinity. In an adapted Hodge--Tate basis the representation matrix near the identity is a product of a fixed change of basis, continuous cyclotomic powers and the Galois translate of the inverse change of basis. Its last column contains the cocycle. The cocycle rule propagates continuity at the identity to the whole group. Consequently \(Z_f(L)\) consists exactly of those elements of \(Z_f(V)\) with all their values in \(L\). Since \(G_K\) is compact, one power of \(p\) carries the values of any given element of \(Z_f(V)\) into \(L\). This proves \(Z_f(L)[1/p]=Z_f(V)\).\formalnote{14.13}{13}

To see finite generation without imposing any topology on a space of functions, choose finitely many elements \(g_j\in G_K\) so that evaluation on the \(g_j\) is injective on \(Z_f(V)\). Such elements exist because the common kernel of all evaluations is zero and each new evaluation can be chosen to reduce its dimension until it vanishes. Evaluation embeds \(Z_f(L)\) into the finite free \(O\)-module \(\prod_j L\). A submodule is finite free over this discrete valuation ring. Also a cocycle in \(Z_f(L)\) reducing to the zero \emph{function} modulo \(p^m\) can be divided pointwise by \(p^m\), and its primitive by the same scalar. Hence its kernel of reduction is \(p^mZ_f(L)\). For the image \(A_m(L)\) in the finite-valued cocycles, this proves
\[
                      |A_m(L)|=p^{m(d+nd^-)}.                   \tag{LC2}
\]
The dual jump paired with \(u\) is \(-u-1\). Exactly one is negative, including the cases \(u=0,-1\). Thus for \(M=L/p^mL\) the two finite images have product \(|A_m(L)|\,|A_m(T)|=|M|^{2+n}\). This argument works with the restricted \(\mathbb Z_p\)-lattices regardless of ramification of \(O\).\formalnote{14.14}{14}

\subsection{Finite duality with ramified coefficients}

Put \(F=\operatorname{Frac}O\). The trace pairing on \(F\) identifies \(\operatorname{Hom}_{\mathbb Z_p}(O,\mathbb Z_p)\) with the fractional ideal
\(\{a\in F:\operatorname{Tr}_{F/\mathbb Q_p}(aO)\subset\mathbb Z_p\}\).
Choose an \(O\)-generator \(a\) of this inverse different and write \(\lambda(z)=\operatorname{Tr}_{F/\mathbb Q_p}(az)\). The map \(t\mapsto\lambda\circ t\) identifies \(T\) with \(\operatorname{Hom}_{\mathbb Z_p}(L,\mathbb Z_p)(1)\), as is seen on any \(O\)-basis of \(L\). It is equivariant because Galois acts \(O\)-linearly on \(L\) and cyclotomically on the target. The generator accounts for the different even when the ordinary trace of \(O\) fails to give a perfect integral form. Reduction and the canonical identification \(\mathbb Z_p(1)/p^m=\mu_{p^m}\) now give
\[
 T/p^mT \ \simeq\ M^D=\operatorname{Hom}(M,\overline K^\times).
\]
Indeed every such map takes values in \(\mu_{p^m}\), and dualizing a finite free \(\mathbb Z_p\)-basis commutes with this reduction. Hilbert 90 and the Kummer sequence identify \(H^2(K,\mu_{p^m})\) with \(\operatorname{Br}(K)[p^m]\). Thus vanishing of a finite cup tested in the multiplicative coefficients is exactly vanishing under the local invariant pairing.\formalnote{14.15}{15}

Take \(m\ge1\); for \(m=0\) all finite groups below are trivial. The continuous modules \(M,M^D\) are finite over a characteristic-zero nonarchimedean local field. Tate duality gives a perfect finite pairing between their two \(H^1\)'s, and identifies \(H^2(K,M)\) with the Pontryagin dual of \(H^0(K,M^D)\). The Euler characteristic is \(|M|^{-n}\), since \(M\) is \(p\)-primary and \(|O_K/kO_K|=k^n\) for \(k\) a power of \(p\). For a subgroup \(J\subset Z^1(K,M^D)\) let \(J^\perp\) be the group of cocycles on \(M\) whose classes annihilate its image. Passing to classes has kernel \(J\cap B^1(K,M^D)\) in \(J\), and full kernel \(B^1(K,M)\) on the annihilator side. As \(|B^1(K,M)|=|M|/|M^{G_K}|\), the two theorems give the exact identity
\[
             |J^\perp|\,|J|
                = |M|^{2+n}
                   \frac{|J\cap B^1(K,M^D)|}{|B^1(K,M^D)|}.      \tag{LC3}
\]
This records the only possible loss in passing from orthogonal cohomology classes to groups of cocycles. Away from \(p\), the Euler characteristic for finite \(p\)-primary coefficients is \(1\). There the full \(Z^1\) has order \(|M|\,|(M^D)^{G_K}|\), which bounds the size of any subgroup of global restrictions used in the prime-to-\(p\) local term.
See \cite[I, Corollary 2.3 and Theorem 2.8]{arith-milne}.\formalnote{14.16}{16}

\subsection{Loss from variation of the lattice}

For the pair \(L,T\) itself, the two reduced crystalline groups are exact annihilators. In fact their rational crystalline classes are orthogonal by \cite[Proposition~3.8]{arith-bloch-kato}. The integral cup under the chosen evaluation is in \(H^2(K,\mathbb Z_p(1))=\mathbb Z_p\), which injects into \(H^2(K,\mathbb Q_p(1))\). Integrally this identification follows by taking the inverse limit of the local Brauer invariants for \(\mu_{p^j}\). The preceding \(H^1(K,\mu_{p^j})\) are finite, so the derived inverse limit in degree one vanishes. The integral cup is therefore zero, and so are its reductions. Every finite coboundary lifts to an integral coboundary by lifting the single defining vector. On setting \(J=A_m(T)\), the fraction in (LC3) is one. The product computed in (LC2) and the inclusion \(A_m(L)\subset J^\perp\) now imply equality.

In the Selmer argument the finite class in \(M\) can instead come from a crystalline cocycle of another lattice \(Y\), through an equivariant map \(Y\to M\). An \(O\)-linear lift \(Q:Y\to L\) exists by freeness. Its evaluation pairing with \(T\) has equivariance error divisible by \(p^m\); dividing that error by \(p^m\) gives an integral cocycle with rational primitive the evaluation form divided by \(p^m\). This is the extra cocycle in the chart argument. Define \(H\subset Z_f(T)\) by annihilation of all such allowed pairings as in \S\ref{sec:crystalline-loss}, and \(J\) by its reduction at this level. Both \(A_m(L)\) and the module \(N\) of allowed restrictions annihilate \(J\). Formula (LC3) gives
\[
 \frac{|N+A_m(L)|}{|A_m(L)|}
 \leq\frac{|J^\perp|}{|A_m(L)|}
 \leq \frac{|A_m(T)|}{|J|}.
\]
The final quotient is a quotient of \(Z_f(T)/H\), by reduction. A finite \(O\)-module \(X\) has \(\operatorname{len}_*(X)=\log_p|X|/[F:\mathbb Q_p]\), since its simple factors have cardinality \(p^{f_O}\) and \(O/pO\) has \(e_O\) such factors. This proves the normalized loss inequality for ramified as well as unramified coefficient rings. If \(p^s Z_f(T)\subset H\), the right side in normalized length is at most \(s\operatorname{rank}_O Z_f(T)\le s(1+n)\operatorname{rank}_O T\). The exponents and the sharper low-jump colength for the varying pairings are the integral estimates of \S\ref{sec:cup-comparisons}.\formalnote{14.17}{17}


\section{Integral cup products at a coefficient prime}
\label{sec:cup-comparisons}

We derive the four comparisons used in Section~\ref{sec:crystalline-loss}. Write \(p=\ell\). Two integral crystalline cocycles can pair to a nontrivial finite cup product even when their rational cup product vanishes. The argument constructs a cochain for the finite pairing. We first describe the integral estimates used to construct it, then count the powers of \(p\) that they require.

\subsection{From Hodge orders to a cochain}

Fix a place \(v\mid p\) and write \(\Gamma\) for its local Galois group, an open subgroup of \(G_{\mathbb Q_p}\). A Kummer chart supplies a coordinate \(z\) from compatible \(p\)-power roots of a local uniformizer. The specialization \(\theta:A_{\inf}\to\mathcal O_{\mathbb C_p}\) sends \(z\) to that uniformizer. Write \(E\) for the chart generator of \(\ker\theta\) in \(A_{\inf}\). Write \(e,f\) for the chart's ramification index and residue degree. Section~\ref{sec:kummer-local} constructs the chart used here. In the unramified chart \(e=1\) and \(E=z+p\). Put \(d=0\) if \(e=1\), and \(d=v_p(e)+1\) otherwise. The chart splits coefficients into components indexed by \(i\in\mathbb F_{p^f}^{\times}\). We call an index \(i\) a \emph{letter}. Write \(\zeta\) for the fixed coefficient root of unity of order dividing \(p^f-1\). An embedding \(\gamma\) contributes to component \(i\) when \(\gamma(\zeta)\) is the Teichmüller lift of \(i\). The allowed jumps at \(i\) are the union over all those embeddings, with \(0\) adjoined. Actual multiplicities at \(i\) are the sums over those embeddings. In this section write \(i[-h]\) for the letter \(h\) steps back under \(i\mapsto i^p\), so \(i[0]=i\). The comparison is made separately for each \(v\), using all letters of its chart.

A lattice \(Y\), a parameter lattice \(C\) with trivial action, and a linear map \(c\mapsto z_c\) to crystalline cocycles give the extension \(\widetilde Y=Y\oplus C\) with \(g(y,c)=(gy+z_c(g),c)\). We use two such extensions, belonging to \(Y\) and \(L\). An integral form on \(Y\otimes L\) that is equivariant modulo \(p^m\) gives a finite pairing \(B_0\) with values in \(\mathbb Z/p^m(1)\) after applying the chosen coefficient functional. Its equivariance error divided by \(p^m\) is required to be crystalline. This produces the third extension on \((Y\otimes L)^\vee(1)\). Its possible jumps at a letter are \(0\) and \(-u-u'-1\), with \(u,u'\) jumps from the two sides. Consequently opposite far windows cancel when both jumps are far.

For \(x\) in the localized Kisin lattice of an extension, an order \(o(h,u)\) means that evaluation of \(\varphi^{-h}x\) on each Hodge form of weight \(u\) at \(i[-h]\) lies in \(t^{o(h,u)}B_{\mathrm{dR}}^+\). A depth bound gives these orders for \(gx-x\) for every chart-group element \(g\). The standard lattice has orders \(-u\) at all levels. For the lattice transported by \(\varphi^l\) they are zero for \(h<l\) and \(-u\) for \(h\ge l\). Say that \(p^kE^s\) kills the cokernel at \(i\) when every constant in that component, multiplied by \(p^kE^s\), lifts to the standard lattice. Dividing its lift by \(E^s\) gives a point over \(p^kc\) with a possible pole of order \(s\) at level zero.

We use the following integral chart estimates. Increasing a required depth from \(y\) to \(y+1\) at level \(h\) costs at most
\[
                     v_p(|y|)+d+3h,\qquad v_p(0):=0
                     \quad\text{in this cost formula}.                 \tag{CP1}
\]
A step shared by several weight lines can be taken in one operation when bounding an exponent. For a length bound one can charge individual lines with multiplicity. On the unramified chart a line of multiplicity \(m_u\) has at most \(f m_u\) independent obstructions to a step. The precise statements use compatible intermediate orders on lower weights. The reason for the factor \(f\) is that among \(f m_u+1\) elements there is a primitive \(\mathbb Z_p\)-linear combination whose leading coefficient on that line can be removed modulo the other raised orders. Applied to the filtration on liftable constants, a step killed by \(p^k\) then has colength at most \(k f m_u\). Section~\ref{sec:kisin-height} proves the one-step estimates and specifies the interval in which the multiplicity argument applies.\formalnote{14.18}{18}

Here is the cochain calculation that explains how the costs combine. Work over \(W(\mathbb C_p^\flat)\) modulo \(p^m\) and call the image of \(W(\mathfrak m^\flat)\) in the twisted coefficient module the small part. Write \(\pi_Y,\pi_L\) for projection onto the first coordinates of the two extensions. Suppose \(x,w\) lift \(p^{k_1}c\) and \(x',w'\) lift \(p^{k_2}c'\). Set \(B(v,v')=p^{k_3}B_0(\pi_Yv,\pi_Lv')\) after scalar extension. Move one coordinate at a time along
\[
 (\varphi x,\varphi x')\,,\ (w,\varphi x')\,,\ (w,w')\,,\
 (x,w')\,,\ (x,x').
\]
A first-coordinate move \(v\) is safe if \(B(v,gx'-x')\) is small for the resting point \(x'\). A second-coordinate move \(v'\) is safe if \(B(gx-x,gv')\) is small for the resting point \(x\). All conditions are for every \(g\) in the chart group. Also require \(B(gx-x,g(hx'-x'))\) small for all \(g,h\).

Solve \((\varphi-1)y=B(w,\varphi x'-w')+B(x,w'-x')\) in the twisted Witt module, using surjectivity of \(\varphi-1\) there. For clarity, put \(B_1(v,\lambda)=p^{k_3}B_0(\pi_Yv,\lambda)\) for a vector \(\lambda\) of the second kernel lattice. On the group form
\[
 F(g)=B_1(gx,g\pi_Lx')-g\bigl(B(x,x')-y\bigr)
                  -B(x,gx'-x')-y .
\]
Since \(g\) commutes with \(\varphi\), expansion gives
\begin{align*}
(\varphi-1)F(g)
 &=B(gw-w,g(\varphi x'-w'))-B(\varphi x-w,g\varphi x'-\varphi x')\\
 &\quad -B(w-x,gw'-w')+B(gx-x,g(w'-x')).
\end{align*}
Each term is small by one of the four moves. Modulo the small part, Frobenius-fixed Witt vectors are constants. Thus \(F(g)\) has a unique constant residue \(q(g)\in\mathbb Z/p^m(1)\). Applying the extension action and bilinearity to \(gF(h)-F(gh)+F(g)\) gives
\[
 p^{k_1+k_2+k_3}B_0\bigl(z_c(g),g z_{c'}(h)\bigr)
                         =gq(h)-q(gh)+q(g).
\]
Indeed the unaccounted pairing is \(B(gx-x,g(hx'-x'))\), which is small, and a small constant is zero. The cocycles vanish modulo \(p^m\) near the identity and the Witt correction is small there, so \(q\) vanishes near the identity. Its coboundary identity gives continuity everywhere. Restriction to the chart subgroup detects the class since its index is prime to \(p\) and corestriction after restriction multiplies by that index. This proves the required coboundary in the local multiplicative coefficients.\formalnote{14.19}{19}

We also use the chart's criterion for safe moves. At a cut whose preceding cut is \(l\) steps back with shift \(s\), add the order of the moving component at weight \(u\) to the depth of the resting component at weight \(u'\). The bound for the change starting at that preceding cut is
\[
 \begin{cases}
  1 & h<l,\\
  s-u-u' & h=l,\\
  -u-u' & h>l .
 \end{cases}
\]
The first line gives strict positivity. The remaining lines give the standard orders after division by the relevant Frobenius translate of \(E^s\). Together with the annihilator of the third cokernel this places the paired difference in the small part. The endpoint move uses the analogous bound at the current cut. One can assemble the cochain globally over the letters by transporting the nearest preceding points by Frobenius and summing their components. At a letter between cuts the transport already agrees with the next point. At a cut an intermediate resting point permits the four moves above.\formalnote{14.20}{20}

\subsection{Sparse far windows}

Suppose the jumps at \(i\) are in \([-\eta,\eta]\) or in \(-\epsilon_i[\beta_i,\beta_i+V]\) on \(Y\), and in the low window or \(+\epsilon_i[\beta_i,\beta_i+V]\) on \(L\), with \(\epsilon_i=\pm1\). Let \(\beta_i=b\) at thick letters and \(G\) elsewhere. Assume \(G+V\le\eta\), \(\eta,J\ge1\), and that three successive visits to thick letters with both gaps in \([1,J]\) are impossible. This exclusion applies to the periodic sequence along each entire Frobenius orbit. Suppose either the chart is unramified with \(p\ge7,\ 10p^J\eta\le b\), or the signs are constant with \(8e\le p^{J+1}\). Choose integers \(\delta,X\) satisfying
\[
 2e(b+5\eta)\le p^{J+1}\delta,\qquad
 20(J+1+d)\eta+\lfloor\log_p(2(b+5\eta))\rfloor\le X.
\]
We show that \(10^3(\delta+(J+1)X)\) is a cup exponent.

Take \(W=5\eta,\ B=b+W,\ \kappa=12((J+1)X+\delta)+1\). Mark a cut at each letter followed by \(J\) thin letters. For any \(i\), first inspect \(i[-J-1]\). If this is not a cut there is a thick letter \(i[-r]\), \(1\le r\le J\). If it too fails to be a cut there is a thick letter within the next \(J\) steps. The preceding \(J\) letters of \(i[-r]\) must then be thin, so \(i[-r-J-1]\) is a cut. This puts a cut at distance \(1\le l\le2J+1\). For the nearest such cut, a thick \(i[-r]\) with \(0<r<l\) has \(r+J+1=l\). In fact the \(J\) letters after the cut are thin, giving \(r+J+1\le l\), and the preceding argument produces a cut at \(r+J+1\), giving the reverse inequality.

All three sides have jumps bounded above by \(\beta_i+W\). At a cut these are at most \(2W\) on the next \(J\) letters and at most \(2B\) everywhere. The weighted Kisin cokernel estimate at a zero-shift cut charges at most
\[
    2Jd(2W)+(2J+4)\frac{2W}{p-1}
                       +8\frac{2eB}{p^{J+1}}\le\kappa .
\]
To see the source of the last term, begin modulo \(p^{\kappa+1}\) with an initial monomial \(z^a\) in the lift ideal, where \((p-1)a\le e(2B+\kappa)\). At a backward Frobenius step with upper bound \(\omega=2W\), an exponent \(n\) becomes \(\lfloor(n+\omega)/p\rfloor\) and the scaling costs at most \(d\omega+v_p(\omega!)\). After \(J\) steps there is a monomial \(p^cz^{n}\) with
\[
      n\le a/p^J+\omega/(p-1),\qquad
      c\le J(d\omega+v_p(\omega!)).
\]
Removing \(z^n\) costs at most \(n+v_p(n!)\). Using \(v_p(q!)\le q/(p-1)\), \(p/(p-1)\le2\), and \(8e\le p^{J+1}\), the term involving \(\kappa\) is at most \(\kappa/2\). The remaining terms are at most half the displayed bound for \(\kappa\). Nakayama gives the integral lift. The ramification inequality is automatic in the unramified alternative.

In raising the depth, the possible steps at each of the first \(2J+2\) levels lie in two intervals of radius \(W\), centered at zero and at the signed far weight. For an interval of length \(D\), at most \(D/p^j+1\) integers are multiples of \(p^j\). Summing over \(j\le\lfloor\log_p(2B)\rfloor\) bounds its valuation sum. Including \(d+3h\) from (CP1), the two intervals at a level cost at most \(30X\). In particular
\[
            \kappa+4\delta+(2J+2)30X\le25\kappa .                    \tag{CP2}
\]
The resulting point above \(p^{25\kappa}c\) has depth at least \(-u+\eta\) through level \(l\) inclusive, and at least \(-u\) thereafter.

If all the significant far windows since the preceding cut have the same sign, the side whose far jumps are negative can rest while the other moves. Indeed its jumps satisfy \(u\le\eta\), so its depth \(\eta-u\) below \(l\) gives the nonnegative orders required of a stop. On the moving side \(u'\ge-\eta\) below \(l\). Adding its orders to the resting depth offsets the possible pole of order \(\eta\) in the move between the two cuts. This move uses the criterion at the current cut (transport distance zero), with required sum \(-u-u'\) in the zero-shift case. For differing signs the only possible significant letters are \(i\) and the pivot \(i[-t]\) with \(t+J+1=l\). In this case \(p\ge7\). Put \(s_0=\lfloor(b+10\eta)/(p-3)\rfloor+1\). The effective order brought to the pivot from the preceding cut is bounded by
\[
           s_0+\eta(1+p+\cdots+p^J)
           \le (b+10\eta)/4+1+2p^J\eta < b .
\]
This strict inequality keeps the transported low orders below the pivot's far window. Here is the construction at the pivot. Write \(l_0=J+1,\ A=2B,\ g=\eta\) and \(N=s_0+g\sum_{h<l_0}p^h\). All \(l_0\) upper jumps on the intervening side with negative far window at the pivot are at most \(g\). Start with a point \(y\) of pole at most \(s_0\) above \(c\), supplied by the unramified annihilator below. Multiply by
\[
        F(E)=E^{s_0}\prod_{h<l_0}(\varphi^h E)^g .
\]
Because \(\varphi^hE=(E-p)^{p^h}+p\), this is a monic polynomial in \(E\) of degree \(N\). The orders of \(F(E)y\) are at least \(g-u\) at levels below \(l_0\), and \(-u\) thereafter. The integral frame theorem, with \(A\le p^{l_0}\delta\), writes \(p^{4\delta}F(E)y\) as a linear combination over \(W(\mathbb F_{p^f})[[z]]\) of elements fixed to depth \(A\) at \(E\). Expand their scalar coefficients in powers of \(E=z+p\). In the resulting expansion the coefficient \(\xi_N\) is fixed to depth \(A\). At a line \(u<-N\) it has order at least \(A\): successively removing terms through degree \(N\) leaves strictly positive order there, and a positive order for a coefficient constant modulo \(E^A\) improves to \(A\). Comparison on the constant quotient, using the leading coefficient of \(F\), gives \(\pi_C\xi_N=p^{4\delta}c+w\) with \(w\) of order at least \(A\). The preceding cut supplies a lift \(v\) of \(p^\kappa w/E^A\) in the transported standard lattice. The point \(p^\kappa\xi_N-E^A v\) now lies exactly over \(p^{\kappa+4\delta}c\) and retains the required orders. Every negative far jump satisfies \(u<-N\).

Carry this point forward \(t\) times and raise its other depths below \(l\) to \(1\) using (CP2). At the pivot a moving order as low as \(1-2B\) can now pair safely. At the other levels below \(l\) the moving orders used with this stop are nonnegative. Thus this intermediate point permits the exchange of which side moves first. The construction at each cut uses \(25\kappa\) on each parameter and \(\kappa\) for the pairing. Their sum is \(51\kappa\le10^3(\delta+(J+1)X)\).\formalnote{14.21}{21}

\subsection{Two sizes in an arbitrary order}

Keep the two signed windows and assume the chart unramified, \(p\ge7,\ \eta\ge1,\ 0<G\le b,\ V+4\eta\le W\le K,\ b<Kp^2,\ 2G<Kp\). No spacing condition is needed. Put \(D=\lfloor b/p\rfloor\) and choose the shift \(s_i=\max(0,D-7K)\) if the successor of \(i\) is thick, and \(0\) otherwise. Mark the cuts \(s_i<\beta_i\). A thick letter is a cut since \(s_i\le D<b\). Every predecessor of a thin letter is a cut since its shift vanishes. Consequently the previous cut is one or two steps back.

Take \(J=1,\ X=3K+\lfloor\log_p(2(b+W))\rfloor,\ \delta=4K,\ s_0=D+12K\), and \(\kappa=12(2X+\delta)+1\). At a cut with thin successor the preceding zero-shift estimate applies with thin bound \(G+W\). At a cut with thick successor, \(s_0\le s_i+19K\le s_i+\kappa\), and every upper jump is at most \((p-3)s_0\). The unramified annihilator proved below therefore supplies \(p^\kappa E^{s_i}\). Here \(s_i+W\le pK\), so the widened level-zero intervals still cost at most \(30X\) each level by the multiple count used above. The shift is absent at positive levels.

The only mixed-sign case has distance two, with the middle letter thin and the endpoint thick. Indeed the middle letter is not a cut, so its successor must be thick and \(G+7K\le D\). The earlier cut has shift zero. Now
\[
             s_0+(G+W)(1+p)
                \le (p+2)D-(6p-6)K < b .
\]
For the last inequality use \(D<pK,\ pD\le b,\ p\ge7\). This is exactly the pivot inequality above with \(J=1\) and upper thin bound \(G+W\). The same stops and the same sum \(25+25+1\) apply. Hence \(10^4(K+\lfloor\log_p(2(b+W))\rfloor)\), which dominates \(51\kappa\), is a cup exponent.\formalnote{14.22}{22}

\subsection{Unramified primes with no nearby multiple of \(p\)}

First we prove the annihilator used in the two unramified cases. Suppose \(p\ge5\) and every upper jump is at most \((p-3)A\). For \(s+t=A\), we claim \(p^sE^t\) kills the cokernel. Work in the ideal of scalars lifting constants modulo \(p^{A+1}\). The initial Kisin height estimate supplies \(z^\nu\) in that ideal with \((p-1)\nu\le(p-2)A\). The chart's monomial descent from exponent \(j\) to \(j-1\) costs \(1+v_p(j)\). By Legendre's formula,
\[
 (p-2)(\nu+v_p(\nu!))\le(p-1)\nu\le(p-2)A.
\]
For \(0\le j<\nu\) the total cost down to \(j\) is
\(\nu-j+v_p(\nu!)-v_p(j!)\le A-j\). Above \(\nu\) the ideal is closed under multiplication by \(z\). Thus \(p^{A-j}z^j\) is in the ideal for \(0\le j\le A\). The binomial expansion of \(p^s(z+p)^t\) uses precisely such monomials. The case \(j=0\) supplies \(p^A\) in the ideal, so Nakayama's lemma removes the modulus. This argument works on each of the three sides.\formalnote{14.23}{23}

Now assume \(\eta\ge1\), let \(W\) denote the actual far-window width and put \(n=6(a+\eta+W)\). At every letter suppose
\[
 \begin{gathered}
    \beta_i+W+\eta\le(p-3)a,\qquad n<\beta_i,\\
    0\ne y\in\mathbb Z,\ 
    \min(|y|,|y-\beta_i|,|y+\beta_i|)\le n
                   \quad\Longrightarrow\quad p\nmid y .
 \end{gathered}                                                       \tag{CP3}
\]
All three side heights are at most \((p-3)a\): two far weights of opposite signs sum to an element of \([-W,W]\), and one far and one low have sum of absolute value at most \(\beta_i+W+\eta\). Two low weights give the required upper bound since \(2\eta<\beta_i\). The Tate shift subtracts one from the negative of the sum. We can therefore lift using \(E^a\) on each side.

If both sets of jumps omit \([-a,0)\), deepen the initial order from \(-u-a\) to \(a+\eta-u\) at level zero. All intermediate steps lie in the three intervals of (CP3), so (CP1) charges zero. The resulting lift is pole-free. In detail, a strictly positive initial order improves to the new depth, and an initially nonpositive order improves at least to the minimum of zero and that depth. For \(u<-a\) the first rule gives at least \(-u\). All other occurring weights have \(u\ge0\), where the second rule also gives at least \(-u\). Use every letter as a cut. According to its sign one side has all jumps \(\le\eta\), the other all jumps \(\ge-\eta\), so the first side gives the resting point for the Frobenius step. Both lifts and the third annihilator cost zero. If instead the entire right side is in \([-\eta,0]\), take \(t_0=a+\eta\). The left lift may have a pole of order \(t_0\) and depth \(a+t_0-u\) at level zero. The right side has a pole-free lift with that depth, since all its upper jumps are nonpositive. Frobenius transports its depth to level one, and free level-zero steps produce a right stop with depths \(a+t_0-u,\ t_0-u\) at levels zero and one. Where the left far window is negative the left side may rest, and elsewhere this right stop permits the move. This proves zero defect under either condition.\formalnote{14.24}{24}

For the general case let \(m_i(u)\) be the multiplicity of a right-hand jump, summed over embeddings above \(i\), and choose \(n_i,r_i\) bounding \(f\sum_{|u|\le\eta}m_i(u)\) and \(f\sum_u m_i(u)\), respectively. We construct a sublattice of right parameters which pairs trivially with every left parameter. At each \(i\) first require a pole-free lift of depth \(2a+\eta-u\) at level zero. The free steps just described give every constant such a depth with pole at most \(\eta\), since there is no jump in \([-a,-\eta)\). Multiplication by \(p^a\) permits a pole-free lift of that depth for every constant, by the annihilator argument with \(s=a,t=0\) and the free steps. Remove the pole one order at a time. For the step from \(t\) to \(t+1,\ -\eta\le t<0\), only the line with jump \(t\) gives an obstruction, killed by \(p^a\); the one-line estimate bounds its \(\mathbb Z_p\)-length by \(afm_i(t)\). Total charge is at most \(a n_i\).

If the right far window is positive that condition already suffices. If it is negative, take the least multiple \(T\) of \(p\) at least \(\beta_i+W+a+\eta\). At level zero require depth at least \(T\) as well. A step \(y\) on a negative far line satisfies \(y\ge2a+\eta-u\ge2a+\eta+\beta_i\). If it were divisible by \(p\) and \(y<T\), we would have \(y\le T-p<\beta_i+W+a+\eta\), which puts it in the forbidden interval around \(\beta_i\). Hence these lines cost zero. For a low line the interval starts at a nonnegative integer. Moreover
\[
 T\le\beta_i+W+a+\eta+p-1\le(p-1)(a+1).
\]
All its steps together cost at most \(v_p(T!)\le T/(p-1)\le a+1\). At level one the order increases from \(-u\) through \(W+1\) steps, and through a further \(2a\) steps if \(|u|\le\eta\). Each step lies within the indicated radius of \(0\) or \(\pm\beta_j\) for the \emph{preceding} letter \(j=i[-1]\), so (CP1) charges exactly \(3\). During this second improvement it suffices to retain depth \(\beta_i+W+a+\eta\) at level zero. The intermediate multiple \(T\) puts no further multiple of \(p\) in the open interval \((T,T+p)\), as required in the one-line estimate. Thus the total local charge is at most
\[
          (2a+1)n_i+3(W+1)r_{j}+6a n_{j}.                              \tag{CP4}
\]
This is the source of the different coefficients for low and far multiplicities. For the sublattice defined by these conditions the right side has pole-free lifts at every letter. The left side has lifts with the improved orders obtained from pole \(a\) and depth \(2a+\eta-u\). If the right far window is positive the left side can rest during its move. With negative right window the extra depths at levels zero and one allow the right side to rest. The construction of the right sublattice uses only \(L\), so this argument holds for all left lattices and finite pairings simultaneously.

The quotient by the intersection of all the letter conditions injects into the direct sum of their quotients. Since the predecessor permutes the letters, (CP4) sums to at most \(11\sum_i((a+1)n_i+(W+1)r_i)\). Take the right parameter lattice \(Z_f(L)\). If there are at most \(\mu\) low jumps at each embedding and \(d_i\) embeddings above \(i\), one may take \(n_i=fd_i\mu,\ r_i=fd_i\operatorname{rank}_O L\). The embedding count gives \(\sum_i d_i\le[O:\mathbb Z_p]\), writing the bracket for the rank of the coefficient ring. Divide the \(\mathbb Z_p\)-length by \([O:\mathbb Z_p]\) to get the normalization \(\operatorname{len}_*(O/pO)=1\). The defect therefore satisfies\formalnote{14.25}{25}
\[
        d_{\mathcal C}(L)\le 11f\bigl((a+1)\mu
                        +(W+1)\operatorname{rank}_O L\bigr).             \tag{CP5}
\]

\subsection{Tame level primes}

Suppose the map \(\Gamma\to G_{\mathbb Q_p}\) is an open embedding and surjective. Assume the permitted type on every inertia element includes \(X^e-1\), with \(e=\operatorname{lcm}_{1\le j\le r}(p^j-1)\), and a monic polynomial over \(\mathbb Z_p\) of degree at most \(N\). The second lattice uses the type at the inverse inertia element. Assume the two far windows start at \(b\), have width \(W\) and opposite signs independent of the embedding. Put \(\kappa=r^2(N+1)^2\) and require \(b+6(W+\eta)+\kappa+2<p\). There is a tame chart of ramification \(e\) where these type conditions imply crystallinity. On it a tame generator sends \(z\) to \(\xi z\) with \(\bar\xi\) of order \(e\). We have \(p\nmid e,\ d\le1\), and \(e\le\prod_{j=1}^r p^j\le p^{r^2}\) for \(r\ge1\). For \(r=0\) the empty least common multiple is \(1\).

For either side let \(q\) be its annihilating polynomial for the generator acting on crystalline periods. For the tensor dual one obtains a monic annihilator of degree at most \(N^2\). Indeed on the finite free algebra \(\mathbb Z_p[X,Y]/(q_1(X),q_2(Y))\), multiplication by \(XY\) has a characteristic polynomial of degree \(\deg q_1\deg q_2\). Cayley--Hamilton gives the tensor annihilator after using the inverse generator on the appropriate duals, and the Tate twist is trivial on the inertial period action here. On an extension by constant parameters use \(Q=q(X)(X-1)\).

Here is why the degree controls the loss. If the upper jumps are at most \(n_0=b+6(W+\eta)\), the initial monomial modulo \(p^{\kappa+1}\) satisfies \((p-1)\nu\le e(n_0+\kappa)\); hence \(\nu<e\). Descend from \(z^\nu\). At exponent \(j>0\) the action of \(Q\) gives a factor \(Q(\xi^j)\) for the descent by one. If its reduction is nonzero this factor is a unit and the descent is free. The residues \(\bar\xi^j,\ 1\le j\le\nu\), are distinct, so at most \(\deg q+1\) steps cost anything. At such a step the general descent costs \(1+v_p(j)\le r^2\), since \(j<e\le p^{r^2}\). All three cokernels are therefore killed by \(p^\kappa\), since \(\deg q+1\le(N+1)^2\) in all cases. When \(r=0\) there are no positive steps.

Use all letters as cuts with shift \(a=0\). The condition on \(p\) excludes every nonzero multiple in the intervals of radius \(n=6(W+\eta)\) about \(0,\pm b\), and the forbidden jump interval \([-a,0)\) is empty. Formula (CP1) at level zero charges at most \(d\) per step. The two signed unions of windows used for the two sides have cardinal at most \(8n+4\) each, so the two lifted parameters and the pairing cost at most \(3(\kappa+(8n+4)d)\) altogether. This is at most
\[
               3r^2(N+1)^2+150(\eta+W+1).
\]
This gives the last exponent used in Section~\ref{sec:crystalline-loss}.\formalnote{14.26}{26}

Section~\ref{sec:kisin-height} derives the initial height, the cost of the depth steps and the long-gap coefficient from the period matrix and the local chart. Section~\ref{sec:hodge-cup} proves the integral pairing criterion on Hodge orders and constructs the tame crystalline chart. Section~\ref{sec:sparse-weights} verifies the weight placement used here, including the exclusion of three thick letters in the sparse range.


\section{The integral pairing criterion and the tame chart}\label{sec:hodge-cup}

The cochain in \S\ref{sec:cup-comparisons} requires certain pairings of
Witt vectors to be small modulo \(p^m\). We prove the implication from
the Hodge orders to that integral assertion. Write \(A=A_{\inf}\),
\(\mathfrak m^\flat\) for the maximal ideal of \(\mathcal O_{\mathbb C_p}^\flat\),
\(\tau=[\epsilon^{1/p}]-1\), and \(\Delta=z\tau\). The small subgroup is
the image of \(W(\mathfrak m^\flat)\) in the Witt module for
\(\mathbb Z/p^m(1)\). Use the chart subgroup \(\Gamma'\), its letters
\(i\), and the notation \(i[-a]\) of \S\ref{sec:cup-comparisons}.

\subsection{An equivariant form that lifts the finite pairing}

Let \(Y,L\) be the two lattices, \(D=\operatorname{Hom}_O(Y\otimes_O L,O(1))\),
and \(\psi\in D\) the chosen integral form. The finite pairing is its
reduction after applying a fixed \(\mathbb Z_p\)-linear map
\(O\to\mathbb Z_p\). By hypothesis \(g\psi-\psi=p^m c(g)\) for the
specified crystalline cocycle \(c\) with values in \(D\).
On the extension \(\widetilde D=D\oplus O\) let
\(g(\lambda,b)=(g\lambda+b\,c(g),b)\). Then
\[
     H((\lambda,b),w)=p^m\lambda(w)+b\psi(w),
       \qquad w\in Y\otimes_O L ,
\]
is equivariant into \(O(1)\). Indeed \(p^m g\lambda+
b(p^m c(g)+\psi)=g(p^m\lambda+b\psi)\) in the Tate dual.
After extending to Witt vectors, reducing modulo \(p^m\) deletes the
first term. If the constant projection of \(\nu\in\widetilde D\) after
scalar extension is \(p^\kappa\) times the idempotent of letter \(i\),
then \(H(\nu,w\otimes y)\) reduces to \(p^\kappa B_0(w,y)\) for
\(w,y\) in that letter. This is an equality in the finite Witt module,
before taking cohomology. A crystalline cocycle splits over \(B_{\max}\),
so \(\widetilde D\) is crystalline on \(\Gamma'\) whenever \(D\) is.
Its possible jumps are \(0\) and \(-u-u'-1\) for jumps \(u,u'\) of
\(Y,L\) at the same letter.\formalnote{14.27}{27}

For an equivariant bilinear map between chart lattices, the following
rule holds at every level \(a\). Suppose its two inputs have lower
orders \(o_1(a,u),o_2(a,u')\) at all jumps of their letter \(i[-a]\).
The output has order at least \(b(a,v)\) at the jump \(v\) if
\[
      u+u'\le v\quad\Longrightarrow\quad
                 o_1(a,u)+o_2(a,u')\ge b(a,v) .                 \tag{HC1}
\]
To see the inequality in this direction, take an output form of type
\(v\). It sends the letter part of the ordinary lattice into
\(\operatorname{Fil}^{-v}B_{\rm dR}\). In a Galois-invariant de Rham basis
of the second factor the change-of-basis matrix is
\(P\operatorname{diag}(t^{h_j})\), with \(P\) invertible over
\(B_{\rm dR}^+\). The row dual to column \(j\) gives a form of type
\(h_j\). Contracting the output form against that column gives a first-factor
form of type \(v-h_j\), since every entry of the column is divisible by
\(t^{h_j}\) in the fractional filtration. Thus evaluation on the two
inputs is a sum of products of these forms. A form of type \(h\) can
be bounded using the actual jumps \(u\le h\). Descend from \(h\)
until a jump is met. If there is none, the type spaces do not
change and the value on an integral multiple of the input belongs to
every filtration step, hence is zero. Applying this in each factor gives
the pairs in (HC1). Multiplication of the two filtration ideals adds
their orders. The same expansion at positive Frobenius translates
preserves the integral tests defining the Kisin lattice. Clearing
products of Frobenius translates of \(E\) therefore proves the rule
also on its localization.\formalnote{14.28}{28}

Suppose \(p^\kappa E^s\) kills the cokernel for \(\widetilde D\) at
\(i[-l]\). Lift \(p^\kappa E^s\) times the constant \(1\) there and
divide the lift by \(E^s\). Its order at a jump \(v\) is at least
\(-v-s\) at level zero and \(-v\) at higher levels. Transport the
result by \(\varphi^l\) and call it \(\nu\). It lies over \(p^\kappa\)
times the idempotent at \(i\), with orders at least
\[
 b_\nu(a,v)=
 \begin{cases}
  0&a<l,\\
  -v-s&a=l,\\
  -v&a>l .
 \end{cases}
\]
These lower bounds decrease as \(v\) increases. For kernel vectors
\(w,y\) with orders \(o_1,o_2\), assume for all levels and allowed jumps
\[
       b_\nu(a,-u-u'-1)+o_1(a,u)+o_2(a,u')\ \ge 1 .             \tag{HC2}
\]
Apply (HC1) first to \(w\otimes y\), taking its order at \(v\)
to be \(1-b_\nu(a,-v-1)\). This is allowed because \(u+u'\le v\)
implies \(b_\nu(a,-v-1)\ge b_\nu(a,-u-u'-1)\).
Apply (HC1) again to \(H(\nu,w\otimes y)\). The only Tate jump is
\(-1\), and each pair of input jumps \(j,v\) with \(j+v\le-1\)
satisfies \(b_\nu(a,j)\ge b_\nu(a,-v-1)\). Thus its value \(\xi\)
has order at least one at every level on the Tate lattice. These are
exactly the standard orders there.\formalnote{14.29}{29}

There is now an integral conclusion. The period frame of
\S\ref{sec:period-matrix} places \(\tau^H\) times the standard Kisin
lattice inside \(A\) tensored with the original lattice when \(H\)
bounds the absolute jumps. Here \(H=1\) suffices, so
\(\tau\xi\in A\otimes_{\mathbb Z_p}O(1)\). The element \(z\) is the
Teichmüller lift of a nonunit in the tilt valuation ring, and hence
\(zA\subset W(\mathfrak m^\flat)\). Reduction through the fixed
coefficient map therefore yields
\[
                    \Delta\,p^\kappa B_0(w,y)
                    \ \in\ \text{the small subgroup}.
\]
If \(x'\) is a resting point whose differences have depths \(o_2\),
put \(y=(gx'-x')/\Delta\). Its constant projection is zero. Divisibility
by \(\Delta\) in the localized frame and the nonzero residues of
\(\varphi^{-a}\Delta\) imply it retains orders \(o_2\).
This proves the first-coordinate move. For a second-coordinate move
take \(w=(gx-x)/\Delta\) and use \(gy'\) for the moving vector.
\(\Gamma'\)-invariance of each order gives the same argument.
Substituting the displayed \(b_\nu\) in (HC2) gives the required sum
\(1\) below level \(l\), \(s-u-u'\) at \(l\), and \(-u-u'\) above.
In particular the bound by \(1\) has an integral role even when
the pole at the preceding cut is zero.\formalnote{14.30}{30}

There are two details in passing to the cochain. For the product of
two cocycle differences the second move may itself be \(hx'-x'\).
For points transported \(d\) steps from a cut with shift \(s\), the
depth at levels \(a<d\) is at least \(1\), since
\(\varphi^d\Delta\) vanishes at those levels. At level \(d\) each
standard depth at a jump \(u\) is at least \(s-u\). At higher levels
it is at least \(-u\). Thus the two depths meet the three bounds just
proved, and \(p^\kappa B_0(gx-x,g(hx'-x'))\) is small as required.
Finally \(B_0\) vanishes between different letters. Acting by the
chosen coefficient root on the left or right gives the same pairing
and multiplies it by the respective Teichmüller lifts. Distinct
lifts differ by a unit. Hence the four small pairing conditions and
the cocycle-difference condition sum over letters without cross terms.
The sum of the three exponents charged to the lifts therefore suffices
in the cochain of \S\ref{sec:cup-comparisons}.\formalnote{14.31}{31}

\subsection{The subgroup on which the tame type becomes crystalline}

At the tame level prime the local group maps onto \(G_{\mathbb Q_p}\)
by an open embedding. Identify the two groups. Write
\(e=\operatorname{lcm}_{1\le j\le r}(p^j-1)\). The numerical hypothesis
of \S\ref{sec:cup-comparisons} gives \(r<p\). Take \(f=r!\), with
\(0!=1\). Then \(p\nmid f\), \(p\nmid e\), and \(e\mid p^f-1\),
because \(j\mid f\) for every \(1\le j\le r\). Let \(K_0/\mathbb Q_p\)
be unramified of degree \(f\), choose \(\varpi^e=p\), and put
\(K=K_0(\varpi)\). The polynomial is Eisenstein over \(K_0\) for
\(e>1\), so the degree of \(K/\mathbb Q_p\) is \(ef\).
All \(e\)-th roots of unity already belong to \(K_0\). Thus \(K\)
is the compositum of the unramified extension with the splitting
field of \(X^e-p\), and \(\Gamma'=G_K\) is normal with prime-to-\(p\)
index. It is exactly the fixer of \(\varpi\) and \(W(\mathbb F_{p^f})\).
These assertions also hold when \(e=1\).\formalnote{14.32}{32}

Choose successive \(p\)-power roots of \(\varpi\), defining a tilt
element \(\varpi^\flat\) with untilt \(\varpi\) and \(z=[\varpi^\flat]\).
At each finite level inertia acts transitively on the roots of
\(X^{ep^n}-p\), since the polynomial is Eisenstein over every finite
unramified extension. Compactness of inertia gives an element realizing
any compatible choice on the whole tower, by the finite intersection
property of its closed cosets. In particular one element \(s\) acts
on \(z\) by \(\xi z\) for the Teichmüller lift \(\xi\) of a primitive
\(e\)-th root in \(\mathbb F_{p^f}\). Since inertia fixes the constants,
\(s^n z=\xi^n z\). The same argument with the compatible cyclotomic
factor \(\epsilon^j\) produces \(g_j z=[\epsilon]^j z\). Its untilt
is unchanged, so \(g_j\in\Gamma'\). These are the translations
required by the Kummer frame calculation.\formalnote{14.33}{33}

Let \(V\) be one of the potentially crystalline representations allowed
by the type. On the spanning set of potentially crystalline periods
the action of inertia induces a homomorphism \(\beta\) on the period
space by linear extension. It has open kernel, since finitely many
spanning periods are fixed near the identity. The type condition
applies to every inertia element \(g\) and gives \(\beta(g)^e=1\).
The image of wild inertia in this finite group is a \(p\)-group and
is therefore trivial. The tame quotient of inertia is procyclic of
orders prime to \(p\). Its cyclic quotient of order \(e\) is detected by
\(g\mapsto g\varpi/\varpi\) over the maximal unramified extension. Indeed the tame inertia groups at finite level are cyclic, and over the maximal unramified extension the prime-to-\(p\) roots of a uniformizer exhaust the tame extensions \cite[Lemma~15.114.5 and Lemma~15.116.7]{arith-stacks}.
Consequently \(\beta\) kills every inertia element fixing \(\varpi\).
All inertia elements of \(\Gamma'\) do so. For the right lattice the type
is supplied at \(g^{-1}\), and inversion permutes inertia.\formalnote{14.34}{34}

Over \(K\), choose a finite Galois extension \(M/K\) on which \(V\) is
crystalline. The full period space over the maximal unramified
subfield \(M_0\) has dimension \(\dim_{\mathbb Q_p}V\). Its semilinear
action of \(\operatorname{Gal}(M/K)\) now has trivial inertia. The
quotient acts on \(M_0\) through
\(\operatorname{Gal}(M_0/K_0)\), the Galois group of the residue
extension. Galois descent gives the same dimension over \(K_0\) for
the fixed periods. These are precisely the crystalline periods over
\(K\), proving crystallinity on \(\Gamma'\). In the project's
\(B_{\max}\) model the same descent uses unramified algebraic
constants, whose Witt lifts belong to \(B_{\max}\); the comparison
with \(B_{\rm cris}\) was described in \S\ref{sec:local-lengths}.
The Hodge jumps at the specified coefficient embeddings persist upon
restriction, as only an open subgroup has been removed.\formalnote{14.35}{35}

We now pass from an annihilating polynomial on periods to an estimate on the integral frame. Let \(q\) be the monic polynomial annihilating the tame operator on the periods of the kernel lattice. Put \(Q=(X-1)q\) and write
\(Q=\sum_j Q_jX^j\). Since the parameter quotient is trivial, \(Q\)
kills the period operator of the extension. For a crystalline form
\(\alpha\), set \(\alpha_j(x)=s^{-j}\alpha(s^jx)\); these are again
crystalline forms by normality of \(\Gamma'\).
Testing on the spanning potentially fixed periods gives
\(\sum_jQ_j\alpha_j=0\). For a point \(x\) of the chart Kisin module,
\[
        \alpha\bigl(Q(s)x\bigr)
        =\sum_jQ_j(s^j-1)\alpha_j(x).
 \]
Clear the common powers of \(E\) in these chart values. Each difference
on the right is divisible by \(z\), because \(s^j\) fixes the constants
and acts on the parameter by \(z\mapsto\xi^jz\). The frame's
divisibility test, using that \(z\mid E^n a\) implies \(z\mid a\)
in the chart calculation, places \(Q(s)x/z\) in the \(A_{\inf}\)-Kisin
lattice. Applied to a lift with projected monomial \(p^b z^j\),
this gives \(Q(\xi^j)p^b z^{j-1}\) in the lift ideal, for \(j\ge1\).
If the lift is known modulo \(p^h\), split its correction into its
constant term and a multiple of \(z\). The former disappears because
\(Q(1)=0\), and division by \(z\) preserves the modulus on the latter.
This proves the descent used in the root count of
\S\ref{sec:cup-comparisons}.\formalnote{14.36}{36}

The local implication uses the period frame, the prescribed bounds on
the orders and, at a tame prime, potential crystallinity with the
stated inertia annihilators over \(\mathbb Q_p\). Verifying those
annihilators and the jump windows for the global family is the
arithmetic input to this local calculation.


\section{From the global weights and tame labels to the local bounds}
\label{sec:sparse-weights}

Fix a place \(v\mid\ell\) of \(E\). The integral estimates of
Section~\ref{sec:cup-comparisons} count jumps along Frobenius orbits.
Their application in Section~\ref{sec:crystalline-loss} requires the
positions of the jumps, their multiplicities on the Tate dual, and the
distances between the letters at which a large weight can occur.
We derive these data for the actual summands. Throughout this section
\(h=r^8,\ \Delta=2h,\ t=5r,\ k=r+1-m,\ Q_0=kt\). The two weights are
\(b\) and \(G_0=a/x_1\), as in \S\ref{sec:data}.

\subsection{The jump set and its low part}

Write \(\tau\) for the infinite place of \(\sigma:E\hookrightarrow\mathbb C\),
\(\epsilon_\sigma=1\) on the chosen CM type and \(-1\) on its conjugate.
The place weight \(\kappa_\tau\) is \(b\) on \(\Phi_b\) and \(G_0=a/x_1\) elsewhere. For index
\(i\) the center of the archimedean weight is
\(w_i=\kappa_\tau+i\Delta+h\). In a rank-\(r\) row \(1\le i\le m\)
the jump of index \(0\le j<r\), counted with multiplicity, is
\[
 u_j=\lfloor t/2\rfloor-\epsilon_\sigma(\lambda_j+w_i),
 \qquad \lambda_j=\lambda_j(\tau).
 \tag{SW1}
\]
If \(j+1<k\), the \(j\)-th entry \(\lambda_j=\lambda_j(\tau)\) of the target parameter satisfies \(|\lambda_j|\le k/2\), so
\(|\epsilon_\sigma u_j+w_i|\le t\). For \(j+1\ge k\), put
\(b'=j+2-k\). Then \(1\le b'\le m\) and the target parameter gives
\(\lambda_j+w_i=(i-b'+1/2)\Delta\). Its absolute value lies between
\(\Delta/2\) and \(m\Delta\). Since \(\lfloor t/2\rfloor\le\Delta/4\),
these last \(m\) jumps satisfy
\(\Delta/4\le |u_j|\le2m\Delta\). The data satisfy
\(2m\Delta\le h^2\) and \(2h^2+t<w_i\) even at a place of weight \(G_0\).
Exactly \(m\) jumps have absolute value at most \(2h^2\). At conjugate
embeddings the two jumps of index \(j\) sum to \(2\lfloor t/2\rfloor\),
so at most one is negative.\formalnote{14.37}{37}

For a rank-one main pair the jump is
\(u=z-1-\epsilon_\sigma c_\tau\), with \(m<z\le r\) and
\(w_i+r-t\le c_\tau\le w_i\). Thus \(|\epsilon_\sigma u+w_i|\le t\)
and \(|u|>2h^2\). Here the conjugate sum is \(2(z-1)\).
For either main summand of rank \(\rho\), the number of jumps of
absolute value at most \(2h^2\) is at most
\(\lfloor \rho m/k\rfloor\), since \(k\le r\).
The lifted type also permits \(t'\) Tate summands of jumps
\(0,1,\ldots,t'-1\), where \(t'=t\) for (SW1) and
\(t'=2z-1\le t\) for the rank-one pair.\formalnote{14.38}{38}

Set \(\delta=i\Delta+h-2t,\ A_\sigma=\kappa_\tau+\delta,\ W=4t,\)
and \(\eta=h^2+1\). Then \(w_i=A_\sigma+2t\). A far jump satisfies
\(-\epsilon_\sigma u\in[A_\sigma+t,A_\sigma+3t]\), while every other
allowed jump has \(|u|\le h^2\). Under Tate duality \(u\mapsto-u-1\).
Consequently the left and right far windows are respectively
\(-\epsilon_\sigma[A_\sigma,A_\sigma+W]\) and
\(\epsilon_\sigma[A_\sigma,A_\sigma+W]\), with a common low window
\([-\eta,\eta]\). On the left any negative jump comes from the main
summand, so it is at most \(-h/2\). On the dual of the main summand,
a negative jump \(-u-1\) comes from a nonnegative main jump, hence is
also at most \(-h/2\). Moreover \(|-u-1|\le\eta\) implies
\(|u|\le h^2+2\le2h^2\), so the same bound
\(\lfloor\rho m/k\rfloor\) counts the dual low jumps with multiplicity.
This uses the jumps that actually occur in the given dual. For the
dual of a single Tate summand \(c\) the exact list is instead
\(\{-c-1\}\subset[-\eta,0]\), which is the other zero-defect
criterion of Section~\ref{sec:crystalline-loss}.\formalnote{14.39}{39}

\subsection{Exclusion of three thick letters on a short Frobenius arc}

Let \(e:E\hookrightarrow\overline{\mathbb Q}_\ell\) realize \(v\),
fix the coefficient identification
\(\iota:\mathbb C\simeq\overline{\mathbb Q}_\ell\), and put
\(\sigma_g=\iota^{-1}g^{-1}e\) for \(g\in G_{\mathbb Q_\ell}\).
The chart for \(G_{E_v}\) has residue degree \(f_v\) and ramification
index \(e_v\), with \(e_v f_v=[E_v:\mathbb Q_\ell]\le[E:\mathbb Q]\le h\).
It provides \(F\in G_{E_v}\) acting on the algebraic residue field
by \(x\mapsto x^{\ell^{f_v}}\). Enlarge the coefficient ring to contain a
primitive root \(\zeta\) of order \(\ell^{f_v}-1\). A coefficient embedding
given by \(g\) has letter \(g\zeta\). Mark a letter thick if it has
\emph{some} representative \(g\) for which the place of \(\sigma_g\)
belongs to \(\Phi_b\).\formalnote{14.40}{40}

Suppose \(\ell\le Q_0\) and write \(J=\lfloor\log_\ell(h^3)\rfloor\).
If three thick occurrences have successive distances \(1\le A,B\le J\),
choose thick representatives \(g_0,g_1,g_2\) with respective letters
\(z,z^{\ell^A},z^{\ell^{A+B}}\). Put \(n=v_\ell(q_1)\),
\(q'_1=q_1/\ell^n\), and \(\xi=e(\zeta_{q_1})^{\ell^n}\), of order \(q'_1\).
The element \(F\) fixes \(\xi\) and acts on prime-to-\(\ell\) roots of unity
through the \(\ell^{f_v}\)-power map, so \(q'_1\mid\ell^{f_v}-1\).
Since the letter represents a primitive root of the latter order, an
identity \(g'\zeta=(g\zeta)^u\) forces \(g'g^{-1}\) to raise every
root of that order or a divisor to the \(u\)-th power. Apply this to
\(\xi\) with \((g',g)=(g_2,g_1),(g_2,g_0)\).

Let \(v_j\in(\mathbb Z/q_1)^\times\) be the label of the infinite place
of \(\sigma_{g_j}\), and put \(\epsilon_j=\epsilon_{\sigma_{g_j}}\).
The place label uses the representative in the CM type.
Replacing an embedding by its conjugate inverts its value on roots of
unity. Comparing the two power maps on \(\xi\) therefore gives
\[
 v_1/v_2\equiv\epsilon_1\epsilon_2\,\ell^B,\qquad
 v_0/v_2\equiv\epsilon_0\epsilon_2\,\ell^{A+B}\pmod{q'_1}.
 \tag{SW2}
\]
In the notation of Section~\ref{sec:data}, \(2^s=m^{22}\) and the
forbidden arc length \(l\) is the least power of three above \(27s\).
The parameter bound \(h^3\le(2^s)^{12}\) gives \(J\le12s\).
Thus \(0<B<A+B\le24s<l\). With denominator \(v_2\), (SW2) is exactly
a triple excluded by the construction of \(\Phi_b\) in that section.
This also rules out a short triple that wraps around a Frobenius orbit,
since the exclusion allows the place labels to coincide.\formalnote{14.41}{41}

There are two ways of reading signs on this chart. When
\(\ell\nmid q_1d\), the field \(E_v\) is unramified and lies in the
unramified extension of degree \(f_v\). Equality of letters forces the
relative automorphism to fix this field, hence the two embeddings
\(\sigma_g\) agree. Both thickness and sign are then constant on each
letter. Every \(\ell>Q_0\) occurs in this case, since the prime factors of
\(q_1d\) divide the fixed integer \(qd\le Q_0\). Moreover every prime
less than seven divides \(q\).

When \(\ell\mid q_1d\), all thick representatives at \(v\) have the same
sign. For \(\ell\mid q_1\), opposite signs would conjugate
\(\beta=(1+\sqrt{-d})/2\) to \(1-\beta\) locally. Their \(\ell\)-adic
absolute values would be equal, and their product \((d+1)/4\) is
divisible by \(\ell\) by the choice of \(q,d\). Both would have absolute
value \(<1\), although their sum is 1. For \(\ell=d\), the conductor
\(q_1\) is prime to \(d\), so every local Galois action on
\(\zeta_q\) is a \(d^a\)-power. The ratio of two thick labels is
\(1\bmod q\), since both are \(-1\bmod q\). Opposite signs would give
\(d^a\equiv-1\bmod q\), impossible since \(d\equiv q/2-1\bmod q\) and
\(d^2\equiv1\bmod q\). We may therefore assign this one sign to every
letter (choose either sign if there are no thick representatives).
Thin representatives of a marked letter need no such sign constraint.
Indeed in this small-prime range one widens the low interval to
\([-2G_0,2G_0]\), where \(G_0=a/x_1\) is the place weight outside \(\Phi_b\). This window contains the entire thin far interval because
\(G_0+\delta+W\le2G_0\). We thus have either an unramified chart with
\(\ell\ge7\) or a chart with a single sign for the far jumps that
survive the widening.\formalnote{14.42}{42}

To check the scale of the cut, take \(K\) as in
Section~\ref{sec:crystalline-loss}, where \(2(b+\delta)<h^2K\) for
\(\ell\le Q_0\). We have
\(\ell^J\le h^3<\ell^{J+1}\) and \(99G_0h^3\le b\). In particular
\(10\ell^J(2G_0)\le b+\delta\) in the unramified chart, while
\(8e_v\le8h\le h^3<\ell^{J+1}\) in the constant-sign case. In both
cases the depth needed for the widened low interval satisfies
\[
 e_v\,2(b+\delta+10G_0)
 \le 4h(b+\delta)<2Kh^3\le 2K\ell^{J+1}.
\]
These inequalities and the exclusion (SW2) give the sparse comparison
with exponent \(10^4 K(J+1)\). This explains the logarithm in the
small-prime local loss.\formalnote{14.43}{43}

\subsection{The two polynomials on inertia at a split level prime}

Finally suppose \(\ell=p_0\) or \(p_1\). These primes are \(1\bmod q_1d\),
so \(E_v=\mathbb Q_\ell\) and all \(\sigma_g\) coincide. It remains to
identify the two annihilators required in Section~\ref{sec:hodge-cup}.
At a chosen place \(w\mid\ell\) let \(c_w(X)\) be the label of degree
\(\rho_0\le r\). Normalization twists the attached representation by the
avatar of the inverse of the character of the descent pair, with the
prescribed weight. On the label this multiplies its roots by the residue
field unit naming the inverse of the tame character at \(w\) under the
chosen faithful multiplicative character. Write \(\widetilde c_w\)
for the resulting polynomial. The label actually imposed on the
normalized lift is \(\widetilde c_w(X)(X-1)^{t'}\). The trivial roots
are the extra blocks of the lift. This identifies the normalization and
the type supplied by the eigensystem used in Section~\ref{sec:congruence-detail}.\formalnote{14.44}{44}

Fix an inertial element \(s_0\). Its multiset \(M_{s_0}\) of possible
eigenvalues is obtained by applying the chosen faithful multiplicative
character, raised to the tame exponent of \(s_0\), to each label root,
and then applying \(\iota\). Its cardinality is \(\rho_0+t'\le r+t\).
A nonzero root of a degree-\(\rho_0\) polynomial over
\(\mathbb F_\ell\) lies in \(\mathbb F_{\ell^j}^\times\) for some
\(1\le j\le\rho_0\). Thus every \(y\in M_{s_0}\) satisfies
\(y^{\ell^j-1}=1\) for some \(1\le j\le r\), also for the extra roots
\(1\). The labelled local condition proved in \S\ref{sec:tame-labels} says that the product of the
\(X-y\) kills the operator on potentially crystalline periods. At the
conjugate place, conjugate self-duality gives the same assertion with
\(s_0\) replaced by its conjugate inverse. Since all inertial elements
were allowed, this covers both places.\formalnote{14.45}{45}

The inertial period operator has finite order. Indeed it acts trivially
on a sufficiently small open inertia subgroup, which contains a positive
power of \(s_0\). Its minimal polynomial therefore has simple roots
in characteristic zero. Every polynomial vanishing on \(M_{s_0}\)
annihilates the operator. Two such polynomials are
\[
 X^{e_*}-1,\qquad
 \prod_{y\in M_{s_0}}\operatorname{minpol}_{\mathbb Q_\ell}(y;X),
 \qquad e_*=\mathop{\rm lcm}_{1\le j\le r}(\ell^j-1).
 \tag{SW3}
\]
A root of unity of order dividing \(\ell^j-1\) lifts a residue unit
from \(\mathbb F_{\ell^j}\) uniquely to the degree-\(j\) unramified
extension. Each minimal polynomial in (SW3) is therefore monic over
\(\mathbb Z_\ell\) of degree at most \(j\le r\). The second annihilator
has degree at most \((r+t)r=N_0\). Both polynomials depend only on the
type and the inertia element and hence work for every member of the
local class. Equivariant inclusion of a lattice after scalar extension
preserves these period relations. The dual uses the inverse inertia
element. This supplies the tame type and the constant window center
used in the last case of Section~\ref{sec:crystalline-loss}.\formalnote{14.46}{46}


\section{Existence of the local Kummer chart}\label{sec:kummer-local}
\subsection{A Kummer chart for the actual local group}
Here we construct the Kummer coordinate \(z\) and the chart generator \(E\) used in the preceding local estimates. The same construction will serve for the varying representations in \S\ref{sec:galois-limits}. Fix a homomorphism \(j:\Gamma\to G_{\mathbb Q_p}\) with open image \(G_K\). The fixed field \(K\) is finite over \(\mathbb Q_p\). Let \(e,f\) be its ramification index and residue degree, \(\Lambda=W(\mathbb F_{p^f})\) the integers of its maximal unramified subfield, and \(\pi\) a uniformizer. Put \(v=-p/\pi^e\in\mathcal O_K^\times\). In the unramified case take \(\pi=-p,\ v=1\). Choose compatible \(p\)-power roots of \(\pi\), write the resulting tilt element as \(\pi^\flat\), and put \(z=[\pi^\flat]\). The Teichmüller digits of the residue field express each \(a\in\mathcal O_K\) as a convergent series in \(\pi\) over \(\Lambda\): choose a digit lifting its residue, subtract, divide by \(\pi\), and repeat. In particular there is \(V\in\Lambda[[X]]^\times\) evaluating to \(v\), since its constant digit is a unit. We take \(V=1\) if \(e=1\).

Embed \(\Lambda\) in \(A_{\inf}\) by the unramified Witt vectors. Evaluation at \(z\) makes sense in \(A_{\inf}\): in Witt coordinates the ideal \([u]^m W(\mathcal O_C^\flat)\) on coordinate \(i\) imposes divisibility by \(u^{mp^i}\); the perfect valuation ring \(\mathcal O_C^\flat\) is complete for powers of the nonzero nonunit \(u=\pi^\flat\). Completeness on all coordinates gives the evaluation homomorphism and \(\theta(z)=\pi\). Injectivity is proved in \S\ref{sec:integral-coordinates}. The element \(E=p+z^e V(z)\) is the chart generator of \(\ker\theta\) in \(A_{\inf}\). In fact it maps to zero under \(\theta\). The standard generator \(\xi_0=[p^\flat]-p\), for compatible roots of \(p\), has reduction with tilt valuation \(v(p)\), where the tilt valuation is transported by sharp. The reduction of \(E\) has exactly that valuation because \(e\,v(\pi)=v(p)\) and \(V(z)\) is a unit. Writing \(E=\xi_0 s\) therefore gives \(s\bmod p\) a unit. The \(p\)-adically complete Witt ring makes \(s\) a unit. This also proves that \(E\) is associated to the cyclotomic generator of \(\ker\theta\) \cite[Lemma~3.10, Remark~3.11 and Example~3.16]{legacy-bms}. The common stabilizer of \(\theta(z)\) and \(\Lambda\) is exactly \(G_K\), since the series expansion gives \(\mathcal O_K=\Lambda[[\pi]]\). Thus \(ef=[G_{\mathbb Q_p}:j(\Gamma)]\), and the residue Frobenius \(x\mapsto x^{p^f}\) has a lift in the image because the corresponding automorphism of the maximal unramified extension of \(K\) extends to its algebraic closure.\formalnote{14.47}{47}

One also needs elements carrying \(z\) to \([\epsilon]^a z\). For fixed \(a\in\mathbb N\) and \(m\), the polynomial \(T^{p^m}-\pi\) is Eisenstein over every finite unramified extension of \(K\) and is therefore irreducible over their union. The inertia of \(K\) acts transitively on its roots. The elements sending the chosen root \(\pi_m\) to \(\zeta_{p^m}^a\pi_m\) form a nonempty closed subset of inertia. The subsets decrease with \(m\). Compactness of inertia gives an element in their intersection. Its action on Teichmüller lifts has precisely the required value. These choices give a single chart for the fixed local group, even when the ranks of the represented spaces subsequently vary. For the separate tame inertia subgroup in the integral cup bound, \S\ref{sec:hodge-cup} constructs the extension of degree \(ef\) and computes the tame action on the same root coordinates.\formalnote{14.48}{48}



\section{Integral coordinates on a Kummer chart}
\label{sec:integral-coordinates}

The period-matrix argument needs to take a coordinate of a period while
retaining its divisibility. The denominator estimate also needs a single
integral approximation that works at every nonnegative Frobenius translate.
We prove the two assertions for the rings used in
Sections~\ref{sec:period-matrix}, \ref{sec:kisin-height}
and~\ref{sec:galois-limits}.

\subsection{Extension of the chart coordinates}

Write \(A=W(\mathcal O_{\mathbb C_p}^{\flat})\),
\(D=A_{\max}=\widehat{A[X]/(pX-\xi)}_p\), and \(B^+=D[1/p]\),
with the cyclotomic generator \(\xi=([\epsilon]-1)/([\epsilon^{1/p}]-1)\) of \(\ker\theta\) in \(A\). A chart gives
\(\Lambda=W(\mathbb F_{p^f})\), \(S=\Lambda[[z]]\subset A\) and
\(E=z^e u(z)+p\), where \(u(z)\in S^\times\) and \(E\) is associated
to \(\xi\) in \(A\). Here \(z\) is a Teichmüller lift of a nonzero
nonunit of the tilt. Evaluation of a series is injective. Indeed, divide
a putative nonzero kernel element by its lowest power of \(z\), then
reduce in the Witt vectors of the residue field of the tilt. This kills
\(z\) and embeds \(\Lambda\), so the remaining nonzero constant
coefficient cannot vanish.\formalnote{14.49}{49}

Set \(I=(p,z)\subset S\). The sequence \(p,z\) is regular on \(A\)
because \(p\) is regular and the image of \(z\) is nonzero in the domain \(A/p\).
Choose lifts \(b_i\), including \(b_0=1\), of a residue basis of \(A/IA\)
over \(S/I\). The lifted finite combinations are dense by successive
reduction modulo \(I\). For completeness, the coefficient bound for this
construction is that a finite combination in \(I^{2n}A\) has all its
coefficients in \(I^n\). Residue independence first detects coefficients
modulo \((p,z)\). Successive cancellation of \(z\) modulo \(p\) detects
them modulo \((p,z^n)\), and cancellation of \(p\) modulo \(z^n\)
detects them modulo \((p^n,z^n)\). The latter regularity follows from
the former and regularity of \(p\) on \(A\). Now
\(I^{2n}\subset(p^n,z^n)\subset I^n\). Completeness of \(S\) extends
each finite coefficient map to \(\lambda_i:A\to S\).
The identity coefficient satisfies \(\lambda_0(s)=s\).\formalnote{14.50}{50}

The equality of ideals
\[
                         A\cap p^nD=(p,\xi)^n A                    \tag{IC1}
\]
explains why the coordinates extend integrally. The right side maps into
\(p^nD\) by \(pX=\xi\). Conversely suppose \(a\in A\cap p^{n+1}D\).
Applying \(\theta\) and lifting the integral quotient gives
\(a=p^{n+1}c+\xi w\) with \(c,w\in A\).
Thus \(Xw\in p^nD\). Since \(D/pD=(A/(p,\xi))[X]\),
multiplication by \(X\) is injective modulo \(p\), and induction using
\(p\)-torsion-freeness gives \(w\in p^nD\). Induction on \(n\) proves
(IC1). Every \(S\)-linear \(f:A\to A\) preserves the right side since
\((p,\xi)=(p,E)\) and \(p,E\in S\). On \(D\cap A[1/p]\) one may
therefore set \(\widehat f(p^{-d}a)=p^{-d}f(a)\). This rule keeps \(D\)
and \(p^nD\). Its domain contains the dense ring \(A[X]\) evaluated at
\(X=\xi/p\). Completion and inversion of \(p\) give a unique extension
to \(B^+\) preserving \(D\).\formalnote{14.51}{51}

Let \(S_B\) consist of multipliers commuting with all \(S\)-linear maps
\(B^+\to B^+\) preserving \(D\). Each such bounded map is continuous
for the neighborhoods \(p^nD\), and is automatically \(S[1/p]\)-linear.
Each extended coordinate \(\widehat\lambda_i\) has image in \(S_B\). To see this,
approximate any element of \(B^+\) by elements of \(A[1/p]\). Their
coordinates lie in \(S[1/p]\), and the identity expressing commutation
with a bounded map passes to the limit after multiplication by a fixed
element of \(B^+\). Thus \(\mathrm{pr}=\widehat\lambda_0\) maps into
\(S_B\). Conversely \(\mathrm{pr}(y)=y\,\mathrm{pr}(1)=y\) for
\(y\in S_B\).

Thus \(S_B\), the fixed locus of the continuous projection \(\mathrm{pr}\),
is closed. For \(y\in S_B\), choose \(a_j\in A[1/p]\) tending to \(y\).
Then \(\mathrm{pr}(a_j)\in S[1/p]\) tends to \(y\). Frobenius preserves
\(S[1/p]\subset S_B\), so continuity gives \(\varphi(y)\in S_B\).
Moreover \(\varphi(E)\in S\) is associated to \(p\) in \(D\). Its
inverse in \(B^+\) therefore commutes with every bounded map as
well.\formalnote{14.52}{52}

For \(y\in B^+\) and any precision \(N\), every sufficiently large
finite set \(F\) of basis indices satisfies
\[
              y-\sum_{i\in F}b_i\widehat\lambda_i(y)\in p^N D.
                                                                    \tag{IC2}
\]
Approximate \(y\) modulo \(p^ND\) by \(p^{-d}a\) with \(a\in A\).
Choose \(F\) such that the expansion of \(a\) is accurate modulo
\(I^{e(N+d)}A\). Since \(z^e\in(p,E)\), this ideal lies in
\((p,E)^{N+d}A\subset p^{N+d}D\). Subtracting a finite coordinate sum
preserves \(D\), proving (IC2). A union of finitely many choices of
\(F\) works for all entries of a period matrix.\formalnote{14.53}{53}

\subsection{A simultaneous approximation for all nonnegative Frobenius powers}

If \(\nu\in A\) satisfies \(\xi=pb+\nu g\) in \(A\), then
\(p^mD\subset A+\nu^mD\) for every \(m\ge0\). Indeed every element
of \(D\) belongs to \(A+(X-b)D\). Divide a polynomial approximation
modulo \(p\) by \(X-b\), repeat on its \(p\)-multiple remainder, and sum
the constant remainders in \(A\) and the quotients in \(D\), using
\(p\)-adic completeness. Multiplication by \(p\) replaces \(X-b\)
by \(\nu g\), so iteration gives the formula. For \(\nu\in S\) it yields
\(p^m\mathrm{pr}(D)\subset S+\nu^m S_B\). Taking \(\nu=E\) gives the
integral approximation used in the radius estimate.\formalnote{14.54}{54}

Put \(\tau=[\epsilon^{1/p}]-1\), \(\mu=[\epsilon]-1=\xi\tau\). The
previous formula applies successively to \(\tau\mid \xi-p\) and
\(\xi\). It gives \(p^{2m}D\subset A+\mu^mD\). In particular, for
every \(y\in B^+\) and \(m,N\ge0\) there exist \(d\ge0,\ a\in A\)
and \(w\in D\) with
\[
                 p^d y=a+p^{d+N}\mu^m w .                         \tag{IC3}
\]
Indeed first approximate \(y\) by \(p^{-d}a_0\) modulo
\(p^{N+2m}D\), then apply the double division to the remaining
\(p^{2m}D\) term. The factor \(\mu\) is killed by \(\theta\varphi^j\)
for every \(j\ge0\), so this single approximation controls all these
tests. Since \(B^+\) is dense in its \(\xi\)-adic completion, (IC3)
also shows \(A[1/p]\) surjects onto each \(B_{\rm dR}^+/\xi^m\).
Only finitely many coefficients
of the fixed local extension need this approximation in the trace
projector of Section~\ref{sec:galois-limits}, permitting a common
denominator there.\formalnote{14.55}{55}

Here is the simultaneous divisibility test used to clear the poles of
the period matrix. In \(A\), let \(J=\bigcap_{j\ge0}\ker(\theta\varphi^j)\)
and use bars for reduction modulo \(p\). We claim \(J=(\mu)\).
Certainly \(\mu\in J\). If \(a\in J\), then \(a=\xi b\), and
\(\varphi b\in J\) since \(\theta\varphi^{j+1}\xi=p\).
If \(\bar\mu\) did not divide \(\bar a\), the valuation ring
\(\mathcal O_{\mathbb C_p}^{\flat}\) would give
\(\bar\mu=\bar a c\) with \(c\) a nonunit. Cancel \(\bar\xi\) to get
\(\bar\tau=\bar b c\) and take \(p\)-th powers to get
\(\bar\mu=\overline{\varphi b}c^p\). Iteration gives
\(c^{p^j}\mid\bar\mu\) for every \(j\), impossible by the rank-one
valuation and \(\bar\mu\ne0\). Thus \(\bar\mu\mid\bar a\).
Lift a quotient and subtract its multiple of \(\mu\).
The remainder is \(pa_1\) with \(a_1\in J\) because each target of
\(\theta\varphi^j\) is \(p\)-torsion-free. Repetition and completeness
give \(a\in\mu A\). For \(y\in B^+\) killed by all tests, (IC3) with
\(m=1,N=0\) gives \(a\in J\), hence \(y\in\mu B^+\).
Conversely every such multiple passes the tests.\formalnote{14.56}{56}

\subsection{Reading a fixed tilt component at a later chart level}

There is one use of the \emph{inverse} Frobenius on \(A\). Write
\(\theta_n=\theta\varphi^{-n}\) and \(E_n=\varphi^nE\).
Then \(\ker\theta_n=E_nA\), and \(\theta_n(a)\bmod p\) is component
\(n\) of \(\bar a\). If \(\theta_n(\alpha)\mid\theta_n(a)\) with
\(\alpha\in S\), surjectivity of \(\theta_n\) supplies \(c,d\in A\)
such that \(a=\alpha c+E_n d\). Taking any \(S\)-coordinate gives
\(\lambda_i(a)\equiv\alpha\lambda_i(c)\bmod E_n S\).\formalnote{14.57}{57}

Let \(H\) be the closed subgroup fixing \(S\) pointwise and let
\(b\in\mathcal O_{\mathbb C_p}^{\flat}\) be fixed by \(H\).
For fixed \(j\), the integral element \(t=\theta_j([b])\) is \(H\)-fixed.
Ax's fixed-field theorem, applied to
\(\overline{\mathbb Q}_p^H\), supplies an \(H\)-fixed algebraic \(x\)
with \(|x-t|<|p|\) \cite[main theorem, p.~417]{legacy-ax}. To locate it in the chart, observe that
\(\theta_n(S)\) is finite over \(\mathbb Z_p\) as a module. Indeed
Weierstrass division by \(E_n\), whose reduction has order \(ep^n\),
leaves a polynomial of degree \(<ep^n\) over \(\Lambda\).
These image rings increase with \(n\) because
\(\theta_{n+1}(\varphi s)=\theta_n(s)\). A Galois element fixing every
one of them fixes \(S\), because the maps \(\theta_n\) distinguish elements of
\(A\), first by their components modulo \(p\), then by successive
division by \(p\). Hence the field generated by the union over
\(\mathbb Q_p\) is \(\overline{\mathbb Q}_p^H\), by Galois
correspondence. Its elements are algebraic, so adjoining the union as
a \(\mathbb Q_p\)-algebra already gives a field.\formalnote{14.58}{58}

Clear a denominator to write \(qx=\theta_n(s)\), \(q\in\mathbb Z_p\)
nonzero, for every sufficiently large \(n\) with \(s=s_n\in S\).
As \(x\) is integral, apply the coordinate divisibility statement to
\(\alpha=q,\ a=s,\ i=0\). The identity coordinate fixes \(s\), so
\(\theta_n(s)=q\theta_n(s')\) for some \(s'\in S\). Cancellation gives
\(\theta_n(s')=x\equiv t\bmod p\).
Thus component \(j\) of \(b\) is component \(n\) of \(\bar s'\) for
every sufficiently large \(n\). A maximum threshold works for a finite
list, as required by the determinant argument.\formalnote{14.59}{59}

The constructions here start from a chart for the prescribed local
group and the basic Witt-vector description of \(\ker\theta\).
Section~\ref{sec:kummer-local} constructs the chart of a specified local field, and Section~\ref{sec:hodge-cup} constructs the chart adapted to the tame type. The subsequent passage to invariant
spans in Section~\ref{sec:galois-limits} uses its own \(p\)-adic basis
of \(D\) over a complete coefficient ring with maximal ideal \((p)\).
Lifting a residue basis there needs just \(p\)-torsion-freeness.
Finite coordinate supports at successive precisions can be made
increasing by taking partial unions. The Galois action on these spans,
the fixed-ratio theorem for \(B_{\max}\), and the choice of de Rham
columns are treated in Sections~\ref{sec:period-matrix} and~\ref{sec:derham-columns}.\formalnote{14.60}{60}


\section{A fixed integral change of columns for de Rham periods}\label{sec:derham-columns}
The matrix used in \S\ref{sec:period-matrix} has columns invariant under the whole chart group. Its change of columns can be chosen over the ring of integers of that group's fixed field. This matters for the projector in \S\ref{sec:galois-limits}, where the same coefficient ring must work at every rank. We give a proof by taking the filtration on the space of constant column combinations.

Write \(B=B_{\rm dR}\), \(B^+=B_{\rm dR}^+\), and \(\theta:B^+\to\mathbb C_p\) for the residue map. Thus \(B^+\) is a discrete valuation ring with uniformizer \(t\), \(g(t)=\chi(g)t\), and the canonical embedding \(\overline{\mathbb Q}_p\to B^+\) reduces to the given embedding in \(\mathbb C_p\). Let \(j:\Gamma\to G_{\mathbb Q_p}\) have open image, and write \(g(x)\) for \(j(g)(x)\). Take \(U\in{\rm GL}_d(B)\) and matrices \(N(g)\) with entries in \(\overline{\mathbb Q}_p\). Suppose
\[
                         N(g)\,g(U)=U.                                      \tag{DC1}
\]
There are two versions. If (DC1) holds for all \(g\in\Gamma\), put \(K=\overline{\mathbb Q}_p^{j(\Gamma)}\). If it holds near the identity, put \(K=\overline{\mathbb Q}_p\). Here \emph{near the identity} means for every \(g\) whose image lies in some fixed identity neighborhood in \(G_{\mathbb Q_p}\). In either case there exist \(Y\in {\rm SL}_d(\mathcal O_K)\), \(W\in{\rm GL}_d(B^+)\) and integers \(k_i\) such that
\[
                  UY=W\,\operatorname{diag}(t^{k_i}).                      \tag{DC2}
\]
The entries of \(Y\) are fixed in the same sense as (DC1). In the second case a single open subgroup fixes all entries since there are finitely many. More generally the source allows (DC1) along a filter \(\mathcal A\) on \(\Gamma\) each of whose sets contains the inverse image of an identity neighborhood. In the proof take \(K\) to be the field of algebraic constants fixed eventually along \(\mathcal A\). Finite intersections show that this is a field and that any finite collection of its elements is simultaneously fixed eventually. The matrices \(N(g)\) arise by applying the chosen coefficient embedding to the representation matrices in the jump calculation.\formalnote{14.61}{61}

For \(r\in\mathbb Z\) form the descending flag of \(K\)-spaces
\[
              D^r=\{a\in K^d:Ua\in t^r(B^+)^d\}.
\]
Every nonzero algebraic constant has \(t\)-valuation zero. A lower bound on the orders of the entries of \(U\) thus gives \(D^r=K^d\) for all sufficiently small \(r\). Similarly \(U^{-1}t^r(B^+)^d\subset t(B^+)^d\) for sufficiently large \(r\), giving \(D^r=0\). Choose a basis \(y_1,\ldots,y_d\) of \(\mathcal O_K^d\) adapted over \(K\) to this finite flag. To see why it can be integral even when \(K=\overline{\mathbb Q}_p\), take a nonzero vector in the smallest nonzero flag space. Divide it by a coordinate of largest \(p\)-adic absolute value, so its entries are integral and its chosen coordinate is one. Projection along this vector onto the hyperplane where that coordinate vanishes preserves the integral lattice and each flag space. Induct on this coordinate hyperplane. The resulting basis matrix has determinant in \(\mathcal O_K^\times\). Multiplying one column by the inverse determinant gives \(Y\in{\rm SL}_d(\mathcal O_K)\) with the same flag. For \(d=0\) take the empty identity matrix.

Give \(y_i\) its largest filtration index \(k_i\) and set \(W_i=t^{-k_i}Uy_i\), \(P_i=\theta(W_i)\). These columns give the equality in (DC2). It remains to prove \(W\) invertible over \(B^+\). For each \(r\), the map
\[
          D^r/D^{r+1}\longrightarrow\mathbb C_p^d,\qquad
                  a\longmapsto\theta(t^{-r}Ua)
\]
is injective and \(K\)-linear because \(\ker\theta=tB^+\). Thus the \(P_i\) with \(k_i=r\) are \(K\)-linearly independent. To obtain independence of all columns over \(\mathbb C_p\), reduce (DC1) applied to \(Uy_i=t^{k_i}W_i\). Take the common set on which (DC1) holds and all the \(y_i\) are fixed. On it we have
\[
                         N(g)\,g(P_i)=\chi(g)^{-k_i}P_i.                    \tag{DC3}
\]
Suppose there is a nonzero relation \(\sum_i c_iP_i=0\) of smallest support over \(\mathbb C_p\), normalized by \(c_h=1\). Conjugating and applying \(N(g)\) gives another relation with coefficients \(\chi(g)^{-k_i}g(c_i)\). Subtract \(\chi(g)^{-k_h}\) times the first relation. The \(h\)-th coordinate disappears, hence every coefficient disappears by minimality. This proves
\[
                             g(c_i)=\chi(g)^{k_i-k_h}c_i
\]
on this common set. Its image contains an open subgroup of \(G_{\mathbb Q_p}\) because \(j(\Gamma)\) is open. The vanishing of the invariants of nonzero Tate twists there \cite[Theorem~2.2.7]{legacy-brinon-conrad} forces \(c_i=0\) whenever \(k_i\ne k_h\). For \(k_i=k_h\), the coefficient \(c_i\) is fixed on the common set. The fixed-field part of the same theorem first makes \(c_i\) algebraic, by restriction to the open subgroup. Fixedness on the entire common set then puts \(c_i\) in \(K\) by its definition. This contradicts independence in \(D^{k_h}/D^{k_h+1}\). Consequently \(P=(P_i)\) is invertible, so \(\det W\) is a unit in \(B^+\). This proves the frame assertion, including fixedness under the full group whenever (DC1) holds there. The scalar fixed-field and twisted-invariant theorems are stated in \S\ref{sec:tate-constants}.
\phantomsection\label{route:24}\leanroute{Lean rescales columns by powers of \(t\), then uses a relation among their residues to lower the valuation of the determinant whenever it is positive. Here an integral basis adapted to \(D^\bullet\) fixes all the exponents at once. The minimal-relation argument above proves that its rescaled columns have independent residues.}\formalnote{14.62}{62}

For the application to jumps, (DC3) gives \(P^{-1}N(g)=\operatorname{diag}(\chi(g)^{-k_i})g(P^{-1})\). Each inverse row is a nonzero Hodge--Tate form of weight \(k_i\) for the chosen embedding. The rows of each weight count its multiplicity by the coefficient comparison of \S\ref{sec:jump-symbol}. In \S\ref{sec:period-matrix} apply (DC2) to each \(U=\varphi^jM\) over \(B\). Frobenius on \(B_{\max}\) commutes with Galois and fixes the \(\mathbb Z_p\)-matrix entries of the representation, so each application uses the full chart group. If its jumps are in \([-a,b]\), the formula \(U=W\operatorname{diag}(t^{k_i})Y^{-1}\) puts all entries in \(t^{-a}B^+\). The inverse formula puts all entries of \(U^{-1}\) in \(t^{-b}B^+\). Its use in \S\ref{sec:galois-limits} also gives \(Y^{-1}=\operatorname{adj}(Y)\) over \(\mathcal O_K\), uniformly in the number of columns. The passage from the simultaneous \(t\)-adic bounds to bounds over \(A_{\max}\) uses the separate period-ring divisibility and approximation arguments of \S\ref{sec:period-matrix}.\formalnote{14.63}{63}


\section{From a crystalline period basis to the integral Kisin frame}
\label{sec:period-matrix}

Fix one local representation and one chart as in Section~\ref{sec:kisin-height}. Let \(T\) denote its underlying free \(\mathbb Z_p\)-lattice, \(d\) its rank, and \(\Gamma'\) the group after restriction to the chart. Crystallinity over this group gives a basis of invariant columns in \(B_{\max}\otimes_{\mathbb Z_p}T\). We will construct a free lattice over \(S=W(\mathbb F_{p^f})[[z]]\) inside \(W(\mathbb C_p^\flat)\otimes_{\mathbb Z_p}T\). If \(H\) bounds the absolute Hodge--Tate jumps of \(T\), multiplying by \(\tau^H\) in either direction compares its \(A_{\inf}\)-span with \(A_{\inf}\otimes T\). This is the starting exponent used in the height calculation. Kisin's theorem associates a finite free Breuil--Kisin module to a Galois-stable \(\mathbb Z_p\)-lattice in a crystalline representation \cite[Theorem~4.4]{legacy-bms}. Here we need the explicit comparison of the two \(A_{\inf}\)-lattices by \(\tau^H\), which we prove below.

\subsection{The chart and the invariant matrix}

Write \(A=A_{\inf}\), \(B^+=A_{\max}[1/p]\), \(B_{\max}=B^+[1/\mu]\) in the period-ring model used here. Put
\[
 \mu=[\epsilon]-1,\qquad \tau=[\epsilon^{1/p}]-1,\qquad
 \varphi(\tau)=\mu=\xi\tau=v E\tau ,
\]
where \(v\in A^\times\), \(E\in S\) generates \(\ker\theta\) in \(A\), and \(z\) is the chosen Teichmüller parameter. The chart must satisfy
\[
 \{g\in G_{\mathbb Q_p}:g|_{W(\mathbb F_{p^f})}=1,\quad
                       g(\theta z)=\theta z\}
      \ \subseteq\ {\rm image}(\Gamma'\longrightarrow G_{\mathbb Q_p}).
\]
This inclusion will allow a lift of \(g\) to act on \(T\). Being able to fix \(z\) as an element of \(A\) is a stronger requirement and enters a different part of the proof below.\formalnote{14.64}{64}

Choose a basis of \(T\), and let \(R(\gamma)\) be its representation matrix. The invariant columns give \(M\in{\rm GL}_d(B_{\max})\) satisfying \(R(\gamma)\gamma(M)=M\). Suppose all jumps at the identity embedding of the underlying \(\mathbb Z_p\)-representation lie in \([-a,b]\), \(a,b\ge0\). For every \(j\ge0\), Frobenius commutes with the action, so \(\varphi^jM\) also has invariant columns. Apply the de Rham frame theorem to these columns. There is a fixed algebraic change of columns \(Y_j\), an invertible \(W_j\) over \(B_{\rm dR}^+\), and integers \(k_{ji}\) such that, in \(B_{\rm dR}\),
\[
       (\varphi^j M)Y_j=W_j\,{\rm diag}(t^{k_{ji}}).
\]
The rows of \(\theta(W_j)^{-1}\) give Hodge--Tate forms of those weights, hence \(-a\le k_{ji}\le b\). The entries of \(\varphi^jM\) have Hodge orders at least \(-a\), and inversion gives order at least \(-b\) for those of \(\varphi^j(M^{-1})\).\formalnote{14.65}{65}

Here is why orders at every translate yield a matrix in \(B^+\). Suppose the entries are reduced to a single \(x\in B_{\max}\) with \(\varphi^j x\in{\rm Fil}^{-a}\) for every \(j\). Choose \(m\) so \(c=\mu^m(E\tau)^a x\in B^+\). For \(m>0\), each \(\varphi^j c\) has order at least \(m\), since every \(\varphi^j\mu\) has positive order and \(E\tau\) is associated to \(\mu\). The period-ring divisibility criterion for the maps \(\theta\varphi^j\), \(j\ge0\), gives \(\mu\mid c\) in \(B^+\). Remove one factor and repeat. Applying this to \(M,M^{-1}\) produces\formalnote{14.66}{66}
\[
       U=(E\tau)^aM,\qquad U'=(E\tau)^bM^{-1}
           \quad\text{over }B^+,\qquad
       UU'=(E\tau)^h I,\qquad h=a+b .
\]

Let \(S_B\subset B^+\) consist of scalars whose multiplication commutes with every \(S\)-linear endomorphism preserving \(A_{\max}\). Both \(p\) and \(\varphi E\) have inverses in \(S_B\). The entries of \(X=M^{-1}\varphi M\) and \(X^{-1}\) are Galois fixed, since the same \(R(\gamma)\) occurs on each set of columns. The fixed-ratio theorem of \S\ref{sec:cohomological-frobenius} for the chart group \(G_K\) puts each entry in \(W(\mathbb F_{p^f})[1/p]\), hence in \(S_B\). Thus, with \(\eta=v^a\) and \(F=(\varphi E)^a X\),
\[
                      \varphi U=\eta U F,\qquad F\in{\rm GL}_d(S_B).
\]
For \(g\) fixing the constants and \(\theta z\), choose a lift \(\gamma\in\Gamma'\). The unit \(r=(E\tau)/g(E\tau)\) gives \(r^aR(\gamma)\,g(U)=U\). This is the integral Galois transformation used below.\formalnote{14.67}{67}

\subsection{Retractions and preservation of the Hodge conditions}

Take the topological \(S\)-basis \((b_i)\) of \(A\) from
\S\ref{sec:integral-coordinates}, with \(b_0=1\) and coordinate
\(\lambda_0:A\to S\) at \(b_0\). By the extension proved there, every
\(S\)-linear \(f:A\to A\) extends to a map \(\widehat f:B^+\to B^+\)
preserving \(A_{\max}\). By the definition of \(S_B\), each such extension is
\(S_B\)-linear.\formalnote{14.68}{68}

We need the following stability. For \(y\in(B^+)^d\) and \(n\ge0\),
\[
        Uy\in E^n(B^+)^d\quad\Longrightarrow\quad
        U\widehat f(y)\in E^n(B^+)^d .
\]
First take finitely many algebraic constants \(c_i\) fixed by \(\Gamma'\). The chart condition and Galois descent give a representative in \(S_B\) for each residue of \(c_i\) at \(\xi\). These representatives lift to every order: for the separable polynomial \(P\) of \(c_i\), choose \(D\) with \(DP'\equiv1\bmod P\). The polynomial \(X-DP\) fixes \(c_i\) and has zero derivative there. Applying it squares the order of approximation. Replace \(c_i\) by the resulting \(S_B\)-approximations modulo \(\xi^m\). The map \(\widehat f\) preserves \(E^m B^+\), and divisibility by this power is tested in \(B_{\rm dR}^+\). It follows for each \(m\) that divisibility by \(\xi^m\) of \(\sum c_i y_i\) carries over to \(\sum c_i\widehat f(y_i)\). In the de Rham valuation ring this says the first sum divides the second, including when it vanishes.\formalnote{14.69}{69}

Apply the divisibility row by row to \(Y_0^{-1}y\). There are \(s_i\in B_{\rm dR}^+\) with
\[
   Y_0^{-1}\widehat f(y)={\rm diag}(s_i)Y_0^{-1}y,\qquad
   U\widehat f(y)= W_0{\rm diag}(s_i)W_0^{-1}Uy .
\]
The diagonal of \(t\)-powers commutes with \({\rm diag}(s_i)\), and the scalar \((E\tau)^a\) commutes with both. The second equality proves stability, since the conjugating matrices are integral in the Hodge valuation and \(E^n\)-divisibility descends to \(B^+\).\formalnote{14.70}{70}

\subsection{The solution lattice and its denominators}

Define the \(S\)-module
\[
              P=\{x\in A^d: Uy=E^h x\text{ for some }y\in S_B^d\}.
\]
The solution \(y\) is unique, since \(UU'\) is a nonzero scalar matrix and \(B^+\) is a domain. A relation over \(A\) among such \(x\)'s is equivalent to the same relation among their solutions. Applying each extended basis coordinate to the latter gives a relation with coefficients in \(S\), since the entries of the solutions lie in \(S_B\). An \(S\)-independent family in \(P\) is consequently \(A\)-independent in \(A^d\) and has at most \(d\) members.

Moreover \(z\) is regular on \(P/pP\). If \(zx=py\) for \(x,y\in P\), then \(x=pw\) in \(A^d\), since \(p\nmid z\). Dividing the solution for \(x\) by \(p\) in \(S_B^d\) puts \(w\) in \(P\). Now \(P\) is \((p,z)\)-adically separated as a submodule of \(A^d\) and has no \(p\)-torsion. The adic basis theorem gives a topological \(S\)-basis. Its independence and the rank bound make it finite, and completeness of \(S\) and separation of \(P\) allow one to sum approximating coefficients to obtain an ordinary basis. Pad to \(d\) columns for the moment, and collect the columns and their solutions as \(Q,Z\). Then\formalnote{14.71}{71}
\[
           P=QS^d,\qquad UZ=E^h Q,\qquad
           Q\in{\rm Mat}_d(A),\quad Z\in{\rm Mat}_d(S_B).
\]

To compare with solutions \(y\) over all of \(B^+\), a single denominator must work for every coordinate. Given \(Uy=E^h x\), \(x\in A^d\), denote by \(w\) the vector with \(U\widehat f(y)=E^h w\). Choose \(k_U,k_{U'}\) so that \(p^{k_U}U,p^{k_{U'}}U'\) have entries in \(A_{\max}\). For \(e_0\mid\mu\), division by \(e_0^h\) inside \(B^+\) costs at most \(2h\) powers of \(p\) in \(A_{\max}\) when the numerator is integral and the quotient exists. Apply this first to \(U'(Uy)=(E\tau)^h y\), then to \(U\widehat f(y)=E^h w\), using that \(\widehat f\) keeps \(A_{\max}\). Thus every \(p^m w_i\) lies in \(A_{\max}\) for the common choice\formalnote{14.72}{72}
\[
                     m=k_U+k_{U'}+4h .
\]

One more descent puts these coordinates in \(A\). Set \(x'=\varphi^{-1}(\eta E^h x)\) in \(A^d\) and \(f'=\varphi^{-1}f\varphi\). The conjugate is \(S\)-linear because \(\varphi S\subseteq S\). Write \(y'=U'x'/\tau^h\). To see \(y'\in(B^+)^d\), apply \(\varphi\) to \(U(U'x')=(E\tau)^h x'\) and cancel \(U,\eta\) using \(\varphi U=\eta UF\). This gives \(F\varphi(U'x')=\varphi(E\tau)^h y\). Multiplying by \(F^{-1}\) and using descent of divisibility by \(\varphi(\tau)^h\) proves the assertion. Thus \(Uy'=E^h x'\). Define \(w'\) by \(U\widehat f'(y')=E^h w'\). Applying \(\varphi\), commuting the extended conjugate with \(\varphi\), and using \(S_B\)-linearity with respect to \(F,\varphi E\) gives\formalnote{14.73}{73}
\[
                          \varphi w'=\eta E^h w .
\]

The set \(V\) of all \(p^m w_i\) therefore lies in \(A_{\max}\), and each \(s\in V\) has \(s'\in V\) with \(\xi^h s=u\varphi s'\), where \(u=(\xi/E)^h\eta^{-1}\in A\). The contraction lemma puts \(p^{h+1}s\) in \(A\). Its extra \(1\) can be seen in the approximation
\[
                 p^{h+1}s=c_j+\tau^{h+1}p^j d_j,\qquad
                          c_j\in A,\quad d_j\in A_{\max}.
\]
The initial approximation is division in \(A_{\max}\) using \(\xi\in(p,\tau)A\). Apply Frobenius to the approximation for \(s'\), multiply by \(u\), and cancel \(\xi^h\). Since \(\varphi\tau=\xi\tau\), the error retains one factor \(\xi=p(\xi/p)\) in \(A_{\max}\); this increases \(j\) by one. The cancellation in the integral term uses \(\xi^h B^+\cap A=\xi^h A\). The \(c_j\) are Cauchy for \((p,\xi)\), and their limit in \(A\) equals \(p^{h+1}s\) by \(p\)-adic separation in \(A_{\max}\). In all, \(p^Nw_i\in A\) with \(N=m+h+1\). This is an auxiliary denominator depending on the chosen period basis.\formalnote{14.74}{74}

Here is how the auxiliary denominator disappears. Suppose \(Uy=E^h x\) with \(y\in(B^+)^d,\ x\in A^d\), and a row \(l\) over \(A\) satisfies \(lQ=\delta g\) with \(p\nmid\delta\). For each extended coordinate \(\widehat\lambda_i\), the integral vector \(p^N E^{-h}U\widehat\lambda_i(y)\) belongs to \(P\), because the entries of \(\widehat\lambda_i(y)\) lie in \(S_B\). Applying \(l\) gives entries divisible by \(\delta\). Finite partial sums \(\sum b_i\widehat\lambda_i(y)\) approximate \(y\) to arbitrary precision in \(p^jA_{\max}\). Multiplying by \(U\) and removing \(E^h\) consumes only the fixed loss \(k_U+2h\). Thus \(l(p^N x)\) can be approximated by multiples of \(\delta\) modulo every \((p,\xi)^j A\). The principal ideal \((\delta)\) is closed for this topology under \(p\nmid\delta\). Cancel \(p^N\), using primality of \(p\) in \(A\), to obtain
\[
                              \delta\mid lx .
\]
In particular if \(lQ=0\), take \(\delta=z^j\) and use the solutions \(U'c\) for \(x=\tau^h c,\ c\in A^d\). Separation for the powers of \(z\) gives \(l\tau^h c=0\) for every \(c\), whence \(l=0\). This proves \(\det Q\ne0\).\formalnote{14.75}{75}

\subsection{The determinant modulo \(p\)}

For the integral inverse one needs the stronger \(p\nmid\det Q\). Applying \(\varphi\) to \(UZ=E^h Q\) gives a solution over \(S_B\) for each column of \(\eta^{-1}\varphi Q\), with solution matrix \((E/\varphi E)^h F\varphi Z\). Hence \(\varphi Q=\eta QB\) with \(B\) over \(S\). If \(p\mid\det B\), the determinant identity and bijectivity of \(\varphi\) on \(A\) imply successively \(p^j\mid\det Q\) for all \(j\), contradicting full rank. Moreover \(Qs=pw\) for \(s\in S^d,\ w\in A^d\) forces \(w\in P\) by the solution \(Zs/p\), so \(s\in pS^d\).\formalnote{14.76}{76}

Write bars for reduction to \(\mathcal O_{\mathbb C_p}^{\flat}\). By the last implication, the chart-coordinate rows \(\lambda_i(Q_{j,*})\) separate vectors over \(S\) modulo \(p\). Indeed if all coordinates of \(Qs\) vanish modulo \(p\), its reduction is divisible by all powers of \(\bar z\), hence is zero. Select \(d\) coordinate rows forming a matrix \(D\) over \(S\) with \(\overline{\det D}\ne0\). Allow rows of the identity alongside the prescribed rows. Among the resulting \(d\)-row matrices of nonzero determinant modulo \(p\), choose one with as many prescribed rows as possible. Such a matrix exists since the identity has determinant one. If row \(k\) still has to be replaced, its adjugate column pairs to zero modulo \(p\) with every prescribed row, by maximality. Separation forces that column to vanish modulo \(p\), contradicting the determinant expansion.\formalnote{14.77}{77}

Suppose now \(\det\bar Q=0\). Over the tilt valuation ring choose a nonzero kernel vector of minimal support and divide by a coordinate of least valuation. The resulting \(c\) has \(c_\nu=1\). Every \(g\) fixing \(S\) pointwise fixes \(Z\), by the definition of \(S_B\). The Galois relation on \(U\) then gives \(Q=R'_g g(Q)\) for an integral matrix \(R'_g\). Thus \(g(c)\) is also in \(\ker\bar Q\) and \(g(c)-c\) has smaller support, giving \(g(c)=c\).\formalnote{14.78}{78}

Let \(\pi_j\) denote component \(j\) of an element of the tilt in \(\mathcal O_{\mathbb C_p}/p\). The chart's fixed-element approximation says that for any prescribed \(j\) there exist \(n\ge0,\ s\in S^d\) with \(\pi_{j+n}(\bar s_i)=\pi_j(c_i)\) simultaneously. The underlying field assertion follows by approximating the fixed untilt by a fixed algebraic integer using Ax's theorem, expressing a nonzero \(\mathbb Z_p\)-multiple in the increasing union of the ranges \(\theta\varphi^{-m} S\), and removing this scalar by the chart coordinate.\formalnote{14.79}{79}

Put \(C_n=B\varphi B\cdots\varphi^{n-1}B\) and \(K=\overline{\det B}\,\overline{\det D}\ne0\). Iterating Frobenius gives \(\bar Q\bar C_n c^{p^n}=0\). Also \(K\overline{\det C_n}\mid K^{p^n}\), since the powers of \(\overline{\det B}\) in \(\overline{\det C_n}\) sum to \(1+p+\cdots+p^{n-1}\le p^n-1\). At component \(j+n\), \(c^{p^n}\) can be replaced by \(\bar s\), so this component kills \(Q C_n s\) after reduction. To pass to its selected coordinates one needs the chart divisibility result: if \(\theta\varphi^{-m}x\) is divisible by \(p\), then \(\lambda_i(x)\) is congruent to a multiple of \(p\) in \(S\) modulo \(\varphi^mE\). Consequently component \(j+n\) also kills \(D C_n s\). Since the \(\nu\)-coordinate of \(\bar s\) maps to \(1\), the adjugate gives \(\pi_{j+n}(\overline{\det D}\,\overline{\det C_n})=0\). The preceding divisibility yields \(\pi_j(K)=0\). This holds for every \(j\), contrary to \(K\ne0\). Thus \(p\nmid\det Q\).\formalnote{14.80}{80}

Take \(\delta=\det Q\) and \(l\) any row of \({\rm adj}(Q)\) in the divisibility result. It gives integrality of \(Q^{-1}x\) for every \(A\)-vector \(x\) admitting any \(B^+\)-solution. Apply this first to \(UU'=E^h\tau^h I\) to get \(QQ'=\tau^h I\) with \(Q'\) over \(A\). To divide \(U\) itself by \(Q\), approximate \(U\) as \(p^jU=C+p^j\mu^h R\), with \(C\) over \(A\) and \(R\) over \(A_{\max}\). Such approximations exist with common \(j\) since \(U\) has finitely many entries. As \(UU'=(E\tau)^h I\), the integral matrix \(C=p^j U(I-v^h U'R)\) admits the solution \(E^h p^j(I-v^h U'R)\). Hence \(C=QL\) integrally. The error is also divisible by \(Q\), since \(\tau^h I=QQ'\). This gives \(U=QZ'\) over \(B^+\) and, on canceling \(Q\) in \(QZ'Z=E^h Q\), gives \(Z'Z=E^h I\). Finally for \(g\) fixing \(\theta z\) and the constants, \(Ug(Z)=E^h R'_g g(Q)\) by the transformation of \(U\) and the fact that \(gE/E\) is a unit. The same integral solution criterion and cancellation of \(U\) now give \(gZ=ZG_g\) with \(G_g\) over \(A\).\formalnote{14.81}{81}

\subsection{The two lattice inclusions and the retraction used in the cup bound}

Specialize to \(a=b=H\), so \(h=2H\). Write \(L_A=A\otimes T\) for the ordinary lattice and \(e_1,\ldots,e_d\) for its basis. The source defines \(N_A\) as the \(x\in L_A[1/\tau]\) on which every chart-equivariant form \(\lambda\) has
\(\lambda(\varphi^j x)\in B_{\rm dR}^+\) for every \(j\ge1\). A form here is \(A\)-linear into \(B_{\rm dR}\) and equivariant on the vectors \(1\otimes T\). Define \(N\subset N_A\) by the additional condition \(\lambda(x)\in\mathcal K_0\) for all such forms, where \(\mathcal K_0\) is the field generated by the fixed local field and \(S_B\) in \(B_{\rm dR}\). Put
\[
           n_i=\tau^{-H}\sum_j Q_{ji}e_j .
\]
These vectors exist in \(W(\mathbb C_p^\flat)\otimes T\) since the reduction of \(\tau\) modulo \(p\) is nonzero. In matrix notation \(\boldsymbol n=\tau^{-H}Q\) in the ordinary basis. Combining \(U=(E\tau)^H M\) with \(UZ=E^{2H}Q\) gives \(E^H\boldsymbol n=MZ\) after passage to \(B_{\rm dR}\). At every \(m\ge0\) its corresponding evaluation formula is
\[
  E_m^H\lambda\!\left(\sum_i s_i\varphi^m n_i\right)
     =\bigl(\lambda(\boldsymbol e)\,\varphi^m M\bigr)
                              \bigl((\varphi^m Z)s\bigr),
       \qquad E_m=\varphi^m E,\quad s\in A^d .
\]
All ring entries here are mapped to \(B_{\rm dR}\) after the displayed Frobenius. The first row on the right is fixed, hence lies in the fixed local field inside \(B_{\rm dR}^+\). For \(m>0\), \(\theta(E_m)\ne0\), so its power can be removed from the order test. This gives \(n_i\in N_A\). The value-field test holds for each \(\varphi^m n_i\), since \(\varphi^m Z\) has entries in \(S_B\) and \(E_m\in S\).\formalnote{14.82}{82}

Conversely, the rows of \(\varphi^m M^{-1}\) give crystalline forms whose coordinate rows in this formula are the coordinate unit vectors. Combining their values with \(\varphi^m Z'\) detects \(E_m^{2H}s_i\). In detail, if the values of all crystalline forms on \(\sum s_i\varphi^m n_i\) belong to a \(B^+\)-submodule \(W\) of \(B_{\rm dR}\), so do all \(E_m^{2H}s_i\). This follows by multiplying the vector of these particular evaluations by \(E_m^H\varphi^m Z'\). In particular \(W=0\) detects linear independence.\formalnote{14.83}{83}

For \(x\in N_A\), some \(\tau^r x\) belongs to \(L_A\). Since \(QQ'=\tau^{2H}I\), multiplying this latter vector by \(\tau^H\) expresses it as a combination of the \(n_i\). It remains to remove powers of \(\tau\) from the resulting equation. Suppose \(\sum s_i n_i=\tau^j x\) with \(j>0\). On every positive Frobenius translate, the values of forms on the right have positive Hodge order, since \(\theta(\varphi^m\tau)=0\) for \(m\ge1\). Apply detection with \(W={\rm Fil}^1\) to the translated coefficients \(\varphi^m s_i\). Since \(\theta(E_m)\ne0\), we get \(\theta(\varphi^m s_i)=0\) for all \(m\ge1\). The criterion for all positive translates is precisely \(\tau\mid s_i\): apply the criterion for \(\mu\) to \(\varphi s_i\), use \(\varphi\tau=\mu\), and invert Frobenius on \(A\). Divide by \(\tau\) and iterate. This proves \(N_A=\bigoplus_i A n_i\). It also gives the exact comparison needed for Section~\ref{sec:kisin-height},
\[
          \tau^H N_A\subseteq L_A,\qquad
          \tau^H L_A\subseteq N_A .
\]
The first inclusion follows from the definition of \(n_i\), the second from \(Q'\). Likewise \(\varphi Q=v^H QB\) and \(\varphi\tau=vE\tau\) give \(E^H\varphi\boldsymbol n=\boldsymbol n B\). Thus \(E^H\varphi\) carries the \(S\)-span into itself. The assertion about Frobenius needed here is this inclusion; a reverse finite-height bound would be a further assertion.\formalnote{14.84}{84}

For descent to \(S\), at level \(m\) let \(\mathcal K_m\subset B_{\rm dR}\) be the field generated over the fixed local field by \(S_B\) and \(\varphi^{-m}S\). For any finite family \(b_i\in\mathcal K_m\), the scalar divisibility argument above implies
\[
  \sum_i b_i\varphi^{-m}(r_i)\in W
       \quad\Longrightarrow\quad
  \sum_i b_i\varphi^{-m}(f(r_i))\in W
\]
for \(r_i\in A\), \(S\)-linear \(f:A\to S\), and every \(B_{\rm dR}^+\)-submodule \(W\) of \(B_{\rm dR}\). Indeed, multiply the \(b_i\)'s by a common denominator to write them as sums of \(c\,b\,\varphi^{-m}(s)\) for fixed algebraic \(c\), \(b\in S_B,\ s\in S\). Apply the fixed-coefficient divisibility to the expansion with the extension of \(\varphi^{-m}f\varphi^m\). It commutes with \(b\), and \(S\)-linearity passes \(s\) through \(f\). After the common factor is removed, the new sum is some element of \(B_{\rm dR}^+\) times the old sum.\formalnote{14.85}{85}

At \(m=0\) and \(W=0\), this implies that \(f\) preserves all linear relations between value-field vectors in the fraction span of \(L_A\). To prove it, apply the scalar implication to every crystalline form on each side of the relation. The resulting difference pairs to zero with every such form. Multiplying by a nonzero scalar puts it in \(L_A\), then multiplying by \(\tau^H\) puts it in the span of the frame, where detection with \(W=0\) proves it vanishes. In particular if \(x\) is a value-field point in \(\varphi^m N_A\), write \(x=\sum_i r_i\varphi^m n_i\). The vectors on both sides have their form values in \(\mathcal K_0\). Apply the relation-preserving map \(f=\lambda_0\) at level zero. It fixes the coefficient of \(x\), so \(x=\sum_i\lambda_0(r_i)\varphi^m n_i\). At \(m=0\) this gives \(N=\bigoplus_i S n_i\), and at every \(m\) it gives the chart span used in the radius argument.\formalnote{14.86}{86}

Finally allow denominators which are finite products of the \(E_j\). Let \(N_A^{\rm loc}\) be the corresponding saturation of \(N_A\) in the ambient Witt module and \(N^{\rm loc}\) its value-field points. Clearing a common product and using the frame gives \(N_A^{\rm loc}=A\,N^{\rm loc}\). The same relation argument makes
\[
       {\mathfrak r}\left(\sum_i r_i w_i\right)
          =\sum_i\lambda_0(r_i)w_i,\qquad w_i\in N^{\rm loc},
\]
well defined. Indeed the right-hand map on \(A\otimes_S N^{\rm loc}\) kills the kernel of the sum on the left. It is a retraction since \(\lambda_0(1)=1\), and commutes with maps of lattices by the formula on sums. Thus it respects the kernel and the quotient of a map in the lifting calculation.

Let \(\lambda_m(y)=\lambda(\varphi^{-m}y)\), with \(\lambda\) equivariant, and \(c\in\mathcal K_m\). The form values \(\lambda_m(w_i)\) for \(w_i\in N^{\rm loc}\) belong to \(\mathcal K_m\). This follows by writing a nonzero scalar multiple of \(w_i\) as a combination of \(\varphi^m\)-translates of localized value-field points and choosing a chart coordinate nonzero on the scalar, then dividing the resulting relation in \(\mathcal K_m\). Apply the scalar implication with \(W={\rm Fil}^j\), the summands \(\lambda_m(w_i)\), and one extra term \(-c\) of coefficient \(1\). We obtain
\[
      \lambda_m(y)-c\in{\rm Fil}^j
        \quad\Longrightarrow\quad
      \lambda_m({\mathfrak r}y)-c\in{\rm Fil}^j .
\]
The retraction costs no power of \(p\) and no Hodge order. Functoriality preserves the letter projection as well. These facts justify its use after the more costly divided-difference operators in Sections~\ref{sec:cup-comparisons} and~\ref{sec:kisin-height}.\formalnote{14.87}{87}

\paragraph{Earlier results used here.}
The construction takes as input the specified Kummer chart, crystallinity on its subgroup, the bounds on the jumps, and the fixed change of columns proved in \S\ref{sec:derham-columns}. Its scalar ingredients include the period-ring divisibility tests for simultaneous Frobenius translates, completeness and the integral approximation properties of \(A_{\max}\), fixed ratios in \(B_{\max}\), Ax approximation of fixed elements, and adic bases for a regular sequence. Section~\ref{sec:integral-coordinates} constructs the extended coordinate maps and proves the simultaneous approximation and divisibility tests. The localized congruence just proved preserves targets already in \(\mathcal K_m\). Obtaining a fixed local-field target from Galois invariance modulo \({\rm Fil}^j\) in the long-gap argument follows from the fixed-class theorem of \S\ref{sec:tate-constants}. Producing the global Galois representations and establishing crystallinity of their restrictions are separate steps in the argument.\formalnote{14.88}{88}


\section{Fixed classes in the de Rham filtration}
\label{sec:tate-constants}

Let \(G_p=\operatorname{Gal}(\overline{\mathbb Q}_p/\mathbb Q_p)\),
let \(U\subset G_p\) be an open subgroup, and let \(K=\overline{\mathbb Q}_p^U\).
Then \(K/\mathbb Q_p\) is finite, \(U=G_K\), and the completion of its
algebraic closure identifies with \(\mathbb C_p\). Tate's fixed-field
and twisted-invariant theorems
\cite[\S3.3, Theorems~1--2]{arith-tate-pdiv}
\cite[Theorem~2.2.7]{legacy-brinon-conrad} give
\begin{equation}\label{eq:tate-twisted-invariants}
  \mathbb C_p^U=K,\qquad
  \bigl(\chi(g)^m g(y)=y\ \text{for all }g\in U,\ m\ne0\bigr)
     \ \Longrightarrow\ y=0,
\end{equation}
where \(\chi\) is the cyclotomic character. For \(m\ne0\), the image of inertia under \(\chi^m\) is infinite, as required by the twisted-invariant theorem. We use this for the full
image of the local group when the hypothesis holds everywhere, and
on an open subgroup contained in the image of the identity
neighborhood when it holds only near \(1\).

\phantomsection\label{route:25}\leanroute{The original formalization proves the two scalar assertions by Ax's polynomial approximation and a normalized trace on the cyclotomic tower. The community formalization invokes the Ax--Sen--Tate fixed-field theorem and Tate's vanishing theorem here.\formalnote{14.89}{89}}

\subsection{Lifting an invariant class}

In \S\ref{sec:kisin-height} an integral frame produces coordinates
fixed modulo a filtration step. The retraction of
\S\ref{sec:period-matrix} needs representatives in \(K\). More
precisely, for any \(B_{\rm dR}^+\)-submodule \(S\) of \(B_{\rm dR}\),
\begin{equation}\label{eq:fixed-de-rham-class}
   gx-x\in S\ \ (g\in U)
   \quad\Longrightarrow\quad x-a\in S\ \text{for some }a\in K.
\end{equation}
Write \(\theta:B_{\rm dR}^+\to\mathbb C_p\) for the residue map and
choose the uniformizer \(t\) with \(g(t)=\chi(g)t\). Suppose \(x\notin S\)
and write \(x=t^n v\) with \(v\in(B_{\rm dR}^+)^\times\). For each \(g\in U\),
\[
          gx-x=t^n\bigl(\chi(g)^n g(v)-v\bigr)\in S.
\]
If the parenthesis were a unit, the \(B_{\rm dR}^+\)-module
property would give \(x\in S\). Its residue therefore vanishes.
The nonzero element \(\theta(v)\) satisfies
\(\chi(g)^n g(\theta(v))=\theta(v)\). By \eqref{eq:tate-twisted-invariants}, \(n=0\) and
\(\theta(v)\in K\). Put \(a=\theta(v)\), viewed in \(B_{\rm dR}^+\)
via the embedding of \(K\) lifting the residue embedding. If
\(x-a\notin S\), it still satisfies the invariance hypothesis, and
the same argument would force its \(t\)-adic order to be zero.
Its residue is zero, so its order is positive. This proves \eqref{eq:fixed-de-rham-class},
in particular for \(S=t^jB_{\rm dR}^+\) for every \(j\in\mathbb Z\)
and for \(S=0\).\formalnote{14.90}{90}

\subsection{Separation of weights}

Suppose \(W\subset\mathbb C_p^d\) is a vector space stable under the
semilinear action \(a_i\mapsto\chi(g)^{u_i}g(a_i)\) on \(U\). Within
the support of any nonzero vector of \(W\), choose a nonempty support
minimal among vectors of \(W\). The subspace with this support
allowed is one-dimensional, since eliminating a common nonzero
coordinate would give a smaller support. Normalize a vector
\(a\) on it by setting a coordinate of maximal norm, say \(a_k\),
equal to \(1\). Every twisted conjugate has that support and is
\(\chi(g)^{u_k}a\), as its \(k\)-th coordinate shows. By \eqref{eq:tate-twisted-invariants},
\(a_i\) can be nonzero only if \(u_i=u_k\), in which case \(a_i\in K\).
Subtracting a multiple of this normalized vector reduces the
support of the original vector. Iteration shows that \(W\) is spanned
by vectors with coordinates in \(K\), each supported in a single
weight. Section~\ref{sec:jump-symbol} identifies the particular
symbols to which this separation applies.\formalnote{14.91}{91}


\section{Hodge jumps and the number of leading coefficients}
\label{sec:jump-symbol}

To raise the order of a vector by one, we must cancel its leading coefficient under each equivariant form of the prescribed order. Fix a component \(i\), indexed by a nonzero element of the chart residue field, and an integer \(u\). Write \(m_{i,u}\) for the multiplicity of the Hodge jump \(u\) in component \(i\). We prove that at most \(m_{i,u}\) independent coefficients need to be tested in the lifting step of \S\ref{sec:kisin-height}. Two fields occur in this count. The multiplicities are computed using forms that transform near the identity. The coefficients used to cancel a leading term must preserve equivariance on the entire chart subgroup.

\subsection{Taking the symbol}

Let \(T\) be the finite free \(O\)-lattice with action \(\rho\), where \(O\) is the ring of integers of a finite extension of \(\mathbb Q_p\). Let \(\Gamma'\) be the subgroup on which the chart requires crystallinity, and let \(K\) be the algebraic fixed field of its image in \(G_p\). This image is open. Write \(g\) for its action on scalars and \(\chi(g)\) for its cyclotomic value. Set \(m=p^f-1\), where \(f\) is the chart residue degree. The coefficient datum includes \(\zeta\in O\) with \(\zeta^m=1\). The component index \(i\in\mathbb F_{p^f}^{\times}\) gives the Teichmüller scalar \(\omega_i=[i]\). On \(W(\mathbb C_p^\flat)\otimes_{\mathbb Z_p}T\), multiplication by \(\zeta\) on \(T\) gives an operator \(Z\). Put
\[
              P_i=\frac1m\sum_{j=0}^{m-1}(\omega_i^{-1}Z)^j,
              \qquad w_i=\theta(\omega_i)\in\mathbb C_p .
\]
The geometric sum gives \(ZP_i=\omega_iP_i\), and \(P_i\) is the identity on that eigenspace. Its coefficients belong to \(A_{\inf}\) and are fixed by \(\Gamma'\). The operator therefore preserves the ordinary \(A_{\inf}\)-lattice and commutes with this group.\formalnote{14.92}{92}

Let \(\mathcal F_{i,u}\) consist of the \(A_{\inf}\)-linear maps \(\mu\) from the displayed tensor product to \(B_{\rm dR}\) that are \(\Gamma'\)-equivariant on \(1\otimes T\) and satisfy
\[
            \mu(P_i(1\otimes v))\in t^{-u}B_{\rm dR}^+
                  \qquad(v\in T).
\]
Write \(\theta_{\rm dR}:B_{\rm dR}^+\to\mathbb C_p\) for reduction modulo \(t\). Define the symbol
\[
                s_\mu(v)=\theta_{\rm dR}
                    \bigl(t^u\mu(P_i(1\otimes v))\bigr).
\]
It is \(\mathbb Z_p\)-linear on \(T\). Equivariance of \(\mu\) extends to the ordinary lattice by semilinearity of the group action. Since \(g(t)=\chi(g)t\), the two transformation rules are
\[
      s_\mu(\rho(g)v)=\chi(g)^{-u}g(s_\mu(v))\quad(g\in\Gamma'),
      \qquad s_\mu(\zeta v)=w_i s_\mu(v).
\]
The kernel of \(\mu\mapsto s_\mu\) is exactly \(\mathcal F_{i,u-1}\). Indeed the residue vanishes for every \(v\) precisely when every value of \(t^u\mu(P_i(1\otimes v))\) lies in \(t B_{\rm dR}^+\). Both spaces of forms are \(K\)-linear through the algebraic constants in \(B_{\rm dR}^+\), and reduction preserves these constants.\formalnote{14.93}{93}

\subsection{The two restrictions on its coordinates}

For each embedding \(\gamma:O\hookrightarrow\mathbb C_p\), write \(j_\gamma(u)\) for the dimension over \(\overline{\mathbb Q}_p\) of the \(\gamma\)-semilinear forms \(b:T\to\mathbb C_p\) satisfying \(b(\rho(g)v)=\chi(g)^{-u}g(b(v))\) near the identity. Multiplication by an algebraic constant preserves this last condition because a sufficiently small open subgroup fixes the constant.\formalnote{14.94}{94}

Here is how crystallinity gives a basis for counting these forms. For \(\gamma\), map a tensor \(b\otimes v\) over \(B_{\max}\) to \(b\) times the \(\gamma\)-coordinates of \(v\) in \(B_{\rm dR}^d\), where \(d=\operatorname{rank}_O T\). This map takes invariant periods to vectors fixed near the identity under the matrix action \(x\mapsto\gamma(\rho(g))g(x)\), because the coefficients of \(v\) are algebraic. The invariant periods span by potential crystallinity and the image contains the coordinate vectors. Thus there is an invertible matrix \(U\) with invariant columns.\formalnote{14.95}{95}

Apply the integral frame of \S\ref{sec:derham-columns} to \(U\) and the subgroup on which its columns are invariant. It gives \(UY=W\operatorname{diag}(t^{k_j})\), with \(Y\) integral algebraic of determinant one and fixed near the identity, and \(W\) invertible over \(B_{\rm dR}^+\). Thus \(P=\theta_{\rm dR}(W)\) is invertible. Invariance gives
\[
       P^{-1}\gamma(\rho(g))
          =\operatorname{diag}(\chi(g)^{-k_j})\,g(P^{-1}).
\]
The rows define a basis \(b_{\gamma,j}\) of the \(\gamma\)-semilinear forms over \(\mathbb C_p\). Across all \(\gamma\) these rows together form a basis of all \(\mathbb Z_p\)-linear forms on \(T\) over \(\mathbb C_p\). Indeed the fraction field of \(O\) is separable over \(\mathbb Q_p\), so \(O\otimes_{\mathbb Z_p}\mathbb C_p=\prod_\gamma\mathbb C_p\), and the dual of \(T\otimes_{\mathbb Z_p}\mathbb C_p\) splits into the corresponding summands.\formalnote{14.96}{96}

Expand a weight-\(u\) form in this basis with coefficients \(d_{\gamma,j}\). Test equivariance on a finite \(\mathbb Z_p\)-basis of \(T\) and intersect the finitely many neighborhoods involved. Comparing coordinates gives
\[
              \chi(g)^{k_j-u}g(d_{\gamma,j})=d_{\gamma,j}
                        \qquad(g\text{ near }1),
\]
where \(k_j\) is the weight in the indicated \(\gamma\)-summand. Tate's vanishing from Section~\ref{sec:tate-constants} kills every coefficient with \(k_j\ne u\). The others are fixed near \(1\), hence algebraic. Taking just the \(\gamma\)-summand shows that exactly \(j_\gamma(u)\) rows there have weight \(u\). For a symbol, evaluating also at \(\zeta v\) gives \((\gamma(\zeta)-w_i)d_{\gamma,j}=0\). Hence the \(\mathbb C_p\)-span of the symbols has dimension at most\formalnote{14.97}{97}
\[
                         m_{i,u}=\sum_{\gamma(\zeta)=w_i}j_\gamma(u).
\]

Finally choose a basis of that span consisting of actual symbols \(s_{\nu_1},\ldots,s_{\nu_r}\). If \(s_\mu=\sum_h c_hs_{\nu_h}\), every \(g\in\Gamma'\) acts on both sides by the same weight law. Uniqueness of coordinates gives \(g(c_h)=c_h\) for every such \(g\). The open-subgroup case of Section~\ref{sec:tate-constants} puts every \(c_h\) in \(K\). Consequently \(\mu-\sum_h c_h\nu_h\) is still chart-equivariant and has zero symbol. This proves
\[
       \dim_K\bigl(\mathcal F_{i,u}/\mathcal F_{i,u-1}\bigr)=r\le m_{i,u}.
\]
If the selected embeddings have no jump at \(u\), every form of type \(u\) already has type \(u-1\). At level \(a\) in Section~\ref{sec:kisin-height}, undoing \(\varphi^a\) replaces \(i\) by \(\sigma^{-a}i\) in this count. The later factor \(f\) in the colength enters upon counting residues over \(\mathbb F_p\).\formalnote{14.98}{98}


\section{Integral height, depth, and the coefficient at a long gap}
\label{sec:kisin-height}

Three estimates in Section~\ref{sec:cup-comparisons} concern an extension of crystalline lattices at one place above \(p\). The first gives a uniform power of the chart coordinate that lifts constants. The second raises the order of Galois differences by one. The third constructs a point which is almost fixed to high order even when some Hodge jumps are far apart. We prove these estimates with their respective denominator bounds.

\subsection{The initial height by contraction}

Write \(A_{\inf}=W(\mathcal O_{\mathbb C_p}^{\flat})\), \(\Lambda=W(\mathbb F_{p^f})\), \(S=\Lambda[[z]]\) on the chosen Kummer chart, and \(E\in S\) for its generator of \(\ker\theta\) in \(A_{\inf}\). The chart has \(E\equiv p\bmod z\) and \(z^e\in(E,p)S\). Write \(E_h=\varphi^hE\), \(\theta_h=\theta\varphi^{-h}\), \(\tau=[\epsilon^{1/p}]-1\), and \(\Delta=z\tau\). The residues \(\theta_h(\Delta)\) are nonzero. In particular division by \(\Delta\), when defined on the integral difference of vectors, consumes no Hodge order.

Let \(q:\widetilde Y\to C\) be the quotient of a pairing side by its kernel, with \(C\) the trivial quotient lattice. For \(c\in C\), write \(c_i\) for the projection of \(1\otimes c\) to component \(i\), \(\sigma i=i^p\), and \(M_l(i)\) for the \(A_{\inf}\)-module of vectors at \(i\) with the standard orders of level \(l\) in the localized Kisin lattice extended to \(A_{\inf}\). These orders are \(0\) at \(E_h\) for \(h<l\) and \(-u\) for \(h\ge l\) on each jump \(u\) at \(\sigma^{-h}i\). For a fixed positive modulus \(r\), define \(J_l(r,i)\) as the set of \(s\in A_{\inf}\) for which there is one \(t\in A_{\inf}\) such that every \((s-rt)c_i\) has a lift in \(M_l(i)\). The vector lifting a given constant may depend on \(c\). Write \(J_l(0,i)\) for exact lifting. Then
\[
 J_l(r,i)=J_l(0,i)+(r),\qquad
 \varphi^l J_0(r,i)=J_l(r,\sigma^l i).
\]
The second equality follows by transporting both the order conditions and the constants. If all jumps at \(i\) satisfy \(u\le\omega\), multiplying a level-zero lift by \(E^\omega\) raises its order \(-u\) at \(E\) to at least zero and preserves its other orders. Thus \(E^\omega J_0(r,i)\subseteq J_1(r,i)\).\formalnote{14.99}{99}

Set \(r=p^{\kappa+1}\). Some common power of \(z\) belongs to all \(J_0(r,i)\). Indeed, if \(H\) bounds the absolute Hodge jumps of the extension as in its frame, \(\tau^H\) takes its ordinary \(A_{\inf}\)-lattice into its Kisin \(A_{\inf}\)-lattice. The projection of \(\tau^H(1\otimes(0,c))\) to component \(i\) lifts \(\tau^H c_i\). Hence the intersection of the ideals contains \(p^{\kappa+1}\) and \(\tau^H\). Modulo \(p\) we are in the rank-one valuation ring \(\mathcal O_{\mathbb C_p}^\flat\), where \(z\bmod p\) has positive valuation and \(\tau\bmod p\) is nonzero. Some power of \(z\bmod p\) is therefore a multiple of \(\tau\bmod p\). This places \(z\) in the radical of the intersection.\formalnote{14.100}{100}

To bound the least such exponent, we need a divided difference. Suppose \(f\in S\cap J_l(r,i)\) and, for \(g\) in the chart group,
\[
             D(s+\beta f)=\alpha(gf-f),\qquad \alpha,\beta\in A_{\inf}.
\]
Take \(D=\Delta\) at level zero, or \(D=\Delta E\) at level one. Then \(s\in J_l(r,i)\). Here are the integral requirements for this calculation. For \(f\in S\) the chart-coordinate retraction gives a common correction \(t\in S\) and lifts \(x\) in the localized \(S\)-lattice with standard orders. There \(D\mid gt-t\), and \((gx-x)/D\) belongs to \(M_l(i)\). For division by \(\Delta\) this uses divisibility in the frame and the nonzero residues above. At level one \(x\) has integral value at \(E\) in the chart's value field, on whose residue the group acts trivially there. Its difference gains an order at \(E\), which permits the additional division. Now \(\alpha(gx-x)/D\) lifts
\[
           \left(s+\beta f-r\alpha(gt-t)/D\right)c_i .
\]
This proves \(s+\beta f\in J_l(r,i)\), and subtracting \(\beta f\) proves the claim.\formalnote{14.101}{101}

Choose \(g\) with \(gz=[\epsilon]z\). For each \(j\ge1\), write \(v=v_p(j)\). The prime-to-\(p\) part of \(j\) gives a unit geometric sum, and successive \(p\)-th powers factor \([\epsilon]^j-1\), up to a unit, as \(\tau E_0E_1\cdots E_v\). Since \(E_h\equiv p\bmod z\), there are \(\alpha,\beta\in A_{\inf}\) with
\[
             \alpha([\epsilon]^j-1)=\tau E(p^v+\beta z).
\]
Apply the divided difference to \(f=p^c z^j\). At level one the resulting element is \(p^{c+v}z^{j-1}\), and at level zero it is \(p^{c+v+1}z^{j-1}\) because \(E=p+z b\). In each case the other term is a multiple of \(f\). In particular one can lower \(j\) to \(p\lfloor j/p\rfloor\) at level one without spending a power of \(p\). Frobenius descent then puts \(p^c z^{\lfloor j/p\rfloor}\) in the preceding component of level zero.\formalnote{14.102}{102}

Let \(b\) be the least exponent common to the level-zero ideals and assume the bound \(u\le\omega\) in every component. At level one each ideal contains \(z^bE^\omega\) and \(p^{\kappa+1}\). Expand \(z^b(z^e)^{\omega+\kappa}\) using \(z^e\in(E,p)\). Every term includes \(z^bE^\omega\) or \(p^{\kappa+1}\), hence lies in the ideal. Rounding and descending as above gives the common level-zero exponent \(\lfloor(b+e(\omega+\kappa))/p\rfloor\). Minimality proves\formalnote{14.103}{103}
\[
                       (p-1)b\le e(\omega+\kappa).
\]

\begin{example}[Exact lifting at small height]\label{ex:small-height-lifting}
Work over an unramified chart at \(p=5\), so \(e=1\). Suppose that every Hodge jump is at most \(\omega=3\), and ask for lifts of constants modulo \(5\), so \(\kappa=0\). Let \(b\) be the least nonnegative exponent such that multiplying the constants by \(z^b\) permits these lifts in every component. The inequality \((p-1)b\le e(\omega+\kappa)\) gives
\(4b\le3\). Since \(b\) is a nonnegative integer, \(b=0\). Thus
\(z^b=1\) belongs to \(J_0(5,i)\) for each component \(i\).
By the definition of the lifting ideal, \(1-5t\in J_0(0,i)\) for
some \(t\in A_{\inf}\). This element is a unit, so \(1\in J_0(0,i)\):
the constants lift exactly. With \(\omega=4\) the same bound would
only give \(b\le1\), which allows a factor of \(z\).
\end{example}

In general, \(\kappa=0\) and \(e\omega\le p-2\) force \(b=0\).
For arbitrary \(\kappa\), once \(p^\kappa\in J_0(p^{\kappa+1},i)\),
the same unit argument gives \(p^\kappa\in J_0(0,i)\). Indeed, the
ideal contains \(p^\kappa(1-pt)\) for some \(t\in A_{\inf}\).
The initial estimate supplies a power of \(z\) modulo \(p^{\kappa+1}\).
The next steps convert this into a power of \(p\).

\subsection{Raising a depth by one}

At level \(a\) define the cost of increasing depth \(n\in\mathbb Z\) by one as
\[
 c(a,n)=v_p(|n|)+d+3a,\qquad v_p(0):=0\ \text{in this formula},
 \qquad
 d=\begin{cases}0&e=1,\\v_p(e)+1&e>1.\end{cases}
\]
We derive the bound for \(n\ne0\). Choose \(g\) with \(gz=[\epsilon]^{p^a}z\), and write \(gE=\omega E\) for a unit \(\omega\) in \(A_{\inf}\). Then \(gE_a=\Omega E_a\), with \(\Omega=\varphi^a\omega\), and \(g\) acts trivially on \(S\bmod E_a\) in \(A_{\inf}/E_a\). Indeed \(\theta_a[\epsilon]\) has order \(p^a\), and the chart group fixes \(\Lambda\). The divisibility \(\Delta E\mid gE-E\) gives \(\Delta\mid\omega-1\). As \(\Delta\mid\varphi^a\Delta\), the element \(X=(\Omega^n-1)/\Delta\) is integral, including for negative \(n\).

The different bounds the denominator of \(X^{-1}\) at \(E_a\). Divide \(z^e\) by \(E\) in \(S\) with remainder \(R(z)\) of degree \(<e\) and take \(P(T)=T^e-R(T)\). Then \(P(z)=EG\) and \(\theta(P'(z))\mid p^d\). For \(e=1\) the derivative is \(1\). For \(e>1\), divisibility of \(P'(\theta z)\) by \(p^{v_p(e)+1}\) already forces divisibility of its degree-\((e-1)\) coefficient \(e\) by this power, using the unique integral remainder in the basis \(1,\ldots,(\theta z)^{e-1}\). This proves the required bound. Write \(gz-z=Ew\). Since \(E\tau\) is associated to \([\epsilon]-1\), the geometric sum in \(([\epsilon]^{p^a}-1)/([\epsilon]-1)\) has \(\theta\)-value \(p^a\), and \(\theta w\) is associated to \(p^a\theta\Delta\). Taylor expansion to first order in \(Ew\), followed by division by \(E\) and application of \(\theta\), yields\formalnote{14.104}{104}
\[
       \theta(\omega-1)\theta G=\theta(P'(z))\theta w,\qquad
       \theta(\omega-1)\mid p^{d+a}\theta\Delta .
\]

Normalize \(v_p(p)=1\). As \(v_p(\theta\tau)=1/(p-1)\) and \(v_p(\theta z)>0\), one has \(v_p(\theta(\omega-1))>1/(p-1)\). In this range the first binomial term of \((1+x)^p-1\) has strictly lowest valuation, giving \(v_p(x)+1\). A prime-to-\(p\) power preserves the valuation, and taking inverses preserves it for a unit. Thus \(\theta(\omega^n-1)\) is associated to \(p^{v_p(|n|)}\theta(\omega-1)\). Finally \(\theta\Delta\mid p^{2a}\theta_a\Delta\). For \(a=0\) this is equality, and for \(a\ge1\) each of \(\theta z,\theta\tau\) divides \(p\). Since
\[
       \theta_a(\Delta)\theta_a(X)=\theta(\omega^n-1),
\]
cancellation of the nonzero \(\theta_a\Delta\) gives \(\theta_a X\mid p^{c(a,n)}\). Surjectivity of \(\theta_a\) supplies \(h\in A_{\inf}\) with \(hX\equiv p^{c(a,n)}\bmod E_a\).\formalnote{14.105}{105}

Let \(x\) be in the localized Kisin \(S\)-lattice. If \(\lambda\) is a chart-equivariant form, write \(\lambda_a(x)=\lambda(\varphi^{-a}x)\) and \(\iota_a\) for the map \(A_{\inf}\to B_{\mathrm{dR}}\) via \(\varphi^{-a}\). Suppose \(y=\lambda_a(x)-k\in\mathrm{Fil}^n\) for a constant \(k\) in the chart group's fixed local field. Write \(\pi=\iota_a E_a\). The integral value \(y/\pi^n\) belongs to the chart's value field, in which \(g\) fixes its residue. Since \(g\pi=\iota_a(\Omega)\pi\), it follows that
\[
                     gy-\iota_a(\Omega^n)y\in\mathrm{Fil}^{n+1}.
\]
For \(w=(gx-x)/\Delta\), equivariance and the fact that \(\iota_a\Delta\) has Hodge order zero give \(\lambda_a(w)-\iota_a(X)y\in\mathrm{Fil}^{n+1}\). Thus \(p^{c(a,n)}x-hw\) evaluates to \(p^{c(a,n)}k\bmod\mathrm{Fil}^{n+1}\). The scalar estimate alone would only construct a vector over \(A_{\inf}\). Applying the \(S\)-linear frame-coordinate retraction keeps the congruence in the localized \(S\)-lattice and keeps the correction in the kernel over a constant quotient. All earlier difference orders persist.\formalnote{14.106}{106}

For \(n=0\), choose \(m=p^{d+3a}\) elements \(g_j\) with \(g_jz=[\epsilon]^jz\), \(0\le j<m\). Modulo \(E_a\), every element of \(S\) is a polynomial in \(z\) of degree \(<ep^a\) with coefficients in \(\Lambda\). On its term of degree \(r\), the sum over \(j\) contributes the sum of \(\theta_a([\epsilon])^{jr}\). This vanishes unless \(p^a\mid r\), since \(p^a\mid m\). The surviving monomials evaluate to powers of \(\theta z\) times fixed coefficients. Consequently \(\sum_jg_j\) gives a fixed residue on \(S\). The value-field residue lemma carries this to each integral \(\lambda_a(x)\), and the fixed-value lemma yields a constant modulo \(\mathrm{Fil}^1\). The constant may depend on \(x,\lambda\). On the trivial quotient a vector lying over \(c_i\) goes to \(m c_i\). This proves the zero step at the same stated cost. A finite sequence of steps needs to charge a given pair \((a,n)\) only once across the different jumps, since the operator at that pair works on all forms simultaneously.\formalnote{14.107}{107}

The cost also gives the second contraction used at sparse letters. At \(a=0\), for \(j>0\) choose the preceding data so that \(\Delta X=\omega^j-1\) and \(hX=p^{c(0,j)}+Eb_1\). Write \((gz/z)^m=1+\tau E b_2\), which follows from \(\Delta E\mid gz-z\). For \(f=p^c z^m E^j\), expansion gives
\[
 zh(gf-f)=\Delta E\bigl(p^{c+c(0,j)}z^{m+1}E^{j-1}
                              +(zb_1+h b_2\omega^j)f\bigr).
\]
The level-one divided difference therefore trades a factor \(E^j\) for \(zE^{j-1}\) at cost \(d+v_p(j)\). Starting from \(p^c z^b\) at level zero, multiply by \(E^\omega\) to go up, make \(\omega\) trades, round to a multiple of \(p\), and descend. At the preceding letter we obtain \(p^{c+d\omega+v_p(\omega!)} z^{\lfloor(b+\omega)/p\rfloor}\). This quantifies how a run of small upper jumps reduces the initial height.\formalnote{14.108}{108}

\subsection{The loss in the frame estimate}

Write \(S_B\) for the subring of \(B^+\) singled out by the chart's extended coordinate maps (its bicommutant). The crystalline frame provides a square matrix \(Z\) of size \(h_0\) over \(S_B\), a matrix \(Z'\) over \(B^+\) with \(Z'Z=E^h I\) for a frame height \(h\), and, for the element \(t\) with \(tz=[\epsilon]z\), a matrix \(G\) over \(A_{\inf}\) with \(tZ=ZG\). Fix \(l,\delta,A\ge0\) with \(3e\le P=p^l\) and \(eA\le P\delta\). We use \(A\) for the order of approximation. The radius argument produces a factorization
\[
                 \varphi^l Z= E_l^b D\,V(z^P/p^2)
\]
for any prescribed \(b\ge0\), where \(D\) has fixed entries under the subgroup fixing \(\Lambda\) and \(V\) is an invertible matrix over \(S[[X]]\). All evaluations at \(X=z^P/p^2\) here factor through \(A_{\max}\). Indeed \(z^e\in p A_{\max}\) as \(E\in pA_{\max}\), so \(z^P/p^2\in pA_{\max}\).\formalnote{14.109}{109}

Here is how the scale \(p^2\) enters. Taking the constant parts of the extended chart coordinates of \(p^{-h}Z'\) gives a fixed matrix \(B\) with \(Y=BZ=I+zW\), \(W\) over \(S_B\). The divisibility \(\Delta\mid G-I\) holds entrywise because the entries of \(tZ-Z\) are divisible by \(\Delta\) in \(B^+\), and in \(E^h(G-I)=Z'(tZ-Z)\) one can cancel \(E^h\) from this divisibility in \(A_{\inf}\). Consider a row remainder \(W_n\) over \(S_B\) such that \(qY_i-z^{n+1}W_n\) has entries in the image of \(A_{\inf}\). At \(n=0\) one may take \(q=1\) and the \(i\)-th row of \(W\). Applying \(t\) to this relation and using \(tY=YG\) cancels the terms in \(Y\). The remainder is built from \(([\epsilon]^{n+1}-1)W_n\), terms divisible by \(\Delta\) from \(G-I\), and \(tW_n-W_n\), likewise divisible by \(\Delta\). The same cyclotomic identity used above and \(E=p\bmod z\) yield, after division by \(\tau z^{n+1}\) and projection to the chart,
\[
     p^{v_p(n+1)+1}W_n=s_n+zW_{n+1},\qquad s_n\in S^{h_0} .
\]
Multiplying the old integral row by \(p^{v_p(n+1)+1}\) and adding \(z^{n+1}s_n\) gives the integral row required for the next step.\formalnote{14.110}{110}

Set \(b_n=n-v_p(n!)+2\), so \(b_n=b_{n+1}-2+v_p(n+1)+1\). Under \(\varphi^l\) the row identity gives
\[
 p^{b_n}\varphi^l W_n
 =p^{b_{n+1}-2}\varphi^l s_n+
       (z^P/p^2)p^{b_{n+1}}\varphi^l W_{n+1}.
\]
All coefficient exponents are nonnegative. Iterating and using separation for powers of \(z^P/p^2\) gives \(p^2\varphi^l W\) as an evaluated matrix of power series. Thus \(\varphi^l(BZ)\) is the evaluation of \(I+X F(X)\), a unit because its constant term is \(I\). Writing \(E=p+zr\) gives \(E_l=p(1+Xp\varphi^lr)\) under evaluation. Absorb the \(b\)-th power of this power-series unit into \(V\), and take \(D=(p^b\varphi^lB)^{-1}\). Invertibility follows from the displayed factorization of \(\varphi^l(BZ)\).\formalnote{14.111}{111}

For a scalar series \(F\in S[[X]]\), truncation at \(\delta\) gives
\[
 p^{2\delta}F(z^P/p^2)=s+z^{P\delta}F_1(z^P/p^2),\qquad s\in S .
\]
The tail belongs to \(p^A A_{\max}\) since \(eA\le P\delta\). To approximate it by an element of \(S\) modulo \(E^A B^+\), one uses the extended coordinate projection \(\mathrm{pr}\) fixing these series values. For \(x\in A_{\max}\) the coordinate satisfies
\[
             p^A\mathrm{pr}(x)\in S+E^A S_B .
\]
This follows by writing \(p^A x=a+E^A x'\), \(a\in A_{\inf}\), using that the period-ring generator \(\xi\) for \(\ker\theta\) belongs to \((p,E)\), and applying \(\mathrm{pr}\). It proves that each \(p^{2\delta}F(z^P/p^2)\) is congruent modulo \(E^A B^+\) to an entry of \(S\).\formalnote{14.112}{112}

Denote the evaluated matrix by \(V_0\). Approximate both \(p^{2\delta}V_0^{-1}\) and \(p^{2\delta}V_0\) by matrices \(M,M'\) over \(S\), respectively, modulo \(E^A\). Call a vector \(y\in S^{h_0}\) admissible here when \(V_0 y\) is congruent to a vector of ordinary nonnegative integers modulo \(E^A B^+\). Each column of \(M\) is admissible because \(V_0M\equiv p^{2\delta}I\). Each column of \(p^{4\delta}I-MM'\) is admissible because its transform is zero modulo \(E^A\). Now for every column vector \(x\) over \(S\),
\[
             p^{4\delta}x=(p^{4\delta}I-MM')x+M(M'x).
\]
Thus \(p^{4\delta}S^{h_0}\) is in the \(S\)-span of admissible vectors. The second scalar approximation is needed here to obtain a statement about the span in \(S^{h_0}\) itself.\formalnote{14.113}{113}

Apply the factorization with \(b\) the absolute jump bound of the lattice. The frame identifies value-field points of its \(l\)-th Frobenius translate with the \(S\)-span of the translated frame vectors. Evaluation of the vector corresponding to \(y\) on a chart-equivariant form is a row of fixed integral Hodge coordinates applied to \(DV_0y\). For an admissible \(y\), the integer part multiplied by the fixed matrix \(D\) is fixed. The remaining part is in \(\mathrm{Fil}^A\) since \(E^A\) has Hodge order \(A\) at level zero. Hence \(p^{4\delta}\) times every such point is an \(S\)-linear combination of vectors whose form values are fixed modulo \(\mathrm{Fil}^A\).\formalnote{14.114}{114}

\subsection{Extracting the coefficient that changes the moving side}

Assume now that the chart is unramified in the chosen normalization \(e=1,\ E=z+p\). Let \(l>0\), \(3\le p^l\), \(A\le p^l\delta\). Assume \(E^s\) kills the level-zero cokernel at \(i\), and that every jump at \(\sigma^{-a}i\), \(0\le a<l\), is at most \(g\), with \(g\ge0\). Define
\[
    N=s+g\sum_{0\le a<l}p^a,\qquad N<A .
\]
We construct the coefficient above \(c_i\) used at the pivot. Take \(y\) over \(c_i\) with the standard level-zero orders and with allowance for a pole of order \(s\) at \(E\). Since \(E_a=(E-p)^{p^a}+p\), the expression \(E^s\prod_{a<l} E_a^g\) is \(F(E)\) for a monic integral polynomial \(F\) of degree \(N\). Thus \(F(E)y\) has orders \(g-u\) below level \(l\) and \(-u\) from level \(l\) onwards. These include the standard level-\(l\) orders. At level zero they give at least \(o(u)=\min(A,g-u)\).\formalnote{14.115}{115}

Project the almost fixed span from the previous subsection to the letter \(i\). Unramified Weierstrass division gives a remainder in \(\Lambda\) on division by \(E\), so successive divisions of its scalar coefficients give, for any truncation order \(m\),
\[
 p^{4\delta}F(E)y=\sum_{0\le j<m}E^j\xi_j+E^m v_m .
\]
The remainder \(v_m\) has the standard level-\(l\) orders. Each \(\xi_j\) has these orders and each difference \(\gamma\xi_j-\xi_j\) additionally has order \(A\) at \(E\), for every jump. Indeed the \(\xi_j\) are \(\Lambda\)-linear combinations of the projected almost fixed vectors, and the chart group fixes \(\Lambda\).\formalnote{14.116}{116}

For the leading-coefficient argument, suppose \(\xi+Ev\) has order \(o(u)\le A\) at \(E\), with \(v\) standard and \(\xi\) almost fixed as above. If \(o(u)>0\), then \(\xi\) has order at least one on this jump. Its form value is fixed modulo \(\mathrm{Fil}^A\) and hence congruent to a field constant there, by the fixed-value lemma. If \(A>0\), reduction modulo \(\mathrm{Fil}^1\) shows that constant is zero, since a nonzero element of the fixed local field has Hodge order zero. It follows that \(\xi\) already has order \(A\) on this jump. For \(o(u)\le0\) its standard order suffices. Subtract \(\xi\) and divide by \(E\) to reduce each target order by one. Induction through an expansion of length \(N+1\) gives order \(A\) for \(\xi_N\) wherever \(N<o(u)\), in particular for \(u<-N\). Its differences retain order \(A\) at all jumps.\formalnote{14.117}{117}

To identify its quotient, subtract \(p^{4\delta}[T^j]F(T)c_i\) from each \(q(\xi_j)\). The new coefficients are still almost fixed, and their expansion, including its remainder, equals zero. We may therefore set \(o=A\) for every jump of the trivial quotient. The strict inequality \(N<A\) and monicity give
\[
               w=q(\xi_N)-p^{4\delta}c_i
\]
of order \(A\) at \(E\), with all other standard orders. If \(p^\kappa\) kills the cokernel at \(\sigma^{-l}i\), Frobenius transport lifts \(p^\kappa E^{-A}w\) to a standard level-\(l\) vector \(v\). The point \(p^\kappa\xi_N-E^A v\) now projects exactly to \(p^{\kappa+4\delta}c_i\). At level zero its order is at least \(A\) for \(u<-N\), and its differences have order at least \(A\) for every \(u\). Below \(l\) the other orders are zero, and at and above \(l\) they are \(-u\). This proves the long-gap estimate used in the cup product.\formalnote{14.118}{118}

\subsection{The obstruction bound at one jump}

For the colength comparison, fix the unramified chart, an order system \(o\), a letter \(i\), a level \(a\), and a weight \(u\) at \(\sigma^{-a}i\). Let \(m_u\) be its multiplicity on the kernel lattice. Assume
\[
                         o(a,v)>o(a,u)\qquad(v<u\ \text{a jump}).
\]
First take any finite list of points \(x_j\) with orders \(o\) in this kernel lattice. The forms of type \(u\) modulo forms of type \(u-1\) are generated by \(r\le m_u\) forms \(\nu_1,\ldots,\nu_r\) with coefficients in the fixed local field, as proved in \S\ref{sec:jump-symbol}. This bound comes from sending a form to the residue of \(t^u\) times its value on the ordinary lattice at this letter. The resulting symbols are Hodge--Tate maps of weight \(u\). Their span has dimension at most the jump multiplicity, and the coordinates in a symbol basis are algebraic by Galois descent. The strict inequality on lower weights means that checking the leading terms of these \(r\) forms suffices to raise the order at \((a,u)\).\formalnote{14.119}{119}

Form a matrix over \(\mathbb C_p\) with one row for each point \(x_j\) and one column for each form \(\nu_m\). Its \((j,m)\)-entry is the residue of \(t^{-o(a,u)}(\nu_m)_a(x_j)\). If its rank is \(d'\), choose \(d'\) independent columns, then take a square minor in these columns with maximal absolute determinant. For \(d'=0\) all leading terms already vanish. Otherwise each row is a linear combination of the \(d'\) chosen rows with coefficients of absolute value at most one by Cramer's rule. This also holds in the full matrix because the chosen columns span the column space. Lift the coefficients along \(\theta\) to \(A_{\inf}\), then apply \(\varphi^a\) to compensate for semilinearity at level \(a\). Subtracting the resulting combinations raises just the prescribed order, and integrality preserves all the other orders. Apply the coordinate retraction to these identities. Its scalar coefficients now belong to \(S\), so write them as \(w_{jm}+Eb_{jm}\) with \(w_{jm}\in\Lambda\).\formalnote{14.120}{120}

If there are \(f m_u+1\) points, the vectors \((w_{jm})_{m=1}^{d'}\) satisfy a primitive relation over \(\mathbb Z_p\). Indeed \(\Lambda^{d'}\) has rank \(f d'\le f m_u\), and one may divide a nonzero relation by its coefficient of least valuation. In this combination the \(w_{jm}\) terms cancel, the \(E b_{jm}\) terms gain one order everywhere at level zero, and the retracted errors already gain one order at \((a,u)\). Thus the combination is \(w+z'\), where \(w\) has all level-zero orders raised by one and \(z'\) has the selected order raised by one.\formalnote{14.121}{121}

For a colength assertion impose two more hypotheses. First \(o(b,v)\ge-v\) on the jumps at all levels \(b\). Second, for every integer \(v\), the open interval \((o(0,v),o(0,v)+p^a)\) contains no multiple of \(p\). Let \(\mathcal X\) be the pole-free points of depth \(o\) over constants, \(\mathcal X^+\) those of depth raised at \((a,u)\), and \(\mathcal M\) the kernel points of orders \(o\). The quotient \(T=\mathcal X/(\mathcal X^++\mathcal M)\) maps onto the quotient of liftable constants. The one-step operator kills \(T\) by \(p^{c(a,o(a,u))}\). We bound the dimension of \(T[p]\) over \(\mathbb F_p\).

Given \(f m_u+1\) classes in \(T[p]\), write \(px_j=x_j^++m_j\). Apply the primitive relation just proved to the \(m_j\). In its combination \(x=\sum_j\beta_jx_j\), we obtain \(p x-w\) of the raised depth with \(w\) in the kernel and having all level-zero orders raised by one. To divide by \(p\), apply a cost-zero operator \(B\) to gain depth at level zero up to \(o(0,v)+p^a\). All intervening costs vanish by the open-interval hypothesis. In particular zero lies in none of those intervals. On each jump, the positive-order improvement for a value congruent to a constant then gives the same higher \emph{order} for \(Bw\). When \(o(0,v)+1>0\) use its depth, and when \(o(0,v)+1\le0\) the desired order is at most zero, so depth suffices up to that order. Hence \(Bw=E^{p^a}w'\) with \(w'\) in the kernel of orders \(o\). Because \(E=z+p\), we have \(E^{p^a}=E_a+p r\) in \(S\). Consequently
\[
                 p(Bx-rw')=B(px-w)+E_a w' .
\]
Both terms on the right have the raised depth. The first uses preservation of existing depths by \(B\), the second the additional order at \(E_a\). The factor \(p\) can be canceled in this depth test because \(p\) is a unit in \(B_{\mathrm{dR}}^+\) and \(Bx-rw'\) is already in the original localized lattice. Moreover \(Bx-rw'\) differs from \(x\) by an element of \(\mathcal M\). Thus the primitive combination is zero in \(T[p]\), giving \(\dim_{\mathbb F_p}T[p]\le f m_u\).\formalnote{14.122}{122}

Each successive quotient \(T[p^j]/T[p^{j-1}]\) injects into \(T[p]\) by multiplication by \(p^{j-1}\). Summation through \(j=c(a,o(a,u))\) bounds the colength lost on this step by \(c(a,o(a,u))\,f m_u\). This explains the need for the open-interval hypothesis when deducing the length bound from a bound on the exponent.\formalnote{14.123}{123}

\paragraph{Inputs to the frame calculation.}
Crystallinity and the two-sided jump bound give the frame via the invariant \(B_{\max}\)-basis, the integral period matrix and its change of basis over the chart bicommutant (\S\ref{sec:period-matrix}). The scalar calculation uses Weierstrass division on that chart, extension of its coordinates to \(B^+\), the residue theorem in its value field, and the fixed classes of \(B_{\mathrm{dR}}\) modulo a filtration step (\S\ref{sec:tate-constants}). The multiplicity bound uses the injection of symbols into the corresponding Hodge--Tate maps.\formalnote{14.124}{124}


\editionchapter{Galois attachment and change of rank}\label{ch:galois}

Start with a Hecke eigenform on a definite unitary group. To apply the Galois congruence argument of Chapter~\ref{ch:deform}, we need a continuous representation of the absolute Galois group whose Frobenius characteristic polynomials approximate the given Hecke polynomials. The coefficients are compared in \(\overline{\mathbb Q}_\ell\), with the precise approximation condition given in \S\ref{sec:galois}. We also need the prescribed local behavior of the representation at places above \(\ell\). We explain how the representations with these properties are first obtained in an odd-rank case and then extended to the ranks and weights used in the determinant comparison.

\section{The Galois representation for the congruence estimate}
\label{sec:galois}
Fix a coefficient prime \(\ell\) and an isomorphism \(\iota:\mathbb C\simeq\overline{\mathbb Q}_\ell\). Let \(G_E\) be the absolute Galois group of the CM field. Use the standard definite unitary group of positive rank \(n\) over the totally real subfield, with the Weil data of the construction. For an eigenform let \(P_w(X)\) be its spherical Hecke polynomial at a finite place \(w\) where this is defined. In the applications to denominators there are two tasks. One must recover a global continuous trace from the \(P_w\), and one must control the action on local decomposition groups. For a cohomological conjugate self-dual cuspidal representation of \(\mathrm{GL}_n\) over a CM field, Chenevier and Harris construct a continuous semisimple Galois representation. At a place above \(\ell\) where the automorphic representation is unramified, it is crystalline and its linear crystalline Frobenius has the predicted characteristic polynomial \cite[Theorem~3.2.3 and the theorem in the introduction]{lit-chenevier-harris}. Our starting data are Hecke eigenvalues on a definite unitary group. We specify below the weaker Frobenius approximation and annihilation properties used for those data.

\subsection{The sense in which Hecke data determine a trace}
A Frobenius datum consists of a place \(w\) of prime norm, a prolongation to the algebraic closure, and a geometric Frobenius element \(s\) there. Call a statement true for most such data if, inside every nonempty open subset of \(G_E\), there is a nonempty open subset \(O\) such that, away from finitely many \(w\), it holds for every datum with \(s\in O\). This condition is stable under finite intersections by choosing the open subsets successively. Chebotarev gives data with \(s\) in the smaller open set outside any specified finite set of places.

The representation \(r\) used in the construction is continuous, with nondegenerate trace pairing on the algebra spanned by its image. Given \(\eta>0\), the coefficients of \(\det(X-r(s))-\iota P_w(X)\) all have \(\ell\)-adic absolute value at most \(\eta\) for most Frobenius data in the sense just defined. The qualifying set may depend on \(\eta\).

Suppose \(r,r'\) are two such representations. Their traces agree everywhere. If their difference at \(g\) had absolute value \(b>0\), continuity and the ultrametric inequality would give a neighborhood on which the difference always has absolute value \(b\). Impose both approximation properties with \(\eta<b\) and choose a Frobenius in the resulting open set. Comparison of the coefficients of \(X^{n-1}\) gives a difference of traces of absolute value at most \(\eta\), a contradiction. The same argument works coefficient by coefficient. Nondegeneracy excludes a radical in the image algebra, since a radical element times any element lies in the nilpotent radical and has trace zero. Thus one can reason about the representation using the multiplicities of its simple constituents over \(\overline{\mathbb Q}_\ell\).\formalnote{15.1}{1}

\subsection{Removing the extra characters from odd-rank lifts}
Assume first that attachment has been proved in odd rank. To deduce it for an eigenform of rank \(n\), take two theta lifts of the same form to rank \(2n+1\) using different twisting characters. Section~\ref{sec:odd-lifts} constructs the lifts and proves their nonvanishing. The added \(n+1\) characters in the two lifts can be chosen with no character shared between the two lists. After the normalizing twists, denote the two attached odd-rank representations by \(\rho_0,\rho_1\) and the extra Galois characters by \(\psi_{ij}\) for \(1\le j\le n+1\). For most Frobenius data the Hecke polynomial of the \(i\)-th lift factors as the original degree-\(n\) polynomial times the linear factors for the \(\psi_{ij}\). Comparing the coefficient of degree \(2n\) and using the preceding continuity argument proves, on all \(G_E\),
\[
      \operatorname{tr}\rho_0-\sum_j\psi_{0j}
       =\operatorname{tr}\rho_1-\sum_j\psi_{1j}.
\]
The disjointness of the character lists turns this virtual character into an actual one. For a simple constituent \(S\), let \(b_i(S)\) be its multiplicity in \(\rho_i\) and \(a_i(S)\) its multiplicity in list \(i\). Add the opposite list to each representation. Equal traces imply equal semisimplifications in characteristic zero, so \(b_0(S)-a_0(S)=b_1(S)-a_1(S)\). At least one \(a_i(S)\) is zero. The common difference is consequently nonnegative. Taking each \(S\) with that multiplicity gives an actual representation \(r\) of dimension \(n\), realized as a sum of constituents inside \(\rho_0\). It is continuous and has the desired trace.

The characteristic polynomial of \(\rho_i(g)\) is the product of that of \(r(g)\) and the polynomials of list \(i\). Division of a monic polynomial by a monic polynomial of fixed degree depends polynomially on the coefficients of the dividend and divisor. Their Frobenius approximations therefore descend to the original polynomial, using the bounded coefficients of the continuous representations on the compact Galois group. This gives attachment in rank \(n\).

Fix a split place \(w\) where the form is spherical, and let \(s\) be any geometric Frobenius lift at a chosen place above \(w\). For \(w\nmid\ell\) the local operator is \(r(s)\). For \(w\mid\ell\), write \(|k_w|=\ell^{f_w}\). For each compatible local embedding let \(j(s)\) denote the local image of \(s\). Use the linearized Frobenius on periods associated with the Weil pair \((j(s),f_w)\). Let \(P\) be the Hecke polynomial of the original form at \(w\), with coefficients transported to \(\overline{\mathbb Q}_\ell\). The two lifts give annihilating polynomials \(P R_0\) and \(P R_1\) for the operator just described. The characters are chosen so that their extra factors \(R_0\) and \(R_1\) at \(w\) are coprime. Take \(A_0R_0+A_1R_1=1\) and multiply by \(P\). Evaluating the identity shows that \(P\) annihilates the operator. The crystalline and Hodge--Tate clauses descend through the corresponding direct-sum and character decompositions.\formalnote{15.2}{2}

\subsection{The input to this descent}
The following calculation explains how the degree of a killing polynomial forces the dimension in the odd-rank seed. Start with a continuous Galois representation of dimension \(h\). Suppose a monic polynomial \(P_w\) of degree \(n\) kills the action of most Frobenii. Suppose also that the linear span of the image in its original basis contains each diagonal matrix. Let \(k\) be an extension of the coefficient field in which we can choose a Hodge--Tate basis, and write \(\rho\) for the matrices of the representation in that basis. The Hodge input gives integers \(\kappa_1,\ldots,\kappa_h\) with the following property. Whenever a polynomial in matrix entries vanishes on the images of all elements sufficiently near \(1\), it vanishes at \(\operatorname{diag}(T^{\kappa_a})\) in Laurent polynomials. Among the \(\kappa_a\) are \(n\) integers \(e_1,\ldots,e_n\) such that two monomials of equal total degree at most \((n+1)^2\) and equal \(e\)-weight have identical exponents. We show that \(h=n\) and that \(P_w\) equals the characteristic polynomial for most Frobenius data, in the sense of density defined above.

For a square matrix \(y\), vectorize \(I,y,\ldots,y^{m-1}\) and take all \(m\)-row minors of the resulting \(h^2\)-by-\(m\) matrix. Call them the minors of order \(m\). Over a field they all vanish exactly when \(y\) has an annihilator of degree less than \(m\). The order-\((n+1)\) minors vanish at the stipulated Frobenii and hence everywhere by continuity and density. Whenever the order-\(m\) minors vanish throughout a neighborhood of \(g\), apply the torus property to the polynomials giving the minors of \(Y\rho(g)\) near \(Y=I\). They vanish at
\[
  \operatorname{diag}(T^{\kappa_a})\rho(g).
\]
For \(g=1\), the number of independent powers of that diagonal matrix over the Laurent fraction field equals the number of distinct diagonal entries, by Vandermonde. Thus there are at most \(n\) distinct \(\kappa_a\). All the distinct \(e_i\) occur, so \(\kappa_a=e_{c(a)}\) for a surjection \(c:\{1,\ldots,h\}\to\{1,\ldots,n\}\).

Fix an invertible matrix \(x\) and \(m\le n+1\) such that the order-\(m\) minors of \(\operatorname{diag}(T^{e_{c(a)}})x\) vanish. Replace \(T^{e_i}\) by independent variables \(X_i\) and set \(y=\operatorname{diag}(X_{c(a)})x\). Each minor is homogeneous in \(X_i\) of total degree \(0+1+\cdots+(m-1)\), since the \(j\)-th power is homogeneous of degree \(j\). At this degree the substitution \(X_i=T^{e_i}\) separates the monomials. Vanishing after substitution thus implies vanishing in the polynomial ring \(S=k[X_1,\ldots,X_n]\). Over its fraction field, the minimal polynomial \(p(t)\) of \(y\) has degree less than \(m\) and divides its characteristic polynomial. Both monic factors \(p,q\) belong to \(S[t]\). Indeed their roots are integral over \(S\), hence so are their coefficients, and \(S\) is integrally closed.

As a polynomial in \(t,X_1,\ldots,X_n\), the characteristic polynomial has total degree \(h\) and \(t\)-degree \(h\). Both degrees add in its factorization \(pq\), and the total degree of each factor is at least its \(t\)-degree. Equality holds in each factor. Also \(p(0)\ne0\), since \(\det y\ne0\). Row \(a\) of \(y\) is divisible by \(X_{c(a)}\), so the \((a,a)\) entry of \(p(y)=0\) implies \(X_{c(a)}\mid p(0)\). Every variable therefore divides \(p(0)\). Its degree is at least \(n\) and at most \(\deg_t p<m\le n+1\). In particular \(m=n+1\) and \(p(0)=uX_1\cdots X_n\) for a nonzero constant \(u\).

Write \(p(t)=t^2s(t)+p_1t+p_0\) and consider \(p(y)x^{-1}=0\). The \((a,b)\) entry of \(y s(y)y x^{-1}\) is divisible by \(X_{c(a)}X_{c(b)}\), since \(yx^{-1}\) is diagonal. If \(a,b\) lie in the same fibre, its squarefree coefficient at \(X_1\cdots X_n\) vanishes. The term \(p_1yx^{-1}\) is diagonal and agrees at diagonal positions from that fibre. Taking the squarefree coefficient shows that the \((a,b)\)-block of \(x^{-1}\) on each fibre is a scalar matrix, since the remaining contribution is \(ux^{-1}\).

Apply this conclusion to \(x=\rho(g^{-1})\) for every \(g\). If a fibre contained two indices \(a,b\), every matrix in the span of the conjugated image would have scalar compression to these two coordinates. This holds in particular for the conjugates of the original diagonal matrix units \(E_{jj}\). Each has rank one, so its two-coordinate compression has rank at most one. A scalar \(2\)-by-\(2\) matrix of rank at most one is zero. Summing in \(j\) would give zero as the compression of the identity, a contradiction. Consequently every fibre is a singleton, and \(h=n\).

Finally the order-\(n\) minors cannot all vanish on any nonempty open set. At an interior point the torus argument with \(m=n\) would give \(n<m\). Their simultaneous vanishing is closed, so every nonempty open set contains an open subset on which some such minor is nonzero. On the corresponding Frobenii the minimal polynomial has degree at least \(n\). Both \(P_w\) and the characteristic polynomial are monic annihilators of degree \(n\), and hence both equal the minimal polynomial. This completes the algebraic step from killing polynomials to characteristic polynomials.\formalnote{15.3}{3}

\subsection{The geometric construction of the initial representation}
In the restricted odd-rank case the eigensystem has a cuspidal type at a split place. Its Hodge jumps satisfy a specific independence condition. Two lists of multiplicities with the same total at most \((n+1)^2\) and the same weighted sum of the jumps must coincide. The eigensystem gives a Hecke eigenquotient of the test-ring cohomology of the Kuga--Sato model, a fibre power of the abelian family with endomorphisms and level structure. Section~\ref{sec:cohomological-frobenius} explains how the special-fibre congruence gives annihilating polynomials on that cohomology. Hodge jumps and trace relations are then used to cut the representation to the predicted rank. The extension from restricted weights uses congruences between eigenforms and limits of traces. Quadratic base change and descent extend the construction over the required CM fields. Section~\ref{sec:seed-proof} derives the Hecke eigenquotient, its Hodge jumps and the cocharacter inclusion from the stated geometric comparisons. Section~\ref{sec:galois-limits} gives the trace-limit argument used to extend this construction. Section~\ref{sec:odd-lifts} gives the two compatible automorphic lifts used in the character subtraction above.\formalnote{15.4}{4}


\section{The odd-rank eigensystem in Shimura cohomology}
\label{sec:seed-proof}
Start with an eigenform on a definite unitary group of odd rank. Its spherical Hecke eigenvalues specify polynomials at the unramified places. We first realize the eigensystem on a quotient of the cohomology used to construct the odd-rank Galois representation. A fixed-point calculation compares traces and detects the eigensystem. An identity of correspondences on the special fibre then gives the annihilating polynomial on Frobenius. The Hodge-class calculation detects the Hodge jumps needed in \S\ref{sec:galois}.

\subsection{The places and the central character}
Let \(E\) be a CM field containing \(K=\mathbb Q(\sqrt{-d})\), where \(d>0\) and \(d\equiv3\bmod4\). Assume that \(E\) contains a primitive \(5040\)-th root of unity. Let the odd rank be \(n\ge1\), let \(\kappa\) be the weight, and fix \(\iota:\mathbb C\simeq\overline{\mathbb Q}_\ell\). The nonzero form \(f\) on the standard definite unitary group is spherical away from a finite set \(S\) in the sense that the spherical eigenvalues \(t_{w,j}\) are specified at split pairs with both places outside \(S\). At a marked \(w_0\in S\) it transforms under the first principal congruence subgroup by a cuspidal character. This local hypothesis means that for each proper nontrivial idempotent \(e\), the character is nontrivial on some group element of the form \(1+eY(1-e)\). The place \(v\) induced by the fixed \(\ell\)-adic embedding and its conjugate lie outside \(S\), with \(v\) over the same place of \(K\) as \(w_0\). For the place of \(K\) below \(w_0\) we require an integer \(b\) such that \(h=(1+\sqrt{-d})/2-b\) belongs to its prime ideal and \(h,\bar h\) are coprime in \(\mathcal O_K\). This is the split-prime condition used in the PEL construction. Fix a further place \(w'\) at which the local conclusion is sought. Its restriction to \(K\) has the same splitting property, and if \(w'\mid\ell\) its restriction to \(K\) must agree with that of \(w_0\). The conclusion at \(w'\) assumes both \(w',\bar w'\notin S\).\formalnote{15.5}{5}

Fix a distinguished complex embedding \(\tau_0:E\hookrightarrow\mathbb C\). Choose a CM type containing \(\tau_0\) in which every other embedding restricts to the conjugate of \(\tau_0|_K\). The construction of polarized abelian varieties with prescribed endomorphisms and level structure (the PEL construction), described in Chapter~\ref{ch:shimura}, gives an auxiliary integer \(N\) avoiding \(w_0,w'\). One chooses a further split place \(w_1\) away from the ramification and this integer, where \(f\) is spherical. The level \(\mathfrak N\) can be chosen neat, deep enough for \(f\), contained in \(w_1\) and outside \(\bar w_1\), with the auxiliary pro-Iwahori condition at \(w_1\). A place \(w\) is called safe for this level when the place of \(K\) below it has the prescribed splitting and \(\mathfrak N\) lies neither in \(w\) nor in any place over the conjugate place of \(K\). The level construction keeps \(v\) safe and also keeps \(w'\) safe under the stated hypotheses. Here neatness gives the freeness needed for the fixed-point counts.\formalnote{15.6}{6}

Extend the coefficient action on the definite unitary group to an algebraic representation \(\Xi_{\rm alg}\) of \(\mathrm{GL}_n(E)\), using the convention that identifies the unitary variable \(u\) with \(u^{-T}\) on frames. On the finite adelic cosets
\[
 (y,u)\longmapsto y\,u^{-T}K(\mathfrak N)
\]
put \(\phi(yu^{-T})=\Xi_{\rm alg}(y)f(u)\), and put \(\phi=0\) off their image. For equal cosets represented by \((z,t)\) and \((zy,u)\), the definite frames and the similitude condition at level imply \(y\) is rational unitary. Its finite representative \(c\) satisfies \(t^{-1}cu\in\operatorname{Stab}(f)\). Thus \(f(t)=f(cu)=\Xi_{\rm alg}(y)f(u)\). This proves independence of representatives. The function is right level-invariant, is rationally equivariant, and recovers \(f\) on the definite group. At a split place with \(\mathfrak N\) prime to \(\bar w\), the local frame identity expresses each translate at a fixed argument as the evaluation of a fixed linear map on the corresponding local translate of \(f\). On the displayed cosets the map is evaluation after a fixed translation followed by \(\Xi_{\rm alg}(y)\); it is zero off the cosets. Local linear relations on translates therefore carry over, including the spherical and cuspidal eigenrelations.

The central translates of \(\phi\) by the finite ideles of \(K\) span a nonzero finite-dimensional space of definite forms. These translations commute, so over \(\mathbb C\) they have a common eigenvector \(\psi\ne0\). They are invertible, and their eigenvalues form a multiplicative character \(\omega\). Writing \(\psi\) as a linear combination of central translates carries the same local linear relations from \(f\) to \(\psi\). Central ideles acting in the level fix \(\psi\), so the kernel of \(\omega\) contains an open subgroup. On principal ideles the eigenvalue is given by the action of the rational central scalar in \(\Xi_{\rm alg}\). This supplies the infinity type of an algebraic Hecke character. The norm of a local unit at a safe place acts at level, hence its base change to \(E\) is unramified there. If \(b_w=\omega(\operatorname{Norm}(\varpi_w))\) with the uniformizer idele convention of the source, the eigenvalues on the frame side are
\[
                   \theta_{w,j}=b_w^{-j}t_{w,j}.
\]
Indeed each local Hecke representative for a \(j\)-plane contains the same central factor with exponent \(-j\), while its remaining factor acts by the unitary Hecke sum.\formalnote{15.7}{7}

\subsection{Detection of a Hodge weight by the definite trace}
Write \(d_\infty=[E:\mathbb Q]/2\). For a highest weight \((\lambda_1,\ldots,\lambda_n)\) the jump list at a chosen embedding consists of \(j-1-\lambda_j\); at its conjugate the entries are \(\lambda_j+n-j\), \(1\le j\le n\). Let \(s_0:K\to\mathbb C\) be the restriction selected by the \(\ell\)-adic embedding. Choose \(m\) so large that all the alternant exponents \(l_j=m+z_j\) are nonnegative, with \(z_j\) the entries of the jump list in the ordering of the Weyl determinant. Set \(\mu=mnd_\infty\). The algebraic representation \(\Xi_{\rm alg}\) is formed using the embedding into \(\mathrm{GL}_n(\mathbb C)\) or its conjugate inverse transpose according to the restriction to \(K\), together with a determinant twist. Weyl's determinant formula gives its trace at \(y^{-1}\) as a product over all embeddings. At embeddings restricting to \(s_0\) the factor is the Schur alternant with exponents \(l_j\); at the other ones it is the \(m\)-th power of the determinant. One verifies the shift by multiplying row \(i\) of the alternant by \(r_i^m\), which adds \(m\) to every exponent. Inversion and reversal of the ordered exponents give the stated conjugate jump list. The central exponent of this extension at \(s_0\) is \(\mu\).

Suppose \(\tau_0|_K=s_0\). For regular \(y\) let \(r_1,\ldots,r_n\) be its distinct eigenvalues at \(\tau_0\). Write
\[
 D(r)=\det(r_i^{j-1})_{i,j=1}^n,\quad
 A(r)=\det(r_i^{l_j})_{i,j=1}^n,\quad S_l(r)=A(r)/D(r).
\]
Marking eigenline \(i\) multiplies \(r_i\) by \(X\) in the Lie contribution at \(\tau_0\). Write \(r^{(i)}=(r_1,\ldots,Xr_i,\ldots,r_n)\). The fixed-point weight also carries the Lagrange factor \(L_i(Xr_i)\) with \(L_i(t)=\prod_{j\ne i}(t-r_j)/(r_i-r_j)\). Since \(D(r^{(i)})/D(r)=L_i(Xr_i)\), expanding the numerator determinant gives
\[
 \sum_i L_i(Xr_i) S_l(r^{(i)})
    =\frac{\sum_i A(r^{(i)})}{D(r)}
    = S_l(r)\sum_{j=1}^n X^{l_j}.
\]
For the last equality, in the determinant summand indexed by \(\pi\), replacement of row \(i\) multiplies that summand by \(X^{l_{\pi(i)}}\). At \(\bar\tau_0\) all but the marked line acquire a factor \(X\). At each other chosen embedding all \(n\) lines acquire it; at their conjugates none do. The contributing determinant factors give the common multiplier \(X^{m(n-1+n(d_\infty-1))}\). Hence taking coefficient \(h=m(n-1+n(d_\infty-1))+l_j=z_j+\mu\) in the sum over markings gives \(\operatorname{tr}\Xi_{\rm alg}(y^{-1})\) times the positive multiplicity of \(l_j\) among the exponents. This identifies the exact shift in the transfer.\formalnote{15.8}{8}

Here is what it means for \(\theta\) to occur in this weight. For every finite polynomial \(p\) in words of the safe Hecke operators with \(p(\theta)=1\), some left Hecke multiple \(b p\) has a nonvanishing marked fixed-point test in weight \(h\). Equivalently it is impossible that \emph{every} \(bp\) be \(h\)-null. For regular test elements, \(h\)-nullity means the degree-\(h\) weighted orbit sums vanish for every permitted even multiplication parameter. The definite trace formula and the equality of marked and definite orbit counts for functions of cuspidal type at \(w_0\) turn \(h\)-nullity into vanishing of the trace on definite forms. In this passage one uses polarization in the independent multiplication parameters to test the Schur and determinant polynomials above. Summing over markings and dividing by the positive multiplicity then gives exactly the definite trace.

Indeed \(p\psi=\psi\). Its operators are supported trivially at the auxiliary \(w_1\). The local test-function construction gives a Hecke element \(b\) with every positive power of \(bp\) regular and of the required cuspidal type, and with \((bp)^k\psi\ne0\) for all \(k\). If every left multiple of \(p\) were \(h\)-null, so would \((bp)^k=((bp)^{k-1}b)p\), \(k\ge1\). All positive power traces of \(bp\) on definite forms would vanish. Newton's identities make its characteristic polynomial a power of the variable; Cayley--Hamilton would then give \((bp)^k\psi=0\) for large \(k\). This contradiction proves occurrence in every required weight \(z+\mu\) at \(\tau_0\). Section~\ref{sec:marked-orbits} derives the equality of weighted orbit counts for the marked and definite frames and explains the local tests used for this occurrence argument.\formalnote{15.9}{9}

\subsection{One polynomial relation on every cohomology group}
Take the smooth proper Kuga--Sato base \(Z\) from the integral model, with the number of auxiliary factors chosen for the preceding weight test, and its level covers. Its unitary similitude group has signature \((1,n-1)\) at the distinguished place and is definite elsewhere. The ring of definition is obtained from \(\mathcal O_E\) by inverting the elements outside every safe prime. Safe primes survive in this localization. Conversely, a nonzero prime \(\mathfrak p\) which is not safe meets the multiplicative set: a positive power of \(\mathfrak p\) is principal by finiteness of the ideal class group, and a generator belongs to no other prime. Thus all the surviving nonzero primes are good for the PEL data. The safe place \(v\) ensures that this ring is not a field. The construction of \(Z\) uses the closure of the theta-null points with its abelian family and the reduced closures of graphs of translations for the covers.\formalnote{15.10}{10}

At safe \(w\) put \(q=\# k_w\). A Hecke leg \(u_A\) is indexed by a subspace \(A\le k_w^n\). Write \(j=\dim A\). Fix a point reducing in the special fibre, passing to an arc over a valuation ring as in the model. Among its labelled \(w\)-torsion sections there is a nonzero vector \(t\) whose section specializes to zero. Put \(W=k_w t\). The local isogeny indexed by \(A\) has \(w\)-part of its characteristic-zero kernel labelled by \(A\). After reduction the section map is constant on cosets of \(W\). For each reduced section value \(a\), let \(N_A(a)\) be the number of vectors in \(A\) with that value. Since \(A\) and \(A+W\) have the same image and each fibre of \(A+W\to (A+W)/W\) has size \(q\), the proportion \(N_A(a)/q^j\) depends only on \(A+W\). This follows directly by comparing the kernels of the section map on \(A\) and on \(A+W\). Consequently the weighted counts \(q^{n-j}N_A(a)\) have the same dependence. The multiplier on the isogeny after the \(q^{n-j}\)-Frobenius twist is \(q^j q^{n-j}=q^n\), also independent of \(A\). The geometric rigidity input identifies the Frobenius-twisted targets, including their level data and marked sections, when these counts and the common multiplier agree. Thus the twisted reduced point depends only on \(A+W\). Section~\ref{sec:pel-covers} proves existence of the vanishing torsion section and the rigidity implication by testing the reduced kernel over local algebras.\formalnote{15.11}{11}

For this point set \(c(A)=(-1)^j q^{j(j-1)/2}\). For every \(D\ge W\), of dimension \(d'>0\),
\[
 \sum_{A+W\le D}c(A)=\sum_{A\le D}c(A)
                     =\prod_{i=0}^{d'-1}(1-q^i)=0.
\]
The middle evaluation follows by expanding \(\prod_{i<d'}(1+q^i T)\) and setting \(T=-1\). Its \(T^j\)-coefficient is \(q^{\binom j2}\) times the number of \(j\)-planes: intersecting with a fixed hyperplane gives the recurrence
\(\genfrac{[}{]}{0pt}{}{d'}{j}_q=\genfrac{[}{]}{0pt}{}{d'-1}{j}_q
+q^{d'-j}\genfrac{[}{]}{0pt}{}{d'-1}{j-1}_q\).
For \(D\not\ge W\) the first sum is empty. Subtracting the sums for proper subspaces by induction gives \(\sum_{A+W=D} c(A)=0\) for every \(D\). Hence the weighted sum of the twisted points vanishes as a zero-cycle, even when several different \(D\)'s give the same point. This explains both the sign and the exponent in the Hecke polynomial.\formalnote{15.12}{12}

Write \(V_i\) for the geometric \(\ell\)-adic cohomology of \(Z\), \(0\le i\le2\dim Z\). Its definition uses the functor of points \(Z_{\bar E}(-)\) on test rings. Such a ring \(D\) is reduced, every nonconstant monic polynomial has a root, and any two elements are comparable for divisibility after partition by an idempotent. Let \(\mathcal W_\ell(D)\) be the \(\ell\)-typical Witt vectors of the locally constant functions \(\operatorname{Spec}D\to\mathbb F_\ell\). In the category of abelian-group-valued functors on these rings, the group is
\[
 V_i=\mathbb Q_\ell\otimes_{\mathbb Z_\ell}
   \operatorname{Ext}^i\!\left(\mathbb Z[Z_{\bar E}(-)],\,\mathcal W_\ell\right).
\]
Functoriality in points gives pullback by correspondences and the action of \(G_E\) on these groups. An identification with conventional \'etale cohomology would require a separate comparison. The argument uses the base-change, period and fixed-point properties of the displayed groups.\formalnote{15.13}{13}

At sufficiently deep level the pullback \(u_1^*\) of the identity leg is injective by transfer, and the Hecke action satisfies
\[
              u_1^*T_{w,j}=\sum_{\dim A=j}u_A^* .
\]
This determines \(T_{w,j}\). Its Galois equivariance follows by checking on the right and using injectivity. Suppose first that \(w\nmid\ell\) and \(s\) is geometric Frobenius there, so \(s^{-1}\) acts on the residue algebraic closure by \(q\)-power Frobenius. The smooth and proper base-change results for the test-ring cohomology carry the pointwise congruence to the relation
\[
    \sum_{j=0}^n(-1)^j q^{\binom j2}T_{w,j}\rho_i(s)^{n-j}=0
                           \quad\hbox{on }V_i.
\]
Indeed one may first pull the proposed relation back along \(u_1\). Substituting the Hecke legs gives the pullbacks of the preceding twisted points. On the special fibre absolute Frobenius acts trivially on the constant \(\ell\)-adic coefficients; the surviving Galois action on the generic fibre is the inverse of arithmetic Frobenius. The comparison and injective pullback then give the displayed relation. Inertia acts trivially by good reduction.

At \(w\mid\ell\), the crystalline period theorem for the same cohomology gives the corresponding identity on \(B_{\max}\otimes V_i\) with the period-theoretic linear Frobenius \(\beta_{s,f}\) instead of \(\rho_i(s)\), where \(q=\ell^f\). Its \(a\)-th power is the operator for the Weil pair \((s^a,af)\). Both comparisons require the integral tower at \(w\). Section~\ref{sec:cohomological-frobenius} derives the linear Frobenius relation from the special-fibre points at the coefficient prime.\formalnote{15.14}{14}

\subsection{Taking the quotient and detecting its jumps}
Put \(k=\overline{\mathbb Q}_\ell\) and \(V=\bigoplus_i k\otimes_{\mathbb Q_\ell}V_i\). With \(a=(w,j)\) running over the safe Hecke generators and \(\vartheta_a=\iota(\theta_{w,j})\), form
\[
 U=V\bigg/\sum_a(T_a-\vartheta_a)V,\qquad F_i:V_i\longrightarrow U .
\]
The subspace in the denominator is stable under Galois because the actions commute. We obtain a continuous representation \(r\) on \(U\). For continuity take a basis from the vectors \(F_i(x)\); their translates \(F_i(\rho_i(g)x)\) depend continuously on \(g\). Every family of maps sending \(T_a\) to \(\vartheta_a\) factors through \(U\), and the images of the \(F_i\) span it. In particular the polynomial
\[
   P^\theta_w(X)=\sum_{j=0}^n(-1)^j q^{\binom j2}
                  \iota(\theta_{w,j})X^{n-j}
\]
annihilates \(r(s)\) at \(w\nmid\ell\): on a spanning vector it is \(F_i\) applied to the preceding relation. At \(w\mid\ell\) the images of the invariant periods of each crystalline \(V_i\) span after base extension because the \(F_i\) are equivariant and their scalar multiples generate. Hence the quotient is crystalline. Naturality of \(\beta_{s,f}\) gives the same annihilation on its period module.

The existence of this quotient alone would allow \(U=0\). To see what forces the prescribed jumps, fix a coefficient embedding \(\gamma:k\hookrightarrow\mathbb C_\ell\) and a nonnegative occurring weight \(h\) corresponding to the geometric embedding of \(E\). Suppose \(h\) is absent from \(U\). On the finite-dimensional space of families of Hodge--Tate forms of weight \(h\) on the \(V_i\), the commuting \(T_a\) act by precomposition. A simultaneous eigenvector for the eigenvalues \(\gamma(\vartheta_a)\) would factor through a \(\gamma\)-semilinear form on \(U\). Equivariance and the spanning property of the \(F_i\) make the resulting form of weight \(h\); absence of the jump forces it to vanish. Thus no such simultaneous eigenvector exists.

There is then a polynomial \(c\) in finitely many Hecke operators which annihilates this space of forms and takes value \(1\) at the character \(\vartheta'\) given by \(T_a\mapsto\gamma(\vartheta_a)\). To see this, use the finite-dimensional commutative algebra generated by the operators over \(k\). Joint generalized eigenspaces are indexed by finitely many characters \(\lambda\); each nonzero such space contains a joint eigenvector by simultaneous triangularization. For each \(\lambda\) choose an index \(a_\lambda\) at which it differs from the prescribed conjugate character \(\vartheta'\). Raising \(T_{a_\lambda}-\lambda_{a_\lambda}\) to the dimension of the space kills its \(\lambda\)-part. Multiply these powers and divide by their nonzero value at \(\vartheta'\). The assertion includes the zero-dimensional case by taking \(c=1\). Coefficient embeddings used here identify \(\gamma|_k\) with the corresponding conjugation of the algebraic eigenvalues.

The same \(c(T)\) kills every cohomology vector of weight \(h\) in \(\mathbb C_\ell\otimes V_i\). Pair such a vector \(y\) with a Hodge--Tate basis of forms. Against a form of weight \(h\) move \(c(T)\) to the form by precomposition. Against a form of weight \(h'\ne h\) the pairing transforms on an open subgroup by a nonzero power of the cyclotomic character. Tate's vanishing of \(\mathbb C_\ell(a)\)-invariants for \(a\ne0\) makes that pairing zero. The basis separates vectors, proving the assertion.

Transport the coefficients of \(c\) back to \(\mathbb C\) along the same embedding. It takes the value \(1\) at \(\theta\). Choose a cover deep enough for each of the finitely many Hecke words of \(c\) and each degree \(i\). Write \(\lambda_h\) for the Hodge--Tate map on Hodge classes, which are defined using square-zero test algebras. For each such class \(x\) of degree \(q'\) on \(Z\) (including the prescribed multiplication pullback), the period theorem supplies a weight-\(h\) vector \(v\) in \(\mathbb C_\ell\otimes V_{q'+h}\) and a map \(\Theta\) to cohomology on pairs of tests, semilinear for the chosen embedding of \(\overline{\mathbb Q}_\ell\) into the coefficient field, satisfying
\[
                         \lambda_h(u^*x)=\Theta(u^*v)
\]
for all the legs \(u\) from this cover. Section~\ref{sec:hodge-vector} constructs both maps and shows that \(v\ne0\) exactly when the tested image \(\lambda_h(x)\ne0\). Move the words of \(c\) across these pullbacks. The sum with coefficients from \(b c\) becomes a sum with coefficients from \(b\) evaluated on \(c(T)v=0\). Thus the \(b c\)-combination kills every tested Hodge class under \(\lambda_h\).\formalnote{15.15}{15}

The Woods Hole implication in \S\ref{sec:woods-hole} now gives vanishing of the weighted sum of local terms at simple coincidences. At a coincidence with tangent polynomial \(P\) the contribution is its correspondence coefficient times \([X^h]P(X)/P(1)\), where \(P(1)\ne0\). Group the coincidences by marked classes. On an integral frame lattice \(\Lambda\), the endomorphism \(B=1-2\alpha y\) is the identity modulo \(2\), so its determinant on the underlying integral lattice is odd and nonzero. Multiplication by \(B\) identifies \(B^{-1}\Lambda/\Lambda\) with \(\Lambda/B\Lambda\). Thus the fibre count is \(|\det_{\mathbb R}B|\), which cancels the determinant in the denominator. The vanishing sum gives \(h\)-nullity for \(b c\), contradicting occurrence. Consequently the required weight appears among the jumps of \(r\) at this embedding.\formalnote{15.16}{16}

\subsection{Untwisting, combining embeddings and obtaining the cocharacter}
We now restore the normalization of the unitary group. At safe \(w\) the two polynomial normalizations satisfy
\[
                P^\theta_w(X)=\iota(b_w)^{-n}\,(\iota P_w)(\iota(b_w)X).
\]
Twist by the \(\ell\)-adic avatar of the base change of \(\omega\) with the reciprocity convention of the seed construction. The scalar on geometric Frobenius away from \(\ell\) is \(\iota(b_w)\). The formula above thus gives the original killing polynomial on the twist. The local avatar comparison gives the corresponding statement for the linearized Frobenius over \(\ell\). At an embedding inducing a safe place over \(\ell\), the conditions on the level \(\mathfrak N\) ensure that its restriction to \(K\) is \(s_0\); the central exponent there is \(\mu\). In the jump convention used here the avatar twist removes \(\mu\). Accordingly the weight \(z+\mu\) gives the required \(z\). In constructing the representation for \(\tau_0\), one imposes this occurrence only when \(\tau_0|_K=s_0\); the geometric list is empty otherwise. This suffices because of the restriction just noted. Away from finitely many places every prime-norm Frobenius datum in the density argument is safe.\formalnote{15.17}{17}

Perform the construction for each complex embedding. The finite direct sum of their representations carries all the required jump occurrences at once, and \(P_w\) kills each appropriate Frobenius. Pass to a semisimple representation with the trace of this direct sum. The local comparison for trace limits, applied to a constant trace family, retains good reduction, the killing polynomials and occurrence of the prescribed weights. Next take each distinct irreducible constituent once. This is a retract of the semisimple sum, so the killing relations persist. A jump occurs in a direct sum precisely when it occurs in one of its summands; dropping repetitions therefore retains occurrence. Over \(k\), the algebra spanned by the image on a multiplicity-free semisimple module is a product of full matrix algebras on the distinct constituents. Indeed its faithful action is semisimple and Wedderburn decomposition gives precisely these blocks. In a basis respecting them the span contains every diagonal matrix. The seed machinery now has both the diagonal-span hypothesis and the prescribed jump occurrences.\formalnote{15.18}{18}

For completeness, the cocharacter assertion used in the rank calculation follows from Tate's theorem on twisted invariants \cite[Theorem~2.2.7]{legacy-brinon-conrad}. Write \(\chi_\ell\) for the cyclotomic character. Let \(x_1,\ldots,x_M\in\mathbb C_\ell\) and assign integer weights \(k_i\). Suppose a polynomial \(H\) over \(\mathbb C_\ell\) vanishes on
\[
           (\chi_\ell(g)^{-k_i}g(x_i))_{i=1}^M
\]
for every \(g\) sufficiently near \(1\) in the local absolute Galois group. We claim that it vanishes on \((T^{k_i}x_i)_i\) over \(\mathbb C_\ell[T,T^{-1}]\). Choose an open subgroup \(U_0\) on which the given vanishing holds. For the finitely many monomials \(\alpha\) in \(H\), put \(u_\alpha=\sum_i k_i\alpha_i\) and \(x^\alpha=\prod_i x_i^{\alpha_i}\). Let \(B\) be the space of coefficient lists \(a\) such that
\[
             \sum_\alpha a_\alpha\chi_\ell(g)^{-u_\alpha}g(x^\alpha)=0
                          \qquad(g\in U_0).
\]
For \(h\in U_0\) apply \(h\) to this equality at \(h^{-1}g\). It gives invariance of \(B\) under \(a_\alpha\mapsto\chi_\ell(h)^{u_\alpha}h(a_\alpha)\). Put a basis of \(B\) in reduced row-echelon form. This action preserves the pivot columns. Divide a transformed row of pivot weight \(u\) by \(\chi_\ell(h)^u\) to restore its pivot entry to one. Uniqueness of reduced row-echelon form shows that every entry of the normalized row is unchanged. Thus an entry \(a_\alpha\) in it satisfies \(h(a_\alpha)=\chi_\ell(h)^{u-u_\alpha}a_\alpha\) on \(U_0\). Tate's vanishing forces \(a_\alpha=0\) unless \(u_\alpha=u\). Every row is therefore supported in one weight. Its defining equation at \(g=1\), multiplied by \(T^u\), proves its vanishing on \((T^{k_i}x_i)_i\). The coefficients of \(H\) belong to \(B\), which proves the claim by linearity.

Apply this to the entries of the inverse Hodge--Tate frame \(Q\) above \(v\), assigning to each entry in row \(i\) the frame weight \(k_i\). The frame identity writes the entries of \(Q r(g)\) as \(\chi_\ell(g)^{-k_i}g(Q_{ij})\) locally at \(1\), with \(g\) on the right acting through the local chart. The local group maps onto an open subgroup in that chart, so a polynomial vanishing on \(Qr(g)\) near \(1\) vanishes at \(\operatorname{diag}(T^{k_i})Q\) by the claim. This proves the required inclusion in the Zariski closure. The step uses existence of a Hodge--Tate frame for these local representations and the specified compatibility with the chart.\formalnote{15.19}{19}

At \(v\) the seed assumes the \(n\) prescribed jumps have no small additive relations in the precise sense of \S\ref{sec:galois}. That section now gives dimension \(n\) and equality with the characteristic polynomial for most Frobenius data. At the requested good \(w'\), a Hodge--Tate frame has \(n\) entries counted with multiplicity. The \(n\) distinct jumps of \(\kappa\) all occur, so their multiplicities are exactly one. This supplies the local conclusion of the seed. Removal of its restrictions on type, weights and field is a further step of the argument.

The geometric inputs used here include existence and good reduction of the abelian family on the theta closure, the deformation and theta-coordinate inputs in the construction of the covers and their isogeny rigidity in \S\ref{sec:pel-covers}, and the equality of the marked and definite orbit counts for the specified test functions. The cohomological passage uses smooth and proper base change for the test-ring groups \(V_i\), the crystalline-period Frobenius relation, the functorial realization of Hodge classes and the point and diagonal calculations in \S\S\ref{sec:hodge-vector}--\ref{sec:point-classes}, the Woods Hole vanishing implication and the local comparison for trace limits.\formalnote{15.20}{20}


\section{Weak base change in the quadratic descent}\label{sec:weak-transfer}
A definite eigenform over \(E\) supplies a polynomial at every spherical place used in the descent. We need a form over a finite extension with those same polynomials at specified split places. In \S\ref{sec:galois-limits} the extension will vary, and each one is reached through a tower of prime-degree steps. For one such extension the desired eigensystem is detected by a trace in which the Galois generator permutes tensor factors cyclically.

\subsection{Local data preserved by transfer}
Write \(E'/E\) for a cyclic extension of CM fields of prime degree \(p\) with compatible CM types and generator \(\sigma\). The step used here has an integral element of trace one and a finite-idele character \(\mu\) trivial on principal ideles with open kernel. It satisfies
\[
       \mu(z\bar z)=1\ \Longleftrightarrow\
             z\in E^\times N_{E'/E}(\mathbb A_{E',f}^{\times}).
\]
Assume the root-of-unity hypotheses on the fields in \S\ref{sec:galois-limits}. A place \(w\) is called matched if \(\mu\) is trivial on \(E_w^\times\), the places above it have the same residue norm as \(w\), and every nonidentity Galois element moves each of them. At such a place the local fields upstairs and downstairs are isomorphic. For the unramified-level condition one only needs the step to be unramified at \(w\) and \(\mu\circ N\) trivial on the local units upstairs. Let \(f\ne0\) be a definite algebraic modular form of weight \(\kappa\). The desired \(f'\) of pulled-back weight \(\kappa'\) is spherical above every such unramified place where \(f\) is spherical. At a matched \(w\) split over \(E^+\), every Hecke polynomial of \(f\) there is a Hecke polynomial of \(f'\) at every place above \(w\). In this application there is no additional exceptional set. In particular \(\mu\) and the splitting of \(w\) matter when we read a polynomial downstairs.\formalnote{15.21}{21}

Write \(j\) for base change of rational or finite adelic unitary points as appropriate. At a transfer level \(L\subset K\) upstairs, \(L\) is open and \(\sigma\)-stable, \(K\) is contained in the integral level and cocycles in \(K\) split: if \(s\,\sigma s\cdots\sigma^{p-1}s=1\) then \(s=k(\sigma k)^{-1}\) for some \(k\in K\). The downstairs levels are their inverse images under \(j\). The level comes from the stabilizer of \(f\). At unramified places where \(f\) is spherical choose the full integral group. At the remaining finitely many needed places choose deep enough principal congruence so the pullback fixes \(f\), the character kills the required local norms, and the tame congruence matching lemma applies. In terms of an integral adele prescribing congruence, one can take it equal to \(1\) on the first set and a sufficiently divisible integer on the other set, invariant under conjugation and \(\sigma\). The local norm-one splitting and matching of the indicators at this level are the arithmetic facts used here.

At each matched split pair fix one \(u\) above a choice of \(w\). A matching family of test-function pairs at these places starts with the two level indicators. In the split product, restrict a test function at the chosen factor to the diagonal copy downstairs; on the remaining upstairs factors use integral-level indicators. The local norm-root identity and the choice of one side of the conjugate pair give matching. If one applies the spherical operator for an \(i\)-plane at \(u\) on the chosen side and at \(w\) on the diagonal restriction, the residue subspace representatives are the same, via \(E_w\simeq E'_u\). On the other factors the translations act trivially on the data being restricted, or preserve their level indicators. Thus the pair remains in the matching family. The disjoint sets of places allow products with the other matching level functions. This proves the stability under products that will be needed to extract a simultaneous eigensystem.\formalnote{15.22}{22}

\subsection{The twisted trace and its finite multiplicities}

\begin{example}[The trace with a swap of factors]\label{ex:transfer-swapped-trace}
Let \(V=\mathbb C^2\), let \(S(v\otimes w)=w\otimes v\), and take
\[
 A=\begin{pmatrix}1&2\\0&1\end{pmatrix},
 \qquad B=\begin{pmatrix}1&0\\3&1\end{pmatrix}.
\]
On \(e_i\otimes e_j\) the diagonal coefficient of \((A\otimes B)S\)
is \(A_{ij}B_{ji}\). Summing gives \(1+6+0+1=8\), the trace of
\(AB\). The ordinary tensor product \(A\otimes B\) has trace
\(\operatorname{tr}(A)\operatorname{tr}(B)=4\).
\end{example}

Let \(V,V'\) be the finite-dimensional spaces of level-fixed forms at the two levels with weights \(\kappa,\kappa'\). Above each archimedean place of \(E\), the coefficient space for \(\kappa'\) is a tensor product of \(p\) copies of the downstairs space, indexed cyclically by the extensions of its embedding. Rotation of the factors defines \(s\) and \(s\rho'(\delta)=\rho'(\sigma\delta)s\). If \(\delta\,\sigma\delta\cdots\sigma^{p-1}\delta=j\gamma\), then
\[
                    \operatorname{tr}(\rho'(\delta)s)
                          =\operatorname{tr}\rho(\gamma).
\]
Indeed the diagonal coefficient of a cyclic shift times a tensor of matrices at a multi-index is the product of their matrix entries along the index cycle. Summing the indices successively contracts them to the trace of the ordered matrix product. Multiply over the archimedean places. The operator \(S\) on forms defined by \(S\psi(x)=s\psi(\sigma^{-1}x)\) preserves \(V'\) and intertwines translation by \(g\) with translation by \(\sigma g\).\formalnote{15.23}{23}

Let \(F,F'\) be a pair of matching test functions with operators \(T,T'\). Expanding the trace on the finite set of double cosets gives a sum over a class representative and a rational element. Each representative \(x\) is weighted by \(1/|\Gamma_x(L)|\), the reciprocal of the rational level stabilizer. The terms downstairs and upstairs respectively are
\[
\begin{aligned}
 H(y,\gamma)&=F(y^{-1}\gamma y)\,\operatorname{tr}\rho(\gamma),\\
 H'(x,\delta)&=F'(x^{-1}\delta\,\sigma(x))\,\operatorname{tr}(\rho'(\delta)s).
\end{aligned}
\]
These are finite sums on support by definiteness and the finite coset supports. A rational move \(a\) from \((x,\delta)\) to \((x_1,\delta_1)\) means \(x_1^{-1}a x\in K\) and \(a^{-1}\delta_1\sigma(a)=\delta\); downstairs drop \(\sigma\). For an endpoint at which every source admitting a move has nonzero term, the weighted count of moves to it in the level-\(L\) sum is \([K:L]\). At the contributing endpoints used below the matching condition and twisted \(K\)-invariance give this support property. In fact the translates of its level-\(K\) class which refine to level \(L\) have total inverse stabilizer mass \([K:L]/|\Gamma_x(K)|\), by orbit-stabilizer on \(K/L\). The rational moves to the same coarse class contribute the factor \(|\Gamma_x(K)|\) and their conjugacy equations specify which summand carries each move. The same argument gives \([j^{-1}K:j^{-1}L]\) downstairs.\formalnote{15.24}{24}

Here is the role of rationality in comparing these two counts. The global norm-root construction, applied to the local roots of the matching functions, assigns a rational unitary \(A(\nu)\) and an adelic correction \(B(\nu)\) to contributing norm elements. They are equivariant under rational conjugation, and for \(\nu=j\gamma\) the twisted norm of \(A(\nu)^{-1}\) is \(j\gamma\). Send the downstairs pair to
\[
          m(y,\gamma)=\bigl(B(j\gamma)^{-1}j(y),\, A(j\gamma)^{-1}\bigr).
\]
The finite factors agree along rational moves to its image by local matching and the weight factors agree by cyclic contraction. For the converse start with a contributing pair \((x,\delta)\) and put \(\nu=\delta\,\sigma\delta\cdots\sigma^{p-1}\delta\); then \(\delta\,\sigma(\nu)=\nu\delta\). The root identity gives \(A(\nu)\delta\) of twisted norm one, and conjugation by \(B(\nu)x\) (with \(\sigma\) on the right) puts it in \(K\). The norm-one splitting in \(K\) gives a \(k\in K\) for which it is the adelic coboundary of \(B(\nu)xk\). The Hasse principle for this unitary descent then gives a rational \(g\) with \(A(\nu)\delta=g(\sigma g)^{-1}\). Now \(g^{-1}B(\nu)xk\) and \(g^{-1}\nu g\) are \(\sigma\)-fixed and hence descend to a pair over \(E\). For each prospective downstairs pair the map \(a\mapsto(gj(a))^{-1}\) is the bijection between the corresponding rational moves. In one direction substitute the two fixed expressions in the level and conjugacy equations. In the other direction a move \(a'\) upstairs conjugates the twisted norms, and the root's conjugation rule gives \((a'g)^{-1}\) fixed, recovering \(a\) uniquely. The arithmetic ingredient in constructing \(A,B\) is the idele norm principle in the CM fields generated by the roots of each irreducible factor of the characteristic polynomial. The construction joins factors by the Chinese remainder theorem in the reduced polynomial algebra; definite rational unitary elements are semisimple.\formalnote{15.25}{25}

Double count moves from an upstairs summand to an image \(m(y,\gamma)\), weighting by both finite sums and their common term. Summing first upstairs gives its move multiplicity times the downstairs trace; summing first downstairs gives its move multiplicity times the upstairs trace. Thus
\[
     \operatorname{tr}(T'S)
         =\frac{[K:L]}{[j^{-1}K:j^{-1}L]}\operatorname{tr}(T).     \tag{BC1}
\]
The positive constant is independent of the test function. This independence permits polynomial combinations of the matching operators.\formalnote{15.26}{26}

\subsection{Extracting one form that works at every matched place}
Use the identity as the initial matching operator. In the finite-dimensional downstairs space consider the joint generalized eigenspace for the matched spherical eigenvalues specified by \(f\), for all operators for which it has such an eigenvalue. It contains \(f\), so the trace of the identity on it is its positive dimension. Upstairs include all split spherical operators where the level is maximal, in particular all places above every matched place at which \(f\) is spherical. These families commute on the respective fixed spaces.

There is a single polynomial \(P\) in finitely many \emph{paired} generators that projects downstairs onto the indicated generalized eigenspace and upstairs onto precisely the joint primary spaces with the specified eigenvalues on the chosen paired generators. For each other distinct tuple occurring in either finite-dimensional space, choose one generator distinguishing it from the specified tuple. On each chosen generator use the Chinese remainder theorem with powers larger than both dimensions to obtain a polynomial which is one on the desired primary part and zero on its other primary parts. Multiply these finitely many polynomials. Primary decomposition and commutativity prove the assertion even though the original set of generators may be infinite. Applying (BC1) to the monomials of \(P\), using stability of matching, gives
\[
                \operatorname{tr}\bigl(P(T')S\bigr)
                 =\frac{[K:L]}{[j^{-1}K:j^{-1}L]}
                     \dim\operatorname{im}P(T)\ne0.
\]
Refine the upstairs decomposition into simultaneous generalized eigenspaces for its full commuting family. At least one allowed block, of joint character \(\Theta\), has nonzero diagonal trace of \(S\). Whenever \(ST'_a=T'_{a'}S\), the image under \(S\) of that block lies in the primary space for \(T'_{a'}\) with value \(\Theta_a\): take a sufficiently high power of \(T'_{a'}-\Theta_a\) and commute it through \(S\). If \(\Theta_{a'}\ne\Theta_a\), the diagonal block is zero. Therefore \(\Theta_{a'}=\Theta_a\) for all such pairs. In particular a spherical value at the chosen \(u\) has the same value at all its \(\sigma\)-translates. This is why matching at a single chosen \(u\) above \(w\) gives matching at every place above \(w\). Inside this nonzero primary space choose a simultaneous eigenvector, by triangularization of a commuting family over \(\mathbb C\). Call it \(f'\). Its level gives the spherical assertion, and the matched eigenvalues with the same residue norm give the asserted equality of the Hecke polynomials. The open level makes the spherical eigenvalues at the remaining almost-all prime-norm places an eigenform system for \(f'\).\formalnote{15.27}{27}

For the quadratic family of \S\ref{sec:galois-limits}, membership of a decomposition group in the specified character kernel supplies the matched-place hypothesis, and membership of inertia gives unramifiedness under the splitting conditions stated there. That is all the transfer of forms required at this step. Sections~\ref{sec:tame-root} and~\ref{sec:global-root} prove the tame splitting, matching of level indicators and idele norm-root construction used in (BC1). The final descent of a norm-one cocycle uses the unitary Hasse principle. Section~\ref{sec:tame-labels} carries out the transfer for the nontrivial local types in the tame inertia argument.\formalnote{15.28}{28}


\section{Matching a congruence subgroup with its norm}\label{sec:tame-root}
The twisted trace in \S\ref{sec:weak-transfer} compares an indicator upstairs with its restriction downstairs. Two questions arise even for this simple test function. Can a norm-one cocycle be split inside the level? Can one choose a root that detects membership of a norm element in the level? The first concerns a particular matrix. The second requires one polynomial that works at every matrix with the given characteristic polynomial.

\subsection{Levels and the congruence in Hilbert's Theorem 90}
Keep the prime-degree step \(E'/E\) of \S\ref{sec:weak-transfer}, with generator \(\sigma\) of order \(p\). Let \(O=\widehat{\mathcal O}_{E'}\), \(R=\mathbb A_{E',f}\), and write \(\tau\) for the action on \(O\). A star on matrices denotes conjugate transpose. Take \(M\in O\) fixed by the Galois action \(\tau\) and by CM conjugation, and invertible in \(R\). To impose level at a union of places closed under both actions, let \(u\) be its characteristic idempotent and put \(m=uM,\ v=uM^{-1}\). Thus \(mv=u,\ uv=v\). The relevant subgroup of adelic unitary matrices is
\[
                L=\{g:v(g-1)\in M_n(O)\}.
\]
At a place in the union this is the integral level with \(g-1\) divisible by \(M\). At the other places the condition is vacuous. In particular \(1+u(g-1)\) is integral unitary for \(g\in L\), and equals \(1+mY\) with \(Y=v(g-1)\).

Here the level satisfies the following tameness condition at \(m\) and at \(0\). At a level \(b\), one has \(\tau b=b,\ \tau^p=1\), some element of \(O\) has trace congruent to \(1\bmod b\), and every prime ideal \(P\) for which \(\tau^i x-x\in P\) for all \(x\in O\), for some \(0<i<p\), must contain \(b\). The zero-level trace condition is an exact equality. We also have \(c\in O^\tau\) with \(c-\bar c\) a unit. The step supplies such a \(c\) from downstairs.\formalnote{15.29}{29}

Suppose \(a=1+mY\in{\rm GL}_n(O)\) and \(a\,\tau(a)\cdots\tau^{p-1}(a)=1\). We claim there is an invertible \(k=1+mW\) with \(a\tau(k)=k\). First reduce in \(O/(N!)\), a finite ring on which \(\tau\) acts. Tameness descends because a prime and all the congruences defining inertia pull back through the quotient, while the trace-one congruence pushes forward. These quotients recover \(O\) as their inverse limit. On columns or matrices set \(\mathcal T(x)=a\tau(x)\), so \(\mathcal T^p=1\). Write \(\mathcal S(x)=\sum_{i=0}^{p-1}\mathcal T^i(x)\). If \(e\) is a trace-one element modulo \(m\), then \(\mathcal S(e I)\) is fixed and is \(I\bmod m\). It is invertible in every local factor whose maximal ideal contains \(m\).

On the remaining local factors \(m\) is a unit. The factors are permuted in orbits of length \(1\) or \(p\). For a length-\(p\) orbit, a matrix supported in one residue factor with value \(I\) has sum \(\mathcal S\) invertible at all of them, since its transported values are invertible. For a fixed factor with residue field \(k_0\), all \(p\) actions on \(k_0\) are distinct by tameness. The invariant sums \(\mathcal S(tw)\) span \(k_0^n\) over \(k_0\). Indeed, if a linear functional \(\lambda\) annihilates them, then for every \(t\in k_0\) and \(w\in k_0^n\),
\[
                 \sum_{i=0}^{p-1}\tau^i(t)\,
                             \lambda(\mathcal T^i w)=0 .
\]
Distinct field automorphisms are linearly independent as functions. The term at \(i=0\) gives \(\lambda(w)=0\) for all \(w\). Select a basis among the sums. Lift its summands to the local ring and sum again. This gives fixed columns whose determinant is a unit. Combine the solutions using the sum of the local idempotents in each orbit, which is fixed by \(\tau\), retaining the first solution wherever \(m\) is a nonunit. The resulting \(k\) is invertible and has the required congruence. In the profinite ring impose the three equations \(a\tau k=k,\ k=1+mW,\ kk'=I\) in finite quotients. Any finitely many are solved in a common refinement. Compactness in the variables \(k,W,k'\) gives an exact solution.\formalnote{15.30}{30}

For unitary \(a\), set \(h=k^*k\). The equation for \(k\) implies \(\tau h=h\), and \(h-I\) is divisible by \(m\). The relative orthonormalization in \S\ref{sec:integral-witt} supplies \(g_0=I+V(h-I)\) with \(g_0^*hg_0=I\), fixed by \(\tau\). To see the fixedness, perform the commutative norm calculation on the characteristic algebra of \(h-I\) over the fixed coefficients. The algebra is profinite by its monic power basis, and the fixed coefficient ring has the separating element \(c\); the trace-one element for \(\tau\) also makes its fixed part an additive retract by \(x\mapsto\sum_i\tau^i(e x)\). The correction is a polynomial in \(h-I\) over fixed coefficients. Thus \(t=kg_0\) is unitary, \(a\tau(t)=t\), and \(t\equiv I\bmod m\). For \(s\in L\) with \(s\,\sigma(s)\cdots\sigma^{p-1}(s)=1\), take \(a=1+u(s-1)\). Its twisted norm is one by evaluating on the support of \(u\) and its complement. The solution \(t\) equals \(I\) off the support because \(m=0\) there. On the support \(s=t\,\sigma(t)^{-1}\), and the congruence \(t=I+mW\) gives \(t\in L\).\formalnote{15.31}{31}

We shall need the commutative consequence that every fixed unit congruent to \(1\bmod m\) is a norm of a unit with that congruence. In a finite quotient write \(H=\{v\in O^\times:v\equiv1\bmod m\}\). The one-dimensional argument gives \(\ker(N:H\to H)=\{v/\tau(v):v\in H\}\). Its size is \(|H|/|H^\tau|\), so the norm image has size \(|H^\tau|\) and equals the fixed subgroup. Apply compactness to the norm equation, an inverse and the congruence as above. It applies as well to each finite quotient of a finite free profinite algebra carrying a tame action, and gives the same norm conclusion in the algebra itself.\formalnote{15.32}{32}

\subsection{One polynomial root for a whole characteristic polynomial}
Choose \(g\in L\) with \(P=\det(XI-g)\) fixed by \(\sigma\). The case \(n=0\) has only the identity matrix and is satisfied by the constant polynomial \(1\). Suppose \(n>0\). Write \(Y=v(g-I)\), \(Q=\det(XI-Y)\), and \(A=O[t]/(Q(t))\), using \(t\) also for its class. The affine change of variables in the characteristic polynomial gives \(\tau Q=Q\): componentwise it is the monic polynomial with roots \(v(\lambda-1)\) for the roots \(\lambda\) of \(P\), with multiplicities. Unitarity gives
\[
       Y^*(I+mY)=-Y,\qquad I+mY=I+u(g-I).
\]
The second matrix is invertible over \(O\). In \(A\) the element \(h=1+mt\) is invertible because the determinant of its multiplication map is \(\det(I+mY)\). Indeed the companion matrix of \(Q\) and \(Y\) have the same characteristic polynomial, so these two determinants agree.

On \(A\) let \(\theta\) act by \(\tau\) on coefficients, fixing \(t\). Define \(\rho\) by conjugating coefficients and sending \(t\) to \(-t/(1+mt)\). To check that the substitution descends, represent this last fraction by a polynomial \(b(t)\) in the monic basis. Then \(b(Y)=Y^*\), whose characteristic polynomial is \(\bar Q\). For any polynomial \(b\), \(Q(X)\) divides \(\operatorname{charpoly}_{b(Y)}(b(X))\). At each field component with roots \(\lambda_j\) of \(Q\), the latter is \(\prod_j(b(X)-b(\lambda_j))\); the \(j\)-th factor is divisible by \(X-\lambda_j\), including repeated roots. Monic division gives the same divisibility over \(O\). Consequently \(\bar Q(b(t))=0\). Since \(\rho(h)=h^{-1}\), applying \(\rho\) twice fixes \(t\). It fixes coefficients after two applications, and commutes with \(\theta\) by the formulas.\formalnote{15.33}{33}

In \(A^\theta\), \(h\rho(h)=1\) and \(h\equiv1\bmod m\). The involution on this ring is tame at \(m\): \(c/(c-\bar c)\) has trace one, and \(c-\bar c\) is a unit in every nonzero residue field. The congruence form of Hilbert 90 in dimension one gives a unit \(T\equiv1\bmod mA^\theta\) such that \(h\rho(T)=T\). The ring \(A\) is finite free over \(O\). For its \(\theta\)-action, \(\theta^p=1\) since \(t\) is fixed. The trace congruence is inherited from \(O\), and any prime with trivial residual action contracts to one with the same property on \(O\). Thus \(\theta\) is tame at \(m\). Now apply norm surjectivity for \(\theta\) to \(T\). Write the result and its inverse in the monic basis as
\[
    Z=1+m z(t),\qquad Z^{-1}=1+m z'(t),\qquad
    \prod_{i=0}^{p-1}\theta^i Z=T,\qquad (1+mt)\rho(T)=T.
    \tag{TM1}
\]
Hilbert 90 for the involution produces the rightmost equation; norm surjectivity for the degree-\(p\) action gives the product preceding it.\formalnote{15.34}{34}

Read \(Z\) in the original variable as \(z_0(X)=1+m z(v(X-1))\). Since \(m=uM\), this polynomial also satisfies \((1-u)z_0=1-u\) coefficientwise. If \(B\) is any matrix of any size with \(P(B)=0\), then \(Q(v(B-I))=0\). At places with \(v\ne0\) this is the affine change of variables. Elsewhere \(v=0\) and \(Q(X)=X^n\), which vanishes at the zero matrix. Thus (TM1) can be evaluated at \(v(B-I)\). When \(B\) is unitary, the adjoint of this argument is its image under the same fractional substitution, so with \(N_z=\prod_i\sigma^i z_0\) we get
\[
                    (I+u(B-I))\,N_z(B)^*=N_z(B).             \tag{TM2}
\]
The required determinant condition also holds for every \(B\). The eigenvalues of \(z(v(B-I))\) are integral at each finite place, since its argument satisfies the monic equation \(Q\). The characteristic coefficients are therefore integral in the local field. Expansion of \(\det(I+m z(v(B-I)))\) gives \(1+M d\) with \(d\in O\), using \(m=uM\). The polynomial for \(Z^{-1}\) gives an inverse determinant of the same form. The choice of \(M\) below puts every such idele in \(\ker(\mu\circ N_{E'/E})\). This determinant statement, the polynomial identity off the support, and (TM2) for each unitary matrix with characteristic polynomial \(P\) give all the norm-root conditions.\formalnote{15.35}{35}

Write \(N_\sigma a=a\,\sigma(a)\cdots\sigma^{p-1}(a)\). Since \(L\) is \(\sigma\)-stable, \(N_\sigma\) maps \(L\) into itself. For \(k\in L\), the two factors \(k^{-1}\) and \(\sigma(k)\) belong to \(L\), so \(\mathbf 1_L(k^{-1}\delta\,\sigma(k))=\mathbf 1_L(\delta)\). The characteristic polynomial required for matching comes from \(j(x)\) with \(x\) in the support of \(\mathbf 1_L\circ j\), or from \(N_\sigma\delta\) with \(\delta\in L\). Both possibilities supply the representative \(g\in L\) used above whenever the polynomial is \(\sigma\)-fixed.

Now take a unitary \(\gamma\) with this polynomial and a unitary \(a\) satisfying the root-quotient equation \(a\,(u z_0(\gamma))^*=u z_0(\gamma)\). If \(\gamma\in L\), then \(v(\gamma-I)\) is integral, and \(z_0(\gamma)\) and its inverse are both \(I\bmod m\). On the support of \(u\), the quotient equation gives \(a=z_0(\gamma)(z_0(\gamma)^*)^{-1}\), so \(v(a-I)\) is integral there and zero elsewhere. Thus \(a\in L\). When \(j(x)=\gamma=N_\sigma a\), this implication and \(N_\sigma(L)\subset L\) give \(\mathbf 1_L(j(x))=\mathbf 1_L(a)\). When \(N_\sigma\delta=\gamma\) and \(\delta\in L\), they give \(\gamma\in L\), then \(a\in L\), and finally \(a^{-1}\delta\in L\). These are precisely the two support conditions.\formalnote{15.36}{36}

Finally, the level element \(M\) can be chosen once for all supports. Write \(H\) for the given open subgroup downstairs. At an unramified place \(w\) where \(H\) contains the full integral subgroup supported at \(w,\bar w\), the construction chooses \(M=1\) above \(w,\bar w\). More precisely, choose a finite set \(S_0\) such that \(H\) contains the integral points trivial on \(S_0\), with \(w,\bar w\) outside \(S_0\) whenever the full supported integral subgroup lies in \(H\). Take \(M=1\) above unramified places outside \(S_0\cup\bar S_0\), and \(M=NC\) elsewhere. Here unramified includes the requirement that \(\mu\) kill the norms of local units. Take \(C\) to be the absolute ideal norm of \((\sigma x-x)\) for an integral \(x\) moved by \(\sigma\). To check the inertia condition, let \(P\) be a prime of \(O\) modulo which \(\tau^i\) acts trivially for some \(0<i<p\). Inertia at its contraction to \(\mathcal O_{E'}\) contains \(\sigma^i\) and hence the entire prime-order group. Thus that contraction contains \(\sigma x-x\) and the nonzero integer \(C\). Since \(O/CO\) is a finite product of local integer quotients above divisors of \(C\), \(P\) is the kernel of reduction at the corresponding finite place. That place has full inertia and lies outside the unramified set, so \(M=NC\) vanishes in its residue field. It follows that \(uM\in P\) for every choice of support idempotent \(u\).

The trace-one element comes from the step itself. Away from the retained unramified places, openness of \(H\) and of the idele character kernel gives one integer \(N>0\) deep enough for the pullback level to lie in \(H\) and for \(\mu\) to kill the required norms. On the retained places the level is the full integral group and the character already kills local unit norms. For full support \(u=1\), the polynomial in (TM2) is the local input to \S\ref{sec:global-root}.\formalnote{15.37}{37}


\section{Making the norm root rational up to an adelic correction}
\label{sec:global-root}
In the trace comparison of \S\ref{sec:weak-transfer}, local matching provides a polynomial whose coefficients are finite adeles. Counting rational conjugacy classes requires a rational factor as well. We explain the adjustment and why it works simultaneously for matrices with the given semisimple characteristic data.

\subsection{The problem in a commutative algebra}
Let \(E'/E\) be a cyclic CM step of prime degree \(p\) with generator \(\sigma\). Write \(I_F=\mathbb A_{F,f}^{\times}\) for the finite ideles of \(F\), write a bar for CM conjugation, and identify global elements with their principal ideles. The step comes with a character \(\mu\) of \(I_E\), trivial on \(E^\times\), and the criterion
\[
         \mu(t\bar t)=1\quad\Longleftrightarrow\quad
                        t\in E^\times N_{E'/E} I_{E'} .
 \tag{NR1}
\]
Fix a \(\sigma\)-stable characteristic polynomial \(P\) of a \emph{rational} matrix \(\gamma_0\) unitary for the positive standard form over \(E'\). Let \(z\in\mathbb A_{E',f}[T]\) be a norm root for full support. We recall the two properties used here. When \(P(A)=0\), the determinant of \(z(A)\), for every matrix size, lies in the group of ideles whose norm to \(E\) is killed by \(\mu\). At the adelic image of \(\gamma_0\), writing \(N_z=\prod_{i=0}^{p-1}\sigma^i z\), the second property gives \(\gamma_0 N_z(\gamma_0)^*=N_z(\gamma_0)\). The first property is stronger than invertibility at \(\gamma_0\); it also applies to the companion matrix of each irreducible factor.\formalnote{15.38}{38}

Let \(P_{\rm red}=\prod_q q\) be the product of the distinct monic irreducible factors of \(P\); set \(L_q=E'[T]/q\), \(\xi_q=T\bmod q\). At every embedding into \(\mathbb C\) the eigenvalues of \(\gamma_0\) have absolute value one by unitarity. Thus conjugation of coefficients and \(\xi_q\mapsto\xi_q^{-1}\) define an involution on \(L_q\) that agrees with complex conjugation under every embedding. Indeed the roots of \(q\) in an embedding are carried by complex conjugation into their reciprocals in the conjugate embedding. The fixed field is totally real, and \(L_q\) contains the CM field \(E'\), so \(L_q\) is CM. Finite adelic evaluation at \(\xi_q\) identifies \(\mathbb A_{E',f}[T]/q\) with \(\mathbb A_{L_q,f}\). For this last identity one can take the power basis, tensor at each finite place, and use that the integral order it spans differs from the maximal order at only finitely many places after denominators are cleared. Chinese remaindering identifies the quotient by \(P_{\rm red}\) with the product of these algebras.\formalnote{15.39}{39}

Set \(u_q=z(\xi_q)\). Its norm to \(E'\) is the determinant of the multiplication matrix of \(u_q\) in the power basis. Hence \(u_q\) is an idele and \(\mu(N_{L_q/E}u_q)=1\). Put \(N_q=N_z(\xi_q)\). We have
\[
                            N_q=\xi_q\bar N_q .                 \tag{NR2}
\]
Choose a polynomial \(v\) representing \(T^{-1}\bmod P\), which exists since \(P(0)\ne0\). Put \(D(T)=T\bar N_z(v(T))-N_z(T)\). Since \(\gamma_0\) is unitary, the adjoint of \(N_z(\gamma_0)\) is \(\bar N_z(v(\gamma_0))\). Thus the norm-root identity gives \(D(\gamma_0)=0\). Write \(n=\deg P\). At each finite place, triangularizing \(\gamma_0\) over an algebraic closure shows that \(D\) vanishes on each root of \(P\). Every root of \(P\) has multiplicity at most \(n\), so \(P\) divides \(D^n\) at each place. Division by the monic \(P\) commutes with projection to each place, so the same divisibility holds over the finite adeles. Evaluating at \(\xi_q\) gives \(D(\xi_q)^n=0\). The finite adele ring of \(L_q\) embeds in the product of its local fields and hence is reduced. It follows that \(D(\xi_q)=0\), which gives \textup{(NR2)}.\formalnote{15.40}{40}

\subsection{A scalar norm calculation in one CM field}
Suppose a factor \(q\) is fixed by \(\sigma\). Then \(\sigma\) extends to \(g\) on \(L=L_q\) fixing \(\xi=\xi_q\) by acting on the coefficients. It has order \(p\), because it has that order on \(E'\), and \(E'\) and \(\xi\) generate \(L\). Put \(K=L^{\langle g\rangle}\), \(u=u_q\) and \(N=N_{L/K}u\). Equation (NR2) descends to \(I_K\) as \(N=\xi\bar N\). We show that there are \(a\in L^\times,\ c\in I_L\) with
\[
                          \bar c\,a u\,g(c)
               \quad\text{fixed by conjugation}.               \tag{NR3}
\]

First, for a CM field \(K\), if \(b=\xi\bar b\) is an idele with \(\xi\in K^\times\), some \(y\in K^\times\) makes \(yb\) self-adjoint. Indeed \(y=\bar t+\bar\xi t\) gives \(yb=\bar t b+t\bar b\) since \(\bar\xi b=\bar b\). Choose \(t=1\) unless \(\xi=-1\); in that case take a nonreal \(t\). Apply this to \(N\). The fixed finite idele \(yN\) comes from \(K^+\), hence can be written \(a_0 w\bar w\) with \(a_0\in K^\times\) and \(w\in I_K\). For this norm assertion, the reciprocity character for the quadratic extension \(K/K^+\) has an archimedean sign at every real place. Given a finite idele, prescribe signs at infinity with the needed product, match the signs by weak approximation and adjust by the resulting principal idele. The quadratic norm subgroup criterion now gives the finite idele norm.\formalnote{15.41}{41}

Take the norm to \(E\) of \(yN=a_0 w\bar w\). As \(\mu\) kills principal ideles and \(N_{L/E}u\), (NR1) puts \(N_{K/E}w\) in \(E^\times N_{E'/E}I_{E'}\). Under global reciprocity its Artin symbol therefore acts trivially on \(E'\). By norm functoriality, a representative of the Artin symbol of \(w\) over \(K\) already acts trivially on \(E'\), and it fixes \(\xi\in K\). It fixes \(L=E'[\xi]\), which is cyclic over \(K\). The cyclic norm subgroup criterion yields \(w=b_0N_{L/K}s\) for \(b_0\in K^\times,\ s\in I_L\). There is no infinite-place sign obstruction over \(K\), which is CM. Substitution gives
\[
     N_{L/K}(s\bar s/u)=y/(a_0 b_0\bar b_0)\in K^\times .
\]
This principal element is a local norm at the finite places by the displayed equation and at the infinite places since \(K\) is totally complex. The cyclic Hasse norm theorem gives \(a\in L^\times\) with this same norm. Then \(a/(s\bar s/u)\) is a finite idele of norm one. Hilbert 90 for finite ideles writes it \(c/g(c)\). Rearranging gives
\(\bar c\,a u\,g(c)=s\bar s\,c\bar c\), establishing (NR3). This calculation uses the cyclic Hasse norm theorem on a principal element only after producing its local norm representatives.\formalnote{15.42}{42}

\subsection{Cycling the factors and lifting the coefficients}
On the product over \(q\) define \(T_0(f)=z\,\sigma f\). A factor orbit has length one or \(p\). At a fixed factor the last subsection supplies rational \(a_q\) and an adelic \(c_q\) for which \(\bar c_q a_q(T_0c)_q\) is fixed. For a length-\(p\) orbit choose a representative \(q\) and put \(c_q=1\). Propagate the entries once around the orbit by the rule \(c_{\sigma q}=u_{\sigma q}\,\sigma(c_q)\), stopping before the return. Set the rational factor equal to one at the nonrepresentatives. At each of these places the desired self-adjoint expression is simply \(\bar c_{\sigma q}c_{\sigma q}\). On returning to the representative the propagated value is \(N_q\). By (NR2) and the elementary scalar adjustment above a rational \(a_q\) makes \(a_q N_q\) self-adjoint. Thus in all factors we have a rational unit \(a_q\), an adelic unit \(c_q\), and
\[
               X_q=a_q u_q(\sigma c)_q\ \text{invertible},\qquad
                         \bar c_q X_q=\bar X_q c_q .
\]
Chinese remaindering lifts the \(a_q\) to a single \(w\in E'[T]\) and the \(c_q\) to \(c\in\mathbb A_{E',f}[T]\). Write \(X=w z\,\sigma c\). Inverting \(X_q\) factor by factor and lifting also their inverses proves that \(X(\gamma)\) is invertible for every unitary \(\gamma\) killed by \(P_{\rm red}\). Set \(a=X(\gamma)(X(\gamma)^*)^{-1}\). All polynomials in \(\gamma,\gamma^{-1}\) commute here, so \(a\) is unitary. The identity \(\bar c_q X_q=\bar X_q c_q\) then gives \(a\,c(\gamma)^*=c(\gamma)\). This proves the simultaneous assertion even if \(P\) has repeated factors.\formalnote{15.43}{43}

For a rational unitary matrix \(\nu\) with characteristic polynomial \(P\), the positive definite form makes \(\nu\) normal under a complex embedding, hence semisimple. Thus \(P_{\rm red}(\nu)=0\). Each commuting factor of the invertible \(X(\nu)\) is invertible. The quotient \(w(\nu)(w(\nu)^*)^{-1}\) is rational, while the quotient for \(X(\nu)\) is adelic; these are the \(A(\nu)\) and \(B(\nu)\) in \S\ref{sec:weak-transfer}. The factors were chosen by characteristic polynomial, and polynomial evaluation and the adjoint commute with unitary conjugation. Both assignments consequently transport under rational conjugacy, as the orbit count requires. Section~\ref{sec:tame-root} constructs \(z\) for the level indicators. The subsequent global descent for unitary norm-one cocycles uses the Hasse principle.\formalnote{15.44}{44}


\section{Transfer of tame labels and Iwahori eigenvalues}\label{sec:tame-labels}
\subsection{Transferring the residue-field label}
Section~\ref{sec:sparse-weights} applies the following assertion. At the split place \(w\) let \(k=\mathbb F_q\). The label is a squarefree monic polynomial \(c\in k[X]\) with nonzero constant term, together with distinct determinant characters \(\chi_i\) in ranks \(t_i\). Write \(x_i\in k^\times\) for the residue-field element labeling \(\chi_i\) under the parametrization determined by the faithful character \(T\) of \(\overline k^\times\). None of the \(x_i\) is a root of \(c\). To read the values of \(T\) in the coefficient field of the Galois representation, apply the comparison embedding \(\iota_\ell\). Fix a place over \(w\) in the algebraic closure, with reduction into \(\bar k\) extending the given map to \(k\). Choose \(g\ge1\) such that all roots of \(c\) lie in \(\mathbb F_{q^g}\), and write \(H=q^g-1\). For an inertia element \(s\), take a uniformizer \(\varpi\) at \(w\) and \(\theta\) with \(\theta^H=\varpi\). The quotient \(s(\theta)/\theta\) is an \(H\)-th root of unity. Its reduction defines the tame exponent \(a\) modulo \(H\) by
\[
      T\bigl(\overline{s(\theta)/\theta}\bigr)=\exp(2\pi i a/H).
\]
Put
\[
 B(X)=c(X)\prod_i(X-x_i)^{t_i},\qquad
 Q_a(X)=\prod_{x\in\operatorname{Roots}(B)}(X-\iota_\ell(T(x)^a)),
\]
counting roots with multiplicity. Then \(Q_a\) annihilates the action of \(s\) on the attached representation at \(w\nmid\ell\), and its linearized action on periods at \(w\mid\ell\). The same representation has finite inertia away from \(\ell\), and is potentially crystalline with the given weight jumps at \(\ell\).

First consider an inert unramified prime-degree step \(E'/E\), with \(u\) the unique place above \(w\), residue extension \(k'/k\), and Frobenius \(\phi\). The Green character for \(c\), parabolically induced together with the \(\chi_i\circ\det\) in their indicated ranks, gives a character \(\Xi\) with a Shintani lift. Write \(\rho'\) for the upstairs representation, with character \(\Xi'\). The function needed on the Frobenius coset is the twisted trace
\[
 S(a\phi)=\operatorname{tr}(\rho'(a)J)=\Xi(b),\qquad
 J\rho'(a)J^{-1}=\rho'(\phi(a)),
\]
where \(b\) is a downstairs conjugate of \((a\phi)^{[k':k]}\). On each irreducible constituent this is Shintani's norm-character identity with its normalized intertwiner \cite[Theorem 1]{arith-shintani}. Shintani writes the norm factors in reverse order. To obtain the displayed convention, apply his formula to the transpose of \(a\) and transpose the resulting product. Explicitly, if \(R(a)I=I R(\phi(a))\) is Shintani's convention, take \(\rho'(a)=R(a^{\mathsf t})^{\mathsf t}\) and \(J=I^{\mathsf t}\). The two transpositions give a representation with the same character as \(R\), since each matrix is conjugate to its transpose. The base-change identity for the Macdonald parameter identifies \(\Xi'\) \cite[identity (1.1)]{arith-silberger-zink}. Restricting the tame parameter leaves the inertia eigencharacters fixed and partitions each orbit of labels into orbits for \(x\mapsto x^{|k'|}\). Thus the polynomial is the image of \(c\) in \(k'[X]\), with each \(x_i\) viewed in \(k'^\times\). For determinant characters the formula is the scalar norm identity. Parabolic induction preserves it by the induced-character formula on invariant flags and their graded pieces.
\phantomsection\label{route:26}\leanroute{For the finite residue group we invoke Shintani's norm-character theorem and its compatibility with restriction of tame parameters. The formalization constructs actual characters on the semidirect product lifting the Green characters. The matching argument uses their twisted traces only on the chosen Frobenius coset. Inflation to the compact subgroup and detection of the global eigensystem still require the matching argument below.}\formalnote{15.45}{45}

Inflate \(\overline{\Xi}\) on the maximal compact below, and \(a\mapsto\overline{S(a\phi)}\) on the maximal compact above, extending by zero elsewhere. Use the matching root for the two compact-subgroup indicators in \S\ref{sec:tame-root}. For a matched integral pair the twisted local norm reduces to \((a\phi)^{[k':k]}\), so the inflated functions agree. Nonintegral pairs have zero on both sides, and twisted conjugation by an integral element is conjugation in the semidirect product. The same root and coset condition therefore work for the new pair. The downstairs trace on a finite level space for the conjugate-character operator is a positive constant times the inner product of the two characters. This is positive on the space containing the given labeled vector by finite-group semisimplicity. The nonzero-trace argument of \S\ref{sec:weak-transfer} gives an upstairs eigenform in the image of the operator for \(\overline{S(a\phi)}\). Expanding \(\operatorname{tr}(\rho'(a)J)\) expresses \(S(a\phi)\) as a sum of matrix coefficients of constituents of \(\Xi'\). By Schur orthogonality convolution with its conjugate kills every other constituent. The central character sum for \(\Xi'\) has nonzero eigenvalue on every constituent of its support, and thus detects this nonzero eigenform. This proves that its label is the required one.\formalnote{15.46}{46}

\paragraph{Prescribed prime-degree steps.}
Let \(p\) be prime, and let \(\chi:E_w^\times\to\mathbb Q/\mathbb Z\)
have order \(p\) and kill the \((q-1)\)-st powers of units. The step
used here comes from a global character \(\nu\) of exponent \(p\),
trivial at \(w\), with \(\nu+\nu^c\) inducing \(\chi\) there by
reciprocity. We need this symmetrized character unramified both
above \(p\) and at places fixed by complex conjugation.

The prime-exponent Grunwald argument gives these conditions as
follows. Let \(\psi\) be the finite-idelic character supported at
\(\bar w\) obtained from \(\chi\) by conjugation. Choose a finite set
\(\mathcal T\) of places containing \(w,\bar w\) and the places
above \(p\), large enough to generate the ideal class group by its
prime classes. Let \(C_{\mathcal T}\) be the finite ideles that are
units outside \(\mathcal T\). Then
\(\mathbb A_{E,f}^{\times}=E^\times C_{\mathcal T}\). The
\(\mathcal T\)-unit group modulo its \(p\)-th powers is finite by
the unit theorem. For each class on which \(\psi\) is nontrivial,
choose a representative \(a\). It is still not a \(p\)-th power
in \(E(\mu_p)\), by taking norms of a putative root and using
\([E(\mu_p):E]\mid p-1\). Chebotarev in \(E(\mu_p,a^{1/p})/E\)
gives a place of prime norm away from \(\mathcal T\) and from ramification in \(E/E^+\) whose Frobenius
fixes \(\mu_p\) and moves \(a^{1/p}\). There \(a\) is not a local
\(p\)-th power. In \(C_{\mathcal T}\) take the open subgroup \(U\)
whose components at all these witness places are local \(p\)-th
powers. Then \(C_{\mathcal T}^{\,p}\subset U\) and \(\psi\) kills
\(E^\times\cap U\). Extend the induced character on the image of
\(U\) in \(\mathbb A_{E,f}^{\times}/E^\times\) by divisibility of
\(\mathbb Q/\mathbb Z\). It is continuous since its kernel contains
the open subgroup \(U\cap\ker\psi\), and it has exponent \(p\)
since it kills \(C_{\mathcal T}^{\,p}\). Global reciprocity gives
\(\nu\). The local restrictions agree with \(\psi\) at every
place in \(\mathcal T\) and on all units outside the witness
places. A nonsplit unramified place over \(E^+\) has even absolute residue degree. The prime-norm witnesses thus split over the real
subfield, so the character is unramified at every place fixed
by conjugation. It is also unramified above \(p\), since an
exponent-\(p\) character killing the \((q-1)\)-st powers of units
is trivial on the units when \(p\mid q\).
\cite[V, Theorems 3.23, 5.3 and 5.5]{arith-cft}.\formalnote{15.47}{47}

Put \(E'=\overline E^{\ker(\nu+\nu^c)}\). The extension
is cyclic of degree \(p\), since \(\chi\ne0\). Its kernel is stable
under the ambient complex conjugation \(c\). For any embedding of
\(E'\) into \(\mathbb C\), extend the embedding to \(\overline E\)
and call the resulting conjugation \(c'\). On \(E\), \(c'=c\)
because \(E\) is CM. Thus \(g=cc'\) fixes \(E\), and \(cgc=g^{-1}\),
so \((\nu+\nu^c)(g)=0\). This proves that \(c'\) and \(c\)
agree on \(E'\) under every embedding. The fixed field of \(c\)
is therefore totally real and \(E'\) is CM.\formalnote{15.48}{48}
There is only
one place \(u\) above \(w\) by the nontrivial local character.
Its norm subgroup is \(\ker\chi\), and \(u\) differs from its
conjugate because \(w\) does. The integral trace is onto,
using \(\operatorname{Tr}(1)=p\) away from \(p\) and the surjective
trace on residue fields at the unramified places above \(p\).
The norm criterion of \S\ref{sec:weak-transfer} is reciprocity for
\(\nu+\nu^c\). The Weil data lift by keeping \(j\), pulling the
splitting character back by the norm, and choosing a new idele
\(\alpha'\) so \(j\alpha'\overline{\alpha'}\) generates the inverse
different. At conjugate pairs prescribe the sum of the two
valuations. At a fixed place this sum can be halved, since the
relative different is trivial there and the valuation difference
downstairs was even. These are the step and Weil hypotheses
needed in the transfer.\formalnote{15.49}{49}

If \(g=p h\), take the unramified order-\(p\) character of \(E_w^\times\)
given by the valuation modulo \(p\), with sign from the reciprocity
convention. The step has local residue degree \(p\), ramification
index one and a single place \(u\) above \(w\). Inertia at \(w\) is
in the kernels of \(\nu\) and \(\nu+\nu^c\). Because \(q_u^h=q^g\), the
original uniformizer and its \((q^g-1)\)-st root compute the same tame
exponent upstairs. The label transfer leaves the roots and
multiplicities unchanged in the common algebraic closure. Apply
induction at level \(h\) to the representation attached upstairs,
then the restriction and descent assertions of
\S\ref{sec:galois-limits} to the original representation. This
descends finite type, weight jumps and the identical annihilator.
Each positive \(g\) reaches \(1\). The level \(n!\) is divisible by
all the residue degrees and works for every root. For an
individual inertial element its action on a \((q^g-1)\)-st root
of the uniformizer determines an exponent by reduction on roots
of unity. It remains to prove the split-label case.\formalnote{15.50}{50}

\subsection{The split-label case and the Iwahori eigenvalues}
At \(g=1\) the Green character for \(c\) is the induction product of its distinct rank-one characters. Finite-group Frobenius reciprocity supplies in the space of translates a vector transforming under the block parahoric by the prescribed determinant characters. Call this vector \(x\) and the local representation \(\rho_w\). The pro-unipotent Iwahori \(I_1\) fixes \(x\). On diagonal units \(x\) has the prescribed character, and permutation matrices that preserve that character act on \(x\) by scalars, since they lie in the block parahoric.

Let \(e\) denote averaging over \(I_1\) with Haar probability measure. Normalize \(|\varpi|=q^{-1}\). For \(t=\operatorname{diag}(t_1,\ldots,t_n)\) call \(t\) dominant when \(|t_1|\le\cdots\le|t_n|\). We seek a nonzero \(I_1\)-fixed vector \(v\) in the span of the translates of \(x\) and characters \(\xi_i\) such that
\[
   e\rho_w(t)v=D_\xi(t)v,\qquad
   D_\xi(t)=\prod_{i=1}^n \xi_i(t_i)|t_i|^{n-i}
\]
for every dominant \(t\), with \(|\xi_i(\varpi)|\) weakly decreasing. The \(I_1\)-fixed part of this span is finite-dimensional: the weight and level away from \(w\) stay fixed, and imposing \(I_1\) gives a finite level space. Its invariant scalar product comes from the automorphic inner product.

The dominant translations \(e\rho_w(t)\) commute on this fixed space and preserve the subspace on which diagonal units act as on \(x\). Choose a simultaneous eigenvector there. A strictly dominant translation is injective for a suitable choice of its diagonal entries, and multiplying by powers of that translation makes any diagonal element dominant. Hence the eigencharacter extends to the whole torus and has the form \(D_\xi\). In particular the restrictions of the \(\xi_i\) to units are the rank-one characters once each and the \(\chi_i\) with multiplicities \(t_i\).

For the ordering, use an eigenfunctional \(l\), so that \(l(\rho_w(t)e z)=D_\xi(t)l(z)\). The invariant scalar product and simultaneous triangularization give such a functional on the cyclic representation. It can be taken nonzero at \(x\). Indeed, a nonzero value on a permutation translate of \(x\) requires the permutation to preserve its unit characters, in which case the translate is a scalar multiple of \(x\). If all those values vanished, the Iwasawa and Bruhat decompositions would force \(l(\rho_w(s^r)z)=0\) for large \(r\) for each \(z\) in the cyclic span. Here \(s\) is the fixed strictly dominant element: conjugating the upper unipotent factor by \(s^r\) puts it in \(I_1\), and the remaining torus factor becomes dominant. On an \(I_1\)-fixed vector on which \(l\) is nonzero, the eigenrelation contradicts this vanishing.

At each stage retain the permutation of \(x\) on which the current functional is nonzero. For an inverted adjacent pair \(i,j=i+1\), let \(h\) exchange those entries, let \(w\) have diagonal entries \(\varpi\) up to position \(i\) and \(1\) thereafter, and put \(U_w=\rho_w(w)e\). The exchanged functional is
\[
        (l\circ\rho_w(h)\circ e)\circ(U_w-D_\xi(w)).
\]
Indeed \(U_w-D_\xi(w)\) takes values in the kernel of \(l\), on which the averaged reflection intertwines the two dominant characters. To check the eigenrelation, put \(U_t=\rho_w(t)e\) for dominant \(t\). Then \(eU_tU_w=eU_wU_t\), since the averaged translations commute on the image of \(e\). Thus one may apply the intertwining identity on \(\ker l\) after \(U_w-D_\xi(w)\). The rank-one double-coset relation shows that this functional is nonzero on the correspondingly permuted vector. If the two unit characters differ, vanishing forces \(\xi_i(\varpi)=\xi_j(\varpi)\). If they coincide, the reflection scales the original block vector, and the relation, with one identity coset and \(q-1\) nonzero residue classes, gives \(\xi_i(\varpi)=q\,\xi_j(\varpi)\). Both possibilities contradict the strict inversion \(|\xi_i(\varpi)|<|\xi_j(\varpi)|\). Repeating the exchanges orders the exponents. The resulting eigenfunctional is nonzero on the finite-dimensional fixed space, so its simultaneous generalized eigenspace contains an eigenvector \(v\) with the same character.

If all \(\xi_i\) are unramified, apply the invariant scalar-product construction once more to the sorted vector \(v\). It gives an eigenfunctional \(l=\langle u,\,\cdot\,\rangle\) of the sorted character, with \(u\) in the cyclic span of \(v\). This functional is nonzero on a spherical vector. To see this, suppose it vanishes on averages over \(\mathrm{GL}_n(\mathcal O_w)\). Successively refine the last diagonal block of this maximal compact by splitting off one more line, ending at the ordinary Iwahori. The parahoric coset relation propagates vanishing of the averages at each refinement: its only possible obstruction for an earlier entry \(m\) and a later entry \(j\) is \(\xi_j(\varpi)=q\,\xi_m(\varpi)\), excluded by the ordering. The functional is already invariant under the ordinary Iwahori, because it is invariant under \(I_1\) and the diagonal unit characters are trivial. This would make it zero. Averaging the vector representing the functional over the maximal compact therefore gives a nonzero spherical vector. Write \(g\) for this average of \(u\). For each spherical \(y\), the spherical Hecke formula gives \(\langle g,T_j y\rangle=\lambda_j\langle g,y\rangle\), where \(\lambda_j=q^{-\binom j2}e_j(\xi_1(\varpi),\ldots,\xi_n(\varpi))\) and \(e_j\) is the elementary symmetric polynomial. The spherical Hecke algebra is commutative and closed under adjoints, so \(T_j\) is normal. Thus \(T_j^*g=\overline{\lambda_j}g\) implies \(T_jg=\lambda_jg\). Its Hecke polynomial has the stated roots.

For comparison, Roche identifies the principal-series type and its Jacquet image \cite[Theorem 7.7 and Remark 7.8]{arith-roche}. Casselman's comparison of translations with the upper Jacquet action applies to ordinary Iwahori-fixed vectors in an admissible representation \cite[Propositions 2.4 and 2.5]{arith-casselman}. In the normalized induction convention \cite[\S1.5]{arith-hkp}, the corresponding unnormalized torus character is \(\delta_B(t)^{1/2}\prod_i\alpha_i(t_i)\), where \(\delta_B\) is the upper Borel modulus. Since \(\prod_i |t_i|^{n-i}=\delta_B(t)^{1/2}|\det t|^{(n-1)/2}\), this gives \(\xi_i=\alpha_i|\cdot|^{-(n-1)/2}\).
\phantomsection\label{route:27}\leanroute{The ordering, including exchanges of unequal unit characters, and the passage to the spherical vector follow the pro-unipotent Iwahori calculation in the formal proof. The cited comparisons explain the conventional normalization.}\formalnote{15.51}{51}

For \(x\in E_w^\times\) with \(\operatorname{ord}_w(x)=k>0\), let \(s\) have its local reciprocity class, with \(s^{-1}\) acting by \(q^k\) on the residue field. Put \(b(y)=(|y|,\xi_1(y),\ldots,\xi_n(y))\) and \(D=b^{-1}(\langle b(x)\rangle)\). The index is finite since the \(\xi_i\) kill principal units and the valuation quotient has size \(k\). Index one forces \(k=1\) and all \(\xi_i\) unramified, giving the spherical case above. Otherwise take a prime-order character of \(E_w^\times/D\). It is tame because \(D\) contains the \((q-1)\)-st powers of units. The preceding prescribed-step construction gives a cyclic extension with one place over \(w\) whose local norm subgroup is the character's kernel. Transfer of the Iwahori character gives \(\xi_i\circ N\) on a form upstairs. The element \(s\) fixes the step and lies in the kernel of the unsymmetrized character at \(w\). A representative \(y\) upstairs satisfies \(b(Ny)=b(x)\) by norm compatibility of local reciprocity. Its valuation is \(k/f>0\), where \(f\) is the local residue degree. The new subgroup \(N^{-1}D\) has index \([N(E'_u{}^\times):D]\), strictly smaller than \([E_w^\times:D]\). By induction and restriction-descent we obtain the finite-type and jump assertions, and \(\prod_i(X-\iota_\ell(\xi_i(x)))\) kills \(s\) on the stipulated space.\formalnote{15.52}{52}

For that transfer write \(I\) for the pro-unipotent Iwahori upstairs. For diagonal \(t\) with strictly ordered absolute values, twisted conjugation by \(I\) carries each \(g\in ItI\) to \(td\) with \(d\) diagonal and congruent to \(1\). Indeed put \(H_j=t^{-j}It^j\) and \(X(k)=t^{-1}k^{-1}g\sigma(k)\). Since \(g\in ItI\), some \(X(k)\) belongs to \(H_0\). If \(X(k)\in\bigcap_{0\le j\le b}H_j\), use upper/lower Iwahori factorization to write \(X(k)=u\,\sigma(z)\) with \(u\) upper unitriangular integral and \(t^{-1}\sigma(z)t\in I\). Then \(z\in H_{-1},\ldots,H_b\), and \(X(kz^{-1})=(t^{-1}zt)u\) belongs to \(H_0,\ldots,H_{b+1}\). To shift an interval \([-a,b+1]\) down by one, take \(y=tX(k)t^{-1}\in I\). Direct multiplication gives \(X(ky)=\sigma(y)\) in the shifted intersection, as \(\sigma\) preserves all root valuations. Compactness gives one \(k\) for all \(j\), and strict ordering forces every off-diagonal entry to vanish in \(\bigcap_jH_j\).

The entries \(N(t_i d_i)\) of the twisted norm are distinguished by valuation. In the downstairs coset for \(N(t)\), the diagonal corrections congruent to \(1\) lift by norm surjectivity on principal units of the tame extension, as in \S\ref{sec:tame-root}. For each resulting characteristic polynomial interpolate a polynomial taking \(N(t_i d_i)\) to \(t_i d_i\). These choices match the indicators of the two strictly dominant double cosets. In fact diagonalization in either coset gives precisely the listed norm eigenvalues; distinct valuations determine their order. The eigenvectors downstairs for the norm matrix may be chosen over the downstairs field: on each eigenline normalize an eigenvector to have \(i\)-th coordinate \(1\) in its designated column. This preserves the upper unitriangular residue condition, and the unique normalized vector is fixed under \(\sigma\). If two roots of this norm belong to the upstairs coset, twice applying the diagonalization shows that their diagonal corrections differ by principal units, so each lies in the same right \(I\)-coset as the interpolated root. The cocycle condition in \(I\) follows from an integral \(c_0\) of trace one. For \(s_0\) of twisted norm one, the sum \(\sum_j\sigma^j(c_0)s_0\sigma(s_0)\cdots\sigma^{j-1}(s_0)\) for one period, with empty product \(1\), is integral upper unitriangular modulo the maximal ideal and fixed by \(s_0\sigma\). This verifies the local conditions of \S\ref{sec:weak-transfer}. The dominant translations commute and multiply as diagonal elements after normalization by their double-coset degrees. Multiplication by a strictly dominant translation makes each dominant one strictly dominant. Its eigenvalue is nonzero, so the common-eigenspace projection in that section supplies all the characters \(\xi_i\circ N\) after cancelling the degrees.\formalnote{15.53}{53}

Finally let \(s\) be inertial and \(x\) a unit with its reciprocity class. Choose a geometric Frobenius \(F\) for \(\varpi\). Positive valuation gives the annihilating polynomial for \(s F^m\) with roots \(\iota_\ell(\xi_i(x)\xi_i(\varpi)^m)\) for every \(m>0\). It also holds at \(m=0\). Indeed Frobenius is annihilated by \(\prod_i(X-\iota_\ell(\xi_i(\varpi)))\) with nonzero constant term. Its powers are thus finite sums of \(m^j\lambda^m\) with nonzero scalar \(\lambda\) and fixed operator coefficients, by Jordan decomposition or scalar partial fractions. The terms in the polynomial evaluated at \(sF^m\) are products of such sequences. They satisfy a linear recurrence with nonzero trailing coefficient, so their sum vanishing at every \(m>0\) forces vanishing at zero. For periods use the multiplicative linearized Weil operators, which commute with the scalar coefficients. For a tame exponent \(a\), compatibility of reciprocity with the action on the \((q-1)\)-st root of \(\varpi\) gives \(\xi_i(x)=T(x_i)^a\) for the corresponding unit character, including the rank-one factors from \(c\). This proves annihilation in the base case and completes the induction.\formalnote{15.54}{54}


\section{The two theta lifts to odd rank}\label{sec:odd-lifts}
To subtract characters in \S\ref{sec:galois}, one needs two actual nonzero eigenforms of rank \(2n+1\). A factorization of a polynomial at each prime alone would give no object to which to apply the odd-rank Galois theorem. We construct the forms, including the characters by which their attached representations must be twisted.

\subsection{A theta integral with prescribed weight and level}
Let \(E\) have the CM type and Weil data of \S\ref{sec:galois}, let \(V,W\) be the standard definite hermitian spaces of ranks \(n\ge1\) and \(m=n+s\) with \(s>n\), and let \(f\ne0\) be an algebraic form on \(U(V)\) of weight \((\kappa,e)\). Choose an automorphic character \(b\) of \(U(1)\) of weight \(c\le e\) and trivial tame part. Thus on a norm-one idele integral at all finite places its value is \(\prod_v u_v^{c_v}\), using the selected complex embedding at each infinite place. Write \(\omega_{1,b}\) for the Weil representation of the pair \(U(V)\times U(W)\) with splitting characters \(1,b\). This is the action on the test function in the theta integral below. Write \(y\) for the adelic form corresponding to \(f\), with the source compact type selected at infinity. The candidate lift, in the conventions of the theta kernel, is
\[
    (L y)(h)=\int_{[U(V)]}
        \theta\bigl(\omega_{1,b}(g,h)(v_\infty\otimes\Psi)\bigr)y(g)\,dg .
\]
Here \(\theta\) is the theta summation functional, \([U(V)]\) is the quotient by the rational points, \(v_\infty\) belongs to the real Fock model and \(\Psi\) is the finite Schwartz datum.\formalnote{15.55}{55}

At each infinite place retain the diagram columns \(\kappa_v(h)\) for \(h<n\), add \(e_v-c_v\) columns of height \(n\), and put \(\kappa'_v(h)=0\) for \(h>n\). The target weight is \((\kappa',0)\). Indeed, write the source highest weight in zero-based coordinates as \(\lambda_p=e_v+\sum_{p<h<n}\kappa_v(h)\). In the target it is \(\lambda'_p=\lambda_p-c_v\) for \(p<n\) and \(0\) for \(n\le p<m\). Adding \(c_v\) to the first \(n\) entries therefore recovers the source weight, while the \(s\) additional slots have zero highest weight. Equivalently, for the doubled Harish-Chandra parameter \(2\lambda_p+n-1-2p\) the old block shifts by \(s-2c_v\) and the new entries are \(m-1-2p\).\formalnote{15.56}{56}

Here is the compact correspondence behind this rule. Remove the determinant characters of the vacuum and write \(a_i=\lambda_{i-1}-c_v\) for \(1\le i\le n\), \(a_{n+1}=0\). This is a partition since \(e_v\ge c_v\). In polynomial Fock coordinates on \(m\)-by-\(n\) matrices the simultaneous extremal vector is the product of nested minors
\[
                     \Phi_a(Z)=\prod_{i=1}^n
                         \det(Z_{[i],[i]})^{a_i-a_{i+1}} .
\]
Order rows and columns by the two extremal flags, with the dual flag on a factor carrying the contragredient action. Each diagonal coordinate in position \(i\) occurs with weight \(a_i\) on its side, since the exponents of minors containing it telescope. Triangular row and column additions in the extremal unipotents preserve the minors. Evaluation on the rectangular identity gives one. The Cauchy decomposition of the polynomial algebra pairs the two Schur modules indexed by \(a\), once each as a tensor product, using the dual on the contragredient side. It follows in degree \(d\) also by decomposing the two tensor powers into Schur modules times Specht modules of \(S_d\) and taking diagonal \(S_d\)-invariants. Two Specht factors contribute a line precisely when their partitions agree. The extremal lines in that summand are therefore the ones selected by \(\Phi_a\). Restoring the vacuum characters and dual flags gives the weight pair computed above \cite[\href{https://www.math.auckland.ac.nz/~obrien/howe.pdf}{Theorem 4.4.1}]{theta-howe}.

By compact Schur orthogonality the normalized projections onto the two extremal lines fix \(\Phi_a\). Its translates under the partner group span the partner type by irreducibility. The real integral of \S\ref{sec:real-pairing} is the squared norm of the nonzero projected vector divided by the source dimension. The product over the infinite places gives \(v_\infty\) and the nonzero factor \(Z_\infty\).\formalnote{15.57}{57}

Let \(H\) be the open right stabilizer of \(f\), \(K_0\) the stabilizer of the standard integral lattice, and \(\widehat{\mathcal O}=\prod_w\mathcal O_{E_w}\). Choose a finite set \(B\) such that integral elements trivial at \(B\) and its conjugates belong to \(H\), and an integer \(N>0\) with \(K(N)\subset H\). Take \(B\) to avoid the pairs of places at which \(H\) contains the local maximal compact. Put \(z=N^{-1}\) on \(B\cup\bar B\), zero elsewhere, and let \(S\) be the places at which \(z\) is nonintegral. There is a matrix \(x\in M_{m,n}(\widehat{\mathcal O})\) with \(x^*x=0\) whose first \(n\) rows are the identity. To construct it take two disjoint blocks of \(n\) rows and put \(I_n,\epsilon I_n\) in them. The integral finite adele \(\epsilon\) of norm \(-1\) furnished by the Weil data gives \(1+\bar\epsilon\epsilon=0\). Set \(\Psi\) equal to the finite Heisenberg translate of the vacuum by \(zx\). It has nonzero norm and is a vacuum off \(S\), since that translation is integral there.

Put \(J=H\cap\operatorname{Stab}(\Psi)\), with the stabilizer taken for the source Weil action. The finite matrix coefficient at \(g\) integral off \(S\) is
\[
   \langle\Psi,\omega_{1,b}(g,1)\Psi\rangle
      =\|\Psi\|^2\,{\bf1}_{J\cap K_0}(g).                         \tag{OL1}
\]
For the support calculation put \(D=g^{-1}-1\). In the integral-lattice model a nonzero vacuum overlap requires an integral matrix \(u\) with \(zxD+ug^{-1}\) integral. Reading the identity block of \(x\) and subtracting the integral term \(u\) gives \(zD_{pj}+\sum_k u_{pk}D_{kj}\) integral for every \(p,j\), using only those rows of \(u\). At a fixed place choose an entry \(D_{pj}\) of largest absolute value in its column. If its absolute value exceeds \(1\), the displayed integrality forces \(|z|\le1\): otherwise the first term strictly dominates all the others. But \(g^{-1}\) is integral at such a place by the assumption off \(S\), a contradiction. Hence \(D\) and then \(zD\) are integral. This gives \(g\in K_0\) by unitarity. Its component at \(B\cup\bar B\) is congruent to the identity modulo \(N\); the complementary component is in \(H\) by the choice of \(B\). Thus \(g\in H\). Finally the difference between the translated labels is \(zxD\), which is integral, and the Heisenberg phase is trivial since \(x^*x=0\). So \(g\) fixes \(\Psi\). On \(J\cap K_0\) the coefficient is its value at the identity, proving (OL1).\formalnote{15.58}{58}

\subsection{The nonvanishing test}
Normalize finite Haar measure by \(\operatorname{vol}K_0=1\), and put \(d=\|\Psi\|^2/[K_0:J\cap K_0]>0\). Write \({\cal R}y=\sum_\gamma b_{\rm vac}(\gamma) R(\gamma)y\), where \(R\) is right translation and \(b_{\rm vac}\) is the rank-\(m\) vacuum coefficient for \(b\). The representatives \(\gamma\) of the finite adelic group modulo the subgroup integral off \(S\) are chosen to be supported away from \(S\). Doubling for the vacuum and the constant form on a definite source gives an integrable matrix coefficient. On each coset of the everywhere integral level, the integral of its absolute value is the absolute vacuum coefficient times a fixed positive integral over the identity coset. The coefficients are therefore summable, also on the subset of cosets off \(S\) and after a unitary twist. For the forms of the chosen type the norm formula of \S\ref{sec:doubling-integral} gives
\[
   \langle Ly',Ly\rangle_W=A d\,\langle y',{\cal R}y\rangle_V,\qquad
   A=\operatorname{vol}[U(W)]\,Z_\infty\ne0 .                    \tag{OL2}
\]
Here \(Z_\infty\) is the nonzero real-place factor for the joint harmonic \(v_\infty\).\formalnote{15.59}{59}

Suppose \(Ly=0\). For a finite set \(T\) off \(S\), pair \(Ly\) with the lift of \(y\) for a second datum whose mixed Weil coefficient cuts out integrality at \(T\). Using the frame \(x\) above, apply the unitary shear \(1+u x x^*\), where \(u\) is skew, supported on \(T\cup\bar T\) and nonintegral there. Translate by the Heisenberg operator for the identity rows of \(x\) at these places and normalize by the vacuum multiplier. For \(g\) integral at \(T\), vacuum projection recovers the original coefficient. Conversely, a nonzero pairing forces the two Heisenberg characters to agree on their common stabilizer. At a place \(w\in T\cup\bar T\), let \(h\) be an entry of largest absolute value in a row of \(g_w^{-1}\), and choose \(\mu\) with \(|\mu|_w=\max(|u_w|_w,|h|_w)\). Testing integral multiples of the corresponding matrix unit scaled by \(\mu^{-1}\), self-duality gives \((1+u_w)h/\mu\) integral. If \(|h|_w>1\), its absolute value is \(\min(|u_w|_w,|h|_w)>1\). Thus \(g^{-1}\) is integral on \(T\cup\bar T\), as required by unitarity.

Doubling now makes each restricted integral zero. Let \(T\) exhaust the places off \(S\). Absolute integrability in the range \(m>2n\) permits passage to the limit. On points integral off \(S\), (OL1) gives the finite factor \(d\); the archimedean pairing gives \(Z_\infty\langle y,y\rangle_V\). Thus \(d Z_\infty\langle y,y\rangle_V=0\), forcing \(y=0\). The lift of our nonzero form is therefore nonzero.\formalnote{15.60}{60}

\subsection{The polynomial at a split place}
Where \(f\) is spherical, the finite datum is a vacuum at \(w,\bar w\). The local maximal compact on the target fixes it: the change of variables preserves the lattice and the splitting is trivial on these units. Equivariance of the theta integral makes \(F\) spherical there as well. Suppose also that \(w\ne\bar w\), put \(q=q_w\), and define \(\alpha=\upsilon(\varpi_w)^{-n}q^{n/2}\) and \(\beta=b(\varpi_w/\bar\varpi_w)\upsilon(\varpi_w)^{-m}q^{m/2}\), where \(\varpi_w\) is the idele with a uniformizer at \(w\) and \(1\) elsewhere. Here \(\upsilon\) is evaluated on the full idele class group. Assume further that \(f\) is a simultaneous eigenvector for the spherical Hecke operators at \(w\), with eigenvalues \(t_j\). Its reciprocal eigenpolynomial \(p(Y)\) has coefficients \((-1)^j q^{\binom j2}t_j\), with \(t_0=1\). Then the target reciprocal eigenpolynomial is
\[
   p'(Y)=\prod_{j=0}^{s-1}(1-q^j\alpha Y)\,
                    p(q^s\alpha\beta^{-1}Y).                      \tag{OL3}
\]
Indeed the split Hecke operator is the sum of its lattice modifications indexed by residue-field subspaces. On a residue matrix of rank \(r\) in the vacuum calculation, the eligible subspaces lie in a kernel of dimension \(m-r\) on the target side and \(n-r\) on the source side. With each dimension-\(j\) term weighted by \((-Y)^j q^{\binom j2}\), the finite-field subspace identity gives \(\prod_{j<d}(1-q^jY)\) for a \(d\)-dimensional kernel. Splitting the product of length \(m-r=s+(n-r)\) leaves the length-\(s\) factor and shifts the argument in the shorter product by \(q^s\). The local determinant factors change \(Y\) to \(\alpha Y\) on the long side and to \(\beta Y\) on the short side. On the Weil datum the short-side operator is at the conjugate place. Moving it across the invariant pairing replaces it by its adjoint at the original place, giving (OL3). Comparing coefficients on \(F\ne0\) identifies the individual Hecke eigenvalues.\formalnote{15.61}{61}

For integers \(k,j\), define the algebraic character \(\chi_{k,j}(x)=b(x/\bar x)\upsilon(x)^{-k}|x|^{j-k/2}\), with \(|\varpi_w|=q^{-1}\). Put \(a_w=\chi_{s,s}(\varpi_w)\) and \(\theta_j=\chi_{m,s-j}\) for \(0\le j<s\). Direct substitution gives
\[
       (q^s\alpha/\beta)a_w=1,\qquad
       \theta_j(\varpi_w)=q^j\alpha a_w .
\]
Let \(P_w(X)=X^np(1/X)\), \(Q_w(X)=X^mp'(1/X)\). The polynomial obtained by multiplying all roots of \(Q_w\) by \(a_w\) is therefore
\[
    a_w^m Q_w(X/a_w)
        =P_w(X)\prod_{j=0}^{s-1}(X-\theta_j(\varpi_w)).            \tag{OL4}
\]
Both sides are monic, so (OL4) includes the constant factor as well as the roots.\formalnote{15.62}{62}

\subsection{Choosing the two lifts and recovering their jumps}
Let \(d_E\) be the number of roots of unity in \(E\), and fix one infinite place \(v_0\). For \(i=0,1\) take
\[
     c_{i,v}=-d_E\bigl(\max(0,-e_v)+i\,{\bf1}_{v=v_0}\bigr).
\]
These weights are \(\le e_v\). A norm-one global unit is a root of unity by Kronecker's theorem \cite[Corollary~5.6]{legacy-milne-ant}, so each weight character is trivial on those units. It defines a character on the product of the archimedean norm-one torus and the integral finite subgroup, modulo rational elements, by making it trivial on the latter subgroup. It extends to an automorphic character \(b_i\): the subgroup is open and the circle is divisible, so one can extend over each generator of the finite quotient by choosing a root of its prescribed power. For these choices all cross-list values \(\chi^{(0)}_{k,j}(\varpi_w)\), \(\chi^{(1)}_{k,j'}(\varpi_w)\) differ at every split \(w\). If two were equal, taking their complex phases cancels the common \(\upsilon^{-k}\) and the positive norm factors. Thus \(b_0/b_1\) would take value \(1\) on \(u=\varpi_w/\bar\varpi_w\). Finiteness of the norm-one class quotient gives \(h>0\) and a rational norm-one element \(\gamma\) such that \(u^h\gamma\) is integral and a unit at all finite places. Evaluating the weight character on it gives \(\Phi_{v_0}(\gamma)^{d_E}=1\), since the other weight differences vanish. Consequently \(\gamma\) is a root of unity. Its valuation cannot cancel the nonzero valuation of \(u^h\) at \(w\). This proves the separation at the individual place that is used for local polynomial cancellation.\formalnote{15.63}{63}

Apply (OL4) to the given eigenform with each twisting character \(b_i\) and \(s=n+1\), so \(m=2n+1\). The odd-rank theorem attaches a representation \(\sigma_i\) to \(F_i\). Twist it by the \(\ell\)-adic avatar of \(\chi^{(i)}_{s,s}\) to obtain \(\rho_i\). The character's value at unramified geometric Frobenius scales each root by \(a_w\), just as in (OL4). At a spherical place the constructed characters are unramified as Hecke characters, hence their avatars have the local unramified or crystalline property at the appropriate residue characteristic. Their infinity types are \(j+c_v-k\) on the chosen embedding and \(j-c_v\) on its conjugate for \(\chi_{k,j}\), and their jumps are the negatives of these exponents.

To identify every jump of the twisted representation, observe that the twist shifts the jumps of the rank-\(m\) weight by \(-c_v\) at the chosen embedding and by \(c_v-s\) at its conjugate. At index \(p<n\) the two resulting jumps are \(p-\lambda_p\) and \(\lambda_p+n-1-p\), exactly the original ones. At \(p=n+a\) for \(0\le a<s\) they are \(n+a-c_v\) and \(c_v-1-a\). These are the jumps of \(\theta_a\) on the chosen embedding and of \(\theta_{s-1-a}\) on its conjugate. The added characters thus account for every additional jump with multiplicity. The two cross-disjoint strings give the trace subtractions and the relatively prime polynomial factors used in \S\ref{sec:galois}.\formalnote{15.64}{64}

The underlying inputs used here are the joint-harmonic compact theta correspondence, the convergent Siegel--Weil doubling identity proved from the reduction and local Weil estimates in \S\ref{sec:doubling-integral}, the dual finite datum cutting out integrality, and class-field-theoretic compatibility for algebraic characters. The odd-rank Galois representation itself enters only after the two nonzero algebraic forms have been constructed.\formalnote{15.65}{65}


\section{Approximation with prescribed weight gaps}\label{sec:weight-density}
\subsection{Approximation with the prescribed weight gaps}
In the approximation step of \S\ref{sec:galois-limits} one must keep the weights fixed at some embeddings while allowing only weights with no short additive relations at the other embeddings. We prove the density assertion at the varying embeddings. Let \(X\) be compact, let \(k=\overline{\mathbb Q}_\ell\), and let \(z:X\to\prod_{\tau}\mathrm{GL}_n(L)\) be continuous for a finite extension \(L/\mathbb Q_\ell\) inside \(k\). The finite product ranges over the varying embeddings. In the application the entries lie in the completions of the fixed number field at the places above \(\ell\), embedded in \(k\). Their compositum is finite. The values of the form in its rational translates also lie in a single finite extension: the values at finitely many level classes and the finitely many embedded coefficient fields generate one. A continuous scalar function factoring through \(z\) gives a continuous function on its compact image by the quotient map property.\formalnote{15.66}{66}

At a varying place use Young diagrams with \(c_h=(n+1)^{4(n-h)}\) columns of each height \(1\le h\le n\) and add any number of height-one columns. The shifted entries up to a common constant are \(-y_i\), \(y_i=\sum_{i<h<n}(c_h+1)\) for \(0\le i<n\), with the increased value used for \(c_1\). They have no nontrivial equal-length sum relation of length \(d\le(n+1)^2\). Indeed when \(j>i\), \(y_i>d\,y_j\). For \(1\le j<n-1\) the latter sum has at most \(n\) terms of size \((n+1)^{4(n-j-1)}+1\), whereas \(y_i\ge(n+1)^{4(n-j)}+1\); for \(j=n-1\) use \(y_j=0\). Compare two multiplicity vectors at their first unequal position. If it occurs before the last, the gain of a single copy exceeds the maximal loss of \(d\) copies at subsequent positions. Equality of lengths determines the last multiplicity. For \(n=1\) equality of lengths already suffices.\formalnote{15.67}{67}

Choose nonzero \(v_1,\ldots,v_n\in k\) independent over \(L\) and let \(Y_{ij}=v_i\) for \(j\le i\), zero otherwise. All leading minors of \(Y\) are nonzero. For a diagram, products of minors using the first \(h\) columns for each column of height \(h\), with arbitrary row sets, span the usual polynomial realization of its Schur representation. Row substitution preserves this span by Cauchy--Binet; the standard monomial theorem identifies the resulting representation with the one defining the weights, counting each full column as a determinant. This is the same realization used in \S\ref{sec:odd-lifts}. In the product of the groups take products across the places. Write \(S\) for the span of the functions \(Z\mapsto F((Z_\tau Y)_\tau)\) for these products of minors and all the allowed extra height-one columns, restricted to \(z(X)\). Each is a matrix coefficient of an allowed weight: apply row substitution in a basis of the finite-dimensional representation and follow it by evaluation at the fixed \(Y\).

The functions \(u_{\tau i}(Z)=(Z_\tau Y)_{i1}=\sum_j Z_{\tau,ij}v_j\) multiply \(S\) into itself, since multiplication adds exactly one height-one column. They distinguish matrices by independence of the \(v_j\) over \(L\). No point of \(z(X)\) is a common zero of \(S\). At a fixed \(Z_0\), take the product \(w\) of the prescribed leading minors and translate its arguments by \(Z_{0,\tau}^{-1}\). This stays in the same weight space, and evaluation at \(Z_0Y\) gives \(w((Y)_\tau)\ne0\).\formalnote{15.68}{68}

Here is the elementary approximation step for these multipliers. They embed the compact image homeomorphically into \(F^s\), with \(F=L(v_1,\ldots,v_n)\) finite over \(\mathbb Q_\ell\). Scale the finitely many coordinates to take values in \(\mathcal O_F\); call the resulting compact set \(D\). On \(D\), indicators of residue boxes to any depth are uniform limits of polynomials over \(F\) in these coordinates. For residue cardinality \(q_F\), the polynomials \(1-t^{(q_F-1)\ell^j}\) tend uniformly on \(\mathcal O_F\) to the indicator of its maximal ideal. On units this follows because the compact group \(1+\mathfrak m_F\) has finite \(\ell\)-group image modulo every \(\mathfrak m_F^a\); on nonunits it follows from the positive valuation. Translation detects each residue class. Inductively, inside a box with indicator \(I\) already in the closure, apply the same formula to \(I(x-c)/\varpi_F^a\) for the next digit of a coordinate and multiply the result by \(I\). The argument is integral everywhere, and the closed algebra is stable under these operations. Products handle the coordinates. Hence \(\overline S\) is stable under every such indicator.

Given continuous \(\phi\) and \(\delta>0\), at each point choose \(s\in S\) scaled to agree with \(\phi\) there, using nonvanishing. It approximates \(\phi\) within \(\delta\) on a neighbourhood. Compactness gives a common finite residue-box partition refining these neighbourhoods. Sum the chosen approximants multiplied by the corresponding indicators. The sum lies in \(\overline S\) and is within \(\delta\) of \(\phi\). This proves the required uniform approximation by linear combinations of allowed-weight coefficients, including on a compact image with empty interior. At the fixed embeddings tensor the approximants with the unchanged weight factors. The finite class averaging and eigenspace decomposition in \S\ref{sec:galois-limits} then produce the stated automorphic models at the required level and type.\formalnote{15.69}{69}



\section{Removing the restrictions on weights, local types and splitting}
\label{sec:galois-limits}

Let \(G=G_E\), put \(k=\overline{\mathbb Q}_\ell\), and fix
\(\iota:\mathbb C\simeq k\). Suppose \(E\) contains a primitive \(5040\)-th
root of unity, \(n>0\) is odd, and \(f\) is a definite unitary eigenform of
rank \(n\), weight \(\kappa\), with Hecke polynomials \(P_w\). Fix a finite
place \(u\). We seek an attached representation \(R\), with the local
properties at \(u\) when \(f\) is spherical there. The seed representations
obtained from the test-ring cohomology in Section~\ref{sec:seed-proof}
require special weights and an auxiliary cuspidal type. An approximation removes these two conditions. A
separate descent removes the splitting hypothesis used in the approximation.

\subsection{The limiting representation}

First suppose \(E\) contains \(\mathbb Q(\sqrt{-D})\), \(D\equiv3\bmod 4\),
and that \(u\) and the places above \(\ell\) lie over split primes of this
quadratic field prime to \(D\). Choose distinct auxiliary places \(v,x\)
over \(\ell\), above the same split prime of the quadratic field, disjoint
from each other and from \(u\) up to conjugation. At \(v\) impose a cuspidal
character on the first congruence subgroup. Choose the level away from
\(v,x\) inside the stabilizer of \(f\). Where \(f\) is spherical at \(u\)
one can retain the full integral subgroup at \(u\). Keep \(\kappa\) at the
infinite places inducing \(u\) or its conjugate under \(\iota\).
At the others allow weights \(\mu\) for which the shifted partial sums
\(-\sum_{j<h<n}(\mu_h+1)\), indexed by \(j\), have no small additive
relations. Explicitly, two multisets of the same size at most
\((n+1)^2\) whose sums of these integers agree must have identical
multiplicities at every index.

The weight approximation proved in \S\ref{sec:weight-density} and the projection below put the \(\ell\)-adic model \(F\) of \(f\)
in the uniform closure of the span of models \(F_i\) of eigenforms with
these weights and this type. The models use matrices intertwining the
weights on the components where they agree. Evaluating a vector-valued
function in each of its coordinates lets us regard the models as scalar
functions on one common set. The density assertion uses approximation by
algebraic matrix coefficients in the varying weight components. To apply
it, fix a representative \(c\) of a rational double coset at the chosen
level \(K\). On \(cK\) the model splits as a matrix from the fixed
components times a vector of continuous functions of the varying
components. This factorization uses that an element of \(K\) with all
varying components trivial fixes \(f\) and has trivial level character.
Approximate the entries of the vector by matrix coefficients of free
weights using the ultrametric Stone--Weierstrass lemma. To recover the
global form, sum the local expressions over rational \(\gamma\) with
\(c^{-1}\gamma g\in K\), divide by the order of the stabilizer of \(c\),
and sum over the finitely many classes \(c\). By equivariance the original
expressions give \(F\). These are bounded operators, since the sets
being summed over have uniformly bounded size and the matrix functions
on \(K\) are bounded. They take the approximants to forms of the
required weights and level character. Decompose the resulting spaces
into simultaneous spherical eigenforms. Each resulting eigenform has
a geometric representation \(r_i\), with polynomial \(Q_{i,w}\)
for most Frobenius data. When \(f\) is spherical at \(u\), the construction keeps
the local compatibility for every \(r_i\). In particular, for \(u\mid\ell\)
it keeps crystallinity and the same jumps at every coefficient embedding.\formalnote{15.70}{70}

Choose a point \(z\) with \(F(z)\ne0\) and normalize \(F(z)=1\). For each
\(\eta>0\) choose a finite sum \(\sum_i a_{\eta i} F_{\eta i}\) at uniform
distance at most \(\eta\) from \(F\). Write \(c_{\eta i}=a_{\eta i}F_{\eta i}(z)\).
If a common Hecke operator \(H\) has norm at most one on the scalar models
and eigenvalues \(h,h_{\eta i}\), applying \(H\) to the approximation and
evaluating at \(z\) gives
\[
   \left|h-\sum_i c_{\eta i}h_{\eta i}\right|_\ell\le\eta .
\]
Every product of the operators is still a contraction. Their paired
eigenvalues therefore generate a monoid of tuples
\((h,(h_{\eta i})_{\eta,i})\) on which the same inequality holds. Let
\(N\) be its closure in the product topology. At fixed \(\eta\) the
inequality is closed because its sum is finite. The identity tuple gives
\(\sum_i c_{\eta i}\to1\). For any polynomial \(B\) and any finite list
of such tuples with first entries \(h_1,\ldots,h_b\), applying the estimate
to each monomial gives
\[
 \left|B(h_1,\ldots,h_b)-
    \sum_i c_{\eta i}B(h_{1,\eta i},\ldots,h_{b,\eta i})\right|_\ell
 \le \eta\max_\alpha|B_\alpha|_\ell .
\]
This is uniform in the tuples even though the individual coefficients
\(c_{\eta i}\) can be large. The Hecke operators furnishing \(N\) give
the pairs of \(j\)-th coefficients of \(\iota P_w,\iota Q_{\eta i,w}\),
multiplied by one common nonzero constant \(b_{wj}\). At eligible split
places away from \(\ell\), \(b_{wj}=-1\). At the target split place a
fixed \(b_{uj}\ne0\) also suffices.\formalnote{15.71}{71}

\subsection{Construction of a continuous representation from the traces}

Write \(s\) for geometric Frobenius. The signed coefficient
\(-[X^{n-1}]\iota P_w\) belongs, at the eligible Frobenii, to a fixed
finite extension \(L/\mathbb Q_\ell\). Indeed the values of the
algebraically normalized form can be taken in a finite coefficient
field, and the eigenvalue is recovered by applying the Hecke operator
and evaluating at a nonzero coordinate. Completeness of \(L\) matters
because \(k\) itself is incomplete.

Fix \(g\in G\) and let eligible Frobenius data tend to \(g\) with the
filter of Section~\ref{sec:galois}, refined by the neighborhoods of \(g\).
For a fixed row \(\eta\), the approximating signed coefficients in degree
\(n-1\) agree on a common member of this filter with
\(\operatorname{tr}r_{\eta i}(s)\). There are only finitely many \(i\)
in the row. Their weighted sum therefore tends to
\(\sum_i c_{\eta i}\operatorname{tr}r_{\eta i}(g)\). The first-coordinate
estimate shows that the signed coefficients of \(P_w\) are Cauchy along
the filter, with arbitrarily small error as \(\eta\) decreases. Denote
their limit in \(L\) by \(T(g)\). The same argument for any finite set
of coordinates, and closedness of \(N\), give
\[
  \big(T(g),(\operatorname{tr}r_{\eta i}(g))_{\eta,i}\big)\in N
  \qquad (g\in G).
\]
Thus finite sums of continuous traces converge uniformly to \(T\) on
\(G\), so \(T\) is continuous.

Every universal polynomial identity in traces at finitely many group elements passes to \(T\) by the polynomial estimate. In particular \(T(1)=n\), \(T(gh)=T(hg)\), and the fundamental trace identity of degree \(n+1\) holds. In the last identity one sums over permutations with sign, taking for each cycle the trace of its ordered product. Multilinearity extends the identity to the group algebra. For \(1\le j\le n\) the signed cycle expression of degree \(j\), evaluated on \(j\) copies of the identity, equals \(n(n-1)\cdots(n-j+1)\ne0\). This follows by inserting the last letter either as a new cycle or into one of the \(j-1\) existing positions in the permutation of \(j-1\) letters. Thus the degree \(n+1\) identity is the first one to vanish identically, as required by Rouquier's definition. It follows that \(T\) is a pseudocharacter of dimension \(n\) over the algebraically closed characteristic-zero field \(k\). By the characteristic-zero realization theorem it is the trace of a semisimple \(n\)-dimensional representation \(R\) over \(k\) \cite[Definition 2.1, Lemmas 4.1, 4.3 and the characteristic-zero statement following Corollary 4.4]{arith-rouquier}.

For continuity, the image algebra \(A=\operatorname{span}_k R(G)\) is a product of matrix algebras. The trace of \(R\) on each simple block is its ordinary trace times the positive multiplicity in \(R\); these integers are nonzero in \(k\). Thus the trace pairing on \(A\) is nondegenerate. Choose \(g_j\in G\) whose images span \(A\). The entries \(T(gg_j)\) determine \(R(g)\) by a fixed linear map. Each is continuous in \(g\), which proves continuity of \(R\).

\phantomsection\label{route:28}\leanroute{The formalization proves realization from all polynomial trace identities by taking the quotient by the trace radical, bounding its dimension by \(n^2\), and identifying its matrix blocks and their multiplicities. The cited theorem replaces this algebraic construction. The uniform Hecke approximation is needed beforehand to get a pseudocharacter, and the continuity and prescribed local Frobenius statements still require their own arguments.\formalnote{15.72}{72}}

Finally, Newton's formula writes each coefficient of
\(\det(X-R(s))\) as a fixed polynomial in \(T(s),\ldots,T(s^n)\). At a
given precision use a single sufficiently accurate row. Apply the
polynomial estimate to these trace tuples and the linear estimate to
the signed \(P_w\)-coefficient tuple. On a common member of the
Frobenius filter, all the \(r_{\eta i}\) of this row have characteristic
polynomials \(\iota Q_{\eta i,w}\). The two approximating sums then
agree with the appropriate sign, and the error is bounded by \(\eta\)
times a constant depending on the Newton polynomial. This proves
attachment in every coefficient. The qualifying Frobenii may depend
on the requested precision.\formalnote{15.73}{73}

\subsection{A matrix relation that survives the limit}

Put \(B=\mathbb Z_\ell[G]\), and use the same letter for a representation
and its linear extension to \(B\). Uniform convergence of a combination
of traces on \(G\) gives the same bound on \(B\). In a finite sum
\(\sum_g b_g[g]\), every \(|b_g|_\ell\le1\). More generally, suppose
that for finitely many indices \(j\) we have the uniform convergences
\[
 \sum_i c_{\eta i} q_{\eta i,j}\operatorname{tr}r_{\eta i}(g)
 \longrightarrow p_j\operatorname{tr}R(g).
\]
Choose \(y_{\eta j}\in B\), allowed to vary with the row, such that
\(\sum_j q_{\eta i,j}r_{\eta i}(y_{\eta j})=0\) for every \(i\).
If \(M=\sum_j p_jR(y_{\eta j})\) is independent of \(\eta\), then \(M=0\).
Indeed test each trace convergence on \(y_{\eta j}[g]\), sum over \(j\),
and take the limit. The approximating side vanishes even before passage
to the limit, whence \(\operatorname{tr}(MR(g))=0\) for every \(g\).
The nondegenerate trace pairing proves the assertion.

Suppose \(f\) is spherical at \(u\nmid\ell\). For inertia \(h\), take \(y=[h]-[1]\). Good
approximants are unramified, so \(R(h)=1\). If \(u\) is split, let
\(s\) be a geometric Frobenius and take \(y_j=[s^j]\), with \(p_j\)
and \(q_{\eta i,j}\) the coefficients at \(u\) of the target and seed
polynomials over \(k\). The required convergences follow by multiplying
the coefficient tuples in \(N\) by the trace tuples and dividing by
the fixed nonzero \(b_{uj}\). The annihilating equations on the seeds
give \((\iota P_u)(R(s))=0\).\formalnote{15.74}{74}

\subsection{The additional work at the coefficient prime}

Let \(u\mid\ell\) and suppose \(f\) is spherical at \(u\). Fix a
prolongation and a compatible embedding into an \(\ell\)-adic closure,
and write \(\Delta\subset G\) for its decomposition group. All
\(r_{\eta i}\) are crystalline there with the same prescribed jumps.
Their jumps are consequently bounded in absolute value by some \(d\).
Enlarge \(d\) to bound the finitely many integers occurring with
nonzero jump multiplicity for \(R\) at any coefficient embedding. This bound precedes crystallinity. Indeed a nonzero Hodge--Tate form on \(R\) pulls back to a nonzero form of the same weight on the integral image lattice \(T\) below, by \(y\mapsto yv\) for some \(v\) on which the form is nonzero. For the finite-rank \(\mathbb Z_\ell\)-module \(T\), the span of all its weight forms in \(\operatorname{Hom}_{\mathbb Z_\ell}(T,\mathbb C_\ell)\) has a finite basis of weight forms. Different weights span independent subspaces by the coefficient comparison in \S\ref{sec:jump-symbol}, so only these finitely many weights occur. The pullback to \(T\) works for each coefficient embedding because all restrict to the same map on \(\mathbb Z_\ell\). Twisting by the cyclotomic character to power \(-d\), in the source's convention, puts all \emph{occurring} jumps in \([0,2d]\). At this stage only the approximants are known to be crystalline.

Let \(T\) be the image of \(B\) in the matrix algebra of \(R\). This is a
finite free \(\mathbb Z_\ell\)-lattice. To see finiteness, a compact subgroup of \(\mathrm{GL}_n(\overline{\mathbb Q}_\ell)\) lies in the matrices over a finite extension: the algebraic closure is a countable union of finite local fields, so by Baire one of the corresponding closed subgroups has interior. Its finitely many cosets require only finitely many more coefficients. Over this field boundedness of the compact image gives a stable lattice by taking the span of the translates of the standard lattice. Thus \(T\) is a submodule of a finite matrix lattice over \(\mathbb Z_\ell\), and is torsion-free. The same proof works for every finite family of representations used below. For each integer \(a\ge1\) there
is an \(\epsilon_a>0\) such that
\[
 t\in T,\quad
 |\operatorname{tr}(tR(g))|_\ell\le\epsilon_a\ \ (\forall g)
 \quad\Longrightarrow\quad t\in\ell^aT .
\]
Indeed finitely many trace tests separate \(T\), and the complement of
\(\ell^aT\) in the compact lattice is compact. Choose a row \(\eta_a\)
within the trace tolerance \(\epsilon_a\), and let \(T'_a\) be the joint
image of \(B\) on all the rows chosen through \(a\). If \(y\) is in
the kernel of \(B\to T'_a\), then so is \(y[g]\) for every \(g\).
The approximation gives the displayed inequalities for \(t=R(y)\).
Thus \(\ker(B\to T'_a)\) maps into \(\ell^aT\). This argument works
for arbitrary sizes of the coefficients \(c_{\eta i}\).

Liu's bounded-weight criterion gives the classical lattice statement behind this step. For a finite free \(\mathbb Z_\ell\)-representation \(T_0\) of the absolute Galois group of a finite extension of \(\mathbb Q_\ell\), suppose that for every \(a\ge1\) the reduction \(T_0/\ell^aT_0\) is a quotient of two stable lattices in a crystalline representation with weights in a fixed interval \([0,r]\), \(r>0\). Then \(T_0[1/\ell]\) is crystalline with weights in that interval \cite[Conjecture~1.0.1, crystalline case proved in \S7.3]{lit-liu-torsion}. Here the group-ring images relate convergence of traces to such integral reductions. The period calculation also keeps track of individual jumps and Frobenius annihilators.

The twisted left action of \(\Delta\) on each \(T'_a\) is crystalline
with weights in \([0,2d]\). Over \(\mathbb Q_\ell\), evaluation on all the columns embeds the joint image into the direct sum of its row representations viewed by restriction of scalars from a common finite coefficient field \(L_a\). After extension to \(\overline{\mathbb Q}_\ell\) this target includes every coefficient conjugate, since \(L_a\otimes_{\mathbb Q_\ell}\overline{\mathbb Q}_\ell=\prod_{\sigma:L_a\hookrightarrow\overline{\mathbb Q}_\ell}\overline{\mathbb Q}_\ell\). Their crystallinity is invariant under conjugation of coefficients, and the jump bound holds at every coefficient embedding. Subrepresentations inherit both properties. The uniform
period bound supplies one \(\delta\in A_{\max}\), invertible in
\(B_{\max}\), for \(\Delta\) and this weight interval. In particular,
in \(A_{\max}\otimes T'_a\), \(\delta\otimes1\) belongs to the
\(A_{\max}\)-span of the invariant periods, for every \(a\). To see
how this gives periods in the limit, lift the finite expressions in
invariant periods to \(A_{\max}\otimes B\) and expand their scalar
coefficients in an adic basis over the unramified coefficient ring
of the local period chart. Extension of scalars on the invariant
module maps injectively into the full period module. This makes
the coordinate expressions for later levels and level \(a\) agree
after projection to level \(a\). By the kernel bound their images in
\(A_{\max}\otimes T\) agree modulo \(\ell^a\). Completeness gives
limits of the coordinate expressions, and the fixed-module equations
put those limits in the invariants. Saturation and the adic expansion
put \(\delta\otimes1\) in the span of the invariants of
\(A_{\max}\otimes T\). Apply the column maps \(t\mapsto t v\), for
arbitrary vectors \(v\). After inverting \(\delta\), the invariant
periods span the full period module, which gives the crystalline
comparison for the twisted \(R\). Untwisting proves crystallinity of
\(R|_\Delta\). The uniform period bound and the integral fixed-module
properties are the period-theory inputs to this argument.\formalnote{15.75}{75}

To see why the same denominator works throughout the limit, use the Kummer chart constructed for the local group in \S\ref{sec:kummer-local}, of ramification index \(e\), with the group fixing its constants and \(\theta(z)\). Choose \(j\) with \(p^j\ge3e\), and let \(\alpha:A_{\inf}\to A_{\max}\) be Frobenius to power \(j\) followed by the structure map. For jumps in \([0,h]\) one may take
\[
                         \delta=\alpha(E\tau)^h .
\]
The element is a unit in \(B_{\max}\) since \(E\tau\) is associated to \(\mu\), and its Frobenius transforms become units upon localization. In the notation of the frame calculation, let \(U\) have invariant columns in \(B^+\), \(UZ=E^hQ\), and \(QL=\tau^h I\) over \(A_{\inf}\) in the last equation. Apply the radius factorization of \S\ref{sec:kisin-height} with its free exponent set to zero. It gives \(\varphi^j Z=CV\), where the group fixes \(C\) and \(V\in{\rm GL}_d(A_{\max})\). The two matrices
\[
            P=\alpha(E^hQ)V^{-1},\qquad P'=V\alpha(L)
\]
are integral. Over \(B^+\) the first is \((\varphi^j U)C\), by applying Frobenius to \(UZ=E^hQ\). Thus its columns are invariant since Frobenius commutes with the group action and fixes the \(\mathbb Z_p\)-coefficients of the action on the lattice. The product \(PP'=\delta I\) proves the assertion about the span of invariants for every lattice, of every rank, with this same jump interval. Neither the radius \(j\) nor \(\delta\) depends on the row of the trace approximation.\formalnote{15.76}{76}

To recover the individual multiplicities, for each row take the
joint image lattice \(T_\eta\) for \(R\) and all \(r_{\eta i}\) together.
It is now crystalline in the same interval after twisting. Fix an
embedding \(\gamma:k\to\mathbb C_\ell\) over \(\mathbb Z_\ell\) and a
jump \(w\). The bounded period projector for the shifted index \(w+d\)
gives a scalar \(b\), independent of \(\eta\), and for each row lifts
\(y_{\eta s}\in B\) of a basis of \(T_\eta\) and scalars
\(|z_{\eta s}|_\ell\le1\). For each representation \(\sigma\) in the row,
including \(R\), the projector reads its original jump multiplicity as
\[
 m_{\gamma,w}(\sigma)=
 b\sum_s z_{\eta s}\,\gamma\big(\operatorname{tr}\sigma(y_{\eta s})\big).
\]
The embedding is isometric. Subtract the formulas, weighting the
seed formulas by \(\gamma(c_{\eta i})\). The absolute error is bounded by
\[
  |b|_\ell\sup_{y\in B}
  \left|\operatorname{tr}R(y)-
    \sum_i c_{\eta i}\operatorname{tr}r_{\eta i}(y)\right|_\ell ,
\]
even though the bases and their lengths vary. This tends to zero.
The seed multiplicities at \(u\) all equal \(m_{\gamma,w}(\kappa)\)
and \(\sum_i c_{\eta i}\to1\), so \(R\) has precisely these multiplicities.

Here are the construction and size of the projector used in this step.
Fix the joint lattice \(T_\eta\) with a \(\mathbb Z_p\)-basis and let \(x\)
be the identity in its group-ring image. Take the invariant integral
matrix \(P\) above, \(PP'=\delta I\). In a de Rham frame
\(PY=W\operatorname{diag}(t^{k_i})\), \(W\) is invertible over
\(B_{\rm dR}^+\), and \(Y\) is a change of basis over the fixed local field with
determinant one and entries in its ring of integers. Enlarge the latter
ring to the integers \(O\) of a fixed finite Galois extension if needed.
Let \(J_w=\operatorname{diag}({\bf1}_{k_i=w})\) for the shifted weight
and put \(a=WJ_w W^{-1}x\). The diagonal powers of \(t\) commute
with \(J_w\). Therefore
\[
               \delta a=P\,YJ_w\operatorname{adj}(Y)P'x .
\]
In particular every entry of \(\delta a\) belongs to the image of
\(A_{\max}\otimes_{\mathbb Z_p}O\to B_{\rm dR}^+\). This is why the
denominator is independent of the dimension of the joint lattice.

For the scalar bound write \(\delta=\xi^m u\), \(u\) a unit in
\(B_{\rm dR}^+\). Truncation modulo \(\xi^{m+1}\) and the density of
\(A_{\inf}[1/p]\) in that finite truncation put the image of each generator
of \(A_{\max}\otimes O\) in
\(p^{-c}A_{\max}+\xi^{m+1}B_{\rm dR}^+\) for some \(c\).
One \(c\) works for all elements, since the domain is generated over
\(A_{\max}\) by a finite basis of \(O\). Thus
\(p^c\delta a_s=b_s+\xi^{m+1}v_s\) with \(b_s\in A_{\max}\).
The integral division estimate for \(\xi\mid\mu\) gives
\(p^{2m}b_s/\xi^m\in A_{\max}\): divisibility in \(B_{\rm dR}^+\)
first gives divisibility in \(B^+\), then each division costs at most
\(p^2\). With \(\kappa=p^{2m+c}u\), reduction by Fontaine's
\(\theta\) gives \(|\theta(\kappa a_s)|_p\le1\) and
\(\theta(\kappa)\ne0\). This gives \(b=1/\theta(\kappa)\) and
\(z_{\eta s}=\theta(\kappa a_s)\).

For the integral division estimate, put \(A=A_{\max}\) and
\(X=\xi/p\). Modulo \(p\) its presentation gives
\(A/p=(A_{\inf}/(p,\xi))[X]\). Indeed the relation is \(pX-\xi\)
and reduction modulo \(p\) is unchanged by completion of this
\(p\)-torsion-free algebra. Every monic polynomial \(H(X)\) is regular
modulo \(p\). If \(H(X)d\in A\) for \(d\in A[1/p]\), clear a minimal
power \(p^a\) in \(d\). For \(a>0\), reduction modulo \(p\) would
annihilate the nonzero class of \(p^a d\) by \(H(X)\). Thus \(d\in A\).
Now \(\mu=\tau\xi\), and
\(\xi=\sum_{i=0}^{p-1}(1+\tau)^i\) gives \(\tau\mid(\xi-p)\).
For any \(e'\mid\mu\) in \(A_{\inf}\), we consequently have
\(e'\mid \xi(\xi-p)=p^2 X(X-1)\) in \(A\). If \(e'^{\,m}y\in A\)
for \(y\in A[1/p]\), then the monic polynomial
\((X(X-1))^m\) times \(p^{2m}y\) belongs to \(A\). The preceding
argument gives \(p^{2m}y\in A\), including \(p=2\). In the application
\(e'=\xi\). Divisibility of \(b_s\) by \(\xi^m\) in the completion
\(B_{\rm dR}^+\) makes its image in \(B^+/\xi^m\) zero, so division
indeed takes place first in \(B^+=A[1/p]\).\formalnote{15.77}{77}

For an equivariant map from \(T_\eta\) to the matrices of a constituent
which sends \(x\) to the identity, each column is a map of representations.
After applying \(\theta\), it intertwines the Hodge--Tate projections
because there are no equivariant maps between scalar lines of distinct
weights \cite[Theorem~2.2.7]{legacy-brinon-conrad}. The matrix obtained by applying these column
maps to the projection of \(x\) is therefore the projection of the
identity onto its weight-\(w\) component in the target. The trace of this projector is the multiplicity of \(w\). Expanding \(a\) in the chosen basis yields exactly the
trace formula above, with the same \(b\) for every row.\formalnote{15.78}{78}

At a split \(u\mid\ell\), let \(s\) be a geometric Frobenius for \(u\). In a compatible local chart \((s,k_0)\) is a Weil pair with \(q_u=\ell^{k_0}\). Write \(\beta\) for its linearized Frobenius. The seed polynomial at \(u\) annihilates this operator for each seed because \(s^{-1}\) acts on the residue field of the algebraic closure by the \(q_u\)-power map. Evaluating the target polynomial at the same period
operator gives the local clause. Take \(m\) strictly greater than the degree of each polynomial
in the relation. Apply the uniform period bound to
the joint lattices \(T_\eta\). For \(0\le j<m\), the multiple
\(\ell^{k_0d(m-j)}\delta\otimes1\) is in the span of their integral
invariants. Computing the \(j\)-th Frobenius power there and untwisting
gives a simultaneous integral representative for the \(j\)-th term,
with common factor \(\delta\ell^{k_0dm}\) in the untwisted equation.
For any \(\mathbb Z_\ell\)-linear coordinate
\(\lambda:A_{\max}\to\mathbb Z_\ell\), apply \(\lambda\) to those
representatives and lift the resulting elements of \(T_\eta\) to
\(y_{\eta j}\in B\). For the seeds,
\(\sum_j q_{\eta i,j} r_{\eta i}(y_{\eta j})=0\), since their
Frobenius polynomials kill \(\beta\). For \(R\) the corresponding matrix
\(\sum_j p_j R(y_{\eta j})\) is independent of \(\eta\). On a vector
\(v\) it reads the same coordinate of the integral representative of
\(\delta\ell^{k_0dm}(\iota P_u)(\beta)(1\otimes v)\).
That representative is unique because extension from integral to
rational periods is injective. The matrix relation of the preceding
subsection now gives zero for every coordinate. Coordinates separate,
and the common factor is invertible in \(B_{\max}\), so
\((\iota P_u)(\beta)=0\).\formalnote{15.79}{79}

\subsection{Descent through varying quadratic extensions}

We finally remove the condition on the auxiliary quadratic subfield.
Choose a finite set \(S\) containing the places at which the level of
\(f\) is prescribed, and write \(q_v\) for the residue field size at
\(v\). For \(T\ge0\), put \(a_T=\ell(T+q_u+\max_{v\in S}q_v)!\),
with the maximum taken as zero for empty \(S\). The construction gives an additive
finite character \(\chi_T\), a field \(L_T\) fixed by
\(\ker(\chi_T+\chi_T^c)\), and an open normal subgroup
\[
                N_T=\ker(\chi_T+\chi_T^c)\cap\ker\chi_T\subset G .
\]
To construct the step, write \(a=a_T\), let \(k_1=N_{E/\mathbb Q}(4)\),
and put \(M=4a(2a)^{k_1}\). We first find an integral \(\alpha\) whose
norm \(\alpha\bar\alpha=D\) is \(-1\bmod M\) and has valuation one
at a place away from \(a\). The field \(E\) contains \(i=\sqrt{-1}\)
and a primitive seventh root \(\zeta\). Put
\(s=\zeta+\zeta^2+\zeta^4\). The nonzero residues modulo \(7\)
are the disjoint union of \(\{1,2,4\}\) and its negative, so
\(s+\bar s=-1\). The six ordered differences between distinct
members of \(\{1,2,4\}\) also give each nonzero residue once. Hence
\(s\bar s=3+\sum_{j=1}^6\zeta^j=2\). Thus \(\omega=1+s=-\bar s\)
satisfies \(\omega+\bar\omega=1\) and \(\omega\bar\omega=2\).
For the odd part \(d\) of \(M\), the norm of \(1+2d\omega\) is
\(C=1+2d+8d^2\), coprime to \(M\) and congruent to \(3\bmod4\).
Choose a large prime \(q\equiv-C^{-1}\bmod M\), avoiding \(C\), \(a\),
and the discriminant of \(E\). Since \(q\equiv1\bmod4\), write
\(q=u_1^2+u_2^2\) by Fermat's two-square theorem. Then
\(\alpha=(1+2d\omega)(u_1+i u_2)\) has \(D=Cq\) as required.
At any \(\mathfrak p\) over \(q\), the rational prime \(q\) is
a uniformizer by unramifiedness. Thus \(-D\) is a uniformizer there.\formalnote{15.80}{80}

\begin{example}[A norm congruent to minus one modulo sixteen]\label{ex:quadratic-norm-congruence}
Seek an algebraic integer \(\alpha\in\mathbb Q(i,\zeta_7)\) such that \(\alpha\bar\alpha\equiv-1\bmod16\) and \(\alpha\bar\alpha\) has valuation one at an odd unramified prime. For a primitive seventh root \(\zeta_7\), take \(\omega=1+\zeta_7+\zeta_7^2+\zeta_7^4\). The six nonzero exponents in the two sums for \(\omega-1\) and \(\bar\omega-1\) exhaust the nonzero residues modulo \(7\). Thus \(\omega+\bar\omega=1\). Also the six differences of distinct elements of \(\{1,2,4\}\) exhaust those residues. Hence \((\omega-1)(\bar\omega-1)=3-1=2\), which gives \(\omega\bar\omega=2\). Therefore
\((1+2\omega)(1+2\bar\omega)=1+2(\omega+\bar\omega)+4\omega\bar\omega=11\).
We need a prime \(q\) with \(11q\equiv-1\bmod16\). The prime
\(13=2^2+3^2\) works. Hence \((1+2\omega)(2+3i)\) has product with
its complex conjugate \(11\cdot13=143\equiv-1\bmod16\).
The prime \(13\) is unramified in this cyclotomic field, so \(143\)
also has valuation one at every place above \(13\), with integral
normalization of the valuation.
\end{example}

For the real factor, use the Chinese remainder theorem to choose
\(r\in\mathcal O_E\) with \(r\equiv1\) at one place of each split
pair over \(a\) and \(r\equiv\alpha\) at the other, modulo the
\(k_1\)-th powers of those primes. At fixed places prescribe
\(r\equiv1\) to that depth. Put \(t=r\bar r\). Then at the split
places \(\alpha t\) is congruent to \(D\equiv i^2\bmod\mathfrak w^{k_1}\)
or to \(\alpha^2\bmod\mathfrak w^{k_1}\). Both are squares of units,
and \(v_{\mathfrak w}(4)<k_1\), since
\(N(\mathfrak w^{k_1})\ge2^{k_1}>N((4))\).
Hensel's lemma makes \(\alpha t\)
a local square there. At fixed places over \(a\), \(\alpha t\) is
a unit and the residue characteristic is odd: the norm-two
element \(\omega\) and its conjugate have sum \(1\), so no prime
over \(2\) is fixed by conjugation. A square root of the unit
thus lies in the maximal unramified extension.\formalnote{15.81}{81}

Take \(b^2=\alpha t\) and define \(\chi_T(g)\) to be \(0\) or
\(\tfrac12\) according as \(g(b)=b\) or \(-b\). Since \(i=\sqrt{-1}\)
lies in \(E\), \(y=(i/t)b\bar b\) satisfies \(y^2=-D\). The signs of \(g\)
on \(b,\bar b\) agree precisely when \(g(y)=y\), so \(L_T=E(y)\).
The congruence for \(D+1\) makes \(-D\) a local square above \(2a\):
its distance from \(1\) is smaller than the local bound for \(4\),
and Hensel's criterion applies to the square root at \(1\).
Inertia at \(\mathfrak p\) contains an element
acting by \(-1\) on \(y\), since \(-D\) has odd valuation there. Its
signs on \(b,\bar b\) are opposite. Conjugating this inertia element
by complex conjugation gives an inertia element over
\(\bar{\mathfrak p}\) with the signs interchanged. Their images
therefore generate both sign coordinates, and hence \(G/N_T\).

The required splitting in \(\mathbb Q(\sqrt{-D})\) follows from
\(\beta(1-\beta)=(D+1)/4\) for
\(\beta=(1+\sqrt{-D})/2\). At a prime over a divisor of \(a_T\), one of
\(\beta\) and \(1-\beta\) vanishes in the residue field, and a fixed
prime would force both to vanish. The prime is also prime to \(D\).
At each place of \(E\) dividing \(a_T\), the decomposition group fixes
\(L_T\), inertia lies in \(N_T\),
and, if the place is split in the CM extension, its full decomposition
group lies in \(N_T\). Section~\ref{sec:weak-transfer} constructs an eigenform by weak base change over
\(L_T\). The split-case theorem gives its representation, whose
restriction \(\rho_T\) to \(N_T\) is attached along \(N_T\hookrightarrow G\)
to the original polynomials. Restriction to the character kernel
removes the character occurring in base change.\formalnote{15.82}{82}

Distinguish \(N_0,\rho_0\). Its step supplies a set \(I\) of inertia
elements at places of bounded norm, away from \(a_0\), whose images
generate \(G/N_0\). Beyond a fixed threshold the \(a_T\) include
those places. They are away from \(\ell\) and the level, and their
conjugates are away from the level. Hence every conjugate of each
element of \(I\) lies in \(N_T\) and is killed by \(\rho_T\), by the
unramified local clause. Use only these large \(T\) in the following
intersection, and let
\[
            K=\operatorname{core}_G\Big(\bigcap_T\ker\rho_T\Big).
\]
Then \(I\subseteq K\), so \(N_0K=G\). On \(N_0\cap N_T\), attachment
and density give identical traces for \(\rho_0,\rho_T\). The restrictions
are semisimple, since the translates of any simple submodule for a
normal subgroup span an irreducible representation of the whole group.
Thus the restrictions are conjugate and \(\rho_0\) kills \(N_0\cap K\).

The isomorphism \(N_0/(N_0\cap K)\simeq G/K\) lets \(\rho_0\) define
\(R\) on \(G\). It agrees with \(\rho_0\) on the open subgroup \(N_0\)
and has the same image, so is continuous with nondegenerate trace.
For \(g=n_0k'\in N_T\), with \(k'\in K\), one has
\(n_0\in N_0\cap N_T\), so \(R|_{N_T}\) is conjugate to \(\rho_T\).
Given a nonempty open set in \(G\), choose a Frobenius of prime norm
there by Chebotarev. Its place is split under the field hypotheses
and its norm eventually divides \(a_T\). Intersect the open set with
such an \(N_T\) around the chosen Frobenius. On a smaller open subset,
attachment for \(\rho_T\) gives the Frobenius estimates for \(R\).
This proves global attachment.

\newpage
\paragraph{Local conditions after descent.}
At \(u\), the base-change representation agrees with \(R\) on
\(N_0\) up to conjugacy. Since inertia lies in \(N_0\), it acts
trivially when \(u\nmid\ell\). At \(u\mid\ell\), conjugacy on an open
subgroup gives potentially crystalline periods and preserves the
jump spaces defined by equivariance near the identity. For an
inertia element \(s\), it identifies \(\beta_{s,0}\), the
\(B_{\max}\)-linear extension of \(s\) acting semilinearly on those
periods. This operator is the identity upstairs by crystallinity
and therefore downstairs. A normal open subgroup and inertia fix
a spanning set of periods, so the periods are already defined
over a finite unramified extension. Its unramified coefficient
field embeds in \(B_{\max}\) by Witt vectors. Finite semilinear
descent gives spanning periods over the original field, proving
crystallinity. At a split \(u\) the full decomposition group lies
in \(N_0\). Conjugacy and the matching polynomial give the local
annihilating relation for the Galois operator away from \(\ell\)
and the linear period operator at \(\ell\). Uniqueness of the
attached semisimple trace gives these conclusions for every
attached representation, as required in \S\ref{sec:galois}.\formalnote{15.83}{83}


\editionchapter{The marked unitary Shimura varieties}\label{ch:shimura}

To relate Hecke eigenvalues to Frobenius, we realize the Hecke operators as correspondences on families of abelian varieties with level structure. We count the fixed points of these correspondences after reduction. Example~\ref{ex:torsion-sections-coalesce} shows why distinct characteristic-zero torsion sections can contribute multiplicity at a single point after reduction.

\section{The integral tower and reduction of its Hecke neighbours}\label{sec:pel-covers}
At a prime \(p\) of good reduction where the group is split over \(\mathbb Q_p\), Wedhorn proves for PEL Shimura varieties that the Hecke polynomial annihilates the Frobenius correspondence \cite[main theorem in the introduction, part~(1)]{lit-wedhorn}. We need a pointwise cancellation on the particular covers constructed below. Each term of the Hecke polynomial in \S\ref{sec:seed-proof} is represented by a correspondence on the base of a family of abelian varieties. To compare all its terms at a point, we give the correspondences a common domain by imposing enough level structure to specify all the relevant torsion subgroups. Their isogenies then have kernels that can be compared after reduction modulo \(p\). The comparison uses the multiplicity with which the characteristic-zero torsion sections meet at a point of a special fibre.

\subsection{Theta coordinates, the base and the covers}
Use the polarized abelian varieties with endomorphisms and level structure of \S\ref{sec:seed-proof}. Write \(\mathcal O=\mathcal O_E\), \(n\) for the rank, \(g=n[E:\mathbb Q]/2\) for the abelian dimension, \(R\) for the localization at the safe places and \(\mathfrak N\) for the neat level. At the distinguished complex place the varying negative or positive line in the form of signature \((1,n-1)\) gives a ball of complex dimension \(n-1\) (interchanging the sign convention interchanges the designated line). At a parameter \(b\), evaluate all the projective theta coordinates of the torus at each of its finitely many level-torsion labels \(\omega\). Let \(S\) be the reduced closure over \(R\) of these tuples in the product of projective spaces. The tensor product of the hyperplane bundles embeds \(S\) by the Segre map. In particular it is projective. Two results about these coordinates used below are the integral theta equations at a good prime, with their prescribed ranks, and the construction of an abelian family by gluing their affine charts. We explain where both enter.

At level \(\mathfrak N\) finitely many analytic ball families represent all the data up to level isomorphism. One can choose one complex parameter in each family whose scheme point has the image of the entire family in its closure. Indeed the base of definition is countable and the functions are holomorphic on a connected domain. For each defining regular function which does not vanish identically its zero set has empty interior; one can avoid the countably many such proper analytic sets in local charts. Thus finitely many scheme points have dense union of closures. Given \(s\in S\), properness gives a specialization to a point \(s_0\) above a maximal ideal of \(R\). A valuation ring of \(\mathbb C\) can be chosen with generic point over one of the dense parameters and centre \(s_0\), by the valuative criterion for specialization. By the choice of \(R\) this is a good prime. In a chart with nonzero level coordinate \(i\), the transition function at the generic point is
\[
       p_{\omega j}=
       \frac{\theta_j(\omega)}{\theta_i(\omega)} .
\]
Along the arc choose a chart containing its closed point. Then the ratios belong to the valuation ring and their residues are the homogeneous coordinates at \(s_0\). For a general argument of theta one may divide by a coordinate of largest nonarchimedean absolute value, so some normalized coordinate is always a unit. The good-reduction theta equations identify the resulting affine loci with the prescribed ranks, and the requisite coordinates for the addition formula can be chosen simultaneously at indices \(o,-o,s'-o\) for every dual-torsion index \(s'\). Since nonvanishing is open these charts also cover the original generalization \(s\). The affine reconstruction theorem now glues the theta family \(\mathcal A/S\), its endomorphism action, line bundle and level sections. The calculation explains why the integral equations are required even though all the original coordinates were specified at complex points.\formalnote{16.1}{1}

For smoothness, the deformation estimate bounds the relative tangent dimension at the good centre by \(n-1\), while the analytic parameter gives an immersed ball of dimension \(n-1\) at the generic point. Here is the immersion argument. Transport the real period coordinates between parameters by \(T_{b'}=Q_{b'}Q_b^{-1}\), and let \(A(y)=\partial_v(T_{b'}y)|_b\) for a direction \(v\) stationary on the algebraic base. For each lattice vector \(l\), the section of the analytic uniformization at \(T_{b'}l/k\) is stationary to first order: multiplying by \(k\) gives the zero section over the stationary base, and multiplication by \(k\) has injective differential in characteristic zero. For any regular test function pulled back to the parameter and fibre, call it \(h\), this gives
\[
       Dh_{(b,l/k)}(v,A(l)/k)=0 .
\]
Subtract \(Dh_{(b,0)}(v,0)=0\), multiply by \(k\) and pass to the limit. One gets \(D_zh_{(b,0)}A(l)=-D^2h_{(b,0)}((0,l),(v,0))\). Both sides are real-linear in \(l\) and the lattice spans. The right side is complex-linear in \(l\), so every \(D_zh\) kills \(A(iy)-iA(y)\). Since the fibre uniformization is an immersion, \(A\) is complex-linear. It follows that \(v=0\). Indeed, a nonzero tangent vector to the ball moves the distinguished line in a complementary direction and hence varies the complex structure, contrary to the complex linearity of \(A\). The smoothness criterion along an arc over a Dedekind base, using this lower bound and the special-fibre tangent bound, gives an open smooth neighbourhood of \(s_0\). It contains \(s\) as well. This yields a smooth projective \(S/R\) of relative dimension \(n-1\). For \(\nu>0\), the base of the tower is \(Z=\mathcal A^\nu_S\), with relative dimension \(n-1+\nu g\). Each new factor is obtained by smooth base change and adds \(g\). Multiplying each marked point by an element of \(\mathcal O\) defines the required multiplication maps on \(Z\).\formalnote{16.2}{2}

Let \(K=K(\mathfrak N)\) be the fixed level in the finite adelic similitude group. For a finite list \(i\) of desired legs \((m,a)\), choose a principal level \(M\ne0\) so deep that
\[
          a^{-1}K(M)a\ \subset\ K\cap K(m)
\]
for every listed pair, also adjoining \((1,1)\) and \((M,1)\). The condition comes from openness of the finite intersection of the conjugates of the levels. Let \(D=K/\operatorname{core}_K(K\cap K(M))\). This quotient is finite since the open subgroup has finite index in \(K\). For an infinite-level analytic datum \(x\), consisting of its finite frame, its distinguished reflection and its \(\nu\) marked vectors modulo the period lattice, take the tuple
\[
           \bigl(\operatorname{datum}_m(x\cdot \kappa a)\bigr)_
                 {\kappa\in D,\ (m,a)\in i\cup\{(1,1),(M,1)\}}.
\]
The map \(\operatorname{datum}_m\) records the level-\(m\) torsion and marked points in a fibre power of \(\mathcal A\). Changing a representative of \(\kappa\) gives the same point by the displayed containment. The reduced closure \(C_i\) of these tuples in that product is projective. Its maps to the factors are the projections. Its maps to the base induced by \(a\) depend only on the right coset \(aK\): representatives in one coset give a level isomorphism on each complex datum, and maps to the separated target agreeing on the dense family agree on its reduced closure. To realize translation \(a\) into an already constructed cover, include every translated factor of that cover in the new list. Projection to those factors lands in the old closure by density. This proves that all the finitely many Hecke maps in one proposed relation are available at a common level.\formalnote{16.3}{3}

The quotient \(D\) acts by permuting the factors, inducing \(x\mapsto x\cdot\kappa^{-1}\) and preserving the identity projection \(C_i\to Z\). That projection is onto on complex points: its image is closed by properness and contains the dense base data. Every complex point of \(C_i\) is an actual tuple. One proves this by taking a valuation arc from a generic member to it. The base object at the complex centre is again uniformized by a torus with frame and reflection. Lift the marked sections to real vectors in this torus. Along each projection the isogeny inducing the prescribed translate on torsion specializes to the corresponding isogeny of these uniformizations. Thus the tuple formed from the centre frame and the lifted markings has every one of the prescribed coordinates. This uses the extension of isogenies and full structures along the arcs and the complex torus classification.

Two tuples in a fibre represent base data isomorphic at level \(K\). The rational isomorphism of frames differs by some \(\kappa\in K\); translating one frame by \(\kappa^{-1}\) makes that isomorphism valid at every deeper level, including one determining the tuple. Hence the deck action is transitive on each fibre. To prove freeness, suppose \(\kappa,\kappa'\) have the same effect on one tuple. Compare their level-\(M\) factors after every permutation indexed by \(q\in K\). The rational frame isomorphism on each of these factors fixes the reflection and, after conjugating by the frame, belongs to \(K\). Neatness forces this isomorphism to be \(1\). The remaining level equality then places the difference of the two elements in \(q(K\cap K(M))q^{-1}\) for every \(q\). This says exactly that their classes in \(D\) are equal. Section~\ref{sec:cover-charts} deduces generic finite étaleness from this action. The characteristic-zero covering criterion transfers the simply transitive action to the geometric points of the characteristic-zero fibres.\formalnote{16.4}{4}

Take a rational level isomorphism \(z\) from the translated datum to the original, with separable characteristic polynomial, and assume the reflection at the point is the one attached to the marked eigenvalue \(\mu\) of \(z\). At this fixed point of identity leg \(c\) and a translated and multiplied leg \(v\), the marked ball family through the point has parameter space of complex dimension \(n-1+\nu g\). Its first jets give tangent vectors on \(C_i\); their composition with \(dc\) is injective by the immersion just proved and the fibre uniformizations. The rational automorphism fixing the point induces a linear map \(H\) in the marked ball coordinates; on the universal first jet the identity of the translated data differentiates to \(dv=dc\circ H\) on this parameter space. The local polynomial is therefore read from \(\det(1-XH)\), up to the nonzero scalar in the tangent-polynomial identity. The block matrices and the precise normalizing scalar are computed in \S\ref{sec:geometry-closure}. The weights in the determinant come from the linear action on the varying line and the \(\nu\) fibre directions, with the marked multipliers on the latter. This is the tangent-polynomial property needed by the trace argument even when the cover is presented as a graph closure.\formalnote{16.5}{5}

\subsection{Detection of a nonreduced kernel by torsion sections}

\begin{example}[Torsion sections coalescing in characteristic \(p\)]\label{ex:torsion-sections-coalesce}
Take a primitive \(p\)-th root of unity \(\zeta_p\) and let \(\mu_p\) be the kernel of \(x\mapsto x^p\) on the multiplicative group. Over the integers of \(\mathbb Q_p(\zeta_p)\), its equation is
\[
             x^p-1=\prod_{j=0}^{p-1}(x-\zeta_p^j).
\]
Its \(p\) distinct generic points give integral sections. In the residue field every section has \(x=1\), since \(x^p-1=(x-1)^p\) there. Writing \(u=x-1\), the special fibre has coordinate ring \(\mathbb F_p[u]/(u^p)\), of length \(p\) at its unique point. A nilpotent direction is visible over the dual numbers \(D=\mathbb F_p[\varepsilon]/(\varepsilon^2)\) as well. The element \(1+\varepsilon\) is a nonidentity point of the kernel in \(D\), since \((1+\varepsilon)^p=1\). The reduced equation \(u=0\) would lose this point.
\end{example}

Let \(w\) be safe of residue size \(q=p^d\), and specialize a marked family along a valuation ring of \(\mathbb C\) to an algebraically closed field \(k\) of characteristic \(p\). Write \(\beta=(1+\sqrt{-d_0})/2\). Goodness gives an integer \(b\) with \(h=\beta-b\in w\) and \((h,\bar h)=1\). Adelic \(w\)-torsion identifies with \(k_w^n\) by \(v\mapsto[\varpi_w^{-1}\tilde v]\), supported at \(w\). The ambiguity in an integral lift \(\tilde v\) is an integral vector after the division. The map is injective and has the order of the full \(w\)-torsion. A matrix integral at \(w\) kills a label precisely when its reduction kills \(v\). Hence the Hecke isogeny \(\psi_A\) belonging to \(A\le k_w^n\) has, on these labels, kernel \(A\); its multiplier is \(q^{\dim A}\). Its other kernel labels and their annihilators are controlled by the conjugate-local and central factors in its finite adelic representative.\formalnote{16.6}{6}

The following one-parameter calculation is what makes section counting possible. At the identity of the special fibre the signature condition gives a line of tangent vectors \(L(a)\) with \(L(1)\ne0\), on which \(\mathcal O\) acts by its residue map. The image of \(\bar h\) on the whole tangent space lies on this line. Indeed in the signature chart its matrix is \(Y\otimes\xi\), with \(Y\) having a unit coordinate and every endomorphism acting on \(Y\) by its residue. Since \(h\) and \(\bar h\) are coprime, \(\bar h\) acts invertibly on every \(h^a\)-kernel. Translating to the identity shows that at any point of such a kernel its relative tangent space has dimension at most one. In a finite affine presentation of that kernel choose a function \(T\) whose derivative detects its tangent space. Given a point \(\gamma\) over a local \(k\)-algebra \(B\) reducing to \(o\), put \(u=o(T)-\gamma(T)\in B\). For a finite family of distinct labeled torsion sections let \(N(o)\) be the number specializing to \(o\).

On this test point the ideal of the union of sections is measured by \(u^{N(o)}\). For the divisibility in one direction, take a section \(\theta_j\) passing through \(o\). In the local ring the function \(t_j=\theta_j(T)-T\) generates its section ideal: its class generates the relative conormal fibre by the tangent test, and Nakayama applies. Equivalently, before localizing, for any element \(a\) of the section ideal a multiplier \(\rho\) nonzero at \(o\) makes \(\rho a=t_j a'\). If \(a\) vanished on the other sections through \(o\), then so does \(a'\). Indeed distinct sections into the valuation domain cannot have equal \(T\)-values here: otherwise the same divisibility would identify their values on every function after cancelling \(\rho\). The image of \(\rho\) in \(B\) is a unit and \(\gamma(t_j)=u\). Removing the sections one at a time proves \(u^{N(o)}\mid\gamma(a)\). In the other direction take the product of the \(t_j\) for the sections through \(o\). For each remaining section choose a function vanishing on that section but nonzero at \(o\), and multiply by it. The resulting function vanishes on them all and evaluates to \(u^{N(o)}\) times a unit. Flatness of the finite kernel over the valuation ring, and density of its characteristic-zero torsion sections, identify vanishing on those kernel sections with vanishing on its defining equations. These two divisibilities can therefore also test whether an isogeny kills a point after a change of the test algebra.\formalnote{16.7}{7}

Some nonzero \(w\)-torsion label specializes to zero. To see this take \(c\in w\setminus\{0\}\), \(t\equiv1\bmod w\) with \(tw\subset(c)\), supplied by goodness. Multiplication by \(t\) carries the labels killed by \(c\) into \(w\)-torsion and fixes a \(w\)-torsion test point. Apply the local equation above to all the \(w\)-labels: at a \(w\)-torsion test point \(z\) over \(B\) it gives \(u^{N(o)}=0\) since \(c z=0\). If \(o=0\), it also implies that killing \(u\) kills \(z\), by applying the section equations to the zero section alone. Now take \(B=k[\varepsilon]/(\varepsilon^2)\) and \(z=L(1)\), whose closed point is zero. If only the zero label specialized there, \(N(0)=1\) would give \(u=0\) already in \(B\), and hence \(z=0\). This contradicts the nonzero unit coordinate of \(Y\). Thus there is a nonzero label, say \(t_0\), reducing to zero. Let \(W=k_wt_0\).\formalnote{16.8}{8}

\subsection{Identification of marked points from multiplicities}
Write \(g_A:X_0\to X_A\) for the reduction of \(\psi_A\) followed by the Frobenius map which raises coordinates to the power \(q^{n-\dim A}\). All \(g_A\) have the same polarization multiplier \(q^n=p^{dn}\) and the same map \(\mathcal O\to k\), since powers of \(q\) fix \(k_w\). If \(A+W=A'+W\), their sections have the same reduction image \(H\), and for a point \(o\) in that image
\[
       q^{n-\dim A} N_A(o)=q^n/|H|
                         =q^{n-\dim A'}N_{A'}(o).
\]
Both counts vanish off \(H\). This already gives equality of the powers in the local equation after Frobenius. To justify comparison of the entire kernels including the component at the conjugate prime, put \(s=dn\), \(c=\bar h^s\). The prime-to-\(w\) calculation on the Hecke representative gives \(p^s\) killing all the kernel labels and \(c\) carrying each of them into its indicated \(w\)-label family. The twisted kernels are also killed by \(p^s\); this uses the common polarization multiplier and the nondegeneracy test for the theta pairing. Since \(p\) divides the integer \(h\bar h\), \(h^s c\) kills each twisted kernel.

Take a point \(u\) in \(\ker g_A\) over an arbitrary finite local test algebra. Apply the common local parameter to \(cu\), which is killed by \(h^s\). The family of \(A\)-labels contains the multiples by \(c\) of the old kernel labels. The vanishing test for \(u\) after the first Frobenius therefore gives \(\theta^{q^{n-\dim A}N_A(o)}=0\) for that parameter \(\theta\) at the closed point of \(cu\). By the displayed equality the power for \(A'\) vanishes too. The \(A'\)-labels are killed by \(\psi_{A'}\), so the converse part of the same local test yields \(c g_{A'}(u)=0\). Finite flat quotient descent now gives homomorphisms \(b_A,b_{A'}\) with
\[
             b_A g_A=c g_{A'},\qquad b_{A'}g_{A'}=c g_A .
\]
The second equality follows by interchanging the two families. Testing on local rings here retains the infinitesimal directions in the kernels.\formalnote{16.9}{9}

To remove \(c\), take \(u\in\ker g_{A'}\). The second factorization gives \(c g_A(u)=0\). The first, together with the common polarization multiplier, gives an identity in the additively written group \(\operatorname{Pic}(T)\). For every base change \(a:T\to\operatorname{Spec}B\) of the test algebra and every \(v\in X_A(T)\), it reads
\[
       \langle c v, g_A(a^*u)\rangle_A
                  =\langle b_A(v),g_{A'}(a^*u)\rangle_{A'}=0 .
\]
It suffices to check the identity after pulling \(v\) back through the faithfully flat isogeny \(g_A\); there both sides follow from the same pullback of the polarization. The adjoint of \(c\) is \(h^s\). The nondegeneracy criterion for the theta coordinates on \(p^s\)-torsion gives \(h^s g_A(u)=0\). B\'ezout for \(c,h^s\) now gives \(g_A(u)=0\). Thus \(g_A=j g_{A'}\) for an additive morphism \(j\). It commutes with \(\mathcal O\) because the isogenies do, and pulling back along \(g_{A'}\) shows that \(j\) preserves the polarization pairing. It preserves the level points because both original isogenies preserve them and the reductions of their labeled sections are compatible with the twists. Rigidity of the theta embedding with spanning coordinates identifies the targets over the base. Finally \(g_A=jg_{A'}\) identifies every marked image of a section from the common source. Hence the entire reduced point is a function of \(A+W\). Subspace M\"obius cancellation in \S\ref{sec:seed-proof} applies to these points and gives the desired cohomological relation.\formalnote{16.10}{10}

Sections~\ref{sec:theta-family} and~\ref{sec:theta-tangent} construct the theta family from the good-reduction models and use its torsion to bound the tangent space. The construction requires the arc models and their uniformizations; the identification of targets uses the polarization pairing and rigidity of spanning coordinates. The comparison of the resulting algebraic correspondences with cohomology is made in \S\ref{sec:seed-proof}.\formalnote{16.11}{11}


\section{A finite list of signed frames at each level}
\label{sec:signed-frames}

The finiteness used in \S\ref{sec:good-arcs} concerns the finite adelic frame. A frame records the Hermitian form and the finite adelic basis. We take frames up to the level equivalence described below. Even with this discrete datum fixed, the distinguished complex line varies in a ball and the labeled torsion points vary with it. We give a bound for the finite list specific to the signatures in the construction.

Let \(E\) be the CM field of the datum, \(F=E^+\), \(d=[E:\mathbb Q]=2s\), and \(n\ge1\) the rank. Let \(v_0\) be the distinguished infinite place. The construction assumes \(s\ge2\). Write \(\widehat O=\prod_{w<\infty}\mathcal O_{E_w}\). A signed frame consists of \(x\in\mathrm{GL}_n(\mathbb A_{E,f})\) admitting a rational hermitian form \(J\) with
\[
         x^*Jx=\nu I,\quad \nu\in\widehat{\mathbb Z}^{\,\times},
       \qquad \Phi_v(J)>0\quad(v\ne v_0),
 \tag{SF1}
\]
and with \(\Phi_{v_0}(J)\ell>0\) for its reflection \(\ell\). Write \(G_f\) for the finite adelic similitudes of the identity form with scalar in \(\mathbb A_{\mathbb Q,f}^{\times}\). For each open subgroup \(V\subseteq\mathrm{GL}_n(\mathbb A_{E,f})\), the set of double cosets
\[
           \mathrm{GL}_n(E)\,x\,(G_f\cap V)
           \quad\text{occupied by signed frames}
 \tag{SF2}
\]
is finite, uniformly over the rational forms \(J\) that can occur. Only the definiteness at one \(v\ne v_0\) will enter the count. We prove the claim for principal congruence levels, which suffice since every such \(V\) contains one.\formalnote{16.12}{12}

\subsection{A bounded rational basis}
At each infinite \(v\) diagonalize the nonsingular hermitian matrix \(J_v=\Phi_v(J)\) and put \(a_v=|J_v|^{-1/2}\), taking absolute values on its real eigenvalues. Then \(a_v^*J_v a_v\) is a hermitian involution of operator norm one. For \(g=(a,x)\) the matrix \(g^*Jg\) has determinant of absolute value one at each infinite place and is a unit scalar matrix at each finite place. The product formula for \(\det J\) and invariance of the adelic norm under conjugation give \(|\det g|_{\mathbb A_E}=1\).\formalnote{16.13}{13}

Put \(M_x=E^n\cap x\widehat O^n\) and let \(\mathcal L=a^{-1}\Phi(M_x)\subset(\mathbb C^n)^s\), where \(\Phi\) takes the chosen complex embedding at each infinite place. Use ordinary Lebesgue measure on the real and imaginary coordinates. This is a real lattice of dimension \(D=dn\) and covolume
\[
       c_{E,n}=(2^{-s}\sqrt{|\operatorname{disc}E|})^n.
 \tag{SF3}
\]
Indeed the Minkowski embedding gives the displayed value for \(\mathcal O_E^n\) from the discriminant determinant. Replacing that module by \(M_x\) multiplies covolume by \(|\det x|_f^{-1}\). To see this, clear the finitely many denominators and use the indices of the local elementary divisors of \(x\) at the remaining primes. The map \(a^{-1}\) contributes \(\prod_v|\det a_v|^{-2}\), so the combined factor is \(|\det g|_{\mathbb A_E}^{-1}=1\).

For \(0\ne u\in M_x\), the scalar \(h=u^*Ju\) is in \(\mathcal O_F\). This follows at finite places by inserting \(u=xz\) with \(z\in\widehat O^n\) into (SF1). At a definite infinite place it is strictly positive, so \(h\ne0\). At every infinite \(v\),
\[
       |\Phi_v(h)|\le\|a_v^{-1}\Phi_v(u)\|^2 .
\]
If \(\|a^{-1}\Phi(u)\|<1\), all these values have modulus less than one. This contradicts \(|N_{F/\mathbb Q}h|\ge1\). Hence the first Euclidean successive minimum of \(\mathcal L\) is at least one. Write \(b_D\) for the volume of the Euclidean unit ball in \(\mathbb R^D\). Minkowski's second theorem \cite[Theorem~1.3]{legacy-henk} gives
\[
         \lambda_D\le\lambda_1\cdots\lambda_D
                   \le 2^D c_{E,n}/b_D .
\]
Take \(R=1+2^D c_{E,n}/b_D\). The vectors of \(\mathcal L\) of length at most \(R\) span its real vector space. Their preimages therefore span \(E^n\) over \(E\), since the real span of a proper rational \(E\)-subspace has dimension at most \(d(n-1)\). Choose \(n\) independent preimages as columns of \(Y\in\mathrm{GL}_n(E)\). This step permits a sublattice of finite index in \(M_x\).
\phantomsection\label{route:32}\leanroute{The formal proof reduces a varying Hermitian form using a rational basis with a distinguished isotropic block. Definiteness at one place forces that block to be empty. Here definiteness gives \(\lambda_1\ge1\) directly. Minkowski's second theorem then bounds the lengths of \(n\) independent columns, which suffice for the finite enumeration below.}

\subsection{Gram matrices and residue classes to enumerate}
Set \(\beta=Y^*JY,\ D_f=x^{-1}Y,\ y=Y^{-1}x=D_f^{-1}\). Since the columns of \(Y\) belong to \(M_x\), \(D_f\) is integral. At finite places \(\beta=D_f^*\nu D_f\), so every entry of \(\beta\) is integral. At each infinite place its entries have absolute value at most \(R^2\) by Cauchy--Schwarz and the operator bound on \(a_v^*J_v a_v\). Thus \(\beta\) belongs to the finite set displayed below.
\[
 \begin{split}
   H&=\{z\in\mathcal O_E: |\Phi_v(z)|\le R^2\ \text{for all }v\},\\
   \mathcal H&=\{\beta\in M_n(H):\beta^*=\beta,\ \det\beta\ne0\}.
 \end{split}
\]
The first enumeration is over a bounded region in the Minkowski lattice of \(\mathcal O_E\). Let \(m\) be the product of the positive integers \(|N_{E/\mathbb Q}(\det\beta)|\) for \(\beta\in\mathcal H\). For each such \(\beta\), its determinant divides \(m\) in \(\mathcal O_E\). The adjugate formula therefore gives \(m\beta^{-1}\in M_n(\mathcal O_E)\). Since
\[
                y=\beta^{-1}D_f^*\nu,\qquad y^{-1}=D_f,
\]
both \(my\) and \(y^{-1}\) are integral. In particular the same denominator works for every frame, including when its initial \(\mathcal O_E\)-module is nonfree.

For a nonzero ideal \(\mathfrak m\) let \(V(\mathfrak m)=\ker[\mathrm{GL}_n(\widehat O)\to\mathrm{GL}_n(\mathcal O_E/\mathfrak m)]\). Associate to a frame the pair \((\beta,my\bmod m\mathfrak m\widehat O)\) arising from the chosen columns. Suppose \(y_1,y_2\) have the same pair. Then \(y_2-y_1\in M_n(\mathfrak m\widehat O)\), and both \(y_i^{-1}\) are integral. It follows that \(y_1^{-1}y_2\) and its inverse are integral and congruent to one modulo \(\mathfrak m\). Also \(y_i^*\beta y_i=\nu_i I\), so \((y_1^{-1}y_2)^*(y_1^{-1}y_2)=(\nu_2/\nu_1)I\). These two assertions place \(y_1^{-1}y_2\) in \(G_f\cap V(\mathfrak m)\). The two original \(x_i=Y_i y_i\) therefore represent the same double coset. Each occurring pair accounts for at most one coset. Consequently
\[
      \#\{\text{occupied cosets at }\mathfrak m\}
          \le |\mathcal H|\,N_{E/\mathbb Q}(m\mathfrak m)^{\,n^2}.
 \tag{SF4}
\]
This is an upper bound from a finite candidate list. The source proves the finiteness it needs by a more general adelic reduction theorem allowing a rational basis with an isotropic block. Its diagonal entries would vanish. Definiteness at a second infinite place makes that block empty. Its remaining compact set of finite frames and finite list of Gram matrices give the same double-coset argument.\formalnote{16.14}{14}

\subsection{The point count on the PEL model}
Take \(V=V(\mathfrak m)\) for the nonzero level in \S\ref{sec:pel-covers}. Choose one analytic datum in each occupied class of moduli points. If its frame is \(x'=z x k\), with \(z\in\mathrm{GL}_n(E)\), \(k\in G_f\cap V\), transport \(\ell\) to \(\Phi_{v_0}(z)\ell\Phi_{v_0}(z)^{-1}\) and the marked vectors by \(z\). The new datum is in the ball family for \(x'\), since that family admits every allowed distinguished line with that frame and arbitrary marked coordinates. Its comparison with the original has level element \(x'^{-1}zx=k^{-1}\). This proves that finitely many analytic families suffice for every fixed number of marks. Neatness, used elsewhere for stabilizers, is unnecessary for the count.\formalnote{16.15}{15}

Passing from these families to the scheme points on the theta closure requires two more arguments. On each connected ball the countable-base specialization argument of \S\ref{sec:good-arcs} gives a generic member. To identify the points of level-isomorphic parameters, one first constructs the theta family on the closure of just those two points, where finiteness of the dense parameters is automatic. Their common level model identifies their theta evaluations there. The generic members of the finitely many balls now have dense union of closures in the full theta base.\formalnote{16.16}{16}

The count supplies the finite set used in the valuative arc argument. At an arc centre the assertions that the affine theta equations cut out the required ideal in each chart with the prescribed ranks \(\sigma_j\) in degrees \(r^j\), and that indices \(o,-o,s'-o\) can be chosen simultaneously with nonvanishing coordinates for each dual-torsion label \(s'\), require good reduction of the theta data. The addition and translation identities then glue the abelian scheme, its level sections and its endomorphisms. Smoothness of the base of dimension \(n-1\) additionally uses the tangent bound and immersion of the complex balls. These are the respective inputs of \S\S\ref{sec:theta-family} and~\ref{sec:theta-tangent}.\formalnote{16.17}{17}


\section{The polarized analytic family}\label{sec:geometry-closure}
\subsection{A polarized ball with all its marked points}
Fix one of the signed frames \(g\) of \S\ref{sec:signed-frames}, with rational Hermitian matrix \(\beta\), and write \(H=\Phi_0(\beta)\) at the distinguished place. Let \(h(y,w)=\bar y^tHw\). At a base point the reflection \(J\) has a one-dimensional \(+1\) eigenspace and \(HJ\) is positive definite. Since \(J\) is \(h\)-self-adjoint and \(J^2=1\), its two eigenspaces are perpendicular. On the \(+1\) line \(h\) is positive, and on the complementary hyperplane \(h\) is negative definite. Thus for \(0\ne u\) on the positive line
\[
                   Jw=2u\,h(u,w)/h(u,u)-w .
\]
Conversely each positive line gives such a reflection. The finite frame and distinguished signature determine the rational form. Indeed, two such forms are carried by \(g\) to scalar matrices with scalars in \(\widehat{\mathbb Z}^{\,\times}\). Their ratio is consequently a finite adelic unit. Comparison of one nonzero matrix entry shows that this ratio belongs to \(E\) and is fixed by complex conjugation. At the finite places the ratio is a single rational local scalar at every component above each rational prime, so it belongs to \(\mathbb Q\): expand the global element in a rational basis for \(E\) starting with \(1\), and compare the coefficients after tensoring with \(\mathbb Q_p\). Positivity of the Hermitian forms after applying the fixed reflection gives positive sign. A positive rational number that is a unit at every prime is \(1\).

Put \(a=h(u,u)>0\). Choose a complex coordinate isomorphism \(e\) on the conjugate vectors with \(e(\bar u)=(1,0)\) and \(e(\bar y)_1=h(y,u)/a\). Its second factor identifies the conjugate of \(u^\perp\) with \(\mathbb C^{n-1}\). All positive lines have a unique representative \(\bar x\) with \(e(x)=(1,b)\), since the kernel of the first coordinate is negative definite. Their domain is \(a+h(\bar x-u,\bar x-u)>0\), an open ellipsoid containing \(0\). For brevity write \(s_y=e(\bar y)_1\) for a vector at the distinguished place. The period map \(Q_b\) on the real Minkowski space and its \(\mathcal O_E\)-lattice \(L_g=E^n\cap g\widehat{\mathcal O}_E^n\) is the map of \S\ref{sec:ball-differential}. Its invertibility there gives compact complex tori \(V/Q_bL_g\). For a fixed real period \(w\), \(Q_bw\) is affine holomorphic in \(b\), so the lattice translations act holomorphically on the product of the ellipsoid and \(V\). They act properly discontinuously, since on each compact set of parameters the inverses \(Q_b^{-1}\) are uniformly bounded. The quotient is a holomorphic family of complex tori.\formalnote{16.18}{18}

We check the polarization as the complex structure varies. Write \(c_v=2\pi\operatorname{Im}\Phi_v(\xi)>0\) for the scalar of the PEL datum. For a fibre vector with distinguished component \((p_1,p')\) and remaining components \(p_v\), use the form
\[
 A(y,p)=c_0\{s_{y_0}p_1-e^{-1}(0,p')^tH y_0\}
              +\sum_{v\ne0}c_v\,\bar y_v^t\Phi_v(\beta)p_v .
\]
It is real-linear in \(y\), complex-linear in \(p\). At \(p=Q_bw\) its distinguished summand divided by \(c_0\) is
\[
 s_{y_0}h(\bar x,w_0)+s_{w_0}h(\bar x,y_0)-h(w_0,y_0).
\]
The first two terms cancel on antisymmetrizing. Summing all places gives \(A(y,Q_bw)-A(w,Q_by)=2\pi i\sum_v2\Re(\Phi_v(\xi)\,\bar w_v^t\Phi_v(\beta)y_v)\). The sum here is the prescribed trace pairing. It is integral on \(L_g\) by the finite frame and codifferent condition in \S\ref{sec:good-arcs}. For \(w=y\), the real part of the distinguished summand divided by \(c_0\) is
\[
 |s_{y_0}|^2 h(\bar x,\bar x)
       -h(y_0-\bar s_{y_0}\bar x,y_0-\bar s_{y_0}\bar x).
\]
The subtracted vector belongs to the negative hyperplane. The expression is positive for \(y_0\ne0\), and the other-place summands are positive for \(y_v\ne0\). These are the Riemann relations. The adjoint of multiplication by \(c\in E\) is multiplication by \(\bar c\), as is seen by substituting in the two summands of \(A\). On the distinguished components of \(Q_b w\), \(c\) acts by \(\Phi_0(c)\) and its conjugate, and elsewhere by \(\Phi_v(c)\). This gives exactly the specified signature and the polarization of the required type at every parameter. Applying the theta coordinates for this Riemann form gives the family used in \S\ref{sec:theta-family}; the local Gaussian convergence justifying the holomorphic coordinates is in \S\ref{sec:ball-differential}.\formalnote{16.19}{19}

The level labels are constant classes in the abstract rational period lattice modulo \(L_g\). Their images \(Q_b w\) therefore give holomorphic torsion sections. An additional marked vector based at \(P_i\) has parameter \(z_i\in V\) and real representative \(P_i+Q_b^{-1}z_i\). This reaches every marked point and gives \(\nu\) independent fibre variables. A finite Hecke translate \(g\gamma\) leaves the periods unchanged and rescales the form by the positive inverse rational similitude factor. This verifies the parametrization used to pass from a finite set of frames to a finite set of families.\formalnote{16.20}{20}

At a coincidence let \(Z\) be the rational isomorphism, let \(Z_0u=\mu u\) for \(Z_0=\Phi_0(Z)\), and write \(Z_0^* H Z_0=cH\). Then \(c=|\mu|^2\). Both the line and \(u^\perp\) are stable, so \(e\overline{Z_0}e^{-1}=\bar\mu\oplus N\). Applying the formula for \(Q_b\) shows that the new base coordinate is \(b'=Nb/\bar\mu\). The map from the \(i\)-th marked fibre to the new fibre has blocks
\[
 \Phi_0(\alpha_i)\mu,\qquad \overline{\Phi_0(\alpha_i)}N,\qquad
       \Phi_v(\alpha_i)\Phi_v(Z)\quad(v\ne0).
\]
Indeed \(\bar x'=(Z_0\bar x)/\mu\), hence \(h(\bar x',Z_0w)=\mu h(\bar x,w)\), and \(e(\overline{Z_0 w})_2-e(\overline{Z_0 w})_1 b'=N(e(\bar w)_2-s_w b)\). The constant marked shifts agree modulo the lattice by the level-isomorphism condition at the base point. Thus the two formulae intertwine the correspondences for nearby parameters, and differentiating gives these very blocks and \(N/\bar\mu\) on the base. The eigenvalues of \(N\) are the conjugates of the remaining roots at the distinguished place. All these roots \(r\) have \(|r|^2=c\): the isomorphism preserves the positive form \(HJ\) up to \(c\), since it preserves the reflection. Consequently \(\mu X-r=-r(1-\bar r X/\bar\mu)\). Multiplying the base factors and the indicated fibre factors gives \(p(X)=p(0)\det(1-X\,dZ)\) for the marked local polynomial used in \S\ref{sec:marked-orbits}. This proves the comparison of the two tangent polynomials in \S\ref{sec:pel-covers}.\formalnote{16.21}{21}



\section{Detecting the ball parameter from moving torsion points}\label{sec:ball-differential}
For the lower dimension bound in \S\ref{sec:theta-tangent}, a family with \(n-1\) complex parameters must still have this dimension after passage to the theta coordinates. We prove injectivity of the differential. It suffices to move torsion points accumulating at the origin of the uniformizing torus.

Fix the frame and the embedding of mixed signature. Choose complex-linear coordinates \(e:\mathbb C^n\simeq\mathbb C\oplus\mathbb C^{n-1}\) with the central positive line represented by \((1,0)\) after conjugation. The nearby lines are represented by \(\bar x\), \(x=e^{-1}(1,b)\), for the parameters \(b\) with \(x^{t}H\bar x>0\). The restriction of the Hermitian form \(H\) to the complementary hyperplane is negative definite, so these parameters contain an open ball. For a real period vector \(w=(w_0,(w_s)_{s\ne0})\), write \((y_1,y_2)=e(\bar w_0)\). In the splitting at the chosen embedding, the period coordinates are
\[
               Q_b(w)=\left(x^{t}Hw_0,\ y_2-y_1 b,\ (w_s)_{s\ne0}\right).
\]
The real map \(Q_b\) is invertible. Indeed the vanishing of the second component forces \(w_0\) into \(\mathbb C\bar x\), where the nonzero first component on \(\bar x\) detects its scalar. All other components are unchanged. The explicit form in \S\ref{sec:geometry-closure} gives the polarization on the transported lattice. The displayed formula is affine in \(b\) for each \(w\).\formalnote{16.22}{22}

Let \(v\) be a tangent vector at \(b\), put \(T_{b'}=Q_{b'}Q_b^{-1}\), and write \(T'\) for its derivative in direction \(v\) at \(b\). This is a real-linear map of the fibre. If it is complex-linear, then \(v=0\). To see this take \(w_0=\bar x\) and all other components zero. The second component of \(Q_b(w)\) vanishes, while that of \(T'(Q_b(w))\) is \(-v\). At \(iw\), the first component of \(Q_b\) is multiplied by \(i\), the second still vanishes, and the derivative of the second is \(iv\), by conjugation in \(y_1\). Hence complex linearity would give \(iv=-iv\).\formalnote{16.23}{23}

Let \(x(b')\) be the point of the base and let \(\psi(b',z)\) be the image of the complex uniformizing vector \(z\) in its abelian fibre. The theta-coordinate family makes \(\psi\) holomorphic in both variables on a neighbourhood of each point in question. In analytic coordinates this is the usual normal convergence of the theta series: on a compact subset of the ball the positive form has a uniform positive lower bound on the unit sphere, while the linear part in the summation variable is bounded by a constant times its norm for bounded \(z\). A summable majorant is then a polynomial times \(e^{-a\|w\|^2+C\|w\|}\) for some \(a>0\), also after differentiation. Local holomorphic trivializations of the theta bundle give the same projective point where a chosen coordinate is nonzero. Composing \(\psi(b',0)\) with projection to the base proves that \(x\) is holomorphic.\formalnote{16.24}{24}

Suppose \(x\) is constant to first order in direction \(v\). For a period \(l\) of the fibre at \(b\) and an integer \(k\ge1\), the point \(\psi(b',T_{b'}l/k)\) is \(k\)-torsion. Its multiple is the zero section above \(x(b')\) and is therefore constant to first order. The multiplication-by-\(k\) map on this abelian scheme in characteristic zero is étale, so the torsion point itself is constant to first order. For any regular function near \(\psi(b,0)\) let \(g\) denote its composite with \(\psi\). For all sufficiently large \(k\) we have
\[
                 dg_{(b,l/k)}(v,T'l/k)=0 .
\]
The function of \(t\) given by \(dg_{(b,tl)}(v,tT'l)\) is holomorphic near \(0\) and vanishes at \(1/k\) for all large \(k\). Differentiating at zero gives
\[
                dg_{(b,0)}(0,T'l)=-d^2g_{(b,0)}((0,l),(v,0)).
\]
The period lattice spans the fibre over \(\mathbb R\), and both sides are real-linear in \(l\), so this equality holds for every fibre vector \(l\). Its right side is complex-linear in \(l\). Consequently \(dg_{(b,0)}(0,T'(il)-iT'l)=0\) for every such \(g\). Regular functions through the smooth fibre point detect its tangent vectors, and the differential of the fibre uniformization at the origin is an isomorphism. Hence \(T'\) is complex-linear. The positive-line test gives \(v=0\), as required.\formalnote{16.25}{25}

For a fibre power with additional marked points, a tangent vector in the kernel first projects to one in the base kernel, so its base component is zero. Fixing \(b\), the differential of the \(i\)-th marked point is the differential of the torus uniformization on its own coordinate. Each remaining component is therefore zero. This proves the immersion for the actual marked family as well. Together with the upper tangent bound in \S\ref{sec:theta-tangent}, it explains why the relative dimension obtained there is exactly \(n-1\) on the base.\formalnote{16.26}{26}


\section{The integral family from theta coordinates}\label{sec:theta-family}

Section~\ref{sec:pel-covers} uses an abelian scheme over a closure of
theta-null data. There are two closures. The evaluations at the finitely
many labeled torsion points define the base \(S\) in a product of
projective spaces over the ring \(R\) of safe primes. Over \(S\), the
evaluations at \emph{all} arguments of theta define a second closure
\(\mathcal A\subset S\times\mathbb P^{|T|-1}\). We prove why this
second closure is flat, why addition extends to it, and why its
geometric fibres are abelian varieties.

\subsection{Theta equations on the special fibre}

Fix a complex fibre with vector space \(V\), period lattice
\(\Lambda\), automorphy form \(B\) and multiplier \(\chi\).
Write \(T=\Lambda_B^\vee/\Lambda\) for the finite group of polarization
torsion and \(\theta_s\), \(s\in T\), for the normalized translates of
the theta section obtained from the chosen Lagrangian. The coordinates
generate the linear system and may repeat. Let \(g=\dim_{\mathbb C}V\),
\(r\ge3\) be the padding integer of the PEL datum, and \(i_0\) its
degree cutoff. Put \(m_j=r^j\). The group
\[
 T_j=\Lambda_{m_jB}^\vee/\Lambda
       =\{u\in(V/\Lambda)_{\rm tors}:m_j u\in T\}
\]
has order \(m_j^{2g}|T|\), since multiplication by \(m_j\) is
surjective on the torsion of a complex torus with kernel of order
\(m_j^{2g}\). The theta dimension formula gives
\(\sigma_j=\dim_{\mathbb C}\Theta(m_jB,\chi^{m_j})
=\sqrt{|T_j|}=m_j^g\sqrt{|T|}\).
The level kills \(T_j\) for \(j\le i_0\). The other property of
\(i_0\) is combinatorial. A monomial ideal with at most \(\sigma_j\)
standard monomials of degree \(r^j\) for \(j\le i_0\) satisfies the
same bounds for every \(j\). Such a cutoff exists by noetherian induction.
If an ideal fails one bound, any larger ideal satisfying it must
contain one of the finitely many monomials left out in that degree.
Apply the induction hypothesis to each ideal obtained by adjoining
one such monomial and take the maximum of the resulting cutoffs.\formalnote{16.27}{27}

Let \(v\) be the multiplicative, nonarchimedean valuation on
\(\mathbb C\) in a valuative arc, with ring \(O\) and residue field
\(k\). Its convention is \(v(O)\le1\). For a theta function \(F\)
of type \((B,\chi)\) set
\[
 h_F(z)=v\bigl(e^{-B(z,z)/2}F(z)\bigr),\qquad
 \|F\|_B=\sup_{z\ {\rm torsion}\bmod\Lambda}h_F(z).
\]
The exponential and the evaluation take place in \(\mathbb C\)
before applying \(v\). The arithmetic good-prime condition makes
\(r,|T|\) units and gives the identity
\[
            \|FG\|_{B+B'}=\|F\|_B\,\|G\|_{B'}.
\]
Here the two suprema need not be attained at the same torsion point.
We give the reason that their product is the norm of the product.

The good prime \(\mathfrak p\) contains \(p,\beta-b\) for a rational
prime \(p\) and \(b\in\mathbb Z\). The elements \(\beta-b\) and
\(\bar\beta-b\) are coprime. A distinct prime \(\mathfrak q\)
also contains \(p,\beta-b\). The definition further requires
\((\mathfrak N)+\mathfrak p(p,\bar\beta-b)=\mathcal O_E\).
Choose a generator \(a\) of a high principal power of
\(\mathfrak q\), large at all archimedean embeddings, and take
\(c=\bar a\). Then \(c\) divides a power of \(p\) and is coprime to
\(\bar c\), since \(\mathfrak q\ne\bar{\mathfrak q}\).
Multiplication by \(c\) has arbitrarily small inverse in the complex
operator topology. Multiplication by \(\bar c=a\) has determinant of
\(v\)-value one on \(V\). At the chosen embedding this follows from an
element of \(\mathfrak q\setminus\mathfrak p\). At its conjugate and
at the remaining embeddings occurring in the determinant it follows
from the values of \(\beta-b\) and the CM type of the PEL datum.
The Hermitian form used in this calculation is
\((B+i^*B)/2\), where \(i\) acts by multiplication on \(V\).
It has the same alternating Riemann form as \(B\).\formalnote{16.28}{28}

For a nonnegative size function \(h\) say that a shift \(w\) has
bound \(C\) in its degree if all such \(h\) satisfy
\[
                  (\sup h)^2\le C\sup_z h(z)h(z+w).
\]
The theta transformation input on the kernel of multiplication is
as follows. After taking powers of \(c,\bar c\) that act trivially on
the relevant dual quotient, put \(d=\bar c\) and denote normalized
translation by \(U\). For the canonical multiplier the identity is
\[
    \sum_{z\in d^{-1}\Lambda/\Lambda}
         U_{dz}(F\circ c\circ d^{-1})=\mu F,\qquad
         \mu^8=(\det\nolimits_V d)^8 .
\]
The scalar has \(v(\mu)=1\). If \(cw\in\Lambda\), then \(p^kw\in
\Lambda\) for some \(k\). A B\'ezout relation between \(c,d\)
writes \(w=dw'\) modulo \(\Lambda\) with \(cw'\in\Lambda\).
Subtract the translated identity from the first one. Each error
term has a factor \(\zeta-1\) with \(\zeta^{p^k}=1\), hence
\(v(\zeta-1)\le v(p)^{1/p^k}=\rho<1\).
For a size \(h\) this gives
\(h(y)\le\max\{h(y+w),\rho\sup h\}\).
Taking points arbitrarily close to the supremum gives bound \(1\)
for these kernel shifts in the degrees covered by the identity.

For the descent, let \(H=c^{-1}\Lambda/\Lambda\), and suppose the
shift \(cw\) has bound \(D_0\) in degree \(B\). For a section \(t\)
of degree \(c^*B\), decompose its pullback along \(c^{-1}\) into
the character sums
\[
 s_\xi=\sum_{k\in H}\xi(-k)\,U_{ck}(t\circ c^{-1}).
\]
Each has degree \(B\) and a finite-order multiplier. Character
orthogonality gives \(\sum_\xi s_\xi=|H|\,t\circ c^{-1}\).
In the product of its two shifted sums, a typical pair of terms
has size \(h_t(y+k)h_t(y+w+l)\). The difference \(l-k\) is a
kernel shift with bound \(1\), so this is at most
\(\sup_z h_t(z)h_t(z+w)\) by the periodic log-convexity argument
below. The ultrametric inequality and the bound for \(cw\) now give
a bound \(D_0C\) for \(w\), where \(C=v(|H|)^{-2}\).
It also applies in every degree \(A\) with \(c^*B-A\) positive
by the comparison of positive theta degrees.

An initial finite bound for each shift comes from evaluations on
the space of degree \(2B\) sections. For \(w\) fixed, consider the
evaluation functionals at \(y\) for which \(2y-w\) has bound \(1\).
Their kernels have zero intersection: such \(y\) include the
midpoints of \(w\) with the points of the lattices \(c^{-1}\Lambda\)
of arbitrarily small mesh. A holomorphic section vanishing on these
points vanishes by continuity. Finite dimensionality expresses
evaluation at \(0\) as a finite linear combination of these normalized
evaluations. Apply this to \(y\mapsto e^{-B(x,x)}t(x-y)t(x+y)\);
the maximum of the \(v\)-values of the coefficients is an initial
shift bound. For the finitely many classes killed by the order of
\(w\) we can take one such bound \(D_1\).

Choose \(c\) with \(c^*B-2B\) positive. In any positive multiple
\(bB\), the same descent gives a bound \(D_1C\) in degree \(2bB\)
on these classes because \(c\) preserves them. Squaring sections
then improves \(D_1\) in degree \(bB\) to \(\sqrt{D_1C}\).
Iteration gives \(C\), independent of \(b\) and of the order of
\(w\). Positive multiples of \(B\) dominate any given degree.
Applying the bound in degree \(2^m A\) to the \(2^m\)-th power of
a theta function in degree \(A\) replaces the constant by its
\(2^m\)-th root. Letting \(m\) tend to infinity gives bound \(1\).
Finally fix independent arguments \(x,y\) for the sizes \(h_F,h_G\).
For \(d_0=y-x\) put
\(R_k=\sup_z h_F(z)h_G(z+kd_0)\). The bound \(1\) on the product
of translated functions gives \(R_{k+1}^2\le R_kR_{k+2}\) by
regrouping its four factors. The sequence is periodic since
\(d_0\) is torsion. Thus \(R_1\le R_0\), which gives
\(h_F(x)h_G(y)\le R_0\). Taking suprema proves the claimed
multiplicativity. The finite-kernel theta identity uses the Gauss
sum calculation.\formalnote{16.29}{29}

All \(\theta_s\) have the same positive norm \(M\). In degree
\(m=r^j\) the image of homogeneous polynomials under substitution
\(X_s=\theta_s\) is exactly
\[
 \{F\in\Theta(mB,\chi^m):\|F\|_{mB}\le M^m\}
 \quad\text{for coefficients in }O .
\]
In degree one the normalized translates of any one \(\theta_s\)
run through the coordinates up to units. Finite theta-group
orthogonality expresses a vector in its norm ball in the span of
these translates by evaluating the averaging identity at a
translate of maximal value. Division by the group order is integral
here. For the power step a translate of \(F^r\) by \(u\) in the
degree-\(rB\) dual lattice has the same norm as
\(F^{r-1}\cdot U_{ru}F\). On the Lagrangian lattice those two
sections have the same multiplier. Translates of the second section
needed for orthogonality are products of \(r\) sections in the
original ball, since the Lagrangian lattice contains
\(r^{-1}\Lambda\). Applying the same span calculation in that
lattice puts the first section in the product ball. Iteration gives
the displayed equality for \(m=r^j\).\formalnote{16.30}{30}

Normalize each vector \((\theta_s(z))_s\) by a coordinate of
largest \(v\)-value. Its entries \(q_s(z)\) are integral and one
is a unit, since the complex theta vector has no common zero.
In chart \(X_c=1\) set
\[
 I_c=\bigcap_{\bar q_c(z)\ne0}
 \ker\left(k[X_s:s\in T]\longrightarrow k,\quad
                 X_s\longmapsto\bar q_s(z)/\bar q_c(z)\right).
\]
For homogeneous \(P\) of degree \(m=r^j\) over \(O\), the criterion is
\[
       \bar P\in I_c
          \quad\Longleftrightarrow\quad
                 \|P(\theta)\|_{mB}<M^m .
\]
Here the needed evaluation lemma says that strict inequality is
equivalent to zero reduction on a set representing
\(\Lambda^\vee_{mB}/\Lambda\), using integral homogeneous
coordinates at those points. The norm ball is a finite
torsion-free module over the valuation ring, hence free. The
theta-group translates test its coefficient norm, and the
strictly smaller ball is its maximal ideal times this lattice.
For the restriction to chart \(c\), multiply \(P\) by
\(X_c^{\,r^{j+1}-r^j}\). The product vanishes at all reduced vectors
off the chart, and multiplicativity scales its norm and its
threshold by the same positive factor. Thus restriction to the
chart gives the stated criterion. Since every homogeneous polynomial
over \(k\) lifts in its own degree, reduction modulo \(I_c\) has
dimension \(\sigma_j\) in the image of degree \(r^j\). For
\(j\le i_0\) its kernel can be computed from the labeled level
points alone, by the inclusion of \(T_j\) in the level torsion.
The evaluations here use the homogeneous vectors of those level
points even if their \(c\)-coordinate is zero.\formalnote{16.31}{31}

\subsection{Flatness of the chart rings}

Take finitely many complex parameters whose images in \(S\) are
dense. They may map to generic scheme points through embeddings of
their residue fields in \(\mathbb C\). At each point of \(S\)
the arc argument of \S\ref{sec:pel-covers} gives a specialization
at which the preceding rank data hold. Flatness at that
specialization implies flatness at its generalizations. We describe
the essential local algebra. Let \(A\) be the local ring at the
specialization, with map \(\phi:A\to k\) to its geometric residue
field. Use charts of the finitely many level sections to represent
their coordinates by vectors \(P(\omega)\) over \(A\). The maps
\(\psi_i:A\to\mathbb C\) from the dense parameters whose images
generalize this point have zero common kernel. Write \(J\) for the ideal of
the reduced closure in the \(c\)-chart, read as
\(f\in J\) exactly when \(\psi_i(f)\in I_i\) in every generic chart.
The assertion is \(\phi(J)k[X]=I_c\) and \(A[X]/J\) flat over \(A\).

Fix a monomial order. Begin with \(J_0=J\cap\phi^{-1}(I_c)\),
where \(\phi\) acts on coefficients. Call a monomial leading if it
leads the nonzero reduction of a homogeneous form in \(J_0\).
Its coefficient lifts to a unit in \(A\). Division eliminates such
terms: the higher terms of a lift whose coefficients reduce to zero
must be accounted for as well. In each fixed degree, form the
matrix of eliminating relations for the leading monomials.
Modulo the maximal ideal its relevant triangular block has
invertible diagonal. This block remains invertible over \(A\),
so the standard monomials span modulo \(J_0\), also after base change.

For \(j\le i_0\) put \(D=\sigma_j\). Choose \(D\) standard monomials
independent in the special quotient. The finite level evaluation
test gives \(D\) level vectors where their evaluation determinant
is nonzero. If there were one further standard monomial \(m\),
append it to the columns and put the monomials themselves in
the last row. Expanding this determinant gives a polynomial
supported on the \(D+1\) chosen monomials. It belongs to
\(J_0\), since the evaluation matrix in each fibre has rank \(D\).
Its coefficient at \(m\) is the chosen nonzero minor up to sign.
Thus its reduction has a leading monomial, although all its
monomials were standard. This proves the bound in the initial
degrees, hence in all degrees \(r^j\) by the cutoff. The rank
in every generic and special quotient is already \(\sigma_j\),
so in each such degree the standard monomials are independent
there as well. A relation over \(A\) modulo \(J\) would have
all its coefficients zero under every \(\psi_i\), hence zero.

The equation \(X_c=1\) belongs to the chart in the form
\(X_c-1\in J_0\). Any finite set of polynomials may be homogenized
to a single degree \(r^j\) using \(X_c\) without changing its
images. Thus each finite set of elements in \(A[X]/J\) lies
in a free submodule spanned by the standard monomials at a
sufficiently high degree. These submodules form a filtered
system, giving flatness. The same division applied first to
a relation of \(J\) and then to one of \(I_c\) shows that their
reductions generate exactly \(I_c\): in each direction the
remainder is a relation among independent standard monomials.
Since \(I_c\) is an intersection of kernels of evaluations into
a field, the selected geometric fibre is reduced.
The chart rings come from the single reduced projective closure,
so their overlap maps are the ordinary changes of homogeneous
coordinates.\formalnote{16.32}{32}

\subsection{Recovering addition and the action}

For pairs \(p=((a_l,b_l))_l\) in \(\mathcal O_E^2\) satisfying
\(\sum_l\bar a_la_l=r,\ \sum_l\bar a_lb_l=0\), put
\[
      P_\mu(w,z)=\prod_l\theta_{\mu_l}(a_lz+b_lw).
\]
Let \(Q_\nu\) correspond to a second such tuple \((a'_l,b'_l)\).
When \(\sum_l\bar a_l\mu_l=\sum_l\bar a'_l\nu_l\), Riemann's relation
in the normalization of these coordinates is
\[
  |T|Q_\nu(w,0)P_\mu(w,z)=
  \sqrt{[T:T[r]]}\sum_{s\in T}
       P_{\mu-as}(w,0)\,Q_{\nu+a's}(w,z).
\]
Indeed both products in \(z\) have the same degree-\(r\) multiplier
on the lattice over \(T[r]\). The first sum condition gives the
degree, the second cancels the cross term in \(w\), and the equal
index sums give equal remaining characters. The dual of this
lattice in degree \(rB\) is the lattice over \(T\). Averaging
rank-one operators over its finite theta translations gives a
scalar by irreducibility. Taking the trace gives the dimension
\(\sqrt{[T:T[r]]}\) relative to the number \(|T|\) of translations
in the sum. Evaluating the resulting identity of matrix entries
at \(0,z\) yields the formula: the phases in the shifts by
\(-s\) and \(s\) cancel.\formalnote{16.33}{33}

A denominator must be a unit at the special point. If
\(M_0=\max_t v(\theta_t(0))\), then for every \(s\) some \(u\)
satisfies
\[
 v(\theta_u(0))=v(\theta_{-u}(0))
                =v(\theta_{s-u}(0))=M_0 .
\]
For the argument choose \(o\) attaining \(M_0\) and maximize
\(\min\{v(\theta_{s-u}(0)),v(\theta_{-u}(0))\}\), calling the
maximum \(D_0\). Apply Riemann's relation to the degree-\(r\)
products \(\theta_{ro-s}\theta_s\theta_0^{r-2}\) and
\(\theta_o^r\). In each shifted term the two distinguished
indices contribute a factor bounded by \(D_0\), while the other
constant factors are bounded by \(M_0\). The ultrametric inequality
and \(v(|T|)=1\) give
\(M_0\,h_{\theta_{ro-s}\theta_s\theta_0^{r-2}}(z)
 \le D_0 M^r\). The supremum on the left is \(M_0M^r\).
Thus \(D_0=M_0\), and the symmetry of coordinates under negation
gives the third equality. Translation identifies the sizes of
\(\theta_k\) at the level point \(l\) with those of
\(\theta_{l+k}(0)\) up to a common factor. Therefore each required
nonvanishing chart persists at the arc centre and its generalizations.\formalnote{16.34}{34}

To turn this identity into maps, add to the factors for \(p\) the
factors at \((0,b'_l)\), and interchange the two tuples on the other
side. Take tuples starting with two copies of \((1,0)\). These
give two constant factors at \(z=0\). To attain the required
sum \(s\) of their labels, put \(u,s-u\) there. Their product
is a unit in the chosen normalized coordinates. The coefficient
of the sum after division is the regular transition function for
the square of the zero-section coordinate bundle, multiplied by
\(\sqrt{[T:T[r]]}/|T|\). Interchanging the tuples gives a local
expression in the reverse direction as well.

Here is the gluing argument. Suppose one side is already given
by a line bundle \(M'\) with generating homogeneous sections
\(m_j\). The other coordinates \(w_l\) are locally regular
combinations of the \(m_j\). On the chart for \(m_k\) these give
regular functions \(w_l/m_k\). On overlaps their transition law
holds on the dense theta points and hence on the reduced scheme.
They glue to sections of \(M'\). The reverse combinations imply
that these sections have no common zero. They define a morphism
to projective space; the equations for the Segre-embedded power
of \(\mathcal A\) hold on the separating theta points and hence
everywhere. Projecting to each factor extends the maps given on
the analytic tori by the individual pairs \((a_l,b_l)\).
The equality of Picard classes of the two product bundles follows
from this construction.\formalnote{16.35}{35}

Concretely use the tuple
\((1,0),(1,0),(d,1),(1,-\bar d),(a_l,0)\), with
\(3+\bar d d+\sum_l\bar a_la_l=r\). For \(d=-1\) the \(a_l\)
can be chosen to span \(\mathcal O_E\) over \(\mathbb Z\).
Pair this with the tuple for \(d=0,\ a_l=1\).
On the known side the maps are projections or the zero section.
On the other side appear \(z+w,\ w-z,\ a_lz\). Subtraction
and the spanning condition give an endomorphism for every
\(a\in\mathcal O_E\). The labeled level coordinates already give
sections by the projective embedding; the section for label zero
is the identity. Flatness over the reduced base and density in
the complex fibres give separating families on both
\(\mathcal A^2_S\) and \(\mathcal A^3_S\). Thus the group and ring
axioms hold as identities of morphisms, by testing on the vector
spaces \(V\).\formalnote{16.36}{36}

At each selected geometric specialization the group fibre is
reduced over an algebraically closed field. Its smooth locus
is dense; translations from any smooth point show that the whole
fibre is smooth. Flatness and finite presentation then give
relative smoothness there. Properness allows a specialization in
the total space above the chosen specialization of a base point.
Openness of the smooth locus gives smoothness at the original
point. Multiplication by every positive integer \(N\) is surjective:
its closed image contains each dense theta point \(e_i(z)\) since
\([N]e_i(z/N)=e_i(z)\). On a smooth proper geometric group fibre
the finite component group therefore admits surjective multiplication
by its own order, forcing it to be trivial. This explains the
connectedness encoded in the source by divisibility. The analytic
uniformisations give relative dimension \(g\); the signature on
the Lie algebra uses invertibility of \(\bar\beta-\beta\). Applying
the same product-bundle comparison to the padding for \(\beta\)
gives the polarization class laws, including adjointness for
\(\beta\). Translation by labels gives the finite theta structure.\formalnote{16.37}{37}

This proves the existence of the theta-embedded abelian family over
the specified closure, assuming the PEL datum, the breaking of the
parameters into finitely many dense generic representatives, the
valuative arc theorem and the theta dimension, translation and finite-kernel identities in \S\ref{sec:theta-chart-gluing}. The source constructs the dense
representatives from finitely many ball families. Smoothness of the
\emph{base} \(S/R\) of dimension \(n-1\) needs an additional
deformation estimate at an arc centre, together with immersion of
the ball. The abelian-family argument supplies smoothness of
\(\mathcal A/S\). For the good primes, the valuation arcs and the comparison of integral sections see \S\ref{sec:good-arcs}. The construction also specifies which complex
fibres have the models above and how the base point is determined by
its level coordinates. A fine-moduli statement for arbitrary PEL
families would additionally require a classification over all test
schemes.\formalnote{16.38}{38}


\section{Relative dimension of the level base}\label{sec:theta-tangent}

Let \(S\) be the closure of the theta coordinates at the labeled level
points, with abelian scheme \(\mathcal A\) constructed in
Section~\ref{sec:theta-family}. Write \(n\) for the rank of the
hermitian space. On the complex side a negative line varies in an
\((n-1)\)-dimensional ball. In residue characteristic \(p\), torsion
sections can coincide. We must show that their collisions introduce
no additional tangent directions on \(S\). The proof measures a
deformation by its action on the \(\mathfrak w\)-torsion of one
abelian variety.

\subsection{The differential of the distinguished endomorphism}

Fix a geometric point \(x_0\) over an algebraically closed field \(k\).
Assume it is the specialization of a complex torus model over a
valuation subring of \(\mathbb C\), with all generic torsion points
labeled by sections. Put \(A_0=x_0^*\mathcal A\), let
\(\rho:\mathcal O_E\to k\) be the structure map and suppose
\(\mathfrak w=\ker\rho\) is a good prime. Put
\(F=\mathcal O_E/\mathfrak w\subset k\), \(q=|F|\) and
\(p=\operatorname{char}k\). Goodness supplies an integer \(b\) with
\[
 h=\beta-b\in\mathfrak w,\qquad a=\bar\beta-b,\qquad
 (h,a)=\mathcal O_E,\qquad
 (\mathfrak N)+\mathfrak w(p,a)=\mathcal O_E,
\]
where \(\mathfrak N\) denotes the level element. Choose
\(0\ne c\in\mathfrak w\) generating
\(\mathfrak w/\mathfrak w^2\). Nakayama's lemma applied to
\(\mathfrak w/(c)\) gives \(t\equiv1\bmod\mathfrak w\) such that
\(t\mathfrak w\subset(c)\).\formalnote{16.39}{39}

Let \(M\) be a finite-dimensional \(k\)-space, \(D=k[\epsilon]/(\epsilon^2)\)
and \(x\) a deformation of \(x_0\) over \(k\oplus M^\vee\) with the
base ring acting through \(k\). For \(\lambda\in M\), let \(x_\lambda\)
be its pullback by \((r,\varphi)\mapsto r+\epsilon\varphi(\lambda)\);
let \(x_{\rm c}\) be the constant deformation over \(D\). Assume that
\(x_\lambda=x_{\rm c}\) as maps to \(S\) implies \(\lambda=0\).
For the universal first-order deformation take \(M\) to be the
dual of the relative cotangent space. On a local affine chart its
ring map sends \(r\) to \((r(x_0),dr)\). A direction annihilating
every \(dr\) is zero since these differentials span. We shall
prove \(\dim_k M\le n-1\).
Write \(A, A_\lambda,A_{\rm c}\) for the corresponding abelian schemes.\formalnote{16.40}{40}

In an étale chart \(z\) at zero, write \(m(s)\) for the tangent
matrix of \([s]\), \(s\in\mathcal O_E\). The integral signature
equation gives a column \(Y\), a row \(\xi\) and an index \(i\) with
\[
 m(a)=Y\xi,\qquad m(s)Y=\rho(s)Y,\qquad Y_i\ \text{a unit}.
\]
Here the scalar term of \(m(a)\) is \(\rho(h)=0\). For a point \(v\)
reducing to zero across a square-zero ideal \(I\), the chart
identifies its coordinates with a vector \(z(v)\) in \(I^g\),
where \(g=\dim A_0\). All quadratic terms vanish. Thus addition
becomes vector addition and \(z([s]v)=m(s)z(v)\). Define
\[
 \omega(v)=z_i([a]v)=Y_i\,\xi z(v).
\]
This scalar is additive, belongs to \(I\), commutes with base
change and satisfies \(\omega([s]v)=\rho(s)\omega(v)\).
If \(\omega(v)=0\), the unit coordinate forces
\(z([a]v)=0\), hence \([a]v=0\).
In particular \([a c']v=0\) for \(c'\in\mathfrak w\), and
\([p]v=0\) by linearization of multiplication by \(p\).\formalnote{16.41}{41}

\subsection{A deformation measured on torsion}

The Serre--Tate theorem identifies deformations of an abelian scheme across a nilpotent thickening, over a ring in which \(p\) is nilpotent, with deformations of its \(p\)-divisible group. In Drinfeld's proof, a multiple of a homomorphism on the closed fibre lifts uniquely. Its vanishing on the corresponding torsion subgroup determines whether the original homomorphism lifts \cite[Lemma~1.1.3 and Theorem~1.2.1]{lit-katz-serre-tate}. Below we carry out this lifting and division over a square-zero extension using the finite \(\mathfrak w\)-torsion. The signature then gives the bound on the dimension of the tangent space.

Let \(G=A_0[\mathfrak w]\) be the group scheme cut out by all
endomorphisms in \(\mathfrak w\). For a local \(k\)-algebra \(C\) put
\(C^+=C\oplus\operatorname{Hom}_k(M,C)\) with square-zero second
summand. Smoothness of the abelian scheme over the level base gives
a lift \(Z(v)\in A(C^+)\) of each \(v\in A_0(C)\), by lifting
coordinates in an étale chart. For \(y\in G(C)\), define
\[
 f_C(y)(\lambda)=\operatorname{pr}_2\bigl(\omega([c]Z(y))\bigr)(\lambda).
\]
The difference between any two lifts is infinitesimal, and the
logarithm of its \(c\)-multiple is zero because \(\rho(c)=0\).
Consequently \(f_C\) is independent of the lifts, natural in \(C\),
\(F\)-linear in \(y\), and \(k\)-linear in \(\lambda\). If
\(y\in G(D)\) reduces to zero over \(k\), modify its lift by
subtracting its value at \(\epsilon=0\) pulled back along \(k\to D\).
The result still lifts \(y\) and is infinitesimal for \(D^+\to k^+\).
Its \(c\)-multiple again has zero logarithm, so \(f_D(y)=0\).\formalnote{16.42}{42}

We claim that if \(f_C(y)(\lambda)=0\) for all finite local
\(k\)-algebras \(C\) and \(y\in G(C)\), then \(\lambda=0\).
First lift \([ac]\) on the common special fibre to an additive map
\(P:A_{\rm c}\to A_\lambda\). For a section \(v\) of the source,
choose locally a section \(v'\) of the target with the same reduction
modulo \(\epsilon\), and put \(P(v)=[ac]v'\).
The existence of local lifts follows from smoothness. Since
\([ac]\) kills every infinitesimal difference, the formula is
independent of choices, compatible with addition, and glues.\formalnote{16.43}{43}

We now divide \(P\) by \([c]\). A point \(v\) of
\(A_{\rm c}[c]\) with values in a finite local \(D\)-algebra \(C\)
is a point of \(A_0(C)\) by constancy. The point \(tv\) lies in
\(G(C)\) since \(t\mathfrak w\subset(c)\). Taking \(tZ(v)\) as its
lift gives
\[
 \operatorname{pr}_2\omega(cZ(v))(\lambda)
 =\operatorname{pr}_2\omega(ctZ(v))(\lambda)=f_C(tv)(\lambda)=0
\]
since \(\rho(t)=1\). The first component of \(\omega(cZ(v))\)
vanishes as \(cv=0\). Under
\(C^+\to C,\ (r,\alpha)\mapsto r+\epsilon\alpha(\lambda)\),
the point \(Z(v)\) therefore gives a lift \(v'\) with
\(\omega(cv')=0\), whence \(P(v)=acv'=0\).
This detects the entire finite kernel scheme, since its coordinate
algebra is a product of finite local \(D\)-algebras. The faithfully
flat isogeny \([c]\) is the quotient by that kernel, so
\(P=j\circ[c]\) for an additive \(j:A_{\rm c}\to A_\lambda\).
On the special fibre we can cancel \([c]\), giving \(j_0=[a]\).\formalnote{16.44}{44}

\subsection{Recovering the marked theta object}

The following division step uses the polarization. Write
\(B_N(y,z)=[N]_{y+z}-[N]_y-[N]_z+[N]_0\) for its bilinear Picard
pairing, and put subscripts \({\rm c},\lambda\) on the two pairings.
Any additive \(j:A_{\rm c}\to A_\lambda\) with \(j_0=[a]\) is
\(\mathcal O_E\)-linear and satisfies
\[
 B_\lambda(jy,jz)=B_{\rm c}(ay,az).
\]
Indeed the defect of linearity has infinitesimal image, so its
\(p\)-multiple is zero. Precomposition by the surjection \([p]\)
then kills the defect only if the defect is zero. For the pairing,
absolute Frobenius on a test scheme multiplies a Picard class by
\(p\) and factors through the reduction modulo \(\epsilon\). It gives
equality of the \(p\)-multiples of the two sides. Bilinearity gives
the claimed formula on first arguments of the form \(py'\). Such a
lift \(y'\) exists on the faithfully flat cover defined by \([p]\).
For descent, take the difference of the universal pairing bundles
on \(A_{\rm c}\times A_{\rm c}\). It has a frame after pullback by
\([p]\times1\). Divide this frame by its value at second argument
zero. On overlaps two frames differ by a unit that is constant
along the second abelian variety, because that variety comes by
base change from \(A_0/k\). The normalized ratios agree, hence
descend. The difference bundle is therefore pulled back from its
restriction at second argument zero, where both pairings vanish.\formalnote{16.45}{45}

The map \(j\) is an isogeny. Its reduction is the isogeny \([a]\),
with \(a\ne0\) by \((a,h)=1\) and \(\rho(h)=0\); the nilpotent
thickening preserves surjectivity and quasi-finiteness, which suffice
for an isogeny between these abelian schemes. Its kernel \(K\) is
killed by \(a^2\). Indeed \([a]\) acts as zero on \(K\bmod\epsilon\).
On the augmentation ideal \(J\) of its coordinate algebra \(B_K\),
the pullback \([a]^*\) therefore maps \(J\) into
\(\epsilon B_K\cap J=\epsilon J\), using
\(B_K=D\oplus J\). A second application gives zero.\formalnote{16.46}{46}

For \(v\in K\), the point \(av\) is infinitesimal, so \(pav=0\).
For every \(z\) after base change the adjoint relation gives
\[
 B_{\rm c}(hav,z)=B_{\rm c}(av,az)=B_\lambda(jv,jz)=0.
\]
To see that \(hav=0\), work over each finite local algebra of \(K\).
Pairing trivially with the universal point lifts translation by
\(hav\) to the theta line bundle. Let \(T\) be its polarization
torsion index group. For a system of theta coordinates spanning
the global sections, the \(|T|\)-th power of that translation lift
preserves the projective coordinates, so \(|T|hav=0\). Since
\(|T|\) is prime to \(p\), this and \(phav=0\) imply \(hav=0\).
Now \(a^2v=hav=0\) and \((h,a)=1\) imply \(av=0\). Thus
\([a]\) factors through \(j\), yielding \(j':A_\lambda\to A_{\rm c}\)
which reduces to the identity. The same kernel calculation with
\(a=1\) makes \(j'\) an isogeny with trivial kernel. Its inverse
\(u:A_{\rm c}\to A_\lambda\) reduces to the identity and preserves
the pairings by the formula above.\formalnote{16.47}{47}

For a level label \(\ell\), the difference
\(u\eta_{\rm c}(\ell)-\eta_\lambda(\ell)\) is killed by the level
and is infinitesimal. The latter condition kills it by both \(p\)
and \(a\mathfrak w\), hence by \(\mathfrak w(p,a)\). The unit-ideal
condition on the level forces the two sections to be equal. The
theta class law \([N]_y=B_N(\beta y,y)+[N]_0\), together with
\(\operatorname{Pic}D=0\), now identifies the line-bundle classes
under \(u\). The matching of the projective coordinates requires
one more argument. A lift to sections of a map preserving the class
intertwines normalized theta translations up to scalars, since the
only global functions on the constant source come from the base.
Writing \(U_s\) for a theta translation and \(I\) for inversion,
the theta laws give \(U_{2s} I U_s=U_s I\) on the coordinate
sections. They span, so the equality holds on all sections.
This relation cancels the scalar ambiguity for each coordinate
indexed by \(2s\). The order of \(T\) is odd, so these are all the
coordinates. Thus \(u\) preserves projective evaluation at every
level point. These evaluations define the embedding of \(S\), and
the maps to the coefficient ring already agree. We conclude that
\(x_\lambda=x_{\rm c}\), so \(\lambda=0\), proving the claim.\formalnote{16.48}{48}

Here spanning is needed on each pullback of \(A_0\) to a
noetherian local \(k\)-algebra. It follows from the valuative torus
model. Choose an integer \(m\) invertible in \(k\) so large that
degree-\(m^2\) monomials in the projective coordinates span the
sections of \(N^{m^2}\), by Serre vanishing applied to the ideal sheaf of the embedding \cite[Tag 0B5U]{arith-stacks}. The class laws identify
\([m]^*N\) with \(N^{m^2}\) up to a line from the local base. On
sections of \([m]^*N\) take the sum of translates over \(A[m]\).
At a complex argument \(z/m\), adding a period to \(z\) permutes
the summands. The result thus descends to a theta function of \(N\).
Over the valuation ring a multiple of this function is a combination
of the coordinates that is nonzero somewhere on each polarization
torsion orbit. A nonunit multiplier clearing a denominator would
annihilate its reduction on an entire orbit, a contradiction. Thus
the trace lies in the pullback of the integral span. The spanning
monomials in degree \(m^2\) carry this containment to the local
base change. On a pullback \([m]^*\sigma\) the trace is
\(m^{2g}[m]^*\sigma\). The scalar is invertible and pullback by the
isogeny is injective, so \(\sigma\) lies in the span of the
coordinates of \(N\). Constant global functions follow from
base change from the abelian variety over \(k\).\formalnote{16.49}{49}

\subsection{Counting the possible first-order measurements}

Let \(d=\dim_F G(k)\). The \(q^n\) generic \(\mathfrak w\)-torsion
labels specialize by an additive map \(\nu\) to \(G(k)\).
Section~\ref{sec:pel-covers} gives, for any point \(y\) of \(G\)
over a local algebra with residue \(k\), a parameter satisfying
\(\tau^{|\nu^{-1}(\bar y)|}=0\). Over the field itself an empty
fibre would therefore give \(1=0\). Thus \(\nu\) is onto and
\(|\ker\nu|=q^{n-d}\).

Take the coordinate algebra \(B\) of the local component of \(G\)
at zero, with universal point \(v\). It is finite, since \(G\) is
contained in \(A_0[c]\). For \(v\), the torsion parameter also
satisfies \(v\bmod\tau=0\). By universality the quotient
\(B\to B/(\tau)\) then factors through the augmentation
\(B\to k\), so \(\tau\) generates the maximal ideal. Repeated
expansion in powers of \(\tau\) gives
\[
 B=k[\tau]/(\tau^e),\qquad e\le q^{n-d}.
\]
Here \(e\) is its least nilpotence exponent. The signature line
supplies a nonzero tangent vector \(L(1)\) killed by \(\mathfrak w\),
with multiples \(L(\alpha)\), \(\alpha\in k\), on which
\(s\in\mathcal O_E\) acts by \(\rho(s)\).
The map \(B\to D\) representing \(L(1)\) sends \(\tau\) to a
nonzero multiple of \(\epsilon\). Hence \(e\ge2\), so \(d<n\).
The endomorphism \(\tau\mapsto\alpha\tau\) of \(B\) gives a
\(B\)-valued point with tangent \(L(\alpha)\).\formalnote{16.50}{50}

Fix \(\lambda\in M\) and put \(F_\lambda=f_B(v)(\lambda)\in B\).
If it is nonzero, let \(b_i\tau^i\), \(b_i\ne0\), be its lowest
term. Evaluation at zero gives \(i>0\) and evaluation at \(L(1)\)
excludes \(i=1\). If \(y\in G(B)\) reduces to the origin and has
tangent \(L(\alpha)\), its endomorphism of the universal algebra
sends \(\tau\) to \(\alpha\tau+O(\tau^2)\). Naturality therefore
gives the \(\tau^i\)-coefficient of \(f_B(y)(\lambda)\) as
\(b_i\alpha^i\). Apply this to sums and \(F\)-multiples of lifts
of \(L(\alpha)\). Their tangents add and scale as their arguments
do. Since \(f_B\) is \(F\)-linear, \(\alpha\mapsto\alpha^i\) is
\(F\)-linear on \(k\).

This forces \(i=q^j\). Indeed \(\alpha\mapsto\alpha^{qi}\) is
also \(F\)-linear. Its fixed space consists of the \(qi\) roots of
\(X^{qi}-X\), all simple since the derivative is \(-1\). Algebraic
closure supplies every root, so \(qi\) is the cardinality of a
finite \(F\)-vector space. As \(i>1\) and \(i<e\le q^{n-d}\),
we have \(1\le j<n-d\). This argument constrains the
\emph{lowest} nonzero power and suffices for injectivity.\formalnote{16.51}{51}

Record the \(F\)-linear function \(y\mapsto f_k(y)(\lambda)\) on
\(G(k)\) and the \(\tau^{q^j}\)-coefficient of \(F_\lambda\) for
\(1\le j<n-d\), taking that coefficient to be zero if \(q^j\ge e\).
If all recorded data vanish, the lowest-term argument gives
\(F_\lambda=0\). Over any finite local \(k\)-algebra an arbitrary
point of \(G\) is the sum of its constant residue point and a
point obtained by base change from \(v\). Naturality makes
\(f_C(y)(\lambda)=0\) for both summands, and the separation
claim gives \(\lambda=0\). We have proved a \(k\)-linear injection
\[
 M\hookrightarrow
 \operatorname{Hom}_F(G(k),k)\oplus k^{\,n-d-1}.
\]
The first summand has dimension \(d\) over \(k\); the two
contributions add to \(n-1\).\formalnote{16.52}{52}

\subsection{From the tangent bound to smoothness}

An upper bound alone allows schemes of smaller dimension. For
this theta closure the other input is the holomorphic immersion
of the complex ball. Choose a valuation arc whose complex endpoint
is generic in its analytic family, so its image as a scheme point
specializes to all nearby parameter values. In an affine
neighbourhood \(\operatorname{Spec} B'\) of the closed point, the
bound on the relative cotangent space gives \(n-1\) functions
whose differentials span it. They define an unramified map to
\(\mathbb A_R^{n-1}\) near that point, with a local presentation
as a quotient of a standard étale algebra. The composite of the
period family and these functions has invertible derivative at
the complex endpoint. Any polynomial over \(R\) vanishing on the
arc vanishes on the nearby complex family by specialization, hence
on a complex open set by the inverse function theorem. The
embedding \(R\hookrightarrow\mathbb C\) makes this polynomial zero.
Thus the polynomial algebra \(H=R[X_1,\ldots,X_{n-1}]\) injects
into the valuation ring. To remove the quotient in the presentation,
let \(y\) be the value on the arc of its standard étale generator,
which satisfies a monic equation \(P(y)=0\) with \(P'(y)\) a unit.
Normality of \(H\) gives \(P=\mu Q\) over \(H\), where \(\mu\) is
the minimal polynomial of \(y\) over \(\operatorname{Frac}H\). Since
\(P'(y)=\mu'(y)Q(y)\), \(Q(y)\) is a unit on the arc. Any element
of the quotient kernel is represented, up to an invertible
denominator, by a polynomial vanishing at \(y\), hence divisible
by \(\mu\). Multiplication by \(Q\) kills it in the standard étale
algebra. Inverting \(Q\) therefore identifies the quotient with
that étale algebra near the arc. Such neighbourhoods, using the finitely many
generic analytic parameters and the good-prime arcs, cover the
closure. This proves smoothness of \(S/R\) of relative dimension
\(n-1\).\formalnote{16.53}{53}

The hypotheses used here distinguish three different assertions.
The group family is already smooth over \(S\), allowing local lifts
of its points. The first-order rigidity argument recovers an actual
map to the particular embedded level base from its theta and level
data. Smoothness of that base over \(R\) then uses the analytic
immersion proved in \S\ref{sec:ball-differential}, normality, and the covering by good-prime arcs. The
argument concerns this closure over the stipulated localization.
Representability of a moduli functor of arbitrary PEL deformations
or smoothness at excluded primes would require additional theorems.


\section{Preservation of theta chart ranks under specialization}
\label{sec:theta-chart-gluing}

The signed-frame calculation of \S\ref{sec:signed-frames} controls how many analytic families occur. Fix a specialization of one of these families from a characteristic-zero torus to a residue field via a valuation ring. Section~\ref{sec:good-arcs} constructs such specializations through the required points of the level base. Write \(O\subset\mathbb C\) for its valuation ring, \(v\) for its multiplicative absolute value, \(\mathfrak m\) for the maximal ideal, and \(k=O/\mathfrak m\). The value group can be nondiscrete. We shall identify the kernel of reduction of a space of theta sections, prove the finite-degree test used in constructing the closure, and write the overlap equations through which Riemann's relation gives a morphism on the special fibre.

\subsection{A theta orbit over the valuation ring}
Let \(V/L\) be the complex torus with period lattice \(L\), \(B\) the positive Appell--Humbert form and \(\chi\) its multiplier. Write \(H=H^0(B,\chi)\) for the complex space of theta functions and \(D_B=L_B^\vee/L\) for its finite group of translation labels, where \(L_B^\vee\) is dual for the alternating form of \(B\). Put \(d=\dim_{\mathbb C}H\). The positive Appell--Humbert data \((B,\chi)\) define an ample line bundle on the complex abelian variety \(V/L\). Its theta group is a central extension of \(D_B\) by \(\mathbb C^\times\). It acts irreducibly on \(H\) with the tautological central character, and \(d^2=|D_B|\) \cite[\S1, Proposition~3 and Theorem~2]{legacy-mumford}. In the model associated with a maximal isotropic subgroup, translations and characters act on the functions on a group of order \(d\).\formalnote{16.54}{54}
\phantomsection\label{route:33}\leanroute{The formal proof shows that the translation operators span \(\operatorname{End}_{\mathbb C}(H)\) and computes \(d^2\) as the index of the period lattice in its dual. Mumford's results supply the dimension and scalar commutant over \(\mathbb C\). The identity (TG1) and the argument over the valuation ring are proved below.}

Choose representatives \(a\) of \(D_B\), write \(U_a\) for the normalized translation of \S\ref{sec:theta-family}, and take \(U_{-a}=U_a^{-1}\) for the opposite lift. For any \(f,g\in H\),
\[
   d\sum_{a\in D_B}(U_{-a}g)(z)\,(U_af)(w)
      =|D_B|\,f(z)g(w).
 \tag{TG1}
\]
Indeed the operator \(g\mapsto\sum_a (U_{-a}g)(z)U_af\) commutes with each translation by reindexing the terms. The two scalar factors from changing a representative cancel. Its trace is \(|D_B|f(z)\), since each rank-one summand has trace \((U_{-a}U_af)(z)\). Schur's lemma gives (TG1).\formalnote{16.55}{55}

At a good prime some \(v\)-unit integer kills \(D_B\), so \(|D_B|\in O^\times\). The torsion translation phases are units, and the norm \(\|f\|=\sup_{z\ {\rm torsion}}v((U_z f)(0))\) is invariant under the \(U_a\). For any nonzero \(f,g\), their \(O\)-spans of translates are homothetic. Here is the integral step. Fix \(z\) with \(f(z)\ne0\). In each of the finite lists \((U_{-a}f)(z)\), \((U_{-a}g)(z)\), choose an entry of largest absolute value and denote the entries by \(q_f,q_g\). The second list has a nonzero entry by (TG1), since \(g\ne0\). Apply (TG1) with its \(f\) replaced by a translate \(f'\) satisfying \(f'(z)=q_f\), and divide by \(|D_B|q_g\). The coefficients
\[
          \frac{d}{|D_B|}\,\frac{(U_{-a}g)(z)}{q_g}
\]
are all in \(O\). Thus \((q_f/q_g)g\) belongs to the orbit span of \(f\). The span is stable under the translations, whose composition changes by a unit. Interchanging \(f,g\) gives the reverse inclusion. The ultrametric inequality in both directions now gives \(\|f\|=v(q_f/q_g)\|g\|\). It follows that the orbit span of \(f\) is its entire closed norm ball. A nonzero \(g\) in the ball is a scalar of absolute value \(\|g\|/\|f\|\) times an element with norm \(\|f\|\). Thus the strictly smaller ball equals \(\mathfrak m\) times the closed ball. This proof only takes maxima of finite lists.\formalnote{16.56}{56}

For the PEL system let \(r\ge3\) be the padding integer, \(v(r)=1\), and let \(M=\|\theta_c\|>0\), independent of \(c\). Apply the orbit argument in degrees \(m=r^j\). The normal-generation step of \S\ref{sec:theta-family} uses the lattice over the chosen Lagrangian, which contains \(r^{-1}L\), and the multiplicativity of the theta norm. It identifies the image of the homogeneous forms of degree \(m\) over \(O\) with the closed ball \(\mathcal H_m\) of radius \(M^m\) in \(H^0(mB,\chi^m)\). In particular \(\mathcal H_m\) is finitely generated, torsion-free, hence free over the valuation ring, and spans that full complex vector space. We have proved\formalnote{16.57}{57}
\[
     \{f\in\mathcal H_m:\|f\|<M^m\}=\mathfrak m\mathcal H_m,\qquad
     \dim_k(\mathcal H_m/\mathfrak m\mathcal H_m)
        =\sqrt{|L_{mB}^\vee/L|}=:\sigma_j .
 \tag{TG2}
\]

\subsection{Detecting a nonzero residue at finite level}
Choose points \(z_t\) representing \(L_{mB}^\vee/L\), with \(z_0\in L\). Normalize the evaluations of \(\mathcal H_m\) at \(z_t\) by a scalar \(s_t\), so that the functional \(e_t(f)=f(z_t)/s_t\) takes integral values and some value is a unit. Such a scalar exists by taking the largest of the values of a finite generating family. The coordinate construction gives the same primitive scaling. If \(s_t^{(1)}\) normalizes the degree-one coordinates, then \((s_t^{(1)})^m\) works in degree \(m\), since a power of a unit coordinate is a unit.

The reductions of the \(e_t\) jointly detect \(\mathcal H_m/\mathfrak m\mathcal H_m\). Indeed, suppose each \(e_t(f)\) is small. For a lift \(a\) of a label choose \(t\) with \(z_0+a=z_t\bmod L\). Quasiperiodicity and the common factor in the translation formula make \(f\mapsto e_0(U_af)\) proportional to \(e_t(f)\). The proportionality is a unit because both functionals are integral on the same ball and each has a unit value, the first because \(U_a\) preserves that ball bijectively. Hence every \(e_0(U_af)\) is small. If \(\|f\|=M^m\), the translates of \(f\) span the ball by (TG1), contradicting primitivity of \(e_0\). Conversely (TG2) shows that every element of smaller norm has small evaluation at all \(t\). It follows that the common kernel of these reduced evaluations is exactly the strict ball.\formalnote{16.58}{58}

Fix a projective chart \(X_c\ne0\). Let \(I_c\) be the vanishing ideal of the specialized theta vectors in that chart, normalized to \(X_c=1\). A homogeneous polynomial \(P\) of degree \(m=r^j\) with integral coefficients reduces into \(I_c\) precisely when \(\|P(\theta)\|<M^m\). In fact multiply \(P\) by \(X_c^{rm-m}\). The new polynomial vanishes at all reduced points with \(X_c=0\), and on the others it vanishes precisely when \(P\) does. The evaluation test just proved can then be applied in degree \(rm\) using all points. Multiplicativity gives the new norm as \(M^{rm-m}\|P(\theta)\|\). This proves the criterion. It also gives the rank \(\sigma_j\) of the degree-\(m\) image modulo \(I_c\), since homogeneous \(k\)-forms lift coefficientwise. Through the specified degree cutoff, the labelled level points already include representatives of \(L_{mB}^\vee/L\), so their homogeneous evaluations alone test \(P\) in these degrees, including labels whose \(c\)-coordinate vanishes.\formalnote{16.59}{59}

Here is why the cutoff is finite before any particular fibre or valuation is chosen. In the fixed polynomial variables let \(N\) be a monomial ideal, and let \(b_j(N)\) count the monomials of degree \(r^j\) outside \(N\). There exists \(i_0\) such that, for every monomial ideal \(N\),
\[
   b_j(N)\le\sigma_j\quad(0\le j\le i_0)
       \quad\Longrightarrow\quad
   b_j(N)\le\sigma_j\quad(j\ge0).
 \tag{TG3}
\]
Prove the same assertion for all ideals containing a fixed monomial ideal \(N_0\), by Noetherian induction on \(N_0\). If all its counts satisfy the bounds, every larger ideal does too. Otherwise choose \(j\) with \(b_j(N_0)>\sigma_j\). Each ideal above \(N_0\) satisfying the bound at \(j\) contains at least one of the finitely many monomials \(u\) of degree \(r^j\) outside \(N_0\). Apply the induction hypothesis to the strictly larger ideals \(N_0+(u)\). Take the maximum of their cutoffs and \(j\), then set \(N_0=0\). This proves (TG3). Applied to the leading monomial ideal of the local relations, it is exactly what the cofactor argument of \S\ref{sec:theta-family} needs to pass from the finitely many evaluations to flatness and equality of the specialized chart ideal. The torsion level used for these evaluations may depend on the cutoff \(i_0\) in \textup{(TG3)}.\formalnote{16.60}{60}

\subsection{Regular transition coefficients for the addition law}
Write \(T=L_B^\vee/L\) in degree one and \(T[r]=\ker(r:T\to T)\) for the PEL data. For every label \(s\) the maximal-coordinate argument in \S\ref{sec:theta-family} gives an index \(o\) such that \(o,-o,s-o\) all attain the maximum of the absolute values of the theta constants. Translations by level labels change the corresponding evaluations by a common scalar and units. Thus this condition puts the centre of the arc in the designated charts for all the translated level sections. These are open conditions, so they also hold at its generalizations.\formalnote{16.61}{61}

More explicitly, let two admissible padded Riemann tuples have coefficient pairs \((a_l,b_l),(a'_l,b'_l)\), and use the doubled products of \S\ref{sec:theta-family}. Two slots of the second tuple are \((1,0)\). Prescribe the other labels and let \(s\) be the difference of the two weighted sums of labels \(\sum_l\bar a_l\mu_l\) still to be made up by those slots. Choose the labels \(o,s-o\) there. If \(\alpha,\beta\) denote the original labels of the corresponding two zero-argument factors, the regular coefficient in the formula expressing the first doubled product in terms of the second is, on the chart just obtained,
\[
     \frac{\sqrt{[T:T[r]]}}{|T|}
        \frac{\theta_\alpha(0)\theta_\beta(0)}
             {\theta_o(0)\theta_{s-o}(0)} .
 \tag{TG4}
\]
The first factor is defined on the base because \(|T|\) is a unit at every good prime. The second is a transition function for the tensor square of the zero-label coordinate line bundle, with denominator a unit after local trivialization. Thus (TG4) defines a regular function even at special points where the numerator vanishes. Substituting it into Riemann's relation gives regular linear combinations in each direction by interchanging the two padded tuples.\formalnote{16.62}{62}

For completeness, suppose \(m_i\) are the coordinates of the known product line bundle and \(w_j\) the new coordinates given on the separating complex theta points. On any chart \(m_k\ne0\), the forward combinations define regular functions \(h_{jk}=w_j/m_k\) there. On an overlap they satisfy \(h_{jk}=(m_{k'}/m_k)h_{jk'}\) on the separating points, hence as regular functions. They therefore glue to sections of the known line bundle. The reverse expression for \(m_k\), divided by \(m_k\), writes \(1\) as a regular combination of the \(h_{jk}\), so the new sections generate that bundle in every characteristic. They define a morphism to projective space. The closed equations of the Segre-embedded product vanish on its separating points, so this morphism factors through the product. Projection gives the separate operations. The tuples and verification of the group and endomorphism laws are in \S\ref{sec:theta-family}.\formalnote{16.63}{63}

\subsection{The unit scalar in the good-reduction norm}
We give the scalar calculation used to prove multiplicativity in \S\ref{sec:theta-family}. Let \(c,d=\bar c\) be distinct nonzero coprime multiplications preserving \(L\), both acting trivially on \(L_B^\vee/L\), and use the canonical multiplier from \(\beta+\bar\beta=1\). Here a bar on a multiplication denotes its adjoint for \(B\). For the alternating form \(E_B=\operatorname{Im}B/\pi\) put \(S_y=U_y/\exp(\pi i E_B(\beta y,y))\). Write \(K_a=a^{-1}L/L\). On canonical theta functions the operator
\[
      A_{c,d}f=\sum_{z\in K_d}S_{dz}(f\circ cd^{-1})
\]
is a scalar \(\mu\). Indeed averaging makes it invariant under the normalized \(L\)-translations, by reindexing \(K_d\). For \(a\in L_B^\vee\), the translations \(S_{ca},S_{da}\) on theta functions each agree with \(S_a\), by their prescribed lattice laws. The first passes through pullback as the second, which commutes with all \(S_{dz}\), since \(E_B(da,dz)=E_B(a,cdz)\in\mathbb Z\). Thus \(A_{c,d}\) commutes with the dual translations and Schur's lemma applies.\formalnote{16.64}{64}

Let \(\delta=\det_{\mathbb C}(d)\) and let \(P(x,y)\) be the Poincaré kernel, given by
\[
 P(x,y)=\sum_{l\in L}\chi(l)^{-1}
       \exp\{B(x,y-x+l)-B(l,y)-B(l,l)/2\}.
\]
Its diagonal has mean one, since every term with \(l\ne0\) is a constant times a nontrivial character of \(V/L\). Write \(m\) for the mean of \(P(dx,cx)\) over \(V/L\). Substituting \(f=P(dx,\cdot)\) and evaluating at \(dx\), the left side becomes \(\sum_{z\in K_d}P(d(x+z),c(x+z))\). Integration therefore gives \(\mu=|K_d|m=|\delta|^2m\). The reversed operator \(A_{d,c}\) has scalar \(|K_c|\bar m\), since \(P(cx,dx)=\overline{P(dx,cx)}\) here. In the composed double sum the inner character sum indexed by \(z\in K_d\) vanishes unless \(cz\in L_B^\vee\). Indeed translation by \(\ell\in L\) fixes the kernel average and multiplies its terms by \(\exp(2\pi i E_B(\ell,cz))\). If \(cz\) is dual, \(dcz-cz\in L\). Also \(dcz\in L\), whence \(cz\in L\). Coprimality of \(c,d\) gives \(z\in L\). The surviving term is \(|K_c|f\), so \(\bar m\mu=1\). Thus \(m\bar\delta\) has complex absolute value one.\formalnote{16.65}{65}

To determine its phase let \(\kappa=c-d,\ h=\bar\beta\kappa-c\), and put \(q(z)=\exp B(z,\kappa h z)\) on \(K_\kappa\). To see that \(q\) is well defined, let \(x,y\in\kappa^{-1}L\) represent the same class \(z\), and put \(l=x-y\). Adjointness and \(\beta+\bar\beta=1\) give \(B(x,\kappa h x)-B(y,\kappa h y)=B(l,l')-B(l',l)\), where \(l'=h(\kappa x-\beta\kappa l)\in L\). Integrality of \(E_B\) puts this difference in \(2\pi i\mathbb Z\). Expanding the kernel and indexing its lattice sum by \(\kappa(z+l)\) gives
\[
      P(dx,cx)=\sum_{l\in L}\sum_{z\in K_\kappa} q(z)
       \exp B(l+x+z,\kappa c(l+x+z)).
\]
The change of indices is bijective. In computing the mean, integration of the translated Gaussian over all \(x+l\) gives its integral on \(V\) divided by the covolume. In complex basis coordinates take Lebesgue measure divided by \(\pi^{\dim V}\), and let \(M_\kappa\) represent \(B(u,\kappa v)\). The Gaussian integral is \(1/\overline{\det(M_\kappa)\delta}\). Indeed the matrix in its negative exponent is the adjoint of the matrix for \(B(u,\kappa d v)\), with positive hermitian part since \(\kappa d+(\kappa d)^*=\bar\kappa\kappa\). The form is positive definite on the image of the invertible \(\kappa\). Thus \(m\bar\delta\) is \(\sum_z q(z)\) divided by the real covolume times \(\overline{\det M_\kappa}\). Since \(M_\kappa\) is skew-Hermitian its determinant has real fourth power.\formalnote{16.66}{66}

Put \(N=|K_\kappa|\). Homogeneity and \(Nz=0\) give \(q(z)^{N^2}=1\). Write \(N^2-1\) as a sum of four integer squares. The matrix \(Q\) of multiplication by the corresponding quaternion satisfies \(Q^tQ=QQ^t=(N^2-1)I\). On \(K_\kappa^4\), \(-Q^t\) is its inverse. Polarizing the quadratic exponents gives
\[
     \prod_{j=1}^4 q((Qz)_j)
        =\prod_{j=1}^4 q(z_j)^{N^2-1}
        =\overline{\prod_{j=1}^4 q(z_j)} .
\]
Summing over quadruples shows that \((\sum_z q(z))^4\) is real. Hence \(m\bar\delta\) has absolute value one and real fourth power, so \((m\bar\delta)^8=1\). Since \(\mu=\delta(m\bar\delta)\), we obtain \(\mu^8=\delta^8\), and \(v(\mu)=1\) at the good place where \(v(\delta)=1\). This is precisely the unit needed in the kernel-shift estimate of \S\ref{sec:theta-family}.\formalnote{16.67}{67}

The chart conclusions use the ample theta datum, the good-prime expanding multiplications for the torsion norm, the padded Lagrangian and level inclusions, and the valuation arcs. Sections~\ref{sec:good-arcs} and~\ref{sec:theta-family} supply those constructions and the descent from the kernel-shift estimate. The separate tangent estimate for smoothness of the base is in \S\ref{sec:theta-tangent}.


\section{Generic covers and torus charts}\label{sec:cover-charts}
\subsection{From the deck action to smooth generic covers}
Let \(W\to Z\) denote the identity leg of a correspondence cover of \S\ref{sec:pel-covers}, after base change to an algebraically closed field of characteristic zero containing the coefficient field. Its target is smooth, its source is projective and reduced, and the finite group \(D\) acts simply transitively on the geometric points of each fibre. Reducedness survives this base change because the original closure is reduced, passage to its generic fibre is localization and fields of characteristic zero are perfect. The characteristic-zero formulation follows from the calculation on complex points: any counterexample to surjectivity, transitivity or freeness can be specified over an algebraically closed subfield of countable transcendence degree containing the coefficients, then embedded compatibly in \(\mathbb C\). Equivalently, test surjectivity of the maps \(W\to Z\) and \(D\times W\to W\times_Z W\) on geometric points and emptiness of the nonidentity fixed loci.

The map is proper with finite geometric fibres, hence finite. The action is free on schemes, since for a constant finite group every nonidentity fixed scheme has empty underlying space. The theorem on quotients of quasi-projective schemes by free finite group actions gives a scheme \(W/D\) and a \(D\)-torsor \(W\to W/D\) \cite[Tag 07S7]{arith-stacks}. In particular it is finite étale, by faithfully flat descent from the trivial torsor. Over an affine open of \(Z\) with finite algebra \(B\) for \(W\), the quotient algebra \(B^D\) is reduced and finite over the base. The map \(W/D\to Z\) is bijective on geometric points. On each connected component of the smooth target there is a dominating irreducible component of the quotient. Its image is all of the target by finiteness, and geometric injectivity then places every point of the quotient on this one component. The quotient is thus integral over that component of the target. Its generic field extension has degree one, since it is separable and there is only one geometric generic point over the target. The finite birational morphism to the normal target is an isomorphism. This makes the original leg finite étale in characteristic zero, and proves the smoothness and relative dimension needed for its source.
\phantomsection\label{route:29}\leanroute{For the smoothness of the correspondence covers, Lean works through a local finite-algebra calculation for Galois covers. Here the projective source admits a quotient by the free finite action. The argument above identifies that quotient with the given normal base in characteristic zero.}\formalnote{16.68}{68}

\subsection{The torus charts in cohomology}
Let \(Y\) be the proper model in \S\ref{sec:seed-proof}. The relative smoothness and the dimension at the safe places were obtained in \S\ref{sec:theta-tangent}. For the correspondence covers use the smooth generic fibres given in \S\ref{sec:pel-covers}. Over the coefficient base where smoothness of relative dimension \(d\) holds, each point has an affine open neighbourhood with an étale map to affine \(d\)-space. This follows from the local form of a smooth morphism \cite[Tag 039P]{arith-stacks}. In the application the base is affine; the coordinates over a principal open of the base also give an étale map to affine \(d\)-space over the full base by the open immersion. Properness gives quasi-compactness, so finitely many of these neighbourhoods cover.\formalnote{16.69}{69}

On such a neighbourhood with coordinates \(x_1,\ldots,x_d\), cover by the \(2^d\) principal opens where all the \(x_j+\epsilon_j\) are units, \(\epsilon_j\in\{0,1\}\). Every point lies in one, since \(x_j\) and \(x_j+1\) generate the unit ideal. Translation by the chosen constants preserves étaleness. We now have affine opens étale over a torus. Their finite intersections are affine: in a scheme separated over an affine base, the intersection of two affine opens is the inverse image of the closed diagonal in their affine product. Each intersection inherits its torus map from either factor. Let \(V\) be the union of their representable point functors inside \(Y\). Given a test point, the inverse images of these opens cover its spectrum, so the point belongs to \(V\) locally at every maximal ideal. The distinction matters for a test with many points. For a finite union build \(\mathbb Z[V]\) by the exact sequences
\[
 0\longrightarrow\mathbb Z[U\cap U']\longrightarrow
 \mathbb Z[U]\oplus\mathbb Z[U']\longrightarrow\mathbb Z[U\cup U']
 \longrightarrow0
\]
of point functors, which are exact by the free basis at each test. Induction on the cover, including its intersections, propagates the class of torus-chart representables through precisely the closure operations in \S\ref{sec:primitive-comparison}. This proves the charting hypothesis both for the proper smooth base at the specified prime and for the smooth correspondence schemes over the algebraic closure.\formalnote{16.70}{70}



\section{Good primes, valuation arcs and integral torus models}\label{sec:good-arcs}

Let \(S\) be the reduced closure of the labeled theta points of the
PEL datum over the ring \(R\) of \S\ref{sec:seed-proof}. To construct
the abelian family over \(S\), one first needs a valuation on
\(\mathbb C\) at which the equations of a complex fibre can be
reduced. To compute the tangent space of \(S\), one later needs
sections of the resulting abelian scheme over that valuation ring.
We prove these statements in the order in which they are used.

\subsection{Arranging good reduction at the prescribed places}

Write \(K=\mathbb Q(\sqrt{-d})\subset E\),
\(\beta=(1+\sqrt{-d})/2\) and \(\mathcal O=\mathcal O_E\).
Here \(d\equiv3\bmod4\) and \(E\) contains a primitive \(5040\)-th
root of unity, as in the odd-rank seed. Given a complex embedding
\(\sigma\) of \(E\), take it as the distinguished member of the CM
type. At every other infinite place choose the one of the two
embeddings whose restriction to \(K\) differs from \(\sigma|_K\).
Let \(\mathcal B\) be a finite set of places to avoid. We construct
a positive integer \(N\) outside every \(\mathfrak w\in\mathcal B\)
such that every nonzero level element \(\mathfrak N\) divisible by
\(N\) admits the polarization data used in \S\ref{sec:theta-family}.\formalnote{16.71}{71}

For each \(b=1,\beta\) we need a list of \(r-3\) elements \(a_i\in\mathcal O\)
spanning \(\mathcal O\) over \(\mathbb Z\) with
\[
               r=3+\bar b b+\sum_i\bar a_i a_i .
\]
Choose an integral basis \(e_1,\ldots,e_D\) and write
\(u=(b,e_1,\ldots,e_D)\). Solve
\(\sum_i n_i e_i=\sum_j(1-\bar u_j u_j)\) in integers \(n_i\).
The right side is fixed by conjugation. Take the list consisting
of \(u\), the \(e_i+n_i\), and \(1+\sum_i n_i^2\) zeros.
Its length is \(2+2D+\sum_i n_i^2\). Its sum of norms plus
\(\bar b b\) equals this length, since the cross terms
\(\sum_i n_i(e_i+\bar e_i)\) equal \(2\sum_j(1-\bar u_j u_j)\).
Append as many ones as necessary. This proves the padding identity
for every sufficiently large \(r\), with the same \(r\) admissible
for both \(b\).\formalnote{16.72}{72}

Choose an even factorial \(s\) large enough that all norms of places
in \(\mathcal B\) divide it and \(r=s+1\) satisfies both padding
bounds. Write \(F=E^+\) and \(\mathfrak D_E\) for the absolute
different. Choose \(\gamma\in\mathcal O\) with \(\gamma^2=-1\)
and \(\bar\gamma=-\gamma\). There is
\(\xi\in\mathfrak D_E^{-1}\), purely imaginary and with positive
imaginary part at every chosen CM embedding, and an
\(m\in\mathfrak D_E\) satisfying \(m\xi=sc-1\) for an integer
\(c\in\mathcal O\). Here are the choices. By approximation of ideals,
choose \(\mu\in\mathfrak D_F\) with \((s)+\mu\mathfrak D_F^{-1}
=\mathcal O_F\), so that \(sc_0+\mu a=1\) for some
\(a\in\mathfrak D_F^{-1}\). The relative different of \(E/F\)
is the unit ideal for these fields: an integral element with
\(b-\bar b\) a unit gives a quadratic order of unit discriminant,
as in \S\ref{sec:integral-witt}. Transitivity of the different now
allows \(m=\gamma\mu\) and
\(\xi=a\gamma+s(z-\bar z)\). The product is \(-\mu a\)
modulo \(s\), hence \(-1\) modulo \(s\). In the Minkowski
embedding the integers form a full lattice. Choose \(z\) within
its covering radius of a vector with each imaginary part
sufficiently positive to make
\(\operatorname{Im}\Phi_v(a\gamma)+2s\operatorname{Im}\Phi_v(z)>0\)
for all \(v\).\formalnote{16.73}{73}

Write \(\psi_f\) for the finite additive character. Composing it with the skew-Hermitian form gives the character pairing on the finite adelic \(E\)-frames
\[
       [x,y]=\psi_f\!\left(r^2\xi\,\bar y^{\,t}x\right).
\]
It is alternating since \(\xi\) is imaginary. Put
\(L_0=\widehat{\mathcal O}^{\,n}\),
\(T=L_0^\perp/L_0\) for its finite polarization group.
The group \(r^{-1}L_0\) is isotropic since \(\xi\) lies in the
codifferent. Extend it to a maximal isotropic subgroup by inclusion.
An element orthogonal to that subgroup could be adjoined while
preserving isotropy, so the maximal subgroup equals its
annihilator. Quotient by \(L_0\) gives the required Lagrangian
in \(T\). Testing against a single coordinate of \(L_0\) shows
\(r^2\xi x_i\in\widehat{\mathfrak D_E^{-1}}\) for \(x\in L_0^\perp\).
Multiplying by \(m\) shows that \(\delta=r^2(sc-1)\) kills \(T\).
At every prime above \(s\) we have \(\delta\equiv-1\). Consequently
\(\gcd(|T|,s)=1\), since the order of \(T\) divides a power of
\(|N_{E/\mathbb Q}(\delta)|\). In particular \(|T|\) is odd.
Take the finite cutoff \(j\) for the ranks of the theta equations
in \S\ref{sec:theta-family} and put
\[
                         N=r^{j+1}|T|.
\]
Its prime factors avoid \(\mathcal B\). For \(0\le i\le j\),
\(r^i t\in T\) implies \(r^i|T|\,t=0\), so any level divisible
by \(N\) contains all torsion labels required in these degrees.\formalnote{16.74}{74}

Let \(\mathfrak w\) be safe for this level, over the split place
\(u\) of \(K\), and \(p=N u\). The safety condition supplies \(b\in
\mathbb Z\) with \(h=\beta-b\in u\) coprime to \(\bar h\). The level
avoids \(\mathfrak w\) and every prime over \(\bar u\).
Thus \(r,|T|\) are units at \(\mathfrak w\).
Any prime containing both \(p\) and \(\bar h\) lies above \(\bar u\),
by coprimality at \(u\). It follows that
\((\mathfrak N)+\mathfrak w(p,\bar h)=\mathcal O\).
There is also a second prime of \(E\) above \(u\). In fact one
can avoid any two specified places and their conjugates over \(p\).
If these four exhausted the primes over \(p\), each cyclotomic
subfield \(\mathbb Q(\zeta_m)\subset E\), \(m=80,63\), would have
at most four such primes. For \(p\nmid m\) their number is
\(\varphi(m)/\operatorname{ord}_m p\). The exponents of the unit
groups modulo \(80\) and \(63\) divide \(4\) and \(6\), while
\(\varphi(80)=32,\ \varphi(63)=36\). Hence \(p\) would divide both
\(80\) and \(63\), an impossibility. Conjugation exchanges the two
places of \(K\), giving a further prime over \(u\) as required.
Both primes above \(u\) contain \(p,h\). These are precisely the
extra prime and ideal condition in the good-reduction calculation.\formalnote{16.75}{75}

For the seed one may enlarge the level at a new split
\(\mathfrak w_1\) of norm \(>(2n)^{[E:\mathbb Q]}\).
Indeed an automorphism \(y\) preserving a signed frame and its
reflection has all eigenvalues of absolute value one at every
embedding. If its framed matrix is congruent to \(1\) modulo
\(\mathfrak w_1\), then \(\operatorname{tr}(y)-n\in\mathfrak w_1\).
Every conjugate has absolute value \(\le2n\), so its absolute norm
forces it to be zero. At a definite embedding the unitary matrix
has trace \(n\), whence all its eigenvalues are \(1\) and \(y=1\).
This is the freeness called neatness here. Choose \(\mathfrak w_1\)
away from the prohibited finite set and \(N\). The stabilizer
of the chosen form, the principal congruence at its marked place,
and the pro-Iwahori subgroup at \(\mathfrak w_1\) are open.
A sufficiently deep principal level subordinate to their intersection
can have exactly the specified prime support, by raising prime
ideals to principal powers and imposing the required local
depths. The seed's choice of one side at the split pairs keeps
the requested good places safe.\formalnote{16.76}{76}

\subsection{A real-valued valuation through the level closure}

The analytic parameter set is infinite. At a fixed nonzero level,
finitely many ball families suffice up to level isomorphism. Indeed
the signed frames meet finitely many double cosets for the rational
group on the left and the level on the right. Choose a frame in
each occupied coset. Multiplication by a rational matrix transports
the reflection and carries every other such frame to the chosen one
up to level. Its distinguished line then gives a point of that
frame's ball. On each connected complex ball a member \(\eta_i\)
can be chosen whose scheme point specializes to all other members.
Locally, over the countable base \(R\), avoid the zero sets of the
countably many regular functions whose holomorphic restrictions
are not identically zero. These zero sets are nowhere dense,
so Baire's theorem gives the generic parameter in that chart.
The local specialization property propagates on the connected
ball by analytic continuation.\formalnote{16.77}{77}

There is an order issue here. In applying the finite-coset result
one needs level-isomorphic analytic parameters to determine the
same projective theta point. First do the arc and chart construction
below on the closure of the two parameters being compared. There
the condition that a finite set of complex points be dense holds
by definition. On that family the level-frame isomorphism identifies
their modeled objects and hence their projective evaluations.
Projecting into each ambient projective factor proves equality in
the original closure. With this step, the finitely many
\(\eta_i:\operatorname{Spec}\mathbb C\to S\) have dense union of
closures. No choice of a global family on \(S\) is needed to
start the finite-parameter construction.\formalnote{16.78}{78}

For a point \(s_1\in S\), choose a maximal ideal of \(R\)
containing its image prime. Properness of \(S/R\) makes the
image of \(\overline{\{s_1\}}\) closed, giving a specialization
\(s_1\leadsto s_0\) above that maximal ideal. Some \(\eta_i\)
specializes to \(s_0\) by finiteness and density. The chosen
\(s_1\) is a generalization of the centre and need not itself lie on
the arc. We realize \(\eta_i\leadsto s_0\) by a map from a valuation ring
\(O\subset\mathbb C\) with a \emph{real}-valued valuation. This
value-group condition allows the norm estimates of
\S\ref{sec:theta-family}. A valuation ring dominating the local
ring obtained directly from the specialization can have higher rank.

Write \(B=\mathcal O_{S,s_0}\) and let \(B\to\mathbb C\) be the
map given by \(\eta_i\). Dominate the local image of \(B\) by a
valuation subring \(W\subset\mathbb C\), with maximal ideal
contracting to \(\mathfrak m_B\). For generators \(a_1,\ldots,a_e\)
of \(\mathfrak m_B\) choose \(x\) among these with image dividing
all the others in \(W\), and adjoin the quotients there to the
image of \(B\). This gives a noetherian domain \(C\subset W\)
with \(\mathfrak m_B C=(x)\). At the good centre the residue
characteristic \(p\) has nonzero image in \(\mathbb C\), so a
nonzero \(x\) can be chosen. Since \(x\) is a nonunit, a prime
\(\mathfrak p\) minimal above \(x\) has height one and contracts
to \(\mathfrak m_B\). This is an affine chart of the blowup
calculation needed at the centre.\formalnote{16.79}{79}

Choose a transcendence basis \(\mathbf t\) of \(\mathbb C\) over
\(\operatorname{Frac}C\), set \(D=C[\mathbf t]\), and put
\(P=\mathfrak pD\). A valuation subring \(U\subset\mathbb C\)
dominates \(D_P\). In \(U\) choose a prime \(Q\) minimal over
\(PU\), and localize \(U\) at \(Q\). We verify that \(O=U_Q\)
has rank one. Any nonzero prime \(J\subset O\) contains a nonzero
element \(b\). Since \(\mathbb C/\operatorname{Frac}D\) is
algebraic, clear denominators in an equation for \(b\) with
nonzero constant coefficient. It follows that \(b\) divides
a nonzero element of \(D\) in \(O\). Thus \(J\cap D\) is a
nonzero prime contained in \(P\), and it equals \(P\): in
finitely many variables this follows from the principal ideal
theorem at \(\mathfrak pD\), and a nonzero polynomial in
\(J\cap D\) involves only finitely many variables. Hence
\(J\cap U\) contains \(PU\) and lies in \(Q\). Minimality of \(Q\)
gives \(J=QO\).

To see that the value group is archimedean, take \(z\in\mathbb C^\times\).
Its denominator ideal \(I_z=\{a\in O:az\in O\}\) is nonzero.
Since the maximal ideal is the only nonzero prime of \(O\), it is
contained in \(\sqrt{I_z}\). If \(y\) has value greater than one,
then \(y^{-1}\) is in the maximal ideal. Thus \(y^{-k}\in I_z\)
for some \(k\ge1\), or equivalently \(v(z)\le v(y)^k\).
Embedding the ordered group in the positive reals gives a
multiplicative valuation \(v:\mathbb C\to\mathbb R_{\ge0}\) with
\(O=\{z:v(z)\le1\}\). The map \(B\to O\) is local and its composite
with \(O\hookrightarrow\mathbb C\) is the prescribed map.
Thus the arc
\[
 \alpha:\operatorname{Spec}O\longrightarrow
      \operatorname{Spec}B\longrightarrow S
\]
has generic point \(\eta_i\), centre \(s_0\), and
\(\mathfrak m_O\cap\mathcal O\) is the good prime below \(s_0\).
Construction from \(\mathbb C\) allows a nondiscrete value group.
Its residue field \(k\) is algebraically closed. A monic
polynomial over \(k\) lifts to \(O\) and splits in \(\mathbb C\).
Integral closedness of \(O\) puts every root in \(O\).\formalnote{16.80}{80}

\subsection{The group scheme and the uniformization on sections}

The arc constructed so far has target \(S\). Its first use is
elementary. For each level label \(t\) choose a homogeneous
coordinate \(i(t)\) whose chart contains its closed value. The
chart then contains the entire local spectrum. The ratios
\(\theta_j(t)/\theta_{i(t)}(t)\) lie in \(O\) and their residues
are the coordinates at \(s_0\). For an arbitrary complex argument
\(z\), divide its vector of theta values by a nonzero entry of
maximal \(v\)-value. Every ratio is integral and one equals one.
The good-reduction equations applied to these vectors give the
required ranks and simultaneous nonvanishing charts at \(s_0\).
Openness gives those charts at \(s_1\) too. As proved in
\S\ref{sec:theta-family}, these are the inputs to construction of
a flat projective family, its group operations, and smoothness of
its abelian fibres. The resulting \(\mathcal A/S\) may now be
pulled back to \(A_O=\alpha^*\mathcal A\).\formalnote{16.81}{81}

We explain why its complex fibre has the \emph{specified} period
lattice. Let \(V,\Lambda\) be the prescribed space and lattice and
\(e:V\to A_{\mathbb C}(\mathbb C)\) the additive map given by the
theta coordinates. The construction gives a dense image and
\(e(\Lambda)=0\). Take a Lie uniformization
\(u:W/\Lambda'\to A_{\mathbb C}(\mathbb C)\).
The holomorphic group map \(e\) lifts to a complex-linear map
\(L:V\to W\) with \(L\Lambda\subseteq\Lambda'\). Let \(E_1,E_2\)
be the alternating real Riemann forms of the two systems.
Appell--Humbert comparison for the same line bundle gives
\(L^*E_2=E_1\).

The theta comparison also gives the following integral implication
for \(w\in W\):
\[
   E_2(L\Lambda,w)\subset\mathbb Z
       \ \Longrightarrow\ E_2(\Lambda',w)\subset\mathbb Z .
\]
Indeed translation of the theta bundle by \(u(w)\) changes its
multiplier on \(\ell\in\Lambda'\) by
\(\exp(2\pi i E_2(\ell,w))\). If it is unchanged on \(L\Lambda\),
the pulled-back translated coordinates, after adjusting the
quadratic exponential for the two automorphy forms, are theta
functions of the original type on \(V\). The space generated by
the coordinates contains the orbit of one nonzero coordinate
under the finite theta group, hence the entire theta space by its
irreducibility. These pulled-back coordinates therefore come from
sections of the original bundle. Their transformation law on
the dense image of \(e\) gives a section of the quotient of the
two bundles. On \(W\) its automorphy form is \(B'-B'=0\), where
\(B'\) is the form for the original and translated theta functions.
The quotient section is therefore an entire function
with unitary multiplier. It is nonzero at \(0\) by the absence
of common zeros of the two coordinate systems. Its absolute
value is bounded by periodicity on a compact fundamental domain,
so Liouville's theorem makes the function constant. Its multiplier
is one on all of \(\Lambda'\), proving the implication.\formalnote{16.82}{82}

Now \(L\) is injective by nondegeneracy of \(E_1\). If \(w\) is
orthogonal to \(L(V)\) for \(E_2\), apply the integral implication
to every real multiple of \(w\). It gives \(E_2(\Lambda',w)=0\),
so \(w=0\). Thus \(L\) is onto. If \(e(z)=0\), then \(Lz\in\Lambda'\).
For every \(y\) in the \(E_1\)-dual of \(\Lambda\), \(Ly\) pairs
integrally with \(L\Lambda\), hence with \(Lz\). It follows that
\(z\) is in the double dual, namely \(\Lambda\).
This proves \(V/\Lambda\simeq A_{\mathbb C}(\mathbb C)\) and
\(\dim A_{\mathbb C}=\dim_{\mathbb C}V=n[E:\mathbb Q]/2\).\formalnote{16.83}{83}

Properness of \(A_O/O\) now supplies a unique integral section
through each point of \(A_{\mathbb C}(\mathbb C)\). In fact the
valuative criterion gives existence and separatedness gives
uniqueness. Restriction is therefore an \(\mathcal O\)-linear
group isomorphism
\[
 j:A_O(O)\xrightarrow{\ \sim\ }A_{\mathbb C}(\mathbb C),
 \qquad \widetilde e(z)=j^{-1}(e(z)).
\]
The map \(\widetilde e\) has kernel \(\Lambda\). Its generic
restriction has the same theta coordinates and is immersive.
If an additive family of sections is holomorphic at zero on its
generic restriction, its restriction factors through \(e\) by a
complex-linear map. Injectivity of \(j\) gives the same factorization
through \(\widetilde e\). This is the universal property of the
uniformization on sections in the source. In particular the
holomorphic assertions concern the complex generic evaluation.
Its special fibre is the abelian variety obtained by reducing
the theta equations at the centre of the arc.\formalnote{16.84}{84}

Explicitly the rational vector \(b\in E^n\) labels the section
\(\widetilde e(Qb)\), where \(Q:(E\otimes\mathbb R)^n\to V\) is
the period isomorphism. Two such vectors label the same section
exactly if \(g^{-1}(b-b')\in\widehat{\mathcal O}^{\,n}\) for the
adelic frame \(g\). Rational approximation supplies lifts of every
finite-adelic torsion class, and these give all generic torsion
points because \(\Lambda\otimes\mathbb Q=Q(E^n)\).
The Weil pairing is the commutator of lifted theta translations,
which along \(e\) is the Riemann phase. The frame identifies it with
the specified finite-adelic character, with the common profinite
unit recording the action on roots of unity for a chosen embedding
of \(O\) in \(\mathbb C\). This gives the full structure used at
the centre of the arc in \S\ref{sec:theta-tangent}. Reduction
can identify distinct labeled sections.\formalnote{16.85}{85}

\subsection{Extension of homomorphisms}

Properness extends individual sections. For a Hecke correspondence
one needs a morphism between the abelian schemes carrying all
sections at once. Let \(g,g'\) be frames for two torus models over
the same \(O\). For \(z\in\mathrm{GL}_n(E)\) assume the rational
conjugation carries the first reflection to the second,
\(k=g'^{-1}zg\) has integral entries, and \(k^{-1}\) is a frame
for the scalar \(\mu>0\) in the conventions of the source. On
the uniformizing vector spaces \(F=Q'zQ^{-1}\) is complex-linear:
the reflection identity is precisely compatibility with the two
complex structures. Integrality of \(k\) gives \(F\Lambda\subset
\Lambda'\). The similitude identity gives \(F^*E_2=\mu E_1\).
The canonical multipliers then give
\(\chi'(F\ell)=\chi(\ell)^\mu\) on \(\Lambda\). An exponential
of a quadratic corrects the pulled-back theta functions for the
possible difference of automorphy forms with this same alternating
part. Thus \(v\mapsto\widetilde e'(Fv)\), as a system of
sections indexed by arguments of the first torus, has theta type
\((\mu B,\chi^\mu)\).\formalnote{16.86}{86}

Here is the graph construction for this extension. Write \(N,M\)
for the two projective theta bundles. Pick \(k_0=1\) or \(2\)
with \(\mu+k_0\) a unit in \(O\). Take the reduced closure
\(\Gamma\subset A_O\times A'_O\) of the pairs
\((\widetilde e(v),\widetilde e'(Fv))\) at the generic point.
Its proper projection onto \(A_O\) is surjective by density. The
product bundle \(p_1^*N^{k_0}\otimes p_2^*M\) on the graph
has the same theta type as \(N^{\mu+k_0}\). In this common type a
single scalar puts all its coordinate functions into integral
sections of \(N^{\mu+k_0}\), with no whole dual-lattice orbit of
common zeros. This uses integral theta translation by polarization
torsion of unit order. To exclude even one base point, choose a
unit integer \(l\) killing the dual quotient for the common type.
Pullback by \([l]\) on the graph multiplies each of the two
polarization classes by \(l^2\), by the parallelogram law.
Comparing their coordinate ratios before and after \([l]\)
on the dense complex graph gives a cross-multiplied regular
identity on \(\Gamma\). Translation by a dual-torsion section
leaves its \([l]\)-argument fixed. On charts at the point and
its translate the remaining factors are units. Hence if every
integral coordinate vanished there, their \(l^2\)-th powers
would vanish at the translate as well. This would give a full
orbit of common zeros. The integral coordinates therefore define
a morphism on \(A_O\). By density and the projective embedding
its graph sections are the prescribed pairs. Projection to
\(A'_O\) gives a homomorphism \(\varphi\) with
\(\varphi^*[M]=\mu[N]\), and its additivity and
\(\mathcal O\)-linearity follow on the dense sections.
Its image is closed and contains the dense target sections, since
\(F\) is surjective. The pullback of the ample class has positive
multiplier, so the fibres are zero-dimensional. It follows that
\(\varphi\) is an isogeny, carrying the torsion label \(x\) to
\(kx\). This stronger extension uses the integral polarized graph
in addition to the valuative criterion for points.\formalnote{16.87}{87}

The ingredients still used in this argument include finiteness of the signed-frame classes (\S\ref{sec:signed-frames}), the polarized ball families constructed in \S\ref{sec:geometry-closure}, complex Lie uniformization and the Appell--Humbert and
theta-group theorems invoked above, including integral scaling
of a tuple under dual-lattice translations. The good-reduction equations
and the integral chart construction are those of
\S\ref{sec:theta-family}. These hypotheses suffice for the specific
arc, the abelian scheme on it and its full labeled torsion.


\section{From marked fixed points to a definite trace}\label{sec:marked-orbits}

The geometric seed singles out one root of a characteristic polynomial.
The definite automorphic trace sums over the roots. To pass between them
we first compare the two counts with an \emph{arbitrary} weight on a
single marked root. In this comparison the rank \(n\) is odd.

\subsection{The two counts}

Let \(K=K(\mathfrak N)\) be the neat similitude level, let
\(x\in\mathrm{GL}_n(\mathbb A_{E,f})\), \(y\in\mathrm{GL}_n(E)\), and put
\[
 P=\det(T-y),\qquad t_x=x^{-1}y^{-1}x .
\]
For each archimedean place of the CM field \(E\) fix the embedding and
numbering of roots used in the construction. A mark \(\mu\) is an
index at the distinguished place \(v_0\). Sum over the simultaneous
rational equivalence classes of triples \((xK,y,\mu)\), where \(z\)
acts by \((zxK,zyz^{-1},\mu)\). A regular Hecke function \(f\) has
\(f(t_xK)=0\) unless \(P\) is separable. For the remaining terms write
\(J_\mu\) for the operator at \(v_0\) with eigenvalue \(+1\) on the
\(\mu\)-line of \(y\) and \(-1\) on the other eigenlines. Lagrange
interpolation expresses \(J_\mu\) as a polynomial in \(y\). In particular
its construction commutes with rational change of basis.

A definite triple admits a Hermitian form \(\beta\) over \(E\) such
that \(x^*\beta x\in\widehat{\mathbb Z}^{\,\times}I_n\) at the finite
adeles and \(\beta_v\) is positive at every infinite place. A marked
triple has the same finite condition and positivity away from \(v_0\),
with \(\beta_{v_0}J_\mu\) positive at \(v_0\). For a weight \(u(P,\mu)\)
let \(S_{\mathrm m}(f,u)\) and \(S_{\mathrm d}(f,u)\) be the sums of
\(f(t_xK)u(P,\mu)\) on the two respective sets of classes. Assume
finite support at weight \(1\) on both sides. In the occurrence argument
finiteness on the marked side is part of the hypothesis that a regular
operator is \(h\)-null. On the definite side the definite class set is
finite. For fixed \(x\) and a coset in the finite support of \(f\),
neatness gives at most one rational return \(y\).\formalnote{16.88}{88}

\subsection{Cancellation on a reducible polynomial}

At a split place \(w_0\) put \(\delta(s)=(s_{w_0}^{-1})^{\mathsf t}\).
There is a unitary embedding \(\iota\) from \(\mathrm{GL}_n(E_{w_0})\)
to the group at the split pair satisfying \(\delta(\iota(a))=a\).
Suppose \(f\) sends level vectors
in every representation to vectors on which \(\iota(J)\) acts by a
character \(\chi\). Applying this to the permutation representation on
right cosets, at the vector of the identity coset, gives
\[
 f(\iota(j)sK)=\chi(j)^{-1}f(sK)\qquad(j\in J).
\]
Cuspidality of the type means that every proper two-block unipotent
radical contains an element \(j\in J\) with \(\chi(j)\ne 1\).

Suppose \(P\) is reducible and the term is nonzero.
The local matrix \(g=\delta(t_x)\) is conjugate to \(y^{\mathsf t}\).
Choose a factorization of its characteristic polynomial into two
coprime monic factors of positive degrees, using the same choice for
each occurrence of the polynomial. It exists because \(P\) was already
reducible over \(E\) and is separable. The Chinese remainder theorem
gives \(p\) congruent to \(1\) on the first factor and \(0\) on the
second. The projection \(e=p(g)\) decomposes the local space into the
summands of \(g_1,g_2\). Write
\[
 U_e=\{1+X:X=eX(1-e)\}.
\]
Products of any two of its off-diagonal blocks vanish, so its group law
adds those blocks. Group the raw triples through \((xK,y,\mu)\) into
the block \(B=\{(x\iota(a)K,y,\mu):a\in U_e\}\). Changing the lift
\(x\) to \(xk\) conjugates the projection and its radical by
\(\delta(k)\), leaving the block unchanged. Starting from \(a\in U_e\)
conjugates the radical by \(a\), which preserves it. Rational
translation carries a block to the block through its translate.
Every member of \(B\) has the same admissibility for its sign
condition, since right multiplication by \(\iota(a)\) preserves the
finite Gram condition.

Choose \(j=1+C\in J\cap U_e\) with \(\chi(j)\ne1\).
We need \(d=1+D\in U_e\) with \(d^{-1}gd=jg\). Its single
off-diagonal equation is
\[
                   g_1D-Dg_2=Cg_2 .
\]
The kernel of the operator on the left consists of intertwiners.
An intertwiner also intertwines \(p(g_1)=1\) and \(p(g_2)=0\), so
the kernel vanishes. Equal finite dimensions give existence.
As \(j\) commutes with \(a\in U_e\),
\((da)^{-1}g(da)=j(a^{-1}ga)\).
The component of \(t_x\) left after removing \(\iota(g)\) commutes
with the embedded local group. This follows from the similitude
relation at the conjugate split place. Thus the actual adelic Hecke
argument for \(da\) equals \(\iota(j)\) times that for \(a\).
Translation by \(d\) permutes the points of \(B\) and multiplies
their terms by \(\chi(j)^{-1}\).

Only finitely many points in this block can contribute. Indeed \(f\)
is supported on cosets of similitudes with integral inverse. Hence
\(a^{-1}ga\) has integral entries at every contributing \(a=1+X\).
The matrix \(p(a^{-1}ga)=a^{-1}ea=e+X\) belongs to the image under
\(p\) of the compact set of integral matrices. The resulting parameter
\(X\) ranges in a compact set after taking its closure. The map from
the additive space of blocks \(X\) to \(B\) has open cosets as fibres,
since \(K\) has open pullback to the split place. A compact set meets
only finitely many fibres. The sum on \(B\) therefore equals
\(\chi(j)^{-1}\) times itself and vanishes.

To descend the cancellation, fix a block meeting the support, and
write \(\Gamma_B\) for its stabilizer in \(\mathrm{GL}_n(E)\).
For any contributing \(s\in B\) the map \(z\mapsto zs\) is a
bijection from \(\Gamma_B\) onto its orbit's intersection with \(B\).
Indeed \(zs\in B\) forces \(zB=B\) by the partition property, and
neatness gives a trivial point stabilizer, since a matrix centralizing
\(y\) fixes its reflection. The preceding finiteness makes
\(|\Gamma_B|\) a positive finite integer. Each contributing rational
orbit meeting \(B\) thus occurs exactly \(|\Gamma_B|\) times in its
zero sum. Divide by this integer and sum over the block orbits.
The reducible terms vanish for each choice of signature separately.\formalnote{16.89}{89}

\subsection{Matching the irreducible terms}

Suppose \(P\) is irreducible and a signed term contributes. There is
a positive \(\nu\in\mathbb Q\) for which
\(\nu t_x^*t_x\in\widehat{\mathbb Z}^{\,\times}I_n\).
It follows that \(y^*\beta y=\nu\beta\). Indeed \(x^*y^*\beta yx\)
and \(\nu x^*\beta x\) are rational-unit-idelic multiples of one
another by the two frame identities. Over \(E\) this gives a
rational proportionality factor which is a unit at all finite
primes, hence \(\pm1\). Positivity on the signed archimedean form
selects \(+1\). If \(r\) is a complex eigenvalue, evaluate the
similitude relation on its eigenvector against the positive form
\(\beta_vJ_v\), using \(J_v=1\) away from \(v_0\) and \(yJ_v=J_vy\).
The positive squared length cancels and gives \(r\bar r=\nu\).
Consequently \(L=E[T]/(P)\) is CM with conjugation
\(\theta\mapsto\nu/\theta\) on \(\theta=T\bmod P\).
This root norm also covers a similitude arising from a Frobenius
return. The signs used in the comparison are signs of Hermitian
eigenlines at infinity.

At \(v_0\) prescribe sign \(+1\) at the root indexed by \(\mu\), the eigenvalue of \(y\) on the \(+1\) line of \(J_\mu\), and \(-1\) on
its \(n-1\) companions, and at every other place prescribe \(+1\).
The product over the real places of \(L^+\) is \((-1)^{n-1}=1\).
The quadratic idele norm theorem gives \(a\in(L^+)^\times\) with
these signs and \(h\in\mathbb A_{L,f}^{\times}\) with
\(a=h\bar h\) at all finite places. In fact, take the idele over
\(L^+\) with the prescribed signs and trivial finite components.
Its Artin symbol in the quadratic extension \(L/L^+\) is trivial
by the sign product. Idelic reciprocity places it in the product
of principal ideles and norms, giving the stated \(a,h\).

Evaluate \(h\) as a finite adelic polynomial in \(y\), writing
\(H=h(y)\). The adjoint for \(\beta\) on \(E[y]\) sends
\(y\) to \(\nu y^{-1}\), so \(H^*\beta H=\beta a(y)\).
Consequently \(\beta\) at \(Hx\) and \(\beta a(y)\) at \(x\)
give the same finite Gram matrix and exactly the prescribed change
of archimedean signs. The reverse operation uses \(H^{-1}\).
Choose the polynomials using \(P,\mu\). Conjugate matrices then give
conjugate changes of frame, and commuting with the centralizer makes
the map well defined on rational classes. Since \(H\) commutes with
\(y\), \(t_{Hx}=t_x\), and the mark is preserved. This bijection
identifies the two sums term by term on irreducibles. Combined with
the cancellation it gives \(S_{\mathrm m}(f,u)=S_{\mathrm d}(f,u)\).\formalnote{16.90}{90}

\subsection{Making every positive power a regular test}

Let \(\psi\ne0\) be a level vector of type \((J,\chi)\) at \(w_0\),
spherical at a second split place \(w_1\). The two split pairs are
disjoint, the level is prime to their conjugates, and its local
pullback \(K_0\) lies in \(J\subset\mathrm{GL}_n(\mathcal O_{w_0})\),
while \(K_1\) lies in the pro-Iwahori \(I\) at \(w_1\).
Let \(c\psi=\psi\), with \(c\) supported on similitudes trivial in
the dual component at \(w_1\). Write \(\iota_i\) for the two split
embeddings. The first local test is
\[
 b_0=\frac{1}{[J:K_0]}\sum_{j\in J/K_0}\chi(j)^{-1}[\iota_0(j)] .
\]
Here and below the brackets denote group-algebra terms for a choice
of coset representatives. On level vectors the average is independent
of that choice. Indeed \(\chi\) is trivial on \(K_0\) because
\(\psi\ne0\) is fixed by \(K_0\). Translation of the sum gives its
\(\chi\)-transformation law, and \(b_0\psi=\psi\).

For \(w_1\) choose a uniformizer \(\varpi\) and put
\(t=\operatorname{diag}(\varpi^{-1},\ldots,\varpi^{-n})\).
Let \(e_I\) denote normalized averaging over \(I\).
The operator \(e_I\pi(t)\) on \(I\)-invariants is injective over
\(\mathbb C\). Here is the argument from the local Iwahori algebra.
Factor \(t\) into the \(n\) diagonal translations with entries
\(\varpi^{-1}\) on one terminal interval and \(1\) elsewhere.
Their averaged operators multiply because they are dominant. For
each such translation a cyclic permutation gives a monomial matrix
\(m\) with \(m^{-1}Im\subset I\). If its translation average
annihilates \(x\), the vector \(v=\pi(m)x\) is \(I\)-fixed and
\(e_I\pi(\sigma^{-1})v=0\) for that inverse cycle \(\sigma\).
For an adjacent inversion the rank-two Iwahori quadratic relation
propagates this equality to the permutation with that inversion
removed. In the relation the extra terms average unit-diagonal
translates of \(\pi(\sigma^{-1})v\), which vanish since unit diagonal
matrices normalize \(I\). Sorting to the identity gives \(e_Iv=v=0\).
Thus each factor and their product have zero kernel.\formalnote{16.91}{91}

Choose a positive integer \(D\) clearing the denominators of
\(\iota_1(t)^{-1}\), put \(z=D^{-1}I_n\) adelically and
\(U=I\cap K_1\cap tK_1t^{-1}\). Take
\[
 b_1=\frac{1}{[I:U]}\sum_{j\in I/U}[\iota_1(j)\,z\,\iota_1(t)] .
\]
The terms have integral inverse and the average sends level vectors
to \(I\)-fixed vectors. Undoing \(z\) and averaging over \(I\)
recovers \(e_I\pi(t)\) on \(I\)-fixed vectors, so \(b_1\) kills
no such nonzero vector. Since \(b_1\) commutes with \(b_0c\) and
\(b_0c\psi=\psi\), \(b=b_1b_0\) satisfies
\[
                      (bc)^m\psi=b_1^m\psi\ne0\qquad(m\ge0).
\]
Thus the iterates stay nonzero throughout the trace argument.
For \(m\ge1\), commute one \(b_1\) past \(b_0c\) to put \(b_0\)
on the far left. Thus \((bc)^m\) sends level vectors into the
\(\chi\)-eigenspace.

The dual component of \(bc\) at \(w_1\) lies termwise in
\(I(Dt)I\). Dominant double cosets multiply into the double coset
of the product: write the intervening element of \(I\) as an
upper-unitriangular factor which passes to the left and a remaining
factor which passes to the right under the dominant translations.
The diagonal valuations remain strictly ordered in each positive
power. To see regularity explicitly, a point of \(IsI\), for such
a diagonal \(s\), is conjugate to \(sk\) for \(k\in I\).
Every principal minor of \(k\) is a unit by triangular reduction.
The degree \(n-j\) coefficient of the characteristic polynomial is,
up to sign,
\[
          \sum_{|A|=j}\Bigl(\prod_{i\in A}s_i\Bigr)\det(k_{A,A}).
\]
The unique lowest additive valuation comes from the \(j\) smallest
valuations of the \(s_i\). The successive coefficient valuations
therefore give \(n\) distinct root valuations by the Newton polygon.
This proves separability, including after right multiplication by the
level, whose local dual component lies in \(I\). Applied to \(t_x\)
it proves regularity of each positive power.\formalnote{16.92}{92}

Finally, if \(\sum_\mu u(P,\mu)=\operatorname{tr}\Xi_{\rm alg}(y^{-1})\)
on the contributing polynomials, evaluating definite forms at
representatives of the finite class set makes the trace a sum of
diagonal blocks. A term returns to its own class precisely when its
coset contains \(x^{-1}y^{-1}x\). The return acts by \(\Xi_{\rm alg}(y^{-1})\), with \(\Xi_{\rm alg}\) the algebraic representation of \S\ref{sec:seed-proof}.
The trace is therefore \(S_{\mathrm d}(f,u)=S_{\mathrm m}(f,u)\).
Applied at each positive power with the weight produced by the
coefficient argument, marked vanishing would force every positive
power trace to vanish. Newton's identities would make \(bc\)
nilpotent on the finite-dimensional definite space, contradicting
\((bc)^m\psi\ne0\). The arithmetic inputs to this count are
finiteness of definite classes, neatness, and the quadratic idele
norm theorem. The local averaging uses the Iwahori relation above.
The passage from vanishing of geometric local terms to vanishing
of the marked count and the construction of its coefficient weight
enter separately.\formalnote{16.93}{93}

\subsection{Detection of a Hodge weight by the marked sum}\label{sec:woods-hole}

For smooth proper schemes over an algebraically closed field, the holomorphic Lefschetz formula for a correspondence \((g,f)\) with transverse coincidences and with \(g\) \'etale there expresses a cohomological trace as a sum of local terms. Each term is the trace on the coefficient fibre divided by \(\det(1-df\,(dg)^{-1})\) \cite[Theorem~1.6]{lit-kondyrev-prikhodko}. For the test-ring cohomology used here we prove the vanishing implication below.
The marked count vanishes if the relevant linear combination of pullbacks annihilates the weight-\(h\) images of the Hodge classes. Work over characteristic-zero fields \(E\subset k\subset\overline{\mathbb Q}_p\), with \(k\) algebraically closed. Let \(Y\) be proper, \(W\) projective over the specified coefficient ring, and suppose their indicated generic fibres have atlases of the same dimension \(d\). Let \(u_i:W\to Y\) be finitely many maps with coefficients \(b_i\in k\), and fix \(h\ge0\). Write \(\lambda_h\) for the weight-\(h\) map from Hodge classes of \S\ref{sec:hodge-vector}. Assume that \(h>d\), or that for each such class \(x\) on \(Y\) of degree at most \(d\),
\[
                            \sum_i b_i\lambda_h(u_i^*x)=0.
\]
For a further map \(c:W\to Y\), consider its coincidences \(u_i(w)=c(w)\) with \(b_i\ne0\). At each take the rank-\(d\) tangent polynomial \(P_{iw}\) and assume \(P_{iw}(1)\ne0\). In tangent coordinates \(du_i=dc\, A_{iw}\) and \(P_{iw}(t)=P_{iw}(0)\det(1-t A_{iw})\), with \(dc\) injective. Thus \(du_i-dc\) is invertible. There is no nonzero first-order deformation of a coincidence, so the coincidence schemes are reduced and zero-dimensional. Separatedness gives closed coincidence loci, and the finite-type charts give finitely many such points. The assertion to prove is\formalnote{16.94}{94}
\[
                       \sum_{i,w} b_i
                      \frac{[t^h]P_{iw}(t)}{P_{iw}(1)}=0.       \tag{WH}
\]

On a square-zero test algebra \(B\oplus M\) let \(s_t(a,m)=(a,tm)\). Pull the diagonal class \(\Delta\) of \(Y\times Y\) back along the pair \((u_i\circ s_t,c)\), after extending to \(k\); call this pullback \(\Phi_i(t)\Delta\). Its support lies at the coincidences. In local coordinates, the differences of the coordinates on the two branches cut out a reduced point. Pullback of the top exterior symbol for these differences scales its coefficient by the ratio of their differential determinants,
\[
   \frac{\det(t\,du_i-dc)}{\det(du_i-dc)}
      =\frac{\det(1-tA_{iw})}{\det(1-A_{iw})}
      =\frac{P_{iw}(t)}{P_{iw}(1)} .
\]
The local symbol for the reduced point represents its point class \([w]\). The diagonal-class formula with supports therefore gives
\(\Phi_i(t)\Delta=\sum_w(P_{iw}(t)/P_{iw}(1))[w]\).
For projectively embedded \(W\), the distinct cohomology classes among the \([w]\) have linearly independent images under \(\lambda_d\). Choose a linear functional \(\varphi\) taking value one on each of these images. If several points have the same class, every one of their summands still contributes its own coefficient. After applying \(\varphi\lambda_d\), coefficient extraction on the last displayed pullback computes exactly the \(i\)-th inner sum in (WH).\formalnote{16.95}{95}

Compute the same coefficient using the K\"unneth generators for the diagonal class. On a cross product \(x\times y\), pullback along the pair is the cup product of the scaled pullback of \(u_i^*x\) with the unscaled pullback of \(c^*y\). Under \(\lambda_d\) the scaling gives the factor \(t^j\) on the component with first weight \(j\). The coefficient of \(t^h\) retains only \(j=h\), and the other factor is independent of \(i\). For a first class of degree \(q\), the fixed scalar in the cup rule is \(\binom dh(-1)^{(q+h)(d-h)}\) when \(h\le d\). The weighted sum is consequently zero by the hypothesis on \(u_i^*x\); if \(h>d\), every such coefficient already vanishes. Formally one can take coefficients after applying any linear functional on the target: evaluation over the infinite field \(k\) embeds \(k[t]\) in the space of functions \(k\to k\), and its coefficient functional extends linearly to that space. Additivity proves the vanishing for the diagonal, a sum of cross products. With the preceding choice of \(\varphi\) this proves (WH).\formalnote{16.96}{96}

For the Hecke correspondence, apply (WH) on a common sufficiently deep level for the finitely many maps. Group the coincidences by their marked classes. On a neat level a fibre has the order of the deck group times the number of solutions in its marked lattice. That number cancels the corresponding lattice denominator in the local coefficient. After division by the common deck degree the result is exactly marked vanishing. Thus a nonzero marked sum entails a surviving weight, which is the direction used in \S\ref{sec:seed-proof}. The cohomological inputs above are the construction and local transformation rule of the diagonal, generation by cross products, and independence of the images of the point classes.\formalnote{16.97}{97}


\editionchapter{Comparison on test rings and cohomological Frobenius}\label{ch:cohom}

The geometric calculation in Chapter~\ref{ch:shimura} gives an identity among points. For each point of a common cover in characteristic \(p\), the images under certain Hecke correspondences, followed by powers of Frobenius, cancel with specified integer coefficients. Chapter~\ref{ch:galois} uses the Galois action on the characteristic-zero cohomology to construct its initial odd-rank representations, starting with a Hecke eigensystem quotient of a direct sum of cohomology groups. We transport the identity to a relation among linear operators on the cohomology groups after extending scalars to the crystalline period ring. The test rings below allow both fibres to be evaluated from the same integral family. In ordinary \(p\)-adic Hodge theory, the crystalline comparison theorem for a smooth proper formal scheme over the integers of a \(p\)-adic field identifies, after extension to \(B_{\mathrm{cris}}\), the \(p\)-adic \'etale cohomology of the geometric generic fibre with the crystalline cohomology of the special fibre \cite[Theorem~1.1(i)]{legacy-bms}. In this chapter the cohomology is defined on test rings. Its crystallinity and the action of Frobenius are proved directly on those rings.

\section{Valuation cells, test envelopes, and continuity}\label{sec:test-foundations}
To study a point above a special-fibre prime, we extend its coefficient ring to a valuation ring. For a homomorphism \(B\to W\) into a valuation ring, its centre is the inverse image in \(B\) of the maximal ideal of \(W\). Prescribing that centre keeps the point over the chosen prime. We first construct such an extension in mixed characteristic. We then construct compatible covers for diagrams of ordinary test rings. These covers identify cohomology at the limiting vertex with the filtered colimit of the cohomologies at the individual stages in the proof of proper base change.

\subsection{Producing a cell with prescribed centre}
Fix the prime \(p\). A cell is a \(p\)-adically complete and separated valuation domain \(W\) in which \(p\ne0\) and every monic polynomial splits. Then \(p\) is a nonunit by separatedness. Given a ring \(B\) in which \(p\) acts injectively and a prime \(\mathfrak q\) containing \(p\), there is a map \(B\to W\) with centre \(\mathfrak q\).

Indeed \(B_{\mathfrak q}[1/p]\) is nonzero. Contract a prime of this ring to \(\mathfrak l\subset B_{\mathfrak q}\), put \(D=B_{\mathfrak q}/\mathfrak l\), and take an algebraic closure of its fraction field. In that field a valuation ring \(U\) dominates \(D\) \cite[\href{https://stacks.math.columbia.edu/tag/00IA}{Tag 00IA}]{arith-stacks}. Thus \(p\) is nonzero and is a nonunit in \(U\), and every monic polynomial splits since \(U\) is integrally closed in that field. To retain these properties on completion the source uses a nonprincipal ultrafilter \(\mathcal U\) on \(\mathbb N\). The ultrapower \(G=U^{\mathbb N}/\mathcal U\) is a valuation domain. For either of two divisibility relations between two sequences one holds on a set in \(\mathcal U\), and the coordinatewise quotients on that set give the indicated division. The same argument handles zero products and shows that monic polynomials split. The diagonal map \(U\to G\) is local because any inverse of a constant has an inverse coordinate.

The map \(G\to\varprojlim_n G/p^nG\) is onto. To see this, write any compatible residues as the partial sums of \(\sum_{j\geq0}p^j c_j\), choosing \(c_j\in G\) inductively. If \(c_j\) is represented by \(c_j(k)\), the sequence \(s(k)=\sum_{j<k}p^jc_j(k)\) agrees with the first \(n\) terms modulo \(p^n\) for all \(k\geq n\). Its class has the prescribed residues. Thus the completion \(W\) is \(G/I\) with \(I=\bigcap_n p^nG\), and \(W/p^nW=G/p^nG\). The ideal \(I\) is prime. If \(a,b\notin I\), then \(a\mid p^r\), \(b\mid p^s\) for some \(r,s\) by total divisibility. If \(ab\in I\), then \(p^{r+s+1}\mid ab\mid p^{r+s}\), contradicting regularity and noninvertibility of \(p\). Also \(p\notin I\) by the same cancellation. The quotient \(W\) inherits total divisibility and splitting of monic polynomials. The proper quotient of the local ring \(G\) is local with the same inverse image of the maximal ideal. Hence \(W\) is the required cell. Section~\ref{sec:primitive-comparison} applies this to \(B=\mathcal O_C\otimes_{A_{\inf}}W(V)\) and the kernel of \(B\to\kappa(V)\).\formalnote{17.1}{1}

\subsection{Products and tilted maps}
A \(p\)-adic test is a nonempty product \(T=\prod_i W_i\) of cells with a map \(\mathcal O_C\to T\). Morphisms are all \(\mathcal O_C\)-algebra maps of products. The ring \(T\) is \(p\)-torsion-free and complete by coordinates. Its tilt is \(T^\flat=\varprojlim_{x\mapsto x^p}T/p=\prod_i W_i^\flat\). In each factor choose \(\varpi\) from a \(p\)-th root of \(p\) in \(\mathcal O_C\). The ring is complete for \(\varpi\), and \(W_i/\varpi\to W_i/p\), \(x\mapsto x^p\), is onto by splitting of monic polynomials. If \(x^p\) is divisible by \(p\), then \((x/\varpi)^p\in W_i\); integral closedness gives \(x/\varpi\in W_i\). Hence this Frobenius is an isomorphism on each factor and on the product. The perfectoid criterion \cite[Lemmas~3.9 and~3.10(ii)]{legacy-bms} now gives surjectivity of \(\theta_T:W(T^\flat)\to T\) and a principal kernel. The chosen generator \(\xi\) over \(\mathcal O_C\) remains distinguished under this map of Witt vectors, so generates \(\ker\theta_T\) \cite[Remark~3.11 and proof of Lemma~3.13]{legacy-bms}. In particular
\[
                  \mathcal O_C\otimes_{A_{\inf}} W(T^\flat)\ \simeq\ T .
\]
The element \(p^\flat=\xi\bmod p\) used in the comparison is regular on \(T^\flat\). Indeed its sharp in each cell divides \(p\), hence is nonzero. The multiplicative description of a tilt by compatible \(p\)-power roots in the cell shows that each tilted factor is a domain and that sharp detects zero there.
\phantomsection\label{route:30}\leanroute{Lean proves \(\ker\theta_T=(\xi)\) by reading divisibility of zeroth Witt coordinates after untilting in each valuation cell. We use the perfectoid criterion after verifying \(\varpi\)-adic completeness, regularity of \(\varpi\), and the Frobenius isomorphism \(T/\varpi\xrightarrow{\sim}T/p\). The test-specific adjunction and its full faithfulness are proved below.}\formalnote{17.2}{2}

For a perfect \(\mathcal O_C^\flat\)-algebra \(b\), the following correspondence is natural:
\[
  \operatorname{Hom}_{\mathcal O_C^\flat}(b,T^\flat)
   \simeq \operatorname{Hom}_{\mathcal O_C}
          (\mathcal O_C\otimes_{A_{\inf}}W(b),T).
\]
A map \(f\) on the left gives \(\theta_T W(f)\). Conversely a map on the right restricts to \(g:W(b)\to T\). Reduction gives \(\bar g:b\to T/p\). The coordinates \(\bar g(x^{1/p^n})\) give the unique lift \(f:b\to T^\flat\) with that initial coordinate. The sequence \(g([x^{1/p^n}])\) is already a compatible system of roots in \(T\), so \(g([x])=f(x)^\sharp\). Every Witt vector is a convergent sum of \(p\)-powers times Teichmüller representatives. Thus the two maps agree on all of \(W(b)\), and the base-algebra condition follows by restriction to \(\mathcal O_C^\flat\). Taking \(b=T_1^\flat\) gives full faithfulness on tests.\formalnote{17.3}{3}

\subsection{A test through which all test points factor}
An ordinary test is reduced, splits monic polynomials, and for \(a,b\) has an idempotent \(e\) with \(a\mid eb\) and \(b\mid(1-e)a\). Products are tests by taking the witnesses coordinatewise. To bound the size of the rings through which points factor, fix a ring \(B\). Consider the following equations for one unknown \(x\) with specified coefficients:
\[
\begin{gathered}
 x^2=x,\quad a\mid xb,\quad b\mid (1-x)a
       \quad\text{(one simultaneous condition)},\\
 a^n\mid b_n-x\quad\text{for all }n\geq0,\qquad P(x)=0 .
\end{gathered}
\]
For each kind of equation form a formal name for a solution whenever each coefficient is written as a polynomial over \(B\) in finitely many earlier names. Iterate this construction through the countable ordinals. Each proposed equation uses at most countably many earlier terms, whose stages have a common countable upper bound. The union is therefore a set of well-founded terms closed under all three operations. Write \(P_B\) for the polynomial ring on these names.

Given \(B\to V\), evaluate each name recursively at a chosen solution in \(V\) if one exists, and at zero otherwise. The image \(D\subset V\) of \(P_B\) then contains a solution of every indicated condition over \(D\) solvable in \(V\), by representing each coefficient as one of the defining polynomials. The linear polynomial \(ax-b\) shows that divisibility of elements of \(D\) in \(V\) is divisibility in \(D\). Thus for an ordinary test \(V\), the image \(D\) is itself a test. For a cell \(V\), \(D\) is a valuation domain and contains the limits of all compatible residues modulo powers of \(p\), by the second equation with \(a=p\). It is separated by injection and reflection of divisibility, and \(p\) stays nonzero. Hence \(D\) is a cell in this case.\formalnote{17.4}{4}

For an \(\mathcal O_C\)-algebra \(B\) admitting at least one test point, take the product of all quotients \(P_B/I\) which are cells. This is indexed by a set of ideals and is nonempty by the preceding construction. The diagonal map from \(B\) to this test covers all points. Given \(B\to\prod_j V_j\), each component factors through one of the chosen quotients. Project to that quotient and map to \(V_j\); the resulting product map gives the desired factorization over \(\mathcal O_C\). Combined with tilted full faithfulness this is the covering point used in (PC1).\formalnote{17.5}{5}

On ordinary tests we need functoriality for a specified small family of rings \(B_i\). Fix the set \(\mathcal F\) of all test quotients of all \(P_{B_i}\). For every ring \(R\) form
\[
              E(R)=\prod_{Q\in\mathcal F,\;\alpha:R\to Q} Q,
              \qquad \epsilon_R(r)=(\alpha(r))_{Q,\alpha}.
\]
Precomposition on indices makes \(E\) a covariant functor and \(\epsilon\) natural. For \(R=B_i\), each map \(R\to V\) to an ordinary test factors through \(\epsilon_R\). Indeed the previous image \(D\subset V\) is such a quotient and its projection from \(E(R)\) followed by inclusion gives the original map. Factorizations can be chosen independently; the maps \(\epsilon_R\) themselves are compatible throughout the diagram.\formalnote{17.6}{6}

\subsection{Continuity on a relatively affine atlas}
Write \(h_B(R)=\operatorname{Hom}(B,R)\). For an abelian group \(\Lambda\) with discrete topology, take the coefficient \(L(T)=\operatorname{LC}(\operatorname{Spec}T,\Lambda)\) on ordinary tests and \(H^n(Z)=\operatorname{Ext}^n(\mathbb Z[Z|_{\rm tests}],L)\) in their category of abelian-group-valued functors. Suppose \(h_A\to Y\) is relatively affine and its image on tests is cohomologically dense for \(L\). Thus the preimage of the image over each representable test has the cohomology of that test. For a small cofiltered inverse system \(Y_j\) of functors relatively affine over \(Y\), let \(Y_*=\varprojlim_jY_j\) be its pointwise limit over \(Y\). Then
\[
                  \varinjlim_j H^n(Y_j)\ \xrightarrow{\ \sim\ }\ H^n(Y_*)
                     \qquad (n\geq0).
\]
On a test representable \(h_T\), evaluation is exact and gives \(L(T)\) in degree zero and zero in positive degrees. A filtered ring colimit of ordinary tests is a test. The finite lists of coefficients of each monic polynomial and of each divisibility comparison lift to one stage, and each nilpotence equation holds at some stage where the rings are reduced. Moreover locally constant functions on an affine spectrum have finite image by quasi-compactness. Their level sets are a finite idempotent partition, whose entries, equations and labels descend to a common stage. Equality of two descended functions can likewise be tested by the finitely many equations on their partitions. This proves the formula for a diagram of test representables in all degrees.\formalnote{17.7}{7}

Pull the diagram back along \(h_A\to Y\), obtaining affines \(h_{B_j}\) and limit \(h_{B_*}\) with \(B_*=\varinjlim B_j\). Apply the ordinary envelope to the single family consisting of all \(B_j\) and \(B_*\). The maps \(p_j:h_{E(B_j)}\to Y_j\) have cohomologically dense image because their image on tests contains the preimage of the original atlas image. This persists after any representable pullback. At the limiting vertex the source is \(h_{\varinjlim E(B_j)}\). By naturality there is a map \(\varinjlim E(B_j)\to E(B_*)\), so the dense cover from \(h_{E(B_*)}\) factors through the limiting cover. This is the reason for including the limiting ring in the family.\formalnote{17.8}{8}

To descend continuity through these covers, order the assertions as injectivity in degree zero, surjectivity in degree zero, injectivity in degree one, and so on. Induct on the number \(r\) of assertions, simultaneously for all systems relatively affine over \(Y\). For a map \(p:E\to Z\) the augmented complex whose \(m\)-th group is free on \(E^{m+1}_Z\), with alternating deletion differential, resolves \(\mathbb Z[\operatorname{im}p]\). It is exact at every test because each fibre over an image point has an element, and prepending it contracts that fibre's complex. Both the stagewise filtered colimit of Ext groups and the Ext at the limiting vertex give long exact sequences in the chain argument, by exactness of filtered colimits of abelian groups. For a short exact sequence the four lemma gives the first \(r\) assertions on the quotient from the first \(r\) on the middle and the first \(r-1\) on the kernel. Applying this successively to the kernels in the augmented complex requires only the first \(r-1\) assertions on all the positive fibre powers and the first \(r\) on \(E\). Indeed the number required on each subsequent kernel drops by one until it reaches zero.

Here the sources of \(p_j\) are test representables, for which all the assertions hold. Each fibre power is relatively affine over \(Y\) by pullback and composition, and its limiting vertex is its own fibre power, since finite limits commute with limits. The induction hypothesis gives the first \(r-1\) assertions on these diagrams. The augmented complexes thus give \(r\) assertions on the images. Cohomological density identifies the cohomology of each image with that of its target, at the stages and at the limiting vertex. This completes the induction. Density here is used in its base-change-stable form from \S\ref{sec:test-proper-base}; restricting dense images to a functor preserves cohomology by resolving that functor by free test representables.\formalnote{17.9}{9}

For the applications, decreasing intersections of relatively affine conditions are pointwise limits. Base changes to a directed union of algebras are also such limits by the universal property of the union. The displayed formula therefore gives both continuity assertions in \S\ref{sec:test-proper-base}. The proof uses the chosen atlas's density. In the proof of the valuative density criterion there, the affine test base has its identity atlas. The preimage of the proposed cover over that base is an image of an affine pullback, which maps pointwise onto this image and so gives it a dense atlas as well. Both sides therefore have continuity before invoking the valuative criterion. The criterion then supplies density for the atlas of a proper functor. The coefficient in the continuity theorem is the locally constant ordinary-test coefficient. The separate almost coefficient on tilted tests uses (PC1) and its descent argument in \S\ref{sec:primitive-comparison}.\formalnote{17.10}{10}


\section{Proper base change and field detection on test rings}\label{sec:test-proper-base}

For a proper morphism of schemes, the higher direct images of a torsion abelian \'etale sheaf commute with arbitrary base change \cite[Expos\'e~XII, Theorem~5.1(iii)]{lit-artin-proper}. Here we establish the corresponding fibre statements for constant \(\mathbb F_\ell\) coefficients in the category of test rings.

The line formula in \S\ref{sec:projective-test} needs two facts about a cohomology class in a family. Its zero set on the residue fields of a test is open, and if this set is the whole spectrum the class vanishes. The principal-open argument needs the same facts with coefficients restricted to the zero locus of a multihomogeneous equation. Both arguments also need invariance under extension of algebraically closed fields. We prove them in the category in which the classes were constructed. Section~\ref{sec:test-foundations} gives the continuity theorem for the relatively affine diagrams used in the proof.

Fix a prime \(\ell\), equal to \(p\) in the application. Let \(\mathcal C\) be the full category of reduced rings in which monic polynomials split and for every \(b,c\) there is an idempotent \(e\) such that \(b\mid ec\) and \(c\mid(1-e)b\). For a covariant functor \(Y\) from rings to sets, write
\[
 L(T)=\operatorname{LC}(\operatorname{Spec}T,\mathbb F_\ell),\qquad
 H^n(Y)=\operatorname{Ext}^n_{\operatorname{Fun}(\mathcal C,\mathbf{Ab})}
                 (\mathbb Z[Y|_{\mathcal C}],L).
\]
Here \(\operatorname{LC}\) denotes locally constant functions. A structure map \(Y\to h_D\) specifies \(D\to T\) at each point of \(Y(T)\). In these sections a proper point functor has the following two properties. Equality of two points over any ring is cut out by an ideal. There are a ring \(S\) and a point \(a\in Y(S)\) such that for every algebraically closed field \(K\), valuation ring \(V\subset K\) and \(y\in Y(K)\) carrying \(D\) into \(V\), some \(S\to V\) carries \(a\) to \(y\) over \(K\). In particular two points over any ring agree if their images agree along jointly injective ring maps. Apply those maps to the ideal cutting out their equality.\formalnote{17.11}{11}

\subsection{A cover by spaces of places}
First suppose \(D\) is an absolutely integrally closed valuation domain with maximal ideal \(\mathfrak m\) and \(\ell\in D^\times\). Write \(Y_0\subset Y\) for the points whose structure maps kill \(\mathfrak m\). We claim that for every proper \(Y\) over \(D\),
\[
                         H^n(Y)\ \xrightarrow{\ \sim\ }\ H^n(Y_0)
                         \qquad(n\ge0).                         \tag{PB1}
\]
We describe the acyclic pieces that make this true. Starting with \(a\in Y(S)\), form
\[
  P=\prod_{\mathfrak p\in\operatorname{Spec}S}\overline{\kappa(\mathfrak p)},
                   \qquad \xi=a|_P.
\]
The ring \(P\) is absolutely flat and all its residue fields are algebraically closed. Indeed \(x^2 y=x\) has a coordinatewise solution. Monic polynomials split in the product, hence in its quotients, and its prime quotients are fields. A place \(j\) is a prime \(q\) of \(P\) with a valuation ring \(V_j\subset\kappa(q)\) containing the image of \(D\). Choose the maps \(S\to V_j\) from properness that represent \(\xi|_{\kappa(q)}\). They give \(z\in Y(\prod_j V_j)\) agreeing with \(\xi\) over \(\prod_j\kappa(q_j)\), by the diagonal property.

For \(t\in P\), represent \(t\bmod q_j\) by the coprime numerator and denominator \((t\bmod q_j,1)\) if it belongs to \(V_j\), and \((1,(t\bmod q_j)^{-1})\) otherwise. Write \(n_t,d_t\) for the resulting product elements. The functor \(E\) of places assigns to \(C\) the maps \(f:\prod_jV_j\to C\) modulo the equivalence of having the same image of \(z\) and the same ratios \([f(n_t):f(d_t)]\) for all \(t\). Each pair generates the unit ideal, so cross multiplication tests equality of ratios. Forgetting them gives \(u:E\to Y\). Its geometric points cover \(Y\). A point given by \(S\to K\) factors through one \(\overline{\kappa(\mathfrak p)}\to K\) after passing to that factor of \(P\). The place with trivial valuation there gives its lift.

Two properties of this construction matter for descent. First \(E\) is proper. Equality has the equations for equality in \(Y\) together with the cross products of the ratios. Over a local ring \(C\) a representative \(f\) factors through the localization of \(\prod_jV_j\) at \(f^{-1}\mathfrak m_C\). This is an absolutely integrally closed valuation ring \(V'\). Its fraction field receives \(\prod_j\kappa(q_j)\), the localization of \(\prod_jV_j\) at its regular elements. Thus \(P\) maps into this field through some \(\kappa(q)\). Pulling \(V'\) back to \(\kappa(q)\) gives a valuation ring \(V\) and a local map \(V\to C\) giving the original point. The ratios agree over the fraction field of \(V'\), as do the two images of \(z\), hence they agree over \(V'\). One may equally represent the point using any smaller valuation ring of \(\kappa(q)\) containing the image of \(D\), by restriction. In particular, for \(C=K\) a field and a specified valuation ring \(U\subset K\), pull \(U\) back along the local map \(V\to K\). The resulting subring of \(V\) is a valuation ring of \(\kappa(q)\): if an element of \(V\) maps outside \(U\) its image is nonzero, so the element is a unit of \(V\) and its inverse maps into \(U\). Projection to this smaller place gives the required lift over \(U\).\formalnote{17.12}{12}

Second, \(u\) has cohomologically dense image. This means that after pullback to every representable test point, inclusion of the image induces isomorphisms on \(H^n\). We record the underlying descent facts for this theory. A relatively affine map lifting all points over valuation rings of algebraically closed fields has dense image. To see the criterion, pull back to \(h_T\). Split repeatedly along complementary idempotents of \(T\). On an idempotent partition the locally constant coefficient is the product of its values on the parts, so restriction detects equality and glues classes. A surviving kernel or cokernel class can be followed to a maximal ideal of equations of vanishing by filtered continuity for affine charts. Such an ideal must contain \(e\) or \(1-e\) for each idempotent. Choose a maximal ring ideal \(\mathfrak n\) containing it. For \(b\notin\mathfrak n\), divisibility on a partition yields an idempotent \(e\notin\mathfrak n\) divisible by \(b\): apply comparability to \(b\) and \(1-rb\), where \(rb=1\bmod\mathfrak n\). On the part where \(b\) divides \(1-rb\), it divides \(1\), and this part contains \(\mathfrak n\). Therefore every point killing the chosen ideal inverts \(b\), since it kills \(1-e\). It factors through \(T_{\mathfrak n}\), an absolutely integrally closed valuation ring, where a lift exists. This eliminates the surviving class. The same criterion applies to a finite union of relatively affine pieces covering these valuation points, by using their product algebra and its idempotents.

For a proper \(Y\), the covering point \(h_S\to Y\) is relatively affine by the diagonal equations. It lifts valuation points by properness and separation. The same argument applied to a proper source over a proper target shows that surjectivity on geometric points gives dense image: first lift the generic point to the source, then use its covering point and separation in the target. This applies to \(u\). We also use the filtered continuity theorem for such an affine atlas. In particular cohomology commutes with passage to a directed union of base algebras and with directed intersections of relatively affine conditions. For this last assertion pull the terms of the inverse system back to the atlas and choose compatible test algebras through which all their test points factor. This is the test-envelope construction. Their representables give covers of the terms. The augmentation of each Čech complex reduces continuity in the next degree to continuity on these test representables and the preceding bounds on the fibre powers. On the representables, evaluation reduces the assertion to continuity of locally constant functions.\formalnote{17.13}{13}

\subsection{The calculation over one residue field}
Fix a maximal ideal \(q\subset P\), put \(K=\kappa(q)\), and let \(\gamma:D\to K\). Let \(J\) be either \(0\) or \(\mathfrak m\), and let \(B\subset K\) be an absolutely integrally closed subring containing \(\gamma(D)\). Define \(E_J(B)\) to consist of the places killing \(q\), finite on \(B\), and killing \(J\) on the base. Here a place has value \(c\) at \(t\) if \(f(n_t)=c f(d_t)\). Coprimality forces \(f(d_t)\) to be a unit in this case. The finite values form a ring of pairs \((t,c)\), and \(q\) having value zero makes the resulting value on \(B\) independent of the lift from \(P\). Evaluation gives \(E_J(B)\to h_{B/JB}\). The assertion needed is
\[
 H^0(E_J(B))=\operatorname{LC}(\operatorname{Spec}(B/JB),\mathbb F_\ell),
               \qquad H^n(E_J(B))=0\quad(n>0),                  \tag{PB2}
\]
where the degree-zero identification is pullback by evaluation.

When \(B\) is a valuation ring of \(K\), evaluation identifies \(E_J(B)\) itself on tests with \(h_{B/JB}\). Its inverse sends a map \(B/JB\to C\) to the projection to the place \((q,B)\) followed by that map. Over a local test the local representative above has valuation ring containing \(B\), so restriction to \(B\) recovers the original place. The equality holds for any test \(C\) because its maps to its localizations at maximal ideals are jointly injective. Also \(JB\) is radical. If \(x^2\in JB\), comparability of elements of \(D\) makes \(x^2\) divisible by \(\gamma(a)\) for a single \(a\in J\). Write \(a=d^2\) in \(D\). Then \(d\in J\) and divisibility of squares in the valuation ring gives \(\gamma(d)\mid x\) (with the zero case immediate). Thus \(B/JB\) is a reduced quotient of a test valuation ring and is a test. Evaluation on a representable is exact, proving (PB2) in this case.\formalnote{17.14}{14}

If \(B\) fails the valuation condition, choose \(t,t^{-1}\notin B\). Write \(B_1,B_2\) for the integral closures in \(K\) of \(B[t], B[t^{-1}]\), and \(B_3\) for the subring they generate. They again split monic polynomials. Moreover \(B_3=B_1[1/t]=B_2[1/t^{-1}]\), by localization of integral closure, and \(B_1\cap B_2=B\). For the last equality every valuation ring of \(K\) containing \(B\) contains \(B_1\) or \(B_2\), and their intersection over all such valuation rings is the integral closure of \(B\) in \(K\), namely \(B\). Indeed, if \(x\in K^\times\) is not integral over \(B\), then \(x^{-1}\) is a nonunit of \(B[x^{-1}]\), since an inverse there would give a monic equation for \(x\). Localize at a maximal ideal containing \(x^{-1}\) and choose a dominating valuation ring of \(K\) \cite[\href{https://stacks.math.columbia.edu/tag/00IA}{Tag 00IA}]{arith-stacks}. This ring excludes \(x\). The place functors for \(B_1,B_2\) cover \(E_J(B)\) cohomologically by the valuation lifting criterion. Their intersection is \(E_J(B_3)\). The choice of \(t\) can be made so that, writing \(C_i=\operatorname{LC}(\operatorname{Spec}(B_i/JB_i),\mathbb F_\ell)\) and \(C_B\) similarly, we have
\[
          0\longrightarrow C_B\longrightarrow C_1\oplus C_2
            \xrightarrow{\,f_1-f_2\,} C_3\longrightarrow0.        \tag{PB3}
\]
We give the algebra underlying both exactness assertions that involve gluing.

At the left the spectra cover: dominate \(B_{\mathfrak p}\) by a valuation ring of \(K\) for any \(\mathfrak p\supset JB\). Its centre on one of the two larger rings lies over \(\mathfrak p\) and contains \(J\). At the middle, compatible clopens lift to idempotents modulo \(J\) represented by \(x\in B_1,y\in B_2\). Their images in \(B_3/JB_3\) agree. Thus \(x(1-y)\) and \(y(1-x)\) are divisible in \(B_3\) by one \(\varpi\in\gamma(J)\), by total divisibility in \(D\). To make the two representatives equal before reduction, replace each \(v\) by \(p_n(v)=1-(1-v^n)^n\), which preserves its idempotent residue. The difference of the new representatives is divisible by \(\varpi^n\), since \(p_n(x)(1-p_n(y))\) is divisible by \((x(1-y))^n\), and likewise with \(x,y\) exchanged.

Here is why the extra powers suffice. Put the two representatives and the two quotients \(x(1-y)/\varpi,y(1-x)/\varpi\) into a finite algebra over \(B[t,t^{-1}]\) inside \(B_3\), using integrality. In the case \(\varpi=0\) the two representatives already agree. Choose finitely many module generators \(s\) for this algebra and a finitely generated coefficient subring \(R\subset B\) containing the coefficients of the finite multiplication table and expansions of the specified elements. The terms \(t^j s\) lie in \(B_1\) for \(j\) sufficiently positive and in \(B_2\) for \(j\) sufficiently negative, by the two localization descriptions. The intervening terms span a finite Noetherian \(R\)-module \(N\). Taking the interval large enough, the algebra generated over \(R\) by the specified elements lies in \(B_1+B_2+N\), by expanding products in this table. We can consequently write \(p_n(x)-p_n(y)=\varpi^n(w+z)\) with \(w\in B_1+B_2,z\in N\). The submodules \(\{z\in N:\varpi^kz\in B_1+B_2\}\) stabilize. Choose \(n=k+1\) after stabilization. Subtracting \(\varpi^nw\) shows \(\varpi^{k+1}z\in B_1+B_2\), hence \(\varpi^kz\in B_1+B_2\). The difference is therefore in \(\varpi(B_1+B_2)\). Correct the representatives by these multiples. Their common value lies in \(B_1\cap B_2=B\), and its basic open has the two original preimages. Applying this to the fibres of the locally constant functions glues them since the two spectra cover.\formalnote{17.15}{15}

For surjectivity at the right choose a primitive \(\ell\)-th root \(\zeta\in B\). Its powers are distinct in every residue field since \(\ell\) is a unit. Thus \(v\in C_3\) gives the element \(\zeta^v\) of \(B_3/JB_3\), using the idempotents of its fibres. Lift it to \(w\in B_3\), so \(w^\ell=1+\varpi G\) for some \(\varpi\in\gamma(J)\), \(G\in B_3\). It suffices to find \(W\equiv w\bmod JB_3\) with \(W^\ell=ab\) for elements of \(B_1,B_2\) respectively congruent to \(1\) modulo \(J\). Indeed write \(a=x^\ell,b=y^\ell\). Since \(B_3\) is a domain, \(W=\zeta^i xy\) for some \(i\) (the zero-quotient case is vacuous). The residues of \(x,y\) define locally constant exponents of \(\zeta\) on the two spectra. Their sum plus \(i\) is \(v\), and changing the sign of the second exponent gives the required difference.

For local \(B\), \(\varpi\) may be assumed a nonunit since otherwise \(B_3/JB_3=0\). Look for
\[
       (w(1+\varpi c))^\ell=(1+\varpi a_1)(1+\varpi a_2),
                \qquad a_i\in B_i,\ c\in B_3 .
\]
Use a finite algebra containing \(G\) over \(B[t,t^{-1}]\), with
module generators \(\mathcal S\). Over a finitely generated coefficient
subring of \(B\), let \(M[a,b]\) be the span of \(t^js\) for \(s\in\mathcal S\)
and \(a\le j\le b\). Choose \(k\) so that \(N=M[-k,k]\) contains
\(1,G\) and all products \(ss'\) of generators. Enlarge \(k\) so the
generators \(t^js\) belong to \(B_1\) for \(j\ge k\) and to \(B_2\)
for \(j\le-k\). The multiplication table gives
\[
 M[a,b]M[c,d]\subseteq M[a+c-k,b+d+k].
\]
Take \(a\le-k\), \(b\ge k\) with \(N^{\ell+1}\subseteq M[a,b]\).
With \(M=M[a,b]\), \(P_1=M[k,b]\), and \(P_2=M[a,-k]\), we have
\(M=P_1+P_2+N\) and \(P_1P_2\subseteq M\) by the same multiplication
bound.

Write a finite generating list of \(M\) as \(k_i=a_i'+b_i'+y_i\)
in these three parts. Use indeterminates \(X_i\), put
\(a_1=\sum X_i a_i'\), \(a_2=\sum X_i b_i'\), and
\(c=-(G+\sum X_i y_i)/\ell\). Expansion of the left side gives
\(1+\varpi(G+\ell c+\varpi\Phi)\), where \(\Phi\) has coefficients
in \(N^{\ell+1}\). The coefficients of \(\Phi-a_1a_2\) lie in \(M\),
so write it as \(\sum p_i(X)k_i\). The desired identity follows if
\(X_i=\varpi p_i(X)\). Modulo the maximal ideal this system has
the solution zero with identity Jacobian. The local absolutely
integrally closed ring is henselian since every monic polynomial
splits and each simple residue root is the reduction of one of
its roots. Inverting the Jacobian determinant makes the algebra defined by these equations étale over \(B\) near the zero residue solution. The lifting characterization of henselian local rings yields the \(X_i\) \cite[\href{https://stacks.math.columbia.edu/tag/04GG}{Tag 04GG}]{arith-stacks}.\formalnote{17.16}{16}

For nonlocal \(B\), choose \(d,1-d\) both nonunits such that either \(d(d-1)\in JB\) or every \(r\in JB\) admits \(\epsilon+r\epsilon^2=1\). Such a choice exists. If no such \(d\) existed, a root \(\rho\) of \(\rho^2-\rho=r\) for each \(r\in JB\) would have \(\rho\) or \(1-\rho\) a unit; its inverse solves the displayed equation. This would make every \(d\) with both \(d,1-d\) nonunits an allowed choice, contradicting nonlocality. Take \(t=1-d^{-1}\); neither \(t\) nor \(t^{-1}\) lies in \(B\). The rings \(B_1,B_2,B_3\) are now the localizations inverting \(d,d-1,d(d-1)\), respectively. If \(d(d-1)\in JB\) the overlap modulo \(J\) is empty. Otherwise write \(G=j/[d^m(d-1)^m]\), and choose \(c_0d^m+c_1(d-1)^m=1\). With \(U=d^{-m}\), \(V=(d-1)^{-m}\), we have \(c_1U+c_0V=UV\). Choose \(\epsilon+c_0c_1j\varpi\epsilon^2=1\). Direct multiplication gives
\[
       1+\varpi G=(1+\varpi\epsilon j c_1 U)(1+\varpi\epsilon j c_0 V).
\]
This proves the required factorization with \(W=w\). Thus (PB3) holds in both cases.\formalnote{17.17}{17}

We finish (PB2) by descending from the valuation rings. Both \(C_B\) and \(H^n(E_J(B))\) commute with directed unions of the allowed \(B\)'s. On the first side clopen partitions are described by finitely many idempotents and equations in \(B/JB\). On the second side the places finite on the union are the intersection of the relatively affine conditions of finiteness on its members, to which filtered continuity applies. Consequently a fixed failure of injectivity or surjectivity persists at a maximal \(B\) by Zorn's lemma: vanishing or the existence of a preimage at a chain union would already be witnessed at a stage.

In degree zero a kernel element dies on both larger charts at a maximal failure, so (PB3) kills it. Degree-zero surjectivity follows because the two lifts of the restricted class agree on the intersection by injectivity, and glue by (PB3) and Mayer--Vietoris. In positive degree use induction on the degree, followed by the same maximal failure argument. A class vanishes on the two charts and must be the boundary of a class on the intersection in the previous degree. That group vanishes in positive degree by induction. In degree zero all its elements are differences by (PB3), so their boundary is zero too. This proves (PB2).\formalnote{17.18}{18}

\subsection{From places to the closed fibre}
Take \(B=\gamma(D)\). Every place killing \(q\) is finite on this ring. For \(J=0\) its quotient is a domain, and for \(J=\mathfrak m\) the quotient is the residue field, since \(\ker\gamma\subseteq\mathfrak m\). Thus (PB2) computes both groups as \(\mathbb F_\ell\) in degree zero, with restriction the identity on constants, and zero in all other degrees. This proves the closed-fibre assertion on the places killing \(q\).

To remove \(q\), use the relatively affine loci \(E(I)\) of places killing an ideal \(I\subset P\). Sums correspond to intersections and killing \(1\) is possible only over the zero ring, where the coefficient vanishes. For an idempotent \(e\), the loci killing \(e\) or \(1-e\) cover cohomologically, using the local description of places. Their intersection contributes zero. If restriction failed, a kernel or cokernel class would again persist by continuity up to a maximal choice of equations \(I\). Splitting by \(e\) shows that \(I\) contains \(e\) or \(1-e\) for every idempotent. A proper ideal of this kind in \(P\) is maximal: given \(t\notin I\), choose \(y\) with \(yt^2=t\). The idempotent \(yt\) is outside \(I\), so \(1-yt\in I\), making \(t\) invertible in \(P/I\). The calculation for maximal \(q\) excludes the failure. The closed-fibre assertion now holds on \(E\).

Finally compare the Čech nerves of \(E\to Y\) and its restriction to \(Y_0\). Both maps cover cohomologically; density on the latter follows by pullback. The closed fibre of every fibre power is the fibre power of the restricted cover, since its coordinates have the same image in \(Y\). Each fibre power is proper over \(D\) by the diagonal and valuative conditions. Induct simultaneously for all proper \(Y\) through injectivity in degree zero, surjectivity in degree zero, injectivity in degree one, and so on. In comparing the augmented Čech complexes, the term \(E\) has the full result just proved. Each higher Čech column is shifted by at least one, so the previous induction bounds on the corresponding fibre powers give the next injectivity or surjectivity on the augmentation. In any fixed total degree only finitely many columns enter. This proves (PB1).\formalnote{17.19}{19}

\subsection{Extensions of fields and the vanishing ideal}
Let \(Y\) be proper over \(R\), with \(\ell\in R^\times\), and write \(Y_B=Y\times_{h_R}h_B\) for an \(R\)-algebra \(B\). Base change preserves the properness hypotheses. For any absolutely integrally closed valuation \(R\)-algebra \(V\), (PB1) implies in particular that \(H^n(Y_V)\to H^n(Y_{\kappa(V)})\) is injective: a structure map killing \(\mathfrak m_V\) factors through the residue field.

For algebraically closed fields \(k\subset K\) over \(R\), write \(K\) as the union of its finitely generated \(k\)-subalgebras \(B\). Filtered continuity gives \(H^n(Y_K)=\varinjlim_B H^n(Y_B)\). Each such \(B\) admits a \(k\)-point \(\beta:B\to k\) by the Nullstellensatz \cite[\href{https://stacks.math.columbia.edu/tag/00FV}{Tag 00FV}]{arith-stacks}. Its section shows injectivity of \(H^n(Y_k)\to H^n(Y_K)\): a class dying over \(K\) dies at some finite stage and then at the section. For surjectivity take a representative \(x\in H^n(Y_B)\). Dominate \(B_{\ker\beta}\) by a valuation ring \(V\) of \(K\) \cite[\href{https://stacks.math.columbia.edu/tag/00IA}{Tag 00IA}]{arith-stacks}. Its fraction field is algebraically closed, so \(V\) is absolutely integrally closed. The maps \(B\hookrightarrow V\) and \(B\xrightarrow{\beta}k\hookrightarrow V\) agree modulo \(\mathfrak m_V\), because \(b-\beta(b)\in\ker\beta\). Their pullbacks of \(x\) coincide on the residue fibre, hence on \(Y_V\) by injectivity. Restricting to \(Y_K\) proves that the given class comes from \(\beta^*x\) on \(Y_k\).\formalnote{17.20}{20}

Take a test \(T\) over \(R\) and \(x\in H^n(Y_T)\). Define
\[
 I_x=\{\,b\in T:\ x\text{ restricts to zero on the points where }b
                   \text{ is a unit}\,\}.
\]
This is an ideal. Closure under multiplication by elements of \(T\) follows by containment of unit loci. For addition choose \(e^2=e\) with \(b\mid ec,\ c\mid(1-e)b\). On the part \(e=1\), if \(b+c\) is a unit then \(b\) is a unit because it divides \(b+c\) there. On \(e=0\) the analogous fact holds for \(c\). The two disjoint idempotent parts compute the direct product on cohomology. Thus if \(b,c\in I_x\) then \(b+c\in I_x\). Zero belongs to \(I_x\) because the coefficient vanishes on the zero ring.

If \(x\) vanishes over \(\kappa(\mathfrak p)\) for \(\mathfrak p\in\operatorname{Spec}T\), it vanishes over the absolutely integrally closed valuation domain \(T_{\mathfrak p}\) by residue injectivity. Inverting all of \(T-\mathfrak p\) is the directed intersection of the conditions inverting finite subsets. By filtered continuity some finite subset already suffices. Its product is an element of \(I_x-\mathfrak p\). Conversely any \(b\in I_x-\mathfrak p\) becomes a unit in \(\kappa(\mathfrak p)\), so \(x\) vanishes there. Therefore
\[
 \{\mathfrak p:x|_{\kappa(\mathfrak p)}=0\}
                     =\bigcup_{b\in I_x}D(b).
\]
This proves openness. If the union is the whole spectrum, \(I_x=T\); taking \(b=1\) gives \(x=0\). Thus field restrictions detect higher-degree classes as required.\formalnote{17.21}{21}

\subsection{Closed coefficients on the marked tests}
Specialize now to \(A=\Omega\) and \(X=(\mathbb P^1)^d\) of \S\ref{sec:projective-test}, with its specified functor of rows. The product is proper in the precise sense above: equality is cut out by the determinants of the rows, and a single point is given by \((1+u_i v_i:v_i)_i\) over the polynomial ring in the \(u_i,v_i\). A row over a valued field can be scaled to \((1:c)\) or \((c:1)\) with \(c\) in the valuation ring. Set \((u_i,v_i)=(0,c)\) or \((c-1,1)\) respectively. Take a family of ideals \(I\) on \(X_A\) commuting with extension under every map of points, and let \(W\subset X_A\) be where \(I\) is zero. If \(a\in X_A(S)\) is the covering point and \(J=I(a)\), the reduction of \(a\) over \(S/J\) covers \(W\). Indeed the value of \(J\) is zero generically for a point of \(W\), hence already zero in the valuation ring used to lift it. The closed diagonal is inherited. Thus \(W\) is proper over \(A\).

Let \(g\) forget the mark on pairs \((T,x\in X(T))\) for tests over \(A\). For the ordinary coefficient \(L\) on marked tests take \(I=0\). For the coefficient \(F\) restricted to the zeros of \(I\), \(F(T,x)\) is the quotient of \(L(T)\) by sections which become zero along every further test killing \(I(T,x)\). On the zero locus the coefficient map \(L\to F\) is an isomorphism. The locus is dense for \(F\): on a representable marked test the sieve of maps killing \(I\) is represented by its quotient by \(\sqrt I\) with the induced mark. Radical quotients are tests. Hence higher Ext vanishes on each side by evaluation. In degree zero restriction has kernel exactly the sections divided out in the definition of \(F\), and is onto because locally constant functions on the closed quotient of a test lift. Indeed, lift an idempotent of the quotient to an element \(b\) of the test and compare \(b\) with \(1-b\) on an idempotent partition. On the part where \(b\mid 1-b\) it is a unit, so its idempotent residue equals one. On the other part that residue equals zero. The partition therefore lifts the given idempotent. Lifting the finitely many fibres of a locally constant function and successively taking their differences as clopen sets gives a partition lifting the function.\formalnote{17.22}{22}

For an \(A\)-test \(c\), these facts identify
\[
 \operatorname{Ext}^n_{\mathrm{marked\ tests}}
        \bigl(\mathbb Z[\operatorname{Hom}_A(c,g(-))],F\bigr)
             \ \simeq\ H^n(W_c).
\]
First restrict the Hom functor to the zero locus, using density pulled back to it. Its points are exactly a test \(T\), a point of \(X_A(T)\) killing \(I\), and a compatible \(c\to T\). These are the data of a point of \(W_c(T)\). Above each such point its already marked test with the identity on \(T\) is universal, so the resolutions by free representables compare by evaluation there, where \(F=L\). Both steps commute with a change \(c\to c'\). The just-proved openness, detection and field invariance of \(H^n(W_c)\) therefore hold for the marked Ext groups. The zero family gives the ordinary line formula's field input. The multihomogeneous ideal gives exactly the closed coefficient needed for the principal-open duality.\formalnote{17.23}{23}


\section{The line class and duality for a principal open}\label{sec:projective-test}

We prove the two projective-line formulas used in \S\ref{sec:primitive-comparison} and identify the relative pair required for an étale chart. Work over \(A=\Omega=\overline{\mathbb Q}\), with its chosen embedding in \(\mathbb C_p\). The two categories are ordinary tests over \(A\), and \(p\)-adic tests \(T\) with generic test \(T[1/p]\). Write \(\mathcal P\) for the functor of projective-line points on the relevant generic rings. An argument with parameters in \(\mathcal P\) takes place on pairs consisting of a test and a chosen point of its line. Taking Ext there, with the parameter forgotten in the first argument, is equivalent to taking Ext on the original category with the first argument tensored with \(\mathbb Z[\mathcal P]\). This is the free-group adjunction for the category of elements.

On ordinary tests use \(F=L\), the locally constant \(\mathbb F_p\)-valued functions. On \(p\)-adic tests use \(F(T)=T^\flat\otimes_O M\), where \(O=\mathcal O_{\mathbb C_p}^\flat\) and \(M\) satisfies the acyclicity on perfect algebras in (PC1). In the application \(M=N\otimes_O\mathfrak m\) with \(p^\flat N=0\). For an abelian-group-valued functor \(S\) write \(H^i(S)=\operatorname{Ext}^i(S,F)\). We will obtain
\[
 H^n(S)\oplus H^{n-2}(S)\ \xrightarrow{\ \sim\ }\
 H^n(S\otimes\mathbb Z[\mathcal P]),\qquad
 (a,b)\longmapsto \pi^*a+\eta\,\pi^*b.                            \tag{PL1}
\]
Negative degrees are zero. A chosen compatible system of roots of unity identifies the Tate coefficient acting in this formula.\formalnote{17.24}{24}

\subsection{A class supported on the diagonal}
In either category the Kummer sequence on tests has the form
\[
 0\longrightarrow\mathcal T\longrightarrow
 \operatorname{Hom}(\mathbb Q,\mathbb G_m)
 \xrightarrow{u\mapsto u(1)}\mathbb G_m\longrightarrow0 .
\]
Indeed each test has compatible roots of every unit: choose successive factorial roots and their powers. The kernel consists of systems trivial at \(1\). The chosen section \(\varepsilon:\mathbb Z\to\mathcal T\) maps \(1\) to the compatible roots of unity. On a coefficient killed by \(p\) the Tate action is normalized so that insertion of this generator acts as the identity. Choose a unimodular row for each projective point. A map of tests carries the chosen row to a unit multiple of the new one. These units obey the cocycle law and give the row extension by \(\mathbb G_m\). Its Yoneda product with the Kummer extension is \(\eta\in\operatorname{Ext}^2(\mathbb Z[\mathcal P],\mathcal T)\).\formalnote{17.25}{25}

For two sections \(q,r\), the determinant of their rows generates an ideal \(J\) independent of the choice of representatives, since changing either row multiplies the determinant by a unit. Where this ideal is the unit ideal the determinant itself trivializes the product of the two row cocycles. Thus \(\eta(q)+\eta(r)\) vanishes on that open part. The long exact sequence lifts this sum to a class relative to the open, supported on \(J=0\). On this locus write the two unimodular rows as \((a,b)\) and \((c,d)\). Choose \(h,k\) with \(ha+kb=1\). Since \(ad=bc\), the scalar \(u=hc+kd\) satisfies \(ua=c\) and \(ub=d\). As \(c,d\) generate the unit ideal, \(u\) is a unit. This proves that \(q=r\) on the zero locus over arbitrary test rings.

If \(x\) is a class on \(S\otimes\mathbb Z[\mathcal P]\), write \(x(q)\) for its evaluation at a section. Multiplication by the relative diagonal class only uses the restriction of \(x\) to the zero locus, where \(q=r\). The theories used here have excision on these closed loci. After forgetting support we obtain
\[
 x(q)\,[\eta(q)+\eta(r)]=x(r)\,[\eta(q)+\eta(r)].
\]
Suppose \(x(o)=0\) at the fixed point \(o=(1:1)\). Since \(\eta(o)=0\) by the constant row, taking \(r=o\) gives \(x(q)\eta(q)=0\). The preceding identity reduces to
\[
                       x(q)\eta(r)=x(r)\eta(q).                    \tag{PL2}
\]
It holds in particular on the category with two independent projective-line parameters. This is where the diagonal enters both (PL1) and the duality for opens.\formalnote{17.26}{26}

\subsection{Two discs in the tilt}
Put \(D=O[\mathbb Q_{\ge0}]\), \(U=O[\mathbb Q]\), using monomials \(z^s\). Denote the point functor of any such algebra \(b\) in the tilts by \(h_b\). In the category of abelian group functors on \(p\)-adic tests put \(W'=\mathbb Z[h_U\amalg h_O]\) and \(W=\mathbb Z[h_D\amalg h_D]\). Define \(i:W'\to W\) by sending a torus point \(t\) to its inverse in disc 0 minus itself in disc 1, and the point to the origin of disc 0. This map is objectwise injective. Indeed the restriction of a torus point to disc 1 determines its inverse powers as well; after looking at disc 1 all coefficients of the torus generators in a putative kernel element vanish. The origin of disc 0 then detects the remaining coefficient.

Under (PC1) these representable tilt functors have only degree-zero cohomology. The degree-zero map of \(i\), before tensoring with \(M\), is
\[
 D\oplus D\longrightarrow O\oplus U,\qquad
       (f,g)\longmapsto \bigl(f(0),f(z^{-1})-g(z)\bigr).
\]
Its inverse sends \((a,w)\) to \((a+w_{<0}(z^{-1}),a-w_{\ge0}(z))\), where the subscripts retain the terms of the indicated exponents of \(w\). Tensoring this explicit inverse with \(M\) still gives an inverse. Adding a parameter represented by a \(p\)-adic test \(c\) here amounts to replacing each of these algebras by \(c^\flat\otimes_O b\), since tilting on the test category is fully faithful. These algebras remain perfect, so the same proof gives an isomorphism in every degree. Resolving by sums of free representables extends the vanishing to \(H^n(S\otimes\operatorname{coker} i)=0\) for every objectwise flat \(S\). Flatness ensures that tensor with \(S\) preserves the exact sequence of \(i\).\formalnote{17.27}{27}

To recover the degree-two operation one must track the rational powers, even though the preceding linear map already has an inverse. Choose a character \(\chi:\mathbb Q\to O^\times\) whose untilts \(\chi(s)^\sharp\) are the chosen roots of unity \(\varepsilon(s)\). It is obtained by taking the compatible roots at \(s,s/p,s/p^2,\ldots\). Twist \(z^s\) by \(\chi(s)\) on disc 1 and on the torus, by \(\chi(-s)\) on disc 0, and trivially on the extra point. Let \(\delta,\delta'\) be one minus this twist on \(W,W'\) respectively. The inverse coordinate in the gluing gives \(\delta i=i\delta'\); both origins are fixed.

Form the pushout \(P=W'\amalg_{W'}W\) along \(\delta'\) and \(i\), and let \(\sigma:P\to W\) equal \(i\) on its \(W'\)-summand and \(\delta\) on its \(W\)-summand. Write \(K_0=\ker\sigma\), \(B_0=\operatorname{coker}\sigma\). The two-extension
\[
             0\to K_0\to P\xrightarrow{\sigma}W\to B_0\to0
\]
defines a class \(\omega\in\operatorname{Ext}^2(B_0,K_0)\). After any exact change of functors that makes \(\operatorname{coker}i\) acyclic for the coefficient, \(\sigma\) induces an isomorphism on cohomology. Indeed the inclusion \(W'\to P\) has the same cokernel as \(i\), and its composite with \(\sigma\) is \(i\). The two exact sequences through \(\operatorname{im}\sigma\) give an isomorphism from total Ext of \(K_0\) to total Ext of \(B_0\) by composition with \(\omega\), raising degree by two. To see this, the map from the cohomology of \(W\) to that of \(\operatorname{im}\sigma\) is injective, and the map from the latter to that of \(P\) is surjective, with their composite bijective. The middle group therefore splits as the image of the first map plus the kernel of the second. The two boundary maps identify the latter summand first with Ext of \(K_0\) shifted by one, then with Ext of \(B_0\) shifted the other way. This is ordinary Yoneda composition in the abelian category of functors \cite[\href{https://stacks.math.columbia.edu/tag/06XP}{Tag 06XP}]{arith-stacks}.\formalnote{17.28}{28}

Here are the maps that compare \(\omega\) with \(\eta\). For the untilt \(X\) of the coordinate on a disc, choose the frame \((X,1)\) on disc 1 and \((1,X)\) on disc 0. Send a disc generator in the row extension to the frame of \((1,0)\) minus this chosen frame. On the overlap the two frames \((X,1)\) and \((1,X^{-1})\) differ by the unit \(X\). It has a lift to the Kummer middle term, the system of untilts of \(z^s\) at the torus point. The twist leaves the ordinary frames fixed since \(\varepsilon(1)=1\). Its effect on the system of powers, by contrast, is multiplication by \(\varepsilon(s)\). With the sign in \(i\), the resulting difference on a torus generator equals the image under \(i\) followed by the map taking disc 1 to \(\varepsilon\) and disc 0 to zero. These identities give a morphism of the above two-extension to the splice defining \(\eta\): the frame difference on \(W\) induces \(j:B_0\to\mathbb Z[\mathcal P]\), and the Kummer lifts on the pushout induce \(u:K_0\to\mathcal T\). The origin of disc 1 gives \(y:\mathbb Z\to K_0\), since \(\delta\) kills that origin. By the computation just made \(u y=\varepsilon\) and \(j^*\eta=u_*\omega\).\formalnote{17.29}{29}

We now get (PL1) on the tilted tests. First, if \(\eta v=0\), restrict along \(j\) and use \(j^*\eta=u_*\omega\). Composition by \(\omega\) is injective, so the Tate action on \(v\) pulled back by \(u\) vanishes. Pull back once more by \(y\); the action of \(u y=\varepsilon\) is the identity, hence \(v=0\). Second, take \(x\) with \(x(o)=0\). By surjectivity of composition by \(\omega\), \(j^*x=\omega w\) for some \(w\). Its proposed coefficient of \(\eta\) is \(y^*w\). With a projective parameter \(r\), substitute \(\eta\,x(r)=x\,\eta(r)\) from (PL2) and pull this equality back on the other line along \(j\). The result is \(\omega\,u^*(\text{Tate action on }x(r))=\omega\, w\eta(r)\). Cancel \(\omega\), then pull back by \(y\), to obtain \(x(r)=(y^*w)\eta(r)\). This computation holds on the category with the universal point \(r\) and hence proves the equality of classes. Its acyclicity requirement there follows by adding the free parameter \(\mathbb Z[\mathcal P]\) to \(S\). Evaluation at \(o\) splits \(\pi^*\), giving (PL1). It suffices in this check to take \(S\) free representable; resolution extends it to all \(S\), since there are only two summands.\formalnote{17.30}{30}

\subsection{Locally constant coefficients}
For the ordinary tests the additional input is proper base change on the test category for the line. The precise consequences are invariance over extensions of algebraically closed fields, an open vanishing locus for the restriction of any class to residue fields of a test, and detection of zero by those restrictions. These imply that a homogeneous basis of cohomology on the point for the theory on pairs \((T,r\in\mathcal P(T))\) gives a Leray--Hirsch basis over the test fields above \(A\). Indeed choose classes with source \(L\) whose values on \(1\) give the chosen basis, using Zariski descent for this torsion coefficient. At a field \(L=\mathbb F_p\) on sections. Scalar extension there preserves the same basis by field invariance. To pass from all residue fields to their test, localize where the finitely many coefficients describing a given relation are constant. Fieldwise vanishing gives injectivity. For surjectivity the class locally has a prescribed finite linear combination of the basis, since the difference has open vanishing locus. On a quasicompact spectrum finitely many such combinations suffice. On an overlap their coefficients agree by injectivity, so the functions giving the coefficients glue. Thus the basis applies also at representables. The affine-cover resolution of the line makes it a basis at the line parameter itself.\formalnote{17.31}{31}

For (PL1) we require that the chosen basis contain \(\eta\). If \(\eta\cdot1\) vanished at a field test \(k\), field invariance and pullback through the structure point over \(k\) would force its vanishing at the generic field of \(A\) (here \(A\) itself). Every generic \(p\)-adic test receives this field map. The Tate action commutes with all maps of torsion coefficients, and pullback along \(1:\mathbb Z\to L\) generates the classes of the point. Thus such vanishing would make multiplication by \(\eta\) zero on the point for every tilted coefficient. Take the particular coefficient \(M=(O/(p^\flat))\otimes_O\mathfrak m\). Its point cohomology is \(M\) in degree zero by (PC1); \(p^\flat\bmod p^\flat\mathfrak m\) is nonzero in \(M\), because \(p^\flat=p^\flat a\) with \(a\in\mathfrak m\) would imply \(a=1\). The tilted (PL1) is injective on this summand times \(\eta\). Therefore the ordinary \(\eta\cdot1\) is nonzero and may be included in the basis.

Over a field test expand a class \(x\) vanishing at \(o\) in this basis. The coefficient at the index of \(\eta\) gives a class \(b\) of the base. On the product of two lines, (PL2) says that multiplying this expansion by the class of the second line is the expansion with \(x\) alone at the index of \(\eta\) in the first line. Uniqueness of coefficients with the additional line parameter implies \(b\eta=x\). The basis at the point gives injectivity of multiplication by \(\eta\); evaluation at \(o\) supplies the constant summand. Field detection and resolution as above prove the ordinary (PL1). Its proof uses (PC1) before applying any global primitive comparison, so the use of (PL1) to prove (PC2) is in this order.\formalnote{17.32}{32}

\subsection{The relative group for a principal open}
For \(f\in A[t_1,\ldots,t_d]\) embed \(D(f)\) in \(P=(\mathbb P^1)^d\) by \(t_i\mapsto (t_i:1)\). On rows \((a_i,b_i)\), put \(e_i=\deg_{t_i}f\) and use the principal ideal
\[
    I_f=\left(\left[\sum_{\alpha} f_\alpha
       \prod_i a_i^{\alpha_i}b_i^{e_i-\alpha_i}\right]\prod_i b_i\right).
\]
Rescaling a row rescales the generator by a unit. This ideal is the unit ideal exactly where every \(b_i\) and \(f(a_1/b_1,\ldots,a_d/b_d)\) are units. For \(f=0\) interpret the ideal as zero. Write \(Q_f=\mathbb Z[P]/\mathbb Z[D(f)]\), using the images as point functors, and call a functor good in the sense of \S\ref{sec:primitive-comparison}. We show that \(Q_f\) is good from goodness at the point.\formalnote{17.33}{33}

Use locally constant coefficients restricted to the zeros of \(I_f\) on tests marked by a point of \(P\). Proper field detection applies to this family of ideals on the proper product. The same basis construction as above gives homogeneous classes \(\theta_i\) with source the ordinary coefficient at the base and target the coefficient on the zeros. Their Leray--Hirsch map is bijective on representables. It is bijective on \(P\) and on the principal open \(D(f)\) by affine descent, hence on \(Q_f\) by the exact sequence of the pair. In particular it is surjective there. This supplies classes in ordinary coefficients; to transfer goodness to the comparison with the other theory takes a further argument.\formalnote{17.34}{34}

Let \(q,r\) be two independent \(d\)-tuples in \(P\). Start with the evaluation class at \(q\) obtained by reducing the constant class modulo \(p\) and restricting to the zeros of \(I_f(q)\). Move it to \(r\), replacing the coordinates one at a time. In coordinate \(j\) multiply by the relative diagonal lift of \(\eta(q_j)+\eta(r_j)\). On the zeros of its determinant the two current tuples are identical, so their ideals \(I_f\) agree. Locally where the new ideal is a unit, either the old ideal or the diagonal determinant generates the unit ideal: reduce to the residue field at a maximal ideal to see this, then localize. Excision consequently extends this product as a relative class at the new tuple. Its composite with any cohomology operation that gave the old evaluation gives the new evaluation times \(\eta(q_j)+\eta(r_j)\). Iterating produces a class with this universal property for the product over \(j\).

The surjection for the \(\theta_i\) just obtained expands the moved evaluation in a finite sum with coefficients \(\beta_i\) from the relative object back to the base coefficient. Given a class \(w\) on \(Q_f/p\), extend its restriction at \(q\) to the zero-locus coefficient; this is the extension furnished by excision, whose pullback recovers precisely the original restriction. Compose with each \(\theta_i\) to obtain an operation \(T_i(w)\) on the base coefficient with parameters \(q\). The preceding universal property identifies, at \(r\), the sum of these operations paired with \(\beta_i\) with
\[
             w(r)\prod_{j=1}^d\bigl(\eta(q_j)+\eta(r_j)\bigr).       \tag{PL3}
\]
Integrate out \(q_1,\ldots,q_d\), where integration by (PL1) sends \(a+\eta b\) to \(b\). At each stage \(\eta(r_j)\) is a base class and the coefficient of \(\eta(q_j)\) is \(1\). Thus integration recovers \(w(r)\). The left side gives a finite sum of operations through the good base coefficient. Pullback at the universal \(r\) determines the relative class, so these operations split the pullback from \(Q_f/p\) to \(Q_f\). For the \(p\)-torsion theories this pullback is surjective: \(Q_f\) is objectwise free, and multiplication by \(p\) acts as zero in the coefficient, so use the exact sequence of multiplication by \(p\) on \(Q_f\). The splitting by operations gives surjectivity and injectivity of the comparison on \(Q_f\) by naturality, and surjectivity of a finite sum of finite ordinary cohomology groups gives finiteness. This proves the required relative assertion. Applying it to \(g\) and \(gf\) and taking the sequence of a triple gives the pair \((D(g),D(gf))\) used for a standard étale chart.\formalnote{17.35}{35}

Section~\ref{sec:test-proper-base} establishes detection and proper base change on these tests, including the identification with coefficients restricted to the closed locus. The two-disc argument also uses the Tate action of the chosen roots of unity and the acyclicity on the tilted valuation tests from (PC1). These calculations take place before identifying \(J_Y\) in \S\ref{sec:primitive-comparison}.


\section{Primitive comparison on the test rings}\label{sec:primitive-comparison}

Let \(Y\) be a proper scheme over a ring that maps to \(\mathcal O_{\mathbb C_p}\). On its generic fibre, the coefficients for the cohomology are Witt vectors of locally constant functions. To study reduction, we evaluate \(Y\) on integral test rings and use the Witt vectors of their tilts. We compare these two cohomologies by first proving descent with coefficients in characteristic \(p\) and then passing to the complete period coefficients. The resulting isomorphism is used in \S\ref{sec:cohomological-frobenius}.

Fix a prime \(p\), write \(C=\mathbb C_p\), let \(O=\mathcal O_C^\flat\) with maximal ideal \(\mathfrak m\), and put \(A_{\inf}=W(O)\). Write \(p^\flat\) for the image modulo \(p\) of the chosen generator \(\xi\) of \(\ker(\theta:A_{\inf}\to\mathcal O_C)\). Let \(Y\) be proper over \(R\). Fix compatible maps \(R\to\mathcal O_C\) and \(R\to\Omega\hookrightarrow C\) with \(\Omega=\overline{\mathbb Q}\), and assume the existence of the charts of \(Y_\Omega\) used in \S\ref{sec:seed-proof}. Section~\ref{sec:cover-charts} constructs these charts for each model and generic cover to which this result is applied. In this section the integral and generic point functors will be denoted by \(X'\) and \(X\).

On an ordinary test ring \(A\), let \(X(A)=Y_\Omega(A)\), with the ring map \(\Omega\to A\) understood as part of the point, and let \(L(A)\) be the ring of locally constant \(\mathbb F_p\)-valued functions on \(\operatorname{Spec} A\). An ordinary test is reduced, splits monic polynomials, and permits a partition by complementary idempotents on which either of any two given elements divides the other. On a \(p\)-adic test \(T\), put \(X'(T)=Y(T)\). Such a \(T\) is a nonempty product of \(p\)-adically complete absolutely integrally closed valuation domains with \(p\ne0\), with a structure map from \(\mathcal O_C\). Its generic test is \(T[1/p]\). For a functor of points \(Z\) and coefficient functor \(F\), cohomology here means
\[
               H^i(Z,F)=\operatorname{Ext}^i_{\rm fun}(\mathbb Z[Z],F).
\]
Section~\ref{sec:test-foundations} proves existence of these tests with a prescribed centre, the tilt adjunction and the envelope property used below. In particular \(V^i(Y)=\mathbb Q_p\otimes_{\mathbb Z_p}H^i(X,W(L))\). Scholze's primitive comparison theorem gives an almost isomorphism \(H^i_{\mathrm{\acute et}}(Z,\mathbb F_p)\otimes_{\mathbb F_p}\mathcal O_C/p \to H^i_{\mathrm{\acute et}}(Z,\mathcal O_Z^+/p)\) for a proper smooth rigid space \(Z\) over a complete algebraically closed extension \(C/\mathbb Q_p\) \cite[Theorem~1.3, constant coefficients]{legacy-scholze}. The functors here use the displayed Ext groups, so we trace the comparison on these functors themselves.\formalnote{17.36}{36}

\subsection{From integral points to generic points}
The first map is given by restriction \(Y(T)\to Y(T[1/p])\), with the \(R\)-structure on the right fixed. It is bijective. Indeed \(T=\prod_j V_j\) injects into \(T[1/p]\), which in turn injects into \(\prod_j\operatorname{Frac}(V_j)\). For the second assertion write an element with a single denominator \(p^a\) and check its numerator at each factor. Since \(Y\) is separated, its diagonal is closed. Two maps agreeing after a jointly injective family of ring maps therefore agree, by pulling back the ideal of the diagonal.

Given a point over \(T[1/p]\), restrict it to each fraction field and apply the valuative criterion of properness to obtain points over \(V_j\) \cite[\href{https://stacks.math.columbia.edu/tag/0BX4}{Tag 0BX4}]{arith-stacks}. To assemble them, cover \(Y\) by finitely many affines and assign each cell to a member containing the image of its closed point. This affine open contains the whole image of \(\operatorname{Spec}V_j\), since an open containing the closed point of a local spectrum is the full spectrum. For each group of cells their maps to the affine open combine by the universal property of the product of rings. The finitely many groups correspond to complementary idempotents of \(T\), so these maps give a single \(T\)-point over \(R\). Its restriction equals the prescribed point by the joint injectivity above. This argument needs no uniform bound on denominators of coordinates chosen separately in the fraction fields.\formalnote{17.37}{37}

A point of \(Y_\Omega\) on an ordinary test includes its \(\Omega\)-algebra structure. Thus the category of tests with points can equivalently remember this structure first and evaluate \(Y\) second. More explicitly, above a point over \(A\) the test \(A\) with its specified map from \(\Omega\) and the identity on \(A\) is universal among lifts. This gives the same cohomology under restriction to tests over \(\Omega\). Combining this observation with the preceding bijection gives the point map \(\alpha\) for the comparison. For the coefficient map, an idempotent giving a clopen part of \(\operatorname{Spec}T[1/p]\) is the partition of the cells into those where it is 1 and those where it is 0. It belongs to \(T\) and gives the same partition in \(T^\flat\). A function with values in \(\mathbb F_p\) on these pieces gives an element of \(T^\flat\). Extending by \(O\)-linearity defines
\[
       u_T:O\otimes_{\mathbb F_p}L(T[1/p])\longrightarrow T^\flat .
\]
The natural map on cohomology first pulls back along \(\alpha\), then pushes the coefficients along \(u_T\).\formalnote{17.38}{38}

\subsection{The precise acyclicity needed on the tilt}
Write \(h_b(D)=\operatorname{Hom}_{O}(b,D)\) for a perfect \(O\)-algebra \(b\). If \(p^\flat N=0\), let \(F(D)=D\otimes_O N\otimes_O\mathfrak m\). The local assertion is
\[
  H^i\bigl(T\mapsto h_b(T^\flat),\,T\mapsto F(T^\flat)\bigr)
     =\begin{cases}F(b)&i=0,\\0&i>0.\end{cases}                      \tag{PC1}
\]
The \(\mathfrak m\) in this formula is a tensor factor. Thus (PC1) is an exact assertion about ordinary modules. We explain how it follows from the covering criterion on the tests.

Suppose \(b\to d\) is a map of perfect \(O\)-algebras through which every \(b\to T^\flat\) factors. On the category of perfect algebras put \(Q=\mathbb Z[h_b]/\mathbb Z[\operatorname{im}h_d]\). The obstruction to descent is a relative class of \(Q\). One proves by induction on degree that no such class survives, allowing an additional perfect \(b\)-algebra factor \(E\) in the coefficients. Choose a nonzero class in the first failing degree. Order the divisibility relations on \(E\) induced by maps \(E\to D\) under which this class survives. A chain has an upper bound: take the product of the rings realizing its relations and invert the idempotents of the tails. If the class died there, some localization multiplier equal to 1 on a tail would kill it in the product, contradicting its nonzero image at an index in that tail. Here one uses the Čech presentation of the relative class to clear the multiplier. Zorn's lemma gives a maximal relation.

For two images \(f,g\), the two perfected charts of the blow-up of
\((f,g)\) impose \(f\mid g\) and \(g\mid f\), respectively. The difference
sequence from the ring to the product of these charts and on to their
intersection stays exact after tensoring with a perfect test algebra
\(D\). To tensor further with \(N\), first tensor with the flat ideal
\(\mathfrak m\). The intersection algebra tensored with \(D\) is
perfect and therefore reduced. In any reduced \(O\)-algebra, the
\(O\)-torsion ideal is killed by \(\mathfrak m\). Indeed, if \(az=0\)
with \(0\ne a\in O\) and \(s\in\mathfrak m\), then \(a\mid s^j\) for
some \(j\) by the rank-one valuation. Thus \((sz)^j=0\), so \(sz=0\).
Since \(\mathfrak m^2=\mathfrak m\), tensoring with \(\mathfrak m\)
kills this torsion ideal. The tensor of the intersection algebra
with \(\mathfrak m\) is consequently torsion-free and flat over \(O\).
This preserves left injectivity on tensoring the sequence with \(N\).

The long exact sequence and vanishing in lower degrees now imply
that the class survives on one chart. Maximality implies that its
corresponding divisibility already held. All the images from \(E\)
are consequently comparable. Repeat this step in the receiving
ring. For the second map the kernel is prime: if the images of two
elements have zero product, they are comparable, so one has square
zero and vanishes in the reduced target. In the quotient domain,
saturate the image of \(E\) by its own nonzero denominators. The
resulting ring \(V\) is perfect, and any two of its elements have
comparable numerators after a common denominator is chosen. It is
a valuation domain. The composite factors through \(V\), so the
class still survives there.\formalnote{17.39}{39}

Here \(p^\flat\) in \(V\) is neither a unit nor zero. More generally if its principal ideal is radical, tensoring \(V\) with \(N\otimes\mathfrak m\) kills every element: for \(m\in\mathfrak m\), write \(m=s m'\) with \(m'\in\mathfrak m\) and \(p^\flat\mid s^k\) for some \(k\), using the rank-one valuation of \(O\). In \(V\), \(p^\flat\mid s\), so the factor kills \(N\). A surviving class excludes this case. Realize a map \(\lambda:V\to T^\flat\) whose kernel is contained in \(p^\flat V\) and for which \(V/\ker\lambda\to T^\flat\) is faithfully flat. The cell construction uses the map \(\mathcal O_C\otimes_{A_{\inf}}W(V)\to\kappa(V)\) given on Witt vectors by their zeroth residue coordinate. This map exists because \(\xi\bmod p=p^\flat\) has residue zero. The tensor product has regular \(p\), and the cell-existence lemma produces a mixed-characteristic valuation cell above the kernel. The induced map to its tilt is local. Its image is a valuation domain, so the injective local map from the image to the tilt is faithfully flat. If an element of \(V\) in the kernel divided \(p^\flat\), the nonzero \(p^\flat\) of the cell's tilt would vanish. Total divisibility gives \(\ker\lambda\subset p^\flat V\).

Passing to \(V/\ker\lambda\) therefore preserves the coefficients because \(p^\flat N=0\). On \(T^\flat\) the chosen point factors through \(d\), so the relative class becomes zero. The relative comparison with this extra scalar factor satisfies flat base change, as seen from its augmented Čech resolution by representables. Faithful flatness detects the vanishing already over \(V/\ker\lambda\), a contradiction. This proves the covering criterion. For (PC1) choose one covering point factoring all test points of \(b\); if there are none, take the zero algebra. Such a point comes from the universal test for points of \(\mathcal O_C\otimes_{A_{\inf}}W(b)\). For its tilt \(b'\), the point functor is represented on \(p\)-adic tests, by full faithfulness of tilting there, so its higher Ext vanishes by evaluation. The fibre powers of \(b'\) over \(b\) are represented by perfect tensor-product algebras. The Čech comparison for this cover, inducting on degree simultaneously for all \(b\), gives (PC1).\formalnote{17.40}{40}

\subsection{From a point to a proper charted scheme}
Put \(M=N\otimes_O\mathfrak m\). The desired mod-\(p\) assertion, in each degree, is the bijectivity of
\[
       M\otimes_{\mathbb F_p}H^i(X,L)
          \longrightarrow H^i(X',\,T^\flat\otimes_O M).             \tag{PC2}
\]
Choose any \(\mathbb F_p\)-basis of \(M\). The source of (PC2) is a direct sum of copies of \(H^i(X,L)\), and the map on the copy for a basis vector \(m\) is restriction followed by the coefficient map \(l\mapsto u_T(1\otimes l)\otimes m\). This agrees on pure tensors with the coefficient comparison by the affine-cover flat base-change property of \(X\). At the point over \(\Omega\), \(L(\Omega)=\mathbb F_p\). The tilted point functor is \(h_O(T^\flat)\), consisting of a single point on every test. Formula (PC1) gives precisely the basis assertion there. Its ordinary mod-\(p\) cohomology has dimension one in degree zero.

Call a free or relative point functor good when its ordinary total mod-\(p\) cohomology is finite and these basis maps together are bijective on total cohomology. Along a short exact sequence with good first term the middle term is good exactly when the quotient is, by the two long exact sequences and exactness of direct sum over a field. A good pair means goodness of \(\mathbb Z[Z]/\mathbb Z[U]\). The boundary functor away from the closed fibre of the base field \(\Omega\) can be ignored: the zero ideal extends to the unit ideal only in the zero ring, invisible by Zariski descent.\formalnote{17.41}{41}

Here is the geometric mechanism extending this calculation to a chart. The locally constant and tilted theories each have a projective-line formula with generators \(1\) and \(\eta=c_1(\mathcal O(1))\) (with its Tate coefficient). Integrating the \(\eta\)-component inverts multiplication by \(\eta\) on the pullback of a class from the base. It follows that goodness of the product of a pair with \((\mathbb P^1,\mathbb P^1-\{(0:1)\})\) descends to the pair: lift \(\eta\) relatively, multiply by this lift, then integrate. The composite is the identity. For a principal open \(D(f)\) in affine \(d\)-space, line duality proves goodness of the relative pair \(((\mathbb P^1)^d,D(f))\) by multihomogenizing \(f\). Section~\ref{sec:projective-test} proves the two line formulas and this relative assertion, using field detection for coefficients on the zero locus.\formalnote{17.42}{42}

To see the étale-chart step, take a standard étale algebra \((P[x]/(f))[1/g]\) over a polynomial algebra \(P\), with \(f\) monic and \(f'\) invertible where required. Write \(Z\) for its points. Over the principal open where \(g(x)\) is a unit, use pairs \((y,x)\) where \(f(y)=0\), \(g(y)\) is a unit and \(f(x)\) is associated to \(x-y\). Map such a pair first to \(x\) and second to \((y,(x-y:1))\in Z\times\mathbb P^1\). The loci where \(f(x)\), respectively the first projective coordinate, are units have identical inverse image. Near a maximal ideal containing \(f(x)\) the test lifts the prescribed residue root \(x\bmod\mathfrak m_T\). The divided difference \((f(x)-f(y))/(x-y)\), read as a polynomial, reduces to the unit \(f'(y)\). Localizing makes it a unit, giving the required association. Two such roots over the same \(x\) agree locally on the complement of the unit locus, by the same divided-difference argument at the two roots. For the second map, a point near the missing section has the coordinate \((t:1)\) with \(t\) in the maximal ideal. Taking \(x=y+t\) gives a local preimage, and both \(y\) and \(x\) are recovered on the image. Excision identifies the two relative pairs. On the \(x\)-side subtract the two relative pairs for \(g\) and \(gf\) in the exact sequence of a triple, obtaining goodness of \((D(g),D(gf))\). On the other side the preceding integration removes \(\mathbb P^1\). This gives goodness of \(Z\).\formalnote{17.43}{43}

Cover each affine torus chart by finitely many principal standard étale pieces. Their further principal localizations have the same form, so intersections are covered by the same argument. Mayer--Vietoris and the defining gluing property of the atlas give goodness on the charting subfunctor. It is locally all of \(X\), hence has the same cohomology. This proves (PC2) and finiteness in every degree of ordinary mod-\(p\) cohomology. Finiteness alone follows even by taking \(M=0\) in the joint goodness property. Affine covers also give flat base change on \(X'\), the other cohomology functor required below.\formalnote{17.44}{44}

\subsection{Removing the almost coefficient after Witt lifting}
Let \(S\) be a \(p\)-complete, \(p\)-torsion-free \(A_{\inf}\)-algebra with \(\xi\in pS\), and let \(K\) be a flat \(S\)-algebra in which \(\tau=[\epsilon^{1/p}]-1\) is invertible. In the application \(S=A_{\max}\) has the twisted structure map of \S\ref{sec:cohomological-frobenius}, and \(K=B_{\max}\) is a \(\mathbb Q_p\)-algebra. Take \(N=O\otimes_{A_{\inf}}S\). Multiplication by \(p^\flat\) on \(N\) is multiplication by the image of \(\xi\), which vanishes since \(pN=0\). Thus (PC2) applies. Put \(I_W=\ker(W(O)\to W(O/\mathfrak m))\). Zeroth coordinate gives \(I_W/pI_W\simeq\mathfrak m\). Indeed every element of \(\mathfrak m\) has a Teichmüller lift in \(I_W\), the kernel of zeroth coordinate in \(W(O)\) is \(pW(O)\), and \(I_W\) is \(p\)-saturated because \(p\) is regular in the Witt ring of the perfect residue field.

For a perfect \(O\)-algebra \(b\) with regular \(p^\flat\), form
\[
 P(b)=\left((S\otimes_{A_{\inf}} I_W)\otimes_{A_{\inf}}W(b)\right)^\wedge_p .
\]
Its reduction modulo \(p\) is \(b\otimes_O N\otimes_O\mathfrak m\): reduce the two Witt factors by zeroth coordinate and move scalars through the tensor product. Multiplication by \(p\) is injective before and after completion. For the first tensor the reduction of \(I_W\) is the flat \(O\)-module \(\mathfrak m\); for the second the reduction of \(W(b)\) is \(b\), torsion-free over \(O\) because \(p^\flat\) acts regularly and any nonzero scalar divides a power of \(p^\flat\). Both assertions use the tensor regularity criterion from reduction modulo a regular element. They apply to \(b=O\otimes L(A)\), which is free as an \(O\)-module, and to each test tilt \(T^\flat\). Consequently the reduction of the cohomology complexes with these \(P\)-coefficients is computed by the ordinary coefficient quotient. The complexes are derived \(p\)-complete: their coefficient values are complete, and applying derived \(\operatorname{Hom}(\mathbb Z[Z],-)\) commutes with derived inverse limits. Their comparison cone has zero derived reduction by (PC2). Derived Nakayama for the principal ideal \((p)\) \cite[\href{https://stacks.math.columbia.edu/tag/0G1U}{Tag 0G1U}]{arith-stacks} gives the integral comparison for \(P\). Flat base change then gives it for \(K\otimes_S P\).\formalnote{17.45}{45}

There is an explicit reason that the residual Witt ideal disappears. The maps
\[
  S\otimes_{A_{\inf}} I_W\ \underset{v}{\stackrel{u}{\rightleftarrows}}\ S,
 \qquad u(s\otimes w)=s w,\qquad v(s)=s\otimes\tau
\]
are defined since \(\epsilon^{1/p}\) has residue \(1\) and hence \(\tau\in I_W\). Both composites are multiplication by \(\tau\). Tensor with \(W(b)\), complete at \(p\), then extend to \(K\). Now \(\tau^{-1}v\) is a two-sided inverse. The comparison just proved applies to the actual coefficients
\[
              b\longmapsto K\otimes_S
                         \left(S\otimes_{A_{\inf}}W(b)\right)^\wedge_p .
\]
This argument also specifies the order of operations; the inverses only have to exist after passing to \(K\).\formalnote{17.46}{46}

It remains to identify the ordinary side of the isomorphism. Write \(H_i=H^i(X,W(L))\). The sequence of coefficient functors \(0\to W(L)\xrightarrow{p}W(L)\to L\to0\) embeds \(H_i/pH_i\) into the finite group \(H^i(X,L)\). Moreover \(H_i\) is derived \(p\)-complete by the same argument with Witt coefficients. Lifting a finite basis of \(H_i/pH_i\) generates \(H_i\): the cokernel of the map from a finite free \(\mathbb Z_p\)-module is derived complete and equals its own multiple by \(p\), so successive approximation (derived Nakayama for surjectivity) kills the cokernel. Thus \(H_i\) is finitely generated, allowing torsion.\formalnote{17.47}{47}

There is a coefficient map from \(S\otimes_{\mathbb Z_p}W(L(A))\) to the completion of \(S\otimes_{A_{\inf}}W(O\otimes_{\mathbb F_p}L(A))\). Modulo \(p\) it identifies \((S/p)\otimes_{\mathbb F_p}L(A)\) on both sides. For its source cohomology, \(S\) is flat over \(\mathbb Z_p\) by \(p\)-torsion-freeness, so flat base change identifies it with \(S\otimes_{\mathbb Z_p}H_i\). These modules are derived \(p\)-complete since the structure theorem for the finite \(\mathbb Z_p\)-module \(H_i\) decomposes them into finite sums of \(S\) and \(S/p^k S\). On the target, the complete \(p\)-torsion-free coefficients give derived completeness as above. The coefficient map is therefore an isomorphism on cohomology, by reduction modulo \(p\) and derived Nakayama. Extending scalars to \(K=B_{\max}\), composing with the comparison of \(b=O\otimes L\) and \(b=T^\flat\), and using \(K\otimes_{\mathbb Z_p}H_i=K\otimes_{\mathbb Q_p}V^i(Y)\), gives
\[
 J_Y:K\otimes_{\mathbb Q_p}V^i(Y)\ \xrightarrow{\ \sim\ }\
 H^i\!\left(X',\,T\longmapsto
 K\otimes_S(S\otimes_{A_{\inf}}W(T^\flat))^\wedge_p\right).
\]
On a pure coefficient tensor \(a\otimes w\) its map on the test \(T\) is \(a\otimes(1\otimes W(l\mapsto u_T(1\otimes l))(w))\), followed by the indicated completion in the inner factor. This makes pullback in \(Y\) transparent. Section~\ref{sec:cohomological-frobenius} uses exactly this \(J_Y\) and the analogous \(J_W\); the spanning assertion there uses in addition the integral root charts and the reduction map. Section~\ref{sec:projective-test} proves the projective-line and principal-open inputs from proper field detection and excision on the tests. Section~\ref{sec:test-proper-base} supplies field detection and the marked-test excision. The valuation argument above also uses the construction and tilting properties of the chosen cells.\formalnote{17.48}{48}


\section{The two root towers in crystalline comparison}\label{sec:crystalline-towers}

On an integral torus chart the coordinates are units. We adjoin compatible roots of a coordinate on each side of the comparison, first over the crystalline rings and then over the \(p\)-adic tests. The two calculations give the comparison map needed in \S\ref{sec:cohomological-frobenius}. Write \(A=A_{\max}\), \(B=B_{\max}\), \(O=\mathcal O_{\mathbb C_p}^{\flat}\), and use the Frobenius-twisted structure \(\iota:A_{\inf}\to A\) of that section. Put \(q=\iota([\epsilon]-1)\), so \(p^3\mid q\) in \(A\), \(q\) is a unit in \(B\), and \(p\mid\iota(\xi)\). For a \(p\)-complete, \(p\)-torsion-free \(A\)-algebra \(D\), put \(\bar D=D/\mathfrak d(D)\), where \(x\in\mathfrak d(D)\) if \(n!\mid x^n\) for all \(n\). We compare \(U(\bar D)\), with coefficients \(B\otimes_A D\), to \(U(S)\) on the \(p\)-adic tests \(S\) (products of valuation cells), with coefficients \(B\otimes_A D_S\), where \(D_S=(A\otimes_{A_{\inf}}W(S^\flat))^\wedge_p\). The map on Ext is pullback by \(S\to\bar D_S\) with the identified coefficients. Here \(U\) is regarded as a functor on \(\mathcal O_{\mathbb C_p}\)-algebras by restriction from the chart's base.\formalnote{17.49}{49}

\subsection{The crystalline tower for one unit}
Take \(t\in D\) whose image in \(\bar D\) is a unit. A compatible root system over \(f:D\to D'\) is a character of the additive monoid \(\mathbb Q_{\ge0}\) in the multiplicative monoid of \(\bar D'\), with value \(\overline{f(t)}\) at \(1\). It is represented by the completed divided-power envelope \(E\) of \(D[v^{\mathbb Q_{\ge0}}]\) along \(z=t-v\). Indeed a character with values modulo \(\mathfrak d\) lifts uniquely as a multiplicative character before reduction. To see this, \(x\mapsto x^p\) gives a map \(\bar D'\to D'/p\), since \(x^p\in pD'\) for \(x\in\mathfrak d(D')\). Composing it with the reduction \(D'/p\to\bar D'\) in either order is Frobenius. On maps from a monoid with unique \(p\)-th roots this identifies the two sets of characters. A compatible system in \(D'/p\) lifts multiplicatively to \(D'\): lift its \(p^j\)-th roots arbitrarily and take their \(p^j\)-th powers, which stabilize modulo every fixed power of \(p\) by the binomial theorem. The same estimate gives uniqueness. The lifted character sends \(z\) into the divided-power ideal. The images of \(z^n/n!\) are uniquely determined since the target is torsion-free over \(\mathbb Z\) (\(p\) is regular and integers prime to \(p\) are units). Adjoining these elements and completing gives the universal property claimed.\formalnote{17.50}{50}

Write \(u_r=v^r\), \(0\le r<1\) rational, and \(y_n=z^n/n!\). The uncompleted envelope has the \(D\)-basis \(e_{r,n}=u_r y_n\). Split an exponent into its fractional and integral parts. In each fractional part, changing the polynomial coordinate \(v\) to \(t-v\) is an involution. Over \(\mathbb Q\) one can then divide the \(n\)-th basis vector by \(n!\) without losing independence. Integrally the new vectors span the envelope because \(y_i y_j=\binom{i+j}{i}y_{i+j}\) and each original \(z^n\) is \(n!y_n\). In \(E\) the coefficients on this basis tend to zero \(p\)-adically, meaning only finitely many survive modulo \(p^k\) for each \(k\). Any assignment of an integral element of \(E\) to each basis vector gives a continuous \(D\)-linear operator. Also \(v\) is a unit of \(E\): it is congruent to \(t\) modulo \(\mathfrak d(E)\) and this ideal lies in the Jacobson radical because \(x^p\in pE\) for its elements.\formalnote{17.51}{51}

Use compatible primitive roots of unity \(\zeta_m\) of every order, extending the chosen system of \(p\)-power roots and satisfying \(\zeta_{mn}^{\,m}=\zeta_n\). For \(a=b/l\) let \(\epsilon^a\) be the tilt element represented by \((\zeta_l^b,\zeta_{lp}^b,\ldots)\). It has untilt \(\zeta_l^b\). The root twist \(\gamma\) acts on \(E\) by \(\gamma(v^a)=\iota([\epsilon^a])v^a\). In particular \(\gamma(v)=(1+q)v\). Two integral estimates allow us to compute its difference \(\Delta=\gamma-1\).

First, for \(r=m/l\) in lowest terms with \(0<r<1\), put \(\eta_r=\iota([\epsilon^r])-1\). There is \(c_r\in A\) with
\[
                              q=p\eta_r c_r.                       \tag{RT1}
\]
Set \(x=[\epsilon^{1/l}]\in A_{\inf}\). Its untilt \(\theta(x)=\zeta_l\) is a primitive \(l\)-th root. The geometric sum \(1+x+\cdots+x^{l-1}\) is divisible by \(\xi\) since \(\theta\) kills it, hence its image in \(A\) is divisible by \(p\). The sum of the first \(m\) terms is a unit after reduction \(A\to\bar A\). It comes there from \(1+\zeta_l+\cdots+\zeta_l^{m-1}\), a unit because a geometric sum with ratio \(\zeta_l^m\) and length \(b\), \(mb\equiv1\bmod l\), is its inverse. A unit modulo \(\mathfrak d(A)\) lifts to a unit in \(A\). Multiplying these two geometric sums separately by \(x-1\) gives \(x^l-1\) and \(x^m-1\), proving (RT1).\formalnote{17.52}{52}

Second, there are integral \(w_n\in E\) such that
\[
                \Delta y_{n+1}=-qv(y_n+p w_n).                   \tag{RT2}
\]
Indeed \(\gamma(z)=z-qv\). In the divided-power binomial expansion the \(k=1\) term is \(-qv y_n\). For \(k\ge2\) the term is \((-qv)^k y_{n+1-k}/k!\). After the factor \(qv\) is removed, \(q^{k-1}/k!\) is in \(pA\), since \(q\in p^3A\) and \(v_p(k!)\le k-1\). The assertion is an equality of integral elements as \(E\) embeds in its rationalization. Put \(L_n=y_n+p w_n\) and \(L_{-1}=0\), so \(\Delta y_n=-qv L_{n-1}\) also for \(n=0\).\formalnote{17.53}{53}

\subsection{A right inverse and the fixed summand}
Let \(m:D\to E\) be inclusion and \(\kappa:E\to D\) the coefficient at \((r,n)=(0,0)\), so \(\kappa m=1\). Define a continuous \(D\)-linear operator \(S_1\) by
\[
       S_1 y_n=-y_{n+1},\qquad S_1 e_{r,n}=p c_r v e_{r,n}\quad(r>0).
\]
On the integral-exponent part (RT2) gives \(\Delta S_1 y_n=qv(y_n+p w_n)\). On a fractional-exponent vector direct multiplication gives
\[
 \Delta S_1 e_{r,n}
     =qv\bigl(e_{r,n}
          +p c_r(1+\eta_r)u_r[\,y_n-(1+q)vL_{n-1}\,]\bigr).
\]
Indeed the coefficient of \(v u_r y_n\) from \(\gamma u_r-u_r\) is \(pc_r\eta_r=q\). The other terms contain the change of \(v\) or the change of \(y_n\). The error terms in parentheses define a single integral continuous endomorphism \(W\) with
\[
                          \Delta S_1=qv(1+pW).                    \tag{RT3}
\]
Since \(W\) preserves the \(p\)-adic filtration, \(\sum_{j\ge0}(-pW)^j\) inverts \(1+pW\) on \(E\). After tensoring over \(A\) with \(B\), multiplication by \(qv\) becomes invertible too. Thus \(\Delta S_1\) is bijective and \(S_1(\Delta S_1)^{-1}\) is a right inverse for \(\Delta\). The geometric series is formed integrally before scalar extension.

For the kernel put \(y_{-1}=0\) and define \(G_1:E\to E\) by
\[
      G_1 y_n=-q[(n+1)y_n-t y_{n-1}],\qquad
      G_1 e_{r,n}=-\eta_r e_{r,n}\quad(r>0).
\]
On integral exponents, \(v y_n=t y_n-(n+1)y_{n+1}\). On fractional exponents apply (RT1). These identities give on every basis vector, hence on \(E\),
\[
                  S_1G_1=m(qt\,\kappa)-qv\,\mathrm{id}.          \tag{RT4}
\]
For \(n=0,r=0\) the left side is \(q y_1=q(t-v)\), which accounts for the constant term. Over \(B\), write any fixed vector \(x\) as \(qv x'\). By (RT4), \(x=m(qt\,\kappa x')-S_1G_1x'\). The first term is fixed. Therefore \(\Delta S_1G_1x'=0\), and (RT3) gives \(G_1x'=0\). The fixed vector lies in the image of \(m\). Conversely that image is fixed, and \(\kappa\) splits \(m\) even after scalar extension. We have proved
\[
 0\longrightarrow B\otimes_A D\xrightarrow{\,m\,}B\otimes_A E
       \xrightarrow{\,\gamma-1\,} B\otimes_A E\longrightarrow0.  \tag{RT5}
\]
The constants here are the entire base \(B\otimes_A D\). Formula (RT1) keeps \(S_1,G_1,W\) integral simultaneously for all rational exponents. The only period inverted in this calculation is \(q\).\formalnote{17.54}{54}

\subsection{The ordinary \(p\)-adic tower}
Fix a \(p\)-adic test \(S\), a unit \(t\in S\), and set \(R'=S^\flat\). Choose a compatible system of roots of \(t\) using absolute integral closedness in the cells. Write \(P(R')\) for the ring of periodic sequences \(\mathbb N\to R'\). It is the union of the finite product algebras of sequences of period \(d!\). It is perfect, since taking \(p\)-th roots termwise preserves a period. With the shift \(\sigma f(k)=f(k+1)\), the functor of a test map out of \(S\) together with roots of its image of \(t\) is corepresented in the tilts by \(P(R')\). To verify this, over a target cell any second root at the index \(1/d!\) differs from the chosen one by a unique power of the primitive \(d!\)-th root. The map from \(\mathcal O_{\mathbb C_p}\) to a cell is injective, because every nonzero element divides a power of \(p\) and \(p\ne0\) there. Thus the exponents give compatible residue classes in \(\mathbb Z/d!\). These select exactly one factor on each finite product algebra. Conversely, any ring map to the domain constituting a target cell's tilt selects such a factor, by its idempotents. Applying this to all cells and using full faithfulness of tilting on test maps gives the asserted parametrization. Under it the root twist translates every selected exponent by one, hence acts by \(\sigma\).\formalnote{17.55}{55}

For \(f\in P(R')\) put \(T f(0)=0\), \(T f(k)=\sum_{j<k}f(j)\). If \(N\) is a period for \(f\), shifting its partial sums by \(pN\) adds \(p\sum_{j<N}f(j)=0\) in characteristic \(p\). Thus \(T f\) is periodic, and
\[
                 (\sigma-1)T f=f,\qquad
                 T(\sigma-1)f=f-f(0).
\]
This splits the exact sequence \(0\to R'\to P(R')\xrightarrow{\sigma-1}P(R')\to0\) on the underlying \(O\)-modules. It remains exact after tensoring with any \(O\)-module \(M\). For the particular almost coefficient of \S\ref{sec:primitive-comparison}, the Ext of the perfect-algebra point functor for \(P(R')\) is concentrated in degree zero by (PC1), where its value is \(P(R')\otimes_O M\). This is where acyclicity enters the ordinary calculation.\formalnote{17.56}{56}

Both exact sequences give descent by the same resolution. For a functor \(Z\) with an endomorphism \(\gamma\), take its action on the product of the test category with the one-object category \(B\mathbb N\) of the monoid \(\mathbb N\). Coefficients have trivial monoid action. In that category resolve \(\mathbb Z[Z]\) by two copies of the free group on the underlying \(Z\) times \(\mathbb N\), with translation in the \(\mathbb N\)-coordinate. The maps on generators are
\[
 [z,k]\longmapsto[\gamma z,k]-[z,k+1],\qquad
 [z,k]\longmapsto[\gamma^k z].
\]
The highest nonzero \(k\)-component proves injectivity of the first map on finite sums. The displayed relations move each \([z,k]\) to \([\gamma^k z,0]\), proving exactness at the middle and surjectivity at the last term. Each copy with its free \(\mathbb N\)-action has the Ext of the underlying functor on tests, by induction and restriction. Over a representable test point of the base that functor is acyclic on each side. On crystalline tests it is represented by \(E\), and on ordinary \(p\)-adic tests use (PC1) as above. Thus degree zero on the category with the action is the fixed group, degree one is the cokernel of \(\gamma-1\), and all higher degrees vanish. On crystalline tests (RT5) computes these groups. On ordinary tests the calculation first applies to the almost coefficient. The period-coefficient lifting of \S\ref{sec:primitive-comparison} gives the same conclusion for \(B\otimes_A D_S\). Its flat-base-change hypothesis holds on the root functor with action forgotten because its functor of tilted points is isomorphic to the functor of points of an \(\mathcal O_{\mathbb C_p}\)-algebra. The two induced terms in the resolution then give flat base change also with the action retained. Consequently projection from the roots with their action recovers the cohomology of the base in each theory. Resolution by representables gives the same result over a varying base functor.\formalnote{17.57}{57}

\subsection{The chart with all roots and the comparison map}
Write \(U=h_C\) on the chart, with units \(t_i\in C\) and an étale map \(R[X_1,\ldots,X_d]\to C\) sending \(X_i\) to \(t_i\). For \(G=\mathbb Q_{\ge0}^d\) set
\[
       C_\infty=\mathcal O_{\mathbb C_p}[G]\otimes_{R[X_1,\ldots,X_d]} C.
\]
Giving a point of \(C_\infty\) means giving a point of \(C\) and a system of compatible roots for every coordinate. The algebra \(C_\infty\) is étale over \(\mathcal O_{\mathbb C_p}[G]\). Lift it across the surjection \(A_{\inf}[G]\to\mathcal O_{\mathbb C_p}[G]\) to an étale algebra \(T\) \cite[\href{https://stacks.math.columbia.edu/tag/04D1}{Tag 04D1}]{arith-stacks} and put \(b=T/pT\). Its base is \(O[G]\), which is perfect since multiplication by \(p\) is bijective on \(G\). Étaleness gives perfection of \(b\) by the relative Frobenius isomorphism \cite[\href{https://stacks.math.columbia.edu/tag/0EBS}{Tag 0EBS}]{arith-stacks}. Also \(b\) is flat over \(O\), so \(p^\flat=\xi\bmod p\) acts regularly on it.\formalnote{17.58}{58}

Make \(W(b)\) an \(A_{\inf}[G]\)-algebra by lifting the monomials by Teichmüller representatives. Formal étaleness lifts \(T\to b\) successively to \(W(b)/p^n\), hence to \(W(b)\). The resulting map induces an isomorphism of completions. Modulo \(p\) it is the identity of \(b\). Successive approximation gives surjectivity modulo \(p^n\), and \(p\)-regularity of \(W(b)\) lifts injectivity modulo \(p\) one power at a time. A point of \(C_\infty\) over \(\bar D\) is therefore equivalently an \(A_{\inf}\)-map \(W(b)\to\bar D\). This factors through the zeroth Witt coefficient in characteristic \(p\), and has a unique lift to \(D\). More explicitly the map \(b\to\bar D\) gives \(b\to D/p\) by first taking a \(p\)-th root in \(b\) and then applying \(x\mapsto x^p\bmod p\) on \(\bar D\). Lift to the inverse Frobenius limit using perfection of \(b\), take Witt vectors, and compose with Fontaine's \(\theta\) for \(D\). Uniqueness follows on Teichmüller representatives from multiplicative lifting and then on all Witt vectors by \(p\)-adic completeness. Thus the points on crystalline tests are represented by \(D_b=(A\otimes_{A_{\inf}}W(b))^\wedge_p\). On the \(p\)-adic tests the corresponding functor is \(\operatorname{Hom}_O(b,S^\flat)\), by the Witt-vector description of the tilts. The comparison sends \(1\otimes w\) to the corresponding period on the ordinary test. Evaluation on the first side and (PC1) with the period-coefficient lifting on the second identify both groups with \(B\otimes_A D_b\) in degree zero and zero above it.\formalnote{17.59}{59}

Now remove the root systems one coordinate at a time. Suppose comparison holds with an additional root system retained and with its action forgotten. It holds on the two induced terms in the monoid resolution, hence on the root functor with its action by the long exact sequence. The fibre calculations on both sides identify the resulting comparison with the one on the functor before this root system was adjoined. Naturality holds because reduction \(S\to\bar D_S\) commutes with the chosen twists. Starting at \(C_\infty\) proves comparison on \(U\). The defining cohomological gluing of the chart atlas then proves it on its charting subfunctor.\formalnote{17.60}{60}

\subsection{The scalar extension yielding spanning}
Consider the final change from \(A\)-tests to \(W_0=W(k)\)-tests of \S\ref{sec:cohomological-frobenius}. Write \(H^n_0\) for cohomology of the reduced chart functor with coefficient \(D\mapsto D\) on \(W_0\)-tests, and \(H^n_L\) for its cohomology with coefficient \(D\mapsto(A\otimes_{W_0}D)^\wedge_p\). The map to the completion induces
\[
                      j_n:A\otimes_{W_0}H^n_0\longrightarrow H^n_L .
\]
Flat base change for the chart gives the displayed source. The uncompleted and completed coefficient algebras are \(p\)-torsion-free and have the same reduction modulo \(p\). The long exact sequences of multiplication by \(p\) make \(\ker j_n\) \(p\)-divisible. Take any element in the kernel, lift it to a \(p\)-division on the source by equality of the reductions, and use the common connecting map in degree \(n-1\) to arrange that the division too lies in the kernel.

The topological \(W_0\)-basis of \(A\) supplies separating continuous \(W_0\)-linear coordinates \(\lambda_i:A\to W_0\). It can be obtained by lifting a basis modulo \(p\), expanding successively by powers of \(p\), and using \(p\)-regularity and completeness for uniqueness. Applying \(\lambda_i\otimes1\) to \(\ker j_n\) gives a \(p\)-divisible submodule of \(H^n_0\), hence zero: the latter group is derived \(p\)-complete since its coefficient values are so. Indeed an element admitting indefinitely compatible divisions by \(p\) supplies a solution to the homogeneous equations \(x_j-px_{j+1}=0\), whose only solution in a derived complete module is zero. It follows that a \(W_0\)-linearly independent family \(x_d\) in \(H^n_0\) gives an \(A\)-linearly independent family \(j_n(1\otimes x_d)\). Each coordinate of a finite relation on the latter is a \(W_0\)-linear relation on the former. Flat extension \(A\to B\) preserves this independence. The comparison just proved, and the adjunction between completed scalar extension and restriction, identify \(B\otimes_A H^n_L\) with the period cohomology of the charting subfunctor. In the proper case of the Frobenius argument this is a finite \(B\)-module by primitive comparison. Thus \(H^n_0\) has finite rank over the complete discrete valuation ring \(W_0\).

Apply the finite-rank lemma below with \(V=W_0\), \(\pi=p\), \(X=H^n_0\), and \(S=A\). It gives derived \(p\)-completeness of the source of \(j_n\), so the \(p\)-divisible kernel vanishes in every degree. The common reduction and injectivity in degree \(n+1\) now give surjectivity modulo \(p\) in degree \(n\). To see that the cokernel itself vanishes, start with \(v_0\in H^n_L\) and choose \(v_{j+1}\) and source elements \(a_j\) successively so that \(v_j=j_n(a_j)+p v_{j+1}\). Solve \(u_j-pu_{j+1}=a_j\) in the source. Then \(v_j-j_n(u_j)\) satisfies the homogeneous equations in the derived complete target, whence \(v_0=j_n(u_0)\). Thus \(j_n\) is an isomorphism before and after further base change to \(B\). The triangular identity of the scalar extension adjunction sends \(1\otimes1\otimes x\) to \({\rm red}_0^*x\), proving the spanning used there.\formalnote{17.61}{61}

\subsection{Finite rank and derived completion}
For a \(V\)-module \(M\), put \(d_M((x_j))=(x_j-\pi x_{j+1})\) on \(P=M^{\mathbb N}\). The kernel and cokernel fit into the exact sequence
\[
0\longrightarrow\operatorname{Hom}_V(V[1/\pi],M)
\longrightarrow P\xrightarrow{d_M}P
\longrightarrow\operatorname{Ext}^1_V(V[1/\pi],M)\longrightarrow 0 .
\]
Indeed, on the free module with basis \(e_j\), \(j\geq0\), the map \(e_j\mapsto e_j-\pi e_{j+1}\) is injective by the least nonzero coefficient. Its cokernel is the direct limit of repeated multiplication by \(\pi\) on \(V\), namely \(V[1/\pi]\), via \(e_j\mapsto\pi^{-j}\). Applying \(\operatorname{Hom}_V(-,M)\) gives the sequence and vanishing of higher Ext. Derived \(\pi\)-completeness means that both end groups vanish, equivalently that \(d_M\) is bijective \cite[\href{https://stacks.math.columbia.edu/tag/091N}{Tag 091N}]{arith-stacks}.

\textit{Finite-rank lemma.} Let \(V\) be a complete discrete valuation ring with uniformizer \(\pi\), \(X\) a derived \(\pi\)-complete \(V\)-module of finite rank, and \(S\) a derived \(\pi\)-complete \(V\)-algebra. Then the ordinary tensor product \(S\otimes_V X\) is derived \(\pi\)-complete.

Choose a maximal independent family in \(X\). It is finite, so its span \(L\) is the image of an injection \(i:F=V^d\to X\). Each \(x\in X\) has some \(\pi^n x\in L\), by maximality and the decomposition of a nonzero scalar as a unit times a power of \(\pi\). We first prove a uniform bound on \(n\). There exist \(m\geq0\) and \(r\in X\) such that
\[
  \forall y\in X\quad\exists z\in X,\qquad
          \pi^m\bigl(r+\pi^m(y+\pi z)\bigr)\in L.
\]
Otherwise start with \(r_0=0\). At step \(j\), choose \(y_j\) violating the displayed condition for \(m=j,r=r_j\) and every \(z\), and put \(r_{j+1}=r_j+\pi^j y_j\). Solve \(u_j-\pi u_{j+1}=y_j\) using surjectivity of \(d_X\). Telescoping gives \(u_0=r_j+\pi^j(y_j+\pi u_{j+1})\). Some \(\pi^n u_0\) belongs to \(L\). Taking \(j=n\) and \(z=u_{n+1}\) contradicts the choice of \(y_n\).

Subtract the displayed inclusions for \(y\) and for \(0\). With \(N=2m\) they give \(\pi^N X\subset L+\pi(\pi^N X)\). For each \(h_0\in\pi^N X\) choose recursively \(f_j\in F\), \(h_{j+1}\in\pi^N X\) with \(h_j=i(f_j)+\pi h_{j+1}\). In the finite free module \(F\), the convergent series \(g_j=\sum_{k\geq0}\pi^k f_{j+k}\) satisfies \(g_j-\pi g_{j+1}=f_j\). Hence \(d_X((h_j-i(g_j))_j)=0\). Injectivity gives \(h_0=i(g_0)\), proving \(\pi^N X\subset L\). In particular \(i\) has zero kernel and its cokernel is killed by the single power \(\pi^N\).

Since \(i:F\to L\) is an isomorphism onto its image, the bound defines a linear map \(b:X\to F\) by \(i(b(x))=\pi^N x\). After extension to \(S\), write \(a=i_S:F_S\to X_S\) and \(b_S:X_S\to F_S\), where \(F_S=S^d\) and \(X_S=S\otimes_V X\). Then \(a b_S=\pi^N\). For a sequence \(v_j\in X_S\), solve \(u_j-\pi u_{j+1}=b_S(v_j)\) coordinatewise in \(F_S\). Direct subtraction shows that
\[
    x_j=\sum_{k=0}^{N-1}\pi^k v_{j+k}+a(u_{j+N})
    \qquad\text{satisfies}\qquad x_j-\pi x_{j+1}=v_j.
\]
If \(d_{X_S}x=0\), then \(d_{F_S}(b_S x)=0\), so \(b_S x=0\) and \(x_j=\pi^N x_{j+N}=a b_S(x_{j+N})=0\). This proves both required vanishings. The only completeness inputs in this scalar-extension argument are ordinary completeness of \(V\) to sum in \(F\), and the sequence criterion for \(X\) and \(S\). It applies to arbitrary torsion in the finite-rank hypothesis and requires no flatness of \(S\).\formalnote{17.62}{62}

The local computation uses the twisted period-ring properties stated at the beginning, full faithfulness on the chosen valuation tests, and the almost acyclicity and period lift in \S\ref{sec:primitive-comparison}. Chart gluing and flat base change enter when passing from one representation by a ring to the atlas. Properness enters the final finite-rank argument through primitive comparison on that atlas. Section~\ref{sec:cohomological-frobenius} uses Fontaine's fixed-period theorem to define the linear Frobenius on this space.


\section{Hodge classes and vectors of prescribed weight}
\label{sec:hodge-vector}

In the odd-rank attachment argument of \S\ref{sec:seed-proof}, a combination of Hecke correspondences is represented by maps \(u:W\to Y\) from a common cover. A Hodge class on \(Y\) can be pulled back along each of these maps. We associate to this class a vector of prescribed Galois weight in the period space of \(Y\), in a way compatible with every one of those pullbacks. The compatibility is what lets us use the same linear combination on periods.

Let \(R\to E\to\Omega=\overline{\mathbb Q}\) be the base maps, with \(E\) a field. Fix a place ring \(A\subset\Omega\) containing the image of \(R\), with place embedding \(e:\Omega\hookrightarrow\overline{\mathbb Q}_p\) and compatible local action \(j:\Gamma\to G_{\mathbb Q_p}\), where \(\Gamma\) is its decomposition group over \(E\). Thus \(e(gx)=j(g)e(x)\). Let \(E\to k\overset{\iota}{\hookrightarrow}\overline{\mathbb Q}_p\) agree on \(E\) with \(e\), and choose \(\varepsilon:\overline{\mathbb Q}_p\to k\) with \(\iota\varepsilon=1\). In the application \(k\) is a copy of the algebraically closed coefficient field with this chosen identification. Assume \(Y/R\) proper, smooth of relative dimension \(d\), and charted over \(\Omega\) for primitive comparison. Assume \(W/R\) proper. Let \(q,h\ge0\), \(n=q+h\). For \(x\) in the degree-\(q\) Hodge cohomology of \(Y_E\), we construct
\[
 v\in\mathbb C_p\otimes_{\mathbb Q_p}V^{n}(Y),\qquad
 \Theta_W:\mathbb C_p\otimes_{\mathbb Q_p}V^{n}(W)
                   \longrightarrow H_{\rm pair}^{n}(W_k),
\]
where \(\Theta_W\) is additive and \(\Theta_W(az)=\varepsilon(a)\Theta_W(z)\) for \(a\in\overline{\mathbb Q}_p\). With \(\rho\) the action on \(V^n(Y)\) and \(\chi\) the cyclotomic character, they satisfy
\[
  (j(g)\otimes\rho(g))v=\chi(j(g))^{-h}v,\qquad
  \Theta_W(u^*v)=\lambda_h\big((u^*x)_k\big)
            \quad\text{for every }u:W\to Y\text{ over }R .             \tag{HV1}
\]
The weight convention in (HV1) acts semilinearly on the tensor product. The target of \(\Theta_W\) is a cohomology group.\formalnote{17.63}{63}

\subsection{The square-zero extension and the contraction}

A \(p\)-adic test has ring \(S\), a nonempty product of the valuation rings called cells, with structure map \(\mathcal O_{\mathbb C_p}\to S\). Set \(T=S[1/p]\), \(D=W(S^\flat)[1/p]\) and \(I=\ker(\theta:D\to T)\). Put \(B_S=D/I^2\). There is a natural exact sequence of additive groups
\[
 0\longrightarrow T\xrightarrow{\iota_S} B_S\xrightarrow{\pi}T\longrightarrow0,
 \qquad \iota_S(\theta b)=\mu b\bmod I^2 .
\]
Indeed \(I=(\xi)\) with \(\xi\) regular, and \(\mu=\xi\tau\) with \(\theta(\tau)\) a unit in \(T\). Multiplication by \(\xi\) identifies \(D/I\) with \(I/I^2\). Replacing \(\xi\) by \(\mu\) composes that isomorphism with multiplication by the unit \(\theta(\tau)\). The product of two elements of this kernel is zero. The coefficient embedding \(\overline{\mathbb Q}_p\to T\) lifts uniquely to \(B_S\), since the algebraic extension of \(\mathbb Q_p\) is separable and every finite subextension is formally étale. The unique lifts on finite subextensions agree.\formalnote{17.64}{64}

An ordinary Hodge test is a \(k\)-algebra \(U\) and a standard finite free module \(P\). Its point algebra is \(U\oplus P\) with \(P^2=0\), and its coefficient functor is \(\bigwedge_U P\), the full exterior algebra. The class \(x_k\) is an Ext class from the free abelian functor of these points of \(Y_k\) to this coefficient functor. For the comparison use pairs \((S,M)\) with \(M=T^r\), the scalar coefficient \(T\), and point algebra
\[
                   B_{S,M}=B_S\times_T(T\oplus M).
\]
The degree-\(n\) Ext of the free abelian functor of \(W_k(B_{S,M})\) against these scalar coefficients is \(H^n_{\rm pair}(W_k)\). A morphism of pairs over a test morphism \(f:S\to S'\) has a semilinear part \((a,b):M\to T'\oplus M'\). It sends \((z,(t,m))\) in the fibre product to \((f_B(z)+\iota_{S'}(a(m)),(f_T(t),b(m)))\). The first component is multiplicative by square-zero multiplication and the linearity of \(a\) over \(f_T\).\formalnote{17.65}{65}

The arrow acts on exterior forms by \(\bigwedge b\). Associate to it the contraction
\[
 D_{a,b}(m_1\wedge\cdots\wedge m_r)
     =\sum_{\nu=1}^r(-1)^{\nu-1}a(m_\nu)
       b(m_1)\wedge\cdots\wedge\widehat{b(m_\nu)}\wedge\cdots\wedge b(m_r).
\]
For consecutive arrows \(f,g\), applying this formula to their composite gives
\(D_{gf}=(\bigwedge b_g)D_f+D_g(\bigwedge b_f)\): the omitted slot can be evaluated in the scalar component at the first arrow or at the second. Thus \(D\) represents a degree-one extension of the exterior coefficient functor by itself. Pull \(x_k\) along \(B_{S,M}\to T\oplus M\), compose with the \(h\)-fold Yoneda power of \([D]\), and project exterior forms onto their degree-zero part. This defines \(\lambda_h(x_k)\) in degree \(q+h\). It commutes with pullback in \(Y\) by associativity of composition. On exterior degree \(r\), scaling \(M\) by \(a_0\in k\) scales by \(a_0^r\). Only degree \(h\) can pass through \(h\) contractions and survive projection to zero. Equivalently, commuting the scaling through each contraction extracts one \(a_0\). The resulting class is multiplied by \(a_0^h\). For the classical Hodge--Tate comparison, Bhatt constructs the degree-one map on differentials from the completed cotangent complex and obtains the higher maps by exterior powers \cite[Lecture~3, \S3.4]{legacy-bhatt}. Here the \(h\)-fold Yoneda product above gives the scaling law on pairs.\formalnote{17.66}{66}

\subsection{The universal lift on a chart}

Let \(r_Z\) denote restriction from the cohomology of the points \(Z_k(T)\) on tests with coefficients \(T\) to \(H_{\rm pair}^{n}(Z_k)\), using \(B_{S,M}\to T\). We need \(r_Y\) bijective degree by degree. Smoothness provides étale coordinate charts. On one such chart, fix a \(T\)-point and lift each of its \(d\) coordinate values \(s_i\) to \(z_i\in B_S\). On \((S,T^d)\) prescribe the values \((z_i,(s_i,e_i))\), where \(e_i\) is the \(i\)-th basis vector. The projection of \(B_{S,T^d}\) to \(T\) has square-zero kernel. Formal étaleness over the coordinate polynomial ring therefore gives a unique lift of the point with these coordinates \cite[\href{https://stacks.math.columbia.edu/tag/00UP}{Tag 00UP}]{arith-stacks}.\formalnote{17.67}{67}

This lift is universal over the \(T\)-point with coordinates \((s_1,\ldots,s_d)\). Given any lift on \((S',M')\) over \(f:S\to S'\), write its \(i\)-th lifted coordinate as \((z'_i,(f_T(s_i),m_i))\). There is exactly one \(t_i\in T'\) with \(\iota_{S'}(t_i)=z'_i-f_B(z_i)\). Send \(e_i\) to \((t_i,m_i)\) to obtain the required and unique arrow of pairs. Formal étaleness gives agreement on the entire chart. Thus over the chart point the lift functor is represented by this one pair and its coefficient value is \(T\), the original coefficient value. For free abelian representable functors, taking Hom into the coefficients is evaluation and is exact. Applying this to a resolution by sums of representables gives the comparison in all degrees over chart points. Finally the chart loci cover locally on tests and on pairs. Their lifts remain in the chart because a square-zero thickening has the same open subsets. The scalar coefficient respects idempotent partitions of \(S\) and has the Zariski descent used for the tests, so restriction to these covering loci preserves cohomology on both sides. This proves the asserted bijectivity. There is consequently a unique test class \(\eta\) with\formalnote{17.68}{68}
\[
                    r_Y(\eta)=\lambda_h(x_k).                          \tag{HV2}
\]

\subsection{Recovering the vector and its character}

Let \(B_Z\) denote pullback from \(Z_k(S[1/p])\) to \(Z(S)\) in the corresponding Ext groups of tests with coefficient \(T\). It is an isomorphism for proper \(Z/R\). Indeed a generic test point extends uniquely on each valuation cell by the valuative criterion for properness \cite[\href{https://stacks.math.columbia.edu/tag/0BX4}{Tag 0BX4}]{arith-stacks}. A finite affine cover groups the cells according to a chart containing their closed points and assembles the local maps on the corresponding idempotent factors of the product. Separatedness tests the resulting equality in the fraction fields of all the cells, which jointly detect equality on \(S[1/p]\).\formalnote{17.69}{69}

The primitive period map \(J_Z\) of \S\ref{sec:primitive-comparison}, here with coefficient pair \((\mathcal O_{\mathbb C_p},\mathbb C_p)\), has the same target as \(B_Z\), is \(\mathbb C_p\)-linear, and commutes with pullbacks. Its coefficient map on a test is given by Fontaine's \(\theta\) to \(S[1/p]\). It is bijective on proper \(Z\) charted over \(\Omega\). Hence define \(v=J_Y^{-1}B_Y\eta\). For \(W\) use the forward map only, and put \(\Theta_W=r_W B_W^{-1}J_W\). Since \(\iota\varepsilon\) is the identity on \(\overline{\mathbb Q}_p\), scalar multiplication by \(a\in\overline{\mathbb Q}_p\) before applying \(J_W\) is multiplication by \(\varepsilon(a)\) in the chosen \(k\)-structure. This proves the required semilinearity. Naturality of the three maps gives the second assertion of (HV1) simultaneously for all the \(u\).\formalnote{17.70}{70}

To compute the character, twist a test by \(j(g)\), which fixes the image of \(E\). The underlying rings remain identified. In the twisted structure the generator of the square-zero kernel changes according to \(j(g)^{-1}\mu\equiv\chi(j(g)^{-1})\mu\bmod\mu^2\); this follows by expanding the cyclotomic action on \(1+\mu=[\epsilon]\) to first order. In terms of the maps above, \(\iota_S(z)=\iota_{S^{j(g)}}(\chi(j(g)^{-1})z)\) under this identification. The functorial change of coordinates on pairs therefore scales \(M\) by \(a_0=\varepsilon(\chi(j(g)^{-1}))\) in the square-zero component. For each arrow, the added first component \(\iota_{S'}(a(m))\) agrees in the two descriptions by this very identity. The Galois pullback of the underlying Hodge class agrees with its pullback by the change of coordinates, since \(x\) is defined over \(E\). The contraction itself commutes with twisting of the underlying tests. Moving that scaling through its \(h\) factors gives \(r_Y(g\eta)=a_0^h r_Y(\eta)\). Injectivity of \(r_Y\) yields \(g\eta=a_0^h\eta\). The identity of underlying test rings commutes with passing to generic points, so \(B_Y\) respects this twist, and equivariance of \(J_Y\) identifies it with \(j(g)\otimes\rho(g)\). Injectivity of \(J_Y\) proves the first assertion of (HV1).\formalnote{17.71}{71}

In particular \(v\ne0\) exactly when \(\lambda_h(x_k)\ne0\), by (HV2) and the two isomorphisms in its construction. To deduce occurrence in the eigenquotient, the seed uses the contrapositive. Absence of \(h\) would give a Hecke polynomial \(c\) taking value \(1\) at the required character and vanishing on the weight-\(h\) vectors. For any test polynomial \(b\), put all its legs and those needed for \(c\) on the common cover. The linear combination of their pullbacks for \(bc\), tested by \(\lambda_h\), can be expressed by (HV1) as the combination for \(b\) tested by \(\Theta_W\) on pullbacks of \(c(T)v=0\), transporting the coefficients by \(\varepsilon\). It would vanish for every Hodge class tested. On the projectively embedded cover, the distinct \emph{classes} of points have linearly independent images under \(\lambda_d\). Consequently there is a linear form taking value \(1\) on each of them. Applied to the diagonal identity of \S\ref{sec:woods-hole}, this gives vanishing of the weighted local terms for \(bc\). The nonzero marked sum at the chosen character contradicts that consequence, as in \S\ref{sec:seed-proof}. This is the nonvanishing assertion used to obtain a jump. It requires no injectivity assertion on all Hodge classes.\formalnote{17.72}{72}


\section{Point classes and the coincidence sum}\label{sec:point-classes}

Let \(k\) be algebraically closed, with a chosen embedding in
\(\overline{\mathbb Q}_p\). Let \(X\) be the functor of \(k\)-algebra points
of a projectively embedded scheme over the base ring, with a separated
finite atlas of étale charts of dimension \(d\). Write \([z]\) for the
degree-\(d\) Hodge class of \(z\in X(k)\). The precise assertion is that
the vectors
\[
       \lambda_d(v)\in H^{2d}_{\rm pair}(X),\qquad
       v\in\{[z]\mid z\in X(k)\},
                                                               \tag{PT1}
\]
are linearly independent over \(k\). The indexing set matters. Points on
one projective irreducible component give the same class. Each coincidence
must still contribute once when we apply the functional in
\S\ref{sec:woods-hole}.

\subsection{The symbol for the diagonal}

On an étale chart with coordinates \(x_1,\ldots,x_d\), put
\(f_i=x_i\otimes1-1\otimes x_i\). Shrink the product chart around its
diagonal until equality of coordinates implies equality of points. This
shrinking removes the other branches of the fibre product over affine
space. The \(f_i\) cut out the diagonal there. In the polynomial coordinate
ring the change of variables from the first coordinates to their differences
makes \(f_1,\ldots,f_d\) part of a coordinate system. Flat base change to
the étale chart preserves this regular sequence wherever its zero locus is
nonempty. The familiar facts behind this calculation are Koszul exactness
for a regular sequence \cite[\href{https://stacks.math.columbia.edu/tag/062F}{Tag 062F}]{arith-stacks} and the identification of the conormal bundle of
the diagonal with the relative cotangent sheaf \cite[\href{https://stacks.math.columbia.edu/tag/08S2}{Tag 08S2}]{arith-stacks}.\formalnote{17.73}{73}

Here is the symbol in the square-zero functors used in (PT1).
On a test \(A\oplus M\), write a tuple as \((f_i,\omega_i)_{i=1}^d\).
Off its zero set, the \(f_i\) generate the unit ideal locally. Use the
cover by witnesses \(l=(l_i)\) to \(\sum_i l_i f_i=1\). On an overlap of
\(d\) witnesses the cochain is
\[
       \det(l^{(a)}_b)_{a,b=1}^d\,
                 \omega_1\wedge\cdots\wedge\omega_d .
\]
The alternating sum on \(d+1\) witnesses is zero, by expanding minors
and using \(l^{(a)}f=1\) for each row. Its relative class is denoted
\([\omega_1\wedge\cdots\wedge\omega_d/f_1,\ldots,f_d]\).
If \(f=M_0g\), a witness \(l\) for \(f\) gives \(lM_0\) for \(g\), so
changing the denominator multiplies the numerator by \(\det M_0\)
when both fractions are written over \(f\). A numerator divisible by
any \(f_i\) gives zero. One can see this by changing the denominator
from the \(i\)-th unit vector using the identity matrix with its \(i\)-th
column replaced by \(f\), whose determinant is \(f_i\). Thus if
\(\omega\equiv N_0\eta\bmod(f)\) and
\(\det N_0\equiv a\det M_0\bmod(f)\), the first symbol is \(a\)
times the symbol of \((g,\eta)\).\formalnote{17.74}{74}

Apply this with \(\omega_i=df_i\), viewed through the universal
square-zero lift on the chart. Where two coordinate differences generate
the same ideal \(I\), write \(f=M_0g\). Differentiating modulo \(I\)
gives \(df\equiv M_0\,dg\), so the multiplier is \(1\). These symbols
glue. Indeed, supported chart cohomology vanishes in degrees \(m<d\).
For the first entry \(r\) of the regular sequence, the long exact
sequence for multiplication by \(r\) and induction on the remaining
sequence give injectivity of \(r\) on this group. Localization at \(r\)
kills the group, since the support disappears there. Each class is
therefore killed by a power of \(r\) and must vanish. This argument on
intersections, followed by the relative Mayer--Vietoris sequence of the
finite cover, glues the compatible degree-\(d\) classes. Call the
result \(\Delta_{\rm supp}\), and its absolute image \(\Delta\).\formalnote{17.75}{75}

Pulling the diagonal back along \((\operatorname{const}_z,1_X)\) gives
\([z]\). For a finite reduced fibre of a map between charted functors of
the same dimension this pullback gives the sum of the point classes of
the fibre. This follows by excising its complement and using the symbol
with coefficient one near each point. For a transverse coincidence of \(u,c\), scaling
the first infinitesimal component by \(t\) leaves the denominator \(u-c\)
in coordinates and replaces its differential by \(t\,du-dc\). The change
of numerator just proved gives \(\det(t\,du-dc)/\det(du-dc)\).\formalnote{17.76}{76}

We also use that two points on which all global functions agree give
the same class. To see this, express \(\Delta\) as a sum of external
products. On a represented Hodge test all positive-degree cohomology
vanishes by Yoneda, and degree zero evaluates to the exterior algebra of
its free module. For the product of two represented tests the wedge
basis of the direct sum consists of products of the two wedge bases.
The chart calculation for external products extends by the long exact
sequences for finite unions of chart loci, first in one variable and
then in the other. Zariski descent takes the result to \(X\times X\).
On the point \((k,0)\) a factor of positive degree restricts to zero,
and a degree-zero factor restricts by evaluating a global function.
It follows that pullback along \((\operatorname{const}_z,1_X)\) depends
only on these evaluations. Regular functions on an integral proper
component are constant over \(k\), since a morphism to \(\mathbb A^1\) has
proper irreducible image in \(\mathbb P^1\) missing infinity, hence
image a point.\formalnote{17.77}{77}

\subsection{One nonzero class on projective space}

Keep \(T=S[1/p]\), the square-zero thickening \(B_S\) and the map
\(r_X\) from test cohomology to pair cohomology of
\S\ref{sec:hodge-vector}. The first step constructs
\(\eta\in H^2_{\rm test}(\mathbb P^1,T)\) such that
\(\lambda_1([1:0])=r_{\mathbb P^1}\eta\) and exterior product with
\(\eta\) is injective on \(H^n(N,T)\) for every abelian test functor \(N\)
with flat values. On the first affine chart the symbol is the relative
boundary of \(dX/X\). Its contraction on a pair is computed by lifting
a unit in \(T\) to a unit in \(B_S\). The difference between two such
lifts is measured by
\[
       0\longrightarrow (T,+)\ \xrightarrow{a\mapsto1-\mu a}\
       B_S^\times\longrightarrow T^\times\longrightarrow1.
\]
Here \(\mu a\) denotes the identification of \(T\) with the square-zero
kernel. If \(X=(s,m)\) with \(s\) a unit, the logarithmic differential
evaluates to \(s^{-1}m\). The contraction evaluates its scalar component
on each arrow, precisely the additive difference of the two lifts in
this sequence. Composing the units extension with the boundary of
the pair \((\mathbb A^1,\mathbb A^1-\{0\})\) gives a degree-two test
class. The other chart avoids the point, so relative excision gives the
asserted \(\eta\) after forgetting supports.\formalnote{17.78}{78}

For injectivity, take the two perfect discs represented on tilts by
\(\mathcal O^\flat[\mathbb Q_{\ge0}]\), and the torus by
\(\mathcal O^\flat[\mathbb Q]\), where \(\mathcal O=\mathcal O_{\mathbb C_p}\).
The gluing map \(i\) from the free abelian functor on the torus and one
point to that on the two discs sends a torus point to the difference
of its reciprocal in the first disc and itself in the second. The extra
point maps to the first origin. In degree zero, pullback sends a pair of
polynomials \((a,b)\) to \((a(0),a(U^{-1})-b(U))\). The constant and
the negative and nonnegative parts of a Laurent polynomial recover the
two inputs. For the almost coefficients on tilts, higher Ext vanishes on
the perfect algebras representing these pieces. The calculation survives
tensoring their rings with the tilt of any test. The Witt-vector lifting
and inversion of \(p\) in \S\ref{sec:primitive-comparison} then give
the same cohomology equivalence for \(T\). Resolving \(N\) by sums of free
representables and using flatness yields
\(\operatorname{Ext}^*(N\otimes\operatorname{coker}i,T)=0\).\formalnote{17.79}{79}

On the overlap, take the Teichmüller lift \([u]\) of the coordinate and
twist it by the compatible roots of unity \(\epsilon\). Its image in \(T\)
is fixed. Its lift changes by \([\epsilon]=1+\mu\), so
\((1-\mu)[\epsilon u]=[u]\). The difference measured in the additive
kernel of the displayed units extension is therefore one. Here is
how this detects a product. Denote one minus the twist on the overlap
and discs by \(\delta'\) and \(\delta\), with \(i\delta'=\delta i\).
Take the pushout of \(\delta'\) and \(i\), and the induced map
\(\sigma=(i,\delta)\) from this pushout to the discs. The inclusion of
the overlap into the discs and its pushout both have the same acyclic
cokernel. Consequently \(\sigma\) induces an isomorphism on Ext into
\(T\) after tensoring with \(N\). The two short exact sequences for
its image give a two-extension from \(\operatorname{coker}\sigma\) to
\(\ker\sigma\); composing their two boundary maps is an isomorphism
on total Ext, by the two long exact sequences. The coordinate lifts
just constructed map these two sequences to the affine relative and
units sequences. On their common kernel the second origin maps to
the computed defect \(1\). Thus if \(x\boxtimes\eta=0\), its pullback
to \(\operatorname{coker}\sigma\) vanishes. Cancel the composite of
the two boundary maps and evaluate at that origin. The result is
\(x\cdot1=0\). In particular \(\eta\ne0\), by taking \(x\) to be the
identity of the nonzero coefficient functor.\formalnote{17.80}{80}

For higher dimension use the map
\(\pi:\mathbb P^d\times\mathbb P^1\to\mathbb P^{d+1}\) given by
\[
 ([v_0:\cdots:v_d],[a:b])\longmapsto
 [\,a{\textstyle\sum_i}v_i:(b-0a)v_0:\cdots:(b-da)v_d\,].
\]
Over \([1:0:\cdots:0]\) the fibre consists of the \(d+1\) reduced points
\((e_j,[1:j])\), \(0\le j\le d\). Indeed \(a\sum v_i\) is a unit on
the affine target chart at a point of the fibre. Over a local ring
some \(v_j\) is a unit, so \(b=ja\). For \(i\ne j\), \(j-i\) is a
nonzero scalar, forcing \(v_i=0\). This works also over the dual numbers.
The product rule for the local symbol concatenates the equations for
the two points, giving the exterior product of their classes. All
projective points in each factor have the same class. Hence the
pullback of the target point class is \(d+1\) times that external product.
The weight \(l\) component of a point in dimension \(d\) is zero for
\(l\ne d\). The scaled identity against the constant map has determinant
\(t^d\), whereas \(\lambda_l\) scales by \(t^l\). Compare \(t=1,2\).
In the weight \(d+1\) product formula only the two weights \(d,1\)
remain. The binomial coefficient \(d+1\) and the sign are invertible
in \(k\). Since \(r_X\) is bijective on charted functors, cancellation
of \(\eta\) proves that the weight of this pullback vanishes exactly
when \(\lambda_d([e_0])\) vanishes. In dimension zero the map \(\pi\)
has a section, so such vanishing would force \(\eta=0\). Induction
proves nonvanishing on every \(\mathbb P^d\).\formalnote{17.81}{81}

\subsection{Separating a class on a projective scheme}

Fix \(z_0\in X(k)\) and a finite set \(D\) of points whose classes differ
from \([z_0]\). Choose an irreducible component through \(z_0\) and
take a homogeneous Noether normalization of its cone by linear forms
\(\Lambda_0,\ldots,\Lambda_d\). On \(\Lambda_0\ne0\) the coordinate
ratios are algebraically independent and give a generically finite
projection. Its function field extension is separable in characteristic
zero. After localization these ratios are étale coordinates on a chart
\(U\). Both they and the original coordinates give bases of the
differentials, accounting for the index \(d\). Choose a point
\(Q=[1:q_1:\cdots:q_d]\) whose component fibre lies in \(U\).
One way to see the choice is to take each nonzero element that must
remain invertible and clear denominators in an algebraic equation for
it over the coordinate ratios. Avoiding the zero of the nonzero constant
coefficient keeps that element invertible. Such choices exist since \(k\) is infinite.\formalnote{17.82}{82}

Form the graph of this projection wherever the coordinates can be
represented by a unimodular tuple and the evaluated forms are unimodular.
In the graph take the open \(O\) omitting the points above \(D\) and
the points over \(Q\) outside \(U\). Write \(j:O\to X\) and
\(\pi:O\to\mathbb P^d\) for its two maps. The first is injective on
points over every algebra, since changing a coordinate tuple multiplies all
evaluated forms by the same unit. The component selected above is
retained, since global functions are constant on it and its \(Q\)-fibre
lies inside \(U\).

The required injectivity of \(\pi^*\) is in the test cohomology.
To prove it, take the finite Galois closure of the normalization on
the cone, with group \(G\). A lift of a projective test point gives a
family of points of \(O\), indexed by the elements of \(G\), above that
same point. To lift at a prime of the generic test ring, use the
algebraically closed fraction field of its valuation localization.
Integrality gives a lift of the finite cone algebra to this field.
The images of finitely many algebra generators are integral over
the image of the test ring, which is absolutely integrally closed,
so these values lift to the test ring. Their finitely many defining
relations hold after localizing at one element outside the prime.
Over the double overlap the two families are permutations of each
other locally. At the geometric point this follows by conjugacy in
the Galois closure, and the finitely many coordinate equalities also
hold after localization. Summing
the family therefore descends along the cover. More explicitly it
gives maps \(l:L\to\mathbb Z[\mathbb P^d(T)]\) and
\(\nu:L\to\mathbb Z[O(T)]\), where \(l\) induces an isomorphism on
Ext into \(T\) by Zariski descent. Their composite satisfies
\(\mathbb Z[\pi]\nu=|G|l\). If \(\pi^*\xi=0\), this implies
\(|G|l^*\xi=0\), hence \(\xi=0\) in characteristic zero.\formalnote{17.83}{83}

The remaining fibre \(F=\pi^{-1}(Q)(k)\) is finite and reduced. In
the étale chart, fixing the coordinate ratios fixes every first-order
lift. Reduced-fibre pullback gives
\[
                 \pi^*[Q]=\sum_{y\in F}[y]
                          =|F|\,j^*[z_0].
\]
Indeed each \(j(y)\) has class \([z_0]\) by constancy of global
functions on the chosen component, and the fibre of \(j\) over \(j(y)\)
is the single reduced point \(y\). Apply \(\lambda_d\). On the left,
the projective-space class comes from a nonzero test class, its
pullback stays nonzero by the transfer, and restriction to pairs is
injective for \(O\) by its étale charts. Consequently
\(\lambda_d(j^*[z_0])\ne0\). For \(z\in D\) the fibre of \(j\) is empty,
so \(j^*[z]=0\). Pulling back any proposed finite linear relation in
(PT1), with \(D\) representing its other classes, shows the coefficient
of \([z_0]\) is zero. This proves (PT1).\formalnote{17.84}{84}

For the \(i\)-th map \(u_i\) at the common Hecke level, let \(w\) range over its coincidences \(u_i(w)=c(w)\).
On the span of their distinct vectors (PT1) define a linear form by
\(\sum_v a_v\lambda_d(v)\mapsto\sum_v a_v\) and extend it to the ambient
vector space. It sends \(\sum_w m_w(t)\lambda_d([w])\) to
\(\sum_w m_w(t)\), including every repeated class with its full
multiplicity. For the \(i\)-th map, \(P_{iw}(t)\) is the tangent polynomial at the coincidence \(w\). Its normalized value \(m_w(t)=P_{iw}(t)/P_{iw}(1)\) is the coefficient of \([w]\) in the diagonal pullback. Taking the coefficient of \(t^h\) after applying the linear form gives the numerical local terms in \textup{(WH)}. The separate passage from these local terms to the marked
lattice count is given in \S\ref{sec:woods-hole}.\formalnote{17.85}{85}


\section{From special-fibre points to linear Frobenius on periods}
\label{sec:cohomological-frobenius}

Let \(u_i:W\to Y\) be maps of proper schemes. Suppose a linear combination of these maps, each composed with a specified power of Frobenius, vanishes on the field-valued points of the special fibre in the free abelian group on the target points. The vanishing in \S\ref{sec:seed-proof} has this form. The aim is to obtain a linear relation between the resulting pullbacks on the period space of the \(p\)-adic cohomology. First we pass from field-valued points to the test rings on which the comparison is defined. Then we use invariant periods to obtain a linear operator from the semilinear Frobenius action. The same comparison gives crystallinity.

Let \(\Omega=\overline{\mathbb Q}\), let \(R\) be noetherian, and let \(R\to K\to\Omega\) be the base maps with \(K\) a field. Fix a \(p\)-adic place of \(\Omega\) over \(K\) with embedding \(e:\Omega\hookrightarrow\overline{\mathbb Q}_p\), decomposition group \(\Gamma\), and compatible open embedding \(j:\Gamma\to G_p\). The image of \(R\) is integral at this place. Write \(O_C=\mathcal O_{\mathbb C_p}\), \(\sigma:R\to O_C\), and
\[
       \mathfrak p=\{r:|\sigma(r)|_p<1\},\qquad k=R/\mathfrak p .
\]
Assume that \(Y,W\) are proper over \(R\) and charted over \(\Omega\), and that \(Y\) is also charted over \(R\). The charts here use étale coordinates that are units. Let \(V^h(Y)\) be the degree-\(h\) rational \(p\)-adic test-ring cohomology of \(Y\), with Galois action \(\rho\), as in \S\ref{sec:seed-proof}. Put \(B=B_{\max}\) and call \(B\otimes_{\mathbb Q_p}V^h(Y)\) its period space. Given integers \(m_i\), nonnegative \(a_i\) and \(s_i\in\Gamma\), the two hypotheses are
\[
 \sum_i m_i[\,u_i(x)\circ{\rm Frob}^{a_i}\,]=0,\qquad
 j(s_i)^{-1}\bmod\mathfrak m_{\mathcal O_{\overline{\mathbb Q}_p}}
                       ={\rm Frob}^{a_i}.
\]
The first equality is required for every field-valued point of \(W\) killing \(\mathfrak p\) in characteristic \(p\). Here precomposition by Frobenius means applying \(\lambda\mapsto\lambda^{p^{a_i}}\) to a point's coordinates. The second equality concerns the action on the full residue field of the algebraic integers at \(p\). A pair \((j(s),a)\) satisfying it is called a Weil pair. For such a pair \(\beta_{s,a}\) denotes the \(B\)-linear operator that agrees with \(\varphi^a j(s)\otimes\rho(s)\) on periods fixed by an open subgroup of \(\Gamma\), where \(\varphi\) is the period-ring Frobenius. The construction below proves its existence and uniqueness. The conclusions are crystallinity of \(\rho|_\Gamma\) and\formalnote{17.86}{86}
\[
      \sum_i m_i(1\otimes u_i^*)\circ\beta_{s_i,a_i}=0
               \quad\text{from } B\otimes_{\mathbb Q_p}V^h(Y)
                           \text{ to } B\otimes_{\mathbb Q_p}V^h(W).
                                                               \tag{CF1}
\]

\subsection{Extending the relation beyond fields}

\begin{example}[Passing from residue fields to a quotient of a valuation ring]\label{ex:comparison-reduction-scalar}
Take \(T=\mathbb Z_p[t]/(t^2-p)\). Two maps from \(\operatorname{Spec} T\)
to the projective line have affine coordinates \(0\) and \(t\).
They agree at every field-valued point of the special fibre, since
\(t^2=0\bmod p\). They still differ in \(T/pT=\mathbb F_p[t]/(t^2)\).
In the valuation ring \(O_C=\mathcal O_{\mathbb C_p}\), evaluate \(t\) at either
square root of \(p\) in \(O_C\). With \(v_p(p)=1\), its value has
valuation \(1/2\), so its image in \(O_C/(p)\) is nonzero.
If \(x_0\) has \(0<v_p(x_0)\le1/2\), both square roots vanish in \(O_C/(x_0)\). Choose \(x_0^3=p\), so \(v_p(x_0)=1/3\). Then both
square roots belong to \((x_0)\). Thus the two maps agree in
\(O_C/(x_0)\) for both evaluations. Here \((p)=(t)^2\) in \(T\). In the argument below \(M\) prime ideals, with repetition allowed, have product contained in the ideal of the special fibre. The root chosen there has degree \(M+1\). Here one can take the two ideals \((t),(t)\), giving \(M=2\).
\end{example}

We seek a single \(x_0\in O_C\) such that the field-point relation holds on the appropriate quotients of all tests. The ring \(S\) of a \(p\)-adic test is a product of valuation rings called cells. The factorization lemma for proper \(W\) gives a single \({\rm Spec}\,T\to W\), with \(T\) noetherian, such that any test point factors through it. Put \(\mathfrak a=\mathfrak pT\). Choose primes \(\mathfrak q_1,\ldots,\mathfrak q_M\) containing \(\mathfrak a\) such that \(\prod_l\mathfrak q_l\subseteq\mathfrak a\). This follows by pulling back a product of primes that vanishes in \(T/\mathfrak a\). In \(O_C\) the image ideal of \(\mathfrak p\) is principal because \(\mathfrak p\) is finitely generated and divisibility is totally ordered. Choose \(x_0^{M+1}\) equal to a generator. The ideal contains \(p\) and lies in the maximal ideal. Thus \(x_0\ne0\), \(|x_0|_p<1\), and \(p\mid x_0^N\) for some \(N\).\formalnote{17.87}{87}

For a factorization \(T\to S\), consider one valuation cell. Some \(\mathfrak q_l\) maps into \((x_0)\) there. Otherwise pick \(b_l\in\mathfrak q_l\) whose image is not divisible by \(x_0\). Each image divides \(x_0\), so their product divides \(x_0^M\). The product is also in \(\mathfrak a\), so its image is divisible by \(x_0^{M+1}\). This is impossible because the image of \(x_0\) in the cell is nonzero and a nonunit, as \(x_0^{M+1}\mid p\) and \(p\mid x_0^N\) there. Assign such an \(l\) to each cell. The idempotents selecting the cells with each label sum to \(1\). At any maximal ideal of \(S\) one of these idempotents, say \(b\), is invertible locally. In any further ring map \(S\to D\) inverting \(b\) and killing \(p,x_0\), the composite \(T\to D\) therefore kills \(\mathfrak q_l\).

To obtain the relation over \(T/\mathfrak q_l\), apply the field-point hypothesis first to its fraction field. If two twisted maps agree at the generic point, their maps to \({\rm Spec}\,R\) agree throughout, by the injection of \(T/\mathfrak q_l\) into its fraction field. Their equalizer over this common base is closed because \(Y\) is separated over \(R\). The generic point is schematically dense in the integral scheme \({\rm Spec}(T/\mathfrak q_l)\), so the two maps agree there. The same terms therefore cancel over \(T/\mathfrak q_l\), and pulling back along \(T/\mathfrak q_l\to D\) gives the required cancellation on the test. The scalar \(x_0\) works for every test.\formalnote{17.88}{88}

\subsection{Test rings supplying spanning classes}

For a \(p\)-complete, \(p\)-torsion-free test algebra \(D\), set
\[
          \mathfrak d(D)=\{x:\ n!\mid x^n\text{ for all }n\ge0\}.
\]
The binomial theorem makes this an ideal. It contains \(p\), since \(v_p(n!)\le n\). In particular \(D/\mathfrak d(D)\) has characteristic dividing \(p\) and carries its \(p\)-power endomorphism. Give \(A_{\max}\) the \(A_{\inf}\)-algebra map obtained by applying \(\varphi^{\,b}\) before its usual structure map, \(b=\max(N,3)\). It commutes with Frobenius and Galois. The usual map sends \((p,\xi)\) into \((p)\), and \(\varphi^b\xi\in(p,\xi)\). Hence the composite into \(A_{\max}/\mathfrak d(A_{\max})\) factors through \(\theta:A_{\inf}\to O_C\). Call the resulting map \(\eta\). For \(x_0=v^\sharp\), the untilt of \(v^{p^b}\) is \(x_0^{p^b}\), divisible by \(p\). The Witt-coordinate test then puts \(\varphi^b[v]\) in \((p,\xi)\), giving \(\eta(x_0)=0\). The same twist sends \(\mu=[\epsilon]-1\) into \(p^3 A_{\max}\), while its image is invertible in \(B\).\formalnote{17.89}{89}

The domain \(k\) embeds in the residue field at the place. If \((j(s),a)\) is a Weil pair, \(s^{-1}\) fixes the image of \(R\) in characteristic zero and acts on its residue by the \(p^a\)-power map. Thus \(z^{p^a}=z\) on \(k\). The open decomposition group contains a geometric Frobenius lift with some positive exponent \(a\), independently of the chosen legs. The polynomial \(X^{p^a}-X\) bounds the size of \(k\), so \(k\) is a finite field. Since \(x_0\) divides all elements of \(\sigma(\mathfrak p)\), \(\eta\sigma\) induces \(k\to A_{\max}/\mathfrak d(A_{\max})\). Its unique Witt lift gives \(W(k)\to A_{\max}\), characterized by agreement on zeroth Witt coordinates after reduction. In \(W(k)\), \(\mathfrak d(W(k))=(p)\), since a unit fails the divisibility test at \(n=p\). Thus for a \(W(k)\)-test \(D\), the map from \(R\) to \(D/\mathfrak d(D)\) sends \(r\) to the image of \([\bar r]\).\formalnote{17.90}{90}

Choose a charting subfunctor \(U\) of the points of \(Y\) over \(R\). It covers locally and is built from étale torus charts. Write \(H^h_{\rm red}(U)\) for the degree-\(h\) cohomology of \(D\mapsto U(D/\mathfrak d(D))\) on the \(p\)-complete torsion-free \(W(k)\)-algebras, with coefficient functor \(D\mapsto D\). For an ordinary \(p\)-adic test \(S\) define
\[
 D_S=\bigl(A_{\max}\otimes_{A_{\inf}}W(S^\flat)\bigr)^{\wedge}_p,
 \qquad \Psi(S)=B\otimes_{A_{\max}}D_S .
\]
There is a reduction map \(S\to D_S/\mathfrak d(D_S)\). It sends \(\theta(w)\) to the class of \(1\otimes w\). For every \(a\) with \(z^{p^a}=z\) on \(k\), compose this reduction with the \(p^a\)-power map to get a map of points over \(R\), denoted \({\rm red}_a\). For each \(r\in R\), the two maps to the quotient agree because the image of \(\eta(\sigma(r))\) equals the image of the Teichmüller lift \([r\bmod\mathfrak p]\in W(k)\). The \(p^a\)-power map fixes this common image since it fixes \(k\). Write \({\rm red}_a^*\) for the resulting pullback on cohomology with its coefficient map.\formalnote{17.91}{91}

There are two comparisons. First the period map \(J_Y\), followed by restriction to \(U\), gives
\[
             E_Y:B\otimes_{\mathbb Q_p}V^h(Y)
                       \ \simeq\ H^h(U(S),\Psi(S)).
\]
The notation on the right means cohomology of the indicated functors as \(S\) varies. On \(Y\) the construction moves \(B\) inside the cohomology by flat change of coefficients, then sends the Witt vectors of locally constant \(\mathbb F_p\)-valued functions on \({\rm Spec}(S[1/p])\) to \(W(S^\flat)\). Properness and the charts over \(\Omega\) give bijectivity by primitive comparison and passage to complete period coefficients. Zariski descent allows restriction to \(U\). The map on \(W\) before any restriction is likewise bijective.\formalnote{17.92}{92}

Second, the map \({\rm red}_0^*\) on \(H^h_{\rm red}(U)\) has \(B\)-spanning image. For this calculation the integral charts of \(Y\) matter. On a chart with étale unit coordinates \(t_1,\ldots,t_d\), adjoining all their compatible roots is the base change from \(R[X_1,\ldots,X_d]\) to \(O_C[\mathbb Q_{\ge0}^{\,d}]\). The resulting chart is étale over the monoid algebra. Its étale lift over \(A_{\inf}[\mathbb Q_{\ge0}^{\,d}]\) has the same \(p\)-completion as the Witt vectors of the corresponding perfect special algebra. Comparison there follows by the Witt lifting calculation. One then descends one root system at a time. On crystalline tests the root coordinate \(v\) over \(t\) has divided-power basis \(v^r(t-v)^n/n!\), \(0\le r<1\) rational. The twist \(\delta\) multiplies \(v^r\) by \([\epsilon^r]\). Here \(t\) is a unit modulo \(\mathfrak d\), so \(v\) is a unit in the complete envelope. After extension to \(B\), \(\delta-1\) is surjective with kernel the constant summand. The calculation uses \(p^3\mid\mu\), invertibility of \(\mu\) in \(B\), and \(p([\epsilon^r]-1)\mid\mu\) for \(0<r<1\), all read through the chosen structure map. Together with the corresponding root-system calculation on the ordinary tests this gives the comparison at each step.\formalnote{17.93}{93}
Section~\ref{sec:crystalline-towers} gives the coordinate calculation and the denominators in this descent.

Passing from \(W(k)\)-tests to \(A_{\max}\)-tests uses flatness, since every nonzero element of \(W(k)\) divides a power of \(p\) and \(A_{\max}\) is \(p\)-torsion-free. Completed scalar extension is left adjoint to restriction. The map from the ordinary to the completed tensor product induces a cohomology isomorphism by a criterion that uses finite generation in all degrees after extending to \(B\), supplied by primitive comparison. Its kernel test applies adic coordinate maps \(A_{\max}\to W(k)\) to a \(p\)-divisible kernel, whose images vanish in derived \(p\)-complete cohomology. Combining this with flat base change and the chart comparison gives the isomorphism
\[
 B\otimes_{A_{\max}}\bigl(A_{\max}\otimes_{W(k)} H^h_{\rm red}(U)\bigr)
                  \ \simeq\ H^h(U(S),\Psi(S)).
\]
It sends \(1\otimes1\otimes z\) to \({\rm red}_0^*z\) because the unit and counit of completed base change compose to the identity. These elements span on the left. For \(z\in H^h_{\rm red}(U)\), define the period \(Cz=E_Y^{-1}({\rm red}_0^*z)\). These periods span the period space \(B\otimes V^h(Y)\).\formalnote{17.94}{94}

\subsection{The semilinear action and its linear extension}

\begin{example}[Linear extension of an action on invariant vectors]\label{ex:period-linear-extension}
Let \(K_0=\mathbb F_{p^2}\), let \(F(c)=c^p\), and let \(F\) act
on \(K_0e_1\oplus K_0e_2\) by acting on the coefficients. Its
invariant vectors form \(W=\mathbb F_pe_1\oplus\mathbb F_pe_2\).
Consider the \(F\)-semilinear operator \(S\) that swaps the two basis vectors and applies \(F\) to their coefficients.
\[
       S(ae_1+be_2)=F(b)e_1+F(a)e_2 .
\]
On \(W\), this is just the interchange of \(e_1,e_2\). Since
\(e_1,e_2\) are also a basis over \(K_0\), the restriction to \(W\)
has a unique \(K_0\)-linear extension \(L\). For \(c\in K_0\setminus
\mathbb F_p\) the two values are \(S(ce_1)=c^p e_2\) and
\(L(ce_1)=ce_2\).
\end{example}

Call a period \emph{potentially invariant} if it is fixed by an open subgroup of \(\Gamma\). To define the \(B\)-linear operator \(\beta_{s,a}\) of \textup{(CF1)}, we must show that \(\varphi^a j(s)\) fixes the coordinates of a potentially invariant period in a basis of potentially invariant periods.

For \(s\in\Gamma\) and \(a\ge0\), write \(F_{a,s}=\varphi^a j(s)\) on period scalars. If the \(p^a\)-power map is the identity on \(k\), then for each cohomology class \(z\in H^h_{\rm red}(U)\) one has
\[
          E_Y\bigl((F_{a,s}\otimes\rho(s))Cz\bigr)={\rm red}_a^* z .
                                                               \tag{CF2}
\]
To verify it, twist the test by \(j(s)\). The map back between its period rings sends a pure tensor \(b\otimes w\) to \(F_{a,s}(b)\otimes W({\rm Frob}^a)w\), then extends to the completion. On a Teichmüller vector it raises the reduced test point to the \(p^a\)-th power. Every test element is an untilt, so this identity determines the map on all reduced points. Also \(F_{a,s}\) fixes the image of \(W(k)\). Indeed it preserves the reduction of the zeroth Witt coordinate since \(s\) fixes \(R\) and \(p^a\) fixes \(k\), and the Witt lift is unique. Hence the map is available on the very tests defining \(z\). On coefficients the two paths send \(1\otimes x\) to the same \(1\otimes F(x)\), with \(F\) the indicated map of test period rings. Finally \(J_Y\) commutes with this twist on pure tensors because Witt vectors of locally constant \(\mathbb F_p\)-valued functions have Frobenius-fixed coefficients. This proves (CF2).\formalnote{17.95}{95}

For \(a=0\), (CF2) and injectivity of \(E_Y\) imply \((j(s)\otimes\rho(s))Cz=Cz\) for every \(s\). These periods \(Cz\) are invariant under the full group \(\Gamma\) and span the period space, proving crystallinity. For a Weil pair \((j(s),a)\), the operator \(\beta_{s,a}\) is \(B\)-linear and agrees with \(F_{a,s}\otimes\rho(s)\) on potentially invariant periods. To construct the extension, choose from the spanning periods \(Cz\) a basis over \({\rm Frac}(B)\) after extension of scalars, where \(B=B_{\max}\). Write a potentially invariant period \(v\) in this basis. An open subgroup fixes \(v\) and the basis vectors, so uniqueness of coordinates forces it to fix each coordinate in \({\rm Frac}(B)\). The fixed-field result below puts these coordinates in a finite unramified field \(L_0\subset B\). This applies in particular to every \(Cz\). Since the \(Cz\) span over \(B\), the selected family is already a \(B\)-basis. On \(L_0\) the inverse residue Frobenius \(j(s)\) cancels \(\varphi^a\), so \(F_{a,s}\) fixes the coordinates of \(v\). Extending the prescription \(B\)-linearly from that basis therefore gives the required value on every potentially invariant period. Spanning gives uniqueness. In particular, (CF2) computes \(\beta_{s,a}Cz\). On an arbitrary multiple \(bCz\) the linear extension uses \(b\) itself.\formalnote{17.96}{96}

\phantomsection\label{route:31}\leanroute{The source proves the fixed-ratio assertion by extending coordinates from a Kummer chart and clearing the poles of the periods. We use the fixed-field corollary of Fontaine's period-ring theorem after identifying the two normalizations of \(B_{\max}\). The comparison on the test rings and the Frobenius calculation (CF2) are still required.}
If \(u/v\), \(v\ne0\), is a ratio in \({\rm Frac}(B_{\max})\) fixed near \(1\), it is fixed by \(G_L\) for some finite \(L/\mathbb Q_p\). Put \(L_0=W(k_L)[1/p]\). In the normalization used here \(A_{\max}\) is the \(p\)-completion obtained by adjoining \(\xi/p\) to \(A_{\inf}\), where \(\xi\) generates \(\ker\theta\). Changing this generator by a unit gives the same algebra. After inverting \(p\), inverting \(\mu\) is equivalent to inverting the crystalline period \(t\), by the period identification in \S\ref{sec:local-lengths}. The corollary to Fontaine's injectivity theorem gives \({\rm Frac}(B_{\max})^{G_L}=L_0\) \cite[\S13, corollary to Theorem~13.1]{legacy-berger}. Thus \(p^j u=w v\) for some \(w\in W(k_L)\) and \(j\ge0\). For a Weil pair of exponent \(a\), inverse residue Frobenius followed by \(\varphi^a\) fixes this \(w\), giving precisely the linearization used above.\formalnote{17.97}{97}

It remains to prove (CF1) on these spanning vectors. In the tests for \(W\), restrict to points that every \(u_i\) sends into \(U\). This restriction preserves cohomology by descent and finite intersection. By naturality of \(J\) and (CF2), the \(i\)-th term on \(Cz\) is now the pullback of the \emph{same} \(z\) along the reduction of \(u_i\) followed by \(p^{a_i}\)-power Frobenius. The preceding local relation applies in \(D_S/\mathfrak d(D_S)\), where \(p,x_0\) vanish and any test unit stays invertible. Thus the integer coefficients sum to zero on each fibre of these maps of points, locally on every test. On the locus of cancellation their weighted sum is zero on free abelian groups, hence also on pullbacks in cohomology. Zariski descent lets us check on that locus. Injectivity of \(J_W\) and spanning by \(Cz\) finish the proof.\formalnote{17.98}{98}

At a safe PEL place of norm \(p^f\), let \(s\) be a geometric Frobenius there. The \(B_{\max}\)-linear operator on periods associated with the Weil pair \((s,f)\) is \(\beta_{s,f}\) of \textup{(CF1)}. The term of exponent \(d\) in the Hecke polynomial uses the pair \((s^d,df)\). Uniqueness on fixed periods gives \(\beta_{s^d,df}=\beta_{s,f}^d\). The integral atlas of the base and the geometric atlases of the base and cover supply precisely the chart hypotheses above. After pulling back by the injective identity leg, (CF1) becomes the Hecke relation of Section~\ref{sec:seed-proof}. The maps to the eigenquotient are Galois-equivariant, so they carry invariant periods to invariant periods and intertwine these linearized operators.\formalnote{17.99}{99}

\paragraph{Scope of the comparison.}
The argument uses rational test-ring cohomology. Sections~\ref{sec:test-proper-base}--\ref{sec:primitive-comparison} give primitive comparison, finiteness on proper charted schemes, and the lift to period coefficients. Section~\ref{sec:crystalline-towers} proves the exactness calculation on the two root towers. The fixed-ratio theorem above justifies the scalar extension of Frobenius. Section~\ref{sec:hodge-vector} realizes the Hodge class used in the trace comparison by a vector of the specified weight.\formalnote{17.100}{100}


\appendix\renewcommand{\currentpart}{Appendices}\renewcommand{\chaptermark}[1]{\markboth{Appendix \thechapter}{}}

\chapter{Comparison with the formal development (optional reference)}\label{sec:outstanding}
The entries below identify places where the argument in the text differs from the readable Lean edition at \texttt{72fead1}. Each entry names the calculation or theorem concerned. Where a comparison uses a characteristic-zero result to obtain an integral bound, the linked passage also gives the argument on lattices.
Each heading links to the passage where the two arguments are compared. Chapter~\ref{ch:volume} combines Theorems A, B and C, and Chapter~\ref{ch:analytic} deduces the conclusion about zeros. Parts II--IV give the proofs of the three estimates.
For the division between established results and the estimates particular to this argument, see the roadmap in \S\ref{map:forward}. The comparisons here record a different distinction, namely where the written proof uses a different derivation of the same input.
\section*{Proofs from the community revision}
The following passages use the community edition at \texttt{120962b4}.
For CM sections, \hyperlink{call.4.18}{the local height and a single evaluation}
give the integral coefficient bound. In the finite-place mass calculation,
\hyperlink{call.6.61}{translation in a hyperbolic plane} bounds the mass of a box on a quadric.
For the real discrete series, \hyperlink{call.7.18}{a trace logarithm} controls each minor.
The \hyperlink{call.10.8}{equal-rank theta period} uses two exchanges in the prescribed signature
and \hyperlink{call.10.46}{descent on explicit determinant vectors}.
The \hyperlink{call.12.51}{effective class-number bound} bounds the number of generators needed
for ray-class detection.
The \hyperlink{call.14.52}{Frobenius-stability argument} uses the coordinate projection directly.
For injectivity of the definite theta lift the community edition adopts the
\hyperlink{call.15.60}{second test function in the doubling integral}
from the authors' compact edition.

\section*{The main theorem and its reduction}
\noindent\hyperref[route:1]{\textbf{Primes for the two level windows} (p.~\pageref*{route:1})}. Siegel--Walfisz gives the required uniformity for moduli bounded by powers of the logarithm. Lean obtains a zero-free estimate for the special moduli in the construction.\par\smallskip
\noindent\hyperref[route:2]{\textbf{The error at the auxiliary cutoff} (p.~\pageref*{route:2})}. The quantitative prime number theorem for the fixed modulus replaces a zero-free region and a quantitative Tauberian theorem. Partial summation still determines the centering constant.\par\smallskip
\noindent\hyperref[route:3]{\textbf{Theta transformation in the zero detector} (p.~\pageref*{route:3})}. The primitive-character theta transformation supplies the identity obtained by Poisson summation in Lean. Only its absolute value enters.\par\smallskip
\noindent\hyperref[route:4]{\textbf{Averaging over twisting primes} (p.~\pageref*{route:4})}. Heath-Brown's quadratic large sieve absorbs the fixed, possibly complex character into its coefficients. It replaces the off-diagonal estimates for Jacobi symbols.\par\smallskip
\noindent\hyperref[route:5]{\textbf{Choosing a twist in the required progression} (p.~\pageref*{route:5})}. A power saving in the number of bad primes suffices when compared directly with the prime number theorem in the progression. Lean makes the selection using a divergent weighted prime sum.\par\smallskip
\noindent\hyperref[route:6]{\textbf{The real singularity of a positive transform} (p.~\pageref*{route:6})}. Landau's theorem supplies the singularity at the abscissa of convergence. The residue and its sign are still calculated for the prime statistic.\par\smallskip
\section*{Short vectors and spanning by translates}
\noindent\hyperref[route:7]{\textbf{The sigma addition formula} (p.~\pageref*{route:7})}. The classical addition theorem replaces evaluation in a theta space of degree two. The change of quadratic gauge and the resulting arithmetic duplication formulas are retained.\par\smallskip
\noindent\hyperref[route:8]{\textbf{The CM lattice zeta constant} (p.~\pageref*{route:8})}. Kronecker's first limit formula supplies the constant term. The area, scale and lattice discriminant are converted to the normalization of the elliptic period.\par\smallskip
\noindent\hyperref[route:34]{\textbf{Rational invariance of the theta sum} (p.~\pageref*{route:34})}. Weil's theorem supplies invariance under the rational splitting. The comparison of the two theta functionals uses the normalized vacuum mean.\par\smallskip
\noindent\hyperref[route:9]{\textbf{Products of flag minors} (p.~\pageref*{route:9})}. Howe's standard monomial basis gives the required independence. The proof still checks integrality and bounds the archimedean norms of the chosen spanning vectors.\par\smallskip
\noindent\hyperref[route:10]{\textbf{Mass of the compact coefficient} (p.~\pageref*{route:10})}. Schur orthogonality and Weyl's dimension formula bound the mass used to normalize the Gaussian integral. The full-group identity-neighborhood bound uses the derivative estimate as in Lean.\par\smallskip
\noindent\hyperref[route:11]{\textbf{Lowest compact type} (p.~\pageref*{route:11})}. The compact averaging identities give the root-cone support. Blattner's formula identifies the Harish--Chandra parameter of the actual coefficient for use in the published character formula.\par\smallskip
\noindent\hyperref[route:12]{\textbf{Lattice transference} (p.~\pageref*{route:12})}. Banaszczyk's theorem bounds the covering radius by the first minimum of the dual. The Hermitian trace-dual normalization is computed locally.\par\smallskip
\noindent\hyperref[route:13]{\textbf{Adelic reduction} (p.~\pageref*{route:13})}. Borel's reduction theory supplies the qualitative volume and compactness input. The estimates depending on the first minimum still have to be established.\par\smallskip
\noindent\hyperref[route:14]{\textbf{Local-to-global matching of Hermitian forms} (p.~\pageref*{route:14})}. Landherr's theorem identifies the rational isometry class once the local invariants agree. Lean proves a restricted Hasse principle by norm and frame constructions.\par\smallskip
\noindent\hyperref[route:15]{\textbf{Finite cuspidal tests} (p.~\pageref*{route:15})}. The classical character construction gives cuspidality at a regular label. Vanishing of the required translated mean is then deduced for the actual finite test.\par\smallskip
\noindent\hyperref[route:16]{\textbf{Elliptic orbital characters} (p.~\pageref*{route:16})}. The elliptic character identity and discrete-series formal degree replace the compact-multiplicity and open-orbit derivation. Their signs and measures are matched to the trace formula.\par\smallskip
\noindent\hyperref[route:17]{\textbf{Compact-dual measure} (p.~\pageref*{route:17})}. A tangent-space comparison determines the Haar normalization required by the published character formula, in place of the integral over an open flag orbit.\par\smallskip
\noindent\hyperref[route:18]{\textbf{Descent to the fibre of a polynomial projector} (p.~\pageref*{route:18})}. Frobenius descent for invariant distributions gives the passage to a stabilizer. Its unimodularity and the measures on the fibre are checked here.\par\smallskip
\section*{The determinant lower bound}
\noindent\hyperref[route:19]{\textbf{Endpoint of the Hecke product} (p.~\pageref*{route:19})}. Hecke's continuation theorem replaces the adelic test-function calculation. The angular type and a positive Euler-factor inequality give nonvanishing at one.\par\smallskip
\noindent\hyperref[route:20]{\textbf{Rank-one Galois avatars} (p.~\pageref*{route:20})}. Conrad's rank-one criterion determines the crystalline periods of the local avatar. The unramified Frobenius convention and character twists are checked explicitly.\par\smallskip
\noindent\hyperref[route:21]{\textbf{Global Selmer relation count} (p.~\pageref*{route:21})}. Greenberg--Wiles replaces the layered unit-lattice count. The dual local conditions, the universal relation and restriction at the auxiliary primes still have to be checked.\par\smallskip
\noindent\hyperref[route:22]{\textbf{Valuations of the Jacobi sums} (p.~\pageref*{route:22})}. Gross--Koblitz gives the degree-one Gauss-sum valuations after a Teichm\"uller normalization. The sign, auxiliary reduction and the prime two are treated separately.\par\smallskip
\clearpage
\section*{Galois representations and integral comparison}
\noindent\hyperref[route:23]{\textbf{Local crystalline cocycle lengths} (p.~\pageref*{route:23})}. Bloch--Kato and local duality give the rational dimension after the local condition is identified. The integral lattice loss is computed separately.\par\smallskip
\noindent\hyperref[route:24]{\textbf{Integral de Rham columns} (p.~\pageref*{route:24})}. A basis adapted to the filtration replaces the successive determinant-lowering procedure and gives the fixed integral change of columns used in the period matrix.\par\smallskip
\noindent\hyperref[route:25]{\textbf{Twisted invariants in the completed local field} (p.~\pageref*{route:25})}. The community formalization also invokes the two scalar theorems. The lifting of a fixed filtered class and the determination of its pure-weight coordinates are proved here.\par\smallskip
\noindent\hyperref[route:26]{\textbf{Transfer of residue-field labels} (p.~\pageref*{route:26})}. Shintani relates the character on the Frobenius coset to its value at the norm. Restriction of the tame parameter identifies the lifted label.\par\smallskip
\noindent\hyperref[route:27]{\textbf{Split Iwahori blocks} (p.~\pageref*{route:27})}. The ordering and the spherical-vector argument follow the formal proof. The classical Jacquet-module comparisons identify the normalization of the exponents.\par\smallskip
\noindent\hyperref[route:28]{\textbf{Realizing a limiting trace} (p.~\pageref*{route:28})}. The characteristic-zero pseudocharacter theorem replaces the matrix-algebra construction. Continuity and the assertions at the prescribed Frobenius elements remain separate steps.\par\smallskip
\noindent\hyperref[route:32]{\textbf{A bounded basis for signed frames} (p.~\pageref*{route:32})}. A definite place gives a lower bound on the first minimum. Minkowski's second theorem then supplies the short columns used to enumerate the Gram matrices.\par\smallskip
\noindent\hyperref[route:33]{\textbf{The theta representation over the complex numbers} (p.~\pageref*{route:33})}. Mumford's dimension and irreducibility theorems give the complex input. The argument establishing the integral orbit span follows in the text.\par\smallskip
\noindent\hyperref[route:29]{\textbf{The generic correspondence cover} (p.~\pageref*{route:29})}. The quotient by the free finite action gives the generic torsor. Finiteness, reducedness, and the identification of its quotient with the specified base are checked.\par\smallskip
\noindent\hyperref[route:30]{\textbf{The period-ring map on valuation tests} (p.~\pageref*{route:30})}. The perfectoid criterion of Bhatt--Morrow--Scholze gives the kernel of the period-ring map after its hypotheses are verified on the product of valuation rings.\par\smallskip
\noindent\hyperref[route:31]{\textbf{Fixed ratios of crystalline periods} (p.~\pageref*{route:31})}. Fontaine's fixed-field theorem replaces clearing poles in chart coordinates. The integral period bounds elsewhere in the proof are still required.\par\smallskip

\chapter{Formal source notes (optional reference)}\label{app:formal}
\noindent Each superscript in the text leads to a note below. The numbers restart in each chapter. Follow the number to return to the mathematical argument, or the description to open the Lean file. Each note gives the source folder before the first link and whenever the folder changes.

\noindent\textsf{[Lib]} marks the library of background results that the authors classify as known and that were missing from their version of Mathlib. It can be read independently of the GRH argument. \textsf{[New]} marks the setup and steps of the argument, including known theorems specialized to its data. \textsf{[src]} marks a file at the top level. Both folders use the Lean namespace \texttt{New}. The authors describe the folder names as a working classification and make no claim of historical priority on the basis of a name. See the \rootlink{grh/GUIDE.md\#4-conventions}{guide to the readable edition}.

For an individual theorem, the \rootlink{grh/INDEX.md}{index of results} gives a finer distinction. Under \emph{Known results} it includes the prime number theorem at the auxiliary cutoff, equation~\eqref{eq:an-fixed-modulus-pnt}, although its statement for the weighted prime sum is in \texttt{New}. Under \emph{Steps of the argument} it includes the \texttt{Lib} bound for the number of Galois cocycles with specified local trivializations. Section~\ref{sec:split-duality} gives the corresponding Selmer count using Greenberg--Wiles. The introductory note in each file distinguishes sources from related work. The \href{https://github.com/leanprover-community/riemann-hypothesis/blob/120962b4fc014406f9c76adc53b5cf64475143d3/REFERENCES.md}{community edition's reference index} lists papers by Lean file and records which theorem or page references its editors checked.
\begingroup\small\setlength{\parskip}{1.8pt}
\section*{Chapter 1. Introduction}
\noteline{1.1}{1}{\src{Statements.lean}{Dirichlet and zeta zero statements}. \src{Main.lean}{The final deductions in the readable edition}.}
\section*{Chapter 2. The arithmetic reduction}
\noteline{2.1}{1}{\src{New/Reduction/Data/AdmissibleData.lean\#L77-L129}{Scale, rank, and weight conditions}.}
\noteline{2.2}{2}{\src{New/Reduction/Data/Existence.lean}{Construction of the set of places}.}
\noteline{2.3}{3}{\src{Lib/Analytic/Primes.lean}{The prime-in-window theorem}.}
\noteline{2.4}{4}{\src{New/Reduction/Data/AdmissibleData.lean}{Definition of the level window}. \src{New/Reduction/Data/Existence.lean}{Existence at every large integer}.}
\noteline{2.5}{5}{\src{New/UpperBound/Objects/Columns.lean}{Definite chain and restriction}. \src{New/UpperBound/Objects/Rows.lean}{Target spaces, weights and good types}. \src{New/UpperBound/Objects/RowChain.lean}{Twisted target chain}. \src{New/UpperBound/Objects/PoincareMap.lean}{The period defining the Poincar\'e map}. \src{New/UpperBound/Size/ShortColumns.lean\#L599-L687}{Local scalars generating the residue-field units}.}
\noteline{2.6}{6}{\src{New/Reduction/PrimeSums.lean}{Bounds on both terms of the direction amount}.}
\noteline{2.7}{7}{\src{New/Reduction/MainInequality.lean}{The comparison on one family}.}
\section*{Chapter 3. Prime sums and zeros of Dirichlet L-functions}
\noteline{3.1}{1}{\src{New/Reduction/Zeros/PrimeError.lean}{The weighted sum and its poles}.}
\noteline{3.2}{2}{\src{New/Reduction/PrimeSums.lean}{The place count}. \src{New/Reduction/OneSidedBound.lean}{The three terms in the bound}.}
\noteline{3.3}{3}{\src{Lib/Analytic/PrimeNumberTheorem/ZeroFreeRegion.lean}{Zero-free region}. \src{Lib/Analytic/PrimeNumberTheorem/Tauberian.lean}{Gamma-kernel estimate}. \src{New/Reduction/Zeros/PrimeNumberTheorem.lean}{Application to the smoothed count}.}
\noteline{3.4}{4}{\src{Lib/Analytic/DirichletTheta.lean}{Mellin formula and theta transformation}.}
\noteline{3.5}{5}{\src{Lib/Analytic/QuadraticTwist.lean}{Three-lines zero detector}.}
\noteline{3.6}{6}{\src{Lib/Analytic/QuadraticTwist.lean}{Expansion of the Jacobi-symbol mean and the twist count}.}
\noteline{3.7}{7}{\src{Lib/Analytic/Landau.lean}{Singularity for positive measures}.}
\noteline{3.8}{8}{\src{Lib/Analytic/Landau.lean}{Application to meromorphic transforms with prescribed residue signs}.}
\noteline{3.9}{9}{\src{New/Reduction/Zeros/Criterion.lean}{The one-sided criterion}.}
\noteline{3.10}{10}{\src{New/Main.lean}{Reflection}. \src{Main.lean}{The five formal conclusions}.}
\section*{Chapter 4. Divisibility and complex multiplication}
\noteline{4.1}{1}{\src{New/Derivatives/FormalGroup/Exponent.lean}{One-curve exponent}. \src{New/Derivatives/FormalGroup/ThetaProducts.lean}{Passage through a resolution}. \src{New/Derivatives/Vacuum/Divisibility.lean}{Normal form for the actual data}.}
\noteline{4.2}{2}{\src{New/Derivatives/Divisibility/NullFrames.lean}{Simultaneous changes of direction}.}
\noteline{4.3}{3}{\src{New/Derivatives/Divisibility/Count.lean}{The uncharged and charged minor estimates}. \src{New/Derivatives/Vacuum/Divisibility.lean}{From minors to Schur derivatives}.}
\noteline{4.4}{4}{\src{New/UpperBound/Divisibility/JetExponent.lean}{Proof of the stage inequality (A1)}.}
\noteline{4.5}{5}{\src{New/UpperBound/Size/Estimates.lean}{Grouping the prime powers}. \src{New/Reduction/PrimeSums.lean}{The two-cutoff identity}.}
\noteline{4.6}{6}{\src{New/UpperBound/Divisibility/ValuationAtNonSplitPrimes.lean}{Hecke transfer through the lifts}. \src{New/UpperBound/Divisibility/LocalDivisibility.lean}{Valuation on the marked image}.}
\noteline{4.7}{7}{\src{New/UpperBound/ElementaryDivisors/ElementaryDivisorTheorem.lean}{Elementary divisors on edges}. \src{New/UpperBound/Divisibility/MarkOperator.lean}{Application to the marked orbit}.}
\noteline{4.8}{8}{\src{New/UpperBound/Divisibility/ValuationAtLevel.lean}{Valuation at the level prime}. \src{New/UpperBound/Divisibility/ColengthAtLevel.lean}{Colength at the level prime}.}
\noteline{4.9}{9}{\src{New/UpperBound/Size/Estimates.lean}{Summation of the losses}. \src{New/UpperBound/Divisibility/LocalDivisibility.lean}{The uniform local theorem}.}
\noteline{4.10}{10}{\src{New/UpperBound/Objects/PoincareMap.lean}{The defining Bessel period}. \src{New/UpperBound/Seesaw/SeesawIdentity.lean}{Unfolding and the controlled chain datum}.}
\noteline{4.11}{11}{\src{New/Reduction/ShortAndEnough.lean}{The maps, the floor and the scale}.}
\noteline{4.12}{12}{\src{New/UpperBound/Kernels/JetConjugatesAndBounds.lean}{Conjugates of the actual jets}. \src{New/UpperBound/Size/ArchimedeanSize.lean}{Multiplication of the size bounds}.}
\noteline{4.13}{13}{\src{New/UpperBound/Size/ShortColumns.lean}{Short spanning vectors and ramification}.}
\noteline{4.14}{14}{\src{New/Reduction/ShortAndEnough.lean}{Theorem A(a), \texttt{span\_\allowbreak{}short\_\allowbreak{}integralRamified\_\allowbreak{}eq\_\allowbreak{}top}}.}
\noteline{4.15}{15}{\src{New/Reduction/Lattices/TwistedLattice.lean}{The determinant definition of index}. \src{New/Reduction/ShortAndEnough.lean}{Theorem A(b), \texttt{indexLe\_\allowbreak{}normalizedPoincare}}.}
\noteline{4.16}{16}{\src{Lib/CM/GoodReduction.lean}{Degree-two isogenies, ideal classes and the unit discriminant}.}
\noteline{4.17}{17}{\comsrc{Lib/CM/Integrality/Doubling.lean\#L127-L176}{Duplication formulas}. \comsrc{Lib/CM/Integrality/TorsionValues.lean\#L183-L267}{Local height and monomial bounds}. \comsrc{Lib/CM/Integrality/TorsionValues.lean\#L306-L451}{Isogenies, torsion translation and integral values}.}
\noteline{4.18}{18}{\comsrc{Lib/CM/Integrality/IntegralSections.lean\#L66-L127}{Polynomial coordinates and even sections}. \comsrc{Lib/CM/Integrality/IntegralSections.lean\#L189-L354}{An integral point attaining the coefficient norms and the integral basis}.}
\noteline{4.19}{19}{\src{Lib/CM/Sigma/ThetaSections.lean}{Period multiplier and the transformed sections}. \comsrc{Lib/CM/Integrality/IsogenyIntegrality.lean\#L50-L115}{Pullbacks and integral translated coordinates}.}
\noteline{4.20}{20}{\src{Lib/CM/Integrality/LogTheta.lean}{Logarithmic second derivative, norm pairs and the quadratic coefficient}. \src{Lib/General/NumberTheory/PadicLimitsValuation.lean}{Norm pairs in the quadratic completion}.}
\noteline{4.21}{21}{\comsrc{Lib/CM/FormalGroup/InseparableDegree.lean\#L193-L255}{A unit coefficient of the even section}. \comsrc{Lib/CM/FormalGroup/InseparableDegree.lean}{Inseparable degree in all three prime types}.}
\noteline{4.22}{22}{\src{Lib/CM/FormalGroup/FormalParameter.lean}{Integral generator, its unit coefficient and scaled series}.}
\noteline{4.23}{23}{\src{Lib/CM/Integrality/LogTheta.lean}{Euler addition in integral coordinates}.}
\noteline{4.24}{24}{\src{New/Derivatives/Sharpness/LeastValuation.lean}{Coefficient pullback and six-copy units}. \src{New/Derivatives/Sharpness/Couplings.lean}{Derivative in the second group of variables}.}
\noteline{4.25}{25}{\src{Lib/CM/ChowlaSelberg.lean}{Choice of ideals, conjugation, product of periods and the limit calculation}. \src{Lib/CM/KroneckerLimit.lean}{Kronecker formula with the stated scale}.}
\noteline{4.26}{26}{\src{Lib/CM/Sigma/Basic.lean\#L400-L556}{Construction of the section and its addition identity}.}
\noteline{4.27}{27}{\src{Lib/CM/Sigma/Basic.lean\#L509-L556}{Addition by evaluation in degree two}. \src{Lib/CM/Sigma/Basic.lean\#L606-L618}{The derivative at the diagonal}.}
\noteline{4.28}{28}{\src{Lib/CM/Sigma/ThetaSections.lean\#L160-L211}{Parity, the square identity and the invariants}. \src{Lib/CM/Sigma/ThetaSections.lean\#L243-L278}{Translated addition formula}. \src{Lib/CM/Integrality/Doubling.lean\#L132-L180}{Both duplication formulas in the integral coordinates}.}
\noteline{4.29}{29}{\src{Lib/CM/Sigma/ThetaSections.lean\#L716-L797}{Spanning by degrees and division by the theta section}.}
\noteline{4.30}{30}{\src{Lib/CM/ChowlaSelberg.lean\#L84-L158}{The prime discriminant and injectivity of squaring}.}
\noteline{4.31}{31}{\src{Lib/CM/ChowlaSelberg.lean\#L254-L268}{Absorbing a principal factor}. \src{Lib/CM/ChowlaSelberg.lean\#L354-L395}{Period, area and denominator}.}
\noteline{4.32}{32}{\src{Lib/CM/ConjugateCurve.lean\#L62-L122}{The formal recursion and the factor by degree}.}
\noteline{4.33}{33}{\src{Lib/CM/ConjugateCurve.lean\#L123-L222}{Matching points on both charts}.}
\noteline{4.34}{34}{\src{Lib/CM/ConjugateCurve.lean\#L259-L335}{The two transported sections and the tests at a period and at the half-period}.}
\noteline{4.35}{35}{\src{Lib/CM/ConjugateCurve.lean\#L349-L394}{The conjugate marked lattice}. \src{Lib/CM/ChowlaSelberg.lean\#L505-L535}{The inverse-ideal argument}.}
\noteline{4.36}{36}{\src{Lib/CM/ConjugateCurve.lean\#L425-L518}{Rescaling and comparison of the first nonconstant coefficient}.}
\noteline{4.37}{37}{\src{Lib/CM/ConjugateCurve.lean\#L527-L598}{Period shifts and their multipliers}. \src{Lib/CM/ConjugateCurve.lean\#L599-L695}{Comparison at the half-shift, including a vanishing value}. \src{Lib/CM/ConjugateCurve.lean\#L696-L764}{Simultaneous finite-set construction}.}
\noteline{4.38}{38}{\src{Lib/CM/ChowlaSelberg.lean\#L270-L327}{Compactness and the explicit congruence for nearby rational shifts}.}
\noteline{4.39}{39}{\src{Lib/CM/Sigma/ThetaSections.lean\#L109-L156}{The marked shift, lattice sign and semicharacter}.}
\noteline{4.40}{40}{\src{Lib/CM/Sigma/ThetaSections.lean\#L440-L462}{The two trivial signs}. \src{Lib/CM/Sigma/ThetaSections.lean\#L484-L527}{Evaluation of the remaining sign}.}
\noteline{4.41}{41}{\src{Lib/CM/Sigma/ThetaSections.lean\#L530-L566}{The isogeny law}. \src{Lib/CM/Sigma/ThetaSections.lean\#L821-L849}{The commutator for torsion}. \comsrc{Lib/CM/Integrality/IsogenyIntegrality.lean\#L82-L159}{Integral pullbacks, the linear coefficient and torsion shifts}.}
\noteline{4.42}{42}{\src{New/Derivatives/Vacuum/IntegralVacuum.lean\#L76-L153}{Finite Fourier projection, trace coordinates and the product of jets}. \src{New/Derivatives/Vacuum/IntegralVacuum.lean\#L160-L209}{Construction on a normalized CM lattice}. \src{New/Derivatives/Vacuum/IntegralVacuum.lean\#L493-L510}{Generation by the finite translates}.}
\noteline{4.43}{43}{\src{New/Derivatives/Vacuum/IntegralVacuum.lean\#L451-L493}{Evaluation of the absolute scale}. \src{New/Derivatives/Vacuum/IntegralVacuum.lean\#L510-L532}{The place permutation on the linear forms}.}
\noteline{4.44}{44}{\src{New/Derivatives/Vacuum/IntegralVacuum.lean\#L536-L582}{Simultaneous rational approximation and substitution}. \src{New/Derivatives/Vacuum/IntegralVacuum.lean\#L595-L649}{The common multiplier in all ranks}.}
\noteline{4.45}{45}{\src{New/Derivatives/Vacuum/TightFrames.lean\#L509-L523}{Faithfulness of the entire jets on the finite part}. \src{New/UpperBound/Kernels/JetConjugatesAndBounds.lean\#L149-L184}{Construction of the semilinear map}.}
\noteline{4.46}{46}{\src{New/UpperBound/Kernels/JetConjugatesAndBounds.lean\#L90-L140}{Recovering the additive character from rank one}. \src{New/UpperBound/Kernels/JetConjugatesAndBounds.lean\#L184-L209}{Intertwining all translates}. \src{New/UpperBound/Kernels/JetConjugatesAndBounds.lean\#L232-L250}{Finite evaluations as constant coefficients}.}
\noteline{4.47}{47}{\src{Lib/Weil/DualPair/GaloisLaw.lean\#L99-L135}{Local units from self-duality}. \src{Lib/Weil/DualPair/GaloisLaw.lean\#L224-L278}{Transport of the normalized translations and the vacuum}.}
\noteline{4.48}{48}{\src{Lib/Weil/DualPair/GaloisLaw.lean\#L58-L96}{The multiplier quotient and its square}. \src{Lib/Weil/DualPair/GaloisLaw.lean\#L279-L312}{Application to the finite projection sum}. \src{Lib/Weil/DualPair/GaloisLaw.lean\#L369-L385}{Removal of the central unitary scalar}. \src{Lib/ClassField/HeckeCharacter/AlgebraicValues.lean\#L362-L380}{Archimedean norms of conjugate character values}.}
\noteline{4.49}{49}{\src{Lib/CM/KroneckerLimit.lean\#L122-L386}{Row subtraction, constant term and the derivative at the prescribed scale}. \src{Lib/CM/Sigma/Basic.lean\#L815-L857}{The lattice discriminant with periods \(a,a\tau\)}.}
\noteline{4.50}{50}{\src{Lib/CM/ChowlaSelberg.lean\#L162-L217}{Euler factors by splitting type}. \src{Lib/CM/ChowlaSelberg.lean\#L541-L630}{The logarithmic derivative after summing ideal classes}.}
\noteline{4.51}{51}{\comsrc{Lib/CM/Integrality/LogTheta.lean}{CM logarithm}. \comsrc{New/Derivatives/FormalGroup/FrobeniusCongruence.lean\#L147-L169}{The even section in the Frobenius congruence}.}
\noteline{4.52}{52}{\src{Lib/CM/FormalGroup/InseparableDegree.lean}{The formal-group height input}. \src{Lib/CM/FormalGroup/FormalParameter.lean}{Construction of the scaled parameter}.}
\noteline{4.53}{53}{\src{New/Derivatives/FormalGroup/Exponent.lean}{Functional equation and its solution}.}
\noteline{4.54}{54}{\src{New/Derivatives/Sharpness/OneCurve.lean}{Sharp one-curve rate}.}
\noteline{4.55}{55}{\src{New/Derivatives/Sharpness/Cases.lean}{Twelve-copy case estimates}. \src{New/Derivatives/Sharpness/NearPairs.lean}{Sparsity and the losses}. \src{New/Derivatives/Sharpness/BoundAttained.lean}{Fourier inversion and the common scalar}.}
\section*{Chapter 5. Integral frames and the definite theta lift}
\noteline{5.1}{1}{\src{New/Derivatives/Field/PrimeDecomposition.lean}{Integral generators and the relative different}. \src{New/Derivatives/Field/WeilData.lean}{The components of the lattice idele}.}
\noteline{5.2}{2}{\src{New/Derivatives/Vacuum/ShortVectors.lean}{Orbit resolutions, pair replacements and the retained member}.}
\noteline{5.3}{3}{\src{Lib/ClassField/HeckeCharacter/AlgebraicValues.lean}{Finite norms of the algebraic Hecke character}.}
\noteline{5.4}{4}{\src{New/Derivatives/Vacuum/TightFrames.lean}{The local period norm from the vacuum root}.}
\noteline{5.5}{5}{\src{New/Derivatives/FormalGroup/ThetaProducts.lean}{The ring congruence and its action on the frame}.}
\noteline{5.6}{6}{\src{New/Derivatives/Sharpness/NearPairs.lean}{Three-point obstruction and incidence at the ends of near pairs}.}
\noteline{5.7}{7}{\src{New/Derivatives/Sharpness/NearPairs.lean}{Cutoff for the tail and the consecutive-integer loss}. \src{New/Derivatives/Sharpness/Cases.lean}{The small-prime allocation and the cost of multinomials}.}
\noteline{5.8}{8}{\src{New/Derivatives/Sharpness/Cases.lean}{The large-prime factorial recurrence and the split case}.}
\noteline{5.9}{9}{\src{New/Derivatives/Vacuum/ShortVectors.lean}{The two coprime frame lengths and unit member}.}
\noteline{5.10}{10}{\src{New/Derivatives/FormalGroup/ThetaProducts.lean}{Frobenius on the frame and the product normal form}. \src{New/Derivatives/Sharpness/Cases.lean}{The period rescaling at an inert place}.}
\noteline{5.11}{11}{\src{New/Derivatives/Sharpness/LeastValuation.lean}{The six scaled theta products and their inverses}.}
\noteline{5.12}{12}{\src{New/Derivatives/Sharpness/Couplings.lean}{Differentiation, unique lowest term and the exact coefficient norm}.}
\noteline{5.13}{13}{\src{New/Derivatives/Sharpness/Cases.lean}{Degree inequalities and the cases giving the twelve-copy bound}. \src{New/Derivatives/Sharpness/NearPairs.lean}{Near-pair uniqueness, ranges and arithmetic losses}.}
\noteline{5.14}{14}{\src{New/Derivatives/Vacuum/TightFrames.lean}{Pullback of the monomial sections, signed periods and integral finite values}. \src{New/Derivatives/Sharpness/BoundAttained.lean}{The twelve-copy Taylor identity at the prescribed scale}.}
\noteline{5.15}{15}{\src{Lib/Weil/Construction/FourierInversion.lean}{The commutator pairing, both translation sums and their integrality}.}
\noteline{5.16}{16}{\src{Lib/General/Algebra/LevelModulus.lean}{A finite vacuum translate in the basis of residue indicators}.}
\noteline{5.17}{17}{\src{New/Derivatives/Sharpness/BoundAttained.lean}{Coefficient inequality and the height of the common scalar}.}
\noteline{5.18}{18}{\src{New/Derivatives/Field/PrimeDecomposition.lean}{Frobenius and the decomposition groups}. \src{New/UpperBound/Divisibility/MarkedPrimes.lean}{Exceptional orders and the marked places}. \src{New/UpperBound/Divisibility/MarkOperator.lean}{Cycle numbering, letter embeddings and operator scale}.}
\noteline{5.19}{19}{\src{New/UpperBound/Divisibility/HeckeEdges.lean}{Block scaling, edge steps and intertwining with the global scaled operator}.}
\noteline{5.20}{20}{\src{New/UpperBound/ElementaryDivisors/InvariantAverage.lean}{The integral trace construction}.}
\noteline{5.21}{21}{\src{New/UpperBound/ElementaryDivisors/FrobeniusFactors.lean}{The weight bound on each surviving full monomial}. \src{New/UpperBound/ElementaryDivisors/Edges.lean}{Adjoint root estimate on all-edge forms}.}
\noteline{5.22}{22}{\src{New/UpperBound/ElementaryDivisors/AveragingMap.lean}{Frobenius contraction, frame change and the pairing congruence}.}
\noteline{5.23}{23}{\src{New/UpperBound/ElementaryDivisors/FrobeniusFactors.lean}{One-factor congruence, telescoping weights and the full-monomial bound}.}
\noteline{5.24}{24}{\src{New/UpperBound/ElementaryDivisors/RootLemma.lean}{Improving deep degree by one third of the defect}. \src{New/UpperBound/ElementaryDivisors/Edges.lean}{Passage to the lattice and its adjoint pairing}.}
\noteline{5.25}{25}{\src{New/UpperBound/ElementaryDivisors/Depth.lean}{The unipotent interpolation points and their equations}. \src{Lib/General/Algebra/MultivariatePolynomials.lean}{Vanishing ideals, powers and weighted monomials}.}
\noteline{5.26}{26}{\src{New/UpperBound/ElementaryDivisors/DefectCount.lean}{The least-height argument, packing and integral translation}.}
\noteline{5.27}{27}{\src{New/UpperBound/ElementaryDivisors/PieriCone.lean}{Interlacing indicator, signed change of variables and one-letter evaluation}.}
\noteline{5.28}{28}{\src{New/UpperBound/ElementaryDivisors/DefectCount.lean}{The flag indices and final dimension estimate}.}
\noteline{5.29}{29}{\src{New/UpperBound/ElementaryDivisors/RootLemma.lean}{The two substitutions, the top coefficient and Stepanov's degree comparison}.}
\noteline{5.30}{30}{\src{New/UpperBound/Objects/LiftChains.lean}{The principal-level norm bound and freeness}. \src{Lib/Adelic/AutomorphicForms/PAdicGauge.lean}{Weak approximation and integral representatives}.}
\noteline{5.31}{31}{\src{New/UpperBound/Divisibility/HeckeEdges.lean}{Rational edge factors, uniqueness at the origin, and the global coset sum}.}
\noteline{5.32}{32}{\src{Lib/GLn/SchurModule/SchurMatrix.lean}{The one-row Schur matrix and the reversed monomial indices}. \src{New/UpperBound/Divisibility/HeckeEdges.lean}{The two letter parities and integrality of the multiplicity action}.}
\noteline{5.33}{33}{\src{Lib/Adelic/AutomorphicForms/HeckeEigenforms.lean}{Evaluation at free finite double cosets}. \src{New/UpperBound/ElementaryDivisors/ElementaryDivisorTheorem.lean}{The edge operator and transport of the colength bound}.}
\noteline{5.34}{34}{\src{New/UpperBound/Divisibility/MarkOperator.lean}{The cycle and its size}. \src{New/UpperBound/Divisibility/HeckeEdges.lean}{Identifying the scaled operator}.}
\noteline{5.35}{35}{\src{New/UpperBound/ElementaryDivisors/Depth.lean}{Interpolation and the depth count}. \src{New/UpperBound/ElementaryDivisors/Edges.lean}{The integral lattice factorization}.}
\noteline{5.36}{36}{\src{New/UpperBound/ElementaryDivisors/RootLemma.lean}{The auxiliary polynomial and its zeros}.}
\noteline{5.37}{37}{\src{New/UpperBound/ElementaryDivisors/DefectCount.lean}{Packing and the dimension inequality}. \src{New/UpperBound/ElementaryDivisors/PieriCone.lean}{The exact relative volume}.}
\noteline{5.38}{38}{\src{New/UpperBound/ElementaryDivisors/AveragingMap.lean}{The integral averaging theorem}.}
\noteline{5.39}{39}{\src{New/UpperBound/ElementaryDivisors/ElementaryDivisorTheorem.lean}{The determinant argument}. \src{New/UpperBound/Divisibility/MarkOperator.lean}{Substitution of the admissible scales}.}
\noteline{5.40}{40}{\src{New/UpperBound/Objects/Rows.lean}{Lower levels and the theta datum}. \src{New/UpperBound/Objects/Columns.lean}{Upper levels, branching and the chain tensor}. \src{New/UpperBound/Seesaw/EdgePairEvaluation.lean}{The gradings and the exceptional last pair}.}
\noteline{5.41}{41}{\src{Lib/Weil/Periods/Unfolding.lean}{Surviving rational classes and the one-step integral}. \src{New/UpperBound/Seesaw/UnfoldingStages.lean}{The translated stages and their recursive unfolding}. \src{New/UpperBound/Seesaw/SeesawIdentity.lean}{Transpose, rapidity and period identities}.}
\noteline{5.42}{42}{\src{New/UpperBound/Seesaw/UnfoldingStages.lean}{Nonzero stages satisfy the graded lattice bounds}. \src{New/UpperBound/Seesaw/Sandwich.lean}{Propagation of upper and lower graded lattice bounds}. \src{Lib/Weil/Periods/LatticeBounds.lean}{Finite representatives for the bounded lattice orbits}. \src{New/UpperBound/Seesaw/Cut.lean}{Construction of a finite cut on the support}.}
\noteline{5.43}{43}{\src{New/UpperBound/Seesaw/IntegralAsCosetSum.lean}{Stabilizer indices, Haar densities and the finite period}. \src{New/UpperBound/Seesaw/SeesawIdentity.lean}{Induction over the outer integrals}.}
\noteline{5.44}{44}{\src{New/UpperBound/Seesaw/UnfoldingStages.lean}{Products of constant terms, unipotent means and the triple theta integral}.}
\noteline{5.45}{45}{\src{New/UpperBound/Seesaw/EdgePairEvaluation.lean}{The two explicit splittings and their finite evaluations}. \src{New/UpperBound/Seesaw/UnfoldingStages.lean}{Induction identifying each even stage with a definite matrix coefficient}.}
\noteline{5.46}{46}{\src{New/UpperBound/Seesaw/SeesawIdentity.lean}{Top Schur factor, equality by pairing and joint level invariance}. \comsrc{Lib/GLn/Branching/NormBound.lean\#L381-L425}{The hook-number bound for the branching vector via an integral left inverse}. \comsrc{Lib/GLn/Branching/IntegralMap.lean\#L578-L603}{Use of the bound in the branching map}.}
\noteline{5.47}{47}{\src{New/UpperBound/Seesaw/AdmissibleEvaluation.lean}{Integral lattice lifts, neighbour lattices and preservation of denominators}. \src{New/UpperBound/Seesaw/GoodDatum.lean}{Fourier coefficients of the level datum and finite evaluation}.}
\noteline{5.48}{48}{\src{New/UpperBound/Seesaw/Cut.lean}{Number of untwisted cut points}. \src{New/UpperBound/Seesaw/GoodDatum.lean}{Pairwise denominators and the integral-vacuum scalar}. \src{New/UpperBound/Seesaw/SeesawIdentity.lean}{Final normalization and supported translations}.}
\noteline{5.49}{49}{\src{New/UpperBound/Objects/Columns.lean}{Finite gluing as a sum of Schrödinger data}. \src{New/UpperBound/Divisibility/ValuationAtLevel.lean}{Marked coordinate changes and their two determinant factors}.}
\noteline{5.50}{50}{\src{New/UpperBound/Objects/LiftLawsModuloN.lean}{Primitive columns, mirror data and the two formulas for the marked operator}.}
\noteline{5.51}{51}{\src{New/UpperBound/Divisibility/ValuationAtLevel.lean}{Iteration of the explicit local scalar and its valuation}.}
\noteline{5.52}{52}{\src{Lib/Weil/Construction/HeisenbergModel.lean\#L266-L306}{The integrable model and its defining kernel}. \src{Lib/Weil/Construction/AdelicRepresentation.lean\#L321-L352}{The adelic vacuum and measure}.}
\noteline{5.53}{53}{\src{Lib/Weil/Construction/HeisenbergModel.lean\#L607-L631}{Evaluation and superposition for a bounded operator}.}
\noteline{5.54}{54}{\src{New/UpperBound/Seesaw/EdgePairEvaluation.lean\#L770-L844}{The string, parity and equivariance hypotheses}.}
\noteline{5.55}{55}{\src{Lib/Weil/DualPair/Defs.lean\#L147-L184}{Rank twists, negation and the pair character}. \src{Lib/Weil/DualPair/TwoEdges.lean\#L70-L86}{The adjoint flip and the three determinants}. \src{Lib/Weil/DualPair/TwoEdges.lean\#L189-L209}{The product of the two actions on joined data}.}
\noteline{5.56}{56}{\src{Lib/Weil/MixedModel/ThetaInvariance.lean\#L231-L262}{Uniqueness from the Heisenberg law}. \src{Lib/Weil/MixedModel/ThetaInvariance.lean\#L305-L390}{The slice Jacobian and the phase of the hyperbolic multiplier}. \src{Lib/Weil/MixedModel/ThetaInvariance.lean\#L405-L443}{The radical and the full parabolic formula}.}
\noteline{5.57}{57}{\src{Lib/Weil/Periods/ConstantTerm.lean\#L167-L205}{Determinant on the positive-weight frame}. \src{New/UpperBound/Seesaw/EdgePairEvaluation.lean\#L877-L979}{The two determinant factors and their cancellation}. \src{Lib/Weil/MixedModel/ZeroTerm.lean\#L211-L230}{Composition of zero terms}.}
\noteline{5.58}{58}{\src{Lib/Weil/DualPair/TwoEdges.lean\#L124-L151}{The global functional in the project's model and the effect of the flip}. \src{Lib/Weil/Construction/CanonicalTheta.lean\#L50-L85}{Change of coordinates for the mean-normalized theta functional}.}
\noteline{5.59}{59}{\src{New/UpperBound/Seesaw/EdgePairEvaluation.lean\#L785-L835}{The grade-one entries in the prescribed splitting}. \src{New/UpperBound/Seesaw/UnfoldingStages.lean\#L307-L363}{The degree-zero displacement and the vanishing Gram matrix}.}
\noteline{5.60}{60}{\src{New/UpperBound/Seesaw/UnfoldingStages.lean\#L388-L443}{The outer character and both Heisenberg phases}. \src{New/UpperBound/Seesaw/UnfoldingStages.lean\#L457-L551}{The two unipotent averages}.}
\noteline{5.61}{61}{\src{New/UpperBound/Seesaw/UnfoldingStages.lean\#L128-L237}{The closed finite-dimensional space of families}. \src{New/UpperBound/Seesaw/UnfoldingStages.lean\#L716-L851}{Separation of the real vector and iteration on the pairs}.}
\noteline{5.62}{62}{\src{New/UpperBound/Seesaw/AdmissibleCombinations.lean\#L445-L488}{The two unit tests}. \src{New/UpperBound/Seesaw/AdmissibleCombinations.lean\#L727-L786}{The lattice coset expansion}. \src{New/UpperBound/Seesaw/GoodDatum.lean\#L237-L261}{The identity using the inverse multiplier}. \src{New/UpperBound/Seesaw/GoodDatum.lean\#L370-L411}{The determinant character of a moved arrow}.}
\noteline{5.63}{63}{\src{Lib/Weil/DualPair/TwoEdges.lean\#L156-L186}{The finite semicharacter correction}. \src{New/UpperBound/Seesaw/GoodDatum.lean\#L267-L310}{The two clearing factors}. \src{New/UpperBound/Seesaw/GoodDatum.lean\#L413-L448}{The three matrix counts}. \src{New/UpperBound/Seesaw/SeesawIdentity.lean\#L584-L654}{Denominator clearing with a common normalization}.}
\noteline{5.64}{64}{\src{Lib/Adelic/AutomorphicForms/SmoothAlgebraicForms.lean\#L120-L135}{The conjugate embedding in the definition of the gauge}. \src{Lib/Adelic/AutomorphicForms/PAdicGauge.lean\#L84-L163}{Weak approximation and both changes of frame}. \src{Lib/Adelic/AutomorphicForms/PAdicGauge.lean\#L187-L203}{Testing the bound in integral coordinates}.}
\noteline{5.65}{65}{\src{Lib/Weil/Periods/BesselIntegral.lean\#L528-L610}{Survival, the Eichler transformation and the Gram-entry test}.}
\noteline{5.66}{66}{\src{Lib/Weil/Periods/BesselIntegral.lean\#L612-L817}{String pivot, stabilizer, row induction and Witt extension}.}
\noteline{5.67}{67}{\src{Lib/Adelic/QuotientMeasure/UnfoldingInStages.lean\#L361-L393}{The quotient mean and the twist argument}.}
\noteline{5.68}{68}{\src{Lib/Weil/Periods/ConstantTerm.lean\#L583-L626}{The two expansions of the same Fourier coefficient}.}
\noteline{5.69}{69}{\src{Lib/Weil/Periods/Unfolding.lean\#L163-L224}{Absolute constant-term bounds}. \src{Lib/Weil/Periods/Unfolding.lean\#L343-L532}{The two orbit sums and the rapid-decay estimate}.}
\noteline{5.70}{70}{\src{Lib/Adelic/Adeles/Basic.lean\#L463-L483}{The integral element with unit difference}. \src{Lib/Adelic/Witt/WittTheorem.lean\#L506-L534}{Witt extension for primitive frames}.}
\noteline{5.71}{71}{\src{Lib/Adelic/Witt/WittTheorem.lean\#L485-L505}{Construction from two integral left inverses}.}
\noteline{5.72}{72}{\src{Lib/Adelic/Witt/WittTheorem.lean\#L442-L486}{Correction to an invertible matrix by residue-field reduction}.}
\noteline{5.73}{73}{\src{Lib/GaloisRep/Hilbert90/UnitaryGroups.lean\#L267-L296}{The commutative algebra of the characteristic polynomial}. \src{Lib/GaloisRep/Hilbert90/UnitaryGroups.lean\#L342-L368}{From an involution norm to the relative frame}.}
\noteline{5.74}{74}{\src{Lib/Adelic/Witt/WittTheorem.lean\#L313-L365}{The two integral orthonormalization statements}.}
\noteline{5.75}{75}{\src{Lib/Weil/Periods/PointLattice.lean\#L294-L325}{A primitive basis of the intersection}. \src{Lib/Weil/Periods/PointLattice.lean\#L425-L467}{Extension and the pointwise stabilizer of the string}.}
\noteline{5.76}{76}{\src{New/UpperBound/Seesaw/StandardLines.lean}{Integral string lines and the shift onto the shorter intersection}.}
\noteline{5.77}{77}{\src{New/UpperBound/Seesaw/Frequency.lean}{Support of the finite constant term and extraction of degree minus one}. \src{New/UpperBound/Seesaw/EdgePairEvaluation.lean}{The integral left inverse and the Heisenberg commutator}.}
\noteline{5.78}{78}{\src{New/UpperBound/Seesaw/Sandwich.lean}{Duality, the one-dimensional quotient, and the exceptional split}. \src{New/UpperBound/Seesaw/UnfoldingStages.lean}{The ramification scalar annihilates the exceptional stage}.}
\noteline{5.79}{79}{\src{Lib/Weil/Periods/GradedLattice.lean}{Adapted integral bases and unipotent Gram reduction}. \src{Lib/Weil/Periods/PointLattice.lean}{Integral Witt extension on a common primitive string}. \src{Lib/Weil/Periods/LatticeBounds.lean}{Finite intermediate submodules and truncation to the conductor}.}
\noteline{5.80}{80}{\src{New/UpperBound/Seesaw/NeighbourLattice.lean}{The rank-one unitary, its coset representatives, and their uniqueness}.}
\noteline{5.81}{81}{\src{New/UpperBound/Seesaw/AdmissibleEvaluation.lean}{The lift on both pairs and its adjoint on the shift}. \src{New/UpperBound/Seesaw/GoodDatum.lean}{The scale identity and the auxiliary level ideal}.}
\noteline{5.82}{82}{\src{New/UpperBound/Seesaw/EdgePairEvaluation.lean}{Nonvanishing of the vacuum zero term and the translated-vacuum formula}. \src{Lib/Weil/Construction/VacuumProjections.lean}{Lattice averaging and the absolute square of the finite multiplier}.}
\noteline{5.83}{83}{\src{New/UpperBound/Seesaw/NeighbourVacuum.lean}{The paired residue shifts, roots of unity, and the normalized average}.}
\noteline{5.84}{84}{\src{Lib/Weil/Construction/VacuumProjections.lean\#L359-L449}{The corrected translations, self-adjoint average and vacuum coefficient}. \src{Lib/Weil/Construction/VacuumProjections.lean\#L480-L563}{Projection on a translate and orthogonality}. \src{Lib/Weil/Construction/VacuumProjections.lean\#L802-L845}{The finite-part basis}.}
\noteline{5.85}{85}{\src{Lib/Weil/Construction/FiniteMultiplier.lean\#L695-L718}{Elementary divisors and the finite character}. \src{Lib/Weil/Construction/FiniteMultiplier.lean\#L892-L936}{The intersection lattice and its index}. \src{Lib/Weil/Construction/WeilOperators.lean\#L291-L316}{Vacuum normalization at zero}. \src{Lib/Weil/Construction/VacuumProjections.lean\#L713-L781}{Conjugating the projection and summing its cosets}.}
\noteline{5.86}{86}{\src{Lib/Weil/MixedModel/ZeroTerm.lean\#L66-L181}{Boundedness, the quotient mean and its translation law}. \src{Lib/Weil/MixedModel/ZeroTerm.lean\#L212-L239}{Successive zero terms}.}
\noteline{5.87}{87}{\src{Lib/Weil/MixedModel/RealTheta.lean\#L336-L396}{The normalized Gaussian integral and real projection}. \src{Lib/Weil/MixedModel/ZeroTerm.lean\#L435-L485}{The two finite cases}. \src{Lib/Weil/MixedModel/ZeroTerm.lean\#L531-L583}{Combining the projections}.}
\noteline{5.88}{88}{\src{Lib/Weil/Construction/WeilOperators.lean\#L414-L502}{The local average and bounded overlap estimate}. \src{Lib/Weil/MixedModel/ZeroTerm.lean\#L681-L746}{Absolute summability and Fourier inversion}. \src{Lib/Weil/MixedModel/ZeroTerm.lean\#L748-L797}{Nonvanishing and the translated-vacuum formula}.}
\noteline{5.89}{89}{\src{New/Derivatives/Divisibility/NullFrames.lean\#L34-L107}{Elimination and moving a congruence across Frobenius}. \src{New/Derivatives/Divisibility/NullFrames.lean\#L221-L288}{The two minimum-cost choices}. \src{New/Derivatives/Divisibility/NullFrames.lean\#L385-L445}{Constructing all matrices from the original directions}.}
\noteline{5.90}{90}{\src{New/Derivatives/Divisibility/Count.lean}{Charged substitution, chain rule, and the factorial estimate}.}
\noteline{5.91}{91}{\src{New/Derivatives/Vacuum/IntegralVacuum.lean\#L201-L276}{The finite average defining the integral multiple}. \src{New/Derivatives/Vacuum/IntegralVacuum.lean\#L376-L425}{Tensor powers and the mean square}.}
\noteline{5.92}{92}{\src{New/UpperBound/Kernels/JetConjugatesAndBounds.lean\#L283-L431}{The trace lattice and its Gaussian theta sum}.}
\noteline{5.93}{93}{\src{New/UpperBound/Kernels/JetConjugatesAndBounds.lean\#L579-L659}{Gaussian evaluation and its integrable majorant}. \src{New/UpperBound/Kernels/JetConjugatesAndBounds.lean\#L783-L894}{The Schur growth exponent and the operator norm}. \src{New/UpperBound/Size/ArchimedeanSize.lean\#L222-L259}{Insertion of the actual scale \(M/\Omega\) and the rank bounds}.}
\noteline{5.94}{94}{\src{New/UpperBound/Size/ShortColumns.lean\#L212-L329}{Finite gauges and the squared Weyl norm}. \src{New/UpperBound/Size/ShortColumns.lean\#L339-L431}{Free-level delta forms and spanning by evaluations}.}
\section*{Chapter 6. The second moment and the orbital sieve}
\noteline{6.1}{1}{\src{New/Spanning/Unfolding/FillingAdmissible.lean\#L123-L128}{Source coordinates after the central translation}. \src{New/Reduction/ShortAndEnough.lean\#L410-L423}{Nontriviality of the full source and application of the spreading argument}.}
\noteline{6.2}{2}{\src{New/Spanning/SecondMoment/FewTranslates.lean}{Dimension gain and greedy iteration}.}
\noteline{6.3}{3}{\src{New/Spanning/Multiplicities/EqualMultiplicities.lean}{Trace comparison on good forms}.}
\noteline{6.4}{4}{\src{New/Spanning/MainTerm/EnergyEstimate.lean}{Open-cell main term}. \src{New/Spanning/RemainingTerms/Sum.lean}{Remaining pairs}. \src{New/Spanning/Unfolding/SquareExpansion.lean}{Expansion into pairs}.}
\noteline{6.5}{5}{\src{New/Spanning/SecondMoment/PositiveType.lean}{Character orthogonality and the determinant test}.}
\noteline{6.6}{6}{\src{New/Spanning/Unfolding/MomentBounds.lean}{Local sizes, terminal estimate and maximization}.}
\noteline{6.7}{7}{\src{New/Spanning/MainTerm/BadTypes.lean}{Green degree bound, weights and product count}. \src{New/UpperBound/Objects/LabelsAndIndexTypes.lean\#L347-L379}{The exponent and the separate condition at each level prime}.}
\noteline{6.8}{8}{\src{New/Spanning/MainTerm/BadTypes.lean}{Finite type count}. \src{New/Spanning/Unfolding/Moments.lean}{Identity block and projection loss}.}
\noteline{6.9}{9}{\src{New/UpperBound/Size/ShortColumns.lean}{Codimension of the ramified source}. \src{New/Reduction/ShortAndEnough.lean}{Passage to the normalized maps}.}
\noteline{6.10}{10}{\src{New/Spanning/MainTerm/SieveSystem.lean}{Positivity and descent}. \src{New/Spanning/MainTerm/GlobalSieve.lean}{One-sided Poisson summation}.}
\noteline{6.11}{11}{\src{New/Spanning/Frames/RealBandLimits.lean}{Complex-place transforms and mass}. \src{New/Spanning/MainTerm/FiniteSieveSystem.lean}{Finite-place transforms}.}
\noteline{6.12}{12}{\src{New/Spanning/MainTerm/Frequencies.lean}{Summable frequency heights}.}
\noteline{6.13}{13}{\src{New/Spanning/MainTerm/GlobalSieve.lean}{Head weights and the assembled bound}.}
\noteline{6.14}{14}{\src{New/Spanning/Frames/FiniteShells.lean}{Finite orbital shell bound}. \src{New/Spanning/Frames/FixedLines.lean}{Stabilizer index and product across places}.}
\noteline{6.15}{15}{\src{New/Spanning/MainTerm/SiegelTopBound.lean}{One-vector class sum and explicit measure factors}. \src{Lib/SiegelWeil/OneVector/SiegelClassSum.lean}{One-vector Siegel formula}.}
\noteline{6.16}{16}{\src{New/Spanning/MainTerm/EnergyEstimate.lean}{Jacobian comparison and normalized energy}.}
\noteline{6.17}{17}{\src{New/Spanning/RemainingTerms/Count.lean}{Count with fixed common-lattice corank}. \src{New/Spanning/RemainingTerms/LatticesOfCorank.lean}{Genus sum of discriminants}. \src{New/Spanning/RemainingTerms/Sum.lean}{Gauge gains and summation over pairs}. \src{New/Spanning/Unfolding/MomentBounds.lean}{Returning to the full kernel}.}
\noteline{6.18}{18}{\src{New/Spanning/Unfolding/Terms.lean\#L220-L290}{Products and the positive middle-coordinate identity}. \src{New/Spanning/Unfolding/Link.lean\#L872-L931}{The two cases of the discriminant criterion}.}
\noteline{6.19}{19}{\src{New/Spanning/Unfolding/Climbing.lean\#L540-L597}{The stabilizer quotient integral and its character}.}
\noteline{6.20}{20}{\src{Lib/Adelic/Witt/WittTheorem.lean\#L736-L751}{The Eichler transformation}. \src{New/Spanning/Unfolding/BesselLevelStructure.lean\#L411-L460}{Its triangular blocks and simple-root entries}. \src{New/Spanning/Unfolding/Climbing.lean\#L696-L710}{Separation by linear forms}. \src{New/Spanning/Unfolding/Climbing.lean\#L731-L802}{Agreement after character normalization}.}
\noteline{6.21}{21}{\src{New/Spanning/Unfolding/Climbing.lean\#L1199-L1233}{Normalization and the least index}. \src{New/Spanning/Unfolding/Link.lean\#L127-L182}{Prescribing the tail in either frame}.}
\noteline{6.22}{22}{\src{New/Spanning/Unfolding/Climbing.lean\#L1230-L1282}{The nonzero leading minor}. \src{New/Spanning/Unfolding/TermStructure.lean\#L673-L730}{The two decompositions and the intrinsic rank}.}
\noteline{6.23}{23}{\src{New/Spanning/RemainingTerms/BorderedDeterminant.lean\#L492-L544}{Inverse terms and the column of the padded inverse}. \src{New/Spanning/RemainingTerms/BorderedDeterminant.lean\#L654-L730}{Replacement of one row and the subframe determinant}.}
\noteline{6.24}{24}{\src{Lib/General/NumberTheory/LatticeCovolumes.lean\#L248-L337}{Jacobi's index identity for a saturated orthogonal complement}. \src{New/Spanning/RemainingTerms/LinkLattice.lean\#L205-L254}{The form on primitive lifts}. \src{New/Spanning/RemainingTerms/LinkLattice.lean\#L255-L315}{The subframe determinant and the two factors of the index}.}
\noteline{6.25}{25}{\src{New/Spanning/Unfolding/Climbing.lean\#L876-L944}{Split-place lifts and integral Eichler transformations}. \src{New/Spanning/Unfolding/Climbing.lean\#L950-L1014}{Comparison of the full level reductions}.}
\noteline{6.26}{26}{\src{New/Spanning/Unfolding/Climbing.lean\#L1016-L1084}{Propagation and the quotient onto residue fields}. \src{New/Spanning/Unfolding/Climbing.lean\#L1092-L1188}{The last-column determinant in lattice coordinates}.}
\noteline{6.27}{27}{\src{New/Spanning/Bessel/Basic.lean\#L368-L393}{Invariance of the successive perpendicular projections}. \src{New/Spanning/Unfolding/ReducedPoints.lean\#L530-L563}{The rational change of triangular coordinates}.}
\noteline{6.28}{28}{\src{Lib/Adelic/Adeles/SemilocalProducts.lean\#L50-L75}{Covering estimate from the different and the Euclidean lattice minimum}. \src{New/Spanning/RemainingTerms/Box.lean\#L145-L181}{The covering radius with the admissible constants}.}
\noteline{6.29}{29}{\src{New/Spanning/Unfolding/ReducedPoints.lean\#L268-L353}{An integral rational dual form and its elementary frame action}. \src{New/Spanning/Unfolding/ReducedPoints.lean\#L599-L616}{The level of the frame entries}.}
\noteline{6.30}{30}{\src{New/Spanning/Unfolding/ReducedPoints.lean\#L465-L525}{Successive reduction of the superdiagonals}. \src{New/Spanning/Unfolding/ReducedPoints.lean\#L617-L637}{The coset of the fixer}.}
\noteline{6.31}{31}{\src{New/Spanning/Bessel/Basic.lean\#L395-L425}{The length bounds for Gram--Schmidt vectors}. \src{New/Spanning/RemainingTerms/Box.lean\#L253-L312}{Basis scales and the frame expansion}. \src{New/Spanning/RemainingTerms/Box.lean\#L314-L331}{The radius after imposing the cut}.}
\noteline{6.32}{32}{\src{New/Spanning/Frames/RealPlaceGauge.lean}{Common positive directions, stable products, and their volume}.}
\noteline{6.33}{33}{\src{New/Spanning/Frames/QuadricVolumes.lean}{Disintegrating in the reserved block and bounding the remaining coordinates}. \src{New/Spanning/Frames/QuadricPlancherel.lean}{Transfer between the floor orbit and quadric measure}.}
\noteline{6.34}{34}{\src{New/Spanning/Frames/Depth.lean}{The difference identity, the centered coordinates, and the two opposite weight integrals}.}
\noteline{6.35}{35}{\src{New/Spanning/Frames/OneFrame.lean}{Disintegration, the one-floor Parseval step and the weighted second moment}.}
\noteline{6.36}{36}{\src{New/Spanning/Frames/Gains.lean}{Cauchy--Schwarz on the head sections and decomposition by deep pivot}.}
\noteline{6.37}{37}{\src{New/Spanning/Frames/Shells.lean}{The shell estimate and its constants}.}
\noteline{6.38}{38}{\src{New/Spanning/Frames/TwoGauges.lean}{Comparison on successive quotients and head duals}. \src{New/Spanning/Frames/Depth.lean}{Normalized stable duals and the size of the relative position}.}
\noteline{6.39}{39}{\src{New/Spanning/Frames/RealShells.lean}{Remote real positions and the dyadic scale}.}
\noteline{6.40}{40}{\src{New/Spanning/RemainingTerms/FrameParts.lean}{Rational lifts, row lattice of the Gram matrix, and the level move}.}
\noteline{6.41}{41}{\src{New/Spanning/RemainingTerms/FrameCount.lean}{The exceptional column and its archimedean length}. \src{New/Spanning/RemainingTerms/FrameParts.lean}{Congruent definite parts and the orthogonal lattice}.}
\noteline{6.42}{42}{\src{New/Spanning/RemainingTerms/FrameCount.lean}{The adjugate, the skew-Hermitian solution and its level}.}
\noteline{6.43}{43}{\src{New/Spanning/RemainingTerms/Box.lean}{The reduced box}. \src{New/Spanning/RemainingTerms/FramesOverLattice.lean}{The successive finite counts and the final exponent}. \src{New/Spanning/RemainingTerms/FrameCount.lean}{Length comparisons at the two points and separation by the label}.}
\noteline{6.44}{44}{\src{New/Spanning/Unfolding/Terms.lean}{Rational and real frame products}.}
\noteline{6.45}{45}{\src{New/Spanning/RemainingTerms/OrthogonalLifts.lean}{Lifts, the two corrected last vectors, and row indices}. \src{New/Spanning/RemainingTerms/BorderedDeterminant.lean}{Complementary minors for a primitive orthogonal basis}. \src{New/Spanning/RemainingTerms/LinkLattice.lean}{The quotient computation and passage to level lifts}.}
\noteline{6.46}{46}{\src{New/Spanning/Unfolding/Climbing.lean}{Propagation of the relation and its last-column determinant}.}
\noteline{6.47}{47}{\src{New/Spanning/RemainingTerms/Cut.lean}{The split filtration and its determinant}. \src{New/Spanning/RemainingTerms/LinkVectors.lean}{The scalar representers and their Gram minors}.}
\noteline{6.48}{48}{\src{New/Spanning/RemainingTerms/Gain.lean}{The mixed exterior-power estimate, crossed torsion indices, and pair discriminant}.}
\noteline{6.49}{49}{\src{New/Spanning/RemainingTerms/Discriminants.lean}{Integral bases, discriminants and the index zeta bound}. \src{Lib/SiegelWeil/Mass.lean}{The ideal zeta estimate}.}
\noteline{6.50}{50}{\src{New/Spanning/RemainingTerms/SiegelLemma.lean}{Saturated orthogonal lattices and the number-field Siegel lemma}.}
\noteline{6.51}{51}{\src{New/Spanning/RemainingTerms/SiegelLemma.lean}{Height zeta bound}.}
\noteline{6.52}{52}{\src{New/Spanning/RemainingTerms/LatticesOfCorank.lean}{Row cosets, averaging and constants}. \src{New/Spanning/Unfolding/Climbing.lean}{Stabilizers at the fine level}.}
\noteline{6.53}{53}{\src{New/Spanning/RemainingTerms/FramesOverLattice.lean}{The three refinements and their exponents}. \src{New/Spanning/RemainingTerms/Gain.lean}{Gain from an agreeing block}.}
\noteline{6.54}{54}{\src{New/Spanning/Unfolding/Moments.lean\#L137-L299}{The unfolded blocks and Bessel equivariance}.}
\noteline{6.55}{55}{\src{New/Spanning/Unfolding/SquareExpansion.lean\#L572-L692}{The restricted square expansion and the meeting norm}.}
\noteline{6.56}{56}{\src{New/Spanning/Unfolding/Moments.lean\#L600-L675}{Restriction of the convolution projection}.}
\noteline{6.57}{57}{\src{New/Spanning/Unfolding/Moments.lean\#L696-L779}{The trace at a single level factor}. \src{New/Spanning/Unfolding/Moments.lean\#L900-L1024}{Reduction of the local integral to the finite conjugacy average}.}
\noteline{6.58}{58}{\src{New/Spanning/Unfolding/Moments.lean\#L1026-L1210}{Trace of the discarded types and the two-factor loss}.}
\noteline{6.59}{59}{\src{New/Spanning/MainTerm/FillingLevel.lean\#L532-L558}{The finite lattice datum and its rank constraints}.}
\noteline{6.60}{60}{\comsrc{New/Spanning/Frames/StandardLevel.lean}{Common integral vectors and their content moments}.}
\noteline{6.61}{61}{\comsrc{Lib/SiegelWeil/OneVector/StandardQuadrics.lean\#L347-L377}{The integral hyperbolic basis}. \comsrc{Lib/SiegelWeil/OneVector/StandardQuadrics.lean\#L236-L290}{Translation by the content of a hyperbolic block}. \comsrc{Lib/SiegelWeil/OneVector/StandardQuadrics.lean\#L428-L456}{The standard quadric in rank at least \(2m+4\)}. \comsrc{New/Spanning/Frames/StandardLevel.lean\#L553-L581}{The bound used at the pivot}.}
\noteline{6.62}{62}{\src{New/Spanning/Frames/FiniteDepth.lean\#L80-L162}{Stable-pair lattices and the coordinate ideal}. \src{New/Spanning/Frames/FiniteDepth.lean\#L445-L543}{A common multiple and propagation through the frame}.}
\noteline{6.63}{63}{\src{New/Spanning/Frames/Shells.lean\#L88-L247}{Per-box losses and the choice of gain on each shell}. \src{New/Spanning/Frames/FiniteShells.lean\#L198-L318}{The ideal-shell cover and its sum}. \src{New/Spanning/Frames/FiniteShells.lean\#L591-L724}{Both pivot gains off the congruence group and the finite mass}.}
\noteline{6.64}{64}{\src{New/Spanning/Frames/FiniteDepth.lean\#L262-L405}{Head products, weights and scale}. \src{New/Spanning/Frames/FiniteShells.lean\#L345-L552}{Replacing a frequency and retaining the denominator saving}.}
\section*{Chapter 7. Archimedean Bessel periods}
\noteline{7.1}{1}{\src{New/Spanning/Bessel/Orthonormality.lean}{Compact-type normalization}. \src{New/Spanning/Bessel/SchurBound.lean}{Discrete-series convolution and coefficient}.}
\noteline{7.2}{2}{\src{New/Spanning/Bessel/CutTailBound.lean}{Decay majorant and cutoff}. \src{New/Spanning/Bessel/FillingData.lean\#L313-L340}{Convolution convention and the decay parameter at a real place}.}
\noteline{7.3}{3}{\src{Lib/DiscreteSeries/Haar/CayleyChart.lean\#L359-L376}{The compact-dual normalizing factor in the Haar density}. \src{New/Spanning/Frames/RealPlaceGauge.lean}{Comparison of group and quotient Haar measures}.}
\noteline{7.4}{4}{\src{New/Spanning/Frames/RealBandLimits.lean}{Near and remote contributions to the mass}.}
\noteline{7.5}{5}{\src{Lib/DiscreteSeries/Bruhat/FourierKernels.lean}{Fourier bound between solved blocks}. \src{New/Spanning/Frames/Gains.lean}{The two frame phases}.}
\noteline{7.6}{6}{\src{New/Spanning/Frames/Shells.lean}{Shell volume and depth estimates}. \src{New/Spanning/Frames/RealShells.lean}{Complex-place specialization}.}
\noteline{7.7}{7}{\src{New/Spanning/Bessel/Basic.lean}{Top orbital normalization}. \src{New/Spanning/Frames/RealBandLimits.lean}{Local descent and top Gram kernel}.}
\noteline{7.8}{8}{\src{New/Spanning/Bessel/FillingData.lean}{Specialization to the parameters of the proof}.}
\noteline{7.9}{9}{\src{New/Appendix/DiscreteSeries/StationaryPosition.lean}{Stationary flag, residue pairings and its size}.}
\noteline{7.10}{10}{\src{New/Appendix/DiscreteSeries/PositionGroup.lean}{Trace calculation for the position group}.}
\noteline{7.11}{11}{\src{New/Appendix/DiscreteSeries/Localization.lean}{The logarithmic phase identity}.}
\noteline{7.12}{12}{\src{New/Spanning/Bessel/FillingData.lean}{Specialization to the two parameters}.}
\noteline{7.13}{13}{\comsrc{New/Appendix/DiscreteSeries/GaussianShape.lean\#L182-L199}{Product of the minors}. \comsrc{New/Appendix/DiscreteSeries/GaussianShape.lean\#L296-L358}{Gaussian estimate for the coefficient}.}
\noteline{7.14}{14}{\src{New/Appendix/DiscreteSeries/Localization.lean}{Gap estimate and translation to the identity}.}
\noteline{7.15}{15}{\src{New/Appendix/DiscreteSeries/Periods.lean}{Period inequality with the small-ball scale}. \src{Lib/DiscreteSeries/Haar/TotalMass.lean}{Haar covering estimates}.}
\noteline{7.16}{16}{\src{New/Spanning/Bessel/LargeFunction.lean}{Controlled position, unfolding and norm of the average}.}
\noteline{7.17}{17}{\src{New/Spanning/Bessel/CutTailBound.lean}{The group normalization and transfer of decay to the average}.}
\noteline{7.18}{18}{\comsrc{New/Appendix/DiscreteSeries/MinorExpansion.lean\#L144-L236}{Logarithm of a compressed matrix and the uniform trace estimate}.}
\noteline{7.19}{19}{\comsrc{New/Appendix/DiscreteSeries/MinorExpansion.lean\#L250-L346}{Triangular factors in the Flensted--Jensen integral}.}
\noteline{7.20}{20}{\comsrc{New/Appendix/DiscreteSeries/MinorExpansion.lean\#L347-L462}{Complementary minors and the estimate for one minor}. \comsrc{New/Appendix/DiscreteSeries/GaussianShape.lean\#L188-L254}{Threshold sum and comparison of weights}.}
\noteline{7.21}{21}{\comsrc{New/Appendix/DiscreteSeries/GaussianShape.lean\#L51-L88}{The compact squared norm and its lower bound}.}
\noteline{7.22}{22}{\src{Lib/DiscreteSeries/FlenstedJensen/SquareIntegrability.lean\#L849-L874}{The derivative bound at the required small-ball scale}.}
\noteline{7.23}{23}{\src{Lib/DiscreteSeries/FlenstedJensen/SquareIntegrability.lean}{Gap, complementary positive minors, and the polar coefficient integral}.}
\noteline{7.24}{24}{\src{New/Appendix/DiscreteSeries/Localization.lean}{Shared compact translation and the two distance bounds}.}
\noteline{7.25}{25}{\src{New/Appendix/DiscreteSeries/PositionGroup.lean\#L124-L238}{The closed group and its unitary quotient action}.}
\noteline{7.26}{26}{\src{New/Appendix/DiscreteSeries/Positions.lean\#L245-L341}{The middle, flag and dual-flag terms}.}
\noteline{7.27}{27}{\src{New/Appendix/DiscreteSeries/Positions.lean\#L422-L482}{The two traces at a stationary position}.}
\noteline{7.28}{28}{\src{New/Appendix/DiscreteSeries/Positions.lean\#L388-L422}{Pairings of higher powers of the weight}. \src{New/Spanning/Bessel/LargeFunction.lean\#L218-L285}{Witt extension and entry bounds}.}
\noteline{7.29}{29}{\src{Lib/DiscreteSeries/Bessel/Basic.lean\#L505-L568}{The frame action and its free dual columns}. \src{New/Spanning/Bessel/Basic.lean\#L338-L361}{Realizing Euclidean orthogonalization inside the Bessel group}.}
\noteline{7.30}{30}{\src{New/Spanning/Bessel/Basic.lean\#L202-L245}{Operator norm from the negative directions}. \src{New/Spanning/Bessel/Basic.lean\#L479-L577}{The dual-frame length estimate}.}
\noteline{7.31}{31}{\src{New/Spanning/Bessel/LargeFunction.lean\#L125-L188}{Flag action and properness of the pair map}. \src{New/Appendix/DiscreteSeries/PositionGroup.lean\#L426-L475}{Transport of the right invariant measure}.}
\noteline{7.32}{32}{\src{Lib/DiscreteSeries/Bessel/UnipotentHaarMeasure.lean\#L572-L650}{Lebesgue measure on the free entries and the Eichler multiplication law}. \src{New/Spanning/Bessel/Average.lean\#L63-L126}{The graph measure and its invariance}.}
\noteline{7.33}{33}{\src{New/Spanning/Bessel/LargeFunction.lean\#L103-L123}{Unfolding the graph}. \src{New/Spanning/Bessel/LargeFunction.lean\#L681-L726}{The norm identity with this Haar normalization}.}
\noteline{7.34}{34}{\src{Lib/DiscreteSeries/FlenstedJensen/SquareIntegrability.lean\#L874-L937}{Layer integration with the actual formal degree}. \src{Lib/DiscreteSeries/Character/FormalDegree.lean\#L473-L479}{The stronger hypothesis for the Weyl polynomial}.}
\noteline{7.35}{35}{\src{Lib/DiscreteSeries/FlenstedJensen/SquareIntegrability.lean\#L285-L309}{The polar integral with the project's normalization}. \src{Lib/DiscreteSeries/FlenstedJensen/SquareIntegrability.lean\#L700-L747}{The unit generator and invariant norm}.}
\noteline{7.36}{36}{\src{Lib/DiscreteSeries/FlenstedJensen/SquareIntegrability.lean\#L132-L174}{The polar determinant bound}.}
\noteline{7.37}{37}{\src{Lib/DiscreteSeries/FlenstedJensen/SquareIntegrability.lean\#L194-L278}{The one-threshold inequality and its product}.}
\noteline{7.38}{38}{\src{Lib/DiscreteSeries/FlenstedJensen/SquareIntegrability.lean\#L279-L308}{Integration of the polar inequality}. \src{Lib/DiscreteSeries/FlenstedJensen/PsiKTypes.lean\#L78-L185}{The value at the identity, inversion, and restriction to the compact group}. \src{Lib/DiscreteSeries/FlenstedJensen/PsiKTypes.lean\#L196-L467}{Circle averaging and the compact-type projection}. \src{Lib/DiscreteSeries/FlenstedJensen/SquareIntegrability.lean\#L533-L577}{Realization by right translations}.}
\noteline{7.39}{39}{\src{Lib/DiscreteSeries/FlenstedJensen/CasimirEigenfunction.lean\#L795-L846}{Casimir eigenvalue of the explicit integral}.}
\noteline{7.40}{40}{\src{Lib/DiscreteSeries/Rigidity/HighestVectorsAndFiltration.lean\#L396-L469}{The Casimir calculation, also valid on higher filtration steps with the corresponding weight bound}. \src{Lib/DiscreteSeries/FlenstedJensen/SquareIntegrability.lean\#L765-L780}{Specialization at the generator}.}
\noteline{7.41}{41}{\src{Lib/DiscreteSeries/FlenstedJensen/SquareIntegrability.lean\#L802-L897}{The path, identity ball, and covering estimate}. \src{Lib/DiscreteSeries/FlenstedJensen/SquareIntegrability.lean\#L909-L952}{The layer integral}.}
\noteline{7.42}{42}{\src{Lib/DiscreteSeries/FlenstedJensen/SquareIntegrability.lean\#L533-L577}{Square integrability and the translated coefficient identity}.}
\noteline{7.43}{43}{\src{Lib/DiscreteSeries/FlenstedJensen/PsiKTypes.lean\#L476-L491}{The two compact averaging identities}.}
\noteline{7.44}{44}{\src{Lib/DiscreteSeries/Character/MultiplicityOrbitalIntegral.lean\#L249-L301}{Highest-vector multiplicity and support for the coefficient}.}
\noteline{7.45}{45}{\src{Lib/Adelic/Reduction/FirstMinimumShortBases.lean\#L652-L714}{Trace integrality and the finite lattice}.}
\noteline{7.46}{46}{\src{Lib/Adelic/Reduction/FirstMinimumShortBases.lean\#L716-L788}{Duality in the real coordinates}.}
\noteline{7.47}{47}{\src{Lib/Adelic/Reduction/FirstMinimumShortBases.lean\#L790-L887}{The covering radius and selection of the columns}.}
\noteline{7.48}{48}{\src{Lib/Adelic/Reduction/DefiniteClassSums.lean\#L64-L115}{Compact representatives from bounded columns}.}
\noteline{7.49}{49}{\src{Lib/Weil/DualPair/ThetaLiftConvergence.lean\#L71-L179}{The short-space dilation and iteration}. \src{Lib/Weil/DualPair/ThetaLiftConvergence.lean\#L181-L215}{The resulting operator bound}.}
\noteline{7.50}{50}{\src{Lib/Adelic/Reduction/FiniteVolume.lean}{Reduction by a rational basis and integration of the matrix count}.}
\noteline{7.51}{51}{\src{Lib/Weil/Periods/ConstantTerm.lean\#L303-L365}{The graded null frame and its opposite projector}. \src{Lib/Weil/Periods/ConstantTerm.lean\#L417-L453}{Passage to rational classes}. \src{Lib/Weil/Periods/ConstantTerm.lean\#L483-L561}{The explicit bijection and its norm bound}.}
\noteline{7.52}{52}{\src{Lib/Adelic/Reduction/DefiniteClassSums.lean\#L64-L115}{Short frames and finitely many Gram matrices}. \src{Lib/Adelic/Reduction/DefiniteClassSums.lean\#L117-L157}{Definiteness and the compact cover}.}
\noteline{7.53}{53}{\src{Lib/Weil/Periods/BesselIntegral.lean\#L65-L129}{The downward recursion on entries}. \src{Lib/Weil/Periods/BesselIntegral.lean\#L131-L193}{The rational isometry in the Bessel group}. \src{Lib/Weil/Periods/BesselIntegral.lean\#L253-L265}{The quotient measure}.}
\noteline{7.54}{54}{\src{Lib/Weil/Periods/Unfolding.lean\#L148-L210}{Equivariance and the first uniform bound}. \src{Lib/Weil/DualPair/ThetaLiftConvergence.lean\#L180-L215}{The height bound for a product of operators}. \src{Lib/Weil/Periods/Unfolding.lean\#L457-L511}{The target test and the sum over translates}. \src{Lib/Weil/Periods/Unfolding.lean\#L513-L578}{Absolute integrability of the final unfolding}.}
\section*{Chapter 8. Signature transfer and the target dimension}
\noteline{8.1}{1}{\src{Lib/Adelic/TraceFormula/CuspidalType.lean}{Cuspidality after inflation}.}
\noteline{8.2}{2}{\src{Lib/Adelic/TraceFormula/TraceComparison.lean}{Elliptic class matching and compact character}.}
\noteline{8.3}{3}{\src{Lib/Adelic/TraceFormula/CuspidalComparison.lean}{Auxiliary split place and cut level}. \src{Lib/Adelic/TraceFormula/LineDeterminant.lean}{Recovering the full-level trace}. \src{New/Spanning/Multiplicities/EqualMultiplicities.lean}{Good-space dimension equality}.}
\noteline{8.4}{4}{\src{Lib/Adelic/TraceFormula/LineStabiliser.lean}{The cosets, adjoints and support of the raising operator}.}
\noteline{8.5}{5}{\src{Lib/Adelic/TraceFormula/Convergence.lean}{Cusp frames and bounds uniform over partial polynomial sums}. \src{Lib/Adelic/TraceFormula/OrbitalUnfolding.lean\#L278-L353}{Reproducing kernel and diagonal trace}.}
\noteline{8.6}{6}{\src{Lib/Adelic/TraceFormula/CoprimeFactors.lean}{Descent to the projector slice and unipotent averaging}. \src{Lib/Adelic/TraceFormula/ReduciblePolynomials.lean}{Elimination of all reducible polynomials}.}
\noteline{8.7}{7}{\src{Lib/Adelic/TraceFormula/OrbitalUnfolding.lean}{Triple rotation and absolute moment estimates}. \src{Lib/Adelic/TraceFormula/TraceComparison.lean\#L102-L215}{Height increments and nonelliptic vanishing}.}
\noteline{8.8}{8}{\src{Lib/Adelic/TraceFormula/ChangingSigns.lean\#L249-L367}{CM norm and change of finite frame}. \src{Lib/ClassField/Global/CMSplittingCharacters.lean\#L284-L303}{The signs in quadratic idelic reciprocity}.}
\noteline{8.9}{9}{\src{Lib/Adelic/TraceFormula/TraceComparison.lean\#L568-L651}{Permutation multiplicities and the signature-independent trace}. \src{Lib/DiscreteSeries/Character/FormalDegree.lean}{The real orbital character identity}.}
\noteline{8.10}{10}{\src{Lib/LocalGL/Green/TauFunctions.lean\#L80-L132}{The exponential label and the sign convention}.}
\noteline{8.11}{11}{\src{Lib/LocalGL/Green/CuspidalFunctions.lean\#L62-L105}{Möbius vanishing on the submodule lattice}. \src{Lib/LocalGL/Green/LabelCharacters.lean\#L229-L306}{Orthogonality, virtuality and the positive character in the chosen convention}.}
\noteline{8.12}{12}{\src{Lib/ClassField/Global/CMSplittingCharacters.lean\#L285-L313}{The sign condition in the finite-idele convention}.}
\noteline{8.13}{13}{\src{Lib/Adelic/Landherr.lean\#L433-L475}{The isometry and its action on the pointed spaces}.}
\noteline{8.14}{14}{\src{Lib/Adelic/Landherr.lean\#L529-L595}{Descent of the finite frame, Gram matrix and rational coboundary}.}
\noteline{8.15}{15}{\src{Lib/LocalGL/Green/CuspidalFunctions.lean\#L208-L254}{The Hall block mean}. \src{Lib/LocalGL/Green/LabelCharacters.lean\#L274-L313}{Positive irreducible characters in the label convention}. \src{Lib/Adelic/TraceFormula/CuspidalType.lean\#L92-L155}{The radical sum in the group algebra}.}
\noteline{8.16}{16}{\src{Lib/Adelic/TraceFormula/CuspidalType.lean\#L167-L207}{Inflation of the coefficient to the compact group}. \src{Lib/Adelic/TraceFormula/CuspidalType.lean\#L211-L270}{The marked factor after lifting and convolution}.}
\noteline{8.17}{17}{\src{Lib/DiscreteSeries/Character/MultiplicityOrbitalIntegral.lean\#L797-L901}{Compact multiplicities and the torus polynomial}. \src{Lib/DiscreteSeries/Character/OpenOrbitIntegral.lean\#L351-L426}{Subtraction of the complement of the flag orbit}. \src{Lib/DiscreteSeries/Character/FormalDegree.lean\#L397-L479}{Identification of the numerator, phase and formal degree}.}
\noteline{8.18}{18}{\src{Lib/DiscreteSeries/Character/FormalDegree.lean\#L481-L568}{Absolute convergence and the sum for the conjugated kernel}. \src{Lib/Adelic/TraceFormula/TraceComparison.lean\#L568-L651}{Its application to the two signatures}.}
\noteline{8.19}{19}{\src{Lib/DiscreteSeries/Haar/CayleyChart.lean\#L359-L384}{Haar density in the signed Cayley chart}. \src{Lib/DiscreteSeries/Haar/TotalMass.lean\#L289-L304}{Mass one in compact signature}.}
\noteline{8.20}{20}{\src{Lib/DiscreteSeries/FlenstedJensen/FlagDensities.lean\#L448-L537}{Finite integrals and the formal open-orbit change of variables}.}
\noteline{8.21}{21}{\src{Lib/Adelic/TraceFormula/Convergence.lean\#L60-L90}{The polynomial and partial kernels}. \src{Lib/Adelic/TraceFormula/Convergence.lean\#L431-L479}{The uniform bound for every set of polynomials}.}
\noteline{8.22}{22}{\src{Lib/Adelic/TraceFormula/KernelBound.lean\#L128-L159}{Lower bound for the volume of the level box}. \src{Lib/Adelic/TraceFormula/KernelBound.lean\#L161-L244}{The short basis and disjointness}. \src{Lib/Adelic/TraceFormula/KernelBound.lean\#L247-L299}{Variation and integrability of the majorant}. \src{Lib/Adelic/TraceFormula/KernelBound.lean\#L301-L346}{Summing by disjoint translates}.}
\noteline{8.23}{23}{\src{Lib/Adelic/TraceFormula/Convergence.lean\#L183-L253}{Saving powers of the norm in the absolute majorant}.}
\noteline{8.24}{24}{\src{Lib/Adelic/AutomorphicForms/Cuspidality.lean\#L228-L266}{Product bound for the basis and the nonzero pairing}. \src{Lib/Adelic/AutomorphicForms/Cuspidality.lean\#L267-L338}{Constructing the two spans from the paired lengths}.}
\noteline{8.25}{25}{\src{Lib/Adelic/AutomorphicForms/Cuspidality.lean\#L340-L384}{The quadratic group formula and the rounding bound}. \src{Lib/Adelic/AutomorphicForms/Cuspidality.lean\#L413-L439}{Rational translation of the radical parameters}. \src{Lib/Adelic/AutomorphicForms/Cuspidality.lean\#L491-L557}{Quantitative covering in the prescribed level}. \src{Lib/Adelic/AutomorphicForms/Cuspidality.lean\#L796-L817}{The relation between the radius and the height}.}
\noteline{8.26}{26}{\src{Lib/Adelic/TraceFormula/Stretch.lean\#L60-L107}{Clearing denominators and multiplying the archimedean inequalities}. \src{Lib/Adelic/TraceFormula/Convergence.lean\#L294-L341}{The exponent \(4n^3\) for the moving elements}. \src{Lib/Adelic/TraceFormula/Convergence.lean\#L343-L372}{The resulting uniform absolute bound in the irreducible case}.}
\noteline{8.27}{27}{\src{Lib/Adelic/TraceFormula/IsotropicStabiliser.lean\#L192-L237}{Invariance of the characteristic polynomial along the rational radical}. \src{Lib/Adelic/TraceFormula/Convergence.lean\#L99-L166}{Level, left invariance and reproduction of the partial kernel}.}
\noteline{8.28}{28}{\src{Lib/DiscreteSeries/CuspKernel.lean\#L139-L195}{Translation and the slowly varying phase}. \src{Lib/DiscreteSeries/CuspKernel.lean\#L197-L307}{Anticommutation and injectivity}. \src{Lib/DiscreteSeries/CuspKernel.lean\#L328-L396}{Small conjugate dilations via orthogonal projections}. \src{Lib/Adelic/AutomorphicForms/UnipotentRadical.lean\#L166-L210}{The logarithmic height and its finite moment}. \src{Lib/Adelic/AutomorphicForms/UnipotentRadical.lean\#L216-L263}{Reproduction and invariance of the radical mean}.}
\noteline{8.29}{29}{\src{Lib/Adelic/AutomorphicForms/Cuspidality.lean\#L559-L585}{Bounding a value by oscillation on representatives}. \src{Lib/Adelic/AutomorphicForms/Cuspidality.lean\#L706-L742}{Using the compact reproducing function}. \src{Lib/Adelic/AutomorphicForms/Cuspidality.lean\#L744-L794}{The weighted oscillation estimate}. \src{Lib/Adelic/TraceFormula/Convergence.lean\#L385-L429}{Application at exponent \(nd\) to the polynomial-restricted sum}.}
\noteline{8.30}{30}{\src{Lib/Adelic/TraceFormula/ReduciblePolynomials.lean\#L138-L184}{The two passages to a limit}. \src{Lib/Adelic/TraceFormula/OrbitalUnfolding.lean\#L360-L429}{Absolute integrability for irreducibles and unfolding}.}
\noteline{8.31}{31}{\src{Lib/Adelic/TraceFormula/OrbitalUnfolding.lean\#L475-L516}{The triples and cyclic rotation}. \src{Lib/Adelic/TraceFormula/OrbitalUnfolding.lean\#L518-L550}{Change of variables on the centralizer quotient}. \src{Lib/Adelic/TraceFormula/OrbitalUnfolding.lean\#L552-L605}{The product of the three integrable factors}.}
\noteline{8.32}{32}{\src{Lib/LocalGL/Green/TauFunctions.lean\#L345-L374}{The rank restriction and the cuspidality of the irreducible label}. \src{Lib/LocalGL/Green/LabelCharacters.lean\#L274-L305}{The character theorem for the representation at the marked place}.}
\noteline{8.33}{33}{\src{Lib/Adelic/TraceFormula/CuspidalType.lean\#L92-L154}{From the finite mean to all translated coefficients}. \src{Lib/Adelic/TraceFormula/CuspidalType.lean\#L211-L269}{Lifting the cancellation to the finite adelic test}. \src{Lib/Adelic/TraceFormula/ReduciblePolynomials.lean\#L221-L249}{The use of both local tests}.}
\noteline{8.34}{34}{\src{Lib/Adelic/TraceFormula/FibreIntegral.lean\#L186-L205}{The support and invariance required of the functional}. \src{Lib/Adelic/TraceFormula/CoprimeFactors.lean\#L755-L775}{The local coprime-factor theorem}.}
\noteline{8.35}{35}{\src{Lib/Adelic/TraceFormula/CoprimeFactors.lean\#L307-L367}{Idempotence and rank, including repeated roots}. \src{Lib/Adelic/TraceFormula/CoprimeFactors.lean\#L57-L85}{The block groups}.}
\noteline{8.36}{36}{\src{Lib/General/LinearAlgebra/IntegralLattices.lean\#L254-L277}{Valuation-ring diagonalization}. \src{Lib/Adelic/TraceFormula/CoprimeFactors.lean\#L241-L305}{The two factorizations}.}
\noteline{8.37}{37}{\src{Lib/Adelic/TraceFormula/CoprimeFactors.lean\#L377-L419}{Openness around the identity and the intertwiner}. \src{Lib/Adelic/TraceFormula/CoprimeFactors.lean\#L474-L524}{The projector slice and its fibres}.}
\noteline{8.38}{38}{\src{Lib/Adelic/TraceFormula/FibreIntegral.lean\#L250-L322}{Haar uniqueness with a parameter}. \src{Lib/Adelic/TraceFormula/FibreIntegral.lean\#L432-L582}{Fibre averaging and the descended functional}.}
\noteline{8.39}{39}{\src{Lib/Adelic/TraceFormula/CoprimeFactors.lean\#L146-L168}{Recovering the radical parameter from the projector}. \src{Lib/Adelic/TraceFormula/CoprimeFactors.lean\#L528-L590}{The compact cutoff and shear of the support}.}
\noteline{8.40}{40}{\src{Lib/Adelic/TraceFormula/CoprimeFactors.lean\#L87-L145}{The two integral diagonal blocks}. \src{Lib/Adelic/TraceFormula/CoprimeFactors.lean\#L593-L647}{The coset, its measure and the image test}.}
\noteline{8.41}{41}{\src{Lib/Adelic/TraceFormula/TestFunction.lean\#L129-L189}{Integral lifts and the cosets used in cuspidality}. \src{Lib/Adelic/TraceFormula/FibreIntegral.lean\#L207-L222}{Passing the identity through compact conjugation averaging}.}
\noteline{8.42}{42}{\src{Lib/General/LinearAlgebra/IdempotentRadical.lean\#L130-L147}{The polynomial intertwiner}. \src{Lib/Adelic/TraceFormula/CoprimeFactors.lean\#L438-L453}{Passing from the split factor to the adelic group}. \src{Lib/Adelic/TraceFormula/CoprimeFactors.lean\#L664-L675}{The conjugating radical element}.}
\noteline{8.43}{43}{\src{Lib/Adelic/TraceFormula/CoprimeFactors.lean\#L650-L662}{Right invariance on the radical, obtained directly in the formal proof}. \src{Lib/Adelic/TraceFormula/CoprimeFactors.lean\#L677-L751}{The pointwise sum and its integral}.}
\noteline{8.44}{44}{\src{Lib/Adelic/TraceFormula/ReduciblePolynomials.lean\#L36-L53}{Splitting a maximal primary factor}. \src{Lib/Adelic/TraceFormula/ReduciblePolynomials.lean\#L186-L219}{Support and conjugation invariance of \(I_R\)}. \src{Lib/Adelic/TraceFormula/ReduciblePolynomials.lean\#L221-L254}{Combining the two finite places}.}
\section*{Chapter 9. Volumes of single lifts and of the full chain}
\noteline{9.1}{1}{\src{New/Determinant/Twist/LongForms.lean}{Theorem C and its normalization}.}
\noteline{9.2}{2}{\src{New/Reduction/Lattices/Heights.lean}{Family heights, rational powers, and the norm proof}.}
\noteline{9.3}{3}{\src{New/Determinant/Twist/FitTwists.lean\#L150-L203}{The two projection bounds in the definition of regularity}. \src{New/Congruences/FitTwistsExist.lean}{Local denominators and existence of regular twists}. \src{New/Derivatives/Field/ChosenPlaces.lean\#L268-L314}{Allowed weights and finiteness of critical pairs}.}
\noteline{9.4}{4}{\src{New/Congruences/Selmer/MainSummands.lean}{The main summands in the twisting argument}. \src{New/Congruences/GaloisRepresentations/SelmerHypotheses.lean}{Local hypotheses for the attached representations}.}
\noteline{9.5}{5}{\src{New/Congruences/Selmer/LValueSystems.lean}{Systems for the Hecke values}. \src{New/Congruences/Selmer/GeneralTheorem.lean}{Simultaneous conclusions for systems, local generators, and twists}.}
\noteline{9.6}{6}{\src{New/Congruences/Selmer/Descent.lean}{The slopes, the decreasing length, and the terminal duality estimate}.}
\noteline{9.7}{7}{\src{New/Congruences/Selmer/RootsOfUnity.lean}{Simultaneous polynomial specialization}. \src{New/Congruences/Selmer/TwistingTheorem.lean}{Application to main summands and critical values}.}
\noteline{9.8}{8}{\src{New/Congruences/FirstOrder/RibetElement.lean}{The polynomial separating element and coefficient bound}. \src{New/Congruences/FirstOrder/TrianglePrimes.lean}{Triangular testing of even classes}. \src{Lib/Cyclotomic/OddSha.lean}{Uniform odd cocycle bound}.}
\noteline{9.9}{9}{\src{New/Congruences/FirstOrder/DenominatorBound.lean}{From eigenvalues to a projection on the integral lattice}. \src{New/Congruences/FirstOrder/TraceRelations.lean}{The terms in the clearing bound}.}
\noteline{9.10}{10}{\src{New/Congruences/LocalLoss/CrystallineDefect.lean}{The crystalline local loss}. \src{New/Congruences/FitTwistsExist.lean}{Summation over primes and assembly of the regular family}.}
\noteline{9.11}{11}{\src{New/OneLift/ComparisonOperator/ThreeExchanges.lean}{Gluing and the comparison of the three exchanges}. \src{New/OneLift/LiftVolumeHeight.lean}{The adjoint identity, tight frame, and bounds on the comparison operator}.}
\noteline{9.12}{12}{\src{New/OneLift/LiftVolumeHeight.lean}{The compression determinant and both absolute-value estimates}.}
\noteline{9.13}{13}{\src{New/OneLift/HeckeValues/AlgebraicPartSize.lean}{The common period, empty row, and Archimedean multiplier}.}
\noteline{9.14}{14}{\src{New/Derivatives/Sharpness/BoundAttained.lean}{Attainment of the local derivative bound}. \src{New/OneLift/HeckeValues/TwistedValueHeights.lean}{Construction of the test forms and their two pairings}.}
\noteline{9.15}{15}{\src{New/OneLift/HeckeValues/TwistedValueHeights.lean}{The three strings and the \(k/2\) main positions}. \src{New/OneLift/LiftVolumeHeight.lean}{The final height of one lift}.}
\noteline{9.16}{16}{\src{New/Spanning/Multiplicities/EqualMultiplicities.lean}{The Hecke eigenspace multiplicity transfer}. \src{New/Determinant/Chain/NormOfChain.lean}{Single-step and full-chain Rallis determinants}.}
\noteline{9.17}{17}{\src{New/Determinant/Twist/LongForms.lean}{The chain height and ramified corrections}.}
\noteline{9.18}{18}{\src{New/Determinant/AuxiliaryForms/SpanningForms.lean}{Denominators on the last path and recovery by delta forms}. \src{New/Determinant/Twist/LongForms.lean}{Projection of the spanning family to the target}.}
\noteline{9.19}{19}{\src{New/Reduction/Lattices/TwistedLattice.lean}{Conjugate coordinates, mixed pairings, and the two index matrices}.}
\noteline{9.20}{20}{\src{New/Determinant/Twist/LongForms.lean}{Completion cost and the final lower bound}.}
\section*{Chapter 10. Theta exchange, local factors and the norm formula}
\noteline{10.1}{1}{\src{Lib/SiegelWeil/Tower/FiniteTateIntegral.lean}{Definition of the scalar product of values}.}
\noteline{10.2}{2}{\src{New/Determinant/Exchange/ExplicitMatching.lean}{The rank-zero partner and determinant data}.}
\noteline{10.3}{3}{\src{New/OneLift/ComparisonOperator/ThreeExchanges.lean}{The two carrier identities and level of the carrier}.}
\noteline{10.4}{4}{\src{New/OneLift/ComparisonOperator/ThreeExchanges.lean}{The invariant trace character and third exchange}.}
\noteline{10.5}{5}{\src{New/OneLift/ComparisonOperator/ThreeExchanges.lean}{The three critical pairs and their product}.}
\noteline{10.6}{6}{\src{New/Determinant/Exchange/MatchingPartners.lean}{Complementary minors and the induced sections of the local group}. \src{New/OneLift/ComparisonOperator/OneExchange.lean}{Product kernels, convergence and class sums at the row stabilizer}.}
\noteline{10.7}{7}{\src{New/Determinant/Exchange/LargeTriple.lean}{The Fourier terms, their support in Gram matrices and the period identity}. \src{New/Determinant/Exchange/RemovingPlanes.lean}{Descent on either side by the constant term}. \src{Lib/SiegelWeil/Tower/LargeSideConstantTerm.lean}{Unfolding the isotropic line and the ramified normalization}. \src{Lib/SiegelWeil/Tower/FiniteTateIntegral.lean}{The finite integral, its domain of convergence and the Euler factors}.}
\noteline{10.8}{8}{\comsrc{New/Determinant/Exchange/VanishingDescends.lean\#L45-L52}{The finite determinant function}. \comsrc{New/Determinant/Exchange/TheoremForPartners.lean\#L391-L423}{The two exchanges and nonvanishing}. \comsrc{New/Determinant/Exchange/TheoremForPartners.lean\#L722-L823}{The constant square and passage to the required signature}. \comsrc{New/Determinant/Exchange/TheoremForPartners.lean}{Extension to matched functions}.}
\noteline{10.9}{9}{\src{New/Determinant/Exchange/ExplicitMatching.lean}{Circle projection, unitary polarization, definite class sum, and normalization}. \src{New/OneLift/ComparisonOperator/ThreeExchanges.lean}{The rank-zero integral and nonzero carrier}.}
\noteline{10.10}{10}{\src{Lib/SiegelWeil/RallisInner/Doubling.lean}{The doubled integral and Weil matrix coefficients}. \comsrc{Lib/SiegelWeil/RallisInner/SharpData.lean\#L436-L496}{Injectivity for the sharp datum, the archimedean pairing and the unramified coset sum}.}
\noteline{10.11}{11}{\src{New/OneLift/HeckeValues/CharacterLiftNorm.lean}{Block means, level indices, and the norm on the chosen rank-one character lift}.}
\noteline{10.12}{12}{\src{Lib/SiegelWeil/Formula/WeilRange.lean\#L216-L247}{Summing the dyadic fibres}. \src{Lib/SiegelWeil/Formula/WeilRange.lean\#L286-L390}{Frame reduction and the hermitian parameter count}. \src{Lib/SiegelWeil/Formula/LagrangianZeroTerms.lean\#L343-L399}{The normalized action and the integrable theta majorant}.}
\noteline{10.13}{13}{\src{Lib/SiegelWeil/Formula/WeilRange.lean\#L435-L461}{The two theta integrals at the normalized translate}. \src{Lib/SiegelWeil/Formula/WeilRange.lean\#L467-L525}{The translation estimate on the finite and infinite components}.}
\noteline{10.14}{14}{\src{Lib/SiegelWeil/Formula/WeilRange.lean\#L538-L585}{Vanishing away from the chosen line}. \src{Lib/SiegelWeil/Formula/WeilRange.lean\#L660-L733}{Induction at a line}.}
\noteline{10.15}{15}{\src{Lib/SiegelWeil/Formula/WeilRange.lean\#L595-L660}{The compact normalization and the finite exceptional box}.}
\noteline{10.16}{16}{\src{Lib/SiegelWeil/Formula/LagrangianZeroTerms.lean\#L661-L710}{Hecke transfer at \(-1\)}. \src{Lib/SiegelWeil/Formula/LagrangianZeroTerms.lean\#L756-L891}{The common phase of the vacuum terms}.}
\noteline{10.17}{17}{\src{Lib/SiegelWeil/Formula/LagrangianZeroTerms.lean\#L631-L658}{Strict comparison of paired factors}. \src{Lib/SiegelWeil/Formula/LagrangianZeroTerms.lean\#L901-L959}{Maximum principle, scaling, and density}. \src{Lib/SiegelWeil/Formula/WeilRange.lean\#L757-L823}{The two invariant functionals}.}
\noteline{10.18}{18}{\src{Lib/SiegelWeil/RallisInner/Doubling.lean\#L116-L126}{The two hypotheses on the datum}. \src{Lib/SiegelWeil/RallisInner/Doubling.lean\#L239-L310}{The shifts and their coordinates at \(g\)}. \src{Lib/Weil/DualPair/ShortIsotropicVectors.lean\#L135-L180}{Vanishing on the supporting cosets}.}
\noteline{10.19}{19}{\src{Lib/SiegelWeil/RallisInner/Doubling.lean\#L194-L235}{The graph zero term as a matrix coefficient}. \src{Lib/SiegelWeil/RallisInner/Doubling.lean\#L315-L416}{Identification of the graphs and the doubled integral}.}
\noteline{10.20}{20}{\src{Lib/SiegelWeil/RallisInner/Doubling.lean\#L422-L505}{Absolute estimates, heights and quotient unfolding}. \src{Lib/SiegelWeil/RallisInner/Doubling.lean\#L520-L562}{Petersson normalization and the resulting integral}.}
\noteline{10.21}{21}{\src{Lib/DiscreteSeries/FlenstedJensen/PsiKTypes.lean\#L468-L480}{The lowest compact type in the coefficient}. \src{Lib/DiscreteSeries/FlenstedJensen/SquareIntegrability.lean\#L396-L424}{The compact projection under the group integral}. \src{Lib/DiscreteSeries/FlenstedJensen/SquareIntegrability.lean\#L530-L579}{Integrability, formal degree and Schur orthogonality}.}
\noteline{10.22}{22}{\src{Lib/Adelic/AutomorphicForms/RealConvolution.lean\#L282-L325}{Right convolution with a translate between two projections}. \src{Lib/Adelic/AutomorphicForms/RealConvolution.lean\#L487-L514}{Insertion in the Petersson product}. \src{Lib/SiegelWeil/RallisInner/SharpData.lean\#L359-L410}{The real factor of the sharp-data identity}.}
\noteline{10.23}{23}{\src{Lib/Weil/Lift.lean\#L590-L618}{Local factors of the product integral}. \src{Lib/Weil/Lift.lean\#L690-L740}{Order of the two groups and change of sign of the parameters}. \src{Lib/DiscreteSeries/Theta/LocalCorrespondence.lean\#L414-L442}{Matching and existence of the joint harmonic}.}
\noteline{10.24}{24}{\src{Lib/DiscreteSeries/FlenstedJensen/PsiKTypes.lean\#L87-L102}{Polar expression in the chosen normalization}. \src{Lib/DiscreteSeries/Rigidity/KernelFixes.lean\#L345-L385}{The two weighted compact projections and their integrability}. \src{Lib/DiscreteSeries/Rigidity/KernelFixes.lean\#L386-L430}{The change of variables and the positive integral}.}
\noteline{10.25}{25}{\src{Lib/Weil/Fock/PolynomialVectors.lean\#L143-L198}{Fischer adjunction and the factorial norms}. \src{Lib/Weil/Fock/OscillatorPair.lean\#L182-L225}{Harmonic recurrence and vacuum coefficient}. \src{Lib/Weil/Fock/OscillatorPair.lean\#L253-L293}{Induction proving the matrix-coefficient formula}. \src{Lib/Weil/Fock/OscillatorPair.lean\#L424-L463}{Positive-definite substitution and its norm square}. \src{Lib/DiscreteSeries/Theta/LocalCorrespondence.lean\#L659-L668}{Application to the actual joint harmonic}.}
\noteline{10.26}{26}{\src{Lib/GLn/SchurModule/SchurOrthogonality.lean\#L319-L328}{Normalization of the highest vector}. \src{Lib/GLn/SchurModule/WeylKernel.lean\#L322-L341}{Polynomial of the evaluation vector}. \src{New/OneLift/HeckeValues/CharacterLift.lean\#L303-L355}{Renaming the variables and the Schur matrix coefficient}.}
\noteline{10.27}{27}{\src{Lib/Weil/CompactLift/ThetaKernelDefs.lean\#L259-L301}{Determinant twists and the source matrix substitution}. \src{New/OneLift/HeckeValues/CharacterLift.lean\#L449-L462}{The one-row hook number}. \src{New/OneLift/HeckeValues/CharacterLiftNorm.lean\#L421-L453}{The real integral with the source volume}. \src{New/OneLift/HeckeValues/CharacterLiftNorm.lean\#L469-L515}{Taking the finite constant outside the integral}.}
\noteline{10.28}{28}{\src{Lib/GLn/SchurModule/SchurOrthogonality.lean\#L278-L317}{Compact averaging and the trace of a rank-one map}. \src{Lib/Adelic/AutomorphicForms/AutomorphicAsAlgebraic.lean\#L369-L410}{Double-coset measure and the conversion to a Petersson product}.}
\noteline{10.29}{29}{\src{New/OneLift/HeckeValues/CharacterLiftNorm.lean\#L63-L92}{The character and the two sets of places in the lift data}.}
\noteline{10.30}{30}{\src{Lib/SiegelWeil/UnramifiedSeries/VacuumCoefficients.lean\#L55-L89}{Definition of the vacuum coefficient and supported representatives}. \src{Lib/SiegelWeil/RallisInner/SharpData.lean\#L118-L153}{The sharp mass and the coefficient on a translate}. \src{Lib/SiegelWeil/RallisInner/SharpData.lean\#L221-L250}{Integration over the finite cosets}.}
\noteline{10.31}{31}{\src{Lib/ClassField/Avatar/HeckeBaseChange.lean\#L125-L154}{Unramifiedness at a pair and the inverse uniformizer value}.}
\noteline{10.32}{32}{\src{Lib/SiegelWeil/RallisInner/SharpData.lean\#L759-L785}{The saturated line lattice and its multiplier}. \src{Lib/SiegelWeil/RallisInner/SharpData.lean\#L787-L835}{Cancellation of the splitting character on the two halves}. \src{Lib/SiegelWeil/RallisInner/SharpData.lean\#L872-L910}{The integer cosets and their geometric sums}.}
\noteline{10.33}{33}{\src{Lib/ClassField/HeckeCharacter/SplitTameCharacters.lean\#L121-L132}{Base change on the full multiplicative group at a fixed place}. \src{Lib/SiegelWeil/RallisInner/SharpData.lean\#L849-L871}{The unique local coset}.}
\noteline{10.34}{34}{\src{Lib/Weil/Construction/WeilData.lean\#L63-L89}{The root of unity and the splitting character in the Weil data}. \src{Lib/Adelic/Adeles/Basic.lean\#L451-L483}{An integral element differing from its conjugate by a unit}.}
\noteline{10.35}{35}{\src{Lib/SiegelWeil/UnramifiedSeries/RallisSeries.lean\#L557-L581}{Multiplicativity and exhaustion by finite sets}. \src{Lib/ClassField/Avatar/HeckeBaseChange.lean\#L156-L177}{Inversion of the character by reindexing its primes}.}
\noteline{10.36}{36}{\src{Lib/SiegelWeil/RallisInner/SharpData.lean\#L736-L757}{Paired signs}. \src{Lib/SiegelWeil/RallisInner/SharpData.lean\#L937-L980}{The partial Hecke product}. \src{New/OneLift/HeckeValues/CharacterLiftNorm.lean\#L524-L553}{Removing both places above each marked prime}.}
\noteline{10.37}{37}{\src{New/OneLift/HeckeValues/CharacterLiftNorm.lean\#L155-L185}{The relation between the class sum and the quotient integral}. \src{New/OneLift/HeckeValues/CharacterLiftNorm.lean\#L308-L340}{Source and target level indices}. \src{New/OneLift/HeckeValues/CharacterLiftNorm.lean\#L383-L419}{Combining the corrections at the two kinds of places}. \src{New/OneLift/HeckeValues/CharacterLiftNorm.lean\#L469-L515}{The source constant and specialization of the real integral}. \src{New/OneLift/HeckeValues/CharacterLiftNorm.lean\#L569-L596}{The class-sum norm}. \src{New/OneLift/HeckeValues/CharacterLiftNorm.lean\#L94-L108}{The subsequent normalization of the critical lift}.}
\noteline{10.38}{38}{\src{Lib/ClassField/HeckeCharacter/SplitTameCharacters.lean\#L211-L229}{The support hypotheses}.}
\noteline{10.39}{39}{\src{Lib/SiegelWeil/Tower/ConstantTerm.lean\#L62-L151}{The central character and zero slice}. \src{Lib/SiegelWeil/Tower/LargeSideConstantTerm.lean\#L137-L159}{The primitive support at zero}.}
\noteline{10.40}{40}{\src{Lib/SiegelWeil/Tower/ConstantTerm.lean\#L143-L179}{The source average at a multiple of the chosen column}. \src{Lib/SiegelWeil/Tower/ConstantTerm.lean\#L562-L605}{Grouping by rational span}.}
\noteline{10.41}{41}{\src{Lib/SiegelWeil/Tower/ConstantTerm.lean\#L299-L359}{The two Jacobians}. \src{Lib/SiegelWeil/Tower/ConstantTerm.lean\#L502-L550}{The unfolded Haar integral}. \src{Lib/SiegelWeil/Tower/ConstantTerm.lean\#L643-L721}{Unfolding the nonzero orbit}.}
\noteline{10.42}{42}{\src{Lib/Weil/MixedModel/ZeroTerm.lean\#L681-L740}{Absolute summability and Fourier coefficients of the rational slices}. \src{Lib/SiegelWeil/Tower/LineTerms.lean\#L193-L273}{Two polarizations of the common zero term}. \src{Lib/SiegelWeil/Tower/LargeSideConstantTerm.lean\#L225-L318}{Levi factors and Poisson summation}.}
\noteline{10.43}{43}{\src{Lib/Weil/MixedModel/FockSections.lean\#L53-L130}{The coherent-state integral and rapid slice estimate}. \src{Lib/SiegelWeil/Tower/RealConeIntegral.lean\#L77-L153}{Restriction to the two archimedean lines}.}
\noteline{10.44}{44}{\src{Lib/SiegelWeil/Tower/RealConeIntegral.lean\#L155-L253}{The integrable moments}. \src{Lib/SiegelWeil/Tower/LargeSideConstantTerm.lean\#L164-L207}{The uniform estimate on the other group}.}
\noteline{10.45}{45}{\src{Lib/SiegelWeil/Tower/LargeSideConstantTerm.lean\#L520-L647}{Factoring the integral and its endpoint}. \src{Lib/SiegelWeil/Tower/LargeSideConstantTerm.lean\#L650-L718}{Mellin convergence and unfolding at the rank exponent}. \src{Lib/SiegelWeil/Tower/LargeSideConstantTerm.lean\#L817-L888}{The two resulting plane identities}.}
\noteline{10.46}{46}{\comsrc{Lib/SiegelWeil/Tower/TestVector.lean\#L342-L381}{Determinant vectors}. \comsrc{Lib/SiegelWeil/Tower/TestVector.lean\#L607-L715}{Their slices and positive cone integral}. \comsrc{New/Determinant/Exchange/VanishingDescends.lean\#L187-L248}{Descent with standard row data}.}
\noteline{10.47}{47}{\src{Lib/ClassField/HeckeCharacter/SplitTameCharacters.lean\#L400-L428}{The nontrivial residue characters at the marked places}. \src{Lib/Weil/MixedModel/Slices.lean\#L291-L325}{Conductor of the additive character in the slice}.}
\noteline{10.48}{48}{\src{Lib/Weil/TameGaussSum.lean\#L147-L213}{The transform, including vanishing outside the unit shell}.}
\noteline{10.49}{49}{\src{Lib/SiegelWeil/Tower/LineGaussSums.lean\#L298-L346}{The integral frame matrices and their determinant contribution}.}
\noteline{10.50}{50}{\src{Lib/SiegelWeil/Tower/LineGaussSums.lean\#L347-L384}{The upper triangular matrix on either side}. \src{Lib/SiegelWeil/Tower/LineGaussSums.lean\#L397-L429}{Expression of the original datum in this matrix}. \src{Lib/Weil/MixedModel/Slices.lean\#L354-L409}{The mixed slice as a product of Fourier integrals}.}
\noteline{10.51}{51}{\src{Lib/SiegelWeil/Tower/LineGaussSums.lean\#L431-L508}{The reduced datum and its exact support}. \src{Lib/SiegelWeil/Tower/LargeSideConstantTerm.lean\#L386-L424}{The line factor in the theta kernel}.}
\noteline{10.52}{52}{\src{Lib/ClassField/HeckeCharacter/TameModulus.lean\#L388-L426}{Unit cancellation and the idele supported on the conjugate places}. \src{Lib/SiegelWeil/Tower/FiniteTateIntegral.lean\#L129-L185}{The translation identity and the resulting valuations}.}
\noteline{10.53}{53}{\src{Lib/ClassField/HeckeL/ContinuationAtOne.lean\#L373-L390}{Tate's integral with the unit volume retained}. \src{Lib/SiegelWeil/Tower/FiniteTateIntegral.lean\#L186-L235}{Application to the unshifted line function}.}
\noteline{10.54}{54}{\src{Lib/SiegelWeil/Tower/FiniteTateIntegral.lean\#L238-L260}{The two Poisson variables}. \src{Lib/SiegelWeil/Tower/LargeSideConstantTerm.lean\#L729-L850}{The character-independent modulus and the Levi hypotheses}.}
\noteline{10.55}{55}{\src{Lib/SiegelWeil/Tower/LargeSideConstantTerm.lean\#L500-L520}{Continuity of the real cone integral}. \src{Lib/SiegelWeil/Tower/LargeSideConstantTerm.lean\#L601-L682}{The endpoint argument on both sides of Poisson summation}.}
\noteline{10.56}{56}{\src{New/Determinant/Norms/AdjointThetaLift.lean\#L525-L586}{The selected row, its unit lift and the action of the level}.}
\noteline{10.57}{57}{\src{Lib/SiegelWeil/RallisInner/SharpData.lean\#L501-L573}{Primitivity and disjointness of the support after translation}. \src{Lib/SiegelWeil/RallisInner/SharpData.lean\#L576-L614}{The matrix coefficient on and off the level}. \src{New/OneLift/HeckeValues/CharacterLiftNorm.lean\#L308-L386}{The exact mass of the rank-one block datum}.}
\noteline{10.58}{58}{\src{Lib/Weil/Unramified/HeckeGL.lean}{Residue kernels and elementary divisors}. \src{Lib/Weil/Unramified/HeckeTransfer.lean\#L226-L257}{Translation of a split vacuum}. \src{New/OneLift/HeckeValues/PeriodAtOne.lean\#L43-L61}{The character and the determinant polynomial}.}
\noteline{10.59}{59}{\src{New/OneLift/HeckeValues/PeriodAtOne.lean\#L208-L254}{Subspace sum on the actual determinant datum}.}
\noteline{10.60}{60}{\src{New/OneLift/HeckeValues/PeriodAtOne.lean\#L133-L206}{Kernel equivariance and the standard level}. \src{New/OneLift/HeckeValues/PeriodAtOne.lean\#L256-L300}{Taking the class integral}. \src{New/Congruences/Selmer/ClassSums.lean}{Specialization to the pair of weight \(c-1\)}.}
\noteline{10.61}{61}{\src{New/UpperBound/Kernels/ThetaLiftGauge.lean\#L417-L473}{Fourier inversion on the finite vacuum datum}.}
\noteline{10.62}{62}{\src{New/UpperBound/Kernels/DefiniteKernelSize.lean\#L116-L187}{Explicit scale and integral jets}.}
\noteline{10.63}{63}{\src{New/OneLift/HeckeValues/AlgebraicPartIntegrality.lean\#L171-L206}{Common integer for the point-kernel gauge}. \src{New/UpperBound/Kernels/ThetaLiftGauge.lean}{Changes of gauge and integral translation}.}
\noteline{10.64}{64}{\src{Lib/SiegelWeil/Tower/FiniteTateIntegral.lean\#L226-L234}{The angular type used in the determinant period}.}
\noteline{10.65}{65}{\src{Lib/ClassField/HeckeL/ContinuationAtOne.lean\#L724-L950}{Continuation by cancellation and the zero-order argument}. \src{Lib/ClassField/HeckeL/EulerProduct.lean\#L539-L568}{The principal prime mean}.}
\noteline{10.66}{66}{\src{Lib/Weil/Fock/WeightedNorms.lean\#L59-L84}{Fock moment and measure}. \src{Lib/Weil/DualPair/KernelDecay.lean\#L364-L410}{Surjective support and both exponents}. \src{Lib/Weil/DualPair/KernelDecay.lean\#L412-L459}{The injective variant}.}
\noteline{10.67}{67}{\src{Lib/Weil/DualPair/ThetaLiftConvergence.lean\#L71-L156}{The height-halving dilation and compact reduction}. \src{Lib/Weil/DualPair/ThetaLiftConvergence.lean\#L184-L219}{The cost of the real factors}. \src{Lib/Weil/DualPair/ThetaLiftConvergence.lean\#L346-L367}{Two-variable moderate growth}.}
\noteline{10.68}{68}{\src{Lib/Weil/DualPair/ShortIsotropicVectors.lean\#L104-L123}{The short isotropic column}. \src{Lib/Adelic/Reduction/FirstMinimumShortBases.lean\#L787-L836}{Transference and spanning over \(E\)}. \src{Lib/Adelic/Reduction/FirstMinimumShortBases.lean\#L843-L885}{Selection and size of the basis}.}
\noteline{10.69}{69}{\src{Lib/Weil/DualPair/KernelDecay.lean\#L136-L195}{Integral coordinates and the length product}. \src{Lib/Weil/DualPair/KernelDecay.lean\#L197-L214}{The rank-one Gram matrix}.}
\noteline{10.70}{70}{\src{Lib/Weil/DualPair/ShortIsotropicVectors.lean\#L60-L102}{The nonzero row after changing frames}. \src{Lib/Weil/DualPair/ShortIsotropicVectors.lean\#L132-L213}{Characters on the marked cosets}. \src{Lib/Weil/DualPair/KernelDecay.lean\#L258-L300}{Identification and vanishing of the isotropic mean}.}
\noteline{10.71}{71}{\src{Lib/Weil/DualPair/ShortIsotropicVectors.lean\#L403-L478}{Replacing adelic shifts by bounded real shifts}. \src{Lib/Weil/DualPair/ShortIsotropicVectors.lean\#L544-L590}{Subtracting the Haar mean and controlling its norm}.}
\noteline{10.72}{72}{\src{Lib/Weil/Fock/WeightedNorms.lean\#L464-L502}{Differentiating the translated polynomial vector}. \src{Lib/Weil/Fock/WeightedNorms.lean\#L504-L606}{Induction through all orders}. \src{Lib/Weil/Fock/WeightedNorms.lean\#L113-L128}{Superposition on finite-moment vectors}.}
\noteline{10.73}{73}{\src{Lib/Weil/DualPair/KernelDecay.lean\#L302-L362}{Combining moderate growth with the short length}.}
\noteline{10.74}{74}{\src{Lib/SiegelWeil/Tower/LineGaussSums.lean\#L68-L90}{The determinant multiplier}. \src{Lib/SiegelWeil/Tower/LineGaussSums.lean\#L197-L212}{Its primitive row support}. \src{Lib/SiegelWeil/Tower/LargeSideConstantTerm.lean\#L164-L209}{The zero-exponent application}.}
\section*{Chapter 11. The comparison period and composition of lifts}
\noteline{11.1}{1}{\src{New/OneLift/ComparisonOperator/ProductsSpan.lean\#L778-L853}{Reproduction at the complementary rank and the single-class test}.}
\noteline{11.2}{2}{\src{Lib/DiscreteSeries/Holomorphic/PeaksAwayFromCusp.lean}{The transvection, triangular factor and disk estimate}.}
\noteline{11.3}{3}{\src{New/OneLift/ComparisonOperator/ThetaKernelEstimates.lean}{The dilated Gaussian estimate and separated-lattice sum}. \src{New/OneLift/ComparisonOperator/ProductsSpan.lean\#L416-L694}{Expansion in the allowed tests, carrier integral and quantitative approximation}.}
\noteline{11.4}{4}{\src{New/OneLift/ComparisonOperator/ThetaKernelEstimates.lean\#L976-L1050}{Circle mean and error after reproduction}.}
\noteline{11.5}{5}{\src{Lib/GLn/SchurModule/GramMatrix.lean}{Hook multiple of the inverse Gram matrix and the integral tight frame}.}
\noteline{11.6}{6}{\src{New/OneLift/ComparisonOperator/GluingMaps.lean}{Gluing in the divided lattice and the estimates for its finite sums}. \src{New/OneLift/LiftVolumeHeight.lean\#L290-L551}{Conjugated frames, class denominators and the two norm estimates}.}
\noteline{11.7}{7}{\src{New/OneLift/FalseLifts/BottomOfTower.lean}{Extension from the block datum and the hyperplane relation}. \src{New/OneLift/FalseLifts/HyperbolicPlaneStep.lean}{Zero term, bounded remainder, transfer and the small-root inequality}.}
\noteline{11.8}{8}{\src{New/OneLift/FalseLifts/PeriodsOfFalseLifts.lean}{Block unfolding, level invariance and the resulting class sums}.}
\noteline{11.9}{9}{\src{New/OneLift/HeckeValues/AlgebraicPartSize.lean\#L284-L385}{The common constant and the two one-sided remainder bounds}. \src{New/OneLift/HeckeValues/AlgebraicPartSize.lean\#L408-L521}{Empty-row character, integral multiplier and jet saving}.}
\noteline{11.10}{10}{\src{New/Determinant/Exchange/OneSquare.lean\#L719-L845}{Splitting the exchange into the two paths}.}
\noteline{11.11}{11}{\src{New/Determinant/Chain/SquareLaw.lean\#L146-L216}{The point identity and its average}. \src{New/Determinant/Chain/ChainNorms.lean\#L142-L250}{The first square and its Hecke values}.}
\noteline{11.12}{12}{\src{New/Determinant/Chain/SquareLaw.lean\#L388-L505}{Fourier transformation, corner Gauss sum, and the reversed square}. \src{New/Determinant/Chain/ChainNorms.lean\#L378-L578}{Kernel reversal and the adjoint identity}.}
\noteline{11.13}{13}{\comsrc{New/Determinant/Chain/TriangleLifts.lean\#L446-L517}{The level and first-edge injectivity}.}
\noteline{11.14}{14}{\src{New/Determinant/Chain/TriangleNorms.lean\#L201-L260}{The two inductions along rows}. \src{New/Determinant/Chain/TriangleNorms.lean\#L398-L494}{Gram recurrence and the induction down the triangle}.}
\noteline{11.15}{15}{\src{New/Determinant/Chain/ChainNorms.lean\#L683-L787}{Flag orders, level indices and the final inner product}. \src{New/Determinant/Chain/NormOfChain.lean\#L145-L246}{The Rallis operators and their normalization}. \src{New/Determinant/Chain/NormOfChain.lean\#L291-L428}{Taking determinants and dividing by the trivial twist}.}
\noteline{11.16}{16}{\src{New/Determinant/AuxiliaryForms/DeepLevelPoincarePeriod.lean}{Separation by the last column and the identity term in the period}. \src{New/Determinant/AuxiliaryForms/DeepLevelVanishing.lean}{The two half-bounds, smoothing and the choice of the real translate}. \src{New/Determinant/AuxiliaryForms/FirstPath.lean}{Surjectivity of the transposed chain and the pairing on the first path}.}
\noteline{11.17}{17}{\src{New/Determinant/AuxiliaryForms/Turns.lean}{Intermediate basis changes and transfer off the diagonal}. \src{New/Determinant/AuxiliaryForms/DiagonalTurns.lean}{Diagonal reversal and composition through the two grids}.}
\noteline{11.18}{18}{\src{New/Determinant/AuxiliaryForms/SpanningForms.lean\#L66-L385}{Clearing one kernel and the flip scalar, iterating on two paths, and counting the last-path terms}.}
\noteline{11.19}{19}{\src{New/Determinant/AuxiliaryForms/SpanningForms.lean\#L426-L773}{The dual hook, class-supported tests, common scaling and the spanning family}. \src{New/UpperBound/Size/ShortColumns.lean\#L363-L429}{Short bases and the spanning evaluations at a free level}. \src{New/Determinant/Twist/LongForms.lean}{Projection onto the good target and the simultaneous basis completion}.}
\section*{Chapter 12. Auxiliary characters and global Selmer bounds}
\noteline{12.1}{1}{\src{New/OneLift/HeckeValues/AlgebraicPartIntegrality.lean\#L302-L389}{Gram expansion and integral coefficients}. \src{New/Congruences/Selmer/ClassSums.lean\#L92-L152}{Uniform normalization and the two Euler factors}.}
\noteline{12.2}{2}{\src{New/Congruences/Selmer/ClassSums.lean\#L162-L300}{Rank-two norm and finite class expansion}. \src{New/OneLift/HeckeValues/PeriodAtOne.lean}{Rank-two insertion at an unramified place}.}
\noteline{12.3}{3}{\src{New/Congruences/Selmer/StickelbergerSystems.lean\#L98-L174}{Evaluation on a norm-one class using the Artin symbol of a preimage}. \src{New/Congruences/Selmer/LValueSystems.lean\#L153-L229}{Finite ideles, integral group rings and evaluation}.}
\noteline{12.4}{4}{\src{New/Congruences/Selmer/ClassSums.lean\#L390-L454}{Identifying the finite Galois and tame characters}. \src{New/Congruences/Selmer/LValueSystems.lean\#L249-L363}{The norm relation and evaluation at twists}.}
\noteline{12.5}{5}{\src{New/Congruences/Selmer/LValueSystems.lean\#L233-L245}{Tame generator and valuation cost}.}
\noteline{12.6}{6}{\src{Lib/PadicHodge/Characters.lean\#L66-L151}{The principal-idele correction and continuous avatar}.}
\noteline{12.7}{7}{\src{Lib/PadicHodge/Characters.lean\#L149-L176}{Local reciprocity away from the correction}. \src{Lib/PadicHodge/Characters.lean\#L1005-L1043}{The geometric Frobenius convention}.}
\noteline{12.8}{8}{\src{New/Congruences/Selmer/LValueSystems.lean\#L56-L151}{Unramifiedness, tame evaluation and the inverse Frobenius}.}
\noteline{12.9}{9}{\src{New/Congruences/Selmer/ClassSums.lean\#L430-L475}{Realizing the admissible class characters by tame twists}.}
\noteline{12.10}{10}{\src{New/Congruences/Selmer/LValueSystems.lean\#L153-L229}{Integral coefficients and evaluation at all admissible characters}.}
\noteline{12.11}{11}{\src{Lib/PadicHodge/Characters.lean\#L984-L1025}{Finite inertia and the scalar of the project's linear period operator}.}
\noteline{12.12}{12}{\src{Lib/ClassField/Avatar/LocalGlobalCompatibility.lean\#L115-L226}{Split-place comparison and descent in Lean}.}
\noteline{12.13}{13}{\src{Lib/ClassField/GreenbergWiles/Relations.lean\#L498-L530}{Avoidance and the split-prime conditions}.}
\noteline{12.14}{14}{\src{Lib/ClassField/GreenbergWiles/Relations.lean\#L498-L582}{Formal construction of the auxiliary set}.}
\noteline{12.15}{15}{\src{Lib/ClassField/GreenbergWiles/Relations.lean\#L70-L140}{The local conditions and the required test at future primes}.}
\noteline{12.16}{16}{\src{Lib/ClassField/GreenbergWiles/SplitPrimeReciprocity.lean\#L267-L319}{The chosen root of unity and the local generators}.}
\noteline{12.17}{17}{\src{Lib/ClassField/GreenbergWiles/Relations.lean\#L584-L615}{The inequality at the interface with descent}.}
\noteline{12.18}{18}{\src{New/Congruences/Selmer/SelmerData.lean}{The Selmer datum, its normalizing scalar, local ranks and the three allowances}.}
\noteline{12.19}{19}{\src{New/Congruences/Selmer/AuxiliaryPrimes.lean}{The Kummer splitting conditions, idele character and choice of Frobenius lift}. \src{New/Congruences/Selmer/Mirror.lean}{Admissible new primes in neighborhoods of the required translates of the commutator group}.}
\noteline{12.20}{20}{\src{New/Congruences/Selmer/Descent.lean}{The slope lemma, comparison at all levels and the bound independent of the character order}.}
\noteline{12.21}{21}{\src{Lib/ClassField/GreenbergWiles/Relations.lean}{The global counting theorem for relations at split primes}. \src{New/Congruences/Selmer/Count.lean}{The local length estimates and cancellation of the global cocycle and coboundary terms}.}
\noteline{12.22}{22}{\src{New/Congruences/Selmer/Mirror.lean}{Paired reciprocity, conjugation by the scalar and the injection of relations into the global Hom module}.}
\noteline{12.23}{23}{\src{New/Congruences/Selmer/RootsOfUnity.lean}{Root selection with a uniform order bound}.}
\noteline{12.24}{24}{\src{Lib/General/NumberTheory/GaussValuation.lean}{Multiplicativity and simultaneous attainment}.}
\noteline{12.25}{25}{\src{New/Congruences/Selmer/RootsOfUnity.lean}{Cyclotomic radii and truncated averaging}.}
\noteline{12.26}{26}{\src{Lib/General/NumberTheory/GaussValuation.lean}{Radicals of \(1+x\) and cyclotomic valuations}. \src{New/Congruences/Selmer/GeneralTheorem.lean}{Specialization simultaneously for powers \(1\) through \(m\)}.}
\noteline{12.27}{27}{\src{New/Congruences/Selmer/OneSystem.lean}{Uniform finite quotient and polynomial relation matrix}.}
\noteline{12.28}{28}{\src{Lib/General/Algebra/SharpMinors.lean}{Maximal minors with a finite torsion level}.}
\noteline{12.29}{29}{\src{New/Congruences/Selmer/Untied.lean}{Affine tame relation and the loss on removing the ties}.}
\noteline{12.30}{30}{\src{New/Congruences/Selmer/Line.lean}{Restriction of the tied minor and passage from finite depth}.}
\noteline{12.31}{31}{\src{New/Congruences/Selmer/StickelbergerLine.lean}{First-prime factors and the torsion-polynomial congruence}.}
\noteline{12.32}{32}{\src{New/Congruences/Selmer/Descent.lean}{Construction and quantitative bound for the tied primes}. \src{New/Congruences/Selmer/Line.lean}{Order of choices and the torsion-level hypothesis}.}
\noteline{12.33}{33}{\src{New/Congruences/Selmer/GeneralTheorem.lean}{The simultaneous anticyclotomic twist}. \src{New/Congruences/Selmer/LValueSystems.lean}{Systems attached to the critical values}.}
\noteline{12.34}{34}{\src{Lib/Cyclotomic/Stickelberger.lean\#L240-L264}{The precise assertions about a Jacobi family}.}
\noteline{12.35}{35}{\src{Lib/Cyclotomic/Stickelberger.lean\#L93-L192}{Formal valuation coordinate and the choice of additive character}.}
\noteline{12.36}{36}{\src{Lib/Cyclotomic/Stickelberger.lean\#L287-L317}{Compatible character parameters at all divisor levels}. \src{Lib/Cyclotomic/Stickelberger.lean\#L336-L379}{The quotient and its order at the chosen prime}. \src{Lib/Cyclotomic/Stickelberger.lean\#L380-L395}{Absolute values and the Jacobi-family assertions}.}
\noteline{12.37}{37}{\src{Lib/Cyclotomic/Stickelberger.lean\#L45-L73}{The distribution of the floor difference, including reindexing}. \src{Lib/Cyclotomic/Stickelberger.lean\#L194-L238}{Vanishing at the conjugate primes and Kronecker's argument}.}
\noteline{12.38}{38}{\src{Lib/Cyclotomic/Stickelberger.lean\#L319-L328}{Reduction of a Gauss sum after inserting a character of order \(q\)}. \src{Lib/Cyclotomic/Stickelberger.lean\#L396-L414}{Passage to ratios of units}.}
\noteline{12.39}{39}{\src{Lib/ClassField/Places/TameInertia.lean\#L64-L109}{Principal units and the inertia character on a root}. \src{Lib/ClassField/Places/TameInertia.lean\#L121-L189}{The integral order of the orbit product and transitivity}.}
\noteline{12.40}{40}{\src{Lib/Cyclotomic/Kolyvagin/DerivativeCocycles.lean\#L465-L493}{The general commutator exponent}. \src{Lib/Cyclotomic/Kolyvagin/DerivativeCocycles.lean\#L719-L755}{The specialized order and the bottom derivative}.}
\noteline{12.41}{41}{\src{Lib/ClassField/Places/TameInertia.lean\#L426-L463}{Surjectivity to the cyclotomic unit classes in inertia}. \src{Lib/Cyclotomic/Kolyvagin/DerivativeCocycles.lean\#L114-L124}{Carrying the \(q\)-adic cyclotomic action to absolute inertia over \(\mathbb Q\)}. \src{Lib/Cyclotomic/Kolyvagin/DerivativeCocycles.lean\#L439-L461}{The action on the two summands of the translated index}.}
\noteline{12.42}{42}{\src{Lib/Cyclotomic/Stickelberger.lean}{The prime-field Gauss sums, valuation parametrization and Jacobi relations}.}
\noteline{12.43}{43}{\src{Lib/Cyclotomic/Kolyvagin/DerivativeCocycles.lean}{Shift independence, reindexing and the ramification of the derivative sums}.}
\noteline{12.44}{44}{\src{Lib/Cyclotomic/Kolyvagin/DerivativeCocycles.lean}{Telescoping at \(q\) and the bottom derivative at \(p\)}.}
\noteline{12.45}{45}{\src{Lib/Cyclotomic/Kolyvagin/ReciprocityComponentBounds.lean}{Artin reciprocity, the Kummer coordinates, and reduction modulo \(\ell\)}.}
\noteline{12.46}{46}{\src{Lib/Cyclotomic/Kolyvagin/ReciprocityComponentBounds.lean}{Simultaneous detecting primes, tight sets and the descending induction}.}
\noteline{12.47}{47}{\src{Lib/Cyclotomic/OddSha.lean}{Character transform, large and small coefficient primes, uniform length}.}
\noteline{12.48}{48}{\src{New/Congruences/FirstOrder/RibetElement.lean}{The odd, Selmer, local and auxiliary-place terms for the actual lattices}. \src{New/Congruences/FirstOrder/TateTwistCocycles.lean}{The triangle cost in terms of crystalline colength and the odd module}. \src{New/Congruences/FirstOrder/RibetLemma.lean}{Valuation of the glued determinant at the lift}.}
\noteline{12.49}{49}{\src{Lib/ClassField/Global/ExistenceTheorem.lean\#L926-L984}{The quotient, its order and the lifted generators}.}
\noteline{12.50}{50}{\src{Lib/ClassField/Global/ExistenceTheorem.lean\#L95-L163}{Local power indices and their product}.}
\noteline{12.51}{51}{\comsrc{Lib/General/NumberTheory/SUnitsAndClassNumber.lean\#L182-L200}{The effective class-number bound}. \src{Lib/Cyclotomic/OddSha.lean\#L89-L145}{The cyclotomic discriminant and the numerical bound}.}
\noteline{12.52}{52}{\src{Lib/Cyclotomic/OddSha.lean\#L176-L267}{The induction and the cocycle equation over \(K\)}.}
\noteline{12.53}{53}{\src{Lib/Cyclotomic/OddSha.lean\#L114-L130}{The polynomial bound}. \src{Lib/Cyclotomic/OddSha.lean\#L269-L302}{The cyclotomic-kernel version}.}
\noteline{12.54}{54}{\src{Lib/Cyclotomic/OddSha.lean\#L550-L581}{From evaluation to integral length}. \src{Lib/Cyclotomic/OddSha.lean\#L634-L679}{One integer for all depths, twists and primes}.}
\section*{Chapter 13. Congruences and integral Hecke operators}
\noteline{13.1}{1}{\src{New/Congruences/GaloisRepresentations/AttachedRepresentations.lean}{Representations attached to automorphic eigenforms}. \src{New/Congruences/GaloisRepresentations/RowsAndColumnsFromLabels.lean}{Type, inertia, and the main summand}.}
\noteline{13.2}{2}{\src{New/Congruences/FirstOrder/TraceRelations.lean}{Block-diagonal lattices and local crystalline colength}.}
\noteline{13.3}{3}{\src{Lib/GaloisRep/Main/SelmerModules.lean}{Exact local conditions and parity of the cocycles}. \src{Lib/Cyclotomic/OddSha.lean}{Class field theory at exponent zero}.}
\noteline{13.4}{4}{\src{Lib/Cyclotomic/Voronoi.lean}{Integral Bernoulli matrix, the floor-sum congruence, and its size}.}
\noteline{13.5}{5}{\src{Lib/Cyclotomic/Stickelberger.lean}{Jacobi families and their valuations}. \src{Lib/Cyclotomic/Kolyvagin/ReciprocityComponentBounds.lean}{The Artin pairing and component annihilation}.}
\noteline{13.6}{6}{\src{Lib/Cyclotomic/OddSha.lean}{The character transform and integral recovery over \(E\)}.}
\noteline{13.7}{7}{\src{Lib/Cyclotomic/Kolyvagin/DerivativeCocycles.lean}{Derivative cocycles and the Frobenius--inertia relation}. \src{Lib/Cyclotomic/Kolyvagin/ReciprocityComponentBounds.lean}{The induction giving \(1+v_\ell(X_e)\) tests}.}
\noteline{13.8}{8}{\src{Lib/Cyclotomic/OddSha.lean}{Detection at small and large primes and the common integer}.}
\noteline{13.9}{9}{\src{New/Congruences/FirstOrder/DenominatorBound.lean}{Cyclotomic separation and the allowable Frobenius restrictions}. \src{New/Congruences/FirstOrder/TrianglePrimes.lean}{The prime-selection theorem with its hypotheses}.}
\noteline{13.10}{10}{\src{New/Congruences/FirstOrder/KummerMaps.lean}{Dimension of the reductions and reciprocity over the real subfield}. \src{New/Congruences/FirstOrder/TrianglePrimes.lean}{Separation of exponents and the simultaneous Frobenius prescription}.}
\noteline{13.11}{11}{\src{New/Congruences/FirstOrder/TrianglePrimes.lean}{The single remaining reciprocity term and total triangular cost}.}
\noteline{13.12}{12}{\src{New/Congruences/FirstOrder/TateTwistCocycles.lean}{From paired generators to the even condition and the odd residual length}.}
\noteline{13.13}{13}{\src{New/Congruences/FirstOrder/RibetLemma.lean}{Gluing the block and character rows into a trace-ideal multiplier}. \src{New/Congruences/FirstOrder/RibetElement.lean}{Hypotheses and the precise bound on the squarefree coefficient}. \src{New/Congruences/FirstOrder/TraceRelations.lean}{The four terms of the residual valuation}.}
\noteline{13.14}{14}{\src{New/Congruences/FirstOrder/FittingBlock.lean}{The triangular Fitting-ideal estimate for one block}.}
\noteline{13.15}{15}{\src{New/Congruences/FirstOrder/OnePrime.lean}{Norms, their first-order deformations, and the squarefree coefficient}. \src{New/Congruences/FirstOrder/DenominatorBound.lean}{Change of level and the diagonal operator clearing the projected forms}.}
\noteline{13.16}{16}{\src{New/Congruences/FirstOrder/RibetLemma.lean}{Finite coordinates and the trace relation for each tied row}.}
\noteline{13.17}{17}{\src{New/Congruences/FirstOrder/RibetMethod.lean}{The integral adjugate and right multiplication in the trace quotient}.}
\noteline{13.18}{18}{\src{New/Congruences/FirstOrder/RibetMatrices.lean}{The main block, cocycle equations and evaluation of the normalizer row}.}
\noteline{13.19}{19}{\src{New/Congruences/FirstOrder/FittingBlock.lean}{The three Fitting-ideal operations with their separate valuation costs}.}
\noteline{13.20}{20}{\src{New/Congruences/FirstOrder/RibetMatrices.lean}{The reduced relations on each character and the simultaneous matrix of tied rows}.}
\noteline{13.21}{21}{\src{New/Congruences/FirstOrder/RibetLemma.lean}{The common element, its value at the distinguished lift and the full valuation budget}.}
\noteline{13.22}{22}{\src{New/Congruences/FirstOrder/AuxiliaryPrimeTie.lean}{The deformed norm, exact divisibility condition and cancellation on the lattice}.}
\noteline{13.23}{23}{\src{New/Congruences/FirstOrder/DenominatorBound.lean}{Type projection, false lifts, simultaneous choice of an operator and its determinant}.}
\noteline{13.24}{24}{\src{New/Congruences/FirstOrder/OnePrime.lean}{Reduction on the torsion-free endomorphism ring}.}
\noteline{13.25}{25}{\src{Lib/GaloisRep/TaylorWiles.lean}{Orbit weights for the auxiliary level and Higman's trace identity}.}
\noteline{13.26}{26}{\src{New/Congruences/FirstOrder/OnePrime.lean}{First-entry transfer, weighted sums and translated logarithms}. \src{New/Congruences/FirstOrder/TiedTraces.lean}{Local Artin choices and the diamond operator}.}
\noteline{13.27}{27}{\src{Lib/GaloisRep/TaylorWiles.lean}{Restriction of a local Hecke polynomial to the principal-level lattice}. \src{New/Congruences/FirstOrder/TiedTraces.lean}{Weighted principal-level operators and parahoric invariance}.}
\noteline{13.28}{28}{\src{Lib/GaloisRep/LocalGlobal/IwahoriCompatibility.lean}{Detection by sorted exponents and the trace identity via polynomial interpolation}.}
\noteline{13.29}{29}{\src{New/Congruences/FirstOrder/TiedTraces.lean}{Character decomposition on the two parahorics and the integral congruence}.}
\noteline{13.30}{30}{\src{Lib/LocalGL/Iwahori/HeckeDegrees.lean\#L491-L582}{The forward hyperplane count}. \src{Lib/LocalGL/SphericalChainLemma.lean\#L94-L205}{The inverse count by dual lattices and the endpoint}.}
\noteline{13.31}{31}{\src{Lib/LocalGL/Iwahori/HeckeDegrees.lean}{The hyperplane incidence counts}. \src{Lib/LocalGL/SphericalChainLemma.lean}{The partial sums and the adjoint cofactor calculation}.}
\noteline{13.32}{32}{\src{New/Congruences/FirstOrder/HeckeStrings.lean}{The Galois root, integral complement and separation from one}. \src{New/Congruences/FirstOrder/DenominatorBound.lean}{Combining the triangular primes with the root-chain condition}.}
\noteline{13.33}{33}{\src{New/Congruences/FirstOrder/TraceRelations.lean}{Passing from the chosen split place to inverse-conjugate traces}.}
\noteline{13.34}{34}{\src{New/Congruences/FirstOrder/AuxiliaryPrimeTie.lean}{Construction at each place, product and norm lift}. \src{New/Congruences/FirstOrder/OnePrime.lean}{Finite logarithm difference, stabilizer cost and coefficient induction}.}
\section*{Chapter 14. Local cohomology and the integral period lattice}
\noteline{14.1}{1}{\src{New/Congruences/LocalLoss/Basic.lean}{Relations among the matrices and the integral graph}. \src{New/Congruences/LocalLoss/TateDualBound.lean}{Crystalline cocycles in the graph lattice}.}
\noteline{14.2}{2}{\src{New/Congruences/LocalLoss/TateDualBound.lean}{Integral pairing and the finite local duality index}. \src{Lib/ClassField/LocalDuality/CrystallineDuality.lean}{The cocycle annihilator count for finite local Tate duality}.}
\noteline{14.3}{3}{\src{New/Congruences/LocalLoss/CrystallineCocycles.lean}{Cocycle rank and the estimate from a cup exponent}.}
\noteline{14.4}{4}{\src{New/Congruences/LocalLoss/CrystallineDefect.lean}{Jump windows, sparse letters, and the choice of local charts}.}
\noteline{14.5}{5}{\src{New/Congruences/Lifting/CupOrthogonalityBounds.lean}{Sparse-letter and unramified cup comparisons}. \src{New/Congruences/LocalLoss/CrystallineCocycles.lean}{Application to the two ranges of primes}.}
\noteline{14.6}{6}{\src{New/Congruences/Lifting/CupOrthogonalityBounds.lean}{Zero-defect and low-multiplicity comparison theorems}.}
\noteline{14.7}{7}{\src{New/Congruences/LocalLoss/LocalLattices.lean}{Tame inertial annihilators}. \src{New/Congruences/LocalLoss/CrystallineDefect.lean}{Tame cup exponent against the level-prime allowance}.}
\noteline{14.8}{8}{\src{New/Congruences/FirstOrder/TraceRelations.lean}{Simultaneous local constraints on all summands}. \src{New/Congruences/FitTwistsExist.lean}{The logarithmic sum of local losses}.}
\noteline{14.9}{9}{\src{New/Congruences/LocalLoss/TateDualBound.lean\#L101-L144}{The open embedding, free lattice, finite reduction and admissibility hypotheses}.}
\noteline{14.10}{10}{\src{Lib/PadicHodge/PeriodRings/CrystallineRings.lean\#L56-L78}{The period-ring construction}. \src{Lib/PadicHodge/PeriodRings/DeRhamPeriods.lean\#L208-L228}{Association with \(p\mu\) and the cyclotomic action}.}
\noteline{14.11}{11}{\src{Lib/PadicHodge/BlochKato.lean\#L153-L179}{Hodge bases in the formal calculation}.}
\noteline{14.12}{12}{\src{Lib/PadicHodge/BlochKato.lean\#L153-L179}{The basis and the filtration convention}.}
\noteline{14.13}{13}{\src{Lib/PadicHodge/BlochKato.lean\#L305-L348}{Continuity of the last column and its finite reductions}.}
\noteline{14.14}{14}{\src{Lib/PadicHodge/BlochKato.lean\#L433-L472}{Reduction as cocycles and its kernel}. \src{Lib/ClassField/LocalDuality/CrystallineDuality.lean\#L446-L486}{The complementary jumps and the product of the indices}.}
\noteline{14.15}{15}{\src{Lib/ClassField/LocalDuality/CrystallineDuality.lean\#L414-L445}{The coefficient functional giving a perfect integral form}. \src{Lib/ClassField/LocalDuality/CrystallineDuality.lean\#L487-L531}{The induced roots-of-unity pairing after reduction}.}
\noteline{14.16}{16}{\src{New/Congruences/Selmer/SelmerData.lean\#L78-L92}{The prime-to-\(p\) term uses actual restrictions}.}
\noteline{14.17}{17}{\src{New/Congruences/LocalLoss/TateDualBound.lean\#L73-L100}{Dividing the equivariance error}. \src{New/Congruences/LocalLoss/TateDualBound.lean\#L139-L211}{Orthogonality of the reductions and the length comparison}.}
\noteline{14.18}{18}{\src{New/Congruences/Lifting/ExponentLattices.lean}{Integral exponents and order improvement}. \src{New/Congruences/Lifting/HodgeExpansion.lean}{Hodge evaluations and primitive relations on one line}. \src{New/Congruences/Lifting/SidePoints.lean}{Depth raising, its cost, and transport by Frobenius}.}
\noteline{14.19}{19}{\src{New/Congruences/Lifting/CupProductBound.lean}{The corrected cochain, four-term identity and descent}.}
\noteline{14.20}{20}{\src{New/Congruences/Lifting/SafeMoves.lean}{Pairing criterion on orders and assembly along Frobenius orbits}.}
\noteline{14.21}{21}{\src{New/Congruences/Lifting/LongGap.lean}{Integral depth at a pivot reached through a long gap}. \src{Lib/PadicHodge/Kisin/FrameTheorem.lean}{The integral frame and the \(p^{4\delta}\) nearly invariant span}. \src{New/Congruences/Lifting/CutExponentBounds.lean}{Sparse cuts, the two allowed thick positions and the constants in (CP2)}.}
\noteline{14.22}{22}{\src{New/Congruences/Lifting/CutExponentBounds.lean}{Shifted cuts for arbitrary patterns of the two sizes}.}
\noteline{14.23}{23}{\src{New/Congruences/Lifting/CupProductBound.lean}{Initial monomial, descent, and unramified cokernel bound}.}
\noteline{14.24}{24}{\src{New/Congruences/Lifting/DeepPoints.lean}{Pole-free lifts and the two vanishing arguments}.}
\noteline{14.25}{25}{\src{New/Congruences/Lifting/FewLowJumps.lean}{Pole filtration, factorial bound and simultaneous annihilator}. \src{New/Congruences/Lifting/CupOrthogonalityBounds.lean}{From letter multiplicities to normalized local defect}.}
\noteline{14.26}{26}{\src{New/Congruences/Lifting/TameCharts.lean}{Tame chart and integral descent by the root count}. \src{New/Congruences/Lifting/CupOrthogonalityBounds.lean}{Tensor annihilator and passage to all allowed pairings}.}
\noteline{14.27}{27}{\src{Lib/ClassField/LocalDuality/CrystallineDuality.lean\#L220-L251}{Construction and equivariance of the integral form}. \src{New/Congruences/Lifting/CupProductBound.lean\#L74-L139}{Crystalline extensions, the tensor lattice and the Tate dual}. \src{New/Congruences/Lifting/SidePoints.lean\#L806-L842}{The idempotent and reduction of the evaluated form}.}
\noteline{14.28}{28}{\src{New/Congruences/Lifting/HodgeExpansion.lean\#L123-L182}{The invariant basis, the dual rows and the type shift}. \src{New/Congruences/Lifting/ExponentLattices.lean\#L854-L885}{Reduction from a type bound to the actual jumps}. \src{New/Congruences/Lifting/HodgeExpansion.lean\#L463-L571}{Expansion of a bilinear form and the integral order rule}.}
\noteline{14.29}{29}{\src{New/Congruences/Lifting/HodgeExpansion.lean\#L321-L353}{Standard orders and the Frobenius-translated lattice}. \src{New/Congruences/Lifting/SafeMoves.lean\#L115-L160}{The successive applications to the tensor and to its evaluation}.}
\noteline{14.30}{30}{\src{New/Congruences/Lifting/CupProductBound.lean\#L533-L551}{The integral denominator of a pure Tate power}. \src{New/Congruences/Lifting/ExponentLattices.lean\#L889-L897}{Division of a difference by \(z\tau\) with all orders preserved}. \src{New/Congruences/Lifting/SafeMoves.lean\#L97-L113}{Passage to the small ideal}. \src{New/Congruences/Lifting/SafeMoves.lean\#L157-L219}{Reduction modulo \(p^m\) and the two moves}. \src{New/Congruences/Lifting/SafeMoves.lean\#L322-L384}{The three bounds from the dual cokernel}.}
\noteline{14.31}{31}{\src{New/Congruences/Lifting/SidePoints.lean\#L329-L368}{Orthogonality of the letter projections}. \src{New/Congruences/Lifting/SidePoints.lean\#L777-L805}{Transported depths below the preceding cut}. \src{New/Congruences/Lifting/SafeMoves.lean\#L451-L478}{The product of two differences}. \src{New/Congruences/Lifting/SafeMoves.lean\#L515-L556}{Passage to the one cochain}.}
\noteline{14.32}{32}{\src{New/Congruences/Lifting/TameCharts.lean\#L447-L508}{The simultaneous fixer and its index}. \src{New/Congruences/Lifting/CupOrthogonalityBounds.lean\#L491-L501}{The choice \(f=r!\)}.}
\noteline{14.33}{33}{\src{New/Congruences/Lifting/TameCharts.lean\#L514-L563}{Transitivity on towers and action on the coordinate}. \src{New/Congruences/Lifting/TameCharts.lean\#L573-L634}{The Kummer elements and the tame generator}.}
\noteline{14.34}{34}{\src{Lib/ClassField/Places/TameInertia.lean\#L329-L353}{The finite-exponent quotient detected by a root of the uniformizer}. \src{New/Congruences/Lifting/CupOrthogonalityBounds.lean\#L434-L465}{The open kernel on periods and the application to \(\Gamma'\)}.}
\noteline{14.35}{35}{\src{Lib/PadicHodge/Crystalline/Regularity.lean\#L716-L784}{Descent using the unramified constants inside \(B_{\max}\)}. \src{New/Congruences/Lifting/CupOrthogonalityBounds.lean\#L467-L489}{Descent in the local class with its original jump sets}.}
\noteline{14.36}{36}{\src{New/Congruences/Lifting/TameCharts.lean\#L91-L174}{The relation on conjugated crystalline forms}. \src{New/Congruences/Lifting/TameCharts.lean\#L177-L250}{Divisibility on the frame and the extra factor for the extension}. \src{Lib/PadicHodge/Kisin/Bicommutant.lean\#L554-L571}{Removing the denominators from divisibility by \(z\)}. \src{New/Congruences/Lifting/TameCharts.lean\#L287-L358}{Descent of the lifted monomial including its correction}.}
\noteline{14.37}{37}{\src{New/Congruences/GaloisRepresentations/RowsFromEigensystems.lean\#L65-L112}{Target parameter in the two ranges of indices}. \src{New/Congruences/GaloisRepresentations/RowsFromEigensystems.lean\#L118-L171}{Exact multiplicity and conjugate jumps}.}
\noteline{14.38}{38}{\src{New/Congruences/GaloisRepresentations/RowsFromEigensystems.lean\#L173-L216}{The rank-one pair}.}
\noteline{14.39}{39}{\src{New/Congruences/LocalLoss/CrystallineDefect.lean\#L49-L110}{Window margins and the dual multiplicity}. \src{New/Congruences/LocalLoss/CrystallineDefect.lean\#L552-L620}{Passing to actual jumps of each individual dual}.}
\noteline{14.40}{40}{\src{New/Congruences/LocalLoss/LocalLattices.lean\#L191-L221}{Reading the actual lattice jumps through the local embedding}. \src{New/Congruences/LocalLoss/CrystallineDefect.lean\#L205-L255}{Construction of the chart and thick-letter predicate}.}
\noteline{14.41}{41}{\src{New/Congruences/LocalLoss/Basic.lean\#L180-L275}{Ratios of CM place labels and the signed short-arc criterion}. \src{New/Congruences/LocalLoss/CrystallineDefect.lean\#L168-L202}{Prime-to-\(\ell\) root and the two distances}.}
\noteline{14.42}{42}{\src{New/Congruences/LocalLoss/Basic.lean\#L278-L340}{Common sign over a prime of the conductor}. \src{New/Congruences/LocalLoss/CrystallineDefect.lean\#L224-L293}{Sign choices and the size of the ramified primes}.}
\noteline{14.43}{43}{\src{New/Congruences/LocalLoss/CrystallineCocycles.lean\#L373-L411}{Use of the widened windows and short-arc bound}.}
\noteline{14.44}{44}{\src{New/Congruences/GaloisRepresentations/SelmerHypotheses.lean\#L563-L584}{The scalar which rescales the roots}. \src{New/Congruences/GaloisRepresentations/BaseLevelOfLabel.lean\#L577-L610}{Normalized representation of the projected eigensystem}.}
\noteline{14.45}{45}{\src{New/Congruences/GaloisRepresentations/RowsAndColumnsFromLabels.lean\#L175-L242}{Evaluation of the normalized type on tame inertia}. \src{New/Congruences/GaloisRepresentations/RowsAndColumnsFromLabels.lean\#L245-L274}{Transport to all places above the level prime}.}
\noteline{14.46}{46}{\src{Lib/PadicHodge/Characters.lean\#L215-L235}{Finite order of the inertial period operator}. \src{New/Congruences/LocalLoss/LocalLattices.lean\#L43-L85}{The two annihilators and the integral degree bound}. \src{New/Congruences/LocalLoss/LocalLattices.lean\#L129-L147}{Descent of annihilation to the lattice}. \src{New/Congruences/FirstOrder/TraceRelations.lean\#L631-L661}{Application to the whole labelled family}.}
\noteline{14.47}{47}{\src{Lib/PadicHodge/Kisin/Charts.lean\#L151-L195}{The series evaluation and associated generators}. \src{Lib/PadicHodge/Kisin/KummerChartExistence.lean\#L338-L447}{The fixed local field and the choice in the unramified case}.}
\noteline{14.48}{48}{\src{Lib/PadicHodge/Kisin/KummerChartExistence.lean\#L201-L253}{From the finite transitive sets to the tilt element}. \src{Lib/PadicHodge/Kisin/KummerChartExistence.lean\#L422-L446}{Action within the chart group}.}
\noteline{14.49}{49}{\src{Lib/PadicHodge/Kisin/Charts.lean\#L108-L127}{The chart data}. \src{Lib/PadicHodge/Kisin/Charts.lean\#L282-L344}{Injectivity and the two expressions for the divisor}.}
\noteline{14.50}{50}{\src{Lib/General/Algebra/AdicCompletion.lean\#L522-L581}{Adic coordinates for two generators}. \src{Lib/PadicHodge/Kisin/Charts.lean\#L425-L475}{Application to the chart including the vector \(1\)}.}
\noteline{14.51}{51}{\src{Lib/PadicHodge/PeriodRings/CrystallineRings.lean\#L365-L420}{Reduction and cancellation of monic polynomials}. \src{Lib/PadicHodge/PeriodRings/CrystallineRings.lean\#L466-L493}{Intersection with \(p^nD\)}. \src{Lib/PadicHodge/Kisin/Bicommutant.lean\#L105-L205}{The bounded extension}.}
\noteline{14.52}{52}{\comsrc{Lib/PadicHodge/Kisin/Bicommutant.lean\#L235-L297}{Coordinates in the bicommutant and its retraction}. \comsrc{Lib/PadicHodge/Kisin/Bicommutant.lean\#L356-L376}{Frobenius and the inverse of \(\varphi(E)\)}.}
\noteline{14.53}{53}{\src{Lib/PadicHodge/Kisin/Bicommutant.lean\#L406-L442}{Common finite coordinate expansions}.}
\noteline{14.54}{54}{\src{Lib/PadicHodge/PeriodRings/CrystallineRings.lean\#L419-L466}{Division by the linear polynomial and iteration}. \src{Lib/PadicHodge/Kisin/Bicommutant.lean\#L467-L490}{Projection of the integral remainder}.}
\noteline{14.55}{55}{\src{Lib/PadicHodge/PeriodRings/CrystallineRings.lean\#L682-L713}{Approximation with a power of \(\mu\) in the error}. \src{Lib/PadicHodge/PeriodRings/CrystallineRings.lean\#L857-L886}{The finite de Rham truncations}.}
\noteline{14.56}{56}{\src{Lib/PadicHodge/PeriodRings/Ainf.lean\#L834-L887}{The intersection in the Witt ring}. \src{Lib/PadicHodge/PeriodRings/CrystallineRings.lean\#L721-L748}{Passage to \(B^+\)}.}
\noteline{14.57}{57}{\src{Lib/PadicHodge/Kisin/Charts.lean\#L636-L713}{Readout of components and divisibility at one level}.}
\noteline{14.58}{58}{\src{Lib/PadicHodge/Kisin/Charts.lean\#L565-L609}{Degree of the remainder at level \(n\)}. \src{Lib/PadicHodge/Kisin/Charts.lean\#L716-L754}{Algebraicity of the values and Galois correspondence}.}
\noteline{14.59}{59}{\src{Lib/PadicHodge/Kisin/Charts.lean\#L755-L821}{Fixed-element approximation and integral cancellation}.}
\noteline{14.60}{60}{\src{Lib/PadicHodge/PeriodRings/CrystallineRings.lean\#L978-L993}{The separate basis on \(D\) for the trace limit}. \src{Lib/PadicHodge/Limit/CrystallineLimit.lean\#L478-L500}{Increasing finite coordinate supports}.}
\noteline{14.61}{61}{\src{Lib/PadicHodge/Crystalline/HodgeTateJumps.lean\#L106-L146}{The two integrality conditions, fixedness and residue convention}. \src{Lib/PadicHodge/Crystalline/HodgeTateJumps.lean\#L202-L220}{The open-image hypothesis and general form of the theorem}.}
\noteline{14.62}{62}{\src{Lib/PadicHodge/TateSen.lean\#L349-L400}{The fixed coefficients in a relation of minimal support}. \src{Lib/PadicHodge/Crystalline/HodgeTateJumps.lean\#L221-L282}{The source's determinant argument}.}
\noteline{14.63}{63}{\src{Lib/PadicHodge/Crystalline/HodgeTateJumps.lean\#L308-L350}{Inverse rows and their weights}. \src{Lib/PadicHodge/Kisin/PeriodMatrixBasis.lean\#L722-L778}{Crystalline columns and their Frobenius translates}. \src{Lib/PadicHodge/Kisin/PeriodMatrixBasis.lean\#L799-L828}{The two filtration bounds}. \src{Lib/PadicHodge/Limit/TraceLimit.lean\#L301-L325}{The full group's fixed integral coefficients for the projector}.}
\noteline{14.64}{64}{\src{Lib/PadicHodge/Kisin/Charts.lean\#L250-L273}{The chart and its required Galois image}. \src{Lib/PadicHodge/Kisin/PeriodMatrixBasis.lean\#L105-L135}{The two cyclotomic factors}.}
\noteline{14.65}{65}{\src{Lib/PadicHodge/Kisin/PeriodMatrixBasis.lean\#L721-L777}{Invariant columns and the weights of each Frobenius translate}. \src{Lib/PadicHodge/Kisin/PeriodMatrixBasis.lean\#L779-L822}{Orders of the matrix and its inverse}.}
\noteline{14.66}{66}{\src{Lib/PadicHodge/Crystalline/FixedRatios.lean\#L44-L82}{Removing the poles one at a time}.}
\noteline{14.67}{67}{\src{Lib/PadicHodge/Crystalline/FixedRatios.lean\#L329-L381}{Fixed ratios have unramified coefficients}. \src{Lib/PadicHodge/Kisin/PeriodMatrixBasis.lean\#L954-L998}{The Frobenius factor and its inverse}. \src{Lib/PadicHodge/Kisin/PeriodMatrixBasis.lean\#L1004-L1026}{Lifting the prescribed Galois elements}.}
\noteline{14.68}{68}{\src{Lib/PadicHodge/Kisin/Charts.lean\#L425-L475}{A coordinate equal to the identity on \(S\)}. \src{Lib/PadicHodge/Kisin/Bicommutant.lean\#L86-L144}{Divisibility and bounded extension}. \src{Lib/PadicHodge/Kisin/Bicommutant.lean\#L145-L205}{Passage through \(A_{\max}\) and inversion of \(p\)}.}
\noteline{14.69}{69}{\src{Lib/PadicHodge/Crystalline/FixedRatios.lean\#L134-L159}{Newton iteration for the fixed constants}. \src{Lib/PadicHodge/Crystalline/FixedRatios.lean\#L161-L224}{Divisibility of a sum with fixed coefficients}.}
\noteline{14.70}{70}{\src{Lib/PadicHodge/Kisin/PeriodMatrixBasis.lean\#L866-L891}{The rows of the inverse fixed frame}. \src{Lib/PadicHodge/Kisin/PeriodMatrixBasis.lean\#L893-L940}{Conjugation of the coordinate multipliers}.}
\noteline{14.71}{71}{\src{Lib/PadicHodge/Kisin/PeriodMatrixBasis.lean\#L293-L314}{Descending a dependence coefficient by coefficient}. \src{Lib/PadicHodge/Kisin/PeriodMatrixBasis.lean\#L471-L539}{Regularity and the spanning matrix}. \src{Lib/General/Algebra/AdicCompletion.lean\#L578-L599}{Finite freeness from the adic rank bound}.}
\noteline{14.72}{72}{\src{Lib/PadicHodge/Kisin/PeriodMatrixBasis.lean\#L249-L265}{The \(2h\) cost of division by a factor of \(\mu^h\)}. \src{Lib/PadicHodge/Kisin/PeriodMatrixBasis.lean\#L353-L379}{The two matrix estimates}.}
\noteline{14.73}{73}{\src{Lib/PadicHodge/Kisin/PeriodMatrixBasis.lean\#L269-L291}{The divisible inverse image}. \src{Lib/PadicHodge/Kisin/PeriodMatrixBasis.lean\#L381-L437}{Transporting the coordinate under Frobenius}.}
\noteline{14.74}{74}{\src{Lib/PadicHodge/PeriodRings/CrystallineRings.lean\#L1007-L1045}{The integral cancellation and first approximation}. \src{Lib/PadicHodge/PeriodRings/CrystallineRings.lean\#L1046-L1086}{Gain of one power and convergence}. \src{Lib/PadicHodge/Kisin/PeriodMatrixBasis.lean\#L439-L469}{Application to the whole family of coordinates}.}
\noteline{14.75}{75}{\src{Lib/PadicHodge/Kisin/PeriodMatrixBasis.lean\#L555-L597}{Approximation by solutions over \(S_B\)}. \src{Lib/PadicHodge/Kisin/PeriodMatrixBasis.lean\#L620-L663}{Cancellation and full rank}.}
\noteline{14.76}{76}{\src{Lib/PadicHodge/Kisin/PeriodMatrixFrame.lean\#L301-L347}{Frobenius on the chart basis}. \src{Lib/PadicHodge/Kisin/PeriodMatrixFrame.lean\#L370-L395}{Saturation on chart coefficients}.}
\noteline{14.77}{77}{\src{Lib/PadicHodge/Kisin/PeriodMatrixFrame.lean\#L47-L69}{Selecting a determinant by row replacement}. \src{Lib/PadicHodge/Kisin/PeriodMatrixFrame.lean\#L102-L141}{Separation by the chart-coordinate rows}.}
\noteline{14.78}{78}{\src{Lib/PadicHodge/Kisin/PeriodMatrixFrame.lean\#L71-L100}{The normalized minimal-support vector}. \src{Lib/PadicHodge/Kisin/PeriodMatrixFrame.lean\#L397-L421}{The fixer of \(S\) preserves its kernel relation}.}
\noteline{14.79}{79}{\src{Lib/PadicHodge/Kisin/Charts.lean\#L755-L804}{Approximation of a fixed tilt component}.}
\noteline{14.80}{80}{\src{Lib/PadicHodge/Kisin/PeriodMatrixFrame.lean\#L171-L220}{The Frobenius product and its determinant}. \src{Lib/PadicHodge/Kisin/PeriodMatrixFrame.lean\#L221-L245}{The component calculation through the coordinate maps}.}
\noteline{14.81}{81}{\src{Lib/PadicHodge/Kisin/PeriodMatrixFrame.lean\#L267-L299}{Integral approximation of \(U\)}. \src{Lib/PadicHodge/Kisin/PeriodMatrixFrame.lean\#L423-L460}{The adjugate and both partial inverses}. \src{Lib/PadicHodge/Kisin/PeriodMatrixFrame.lean\#L473-L506}{Assembly of the frame and its Galois transform}.}
\noteline{14.82}{82}{\src{Lib/PadicHodge/Kisin/Modules.lean\#L216-L266}{Definitions of the lattices and forms}. \src{Lib/PadicHodge/Kisin/FrameTheorem.lean\#L213-L234}{Construction in the ambient Witt module}. \src{Lib/PadicHodge/Kisin/FrameTheorem.lean\#L265-L299}{Fixed rows and the evaluation formula}. \src{Lib/PadicHodge/Kisin/FrameTheorem.lean\#L321-L365}{Hodge orders and value-field points of the frame}.}
\noteline{14.83}{83}{\src{Lib/PadicHodge/Kisin/FrameTheorem.lean\#L367-L405}{Crystalline forms reading the frame coordinates}.}
\noteline{14.84}{84}{\src{Lib/PadicHodge/Kisin/FrameTheorem.lean\#L458-L514}{Removing the powers of \(\tau\)}. \src{Lib/PadicHodge/PeriodRings/Ainf.lean\#L889-L895}{The positive-translate divisibility test}. \src{Lib/PadicHodge/Kisin/FrameTheorem.lean\#L651-L678}{The two inclusions and normalized Frobenius}.}
\noteline{14.85}{85}{\src{Lib/PadicHodge/Kisin/KummerChartDeRham.lean\#L247-L256}{The field at the prescribed level}. \src{Lib/PadicHodge/Kisin/KummerChartDeRham.lean\#L370-L435}{Clearing denominators and conjugating the coordinate}.}
\noteline{14.86}{86}{\src{Lib/PadicHodge/Kisin/FrameTheorem.lean\#L418-L442}{From scalar relations to vector relations}. \src{Lib/PadicHodge/Kisin/FrameTheorem.lean\#L517-L555}{Spanning after descent and linear independence}.}
\noteline{14.87}{87}{\src{Lib/PadicHodge/Kisin/Modules.lean\#L268-L279}{Saturation by products of the chart primes}. \src{New/Congruences/Lifting/ExponentLattices.lean\#L754-L786}{Locating the form values in \(\mathcal K_m\)}. \src{New/Congruences/Lifting/ExponentLattices.lean\#L788-L852}{Factoring the retraction through the tensor product}. \src{New/Congruences/Lifting/HodgeExpansion.lean\#L188-L228}{Preservation of the target congruence and the letter}.}
\noteline{14.88}{88}{\src{Lib/PadicHodge/Crystalline/HodgeTateJumps.lean\#L208-L240}{Existence of the de Rham frame}. \src{Lib/PadicHodge/TateSen.lean\#L490-L523}{Passage from invariance modulo a filtration step to a constant}.}
\noteline{14.89}{89}{\src{Lib/PadicHodge/PeriodRings/AxTheorem.lean\#L376-L527}{Ax's approximation in the original formalization}. \src{Lib/PadicHodge/TateSen.lean\#L105-L274}{Its normalized trace}. \comsrc{Lib/PadicHodge/PeriodRings/AxTheorem.lean\#L335-L359}{Fixed fields in the community formalization}. \comsrc{Lib/PadicHodge/TateSen.lean\#L48-L74}{Twisted vanishing for open subgroups}.}
\noteline{14.90}{90}{\src{Lib/PadicHodge/TateSen.lean\#L416-L480}{Residues and the filtration}. \src{Lib/PadicHodge/TateSen.lean\#L487-L550}{Lifting a fixed class by subtracting its residue}.}
\noteline{14.91}{91}{\src{Lib/PadicHodge/TateSen.lean\#L291-L402}{Support elimination and algebraic coordinates}.}
\noteline{14.92}{92}{\src{Lib/PadicHodge/Kisin/Modules.lean\#L127-L183}{Coefficient roots, letters and the crystallinity hypothesis}. \src{Lib/PadicHodge/Kisin/Modules.lean\#L291-L317}{The projector, type of a functional and jump multiplicity}. \src{Lib/PadicHodge/Kisin/Modules.lean\#L458-L483}{Image and fixed vectors of the projector}. \src{Lib/PadicHodge/Kisin/KummerChartDeRham.lean\#L150-L155}{Openness of the chart subgroup's image}.}
\noteline{14.93}{93}{\src{Lib/PadicHodge/Kisin/Modules.lean\#L599-L614}{Equivariance after applying the projector}. \src{Lib/PadicHodge/Kisin/Modules.lean\#L648-L688}{The symbol and its two transformation rules}.}
\noteline{14.94}{94}{\src{Lib/PadicHodge/Crystalline/Basic.lean\#L193-L214}{The coefficient field, weight convention and definition of \(j_\gamma\)}.}
\noteline{14.95}{95}{\src{Lib/PadicHodge/Crystalline/HodgeTateJumps.lean\#L355-L394}{The spanning invariant columns after applying \(\gamma\)}.}
\noteline{14.96}{96}{\src{Lib/PadicHodge/Crystalline/HodgeTateJumps.lean\#L308-L350}{Weights of the inverse rows}. \src{Lib/PadicHodge/Crystalline/HodgeTateJumps.lean\#L812-L854}{Combining the bases for the coefficient embeddings}.}
\noteline{14.97}{97}{\src{Lib/PadicHodge/Crystalline/HodgeTateForms.lean\#L102-L148}{The coordinate comparison for two weights}. \src{Lib/PadicHodge/Crystalline/HodgeTateJumps.lean\#L853-L878}{The second restriction from the coefficient root}.}
\noteline{14.98}{98}{\src{Lib/PadicHodge/Kisin/Modules.lean\#L689-L733}{Descent of the coordinates and the zero-multiplicity case}.}
\noteline{14.99}{99}{\src{New/Congruences/Lifting/CupProductBound.lean\#L606-L665}{The lifting ideal and the change at \(E\)}. \src{New/Congruences/Lifting/CupProductBound.lean\#L668-L731}{Exact lifting and transport of the ideal}.}
\noteline{14.100}{100}{\src{Lib/PadicHodge/Kisin/FrameTheorem.lean\#L651-L666}{The two inclusions after multiplication by \(\tau^H\)}. \src{New/Congruences/Lifting/CupProductBound.lean\#L1064-L1095}{A uniform starting exponent}.}
\noteline{14.101}{101}{\src{New/Congruences/Lifting/CupProductBound.lean\#L738-L782}{Integral lifts using chart coordinates}. \src{New/Congruences/Lifting/CupProductBound.lean\#L798-L893}{Dividing the difference at levels zero and one}.}
\noteline{14.102}{102}{\src{Lib/PadicHodge/Kisin/PeriodMatrixBasis.lean\#L152-L193}{The cyclotomic factorization}. \src{New/Congruences/Lifting/CupProductBound.lean\#L895-L932}{Rounding a power down before Frobenius descent}. \src{New/Congruences/Lifting/CupProductBound.lean\#L991-L1037}{The extra factor \(p\) at level zero}.}
\noteline{14.103}{103}{\src{New/Congruences/Lifting/CupProductBound.lean\#L1091-L1124}{Minimality and the uniform height inequality}.}
\noteline{14.104}{104}{\src{Lib/PadicHodge/Kisin/KummerChartDeRham.lean\#L82-L149}{The different from Weierstrass division}. \src{New/Congruences/Lifting/ExponentLattices.lean\#L83-L134}{The Taylor identity for the moved uniformizer}.}
\noteline{14.105}{105}{\src{New/Congruences/Lifting/ExponentLattices.lean\#L135-L245}{The three contributions \(v_p(|n|)\), \(d+a\), and \(2a\)}.}
\noteline{14.106}{106}{\src{New/Congruences/Lifting/ExponentLattices.lean\#L1040-L1138}{The first-order action and its cancellation}. \src{New/Congruences/Lifting/SidePoints.lean\#L596-L679}{Passage back to the integral chart lattice}.}
\noteline{14.107}{107}{\src{New/Congruences/Lifting/ExponentLattices.lean\#L247-L312}{The residue sum on the powers of \(z\)}. \src{New/Congruences/Lifting/ExponentLattices.lean\#L1140-L1181}{The zero step on the value field}. \src{New/Congruences/Lifting/SidePoints.lean\#L681-L728}{Composition for a finite set of pairs}.}
\noteline{14.108}{108}{\src{New/Congruences/Lifting/CupProductBound.lean\#L933-L989}{The one-factor trade}. \src{New/Congruences/Lifting/CupProductBound.lean\#L1038-L1062}{The weighted contraction at one letter}.}
\noteline{14.109}{109}{\src{Lib/PadicHodge/Kisin/PeriodMatrixFrame.lean\#L445-L470}{The integral change-of-frame matrices}. \src{Lib/PadicHodge/RadiusLemma.lean\#L402-L426}{Evaluation and normalization at \(z=0\)}.}
\noteline{14.110}{110}{\src{Lib/PadicHodge/RadiusLemma.lean\#L58-L101}{Cancellation of the vertical and coordinate divisors}. \src{Lib/PadicHodge/RadiusLemma.lean\#L285-L340}{A divided difference for a row remainder}.}
\noteline{14.111}{111}{\src{Lib/PadicHodge/RadiusLemma.lean\#L342-L401}{Successive row approximation and its factorial exponents}. \src{Lib/PadicHodge/RadiusLemma.lean\#L471-L540}{Factorization of the Frobenius translate}.}
\noteline{14.112}{112}{\src{Lib/PadicHodge/Kisin/Bicommutant.lean\#L467-L489}{Projection of the \(A_{\max}\) remainder}. \src{Lib/PadicHodge/RadiusLemma.lean\#L177-L215}{Truncation at the chosen evaluation point}.}
\noteline{14.113}{113}{\src{Lib/PadicHodge/RadiusLemma.lean\#L217-L283}{Two approximations give an integral spanning identity}.}
\noteline{14.114}{114}{\src{Lib/PadicHodge/Kisin/FrameTheorem.lean\#L555-L600}{Application to all chart-equivariant forms}.}
\noteline{14.115}{115}{\src{New/Congruences/Lifting/LongGap.lean\#L275-L305}{The monic polynomial and its degree}.}
\noteline{14.116}{116}{\src{New/Congruences/Lifting/LongGap.lean\#L177-L270}{Projecting the frame span and expanding its coefficients}.}
\noteline{14.117}{117}{\src{New/Congruences/Lifting/LongGap.lean\#L118-L175}{Peeling the coefficients and improving a positive order}.}
\noteline{14.118}{118}{\src{New/Congruences/Lifting/LongGap.lean\#L306-L364}{Monicity and transport of the quotient lift}. \src{New/Congruences/Lifting/LongGap.lean\#L365-L473}{The coefficient and the exact corrected point}.}
\noteline{14.119}{119}{\src{Lib/PadicHodge/Kisin/Modules.lean\#L649-L712}{Forms of one type modulo forms of lower type}.}
\noteline{14.120}{120}{\src{New/Congruences/Lifting/HodgeExpansion.lean\#L597-L669}{Integral row reduction on the leading residues}.}
\noteline{14.121}{121}{\src{New/Congruences/Lifting/HodgeExpansion.lean\#L671-L765}{The primitive relation and its two terms}.}
\noteline{14.122}{122}{\src{New/Congruences/Lifting/FewLowJumps.lean\#L109-L208}{A zero-cost improvement and division by \(p\)}. \src{New/Congruences/Lifting/FewLowJumps.lean\#L220-L264}{Applying the relation to the kernel points}.}
\noteline{14.123}{123}{\src{New/Congruences/Lifting/FewLowJumps.lean\#L265-L310}{The colength of the liftable constants}.}
\noteline{14.124}{124}{\src{Lib/PadicHodge/Kisin/PeriodMatrixBasis.lean\#L950-L977}{Crystallinity gives the period matrix}. \src{Lib/PadicHodge/Kisin/FrameTheorem.lean\#L176-L242}{The frame data for a bounded crystalline lattice}. \src{Lib/PadicHodge/Kisin/KummerChartDeRham.lean\#L434-L471}{Residues in the prescribed value field}. \src{Lib/PadicHodge/TateSen.lean\#L490-L523}{Fixed values modulo the Hodge filtration}.}
\section*{Chapter 15. Galois attachment and change of rank}
\noteline{15.1}{1}{\src{Lib/GaloisRep/Main/Defs.lean}{Frobenius data and density}. \src{Lib/GaloisRep/Main/AttachedRepresentations.lean}{Attachment and its local clauses}.}
\noteline{15.2}{2}{\src{Lib/GaloisRep/Main/AllRanks.lean}{Construction of the two twists, cancellation and local compatibility}.}
\noteline{15.3}{3}{\src{Lib/GaloisRep/KillingPolynomials.lean}{Power minors, rank and characteristic polynomials}.}
\noteline{15.4}{4}{\src{Lib/GaloisRep/Main/SeedRepresentation.lean}{Odd-rank seed}. \src{Lib/GaloisRep/Main/SplitCase.lean}{Removing the weight and type restrictions}. \src{Lib/GaloisRep/Main/OddRank.lean}{Descent over CM fields}. \src{Lib/GaloisRep/Main/ThetaLift.lean}{Odd-rank theta lifts}.}
\noteline{15.5}{5}{\src{Lib/GaloisRep/Main/SeedRepresentation.lean}{The seed and its place hypotheses}.}
\noteline{15.6}{6}{\src{Lib/Shimura/IntegralModel/PELDataExistence.lean}{Choice of PEL data}. \src{Lib/GaloisRep/Main/SeedRepresentation.lean}{Auxiliary place and level choices}.}
\noteline{15.7}{7}{\src{Lib/GaloisRep/Main/SeedRepresentation.lean}{Frame extension, central eigenform and central twist}.}
\noteline{15.8}{8}{\src{Lib/Shimura/Transfer/WeightOccurs.lean}{Alternant calculation, coefficient representation and weight shift}.}
\noteline{15.9}{9}{\src{Lib/Shimura/TraceFormula/TestFunction.lean}{Test function and its powers}. \src{Lib/Shimura/TraceFormula/MarkedDefiniteSums.lean}{Comparison of the two counts}. \src{Lib/Shimura/Transfer/WeightOccurs.lean}{Passage from traces to occurrence}.}
\noteline{15.10}{10}{\src{Lib/Shimura/ThetaModuli/UniversalFamily.lean}{Abelian family over the theta closure}. \src{Lib/Shimura/IntegralModel/TowerExistence.lean}{The integral tower and its Hecke legs}.}
\noteline{15.11}{11}{\src{Lib/Shimura/CongruenceRelation/KernelReduction.lean}{Section multiplicities and reduction}. \src{Lib/Shimura/CongruenceRelation/HeckeNeighbours.lean}{From torsion counts to equal reduced points}.}
\noteline{15.12}{12}{\src{Lib/Shimura/IntegralModel/TowerExistence.lean}{Cancellation on the fibres over subspace sums}.}
\noteline{15.13}{13}{\src{Lib/Cohomology/Etale/Basic.lean\#L356-L374}{The category of test rings}. \src{Lib/Cohomology/Etale/EllAdicCohomology.lean\#L44-L92}{The coefficient functor, geometric cohomology and Galois action}. \src{Lib/PadicComparison/CrystallineFrobenius.lean\#L20-L50}{Scope of the comparisons}.}
\noteline{15.14}{14}{\src{Lib/Shimura/CongruenceRelation/EichlerShimura.lean}{Congruence and geometric Frobenius away from \(\ell\)}. \src{Lib/PadicComparison/CrystallineFrobenius.lean\#L713-L752}{Congruence over the period ring}.}
\noteline{15.15}{15}{\src{Lib/PadicComparison/CrystallineFrobenius.lean\#L37-L47}{Precisely the Hodge--Tate realization needed here}. \src{Lib/Shimura/Eigenspace.lean\#L443-L472}{A common cover and the vanishing on Hodge classes}.}
\noteline{15.16}{16}{\src{Lib/Cohomology/WoodsHole.lean\#L205-L242}{From vanishing on Hodge classes to vanishing of local terms}. \src{Lib/Shimura/Transfer/Lefschetz.lean\#L460-L478}{Grouping coincidences by marked classes}.}
\noteline{15.17}{17}{\src{Lib/ClassField/Avatar/HeckeBaseChange.lean}{Twisting and base change}. \src{Lib/GaloisRep/Main/SeedRepresentation.lean}{Untwisting on the prescribed places}.}
\noteline{15.18}{18}{\src{Lib/GaloisRep/Main/CuspidalTypes.lean}{Combining the embeddings and removing multiplicities}.}
\noteline{15.19}{19}{\src{Lib/GaloisRep/Main/CuspidalTypes.lean}{The Hodge--Tate cocharacter}. \src{Lib/GaloisRep/Main/SeedRepresentation.lean}{Application to the inverse local frame}.}
\noteline{15.20}{20}{\src{Lib/Shimura/IntegralModel/TowerExistence.lean}{Geometric hypotheses and construction}. \src{Lib/Shimura/Eigenspace.lean\#L410-L418}{The implication used to detect a jump}.}
\noteline{15.21}{21}{\src{Lib/GaloisRep/BaseChange/PrimeDegreeTransfer.lean\#L46-L76}{Definition of weak base change}. \src{Lib/GaloisRep/Main/OddRank.lean\#L267-L295}{The transfer used in the quadratic descent}.}
\noteline{15.22}{22}{\src{Lib/GaloisRep/BaseChange/CongruenceLevelMatching.lean}{Integral levels, cocycle splitting and matching of their indicators}. \src{Lib/GaloisRep/BaseChange/StepMatchingPairs.lean}{The local root identity, representatives and Hecke-stable pairs}.}
\noteline{15.23}{23}{\src{Lib/GaloisRep/BaseChange/TwistedTraceIdentity.lean\#L330-L395}{Cyclic factors of the coefficient representation}.}
\noteline{15.24}{24}{\src{Lib/GaloisRep/BaseChange/TwistedTraceIdentity.lean\#L171-L198}{Multiplicity of the rational moves}. \src{Lib/GaloisRep/BaseChange/TwistedTraceIdentity.lean\#L215-L260}{Double-coset expansion}.}
\noteline{15.25}{25}{\src{Lib/GaloisRep/BaseChange/GlobalRoot.lean}{Norm roots and their rational adjustment}. \src{Lib/GaloisRep/BaseChange/TwistedTraceIdentity.lean\#L850-L927}{Bijection on moves}.}
\noteline{15.26}{26}{\src{Lib/GaloisRep/BaseChange/TwistedTraceIdentity.lean\#L930-L1020}{The coefficient and twisted trace identity}.}
\noteline{15.27}{27}{\src{Lib/GaloisRep/BaseChange/PrimeDegreeTransfer.lean\#L693-L740}{Choice of the commuting spherical families}. \src{Lib/GaloisRep/BaseChange/PrimeDegreeTransfer.lean\#L802-L854}{Nonzero block, twist and eigenvalues}.}
\noteline{15.28}{28}{\src{Lib/GaloisRep/Main/OddRank.lean}{Using weak base change on the kernels of the quadratic characters}. \src{Lib/GaloisRep/BaseChange/ShintaniLift.lean}{Transfer when a finite residue-field type is prescribed}.}
\noteline{15.29}{29}{\src{Lib/GaloisRep/Hilbert90/TameActions.lean\#L48-L72}{Tameness at a level}. \src{Lib/GaloisRep/BaseChange/CongruenceLevelMatching.lean\#L622-L665}{Integral closure and the element separating the two conjugations}.}
\noteline{15.30}{30}{\src{Lib/GaloisRep/Hilbert90/TameActions.lean\#L301-L346}{Finite quotients, inherited tameness and residue descent}. \src{Lib/GaloisRep/Hilbert90/TameActions.lean\#L463-L493}{The congruent fixed matrix}. \src{Lib/GaloisRep/Hilbert90/TameActions.lean\#L502-L584}{Making it invertible and passing to the profinite ring}.}
\noteline{15.31}{31}{\src{Lib/GaloisRep/Hilbert90/UnitaryGroups.lean\#L346-L402}{Fixed orthonormalization and the unitary congruence}.}
\noteline{15.32}{32}{\src{Lib/GaloisRep/Hilbert90/TameActions.lean\#L637-L712}{Counting the image of the norm and lifting through the quotients}.}
\noteline{15.33}{33}{\src{Lib/GaloisRep/Hilbert90/UnitaryGroups.lean\#L89-L169}{The characteristic algebra and its commuting actions}.}
\noteline{15.34}{34}{\src{Lib/GaloisRep/Hilbert90/UnitaryGroups.lean\#L171-L248}{A fixed unit in the characteristic algebra}. \src{Lib/GaloisRep/Hilbert90/UnitaryGroups.lean\#L249-L342}{The twisted polynomial root}.}
\noteline{15.35}{35}{\src{Lib/GaloisRep/BaseChange/Defs.lean\#L145-L161}{The three conditions on a polynomial root}. \src{Lib/GaloisRep/BaseChange/CongruenceLevelMatching.lean\#L52-L111}{The affine substitution and determinants of integral elements}. \src{Lib/GaloisRep/BaseChange/CongruenceLevelMatching.lean\#L181-L260}{The determinant condition for arbitrary matrix size}.}
\noteline{15.36}{36}{\src{Lib/GaloisRep/BaseChange/Defs.lean\#L163-L184}{The matching quantifiers and the cocycle condition}. \src{Lib/GaloisRep/BaseChange/CongruenceLevelMatching.lean\#L113-L125}{The indicator and the second support condition}. \src{Lib/GaloisRep/BaseChange/CongruenceLevelMatching.lean\#L269-L283}{Integrality of the quotient on the support}. \src{Lib/GaloisRep/BaseChange/CongruenceLevelMatching.lean\#L347-L383}{Cocycle splitting and choice of a representative in the level}.}
\noteline{15.37}{37}{\src{Lib/GaloisRep/BaseChange/CongruenceLevelMatching.lean\#L442-L456}{The exact maximality and pullback conditions}. \src{Lib/GaloisRep/BaseChange/CongruenceLevelMatching.lean\#L784-L833}{Inertia and tameness for the choice of level}. \src{Lib/GaloisRep/BaseChange/CongruenceLevelMatching.lean\#L840-L896}{Choosing the level below the given open subgroup}.}
\noteline{15.38}{38}{\src{Lib/GaloisRep/BaseChange/Defs.lean\#L146-L164}{Determinants for every matrix size and the norm-root identity}.}
\noteline{15.39}{39}{\src{Lib/GaloisRep/BaseChange/GlobalRoot.lean\#L920-L1001}{The CM fields, their involutions, and the finite adelic presentations}.}
\noteline{15.40}{40}{\src{Lib/GaloisRep/BaseChange/GlobalRoot.lean\#L578-L595}{The characteristic-polynomial divisibility}. \src{Lib/GaloisRep/BaseChange/GlobalRoot.lean\#L728-L753}{Passage to the reduced algebra}.}
\noteline{15.41}{41}{\src{Lib/GaloisRep/BaseChange/GlobalRoot.lean\#L50-L79}{Making an idele fixed under CM conjugation}. \src{Lib/ClassField/Global/CMSplittingCharacters.lean\#L259-L285}{The norm criterion with signs at infinity}.}
\noteline{15.42}{42}{\src{Lib/GaloisRep/BaseChange/GlobalRoot.lean\#L92-L114}{Hasse's norm theorem followed by Hilbert 90}. \src{Lib/GaloisRep/BaseChange/GlobalRoot.lean\#L129-L230}{The norm character, the Artin restriction and the scalar correction}.}
\noteline{15.43}{43}{\src{Lib/GaloisRep/BaseChange/GlobalRoot.lean\#L516-L557}{Passing a unit around each orbit}. \src{Lib/GaloisRep/BaseChange/GlobalRoot.lean\#L597-L629}{Recovering the matrix quotient from self-adjointness}.}
\noteline{15.44}{44}{\src{Lib/GaloisRep/BaseChange/TwistedTraceIdentity.lean\#L508-L579}{The rational and adelic assignments in the trace calculation}.}
\noteline{15.45}{45}{\src{Lib/LocalGL/Shintani/LiftingGreenCharacters.lean\#L649-L725}{The finite-field theorem for these Green characters}. \src{Lib/GaloisRep/LocalGlobal/LabelCompatibility.lean\#L75-L107}{Reduction of the inertia quotient defining the tame exponent}.}
\noteline{15.46}{46}{\src{Lib/GaloisRep/BaseChange/ShintaniLift.lean\#L244-L338}{Matching the inflated functions}. \src{Lib/GaloisRep/BaseChange/ShintaniLift.lean\#L342-L400}{Trace downstairs and detection of the upstairs type}.}
\noteline{15.47}{47}{\src{Lib/ClassField/Grunwald.lean\#L85-L161}{The Kummer detection and extension steps}. \src{Lib/ClassField/Grunwald.lean\#L171-L271}{Prime-exponent Grunwald with prime-norm witnesses}. \src{Lib/ClassField/Grunwald.lean\#L769-L847}{Prescribing the character and its conjugate}.}
\noteline{15.48}{48}{\src{Lib/ClassField/Grunwald.lean\#L318-L356}{Conjugation at every embedding of the symmetrized extension}.}
\noteline{15.49}{49}{\src{Lib/ClassField/Grunwald.lean\#L398-L494}{Prime degree, norm criterion and integral trace one}. \src{Lib/ClassField/Grunwald.lean\#L652-L719}{Lifting the Weil data}. \src{Lib/ClassField/Grunwald.lean\#L850-L926}{The resulting prescribed completion}.}
\noteline{15.50}{50}{\src{Lib/GaloisRep/LocalGlobal/LabelCompatibility.lean\#L351-L386}{The inert global step}. \src{Lib/GaloisRep/LocalGlobal/LabelCompatibility.lean\#L410-L476}{Removing a prime from the tame exponent level}.}
\noteline{15.51}{51}{\src{Lib/LocalGL/Iwahori/SortedExponents.lean\#L316-L447}{The adjacent exchange and its two obstructions}. \src{Lib/LocalGL/Iwahori/SortedExponents.lean\#L514-L626}{The eigenfunctional and sorted eigenvector}. \src{Lib/LocalGL/Iwahori/SortedExponents.lean\#L841-L879}{Application to the prescribed determinant blocks}. \src{Lib/LocalGL/Iwahori/SphericalExponents.lean\#L518-L634}{Descent from the maximal compact to the Iwahori}. \src{Lib/GaloisRep/LocalGlobal/IwahoriCompatibility.lean\#L145-L190}{The spherical Hecke polynomial}.}
\noteline{15.52}{52}{\src{Lib/GaloisRep/LocalGlobal/IwahoriCompatibility.lean\#L354-L444}{The finite index and its decrease under a prescribed step}.}
\noteline{15.53}{53}{\src{Lib/LocalGL/Iwahori/ParahoricStructure.lean\#L698-L811}{Twisted diagonalization by interval induction}. \src{Lib/GaloisRep/BaseChange/IwahoriEigenforms.lean\#L343-L411}{The interpolated root and coset identities}. \src{Lib/GaloisRep/BaseChange/IwahoriEigenforms.lean\#L755-L834}{Transfer of the dominant eigenvalues}.}
\noteline{15.54}{54}{\src{Lib/GaloisRep/LocalGlobal/IwahoriCompatibility.lean\#L105-L141}{Passage to inertia on both spaces}. \src{Lib/GaloisRep/LocalGlobal/LabelCompatibility.lean\#L289-L349}{Artin values and the tame labels}.}
\noteline{15.55}{55}{\src{Lib/Weil/DualPair/ThetaLiftConvergence.lean\#L287-L296}{The integral and the quotient measure}.}
\noteline{15.56}{56}{\src{Lib/GaloisRep/Main/ThetaLift.lean\#L295-L322}{The weight rule}.}
\noteline{15.57}{57}{\src{Lib/Weil/Lift.lean\#L594-L624}{The normalized coefficient pairing}. \src{Lib/Weil/Lift.lean\#L727-L741}{The joint harmonic in the oscillator representation}.}
\noteline{15.58}{58}{\src{Lib/SiegelWeil/RallisInner/Doubling.lean\#L630-L652}{The finite isotropic frame}. \src{Lib/GaloisRep/Main/ThetaLift.lean\#L80-L277}{Integral support, stabilizer and sharp datum}.}
\noteline{15.59}{59}{\comsrc{Lib/SiegelWeil/RallisInner/VacuumSummability.lean\#L55-L130}{Absolute summability from doubling}. \comsrc{Lib/SiegelWeil/RallisInner/SharpData.lean\#L436-L496}{The sharp norm identity and the real factor}.}
\noteline{15.60}{60}{\comsrc{Lib/SiegelWeil/RallisInner/DualDatum.lean\#L204-L265}{The shear, the Heisenberg translation and the coefficient at integral points}. \comsrc{Lib/SiegelWeil/RallisInner/DualDatum.lean\#L266-L358}{Integrality from the additive character}. \comsrc{Lib/SiegelWeil/RallisInner/SharpData.lean\#L294-L407}{Exhaustion by finite sets of places and the sharp mass}.}
\noteline{15.61}{61}{\comsrc{Lib/GaloisRep/Main/ThetaLift.lean\#L436-L454}{Sphericity and the eigenpolynomial hypothesis}. \src{Lib/Weil/Unramified/HeckeTransfer.lean\#L336-L451}{The two residue kernels and the product identity}. \src{Lib/Weil/Unramified/HeckeTransfer.lean\#L608-L653}{Transfer through the theta integral}.}
\noteline{15.62}{62}{\src{Lib/GaloisRep/Main/ThetaLift.lean\#L276-L298}{Conversion from reciprocal to monic polynomials}. \src{Lib/GaloisRep/Main/ThetaLift.lean\#L323-L350}{The two character normalizations}.}
\noteline{15.63}{63}{\src{Lib/GaloisRep/Main/AllRanks.lean\#L108-L161}{The two weights and the split-place contradiction}.}
\noteline{15.64}{64}{\src{Lib/ClassField/Avatar/HeckeBaseChange.lean\#L190-L215}{Infinity types and unramified character values}. \src{Lib/GaloisRep/Main/AllRanks.lean\#L384-L454}{The automorphic lifts, twists and jump calculation}. \src{Lib/GaloisRep/Main/AllRanks.lean\#L466-L539}{The two strings and local cancellation}.}
\noteline{15.65}{65}{\src{Lib/SiegelWeil/RallisInner/Doubling.lean\#L520-L553}{Doubling before imposing the local support condition}. \src{Lib/ClassField/HeckeCharacter/Basic.lean\#L590-L618}{Extension of a unitary character with specified weight}.}
\noteline{15.66}{66}{\src{Lib/GaloisRep/WeightCongruence.lean\#L125-L157}{A coefficient field for the models}.}
\noteline{15.67}{67}{\src{Lib/GaloisRep/WeightCongruence.lean\#L275-L333}{The allowed lacunary weights}.}
\noteline{15.68}{68}{\src{Lib/GaloisRep/WeightCongruence.lean\#L338-L406}{The separating multipliers and the translated nonvanishing section}.}
\noteline{15.69}{69}{\src{Lib/GaloisRep/WeightCongruence.lean\#L577-L596}{The decomposition into fixed and varying weight factors}. \src{Lib/GaloisRep/WeightCongruence.lean\#L625-L719}{Projection to the level character and recovery of the original model}.}
\noteline{15.70}{70}{\src{Lib/GaloisRep/WeightCongruence.lean}{Density of the weight and type models}. \src{Lib/GaloisRep/Main/SplitCase.lean}{Choice of the places and level and application of the free-weight construction}.}
\noteline{15.71}{71}{\src{Lib/GaloisRep/Pseudocharacter/Pseudorepresentations.lean}{Uniform approximation on the multiplicative monoid and its polynomial consequences}. \src{Lib/GaloisRep/Main/SplitCase.lean}{Normalized Hecke contractions and coefficient pairs}.}
\noteline{15.72}{72}{\src{Lib/GaloisRep/Pseudocharacter/Pseudorepresentations.lean}{The trace-identity realization and the approximation theorem in Lean}.}
\noteline{15.73}{73}{\src{Lib/GaloisRep/Pseudocharacter/Pseudorepresentations.lean}{Realization of trace identities in characteristic zero}. \src{Lib/GaloisRep/Main/SplitCase.lean}{Frobenius filters, coefficient field, continuity and attachment}.}
\noteline{15.74}{74}{\src{Lib/PadicHodge/Limit/TraceLimit.lean}{Passage of matrix relations to the limit over the integral group ring}.}
\noteline{15.75}{75}{\src{Lib/PadicHodge/Limit/CrystallineLimit.lean}{The period bound, integral kernel estimate, and adic passage of invariant spans}.}
\noteline{15.76}{76}{\src{Lib/PadicHodge/Limit/CrystallineLimit.lean\#L98-L139}{The uniform choice of denominator}. \src{Lib/PadicHodge/Limit/CrystallineLimit.lean\#L145-L195}{The two integral matrices and the fixed columns}.}
\noteline{15.77}{77}{\src{Lib/PadicHodge/PeriodRings/CrystallineRings.lean\#L353-L421}{The special fibre and monic polynomials}. \src{Lib/PadicHodge/PeriodRings/CrystallineRings.lean\#L639-L677}{The integral loss for every divisor of \(\mu\)}.}
\noteline{15.78}{78}{\src{Lib/PadicHodge/Limit/TraceLimit.lean\#L95-L145}{The uniform bound after integral division}. \src{Lib/PadicHodge/Limit/TraceLimit.lean\#L148-L221}{Compatibility with weight projections and their traces}. \src{Lib/PadicHodge/Limit/TraceLimit.lean\#L222-L333}{The adjugate and the finite-coefficient bound}.}
\noteline{15.79}{79}{\src{Lib/PadicHodge/Limit/TraceLimit.lean}{Uniform jump projectors, integral representatives for powers of Frobenius, and their coordinate tests}.}
\noteline{15.80}{80}{\src{Lib/GaloisRep/Main/QuadraticStep.lean\#L130-L181}{Construction of the norm and the uniformizer}.}
\noteline{15.81}{81}{\src{Lib/GaloisRep/Main/QuadraticStep.lean\#L212-L280}{The square-root test and the real Chinese-remainder factor}.}
\noteline{15.82}{82}{\src{Lib/GaloisRep/Main/QuadraticStep.lean}{Construction of the quadratic step, local conditions, and generating inertia}. \src{Lib/GaloisRep/BaseChange/PrimeDegreeTransfer.lean}{Transfer of the eigenform through the cyclic step}.}
\noteline{15.83}{83}{\src{Lib/GaloisRep/Main/OddRank.lean}{Weak base change on the character kernels, gluing, and descent of the local clauses}. \src{Lib/PadicHodge/Crystalline/Regularity.lean\#L709-L786}{Crystalline descent using the exponent-zero period operator}.}
\section*{Chapter 16. The marked unitary Shimura varieties}
\noteline{16.1}{1}{\src{Lib/Shimura/ThetaModuli/UniversalFamily.lean\#L126-L189}{Integral charts at the centre of an arc}. \src{Lib/Shimura/ThetaModuli/UniversalFamily.lean\#L278-L360}{Construction of the family over the theta closure}.}
\noteline{16.2}{2}{\src{Lib/Shimura/IntegralModel/BallImmersion.lean\#L100-L192}{Differentiation along torsion sections}. \src{Lib/Shimura/ThetaModuli/UniversalFamily.lean\#L460-L534}{Tangent bound and smoothness of the closure}.}
\noteline{16.3}{3}{\src{Lib/Shimura/IntegralModel/TowerExistence.lean\#L239-L276}{Legs determined by cosets}. \src{Lib/Shimura/IntegralModel/TowerExistence.lean\#L297-L349}{The finite tuple and permutations of its factors}.}
\noteline{16.4}{4}{\src{Lib/Shimura/IntegralModel/BallImmersion.lean\#L399-L452}{Recovering an actual tuple at a complex centre}. \src{Lib/Shimura/IntegralModel/TowerExistence.lean\#L350-L378}{Freeness from neatness and the normal core}.}
\noteline{16.5}{5}{\src{Lib/Shimura/IntegralModel/TowerExistence.lean\#L138-L217}{Tangent comparison through the two jets}.}
\noteline{16.6}{6}{\src{Lib/Shimura/CongruenceRelation/HeckeNeighbours.lean\#L57-L99}{The residue torsion and its natural frame}. \src{Lib/Shimura/CongruenceRelation/HeckeNeighbours.lean\#L388-L441}{Kernel labels at the split prime}.}
\noteline{16.7}{7}{\src{Lib/Shimura/CongruenceRelation/KernelReduction.lean\#L349-L419}{The tangent line from the signature condition}. \src{Lib/Shimura/CongruenceRelation/KernelReduction.lean\#L530-L601}{Divisibility in both directions}. \src{Lib/Shimura/CongruenceRelation/KernelReduction.lean\#L689-L763}{One local parameter for all label families}.}
\noteline{16.8}{8}{\src{Lib/Shimura/CongruenceRelation/KernelReduction.lean\#L803-L866}{Applying the local equation over the dual numbers}.}
\noteline{16.9}{9}{\src{Lib/Shimura/CongruenceRelation/HeckeNeighbours.lean\#L674-L745}{Factorization using the Frobenius-weighted local equation}. \src{Lib/Shimura/CongruenceRelation/HeckeNeighbours.lean\#L905-L941}{The annihilators and the two factorizations}.}
\noteline{16.10}{10}{\src{Lib/Shimura/CongruenceRelation/HeckeNeighbours.lean\#L447-L528}{Removing the conjugate scalar and identifying targets}. \src{Lib/Shimura/CongruenceRelation/HeckeNeighbours.lean\#L950-L1006}{Application to safe Hecke neighbours}.}
\noteline{16.11}{11}{\src{Lib/Shimura/ThetaModuli/FamilyFromPoints.lean}{Construction from fibre equations}. \src{Lib/Shimura/Deformation/TangentBound.lean}{Deformations at a good prime}. \src{Lib/Shimura/ThetaModuli/ThetaFamilyRigidity.lean}{Pairings, level data and the theta embedding}.}
\noteline{16.12}{12}{\src{Lib/Shimura/Uniformisation/AdelicFrames.lean\#L58-L87}{The place hypothesis and the integral unit in the frame}. \src{Lib/Shimura/Uniformisation/AdelicFrames.lean\#L118-L169}{The similitude group and the signs of the form}.}
\noteline{16.13}{13}{\src{Lib/Shimura/TraceFormula/DefiniteTrace.lean\#L194-L250}{Choosing the archimedean factors in a common compact Gram set}.}
\noteline{16.14}{14}{\src{Lib/Adelic/Reduction/FiniteVolume.lean\#L228-L248}{The isotropic and bounded parts of a rational basis}. \src{Lib/Adelic/Reduction/FiniteVolume.lean\#L271-L303}{Reduction for a varying rational form}. \src{Lib/Shimura/TraceFormula/DefiniteTrace.lean\#L254-L321}{Eliminating the isotropic block and counting occupied cosets}.}
\noteline{16.15}{15}{\src{Lib/Shimura/IntegralModel/BallImmersion.lean\#L194-L233}{Transporting the reflection and marks to a chosen ball}.}
\noteline{16.16}{16}{\src{Lib/Shimura/ThetaModuli/UniversalFamily.lean\#L411-L467}{The two-point comparison precedes the passage to a finite dense family}.}
\noteline{16.17}{17}{\src{Lib/Shimura/ThetaModuli/UniversalFamily.lean\#L126-L189}{Specialized chart ranks and nonvanishing}. \src{Lib/Shimura/ThetaModuli/UniversalFamily.lean\#L309-L349}{Application of the chart reconstruction}. \src{Lib/Shimura/ThetaModuli/UniversalFamily.lean\#L469-L530}{The further tangent and smoothness step}.}
\noteline{16.18}{18}{\src{Lib/Shimura/Uniformisation/BallModel.lean\#L312-L362}{The periods and their reflection}. \src{Lib/Shimura/Uniformisation/BallModel.lean\#L664-L701}{The connected positive domain}.}
\noteline{16.19}{19}{\src{Lib/Shimura/Uniformisation/BallModel.lean\#L392-L507}{The form on periods and its alternating part}. \src{Lib/Shimura/Uniformisation/BallModel.lean\#L572-L643}{The induced endomorphism action and polarized tori}.}
\noteline{16.20}{20}{\src{Lib/Shimura/Uniformisation/BallModel.lean\#L718-L802}{Level, marks, similitude and coverage}.}
\noteline{16.21}{21}{\src{Lib/Shimura/Uniformisation/BallModel.lean\#L803-L876}{The block determinant and its normalization}. \src{Lib/Shimura/Uniformisation/BallModel.lean\#L908-L990}{The isomorphism in a neighbourhood of the marked point}.}
\noteline{16.22}{22}{\src{Lib/Shimura/Uniformisation/BallModel.lean\#L533-L609}{The period coordinates, positivity domain, and derivative}.}
\noteline{16.23}{23}{\src{Lib/Shimura/Uniformisation/BallModel.lean\#L639-L661}{The positive-line test for the derivative}.}
\noteline{16.24}{24}{\src{Lib/Shimura/ThetaModuli/UniversalFamily.lean\#L369-L408}{Local holomorphic coordinates of the family}.}
\noteline{16.25}{25}{\src{Lib/Shimura/IntegralModel/BallImmersion.lean\#L107-L191}{Torsion sections, the second derivative and the lattice span}.}
\noteline{16.26}{26}{\src{Lib/Shimura/IntegralModel/BallImmersion.lean\#L238-L295}{Projections to the base, torsion and individual marks}.}
\noteline{16.27}{27}{\src{Lib/Shimura/ThetaModuli/PELThetaTori.lean\#L89-L138}{The level, Lagrangian and padding hypotheses}. \src{Lib/Shimura/ThetaModuli/FlatClosure.lean\#L73-L138}{The cutoff and the precise meaning of the rank equations}.}
\noteline{16.28}{28}{\src{Lib/Shimura/ThetaModuli/PELThetaTori.lean\#L172-L190}{Good primes}. \src{Lib/Shimura/GoodReduction/MultiplicativeNorm.lean\#L883-L998}{The second prime, expanding multiplications and their determinant}.}
\noteline{16.29}{29}{\src{Lib/Shimura/GoodReduction/MultiplicativeNorm.lean\#L376-L436}{The kernel average and its eighth-power normalization}. \src{Lib/Shimura/GoodReduction/MultiplicativeNorm.lean\#L493-L571}{The strict error estimate for a kernel shift}. \src{Lib/Shimura/GoodReduction/MultiplicativeNorm.lean\#L590-L816}{Descent of bounds and the expanding criterion}. \src{Lib/Shimura/GoodReduction/TorsionSizes.lean\#L345-L394}{Log-convexity on the finite shift orbit}.}
\noteline{16.30}{30}{\src{Lib/Shimura/ThetaNormality/NormalGeneration.lean\#L147-L237}{Integral spans of translated theta sections}. \src{Lib/Shimura/ThetaNormality/NormalGeneration.lean\#L442-L545}{The power step}. \src{Lib/Shimura/ThetaModuli/PELThetaTori.lean\#L1063-L1106}{Its application to the PEL coordinates}.}
\noteline{16.31}{31}{\src{Lib/Shimura/ThetaNormality/NormalGeneration.lean\#L558-L669}{Rank of ball reduction and the evaluation test}. \src{Lib/Shimura/ThetaNormality/RiemannRelation.lean\#L537-L710}{The affine ideal, its ranks and the finite level test}.}
\noteline{16.32}{32}{\src{Lib/Shimura/ThetaModuli/FlatClosure.lean\#L150-L205}{Division by relations with unit leading minors}. \src{Lib/Shimura/ThetaModuli/FlatClosure.lean\#L268-L485}{Cofactors, the cutoff and flatness}. \src{Lib/Shimura/ThetaModuli/FlatClosure.lean\#L784-L934}{Application over the local rings of the projective closure}.}
\noteline{16.33}{33}{\src{Lib/Shimura/ThetaNormality/RiemannRelation.lean\#L224-L329}{The translation representation, factors and scalar}.}
\noteline{16.34}{34}{\src{Lib/Shimura/ThetaNormality/RiemannRelation.lean\#L334-L497}{The maximal theta constants}. \src{Lib/Shimura/ThetaModuli/UniversalFamily.lean\#L76-L189}{Transfer of the unit charts along an arc}.}
\noteline{16.35}{35}{\src{Lib/Shimura/ThetaModuli/FamilyFromPoints.lean\#L53-L150}{Extension of the linear system from its two coordinate expressions}. \src{Lib/Shimura/ThetaModuli/FamilyFromPoints.lean\#L586-L704}{Regular coefficients and projection from the Segre embedding}.}
\noteline{16.36}{36}{\src{Lib/Shimura/ThetaModuli/FamilyFromPoints.lean\#L710-L774}{The two explicit padding tuples}. \src{Lib/Shimura/ThetaModuli/FamilyFromPoints.lean\#L248-L329}{Extension of the group and module laws}.}
\noteline{16.37}{37}{\src{Lib/Shimura/ThetaModuli/FamilyFromPoints.lean\#L185-L247}{Smoothness from specialized geometric group fibres}. \src{Lib/Shimura/ThetaModuli/FamilyFromPoints.lean\#L776-L907}{Polarization and integral translation laws}. \src{Lib/Shimura/ThetaModuli/FamilyFromPoints.lean\#L952-L1110}{Uniformisation, signature and the assembled family}.}
\noteline{16.38}{38}{\src{Lib/Shimura/ThetaModuli/UniversalFamily.lean\#L437-L530}{Dense parameters, tangent estimate and the precise existence theorem}.}
\noteline{16.39}{39}{\src{Lib/Shimura/CongruenceRelation/KernelReduction.lean\#L205-L239}{The good-prime identities}.}
\noteline{16.40}{40}{\src{Lib/General/Geometry/SmoothnessCriterion.lean\#L158-L190}{The universal first-order test}.}
\noteline{16.41}{41}{\src{Lib/Shimura/ThetaModuli/ThetaFamilyRigidity.lean\#L345-L397}{The rank-one logarithm in an étale chart}. \src{Lib/Shimura/Deformation/TangentBound.lean\#L119-L163}{The two annihilators of an infinitesimal point}.}
\noteline{16.42}{42}{\src{Lib/Shimura/Deformation/TangentBound.lean\#L462-L508}{Lifting points to the square-zero algebra}. \src{Lib/Shimura/Deformation/TangentBound.lean\#L800-L847}{Independence, linearity and vanishing on tangent vectors}.}
\noteline{16.43}{43}{\src{Lib/Shimura/Deformation/TangentBound.lean\#L165-L238}{Drinfeld's lifting argument}.}
\noteline{16.44}{44}{\src{Lib/Shimura/Deformation/TangentBound.lean\#L681-L762}{Killing the scheme-theoretic kernel and dividing by \(c\)}.}
\noteline{16.45}{45}{\src{Lib/Shimura/Deformation/TangentBound.lean\#L264-L328}{Linearity, Frobenius and faithfully flat comparison of the pairings}. \src{Lib/Shimura/AbelianScheme/Pairings.lean\#L317-L387}{Descent of the frame normalized at zero}.}
\noteline{16.46}{46}{\src{Lib/Shimura/Deformation/TangentBound.lean\#L335-L414}{The kernel and descent through a lifted isogeny}.}
\noteline{16.47}{47}{\src{Lib/Shimura/ThetaModuli/ThetaFamilyRigidity.lean\#L543-L582}{Detecting \(p\)-primary points by the pairing}. \src{Lib/Shimura/ThetaGroup.lean\#L499-L534}{The power of a translation lift that fixes the embedding}. \src{Lib/Shimura/Deformation/TangentBound.lean\#L527-L579}{The two isogeny divisions}.}
\noteline{16.48}{48}{\src{Lib/Shimura/ThetaGroup.lean\#L443-L497}{The translation and inversion identity}. \src{Lib/Shimura/ThetaModuli/ThetaFamilyRigidity.lean\#L584-L618}{Recovering the base map from the marked coordinates}. \src{Lib/Shimura/Deformation/TangentBound.lean\#L579-L603}{Keeping all level points}.}
\noteline{16.49}{49}{\src{Lib/Shimura/IntegralModel/UniversalFamily.lean\#L251-L331}{Integral descent of a periodic section}. \src{Lib/Shimura/IntegralModel/UniversalFamily.lean\#L334-L359}{Periodicity of the torsion trace}. \src{Lib/Shimura/IntegralModel/UniversalFamily.lean\#L411-L478}{Descent of spanning through multiplication}. \src{Lib/Shimura/IntegralModel/UniversalFamily.lean\#L482-L527}{The Serre bound on local pullbacks}.}
\noteline{16.50}{50}{\src{Lib/Shimura/Deformation/TangentCount.lean\#L59-L123}{The universal point of the local torsion scheme}. \src{Lib/Shimura/CongruenceRelation/KernelReduction.lean\#L803-L835}{The torsion parameter from labeled sections}. \src{Lib/Shimura/Deformation/TangentCount.lean\#L212-L274}{A single parameter and lifts of the tangent line}. \src{Lib/Shimura/Deformation/TangentCount.lean\#L438-L481}{The number of labels through zero}.}
\noteline{16.51}{51}{\src{Lib/Shimura/Deformation/TangentCount.lean\#L29-L57}{The finite-field monomial lemma}. \src{Lib/Shimura/Deformation/TangentCount.lean\#L374-L414}{The transformation of the lowest coefficient}.}
\noteline{16.52}{52}{\src{Lib/Shimura/Deformation/TangentCount.lean\#L475-L499}{The injective map and dimension bound}.}
\noteline{16.53}{53}{\src{Lib/General/Geometry/SmoothnessCriterion.lean\#L369-L433}{The relative cotangent and standard étale presentation}. \src{Lib/General/Geometry/SmoothnessCriterion.lean\#L440-L462}{Dominance along the analytic family and normality}. \src{Lib/General/Geometry/SmoothnessCriterion.lean\#L32-L62}{Removing the quotient kernel over the normal polynomial ring}. \src{Lib/Shimura/ThetaModuli/UniversalFamily.lean\#L435-L492}{Generic analytic parameters and the tangent estimate on an arc}. \src{Lib/Shimura/ThetaModuli/UniversalFamily.lean\#L496-L531}{The open cover of the good-prime base}.}
\noteline{16.54}{54}{\src{Lib/Shimura/ComplexTheta/DimensionFormula.lean\#L546-L609}{The dimension and scalar commutant for the analytic realization}.}
\noteline{16.55}{55}{\src{Lib/Shimura/ThetaNormality/NormalGeneration.lean\#L72-L137}{The trace calculation and the unit condition on the translation group}.}
\noteline{16.56}{56}{\src{Lib/Shimura/ThetaNormality/NormalGeneration.lean\#L142-L235}{Comparison of the two integral orbit spans}. \src{Lib/Shimura/ThetaNormality/NormalGeneration.lean\#L355-L432}{Exact scaling of the norm, also for a sublattice of dual translations}.}
\noteline{16.57}{57}{\src{Lib/Shimura/ThetaNormality/NormalGeneration.lean\#L442-L535}{Powering the norm ball using the larger lattice}. \src{Lib/Shimura/ThetaModuli/PELThetaTori.lean\#L1063-L1106}{Identification with the actual PEL coordinate polynomials}. \src{Lib/Shimura/ThetaNormality/NormalGeneration.lean\#L548-L635}{The quotient by the strict ball}.}
\noteline{16.58}{58}{\src{Lib/Shimura/ThetaNormality/NormalGeneration.lean\#L637-L705}{Comparison of evaluations using a unit coordinate}.}
\noteline{16.59}{59}{\src{Lib/Shimura/ThetaNormality/RiemannRelation.lean\#L537-L628}{The test on points and restriction to one chart}. \src{Lib/Shimura/ThetaNormality/RiemannRelation.lean\#L641-L710}{The ranks and the finite labelled evaluation criterion}.}
\noteline{16.60}{60}{\src{Lib/Shimura/ThetaModuli/FlatClosure.lean\#L70-L138}{Uniform finite cutoff and the conditions on a fibre ideal}. \src{Lib/Shimura/ThetaModuli/FlatClosure.lean\#L316-L374}{The evaluation determinants bound the initial monomial counts}. \src{Lib/Shimura/ThetaModuli/FlatClosure.lean\#L398-L485}{Flatness and equality of the specialized ideal}.}
\noteline{16.61}{61}{\src{Lib/Shimura/ThetaNormality/RiemannRelation.lean\#L334-L497}{The three simultaneous maximal coordinates}. \src{Lib/Shimura/ThetaModuli/UniversalFamily.lean\#L76-L189}{Passage from the primitive coordinates of the arc to the charts}.}
\noteline{16.62}{62}{\src{Lib/Shimura/ThetaModuli/FamilyFromPoints.lean\#L586-L644}{The two indices, the exact coefficient and its regularity}.}
\noteline{16.63}{63}{\src{Lib/Shimura/ThetaModuli/FamilyFromPoints.lean\#L53-L150}{Gluing the ratios and excluding common zeros by the reverse formula}. \src{Lib/Shimura/ThetaModuli/FamilyFromPoints.lean\#L645-L703}{Factoring the product map}.}
\noteline{16.64}{64}{\src{Lib/Shimura/GoodReduction/MultiplicativeNorm.lean\#L135-L247}{Covariance of the averaged pullback}.}
\noteline{16.65}{65}{\src{Lib/Shimura/ComplexTheta/DimensionFormula.lean\#L150-L198}{The diagonal mean from character orthogonality}. \src{Lib/Shimura/GoodReduction/MultiplicativeNorm.lean\#L249-L302}{Testing the average on the kernel}. \src{Lib/Shimura/GoodReduction/MultiplicativeNorm.lean\#L303-L375}{The off-diagonal character sums and conjugate kernel}.}
\noteline{16.66}{66}{\src{Lib/Shimura/GoodReduction/GaussSums.lean\#L345-L393}{Completing the exponent after reindexing the kernel}. \src{Lib/Shimura/GoodReduction/GaussSums.lean\#L395-L477}{The unfolded Gaussian and its determinant}.}
\noteline{16.67}{67}{\src{Lib/Shimura/GoodReduction/GaussSums.lean\#L264-L304}{Four squares and the phase of the finite sum}. \src{Lib/Shimura/GoodReduction/MultiplicativeNorm.lean\#L378-L435}{The eighth power of the determinant ratio}.}
\noteline{16.68}{68}{\src{Lib/General/Geometry/GaloisCovers.lean\#L481-L512}{Transport of the pointwise conditions}. \src{Lib/General/Geometry/GaloisCovers.lean\#L779-L824}{Smoothness after base change to the algebraic closure}.}
\noteline{16.69}{69}{\src{Lib/General/Geometry/PointsOfProperSchemes.lean\#L348-L447}{The finite atlas for the functor of points}.}
\noteline{16.70}{70}{\src{Lib/General/Geometry/EtaleCharts.lean\#L924-L1008}{The pointwise union, change of coordinates and charting subfunctor}. \src{Lib/Shimura/Eigenspace.lean\#L569-L583}{The base and the generic covers to which primitive comparison is applied}.}
\noteline{16.71}{71}{\src{Lib/Shimura/IntegralModel/PELDataExistence.lean\#L168-L189}{Choosing the CM type}.}
\noteline{16.72}{72}{\src{Lib/Shimura/IntegralModel/PELDataExistence.lean\#L54-L86}{The padding lists}.}
\noteline{16.73}{73}{\src{Lib/Shimura/IntegralModel/PELDataExistence.lean\#L88-L164}{The signed generator of the codifferent at the primes of \(s\)}.}
\noteline{16.74}{74}{\src{Lib/Shimura/IntegralModel/PELDataExistence.lean\#L238-L319}{The Lagrangian and its annihilator bound}. \src{Lib/Shimura/IntegralModel/PELDataExistence.lean\#L365-L451}{The integer avoided by the prescribed places}.}
\noteline{16.75}{75}{\src{Lib/GaloisRep/Main/QuadraticStep.lean\#L70-L115}{Counting the places above \(p\)}. \src{Lib/Shimura/IntegralModel/PELDataExistence.lean\#L197-L236}{Why each safe prime is good}.}
\noteline{16.76}{76}{\src{Lib/Shimura/HeckeTower/HeckeData.lean\#L855-L926}{The trace, norm bound and neatness}. \src{Lib/GaloisRep/Main/SeedRepresentation.lean\#L609-L634}{The supported principal level}. \src{Lib/GaloisRep/Main/SeedRepresentation.lean\#L713-L745}{Choosing the auxiliary place and retaining the safe places}.}
\noteline{16.77}{77}{\src{Lib/Shimura/TraceFormula/DefiniteTrace.lean\#L268-L378}{Finiteness of the signed-frame double cosets}. \src{Lib/Shimura/IntegralModel/BallImmersion.lean\#L194-L232}{Reduction to finitely many signed frames}. \src{Lib/General/Geometry/GaloisCovers.lean\#L294-L314}{A generic member of each ball}.}
\noteline{16.78}{78}{\src{Lib/Shimura/ThetaModuli/UniversalFamily.lean\#L404-L458}{The two-parameter closure and density for the full base}.}
\noteline{16.79}{79}{\src{Lib/Shimura/ThetaModuli/UniversalFamily.lean\#L273-L301}{The two specializations and contraction to a good prime}. \src{Lib/General/Geometry/DenseComplexPoints.lean\#L291-L336}{Making the centre principal in a noetherian subring}.}
\noteline{16.80}{80}{\src{Lib/General/Geometry/DenseComplexPoints.lean\#L205-L239}{Height one after adjoining indeterminates}. \src{Lib/General/Geometry/DenseComplexPoints.lean\#L241-L289}{An archimedean value group from the prime criterion}. \src{Lib/General/Geometry/DenseComplexPoints.lean\#L338-L429}{The rank-one localization and the arc}.}
\noteline{16.81}{81}{\src{Lib/Shimura/ThetaModuli/UniversalFamily.lean\#L126-L188}{Passing from a base arc to the reduced theta vectors}. \src{Lib/Shimura/ThetaModuli/UniversalFamily.lean\#L308-L355}{Applying the equations to reconstruct the family}.}
\noteline{16.82}{82}{\src{Lib/Shimura/Uniformisation/TorusModel.lean\#L106-L156}{Translation multipliers and irreducibility of the theta space}. \src{Lib/Shimura/Uniformisation/TorusModel.lean\#L156-L191}{The quotient section of zero Riemann form}.}
\noteline{16.83}{83}{\src{Lib/Shimura/Uniformisation/TorusModel.lean\#L193-L285}{The period map is bijective and its kernel is the stated lattice}. \src{Lib/Shimura/ThetaModuli/FamilyFromPoints.lean\#L949-L977}{Application on the complex fibre of the reconstructed family}.}
\noteline{16.84}{84}{\src{Lib/Shimura/Uniformisation/HolomorphicFamilies.lean\#L106-L130}{Existence and uniqueness of sections over the valuation ring}. \src{Lib/Shimura/IntegralModel/UniversalFamily.lean\#L68-L149}{Extension of the model, its theta coordinates and its level sections}.}
\noteline{16.85}{85}{\src{Lib/Shimura/Uniformisation/TorusModel.lean\#L304-L363}{The common unit on roots of unity}. \src{Lib/Shimura/Uniformisation/TorusModel.lean\#L365-L410}{Labeling torsion and comparing the commutator pairing}.}
\noteline{16.86}{86}{\src{Lib/Shimura/Uniformisation/TorusModel.lean\#L687-L753}{Complex linearity, the lattices and the multiplier}.}
\noteline{16.87}{87}{\src{Lib/Shimura/ArcExtension.lean\#L432-L486}{Integral sections and the closure of the graph}. \src{Lib/Shimura/ArcExtension.lean\#L486-L542}{Propagation of common zeros along the polarization torsion}. \src{Lib/Shimura/ArcExtension.lean\#L555-L615}{The graph section and identity of polarization classes}. \src{Lib/Shimura/Uniformisation/TorusModel.lean\#L754-L783}{The resulting integral isogeny and its torsion labels}.}
\noteline{16.88}{88}{\src{Lib/Shimura/Uniformisation/AdelicFrames.lean}{Signed frames and finite Gram matrices}. \src{Lib/Shimura/Transfer/Lefschetz.lean}{The marked finiteness condition}. \src{Lib/Shimura/TraceFormula/DefiniteTrace.lean}{Definite finiteness and rational returns}.}
\noteline{16.89}{89}{\src{Lib/Shimura/TraceFormula/MarkedDefiniteSums.lean}{Splitting projections, unipotent translations and block multiplicities}.}
\noteline{16.90}{90}{\src{Lib/Shimura/HeckeTower/HeckeData.lean}{Similitude and root modulus from a signed frame}. \src{Lib/Adelic/TraceFormula/ChangingSigns.lean}{Idelic norms and the equivariant frame permutation}.}
\noteline{16.91}{91}{\src{Lib/LocalGL/Iwahori/QuadraticRelation.lean}{The rank-two relation and injective dominant translation}.}
\noteline{16.92}{92}{\src{Lib/Shimura/TraceFormula/TestFunction.lean}{The two local averages and all their powers}.}
\noteline{16.93}{93}{\src{Lib/Shimura/TraceFormula/DefiniteTrace.lean}{Evaluation at definite classes and the trace identity}.}
\noteline{16.94}{94}{\src{Lib/Cohomology/Etale/Basic.lean\#L981-L1008}{Tangent polynomial and coincidence condition}. \src{Lib/Cohomology/WoodsHole.lean\#L38-L53}{The geometric vanishing implication}.}
\noteline{16.95}{95}{\src{Lib/Cohomology/Hodge/FixedPointFormula.lean\#L683-L711}{Pullback of the supported diagonal in local coordinates}. \src{Lib/Cohomology/PointClass/LinearIndependence.lean\#L626-L646}{Independence for the set of point classes}.}
\noteline{16.96}{96}{\src{Lib/Cohomology/WoodsHole.lean\#L81-L133}{Coefficient extraction and pullback on a cross product}. \src{Lib/Cohomology/WoodsHole.lean\#L174-L211}{Extraction from the scaled cup product}. \src{Lib/Cohomology/WoodsHole.lean\#L221-L293}{Pairing with the point classes}.}
\noteline{16.97}{97}{\src{Lib/Shimura/Transfer/Lefschetz.lean\#L459-L499}{Grouping coincidences on a common level}. \src{Lib/Shimura/Transfer/Lefschetz.lean\#L508-L546}{Fibres, local coefficients and the resulting marked sum}.}
\section*{Chapter 17. Comparison on test rings and cohomological Frobenius}
\noteline{17.1}{1}{\src{Lib/PadicComparison/Descent/TestCovers.lean\#L330-L429}{Ultrapowers, completion and the prescribed centre}.}
\noteline{17.2}{2}{\src{Lib/PadicComparison/Descent/TestCovers.lean\#L82-L167}{Product tilts, divisibility, regularity and the kernel of \(\theta_T\)}.}
\noteline{17.3}{3}{\src{Lib/PadicComparison/Descent/TestCovers.lean\#L195-L317}{The tensor isomorphism, Witt-vector adjunction and full faithfulness}.}
\noteline{17.4}{4}{\src{Lib/Cohomology/Etale/FilteredLimits.lean\#L462-L547}{Solution names and reflection of divisibility}. \src{Lib/PadicComparison/Descent/TestCovers.lean\#L430-L454}{Descent of the cell conditions to the image}.}
\noteline{17.5}{5}{\src{Lib/PadicComparison/Descent/TestCovers.lean\#L455-L490}{The covering point for a mixed-characteristic test}.}
\noteline{17.6}{6}{\src{Lib/Cohomology/Etale/FilteredLimits.lean\#L548-L584}{Construction of the functorial ordinary envelope}.}
\noteline{17.7}{7}{\src{Lib/Cohomology/Etale/FilteredLimits.lean\#L237-L272}{The finite idempotent calculation}. \src{Lib/Cohomology/Etale/FilteredLimits.lean\#L352-L425}{Filtered colimits and evaluation on tests}.}
\noteline{17.8}{8}{\src{Lib/Cohomology/Etale/FilteredLimits.lean\#L843-L910}{Compatible covers at the stages and the limiting vertex}.}
\noteline{17.9}{9}{\src{Lib/Cohomology/Etale/BijectiveComparisons.lean\#L429-L484}{The finite induction on the kernels of the augmented complex}. \src{Lib/Cohomology/Etale/FilteredLimits.lean\#L123-L215}{The filtered-colimit comparison of long exact sequences}. \src{Lib/Cohomology/Etale/FilteredLimits.lean\#L907-L928}{The induction for the fibre powers}.}
\noteline{17.10}{10}{\src{Lib/Cohomology/Etale/FilteredLimits.lean\#L735-L823}{Directed intersections and increasing unions of scalar algebras}.}
\noteline{17.11}{11}{\src{Lib/Cohomology/Etale/Basic.lean\#L355-L390}{Test rings and locally constant coefficients}. \src{Lib/Cohomology/ProperBaseChange/PlaceSpace.lean\#L40-L77}{The single valuative covering point and the diagonal}.}
\noteline{17.12}{12}{\src{Lib/Cohomology/ProperBaseChange/PlaceSpace.lean\#L196-L225}{The space of places as a quotient}. \src{Lib/Cohomology/ProperBaseChange/PlaceSpace.lean\#L373-L508}{Local representatives and properness}. \src{Lib/Cohomology/Etale/ProperFunctors.lean\#L758-L827}{Constructing the lift and covering the geometric points}.}
\noteline{17.13}{13}{\src{Lib/Cohomology/Etale/ProperFunctors.lean\#L383-L450}{Reduction by complementary idempotents}. \src{Lib/Cohomology/Etale/ProperFunctors.lean\#L511-L574}{The valuation lifting criterion for density}. \src{Lib/Cohomology/Etale/ProperFunctors.lean\#L729-L757}{Density of a proper geometric cover}. \src{Lib/Cohomology/Etale/FilteredLimits.lean\#L843-L908}{Filtered continuity from an affine atlas}.}
\noteline{17.14}{14}{\src{Lib/Cohomology/ProperBaseChange/Theorem.lean\#L138-L204}{The value map on the residue subring}. \src{Lib/Cohomology/ProperBaseChange/Theorem.lean\#L205-L275}{Representability in the valuation case}.}
\noteline{17.15}{15}{\src{Lib/Cohomology/ProperBaseChange/LineCharts.lean\#L130-L219}{Finite spans of Laurent powers}. \src{Lib/Cohomology/ProperBaseChange/LineCharts.lean\#L220-L292}{The correction of two idempotent representatives}.}
\noteline{17.16}{16}{\src{Lib/Cohomology/ProperBaseChange/LineCharts.lean\#L295-L371}{The polynomial system for the local case}.}
\noteline{17.17}{17}{\src{Lib/Cohomology/ProperBaseChange/SubringSplitting.lean\#L157-L207}{Choosing the two charts}. \src{Lib/Cohomology/ProperBaseChange/LineCharts.lean\#L411-L479}{Recovering locally constant exponents}. \src{Lib/Cohomology/ProperBaseChange/LineCharts.lean\#L484-L526}{Factorization on the localization charts}.}
\noteline{17.18}{18}{\src{Lib/Cohomology/ProperBaseChange/Theorem.lean\#L319-L429}{The covering and the boundary map for the split}. \src{Lib/Cohomology/ProperBaseChange/Theorem.lean\#L443-L512}{Continuity and induction on classes}.}
\noteline{17.19}{19}{\src{Lib/Cohomology/ProperBaseChange/Theorem.lean\#L520-L562}{Restriction at a fixed maximal ideal}. \src{Lib/Cohomology/ProperBaseChange/Theorem.lean\#L584-L671}{Idempotent reduction and descent from the space of places}. \src{Lib/Cohomology/Etale/BijectiveComparisons.lean\#L729-L762}{The one-step comparison of the augmented Čech complexes}.}
\noteline{17.20}{20}{\src{Lib/Cohomology/ProperBaseChange/Theorem.lean\#L717-L824}{Residue injectivity and invariance under field extension}.}
\noteline{17.21}{21}{\src{Lib/Cohomology/ProperBaseChange/Theorem.lean\#L844-L905}{Addition in the vanishing ideal}. \src{Lib/Cohomology/ProperBaseChange/Theorem.lean\#L907-L992}{Finite localization, openness and detection}.}
\noteline{17.22}{22}{\src{Lib/Cohomology/ProperBaseChange/Theorem.lean\#L1006-L1035}{The covering point for products of lines}. \src{Lib/Cohomology/ProperBaseChange/Theorem.lean\#L1080-L1108}{Properness on the zeros of the family}. \src{Lib/Cohomology/TorsionTheories.lean\#L686-L711}{Restriction onto a radical quotient}. \src{Lib/Cohomology/TorsionTheories.lean\#L786-L825}{Density for zeros coefficients}.}
\noteline{17.23}{23}{\src{Lib/Cohomology/ProperBaseChange/Theorem.lean\#L1119-L1198}{The natural identification on the category of elements}. \src{Lib/Cohomology/ProjectiveLine/TiltedLineFormula.lean\#L357-L380}{Application to the ordinary line}. \src{Lib/PadicComparison/Duality/GoodPairs.lean\#L509-L532}{Application to the multihomogeneous closed locus}.}
\noteline{17.24}{24}{\src{Lib/Cohomology/ProjectiveLine/ChernClass.lean\#L1043-L1064}{The degrees, operations, and quantification over \(S\)}. \src{Lib/Cohomology/ProjectiveLine/TiltedLineFormula.lean\#L591-L622}{The category of points and multiplication by a class}.}
\noteline{17.25}{25}{\src{Lib/Cohomology/ProjectiveLine/ChernClass.lean\#L72-L119}{The exact Kummer sequence on tests}. \src{Lib/Cohomology/ProjectiveLine/ChernClass.lean\#L138-L167}{The Tate generator}. \src{Lib/Cohomology/ProjectiveLine/ChernClass.lean\#L690-L751}{The cocycle of the row and its two-extension}.}
\noteline{17.26}{26}{\src{Lib/Cohomology/ProjectiveLine/ChernClass.lean\#L850-L898}{The determinant, its trivialization, and the supported lift}. \src{Lib/Cohomology/ProjectiveLine/TiltedLineFormula.lean\#L625-L662}{The swap identity with a zero base-point value}.}
\noteline{17.27}{27}{\src{Lib/Cohomology/ProjectiveLine/PerfectLine.lean\#L86-L115}{The two discs, torus, base point and gluing}. \src{Lib/Cohomology/ProjectiveLine/PerfectLine.lean\#L189-L214}{Objectwise injectivity}. \src{Lib/Cohomology/ProjectiveLine/PerfectLine.lean\#L260-L306}{The inverse on monomials}. \src{Lib/Cohomology/ProjectiveLine/PerfectLine.lean\#L337-L415}{Base change on perfect algebras and their acyclicity}. \src{Lib/Cohomology/ProjectiveLine/PerfectLine.lean\#L443-L473}{Resolving the additional parameter}.}
\noteline{17.28}{28}{\src{Lib/Cohomology/ProjectiveLine/TiltedLineFormula.lean\#L405-L435}{The two exact triangles}. \src{Lib/Cohomology/ProjectiveLine/TiltedLineFormula.lean\#L473-L528}{The pushout and its two-extension}. \src{Lib/Cohomology/ProjectiveLine/TiltedLineFormula.lean\#L535-L574}{Composition with the extension class}.}
\noteline{17.29}{29}{\src{Lib/Cohomology/ProjectiveLine/TiltedLineFormula.lean\#L530-L556}{Maps on the kernel and cokernel}. \src{Lib/Cohomology/ProjectiveLine/TiltedLineFormula.lean\#L853-L904}{Lifting the character and the required squares}. \src{Lib/Cohomology/ProjectiveLine/TiltedLineFormula.lean\#L914-L962}{The two chosen frames and their untilts}. \src{Lib/Cohomology/ProjectiveLine/TiltedLineFormula.lean\#L966-L1030}{The transition function, its roots and the Tate generator}.}
\noteline{17.30}{30}{\src{Lib/Cohomology/ProjectiveLine/TiltedLineFormula.lean\#L674-L724}{Evaluation at the base point and the universal swap}. \src{Lib/Cohomology/ProjectiveLine/TiltedLineFormula.lean\#L770-L849}{Cancellation on the test and on the category with a parameter}. \src{Lib/Cohomology/LerayHirsch.lean\#L290-L319}{The reduction to representables}.}
\noteline{17.31}{31}{\src{Lib/Cohomology/ProperBaseChange/Theorem.lean\#L958-L1008}{The open vanishing locus and detection on tests}. \src{Lib/Cohomology/ProjectiveLine/TiltedLineFormula.lean\#L122-L208}{Extending a field basis to a test}. \src{Lib/Cohomology/ProjectiveLine/TiltedLineFormula.lean\#L218-L285}{Selecting a basis of values on \(1\)}. \src{Lib/Cohomology/ProjectiveLine/TiltedLineFormula.lean\#L344-L380}{The affine cover and properness for the line}.}
\noteline{17.32}{32}{\src{Lib/Cohomology/ProjectiveLine/ConstantLineFormula.lean\#L60-L99}{Recovering \(1,\eta\) from a basis containing \(\eta\)}. \src{Lib/Cohomology/ProjectiveLine/ConstantLineFormula.lean\#L299-L361}{Testing the vanishing of \(\eta\) at the generic field}. \src{Lib/Cohomology/ProjectiveLine/ConstantLineFormula.lean\#L378-L447}{Propagation to the point and to the line}. \src{Lib/Cohomology/ProjectiveLine/ConstantLineFormula.lean\#L454-L496}{The nonzero element of the tilted coefficient}.}
\noteline{17.33}{33}{\src{Lib/PadicComparison/Duality/GoodPairs.lean\#L59-L84}{The ideal and its change under rescaling}. \src{Lib/PadicComparison/Duality/GoodPairs.lean\#L147-L196}{The unit locus}.}
\noteline{17.34}{34}{\src{Lib/Cohomology/ProperBaseChange/Theorem.lean\#L1103-L1147}{Field detection with coefficients on a closed locus}. \src{Lib/PadicComparison/Duality/GoodPairs.lean\#L508-L560}{Choosing classes for the zero locus}. \src{Lib/PadicComparison/Duality/GoodPairs.lean\#L567-L613}{The basis on the principal open, projective product, and relative group}.}
\noteline{17.35}{35}{\src{Lib/PadicComparison/Duality/LineIntegration.lean\#L204-L285}{Moving one coordinate using its determinant}. \src{Lib/PadicComparison/Duality/LineIntegration.lean\#L289-L363}{Iterating over coordinates}. \src{Lib/PadicComparison/Duality/LineIntegration.lean\#L683-L776}{The retract by operations and its consequences}. \src{Lib/PadicComparison/Duality/LineIntegration.lean\#L779-L871}{Integrating the parameters}. \src{Lib/PadicComparison/Duality/LineIntegration.lean\#L889-L978}{The moved evaluation, expansion, and identification with the relative class}.}
\noteline{17.36}{36}{\src{Lib/Cohomology/Etale/Basic.lean\#L355-L374}{Ordinary tests}. \src{Lib/PadicComparison/Descent/TestCategory.lean\#L243-L264}{Cells and \(p\)-adic tests}. \src{Lib/PadicComparison/Primitive/ComparisonTheorem.lean\#L25-L45}{The two cohomology theories in the comparison}.}
\noteline{17.37}{37}{\src{Lib/PadicComparison/Primitive/ComparisonTheorem.lean\#L59-L133}{The cellwise extension and its uniqueness}.}
\noteline{17.38}{38}{\src{Lib/Cohomology/SmoothBaseChange.lean\#L56-L97}{The universal lift after base change}. \src{Lib/PadicComparison/Primitive/PeriodCoefficients.lean\#L469-L480}{Extending the locally constant functions}. \src{Lib/PadicComparison/Primitive/PeriodCoefficients.lean\#L511-L555}{From a partition to the map of tilt algebras}.}
\noteline{17.39}{39}{\src{Lib/PadicComparison/Descent/AlmostDescent.lean\#L334-L389}{The Čech presentation and clearing a localization multiplier}. \src{Lib/PadicComparison/Descent/AlmostDescent.lean\#L423-L504}{The chain construction}. \src{Lib/PadicComparison/Descent/AlmostDescent.lean\#L505-L577}{Factoring through the saturated valuation domain}. \src{Lib/PadicComparison/Descent/AlmostDescent.lean\#L667-L728}{The two blow-up charts and the tensor exactness}.}
\noteline{17.40}{40}{\src{Lib/PadicComparison/Descent/AlmostDescent.lean\#L189-L258}{Realization by a valuation cell and faithful flatness}. \src{Lib/PadicComparison/Descent/AlmostDescent.lean\#L579-L655}{Passing to the quotient and detecting vanishing}. \src{Lib/PadicComparison/Descent/AlmostDescent.lean\#L781-L824}{The covering criterion}. \src{Lib/PadicComparison/Descent/AlmostDescent.lean\#L862-L935}{The covering point and induction on its Čech powers}.}
\noteline{17.41}{41}{\src{Lib/PadicComparison/Primitive/ComparisonTheorem.lean\#L152-L176}{Removing the empty boundary over a field}. \src{Lib/PadicComparison/Primitive/ComparisonTheorem.lean\#L247-L324}{The basis map and its verification at a point}. \src{Lib/PadicComparison/Duality/LineIntegration.lean\#L601-L656}{Exact sequences and the finite total cohomology condition}.}
\noteline{17.42}{42}{\src{Lib/Cohomology/ProjectiveLine/TiltedLineFormula.lean\#L1037-L1053}{The tilted projective-line formula}. \src{Lib/Cohomology/ProjectiveLine/ConstantLineFormula.lean}{Its locally constant counterpart}. \src{Lib/PadicComparison/Duality/GoodPairs.lean\#L289-L340}{Multiplication and integration over the line}. \src{Lib/PadicComparison/Duality/GoodPairs.lean\#L492-L532}{Principal opens via the zero locus and the projective product}.}
\noteline{17.43}{43}{\src{Lib/PadicComparison/Duality/GoodPairs.lean\#L345-L419}{The pairs of roots and the divided difference}. \src{Lib/PadicComparison/Duality/GoodPairs.lean\#L433-L488}{The second map and its image near the missing section}. \src{Lib/PadicComparison/Duality/GoodPairs.lean\#L615-L665}{The two principal opens and the product pair}. \src{Lib/PadicComparison/Duality/GoodPairs.lean\#L671-L721}{Excision at a standard étale chart}.}
\noteline{17.44}{44}{\src{Lib/Cohomology/SmoothBaseChange.lean\#L113-L189}{The finite cover, intersections, and relative excision}. \src{Lib/PadicComparison/Primitive/ComparisonTheorem.lean\#L325-L339}{Return from the charted functor to the actual point map}. \src{Lib/PadicComparison/HodgeTate/PeriodMap.lean\#L239-L250}{Flat base change for both functors of points}.}
\noteline{17.45}{45}{\src{Lib/PadicComparison/Primitive/PeriodCoefficients.lean\#L121-L137}{The two period rings and their divisibility conditions}. \src{Lib/PadicComparison/Primitive/PeriodCoefficients.lean\#L174-L182}{Why \(p^\flat N=0\)}. \src{Lib/PadicComparison/Primitive/PeriodCoefficients.lean\#L192-L223}{The reduction of the Witt ideal}. \src{Lib/PadicComparison/Primitive/PeriodCoefficients.lean\#L242-L345}{Reduction and regularity for the completed tensor product}.}
\noteline{17.46}{46}{\src{Lib/PadicComparison/Primitive/PeriodCoefficients.lean\#L225-L240}{The two multiplication maps}. \src{Lib/PadicComparison/Primitive/PeriodCoefficients.lean\#L347-L390}{Their composites after completion}. \src{Lib/PadicComparison/Primitive/PeriodCoefficients.lean\#L391-L437}{Applying flat base change to the comparison}.}
\noteline{17.47}{47}{\src{Lib/Cohomology/Etale/EllAdicCohomology.lean\#L142-L186}{The exact sequence, completeness and the finite generating set}.}
\noteline{17.48}{48}{\src{Lib/PadicComparison/Primitive/Periods.lean\#L114-L153}{Factorization of the actual map on a test}. \src{Lib/PadicComparison/Primitive/Periods.lean\#L163-L221}{Its reduction on the ordinary side}. \src{Lib/PadicComparison/Primitive/Periods.lean\#L233-L294}{Finite cohomology, completion, and composition}. \src{Lib/PadicComparison/HodgeTate/PeriodMap.lean\#L185-L216}{Identification with the period map on the variety}.}
\noteline{17.49}{49}{\src{Lib/PadicComparison/Crystalline/ComparisonMap.lean\#L94-L123}{The period scalars and reduction ring}. \src{Lib/PadicComparison/Crystalline/ComparisonMap.lean\#L193-L249}{The morphism of the two cohomology theories}.}
\noteline{17.50}{50}{\src{Lib/PadicComparison/Crystalline/DividedPowerEnvelope.lean\#L92-L154}{Multiplicative lifting and its Witt-vector counterpart}. \src{Lib/PadicComparison/Crystalline/DividedPowerEnvelope.lean\#L209-L264}{The envelope inside the rational algebra and its universal property}. \src{Lib/PadicComparison/Crystalline/ComparisonMap.lean\#L551-L608}{Corepresenting the root systems over a crystalline test}.}
\noteline{17.51}{51}{\src{Lib/PadicComparison/Crystalline/DividedPowerEnvelope.lean\#L294-L366}{Independence and span of the divided-power basis}.}
\noteline{17.52}{52}{\src{Lib/PadicComparison/Crystalline/ComparisonMap.lean\#L490-L523}{The bound on every fractional exponent}.}
\noteline{17.53}{53}{\src{Lib/PadicComparison/Crystalline/DividedPowerEnvelope.lean\#L367-L408}{First-order Taylor estimate in the complete envelope}.}
\noteline{17.54}{54}{\src{Lib/PadicComparison/Crystalline/DividedPowerEnvelope.lean\#L413-L485}{Passage from the integral identities to kernel and cokernel}. \src{Lib/PadicComparison/Crystalline/DividedPowerEnvelope.lean\#L487-L578}{The basis operators and their two identities}.}
\noteline{17.55}{55}{\src{Lib/PadicComparison/ToricTower.lean\#L404-L436}{The periodic sequence algebra}. \src{Lib/PadicComparison/ToricTower.lean\#L499-L572}{Compatible exponents and evaluation at a residue class}. \src{Lib/PadicComparison/ToricTower.lean\#L745-L829}{The representation on the tilts of products of cells}.}
\noteline{17.56}{56}{\src{Lib/PadicComparison/ToricTower.lean\#L443-L490}{The two partial-sum identities and tensoring}.}
\noteline{17.57}{57}{\src{Lib/PadicComparison/ToricTower.lean\#L151-L194}{The resolution on the category with a monoid action}. \src{Lib/PadicComparison/ToricTower.lean\#L228-L256}{The fibre criterion using its long exact sequence}. \src{Lib/PadicComparison/Crystalline/ComparisonMap.lean\#L810-L834}{Passage from almost coefficients to ordinary period coefficients}.}
\noteline{17.58}{58}{\src{Lib/PadicComparison/Crystalline/ComparisonMap.lean\#L438-L478}{The lift of the complete root chart}.}
\noteline{17.59}{59}{\src{Lib/PadicComparison/Crystalline/ComparisonMap.lean\#L250-L350}{Representatives on both test categories and their coefficient map}. \src{Lib/PadicComparison/Crystalline/ComparisonMap.lean\#L352-L435}{The map \(T\to W(b)\) on \(p\)-adic completions}.}
\noteline{17.60}{60}{\src{Lib/PadicComparison/Crystalline/ComparisonMap.lean\#L747-L802}{Induction after adjoining one coordinate's roots}. \src{Lib/PadicComparison/Crystalline/ComparisonMap.lean\#L843-L882}{The tensor-product chart and its universal property}.}
\noteline{17.61}{61}{\src{Lib/PadicComparison/Completion.lean\#L119-L193}{Finite rank and derived completeness after scalar extension}. \src{Lib/PadicComparison/Completion.lean\#L265-L319}{The kernel and surjectivity argument using the common reduction}. \src{Lib/PadicComparison/Crystalline/ComparisonMap.lean\#L1109-L1174}{Adic coordinates and preservation of independence}. \src{Lib/PadicComparison/Crystalline/ComparisonMap.lean\#L1200-L1268}{Identifying the spanning classes}.}
\noteline{17.62}{62}{\src{Lib/PadicComparison/Completion.lean\#L88-L112}{The sequence solution for a factorization of multiplication by \(\pi^N\)}. \src{Lib/PadicComparison/Completion.lean\#L119-L184}{The uniform bound and scalar-extension lemma}.}
\noteline{17.63}{63}{\src{Lib/PadicComparison/HodgeTate/PeriodMap.lean\#L259-L287}{The fields, schemes, coefficient section and maps in the comparison}.}
\noteline{17.64}{64}{\src{Lib/PadicComparison/HodgeTate/GenericFibre.lean\#L155-L201}{Unique lifting of coefficients}. \src{Lib/PadicComparison/HodgeTate/GenericFibre.lean\#L437-L481}{The two identifications of the square-zero kernel}.}
\noteline{17.65}{65}{\src{Lib/Cohomology/Hodge/Site.lean\#L226-L252}{The Hodge Ext groups and their pullback}. \src{Lib/PadicComparison/HodgeTate/Map.lean\#L92-L143}{The pairs, their arrows and their rings}. \src{Lib/PadicComparison/HodgeTate/Map.lean\#L154-L231}{The action on the two components}.}
\noteline{17.66}{66}{\src{Lib/PadicComparison/HodgeTate/Map.lean\#L253-L283}{The cocycle and its Yoneda powers}. \src{Lib/PadicComparison/HodgeTate/Map.lean\#L317-L366}{Scaling through the contraction}. \src{Lib/PadicComparison/HodgeTate/Map.lean\#L379-L447}{The two natural maps into cohomology on pairs}.}
\noteline{17.67}{67}{\src{Lib/PadicComparison/HodgeTate/Map.lean\#L728-L749}{The square-zero kernel for a pair}.}
\noteline{17.68}{68}{\src{Lib/PadicComparison/HodgeTate/Map.lean\#L774-L825}{Descent and the universal pair on a chart}. \src{Lib/PadicComparison/HodgeTate/Map.lean\#L823-L853}{Uniqueness on coordinates and passage to the open chart}. \src{Lib/Cohomology/Etale/BijectiveComparisons.lean\#L847-L884}{Comparison from represented fibres with equal coefficients}.}
\noteline{17.69}{69}{\src{Lib/PadicComparison/Primitive/ComparisonTheorem.lean\#L58-L135}{Passage from a product of valuation rings to its generic points}.}
\noteline{17.70}{70}{\src{Lib/PadicComparison/HodgeTate/PeriodMap.lean\#L102-L128}{Naturality, equivariance and the bijectivity hypotheses on \(J_Z\)}. \src{Lib/PadicComparison/HodgeTate/PeriodMap.lean\#L292-L335}{The coefficient pair and the definitions of \(v\) and \(\Theta_W\)}. \src{Lib/PadicComparison/HodgeTate/PeriodMap.lean\#L361-L373}{Naturality for every map out of the same cover}.}
\noteline{17.71}{71}{\src{Lib/PadicComparison/HodgeTate/GaloisWeight.lean\#L70-L119}{Twisting the generator of the infinitesimal kernel}. \src{Lib/PadicComparison/HodgeTate/GaloisWeight.lean\#L183-L262}{The arrow on pairs and the weight of the Yoneda power}. \src{Lib/PadicComparison/HodgeTate/PeriodMap.lean\#L332-L362}{Injectivity at the two steps of the weight argument}.}
\noteline{17.72}{72}{\src{Lib/Cohomology/PointClass/LinearIndependence.lean\#L625-L646}{Independence on the set of point classes}. \src{Lib/Cohomology/WoodsHole.lean\#L205-L242}{The diagonal and the functional on point classes}. \src{Lib/Shimura/Eigenspace.lean\#L443-L477}{Moving the Hecke words across all pullbacks}.}
\noteline{17.73}{73}{\src{Lib/Cohomology/Hodge/DiagonalClass.lean\#L102-L150}{Shrinking the product chart around the diagonal}.}
\noteline{17.74}{74}{\src{Lib/Cohomology/Hodge/GrothendieckSymbol.lean\#L157-L260}{Witnesses and the determinant cocycle}. \src{Lib/Cohomology/Hodge/GrothendieckSymbol.lean\#L323-L425}{Change of denominator and annihilation of the ideal}.}
\noteline{17.75}{75}{\src{Lib/Cohomology/Hodge/DepthVanishing.lean\#L63-L88}{Depth induction}. \src{Lib/Cohomology/Hodge/DepthVanishing.lean\#L273-L319}{Its application to the coordinate differences}. \src{Lib/Cohomology/Hodge/DiagonalClass.lean\#L151-L280}{Agreement on overlaps, excision and gluing}.}
\noteline{17.76}{76}{\src{Lib/Cohomology/Hodge/FixedPointFormula.lean\#L650-L711}{The two differentials at a coincidence}. \src{Lib/Cohomology/Hodge/FixedPointFormula.lean\#L781-L865}{Absolute classes and finite reduced pullback}.}
\noteline{17.77}{77}{\src{Lib/Cohomology/Hodge/Kunneth.lean\#L96-L141}{Exterior bases on a product of represented tests}. \src{Lib/Cohomology/Hodge/Kunneth.lean\#L229-L330}{Passage from these products to the atlas}. \src{Lib/Cohomology/Hodge/FixedPointFormula.lean\#L866-L935}{Global functions determine restriction to a point}.}
\noteline{17.78}{78}{\src{Lib/Cohomology/PointClass/ProjectiveNonvanishing.lean\#L80-L151}{Units in the thickening and the lifted coordinate}. \src{Lib/Cohomology/PointClass/ProjectiveNonvanishing.lean\#L202-L230}{The contraction of the logarithmic differential}. \src{Lib/Cohomology/PointClass/ProjectiveNonvanishing.lean\#L283-L355}{The two affine charts and the relative class}.}
\noteline{17.79}{79}{\src{Lib/Cohomology/ProjectiveLine/PerfectLine.lean\#L256-L306}{The inverse on rational Laurent monomials}. \src{Lib/Cohomology/ProjectiveLine/PerfectLine.lean\#L337-L411}{The calculation on perfect algebras}. \src{Lib/Cohomology/ProjectiveLine/PerfectLine.lean\#L415-L466}{Witt lifting and the acyclic cokernel}.}
\noteline{17.80}{80}{\src{Lib/Cohomology/ProjectiveLine/TiltedLineFormula.lean\#L472-L580}{Pushout, kernel and the two connecting maps}. \src{Lib/Cohomology/PointClass/ProjectiveNonvanishing.lean\#L588-L641}{Cancellation applied to an exterior product}. \src{Lib/Cohomology/PointClass/ProjectiveNonvanishing.lean\#L648-L726}{The Teichmüller lifts and the unit at the origin}.}
\noteline{17.81}{81}{\src{Lib/Cohomology/Hodge/GrothendieckSymbol.lean\#L709-L790}{The product of supported symbols}. \src{Lib/Cohomology/PointClass/ProjectiveNonvanishing.lean\#L784-L868}{The weights and the finite reduced fibre}. \src{Lib/Cohomology/PointClass/ProjectiveNonvanishing.lean\#L871-L951}{The comparison and the induction}.}
\noteline{17.82}{82}{\src{Lib/Cohomology/PointClass/LinearIndependence.lean\#L83-L139}{The component through the chosen point and its normalization}. \src{Lib/Cohomology/PointClass/LinearIndependence.lean\#L255-L327}{The étale chart and the choice of fibre}.}
\noteline{17.83}{83}{\src{Lib/Cohomology/PointClass/NoetherNormalization.lean\#L456-L549}{The homogeneous Galois closure and conjugacy}. \src{Lib/Cohomology/PointClass/ProjectiveSpace.lean\#L309-L385}{Lifting near a prime of a test ring}. \src{Lib/Cohomology/PointClass/ProjectiveSpace.lean\#L482-L543}{Local conjugacy and the transfer identity}. \src{Lib/Cohomology/PointClass/ProjectiveSpace.lean\#L547-L622}{Existence of the two maps and descent on the tests}.}
\noteline{17.84}{84}{\src{Lib/Cohomology/PointClass/LinearIndependence.lean\#L450-L512}{Pullback along an injection and detection on the fibre}. \src{Lib/Cohomology/PointClass/LinearIndependence.lean\#L540-L646}{The separating open, its fibre and linear independence}.}
\noteline{17.85}{85}{\src{Lib/Cohomology/WoodsHole.lean\#L207-L243}{The linear form used for coefficient extraction}.}
\noteline{17.86}{86}{\src{Lib/PadicComparison/CrystallineFrobenius.lean\#L703-L732}{Hypotheses on the two schemes, the place and the legs}. \src{Lib/PadicHodge/Crystalline/Basic.lean\#L129-L136}{The inverse Frobenius convention for a Weil pair}.}
\noteline{17.87}{87}{\src{Lib/PadicComparison/CrystallineFrobenius.lean\#L114-L128}{The product of primes}. \src{Lib/PadicComparison/CrystallineFrobenius.lean\#L202-L236}{Uniform choice of an affine scheme through which the test points factor}. \src{Lib/PadicComparison/CrystallineFrobenius.lean\#L388-L402}{Choosing the root of the generator}.}
\noteline{17.88}{88}{\src{Lib/PadicComparison/CrystallineFrobenius.lean\#L94-L112}{The valuation argument}. \src{Lib/PadicComparison/CrystallineFrobenius.lean\#L138-L171}{Extending equality from the generic point of a prime quotient}. \src{Lib/PadicComparison/CrystallineFrobenius.lean\#L245-L285}{Cell idempotents and pullback to another ring}.}
\noteline{17.89}{89}{\src{Lib/PadicComparison/Descent/TestCategory.lean\#L57-L80}{The functorial divided-power ideal}. \src{Lib/PadicComparison/Crystalline/ComparisonMap.lean\#L98-L123}{The coefficient structure and its map out of \(O_C\)}. \src{Lib/PadicComparison/Crystalline/ComparisonMap.lean\#L130-L172}{The Frobenius twist and vanishing of \(x_0\)}.}
\noteline{17.90}{90}{\src{Lib/PadicComparison/CrystallineFrobenius.lean\#L404-L427}{The residue action and the Teichmüller reduction map}. \src{Lib/PadicComparison/CrystallineFrobenius.lean\#L735-L779}{Finiteness, the open subgroup and the Witt lift}. \src{Lib/PadicComparison/Crystalline/DividedPowerEnvelope.lean\#L116-L147}{The unique Witt lift for a complete torsion-free test}.}
\noteline{17.91}{91}{\src{Lib/PadicComparison/Crystalline/ComparisonMap.lean\#L211-L240}{Reduction of a test and its action on Witt vectors}. \src{Lib/PadicComparison/CrystallineFrobenius.lean\#L429-L477}{Both coefficient functors and the \(R\)-linearity of the power map}.}
\noteline{17.92}{92}{\src{Lib/PadicComparison/HodgeTate/PeriodMap.lean\#L107-L126}{Naturality, equivariance and the schemes on which \(J\) is bijective}. \src{Lib/PadicComparison/HodgeTate/PeriodMap.lean\#L163-L202}{Construction by changing coefficients and pulling back}. \src{Lib/PadicComparison/CrystallineFrobenius.lean\#L780-L797}{Restriction to chart points}.}
\noteline{17.93}{93}{\src{Lib/PadicComparison/Crystalline/ComparisonMap.lean\#L445-L475}{The étale Witt lift after adjoining roots}. \src{Lib/PadicComparison/Crystalline/ComparisonMap.lean\#L610-L668}{The divided-power basis and divisibility for a nonintegral exponent}. \src{Lib/PadicComparison/Crystalline/DividedPowerEnvelope.lean\#L487-L505}{The exact sequence for the root coordinate}. \src{Lib/PadicComparison/Crystalline/ComparisonMap.lean\#L810-L861}{Descent for root systems and the integral torus chart}.}
\noteline{17.94}{94}{\src{Lib/PadicComparison/Crystalline/ComparisonMap.lean\#L1103-L1129}{The completion criterion}. \src{Lib/PadicComparison/Crystalline/ComparisonMap.lean\#L1139-L1178}{Adic coordinates, independence and finite rank}. \src{Lib/PadicComparison/Crystalline/ComparisonMap.lean\#L1189-L1235}{Flat base change, the adjunction and passage to charts}. \src{Lib/PadicComparison/Crystalline/ComparisonMap.lean\#L1236-L1268}{The two changes of coefficients and the spanning elements}.}
\noteline{17.95}{95}{\src{Lib/PadicComparison/CrystallineFrobenius.lean\#L342-L386}{The coefficient twist on pure tensors and reduced test points}. \src{Lib/PadicComparison/CrystallineFrobenius.lean\#L479-L515}{Why the same twist fixes the Witt coefficient ring}. \src{Lib/PadicComparison/CrystallineFrobenius.lean\#L545-L609}{Compatibility on points, coefficients and period maps}.}
\noteline{17.96}{96}{\src{Lib/PadicComparison/CrystallineFrobenius.lean\#L798-L825}{Full-group invariance and crystallinity}. \src{Lib/PadicHodge/Crystalline/Regularity.lean\#L325-L373}{Fixed Witt coordinates and the \(B_{\max}\)-linear extension}.}
\noteline{17.97}{97}{\src{Lib/PadicHodge/PeriodRings/DeRhamPeriods.lean\#L146-L210}{The period and its unit factor}. \src{Lib/PadicHodge/Crystalline/FixedRatios.lean\#L386-L401}{The denominator-cleared fixed-ratio assertion}.}
\noteline{17.98}{98}{\src{Lib/PadicComparison/CrystallineFrobenius.lean\#L652-L687}{Common restriction and pullback of one class}. \src{Lib/PadicComparison/CrystallineFrobenius.lean\#L688-L701}{Cancellation on the local divided-power quotient}.}
\noteline{17.99}{99}{\src{Lib/PadicHodge/Crystalline/Regularity.lean\#L407-L440}{Powers and functoriality of linear Frobenius}. \src{Lib/Shimura/Eigenspace.lean\#L645-L679}{The three atlases, residue degree, and injective leg}.}
\noteline{17.100}{100}{\src{Lib/PadicComparison/Primitive/ComparisonTheorem.lean\#L219-L247}{The primitive comparison input and its finiteness hypotheses}. \src{Lib/PadicComparison/Primitive/Periods.lean\#L299-L319}{Lifting primitive comparison to period coefficients}. \src{Lib/PadicComparison/CrystallineFrobenius.lean\#L20-L54}{The comparison statements used for the seed}.}
\endgroup

\begin{thebibliography}{99}\markboth{Bibliography}{Bibliography}\phantomsection\addcontentsline{toc}{chapter}{Bibliography}

\bibitem{rubinstein} Claude Fable 5.2, Codex Astra 6.2 and Gemini Carbon, at the behest of M.~O. Rubinstein.
\href{https://www.math.uwaterloo.ca/~mrubinst/alpoge_claude_grh/zeta-paper_v2.pdf}{\emph{The Alp\"oge--Claude proof of the Riemann hypothesis explained for the Riemann zeta function}}. Version 2, 22 November 2026.

\bibitem{montgomery-vaughan} H.~L. Montgomery and R.~C. Vaughan, \href{https://doi.org/10.1017/CBO9780511618314}{\emph{Multiplicative Number Theory I. Classical Theory}}, Cambridge Studies in Advanced Mathematics 97, Cambridge University Press, 2006.

\bibitem{pritsker} I.~E. Pritsker, \href{https://doi.org/10.4153/CJM-2005-041-8}{The Gelfond--Schnirelman method in prime number theory}, \emph{Canad. J. Math.} \textbf{57} (2005), 1080--1101.

\bibitem{oldguide} The \emph{riemann-rev} project.
\href{https://nasqret.github.io/riemann-rev/guide/proof-guide.pdf}{\emph{The proof of the Riemann hypothesis announced by L.~Alp\"oge. Its architecture: the statements, the flow, and why each idea calls for the next}}.

\bibitem{katz-cartier} N.~M. Katz, \href{https://web.math.princeton.edu/~nmk/old/divcongCartier.pdf}{Divisibilities, congruences, and Cartier duality}, \emph{J. Fac. Sci. Univ. Tokyo Sect. IA Math.} \textbf{28} (1981), 667--678.

\bibitem{bannai-integral} K. Bannai and S. Kobayashi, \href{https://www.numdam.org/articles/10.5802/aif.3018/}{Integral structures on \(p\)-adic Fourier theory}, \emph{Ann. Inst. Fourier (Grenoble)} \textbf{66} (2016), 521--550.

\bibitem{upper-moore-wolf} C.~C.~Moore and J.~A.~Wolf, \href{https://math.berkeley.edu/~jawolf/publications.pdf/paper_056.pdf}{Square integrable representations of nilpotent groups}, \emph{Trans. Amer. Math. Soc.} \textbf{185} (1973), 445--462.

\bibitem{gan-qiu-takeda-rallis} W.~T. Gan, Y.~Qiu and S.~Takeda, \href{https://arxiv.org/pdf/1207.4709}{The regularized Siegel--Weil formula (the second term identity) and the Rallis inner product formula}, \emph{Invent. Math.} \textbf{198} (2014), 739--831.

\bibitem{behrend} F.~A. Behrend, \href{https://doi.org/10.1073/pnas.32.12.331}{On sets of integers which contain no three terms in arithmetical progression}, \emph{Proc. Nat. Acad. Sci. U.S.A.} \textbf{32} (1946), 331--332.

\bibitem{shapira-behrend} A. Shapira, \href{https://www.math.tau.ac.il/~asafico/behrend.pdf}{Behrend-type constructions for sets of linear equations}, \emph{Acta Arith.} \textbf{122} (2006), 17--33. The link is to the author's preprint.

\bibitem{alon-shapira} N. Alon and A. Shapira, \href{https://www.math.tau.ac.il/~asafico/test_induced.pdf}{A characterization of easily testable induced subgraphs}, \emph{Combin. Probab. Comput.} \textbf{15} (2006), 791--805. The link is to the authors' preprint.

\bibitem{newman} D.~J. Newman, \href{https://www.jstor.org/stable/2321853}{Simple analytic proof of the prime number theorem}, \emph{Amer. Math. Monthly} \textbf{87} (1980), 693--696.

\bibitem{jutila} M. Jutila, \href{https://matwbn.icm.edu.pl/ksiazki/aa/aa27/aa27116.pdf}{On mean values of Dirichlet polynomials with real characters}, \emph{Acta Arith.} \textbf{27} (1975), 191--198.

\bibitem{heath-brown-sieve} D.~R. Heath-Brown, \href{https://doi.org/10.4064/aa-72-3-235-275}{A mean value estimate for real character sums}, \emph{Acta Arith.} \textbf{72} (1995), 235--275.

\bibitem{widder-laplace} D.~V. Widder, \href{https://doi.org/10.1515/9781400876457}{\emph{The Laplace Transform}}, Princeton Mathematical Series 6, Princeton University Press, 1941.

\bibitem{suzuki} M. Suzuki, \href{https://doi.org/10.1007/s11139-025-01238-9}{On variants of Chebyshev's conjecture}, \emph{Ramanujan J.} \textbf{68} (2025), article 95.

\bibitem{blake-deuring} C. Blake, \href{https://doi.org/10.1112/blms/bdu079}{A Deuring criterion for abelian varieties}, \emph{Bull. Lond. Math. Soc.} \textbf{46} (2014), 1256--1263. \href{https://arxiv.org/pdf/1311.5454}{Author's preprint}.

\bibitem{theta-dlmf} F. W. J. Olver et al., eds.,
\href{https://dlmf.nist.gov/}{NIST Digital Library of Mathematical Functions},
Chs.~20 and~23.

\bibitem{upper-milne-cm} J.~S.~Milne, \href{https://www.jmilne.org/math/CourseNotes/CM.pdf}{\emph{Complex Multiplication}}, version 0.10, 2020.

\bibitem{legacy-bannai-kobayashi} K.~Bannai and S.~Kobayashi, \href{https://doi.org/10.1215/00127094-2010-024}{Algebraic theta functions and the \(p\)-adic interpolation of Eisenstein--Kronecker numbers}, \emph{Duke Math. J.} \textbf{153} (2010), 229--295. Numbering follows the \href{https://arxiv.org/pdf/math/0610163}{authors' preprint}.

\bibitem{theta-siegel} C. L. Siegel,
\href{https://mathweb.tifr.res.in/sites/default/files/publications/ln/tifr23.pdf}{Lectures on advanced analytic number theory},
Tata Institute of Fundamental Research Lectures on Mathematics 23, Ch.~I, \S1.

\bibitem{upper-gross-periods} B.~H.~Gross, \href{https://arxiv.org/pdf/2005.04194}{On the periods of abelian varieties}, \emph{Notices of the ICCM} \textbf{8} (2020), no.~2, 10--18. Section numbers refer to the linked author's version.

\bibitem{upper-fourier} P.~Schneider and J.~Teitelbaum, \href{https://ems.press/journals/dm/articles/8965041}{\(p\)-adic Fourier theory}, \emph{Doc. Math.} \textbf{6} (2001), 447--481.

\bibitem{upper-benson-doty} D.~Benson and S.~Doty, \href{https://arxiv.org/pdf/0805.1235}{Schur--Weyl duality over finite fields}, \emph{Arch. Math. (Basel)} \textbf{93} (2009), 425--435.

\bibitem{upper-greenberg-voight} M.~Greenberg and J.~Voight, \href{https://jvoight.github.io/articles/orthohecke-240123.pdf}{Lattice methods for algebraic modular forms on classical groups}, in \emph{Computations with Modular Forms}, Contributions in Mathematical and Computational Sciences 6, Springer, 2014, 147--179.

\bibitem{upper-kudla} S.~S.~Kudla, \href{https://wiki.epfl.ch/tan/documents/1984\%20kudla-1.pdf}{Seesaw dual reductive pairs}, in \emph{Automorphic Forms of Several Variables}, Progress in Mathematics 46, Birkh\"auser, 1984, 244--268.

\bibitem{legacy-weil} A.~Weil, \href{https://doi.org/10.1007/BF02391012}{Sur certains groupes d'op\'erateurs unitaires}, \emph{Acta Math.} \textbf{111} (1964), 143--211.

\bibitem{upper-wang} B.~P.~J.~Wang, \href{https://arxiv.org/pdf/2311.14234}{Global theta lifting and automorphic periods associated to nilpotent orbits}, \emph{J. Number Theory} \textbf{271} (2025), 122--149.

\bibitem{upper-furusawa-morimoto} M.~Furusawa and K.~Morimoto, \href{https://arxiv.org/pdf/2205.09471}{On the Gan--Gross--Prasad conjecture and its refinement for \(({\rm U}(2n),{\rm U}(1))\)}, \emph{Math. Ann.} \textbf{391} (2025), 3799--3862.

\bibitem{theta-howe} R. Howe,
\href{https://www.math.auckland.ac.nz/~obrien/howe.pdf}{Standard monomial theory for \(\mathrm{GL}_n\)}, notes for the 2010 NZMRI Summer Workshop, \S\S4--5.

\bibitem{upper-bombieri-vaaler} E.~Bombieri and J.~Vaaler, \href{https://doi.org/10.1007/BF01393823}{On Siegel's lemma}, \emph{Invent. Math.} \textbf{73} (1983), 11--32. For the formulation in terms of subspace height see L.~Fukshansky, \href{https://www1.cmc.edu/pages/faculty/Lenny/papers/CIMPA_2024-lecture_notes.pdf}{\emph{Search bounds for Diophantine equations}}, CIMPA, 2024, Theorem~9.1.

\bibitem{theta-knapp} A. W. Knapp,
\href{https://www.math.stonybrook.edu/~aknapp/download/Beyond2-clickable.pdf}{Lie Groups Beyond an Introduction},
2nd ed., Progress in Mathematics 140, Birkh\"auser, 2002. The linked digital second edition appeared in 2023.

\bibitem{theta-flensted} M. Flensted-Jensen,
\href{https://annals.math.princeton.edu/1980/111-2/p05}{Discrete series for semisimple symmetric spaces},
\emph{Ann. of Math. (2)} \textbf{111} (1980), 253--311.

\bibitem{theta-baldoni-vergne} V. Baldoni and M. Vergne,
\href{https://arxiv.org/abs/1008.5360}{Discrete series representations and \(K\)-multiplicities for \(U(p,q)\). User's guide}, arXiv:1008.5360 (2010), \S2.

\bibitem{theta-banaszczyk} W. Banaszczyk,
\href{https://doi.org/10.1007/BF01445125}{New bounds in some transference theorems in the geometry of numbers},
\emph{Math. Ann.} \textbf{296} (1993), 625--635.

\bibitem{theta-borel} A. Borel,
\href{https://numdam.org/articles/10.1007/BF02684289/}{Some finiteness properties of adele groups over number fields},
\emph{Publ. Math. IH\'ES} \textbf{16} (1963), 5--30.

\bibitem{theta-green} J. A. Green,
\href{https://www.ams.org/tran/1955-080-02/S0002-9947-1955-0072878-2/S0002-9947-1955-0072878-2.pdf}{The characters of the finite general linear groups},
\emph{Trans. Amer. Math. Soc.} \textbf{80} (1955), 402--447.

\bibitem{upper-labesse} J.-P.~Labesse, \href{https://webusers.imj-prg.fr/~michael.harris/CBCM.pdf}{Changement de base CM et s\'eries discr\`etes}, in \emph{On the Stabilization of the Trace Formula}, vol.~1, International Press, 2011, 429--470.

\bibitem{arith-cft} J. S. Milne. \emph{Class Field Theory}, version 4.03, 2020. \href{https://www.jmilne.org/math/CourseNotes/CFTc.pdf}{Author's edition.}

\bibitem{legacy-kneser} M.~Kneser, \href{https://mathweb.tifr.res.in/sites/default/files/publications/ln/tifr47.pdf}{\emph{Lectures on Galois Cohomology of Classical Groups}}, Tata Institute of Fundamental Research Lectures on Mathematics and Physics 47, Bombay, 1969.

\bibitem{theta-deligne-lusztig} P. Deligne and G. Lusztig,
\href{https://publications.ias.edu/sites/default/files/Number27.pdf}{Representations of reductive groups over finite fields},
\emph{Ann. of Math. (2)} \textbf{103} (1976), 103--161.

\bibitem{theta-arthur} J. Arthur,
\href{https://www.claymath.org/library/cw/arthur/pdf/4.pdf}{The characters of discrete series as orbital integrals},
\emph{Invent. Math.} \textbf{32} (1976), 205--261.

\bibitem{theta-hc-orbital} Harish-Chandra,
\href{https://doi.org/10.1007/BF02392813}{Discrete series for semisimple Lie groups. II. Explicit determination of the characters},
\emph{Acta Math.} \textbf{116} (1966), 1--111.

\bibitem{theta-atiyah-schmid} M. F. Atiyah and W. Schmid,
\href{https://www.math.utah.edu/~ptrapa/math-library/schmid/atiyah-schmid.pdf}{A geometric construction of the discrete series for semisimple Lie groups},
\emph{Invent. Math.} \textbf{42} (1977), 1--62. \href{https://doi.org/10.1007/BF01408936}{Erratum}, \textbf{54} (1979), 189--192.

\bibitem{theta-bernstein} J. N. Bernstein,
\href{https://www.math.tau.ac.il/~bernstei/Publication_list/publication_texts/Bernstein-P-invar-SLN.pdf}{\(P\)-invariant distributions on \(\mathrm{GL}(N)\) and the classification of unitary representations of \(\mathrm{GL}(N)\) (non-Archimedean case)},
in \emph{Lie Group Representations, II}, Lecture Notes in Mathematics 1041, Springer, 1984, 50--102.

\bibitem{prasanna-norms} K.~Prasanna, \href{https://sites.lsa.umich.edu/kartikp/wp-content/uploads/sites/1326/2024/07/prasanna_final.pdf}{Integrality of a ratio of Petersson norms and level-lowering congruences}, \emph{Ann. of Math.} \textbf{163} (2006), 901--967.

\bibitem{gaudron} \'E. Gaudron, \href{https://www.numdam.org/item/RSMUP_2008__119__21_0/}{Pentes des fibr\'es vectoriels ad\'eliques sur un corps global}, \emph{Rend. Semin. Mat. Univ. Padova} \textbf{119} (2008), 21--95. \href{https://www-fourier.univ-grenoble-alpes.fr/sites/default/files/REF_688.pdf}{Author's preprint}.

\bibitem{bost} J.-B. Bost, \href{https://numdam.org/articles/10.1007/s10240-001-8191-3/}{Algebraic leaves of algebraic foliations over number fields}, \emph{Publ. Math. Inst. Hautes \'Etudes Sci.} \textbf{93} (2001), 161--221.

\bibitem{yang-hecke} T.~Yang, \href{https://people.math.wisc.edu/~tonghaiyang/thetashort.pdf}{Theta liftings and Hecke \(L\)-functions}, \emph{J. Reine Angew. Math.} \textbf{485} (1997), 25--53.

\bibitem{snitz-matching} K.~Snitz, \href{https://arxiv.org/pdf/math/0509486}{Theta correspondence of automorphic characters}, \emph{J. Number Theory} \textbf{122} (2007), 429--465.

\bibitem{weil-siegel-formula} A.~Weil, \href{https://doi.org/10.1007/BF02391774}{Sur la formule de Siegel dans la th\'eorie des groupes classiques}, \emph{Acta Math.} \textbf{113} (1965), 1--87.

\bibitem{wu-unitary-tower} C.~Wu, \href{https://arxiv.org/pdf/1409.0770}{Irreducibility of theta lifting for unitary groups}, \emph{J. Number Theory} \textbf{133} (2013), 3296--3318.

\bibitem{hecke1920} E.~Hecke.
\href{https://geodesic.mathdoc.fr/item/MZ_1920__6_167556/}{Eine neue Art von Zetafunktionen und ihre Beziehungen zur Verteilung der Primzahlen. Zweite Mitteilung}.
\emph{Math. Z.} \textbf{6} (1920), 11--51.

\bibitem{tate1950} J.~T. Tate. Fourier analysis in number fields and Hecke's zeta-functions. \href{https://sites.math.rutgers.edu/~alexk/2023S572/Tate1950.pdf}{Princeton dissertation}, 1950. Reprinted in J.~W.~S. Cassels and A.~Fr\"ohlich (eds.), \emph{Algebraic Number Theory}. Academic Press, 1967, 305--347.

\bibitem{katz-cm-lfunctions} N.~M. Katz, \href{https://web.math.princeton.edu/~nmk/old/padicLCM.pdf}{\(p\)-adic \(L\)-functions for CM fields}, \emph{Invent. Math.} \textbf{49} (1978), 199--297.

\bibitem{chida-hsieh-theta} M.~Chida and M.-L.~Hsieh, \href{https://www.math.ntu.edu.tw/~mlhsieh/research/central.pdf}{Special values of anticyclotomic \(L\)-functions for modular forms}, \emph{J. Reine Angew. Math.} \textbf{741} (2018), 87--131.

\bibitem{arith-conrad} B. Conrad. \emph{Lifting global representations with local properties}. 2011, Appendix~B. \href{https://math.stanford.edu/~conrad/papers/locchar.pdf}{Author's text.}

\bibitem{arith-ddt} H. Darmon, F. Diamond, and R. Taylor. \href{https://www.math.mcgill.ca/darmon/pub/Articles/Expository/05.DDT/paper.pdf}{Fermat's Last Theorem}. \emph{Current Developments in Mathematics} 1995, 1--154, International Press. The theorem numbers refer to the linked authors' version dated September 9, 2007.

\bibitem{arith-milne} J. S. Milne. \emph{Arithmetic Duality Theorems}, 2nd ed. BookSurge, 2006. \href{https://www.jmilne.org/math/Books/ADTnot.pdf}{Author's edition.}

\bibitem{mazur-rubin-stability} B.~Mazur and K.~Rubin, with an appendix by M.~Larsen, \href{https://arxiv.org/pdf/1503.04642}{Diophantine stability}, \emph{Amer. J. Math.} \textbf{140} (2018), 571--616.

\bibitem{arith-gross-koblitz} B. H. Gross and N. Koblitz. Gauss sums and the \(p\)-adic \(\Gamma\)-function. \emph{Ann. of Math.} 109 (1979), 569--581. \href{https://annals.math.princeton.edu/1979/109-3/p06}{Publisher.}

\bibitem{rubin-euler-systems} K.~Rubin, \href{https://swc-math.github.io/aws/1999/99RubinES.pdf}{\emph{Euler Systems}}, Annals of Mathematics Studies 147, Princeton University Press, 2000.

\bibitem{legacy-milne-ant} J.~S.~Milne, \href{https://www.jmilne.org/math/CourseNotes/ANT301.pdf}{\emph{Algebraic Number Theory}}, version 3.01, September 28, 2008.

\bibitem{dksw-ribet} S.~Dasgupta, M.~Kakde, J.~Silliman and J.~Wang, \href{https://arxiv.org/pdf/2310.16396}{The residually indistinguishable case of Ribet's method for \(\mathrm{GL}_2\)}, arXiv:2310.16396 (2023).

\bibitem{arith-stacks} The Stacks Project Authors. \href{https://stacks.math.columbia.edu/}{\emph{The Stacks Project}}.

\bibitem{arith-colmez} P. Colmez. Th\'eorie d'Iwasawa des repr\'esentations de de Rham d'un corps local. \emph{Ann. of Math.} 148 (1998), 485--571. \href{https://webusers.imj-prg.fr/~pierre.colmez/iwasawa-deRham.pdf}{Author's text.}

\bibitem{arith-bloch-kato} S. Bloch and K. Kato. \(L\)-functions and Tamagawa numbers of motives. In \emph{The Grothendieck Festschrift}, Vol.~I, Progress in Mathematics 86, Birkh\"auser, 1990, 333--400. \href{https://math.stanford.edu/~conrad/BSDseminar/refs/BKTamagawa.pdf}{Text.}

\bibitem{legacy-bms} B.~Bhatt, M.~Morrow and P.~Scholze, \href{https://numdam.org/articles/10.1007/s10240-019-00102-z/}{Integral \(p\)-adic Hodge theory}, \emph{Publ. Math. Inst. Hautes \'Etudes Sci.} \textbf{128} (2018), 219--397.

\bibitem{legacy-ax} J.~Ax, \href{https://doi.org/10.1016/0021-8693(70)90069-4}{Zeros of polynomials over local fields---the Galois action}, \emph{J. Algebra} \textbf{15} (1970), 417--428.

\bibitem{legacy-brinon-conrad} O.~Brinon and B.~Conrad, \href{https://math.stanford.edu/~conrad/papers/notes.pdf}{\emph{CMI summer school notes on \(p\)-adic Hodge theory}}, preliminary version, June 24, 2009.

\bibitem{arith-tate-pdiv} J.~T. Tate.
\href{https://www.math.purdue.edu/~tongliu/teaching/598/p-divisible.pdf}{\(p\)-Divisible groups}.
In \emph{Proceedings of a Conference on Local Fields}, Springer, 1967, 158--183.

\bibitem{lit-chenevier-harris} G.~Chenevier and M.~Harris, \href{https://intlpress.com/site/pub/files/_fulltext/journals/cjm/2013/0001/0001/CJM-2013-0001-0001-a002.pdf}{Construction of automorphic Galois representations, II}, \emph{Cambridge J. Math.} \textbf{1} (2013), 53--73.

\bibitem{arith-shintani} T. Shintani. Two remarks on irreducible characters of finite general linear groups. \emph{J. Math. Soc. Japan} 28 (1976), 396--414. \href{https://doi.org/10.2969/jmsj/02820396}{Publisher.}

\bibitem{arith-silberger-zink} A. J. Silberger and E.-W. Zink. Explicit Shintani base change and the Macdonald correspondence for characters of \(\mathrm{GL}_n(k)\). \emph{J. Algebra} 319 (2008), 4147--4176. \href{https://doi.org/10.1016/j.jalgebra.2008.02.010}{Publisher.}

\bibitem{arith-roche} A. Roche. Types and Hecke algebras for principal series representations of split reductive \(p\)-adic groups. \emph{Ann. Sci. \'Ecole Norm. Sup.} 31 (1998), 361--413. \href{https://www.numdam.org/item/ASENS_1998_4_31_3_361_0/}{Text.}

\bibitem{arith-casselman} W. Casselman. The unramified principal series of \(p\)-adic groups I. The spherical function. \emph{Compositio Math.} 40 (1980), 387--406. \href{https://www.numdam.org/item/CM_1980__40_3_387_0/}{Text.}

\bibitem{arith-hkp} T. J. Haines, R. E. Kottwitz and A. Prasad. Iwahori--Hecke algebras. \emph{J. Ramanujan Math. Soc.} 25 (2010), 113--145. \href{https://arxiv.org/abs/math/0309168}{Author version.}

\bibitem{arith-rouquier} R. Rouquier. Caract\'erisation des caract\`eres et pseudo-caract\`eres. \emph{J. Algebra} 180 (1996), 571--586. \href{https://www.math.ucla.edu/~rouquier/papers/carcar.pdf}{Author's text.}

\bibitem{legacy-henk} M.~Henk, \href{https://arxiv.org/abs/math/0204158}{Successive minima and lattice points}, \emph{Rend. Circ. Mat. Palermo (2), Suppl.} \textbf{70} (2002), 377--384.

\bibitem{lit-liu-torsion} T.~Liu, \href{https://www.numdam.org/item/ASENS_2007_4_40_4_633_0.pdf}{Torsion \(p\)-adic Galois representations and a conjecture of Fontaine}, \emph{Ann. Sci. \'Ecole Norm. Sup.} (4) \textbf{40} (2007), 633--674.

\bibitem{lit-wedhorn} T.~Wedhorn, \href{https://doi.org/10.1515/crll.2000.060}{Congruence relations on some Shimura varieties}, \emph{J. reine angew. Math.} \textbf{524} (2000), 43--71.

\bibitem{lit-katz-serre-tate} N.~M.~Katz, \href{https://web.math.princeton.edu/~nmk/old/serretatelocmod.pdf}{Serre--Tate local moduli}, in \emph{Surfaces alg\'ebriques} (Orsay, 1976--78), Lecture Notes in Mathematics \textbf{868}, Springer, 1981, 138--202.

\bibitem{legacy-mumford} D.~Mumford, \href{https://pazuki.perso.math.cnrs.fr/index_fichiers/Mumford66.pdf}{On the equations defining abelian varieties. I}, \emph{Invent. Math.} \textbf{1} (1966), 287--354.

\bibitem{lit-kondyrev-prikhodko} G.~Kondyrev and A.~Prikhodko, \href{https://armj.math.stonybrook.edu/pdf-Springer-final/022-0206.pdf}{Holomorphic Atiyah--Bott formula for correspondences}, \emph{Arnold Math. J.} \textbf{8} (2022), 497--511.

\bibitem{lit-artin-proper} M.~Artin, \href{https://link.springer.com/chapter/10.1007/BFb0070718}{Th\'eor\`eme de changement de base pour un morphisme propre}, in \emph{Th\'eorie des topos et cohomologie \'etale des sch\'emas} (SGA~4), tome~3, expos\'e~XII, Lecture Notes in Mathematics \textbf{305}, Springer, 1973, 79--131.

\bibitem{legacy-scholze} P.~Scholze, \href{https://www.math.uni-bonn.de/people/scholze/pAdicHodgeTheory.pdf}{\(p\)-adic Hodge theory for rigid-analytic varieties}, \emph{Forum Math. Pi} \textbf{1} (2013), e1.

\bibitem{legacy-bhatt} B.~Bhatt, \href{https://swc-math.github.io/aws/2017/2017BhattNotes.pdf}{\emph{The Hodge--Tate decomposition via perfectoid spaces}}, Arizona Winter School lecture notes, 2017.

\bibitem{legacy-berger} L.~Berger, \href{https://perso.ens-lyon.fr/laurent.berger/autrestextes/barcelone.pdf}{On \(p\)-adic Galois representations}, in \emph{Elliptic Curves, Hilbert Modular Forms and Galois Deformations}, Advanced Courses in Mathematics, CRM Barcelona, Birkh\"auser, 2013, 3--19.

\end{thebibliography}\end{document}
